% Generated by matins/tools/build_latex.py from data/corpus.json — do not edit.
\documentclass[10pt,twoside,openright]{book}
\usepackage{matins}
\begin{document}
\BodySize
\frontmatter
\thispagestyle{empty}
\vspace*{0.25\textheight}
{\centering
{\Huge\rubr{Lectiones ad Matutinum}}\par\vspace{0.4em}
{\large\itshape ex Breviario Romano anni MCMLIV}\par\vspace{2em}
{\LARGE The Lessons of Matins}\par\vspace{0.4em}
{\large\itshape from the Roman Breviary of 1954, in Latin and English}\par\vspace{3em}
{\Large\rubr{Pars Hiemalis}}\par\vspace{0.3em}{\large\itshape Winter}\par\vspace{2em}
{\small }\par}
\clearpage
\thispagestyle{empty}
{\footnotesize\setlength{\parindent}{0pt}\setlength{\parskip}{0.35em}\raggedright
\par\addvspace{0.6em}\textsc{Divinum Officium}\par
The texts, and the program that arranges them according to the rubrics, are the work of the Divinum Officium Project (divinumofficium.com), begun by the late Laszlo Kiss and continued by the Project's volunteers. Its code and data are available at github.com/DivinumOfficium/divinum-officium under the MIT License.\par
\par\addvspace{0.6em}\textsc{Latin}\par
The Latin text derives from the Ratisbon 1888 edition of the Breviarium Romanum, scanned by the University of St Michael's College, Toronto, amended from the 1906 edition, and corrected by Laszlo Kiss and the Project team against later printings: Pustet (Regensburg, 1943), Benziger Brothers (New York, 1945), Burns Oates \& Washbourne (London, 1946), La Presse Catholique Panaméricaine (Montréal, 1943), and Desclée and Mame (1962).\par
The accented text of the Sixto-Clementine Vulgate was supplied by Frère Romain-Marie of the Abbaye Saint-Joseph de Clairval, Flavigny-sur-Ozerain.\par
\par\addvspace{0.6em}\textsc{English}\par
Holy Scripture is given in the Douay-Rheims translation.\par
The other lessons and the responsories are from the translation of the Roman Breviary by John, Marquess of Bute (edition of 1908), made available by the University of St Michael's College, Toronto.\par
\par\addvspace{0.6em}\textsc{This edition}\par
The arrangement of this edition follows the Divinum Officium program's choice of lessons for each day under the rubrics of 1954. Any errors in the arrangement are the editor's, not the Project's.\par
\par\addvspace{0.6em}\textsc{Typeface}\par
Set in EB Garamond, designed by Georg Duffner and Octavio Pardo, under the SIL Open Font License 1.1.\par
}
\mainmatter
\PartTitle{Proprium de Tempore}{Proper of Time}

\addcontentsline{toc}{chapter}{Proprium de Tempore\enspace\textit{Proper of Time}}

\SeasonTitle{Tempus Adventus}{Advent}

\addcontentsline{toc}{section}{Tempus Adventus\enspace\textit{Advent}}

\Office{Dominica I Adventus}{First Sunday of Advent}{Semiduplex I classis}
\addcontentsline{toc}{subsection}{Dominica I Adventus\enspace\textit{First Sunday of Advent}}
\begin{paracol}{2}
\LA
\LessonHead{Lectio I}
\switchcolumn
\EN
\LessonHead{Lesson I}
\switchcolumn*

\LA
\LessonTitle{Incipit liber Isaíæ Prophétæ}
\switchcolumn
\EN
\LessonTitle{Lesson from the book of Isaias}
\switchcolumn*

\LA
\Cite{Isa 1:1-3}
\switchcolumn
\EN
\Cite{Isa 1:1-3}
\switchcolumn*

\LA
\noindent Vísio Isaíæ fílii Amos, quam vidit super Judam et Jerúsalem, in diébus Ozíæ, Jóatham, Achaz, et Ezechíæ, regum Juda. Audíte, cæli, et áuribus pércipe, terra, quóniam Dóminus locútus est: Fílios enutrívi, et exaltávi: ipsi autem sprevérunt me. Cognóvit bos possessórem suum, et ásinus præsépe dómini sui: Israël autem me non cognóvit, et pópulus meus non intelléxit.\par
\switchcolumn
\EN
\noindent The vision of Isaias the son of Amos, which he saw concerning Juda and Jerusalem in the days of Ozias, Joathan, Achaz, and Ezechias, kings of Juda Hear, O ye heavens, and give ear, O earth, for the Lord hath spoken. I have brought up children, and exalted them: but they have despised me. The ox knoweth his owner, and the ass his master's crib: but Israel hath not known me, and my people hath not understood.\par
\switchcolumn*

\LA
\begin{Resp}\Rsign Aspíciens a longe, ecce vídeo Dei poténtiam veniéntem, et nébulam totam terram tegéntem. \Star{} Ite óbviam ei, et dícite: * Núntia nobis, si tu es ipse, * Qui regnatúrus es in pópulo Israël.\end{Resp}
\begin{Resp}\Vsign Quique terrígenæ, et fílii hóminum, simul in unum dives et pauper. Ite óbviam ei, et dícite.\end{Resp}
\begin{Resp}\Vsign Qui regis Israël, inténde, qui dedúcis velut ovem Joseph. Núntia nobis, si tu es ipse.\end{Resp}
\begin{Resp}\Vsign Tóllite portas, príncipes, vestras, et elevámini portæ æternáles, et introíbit Rex glóriæ. Qui regnatúrus es in pópulo Israël.\end{Resp}
\begin{Resp}\Vsign Glória Patri, et Fílio, \Star{} et Spirítui Sancto.\end{Resp}
\begin{Resp}\Rsign Aspíciens a longe, ecce vídeo Dei poténtiam veniéntem, et nébulam totam terram tegéntem. * Ite óbviam ei, et dícite: * Núntia nobis, si tu es ipse, * Qui regnatúrus es in pópulo Israël.\end{Resp}
\switchcolumn
\EN
\begin{Resp}\Rsign I look from afar, and behold I see the Power of God, coming like as a cloud to cover the land with the hosts of his People: \Star{} Go ye out to meet him and say: * Tell us if thou art he, * That shalt reign over God's people Israel.\end{Resp}
\begin{Resp}\Vsign All ye that dwell in the world, all ye children of men, high and low, rich and poor, one with another. Go ye out to meet him and say.\end{Resp}
\begin{Resp}\Vsign Hear, O thou Shepherd of Israel, thou that leadest Joseph like a sheep. Tell us if thou art he.\end{Resp}
\begin{Resp}\Vsign Lift up your heads, O ye gates, and be ye lift up, ye everlasting doors, and the King of Glory shall come in. That shalt reign over God's people Israel.\end{Resp}
\begin{Resp}\Vsign Glory be to the Father, and to the Son, \Star{} and to the Holy Ghost.\end{Resp}
\begin{Resp}\Rsign I look from afar, and behold I see the Power of God, coming like as a cloud to cover the land with the hosts of his People: * Go ye out to meet him and say: * Tell us if thou art he, * That shalt reign over God's people Israel.\end{Resp}
\switchcolumn*

\LA
\LessonHead{Lectio II}
\switchcolumn
\EN
\LessonHead{Lesson II}
\switchcolumn*

\LA
\Cite{Isa 1:4-6}
\switchcolumn
\EN
\Cite{Isa 1:4-6}
\switchcolumn*

\LA
\noindent Væ genti peccatríci, pópulo gravi iniquitáte, sémini nequam, fíliis scelerátis: dereliquérunt Dóminum, blasphemavérunt Sanctum Israël, abalienáti sunt retrórsum. Super quo percútiam vos ultra, addéntes prævaricatiónem? omne caput lánguidum, et omne cor mærens. A planta pedis usque ad vérticem non est in eo sánitas: vulnus, et livor, et plaga tumens non est circumligáta, nec curáta medicámine, neque fota óleo.\par
\switchcolumn
\EN
\noindent Woe to the sinful nation, a people laden with iniquity, a wicked seed, ungracious children: they have forsaken the Lord, they have blasphemed the Holy One of Israel, they are gone away backwards. For what shall I strike you any more, you that increase transgression? the whole head is sick, and the whole heart is sad. From the sole of the foot unto the top of the head, there is no soundness therein: wounds and bruises and swelling sores: they are not bound up, nor dressed, nor fomented with oil.\par
\switchcolumn*

\LA
\begin{Resp}\Rsign Aspiciébam in visu noctis, et ecce in núbibus cæli Fílius hóminis veniébat: et datum est ei regnum, et honor: \Star{} Et omnis pópulus, tribus, et linguæ sérvient ei.\end{Resp}
\begin{Resp}\Vsign Potéstas ejus, potéstas ætérna, quæ non auferétur: et regnum ejus, quod non corrumpétur.\end{Resp}
\begin{Resp}\Rsign Et omnis pópulus, tribus, et linguæ sérvient ei.\end{Resp}
\switchcolumn
\EN
\begin{Resp}\Rsign I saw in the night visions, and, behold, the Son of man came with the clouds of heaven, and there was given Him a Kingdom, and glory; \Star{} And all people, nations, and languages shall serve Him.\end{Resp}
\begin{Resp}\Vsign His dominion is an everlasting dominion which shall not pass away, and His Kingdom that which shall not be destroyed.\end{Resp}
\begin{Resp}\Rsign And all people, nations, and languages shall serve Him.\end{Resp}
\switchcolumn*

\LA
\LessonHead{Lectio III}
\switchcolumn
\EN
\LessonHead{Lesson III}
\switchcolumn*

\LA
\Cite{Isa 1:7-9}
\switchcolumn
\EN
\Cite{Isa 1:7-9}
\switchcolumn*

\LA
\noindent Terra vestra desérta, civitátes vestræ succénsæ igni: regiónem vestram coram vobis aliéni dévorant, et desolábitur sicut in vastitáte hostíli. Et derelinquétur fília Sion ut umbráculum in vínea, et sicut tugúrium in cucumerário, et sicut cívitas quæ vastátur. Nisi Dóminus exercítuum reliquísset nobis semen, quasi Sódoma fuissémus, et quasi Gomórrha símiles essémus.\par
\switchcolumn
\EN
\noindent Your land is desolate, your cities are burnt with fire: your country strangers devour before your face, and it shall be desolate as when wasted by enemies. And the daughter of Sion shall be left as a covert in a vineyard, and as a lodge in a garden of cucumbers, and as a city that is laid waste. Except the Lord of hosts had left us seed, we had been as Sodom, and we should have been like to Gomorrha.\par
\switchcolumn*

\LA
\begin{Resp}\Rsign Missus est Gábriel Angelus ad Maríam Vírginem desponsátam Joseph, núntians ei verbum; et expavéscit Virgo de lúmine: ne tímeas, María, invenísti grátiam apud Dóminum: \Star{} Ecce concípies, et páries, et vocábitur Altíssimi Fílius.\end{Resp}
\begin{Resp}\Vsign Dabit ei Dóminus Deus sedem David, patris ejus, et regnábit in domo Jacob in ætérnum.\end{Resp}
\begin{Resp}\Rsign Ecce concípies, et páries, et vocábitur Altíssimi Fílius.\end{Resp}
\begin{Resp}\Vsign Glória Patri, et Fílio, \Star{} et Spirítui Sancto.\end{Resp}
\begin{Resp}\Rsign Ecce concípies, et páries, et vocábitur Altíssimi Fílius.\end{Resp}
\switchcolumn
\EN
\begin{Resp}\Rsign The Angel Gabriel was sent to Mary, a Virgin espoused to Joseph, to bring unto her the word of the Lord, and when the Virgin saw the light, she was afraid. Fear not, Mary, for thou hast found grace from the Lord. \Star{} Behold, thou shalt conceive and bring forth a son, and He shall be called the Son of the Highest.\end{Resp}
\begin{Resp}\Vsign The Lord God shall give unto Him the throne of His father David, and He shall reign over the house of Jacob for ever.\end{Resp}
\begin{Resp}\Rsign Behold, thou shalt conceive, and bring forth a son, and He shall be called the Son of the Highest.\end{Resp}
\begin{Resp}\Vsign Glory be to the Father, and to the Son, \Star{} and to the Holy Ghost.\end{Resp}
\begin{Resp}\Rsign Behold, thou shalt conceive and bring forth a son, and He shall be called the Son of the Highest.\end{Resp}
\switchcolumn*

\LA
\LessonHead{Lectio IV}
\switchcolumn
\EN
\LessonHead{Lesson IV}
\switchcolumn*

\LA
\LessonTitle{Sermo sancti Leónis Papæ}
\switchcolumn
\EN
\LessonTitle{From the Sermons of Pope St. Leo (the Great,)}
\switchcolumn*

\LA
\Source{Sermo 8 de jejunio decimi mensis et eleemosynis}
\switchcolumn
\EN
\Source{8th on the December Fast, and almsgiving.}
\switchcolumn*

\LA
\noindent Cum de advéntu regni Dei, et de mundi fine ac témporum, discípulos suos Salvátor instrúeret, totámque Ecclésiam suam in Apóstolis erudíret: Cavéte, inquit, ne forte gravéntur corda vestra in crápula, et ebrietáte, et cogitatiónibus sæculáribus. Quod útique præcéptum, dilectíssimi, ad nos speciálius pertinére cognóscimus, quibus denuntiátus dies, etiámsi est occúltus, non dubitátur esse vicínus.\par
\switchcolumn
\EN
\noindent Our Saviour Himself instructed His disciples concerning the times and seasons of the coming of the Kingdom of God and the end of the world, and He hath given the same teaching to the Church by the mouth of His Apostles. In connection with this subject then, Our Lord biddeth us beware lest we let our hearts grow heavy through excess of meat and drink, and worldly thoughts. Dearly beloved brethren, we know how that this warning applieth particularly to us. We know that that day is coming, and though for a season we know not the very hour, yet this we know, that it is near.\par
\switchcolumn*

\LA
\begin{Resp}\Rsign Ave, María, grátia plena, Dóminus tecum: \Star{} Spíritus Sanctus supervéniet in te, et virtus Altíssimi obumbrábit tibi: quod enim ex te nascétur Sanctum, vocábitur Fílius Dei.\end{Resp}
\begin{Resp}\Vsign Quómodo fiet istud, quóniam virum non cognósco? Et respóndens Angelus, dixit ei.\end{Resp}
\begin{Resp}\Rsign Spíritus Sanctus supervéniet in te, et virtus Altíssimi obumbrábit tibi: quod enim ex te nascétur Sanctum, vocábitur Fílius Dei.\end{Resp}
\switchcolumn
\EN
\begin{Resp}\Rsign Hail, Mary, full of grace; the Lord is with thee; \Star{} The Holy Ghost shall come upon thee, and the power of the Highest shall overshadow thee; therefore, also that Holy Thing Which shall be born of thee shall be called the Son of God.\end{Resp}
\begin{Resp}\Vsign How shall this be, seeing I know not a man? And the Angel answered and said unto her,\end{Resp}
\begin{Resp}\Rsign The Holy Ghost shall come upon thee, and the power of the Highest shall overshadow thee; therefore, also that Holy Thing Which shall be born of thee shall be called the Son of God.\end{Resp}
\switchcolumn*

\LA
\LessonHead{Lectio V}
\switchcolumn
\EN
\LessonHead{Lesson V}
\switchcolumn*

\LA
\noindent Ad cujus advéntum omnem hóminem cónvenit præparári: ne quem aut ventri déditum, aut curis sæculáribus invéniat implicátum. Quotidiáno enim, dilectíssimi, experiménto probátur, potus satietáte áciem mentis obtúndi, et cibórum nimietáte vigórem cordis hebetári; ita ut delectátio edéndi étiam córporum contrária sit salúti, nisi rátio temperántiæ obsístat illécebræ, et quod futúrum est óneri, súbtrahat voluptáti.\par
\switchcolumn
\EN
\noindent Let every man then make himself ready against the coming of the Lord, so that He may not find him making his belly his god, or the world his chief care. Dearly beloved brethren, it is a matter of every day experience that fulness of drink dulleth the keenness of the mind, and that excess of eating unnerveth the strength of the will. The very stomach protesteth that gluttony doth harm to the bodily health, unless temperance get the better of desire, and the thought of the indigestion afterward check the indulgence of the moment.\par
\switchcolumn*

\LA
\begin{Resp}\Rsign Salvatórem exspectámus, Dóminum Jesum Christum, \Star{} Qui reformábit corpus humilitátis nostræ configurátum córpori claritátis suæ.\end{Resp}
\begin{Resp}\Vsign Sóbrie, et juste, et pie vivámus in hoc sǽculo, exspectántes beátam spem, et advéntum glóriæ magni Dei.\end{Resp}
\begin{Resp}\Rsign Qui reformábit corpus humilitátis nostræ configurátum córpori claritátis suæ.\end{Resp}
\switchcolumn
\EN
\begin{Resp}\Rsign We look for the Saviour, the Lord Jesus Christ; \Star{} Who shall change our vile body, that it may be fashioned like unto His glorious Body.\end{Resp}
\begin{Resp}\Vsign We should live soberly, and righteously, and godly in this present world, looking for that blessed hope, and the glorious appearing of the great God.\end{Resp}
\begin{Resp}\Rsign Who shall change our vile body, that it may be fashioned like unto His glorious Body.\end{Resp}
\switchcolumn*

\LA
\LessonHead{Lectio VI}
\switchcolumn
\EN
\LessonHead{Lesson VI}
\switchcolumn*

\LA
\noindent Quamvis enim sine ánima nihil caro desíderet, et inde accípiat sensus, unde sumit et motus: ejúsdem tamen est ánimæ, quædam sibi súbditæ negáre substántiæ, et interióri judício ab inconveniéntibus exterióra frenáre: ut a corpóreis cupiditátibus sǽpius líbera, in aula mentis possit divínæ vacáre sapiéntiæ: ubi omni strépitu terrenárum silénte curárum, in meditatiónibus sanctis, et in delíciis lætétur ætérnis.\par
\switchcolumn
\EN
\noindent The body without the soul hath no desires; its sensibility cometh from the same source as its movements. And it is the duty of a man with a reasonable soul to deny something to his lower nature and to keep back the outer man from things unseemly. Then will his soul, free from fleshly cravings, sit often at leisure in the palace of the mind, dwelling on the wisdom of God. There, when the roar and rattle of earthly cares are stilled, will she feed on holy thoughts and entertain herself with the expectation of the everlasting joy.\par
\switchcolumn*

\LA
\begin{Resp}\Rsign Obsecro, Dómine, mitte quem missúrus es: vide afflictiónem pópuli tui: \Star{} Sicut locútus es, veni, * Et líbera nos.\end{Resp}
\begin{Resp}\Vsign Qui regis Israël, inténde, qui dedúcis velut ovem Joseph, qui sedes super Chérubim.\end{Resp}
\begin{Resp}\Rsign Sicut locútus es, veni.\end{Resp}
\begin{Resp}\Vsign Glória Patri, et Fílio, \Star{} et Spirítui Sancto.\end{Resp}
\begin{Resp}\Rsign Et líbera nos.\end{Resp}
\switchcolumn
\EN
\begin{Resp}\Rsign O my Lord, send I pray thee, Him Whom Thou wilt send; see the affliction of thy people. As Thou hast promised, come; \Star{} And deliver us.\end{Resp}
\begin{Resp}\Vsign Give ear, O Shepherd of Israel, Thou That leadest Joseph like a flock, Thou That sittest upon the Cherubim!\end{Resp}
\begin{Resp}\Rsign As Thou hast promised, come.\end{Resp}
\begin{Resp}\Vsign Glory be to the Father, and to the Son, \Star{} and to the Holy Ghost.\end{Resp}
\begin{Resp}\Rsign And deliver us.\end{Resp}
\switchcolumn*

\LA
\LessonHead{Lectio VII}
\switchcolumn
\EN
\LessonHead{Lesson VII}
\switchcolumn*

\LA
\LessonTitle{Léctio sancti Evangélii secúndum Lucam}
\switchcolumn
\EN
\LessonTitle{From the Holy Gospel according to Luke}
\switchcolumn*

\LA
\Cite{Luc 21:25-33}
\switchcolumn
\EN
\Cite{Luke 21:25-33}
\switchcolumn*

\LA
\noindent In illo témpore: Dixit Jesus discípulis suis: Erunt signa in sole, et luna, et stellis, et in terris pressúra géntium. Et réliqua.\par
\switchcolumn
\EN
\noindent In that time, Jesus said to his disciples: And there shall be signs in the sun, and in the moon, and in the stars; and upon the earth distress of nations. And so on.\par
\switchcolumn*

\LA
\LessonTitle{Homilía sancti Gregórii Papæ}
\switchcolumn
\EN
\LessonTitle{Homily by Pope St. Gregory (the Great,)}
\switchcolumn*

\LA
\Source{Homilia 1 in Evangelia}
\switchcolumn
\EN
\Source{1st on the Gospels}
\switchcolumn*

\LA
\noindent Dóminus ac Redémptor noster parátos nos inveníre desíderans, senescéntem mundum quæ mala sequántur denúntiat, ut nos ab ejus amóre compéscat. Appropinquántem ejus términum quantæ percussiónes prævéniant, innotéscit: ut, si Deum metúere in tranquillitáte nólumus, saltem vicínum ejus judícium vel percussiónibus attríti timeámus.\par
\switchcolumn
\EN
\noindent Our Lord and Saviour wisheth to find us ready at His second coming. Therefore He telleth us what will be the evils of the world as it groweth old, that He may wean our hearts from worldly affections. Here we read what great convulsions will go before the end, that, if we will not fear God in our prosperity, we may at least be scourged into fearing His judgment when it is at hand.\par
\switchcolumn*

\LA
\begin{Resp}\Rsign Ecce virgo concípiet, et páriet fílium, dicit Dóminus: \Star{} Et vocábitur nomen ejus Admirábilis, Deus, Fortis.\end{Resp}
\begin{Resp}\Vsign Super sólium David, et super regnum ejus sedébit in ætérnum.\end{Resp}
\begin{Resp}\Rsign Et vocábitur nomen ejus Admirábilis, Deus, Fortis.\end{Resp}
\switchcolumn
\EN
\begin{Resp}\Rsign Behold, the Virgin shall conceive, and bear a son, saith the Lord; \Star{} And His name shall be called Wonderful, the Mighty God.\end{Resp}
\begin{Resp}\Vsign He shall sit upon the throne of David, and upon his kingdom for ever.\end{Resp}
\begin{Resp}\Rsign And His name shall be called Wonderful, the Mighty God.\end{Resp}
\switchcolumn*

\LA
\LessonHead{Lectio VIII}
\switchcolumn
\EN
\LessonHead{Lesson VIII}
\switchcolumn*

\LA
\noindent Huic étenim lectióni sancti Evangélii, quam modo vestra fratérnitas audívit, paulo supérius Dóminus præmísit, dicens: Exsúrget gens contra gentem, et regnum advérsus regnum: et erunt terræmótus magni per loca, et pestiléntiæ, et fames. Et quibúsdam interpósitis, hoc, quod modo audístis, adjúnxit: Erunt signa in sole, et luna, et stellis, et in terris pressúra géntium præ confusióne sónitus maris, et flúctuum. Ex quibus profécto ómnibus ália jam facta cérnimus, ália in próximo ventúra formidámus.\par
\switchcolumn
\EN
\noindent Immediately before the passage which hath just been read from the Gospel, are found the following words of our Lord, Nation shall rise against nation, and kingdom against kingdom, and great earthquakes shall be in diverse places, and pestilences and famines. Then, after a few more verses, cometh today's Gospel. There shall be signs in the sun, and in the moon, and in the stars; and upon the earth distress of nations with perplexity, the sea and the waves roaring. Now some of these things are come to pass already, and we fear the others are not far off.\par
\switchcolumn*

\LA
\begin{Resp}\Rsign Audíte verbum Dómini, gentes, et annuntiáte illud in fínibus terræ: \Star{} Et ínsulis, quæ procul sunt, dícite: Salvátor noster advéniet.\end{Resp}
\begin{Resp}\Vsign Annuntiáte, et audítum fácite: loquímini, et clamáte.\end{Resp}
\begin{Resp}\Rsign Et ínsulis, quæ procul sunt, dícite: Salvátor noster advéniet.\end{Resp}
\switchcolumn
\EN
\begin{Resp}\Rsign Hear the word of the Lord, O ye nations, and declare it in the ends of the earth; \Star{} And in the isles afar off, and say: Our Saviour shall come.\end{Resp}
\begin{Resp}\Vsign Declare it and make it known, lift up your voice and cry aloud.\end{Resp}
\begin{Resp}\Rsign And in the isles afar off, and say: Our Saviour shall come.\end{Resp}
\switchcolumn*

\LA
\LessonHead{Lectio IX}
\switchcolumn
\EN
\LessonHead{Lesson IX}
\switchcolumn*

\LA
\noindent Nam gentem contra gentem exsúrgere, earúmque pressúram terris insístere, plus jam in nostris tempóribus cérnimus, quam in codícibus légimus. Quod terræmótus urbes innúmeras óbruat, ex áliis mundi pártibus scitis quam frequénter audívimus. Pestiléntias sine cessatióne pátimur. Signa vero in sole, et luna, et stellis, adhuc apérte mínime vídimus: sed quia et hæc non longe sint, ex ipsa jam áëris immutatióne collígimus.\par
\switchcolumn
\EN
\noindent In these our days we see nation rise against nation, and their distress over all the earth, more than we read in books hath ever come to pass of old time. Ye know also how often we hear of earthquakes overwhelming countless cities in other parts of the world. As for pestilences, we suffer from them ourselves, with hardly any intermission. As yet we do not see signs in the sun, and in the moon, and in the stars; but the changes of seasons and climates warn us that we may look for these also before long.\par
\switchcolumn*

\LA
\begin{Resp}\Rsign Ecce dies véniunt, dicit Dóminus, et suscitábo David germen justum: et regnábit rex, et sápiens erit, et fáciet judícium et justítiam in terra: \Star{} Et hoc est nomen quod vocábunt eum: * Dóminus justus noster.\end{Resp}
\begin{Resp}\Vsign In diébus illis salvábitur Juda, et Israël habitábit confidénter.\end{Resp}
\begin{Resp}\Rsign Et hoc est nomen quod vocábunt eum:\end{Resp}
\begin{Resp}\Vsign Glória Patri, et Fílio, \Star{} et Spirítui Sancto.\end{Resp}
\begin{Resp}\Rsign Dóminus justus noster.\end{Resp}
\switchcolumn
\EN
\begin{Resp}\Rsign Behold, the days come, saith the Lord, that I will raise unto David a righteous Branch; and a King shall reign in wisdom and shall execute judgment and justice in the earth: \Star{} And this is His name whereby He shall be called: * The Lord our Righteous one.\end{Resp}
\begin{Resp}\Vsign In His days Judah shall be saved, and Israel shall dwell safely.\end{Resp}
\begin{Resp}\Rsign And this is His name whereby He shall be called.\end{Resp}
\begin{Resp}\Vsign Glory be to the Father, and to the Son, \Star{} and to the Holy Ghost.\end{Resp}
\begin{Resp}\Rsign The Lord our Righteous one.\end{Resp}
\switchcolumn*

\end{paracol}

\Office{Feria II infra Hebdomadam I Adventus}{Monday of the First Week of Advent}{Feria major}
\addcontentsline{toc}{subsection}{Feria II infra Hebdomadam I Adventus\enspace\textit{Monday of the First Week of Advent}}
\begin{paracol}{2}
\LA
\LessonHead{Lectio I}
\switchcolumn
\EN
\LessonHead{Lesson I}
\switchcolumn*

\LA
\LessonTitle{De Isaía Prophéta}
\switchcolumn
\EN
\LessonTitle{Lesson from the book of Isaias}
\switchcolumn*

\LA
\Cite{Isa 1:16-18}
\switchcolumn
\EN
\Cite{Isa 1:16-18}
\switchcolumn*

\LA
\noindent Lavámini, mundi estóte, auférte malum cogitatiónum vestrárum ab óculis meis: quiéscite ágere pervérse, díscite benefácere: quǽrite judícium, subveníte opprésso, judicáte pupíllo, deféndite víduam. Et veníte, et argúite me, dicit Dóminus: si fúerint peccáta vestra ut cóccinum, quasi nix dealbabúntur: et si fúerint rubra quasi vermículus, velut lana alba erunt.\par
\switchcolumn
\EN
\noindent Wash yourselves, be clean, take away the evil of your devices from my eyes: cease to do perversely, Learn to do well: seek judgment, relieve the oppressed, judge for the fatherless, defend the widow. And then come, and accuse me, saith the Lord: if your sins be as scarlet, they shall be made as white as snow: and if they be red as crimson, they shall be white as wool.\par
\switchcolumn*

\LA
\begin{Resp}\Rsign Súscipe verbum, Virgo María, quod tibi a Dómino per Angelum transmíssum est: concípies et páries Deum páriter et hóminem, \Star{} Ut benedícta dicáris inter omnes mulíeres.\end{Resp}
\begin{Resp}\Vsign Páries quidem fílium, et virginitátis non patiéris detriméntum: efficiéris grávida, et eris mater semper intácta.\end{Resp}
\begin{Resp}\Rsign Ut benedícta dicáris inter omnes mulíeres.\end{Resp}
\switchcolumn
\EN
\begin{Resp}\Rsign Receive, O Virgin Mary, receive the word of the Lord, which is sent thee by His Angel; thou shalt conceive, and shalt bring forth God and Man together. \Star{} And thou shalt be called blessed among all women.\end{Resp}
\begin{Resp}\Vsign Thou shalt bring forth a son, and remain a maiden undefiled thou shalt conceive and be a Mother, still Virgin unspotted.\end{Resp}
\begin{Resp}\Rsign And thou shalt be called blessed among all women.\end{Resp}
\switchcolumn*

\LA
\LessonHead{Lectio II}
\switchcolumn
\EN
\LessonHead{Lesson II}
\switchcolumn*

\LA
\Cite{Isa 1:19-23}
\switchcolumn
\EN
\Cite{Isa 1:19-23}
\switchcolumn*

\LA
\noindent Si voluéritis, et audiéritis me, bona terræ comedétis. Quod si noluéritis, et me ad iracúndiam provocavéritis, gládius devorábit vos, quia os Dómini locútum est. Quómodo facta est méretrix cívitas fidélis, plena judícii? Justítia habitávit in ea, nunc autem homicídæ. Argéntum tuum versum est in scóriam: vinum tuum mixtum est aqua. Príncipes tui infidéles, sócii furum: omnes díligunt múnera, sequúntur retributiónes. Pupíllo non júdicant: et causa víduæ non ingréditur ad illos.\par
\switchcolumn
\EN
\noindent if you be willing, and will hearken to me, you shall eat the good things of the land. But if you will not, and will provoke me to wrath: the sword shall devour you because the mouth of the Lord hath spoken it. How is the faithful city, that was full of judgment, become a harlot? justice dwelt in it, but now murderers. thy silver is turned into dross: thy wine is mingled with water. thy princes are faithless, companions of thieves: they all love bribes, they run after rewards. They judge not for the fatherless: and the widow's cometh not in to them.\par
\switchcolumn*

\LA
\begin{Resp}\Rsign Læténtur cæli, et exsúltet terra, jubiláte, montes, laudem: quia Dóminus noster véniet, \Star{} Et páuperum suórum miserébitur.\end{Resp}
\begin{Resp}\Vsign Oriétur in diébus ejus justítia, et abundántia pacis.\end{Resp}
\begin{Resp}\Rsign Et páuperum suórum miserébitur.\end{Resp}
\switchcolumn
\EN
\begin{Resp}\Rsign Sing, O heavens; and be joyful, O earth; and break forth into singing, O mountains, for our Lord will come; \Star{} And will have mercy on His afflicted.\end{Resp}
\begin{Resp}\Vsign In His days shall righteousness flourish and abundance of peace.\end{Resp}
\begin{Resp}\Rsign And will have mercy upon His afflicted.\end{Resp}
\switchcolumn*

\LA
\LessonHead{Lectio III}
\switchcolumn
\EN
\LessonHead{Lesson III}
\switchcolumn*

\LA
\Cite{Isa 1:24-28}
\switchcolumn
\EN
\Cite{Isa 1:24-28}
\switchcolumn*

\LA
\noindent Propter hoc ait Dóminus Deus exercítuum fortis Israël: Heu! consolábor super hóstibus meis, et vindicábor de inimícis meis. Et convértam manum meam ad te, et éxcoquam ad purum scóriam tuam, et áuferam omne stannum tuum. Et restítuam júdices tuos ut fuérunt prius, et consiliários tuos sicut antíquitus: post hæc vocáberis cívitas justi, urbs fidélis. Sion in judício redimétur, et redúcent eam in justítia: et cónteret sceléstos, et peccatóres simul: et qui dereliquérunt Dóminum, consuméntur.\par
\switchcolumn
\EN
\noindent Therefore saith the Lord the God of hosts, the mighty one of Israel: Ah! I will comfort myself over my adversaries: and I will be revenged of my enemies. And I will turn my hand to thee, and I will clean purge away thy dross, and I will take away all thy tin. And I will restore thy judges as they were before, and thy counsellors as of old. After this thou shalt be called the city of the just, a faithful city. Sion shall be redeemed in judgment, and they shall bring her back in justice. And he shall destroy the wicked, and the sinners together: and they that have forsaken the Lord, shall be consumed.\par
\switchcolumn*

\LA
\begin{Resp}\Rsign Aliéni non transíbunt per Jerúsalem ámplius: \Star{} Nam in illa die stillábunt montes dulcédinem, et colles fluent lac et mel, dicit Dóminus.\end{Resp}
\begin{Resp}\Vsign Deus a Líbano véniet, et Sanctus de monte umbróso et condénso.\end{Resp}
\begin{Resp}\Rsign Nam in illa die stillábunt montes dulcédinem, et colles fluent lac et mel, dicit Dóminus.\end{Resp}
\begin{Resp}\Vsign Glória Patri, et Fílio, \Star{} et Spirítui Sancto.\end{Resp}
\begin{Resp}\Rsign Nam in illa die stillábunt montes dulcédinem, et colles fluent lac et mel, dicit Dóminus.\end{Resp}
\switchcolumn
\EN
\begin{Resp}\Rsign There shall no strangers pass through Jerusalem any more \Star{} For in that day the mountains shall drop down sweet wine, and the hills shall flow with milk and honey, saith the Lord.\end{Resp}
\begin{Resp}\Vsign God shall come from Lebanon, and the Holy One from the thick and shady mountain.\end{Resp}
\begin{Resp}\Rsign For in that day the mountains shall drop down sweet wine, and the hills shall flow with milk and honey, saith the Lord.\end{Resp}
\begin{Resp}\Vsign Glory be to the Father, and to the Son, \Star{} and to the Holy Ghost.\end{Resp}
\begin{Resp}\Rsign For in that day the mountains shall drop down sweet wine, and the hills shall flow with milk and honey, saith the Lord.\end{Resp}
\switchcolumn*

\end{paracol}

\Office{Feria III infra Hebdomadam I Adventus}{Tuesday of the First Week of Advent}{Feria major}
\addcontentsline{toc}{subsection}{Feria III infra Hebdomadam I Adventus\enspace\textit{Tuesday of the First Week of Advent}}
\begin{paracol}{2}
\LA
\LessonHead{Lectio I}
\switchcolumn
\EN
\LessonHead{Lesson I}
\switchcolumn*

\LA
\LessonTitle{De Isaía Prophéta}
\switchcolumn
\EN
\LessonTitle{Lesson from the book of Isaias}
\switchcolumn*

\LA
\Cite{Isa 2:1-3}
\switchcolumn
\EN
\Cite{Isa 2:1-3}
\switchcolumn*

\LA
\noindent Verbum, quod vidit Isaías, fílius Amos, super Judam et Jerúsalem. Et erit in novíssimis diébus præparátus mons domus Dómini in vértice móntium, et elevábitur super colles; et fluent ad eum omnes gentes. Et ibunt pópuli multi, et dicent: Veníte, et ascendámus ad montem Dómini, et ad domum Dei Jacob, et docébit nos vias suas, et ambulábimus in sémitis ejus: quia de Sion exíbit lex, et verbum Dómini de Jerúsalem.\par
\switchcolumn
\EN
\noindent The word that Isaias the son of Amos saw, concerning Juda and Jerusalem. And in the last days the mountain of the house of the Lord shall be prepared on the top of mountains, and it shall be exalted above the hills, and all nations shall flow unto it. And many people shall go, and say: Come and let us go up to the mountain of the Lord, and to the house of the God of Jacob, and he will teach us his ways, and we will walk in his paths: for the law shall come forth from Sion, and the word of the Lord from Jerusalem.\par
\switchcolumn*

\LA
\begin{Resp}\Rsign Montes Israël, ramos vestros expándite, et floréte, et fructus fácite: \Star{} Prope est ut véniat dies Dómini.\end{Resp}
\begin{Resp}\Vsign Roráte, cæli, désuper, et nubes pluant justum: aperiátur terra, et gérminet Salvatórem.\end{Resp}
\begin{Resp}\Rsign Prope est ut véniat dies Dómini.\end{Resp}
\switchcolumn
\EN
\begin{Resp}\Rsign O ye mountains of Israel, shoot forth your branches and blossom and bring forth fruit. \Star{} The day of the Lord is at hand to come.\end{Resp}
\begin{Resp}\Vsign Drop down, ye heavens, from above, and let the skies pour down the Righteous One; let the earth open, and let her bring forth the Saviour.\end{Resp}
\begin{Resp}\Rsign The day of the Lord is at hand to come.\end{Resp}
\switchcolumn*

\LA
\LessonHead{Lectio II}
\switchcolumn
\EN
\LessonHead{Lesson II}
\switchcolumn*

\LA
\Cite{Isa 2:4-6}
\switchcolumn
\EN
\Cite{Isa 2:4-6}
\switchcolumn*

\LA
\noindent Et judicábit gentes, et árguet pópulos multos: et conflábunt gládios suos in vómeres, et lánceas suas in falces: non levábit gens contra gentem gládium, nec exercebúntur ultra ad prǽlium. Domus Jacob, veníte, et ambulémus in lúmine Dómini. Projecísti enim pópulum tuum, domum Jacob: quia repléti sunt ut olim, et áugures habuérunt ut Philísthiim, et púeris aliénis adhæsérunt.\par
\switchcolumn
\EN
\noindent And he shall judge the Gentiles, and rebuke many people: and they shall turn their swords into ploughshares, and their spears into sickles: nation shall not lift up sword against nation, neither shall they be exercised any more to war. O house of Jacob, come ye, and let us walk in the light of the Lord. For thou hast cast off thy people, the house of Jacob: because they are filled as in times past, and have had soothsayers as the Philistines, and have adhered to strange children.\par
\switchcolumn*

\LA
\begin{Resp}\Rsign Erúmpant montes jucunditátem, et colles justítiam: \Star{} Quia lux mundi Dóminus cum poténtia venit.\end{Resp}
\begin{Resp}\Vsign De Sion exíbit lex, et verbum Dómini de Jerúsalem.\end{Resp}
\begin{Resp}\Rsign Quia lux mundi Dóminus cum poténtia venit.\end{Resp}
\switchcolumn
\EN
\begin{Resp}\Rsign Let the mountains break forth into singing, and the hills bring forth righteousness \Star{} For the Lord, the Light of the world, cometh with power.\end{Resp}
\begin{Resp}\Vsign Out of Zion shall go forth the law, and the word of the Lord from Jerusalem.\end{Resp}
\begin{Resp}\Rsign For the Lord, the Light of the world, cometh with power.\end{Resp}
\switchcolumn*

\LA
\LessonHead{Lectio III}
\switchcolumn
\EN
\LessonHead{Lesson III}
\switchcolumn*

\LA
\Cite{Isa 2:7-9}
\switchcolumn
\EN
\Cite{Isa 2:7-9}
\switchcolumn*

\LA
\noindent Repléta est terra argénto et auro: et non est finis thesaurórum ejus: et repléta est terra ejus equis: et innumerábiles quadrígæ ejus. Et repléta est terra ejus idólis: opus mánuum suárum adoravérunt, quod fecérunt dígiti eórum. Et incurvávit se homo, et humiliátus est vir, ne ergo dimíttas eis.\par
\switchcolumn
\EN
\noindent Their land is filled with silver and gold: and there is no end of their treasures. And their land is filled with horses: and their chariots are innumerable. Their land also is full of idols: they have adored the work of their own hands, which their own fingers have made. And man hath bowed himself down, and man hath been debased: therefore forgive them not.\par
\switchcolumn*

\LA
\begin{Resp}\Rsign Ecce ab Austro vénio, ego Dóminus Deus vester, \Star{} Visitáre vos in pace.\end{Resp}
\begin{Resp}\Vsign Aspíciam vos, et créscere fáciam: multiplicabímini, et firmábo pactum meum vobíscum.\end{Resp}
\begin{Resp}\Rsign Visitáre vos in pace.\end{Resp}
\begin{Resp}\Vsign Glória Patri, et Fílio, \Star{} et Spirítui Sancto.\end{Resp}
\begin{Resp}\Rsign Visitáre vos in pace.\end{Resp}
\switchcolumn
\EN
\begin{Resp}\Rsign Behold, I, the Lord your God, come from the South \Star{} To visit you in peace.\end{Resp}
\begin{Resp}\Vsign I will look again upon you and make you to increase ye shall be multiplied, and I will establish My covenant with you.\end{Resp}
\begin{Resp}\Rsign To visit you in peace.\end{Resp}
\begin{Resp}\Vsign Glory be to the Father, and to the Son, \Star{} and to the Holy Ghost.\end{Resp}
\begin{Resp}\Rsign To visit you in peace.\end{Resp}
\switchcolumn*

\end{paracol}

\Office{Feria IV infra Hebdomadam I Adventus}{Wednesday of the First Week of Advent}{Feria major}
\addcontentsline{toc}{subsection}{Feria IV infra Hebdomadam I Adventus\enspace\textit{Wednesday of the First Week of Advent}}
\begin{paracol}{2}
\LA
\LessonHead{Lectio I}
\switchcolumn
\EN
\LessonHead{Lesson I}
\switchcolumn*

\LA
\LessonTitle{De Isaía Prophéta}
\switchcolumn
\EN
\LessonTitle{Lesson from the book of Isaias}
\switchcolumn*

\LA
\Cite{Isa 3:1-4}
\switchcolumn
\EN
\Cite{Isa 3:1-4}
\switchcolumn*

\LA
\noindent Ecce enim Dominátor Dóminus exercítuum áuferet a Jerúsalem et a Juda válidum et fortem, omne robur panis, et omne robur aquæ: fortem, et virum bellatórem, Júdicem, et prophétam, et aríolum, et senem: príncipem super quinquagínta, et honorábilem vultu, et consiliárium, et sapiéntem de architéctis, et prudéntem elóquii mýstici. Et dabo púeros príncipes eórum, et effemináti dominabúntur eis.\par
\switchcolumn
\EN
\noindent For behold, the sovereign the Lord of hosts shall take away from Jerusalem, and from Juda the valiant and the strong, the whole strength of bread, and the whole strength of water. The strong man, and the man of war, the judge, and the prophet, and the cunning man, and the ancient. The captain over fifty, and the honourable in countenance, and the counsellor, and the architect, and the skilful in eloquent speech. And I will give children to be their princes, and the effeminate shall rule over them.\par
\switchcolumn*

\LA
\begin{Resp}\Rsign Rex noster advéniet Christus, \Star{} Quem Joánnes prædicávit Agnum esse ventúrum.\end{Resp}
\begin{Resp}\Vsign Super ipsum continébunt reges os suum, ipsum gentes deprecabúntur.\end{Resp}
\begin{Resp}\Rsign Quem Joánnes prædicávit Agnum esse ventúrum.\end{Resp}
\switchcolumn
\EN
\begin{Resp}\Rsign Christ our King cometh. \Star{} And John hath testified of Him, that He is the Lamb that should come!\end{Resp}
\begin{Resp}\Vsign The kings shall shut their mouths at Him, all nations shall serve Him.\end{Resp}
\begin{Resp}\Rsign And John hath testified of Him, that He is the Lamb that should come!\end{Resp}
\switchcolumn*

\LA
\LessonHead{Lectio II}
\switchcolumn
\EN
\LessonHead{Lesson II}
\switchcolumn*

\LA
\Cite{Isa 3:5-7}
\switchcolumn
\EN
\Cite{Isa 3:5-7}
\switchcolumn*

\LA
\noindent Et írruet pópulus vir ad virum, et unusquísque ad próximum suum: tumultuábitur puer contra senem, et ignóbilis contra nóbilem. Apprehéndet enim vir fratrem suum domésticum patris sui: Vestiméntum tibi est, princeps esto noster, ruína autem hæc sub manu tua. Respondébit in die illa, dicens: Non sum médicus, et in domo mea non est panis, neque vestiméntum: nolíte constitúere me príncipem pópuli.\par
\switchcolumn
\EN
\noindent And the people shall rush one upon another, and every man against his neighbour: the child shall make it tumult against the ancient, and the base against the honourable. For a man shall take hold of his brother, one of the house of his father, saying: Thou hast a garment, be thou our ruler, and let this ruin be under thy hand. In that day he shall answer, saying: I am no healer, and in my house there is no bread, nor clothing: make me not ruler of the people.\par
\switchcolumn*

\LA
\begin{Resp}\Rsign Ante multum tempus prophetávit Ezéchiel: Vidi portam clausam; ecce Deus ante sǽcula ex ea procedébat pro salúte mundi: \Star{} Et erat íterum clausa, demónstrans Vírginem, quia post partum permánsit virgo.\end{Resp}
\begin{Resp}\Vsign Porta quam vidísti, Dóminus solus transíbit per illam.\end{Resp}
\begin{Resp}\Rsign Et erat íterum clausa, demónstrans Vírginem, quia post partum permánsit virgo.\end{Resp}
\switchcolumn
\EN
\begin{Resp}\Rsign Of a long time, said Ezekiel the Prophet, I saw the gate shut behold, God went forth from it before the ages for the salvation of the world. \Star{} And it was shut again, for it is a figure of the Virgin, in that after child-birth she remained a Virgin still.\end{Resp}
\begin{Resp}\Vsign The Lord alone shall enter by the gate that thou savest.\end{Resp}
\begin{Resp}\Rsign And it was shut again, for it is a figure of the Virgin, in that after child-birth she remained a Virgin still.\end{Resp}
\switchcolumn*

\LA
\LessonHead{Lectio III}
\switchcolumn
\EN
\LessonHead{Lesson III}
\switchcolumn*

\LA
\Cite{Isa 3:8-11}
\switchcolumn
\EN
\Cite{Isa 3:8-11}
\switchcolumn*

\LA
\noindent Ruit enim Jerúsalem, et Judas cóncidit: quia lingua eórum et adinventiónes eórum contra Dóminum, ut provocárent óculos majestátis ejus. Agnítio vultus eórum respóndit eis: et peccátum suum quasi Sódoma prædicavérunt, nec abscondérunt: væ ánimæ eórum, quóniam réddita sunt eis mala. Dícite justo quóniam bene, quóniam fructum adinventiónum suárum cómedet. Væ ímpio in malum: retribútio enim mánuum ejus fiet ei.\par
\switchcolumn
\EN
\noindent For Jerusalem is ruined, and Juda is fallen: because their tongue, and their devices are against the Lord, to provoke the eyes of his majesty. The shew of their countenance hath answered them: and they have proclaimed abroad their sin as Sodom, and they have not hid it: woe to their souls, for evils are rendered to them. Say to the just man that it is well, for he shall eat the fruit of his doings. Woe to the wicked unto evil: for the reward of his hands shall be given him.\par
\switchcolumn*

\LA
\begin{Resp}\Rsign Ecce dies véniunt, dicit Dóminus, et suscitábo David germen justum: et regnábit rex, et sápiens erit, et fáciet judícium et justítiam in terra: \Star{} Et hoc est nomen quod vocábunt eum: * Dóminus justus noster.\end{Resp}
\begin{Resp}\Vsign In diébus illis salvábitur Juda, et Israël habitábit confidénter.\end{Resp}
\begin{Resp}\Rsign Et hoc est nomen quod vocábunt eum:\end{Resp}
\begin{Resp}\Vsign Glória Patri, et Fílio, \Star{} et Spirítui Sancto.\end{Resp}
\begin{Resp}\Rsign Dóminus justus noster.\end{Resp}
\switchcolumn
\EN
\begin{Resp}\Rsign Behold, the days come, saith the Lord, that I will raise unto David a righteous Branch; and a King shall reign in wisdom and shall execute judgment and justice in the earth: \Star{} And this is His name whereby He shall be called: * The Lord our Righteous one.\end{Resp}
\begin{Resp}\Vsign In His days Judah shall be saved, and Israel shall dwell safely.\end{Resp}
\begin{Resp}\Rsign And this is His name whereby He shall be called.\end{Resp}
\begin{Resp}\Vsign Glory be to the Father, and to the Son, \Star{} and to the Holy Ghost.\end{Resp}
\begin{Resp}\Rsign The Lord our Righteous one.\end{Resp}
\switchcolumn*

\end{paracol}

\Office{Feria V infra Hebdomadam I Adventus}{Thursday of the First Week of Advent}{Feria major}
\addcontentsline{toc}{subsection}{Feria V infra Hebdomadam I Adventus\enspace\textit{Thursday of the First Week of Advent}}
\begin{paracol}{2}
\LA
\LessonHead{Lectio I}
\switchcolumn
\EN
\LessonHead{Lesson I}
\switchcolumn*

\LA
\LessonTitle{De Isaía Prophéta}
\switchcolumn
\EN
\LessonTitle{Lesson from the book of Isaias}
\switchcolumn*

\LA
\Cite{Isa 4:1-3}
\switchcolumn
\EN
\Cite{Isa 4:1-3}
\switchcolumn*

\LA
\noindent Et apprehéndent septem mulíeres virum unum in die illa, dicéntes: Panem nostrum comedémus, et vestiméntis nostris operiémur: tantúmmodo invocétur nomen tuum super nos, aufer oppróbrium nostrum. In die illa erit germen Dómini in magnificéntia, et glória, et fructus terræ sublímis, et exsultátio his, qui salváti fúerint de Israël. Et erit: Omnis qui relíctus fúerit in Sion, et resíduus in Jerúsalem, sanctus vocábitur, omnis qui scriptus est in vita in Jerúsalem.\par
\switchcolumn
\EN
\noindent And in that day seven women shall take hold of one man, saying: We will eat our own bread, and wear our own apparel: only let us be called by thy name, take away our reproach. In that day the bud of the Lord shall be in magnificence and glory, and the fruit of the earth shall be high, and a great joy to them that shall have escaped of Israel. And it shall come to pass, that every one that shall be left in Sion, and that shall remain in Jerusalem, shall be called holy, every one that is written in life in Jerusalem.\par
\switchcolumn*

\LA
\begin{Resp}\Rsign Súscipe verbum, Virgo María, quod tibi a Dómino per Angelum transmíssum est: concípies et páries Deum páriter et hóminem, \Star{} Ut benedícta dicáris inter omnes mulíeres.\end{Resp}
\begin{Resp}\Vsign Páries quidem fílium, et virginitátis non patiéris detriméntum: efficiéris grávida, et eris mater semper intácta.\end{Resp}
\begin{Resp}\Rsign Ut benedícta dicáris inter omnes mulíeres.\end{Resp}
\switchcolumn
\EN
\begin{Resp}\Rsign Receive, O Virgin Mary, receive the word of the Lord, which is sent thee by His Angel; thou shalt conceive, and shalt bring forth God and Man together. \Star{} And thou shalt be called blessed among all women.\end{Resp}
\begin{Resp}\Vsign Thou shalt bring forth a son, and remain a maiden undefiled thou shalt conceive and be a Mother, still Virgin unspotted.\end{Resp}
\begin{Resp}\Rsign And thou shalt be called blessed among all women.\end{Resp}
\switchcolumn*

\LA
\LessonHead{Lectio II}
\switchcolumn
\EN
\LessonHead{Lesson II}
\switchcolumn*

\LA
\Cite{Isa 5:1-4}
\switchcolumn
\EN
\Cite{Isa 5:1-4}
\switchcolumn*

\LA
\noindent Cantábo dilécto meo cánticum patruélis mei víneæ suæ: Vínea facta est dilécto meo in cornu fílio ólei. Et sepívit eam, et lápides elégit ex illa, et plantávit eam eléctam, et ædificávit turrem in médio ejus, et tórcular exstrúxit in ea: et exspectávit ut fáceret uvas, et fecit labrúscas. Nunc ergo, habitatóres Jerúsalem, et viri Juda, judicáte inter me et víneam meam. Quid est quod débui ultra fácere víneæ meæ, et non feci ei? An quod exspectávi, ut fáceret uvas, et fecit labrúscas?\par
\switchcolumn
\EN
\noindent I will sing to my beloved the canticle of my cousin concerning his vineyard. My beloved had a vineyard on a hill in a fruitful place. And he fenced it in, and picked the stones out of it, and planted it with the choicest vines, and built a tower in the midst thereof, and set up a winepress therein: and he looked that it should bring forth grapes, and it brought forth wild grapes. And now, O ye inhabitants of Jerusalem, and ye men of Juda, judge between me and my vineyard. What is there that I ought to do more to my vineyard, that I have not done to it? was it that I looked that it should bring forth grapes, and it hath brought forth wild grapes?\par
\switchcolumn*

\LA
\begin{Resp}\Rsign Aspiciébam in visu noctis, et ecce in núbibus cæli Fílius hóminis veniébat: et datum est ei regnum, et honor: \Star{} Et omnis pópulus, tribus, et linguæ sérvient ei.\end{Resp}
\begin{Resp}\Vsign Potéstas ejus, potéstas ætérna, quæ non auferétur: et regnum ejus, quod non corrumpétur.\end{Resp}
\begin{Resp}\Rsign Et omnis pópulus, tribus, et linguæ sérvient ei.\end{Resp}
\switchcolumn
\EN
\begin{Resp}\Rsign I saw in the night visions, and, behold, the Son of man came with the clouds of heaven, and there was given Him a Kingdom, and glory; \Star{} And all people, nations, and languages shall serve Him.\end{Resp}
\begin{Resp}\Vsign His dominion is an everlasting dominion which shall not pass away, and His Kingdom that which shall not be destroyed.\end{Resp}
\begin{Resp}\Rsign And all people, nations, and languages shall serve Him.\end{Resp}
\switchcolumn*

\LA
\LessonHead{Lectio III}
\switchcolumn
\EN
\LessonHead{Lesson III}
\switchcolumn*

\LA
\Cite{Isa 5:5-7}
\switchcolumn
\EN
\Cite{Isa 5:5-7}
\switchcolumn*

\LA
\noindent Et nunc osténdam vobis quid ego fáciam víneæ meæ: áuferam sepem ejus, et erit in direptiónem: díruam macériam ejus, et erit in conculcatiónem. Et ponam eam desértam: non putábitur, et non fodiétur. Et ascéndent vepres et spinæ; et núbibus mandábo ne pluant super eam imbrem. Vínea enim Dómini exercítuum domus Israël est: et vir Juda germen ejus delectábile: et exspectávi ut fáceret judícium, et ecce iníquitas: et justítiam, et ecce clamor.\par
\switchcolumn
\EN
\noindent And now I will shew you what I will do to my vineyard. I will take away the hedge thereof, and it shall be wasted: I will break down the wall thereof, and it shall be trodden down. And I will make it desolate: it shall not be pruned, and it shall not be digged: but briers and thorns shall come up: and I will command the clouds to rain no rain upon it. For the vineyard of the Lord of hosts is the house of Israel: and the man of Juda, his pleasant plant: and I looked that he should do judgment, and behold iniquity: and do justice, and behold a cry.\par
\switchcolumn*

\LA
\begin{Resp}\Rsign Missus est Gábriel Angelus ad Maríam Vírginem desponsátam Joseph, núntians ei verbum; et expavéscit Virgo de lúmine: ne tímeas, María, invenísti grátiam apud Dóminum: \Star{} Ecce concípies, et páries, et vocábitur Altíssimi Fílius.\end{Resp}
\begin{Resp}\Vsign Dabit ei Dóminus Deus sedem David, patris ejus, et regnábit in domo Jacob in ætérnum.\end{Resp}
\begin{Resp}\Rsign Ecce concípies, et páries, et vocábitur Altíssimi Fílius.\end{Resp}
\begin{Resp}\Vsign Glória Patri, et Fílio, \Star{} et Spirítui Sancto.\end{Resp}
\begin{Resp}\Rsign Ecce concípies, et páries, et vocábitur Altíssimi Fílius.\end{Resp}
\switchcolumn
\EN
\begin{Resp}\Rsign The Angel Gabriel was sent to Mary, a Virgin espoused to Joseph, to bring unto her the word of the Lord, and when the Virgin saw the light, she was afraid. Fear not, Mary, for thou hast found grace from the Lord. \Star{} Behold, thou shalt conceive and bring forth a son, and He shall be called the Son of the Highest.\end{Resp}
\begin{Resp}\Vsign The Lord God shall give unto Him the throne of His father David, and He shall reign over the house of Jacob for ever.\end{Resp}
\begin{Resp}\Rsign Behold, thou shalt conceive, and bring forth a son, and He shall be called the Son of the Highest.\end{Resp}
\begin{Resp}\Vsign Glory be to the Father, and to the Son, \Star{} and to the Holy Ghost.\end{Resp}
\begin{Resp}\Rsign Behold, thou shalt conceive and bring forth a son, and He shall be called the Son of the Highest.\end{Resp}
\switchcolumn*

\end{paracol}

\Office{Feria VI infra Hebdomadam I Adventus}{Friday of the First Week of Advent}{Feria major}
\addcontentsline{toc}{subsection}{Feria VI infra Hebdomadam I Adventus\enspace\textit{Friday of the First Week of Advent}}
\begin{paracol}{2}
\LA
\LessonHead{Lectio I}
\switchcolumn
\EN
\LessonHead{Lesson I}
\switchcolumn*

\LA
\LessonTitle{De Isaía Prophéta}
\switchcolumn
\EN
\LessonTitle{Lesson from the book of Isaias}
\switchcolumn*

\LA
\Cite{Isa 6:1-3}
\switchcolumn
\EN
\Cite{Isa 6:1-3}
\switchcolumn*

\LA
\noindent In anno, quo mórtuus est rex Ozías, vidi Dóminum sedéntem super sólium excélsum et elevátum: et ea, quæ sub ipso erant, replébant templum. Séraphim stabant super illud: sex alæ uni, et sex alæ álteri, duábus velábant fáciem ejus, et duábus velábant pedes ejus, et duábus volábant. Et clamábant alter ad álterum, et dicébant: Sanctus, sanctus, sanctus, Dóminus exercítuum, plena est omnis terra glória ejus.\par
\switchcolumn
\EN
\noindent In the year that king Ozias died, I saw the Lord sitting upon a throne high and elevated: and his train filled the temple. Upon it stood the seraphims: the one had six wings, and the other had six wings: with two they covered his face, and with two they covered his feet, and with two they flew. And they cried one to another, and said: Holy, holy, holy, the Lord God of hosts, all the earth is full of his glory.\par
\switchcolumn*

\LA
\begin{Resp}\Rsign Ave, María, grátia plena, Dóminus tecum: \Star{} Spíritus Sanctus supervéniet in te, et virtus Altíssimi obumbrábit tibi: quod enim ex te nascétur Sanctum, vocábitur Fílius Dei.\end{Resp}
\begin{Resp}\Vsign Quómodo fiet istud, quóniam virum non cognósco? Et respóndens Angelus, dixit ei.\end{Resp}
\begin{Resp}\Rsign Spíritus Sanctus supervéniet in te, et virtus Altíssimi obumbrábit tibi: quod enim ex te nascétur Sanctum, vocábitur Fílius Dei.\end{Resp}
\switchcolumn
\EN
\begin{Resp}\Rsign Hail, Mary, full of grace; the Lord is with thee; \Star{} The Holy Ghost shall come upon thee, and the power of the Highest shall overshadow thee; therefore, also that Holy Thing Which shall be born of thee shall be called the Son of God.\end{Resp}
\begin{Resp}\Vsign How shall this be, seeing I know not a man? And the Angel answered and said unto her,\end{Resp}
\begin{Resp}\Rsign The Holy Ghost shall come upon thee, and the power of the Highest shall overshadow thee; therefore, also that Holy Thing Which shall be born of thee shall be called the Son of God.\end{Resp}
\switchcolumn*

\LA
\LessonHead{Lectio II}
\switchcolumn
\EN
\LessonHead{Lesson II}
\switchcolumn*

\LA
\Cite{Isa 6:4-7}
\switchcolumn
\EN
\Cite{Isa 6:4-7}
\switchcolumn*

\LA
\noindent Et commóta sunt superliminária cárdinum a voce clamántis, et domus repléta est fumo. Et dixi: Væ mihi, quia tácui, quia vir pollútus lábiis ego sum, et in médio pópuli pollúta lábia habéntis ego hábito, et Regem Dóminum exercítuum vidi óculis meis. Et volávit ad me unus de Séraphim, et in manu ejus cálculus, quem fórcipe túlerat de altári. Et tétigit os meum, et dixit: Ecce tétigit hoc lábia tua, et auferétur iníquitas tua, et peccátum tuum mundábitur.\par
\switchcolumn
\EN
\noindent And the lintels of the doors were moved at the voice of him that cried, and the house was filled with smoke. And I said: Woe is me, because I have held my peace; because I am a man of unclean lips, and I dwell in the midst of a people that hath unclean lips, and I have seen with my eyes the King the Lord of hosts. And one of the seraphims flew to me, and in his hand was a live coal, which he had taken with the tongs off the altar. And he touched my mouth, and said: Behold this hath touched thy lips, and thy iniquities shall be taken away, and thy sin shall be cleansed.\par
\switchcolumn*

\LA
\begin{Resp}\Rsign Salvatórem exspectámus, Dóminum Jesum Christum, \Star{} Qui reformábit corpus humilitátis nostræ configurátum córpori claritátis suæ.\end{Resp}
\begin{Resp}\Vsign Sóbrie, et juste, et pie vivámus in hoc sǽculo, exspectántes beátam spem, et advéntum glóriæ magni Dei.\end{Resp}
\begin{Resp}\Rsign Qui reformábit corpus humilitátis nostræ configurátum córpori claritátis suæ.\end{Resp}
\switchcolumn
\EN
\begin{Resp}\Rsign We look for the Saviour, the Lord Jesus Christ; \Star{} Who shall change our vile body, that it may be fashioned like unto His glorious Body.\end{Resp}
\begin{Resp}\Vsign We should live soberly, and righteously, and godly in this present world, looking for that blessed hope, and the glorious appearing of the great God.\end{Resp}
\begin{Resp}\Rsign Who shall change our vile body, that it may be fashioned like unto His glorious Body.\end{Resp}
\switchcolumn*

\LA
\LessonHead{Lectio III}
\switchcolumn
\EN
\LessonHead{Lesson III}
\switchcolumn*

\LA
\Cite{Isa 6:8-10}
\switchcolumn
\EN
\Cite{Isa 6:8-10}
\switchcolumn*

\LA
\noindent Et audívi vocem Dómini dicéntis: Quem mittam? et quis ibit nobis? Et dixi: Ecce ego, mitte me. Et dixit: Vade, et dices pópulo huic: Audíte audiéntes, et nolíte intellégere: et vidéte visiónem, et nolíte cognóscere. Excǽca cor pópuli hujus, et aures ejus ággrava, et óculos ejus claude, ne forte vídeat óculis suis, et áuribus suis áudiat, et corde suo intéllegat, et convertátur, et sanem eum.\par
\switchcolumn
\EN
\noindent And I heard the voice of the Lord, saying: Whom shall I send? and who shall go for us? And I said: Lo, here am I, send me. And he said: Go, and thou shalt say to this people: Hearing, hear, and understand not: and see the vision, and know it not. Blind the heart of this people, and make their ears heavy, and shut their eyes: lest they see with their eyes, and hear with their ears, and understand with their heart, and be converted and I heal them.\par
\switchcolumn*

\LA
\begin{Resp}\Rsign Obsecro, Dómine, mitte quem missúrus es: vide afflictiónem pópuli tui: \Star{} Sicut locútus es, veni, * Et líbera nos.\end{Resp}
\begin{Resp}\Vsign Qui regis Israël, inténde, qui dedúcis velut ovem Joseph, qui sedes super Chérubim.\end{Resp}
\begin{Resp}\Rsign Sicut locútus es, veni.\end{Resp}
\begin{Resp}\Vsign Glória Patri, et Fílio, \Star{} et Spirítui Sancto.\end{Resp}
\begin{Resp}\Rsign Et líbera nos.\end{Resp}
\switchcolumn
\EN
\begin{Resp}\Rsign O my Lord, send I pray thee, Him Whom Thou wilt send; see the affliction of thy people. As Thou hast promised, come; \Star{} And deliver us.\end{Resp}
\begin{Resp}\Vsign Give ear, O Shepherd of Israel, Thou That leadest Joseph like a flock, Thou That sittest upon the Cherubim!\end{Resp}
\begin{Resp}\Rsign As Thou hast promised, come.\end{Resp}
\begin{Resp}\Vsign Glory be to the Father, and to the Son, \Star{} and to the Holy Ghost.\end{Resp}
\begin{Resp}\Rsign And deliver us.\end{Resp}
\switchcolumn*

\end{paracol}

\Office{Sabbato infra Hebdomadam I Adventus}{Saturday of the First Week of Advent}{Feria major}
\addcontentsline{toc}{subsection}{Sabbato infra Hebdomadam I Adventus\enspace\textit{Saturday of the First Week of Advent}}
\begin{paracol}{2}
\LA
\LessonHead{Lectio I}
\switchcolumn
\EN
\LessonHead{Lesson I}
\switchcolumn*

\LA
\LessonTitle{De Isaía Prophéta}
\switchcolumn
\EN
\LessonTitle{Lesson from the book of Isaias}
\switchcolumn*

\LA
\Cite{Isa 7:1-3}
\switchcolumn
\EN
\Cite{Isa 7:1-3}
\switchcolumn*

\LA
\noindent Et factum est in diébus Achaz, fílii Jóatham, fílii Ozíæ, regis Juda, ascéndit Rasin, rex Sýriæ, et Phácee, fílius Romelíæ, rex Israël, in Jerúsalem, ad præliándum contra eam: et non potuérunt debelláre eam. Et nuntiavérunt dómui David, dicéntes: Requiévit Sýria super Ephraim, et commótum est cor ejus, et cor pópuli ejus, sicut movéntur ligna silvárum a fácie venti. Et dixit Dóminus ad Isaíam: Egrédere in occúrsum Achaz tu, et qui derelíctus est Jasub fílius tuus, ad extrémum aquædúctus piscínæ superióris in via Agri fullónis.\par
\switchcolumn
\EN
\noindent And it came to pass in the days of Achaz the son of Joathan, the son of Ozias, king of Juda, that Basin king of Syria, and Phacee the son of Romelia king of Israel, came up to Jerusalem, to fight against it: but they could not prevail over it. And they told the house of David, saying: Syria hath rested upon Ephraim, and his heart was moved, and the heart of his people, as the trees of the woods are moved with the wind. And the Lord said to Isaias: Go forth to meet Achaz, thou and Jasub thy son that is left, to the conduit in the upper pool, in the way of the fuller's field.\par
\switchcolumn*

\LA
\begin{Resp}\Rsign Ecce virgo concípiet, et páriet fílium, dicit Dóminus: \Star{} Et vocábitur nomen ejus Admirábilis, Deus, Fortis.\end{Resp}
\begin{Resp}\Vsign Super sólium David, et super regnum ejus sedébit in ætérnum.\end{Resp}
\begin{Resp}\Rsign Et vocábitur nomen ejus Admirábilis, Deus, Fortis.\end{Resp}
\switchcolumn
\EN
\begin{Resp}\Rsign Behold, the Virgin shall conceive, and bear a son, saith the Lord; \Star{} And His name shall be called Wonderful, the Mighty God.\end{Resp}
\begin{Resp}\Vsign He shall sit upon the throne of David, and upon his kingdom for ever.\end{Resp}
\begin{Resp}\Rsign And His name shall be called Wonderful, the Mighty God.\end{Resp}
\switchcolumn*

\LA
\LessonHead{Lectio II}
\switchcolumn
\EN
\LessonHead{Lesson II}
\switchcolumn*

\LA
\Cite{Isa 7:4-6}
\switchcolumn
\EN
\Cite{Isa 7:4-6}
\switchcolumn*

\LA
\noindent Et dices ad eum: Vide ut síleas: noli timére, et cor tuum ne formídet a duábus caudis titiónum fumigántium istórum, in ira furóris Rasin, regis Sýriæ, et fílii Romelíæ: eo quod consílium iníerit contra te Sýria in malum Ephraim, et fílius Romelíæ, dicéntes: Ascendámus ad Judam, et suscitémus eum, et avellámus eum ad nos, et ponámus regem in médio ejus fílium Tábeel.\par
\switchcolumn
\EN
\noindent And thou shalt say to him: See thou be quiet: fear not, and let not thy heart be afraid of the two tails of these fire brands, smoking with the wrath of the fury of Rasin king of Syria, and of the son of Romelia. Because Syria hath taken counsel against thee, unto the evil of Ephraim and the son of Romelia, saying: Let us go up to Juda, and rouse it up, and draw it away to us, and make the son of Tabeel king in the midst thereof.\par
\switchcolumn*

\LA
\begin{Resp}\Rsign Audíte verbum Dómini, gentes, et annuntiáte illud in fínibus terræ: \Star{} Et ínsulis, quæ procul sunt, dícite: Salvátor noster advéniet.\end{Resp}
\begin{Resp}\Vsign Annuntiáte, et audítum fácite: loquímini, et clamáte.\end{Resp}
\begin{Resp}\Rsign Et ínsulis, quæ procul sunt, dícite: Salvátor noster advéniet.\end{Resp}
\switchcolumn
\EN
\begin{Resp}\Rsign Hear the word of the Lord, O ye nations, and declare it in the ends of the earth; \Star{} And in the isles afar off, and say: Our Saviour shall come.\end{Resp}
\begin{Resp}\Vsign Declare it and make it known, lift up your voice and cry aloud.\end{Resp}
\begin{Resp}\Rsign And in the isles afar off, and say: Our Saviour shall come.\end{Resp}
\switchcolumn*

\LA
\LessonHead{Lectio III}
\switchcolumn
\EN
\LessonHead{Lesson III}
\switchcolumn*

\LA
\Cite{Isa 7:10-15}
\switchcolumn
\EN
\Cite{Isa 7:10-15}
\switchcolumn*

\LA
\noindent Et adjécit Dóminus loqui ad Achaz, dicens: Pete tibi signum a Dómino Deo tuo in profúndum inférni, sive in excélsum supra. Et dixit Achaz: Non petam, et non tentábo Dóminum. Et dixit: Audíte ergo, domus David: Numquid parum vobis est, moléstos esse homínibus, quia molésti estis et Deo meo? Propter hoc dabit Dóminus ipse vobis signum. Ecce virgo concípiet, et páriet fílium, et vocábitur nomen ejus Emmánuel. Butýrum et mel cómedet, ut sciat reprobáre malum, et elígere bonum.\par
\switchcolumn
\EN
\noindent And the Lord spoke again to Achaz, saying: Ask thee a sign of the Lord thy God either unto the depth of hell, or unto the height above. And Achaz said: I will not ask, and I will not tempt the Lord. And he said: Hear ye therefore, O house of David: Is it a small thing for you to be grievous to men, that you are grievous to my God also? Therefore the Lord himself shall give you a sign. Behold a virgin shall conceive, and bear a son, and his name shall be called Emmanuel. He shall eat butter and honey, that he may know to refuse the evil, and to choose the good.\par
\switchcolumn*

\LA
\begin{Resp}\Rsign Ecce dies véniunt, dicit Dóminus, et suscitábo David germen justum: et regnábit rex, et sápiens erit, et fáciet judícium et justítiam in terra: \Star{} Et hoc est nomen quod vocábunt eum: * Dóminus justus noster.\end{Resp}
\begin{Resp}\Vsign In diébus illis salvábitur Juda, et Israël habitábit confidénter.\end{Resp}
\begin{Resp}\Rsign Et hoc est nomen quod vocábunt eum:\end{Resp}
\begin{Resp}\Vsign Glória Patri, et Fílio, \Star{} et Spirítui Sancto.\end{Resp}
\begin{Resp}\Rsign Dóminus justus noster.\end{Resp}
\switchcolumn
\EN
\begin{Resp}\Rsign Behold, the days come, saith the Lord, that I will raise unto David a righteous Branch; and a King shall reign in wisdom and shall execute judgment and justice in the earth: \Star{} And this is His name whereby He shall be called: * The Lord our Righteous one.\end{Resp}
\begin{Resp}\Vsign In His days Judah shall be saved, and Israel shall dwell safely.\end{Resp}
\begin{Resp}\Rsign And this is His name whereby He shall be called.\end{Resp}
\begin{Resp}\Vsign Glory be to the Father, and to the Son, \Star{} and to the Holy Ghost.\end{Resp}
\begin{Resp}\Rsign The Lord our Righteous one.\end{Resp}
\switchcolumn*

\end{paracol}

\Office{Dominica II Adventus}{Second Sunday of Advent}{Semiduplex II classis}
\addcontentsline{toc}{subsection}{Dominica II Adventus\enspace\textit{Second Sunday of Advent}}
\begin{paracol}{2}
\LA
\LessonHead{Lectio I}
\switchcolumn
\EN
\LessonHead{Lesson I}
\switchcolumn*

\LA
\LessonTitle{De Isaía Prophéta}
\switchcolumn
\EN
\LessonTitle{Lesson from the book of Isaias}
\switchcolumn*

\LA
\Cite{Isa 11:1-4}
\switchcolumn
\EN
\Cite{Isa 11:1-4}
\switchcolumn*

\LA
\noindent Et egrediétur virga de radíce Jesse, et flos de radíce ejus ascéndet. Et requiéscet super eum Spíritus Dómini: spíritus sapiéntiæ et intelléctus, spíritus consílii et fortitúdinis, spíritus sciéntiæ et pietátis; Et replébit eum spíritus timóris Dómini: non secúndum visiónem oculórum judicábit, neque secúndum audítum áurium árguet: Sed judicábit in justítia páuperes, et árguet in æquitáte pro mansuétis terræ:\par
\switchcolumn
\EN
\noindent And there shall come forth a rod out of the root of Jesse, and a flower shall rise up out of his root. And the spirit of the Lord shall rest upon him: the spirit of wisdom, and of understanding, the spirit of counsel, and of fortitude, the spirit of knowledge, and of godliness. And he shall be filled with the spirit of the fear of the Lord. He shall not judge according to the sight of the eyes, nor reprove according to the hearing of the ears. But he shall judge the poor with justice, and shall reprove with equity for the meek of the earth.\par
\switchcolumn*

\LA
\begin{Resp}\Rsign Jerúsalem, cito véniet salus tua: Quare mœróre consúmeris? numquid consiliárius non est tibi, quia innovávit te dolor? \Star{} Salvábo te, et liberábo te, noli timére.\end{Resp}
\begin{Resp}\Vsign Ego enim sum Dóminus, Deus tuus, Sanctus Israël, Redémptor tuus.\end{Resp}
\begin{Resp}\Rsign Salvábo te, et liberábo te, noli timére.\end{Resp}
\switchcolumn
\EN
\begin{Resp}\Rsign Thy salvation cometh quickly, O Jerusalem; why art thou wasted with sorrow? Is there no counselor in thee, that pangs have taken thee? \Star{} Fear not, for I will save thee and deliver thee.\end{Resp}
\begin{Resp}\Vsign For I am the Lord, thy God, the Holy One of Israel, thy Saviour.\end{Resp}
\begin{Resp}\Rsign Fear not, for I will save thee, and deliver thee.\end{Resp}
\switchcolumn*

\LA
\LessonHead{Lectio II}
\switchcolumn
\EN
\LessonHead{Lesson II}
\switchcolumn*

\LA
\Cite{Isa 11:4-7}
\switchcolumn
\EN
\Cite{Isa 11:4-7}
\switchcolumn*

\LA
\noindent Et percútiet terram virga oris sui, et spíritu labiórum suórum interfíciet ímpium. Et erit justítia cíngulum lumbórum ejus: et fides cinctórium renis ejus. Habitábit lupus cum agno, et pardus cum hædo accubábit, vitúlus et leo et ovis simul morabúntur, et puer párvulus minábit eos. Vítulus et ursus pascéntur, simul requiéscent cátuli eórum, et leo quasi bos cómedet páleas.\par
\switchcolumn
\EN
\noindent And he shall strike the earth with the rod of his mouth, and with the breath of his lips he shall slay the wicked. And justice shall be the girdle of his loins: and faith the girdle of his reins. The wolf shall dwell with the lamb: and the leopard shall lie down with the kid: the calf and the lion, and the sheep shall abide together, and a little child shall lead them. The calf and the bear shall feed: their young ones shall rest together: and the lion shall eat straw like the ox.\par
\switchcolumn*

\LA
\begin{Resp}\Rsign Ecce Dóminus véniet, et omnes Sancti ejus cum eo, et erit in die illa lux magna: et exíbunt de Jerúsalem sicut aqua munda: et regnábit Dóminus in ætérnum \Star{} Super omnes gentes.\end{Resp}
\begin{Resp}\Vsign Ecce Dóminus cum virtúte véniet: et regnum in manu ejus, et potéstas, et impérium.\end{Resp}
\begin{Resp}\Rsign Super omnes gentes.\end{Resp}
\switchcolumn
\EN
\begin{Resp}\Rsign Behold, the Lord shall come, and all His saints with Him, and it shall come to pass in that day that the light shall be great; and they shall go out from Jerusalem like clean water; and the Lord shall be King for ever, \Star{} Over all the earth.\end{Resp}
\begin{Resp}\Vsign Behold, the Lord cometh with an host, and in His hand are the kingdom, and power, and dominion.\end{Resp}
\begin{Resp}\Rsign Over all the earth.\end{Resp}
\switchcolumn*

\LA
\LessonHead{Lectio III}
\switchcolumn
\EN
\LessonHead{Lesson III}
\switchcolumn*

\LA
\Cite{Isa 11:8-10}
\switchcolumn
\EN
\Cite{Isa 11:8-10}
\switchcolumn*

\LA
\noindent Et delectábitur infans ab úbere super forámine áspidis, et in cavérna réguli, qui ablactátus fúerit, manum suam mittet. Non nocébunt, et non occídent in univérso monte sancto meo: quia repléta est terra sciéntia Dómini, sicut aquæ maris operiéntes. In die illa radix Jesse, qui stat in signum populórum, ipsum gentes deprecabúntur, et erit sepúlcrum ejus gloriósum.\par
\switchcolumn
\EN
\noindent And the sucking child shall play on the hole of the asp: and the weaned child shall thrust his hand into the den of the basilisk. They shall not hurt, nor shall they kill in all my holy mountain, for the earth is filled with the knowledge of the Lord, as the covering waters of the sea. In that day the root of Jesse, who standeth for an ensign of the people, him the Gentiles shall beseech, and his sepulchre shall be glorious.\par
\switchcolumn*

\LA
\begin{Resp}\Rsign Cívitas Jerúsalem, noli flere: quóniam dóluit Dóminus super te: \Star{} Et áuferet a te omnem tribulatiónem.\end{Resp}
\begin{Resp}\Vsign Ecce Dóminus in fortitúdine véniet: et brácchium ejus dominábitur.\end{Resp}
\begin{Resp}\Rsign Et áuferet a te omnem tribulatiónem.\end{Resp}
\begin{Resp}\Vsign Glória Patri, et Fílio, \Star{} et Spirítui Sancto.\end{Resp}
\begin{Resp}\Rsign Et áuferet a te omnem tribulatiónem.\end{Resp}
\switchcolumn
\EN
\begin{Resp}\Rsign O, thou city of Jerusalem, weep not, for the Lord hath repented Him concerning thee. \Star{} And He will take away from thee all distress.\end{Resp}
\begin{Resp}\Vsign Behold, the Lord shall come with might, and His arm shall rule.\end{Resp}
\begin{Resp}\Rsign And He will take away from thee all distress.\end{Resp}
\begin{Resp}\Vsign Glory be to the Father, and to the Son, \Star{} and to the Holy Ghost.\end{Resp}
\begin{Resp}\Rsign And He will take away from thee all distress.\end{Resp}
\switchcolumn*

\LA
\LessonHead{Lectio IV}
\switchcolumn
\EN
\LessonHead{Lesson IV}
\switchcolumn*

\LA
\LessonTitle{De Expositióne sancti Hierónymi Presbýteri in Isaíam Prophétam}
\switchcolumn
\EN
\LessonTitle{From the Commentary on the Prophecies of Isaiah made by St. Jerome, Priest (at Bethlehem)}
\switchcolumn*

\LA
\Source{Liber 4. in cap. 11 Isaíæ.}
\switchcolumn
\EN
\Source{Book iv, c. xi}
\switchcolumn*

\LA
\noindent Et egrediétur virga de radíce Jesse. Usque ad principium visiónis, vel pónderis Babylónis, quod vidit Isaías, fílius Amos, omnis hæc prophetía de Christo est: quam per partes vólumus explanare, ne simul propósita atque disserta lectóris confundat memóriam. Virgam et florem de radíce Jesse ipsum Dóminum Judæi interpretántur: quod scilicet in virga regnántis poténtia, in flore pulchritúdo monstrétur.\par
\switchcolumn
\EN
\noindent And there shall come forth a rod out of the stem of Jesse. From the beginning of the Book of this Prophet till the xiiith chapter, where commenceth the vision, or burden of Babylon, the whole of the vision of Isaiah, the son of Amoz, is one continual prophecy of Christ. We must explain it part by part, for if we were to take it all at once, the memory of the reader would be confused. According to the Jewish commentators, the rod and the flower would both relate to the Lord Himself. They take the rod to mean the sceptre of His Royal dominion, and the flower the loveliness of His beauty.\par
\switchcolumn*

\LA
\begin{Resp}\Rsign Ecce véniet Dóminus, protéctor noster, Sanctus Israël, \Star{} Corónam regni habens in cápite suo.\end{Resp}
\begin{Resp}\Vsign Et dominábitur a mari usque ad mare, et a flúmine usque ad términos orbis terrárum.\end{Resp}
\begin{Resp}\Rsign Corónam regni habens in cápite suo.\end{Resp}
\switchcolumn
\EN
\begin{Resp}\Rsign Behold, there cometh the Lord, our defender, the Holy One of Israel, \Star{} Wearing a royal crown upon His head.\end{Resp}
\begin{Resp}\Vsign And His dominion shall be from sea even to sea, and from the river even to the ends of the earth.\end{Resp}
\begin{Resp}\Rsign Wearing a royal crown upon His head.\end{Resp}
\switchcolumn*

\LA
\LessonHead{Lectio V}
\switchcolumn
\EN
\LessonHead{Lesson V}
\switchcolumn*

\LA
\noindent Nos autem virgam de radíce Jesse sanctam Mariam Vírginem intelligámus, quæ nullum hábuit sibi frúticem cohæréntem, de qua et supra légimus: Ecce virgo concípiet et pariet fílium. Et florem, Dóminum Salvatórem, qui dicit in Cántico canticórum: Ego flos campi, et lílium convallium.\par
\switchcolumn
\EN
\noindent We, however, understand that the rod out of the root of Jesse signifieth the holy Virgin Mary. She was a clean stem that had as yet put forth no shoot; as we have read above Behold, the Virgin shall conceive and bear a son. (Isa. vii. 14.) And the flower we believe to mean the Lord our Redeemer, Who hath elsewhere compared Himself to a flower; I am a flower of the plain, and a lily of the valleys. (Cant. ii. 1.)\par
\switchcolumn*

\LA
\begin{Resp}\Rsign Sicut mater consolátur fílios suos, ita consolábor vos, dicit Dóminus: et de Jerúsalem civitáte quam elégi, véniet vobis auxílium: \Star{} Et vidébitis et gaudébit cor vestrum.\end{Resp}
\begin{Resp}\Vsign Dabo in Sion salútem, et in Jerúsalem glóriam meam.\end{Resp}
\begin{Resp}\Rsign Et vidébitis et gaudébit cor vestrum.\end{Resp}
\switchcolumn
\EN
\begin{Resp}\Rsign As a mother comforteth her children, so will I comfort you, saith the Lord; My help also cometh unto you out of Jerusalem, the city which I have chosen. \Star{} And when ye see this, your heart shall rejoice.\end{Resp}
\begin{Resp}\Vsign I will place salvation in Zion and in Jerusalem My glory.\end{Resp}
\begin{Resp}\Rsign And when ye shall see this, your heart shall rejoice.\end{Resp}
\switchcolumn*

\LA
\LessonHead{Lectio VI}
\switchcolumn
\EN
\LessonHead{Lesson VI}
\switchcolumn*

\LA
\noindent Super hunc ígitur florem, qui de trunco et radíce Jesse per Mariam Vírginem repénte consúrget, requiéscet Spíritus Dómini: quia in ipso complácuit omnem plenitúdinem divinitátis habitáre corporáliter: nequáquam per partes, ut in ceteris Sanctis: sed juxta Evangélium eórum, quod Hebræo sermóne conscríptum legunt Nazaræi: Descéndet super eum omnis fons Spíritus Sancti. Dóminus autem Spíritus est: et ubi Spíritus Dómini, ibi libertas.\par
\switchcolumn
\EN
\noindent The Spirit of the Lord then shall rest upon this flower; this flower which shall come forth from the stem and roots of Jesse by means of the Virgin Mary. And truly the Spirit of the Lord did rest upon our Redeemer. It is written that In Him dwelleth all the fulness of the Godhead bodily. (Col. ii. 9.) The Spirit was not shed on Him by measure, as it is upon the (4 Isa. Ixvi. 13, 14. Isa. xlvi. 13.) Saints. To Him we may apply the words of the Hebrew Gospel used by the Nazarenes; The whole fountain of the Holy Ghost shall be poured forth upon Him The Lord is a spirit, and where the Spirit of the Lord is, there is liberty. (2 Cor. iii. 17.)\par
\switchcolumn*

\LA
\begin{Resp}\Rsign Jerúsalem, plantábis víneam in móntibus tuis: exsultábis, quóniam dies Dómini véniet: surge, Sion, convértere ad Dóminum Deum tuum: gaude et lætáre, Jacob: \Star{} Quia de médio géntium Salvátor tuus véniet.\end{Resp}
\begin{Resp}\Vsign Exsúlta satis, fília Sion: júbila, fília Jerúsalem.\end{Resp}
\begin{Resp}\Rsign Quia de médio géntium Salvátor tuus véniet.\end{Resp}
\begin{Resp}\Vsign Glória Patri, et Fílio, \Star{} et Spirítui Sancto.\end{Resp}
\begin{Resp}\Rsign Quia de médio géntium Salvátor tuus véniet.\end{Resp}
\switchcolumn
\EN
\begin{Resp}\Rsign Thou shalt yet plant vines upon thy mountains, O Jerusalem thou shalt sing for joy, for the day of the Lord cometh; arise, O Zion, and turn unto the Lord thy God; rejoice and be glad, O Jacob. \Star{} For thy Saviour cometh from the midst of the nations.\end{Resp}
\begin{Resp}\Vsign Sing aloud for joy, O daughter of Zion; shout with gladness, O daughter of Jerusalem.\end{Resp}
\begin{Resp}\Rsign For thy Saviour cometh from the midst of the nations.\end{Resp}
\begin{Resp}\Vsign Glory be to the Father, and to the Son, \Star{} and to the Holy Ghost.\end{Resp}
\begin{Resp}\Rsign For thy Saviour cometh from the midst of the nations.\end{Resp}
\switchcolumn*

\LA
\LessonHead{Lectio VII}
\switchcolumn
\EN
\LessonHead{Lesson VII}
\switchcolumn*

\LA
\LessonTitle{Léctio sancti Evangélii secúndum Matthǽum}
\switchcolumn
\EN
\LessonTitle{From the Holy Gospel according to Matthew}
\switchcolumn*

\LA
\Cite{Matt 11:2-10}
\switchcolumn
\EN
\Cite{Matt 11:2-10}
\switchcolumn*

\LA
\noindent In illo témpore: Cum audísset Joánnes in vínculis ópera Christi, mittens duos de discípulis suis, ait illi: Tu es qui ventúrus es, an álium exspectámus? Et réliqua.\par
\switchcolumn
\EN
\noindent In that time, when John had heard in prison the works of Christ: sending two of his disciples he said to him: Art thou he that art to come, or look we for another? And so on.\par
\switchcolumn*

\LA
\LessonTitle{Homilía sancti Gregórii Papæ}
\switchcolumn
\EN
\LessonTitle{Homily by Pope St. Gregory (the Great.)}
\switchcolumn*

\LA
\Source{Homilia 6 in Evangelia, post initium}
\switchcolumn
\EN
\Source{6th Homily on the Gospels}
\switchcolumn*

\LA
\noindent Visis tot signis tantísque virtútibus, non scandalizári quisque pótuit, sed admirári. Sed infidélium mens grave in illo scándalum pértulit, cum eum post tot mirácula moriéntem vidit. Unde et Paulus dicit: Nos autem prædicámus Christum crucifíxum, Judǽis quidem scándalum, géntibus autem stultítiam. Stultum quippe homínibus visum est, ut pro homínibus Auctor vitæ morerétur: et inde contra eum homo scándalum sumpsit, unde ei ámplius débitor fíeri débuit. Nam tanto Deus ab homínibus dígnius honorándus est, quanto pro homínibus et indígna suscépit.\par
\switchcolumn
\EN
\noindent The sight of so many signs and so many mighty works should have been a source of wonder, and not a stumbling-block. And yet the unfaithful (Jer. xxxi. 5.) found these very works a rock of offence, when they afterwards saw Him Who had worked so many miracles dying on the Cross. Hence Paul saith: “We preach Christ crucified, unto the Jews a stumbling-block and unto the Gentiles foolishness.” (1 Cor. i. 23.) It is indeed folly in the eyes of men to say that the Author of life died for men and thus men put as a stumbling-block to hinder them from coming to Jesus, the very thing that doth oblige them the most unto Him. For the more humbling God hath undergone for man's sake, the more worthy is He that man should worship Him.\par
\switchcolumn*

\LA
\begin{Resp}\Rsign Egrediétur Dóminus de Samaría ad portam, quæ réspicit ad Oriéntem: et véniet in Béthlehem ámbulans super aquas redemptiónis Judæ: \Star{} Tunc salvus erit omnis homo: quia ecce véniet.\end{Resp}
\begin{Resp}\Vsign Et præparábitur in misericórdia sólium ejus, et sedébit super illud in veritáte.\end{Resp}
\begin{Resp}\Rsign Tunc salvus erit omnis homo: quia ecce véniet.\end{Resp}
\switchcolumn
\EN
\begin{Resp}\Rsign The Lord shall go forth out of Samaria unto the gate that looketh toward the East; and He shall come into Bethlehem, walking upon the waters of the redemption of Judah. \Star{} Then shall every one be saved for, behold, He cometh.\end{Resp}
\begin{Resp}\Vsign And in mercy shall His throne be established, and He shall sit upon it in truth.\end{Resp}
\begin{Resp}\Rsign Then shall every one be saved for, behold, He cometh.\end{Resp}
\switchcolumn*

\LA
\LessonHead{Lectio VIII}
\switchcolumn
\EN
\LessonHead{Lesson VIII}
\switchcolumn*

\LA
\noindent Quid est ergo dicere: Beátus qui non fúerit scandalizátus in me; nisi apérta voce abjectiónem mortis suæ humilitátemque signare? Ac si patenter dicat: Mira quidem facio, sed abjecta pérpeti non dedignor. Quia ergo moriéndo te súbsequor, cavéndum valde est homínibus, ne in me mortem despíciant, qui signa venerántur.\par
\switchcolumn
\EN
\noindent And blessed is he, whosoever shall not be offended in Me. Now what is this, but a plain mention of that time, when He afterwards humbled Himself, becoming obedient unto death, even the death of the Cross? It is as if He said: I indeed do wonderful works, but the day will come when I shall not refuse to suffer shame and evil treatment. Take heed then, ye who now worship Me for the works' sake, that when I come to die ye despise Me not for My death's sake.\par
\switchcolumn*

\LA
\begin{Resp}\Rsign Festína, ne tardáveris, Dómine: \Star{} Et líbera pópulum tuum.\end{Resp}
\begin{Resp}\Vsign Veni, Dómine, et noli tardáre: reláxa facínora plebi tuæ.\end{Resp}
\begin{Resp}\Rsign Et líbera pópulum tuum.\end{Resp}
\switchcolumn
\EN
\begin{Resp}\Rsign Make haste, O Lord, make no tarrying. \Star{} And deliver thy people.\end{Resp}
\begin{Resp}\Vsign O Lord, come and make no tarrying loose the bonds of thy people.\end{Resp}
\begin{Resp}\Rsign And deliver thy people.\end{Resp}
\switchcolumn*

\LA
\LessonHead{Lectio IX}
\switchcolumn
\EN
\LessonHead{Lesson IX}
\switchcolumn*

\LA
\noindent Sed dimissis Joánnis discípulis, quid de eodem Joánne turbis dicat, audiámus: Quid exístis in desértum vidére? Arúndinem vento agitatam? Quod videlicet non asseréndo , sed negando intulit. Arúndinem quippe mox ut aura contígerit, in partem álteram inflectit. Et quid per arúndinem, nisi carnalis animus designátur? Qui mox ut favore vel detractióne tangitur, statim in partem quámlibet inclinátur?\par
\switchcolumn
\EN
\noindent And, as the disciples of John departed, what did Jesus say unto the multitudes concerning this same John? Let us hear. What went ye out into the wilderness to see? A reed shaken with the wind? Here our Lord teacheth not by assertion, but by negation. Now a reed is a thing so made that as soon as the wind bloweth upon it, it bendeth it over toward the opposite quarter. And the fleshly-minded man is like a human reed. As he is praised or blamed so he bendeth himself in the one direction or the other.\par
\switchcolumn*

\LA
\begin{Resp}\Rsign Ecce Dóminus véniet cum splendóre descéndens, et virtus ejus cum eo, \Star{} Visitáre pópulum suum in pace, et constitúere super eum vitam sempitérnam.\end{Resp}
\begin{Resp}\Vsign Ecce Dóminus noster cum virtúte véniet.\end{Resp}
\begin{Resp}\Rsign Visitáre pópulum suum in pace, et constitúere super eum vitam sempitérnam.\end{Resp}
\begin{Resp}\Vsign Glória Patri, et Fílio, \Star{} et Spirítui Sancto.\end{Resp}
\begin{Resp}\Rsign Visitáre pópulum suum in pace, et constitúere super eum vitam sempitérnam.\end{Resp}
\switchcolumn
\EN
\begin{Resp}\Rsign Behold, the Lord cometh down with glory, and His host is with Him. \Star{} To visit His people in peace, and to establish them in life everlasting.\end{Resp}
\begin{Resp}\Vsign Behold, our Lord cometh with an host.\end{Resp}
\begin{Resp}\Rsign To visit His people in peace, and to establish them in life everlasting.\end{Resp}
\begin{Resp}\Vsign Glory be to the Father, and to the Son, \Star{} and to the Holy Ghost.\end{Resp}
\begin{Resp}\Rsign To visit His people in peace, and to establish them in life everlasting.\end{Resp}
\switchcolumn*

\end{paracol}

\Office{Feria II infra Hebdomadam II Adventus}{Monday of the Second Week of Advent}{Feria major}
\addcontentsline{toc}{subsection}{Feria II infra Hebdomadam II Adventus\enspace\textit{Monday of the Second Week of Advent}}
\begin{paracol}{2}
\LA
\LessonHead{Lectio I}
\switchcolumn
\EN
\LessonHead{Lesson I}
\switchcolumn*

\LA
\LessonTitle{De Isaía Prophéta}
\switchcolumn
\EN
\LessonTitle{Lesson from the book of Isaias}
\switchcolumn*

\LA
\Cite{Isa 13:1-4}
\switchcolumn
\EN
\Cite{Isa 13:1-4}
\switchcolumn*

\LA
\noindent Onus Babylónis, quod vidit Isaías fílius Amos. Super montem caliginósum leváte signum, exaltáte vocem, leváte manum, et ingrediántur portas duces. Ego mandávi sanctificátis meis, et vocávi fortes meos in ira mea, exsultántes in glória mea. Vox multitúdinis in móntibus, quasi populórum frequéntium: vox sónitus regum, géntium congregatárum.\par
\switchcolumn
\EN
\noindent The burden of Babylon, which Isaias the son of Amos saw. Upon the dark mountain lift ye up a banner, exalt the voice, lift up the hand, and let the rulers go into the gates. I have commanded my sanctified ones, and have called my strong ones in my wrath, them that rejoice in my glory. The noise of a multitude in the mountains, as it were of many people, the noise of the sound of kings, of nations gathered together.\par
\switchcolumn*

\LA
\begin{Resp}\Rsign Súscipe verbum, Virgo María, quod tibi a Dómino per Angelum transmíssum est: concípies et páries Deum páriter et hóminem, \Star{} Ut benedícta dicáris inter omnes mulíeres.\end{Resp}
\begin{Resp}\Vsign Páries quidem fílium, et virginitátis non patiéris detriméntum: efficiéris grávida, et eris mater semper intácta.\end{Resp}
\begin{Resp}\Rsign Ut benedícta dicáris inter omnes mulíeres.\end{Resp}
\switchcolumn
\EN
\begin{Resp}\Rsign Receive, O Virgin Mary, receive the word of the Lord, which is sent thee by His Angel; thou shalt conceive, and shalt bring forth God and Man together. \Star{} And thou shalt be called blessed among all women.\end{Resp}
\begin{Resp}\Vsign Thou shalt bring forth a son, and remain a maiden undefiled thou shalt conceive and be a Mother, still Virgin unspotted.\end{Resp}
\begin{Resp}\Rsign And thou shalt be called blessed among all women.\end{Resp}
\switchcolumn*

\LA
\LessonHead{Lectio II}
\switchcolumn
\EN
\LessonHead{Lesson II}
\switchcolumn*

\LA
\Cite{Isa 13:4-8}
\switchcolumn
\EN
\Cite{Isa 13:4-8}
\switchcolumn*

\LA
\noindent Dóminus exercítuum præcépit milítiæ belli, veniéntibus de terra procul a summitáte cæli: Dóminus, et vasa furóris ejus, ut dispérdat omnem terram. Ululáte, quia prope est dies Dómini: quasi vástitas a Dómino véniet. Propter hoc omnes manus dissolvéntur, et omne cor hóminis contabéscet, et conterétur. Torsiónes et dolóres tenébunt; quasi partúriens dolébunt: unusquísque ad próximum suum stupébit, fácies combústæ vultus eórum.\par
\switchcolumn
\EN
\noindent The Lord of hosts hath given charge to the troops of war. To them that come from a country afar off, from the end of heaven: the Lord and the instruments of his wrath, to destroy the whole land. Howl ye, for the day of the Lord is near: it shall come as a destruction from the Lord. Therefore shall all hands be faint, and every heart of man shall melt, And shall be broken. Gripings and pains shall take hold of them, they shall be in pain as a woman in labour. Every one shall be amazed at his neighbour, their countenances shall be as faces burnt.\par
\switchcolumn*

\LA
\begin{Resp}\Rsign Læténtur cæli, et exsúltet terra, jubiláte, montes, laudem: quia Dóminus noster véniet, \Star{} Et páuperum suórum miserébitur.\end{Resp}
\begin{Resp}\Vsign Oriétur in diébus ejus justítia, et abundántia pacis.\end{Resp}
\begin{Resp}\Rsign Et páuperum suórum miserébitur.\end{Resp}
\switchcolumn
\EN
\begin{Resp}\Rsign Sing, O heavens; and be joyful, O earth; and break forth into singing, O mountains, for our Lord will come; \Star{} And will have mercy on His afflicted.\end{Resp}
\begin{Resp}\Vsign In His days shall righteousness flourish and abundance of peace.\end{Resp}
\begin{Resp}\Rsign And will have mercy upon His afflicted.\end{Resp}
\switchcolumn*

\LA
\LessonHead{Lectio III}
\switchcolumn
\EN
\LessonHead{Lesson III}
\switchcolumn*

\LA
\Cite{Isa 13:9-11}
\switchcolumn
\EN
\Cite{Isa 13:9-11}
\switchcolumn*

\LA
\noindent Ecce dies Dómini véniet, crudélis, et indignatiónis plenus, et iræ furorísque, ad ponéndam terram in solitúdinem, et peccatóres ejus conteréndos de ea. Quóniam stellæ cæli, et splendor eárum, non expándent lumen suum: obtenebrátus est sol in ortu suo, et luna non splendébit in lúmine suo. Et visitábo super orbis mala, et contra ímpios iniquitátem eórum, et quiéscere fáciam supérbiam infidélium, et arrogántiam fórtium humiliábo.\par
\switchcolumn
\EN
\noindent Behold, the day of the Lord shall come, a cruel day, and full of indignation, and of wrath, and fury, to lay the land desolate, and to destroy the sinners thereof out of it. For the stars of heaven, and their brightness shall not display their light: the sun shall be darkened in his rising, and the moon shall not shine with her light. And I will visit the evils of the world, and against the wicked for their iniquity: and I will make the pride of infidels to cease, and will bring down the arrogancy of the mighty.\par
\switchcolumn*

\LA
\begin{Resp}\Rsign Aliéni non transíbunt per Jerúsalem ámplius: \Star{} Nam in illa die stillábunt montes dulcédinem, et colles fluent lac et mel, dicit Dóminus.\end{Resp}
\begin{Resp}\Vsign Deus a Líbano véniet, et Sanctus de monte umbróso et condénso.\end{Resp}
\begin{Resp}\Rsign Nam in illa die stillábunt montes dulcédinem, et colles fluent lac et mel, dicit Dóminus.\end{Resp}
\begin{Resp}\Vsign Glória Patri, et Fílio, \Star{} et Spirítui Sancto.\end{Resp}
\begin{Resp}\Rsign Nam in illa die stillábunt montes dulcédinem, et colles fluent lac et mel, dicit Dóminus.\end{Resp}
\switchcolumn
\EN
\begin{Resp}\Rsign There shall no strangers pass through Jerusalem any more \Star{} For in that day the mountains shall drop down sweet wine, and the hills shall flow with milk and honey, saith the Lord.\end{Resp}
\begin{Resp}\Vsign God shall come from Lebanon, and the Holy One from the thick and shady mountain.\end{Resp}
\begin{Resp}\Rsign For in that day the mountains shall drop down sweet wine, and the hills shall flow with milk and honey, saith the Lord.\end{Resp}
\begin{Resp}\Vsign Glory be to the Father, and to the Son, \Star{} and to the Holy Ghost.\end{Resp}
\begin{Resp}\Rsign For in that day the mountains shall drop down sweet wine, and the hills shall flow with milk and honey, saith the Lord.\end{Resp}
\switchcolumn*

\end{paracol}

\Office{Feria III infra Hebdomadam II Adventus}{Tuesday of the Second Week of Advent}{Feria major}
\addcontentsline{toc}{subsection}{Feria III infra Hebdomadam II Adventus\enspace\textit{Tuesday of the Second Week of Advent}}
\begin{paracol}{2}
\LA
\LessonHead{Lectio I}
\switchcolumn
\EN
\LessonHead{Lesson I}
\switchcolumn*

\LA
\LessonTitle{De Isaía Prophéta}
\switchcolumn
\EN
\LessonTitle{Lesson from the book of Isaias}
\switchcolumn*

\LA
\Cite{Isa 14:1-2}
\switchcolumn
\EN
\Cite{Isa 14:1-2}
\switchcolumn*

\LA
\noindent Prope est ut véniat tempus ejus, et dies ejus non elongabúntur. Miserébitur enim Dóminus Jacob, et éliget adhuc de Israël, et requiéscere eos fáciet super humum suam: adjungétur ádvena ad eos, et adhærébit dómui Jacob. Et tenébunt eos pópuli, et addúcent eos in locum suum: et possidébit eos domus Israël super terram Dómini in servos et ancíllas: et erunt capiéntes eos, qui se céperant, et subícient exactóres suos.\par
\switchcolumn
\EN
\noindent Her time is near at hand, and her days shall not be prolonged. For the Lord will have mercy on Jacob, and will yet choose out of Israel, and will make them rest upon their own ground: and the stranger shall be joined with them, and shall adhere to the house of Jacob. And the people shall take them, and bring them into their place: and the house of Israel shall possess them in the land of the Lord for servants and handmaids: and they shall make them captives that had taken them, and shall subdue their oppressors.\par
\switchcolumn*

\LA
\begin{Resp}\Rsign Montes Israël, ramos vestros expándite, et floréte, et fructus fácite: \Star{} Prope est ut véniat dies Dómini.\end{Resp}
\begin{Resp}\Vsign Roráte, cæli, désuper, et nubes pluant justum: aperiátur terra, et gérminet Salvatórem.\end{Resp}
\begin{Resp}\Rsign Prope est ut véniat dies Dómini.\end{Resp}
\switchcolumn
\EN
\begin{Resp}\Rsign O ye mountains of Israel, shoot forth your branches and blossom and bring forth fruit. \Star{} The day of the Lord is at hand to come.\end{Resp}
\begin{Resp}\Vsign Drop down, ye heavens, from above, and let the skies pour down the Righteous One; let the earth open, and let her bring forth the Saviour.\end{Resp}
\begin{Resp}\Rsign The day of the Lord is at hand to come.\end{Resp}
\switchcolumn*

\LA
\LessonHead{Lectio II}
\switchcolumn
\EN
\LessonHead{Lesson II}
\switchcolumn*

\LA
\Cite{Isa 14:3-6}
\switchcolumn
\EN
\Cite{Isa 14:3-6}
\switchcolumn*

\LA
\noindent Et erit in die illa, cum réquiem déderit tibi Deus a labóre tuo, et a concussióne tua, et a servitúte dura, qua ante servísti: sumes parábolam istam contra regem Babylónis, et dices: Quómodo cessávit exáctor, quiévit tribútum? Contrívit Dóminus báculum impiórum, virgam dominántium, cædéntem pópulos in indignatióne, plaga insanábili, subiciéntem in furóre gentes, persequéntem crudéliter.\par
\switchcolumn
\EN
\noindent And it shall come to pass in that day, that when God shall give thee rest from thy labour, and from thy vexation, and from the hard bondage, wherewith thou didst serve before, Thou shalt take up this parable against the king of Babylon, and shalt say: How is the oppressor come to nothing, the tribute hath ceased? The Lord hath broken the staff of the wicked, the rod of the rulers, That struck the people in wrath with an incurable wound, that brought nations under in fury, that persecuted in a cruel manner.\par
\switchcolumn*

\LA
\begin{Resp}\Rsign Erúmpant montes jucunditátem, et colles justítiam: \Star{} Quia lux mundi Dóminus cum poténtia venit.\end{Resp}
\begin{Resp}\Vsign De Sion exíbit lex, et verbum Dómini de Jerúsalem.\end{Resp}
\begin{Resp}\Rsign Quia lux mundi Dóminus cum poténtia venit.\end{Resp}
\switchcolumn
\EN
\begin{Resp}\Rsign Let the mountains break forth into singing, and the hills bring forth righteousness \Star{} For the Lord, the Light of the world, cometh with power.\end{Resp}
\begin{Resp}\Vsign Out of Zion shall go forth the law, and the word of the Lord from Jerusalem.\end{Resp}
\begin{Resp}\Rsign For the Lord, the Light of the world, cometh with power.\end{Resp}
\switchcolumn*

\LA
\LessonHead{Lectio III}
\switchcolumn
\EN
\LessonHead{Lesson III}
\switchcolumn*

\LA
\Cite{Isa 14:12-15}
\switchcolumn
\EN
\Cite{Isa 14:12-15}
\switchcolumn*

\LA
\noindent Quómodo cecidísti de cælo, lúcifer, qui mane oriebáris? corruísti in terram, qui vulnerábas gentes? qui dicébas in corde tuo: In cælum conscéndam, super astra Dei exaltábo sólium meum, sedébo in monte testaménti, in latéribus Aquilónis. Ascéndam super altitúdinem núbium, símilis ero Altíssimo. Verúmtamen ad inférnum detrahéris in profúndum laci.\par
\switchcolumn
\EN
\noindent How art thou fallen from heaven, O Lucifer, who didst rise in the morning? how art thou fallen to the earth, that didst wound the nations? And thou saidst in thy heart: I will ascend into heaven, I will exalt my throne above the stars of God, I will sit in the mountain of the covenant, in the sides of the north. I will ascend above the height of the clouds, I will be like the most High. But yet thou shalt be brought down to hell, into the depth of the pit.\par
\switchcolumn*

\LA
\begin{Resp}\Rsign Ecce ab Austro vénio, ego Dóminus Deus vester, \Star{} Visitáre vos in pace.\end{Resp}
\begin{Resp}\Vsign Aspíciam vos, et créscere fáciam: multiplicabímini, et firmábo pactum meum vobíscum.\end{Resp}
\begin{Resp}\Rsign Visitáre vos in pace.\end{Resp}
\begin{Resp}\Vsign Glória Patri, et Fílio, \Star{} et Spirítui Sancto.\end{Resp}
\begin{Resp}\Rsign Visitáre vos in pace.\end{Resp}
\switchcolumn
\EN
\begin{Resp}\Rsign Behold, I, the Lord your God, come from the South \Star{} To visit you in peace.\end{Resp}
\begin{Resp}\Vsign I will look again upon you and make you to increase ye shall be multiplied, and I will establish My covenant with you.\end{Resp}
\begin{Resp}\Rsign To visit you in peace.\end{Resp}
\begin{Resp}\Vsign Glory be to the Father, and to the Son, \Star{} and to the Holy Ghost.\end{Resp}
\begin{Resp}\Rsign To visit you in peace.\end{Resp}
\switchcolumn*

\end{paracol}

\Office{Feria IV infra Hebdomadam II Adventus}{Wednesday of the Second Week of Advent}{Feria major}
\addcontentsline{toc}{subsection}{Feria IV infra Hebdomadam II Adventus\enspace\textit{Wednesday of the Second Week of Advent}}
\begin{paracol}{2}
\LA
\LessonHead{Lectio I}
\switchcolumn
\EN
\LessonHead{Lesson I}
\switchcolumn*

\LA
\LessonTitle{De Isaía Prophéta}
\switchcolumn
\EN
\LessonTitle{Lesson from the book of Isaias}
\switchcolumn*

\LA
\Cite{Isa 16:1-4}
\switchcolumn
\EN
\Cite{Isa 16:1-4}
\switchcolumn*

\LA
\noindent Emítte Agnum, Dómine, Dominatórem terræ de Petra desérti ad montem fíliæ Sion. Et erit: Sicut avis fúgiens, et pulli de nido avolántes, sic erunt fíliæ Moab in transcénsu Arnon. Ini consílium, coge concílium: pone quasi noctem umbram tuam in merídie: abscónde fugiéntes, et vagos ne prodas. Habitábunt apud te prófugi mei: Moab esto latíbulum eórum a fácie vastatóris.\par
\switchcolumn
\EN
\noindent Send forth, O Lord, the lamb, the ruler of the earth, from Petra of the desert, to the mount of the daughter of Sion. And it shall come to pass, that as a bird fleeing away, and as young ones flying out of the nest, so shall the daughters of Moab be in the passage of Arnon. Take counsel, gather a council: make thy shadow as the night in the midday: hide them that flee, and betray not them that wander about. My fugitives shall dwell with thee: O Moab, be thou a covert to them from the face of the destroyer.\par
\switchcolumn*

\LA
\begin{Resp}\Rsign Rex noster advéniet Christus, \Star{} Quem Joánnes prædicávit Agnum esse ventúrum.\end{Resp}
\begin{Resp}\Vsign Super ipsum continébunt reges os suum, ipsum gentes deprecabúntur.\end{Resp}
\begin{Resp}\Rsign Quem Joánnes prædicávit Agnum esse ventúrum.\end{Resp}
\switchcolumn
\EN
\begin{Resp}\Rsign Christ our King cometh. \Star{} And John hath testified of Him, that He is the Lamb that should come!\end{Resp}
\begin{Resp}\Vsign The kings shall shut their mouths at Him, all nations shall serve Him.\end{Resp}
\begin{Resp}\Rsign And John hath testified of Him, that He is the Lamb that should come!\end{Resp}
\switchcolumn*

\LA
\LessonHead{Lectio II}
\switchcolumn
\EN
\LessonHead{Lesson II}
\switchcolumn*

\LA
\Cite{Isa 16:4-6}
\switchcolumn
\EN
\Cite{Isa 16:4-6}
\switchcolumn*

\LA
\noindent Fínitus est enim pulvis, consummátus est miser, defécit qui conculcábat terram. Et præparábitur in misericórdia sólium, et sedébit super illud in veritáte in tabernáculo David, júdicans et quærens judícium, et velóciter reddens quod justum est. Audívimus supérbiam Moab, supérbus est valde: supérbia ejus et arrogántia ejus, et indignátio ejus, plus quam fortitúdo ejus.\par
\switchcolumn
\EN
\noindent For the dust is at an end, the wretch is consumed: he hath failed, that trod the earth under foot. And a throne shall be prepared in mercy, and one shall sit upon it in truth in the tabernacle of David, judging and seeking judgment and quickly rendering that which is just. We have heard of the pride of Moab, he is exceeding proud: his pride and his arrogancy, and his indignation is more than his strength.\par
\switchcolumn*

\LA
\begin{Resp}\Rsign Ante multum tempus prophetávit Ezéchiel: Vidi portam clausam; ecce Deus ante sǽcula ex ea procedébat pro salúte mundi: \Star{} Et erat íterum clausa, demónstrans Vírginem, quia post partum permánsit virgo.\end{Resp}
\begin{Resp}\Vsign Porta quam vidísti, Dóminus solus transíbit per illam.\end{Resp}
\begin{Resp}\Rsign Et erat íterum clausa, demónstrans Vírginem, quia post partum permánsit virgo.\end{Resp}
\switchcolumn
\EN
\begin{Resp}\Rsign Of a long time, said Ezekiel the Prophet, I saw the gate shut behold, God went forth from it before the ages for the salvation of the world. \Star{} And it was shut again, for it is a figure of the Virgin, in that after child-birth she remained a Virgin still.\end{Resp}
\begin{Resp}\Vsign The Lord alone shall enter by the gate that thou savest.\end{Resp}
\begin{Resp}\Rsign And it was shut again, for it is a figure of the Virgin, in that after child-birth she remained a Virgin still.\end{Resp}
\switchcolumn*

\LA
\LessonHead{Lectio III}
\switchcolumn
\EN
\LessonHead{Lesson III}
\switchcolumn*

\LA
\Cite{Isa 16:7-8}
\switchcolumn
\EN
\Cite{Isa 16:7-8}
\switchcolumn*

\LA
\noindent Idcírco ululábit Moab ad Moab, univérsus ululábit: his, qui lætántur super muros cocti láteris, loquímini plagas suas. Quóniam suburbána Hésebon desérta sunt, et víneam Sábama dómini géntium excidérunt: flagélla ejus usque ad Jazer pervenérunt: erravérunt in desérto, propágines ejus relictæ sunt, transiérunt mare.\par
\switchcolumn
\EN
\noindent Therefore shall Moab howl to Moab, every one shall howl: to them that rejoice upon the brick walls, tell ye their stripes. For the suburbs of Hesebon are desolate, and the lords of the nations have destroyed the vineyard of Sabama: the branches thereof have reached even to Jazer: they have wandered in the wilderness, the branches thereof are left, they are gone over the sea.\par
\switchcolumn*

\LA
\begin{Resp}\Rsign Ecce Dóminus véniet cum splendóre descéndens, et virtus ejus cum eo, \Star{} Visitáre pópulum suum in pace, et constitúere super eum vitam sempitérnam.\end{Resp}
\begin{Resp}\Vsign Ecce Dóminus noster cum virtúte véniet.\end{Resp}
\begin{Resp}\Rsign Visitáre pópulum suum in pace, et constitúere super eum vitam sempitérnam.\end{Resp}
\begin{Resp}\Vsign Glória Patri, et Fílio, \Star{} et Spirítui Sancto.\end{Resp}
\begin{Resp}\Rsign Visitáre pópulum suum in pace, et constitúere super eum vitam sempitérnam.\end{Resp}
\switchcolumn
\EN
\begin{Resp}\Rsign Behold, the Lord cometh down with glory, and His host is with Him. \Star{} To visit His people in peace, and to establish them in life everlasting.\end{Resp}
\begin{Resp}\Vsign Behold, our Lord cometh with an host.\end{Resp}
\begin{Resp}\Rsign To visit His people in peace, and to establish them in life everlasting.\end{Resp}
\begin{Resp}\Vsign Glory be to the Father, and to the Son, \Star{} and to the Holy Ghost.\end{Resp}
\begin{Resp}\Rsign To visit His people in peace, and to establish them in life everlasting.\end{Resp}
\switchcolumn*

\end{paracol}

\Office{Feria V infra Hebdomadam II Adventus}{Thursday of the Second Week of Advent}{Feria major}
\addcontentsline{toc}{subsection}{Feria V infra Hebdomadam II Adventus\enspace\textit{Thursday of the Second Week of Advent}}
\begin{paracol}{2}
\LA
\LessonHead{Lectio I}
\switchcolumn
\EN
\LessonHead{Lesson I}
\switchcolumn*

\LA
\LessonTitle{De Isaía Prophéta}
\switchcolumn
\EN
\LessonTitle{Lesson from the book of Isaias}
\switchcolumn*

\LA
\Cite{Isa 19:1-2}
\switchcolumn
\EN
\Cite{Isa 19:1-2}
\switchcolumn*

\LA
\noindent Onus Ægýpti. Ecce Dóminus ascéndet super nubem levem, et ingrediétur Ægýptum, et commovebúntur simulácra Ægýpti a fácie ejus, et cor Ægýpti tabéscet in médio ejus. Et concúrrere fáciam Ægýptios advérsus Ægýptios: et pugnábit vir contra fratrem suum, et vir contra amícum suum, cívitas advérsus civitátem, regnum advérsus regnum.\par
\switchcolumn
\EN
\noindent The burden of Egypt. Behold the Lord will ascend upon a swift cloud, and will enter into Egypt, and the idols of Egypt shall be moved at his presence, and the heart of Egypt shall melt in the midst thereof. And I will set the Egyptians to fight against the Egyptians: and they shall fight brother against brother, and friend against friend, city against city, kingdom against kingdom.\par
\switchcolumn*

\LA
\begin{Resp}\Rsign Jerúsalem, cito véniet salus tua: Quare mœróre consúmeris? numquid consiliárius non est tibi, quia innovávit te dolor? \Star{} Salvábo te, et liberábo te, noli timére.\end{Resp}
\begin{Resp}\Vsign Ego enim sum Dóminus, Deus tuus, Sanctus Israël, Redémptor tuus.\end{Resp}
\begin{Resp}\Rsign Salvábo te, et liberábo te, noli timére.\end{Resp}
\switchcolumn
\EN
\begin{Resp}\Rsign Thy salvation cometh quickly, O Jerusalem; why art thou wasted with sorrow? Is there no counselor in thee, that pangs have taken thee? \Star{} Fear not, for I will save thee and deliver thee.\end{Resp}
\begin{Resp}\Vsign For I am the Lord, thy God, the Holy One of Israel, thy Saviour.\end{Resp}
\begin{Resp}\Rsign Fear not, for I will save thee, and deliver thee.\end{Resp}
\switchcolumn*

\LA
\LessonHead{Lectio II}
\switchcolumn
\EN
\LessonHead{Lesson II}
\switchcolumn*

\LA
\Cite{Isa 19:3-6}
\switchcolumn
\EN
\Cite{Isa 19:3-6}
\switchcolumn*

\LA
\noindent Et dirumpétur spíritus Ægýpti in viscéribus ejus, et consílium ejus præcipitábo: et interrogábunt simulácra sua, et divínos suos, et pythónes, et aríolos. Et tradam Ægýptum in manu dominórum crudélium, et rex fortis dominábitur eórum, ait Dóminus Deus exercítuum. Et aréscet aqua de mari, et flúvius desolábitur, atque siccábitur. Et defícient flúmina: attenuabúntur, et siccabúntur rivi ággerum.\par
\switchcolumn
\EN
\noindent And the spirit of Egypt shall be broken in the bowels thereof, and I will cast down their counsel: and they shall consult their idols, and their diviners, and their wizards, and soothsayers. And I will deliver Egypt into the hand of cruel masters, and a strong king shall rule over them, saith the Lord the God of hosts. And the water of the sea shall be dried up, and the river shall be wasted and dry. And the rivers shall fail: the streams of the banks shall be diminished, and be dried up. The reed and the bulrush shall wither away.\par
\switchcolumn*

\LA
\begin{Resp}\Rsign Ecce Dóminus véniet, et omnes Sancti ejus cum eo, et erit in die illa lux magna: et exíbunt de Jerúsalem sicut aqua munda: et regnábit Dóminus in ætérnum \Star{} Super omnes gentes.\end{Resp}
\begin{Resp}\Vsign Ecce Dóminus cum virtúte véniet: et regnum in manu ejus, et potéstas, et impérium.\end{Resp}
\begin{Resp}\Rsign Super omnes gentes.\end{Resp}
\switchcolumn
\EN
\begin{Resp}\Rsign Behold, the Lord shall come, and all His saints with Him, and it shall come to pass in that day that the light shall be great; and they shall go out from Jerusalem like clean water; and the Lord shall be King for ever, \Star{} Over all the earth.\end{Resp}
\begin{Resp}\Vsign Behold, the Lord cometh with an host, and in His hand are the kingdom, and power, and dominion.\end{Resp}
\begin{Resp}\Rsign Over all the earth.\end{Resp}
\switchcolumn*

\LA
\LessonHead{Lectio III}
\switchcolumn
\EN
\LessonHead{Lesson III}
\switchcolumn*

\LA
\Cite{Isa 19:11-13}
\switchcolumn
\EN
\Cite{Isa 19:11-13}
\switchcolumn*

\LA
\noindent Stulti príncipes Táneos, sapiéntes consiliárii Pharaónis dedérunt consílium insípiens. Quómodo dicétis Pharaóni: Fílius sapiéntium ego, fílius regum antiquórum? Ubi nunc sunt sapiéntes tui? annúntient tibi, et índicent quid cogitáverit Dóminus exercítuum super Ægýptum. Stulti facti sunt príncipes Táneos, emarcuérunt príncipes Mémpheos, decepérunt Ægýptum, ángulum populórum ejus.\par
\switchcolumn
\EN
\noindent The princes of Tanis are become fools, the wise counsellors of Pharao have given foolish counsel: how will you say to Pharao: I am the son of the wise, the son of ancient kings? Where are now thy wise men? let them tell thee, and shew what the Lord of hosts hath purposed upon Egypt. The princes of Tanis are become fools, the princes of Memphis are gone astray, they have deceived Egypt, the stay of the people thereof.\par
\switchcolumn*

\LA
\begin{Resp}\Rsign Cívitas Jerúsalem, noli flere: quóniam dóluit Dóminus super te: \Star{} Et áuferet a te omnem tribulatiónem.\end{Resp}
\begin{Resp}\Vsign Ecce Dóminus in fortitúdine véniet: et brácchium ejus dominábitur.\end{Resp}
\begin{Resp}\Rsign Et áuferet a te omnem tribulatiónem.\end{Resp}
\begin{Resp}\Vsign Glória Patri, et Fílio, \Star{} et Spirítui Sancto.\end{Resp}
\begin{Resp}\Rsign Et áuferet a te omnem tribulatiónem.\end{Resp}
\switchcolumn
\EN
\begin{Resp}\Rsign O, thou city of Jerusalem, weep not, for the Lord hath repented Him concerning thee. \Star{} And He will take away from thee all distress.\end{Resp}
\begin{Resp}\Vsign Behold, the Lord shall come with might, and His arm shall rule.\end{Resp}
\begin{Resp}\Rsign And He will take away from thee all distress.\end{Resp}
\begin{Resp}\Vsign Glory be to the Father, and to the Son, \Star{} and to the Holy Ghost.\end{Resp}
\begin{Resp}\Rsign And He will take away from thee all distress.\end{Resp}
\switchcolumn*

\end{paracol}

\Office{Feria VI infra Hebdomadam II Adventus}{Friday of the Second Week of Advent}{Feria major}
\addcontentsline{toc}{subsection}{Feria VI infra Hebdomadam II Adventus\enspace\textit{Friday of the Second Week of Advent}}
\begin{paracol}{2}
\LA
\LessonHead{Lectio I}
\switchcolumn
\EN
\LessonHead{Lesson I}
\switchcolumn*

\LA
\LessonTitle{De Isaía Prophéta}
\switchcolumn
\EN
\LessonTitle{Lesson from the book of Isaias}
\switchcolumn*

\LA
\Cite{Isa 24:1-3}
\switchcolumn
\EN
\Cite{Isa 24:1-3}
\switchcolumn*

\LA
\noindent Ecce Dóminus dissipábit terram, et nudábit eam, et afflíget fáciem ejus, et dispérget habitatóres ejus. Et erit sicut pópulus, sic sacérdos: et sicut servus, sic dóminus ejus: sicut ancílla, sic dómina ejus: sicut emens, sic ille qui vendit: sicut fœnerátor, sic is qui mútuum áccipit: sicut qui répetit, sic qui debet. Dissipatióne dissipábitur terra, et direptióne prædábitur. Dóminus enim locútus est verbum hoc.\par
\switchcolumn
\EN
\noindent Behold the Lord shall lay waste the earth, and shall strip it, and shall afflict the face thereof, and scatter abroad the inhabitants thereof. And it shall be as with the people, so with the priest: and as with the servant, so with his master: as with the handmaid, so with her mistress: as with the buyer, so with the seller: as with the lender, so with the borrower: as with him that calleth for his money, so with him that oweth. With desolation shall the earth be laid waste, and it shall be utterly spoiled: for the Lord hath spoken this word.\par
\switchcolumn*

\LA
\begin{Resp}\Rsign Ecce véniet Dóminus, protéctor noster, Sanctus Israël, \Star{} Corónam regni habens in cápite suo.\end{Resp}
\begin{Resp}\Vsign Et dominábitur a mari usque ad mare, et a flúmine usque ad términos orbis terrárum.\end{Resp}
\begin{Resp}\Rsign Corónam regni habens in cápite suo.\end{Resp}
\switchcolumn
\EN
\begin{Resp}\Rsign Behold, there cometh the Lord, our defender, the Holy One of Israel, \Star{} Wearing a royal crown upon His head.\end{Resp}
\begin{Resp}\Vsign And His dominion shall be from sea even to sea, and from the river even to the ends of the earth.\end{Resp}
\begin{Resp}\Rsign Wearing a royal crown upon His head.\end{Resp}
\switchcolumn*

\LA
\LessonHead{Lectio II}
\switchcolumn
\EN
\LessonHead{Lesson II}
\switchcolumn*

\LA
\Cite{Isa 24:4-6}
\switchcolumn
\EN
\Cite{Isa 24:4-6}
\switchcolumn*

\LA
\noindent Luxit et deflúxit terra, et infirmáta est: deflúxit orbis, infirmáta est altitúdo pópuli terræ. Et terra infécta est ab habitatóribus suis: quia transgréssi sunt leges, mutavérunt jus, dissipavérunt fœdus sempitérnum. Propter hoc maledíctio vorábit terram, et peccábunt habitatóres ejus: ideóque insánient cultóres ejus, et relinquéntur hómines pauci.\par
\switchcolumn
\EN
\noindent The earth mourned, and faded away, and is weakened: the world faded away, the height of the people of the earth is weakened. And the earth is infected by the inhabitants thereof: because they have transgressed the laws, they have changed the ordinance, they have broken the everlasting covenant. Therefore shall a curse devour the earth, and the inhabitants thereof shall sin: and therefore they that dwell therein shall be mad, and few men shall be left.\par
\switchcolumn*

\LA
\begin{Resp}\Rsign Sicut mater consolátur fílios suos, ita consolábor vos, dicit Dóminus: et de Jerúsalem civitáte quam elégi, véniet vobis auxílium: \Star{} Et vidébitis et gaudébit cor vestrum.\end{Resp}
\begin{Resp}\Vsign Dabo in Sion salútem, et in Jerúsalem glóriam meam.\end{Resp}
\begin{Resp}\Rsign Et vidébitis et gaudébit cor vestrum.\end{Resp}
\switchcolumn
\EN
\begin{Resp}\Rsign As a mother comforteth her children, so will I comfort you, saith the Lord; My help also cometh unto you out of Jerusalem, the city which I have chosen. \Star{} And when ye see this, your heart shall rejoice.\end{Resp}
\begin{Resp}\Vsign I will place salvation in Zion and in Jerusalem My glory.\end{Resp}
\begin{Resp}\Rsign And when ye shall see this, your heart shall rejoice.\end{Resp}
\switchcolumn*

\LA
\LessonHead{Lectio III}
\switchcolumn
\EN
\LessonHead{Lesson III}
\switchcolumn*

\LA
\Cite{Isa 24:7-16}
\switchcolumn
\EN
\Cite{Isa 24:7-16}
\switchcolumn*

\LA
\noindent Luxit vindémia, infirmáta est vitis, ingemuérunt omnes qui lætabántur corde. Cessávit gáudium tympanórum, quiévit sónitus lætántium, contícuit dulcédo cítharæ. Cum cántico non bibent vinum, amára erit pótio bibéntibus illam. Attríta est cívitas vanitátis, clausa est omnis domus, nullo introëúnte. Clamor erit super vino in platéis: desérta est omnis lætítia: translátum est gáudium terræ. Relícta est in urbe solitudo, et calámitas ópprimet portas. Quia hæc erunt in médio terræ, in médio populórum: quómodo si paucæ olívæ, quæ remansérunt, excutiántur ex ólea: et racémi, cum fúerit fínita vindémia. Hi levábunt vocem suam, atque laudábunt: cum glorificátus fúerit Dóminus, hínnient de mari. Propter hoc in doctrínis glorificáte Dóminum, in ínsulis maris nomen Dómini Dei Israël. A fínibus terræ laudes audívimus, glóriam justi.\par
\switchcolumn
\EN
\noindent The vintage hath mourned, the vine hath languished away, all the merryhearted have sighed. The mirth of timbrels hath ceased, the noise of them that rejoice is ended, the melody of the harp is silent. They shall not drink wine with a song: the drink shall be bitter to them that drink it. The city of vanity is broken down, every house is shut up, no man cometh in. There shall be a crying for wine in the streets: all mirth is forsaken: the joy of the earth is gone away. Desolation is left in the city, and calamity shall oppress the gates. For it shall be thus in the midst of the earth, in the midst of the people, as if a few olives, that remain, should be shaken out of the olive tree: or grapes, when the vintage is ended. These shall lift up their voice, and shall give praise: when the Lord shall be glorified, they shall make a joyful noise from the sea. Therefore glorify ye the Lord in instruction: the name of the Lord God of Israel in the islands of the sea. From the ends of the earth we have heard praises, the glory of the Just one.\par
\switchcolumn*

\LA
\begin{Resp}\Rsign Jerúsalem, plantábis víneam in móntibus tuis: exsultábis, quóniam dies Dómini véniet: surge, Sion, convértere ad Dóminum Deum tuum: gaude et lætáre, Jacob: \Star{} Quia de médio géntium Salvátor tuus véniet.\end{Resp}
\begin{Resp}\Vsign Exsúlta satis, fília Sion: júbila, fília Jerúsalem.\end{Resp}
\begin{Resp}\Rsign Quia de médio géntium Salvátor tuus véniet.\end{Resp}
\begin{Resp}\Vsign Glória Patri, et Fílio, \Star{} et Spirítui Sancto.\end{Resp}
\begin{Resp}\Rsign Quia de médio géntium Salvátor tuus véniet.\end{Resp}
\switchcolumn
\EN
\begin{Resp}\Rsign Thou shalt yet plant vines upon thy mountains, O Jerusalem thou shalt sing for joy, for the day of the Lord cometh; arise, O Zion, and turn unto the Lord thy God; rejoice and be glad, O Jacob. \Star{} For thy Saviour cometh from the midst of the nations.\end{Resp}
\begin{Resp}\Vsign Sing aloud for joy, O daughter of Zion; shout with gladness, O daughter of Jerusalem.\end{Resp}
\begin{Resp}\Rsign For thy Saviour cometh from the midst of the nations.\end{Resp}
\begin{Resp}\Vsign Glory be to the Father, and to the Son, \Star{} and to the Holy Ghost.\end{Resp}
\begin{Resp}\Rsign For thy Saviour cometh from the midst of the nations.\end{Resp}
\switchcolumn*

\end{paracol}

\Office{Sabbato infra Hebdomadam II Adventus}{Saturday of the Second Week of Advent}{Feria major}
\addcontentsline{toc}{subsection}{Sabbato infra Hebdomadam II Adventus\enspace\textit{Saturday of the Second Week of Advent}}
\begin{paracol}{2}
\LA
\LessonHead{Lectio I}
\switchcolumn
\EN
\LessonHead{Lesson I}
\switchcolumn*

\LA
\LessonTitle{De Isaía Prophéta}
\switchcolumn
\EN
\LessonTitle{Lesson from the book of Isaias}
\switchcolumn*

\LA
\Cite{Isa 25:1-4}
\switchcolumn
\EN
\Cite{Isa 25:1-4}
\switchcolumn*

\LA
\noindent Dómine, Deus meus es tu; exaltábo te, et confitébor nómini tuo: quóniam fecísti mirabília, cogitatiónes antíquas fidéles. Amen. Quia posuísti civitátem in túmulum, urbem fortem in ruínam, domum alienórum: ut non sit cívitas, et in sempitérnum non ædificétur. Super hoc laudábit te pópulus fortis; cívitas géntium robustárum timébit te: quia factus es fortitúdo páuperi, fortitúdo egéno in tribulatióne sua, spes a túrbine, umbráculum ab æstu.\par
\switchcolumn
\EN
\noindent O Lord, thou art my God, I will exalt thee, and give glory to thy name: for thou hast done wonderful things, thy designs of old faithful, amen. For thou hast reduced the city to a heap, the strong city to ruin, the house of strangers, to be no city, and to be no more built up for ever. Therefore shall a strong people praise thee, the city of mighty nations shall fear thee. Because thou hast been a strength to the poor, a strength to the needy in his distress: a refuge from the whirlwind, a shadow from the heat.\par
\switchcolumn*

\LA
\begin{Resp}\Rsign Egrediétur Dóminus de Samaría ad portam, quæ réspicit ad Oriéntem: et véniet in Béthlehem ámbulans super aquas redemptiónis Judæ: \Star{} Tunc salvus erit omnis homo: quia ecce véniet.\end{Resp}
\begin{Resp}\Vsign Et præparábitur in misericórdia sólium ejus, et sedébit super illud in veritáte.\end{Resp}
\begin{Resp}\Rsign Tunc salvus erit omnis homo: quia ecce véniet.\end{Resp}
\switchcolumn
\EN
\begin{Resp}\Rsign The Lord shall go forth out of Samaria unto the gate that looketh toward the East; and He shall come into Bethlehem, walking upon the waters of the redemption of Judah. \Star{} Then shall every one be saved for, behold, He cometh.\end{Resp}
\begin{Resp}\Vsign And in mercy shall His throne be established, and He shall sit upon it in truth.\end{Resp}
\begin{Resp}\Rsign Then shall every one be saved for, behold, He cometh.\end{Resp}
\switchcolumn*

\LA
\LessonHead{Lectio II}
\switchcolumn
\EN
\LessonHead{Lesson II}
\switchcolumn*

\LA
\Cite{Isa 25:4-7}
\switchcolumn
\EN
\Cite{Isa 25:4-7}
\switchcolumn*

\LA
\noindent Spíritus enim robustórum quasi turbo impéllens paríetem. Sicut æstus in siti, tumúltum alienórum humiliábis; et quasi calóre sub nube torrénte, propáginem fórtium marcéscere fácies. Et fáciet Dóminus exercítuum ómnibus pópulis in monte hoc convívium píngujum, convívium vindémiæ, píngujum medullatórum, vindémiæ defæcátæ. Et præcipitábit in monte isto fáciem vínculi colligáti super omnes pópulos, et telam quam ordítus est super omnes natiónes.\par
\switchcolumn
\EN
\noindent For the blast of the mighty is like a whirlwind beating against a wall. Thou shalt bring down the tumult of strangers, as heat in thirst: and as with heat under a burning cloud, thou shalt make the branch of the mighty to wither away. And the Lord of hosts shall make unto all people in this mountain, a feast of fat things, a feast of wine, of fat things full of marrow, of wine purified from the lees. And he shall destroy in this mountain the face of the bond with which all people were tied, and the web that he began over all nations.\par
\switchcolumn*

\LA
\begin{Resp}\Rsign Festína, ne tardáveris, Dómine: \Star{} Et líbera pópulum tuum.\end{Resp}
\begin{Resp}\Vsign Veni, Dómine, et noli tardáre: reláxa facínora plebi tuæ.\end{Resp}
\begin{Resp}\Rsign Et líbera pópulum tuum.\end{Resp}
\switchcolumn
\EN
\begin{Resp}\Rsign Make haste, O Lord, make no tarrying. \Star{} And deliver thy people.\end{Resp}
\begin{Resp}\Vsign O Lord, come and make no tarrying loose the bonds of thy people.\end{Resp}
\begin{Resp}\Rsign And deliver thy people.\end{Resp}
\switchcolumn*

\LA
\LessonHead{Lectio III}
\switchcolumn
\EN
\LessonHead{Lesson III}
\switchcolumn*

\LA
\Cite{Isa 25:8-12}
\switchcolumn
\EN
\Cite{Isa 25:8-12}
\switchcolumn*

\LA
\noindent Præcipitábit mortem in sempitérnum; et áuferet Dóminus Deus lácrimam ab omni fácie, et oppróbrium pópuli sui áuferet de univérsa terra: quia Dóminus locútus est. Et dicet in die illa: Ecce Deus noster iste; exspectávimus eum, et salvábit nos; iste Dóminus, sustinúimus eum: exsultábimus, et lætábimur in salutári ejus. Quia requiéscet manus Dómini in monte isto; et triturábitur Moab sub eo, sícuti terúntur páleæ in plaustro. Et exténdet manus suas sub eo sicut exténdit natans ad natándum; et humiliábit glóriam ejus cum allisióne mánuum ejus. Et muniménta sublímium murórum tuórum cóncident, et humiliabúntur, et detrahéntur in terram usque ad púlverem.\par
\switchcolumn
\EN
\noindent He shall cast death down headlong forever: and the Lord God shall wipe away tears from every face, and the reproach of his people he shall take away from off the whole earth: for the Lord hath spoken it. And they shall say in that day: Lo, this is our God, we have waited for him, and he will save us: this is the Lord, we have patiently waited for him, we shall rejoice and be joyful in his salvation. For the hand of the Lord shall rest in this mountain: and Moab shall be trodden down under him, as straw is broken in pieces with the wain. And he shall stretch forth his hands under him, as he that swimmeth stretcheth forth his hands to swim: and he shall bring down his glory with the dashing of his hands. And the bulwarks of thy high walls shall fall, and be brought low, and shall be pulled down to the ground, even to the dust.\par
\switchcolumn*

\LA
\begin{Resp}\Rsign Ecce Dóminus véniet cum splendóre descéndens, et virtus ejus cum eo, \Star{} Visitáre pópulum suum in pace, et constitúere super eum vitam sempitérnam.\end{Resp}
\begin{Resp}\Vsign Ecce Dóminus noster cum virtúte véniet.\end{Resp}
\begin{Resp}\Rsign Visitáre pópulum suum in pace, et constitúere super eum vitam sempitérnam.\end{Resp}
\begin{Resp}\Vsign Glória Patri, et Fílio, \Star{} et Spirítui Sancto.\end{Resp}
\begin{Resp}\Rsign Visitáre pópulum suum in pace, et constitúere super eum vitam sempitérnam.\end{Resp}
\switchcolumn
\EN
\begin{Resp}\Rsign Behold, the Lord cometh down with glory, and His host is with Him. \Star{} To visit His people in peace, and to establish them in life everlasting.\end{Resp}
\begin{Resp}\Vsign Behold, our Lord cometh with an host.\end{Resp}
\begin{Resp}\Rsign To visit His people in peace, and to establish them in life everlasting.\end{Resp}
\begin{Resp}\Vsign Glory be to the Father, and to the Son, \Star{} and to the Holy Ghost.\end{Resp}
\begin{Resp}\Rsign To visit His people in peace, and to establish them in life everlasting.\end{Resp}
\switchcolumn*

\end{paracol}

\Office{Dominica III Adventus}{Third Sunday of Advent}{Semiduplex II classis}
\addcontentsline{toc}{subsection}{Dominica III Adventus\enspace\textit{Third Sunday of Advent}}
\begin{paracol}{2}
\LA
\LessonHead{Lectio I}
\switchcolumn
\EN
\LessonHead{Lesson I}
\switchcolumn*

\LA
\LessonTitle{De Isaía Prophéta}
\switchcolumn
\EN
\LessonTitle{Lesson from the book of Isaias}
\switchcolumn*

\LA
\Cite{Isa 26:1-6}
\switchcolumn
\EN
\Cite{Isa 26:1-6}
\switchcolumn*

\LA
\noindent In die illa cantábitur cánticum istud in terra Juda: Urbs fortitúdinis nostræ Sion Salvátor, ponétur in ea murus et antemurále. Aperíte portas et ingrediátur gens justa, custódiens veritátem. Vetus error ábiit: servábis pacem: pacem, quia in te sperávimus. Sperástis in Dómino in sǽculis ætérnis, in Dómino Deo forti, in perpétuum. Quia incurvábit habitántes in excélso, civitátem sublímem humiliábit. Humiliábit eam usque ad terram, détrahet eam usque ad púlverem. Conculcábit eam pes, pedes páuperis, gressus egenórum.\par
\switchcolumn
\EN
\noindent In that day shall this canticle be sung in the land of Juda. Sion the city of our strength a saviour, a wall and a bulwark shall be set therein. Open ye the gates, and let the just nation, that keepeth the truth, enter in. The old error is passed away: thou wilt keep peace: peace, because we have hoped in thee. You have hoped in the Lord for evermore, in the Lord God mighty for ever. For he shall bring down them that dwell on high, the high city he shall lay low. He shall bring it down even to the ground, he shall pull it down even to the dust. The foot shall tread it down, the feet of the poor, the steps of the needy.\par
\switchcolumn*

\LA
\begin{Resp}\Rsign Ecce apparébit Dóminus super nubem cándidam, \Star{} Et cum eo Sanctórum míllia: et habébit in vestiménto, et in fémore suo scriptum: Rex regum, et Dóminus dominántium.\end{Resp}
\begin{Resp}\Vsign Apparébit in finem, et non mentiétur; si moram fécerit, exspécta eum, quia véniens véniet.\end{Resp}
\begin{Resp}\Rsign Et cum eo Sanctórum míllia: et habébit in vestiménto, et in fémore suo scriptum: Rex regum, et Dóminus dominántium.\end{Resp}
\switchcolumn
\EN
\begin{Resp}\Rsign Behold, the Lord shall appear upon a white cloud, \Star{} And ten thousand of His saints with Him; and He shall have on His vesture, and on His thigh a name written King of kings, and Lord of lords.\end{Resp}
\begin{Resp}\Vsign He shall appear and not lie; though He tarry, wait for Him, because He will surely come.\end{Resp}
\begin{Resp}\Rsign And ten thousand of His saints with Him; and He shall have on His vesture, and on His thigh a name written King of kings, and Lord of lords.\end{Resp}
\switchcolumn*

\LA
\LessonHead{Lectio II}
\switchcolumn
\EN
\LessonHead{Lesson II}
\switchcolumn*

\LA
\Cite{Isa 26:7-10}
\switchcolumn
\EN
\Cite{Isa 26:7-10}
\switchcolumn*

\LA
\noindent Sémita justi recta est, rectus callis justi ad ambulándum. Et in sémita judiciórum tuórum, Dómine, sustinúimus te: nomen tuum, et memoriále tuum in desidério ánimæ. Anima mea desiderávit te in nocte, sed et spíritu meo in præcórdiis meis de mane vigilábo ad te. Cum féceris judícia tua in terra, justítiam discent habitatóres orbis. Misereámur ímpio, et non discet justítiam: in terra sanctórum iníqua gessit, et non vidébit glóriam Dómini.\par
\switchcolumn
\EN
\noindent The way of the just is right, the path of the just is right to walk in. And in the way of thy judgments, O Lord, we have patiently waited for thee: thy name, and thy remembrance are the desire of the soul. My soul hath desired thee in the night: yea, and with my spirit within me in the morning early I will watch to thee. When thou shalt do thy judgments on the earth, the inhabitants of the world shall learn justice. Let us have pity on the wicked, but he will not learn justice: in the land of the saints he hath done wicked things, and he shall not see the glory of the Lord.\par
\switchcolumn*

\LA
\begin{Resp}\Rsign Béthlehem cívitas Dei summi, ex te éxiet Dominátor Israël, et egréssus ejus sicut a princípio diérum æternitátis, et magnificábitur in médio univérsæ terræ: \Star{} Et pax erit in terra nostra, dum vénerit.\end{Resp}
\begin{Resp}\Vsign Loquétur pacem in géntibus, et potéstas ejus a mari usque ad mare.\end{Resp}
\begin{Resp}\Rsign Et pax erit in terra nostra, dum vénerit.\end{Resp}
\switchcolumn
\EN
\begin{Resp}\Rsign Thou, Bethlehem, art the city of the Most High God, out of thee shall He come forth That is to be Ruler in Israel; Whose goings forth have been from of old, from everlasting, and now shall He be great unto the ends of the earth. \Star{} And this Man shall be the peace in our land, when He shall come.\end{Resp}
\begin{Resp}\Vsign He shall speak peace unto the Gentiles, and shall have dominion from sea to sea.\end{Resp}
\begin{Resp}\Rsign And this Man shall be the peace in our land, when He shall come.\end{Resp}
\switchcolumn*

\LA
\LessonHead{Lectio III}
\switchcolumn
\EN
\LessonHead{Lesson III}
\switchcolumn*

\LA
\Cite{Isa 26:11-14}
\switchcolumn
\EN
\Cite{Isa 26:11-14}
\switchcolumn*

\LA
\noindent Dómine, exaltétur manus tua, et non vídeant: vídeant, et confundántur zelántes pópuli: et ignis hostes tuos dévoret. Dómine, dabis pacem nobis: ómnia enim ópera nostra operátus es nobis. Dómine Deus noster, possedérunt nos dómini absque te, tantum in te recordémur nóminis tui. Moriéntes non vivant, gigántes non resúrgant: proptérea visitásti et contrivísti eos, et perdidísti omnem memóriam eórum.\par
\switchcolumn
\EN
\noindent Lord, let thy hand be exalted, and let them not see: let the envious people see, and be confounded: and let fire devour thy enemies. Lord, thou wilt give us peace: for thou hast wrought all our works for us. O Lord our God, other lords besides thee have had dominion over us, only in thee let us remember thy name. Let not the dead live, let not the giants rise again: therefore hast thou visited and destroyed them, and hast destroyed all their memory.\par
\switchcolumn*

\LA
\begin{Resp}\Rsign Qui ventúrus est, véniet, et non tardábit: et jam non erit timor in fínibus nostris: \Star{} Quóniam ipse est Salvátor noster.\end{Resp}
\begin{Resp}\Vsign Depónet omnes iniquitátes nostras, et proíciet in profúndum maris ómnia peccáta nostra.\end{Resp}
\begin{Resp}\Rsign Quóniam ipse est Salvátor noster.\end{Resp}
\begin{Resp}\Vsign Glória Patri, et Fílio, \Star{} et Spirítui Sancto.\end{Resp}
\begin{Resp}\Rsign Quóniam ipse est Salvátor noster.\end{Resp}
\switchcolumn
\EN
\begin{Resp}\Rsign He That shall come, will come, and will not tarry; and there shall no more be fear in our borders. \Star{} For He is our Saviour.\end{Resp}
\begin{Resp}\Vsign He shall tread down all our iniquities, and cast all our sins into the depths of the sea.\end{Resp}
\begin{Resp}\Rsign For He is our Saviour.\end{Resp}
\begin{Resp}\Vsign Glory be to the Father, and to the Son, \Star{} and to the Holy Ghost.\end{Resp}
\begin{Resp}\Rsign For He is our Saviour.\end{Resp}
\switchcolumn*

\LA
\LessonHead{Lectio IV}
\switchcolumn
\EN
\LessonHead{Lesson IV}
\switchcolumn*

\LA
\LessonTitle{Sermo sancti Leónis Papæ}
\switchcolumn
\EN
\LessonTitle{From the Sermons of Pope St. Leo (the Great)}
\switchcolumn*

\LA
\Source{Sermo 2 de jejúnio décimi mensis et collectis}
\switchcolumn
\EN
\Source{Second on the December Fast, and Almsgiving}
\switchcolumn*

\LA
\noindent Quod témporis rátio, et devotiónis nostræ ádmonet consuetúdo, pastoráli vobis, dilectíssimi, sollicitúdine prædicámus, décimi mensis celebrándum esse jejúnium, quo pro consummáta perceptióne ómnium frúctuum, digníssime largitóri eórum Deo continéntiæ libámen offértur. Quid enim potest efficácius esse jejúnio? cujus observántia appropinquámus Deo, et resisténtes diábolo, vítia blanda superámus.\par
\switchcolumn
\EN
\noindent Dearly beloved brethren, with the care which becometh us as the shepherd of your souls, we urge upon you the rigid observance of this December Fast. The month of December hath come round again, and with it this devout custom of the Church. The fruits of the year, which is drawing to a close, are now all gathered in, and we most meetly offer our abstinence to God as a sacrifice of thanksgiving. And what can be more useful than fasting, that exercise by which we draw nigh to God, make a stand against the devil, and overcome the softer enticements of sin?\par
\switchcolumn*

\LA
\begin{Resp}\Rsign Ægýpte, noli flere, quia Dominátor tuus véniet tibi, ante cujus conspéctum movebúntur abýssi, \Star{} Liberáre pópulum suum de manu poténtiæ.\end{Resp}
\begin{Resp}\Vsign Ecce véniet Dóminus exercítuum, Deus tuus cum potestáte magna.\end{Resp}
\begin{Resp}\Rsign Liberáre pópulum suum de manu poténtiæ.\end{Resp}
\switchcolumn
\EN
\begin{Resp}\Rsign Weep not, O Egypt, for the Ruler cometh unto thee, and the depths shall be moved at His presence. \Star{} To deliver His people out of the hand of the mighty.\end{Resp}
\begin{Resp}\Vsign Behold, the Lord of hosts, thy God, cometh with great power.\end{Resp}
\begin{Resp}\Rsign To deliver His people out of the hand of the mighty.\end{Resp}
\switchcolumn*

\LA
\LessonHead{Lectio V}
\switchcolumn
\EN
\LessonHead{Lesson V}
\switchcolumn*

\LA
\noindent Semper enim virtúti cibus jejúnium fuit. De abstinéntia dénique pródeunt castæ cogitatiónes, rationábiles voluntátes, salubrióra consília: et per voluntárias afflictiónes caro concupiscéntiis móritur, virtútibus spíritus innovátur. Sed quia non solo jejúnio animárum nostrárum salus acquíritur, jejúnium nostrum misericórdiis páuperum suppleámus. Impendámus virtúti, quod subtráhimus voluptáti. Fiat reféctio páuperis abstinéntia jejunántis.\par
\switchcolumn
\EN
\noindent Fasting hath ever been the bread of strength. From abstinence proceed pure thoughts, reasonable desires, and healthy counsels. By voluntary mortifications the flesh dieth to lust, and the soul is renewed in might. But since fasting is not the only mean whereby we get health for our souls, let us add to our fasting works of mercy. Let us spend in good deeds what we take from indulgence. Let our fast become the banquet of the poor.\par
\switchcolumn*

\LA
\begin{Resp}\Rsign Prope est ut véniat tempus ejus, et dies ejus non elongabúntur: \Star{} Miserébitur Dóminus Jacob, et Israël salvábitur.\end{Resp}
\begin{Resp}\Vsign Revértere, virgo Israël, revértere ad civitátes tuas.\end{Resp}
\begin{Resp}\Rsign Miserébitur Dóminus Jacob, et Israël salvábitur.\end{Resp}
\switchcolumn
\EN
\begin{Resp}\Rsign Her time is near to come, and her days shall not be prolonged. \Star{} Lord shall have mercy on Jacob, and Israel shall be saved.\end{Resp}
\begin{Resp}\Vsign Turn again, O Virgin of Israel, turn again to thy cities.\end{Resp}
\begin{Resp}\Rsign For the Lord shall have mercy on Judah, and Israel shall be saved.\end{Resp}
\switchcolumn*

\LA
\LessonHead{Lectio VI}
\switchcolumn
\EN
\LessonHead{Lesson VI}
\switchcolumn*

\LA
\noindent Studeámus viduárum defensióni, pupillórum utilitáti, lugéntium consolatióni, dissidéntium paci. Suscipiátur peregrínus, adjuvétur oppréssus, vestiátur nudus, foveátur ægrótus: ut quicúmque nostrum de justis labóribus auctóri bonórum ómnium Deo sacrifícium hujus pietátis obtúlerit, ab eodem regni cæléstis prǽmium percípere mereátur. Quarta ígitur et sexta Feria jejunémus; Sábbato autem apud beátum Petrum Apóstolum páriter vigilémus: cujus suffragántibus méritis, quæ póscimus, impetráre possímus per Dóminum nostrum Jesum Christum, qui cum Patre et Sancto Spíritu vivit et regnat in sǽcula sæculórum. Amen.\par
\switchcolumn
\EN
\noindent Let us defend the widow and serve the orphan; let us comfort the afflicted and reconcile the estranged; let us take in the wanderer and succour the oppressed; let us clothe the naked and cherish the sick. And may every one of us that shall offer to the God of all goodness this Advent sacrifice of fasting and alms be by Him fitted to receive an eternal reward in His heavenly kingdom! We fast on Wednesday and Friday; and there is likewise a Vigil on Saturday at the Church of St. Peter, that by his good prayers we may the more effectually obtain what we ask for, through our Lord Jesus Christ, Who with the Father and the Holy Ghost, liveth and reigneth, God, world without end. Amen.\par
\switchcolumn*

\LA
\begin{Resp}\Rsign Descéndet Dóminus sicut plúvia in vellus: \Star{} Oriétur in diébus ejus justítia, et abundántia pacis.\end{Resp}
\begin{Resp}\Vsign Et adorábunt eum omnes reges, omnes gentes sérvient ei.\end{Resp}
\begin{Resp}\Rsign Oriétur in diébus ejus justítia, et abundántia pacis.\end{Resp}
\begin{Resp}\Vsign Glória Patri, et Fílio, \Star{} et Spirítui Sancto.\end{Resp}
\begin{Resp}\Rsign Oriétur in diébus ejus justítia, et abundántia pacis.\end{Resp}
\switchcolumn
\EN
\begin{Resp}\Rsign The Lord shall come down like rain upon a fleece. \Star{} In His days shall righteousness flourish, and abundance of peace.\end{Resp}
\begin{Resp}\Vsign All the kings of the earth shall fall down before Him, all nations shall serve Him.\end{Resp}
\begin{Resp}\Rsign In His days shall righteousness flourish, and abundance of peace.\end{Resp}
\begin{Resp}\Vsign Glory be to the Father, and to the Son, \Star{} and to the Holy Ghost.\end{Resp}
\begin{Resp}\Rsign In His days shall righteousness flourish, and abundance of peace.\end{Resp}
\switchcolumn*

\LA
\LessonHead{Lectio VII}
\switchcolumn
\EN
\LessonHead{Lesson VII}
\switchcolumn*

\LA
\LessonTitle{Léctio sancti Evangélii secúndum Joánnem}
\switchcolumn
\EN
\LessonTitle{From the Holy Gospel according to John}
\switchcolumn*

\LA
\Cite{Joannes 1:19-28}
\switchcolumn
\EN
\Cite{John 1:19-28}
\switchcolumn*

\LA
\noindent In illo témpore: Misérunt Judǽi ab Jerosólymis sacerdótes et levítas ad Joánnem, ut interrogárent eum: Tu quis es? Et réliqua.\par
\switchcolumn
\EN
\noindent In that time were sent from Jerusalem priests and Levites to John, to ask him: Who art thou? And so on.\par
\switchcolumn*

\LA
\LessonTitle{Homilía sancti Gregórii Papæ}
\switchcolumn
\EN
\LessonTitle{Homily by Pope St. Gregory (the Great)}
\switchcolumn*

\LA
\Source{Homilia 7 in Evangelia}
\switchcolumn
\EN
\Source{7th on the Gospels}
\switchcolumn*

\LA
\noindent Ex hujus nobis lectiónis verbis, fratres caríssimi, Joánnis humílitas commendátur: qui cum tantæ virtútis esset, ut Christus credi potuísset, elégit sólide subsístere in se, ne humána opinióne raperétur inániter super se. Nam conféssus est, et non negávit: et conféssus est, Quia non sum ego Christus. Sed qui dixit, Non sum; negávit plane quod non erat, sed non negávit quod erat: ut veritátem loquens, ejus membrum fíeret, cujus sibi nomen falláciter non usurpáret. Cum ergo non vult appétere nomen Christi, factus est membrum Christi: quia dum infirmitátem suam stúduit humíliter agnóscere, illíus celsitúdinem méruit veráciter obtinére.\par
\switchcolumn
\EN
\noindent Dearly beloved brethren, the first thing which striketh us in today's Gospel is the lowly-mindedness of John. He was so great that it was thought he might be the Christ; yet he soberly chose rather to seem only what he really was, than to let the belief of men invest him with a dignity which did not belong to him; for he confessed, and denied not, but confessed, I am not the Christ, at the same time he would not deny what he was in reality; and thus his very truth-speaking made him a member of Him Whose title he would not by falsehood take. In that he arrogated not to himself the name of Christ, he became a member of Christ. While he humbly strove to confess his own weakness, he earned by his simplicity a part in the grandeur of his Master.\par
\switchcolumn*

\LA
\begin{Resp}\Rsign Veni, Dómine, et noli tardáre: reláxa facínora plebi tuæ, \Star{} Et révoca dispérsos in terram suam.\end{Resp}
\begin{Resp}\Vsign Excita, Dómine, poténtiam tuam, et veni, ut salvos fácias nos.\end{Resp}
\begin{Resp}\Rsign Et révoca dispérsos in terram suam.\end{Resp}
\switchcolumn
\EN
\begin{Resp}\Rsign O Lord, come and make no tarrying; loosen the bonds of thy people. \Star{} And gather again into their own land them that are scattered abroad.\end{Resp}
\begin{Resp}\Vsign O Lord, stir up thy strength, and come and save us.\end{Resp}
\begin{Resp}\Rsign And gather again into their own land them that are scattered abroad.\end{Resp}
\switchcolumn*

\LA
\LessonHead{Lectio VIII}
\switchcolumn
\EN
\LessonHead{Lesson VIII}
\switchcolumn*

\LA
\noindent Sed cum ex lectióne alia, Redemptóris nostri senténtia ad mentem redúcitur, ex hujus lectiónis verbis nobis quǽstio valde impléxa generátur. Alio quippe in loco inquisítus a discípulis Dóminus de Elíæ advéntu, respóndit: Elías jam venit, et non cognovérunt eum, sed fecérunt in eum quæcúmque voluérunt: et, si vultis scire, Joánnes ipse est Elías. Requisítus autem Joánnes dicit: Non sum Elías. Quid est hoc, fratres caríssimi, quia quod Véritas affírmat, hoc prophéta Veritátis negat? Valde namque inter se divérsa sunt: Ipse est: et, Non sum. Quómodo ergo prophéta veritátis est, si ejusdem Veritátis sermónibus concors non est?\par
\switchcolumn
\EN
\noindent In considering this subject we find an apparent contradiction between one of John's statements, and the saying of our Redeemer recorded in another part of the Gospel. (Matth. xvii. 10-12.) When His disciples asked our Lord regarding the coming of Elias, He answered: Elias is come already, and they knew him not, but have done unto him whatsoever they listed. And if ye will receive it, this (that is, John) is Elias. (Matth. xi. 14.) But when John was asked if he was Elias, he answered: I am not. How comes it then, dearly beloved brethren, that we find the Truth Itself asserting what the prophet of the Truth denied? It must evidently be that our Lord meant one thing and John another, when the Lord said, This is, and John, I am not. For how can he be the prophet of truth, if he speak not according to the word of Him Who is the Eternal Truth?\par
\switchcolumn*

\LA
\begin{Resp}\Rsign Ecce radix Jesse descéndet in salútem populórum, ipsum gentes deprecabúntur: \Star{} Et erit nomen ejus gloriósum.\end{Resp}
\begin{Resp}\Vsign Dabit ei Dóminus Deus sedem David, patris ejus, et regnábit in domo Jacob in ætérnum.\end{Resp}
\begin{Resp}\Rsign Et erit nomen ejus gloriósum.\end{Resp}
\switchcolumn
\EN
\begin{Resp}\Rsign Behold, there shall be a root of Jesse, which shall come for salvation unto the people, to it shall the Gentiles seek; \Star{} And His name shall be glorious.\end{Resp}
\begin{Resp}\Vsign The Lord God shall give unto Him the throne of His father David, and He shall reign over the house of Jacob for ever.\end{Resp}
\begin{Resp}\Rsign And His name shall be glorious.\end{Resp}
\switchcolumn*

\LA
\LessonHead{Lectio IX}
\switchcolumn
\EN
\LessonHead{Lesson IX}
\switchcolumn*

\LA
\noindent Sed si subtíliter véritas ipsa requirátur, hoc quod inter se contrárium sonat, quómodo contrárium non sit, invenítur. Ad Zacharíam namque de Joánne Angelus dicit: Ipse præcédet ante illum in spíritu et virtúte Elíæ. Qui idcírco ventúrus in spíritu et virtúte Elíæ dícitur, quia sicut Elías secúndum Dómini advéntum prævéniet, ita Joánnes prævénit primum. Sicut ille præcúrsor ventúrus est Júdicis, ita iste præcúrsor est factus Redemptóris. Joánnes ígitur in spíritu Elías erat, in persóna Elías non erat. Quod ergo Dóminus fatétur de spíritu, hoc Joánnes dénegat de persóna.\par
\switchcolumn
\EN
\noindent Let us then more minutely examine these words, and we shall find that there is no real contradiction. When the Angel announced to Zacharias the coming birth of John he said He shall go before Him in the spirit and power of Elias, (Luke i. 17.) As the old Elias will come again before the Second Advent of the Lord, so did John, as the new Elias, go before the First Advent, in the spirit and power of Elias. As the old Elias will be the Fore-runner of the Judge, so the new Elias was the Forerunner of the Saviour. John then was Elias in spirit, but not in person; and our Lord asserteth of the spirit what John denieth of the person.\par
\switchcolumn*

\LA
\begin{Resp}\Rsign Docébit nos Dóminus vias suas, et ambulábimus in sémitis ejus: \Star{} Quia de Sion exíbit lex, et verbum Dómini de Jerúsalem.\end{Resp}
\begin{Resp}\Vsign Veníte, ascendámus ad montem Dómini, et ad domum Dei Jacob.\end{Resp}
\begin{Resp}\Rsign Quia de Sion exíbit lex, et verbum Dómini de Jerúsalem.\end{Resp}
\begin{Resp}\Vsign Glória Patri, et Fílio, \Star{} et Spirítui Sancto.\end{Resp}
\begin{Resp}\Rsign Quia de Sion exíbit lex, et verbum Dómini de Jerúsalem.\end{Resp}
\switchcolumn
\EN
\begin{Resp}\Rsign The Lord will teach us of His ways, and we will walk in His paths. \Star{} For out of Zion shall go forth the law, and the word of the Lord from Jerusalem.\end{Resp}
\begin{Resp}\Vsign Come ye, and let us go up to the mountain of the Lord, and to the house of the God of Jacob.\end{Resp}
\begin{Resp}\Rsign For out of Zion shall go forth the law, and the word of the Lord from Jerusalem.\end{Resp}
\begin{Resp}\Vsign Glory be to the Father, and to the Son, \Star{} and to the Holy Ghost.\end{Resp}
\begin{Resp}\Rsign For out of Zion shall go forth the law, and the word of the Lord from Jerusalem.\end{Resp}
\switchcolumn*

\end{paracol}

\Office{Feria II infra Hebdomadam III Adventus}{Monday of the Third Week of Advent}{Feria major}
\addcontentsline{toc}{subsection}{Feria II infra Hebdomadam III Adventus\enspace\textit{Monday of the Third Week of Advent}}
\begin{paracol}{2}
\LA
\LessonHead{Lectio I}
\switchcolumn
\EN
\LessonHead{Lesson I}
\switchcolumn*

\LA
\LessonTitle{De Isaía Prophéta}
\switchcolumn
\EN
\LessonTitle{Lesson from the book of Isaias}
\switchcolumn*

\LA
\Cite{Isa 28:1-3}
\switchcolumn
\EN
\Cite{Isa 28:1-3}
\switchcolumn*

\LA
\noindent Væ corónæ supérbiæ, ébriis Ephraim, et flori decidénti, glóriæ exsultatiónis ejus, qui erant in vértice vallis pinguíssimæ, errántes a vino. Ecce válidus et fortis Dóminus, sicut ímpetus grándinis: turbo confríngens, sicut ímpetus aquárum multárum inundántium, et emissárum super terram spatiósam. Pédibus conculcábitur coróna supérbiæ ebriórum Ephraim.\par
\switchcolumn
\EN
\noindent Woe to the crown of pride, to the drunkards of Ephraim, and to the fading flower the glory of his joy, who were on the head of the fat valley, staggering with wine. Behold the Lord is mighty and strong, as a storm of hail: a destroying whirlwind, as the violence of many waters overflowing, and sent forth upon a spacious land. The crown of pride of the drunkards of Ephraim shall be trodden under feet.\par
\switchcolumn*

\LA
\begin{Resp}\Rsign Ecce apparébit Dóminus super nubem cándidam, \Star{} Et cum eo Sanctórum míllia: et habébit in vestiménto, et in fémore suo scriptum: Rex regum, et Dóminus dominántium.\end{Resp}
\begin{Resp}\Vsign Apparébit in finem, et non mentiétur; si moram fécerit, exspécta eum, quia véniens véniet.\end{Resp}
\begin{Resp}\Rsign Et cum eo Sanctórum míllia: et habébit in vestiménto, et in fémore suo scriptum: Rex regum, et Dóminus dominántium.\end{Resp}
\switchcolumn
\EN
\begin{Resp}\Rsign Behold, the Lord shall appear upon a white cloud, \Star{} And ten thousand of His saints with Him; and He shall have on His vesture, and on His thigh a name written King of kings, and Lord of lords.\end{Resp}
\begin{Resp}\Vsign He shall appear and not lie; though He tarry, wait for Him, because He will surely come.\end{Resp}
\begin{Resp}\Rsign And ten thousand of His saints with Him; and He shall have on His vesture, and on His thigh a name written King of kings, and Lord of lords.\end{Resp}
\switchcolumn*

\LA
\LessonHead{Lectio II}
\switchcolumn
\EN
\LessonHead{Lesson II}
\switchcolumn*

\LA
\Cite{Isa 28:4-7}
\switchcolumn
\EN
\Cite{Isa 28:4-7}
\switchcolumn*

\LA
\noindent Et erit flos décidens glóriæ exsultatiónis ejus, qui est super vérticem vallis píngujum, quasi temporáneum ante maturitátem autúmni: quod cum aspéxerit videns, statim ut manu tenúerit, devorábit illud. In die illa erit Dóminus exercítuum coróna glóriæ, et sertum exsultatiónis resíduo pópuli sui; et spíritus judícii sedénti super judícium, et fortitúdo reverténtibus de bello ad portam. Verum hi quoque præ vino nesciérunt, et præ ebrietáte erravérunt: sacérdos et prophéta nesciérunt præ ebrietáte, absórpti sunt a vino.\par
\switchcolumn
\EN
\noindent And the fading flower the glory of his joy, who is on the head of the fat valley, shall be as a hasty fruit before the ripeness of autumn: which when he that seeth it shall behold, as soon as he taketh it in his hand, he will eat it up. In that day the Lord of hosts shall be a crown of glory, and a garland of joy to the residue of his people: And a spirit of judgment to him that sitteth in judgment, and strength to them that return out of the battle to the gate. But these also have been ignorant through wine, and through drunkenness have erred: the priest and the prophet have been ignorant through drunkenness, they are swallowed up with wine, they have gone astray in drunkenness, they have not known him that seeth, they have been ignorant of judgment.\par
\switchcolumn*

\LA
\begin{Resp}\Rsign Béthlehem cívitas Dei summi, ex te éxiet Dominátor Israël, et egréssus ejus sicut a princípio diérum æternitátis, et magnificábitur in médio univérsæ terræ: \Star{} Et pax erit in terra nostra, dum vénerit.\end{Resp}
\begin{Resp}\Vsign Loquétur pacem in géntibus, et potéstas ejus a mari usque ad mare.\end{Resp}
\begin{Resp}\Rsign Et pax erit in terra nostra, dum vénerit.\end{Resp}
\switchcolumn
\EN
\begin{Resp}\Rsign Thou, Bethlehem, art the city of the Most High God, out of thee shall He come forth That is to be Ruler in Israel; Whose goings forth have been from of old, from everlasting, and now shall He be great unto the ends of the earth. \Star{} And this Man shall be the peace in our land, when He shall come.\end{Resp}
\begin{Resp}\Vsign He shall speak peace unto the Gentiles, and shall have dominion from sea to sea.\end{Resp}
\begin{Resp}\Rsign And this Man shall be the peace in our land, when He shall come.\end{Resp}
\switchcolumn*

\LA
\LessonHead{Lectio III}
\switchcolumn
\EN
\LessonHead{Lesson III}
\switchcolumn*

\LA
\Cite{Isa 28:16-18}
\switchcolumn
\EN
\Cite{Isa 28:16-18}
\switchcolumn*

\LA
\noindent Idcírco hæc dicit Dóminus Deus: Ecce ego mittam in fundaméntis Sion lápidem, lápidem probátum, angulárem, pretiósum, in fundaménto fundátum. Qui credíderit, non festínet. Et ponam in póndere judícium, et justítiam in mensúra: et subvértet grando spem mendácii, et protectiónem aquæ inundábunt. Et delébitur fœdus vestrum cum morte, et pactum vestrum cum inférno non stabit.\par
\switchcolumn
\EN
\noindent Therefore thus saith the Lord God: Behold I will lay a stone in the foundations of Sion, a tried stone, a corner stone, a precious stone, founded in the foundation. He that believeth, let him not hasten. And I will set judgment in weight, and justice in measure: and hail shall overturn the hope of falsehood: and waters shall overflow its protection. And your league with death shall be abolished, and your covenant with hell shall not stand: when the overflowing scourge shall pass, you shall be trodden down by it.\par
\switchcolumn*

\LA
\begin{Resp}\Rsign Qui ventúrus est, véniet, et non tardábit: et jam non erit timor in fínibus nostris: \Star{} Quóniam ipse est Salvátor noster.\end{Resp}
\begin{Resp}\Vsign Depónet omnes iniquitátes nostras, et proíciet in profúndum maris ómnia peccáta nostra.\end{Resp}
\begin{Resp}\Rsign Quóniam ipse est Salvátor noster.\end{Resp}
\begin{Resp}\Vsign Glória Patri, et Fílio, \Star{} et Spirítui Sancto.\end{Resp}
\begin{Resp}\Rsign Quóniam ipse est Salvátor noster.\end{Resp}
\switchcolumn
\EN
\begin{Resp}\Rsign He That shall come, will come, and will not tarry; and there shall no more be fear in our borders. \Star{} For He is our Saviour.\end{Resp}
\begin{Resp}\Vsign He shall tread down all our iniquities, and cast all our sins into the depths of the sea.\end{Resp}
\begin{Resp}\Rsign For He is our Saviour.\end{Resp}
\begin{Resp}\Vsign Glory be to the Father, and to the Son, \Star{} and to the Holy Ghost.\end{Resp}
\begin{Resp}\Rsign For He is our Saviour.\end{Resp}
\switchcolumn*

\end{paracol}

\Office{Feria III infra Hebdomadam III Adventus}{Tuesday of the Third Week of Advent}{Feria major}
\addcontentsline{toc}{subsection}{Feria III infra Hebdomadam III Adventus\enspace\textit{Tuesday of the Third Week of Advent}}
\begin{paracol}{2}
\LA
\LessonHead{Lectio I}
\switchcolumn
\EN
\LessonHead{Lesson I}
\switchcolumn*

\LA
\LessonTitle{De Isaía Prophéta}
\switchcolumn
\EN
\LessonTitle{Lesson from the book of Isaias}
\switchcolumn*

\LA
\Cite{Isa 30:18-20}
\switchcolumn
\EN
\Cite{Isa 30:18-20}
\switchcolumn*

\LA
\noindent Exspéctat Dóminus ut misereátur vestri, et ídeo exaltábitur parcens vobis: quia Deus judícii Dóminus: beáti omnes qui exspéctant eum. Pópulus enim Sion habitábit in Jerúsalem: plorans nequáquam plorábis, míserans miserébitur tui: ad vocem clamóris tui statim ut audíerit, respondébit tibi. Et dabit vobis Dóminus panem arctum, et aquam brevem: et non fáciet avoláre a te ultra doctórem tuum: et erunt óculi tui vidéntes præceptórem tuum.\par
\switchcolumn
\EN
\noindent The Lord waiteth that he may have mercy on you: and therefore shall he be exalted sparing you: because the Lord is the God of judgment: blessed are all they that wait for him. For the people of Sion shall dwell in Jerusalem: weeping thou shalt not weep, he will surely have pity on thee: at the voice of thy cry, as soon as he shall hear, he will answer thee. And the Lord will give you spare bread, and short water: and will not cause thy teacher to flee away from thee any more, and thy eyes shall see thy teacher.\par
\switchcolumn*

\LA
\begin{Resp}\Rsign Ægýpte, noli flere, quia Dominátor tuus véniet tibi, ante cujus conspéctum movebúntur abýssi, \Star{} Liberáre pópulum suum de manu poténtiæ.\end{Resp}
\begin{Resp}\Vsign Ecce véniet Dóminus exercítuum, Deus tuus cum potestáte magna.\end{Resp}
\begin{Resp}\Rsign Liberáre pópulum suum de manu poténtiæ.\end{Resp}
\switchcolumn
\EN
\begin{Resp}\Rsign Weep not, O Egypt, for the Ruler cometh unto thee, and the depths shall be moved at His presence. \Star{} To deliver His people out of the hand of the mighty.\end{Resp}
\begin{Resp}\Vsign Behold, the Lord of hosts, thy God, cometh with great power.\end{Resp}
\begin{Resp}\Rsign To deliver His people out of the hand of the mighty.\end{Resp}
\switchcolumn*

\LA
\LessonHead{Lectio II}
\switchcolumn
\EN
\LessonHead{Lesson II}
\switchcolumn*

\LA
\Cite{Isa 30:22-25}
\switchcolumn
\EN
\Cite{Isa 30:22-25}
\switchcolumn*

\LA
\noindent Egrédere, dices ei: Et dábitur plúvia sémini tuo, ubicúmque semináveris in terra: et panis frugum terræ erit ubérrimus et pinguis. Pascétur in possessióne tua in die illo agnus spatióse: et tauri tui, et pulli asinórum, qui operántur terram, commístum migma cómedent sicut in área ventilátum est. Et erunt super omnem montem excélsum, et super omnem collem elevátum, rivi curréntium aquárum, in die interfectiónis multórum, cum cecíderint turres.\par
\switchcolumn
\EN
\noindent And thou shalt defile the plates of thy graven things of silver, and the garment of thy molten things of gold, and shalt cast them away as the uncleanness of a menstruous woman. Thou shalt say to it: Get thee hence. And rain shall be given to thy seed, wheresoever thou shalt sow in the land: and the bread of the corn of the land shall be most plentiful, and fat. The lamb in that day shall feed at large in thy possession: And thy oxen, and the ass colts that till the ground, shall eat mingled provender as it was winnowed in the floor. And there shall be upon every high mountain, and upon every elevated hill rivers of running waters in the day of the slaughter of many, when the tower shall fall.\par
\switchcolumn*

\LA
\begin{Resp}\Rsign Prope est ut véniat tempus ejus, et dies ejus non elongabúntur: \Star{} Miserébitur Dóminus Jacob, et Israël salvábitur.\end{Resp}
\begin{Resp}\Vsign Revértere, virgo Israël, revértere ad civitátes tuas.\end{Resp}
\begin{Resp}\Rsign Miserébitur Dóminus Jacob, et Israël salvábitur.\end{Resp}
\switchcolumn
\EN
\begin{Resp}\Rsign Her time is near to come, and her days shall not be prolonged. \Star{} Lord shall have mercy on Jacob, and Israel shall be saved.\end{Resp}
\begin{Resp}\Vsign Turn again, O Virgin of Israel, turn again to thy cities.\end{Resp}
\begin{Resp}\Rsign For the Lord shall have mercy on Judah, and Israel shall be saved.\end{Resp}
\switchcolumn*

\LA
\LessonHead{Lectio III}
\switchcolumn
\EN
\LessonHead{Lesson III}
\switchcolumn*

\LA
\Cite{Isa 30:26-28}
\switchcolumn
\EN
\Cite{Isa 30:26-28}
\switchcolumn*

\LA
\noindent Et erit lux lunæ sicut lux solis, et lux solis erit septemplíciter sicut lux septem diérum in die, qua alligáverit Dóminus vulnus pópuli sui, et percussúram plagæ ejus sanáverit. Ecce nomen Dómini venit de longínquo, ardens furor ejus, et gravis ad portándum: lábia ejus repléta sunt indignatióne, et lingua ejus quasi ignis dévorans. Spíritus ejus velut torrens inúndans usque ad médium colli, ad perdéndas gentes in níhilum, et frenum erróris, quod erat in maxíllis populórum.\par
\switchcolumn
\EN
\noindent And the light of the moon shall be as the light of the sun, and the light of the sun shall be sevenfold, as the light of seven days: in the day when the Lord shall bind up the wound of his people, and shall heal the stroke of their wound. Behold the name of the Lord cometh from afar, his wrath burneth, and is heavy to bear: his lips are filled with indignation, and his tongue as a devouring fire. His breath as a torrent overflowing even to the midst of the neck, to destroy the nations unto nothing, and the bridle of error that was in the jaws of the people.\par
\switchcolumn*

\LA
\begin{Resp}\Rsign Descéndet Dóminus sicut plúvia in vellus: \Star{} Oriétur in diébus ejus justítia, et abundántia pacis.\end{Resp}
\begin{Resp}\Vsign Et adorábunt eum omnes reges, omnes gentes sérvient ei.\end{Resp}
\begin{Resp}\Rsign Oriétur in diébus ejus justítia, et abundántia pacis.\end{Resp}
\begin{Resp}\Vsign Glória Patri, et Fílio, \Star{} et Spirítui Sancto.\end{Resp}
\begin{Resp}\Rsign Oriétur in diébus ejus justítia, et abundántia pacis.\end{Resp}
\switchcolumn
\EN
\begin{Resp}\Rsign The Lord shall come down like rain upon a fleece. \Star{} In His days shall righteousness flourish, and abundance of peace.\end{Resp}
\begin{Resp}\Vsign All the kings of the earth shall fall down before Him, all nations shall serve Him.\end{Resp}
\begin{Resp}\Rsign In His days shall righteousness flourish, and abundance of peace.\end{Resp}
\begin{Resp}\Vsign Glory be to the Father, and to the Son, \Star{} and to the Holy Ghost.\end{Resp}
\begin{Resp}\Rsign In His days shall righteousness flourish, and abundance of peace.\end{Resp}
\switchcolumn*

\end{paracol}

\Office{Feria IV Quattuor Temporum in Adventu}{Ember Wednesday of Advent}{Feria major}
\addcontentsline{toc}{subsection}{Feria IV Quattuor Temporum in Adventu\enspace\textit{Ember Wednesday of Advent}}
\begin{paracol}{2}
\LA
\LessonHead{Lectio I}
\switchcolumn
\EN
\LessonHead{Lesson I}
\switchcolumn*

\LA
\LessonTitle{Léctio sancti Evangélii secúndum Lucam}
\switchcolumn
\EN
\LessonTitle{From the Holy Gospel according to Luke}
\switchcolumn*

\LA
\Cite{Luc 1:26-38}
\switchcolumn
\EN
\Cite{Luke 1:26-38}
\switchcolumn*

\LA
\noindent In illo témpore: Missus est Angelus Gábriel a Deo in civitátem Galilǽæ, cui nomen Názareth, ad Vírginem desponsátam viro, cui nomen erat Joseph, de domo David, et nomen Vírginis María. Et réliqua.\par
\switchcolumn
\EN
\noindent At that time the angel Gabriel was sent from God into a city of Galilee, called Nazareth, To a virgin espoused to a man whose name was Joseph, of the house of David; and the virgin's name was Mary. And so on.\par
\switchcolumn*

\LA
\LessonTitle{Homilía sancti Ambrósii Epíscopi}
\switchcolumn
\EN
\LessonTitle{Homily by St. Ambrose, Bishop (of Milan.)}
\switchcolumn*

\LA
\Source{Liber 2 in Lucam}
\switchcolumn
\EN
\Source{Bk. ii on Luke}
\switchcolumn*

\LA
\noindent Latent quidem divína mystéria, nec fácile, juxta prophéticum dictum, quisquam hóminum potest scire consílium Dei. Sed tamen ex céteris factis, atque præcéptis Dómini Salvatóris póssumus intellégere, et hoc propensióris fuísse consílii, quod ea potíssimum elécta est, ut Dóminum páreret, quæ erat desponsáta viro. Cur autem non ántequam desponsarétur, impléta est? Fortásse ne dicerétur quod concéperat ex adultério.\par
\switchcolumn
\EN
\noindent The mysteries of God are unsearchable, and it is especially declared by a Prophet, that a man can hardly know His counsels. (Wisd. ix. 13.) Nevertheless, some things have been revealed to us, and we may gather from some of the words and works of the Lord our Saviour, that there was a special purpose of God, in the fact that she who was chosen to be the mother of the Lord was espoused to a man. Why did not the power of the Highest overshadow her before she was so espoused? Perhaps it was lest any might blasphemously say that she had conceived in fornication the Holy One.\par
\switchcolumn*

\LA
\begin{Resp}\Rsign Clama in fortitúdine, qui annúntias pacem in Jerúsalem: \Star{} Dic civitátibus Judæ, et habitatóribus Sion: Ecce Deus noster, quem exspectabámus, advéniet.\end{Resp}
\begin{Resp}\Vsign Supra montem excélsum ascénde tu, qui evangelízas Sion, exálta in fortitúdine vocem tuam.\end{Resp}
\begin{Resp}\Rsign Dic civitátibus Judæ, et habitatóribus Sion: Ecce Deus noster, quem exspectabámus, advéniet.\end{Resp}
\switchcolumn
\EN
\begin{Resp}\Rsign O thou that bringest good tidings of peace to Jerusalem, lift up thy voice with strength! \Star{} Say unto the cities of Judah, and to the inhabitants of Jerusalem: Behold, our God will come, for Whom we waited.\end{Resp}
\begin{Resp}\Vsign O thou that tellest good tidings to Zion, get thee up into the high mountain, lift up thy voice with strength.\end{Resp}
\begin{Resp}\Rsign Say unto the cities of Judah, and to the inhabitants of Jerusalem: Behold, our God will come, for Whom we waited.\end{Resp}
\switchcolumn*

\LA
\LessonHead{Lectio II}
\switchcolumn
\EN
\LessonHead{Lesson II}
\switchcolumn*

\LA
\noindent Et ingréssus ad eam Angelus. Disce vírginem móribus, disce vírginem verecúndia, disce oráculo, disce mystério. Trepidáre vírginum est, et ad omnes viri ingréssus pavére, omnes viri affátus veréri. Discant mulíeres propósitum pudóris imitári. Sola in penetrálibus, quam nemo virórum víderit, solus Angelus repérerit: sola sine cómite, sola sine teste, ne quo degénere depravarétur affátu, ab Angelo salutátur.\par
\switchcolumn
\EN
\noindent And the angel came in unto her. Let us learn from this Virgin how to bear ourselves, let us learn her modesty, let us learn by her devout utterance, above all let us learn by the holy mystery enacted. It is the part of a maiden to be timid, to avoid the advances of men, and to shrink from men's addresses. Would that our women would learn from the example of modesty here set before us. She upon whom the stare of men had never been fixed was alone in her chamber, and was found only by an angel. There was neither companion nor witness there, that what passed might not be debased in gossip and the angel saluted her.\par
\switchcolumn*

\LA
\begin{Resp}\Rsign Oriétur stella ex Jacob, et exsúrget homo de Israël, et confrínget omnes duces alienigenárum: \Star{} Et erit omnis terra posséssio ejus.\end{Resp}
\begin{Resp}\Vsign Adorábunt eum omnes reges terræ, omnes gentes sérvient ei.\end{Resp}
\begin{Resp}\Rsign Et erit omnis terra posséssio ejus.\end{Resp}
\switchcolumn
\EN
\begin{Resp}\Rsign There shall come a Star out of Jacob, and a Man shall rise out of Israel, and shall smite through all the princes of the aliens. \Star{} And all the earth shall be His possession.\end{Resp}
\begin{Resp}\Vsign All kings shall fall down before Him, all nations shall serve Him.\end{Resp}
\begin{Resp}\Rsign And all the earth shall be His possession.\end{Resp}
\switchcolumn*

\LA
\LessonHead{Lectio III}
\switchcolumn
\EN
\LessonHead{Lesson III}
\switchcolumn*

\LA
\noindent Tanti namque mandáti mystérium non hóminis fuit, sed Angeli ore proméndum. Hódie primum audítur: Spíritus Sanctus supervéniet in te. Et audítur, et créditur. Dénique, Ecce, inquit, ancílla Dómini: contíngat mihi secúndum verbum tuum. Vide humilitátem, vide devotiónem. Ancíllam se dicit Dómini, quæ mater elígitur: nec repentíno exaltáta promísso est.\par
\switchcolumn
\EN
\noindent The message of God to the Virgin was a mystery, which it was not lawful for the mouth of men, but only of angels, to utter. For the first time on earth the words are spoken The Holy Ghost shall come upon thee. The holy maiden heareth, and believeth. At length she saith Behold the handmaid of the Lord be it unto me according to thy word. Here is an example of lowliness, here is a pattern of true devotion. At the very moment that she is told she is chosen to be the mother of the Lord she at once declareth herself His handmaid. The knowledge that she was Mother of God caused in the heart of Mary only an act of humility.\par
\switchcolumn*

\LA
\begin{Resp}\Rsign Modo véniet Dominátor Dóminus: \Star{} Et nomen ejus Emmánuel vocábitur.\end{Resp}
\begin{Resp}\Vsign Oriétur in diébus ejus justítia, et abundántia pacis.\end{Resp}
\begin{Resp}\Rsign Et nomen ejus Emmánuel vocábitur.\end{Resp}
\begin{Resp}\Vsign Glória Patri, et Fílio, \Star{} et Spirítui Sancto.\end{Resp}
\begin{Resp}\Rsign Et nomen ejus Emmánuel vocábitur.\end{Resp}
\switchcolumn
\EN
\begin{Resp}\Rsign The Lord, the Ruler, cometh quickly. \Star{} And His name shall be called Emmanuel.\end{Resp}
\begin{Resp}\Vsign In His days shall righteousness flourish, and abundance of peace.\end{Resp}
\begin{Resp}\Rsign And His name shall be called Emmanuel.\end{Resp}
\begin{Resp}\Vsign Glory be to the Father, and to the Son, \Star{} and to the Holy Ghost.\end{Resp}
\begin{Resp}\Rsign And His name shall be called Emmanuel.\end{Resp}
\switchcolumn*

\end{paracol}

\Office{Feria V infra Hebdomadam III Adventus}{Thursday of the Third Week of Advent}{Feria major}
\addcontentsline{toc}{subsection}{Feria V infra Hebdomadam III Adventus\enspace\textit{Thursday of the Third Week of Advent}}
\begin{paracol}{2}
\LA
\LessonHead{Lectio I}
\switchcolumn
\EN
\LessonHead{Lesson I}
\switchcolumn*

\LA
\LessonTitle{De Isaía Prophéta}
\switchcolumn
\EN
\LessonTitle{Lesson from the book of Isaias}
\switchcolumn*

\LA
\Cite{Isa 33:1-2}
\switchcolumn
\EN
\Cite{Isa 33:1-2}
\switchcolumn*

\LA
\noindent Væ qui prædáris, nonne et ipse prædáberis? et qui spernis, nonne et ipse spernéris? cum consummáveris deprædatiónem, deprædáberis: cum fatigátus desíeris contémnere, contemnéris. Dómine, miserére nostri, te enim exspectávimus: esto brácchium nostrum in mane, et salus nostra in témpore tribulatiónis.\par
\switchcolumn
\EN
\noindent Woe to thee that spoilest, shalt not thou thyself also be spoiled? and thou that despisest, shalt not thyself also be despised? when thou shalt have made an end of spoiling, thou shalt be spoiled: when being wearied thou shalt cease to despise, thou shalt be despised. O Lord, have mercy on us: for we have waited for thee: be thou our arm in the morning, and our salvation in the time of trouble.\par
\switchcolumn*

\LA
\begin{Resp}\Rsign Egrediétur Dóminus, et præliábitur contra gentes: \Star{} Et stabunt pedes ejus supra montes olivárum ad Oriéntem.\end{Resp}
\begin{Resp}\Vsign Et elevábitur supra omnes colles, et fluent ad eum omnes gentes.\end{Resp}
\begin{Resp}\Rsign Et stabunt pedes ejus supra montes olivárum ad Oriéntem.\end{Resp}
\switchcolumn
\EN
\begin{Resp}\Rsign The Lord shall go forth and fight against the nations. \Star{} And His feet shall stand upon the mount of Olives on the east.\end{Resp}
\begin{Resp}\Vsign And it shall be exalted above the hills, and all nations shall flow unto it.\end{Resp}
\begin{Resp}\Rsign And His feet shall stand upon the mount of Olives on the east.\end{Resp}
\switchcolumn*

\LA
\LessonHead{Lectio II}
\switchcolumn
\EN
\LessonHead{Lesson II}
\switchcolumn*

\LA
\Cite{Isa 33:3-6}
\switchcolumn
\EN
\Cite{Isa 33:3-6}
\switchcolumn*

\LA
\noindent A voce Angeli fugérunt pópuli, et ab exaltatióne tua dispérsæ sunt gentes. Et congregabúntur spólia vestra sicut collígitur bruchus, velut cum fossæ plenæ fúerint de eo. Magnificátus est Dóminus, quóniam habitávit in excélso: implévit Sion judício et justítia. Et erit fides in tempóribus tuis: divítiæ salútis, sapiéntia, et sciéntia: timor Dómini ipse est thesáurus ejus.\par
\switchcolumn
\EN
\noindent At the voice of the angel the people fled, and at the lifting up thyself the nations are scattered. And your spoils shall be gathered together as the locusts are gathered, as when the ditches are full of them. The Lord is magnified, for he hath dwelt on high: he hath filled Sion with judgment and justice. And there shall be faith in thy times: riches of salvation, wisdom and knowledge: the fear of the Lord is his treasure.\par
\switchcolumn*

\LA
\begin{Resp}\Rsign Præcúrsor pro nobis ingréditur Agnus sine mácula, \Star{} Secúndum órdinem Melchísedech Póntifex factus in ætérnum et in sǽculum sǽculi.\end{Resp}
\begin{Resp}\Vsign Ipse est rex justítiæ, cujus generátio non habet finem.\end{Resp}
\begin{Resp}\Rsign Secúndum órdinem Melchísedech Póntifex factus in ætérnum et in sǽculum sǽculi.\end{Resp}
\switchcolumn
\EN
\begin{Resp}\Rsign The Fore-runner is for us entered, even the Lamb without spot \Star{} Made an High Priest for ever after the order of Melchisedek.\end{Resp}
\begin{Resp}\Vsign This is that King of Righteousness without descent, nor end of life.\end{Resp}
\begin{Resp}\Rsign Made an High Priest for ever after the order of Melchisedek.\end{Resp}
\switchcolumn*

\LA
\LessonHead{Lectio III}
\switchcolumn
\EN
\LessonHead{Lesson III}
\switchcolumn*

\LA
\Cite{Isa 33:14-17}
\switchcolumn
\EN
\Cite{Isa 33:14-17}
\switchcolumn*

\LA
\noindent Contérriti sunt in Sion peccatóres, possédit tremor hypócritas. Quis póterit habitáre de vobis cum igne devoránte? quis habitábit ex vobis cum ardóribus sempitérnis? Qui ámbulat in justítiis, et lóquitur veritátem, qui próicit avarítiam ex calúmnia, et éxcutit manus suas ab omni múnere, qui obtúrat aures suas ne áudiat sánguinem, et claudit óculos suos ne vídeat malum. Iste in excélsis habitábit muniménta saxórum sublímitas ejus: panis ei datus est, aquæ ejus fidéles sunt. Regem in decóre suo vidébunt óculi ejus, cernent terram de longe.\par
\switchcolumn
\EN
\noindent The sinners in Sion are afraid, trembling hath seized upon the hypocrites. Which of you can dwell with devouring fire? which of you shall dwell with everlasting burnings? He that walketh in justices, and speaketh truth, that casteth away avarice by oppression, and shaketh his hands from all bribes, that stoppeth his ears lest he hear blood, and shutteth his eyes that he may see no evil. He shall dwell on high, the fortifications of rocks shall be his highness: bread is given him, his waters are sure. His eyes shall see the king in his beauty, they shall see the land far off.\par
\switchcolumn*

\LA
\begin{Resp}\Rsign Vidébunt gentes justum tuum, et cuncti reges ínclitum tuum: \Star{} Et vocábitur tibi nomen novum, quod os Dómini nominávit.\end{Resp}
\begin{Resp}\Vsign Et eris coróna glóriæ in manu Dómini, et diadéma regni in manu Dei tui.\end{Resp}
\begin{Resp}\Rsign Et vocábitur tibi nomen novum, quod os Dómini nominávit.\end{Resp}
\begin{Resp}\Vsign Glória Patri, et Fílio, \Star{} et Spirítui Sancto.\end{Resp}
\begin{Resp}\Rsign Et vocábitur tibi nomen novum, quod os Dómini nominávit.\end{Resp}
\switchcolumn
\EN
\begin{Resp}\Rsign The Gentiles shall see thy Righteous One, and all kings thy Glorious one. \Star{} And thou shalt be called by a new name, which the mouth of the Lord hath named.\end{Resp}
\begin{Resp}\Vsign Thou shalt also be a crown of glory in the hand of the Lord, and a royal diadem in the hand of thy God.\end{Resp}
\begin{Resp}\Rsign And thou shalt be called by a new name, which the mouth of the Lord hath named.\end{Resp}
\begin{Resp}\Vsign Glory be to the Father, and to the Son, \Star{} and to the Holy Ghost.\end{Resp}
\begin{Resp}\Rsign And thou shalt be called by a new name, which the mouth of the Lord hath named.\end{Resp}
\switchcolumn*

\end{paracol}

\Office{Feria VI Quattuor Temporum in Adventu}{Ember Friday of Advent}{Feria major}
\addcontentsline{toc}{subsection}{Feria VI Quattuor Temporum in Adventu\enspace\textit{Ember Friday of Advent}}
\begin{paracol}{2}
\LA
\LessonHead{Lectio I}
\switchcolumn
\EN
\LessonHead{Lesson I}
\switchcolumn*

\LA
\LessonTitle{Léctio sancti Evangélii secúndum Lucam}
\switchcolumn
\EN
\LessonTitle{Continuation of the Holy Gospel according to Luke}
\switchcolumn*

\LA
\Cite{Luc 1:39-47}
\switchcolumn
\EN
\Cite{Luke 1:39-47}
\switchcolumn*

\LA
\noindent In illo témpore: Exsúrgens María ábiit in montána cum festinatióne in civitátem Juda: et intrávit in domum Zacharíæ, et salutávit Elísabeth. Et réliqua.\par
\switchcolumn
\EN
\noindent In that time, Mary, rising up, went into the hill country with haste into a city of Juda. And she entered into the house of Zachary, and saluted Elizabeth. And so on.\par
\switchcolumn*

\LA
\LessonTitle{Homilía sancti Ambrósii Epíscopi}
\switchcolumn
\EN
\LessonTitle{Homily by St. Ambrose, Bishop (of Milan)}
\switchcolumn*

\LA
\Source{Lib. 2 in Luc. cap. 1 post init.}
\switchcolumn
\EN
\Source{Commentary on Luke, Bk. ii. c}
\switchcolumn*

\LA
\noindent Morále est ómnibus, ut qui fidem éxigunt, fidem ástruant. Et ídeo Angelus, cum abscóndita nuntiáret, ut fidem astrúeret exémplo, senióris féminæ sterilísque concéptum Vírgini Maríæ nuntiávit: ut possíbile Deo omne, quod ei placúerit, asséreret. Ubi audívit hoc María, non quasi incrédula de oráculo, nec quasi incérta de núntio, nec quasi dúbitans de exémplo, sed quasi læta pro voto, religiósa pro offício, festína pro gáudio, in montána perréxit. Quo enim jam Deo plena, nisi ad superióra cum festinatióne conténderet? Nescit tarda molímina Sancti Spíritus grátia.\par
\switchcolumn
\EN
\noindent When any one asketh another for credence, he is bound to give some reasonable ground. And so the Angel, when he announced to Mary the counsel of God, gave, as a proof, the conception of Elizabeth, then aged and barren, that Mary might perceive, by this example, that with God nothing is impossible. When the holy Virgin had heard it, she arose and went to visit her cousin. She did not go to see if what she had heard was true, because she did not believe God, or because she knew not who the messenger had been, or yet because she doubted the fact adduced in proof. She went joyfully as one who hath received a mercy in answer to his vow goeth to pay the same. She went with devotion, as a godly person goeth to execute a religious duty. She went into the hill country in joyful haste. And is it not something that she went up into the hills? God was already in her womb, and her feeling bore her continually upward. The grace of the Holy Spirit knoweth no slow working.\par
\switchcolumn*

\LA
\begin{Resp}\Rsign Emítte Agnum, Dómine, Dominatórem terræ, \Star{} De Petra desérti ad montem fíliæ Sion.\end{Resp}
\begin{Resp}\Vsign Osténde nobis, Dómine, misericórdiam tuam, et salutáre tuum da nobis.\end{Resp}
\begin{Resp}\Rsign De Petra desérti ad montem fíliæ Sion.\end{Resp}
\switchcolumn
\EN
\begin{Resp}\Rsign Send forth the Lamb, O Lord, the Ruler of the land; \Star{} From the rock in the wilderness unto the mount of the daughter of Zion.\end{Resp}
\begin{Resp}\Vsign Show us thy mercy, O Lord, and grant us thy salvation.\end{Resp}
\begin{Resp}\Rsign From the rock in the wilderness unto the mount of the daughter of Zion.\end{Resp}
\switchcolumn*

\LA
\LessonHead{Lectio II}
\switchcolumn
\EN
\LessonHead{Lesson II}
\switchcolumn*

\LA
\noindent Díscite et vos, sanctæ mulíeres, sedulitátem, quam prægnántibus debeátis exhibére cognátis. Maríam, quæ ante sola in íntimis penetrálibus versabátur, non a público virginitátis pudor, non ab stúdio aspéritas móntium, non ab offício prolíxitas itíneris retardávit. In montána Virgo cum festinatióne, Virgo offícii memor, injúriæ ímmemor, afféctu urgénte non sexu, relícta perréxit domo. Díscite, vírgines, non circumcursáre per aliénas ædes, non demorári in platéis, non áliquos in público miscére sermónes. María in domo sera, festína in público, mansit apud cognátam suam tribus ménsibus.\par
\switchcolumn
\EN
\noindent Godly women will learn from the example of the Mother of God to take a tender care of their kinswomen who are with child. In pursuance of this charity, Mary, who had hitherto remained alone at home, was not deterred by her maidenly shyness from entering on a public journey; she faced for this end the hardships of mountain travelling; and encountered with a sense of duty the weary length of the way. The Virgin left her home, and went into the hill country with haste, unmindful of the trouble, and remembering only the office to which her cousinly love prompted her, in spite of the delicacy of her sex. Maidens will learn from her not to idle about from house to house, to loiter in the streets, nor to take part in conversations in public. Mary, as she was hasteful to pass through the public roads, so was she slow again to enter on them she abode with her cousin about three months.\par
\switchcolumn*

\LA
\begin{Resp}\Rsign Roráte, cæli, désuper, et nubes pluant justum: \Star{} Aperiátur terra, et gérminet Salvatórem.\end{Resp}
\begin{Resp}\Vsign Emítte Agnum, Dómine, Dominatórem terræ, de Petra desérti ad montem fíliæ Sion.\end{Resp}
\begin{Resp}\Rsign Aperiátur terra, et gérminet Salvatórem.\end{Resp}
\switchcolumn
\EN
\begin{Resp}\Rsign Drop down, ye heavens, from above, and let the skies pour down the Righteous One. \Star{} Let the earth open, and let her bring forth the Saviour.\end{Resp}
\begin{Resp}\Vsign Send forth the Lamb, O Lord, the Ruler of the land, from the rock in the wilderness unto the mount of the daughter of Zion.\end{Resp}
\begin{Resp}\Rsign Let the earth open, and let her bring forth the Saviour.\end{Resp}
\switchcolumn*

\LA
\LessonHead{Lectio III}
\switchcolumn
\EN
\LessonHead{Lesson III}
\switchcolumn*

\LA
\noindent Didicístis, vírgines, pudórem Maríæ: díscite humilitátem. Venit propínqua ad próximam, júnior ad seniórem: nec solum venit, sed étiam prior salutávit. Decet enim, ut quanto cástior virgo, tanto humílior sit. Nóverit deférre senióribus. Sit magístra humilitátis, in qua est proféssio castitátis. Est et causa pietátis, est étiam norma doctrínæ. Contuéndum est enim, quia supérior venit ad inferiórem, ut inférior adjuvétur: María ad Elísabeth, Christus ad Joánnem.\par
\switchcolumn
\EN
\noindent As the modesty of Mary is a pattern for the imitation of all maidens, so also is her humility. She went to see Elizabeth, like one cousin going to visit another, and as the younger to the elder. Not only did she first go, but she first saluted Elizabeth. Now, the purer a virgin is, the humbler ought she to be. She will know how to submit herself to her elders. She that professeth chastity ought to be a very mistress of humility. Lowly-mindedness is at once the very ground in which devotion groweth, and the first and principal rule of its teaching. In this act of the Virgin then we see the greater going to visit and to succour the lesser: Mary to Elizabeth, Christ to John.\par
\switchcolumn*

\LA
\begin{Resp}\Rsign Germinavérunt campi erémi germen odóris Israël: quia ecce Deus noster cum virtúte véniet, \Star{} Et splendor ejus cum eo.\end{Resp}
\begin{Resp}\Vsign Ex Sion spécies decóris ejus: Deus noster maniféste véniet.\end{Resp}
\begin{Resp}\Rsign Et splendor ejus cum eo.\end{Resp}
\begin{Resp}\Vsign Glória Patri, et Fílio, \Star{} et Spirítui Sancto.\end{Resp}
\begin{Resp}\Rsign Et splendor ejus cum eo.\end{Resp}
\switchcolumn
\EN
\begin{Resp}\Rsign The waste places have brought forth sweet-smelling buds for Israel; for, behold, our God will come with power. \Star{} And His brightness is with Him.\end{Resp}
\begin{Resp}\Vsign Out of Zion the perfection of beauty, our God shall come manifestly.\end{Resp}
\begin{Resp}\Rsign And His brightness is with Him.\end{Resp}
\begin{Resp}\Vsign Glory be to the Father, and to the Son, \Star{} and to the Holy Ghost.\end{Resp}
\begin{Resp}\Rsign And His brightness is with Him.\end{Resp}
\switchcolumn*

\end{paracol}

\Office{Sabbato Quattuor Temporum in Adventu}{Ember Saturday of Advent}{Feria major}
\addcontentsline{toc}{subsection}{Sabbato Quattuor Temporum in Adventu\enspace\textit{Ember Saturday of Advent}}
\begin{paracol}{2}
\LA
\LessonHead{Lectio I}
\switchcolumn
\EN
\LessonHead{Lesson I}
\switchcolumn*

\LA
\LessonTitle{Léctio sancti Evangélii secúndum Lucam}
\switchcolumn
\EN
\LessonTitle{Continuation of the Holy Gospel according to Luke}
\switchcolumn*

\LA
\Cite{Luc 3:1-6}
\switchcolumn
\EN
\Cite{Luke 3:1-6}
\switchcolumn*

\LA
\noindent Anno quintodécimo impérii Tibérii Cǽsaris, procuránte Póntio Piláto Judǽam. Et réliqua.\par
\switchcolumn
\EN
\noindent In the fifteenth year of the reign of Tiberius Caesar, Pontius Pilate being governor of Judea. And so on.\par
\switchcolumn*

\LA
\LessonTitle{Homilía sancti Gregórii Papæ}
\switchcolumn
\EN
\LessonTitle{Homily by Pope St. Gregory (the Great.)}
\switchcolumn*

\LA
\Source{Homilia 20 in Evangelia}
\switchcolumn
\EN
\Source{10th on the Gospels.}
\switchcolumn*

\LA
\noindent Redemptóris nostri Præcúrsor, quo témpore prædicatiónis offícium accéperit, memoráto Románæ reipúblicæ príncipe, et Judǽæ régibus, designátur. Quia enim illum prædicáre veniébat, qui et ex Judǽa quosdam, et multos ex géntibus redemptúrus erat: per regem géntium et príncipes Judæórum prædicatiónis ejus témpora designántur. Quia autem gentílitas colligénda erat, et Judǽa pro culpa perfídiæ dispergénda, ipsa quoque descríptio terréni principátus osténdit: quóniam et in Romána república unus præfuísse descríbitur, et in Judǽæ regno per quartam partem plúrimi principabántur.\par
\switchcolumn
\EN
\noindent The date, at which the Fore-runner of our Redeemer entered on his public office of preaching, is indicated to us by the name of the ruler of the Roman Commonwealth, and by those of the princes of Palestine. The time of his preaching is indicated by these names, because he came as the Fore-runner of Him Who was to be the Redeemer of some Jews and many Gentiles. Moreover in the enumeration of these worldly monarchs there is a foreshadowing of the fact, that the Gentiles were about to be gathered into one, and the Jews to be scattered abroad in punishment of their unbelief; in the whole heathen Commonwealth we find the title of one Emperor, but in the small kingdom of Judaea are mentioned four masters.\par
\switchcolumn*

\LA
\begin{Resp}\Rsign Egrediétur virga de radíce Jesse, et flos de radíce ejus ascéndet: \Star{} Et erit justítia cíngulum lumbórum ejus, et fides cinctórium renum ejus.\end{Resp}
\begin{Resp}\Vsign Et requiéscet super eum spíritus Dómini: spíritus sapiéntiæ, et intelléctus: spíritus consílii, et fortitúdinis.\end{Resp}
\begin{Resp}\Rsign Et erit justítia cíngulum lumbórum ejus, et fides cinctórium renum ejus.\end{Resp}
\switchcolumn
\EN
\begin{Resp}\Rsign There shall come forth a rod out of the stem of Jesse, and a flower shall grow out of his roots. \Star{} And righteousness shall be the girdle of his loins, and faithfulness the girdle of his reins.\end{Resp}
\begin{Resp}\Vsign And the Spirit of the Lord shall rest upon him the spirit of wisdom and understanding the spirit of counsel and might.\end{Resp}
\begin{Resp}\Rsign And righteousness shall be the girdle of his loins, and faithfulness the girdle of his reins.\end{Resp}
\switchcolumn*

\LA
\LessonHead{Lectio II}
\switchcolumn
\EN
\LessonHead{Lesson II}
\switchcolumn*

\LA
\noindent Voce étenim nostri Redemptóris dícitur: Omne regnum in seípsum divisum desolábitur. Liquet ergo, quod ad finem regni Judǽa pervénerat, quæ tot régibus divísa subjacébat. Apte quoque non solum quibus régibus, sed étiam quibus sacerdótibus actum sit, demonstrátur: et quia illum Joánnes Baptísta prædicáret, qui simul Rex et Sacérdos exsísteret, Lucas Evangelísta prædicatiónis ejus témpora per regnum et sacerdótium designávit.\par
\switchcolumn
\EN
\noindent The blessed voice of the Saviour itself hath said, Every kingdom divided against itself is brought to desolation (Luke xi. 17.) And we may well look for the ruin of the Jewish state when we see it divided among so many rulers. We observe likewise that the names of the reigning priests as well as kings are given. The Evangelist Luke hath left on record the chiefs both of the monarchy and of the priesthood who held office when John the Baptist began to preach, because John preached Him Who is at once our Priest and our King.\par
\switchcolumn*

\LA
\begin{Resp}\Rsign Radix Jesse, qui exsúrget judicáre gentes, in eum gentes sperábunt: \Star{} Et erit nomen ejus benedíctum in sǽcula.\end{Resp}
\begin{Resp}\Vsign Super ipsum continébunt reges os suum, ipsum gentes deprecabúntur.\end{Resp}
\begin{Resp}\Rsign Et erit nomen ejus benedíctum in sǽcula.\end{Resp}
\switchcolumn
\EN
\begin{Resp}\Rsign Behold, the root of Jesse that shall arise to bring forth judgment to the Gentiles, in him shall the Gentiles trust. \Star{} And his name shall be blessed for ever.\end{Resp}
\begin{Resp}\Vsign The Kings shall shut their mouths at him, to him shall the Gentiles seek.\end{Resp}
\begin{Resp}\Rsign And his name shall be blessed for ever.\end{Resp}
\switchcolumn*

\LA
\LessonHead{Lectio III}
\switchcolumn
\EN
\LessonHead{Lesson III}
\switchcolumn*

\LA
\noindent Et venit in omnem regiónem Jordánis, prǽdicans baptísmum pœniténtiæ in remissiónem peccatórum. Cunctis legéntibus liquet, quia Joánnes non solum baptísmum pœniténtiæ prædicávit, verum étiam quibúsdam dedit: sed tamen baptísmum suum in remissiónem peccatórum dare non pótuit. Remíssio étenim peccatórum in solo nobis baptísmo Christi tribúitur. Notándum ítaque, quod dícitur: Prǽdicans baptísmum pœniténtiæ in remissiónem peccatórum: quóniam baptísmum, quod peccáta sólveret, quia dare non póterat, prædicábat: ut sicut incarnátum Verbum Patris præcurrébat verbo prædicatiónis, ita baptísmum pœniténtiæ, quo peccáta solvúntur, præcúrreret suo baptísmate, quo peccáta solvi non possunt.\par
\switchcolumn
\EN
\noindent And he came into all the country about Jordan, preaching the baptism of repentance for the remission of sins. It is evident from these words that John the Baptist not only preached, but also administered the baptism of repentance, and yet that baptism of repentance which he gave, was not really a baptism for the remission of sins. For there is only one baptism for the remission of sins, and that is our Christian baptism. It is worthy of note here that the words used are, preaching the baptism of repentance for the remission of sins, for he himself owned that his baptism was not the true baptism that washes away sin. Even as the Eternal Word of God made Flesh was greater than the preacher that went before Him, so was His holy baptism, by which our sins are washed away, far greater than that baptism of repentance which the Fore-runner preached, and which could never wash away sin.\par
\switchcolumn*

\LA
\begin{Resp}\Rsign Veni, Dómine, et noli tardáre: reláxa facínora plebi tuæ, \Star{} Et révoca dispérsos in terram suam.\end{Resp}
\begin{Resp}\Vsign Excíta, Dómine, poténtiam tuam, et veni, ut salvos fácias nos.\end{Resp}
\begin{Resp}\Rsign Et révoca dispérsos in terram suam.\end{Resp}
\begin{Resp}\Vsign Glória Patri, et Fílio, \Star{} et Spirítui Sancto.\end{Resp}
\begin{Resp}\Rsign Et révoca dispérsos in terram suam.\end{Resp}
\switchcolumn
\EN
\begin{Resp}\Rsign O Lord, come, and make no tarrying loosen the bonds of thy people. \Star{} And gather together into their own land them that are scattered abroad.\end{Resp}
\begin{Resp}\Vsign Stir up, O Lord, thy power, and come among us, to save us.\end{Resp}
\begin{Resp}\Rsign And gather together into their own land them that are scattered abroad.\end{Resp}
\begin{Resp}\Vsign Glory be to the Father, and to the Son, \Star{} and to the Holy Ghost.\end{Resp}
\begin{Resp}\Rsign And gather together into their own land them that are scattered abroad.\end{Resp}
\switchcolumn*

\end{paracol}

\Office{Dominica IV Adventus}{Fourth Sunday of Advent}{Semiduplex II classis}
\addcontentsline{toc}{subsection}{Dominica IV Adventus\enspace\textit{Fourth Sunday of Advent}}
\begin{paracol}{2}
\LA
\LessonHead{Lectio I}
\switchcolumn
\EN
\LessonHead{Lesson I}
\switchcolumn*

\LA
\LessonTitle{De Isaía Prophéta}
\switchcolumn
\EN
\LessonTitle{Lesson from the book of Isaias}
\switchcolumn*

\LA
\Cite{Isa 35:1-7}
\switchcolumn
\EN
\Cite{Isa 35:1-7}
\switchcolumn*

\LA
\noindent Lætábitur desérta et ínvia, et exsultábit solitúdo, et florébit quasi lílium. Gérminans germinábit, et exsultábit lætabúnda et laudans: glória Líbani data est ei: decor Carméli et Saron, ipsi vidébunt glóriam Dómini, et decórem Dei nostri. Confortáte manus dissolútas, et génua debília roboráte. Dícite pusillánimis: Confortámini, et nolíte timére: ecce Deus vester ultiónem addúcet retributiónis: Deus ipse véniet, et salvábit vos. Tunc aperiéntur óculi cæcórum, et aures surdórum patébunt. Tunc sáliet sicut cervus claudus, et apérta erit lingua mutórum: quia scissæ sunt in desérto aquæ, et torréntes in solitúdine. Et quæ erat árida, erit in stagnum, et sítiens in fontes aquárum.\par
\switchcolumn
\EN
\noindent The land that was desolate and impassable shall be glad, and the wilderness shall rejoice, and shall flourish like the lily. It shall bud forth and blossom, and shall rejoice with joy and praise: the glory of Libanus is given to it: the beauty of Carmel, and Saron, they shall see the glory of the Lord, and the beauty of our God. Strengthen ye the feeble hands, and confirm the weak knees. Say to the fainthearted: Take courage, and fear not: behold your God will bring the revenge of recompense: God himself will come and will save you. Then shall the eyes of the blind be opened, and the ears of the deaf shall be unstopped. Then shall the lame man leap as a hart, and the tongue of the dumb shall be free: for waters are broken out in the desert, and streams in the wilderness. And that which was dry land, shall become a pool, and the thirsty land springs of water.\par
\switchcolumn*

\LA
\begin{Resp}\Rsign Cánite tuba in Sion, vocáte gentes, annuntiáte pópulis, et dícite: \Star{} Ecce Deus Salvátor noster advéniet.\end{Resp}
\begin{Resp}\Vsign Annuntiáte, et audítum fácite: loquímini, et clamáte.\end{Resp}
\begin{Resp}\Rsign Ecce Deus Salvátor noster advéniet.\end{Resp}
\switchcolumn
\EN
\begin{Resp}\Rsign Blow ye the trumpet in Zion, call together the nations, tell it out among the people, and say: \Star{} Behold, God our Saviour cometh.\end{Resp}
\begin{Resp}\Vsign Tell it out and make it to be heard; speak aloud and cry.\end{Resp}
\begin{Resp}\Rsign Behold, God our Saviour cometh.\end{Resp}
\switchcolumn*

\LA
\LessonHead{Lectio II}
\switchcolumn
\EN
\LessonHead{Lesson II}
\switchcolumn*

\LA
\Cite{Isa 35:7-10}
\switchcolumn
\EN
\Cite{Isa 35:7-10}
\switchcolumn*

\LA
\noindent In cubílibus, in quibus prius dracónes habitábant, oriétur viror cálami et junci. Et erit ibi sémita et via, et via sancta vocábitur: non transíbit per eam pollútus, et hæc erit vobis dirécta via, ita ut stulti non errent per eam. Non erit ibi leo, et mala béstia non ascéndet per eam, nec inveniétur ibi: et ambulábunt, qui liberáti fúerint. Et redémpti a Dómino converténtur, et vénient in Sion cum laude: et lætítia sempitérna super caput eórum: gáudium et lætítiam obtinébunt, et fúgiet dolor et gémitus.\par
\switchcolumn
\EN
\noindent In the dens where dragons dwell before, shall rise up the verdure of the reed and the bulrush. And a path and a way shall be there, and it shall be called the holy way: the unclean shall not pass over it, and this shall be unto you a straight way, so that fools shall not err therein. No lion shall be there, nor shall any mischievous beast go up by it, nor be found there: but they shall walk there that shall be delivered. And the redeemed of the Lord shall return, and shall come into Sion with praise, and everlasting joy shall be upon their heads: they shall obtain joy and gladness, and sorrow and mourning shall flee away.\par
\switchcolumn*

\LA
\begin{Resp}\Rsign Non auferétur sceptrum de Juda, et dux de fémore ejus, donec véniat qui mitténdus est: \Star{} Et ipse erit exspectátio géntium.\end{Resp}
\begin{Resp}\Vsign Pulchrióres sunt óculi ejus vino, et dentes ejus lacte candidióres.\end{Resp}
\begin{Resp}\Rsign Et ipse erit exspectátio géntium.\end{Resp}
\switchcolumn
\EN
\begin{Resp}\Rsign The sceptre shall not depart from Judah, nor the law-giver from his loins, until he that shall be sent cometh. \Star{} And unto him shall the longing of the Gentiles be.\end{Resp}
\begin{Resp}\Vsign His eyes shall be bright with wine, and his teeth white with milk.\end{Resp}
\begin{Resp}\Rsign And unto him shall the desire of the Gentiles be.\end{Resp}
\switchcolumn*

\LA
\LessonHead{Lectio III}
\switchcolumn
\EN
\LessonHead{Lesson III}
\switchcolumn*

\LA
\Cite{Isa 41:1-4}
\switchcolumn
\EN
\Cite{Isa 41:1-4}
\switchcolumn*

\LA
\noindent Táceant ad me ínsulæ, et gentes mutent fortitúdinem: accédant, et tunc loquántur, simul ad judícium propinquémus. Quis suscitávit ab Oriénte justum, vocávit eum ut sequerétur se? dabit in conspéctu ejus gentes et reges obtinébit: dabit quasi púlverem gládio ejus, sicut stípulam vento raptam árcui ejus. Persequétur eos, transíbit in pace, sémita in pédibus ejus non apparébit. Quis hæc operátus est, et fecit, vocans generatiónes ab exórdio? Ego Dóminus, primus et novíssimus ego sum.\par
\switchcolumn
\EN
\noindent Let the islands keep silence before me, and the nations take new strength: let them come near, and then speak, let us come near to judgment together. Who hath raised up the just one from the east, hath called him to follow him? he shall give the nations in his sight, and he shall rule over kings: he shall give them as the dust to his sword, as stubble driven by the wind, to his bow. He shall pursue them, he shall pass in peace, no path shall appear after his feet. Who hath wrought and done these things, calling the generations from the beginning? I the Lord, I am the first and the last.\par
\switchcolumn*

\LA
\begin{Resp}\Rsign Me opórtet mínui, illum autem créscere: qui autem post me venit, ante me factus est: \Star{} Cujus non sum dignus corrígiam calceamentórum sólvere.\end{Resp}
\begin{Resp}\Vsign Ego baptizávi vos aqua: ille autem baptizábit vos Spíritu Sancto.\end{Resp}
\begin{Resp}\Rsign Cujus non sum dignus corrígiam calceamentórum sólvere.\end{Resp}
\begin{Resp}\Vsign Glória Patri, et Fílio, \Star{} et Spirítui Sancto.\end{Resp}
\begin{Resp}\Rsign Cujus non sum dignus corrígiam calceamentórum sólvere.\end{Resp}
\switchcolumn
\EN
\begin{Resp}\Rsign I must decrease, but He must increase; He it is Who, coming after me is preferred before me; \Star{} Whose shoe's latchet I am not worthy to unloose.\end{Resp}
\begin{Resp}\Vsign I baptize you with water; but He shall baptize you with the Holy Ghost.\end{Resp}
\begin{Resp}\Rsign Whose shoe's latchet I am not worthy to unloose.\end{Resp}
\begin{Resp}\Vsign Glory be to the Father, and to the Son, \Star{} and to the Holy Ghost.\end{Resp}
\begin{Resp}\Rsign Whose shoe's latchet I am not worthy to unloose.\end{Resp}
\switchcolumn*

\LA
\LessonHead{Lectio IV}
\switchcolumn
\EN
\LessonHead{Lesson IV}
\switchcolumn*

\LA
\LessonTitle{Sermo sancti Leónis Papæ}
\switchcolumn
\EN
\LessonTitle{From the Sermons of Pope St. Leo (the Great)}
\switchcolumn*

\LA
\Source{Sermo 1 de jejúnio décimi mensis, et collectis}
\switchcolumn
\EN
\Source{1st on the December Fast, and Almsgiving}
\switchcolumn*

\LA
\noindent Si fidéliter, dilectíssimi, atque sapiénter creatiónis nostræ intelligámus exórdium, inveniémus hóminem ideo ad imáginem Dei cónditum, ut imitátor sui esset auctóris: et hanc esse naturálem nostri géneris dignitátem, si in nobis, quasi in quodam spéculo, divínæ benignitátis forma respléndeat. Ad quam quotídie nos útique réparat grátia Salvatóris, dum quod cécidit in Adam primo, erígitur in secúndo.\par
\switchcolumn
\EN
\noindent Dearly beloved brethren, if we study attentively the history of the creation of our race, we shall find that man was made in the image of God, that his ways also might be an imitation of the ways of his Maker. This is the natural, real, and highest dignity to which we are capable of attaining, that the goodness of the Divine nature should have a reflection in us, as in a glass. As a mean of reaching this dignity, we are daily offered the grace of our Saviour, for as in the first Adam all men are fallen, so in the Second Adam can all men be raised up again (i Cor. xv. 22).\par
\switchcolumn*

\LA
\begin{Resp}\Rsign Nascétur nobis párvulus, et vocábitur Deus, Fortis: \Star{} Ipse sedébit super thronum David patris sui, et imperábit: cujus potéstas super húmerum ejus.\end{Resp}
\begin{Resp}\Vsign In ipso benedicéntur omnes tribus terræ, omnes gentes sérvient ei.\end{Resp}
\begin{Resp}\Rsign Ipse sedébit super thronum David patris sui, et imperábit: cujus potéstas super húmerum ejus.\end{Resp}
\switchcolumn
\EN
\begin{Resp}\Rsign Unto us shall a Child be born, and His name shall be called the Mighty God. \Star{} He shall sit upon the throne of His father David, and shall reign, and the government shall be upon His shoulder.\end{Resp}
\begin{Resp}\Vsign In Him shall all the kindreds of the earth be blessed; all nations shall serve Him.\end{Resp}
\begin{Resp}\Rsign He shall sit upon the throne of His father David, and shall reign, and the government shall be upon His shoulder.\end{Resp}
\switchcolumn*

\LA
\LessonHead{Lectio V}
\switchcolumn
\EN
\LessonHead{Lesson V}
\switchcolumn*

\LA
\noindent Causa autem reparatiónis nostræ non est nisi misericórdia Dei: quem non diligerémus, nisi prius nos ipse dilígeret, et ténebras ignorántiæ nostræ, suæ veritátis luce discúteret. Quod per sanctum Isaíam Dóminus denúntians, ait: Addúcam cæcos in viam quam ignorábant, et sémitas quas nesciébant, fáciam illos calcáre: fáciam illis ténebras in lucem, et prava in dirécta. Hæc verba fáciam illis, et non relínquam eos. Et íterum: Invéntus sum, inquit, a non quæréntibus me, et palam appárui iis qui me non interrogábant.\par
\switchcolumn
\EN
\noindent Our restoration from the consequences of Adam's fall is sheer mercy of God, and nothing else; we should not have loved Him unless He had first loved us, (1 John iv. 19,) and scattered the darkness of our ignorance by the light of His truth. This the Lord promised by the mouth of Isaiah, where He saith: (Isa. xlii. 16,) I will bring the blind by a way that they knew not, and I will lead them in paths that they have not known I will make darkness light before them, and crooked things straight. These things will I do unto them and not forsake them. And again: (Isa. lxv. 1, 2; Rom. x. 20,) I was found of them that sought Me not; I was made manifest unto them that asked not after Me.\par
\switchcolumn*

\LA
\begin{Resp}\Rsign Ecce jam venit plenitúdo témporis, in quo misit Deus Fílium suum in terras, natum de Vírgine, factum sub lege: \Star{} Ut eos, qui sub lege erant, redímeret.\end{Resp}
\begin{Resp}\Vsign Propter nímiam caritátem suam, qua diléxit nos Deus, Fílium suum misit in similitúdinem carnis peccáti.\end{Resp}
\begin{Resp}\Rsign Ut eos, qui sub lege erant, redímeret.\end{Resp}
\switchcolumn
\EN
\begin{Resp}\Rsign Behold, the fullness of the time is come, wherein God hath sent forth His Son into the world, born of a Virgin, made under the law; \Star{} To redeem them that were under the law.\end{Resp}
\begin{Resp}\Vsign God, for His great love wherewith He loved us, hath sent His own Son in the likeness of sinful flesh.\end{Resp}
\begin{Resp}\Rsign To redeem them that were under the law.\end{Resp}
\switchcolumn*

\LA
\LessonHead{Lectio VI}
\switchcolumn
\EN
\LessonHead{Lesson VI}
\switchcolumn*

\LA
\noindent Quod quómodo implétum sit, Joánnes Apóstolus docet, dicens: Scimus quóniam Fílius Dei venit, et dedit nobis sensum, ut cognoscámus verum, et simus in vero Fílio ejus. Et íterum: Nos ergo diligámus Deum, quóniam ipse prior diléxit nos. Diligéndo ítaque nos Deus, ad imáginem suam nos réparat: et ut in nobis formam suæ bonitátis invéniat, dat unde ipsi quoque quod operátur operémur, accéndens scílicet méntium nostrárum lucérnas, et igne nos suæ caritátis inflámmans, ut non solum ipsum, sed étiam quidquid díligit, diligámus.\par
\switchcolumn
\EN
\noindent And we know from the Apostle John how God fulfilled His promise, (1 John v. 20.) We know that the Son of God is come, and hath given us an understanding, that we may know Him That is True, and be in Him That is True, even in His Son. And again: (iv. 19,) Let us therefore love God, because He first loved us. For His great love then wherewith he hath loved us, (Eph. ii. 4,) God reneweth His likeness in us. And, moreover, in order that He may find in us the reflection of His goodness, He giveth us that whereby to work along with Himself, (Who worketh all in all,) lighting, as it were, candles in our dark minds, and kindling in us the fire of His love, to make us love not Himself only, but likewise, in Him, whatsoever He loveth.\par
\switchcolumn*

\LA
\begin{Resp}\Rsign Virgo Israël, revértere ad civitátes tuas: \Star{} Usquequo dolens avertéris? generábis Dóminum Salvatórem, oblatiónem novam in terra: * Ambulábunt hómines in salvatiónem.\end{Resp}
\begin{Resp}\Vsign In caritáte perpétua diléxi te: ídeo attráxi te míserans tui.\end{Resp}
\begin{Resp}\Rsign Usquequo dolens avertéris? generábis Dóminum Salvatórem, oblatiónem novam in terra.\end{Resp}
\begin{Resp}\Vsign Glória Patri, et Fílio, \Star{} et Spirítui Sancto.\end{Resp}
\begin{Resp}\Rsign Ambulábunt hómines in salvatiónem.\end{Resp}
\switchcolumn
\EN
\begin{Resp}\Rsign O virgin of Israel, turn again to thy cities. \Star{} How long wilt thou go about sorrowing? Thou shalt bring forth the Lord thy Saviour, a new offering in the earth; men shall walk in paths of salvation.\end{Resp}
\begin{Resp}\Vsign I have loved thee with an everlasting love; therefore with loving-kindness have I drawn thee.\end{Resp}
\begin{Resp}\Rsign How long wilt thou go about sorrowing? Thou shalt bring forth the Lord thy Saviour, a new offering in the earth.\end{Resp}
\begin{Resp}\Vsign Glory be to the Father, and to the Son, \Star{} and to the Holy Ghost.\end{Resp}
\begin{Resp}\Rsign Men shall walk in paths of salvation.\end{Resp}
\switchcolumn*

\LA
\LessonHead{Lectio VII}
\switchcolumn
\EN
\LessonHead{Lesson VII}
\switchcolumn*

\LA
\LessonTitle{Léctio sancti Evangélii secúndum Lucam}
\switchcolumn
\EN
\LessonTitle{From the Holy Gospel according to Luke}
\switchcolumn*

\LA
\Cite{Luc 3:1-6}
\switchcolumn
\EN
\Cite{Luke 3:1-6}
\switchcolumn*

\LA
\noindent Anno quintodécimo impérii Tibérii Cǽsaris, procuránte Póntio Piláto Judǽam. Et réliqua.\par
\switchcolumn
\EN
\noindent In the fifteenth year of the reign of Tiberius Caesar, Pontius Pilate being governor of Judea. And so on.\par
\switchcolumn*

\LA
\LessonTitle{De Homilía sancti Gregórii Papæ}
\switchcolumn
\EN
\LessonTitle{Homily by Pope St. Gregory (the Great)}
\switchcolumn*

\LA
\Source{Homilia 20 in Evangelia ante medium}
\switchcolumn
\EN
\Source{20th on the Gospels}
\switchcolumn*

\LA
\noindent Dicébat Joánnes ad turbas, quæ éxibant ut baptizaréntur ab eo: Genímina viperárum, quis osténdit vobis fúgere a ventúra ira? Ventúra enim ira est animadvérsio ultiónis extrémæ: quam tunc fúgere peccátor non valet, qui nunc ad laménta pœniténtiæ non recúrrit. Et notándum, quod malæ sóboles, malórum paréntum actiónem imitántes, genímina viperárum vocántur: quia per hoc quod bonis ínvident, eósque persequúntur, quod quibúsdam mala retríbuunt, quod læsiónes próximis exquírunt: quóniam in his ómnibus patrum suórum carnálium vias sequúntur, quasi venenáti fílii de venenátis paréntibus nati sunt.\par
\switchcolumn
\EN
\noindent John said unto the multitude, that came forth to be baptised of him: O generation of vipers, who hath warned you to flee from the wrath to come? The wrath to come in one sense signifieth the great vengeance of the Latter Day the sinner that repenteth not of his sin now, will have no mean whereby to flee from punishment then. Let us remark that addressing evil children copying the example of evil parents, the Baptist calleth them a generation of vipers in that they were envious at the righteous, and persecuted them; that they repaid evil for evil; that they hunted out ways of harming their neighbours, in all these things following the pattern of carnal parents, the prophet likeneth them to a venomous brood hatched from a venomous stock.\par
\switchcolumn*

\LA
\begin{Resp}\Rsign Jurávi, dicit Dóminus, ut ultra jam non iráscar super terram: montes enim et colles suscípient justítiam meam, \Star{} Et testaméntum pacis erit in Jerúsalem.\end{Resp}
\begin{Resp}\Vsign Juxta est salus mea, ut véniat: et justítia mea, ut revelétur.\end{Resp}
\begin{Resp}\Rsign Et testaméntum pacis erit in Jerúsalem.\end{Resp}
\switchcolumn
\EN
\begin{Resp}\Rsign I have sworn, saith the Lord, that I will not be wroth any more with the earth; for the mountains and the hills shall receive My righteousness. \Star{} And the covenant of My peace shall be in Jerusalem.\end{Resp}
\begin{Resp}\Vsign My salvation is near to come, and My righteousness to be revealed.\end{Resp}
\begin{Resp}\Rsign And the covenant of My peace shall be in Jerusalem.\end{Resp}
\switchcolumn*

\LA
\LessonHead{Lectio VIII}
\switchcolumn
\EN
\LessonHead{Lesson VIII}
\switchcolumn*

\LA
\noindent Sed quia jam peccávimus, quia usu malæ consuetúdinis involúti sumus: dicat quid nobis faciéndum sit, ut fúgere a ventúra ira valeámus. Séquitur: Fácite ergo fructus dignos pœniténtiæ. In quibus verbis notándum est, quod amícus Sponsi non solum fructus pœniténtiæ, sed dignos pœniténtiæ ádmonet esse faciéndos. Aliud namque est fructum fácere pœniténtiæ, aliud, dignum pœniténtiæ fructum fácere. Ut enim secúndum dignos pœniténtiæ fructus loquámur, sciéndum est, quia quisquis illícita nulla commísit, huic jure concéditur, ut lícitis utátur: sicque pietátis ópera fáciat, ut tamen si volúerit, ea quæ mundi sunt, non relínquat.\par
\switchcolumn
\EN
\noindent We also have sinned, we have fallen into wicked habits. What must we do, if we would flee from the wrath to come? Let us hear John. Bring forth fruits worthy of repentance. In which words let us remark that the Friend of the Bridegroom demandeth not only fruits of repentance, but fruits worthy of repentance. The former are one thing, and the latter another. In considering then what are fruits worthy of repentance, we may remark that if we had done nothing unlawful we might have had free use of things which are lawful, and been able to sanctify ourselves without abstaining from indulgence in the things of the world.\par
\switchcolumn*

\LA
\begin{Resp}\Rsign Non discédimus a te, vivificábis nos, Dómine, et nomen tuum invocábimus: \Star{} Osténde nobis fáciem tuam, et salvi érimus.\end{Resp}
\begin{Resp}\Vsign Meménto nostri, Dómine, in beneplácito pópuli tui: vísita nos in salutári tuo.\end{Resp}
\begin{Resp}\Rsign Osténde nobis fáciem tuam, et salvi érimus.\end{Resp}
\switchcolumn
\EN
\begin{Resp}\Rsign We will not go back from thee. Thou, O Lord, shalt quicken us, and we will call upon thy name. \Star{} Cause thy face to shine upon us, and we shall be saved.\end{Resp}
\begin{Resp}\Vsign Remember us, O Lord, with the favour that Thou showest unto thy people; O visit us with thy salvation.\end{Resp}
\begin{Resp}\Rsign Cause thy face to shine upon us, and we shall be saved.\end{Resp}
\switchcolumn*

\LA
\LessonHead{Lectio IX}
\switchcolumn
\EN
\LessonHead{Lesson IX}
\switchcolumn*

\LA
\noindent At si quis in fornicatiónis culpam, vel fortasse, quod est grávius, in adultérium lapsus est: tanto a se lícita debet abscíndere, quanto se méminit et illícita perpetrásse. Neque enim par fructus boni óperis esse debet, ejus qui minus, et ejus qui ámplius delíquit: aut ejus qui in nullis, et ejus qui in quibúsdam facinóribus cécidit, et ejus qui in multis est lapsus. Per hoc ergo quod dícitur: Fácite fructus dignos pœniténtiæ: uniuscujúsque consciéntia convenítur, ut tanto majóra acquírat bonórum óperum lucra per pœniténtiam, quanto gravióra sibi íntulit damna per culpam.\par
\switchcolumn
\EN
\noindent But if any one, for example, hath fallen into fornication, or perhaps, into what is much worse, adultery, he ought to make up for his lawless pleasure by abstaining in some degree from lawful enjoyments. He that hath sinned less is not bound to mortify himself as much as he that hath sinned more, nor he that is innocent like him that is guilty. Let every one hearing these words bring forth fruits worthy of repentance, proceed to judge himself by his own conscience, and the more he perceiveth that he hath sinned, the greater penance let him do.\par
\switchcolumn*

\LA
\begin{Resp}\Rsign Intuémini, quantus sit iste, qui ingréditur ad salvándas gentes: ipse est Rex justítiæ, \Star{} Cujus generátio non habet finem.\end{Resp}
\begin{Resp}\Vsign Præcúrsor pro nobis ingréditur, secúndum órdinem Melchísedech Póntifex factus in ætérnum.\end{Resp}
\begin{Resp}\Rsign Cujus generátio non habet finem.\end{Resp}
\begin{Resp}\Vsign Glória Patri, et Fílio, \Star{} et Spirítui Sancto.\end{Resp}
\begin{Resp}\Rsign Cujus generátio non habet finem.\end{Resp}
\switchcolumn
\EN
\begin{Resp}\Rsign Consider how great this man is, who is entered in for the salvation of the nations; he is King of Righteousness; \Star{} Without descent, nor end of life.\end{Resp}
\begin{Resp}\Vsign The Fore-runner is for us entered, made an High Priest for ever after the order of Melchisedek.\end{Resp}
\begin{Resp}\Rsign Without descent, nor end of life.\end{Resp}
\begin{Resp}\Vsign Glory be to the Father, and to the Son, \Star{} and to the Holy Ghost.\end{Resp}
\begin{Resp}\Rsign Without descent, nor end of life.\end{Resp}
\switchcolumn*

\end{paracol}

\Office{Feria II infra Hebdomadam IV Adventus}{Monday of the Fourth Week of Advent}{Feria major}
\addcontentsline{toc}{subsection}{Feria II infra Hebdomadam IV Adventus\enspace\textit{Monday of the Fourth Week of Advent}}
\begin{paracol}{2}
\LA
\LessonHead{Lectio I}
\switchcolumn
\EN
\LessonHead{Lesson I}
\switchcolumn*

\LA
\LessonTitle{De Isaía Prophéta}
\switchcolumn
\EN
\LessonTitle{Lesson from the book of Isaias}
\switchcolumn*

\LA
\Cite{Isa 41:8-10}
\switchcolumn
\EN
\Cite{Isa 41:8-10}
\switchcolumn*

\LA
\noindent Et tu, Israël, serve meus, Jacob quem elégi, semen Abraham amíci mei: in quo apprehéndi te ab extrémis terræ, et a longinquis ejus vocávi te, et dixi tibi: Servus meus es tu, elégi te, et non abjéci te. Ne tímeas, quia ego tecum sum: ne declínes, quia ego Deus tuus: confortávi te, et auxiliátus sum tibi, et suscépit te déxtera justi mei.\par
\switchcolumn
\EN
\noindent But thou Israel, art my servant, Jacob whom I have chosen, the seed of Abraham my friend: In whom I have taken thee from the ends of the earth, and from the remote parts thereof have called thee, and said to thee: Thou art my servant, I have chosen thee, and have not cast thee away. Fear not, for I am with thee: turn not aside, for I am thy God: I have strengthened thee, and have helped thee, and the right hand of my just one hath upheld thee.\par
\switchcolumn*

\LA
\begin{Resp}\Rsign Cánite tuba in Sion, vocáte gentes, annuntiáte pópulis, et dícite: \Star{} Ecce Deus Salvátor noster advéniet.\end{Resp}
\begin{Resp}\Vsign Annuntiáte, et audítum fácite: loquímini, et clamáte.\end{Resp}
\begin{Resp}\Rsign Ecce Deus Salvátor noster advéniet.\end{Resp}
\switchcolumn
\EN
\begin{Resp}\Rsign Blow ye the trumpet in Zion, call together the nations, tell it out among the people, and say: \Star{} Behold, God our Saviour cometh.\end{Resp}
\begin{Resp}\Vsign Tell it out and make it to be heard; speak aloud and cry.\end{Resp}
\begin{Resp}\Rsign Behold, God our Saviour cometh.\end{Resp}
\switchcolumn*

\LA
\LessonHead{Lectio II}
\switchcolumn
\EN
\LessonHead{Lesson II}
\switchcolumn*

\LA
\Cite{Isa 41:11-13}
\switchcolumn
\EN
\Cite{Isa 41:11-13}
\switchcolumn*

\LA
\noindent Ecce confundéntur et erubéscent omnes, qui pugnant advérsum te: erunt quasi non sint, et períbunt viri, qui contradícunt tibi. Quæres eos, et non invénies, viros rebélles tuos: erunt quasi non sint, et véluti consúmptio hómines bellántes advérsum te. Quia ego Dóminus Deus tuus apprehéndens manum tuam, dicénsque tibi: Ne tímeas, ego adjúvi te.\par
\switchcolumn
\EN
\noindent Behold all that fight against thee shall be confounded and ashamed, they shall be as nothing, and the men shall perish that strive against thee. Thou shalt seek them, and shalt not find the men that resist thee: they shall be as nothing: and as a thing consumed the men that war against thee. For I am the Lord thy God, who take thee by the hand, and say to thee: Fear not, I have helped thee.\par
\switchcolumn*

\LA
\begin{Resp}\Rsign Non auferétur sceptrum de Juda, et dux de fémore ejus, donec véniat qui mitténdus est: \Star{} Et ipse erit exspectátio géntium.\end{Resp}
\begin{Resp}\Vsign Pulchrióres sunt óculi ejus vino, et dentes ejus lacte candidióres.\end{Resp}
\begin{Resp}\Rsign Et ipse erit exspectátio géntium.\end{Resp}
\switchcolumn
\EN
\begin{Resp}\Rsign The sceptre shall not depart from Judah, nor the law-giver from his loins, until he that shall be sent cometh. \Star{} And unto him shall the longing of the Gentiles be.\end{Resp}
\begin{Resp}\Vsign His eyes shall be bright with wine, and his teeth white with milk.\end{Resp}
\begin{Resp}\Rsign And unto him shall the desire of the Gentiles be.\end{Resp}
\switchcolumn*

\LA
\LessonHead{Lectio III}
\switchcolumn
\EN
\LessonHead{Lesson III}
\switchcolumn*

\LA
\Cite{Isa 41:14-16}
\switchcolumn
\EN
\Cite{Isa 41:14-16}
\switchcolumn*

\LA
\noindent Noli timére, vermis Jacob, qui mórtui estis ex Israël: ego auxiliátus sum tibi, dicit Dóminus: et redémptor tuus Sanctus Israël. Ego pósui te quasi plaustrum tritúrans novum, habens rostra serrántia: triturábis montes, et commínues: et colles quasi púlverem pones. Ventilábis eos, et ventus tollet, et turbo dispérget eos: et tu exsultábis in Dómino, in Sancto Israël lætáberis.\par
\switchcolumn
\EN
\noindent Fear not, thou worm of Jacob, you that are dead of Israel: I have helped thee, saith the Lord: and thy Redeemer the Holy One of Israel. I have made thee as a new thrashing wain, with teeth like a saw: thou shall thrash the mountains, and break them in pieces: and shalt make the hills as chaff. Thou shalt fan them, and the wind shall carry them away, and the whirlwind shall scatter them: and thou shalt rejoice in the Lord, in the Holy One of Israel thou shalt be joyful.\par
\switchcolumn*

\LA
\begin{Resp}\Rsign Me opórtet mínui, illum autem créscere: qui autem post me venit, ante me factus est: \Star{} Cujus non sum dignus corrígiam calceamentórum sólvere.\end{Resp}
\begin{Resp}\Vsign Ego baptizávi vos aqua: ille autem baptizábit vos Spíritu Sancto.\end{Resp}
\begin{Resp}\Rsign Cujus non sum dignus corrígiam calceamentórum sólvere.\end{Resp}
\begin{Resp}\Vsign Glória Patri, et Fílio, \Star{} et Spirítui Sancto.\end{Resp}
\begin{Resp}\Rsign Cujus non sum dignus corrígiam calceamentórum sólvere.\end{Resp}
\switchcolumn
\EN
\begin{Resp}\Rsign I must decrease, but He must increase; He it is Who, coming after me is preferred before me; \Star{} Whose shoe's latchet I am not worthy to unloose.\end{Resp}
\begin{Resp}\Vsign I baptize you with water; but He shall baptize you with the Holy Ghost.\end{Resp}
\begin{Resp}\Rsign Whose shoe's latchet I am not worthy to unloose.\end{Resp}
\begin{Resp}\Vsign Glory be to the Father, and to the Son, \Star{} and to the Holy Ghost.\end{Resp}
\begin{Resp}\Rsign Whose shoe's latchet I am not worthy to unloose.\end{Resp}
\switchcolumn*

\end{paracol}

\Office{Feria III infra Hebdomadam IV Adventus}{Tuesday of the Fourth Week of Advent}{Feria major}
\addcontentsline{toc}{subsection}{Feria III infra Hebdomadam IV Adventus\enspace\textit{Tuesday of the Fourth Week of Advent}}
\begin{paracol}{2}
\LA
\LessonHead{Lectio I}
\switchcolumn
\EN
\LessonHead{Lesson I}
\switchcolumn*

\LA
\LessonTitle{De Isaía Prophéta}
\switchcolumn
\EN
\LessonTitle{Lesson from the book of Isaias}
\switchcolumn*

\LA
\Cite{Isa 42:1-4}
\switchcolumn
\EN
\Cite{Isa 42:1-4}
\switchcolumn*

\LA
\noindent Ecce servus meus, suscípiam eum: eléctus meus, complácuit sibi in illo ánima mea: dedi spíritum meum super eum, judícium géntibus próferet. Non clamábit, neque accípiet persónam, nec audiétur vox ejus foris. Cálamum quassátum non cónteret, et linum fúmigans non exstínguet: in veritáte edúcet judícium. Non erit tristis, neque turbuléntus, donec ponat in terra judícium: et legem ejus ínsulæ exspectábunt.\par
\switchcolumn
\EN
\noindent Behold my servant, I will uphold him: my elect, my soul delighteth in him: I have given my spirit upon him, he shall bring forth judgment to the Gentiles. He shall not cry, nor have respect to person, neither shall his voice be heard abroad. The bruised reed he shall not break, and smoking flax he shall not quench: he shall bring forth judgment unto truth. He shall not be sad, nor troublesome, till he set judgment in the earth: and the islands shall wait for his law.\par
\switchcolumn*

\LA
\begin{Resp}\Rsign Nascétur nobis párvulus, et vocábitur Deus, Fortis: \Star{} Ipse sedébit super thronum David patris sui, et imperábit: cujus potéstas super húmerum ejus.\end{Resp}
\begin{Resp}\Vsign In ipso benedicéntur omnes tribus terræ, omnes gentes sérvient ei.\end{Resp}
\begin{Resp}\Rsign Ipse sedébit super thronum David patris sui, et imperábit: cujus potéstas super húmerum ejus.\end{Resp}
\switchcolumn
\EN
\begin{Resp}\Rsign Unto us shall a Child be born, and His name shall be called the Mighty God. \Star{} He shall sit upon the throne of His father David, and shall reign, and the government shall be upon His shoulder.\end{Resp}
\begin{Resp}\Vsign In Him shall all the kindreds of the earth be blessed; all nations shall serve Him.\end{Resp}
\begin{Resp}\Rsign He shall sit upon the throne of His father David, and shall reign, and the government shall be upon His shoulder.\end{Resp}
\switchcolumn*

\LA
\LessonHead{Lectio II}
\switchcolumn
\EN
\LessonHead{Lesson II}
\switchcolumn*

\LA
\Cite{Isa 42:5-7}
\switchcolumn
\EN
\Cite{Isa 42:5-7}
\switchcolumn*

\LA
\noindent Hæc dicit Dóminus Deus creans cælos, et exténdens eos: firmans terram, et quæ gérminant ex ea: dans flatum pópulo, qui est super eam, et spíritum calcántibus eam. Ego Dóminus vocávi te in justítia, et apprehéndi manum tuam, et servávi te. Et dedi te in fœdus pópuli, in lucem géntium: ut aperíres óculos cæcórum, et edúceres de conclusióne vinctum, de domo cárceris sedéntes in ténebris.\par
\switchcolumn
\EN
\noindent Thus saith the Lord God that created the heavens, and stretched them out: that established the earth, and the things that spring out of it: that giveth breath to the people upon it, and spirit to them that tread thereon. I the Lord have called thee in justice, and taken thee by the hand, and preserved thee. And I have given thee for a covenant of the people, for a light of the Gentiles: That thou mightest open the eyes of the blind, and bring forth the prisoner out of prison, and them that sit in darkness out of the prison house.\par
\switchcolumn*

\LA
\begin{Resp}\Rsign Ecce jam venit plenitúdo témporis, in quo misit Deus Fílium suum in terras, natum de Vírgine, factum sub lege: \Star{} Ut eos, qui sub lege erant, redímeret.\end{Resp}
\begin{Resp}\Vsign Propter nímiam caritátem suam, qua diléxit nos Deus, Fílium suum misit in similitúdinem carnis peccáti.\end{Resp}
\begin{Resp}\Rsign Ut eos, qui sub lege erant, redímeret.\end{Resp}
\switchcolumn
\EN
\begin{Resp}\Rsign Behold, the fullness of the time is come, wherein God hath sent forth His Son into the world, born of a Virgin, made under the law; \Star{} To redeem them that were under the law.\end{Resp}
\begin{Resp}\Vsign God, for His great love wherewith He loved us, hath sent His own Son in the likeness of sinful flesh.\end{Resp}
\begin{Resp}\Rsign To redeem them that were under the law.\end{Resp}
\switchcolumn*

\LA
\LessonHead{Lectio III}
\switchcolumn
\EN
\LessonHead{Lesson III}
\switchcolumn*

\LA
\Cite{Isa 42:10-13}
\switchcolumn
\EN
\Cite{Isa 42:10-13}
\switchcolumn*

\LA
\noindent Cantáte Dómino cánticum novum, laus ejus ab extrémis terræ: qui descénditis in mare, et plenitúdo ejus, ínsulæ, et habitatóres eárum. Sublevétur desértum, et civitátes ejus; in dómibus habitábit Cedar: laudáte habitatóres Petræ, de vértice móntium clamábunt. Ponent Dómino glóriam, et laudem ejus in ínsulis nuntiábunt. Dóminus sicut fortis egrediétur, sicut vir præliátor suscitábit zelum: vociferábitur, et clamábit: super inimícos suos confortábitur.\par
\switchcolumn
\EN
\noindent Sing ye to the Lord a new song, his praise is from the ends of the earth: you that go down to the sea, and all that are therein: ye islands, and ye inhabitants of them. Let the desert and the cities thereof be exalted: Cedar shall dwell in houses: ye inhabitants of Petra, give praise, they shall cry from the top of the mountains. They shall give glory to the Lord, and shall declare his praise in the islands. The Lord shall go forth as a mighty man, as a man of war shall he stir up zeal: he shall shout and cry: he shall prevail against his enemies.\par
\switchcolumn*

\LA
\begin{Resp}\Rsign Virgo Israël, revértere ad civitátes tuas: \Star{} Usquequo dolens avertéris? generábis Dóminum Salvatórem, oblatiónem novam in terra: * Ambulábunt hómines in salvatiónem.\end{Resp}
\begin{Resp}\Vsign In caritáte perpétua diléxi te: ídeo attráxi te míserans tui.\end{Resp}
\begin{Resp}\Rsign Usquequo dolens avertéris? generábis Dóminum Salvatórem, oblatiónem novam in terra.\end{Resp}
\begin{Resp}\Vsign Glória Patri, et Fílio, \Star{} et Spirítui Sancto.\end{Resp}
\begin{Resp}\Rsign Ambulábunt hómines in salvatiónem.\end{Resp}
\switchcolumn
\EN
\begin{Resp}\Rsign O virgin of Israel, turn again to thy cities. \Star{} How long wilt thou go about sorrowing? Thou shalt bring forth the Lord thy Saviour, a new offering in the earth; men shall walk in paths of salvation.\end{Resp}
\begin{Resp}\Vsign I have loved thee with an everlasting love; therefore with loving-kindness have I drawn thee.\end{Resp}
\begin{Resp}\Rsign How long wilt thou go about sorrowing? Thou shalt bring forth the Lord thy Saviour, a new offering in the earth.\end{Resp}
\begin{Resp}\Vsign Glory be to the Father, and to the Son, \Star{} and to the Holy Ghost.\end{Resp}
\begin{Resp}\Rsign Men shall walk in paths of salvation.\end{Resp}
\switchcolumn*

\end{paracol}

\Office{Feria IV infra Hebdomadam IV Adventus}{Wednesday of the Fourth Week of Advent}{Feria major}
\addcontentsline{toc}{subsection}{Feria IV infra Hebdomadam IV Adventus\enspace\textit{Wednesday of the Fourth Week of Advent}}
\begin{paracol}{2}
\LA
\LessonHead{Lectio I}
\switchcolumn
\EN
\LessonHead{Lesson I}
\switchcolumn*

\LA
\LessonTitle{De Isaía Prophéta}
\switchcolumn
\EN
\LessonTitle{Lesson from the book of Isaias}
\switchcolumn*

\LA
\Cite{Isa 51:1-3}
\switchcolumn
\EN
\Cite{Isa 51:1-3}
\switchcolumn*

\LA
\noindent Audíte me, qui sequímini quod justum est, et quǽritis Dóminum: atténdite ad petram unde excísi estis, et ad cavérnam laci, de qua præcísi estis. Atténdite ad Abraham patrem vestrum, et ad Saram, quæ péperit vos, quia unum vocávi eum, et benedíxi ei, et multiplicávi eum. Consolábitur ergo Dóminus Sion, et consolábitur omnes ruínas ejus: et ponet desértum ejus quasi delícias, et solitúdinem ejus quasi hortum Dómini. Gáudium et lætítia inveniétur in ea, gratiárum áctio et vox laudis.\par
\switchcolumn
\EN
\noindent Give ear to me, you that follow that which is just, and you that seek the Lord: look unto the rock whence you are hewn, and to the hole of the pit from which you are dug out. Look unto Abraham your father, and to Sara that bore you: for I called him alone, and blessed him, and multiplied him. The Lord therefore will comfort Sion, and will comfort all the ruins thereof: and he will make her desert as a place of pleasure, and her wilderness as the garden of the Lord. Joy and gladness shall be found therein, thanksgiving, and the voice of praise.\par
\switchcolumn*

\LA
\begin{Resp}\Rsign Jurávi, dicit Dóminus, ut ultra jam non iráscar super terram: montes enim et colles suscípient justítiam meam, \Star{} Et testaméntum pacis erit in Jerúsalem.\end{Resp}
\begin{Resp}\Vsign Juxta est salus mea, ut véniat: et justítia mea, ut revelétur.\end{Resp}
\begin{Resp}\Rsign Et testaméntum pacis erit in Jerúsalem.\end{Resp}
\switchcolumn
\EN
\begin{Resp}\Rsign I have sworn, saith the Lord, that I will not be wroth any more with the earth; for the mountains and the hills shall receive My righteousness. \Star{} And the covenant of My peace shall be in Jerusalem.\end{Resp}
\begin{Resp}\Vsign My salvation is near to come, and My righteousness to be revealed.\end{Resp}
\begin{Resp}\Rsign And the covenant of My peace shall be in Jerusalem.\end{Resp}
\switchcolumn*

\LA
\LessonHead{Lectio II}
\switchcolumn
\EN
\LessonHead{Lesson II}
\switchcolumn*

\LA
\Cite{Isa 51:4-6}
\switchcolumn
\EN
\Cite{Isa 51:4-6}
\switchcolumn*

\LA
\noindent Atténdite ad me, pópule meus, et tribus mea, me audíte: quia lex a me éxiet, et judícium meum in lucem populórum requiéscet. Prope est justus meus, egréssus est salvátor meus, et brácchia mea pópulos judicábunt: me ínsulæ exspectábunt, et brácchium meum sustinébunt. Leváte in cælum óculos vestros, et vidéte sub terra deórsum: quia cæli sicut fumus liquéscent, et terra sicut vestiméntum atterétur, et habitatóres ejus sicut hæc interíbunt: Salus autem mea in sempitérnum erit, et justítia mea non defíciet.\par
\switchcolumn
\EN
\noindent Hearken unto me, O my people, and give ear to me, O my tribes: for a law shall go forth from me, and my judgment shall rest to be a light of the nations. My just one is near at hand, my saviour is gone forth, and my arms shall judge the people: the islands shall look for me, and shall patiently wait for my arm. Lift up your eyes to heaven, and look down to the earth beneath: for the heavens shall vanish like smoke, and the earth shall be worn away like a garment, and the inhabitants thereof shall perish in like manner: but my salvation shall be for ever, and my justice shall not fail.\par
\switchcolumn*

\LA
\begin{Resp}\Rsign Non discédimus a te, vivificábis nos, Dómine, et nomen tuum invocábimus: \Star{} Osténde nobis fáciem tuam, et salvi érimus.\end{Resp}
\begin{Resp}\Vsign Meménto nostri, Dómine, in beneplácito pópuli tui: vísita nos in salutári tuo.\end{Resp}
\begin{Resp}\Rsign Osténde nobis fáciem tuam, et salvi érimus.\end{Resp}
\switchcolumn
\EN
\begin{Resp}\Rsign We will not go back from thee. Thou, O Lord, shalt quicken us, and we will call upon thy name. \Star{} Cause thy face to shine upon us, and we shall be saved.\end{Resp}
\begin{Resp}\Vsign Remember us, O Lord, with the favour that Thou showest unto thy people; O visit us with thy salvation.\end{Resp}
\begin{Resp}\Rsign Cause thy face to shine upon us, and we shall be saved.\end{Resp}
\switchcolumn*

\LA
\LessonHead{Lectio III}
\switchcolumn
\EN
\LessonHead{Lesson III}
\switchcolumn*

\LA
\Cite{Isa 51:7-8}
\switchcolumn
\EN
\Cite{Isa 51:7-8}
\switchcolumn*

\LA
\noindent Audíte me, qui scitis justum, pópulus meus, lex mea in corde eórum: nolíte timére oppróbrium hóminum, et blasphémias eórum ne metuátis. Sicut enim vestiméntum, sic cómedet eos vermis: et sicut lanam, sic devorábit eos tínea: Salus autem mea in sempitérnum erit, et justítia mea in generatiónes generatiónum.\par
\switchcolumn
\EN
\noindent Hearken to me, you that know what is just, my people who have my law in your heart: fear ye not the reproach of men, and be not afraid of their blasphemies. For the worm shall eat them up as a garment: and the moth shall consume them as wool: but my salvation shall be for ever, and my justice from generation to generation,\par
\switchcolumn*

\LA
\begin{Resp}\Rsign Intuémini, quantus sit iste, qui ingréditur ad salvándas gentes: ipse est Rex justítiæ, \Star{} Cujus generátio non habet finem.\end{Resp}
\begin{Resp}\Vsign Præcúrsor pro nobis ingréditur, secúndum órdinem Melchísedech Póntifex factus in ætérnum.\end{Resp}
\begin{Resp}\Rsign Cujus generátio non habet finem.\end{Resp}
\begin{Resp}\Vsign Glória Patri, et Fílio, \Star{} et Spirítui Sancto.\end{Resp}
\begin{Resp}\Rsign Cujus generátio non habet finem.\end{Resp}
\switchcolumn
\EN
\begin{Resp}\Rsign Consider how great this man is, who is entered in for the salvation of the nations; he is King of Righteousness; \Star{} Without descent, nor end of life.\end{Resp}
\begin{Resp}\Vsign The Fore-runner is for us entered, made an High Priest for ever after the order of Melchisedek.\end{Resp}
\begin{Resp}\Rsign Without descent, nor end of life.\end{Resp}
\begin{Resp}\Vsign Glory be to the Father, and to the Son, \Star{} and to the Holy Ghost.\end{Resp}
\begin{Resp}\Rsign Without descent, nor end of life.\end{Resp}
\switchcolumn*

\end{paracol}

\Office{Feria V infra Hebdomadam IV Adventus}{Thursday of the Fourth Week of Advent}{Feria major}
\addcontentsline{toc}{subsection}{Feria V infra Hebdomadam IV Adventus\enspace\textit{Thursday of the Fourth Week of Advent}}
\begin{paracol}{2}
\LA
\LessonHead{Lectio I}
\switchcolumn
\EN
\LessonHead{Lesson I}
\switchcolumn*

\LA
\LessonTitle{De Isaía Prophéta}
\switchcolumn
\EN
\LessonTitle{Lesson from the book of Isaias}
\switchcolumn*

\LA
\Cite{Isa 64:1-4}
\switchcolumn
\EN
\Cite{Isa 64:1-4}
\switchcolumn*

\LA
\noindent Utinam dirúmperes cælos, et descénderes: a fácie tua montes deflúerent. Sicut exústio ignis tabéscerent, aquæ ardérent igni, ut notum fíeret nomen tuum inimícis tuis: a fácie tua gentes turbaréntur. Cum féceris mirabília, non sustinébimus: descendísti, et a fácie tua montes defluxérunt. A sǽculo non audiérunt, neque áuribus percepérunt: óculus non vidit, Deus absque te, quæ præparásti exspectántibus te.\par
\switchcolumn
\EN
\noindent That thou wouldst rend the heavens, and wouldst come down: the mountains would melt away at thy presence. They would melt as at the burning of fire, the waters would burn with fire, that thy name might be made known to thy enemies: that the nations might tremble at thy presence. When thou shalt do wonderful things, we shall not bear them: thou didst come down, and at thy presence the mountains melted away. From the beginning of the world they have not heard, nor perceived with the ears: the eye hath not seen, O God, besides thee, what things thou hast prepared for them that wait for thee.\par
\switchcolumn*

\LA
\begin{Resp}\Rsign Cánite tuba in Sion, vocáte gentes, annuntiáte pópulis, et dícite: \Star{} Ecce Deus Salvátor noster advéniet.\end{Resp}
\begin{Resp}\Vsign Annuntiáte, et audítum fácite: loquímini, et clamáte.\end{Resp}
\begin{Resp}\Rsign Ecce Deus Salvátor noster advéniet.\end{Resp}
\switchcolumn
\EN
\begin{Resp}\Rsign Blow ye the trumpet in Zion, call together the nations, tell it out among the people, and say: \Star{} Behold, God our Saviour cometh.\end{Resp}
\begin{Resp}\Vsign Tell it out and make it to be heard; speak aloud and cry.\end{Resp}
\begin{Resp}\Rsign Behold, God our Saviour cometh.\end{Resp}
\switchcolumn*

\LA
\LessonHead{Lectio II}
\switchcolumn
\EN
\LessonHead{Lesson II}
\switchcolumn*

\LA
\Cite{Isa 64:5-7}
\switchcolumn
\EN
\Cite{Isa 64:5-7}
\switchcolumn*

\LA
\noindent Occurrísti lætánti, et faciénti justítiam: in viis tuis recordabúntur tui: ecce tu irátus es, et peccávimus: in ipsis fúimus semper, et salvábimur. Et facti sumus ut immúndus omnes nos, et quasi pannus menstruátæ univérsæ justítiæ nostræ: et cecídimus quasi fólium univérsi, et iniquitátes nostræ quasi ventus abstulérunt nos. Non est qui ínvocet nomen tuum: qui consúrgat, et téneat te: abscondísti fáciem tuam a nobis, et allisísti nos in manu iniquitátis nostræ.\par
\switchcolumn
\EN
\noindent Thou hast met him that rejoiceth, and doth justice: in thy ways they shall remember thee: behold thou art angry, and we have sinned: in them we have been always, and we shall be saved. And we are all become as one unclean, and all our justices as the rag of a menstruous woman: and we have all fallen as a leaf, and our iniquities, like the wind, have taken us away. There is none that calleth upon thy name: that riseth up, and taketh hold of thee: thou hast hid thy face from us, and hast crushed us in the hand of our iniquity.\par
\switchcolumn*

\LA
\begin{Resp}\Rsign Non auferétur sceptrum de Juda, et dux de fémore ejus, donec véniat qui mitténdus est: \Star{} Et ipse erit exspectátio géntium.\end{Resp}
\begin{Resp}\Vsign Pulchrióres sunt óculi ejus vino, et dentes ejus lacte candidióres.\end{Resp}
\begin{Resp}\Rsign Et ipse erit exspectátio géntium.\end{Resp}
\switchcolumn
\EN
\begin{Resp}\Rsign The sceptre shall not depart from Judah, nor the law-giver from his loins, until he that shall be sent cometh. \Star{} And unto him shall the longing of the Gentiles be.\end{Resp}
\begin{Resp}\Vsign His eyes shall be bright with wine, and his teeth white with milk.\end{Resp}
\begin{Resp}\Rsign And unto him shall the desire of the Gentiles be.\end{Resp}
\switchcolumn*

\LA
\LessonHead{Lectio III}
\switchcolumn
\EN
\LessonHead{Lesson III}
\switchcolumn*

\LA
\Cite{Isa 64:8-11}
\switchcolumn
\EN
\Cite{Isa 64:8-11}
\switchcolumn*

\LA
\noindent Et nunc, Dómine, pater noster es tu, nos vero lutum: et fictor noster tu, et ópera mánuum tuárum omnes nos. Ne irascáris, Dómine, satis, et ne ultra memíneris iniquitátis nostræ: ecce réspice, pópulus tuus omnes nos. Cívitas sancti tui facta est desérta, Sion desérta facta est, Jerúsalem desoláta est. Domus sanctificatiónis nostræ, et glóriæ nostræ, ubi laudavérunt te patres nostri, facta est in exustiónem ignis, et ómnia desiderabília nostra versa sunt in ruínas.\par
\switchcolumn
\EN
\noindent And now, O Lord, thou art our father, and we are clay: and thou art our maker, and we all are the works of thy hands. Be not very angry, O Lord, and remember no longer our iniquity: behold, see we are all thy people. The city of thy sanctuary is become a desert, Sion is made a desert, Jerusalem is desolate. The house of our holiness, and of our glory, where our fathers praised thee, is burnt with fire, and all our lovely things are turned into ruins.\par
\switchcolumn*

\LA
\begin{Resp}\Rsign Me opórtet mínui, illum autem créscere: qui autem post me venit, ante me factus est: \Star{} Cujus non sum dignus corrígiam calceamentórum sólvere.\end{Resp}
\begin{Resp}\Vsign Ego baptizávi vos aqua: ille autem baptizábit vos Spíritu Sancto.\end{Resp}
\begin{Resp}\Rsign Cujus non sum dignus corrígiam calceamentórum sólvere.\end{Resp}
\begin{Resp}\Vsign Glória Patri, et Fílio, \Star{} et Spirítui Sancto.\end{Resp}
\begin{Resp}\Rsign Cujus non sum dignus corrígiam calceamentórum sólvere.\end{Resp}
\switchcolumn
\EN
\begin{Resp}\Rsign I must decrease, but He must increase; He it is Who, coming after me is preferred before me; \Star{} Whose shoe's latchet I am not worthy to unloose.\end{Resp}
\begin{Resp}\Vsign I baptize you with water; but He shall baptize you with the Holy Ghost.\end{Resp}
\begin{Resp}\Rsign Whose shoe's latchet I am not worthy to unloose.\end{Resp}
\begin{Resp}\Vsign Glory be to the Father, and to the Son, \Star{} and to the Holy Ghost.\end{Resp}
\begin{Resp}\Rsign Whose shoe's latchet I am not worthy to unloose.\end{Resp}
\switchcolumn*

\end{paracol}

\Office{Feria VI infra Hebdomadam IV Adventus}{Friday of the Fourth Week of Advent}{Feria major}
\addcontentsline{toc}{subsection}{Feria VI infra Hebdomadam IV Adventus\enspace\textit{Friday of the Fourth Week of Advent}}
\begin{paracol}{2}
\LA
\LessonHead{Lectio I}
\switchcolumn
\EN
\LessonHead{Lesson I}
\switchcolumn*

\LA
\LessonTitle{De Isaía Prophéta}
\switchcolumn
\EN
\LessonTitle{Lesson from the book of Isaias}
\switchcolumn*

\LA
\Cite{Isa 66:5-8}
\switchcolumn
\EN
\Cite{Isa 66:5-8}
\switchcolumn*

\LA
\noindent Audíte verbum Dómini, qui trémitis ad verbum ejus: dixérunt fratres vestri odiéntes vos, et abiciéntes propter nomen meum: glorificétur Dóminus, et vidébimus in lætítia vestra: ipsi autem confundéntur. Vox pópuli de civitáte, vox de templo, vox Dómini reddéntis retributiónem inimícis suis. Antequam parturíret, péperit: ántequam veníret partus ejus, péperit másculum. Quis audívit umquam tale? et quis vidit huic símile? Numquid partúriet terra in die una? aut pariétur gens simul, quia parturívit et péperit Sion fílios suos?\par
\switchcolumn
\EN
\noindent Hear the word of the Lord, you that tremble at his word: Your brethren that hate you, and cast you out for my name's sake, have said: Let the Lord be glorified, and we shall see in your joy: but they shall be confounded. A voice of the people from the city, a voice from the temple, the voice of the Lord that rendereth recompense to his enemies. Before she was in labour, she brought forth; before her time came to be delivered, she brought forth a man child. Who hath ever heard such a thing? and who hath seen the like to this? shall the earth bring forth in one day? or shall a nation be brought forth at once, because Sion hath been in labour, and hath brought forth her children?\par
\switchcolumn*

\LA
\begin{Resp}\Rsign Nascétur nobis párvulus, et vocábitur Deus, Fortis: \Star{} Ipse sedébit super thronum David patris sui, et imperábit: cujus potéstas super húmerum ejus.\end{Resp}
\begin{Resp}\Vsign In ipso benedicéntur omnes tribus terræ, omnes gentes sérvient ei.\end{Resp}
\begin{Resp}\Rsign Ipse sedébit super thronum David patris sui, et imperábit: cujus potéstas super húmerum ejus.\end{Resp}
\switchcolumn
\EN
\begin{Resp}\Rsign Unto us shall a Child be born, and His name shall be called the Mighty God. \Star{} He shall sit upon the throne of His father David, and shall reign, and the government shall be upon His shoulder.\end{Resp}
\begin{Resp}\Vsign In Him shall all the kindreds of the earth be blessed; all nations shall serve Him.\end{Resp}
\begin{Resp}\Rsign He shall sit upon the throne of His father David, and shall reign, and the government shall be upon His shoulder.\end{Resp}
\switchcolumn*

\LA
\LessonHead{Lectio II}
\switchcolumn
\EN
\LessonHead{Lesson II}
\switchcolumn*

\LA
\Cite{Isa 66:9-12}
\switchcolumn
\EN
\Cite{Isa 66:9-12}
\switchcolumn*

\LA
\noindent Numquid ego, qui álios párere fácio, ipse non páriam, dicit Dóminus? si ego, qui generatiónem céteris tríbuo, stérilis ero, ait Dóminus, Deus tuus? Lætámini cum Jerúsalem, et exsultáte in ea, omnes qui dilígitis eam: gaudéte cum ea gáudio, univérsi qui lugétis super eam, ut sugátis, et repleámini ab úbere consolatiónis ejus: ut mulgeátis, et delíciis affluátis ab omnímoda glória ejus. Quia hæc dicit Dóminus: Ecce ego declinábo super eam quasi flúvium pacis, et quasi torréntem inundántem glóriam géntium, quam sugétis: ad úbera portabímini, et super génua blandiéntur vobis.\par
\switchcolumn
\EN
\noindent Shall not I that make others to bring forth children, myself bring forth, saith the Lord? shall I, that give generation to others, be barren, saith the Lord thy God? Rejoice with Jerusalem, and be glad with her, all you that love her: rejoice for joy with her, all you that mourn for her. That you may suck, and be filled with the breasts of her consolations: that you may milk out, and flow with delights, from the abundance of her glory. For thus saith the Lord: Behold I will bring upon her as it were a river of peace, and as an overflowing torrent the glory of the Gentiles, which you shall suck; you shall be carried at the breasts, end upon the knees they shall caress you.\par
\switchcolumn*

\LA
\begin{Resp}\Rsign Ecce jam venit plenitúdo témporis, in quo misit Deus Fílium suum in terras, natum de Vírgine, factum sub lege: \Star{} Ut eos, qui sub lege erant, redímeret.\end{Resp}
\begin{Resp}\Vsign Propter nímiam caritátem suam, qua diléxit nos Deus, Fílium suum misit in similitúdinem carnis peccáti.\end{Resp}
\begin{Resp}\Rsign Ut eos, qui sub lege erant, redímeret.\end{Resp}
\switchcolumn
\EN
\begin{Resp}\Rsign Behold, the fullness of the time is come, wherein God hath sent forth His Son into the world, born of a Virgin, made under the law; \Star{} To redeem them that were under the law.\end{Resp}
\begin{Resp}\Vsign God, for His great love wherewith He loved us, hath sent His own Son in the likeness of sinful flesh.\end{Resp}
\begin{Resp}\Rsign To redeem them that were under the law.\end{Resp}
\switchcolumn*

\LA
\LessonHead{Lectio III}
\switchcolumn
\EN
\LessonHead{Lesson III}
\switchcolumn*

\LA
\Cite{Isa 66:13-16}
\switchcolumn
\EN
\Cite{Isa 66:13-18}
\switchcolumn*

\LA
\noindent Quómodo si cui mater blandiátur, ita ego consolábor vos, et in Jerúsalem consolabímini. Vidébitis, et gaudébit cor vestrum, et ossa vestra quasi herba germinábunt, et cognoscétur manus Dómini servis ejus, et indignábitur inimícis suis. Quia ecce Dóminus in igne véniet, et quasi turbo quadrígæ ejus: réddere in indignatióne furórem suum, et increpatiónem suam in flamma ignis, quia in igne Dóminus dijudicábit, et in gládio suo ad omnem carnem, et multiplicabúntur interfécti a Dómino.\par
\switchcolumn
\EN
\noindent As one whom the mother caresseth, so will I comfort you, and you shall be comforted in Jerusalem. You shall see and your heart shall rejoice, and your bones shall flourish like an herb, and the hand of the Lord shall be known to his servants, and he shall be angry with his enemies. For behold the Lord will come with fire, and his chariots are like a whirlwind, to render his wrath in indignation, and his rebuke with flames of fire. For the Lord shall judge by fire, and by his sword unto all flesh, and the slain of the Lord shall be many. They that were sanctified, and thought themselves clean in the gardens behind the gate within, they that did eat swine's flesh, and the abomination, and the mouse: they shall be consumed together, saith the Lord. But I know their works, and their thoughts: I come that I may gather them together with all nations and tongues: and they shall come and shall see my glory.\par
\switchcolumn*

\LA
\begin{Resp}\Rsign Virgo Israël, revértere ad civitátes tuas: \Star{} Usquequo dolens avertéris? generábis Dóminum Salvatórem, oblatiónem novam in terra: * Ambulábunt hómines in salvatiónem.\end{Resp}
\begin{Resp}\Vsign In caritáte perpétua diléxi te: ídeo attráxi te míserans tui.\end{Resp}
\begin{Resp}\Rsign Usquequo dolens avertéris? generábis Dóminum Salvatórem, oblatiónem novam in terra.\end{Resp}
\begin{Resp}\Vsign Glória Patri, et Fílio, \Star{} et Spirítui Sancto.\end{Resp}
\begin{Resp}\Rsign Ambulábunt hómines in salvatiónem.\end{Resp}
\switchcolumn
\EN
\begin{Resp}\Rsign O virgin of Israel, turn again to thy cities. \Star{} How long wilt thou go about sorrowing? Thou shalt bring forth the Lord thy Saviour, a new offering in the earth; men shall walk in paths of salvation.\end{Resp}
\begin{Resp}\Vsign I have loved thee with an everlasting love; therefore with loving-kindness have I drawn thee.\end{Resp}
\begin{Resp}\Rsign How long wilt thou go about sorrowing? Thou shalt bring forth the Lord thy Saviour, a new offering in the earth.\end{Resp}
\begin{Resp}\Vsign Glory be to the Father, and to the Son, \Star{} and to the Holy Ghost.\end{Resp}
\begin{Resp}\Rsign Men shall walk in paths of salvation.\end{Resp}
\switchcolumn*

\end{paracol}

\SeasonTitle{Tempus Nativitatis}{Christmastide}

\addcontentsline{toc}{section}{Tempus Nativitatis\enspace\textit{Christmastide}}

\Office{Dominica Infra Octavam Nativitatis}{Sunday Within The Octave of Christmas}{Semiduplex}
\addcontentsline{toc}{subsection}{Dominica Infra Octavam Nativitatis\enspace\textit{Sunday Within The Octave of Christmas}}
\begin{paracol}{2}
\LA
\RefLine{Lectiones I–III: de Scriptura occurrente.}
\switchcolumn
\EN
\RefLine{Lessons I–III: from the Scripture of the day.}
\switchcolumn*

\LA
\LessonHead{Lectio IV}
\switchcolumn
\EN
\LessonHead{Lesson IV}
\switchcolumn*

\LA
\LessonTitle{Sermo sancti Leónis Papæ}
\switchcolumn
\EN
\LessonTitle{From the Sermons of Pope St. Leo (the Great.)}
\switchcolumn*

\LA
\Source{Sermo 9 de Nativ. Dom.}
\switchcolumn
\EN
\Source{9th on Christmas.}
\switchcolumn*

\LA
\noindent Excédit quidem, dilectíssimi, multúmque superéminet humáni elóquii facultátem divini operis magnitúdo: et inde óritur difficultas fandi, unde adest rátio non tacendi. Quia in Christo Jesu Fílio Dei non solum ad divinam esséntiam, sed étiam ad humanam spectat natúram quod dictum est per prophétam: Generatiónem ejus quis enarrábit? Utrámque enim substántiam in unam convenísse persónam, nisi fides credat, sermo non explicat. Et ideo numquam materia déficit laudis, quia numquam sufficit copia laudatóris.\par
\switchcolumn
\EN
\noindent Dearly beloved brethren, the greatness of God's work, in its breadth and height, passeth the power of man's utterance; and, therefore, when we must needs not keep silence, we find it hard to know what to say. The words of the Prophet: Who shall declare His generation? (Isa. liii. 8) look not only to the Divine, but also to the human birth of Jesus Christ, the Son of God. Faith believeth, but words cannot explain, how the two natures were joined in one Person, and therein we find that we shall never lack matter of praise in Him, Whose abundance ever outrunneth the power of our expression.\par
\switchcolumn*

\LA
\begin{Resp}\Rsign Beáta Dei Génetrix María, cujus víscera intácta pérmanent: \Star{} Hódie génuit Salvatórem sǽculi.\end{Resp}
\begin{Resp}\Vsign Beáta, quæ crédidit: quóniam perfécta sunt ómnia, quæ dicta sunt ei a Dómino.\end{Resp}
\begin{Resp}\Rsign Hódie génuit Salvatórem sǽculi.\end{Resp}
\switchcolumn
\EN
\begin{Resp}\Rsign Blessed is God's holy Mother, Mary, maiden undefiled. \Star{} This day hath she brought forth the Saviour of the world.\end{Resp}
\begin{Resp}\Vsign Blessed is she that believed; for there is a performance of all those things which were told her from the Lord.\end{Resp}
\begin{Resp}\Rsign This day hath she brought forth the Saviour of the world.\end{Resp}
\switchcolumn*

\LA
\LessonHead{Lectio V}
\switchcolumn
\EN
\LessonHead{Lesson V}
\switchcolumn*

\LA
\noindent Gaudeámus ígitur, quod ad eloquéndum tantæ misericórdiæ sacraméntum ímpares sumus; et cum salútis nostræ altitúdinem prómere non valeamus, sentiámus nobis bonum esse, quod vincimur. Nemo enim ad cognitiónem veritátis magis propinquat, quam qui intélligit in rebus divinis, étiam si multum profíciat, semper sibi superesse quod quærat. Nam qui se ad id, in quod tendit, pervenísse præsúmit, non quæsíta réperit, sed in inquisitióne déficit.\par
\switchcolumn
\EN
\noindent Therefore let us rejoice, that this mystery of mercy is greater than we can ever speak; and let us feel that it is good for us to fail if we try to express the height and depth of redeeming love. He cometh nearest to the knowledge of the truth, who, the farther he advanceth, seeth all the more clearly that he can never overtake that whereafter he searcheth. For he that imagineth therein that he hath ever attained unto the goal, hath not found that which he seeketh, but hath altogether missed.\par
\switchcolumn*

\LA
\begin{Resp}\Rsign Sancta et immaculáta virgínitas, quibus te láudibus éfferam, néscio: \Star{} Quia quem cæli cápere non póterant, tuo grémio contulísti.\end{Resp}
\begin{Resp}\Vsign Benedícta tu in muliéribus, et benedíctus fructus ventris tui.\end{Resp}
\begin{Resp}\Rsign Quia quem cæli cápere non póterant, tuo grémio contulísti.\end{Resp}
\switchcolumn
\EN
\begin{Resp}\Rsign O Mary, how holy and how spotless is thy virginity! I am too dull to praise thee! \Star{} For thou hast borne in thy breast Him Whom the heavens cannot contain.\end{Resp}
\begin{Resp}\Vsign Blessed art thou among women, and blessed is the fruit of thy womb.\end{Resp}
\begin{Resp}\Rsign For thou hast borne in thy breast Him Whom the heavens cannot contain.\end{Resp}
\switchcolumn*

\LA
\LessonHead{Lectio VI}
\switchcolumn
\EN
\LessonHead{Lesson VI}
\switchcolumn*

\LA
\noindent Ne autem infirmitátis nostræ perturbemur angústiis, evangélicæ nos et prophéticæ ádjuvant voces; quibus ita accendimur, et docemur, ut nos nativitátem Dómini, qua Verbum caro factum est, non tam prætéritam recólere, quam præséntem videámur inspícere. Quod enim pastóribus pro gregum suórum custódia vigilántibus nuntiávit Angelus, étiam nostrum implévit audítum: et ideo Dominicis ovibus præsumus, quia verba divínitus édita cordis aure servámus: tamquam et in hodiérna festivitáte dicátur: Evangelizo vobis gáudium magnum, quod erit omni pópulo, quia natus est vobis hódie Salvator, qui est Christus Dóminus, in civitáte David.\par
\switchcolumn
\EN
\noindent But lest we should be confounded at the weakness of our mortality, we have help in the words of the Prophets and Evangelists; and they are able so to inflame and teach us that we may see the Birth of the Lord, wherein the Word was made Flesh, not so much as a thing past, as a thing present. The proclamation of the angel to the shepherds who watched their flocks by night, ringeth in our ears also; and for this end are we appointed to rule the Lord's flock, that we may ever keep in our heart the word revealed from heaven, and say unto you, as we do this day Behold, I bring you good tidings of great joy, which shall be to all people; for unto you is born this day, in the city of David, a Saviour, Which is Christ the Lord!\par
\switchcolumn*

\LA
\begin{Resp}\Rsign Beáta víscera Maríæ Vírginis, quæ portavérunt ætérni Patris Fílium: et beáta úbera, quæ lactavérunt Christum Dóminum: \Star{} Qui hódie pro salúte mundi de Vírgine nasci dignátus est.\end{Resp}
\begin{Resp}\Vsign Dies sanctificátus illúxit nobis: veníte, gentes, et adoráte Dóminum.\end{Resp}
\begin{Resp}\Rsign Qui hódie pro salúte mundi de Vírgine nasci dignátus est.\end{Resp}
\begin{Resp}\Vsign Glória Patri, et Fílio, \Star{} et Spirítui Sancto.\end{Resp}
\begin{Resp}\Rsign Qui hódie pro salúte mundi de Vírgine nasci dignátus est.\end{Resp}
\switchcolumn
\EN
\begin{Resp}\Rsign Blessed be the womb of the Virgin Mary, which bore the Son of the Eternal Father, and blessed be the paps which give suck to Christ our Lord. \Star{} This day hath He been pleased for the salvation of the world to be born of a Virgin.\end{Resp}
\begin{Resp}\Vsign This day which is breaking is holy; O come, ye Gentiles, and worship the Lord.\end{Resp}
\begin{Resp}\Rsign This day hath He been pleased for the salvation of the world to be born of a Virgin.\end{Resp}
\begin{Resp}\Vsign Glory be to the Father, and to the Son, \Star{} and to the Holy Ghost.\end{Resp}
\begin{Resp}\Rsign This day hath He been pleased for the salvation of the world to be born of a Virgin.\end{Resp}
\switchcolumn*

\LA
\LessonHead{Lectio VII}
\switchcolumn
\EN
\LessonHead{Lesson VII}
\switchcolumn*

\LA
\LessonTitle{Léctio sancti Evangélii secúndum Lucam}
\switchcolumn
\EN
\LessonTitle{From the Holy Gospel according to Luke}
\switchcolumn*

\LA
\Cite{Luc 2:33-40}
\switchcolumn
\EN
\Cite{Luke 2:33-40}
\switchcolumn*

\LA
\noindent In illo témpore: Erat Joseph, et María mater Jesu, mirántes super his, quæ dicebántur de illo. Et réliqua.\par
\switchcolumn
\EN
\noindent In that time, his father and mother were wondering at those things which were spoken concerning him. And so on.\par
\switchcolumn*

\LA
\LessonTitle{Homilía sancti Ambrósii Epíscopi}
\switchcolumn
\EN
\LessonTitle{Homily by St. Ambrose, Bishop (of Milan.)}
\switchcolumn*

\LA
\Source{Lib. 2 in cap. 2 Lucæ, prope finem}
\switchcolumn
\EN
\Source{Bk. ii on Luke ii}
\switchcolumn*

\LA
\noindent Vides úberem in omnes grátiam, Dómini generatióne diffúsam, et prophetíam incrédulis negátam esse, non justis. Ecce et Símeon prophétat, in ruínam et resurrectiónem plurimórum venísse Dóminum Jesum Christum, ut justórum, iniquorúmque mérita discérnat; et pro nostrórum qualitáte factórum, judex verus et justus aut supplícia decérnat, aut prǽmia.\par
\switchcolumn
\EN
\noindent We see that God's abounding grace is poured forth on all by the birth of the Lord, and that the gift of prophecy is not denied to the righteous, but to the unbelieving. Simeon prophesieth that our Lord Jesus Christ is set for the fall and rising again of many in Israel, setting forth that the just and the unjust reap different fruits from the coming of the Saviour; so will it be with us; according to our individual works will the True and Just Judge apportion to us punishment or reward.\par
\switchcolumn*

\LA
\begin{Resp}\Rsign Verbum caro factum est, et habitávit in nobis: \Star{} Et vídimus glóriam ejus, glóriam quasi Unigéniti a Patre, plenum grátiæ et veritátis.\end{Resp}
\begin{Resp}\Vsign Omnia per ipsum facta sunt, et sine ipso factum est nihil.\end{Resp}
\begin{Resp}\Rsign Et vídimus glóriam ejus, glóriam quasi Unigéniti a Patre, plenum grátiæ et veritátis.\end{Resp}
\switchcolumn
\EN
\begin{Resp}\Rsign The Word was made flesh, and dwelt among us. \Star{} And we beheld His glory, the glory as of the Only-begotten of the Father, full of grace and truth.\end{Resp}
\begin{Resp}\Vsign All things were made by Him, and without Him was not anything made.\end{Resp}
\begin{Resp}\Rsign And we beheld His glory, the glory as of the Only-begotten of the Father, full of grace and truth.\end{Resp}
\switchcolumn*

\LA
\LessonHead{Lectio VIII}
\switchcolumn
\EN
\LessonHead{Lesson VIII}
\switchcolumn*

\LA
\noindent Et tuam ipsíus ánimam pertransíbit gládius. Nec littera, nec historia docet, ex hac vita Mariam corporalis necis passióne migrasse. Non enim anima, sed corpus materiali gládio transverberátur. Et ideo prudéntiam Maríæ haud ignaram mysterii cæléstis osténdit. Vivum enim verbum Dei, et válidum, et acutius omni gládio acutíssimo, pénetrans usque ad divisiónem ánimæ et spíritus, artuum quoque et medullárum, cogitatiónes cordis et secreta scrutátur animórum, quia nuda et apérta sunt ómnia Dei Fílio, quem consciéntiæ secreta non fallunt.\par
\switchcolumn
\EN
\noindent See, a sword shall pierce through thine own soul also. We have no record or tradition that Mary left this world by suffering a violent death, and the material sword can pierce the body only, and not the soul. Wherefore here we see the wisdom of Mary in that she was not ignorant of the heavenly mysteries. For, the word of God is quick, and powerful, and sharper than any two-edged sword, piercing even to the dividing asunder of soul and spirit, and of the joints and marrow, and is a discerner of the thoughts and intents of the heart for all things are naked and opened unto the eyes of the Son of God, from Whom also the secret things of our conscience are not hidden.\par
\switchcolumn*

\LA
\begin{Resp}\Rsign O Regem cæli, cui talia famulántur obsequia: stábulo pónitur, qui cóntinet mundum: \Star{} Jacet in præsépio, et in cælis regnat.\end{Resp}
\begin{Resp}\Vsign Natus est nobis hódie Salvator, qui est Christus Dóminus, in civitáte David.\end{Resp}
\begin{Resp}\Rsign Jacet in præsépio, et in cælis regnat.\end{Resp}
\begin{Resp}\Vsign Glória Patri, et Fílio, \Star{} et Spirítui Sancto.\end{Resp}
\begin{Resp}\Rsign Jacet in præsépio, et in cælis regnat.\end{Resp}
\switchcolumn
\EN
\begin{Resp}\Rsign How is the King of heaven attended? He that containeth the world is laid in a stable; \Star{} Lying in a manger, reigning in heaven.\end{Resp}
\begin{Resp}\Vsign Unto us is born, this day, in the city of David, a Saviour, Which is Christ the Lord.\end{Resp}
\begin{Resp}\Rsign Lying in a manger, reigning in heaven.\end{Resp}
\begin{Resp}\Vsign Glory be to the Father, and to the Son, \Star{} and to the Holy Ghost.\end{Resp}
\begin{Resp}\Rsign Lying in a manger, reigning in heaven.\end{Resp}
\switchcolumn*

\LA
\LessonHead{Lectio IX}
\switchcolumn
\EN
\LessonHead{Lesson IX}
\switchcolumn*

\LA
\noindent Prophetávit ítaque Simeon, prophetáverat virgo, prophetáverat copuláta conjugio; prophétare debuit étiam vidua, ne qua aut proféssio deésset, aut sexus. Et ideo Anna et stipendiis viduitátis, et móribus talis inducitur, ut digna plane fuísse credátur, quæ Redemptórem ómnium venisse nuntiáret. Cujus mérita cum álibi descripsérimus, cum víduas hortarémur, hoc loco, quóniam ad ália properámus, non putámus iteránda.\par
\switchcolumn
\EN
\noindent There had been a triple prophecy; the prophecy of Simeon had followed the prophecy of the virgin, and the prophecy of the wife; those, namely, of Mary and Elizabeth. And now ought the widow also to prophesy, that no sex nor state might be wanting. And Anna is brought before us with such a title from her widowhood and her life, that we may well believe that she received the grace to announce the Advent of the Redeemer. In our exhortation addressed to widows we have already treated of her gifts at length, and, as we have much matter before us, we will not now again enter on the subject.\par
\switchcolumn*

\LA
\TeDeum{Te Deum}
\switchcolumn
\EN
\TeDeum{Te Deum}
\switchcolumn*

\end{paracol}

\Office{Die 29 Decembris: Diei V infra Octavam Nativitatis}{29 December}{Semiduplex}
\addcontentsline{toc}{subsection}{Die 29 Decembris: Diei V infra Octavam Nativitatis\enspace\textit{29 December}}
\begin{paracol}{2}
\LA
\LessonHead{Lectio I}
\switchcolumn
\EN
\LessonHead{Lesson I}
\switchcolumn*

\LA
\LessonTitle{Incipit Epístola beáti Pauli Apóstoli ad Romános}
\switchcolumn
\EN
\LessonTitle{Beginning of the letter of St. Paul the Apostle to the Romans}
\switchcolumn*

\LA
\Cite{Rom 1:1-7}
\switchcolumn
\EN
\Cite{Rom 1:1-7}
\switchcolumn*

\LA
\noindent Paulus, servus Jesu Christi, vocátus Apóstolus, segregátus in Evangélium Dei, quod ante promíserat per prophétas suos in Scriptúris sanctis de Fílio suo, qui factus est ei ex sémine David secúndum carnem, qui prædestinátus est Fílius Dei in virtúte secúndum spíritum sanctificatiónis ex resurrectióne mortuórum Jesu Christi Dómini nostri: per quem accépimus grátiam, et apostolátum ad obediéndum fídei in ómnibus géntibus pro nómine ejus, in quibus estis et vos vocáti Jesu Christi: ómnibus qui sunt Romæ, diléctis Dei, vocátis sanctis. Grátia vobis, et pax a Deo Patre nostro, et Dómino Jesu Christo.\par
\switchcolumn
\EN
\noindent Paul, a servant of Jesus Christ, called to be an apostle, separated unto the gospel of God, Which he had promised before, by his prophets, in the holy scriptures, Concerning his Son, who was made to him of the seed of David, according to the flesh, Who was predestinated the Son of God in power, according to the spirit of sanctification, by the resurrection of our Lord Jesus Christ from the dead; By whom we have received grace and apostleship for obedience to the faith, in all nations, for his name; Among whom are you also the called of Jesus Christ: To all that are at Rome, the beloved of God, called to be saints. Grace to you, and peace from God our Father, and from the Lord Jesus Christ.\par
\switchcolumn*

\LA
\begin{Resp}\Rsign Hódie nobis de cælo pax vera descéndit: \Star{} Hódie per totum mundum mellíflui facti sunt cæli.\end{Resp}
\begin{Resp}\Vsign Hódie illúxit nobis dies redemptiónis novæ, reparatiónis antíquæ, felicitátis ætérnæ.\end{Resp}
\begin{Resp}\Rsign Hódie per totum mundum mellíflui facti sunt cæli.\end{Resp}
\switchcolumn
\EN
\begin{Resp}\Rsign This day is the true peace come down unto us from heaven. \Star{} This day throughout the whole world the skies drop down sweetness.\end{Resp}
\begin{Resp}\Vsign This day is the daybreak of our new redemption, of the restoring of the old, of everlasting joy.\end{Resp}
\begin{Resp}\Rsign This day throughout the whole world the skies drop down sweetness.\end{Resp}
\switchcolumn*

\LA
\LessonHead{Lectio II}
\switchcolumn
\EN
\LessonHead{Lesson II}
\switchcolumn*

\LA
\Cite{Rom 1:8-12}
\switchcolumn
\EN
\Cite{Rom 1:8-12}
\switchcolumn*

\LA
\noindent Primum quidem grátias ago Deo meo per Jesum Christum pro ómnibus vobis: quia fides vestra annuntiátur in univérso mundo. Testis enim mihi est Deus, cui sérvio in spíritu meo in Evangélio Fílii ejus, quod sine intermissióne memóriam vestri fácio semper in oratiónibus meis: óbsecrans, si quómodo tandem aliquándo prósperum iter hábeam in voluntáte Dei veniéndi ad vos. Desídero enim vidére vos, ut áliquid impértiar vobis grátiæ spirituális ad confirmándos vos: id est, simul consolári in vobis per eam quæ ínvicem est, fidem vestram atque meam.\par
\switchcolumn
\EN
\noindent First I give thanks to my God, through Jesus Christ, for you all, because your faith is spoken of in the whole world. For God is my witness, whom I serve in my spirit in the gospel of his Son, that without ceasing I make a commemoration of you; Always in my prayers making request, if by any means now at length I may have a prosperous journey, by the will of God, to come unto you. For I long to see you, that I may impart unto you some spiritual grace, to strengthen you: That is to say, that I may be comforted together in you, by that which is common to us both, your faith and mine.\par
\switchcolumn*

\LA
\begin{Resp}\Rsign Quem vidístis, pastóres? dícite, annuntiáte nobis, in terris quis appáruit? \Star{} Natum vídimus, et choros Angelórum collaudántes Dóminum.\end{Resp}
\begin{Resp}\Vsign Dícite, quidnam vidístis? et annuntiáte Christi nativitátem.\end{Resp}
\begin{Resp}\Rsign Natum vídimus, et choros Angelórum collaudántes Dóminum.\end{Resp}
\switchcolumn
\EN
\begin{Resp}\Rsign O ye shepherds, speak, and tell us what ye have seen; who is appeared in the earth? \Star{} We saw the new-born Child, and Angels singing praise to the Lord.\end{Resp}
\begin{Resp}\Vsign Speak; what have ye seen? And tell us of the Birth of Christ.\end{Resp}
\begin{Resp}\Rsign We saw the new-born Child, and Angels singing praise to the Lord.\end{Resp}
\switchcolumn*

\LA
\LessonHead{Lectio III}
\switchcolumn
\EN
\LessonHead{Lesson III}
\switchcolumn*

\LA
\Cite{Rom 1:13-19}
\switchcolumn
\EN
\Cite{Rom 1:13-19}
\switchcolumn*

\LA
\noindent Nolo autem vos ignoráre fratres: quia sæpe propósui veníre ad vos (et prohíbitus sum usque adhuc) ut áliquem fructum hábeam et in vobis, sicut et in céteris géntibus. Græcis ac bárbaris, sapiéntibus, et insipiéntibus débitor sum: ita (quod in me) promptum est et vobis, qui Romæ estis, evangelizáre. Non enim erubésco Evangélium. Virtus enim Dei est in salútem omni credénti, Judǽo primum, et Græco. Justítia enim Dei in eo revelátur ex fide in fidem: sicut scriptum est: Justus autem ex fide vivit. Revelátur enim ira Dei de cælo super omnem impietátem, et injustítiam hóminum eórum, qui veritátem Dei in injustítia détinent: quia quod notum est Dei, maniféstum est in illis. Deus enim illis manifestávit.\par
\switchcolumn
\EN
\noindent And I would not have you ignorant, brethren, that I have often purposed to come unto you, (and have been hindered hitherto,) that I might have some fruit among you also, even as among other Gentiles. To the Greeks and to the barbarians, to the wise and to the unwise, I am a debtor; So (as much as is in me) I am ready to preach the gospel to you also that are at Rome. For I am not ashamed of the gospel. For it is the power of God unto salvation to every one that believeth, to the Jew first, and to the Greek. For the justice of God is revealed therein, from faith unto faith, as it is written: The just man liveth by faith. For the wrath of God is revealed from heaven against all ungodliness and injustice of those men that detain the truth of God in injustice: Because that which is known of God is manifest in them. For God hath manifested it unto them.\par
\switchcolumn*

\LA
\begin{Resp}\Rsign O magnum mystérium, et admirábile sacraméntum, ut animália vidérent Dóminum natum, jacéntem in præsépio: \Star{} Beáta Virgo, cujus víscera meruérunt portáre Dóminum Christum.\end{Resp}
\begin{Resp}\Vsign Ave, María, grátia plena: Dóminus tecum.\end{Resp}
\begin{Resp}\Rsign Beáta Virgo, cujus víscera meruérunt portáre Dóminum Christum.\end{Resp}
\begin{Resp}\Vsign Glória Patri, et Fílio, \Star{} et Spirítui Sancto.\end{Resp}
\begin{Resp}\Rsign Beáta Virgo, cujus víscera meruérunt portáre Dóminum Christum.\end{Resp}
\switchcolumn
\EN
\begin{Resp}\Rsign How great is this mystery, how wonderful is the teaching of the faith! The beasts saw the new-born Lord lying in a manger. \Star{} Blessed is that Virgin whose womb was made meet to bear our Lord Christ.\end{Resp}
\begin{Resp}\Vsign Hail, Mary, full of grace the Lord is with thee.\end{Resp}
\begin{Resp}\Rsign Blessed is that Virgin whose womb was made meet to bear our Lord Christ.\end{Resp}
\begin{Resp}\Vsign Glory be to the Father, and to the Son, \Star{} and to the Holy Ghost.\end{Resp}
\begin{Resp}\Rsign Blessed is that Virgin whose womb was made meet to bear our Lord Christ.\end{Resp}
\switchcolumn*

\end{paracol}

\Office{Die 30 Decembris: Diei VI infra Octavam Nativitatis}{30 December}{Semiduplex}
\addcontentsline{toc}{subsection}{Die 30 Decembris: Diei VI infra Octavam Nativitatis\enspace\textit{30 December}}
\begin{paracol}{2}
\LA
\LessonHead{Lectio I}
\switchcolumn
\EN
\LessonHead{Lesson I}
\switchcolumn*

\LA
\LessonTitle{De Epístola ad Romános}
\switchcolumn
\EN
\LessonTitle{Lesson from the letter of St. Paul the Apostle to the Romans}
\switchcolumn*

\LA
\Cite{Rom 2:1-4}
\switchcolumn
\EN
\Cite{Rom 2:1-4}
\switchcolumn*

\LA
\noindent Propter quod inexcusábilis es, o homo omnis qui júdicas. In quo enim júdicas álterum, teípsum condémnas: éadem enim agis quæ júdicas. Scimus enim quóniam judícium Dei est secúndum veritátem in eos qui tália agunt. Exístimas autem hoc, o homo, qui júdicas eos qui tália agunt, et facis ea, quia tu effúgies judícium Dei? an divítias bonitátis ejus, et patiéntiæ, et longanimitátis contémnis? Ignóras quóniam benígnitas Dei ad pœniténtiam te addúcit?\par
\switchcolumn
\EN
\noindent Wherefore thou art inexcusable, O man, whosoever thou art that judgest. For wherein thou judgest another, thou condemnest thyself. For thou dost the same things which thou judgest. For we know that the judgment of God is, according to truth, against them that do such things. And thinkest thou this, O man, that judgest them who do such things, and dost the same, that thou shalt escape the judgment of God? Or despisest thou the riches of his goodness, and patience, and longsuffering? Knowest thou not, that the benignity of God leadeth thee to penance?\par
\switchcolumn*

\LA
\begin{Resp}\Rsign Hódie nobis de cælo pax vera descéndit: \Star{} Hódie per totum mundum mellíflui facti sunt cæli.\end{Resp}
\begin{Resp}\Vsign Hódie illúxit nobis dies redemptiónis novæ, reparatiónis antíquæ, felicitátis ætérnæ.\end{Resp}
\begin{Resp}\Rsign Hódie per totum mundum mellíflui facti sunt cæli.\end{Resp}
\switchcolumn
\EN
\begin{Resp}\Rsign This day is the true peace come down unto us from heaven. \Star{} This day throughout the whole world the skies drop down sweetness.\end{Resp}
\begin{Resp}\Vsign This day is the daybreak of our new redemption, of the restoring of the old, of everlasting joy.\end{Resp}
\begin{Resp}\Rsign This day throughout the whole world the skies drop down sweetness.\end{Resp}
\switchcolumn*

\LA
\LessonHead{Lectio II}
\switchcolumn
\EN
\LessonHead{Lesson II}
\switchcolumn*

\LA
\Cite{Rom 2:5-8}
\switchcolumn
\EN
\Cite{Rom 2:5-8}
\switchcolumn*

\LA
\noindent Secúndum autem durítiam tuam, et impǽnitens cor, thesaurízas tibi iram in die iræ, et revelatiónis justi judícii Dei, qui reddet unicuíque secúndum ópera ejus: iis quidem qui secúndum patiéntiam boni óperis, glóriam, et honórem, et incorruptiónem quærunt, vitam ætérnam: iis autem qui sunt ex contentióne, et qui non acquiéscunt veritáti, credunt autem iniquitáti, ira et indignátio.\par
\switchcolumn
\EN
\noindent But according to thy hardness and impenitent heart, thou treasurest up to thyself wrath, against the day of wrath, and revelation of the just judgment of God. Who will render to every man according to his works. To them indeed, who according to patience in good work, seek glory and honour and incorruption, eternal life: But to them that are contentious, and who obey not the truth, but give credit to iniquity, wrath and indignation.\par
\switchcolumn*

\LA
\begin{Resp}\Rsign Quem vidístis, pastóres? dícite, annuntiáte nobis, in terris quis appáruit? \Star{} Natum vídimus, et choros Angelórum collaudántes Dóminum.\end{Resp}
\begin{Resp}\Vsign Dícite, quidnam vidístis? et annuntiáte Christi nativitátem.\end{Resp}
\begin{Resp}\Rsign Natum vídimus, et choros Angelórum collaudántes Dóminum.\end{Resp}
\switchcolumn
\EN
\begin{Resp}\Rsign O ye shepherds, speak, and tell us what ye have seen; who is appeared in the earth? \Star{} We saw the new-born Child, and Angels singing praise to the Lord.\end{Resp}
\begin{Resp}\Vsign Speak; what have ye seen? And tell us of the Birth of Christ.\end{Resp}
\begin{Resp}\Rsign We saw the new-born Child, and Angels singing praise to the Lord.\end{Resp}
\switchcolumn*

\LA
\LessonHead{Lectio III}
\switchcolumn
\EN
\LessonHead{Lesson III}
\switchcolumn*

\LA
\Cite{Rom 2:9-13}
\switchcolumn
\EN
\Cite{Rom 2:9-13}
\switchcolumn*

\LA
\noindent Tribulátio et angústia in omnem ánimam hóminis operántis malum, Judǽi primum, et Græci: glória autem, et honor, et pax omni operánti bonum, Judǽo primum, et Græco: non enim est accéptio personárum apud Deum. Quicúmque enim sine lege peccavérunt, sine lege períbunt: et quicúmque in lege peccavérunt, per legem judicabúntur. Non enim auditóres legis justi sunt apud Deum, sed factóres legis justificabúntur.\par
\switchcolumn
\EN
\noindent Tribulation and anguish upon every soul of man that worketh evil, of the Jew first, and also of the Greek. But glory, and honour, and peace to every one that worketh good, to the Jew first, and also to the Greek. For there is no respect of persons with God. For whosoever have sinned without the law, shall perish without the law; and whosoever have sinned in the law, shall be judged by the law. For not the hearers of the law are just before God, but the doers of the law shall be justified.\par
\switchcolumn*

\LA
\begin{Resp}\Rsign O magnum mystérium, et admirábile sacraméntum, ut animália vidérent Dóminum natum, jacéntem in præsépio: \Star{} Beáta Virgo, cujus víscera meruérunt portáre Dóminum Christum.\end{Resp}
\begin{Resp}\Vsign Ave, María, grátia plena: Dóminus tecum.\end{Resp}
\begin{Resp}\Rsign Beáta Virgo, cujus víscera meruérunt portáre Dóminum Christum.\end{Resp}
\begin{Resp}\Vsign Glória Patri, et Fílio, \Star{} et Spirítui Sancto.\end{Resp}
\begin{Resp}\Rsign Beáta Virgo, cujus víscera meruérunt portáre Dóminum Christum.\end{Resp}
\switchcolumn
\EN
\begin{Resp}\Rsign How great is this mystery, how wonderful is the teaching of the faith! The beasts saw the new-born Lord lying in a manger. \Star{} Blessed is that Virgin whose womb was made meet to bear our Lord Christ.\end{Resp}
\begin{Resp}\Vsign Hail, Mary, full of grace the Lord is with thee.\end{Resp}
\begin{Resp}\Rsign Blessed is that Virgin whose womb was made meet to bear our Lord Christ.\end{Resp}
\begin{Resp}\Vsign Glory be to the Father, and to the Son, \Star{} and to the Holy Ghost.\end{Resp}
\begin{Resp}\Rsign Blessed is that Virgin whose womb was made meet to bear our Lord Christ.\end{Resp}
\switchcolumn*

\LA
\LessonHead{Lectio IV}
\switchcolumn
\EN
\LessonHead{Lesson IV}
\switchcolumn*

\LA
\LessonTitle{Sermo sancti Leónis Papæ}
\switchcolumn
\EN
\LessonTitle{Sermon by St. Leo the Pope}
\switchcolumn*

\LA
\Source{Sermo 6 de Nativitáte Dómini}
\switchcolumn
\EN
\Source{Sermon 6 on the Nativity of the Lord}
\switchcolumn*

\LA
\noindent Omnibus quidem diébus, dilectíssimi, atque tempóribus, ánimis fidélium divína meditántium Dómini et Salvatóris nostri ex Matre Vírgine ortus occúrrit, ut mens, ad confessiónem sui auctóris erécta, sive in gémitu supplicatiónis, sive in exsultatióne laudis, sive in sacrifícii oblatióne versétur, nihil crébrius nihílque fidéntius spiritáli attíngat intúitu, quam quod Deus, Dei Fílius, génitus de Patre coætérno, idem étiam partu est natus humáno. Sed hanc adorándam in cælo et in terra nativitátem nullus nobis dies magis quam hodiérnus insínuat, et nova étiam in eleméntis luce radiánte, coram sénsibus nostris mirábilis sacraménti íngerit claritátem. Non solum enim in memóriam, sed in conspéctum quodámmodo redit Angeli Gabriélis cum María stupénte collóquium, et concéptio de Spíritu Sancto tam mire promíssa quam crédita.\par
\switchcolumn
\EN
\noindent In any day of the year, dearly beloved, whenever we make our meditations, we are mindful of the birth from a Virgin Mother of our Lord and Saviour. Whenever our souls are uplifted in the worship of our Maker, whether we sigh in supplication, rejoice in praise, or offer sacrifice, there is nothing which we more frequently or more confidently set our minds upon than the fact that God, the Son of God, begotten of the co-eternal Father, was also born by a human birth. But on this day his Nativity, which is to be adored both in heaven and on earth, is brought before us as at no other time. For today it is as though a new and radiant light is shining forth in the heavens, in such wise that the brightness of this wondrous mystery is perceived even by our senses. And not only do we call to mind what the Angel Gabriel said to the awe-stricken Mary, but in some sort we seem even to be present at that colloquy when she conceived of the Holy Ghost. And we marvel both at the promise made to her, and at her faith in that promise.\par
\switchcolumn*

\LA
\begin{Resp}\Rsign Beáta Dei Génetrix María, cujus víscera intácta pérmanent: \Star{} Hódie génuit Salvatórem sǽculi.\end{Resp}
\begin{Resp}\Vsign Beáta, quæ crédidit: quóniam perfécta sunt ómnia, quæ dicta sunt ei a Dómino.\end{Resp}
\begin{Resp}\Rsign Hódie génuit Salvatórem sǽculi.\end{Resp}
\switchcolumn
\EN
\begin{Resp}\Rsign Blessed is God's holy Mother, Mary, maiden undefiled. \Star{} This day hath she brought forth the Saviour of the world.\end{Resp}
\begin{Resp}\Vsign Blessed is she that believed; for there is a performance of all those things which were told her from the Lord.\end{Resp}
\begin{Resp}\Rsign This day hath she brought forth the Saviour of the world.\end{Resp}
\switchcolumn*

\LA
\LessonHead{Lectio V}
\switchcolumn
\EN
\LessonHead{Lesson V}
\switchcolumn*

\LA
\noindent Hódie enim auctor mundi éditus est útero virgináli, et qui omnes natúras cóndidit, ejus est factus fílius, quam creávit. Hódie Verbum Dei carne appáruit vestítum, et quod numquam fuit humánis óculis visíbile, cœpit étiam mánibus esse tractábile. Hódie génitum in nostræ carnis animǽque substántia Salvatórem angélicis vócibus didicére pastóres, et apud Dominicórum prǽsules gregum hódie evangelizándi forma præcóndita est; ut nos quoque cum cæléstis milítiæ dicámus exércitu: Glória in excélsis Deo, et in terra pax homínibus bonæ voluntátis.\par
\switchcolumn
\EN
\noindent For as of today the Maker of the world was brought forth from a virginal womb, and he who made all things became the Son of her whom he had made. As of today the Word of God appeared in a garment of flesh, and that which was never beheld by men's eyes can now be even touched by their hands. As of today the shepherds learned from angelic voices that a Saviour was born in substance of our flesh and soul. And this same angelic message was a pattern to the pastors of the Lord's flock to preach the Gospel on this day, and to do it in such wise that we too may say with the heavenly hosts: Glory to God in the highest, and on earth peace to men of good will.\par
\switchcolumn*

\LA
\begin{Resp}\Rsign Sancta et immaculáta virgínitas, quibus te láudibus éfferam, néscio: \Star{} Quia quem cæli cápere non póterant, tuo grémio contulísti.\end{Resp}
\begin{Resp}\Vsign Benedícta tu in muliéribus, et benedíctus fructus ventris tui.\end{Resp}
\begin{Resp}\Rsign Quia quem cæli cápere non póterant, tuo grémio contulísti.\end{Resp}
\switchcolumn
\EN
\begin{Resp}\Rsign O Mary, how holy and how spotless is thy virginity! I am too dull to praise thee! \Star{} For thou hast borne in thy breast Him Whom the heavens cannot contain.\end{Resp}
\begin{Resp}\Vsign Blessed art thou among women, and blessed is the fruit of thy womb.\end{Resp}
\begin{Resp}\Rsign For thou hast borne in thy breast Him Whom the heavens cannot contain.\end{Resp}
\switchcolumn*

\LA
\LessonHead{Lectio VI}
\switchcolumn
\EN
\LessonHead{Lesson VI}
\switchcolumn*

\LA
\noindent Unde ipsa colláti múneris magnitúdo dignam a nobis éxigit suo splendóre reveréntiam. Ideo enim, sicut beátus Apóstolus docet, non spíritum hujus mundi accépimus, sed spíritum qui ex Deo est, ut sciámus quæ a Deo donáta sunt nobis; qui non áliter pie cólitur, nisi id ei, quod ipse tríbuit, offerátur. Quid autem in thesáuro Domínicæ largitátis ad honórem præséntis festi tam cóngruum póssumus inveníre, quam pacem, quæ in nativitáte Dómini prima est Angélico prædicáta concéntu? Ipsa enim est quæ parit fílios Dei, nutrix dilectiónis et génitrix unitátis, réquies beatórum et æternitátis habitáculum; cujus hoc opus próprium et speciále benefícium est, ut jungat Deo quos secérnit de mundo.\par
\switchcolumn
\EN
\noindent Verily the greatness of the gift bestowed upon us demandeth a reverence worthy of its splendour. For, as the blessed Apostle saith, we have received, not the spirit of the world, but the Spirit which is of God, that we may know the things that are given us from God. And him we can devoutly worship only by offering to him that which he bestoweth. And in the treasury of the Lord's bounty, what can we find so appropriate to the honour of the present Feast, as that peace which at the birth of the Lord was first proclaimed by the angelic choir? For peace it is which bringeth forth the children of God. Peace it is also which is the nurse of affection, the mother of unity, the rest of the blessed, and our eternal home. Peace it is whose proper work and special benefit is to join unto God those whom it separateth from the world.\par
\switchcolumn*

\LA
\begin{Resp}\Rsign Beáta víscera Maríæ Vírginis, quæ portavérunt ætérni Patris Fílium: et beáta úbera, quæ lactavérunt Christum Dóminum: \Star{} Qui hódie pro salúte mundi de Vírgine nasci dignátus est.\end{Resp}
\begin{Resp}\Vsign Dies sanctificátus illúxit nobis: veníte, gentes, et adoráte Dóminum.\end{Resp}
\begin{Resp}\Rsign Qui hódie pro salúte mundi de Vírgine nasci dignátus est.\end{Resp}
\begin{Resp}\Vsign Glória Patri, et Fílio, \Star{} et Spirítui Sancto.\end{Resp}
\begin{Resp}\Rsign Qui hódie pro salúte mundi de Vírgine nasci dignátus est.\end{Resp}
\switchcolumn
\EN
\begin{Resp}\Rsign Blessed be the womb of the Virgin Mary, which bore the Son of the Eternal Father, and blessed be the paps which give suck to Christ our Lord. \Star{} This day hath He been pleased for the salvation of the world to be born of a Virgin.\end{Resp}
\begin{Resp}\Vsign This day which is breaking is holy; O come, ye Gentiles, and worship the Lord.\end{Resp}
\begin{Resp}\Rsign This day hath He been pleased for the salvation of the world to be born of a Virgin.\end{Resp}
\begin{Resp}\Vsign Glory be to the Father, and to the Son, \Star{} and to the Holy Ghost.\end{Resp}
\begin{Resp}\Rsign This day hath He been pleased for the salvation of the world to be born of a Virgin.\end{Resp}
\switchcolumn*

\LA
\LessonHead{Lectio VII}
\switchcolumn
\EN
\LessonHead{Lesson VII}
\switchcolumn*

\LA
\LessonTitle{Léctio sancti Evangélii secúndum Lucam}
\switchcolumn
\EN
\LessonTitle{From the Holy Gospel according to Luke}
\switchcolumn*

\LA
\Cite{Luc 2:15-20}
\switchcolumn
\EN
\Cite{Luke 2:15-20}
\switchcolumn*

\LA
\noindent In illo témpore: Pastóres loquebántur ad ínvicem: Transeámus usque Béthlehem, et videámus hoc verbum quod factum est, quod Dóminus osténdit nobis. Et venérunt festinántes. Et réliqua.\par
\switchcolumn
\EN
\noindent In that time, the shepherds said one to another: Let us go over to Bethlehem, and let us see this word that is come to pass, which the Lord hath shewed to us. And they came with haste. And so on.\par
\switchcolumn*

\LA
\LessonTitle{Homilía sancti Ambrósii Epíscopi}
\switchcolumn
\EN
\LessonTitle{Homily by St. Ambrose, Bishop (of Milan.)}
\switchcolumn*

\LA
\Source{Liber 2 in cap. 2. Lucæ, circa med.}
\switchcolumn
\EN
\Source{Bk. ii. on Luke ii.}
\switchcolumn*

\LA
\noindent Vides festináre pastóres; nemo enim cum desídia Christum requírit. Vides pastóres Angelo credidísse: et tu Patri, Fílio, Spirítui Sancto, Angelis, Prophétis et Apóstolis crédere non vis? Vide, quam signánter Scriptúra singulórum libret moménta verbórum. Festínant, inquit, Verbum vidére. Etenim cum caro Dómini vidétur, Verbum vidétur, quod est Fílius.\par
\switchcolumn
\EN
\noindent The shepherds came with haste. This is how every one cometh who is really earnestly seeking Christ. The shepherds believed the angel. Wilt not thou believe Father, Son, and Holy Ghost, Angels, Prophets, and Apostles? Here also remark how carefully every word in the Scripture is chosen. They came with haste to see this Word, (as the original text hath it. ) A Word, indeed; the Word of God. He that saw the Lord's Flesh, saw the Word, that is, God the Son.\par
\switchcolumn*

\LA
\begin{Resp}\Rsign Verbum caro factum est, et habitávit in nobis: \Star{} Et vídimus glóriam ejus, glóriam quasi Unigéniti a Patre, plenum grátiæ et veritátis.\end{Resp}
\begin{Resp}\Vsign Omnia per ipsum facta sunt, et sine ipso factum est nihil.\end{Resp}
\begin{Resp}\Rsign Et vídimus glóriam ejus, glóriam quasi Unigéniti a Patre, plenum grátiæ et veritátis.\end{Resp}
\switchcolumn
\EN
\begin{Resp}\Rsign The Word was made flesh, and dwelt among us. \Star{} And we beheld His glory, the glory as of the Only-begotten of the Father, full of grace and truth.\end{Resp}
\begin{Resp}\Vsign All things were made by Him, and without Him was not anything made.\end{Resp}
\begin{Resp}\Rsign And we beheld His glory, the glory as of the Only-begotten of the Father, full of grace and truth.\end{Resp}
\switchcolumn*

\LA
\LessonHead{Lectio VIII}
\switchcolumn
\EN
\LessonHead{Lesson VIII}
\switchcolumn*

\LA
\noindent Non medíocre fídei tibi hoc videátur exémplum, quod vilis sit persóna pastórum. Certe quo vílior ad prudéntiam, eo pretiósior ad fidem. Non gymnásia choris reférta sapiéntium, sed plebem Dóminus símplicem requisívit, quæ phaleráre audíta et fucáre nescíret. Simplícitas enim quǽritur, non ambítio desiderátur.\par
\switchcolumn
\EN
\noindent Because the office of a shepherd is mean, think not meanly of the example of their faith. Verily, that which is poorest in learning is richest in faith. The Lord seeketh not for schools crowded with wise men, but for a people of a single heart unused to overlay and to disguise what they learn, by vain and superfluous adornments. He will have straightforwardness rather than vain-glory.\par
\switchcolumn*

\LA
\begin{Resp}\Rsign O Regem cæli, cui talia famulántur obsequia: stábulo pónitur, qui cóntinet mundum: \Star{} Jacet in præsépio, et in cælis regnat.\end{Resp}
\begin{Resp}\Vsign Natus est nobis hódie Salvator, qui est Christus Dóminus, in civitáte David.\end{Resp}
\begin{Resp}\Rsign Jacet in præsépio, et in cælis regnat.\end{Resp}
\begin{Resp}\Vsign Glória Patri, et Fílio, \Star{} et Spirítui Sancto.\end{Resp}
\begin{Resp}\Rsign Jacet in præsépio, et in cælis regnat.\end{Resp}
\switchcolumn
\EN
\begin{Resp}\Rsign How is the King of heaven attended? He that containeth the world is laid in a stable; \Star{} Lying in a manger, reigning in heaven.\end{Resp}
\begin{Resp}\Vsign Unto us is born, this day, in the city of David, a Saviour, Which is Christ the Lord.\end{Resp}
\begin{Resp}\Rsign Lying in a manger, reigning in heaven.\end{Resp}
\begin{Resp}\Vsign Glory be to the Father, and to the Son, \Star{} and to the Holy Ghost.\end{Resp}
\begin{Resp}\Rsign Lying in a manger, reigning in heaven.\end{Resp}
\switchcolumn*

\LA
\LessonHead{Lectio IX}
\switchcolumn
\EN
\LessonHead{Lesson IX}
\switchcolumn*

\LA
\noindent Nec contemnénda putes quasi vília verba pastórum. A pastóribus enim María fidem cólligit, a pastóribus pópulus ad Dei reveréntiam congregátur. Miráti étiam sunt omnes qui audiérunt, de iis, quæ dicebántur a pastóribus ad ipsos. Mária autem conservábat ómnia verba hæc, cónferens in corde suo. Discámus sanctæ Vírginis in ómnibus castitátem, quæ non minus ore pudíca quam córpore, arguménta fídei conferébat in corde.\par
\switchcolumn
\EN
\noindent Think not meanly either of the shepherds' words. The shepherds strengthen the faith even of Mary; the shepherds lead God's people to His worship. For, all they that heard it, wondered at those things which were told them by the shepherds. But Mary kept all these things and pondered them in her heart. Let us learn the modesty of the Holy Virgin, that modesty of speech as of body, whereby she laid up in her heart the evidences of her faith.\par
\switchcolumn*

\LA
\TeDeum{Te Deum}
\switchcolumn
\EN
\TeDeum{Te Deum}
\switchcolumn*

\end{paracol}

\Office{Die 31 Decembris: Diei VII infra Octavam Nativitatis}{31 December}{Semiduplex}
\addcontentsline{toc}{subsection}{Die 31 Decembris: Diei VII infra Octavam Nativitatis\enspace\textit{31 December}}
\begin{paracol}{2}
\LA
\LessonHead{Lectio I}
\switchcolumn
\EN
\LessonHead{Lesson I}
\switchcolumn*

\LA
\LessonTitle{De Epístola ad Romános}
\switchcolumn
\EN
\LessonTitle{Lesson from the letter of St. Paul the Apostle to the Romans}
\switchcolumn*

\LA
\Cite{Rom 3:19-22}
\switchcolumn
\EN
\Cite{Rom 3:19-22}
\switchcolumn*

\LA
\noindent Scimus autem quóniam quæcúmque lex lóquitur, iis, qui in lege sunt, lóquitur: ut omne os obstruátur, et súbditus fiat omnis mundus Deo: quia ex opéribus legis non justificábitur omnis caro coram illo. Per legem enim cognítio peccáti. Nunc autem sine lege justítia Dei manifestáta est: testificáta a lege et prophétis. Justítia autem Dei per fidem Jesu Christi in omnes, et super omnes, qui credunt in eum: non enim est distínctio.\par
\switchcolumn
\EN
\noindent Now we know, that what things soever the law speaketh, it speaketh to them that are in the law; that every mouth may be stopped, and all the world may be made subject to God. Because by the works of the law no flesh shall be justified before him. For by the law is the knowledge of sin. But now without the law the justice of God is made manifest, being witnessed by the law and the prophets. Even the justice of God, by faith of Jesus Christ, unto all and upon all them that believe in him: for there is no distinction:\par
\switchcolumn*

\LA
\begin{Resp}\Rsign Hódie nobis de cælo pax vera descéndit: \Star{} Hódie per totum mundum mellíflui facti sunt cæli.\end{Resp}
\begin{Resp}\Vsign Hódie illúxit nobis dies redemptiónis novæ, reparatiónis antíquæ, felicitátis ætérnæ.\end{Resp}
\begin{Resp}\Rsign Hódie per totum mundum mellíflui facti sunt cæli.\end{Resp}
\switchcolumn
\EN
\begin{Resp}\Rsign This day is the true peace come down unto us from heaven. \Star{} This day throughout the whole world the skies drop down sweetness.\end{Resp}
\begin{Resp}\Vsign This day is the daybreak of our new redemption, of the restoring of the old, of everlasting joy.\end{Resp}
\begin{Resp}\Rsign This day throughout the whole world the skies drop down sweetness.\end{Resp}
\switchcolumn*

\LA
\LessonHead{Lectio II}
\switchcolumn
\EN
\LessonHead{Lesson II}
\switchcolumn*

\LA
\Cite{Rom 3:23-26}
\switchcolumn
\EN
\Cite{Rom 3:23-26}
\switchcolumn*

\LA
\noindent Omnes enim peccavérunt et egent glória Dei. Justificáti gratis per grátiam ipsíus, per redemptiónem, quæ est in Christo Jesu, quem propósuit Deus propitiatiónem per fidem in sánguine ipsíus, ad ostensiónem justítiæ suæ propter remissiónem præcedéntium delictórum in sustentatióne Dei, ad ostensiónem justítiæ ejus in hoc témpore: ut sit ipse justus, et justíficans eum, qui est ex fide Jesu Christi.\par
\switchcolumn
\EN
\noindent For all have sinned, and do need the glory of God. Being justified freely by his grace, through the redemption, that is in Christ Jesus, Whom God hath proposed to be a propitiation, through faith in his blood, to the shewing of his justice, for the remission of former sins, Through the forbearance of God, for the shewing of his justice in this time; that he himself may be just, and the justifier of him, who is of the faith of Jesus Christ.\par
\switchcolumn*

\LA
\begin{Resp}\Rsign Quem vidístis, pastóres? dícite, annuntiáte nobis, in terris quis appáruit? \Star{} Natum vídimus, et choros Angelórum collaudántes Dóminum.\end{Resp}
\begin{Resp}\Vsign Dícite, quidnam vidístis? et annuntiáte Christi nativitátem.\end{Resp}
\begin{Resp}\Rsign Natum vídimus, et choros Angelórum collaudántes Dóminum.\end{Resp}
\switchcolumn
\EN
\begin{Resp}\Rsign O ye shepherds, speak, and tell us what ye have seen; who is appeared in the earth? \Star{} We saw the new-born Child, and Angels singing praise to the Lord.\end{Resp}
\begin{Resp}\Vsign Speak; what have ye seen? And tell us of the Birth of Christ.\end{Resp}
\begin{Resp}\Rsign We saw the new-born Child, and Angels singing praise to the Lord.\end{Resp}
\switchcolumn*

\LA
\LessonHead{Lectio III}
\switchcolumn
\EN
\LessonHead{Lesson III}
\switchcolumn*

\LA
\Cite{Rom 3:27-31}
\switchcolumn
\EN
\Cite{Rom 3:27-31}
\switchcolumn*

\LA
\noindent Ubi est ergo gloriátio tua? Exclúsa est. Per quam legem? Factórum? Non: sed per legem fídei. Arbitrámur enim justificári hóminem per fidem sine opéribus legis. An Judæórum Deus tantum? nonne et géntium? Immo et géntium: quóniam quidem unus est Deus, qui justíficat circumcisiónem ex fide, et præpútium per fidem. Legem ergo destrúimus per fidem? Absit: sed legem statúimus.\par
\switchcolumn
\EN
\noindent Where is then thy boasting? It is excluded. By what law? Of works? No, but by the law of faith. For we account a man to be justified by faith, without the works of the law. Is he the God of the Jews only? Is he not also of the Gentiles? Yes, of the Gentiles also. For it is one God, that justifieth circumcision by faith, and uncircumcision through faith. Do we, then, destroy the law through faith? God forbid: but we establish the law.\par
\switchcolumn*

\LA
\begin{Resp}\Rsign O magnum mystérium, et admirábile sacraméntum, ut animália vidérent Dóminum natum, jacéntem in præsépio: \Star{} Beáta Virgo, cujus víscera meruérunt portáre Dóminum Christum.\end{Resp}
\begin{Resp}\Vsign Ave, María, grátia plena: Dóminus tecum.\end{Resp}
\begin{Resp}\Rsign Beáta Virgo, cujus víscera meruérunt portáre Dóminum Christum.\end{Resp}
\begin{Resp}\Vsign Glória Patri, et Fílio, \Star{} et Spirítui Sancto.\end{Resp}
\begin{Resp}\Rsign Beáta Virgo, cujus víscera meruérunt portáre Dóminum Christum.\end{Resp}
\switchcolumn
\EN
\begin{Resp}\Rsign How great is this mystery, how wonderful is the teaching of the faith! The beasts saw the new-born Lord lying in a manger. \Star{} Blessed is that Virgin whose womb was made meet to bear our Lord Christ.\end{Resp}
\begin{Resp}\Vsign Hail, Mary, full of grace the Lord is with thee.\end{Resp}
\begin{Resp}\Rsign Blessed is that Virgin whose womb was made meet to bear our Lord Christ.\end{Resp}
\begin{Resp}\Vsign Glory be to the Father, and to the Son, \Star{} and to the Holy Ghost.\end{Resp}
\begin{Resp}\Rsign Blessed is that Virgin whose womb was made meet to bear our Lord Christ.\end{Resp}
\switchcolumn*

\end{paracol}

\Office{Die 1 Januarii: Ad Romános, cap. 4}{1 January: Epistle to the Romans, chapter 4}{Feria}
\addcontentsline{toc}{subsection}{Die 1 Januarii: Ad Romános, cap. 4\enspace\textit{1 January: Epistle to the Romans, chapter 4}}
\begin{paracol}{2}
\LA
\LessonHead{Lectio I}
\switchcolumn
\EN
\LessonHead{Lesson I}
\switchcolumn*

\LA
\LessonTitle{De Epístola beati Pauli Apóstoli ad Romános}
\switchcolumn
\EN
\LessonTitle{Lesson from the letter of St. Paul the Apostle to the Romans}
\switchcolumn*

\LA
\Cite{Rom 4:1-8}
\switchcolumn
\EN
\Cite{Rom 4:1-8}
\switchcolumn*

\LA
\noindent Quid ergo dicémus invenísse Abraham patrem nostrum secúndum carnem? Si enim Abraham ex opéribus justificátus est, habet glóriam, sed non apud Deum. Quid enim dicit Scriptúra? Crédidit Abraham Deo: et reputátum est illi ad justítiam. Ei autem, qui operátur, merces non imputátur secúndum grátiam, sed secúndum débitum. Ei vero, qui non operátur, credénti autem in eum, qui justíficat ímpium, reputátur fides ejus ad justítiam secúndum propósitum grátiæ Dei. Sicut et David dicit beatitúdinem hóminis, cui Deus accépto fert justítiam sine opéribus: Beáti, quorum remíssæ sunt iniquitátes, et quorum tecta sunt peccáta. Beátus vir, cui non imputávit Dóminus peccátum.\par
\switchcolumn
\EN
\noindent What shall we say then that Abraham hath found, who is our father according to the flesh. For if Abraham were justified by works, he hath whereof to glory, but not before God. For what saith the scripture? Abraham believed God, and it was reputed to him unto justice. Now to him that worketh, the reward is not reckoned according to grace, but according to debt. But to him that worketh not, yet believeth in him that justifieth the ungodly, his faith is reputed to justice, according to the purpose of the grace of God. As David also termeth the blessedness of a man, to whom God reputeth justice without works: Blessed are they whose iniquities are forgiven, and whose sins are covered. Blessed is the man to whom the Lord hath not imputed sin.\par
\switchcolumn*

\LA
\begin{Resp}\Rsign Ecce Agnus Dei, ecce qui tollit peccáta mundi: ecce de quo dicébam vobis: Qui post me venit, ante me factus est: \Star{} Cujus non sum dignus corrígiam calceaménti sólvere.\end{Resp}
\begin{Resp}\Vsign Qui de terra est, de terra lóquitur: qui de cælo venit, super omnes est.\end{Resp}
\begin{Resp}\Rsign Cujus non sum dignus corrígiam calceaménti sólvere.\end{Resp}
\switchcolumn
\EN
\begin{Resp}\Rsign Behold the Lamb of God, behold Him Which taketh away the sins of the world; behold Him of Whom I said unto you: He That cometh after me is preferred before me \Star{} Whose shoe's latchet I am not worthy to unloose.\end{Resp}
\begin{Resp}\Vsign He that is of the earth speaketh of the earth; He That cometh from heaven is above all.\end{Resp}
\begin{Resp}\Rsign Whose shoe's latchet I am not worthy to unloose.\end{Resp}
\switchcolumn*

\LA
\LessonHead{Lectio II}
\switchcolumn
\EN
\LessonHead{Lesson II}
\switchcolumn*

\LA
\Cite{Rom 4:9-12}
\switchcolumn
\EN
\Cite{Rom 4:9-12}
\switchcolumn*

\LA
\noindent Beatitúdo ergo hæc in circumcisióne tantum manet, an étiam in præpútio? Dícimus enim quia reputáta est Abrahæ fides ad justítiam. Quómodo ergo reputáta est? in circumcisióne, an in præpútio? Non in circumcisióne, sed in præpútio. Et signum accépit circumcisiónis, signáculum justítiæ fídei, quæ est in præpútio: ut sit pater ómnium credéntium per præpútium, ut reputétur et illis ad justítiam: et sit pater circumcisiónis non iis tantum, qui sunt ex circumcisióne, sed et iis, qui sectántur vestígia fídei, quæ est in præpútio patris nostri Abrahæ.\par
\switchcolumn
\EN
\noindent This blessedness then, doth it remain in the circumcision only, or in the uncircumcision also? For we say that unto Abraham faith was reputed to justice. How then was it reputed? When he was in circumcision, or in uncircumcision? Not in circumcision, but in uncircumcision. And he received the sign of circumcision, a seal of the justice of the faith, which he had, being uncircumcised; that he might be the father of all them that believe, being uncircumcised, that unto them also it may be reputed to justice: And might be the father of circumcision; not to them only, that are of the circumcision, but to them also that follow the steps of the faithful, that is in the uncircumcision of our father Abraham.\par
\switchcolumn*

\LA
\begin{Resp}\Rsign Dies sanctificátus illúxit nobis: veníte, gentes, et adoráte Dóminum: \Star{} Quia hódie descéndit lux magna in terris.\end{Resp}
\begin{Resp}\Vsign Hæc dies quam fecit Dóminus, exsultémus et lætémur in ea.\end{Resp}
\begin{Resp}\Rsign Quia hódie descéndit lux magna in terris.\end{Resp}
\switchcolumn
\EN
\begin{Resp}\Rsign This day which is breaking is holy; O come, ye Gentiles, and worship the Lord. \Star{} For this day is much light come down unto us from heaven.\end{Resp}
\begin{Resp}\Vsign This is the day which the Lord hath made, let us rejoice and be glad in it.\end{Resp}
\begin{Resp}\Rsign For this day is much light come down unto us from heaven.\end{Resp}
\switchcolumn*

\LA
\LessonHead{Lectio III}
\switchcolumn
\EN
\LessonHead{Lesson III}
\switchcolumn*

\LA
\Cite{Rom 4:13-17}
\switchcolumn
\EN
\Cite{Rom 4:13-17}
\switchcolumn*

\LA
\noindent Non enim per legem promíssio Abrahæ, aut sémini ejus, ut heres esset mundi: sed per justítiam fídei. Si enim qui ex lege, herédes sunt: exinaníta est fides, abólita est promíssio. Lex enim iram operátur. Ubi enim non est lex: nec prævaricátio. Ideo ex fide, ut secúndum grátiam firma sit promíssio omni sémini, non ei, qui ex lege est solum, sed et ei, qui ex fide est Abrahæ, qui pater est ómnium nostrum (sicut scriptum est: Quia patrem multárum géntium pósui te) ante Deum, cui crédidit, qui vivíficat mórtuos, et vocat ea quæ non sunt, tamquam ea quæ sunt.\par
\switchcolumn
\EN
\noindent For not through the law was the promise to Abraham, or to his seed, that he should be heir of the world; but through the justice of faith. For if they who are of the law be heirs, faith is made void, the promise is made of no effect. For the law worketh wrath. For where there is no law, neither is there transgression. Therefore is it of faith, that according to grace the promise might be firm to all the seed; not to that only which is of the law, but to that also which is of the faith of Abraham, who is the father of us all, As it is written: I have made thee a father of many nations, before God, whom he believed, who quickeneth the dead; and calleth those things that are not, as those that are.\par
\switchcolumn*

\LA
\begin{Resp}\Rsign Benedíctus qui venit in nómine Dómini, Deus Dóminus, et illúxit nobis: \Star{} Allelúja, allelúja.\end{Resp}
\begin{Resp}\Vsign Hæc dies quam fecit Dóminus, exsultémus et lætémur in ea.\end{Resp}
\begin{Resp}\Rsign Allelúja, allelúja.\end{Resp}
\begin{Resp}\Vsign Glória Patri, et Fílio, \Star{} et Spirítui Sancto.\end{Resp}
\begin{Resp}\Rsign Allelúja, allelúja.\end{Resp}
\switchcolumn
\EN
\begin{Resp}\Rsign Blessed be He That cometh in the name of the Lord! God is the Lord Who hath showed us light. \Star{} Alleluia, Alleluia.\end{Resp}
\begin{Resp}\Vsign This is the day which the Lord hath made, let us rejoice and be glad in it.\end{Resp}
\begin{Resp}\Rsign Alleluia, Alleluia.\end{Resp}
\begin{Resp}\Vsign Glory be to the Father, and to the Son, \Star{} and to the Holy Ghost.\end{Resp}
\begin{Resp}\Rsign Alleluia, Alleluia.\end{Resp}
\switchcolumn*

\end{paracol}

\Office{Sanctissimi Nominis Jesu}{Most Holy Name of Jesus}{Duplex II classis}
\addcontentsline{toc}{subsection}{Sanctissimi Nominis Jesu\enspace\textit{Most Holy Name of Jesus}}
\begin{paracol}{2}
\LA
\RefLine{Lectio IX: de In Octava S. Stephani Protomartyris (01-02).}
\switchcolumn
\EN
\RefLine{Lesson IX: of Octave of St. Stephen (01-02).}
\switchcolumn*

\LA
\LessonHead{Lectio I}
\switchcolumn
\EN
\LessonHead{Lesson I}
\switchcolumn*

\LA
\LessonTitle{De Actibus Apostolórum}
\switchcolumn
\EN
\LessonTitle{Lesson from the Acts of the Apostles}
\switchcolumn*

\LA
\Cite{Act 3:1-8}
\switchcolumn
\EN
\Cite{Acts 3:1-8}
\switchcolumn*

\LA
\noindent Petrus autem et Joánnes ascendébant in templum ad horam oratiónis nonam. Et quidam vir, qui erat claudus ex útero matris suæ, bajulabátur; quem ponébant cotídie ad portam templi, quæ dícitur Speciósa, ut péteret eleemósynam ab introëúntibus in templum. Is, cum vidísset Petrum et Joánnem incipiéntes introíre in templum, rogábat ut eleemósynam accíperet. Intuens autem in eum Petrus cum Joánne, dixit: Réspice in nos. At ille intendébat in eos, sperans se áliquid acceptúrum ab eis. Petrus autem dixit: Argéntum et aurum non est mihi: quod autem hábeo, hoc tibi do: In nómine Jesu Christi Nazaréni surge, et ámbula. Et, apprehénsa manu ejus déxtera, allevávit eum, et prótinus consolidátæ sunt bases ejus, et plantæ. Et exsíliens stetit, et ambulábat; et intrávit cum illis in templum ámbulans, et exsíliens, et laudans Deum.\par
\switchcolumn
\EN
\noindent Now Peter and John went up into the temple at the ninth hour of prayer. And a certain man who was lame from his mother's womb, was carried: whom they laid every day at the gate of the temple, which is called Beautiful, that he might ask alms of them that went into the temple. He, when he had seen Peter and John about to go into the temple, asked to receive an alms. But Peter with John fastening his eyes upon him, said: Look upon us. But he looked earnestly upon them, hoping that he should receive something of them. But Peter said: Silver and gold I have none; but what I have, I give thee: In the name of Jesus Christ of Nazareth, arise, and walk. And taking him by the right hand, he lifted him up, and forthwith his feet and soles received strength. And he leaping up, stood, and walked, and went in with them into the temple, walking, and leaping, and praising God.\par
\switchcolumn*

\LA
\begin{Resp}\Rsign Ecce concípies, et paries fílium, et vocábis nomen ejus Jesum: \Star{} Ipse enim salvum fáciet pópulum suum a peccátis eórum.\end{Resp}
\begin{Resp}\Vsign Vocátum est nomen ejus Jesus, quod vocátum est ab Angelo, priúsquam in útero conciperétur.\end{Resp}
\begin{Resp}\Rsign Ipse enim salvum fáciet pópulum suum a peccátis eórum.\end{Resp}
\switchcolumn
\EN
\begin{Resp}\Rsign Behold, thou shalt conceive, and bring forth a Son, and shalt call His Name Jesus \Star{} For He shall save His people from their sins.\end{Resp}
\begin{Resp}\Vsign His Name was called Jesus, which was so named of the Angel, before He was conceived in the womb.\end{Resp}
\begin{Resp}\Rsign For He shall save His people from their sins.\end{Resp}
\switchcolumn*

\LA
\LessonHead{Lectio II}
\switchcolumn
\EN
\LessonHead{Lesson II}
\switchcolumn*

\LA
\Cite{Act 3:9-16}
\switchcolumn
\EN
\Cite{Acts 3:9-16}
\switchcolumn*

\LA
\noindent Et vidit omnis pópulus eum ambulántem, et laudántem Deum. Cognoscébant autem illum, quod ipse erat, qui ad eleemósynam sedébat ad Speciósam portam templi; et impléti sunt stupóre et éxtasi in eo, quod contígerat illi. Cum tenéret autem Petrum et Joánnem, cucúrrit omnis pópulus ad eos ad pórticum, quæ appellátur Salomónis, stupéntes. Videns autem Petrus, respóndit ad pópulum: Viri Israëlítæ, quid mirámini in hoc, aut nos quid intuémini, quasi nostra virtúte aut potestáte fecérimus hunc ambuláre? Deus Abraham, et Deus Isaac, et Deus Jacob, Deus patrum nostrórum glorificávit Fílium suum Jesum, quem vos quidem tradidístis et negástis ante fáciem Piláti, judicánte illo dimítti. Vos autem sanctum et justum negástis, et petístis virum homicídam donári vobis: Auctórem vero vitæ interfecístis, quem Deus suscitávit a mórtuis, cujus nos testes sumus. Et in fide nóminis ejus, hunc, quem vos vidístis et nostis, confirmávit nomen ejus: et fides, quæ per eum est, dedit íntegram sanitátem istam in conspéctu ómnium vestrum.\par
\switchcolumn
\EN
\noindent And all the people saw him walking and praising God. And they knew him, that it was he who sat begging alms at the Beautiful gate of the temple: and they were filled with wonder and amazement at that which had happened to him. And as he held Peter and John, all the people ran to them to the porch which is called Solomon's, greatly wondering. But Peter seeing, made answer to the people: Ye men of Israel, why wonder you at this? or why look you upon us, as if by our strength or power we had made this man to walk? The God of Abraham, and the God of Isaac, and the God of Jacob, the God of our fathers, hath glorified his Son Jesus, whom you indeed delivered up and denied before the face of Pilate, when he judged he should be released. But you denied the Holy One and the Just, and desired a murderer to be granted unto you. But the author of life you killed, whom God hath raised from the dead, of which we are witnesses. And in the faith of his name, this man, whom you have seen and known, hath his name strengthened; and the faith which is by him, hath given this perfect soundness in the sight of you all.\par
\switchcolumn*

\LA
\begin{Resp}\Rsign Benedíctum est nomen tuum, Deus patrum nostrórum, qui cum irátus fúeris, misericórdiæ recordáberis, \Star{} Et in témpore tribulatiónis peccáta dimíttis.\end{Resp}
\begin{Resp}\Vsign Et benedíctum nomen majestátis tuæ in ætérnum, qui facis mirabília solus.\end{Resp}
\begin{Resp}\Rsign Et in témpore tribulatiónis peccáta dimíttis.\end{Resp}
\switchcolumn
\EN
\begin{Resp}\Rsign Blessed is thy Name, O God of our fathers, for in wrath Thou wilt remember mercy. \Star{} And in the time of tribulation Thou forgivest the sins of them that call upon thee.\end{Resp}
\begin{Resp}\Vsign And blessed be thy glorious Name for ever, O Thou Who only doest wondrous things!\end{Resp}
\begin{Resp}\Rsign And in the time of tribulation Thou forgivest the sins of them that call upon thee.\end{Resp}
\switchcolumn*

\LA
\LessonHead{Lectio III}
\switchcolumn
\EN
\LessonHead{Lesson III}
\switchcolumn*

\LA
\Cite{Act 4:5-12}
\switchcolumn
\EN
\Cite{Acts 4:5-12}
\switchcolumn*

\LA
\noindent Factum est autem in crástinum, ut congregaréntur príncipes eórum, et senióres, et scribæ, in Jerúsalem: et Annas princeps sacerdótum, et Cáiphas, et Joánnes, et Alexánder, et quotquot erant de génere sacerdotáli. Et statuéntes eos in médio, interrogábant: In qua virtúte, aut in quo nómine fecístis hoc vos? Tunc replétus Spíritu Sancto Petrus, dixit ad eos: Príncipes pópuli, et senióres, audíte: si nos hódie diiudicámur in benefácto hóminis infírmi, in quo iste salvus factus est, notum sit ómnibus vobis, et omni plebi Israël, quia in nómine Dómini nostri Jesu Christi Nazaréni, quem vos crucifixístis, quem Deus suscitávit a mórtuis, in hoc iste adstat coram vobis sanus. Hic est lapis, qui reprobátus est a vobis ædificántibus, qui factus est in caput ánguli: et non est in álio áliquo salus. Nec enim áliud nomen est sub cælo datum homínibus, in quo opórteat nos salvos fíeri.\par
\switchcolumn
\EN
\noindent And it came to pass on the morrow, that their princes, and ancients, and scribes, were gathered together in Jerusalem; And Annas the high priest, and Caiphas, and John, and Alexander, and as many as were of the kindred of the high priest. And setting them in the midst, they asked: By what power, or by what name, have you done this? Then Peter, filled with the Holy Ghost, said to them: Ye princes of the people, and ancients, hear: If we this day are examined concerning the good deed done to the infirm man, by what means he hath been made whole: Be it known to you all, and to all the people of Israel, that by the name of our Lord Jesus Christ of Nazareth, whom you crucified, whom God hath raised from the dead, even by him this man standeth here before you whole. This is the stone which was rejected by you the builders, which is become the head of the corner. Neither is there salvation in any other. For there is no other name under heaven given to men, whereby we must be saved.\par
\switchcolumn*

\LA
\begin{Resp}\Rsign Laudábo nomen tuum assídue, \Star{} Et collaudábo illud in confessióne.\end{Resp}
\begin{Resp}\Vsign Lætábor et exsultábo in te: psallam nómini tuo, Altíssime.\end{Resp}
\begin{Resp}\Rsign Et collaudábo illud in confessióne.\end{Resp}
\begin{Resp}\Vsign Glória Patri, et Fílio, \Star{} et Spirítui Sancto.\end{Resp}
\begin{Resp}\Rsign Et collaudábo illud in confessióne.\end{Resp}
\switchcolumn
\EN
\begin{Resp}\Rsign I will praise thy Name continually; \Star{} Yea, I will praise it with thanksgiving.\end{Resp}
\begin{Resp}\Vsign I will be glad and rejoice in thee I will sing praise to thy Name, O Thou Most High.\end{Resp}
\begin{Resp}\Rsign Yea, I will praise it with thanksgiving.\end{Resp}
\begin{Resp}\Vsign Glory be to the Father, and to the Son, \Star{} and to the Holy Ghost.\end{Resp}
\begin{Resp}\Rsign Yea, I will praise it with thanksgiving.\end{Resp}
\switchcolumn*

\LA
\LessonHead{Lectio IV}
\switchcolumn
\EN
\LessonHead{Lesson IV}
\switchcolumn*

\LA
\LessonTitle{Sermo sancti Bernárdi Abbátis}
\switchcolumn
\EN
\LessonTitle{From the Sermons of St. Bernard, Abbot (of Clairvaux)}
\switchcolumn*

\LA
\Source{Sermo 15 super Cantica, circa medium}
\switchcolumn
\EN
\Source{1st on the Song of Songs}
\switchcolumn*

\LA
\noindent Non otióse Spíritus Sanctus nomen Sponsi óleo cómparat, cum ita dóceat sponsam ad Sponsum clamáre: Oleum effúsum nomen tuum. Oleum enim lucet, pascit et ungit. Fovet ignem, nutrit carnem, lenit dolórem: lux, cibus, medicína. Vide idem nunc et de Sponsi nómine: Lucet prædicátum, pascit recogitátum, invocátum lenit et ungit. Et percurrámus síngula. Unde putas in toto orbe tanta et tam súbita fídei lux, nisi de prædicáto nómine Jesu? Nonne in hujus nóminis luce Deus nos vocávit in admirábile lumen suum; quibus illuminátis, et in lúmine isto vidéntibus lumen, dicat mérito Paulus: Fuístis aliquándo ténebræ, nunc autem lux in Dómino?\par
\switchcolumn
\EN
\noindent It is not idly that the Holy Ghost likeneth the Name of the Bridegroom to oil, when He maketh the Bride say to the Bridegroom: thy Name is as oil poured forth. Oil indeed giveth light, meat, and unction. It feedeth fire, it nourisheth the flesh, it sootheth pain; it is light, food, and healing. Behold, Thus also is the Name of the Bridegroom. To preach it, is to give light; to think of it, is to feed the soul; to call on it, is to win grace and unction. Let us take it point by point. What, thinkest thou, hath made the light of faith so suddenly and so brightly to shine in the whole world but the preaching of the Name of Jesus? Is it not in the light of this Name that God hath called us into His marvellous light, even that light wherewith we being enlightened, and in His light seeing light, Paul saith truly of us: Ye were sometimes darkness, but now are ye light in the Lord.\par
\switchcolumn*

\LA
\begin{Resp}\Rsign Sperent in te, qui novérunt nomen tuum: \Star{} Quóniam non dereliquísti quæréntes te, Dómine.\end{Resp}
\begin{Resp}\Vsign Exsúrge, Dómine, ádjuva nos, et líbera nos propter nomen tuum.\end{Resp}
\begin{Resp}\Rsign Quóniam non dereliquísti quæréntes te, Dómine.\end{Resp}
\switchcolumn
\EN
\begin{Resp}\Rsign Let them that know thy Name put their trust in thee \Star{} For Thou, Lord, hast not forsaken them that seek thee.\end{Resp}
\begin{Resp}\Vsign Arise, O Lord, help us, and redeem us for thy Name's sake.\end{Resp}
\begin{Resp}\Rsign For Thou, Lord, hast not forsaken them that seek thee.\end{Resp}
\switchcolumn*

\LA
\LessonHead{Lectio V}
\switchcolumn
\EN
\LessonHead{Lesson V}
\switchcolumn*

\LA
\noindent Hoc dénique nomen coram régibus, et géntibus, et fíliis Israël portáre jussus est idem Apóstolus; et portábat nomen tamquam lumen, et illuminábat pátriam, et clamábat ubíque: Nox præcéssit, dies autem appropinquávit. Abiciámus ergo ópera tenebrárum, et induámur arma lucis, sicut in die honéste ambulémus. Et monstrábat ómnibus lucérnam super candelábrum, annúntians in omni loco Jesum, et hunc crucifíxum. Quómodo lux ista resplénduit ac perstrínxit cunctórum intuéntium óculos, quando de ore Petri, tamquam fulgur egrédiens, claudi uníus corporáles plantas solidávit et bases, multósque spirituáliter cæcos illuminávit? Numquid non ignem sparsit, cum ait: In nómine Jesu Christi Nazaréni surge et ámbula?\par
\switchcolumn
\EN
\noindent This is the Name which the Apostle was commanded to bear before Gentiles, and kings, and the children of Israel, the Name which he bore as a light to enlighten his people, crying everywhere The night is far spent, the day is at hand; let us therefore cast off the works of darkness, and let us put on the armour of light, let us walk honestly as in the dayligth, He pointed out to all that candle set upon a candlestick, preaching in every place Jesus and Him crucified. How did that Name shine forth and dazzle every eye that beheld it, when it came like lightning out of the mouth of Peter to give bodily strength to the feet of the lame man, and to clear the sight of many a blind soul? Cast he not fire when he said: In the Name of Jesus Christ of Nazareth, rise up and walk?\par
\switchcolumn*

\LA
\begin{Resp}\Rsign Confiteámur nómini tuo magno, \Star{} Quóniam terríbile et sanctum est.\end{Resp}
\begin{Resp}\Vsign Hi in cúrribus, et hi in equis, nos autem in nómine Dómini, Dei nostri, invocábimus.\end{Resp}
\begin{Resp}\Rsign Quóniam terríbile et sanctum est.\end{Resp}
\switchcolumn
\EN
\begin{Resp}\Rsign Let us praise thy great Name, \Star{} For it is terrible and holy.\end{Resp}
\begin{Resp}\Vsign Some trust in chariots, and some in horses, but we will call upon the Name of the Lord our God.\end{Resp}
\begin{Resp}\Rsign For it is terrible and holy.\end{Resp}
\switchcolumn*

\LA
\LessonHead{Lectio VI}
\switchcolumn
\EN
\LessonHead{Lesson VI}
\switchcolumn*

\LA
\noindent Nec tantum lux est nomen Jesu, sed est et cibus. An non tóties confortáris, quóties recordáris? Quid æque mentem cogitántis impínguat? Quid ita exercitátos réparat sensus, virtútes róborat, végetat mores bonos atque honéstos, castas fovet affectiónes? Aridus est omnis ánimæ cibus, si non óleo isto infúnditur; insípidus est, si non hoc sale condítur. Si scribas, non sapit mihi, nisi légero ibi Jesum. Si dísputes aut cónferas, non sapit mihi, nisi sonúerit ibi Jesus. Jesus mel in ore, in aure melos, in corde júbilus. Sed est et medicína. Tristátur áliquis nostrum? Véniat in cor ejus Jesus, et inde sáliat in os. Et ecce ad exórtum nóminis lumen, núbilum omne díffugit, redit serénum. Lábitur quis in crimen? currit ínsuper ad láqueum mortis desperándo? Nonne, si ínvocet nomen vitæ, conféstim respirábit ad vitam?\par
\switchcolumn
\EN
\noindent The Name of Jesus is not a Name of light only, but it is meat also. Dost thou ever call it to mind, and remain unstrengthened? Is there anything like it to enrich the soul of him that thinketh of it? What is there like it to restore the weakened senses, to fortify strength, to give birth to good lives and pure affections? The soul is fed on husks if that whereon it feedeth lack seasoning with this salt. If thou writest, thou hast no meaning for me if I read not of Jesus there. If thou preach, or dispute, thou hast no meaning for me if I hear not of Jesus there. The mention of Jesus is honey in the mouth, music in the ear, and gladness in the heart. It is our healing too. Is any sorrowful among us? Let the thought of Jesus come into his heart, and spring to his mouth. Behold, when the day of that Name beginneth to break, every cloud will flee away, and there will be a great calm. Doth any fall into sin? Doth any draw nigh to an hopeless death? And if he but call on the life-giving Name of Jesus, will he not draw the breath of a new life again?\par
\switchcolumn*

\LA
\begin{Resp}\Rsign Læténtur omnes, qui sperant in te, Dómine, in ætérnum exsultábunt, et habitábis in eis, et gloriabúntur in te omnes, \Star{} Qui díligunt nomen tuum.\end{Resp}
\begin{Resp}\Vsign Dómine, in lúmine vultus tui ambulábunt, et in nómine tuo exsultábunt tota die.\end{Resp}
\begin{Resp}\Rsign Qui díligunt nomen tuum.\end{Resp}
\begin{Resp}\Vsign Glória Patri, et Fílio, \Star{} et Spirítui Sancto.\end{Resp}
\begin{Resp}\Rsign Qui díligunt nomen tuum.\end{Resp}
\switchcolumn
\EN
\begin{Resp}\Rsign O Lord, let all those that put their trust in thee rejoice, let them ever shout for joy because Thou dwellest in them \Star{} Let them also that love thy Name be joyful in thee.\end{Resp}
\begin{Resp}\Vsign They shall walk, O Lord, in the light of thy countenance; and in thy Name shall they rejoice all the day.\end{Resp}
\begin{Resp}\Rsign Let them also that love thy Name be joyful in thee.\end{Resp}
\begin{Resp}\Vsign Glory be to the Father, and to the Son, \Star{} and to the Holy Ghost.\end{Resp}
\begin{Resp}\Rsign Let them also that love thy Name be joyful in thee.\end{Resp}
\switchcolumn*

\LA
\LessonHead{Lectio VII}
\switchcolumn
\EN
\LessonHead{Lesson VII}
\switchcolumn*

\LA
\LessonTitle{Léctio sancti Evangélii secúndum Lucam}
\switchcolumn
\EN
\LessonTitle{From the Holy Gospel according to Luke}
\switchcolumn*

\LA
\Cite{Luc 2:21}
\switchcolumn
\EN
\Cite{Luke 2:21}
\switchcolumn*

\LA
\noindent In illo témpore: Postquam consummáti sunt dies octo, ut circumciderétur Puer: vocátum est nomen ejus Jesus. Et réliqua.\par
\switchcolumn
\EN
\noindent In that time: when eight days were accomplished, that the child should be circumcised, his name was called Jesus, which was called by the angel, before he was conceived in the womb. And so on.\par
\switchcolumn*

\LA
\LessonTitle{Homilía sancti Bernárdi Abbátis}
\switchcolumn
\EN
\LessonTitle{Homily by St. Bernard, Abbot (of Clairvaux)}
\switchcolumn*

\LA
\Source{Sermo 1 de Circumcisione}
\switchcolumn
\EN
\Source{1st on the Circumcision}
\switchcolumn*

\LA
\noindent Magnum et mirábile sacraméntum! Circumcíditur Puer, et vocátur Jesus. Quid sibi vult ista connéxio? Circumcísio nempe magis salvándi, quam Salvatóris esse vidétur; et Salvatórem circumcídere pótius decet, quam circumcídi. Sed agnósce mediatórem Dei et hóminum, qui ab ipso nativitátis suæ exórdio divínis humána sóciat, ima summis. Náscitur ex mulíere, sed cui fecunditátis fructus sic accédat, ut non décidat flos virginitátis. Pannis invólvitur; sed panni ipsi angélicis láudibus honorántur. Abscónditur in præsépio; sed próditur radiánte stella de cælo. Sic et circumcísio veritátem suscéptæ probat humanitátis; et nomen, quod est super omne nomen, glóriam índicat majestátis. Circumcíditur tamquam verus Abrahæ fílius; Jesus vocátur tamquam verus Fílius Dei.\par
\switchcolumn
\EN
\noindent Behold a mystery, great and full of wonder! The Child is circumcised, and His Name is called Jesus. Why are these two things thus mentioned together? It would seem that circumcision should rather be for the saved than for the Saviour; that the Saviour ought rather to be Circumciser than circumcised. But behold here the Mediator between God and men, how even from His childhood He joineth the things of the Highest to the things of the lowest, the things of God to the things of men. He is born of a woman, but her womb is made fruitful without the loss of the flower of her virginity. He is wrapped in swaddling-bands, but these swaddling-bands are a theme for the jubilation of angels. He is laid in a manger, but a bright star standeth in heaven over the place. So also in His circumcision, the ceremony gave proof of the reality of the Manhood which He had taken, and that Name which is above every name proclaimed the glory of His Blessed Majesty. As very son of Abraham He underwent circumcision; He assumed the Name of Jesus as very Son of God.\par
\switchcolumn*

\LA
\begin{Resp}\Rsign Tribulatiónem et dolórem invéni: \Star{} Et nomen Dómini invocávi.\end{Resp}
\begin{Resp}\Vsign Turris fortíssima nomen Dómini, in ipso sperávi, et adjútus sum.\end{Resp}
\begin{Resp}\Rsign Et nomen Dómini invocávi.\end{Resp}
\switchcolumn
\EN
\begin{Resp}\Rsign Sorrow and trouble did I find. \Star{} Then called I upon the Name of the Lord.\end{Resp}
\begin{Resp}\Vsign The Name of the Lord is a strong tower; I trusted in Him and I am helped.\end{Resp}
\begin{Resp}\Rsign Then called I upon the Name of the Lord.\end{Resp}
\switchcolumn*

\LA
\LessonHead{Lectio VIII}
\switchcolumn
\EN
\LessonHead{Lesson VIII}
\switchcolumn*

\LA
\noindent Neque enim ad instar priórum meus iste Jesus nomen vácuum aut ináne portat: non est in eo magni nóminis umbra, sed véritas. Cǽlitus síquidem índitum nomen Evangelísta testátur, quod vocátum est ab Angelo, priúsquam in útero conciperétur. Et atténde verbi profunditátem: Postquam natus est Jesus, Jesus vocátur ab homínibus, qui vocátus est ab Angelo, priúsquam in útero conciperétur. Idem quippe et Angeli Salvátor et hóminis; sed hóminis ab incarnatióne, Angeli ab inítio creatúræ. Vocátum est, ait, nomen ejus Jesus, quod vocátum est ab Angelo. In ore ergo duórum vel trium téstium stat omne verbum; et ipsum, quod in Prophéta abbreviátum, maniféstius in Evangélio légitur caro factum.\par
\switchcolumn
\EN
\noindent Why Jesus beareth not that Name as others have borne it before Him, as a vain and empty title. It is not in Him the shadow of a great Name, but the very meaning of that Name. That His Name was revealed from heaven, is attested by the Evangelist, where it is written, Which was so named of the Angel before He was conceived in the womb. After Jesus was born, men called Him Jesus, but angels called Him Jesus, before He was conceived in the womb. The One Lord is the Saviour of angels and of men; of men, since His Incarnation; of angels, from the beginning of their creation. His Name, saith the Evangelist, was called Jesus, which was so named of the Angel before He was conceived in the womb. In the mouth therefore of two or three witnesses is every word established; and that word whereof the Prophet spoke as cut short, is set forth at length in the Gospel: the Word made Flesh.\par
\switchcolumn*

\LA
\begin{Resp}\Rsign Exspectábo nomen tuum, Dómine, \Star{} Quóniam bonum est in conspéctu sanctórum tuórum.\end{Resp}
\begin{Resp}\Vsign Ut confiteámur nómini sancto tuo, et gloriémur in laude tua.\end{Resp}
\begin{Resp}\Rsign Quóniam bonum est in conspéctu sanctórum tuórum.\end{Resp}
\begin{Resp}\Vsign Glória Patri, et Fílio, \Star{} et Spirítui Sancto.\end{Resp}
\begin{Resp}\Rsign Quóniam bonum est in conspéctu sanctórum tuórum.\end{Resp}
\switchcolumn
\EN
\begin{Resp}\Rsign O Lord, I will wait on thy Name. \Star{} For it is good before thy Saints.\end{Resp}
\begin{Resp}\Vsign To give thanks unto thy Holy Name, and to triumph in thy praise.\end{Resp}
\begin{Resp}\Rsign For it is good before thy Saints.\end{Resp}
\begin{Resp}\Vsign Glory be to the Father, and to the Son, \Star{} and to the Holy Ghost.\end{Resp}
\begin{Resp}\Rsign For it is good before thy Saints.\end{Resp}
\switchcolumn*

\end{paracol}

\Office{Die 2 Januarii: Ad Romános, cap. 5}{2 January: Epistle to the Romans, chapter 5}{Feria}
\addcontentsline{toc}{subsection}{Die 2 Januarii: Ad Romános, cap. 5\enspace\textit{2 January: Epistle to the Romans, chapter 5}}
\begin{paracol}{2}
\LA
\LessonHead{Lectio I}
\switchcolumn
\EN
\LessonHead{Lesson I}
\switchcolumn*

\LA
\LessonTitle{De Epístola ad Romános}
\switchcolumn
\EN
\LessonTitle{Lesson from the letter of St. Paul the Apostle to the Romans}
\switchcolumn*

\LA
\Cite{Rom 5:1-5}
\switchcolumn
\EN
\Cite{Rom 5:1-5}
\switchcolumn*

\LA
\noindent Justificáti ergo ex fide, pacem habeámus ad Deum per Dóminum nostrum Jesum Christum: per quem et habémus accéssum per fidem in grátiam istam, in qua stamus, et gloriámur in spe glóriæ filiórum Dei. Non solum autem, sed et gloriámur in tribulatiónibus: sciéntes quod tribulátio patiéntiam operátur: patiéntia autem probatiónem, probátio vero spem: spes autem non confúndit: quia cáritas Dei diffúsa est in córdibus nostris per Spíritum Sanctum, qui datus est nobis.\par
\switchcolumn
\EN
\noindent Being justified therefore by faith, let us have peace with God, through our Lord Jesus Christ: By whom also we have access through faith into this grace, wherein we stand, and glory in the hope of the glory of the sons of God. And not only so; but we glory also in tribulations, knowing that tribulation worketh patience; And patience trial; and trial hope; And hope confoundeth not: because the charity of God is poured forth in our hearts, by the Holy Ghost, who is given to us.\par
\switchcolumn*

\LA
\begin{Resp}\Rsign Ecce Agnus Dei, ecce qui tollit peccáta mundi: ecce de quo dicébam vobis: Qui post me venit, ante me factus est: \Star{} Cujus non sum dignus corrígiam calceaménti sólvere.\end{Resp}
\begin{Resp}\Vsign Qui de terra est, de terra lóquitur: qui de cælo venit, super omnes est.\end{Resp}
\begin{Resp}\Rsign Cujus non sum dignus corrígiam calceaménti sólvere.\end{Resp}
\switchcolumn
\EN
\begin{Resp}\Rsign Behold the Lamb of God, behold Him Which taketh away the sins of the world; behold Him of Whom I said unto you: He That cometh after me is preferred before me \Star{} Whose shoe's latchet I am not worthy to unloose.\end{Resp}
\begin{Resp}\Vsign He that is of the earth speaketh of the earth; He That cometh from heaven is above all.\end{Resp}
\begin{Resp}\Rsign Whose shoe's latchet I am not worthy to unloose.\end{Resp}
\switchcolumn*

\LA
\LessonHead{Lectio II}
\switchcolumn
\EN
\LessonHead{Lesson II}
\switchcolumn*

\LA
\Cite{Rom 5:6-9}
\switchcolumn
\EN
\Cite{Rom 5:6-9}
\switchcolumn*

\LA
\noindent Ut quid enim Christus, cum adhuc infírmi essémus, secúndum tempus, pro ímpiis mórtuus est? Vix enim pro justo quis móritur: nam pro bono fórsitan quis áudeat mori. Comméndat autem caritátem suam Deus in nobis: quóniam cum adhuc peccatóres essémus, secúndum tempus, Christus pro nobis mórtuus est: multo ígitur magis nunc justificáti in sánguine ipsíus, salvi érimus ab ira per ipsum.\par
\switchcolumn
\EN
\noindent For why did Christ, when as yet we were weak, according to the time, die for the ungodly? For scarce for a just man will one die; yet perhaps for a good man some one would dare to die. But God commendeth his charity towards us; because when as yet we were sinners, according to the time, Christ died for us; much more therefore, being now justified by his blood, shall we be saved from wrath through him.\par
\switchcolumn*

\LA
\begin{Resp}\Rsign Dies sanctificátus illúxit nobis: veníte, gentes, et adoráte Dóminum: \Star{} Quia hódie descéndit lux magna in terris.\end{Resp}
\begin{Resp}\Vsign Hæc dies quam fecit Dóminus, exsultémus et lætémur in ea.\end{Resp}
\begin{Resp}\Rsign Quia hódie descéndit lux magna in terris.\end{Resp}
\begin{Resp}\Vsign Glória Patri, et Fílio, \Star{} et Spirítui Sancto.\end{Resp}
\begin{Resp}\Rsign Quia hódie descéndit lux magna in terris.\end{Resp}
\switchcolumn
\EN
\begin{Resp}\Rsign This day which is breaking is holy; O come, ye Gentiles, and worship the Lord. \Star{} For this day is much light come down unto us from heaven.\end{Resp}
\begin{Resp}\Vsign This is the day which the Lord hath made, let us rejoice and be glad in it.\end{Resp}
\begin{Resp}\Rsign For this day is much light come down unto us from heaven.\end{Resp}
\begin{Resp}\Vsign Glory be to the Father, and to the Son, \Star{} and to the Holy Ghost.\end{Resp}
\begin{Resp}\Rsign For this day is much light come down unto us from heaven.\end{Resp}
\switchcolumn*

\end{paracol}

\Office{Die 3 Januarii: Ad Romános, cap. 6}{3 January: Epistle to the Romans, chapter 6}{Feria}
\addcontentsline{toc}{subsection}{Die 3 Januarii: Ad Romános, cap. 6\enspace\textit{3 January: Epistle to the Romans, chapter 6}}
\begin{paracol}{2}
\LA
\LessonHead{Lectio I}
\switchcolumn
\EN
\LessonHead{Lesson I}
\switchcolumn*

\LA
\LessonTitle{De Epístola ad Romános}
\switchcolumn
\EN
\LessonTitle{Lesson from the letter of St. Paul the Apostle to the Romans}
\switchcolumn*

\LA
\Cite{Rom 6:1-5}
\switchcolumn
\EN
\Cite{Rom 6:1-5}
\switchcolumn*

\LA
\noindent Quid ergo dicémus? permanébimus in peccáto ut grátia abúndet? Absit. Qui enim mórtui sumus peccáto, quómodo adhuc vivémus in illo? an ignorátis quia quicúmque baptizáti sumus in Christo Jesu, in morte ipsíus baptizáti sumus? Consepúlti enim sumus cum illo per baptísmum in mortem: ut quómodo Christus surréxit a mórtuis per glóriam Patris, ita et nos in novitáte vitæ ambulémus. Si enim complantáti facti sumus similitúdini mortis ejus: simul et resurrectiónis érimus.\par
\switchcolumn
\EN
\noindent What shall we say, then? shall we continue in sin, that grace may abound? God forbid. For we that are dead to sin, how shall we live any longer therein? Know you not that all we, who are baptized in Christ Jesus, are baptized in his death? For we are buried together with him by baptism into death; that as Christ is risen from the dead by the glory of the Father, so we also may walk in newness of life. For if we have been planted together in the likeness of his death, we shall be also in the likeness of his resurrection.\par
\switchcolumn*

\LA
\begin{Resp}\Rsign Congratulámini mihi, omnes qui dilígitis Dóminum: \Star{} Quia, cum essem párvula, plácui Altíssimo, et de meis viscéribus génui Deum et hóminem.\end{Resp}
\begin{Resp}\Vsign Beátam me dicent omnes generatiónes, quia ancíllam húmilem respéxit Deus.\end{Resp}
\begin{Resp}\Rsign Quia, cum essem párvula, plácui Altíssimo, et de meis viscéribus génui Deum et hóminem.\end{Resp}
\switchcolumn
\EN
\begin{Resp}\Rsign Rejoice with me, all ye that love the Lord. \Star{} For while I was yet little I pleased the Most High, and from my womb have I brought forth God and man.\end{Resp}
\begin{Resp}\Vsign All generations shall call me blessed, for God hath regarded the lowliness of His hand-maiden.\end{Resp}
\begin{Resp}\Rsign For while I was yet little I pleased the Most High, and from my womb have I brought forth God and man.\end{Resp}
\switchcolumn*

\LA
\LessonHead{Lectio II}
\switchcolumn
\EN
\LessonHead{Lesson II}
\switchcolumn*

\LA
\Cite{Rom 6:6-11}
\switchcolumn
\EN
\Cite{Rom 6:6-11}
\switchcolumn*

\LA
\noindent Hoc sciéntes, quia vetus homo noster simul crucifíxus est, ut destruátur corpus peccáti, et ultra non serviámus peccáto. Qui enim mórtuus est, justificátus est a peccáto. Si autem mórtui sumus cum Christo, crédimus quia simul étiam vivémus cum Christo, sciéntes quod Christus resúrgens ex mórtuis jam non móritur: mors illi ultra non dominábitur. Quod enim mórtuus est peccáto, mórtuus est semel: quod autem vivit, vivit Deo. Ita et vos existimáte vos mórtuos quidem esse peccáto, vivéntes autem Deo, in Christo Jesu Dómino nostro.\par
\switchcolumn
\EN
\noindent Knowing this, that our old man is crucified with him, that the body of sin may be destroyed, to the end that we may serve sin no longer. For he that is dead is justified from sin. Now if we be dead with Christ, we believe that we shall live also together with Christ: Knowing that Christ rising again from the dead, dieth now no more, death shall no more have dominion over him. For in that he died to sin, he died once; but in that he liveth, he liveth unto God: So do you also reckon, that you are dead to sin, but alive unto God, in Christ Jesus our Lord.\par
\switchcolumn*

\LA
\begin{Resp}\Rsign Confirmátum est cor Vírginis, in quo divína mystéria, Angelo nuntiánte, concépit: tunc speciósum forma præ fíliis hóminum castis suscépit viscéribus: \Star{} Et benedícta in ætérnum, Deum nobis prótulit et hóminem.\end{Resp}
\begin{Resp}\Vsign Domus pudíci péctoris templum repénte fit Dei: intácta nésciens virum, verbo concépit Fílium.\end{Resp}
\begin{Resp}\Rsign Et benedícta in ætérnum, Deum nobis prótulit et hóminem.\end{Resp}
\begin{Resp}\Vsign Glória Patri, et Fílio, \Star{} et Spirítui Sancto.\end{Resp}
\begin{Resp}\Rsign Et benedícta in ætérnum, Deum nobis prótulit et hóminem.\end{Resp}
\switchcolumn
\EN
\begin{Resp}\Rsign The heart of the Virgin was fixed, when the Angel declared unto her the mystery of God, and she conceived, then did she receive in her pure womb Him That is fairer than the children of men. \Star{} And, she that is blessed for ever, brought forth for us God and man.\end{Resp}
\begin{Resp}\Vsign Soon rises, in that modest shrine, The Temple of the Lord Divine The stainless and unwedded one, Within her womb conceived the Son.\end{Resp}
\begin{Resp}\Rsign And, she that is blessed for ever, brought forth for us God and man.\end{Resp}
\begin{Resp}\Vsign Glory be to the Father, and to the Son, \Star{} and to the Holy Ghost.\end{Resp}
\begin{Resp}\Rsign And, she that is blessed for ever, brought forth for us God and man.\end{Resp}
\switchcolumn*

\end{paracol}

\Office{Die 4 Januarii: Ad Romános, cap. 7}{4 January: Epistle to the Romans, chapter 7}{Feria}
\addcontentsline{toc}{subsection}{Die 4 Januarii: Ad Romános, cap. 7\enspace\textit{4 January: Epistle to the Romans, chapter 7}}
\begin{paracol}{2}
\LA
\LessonHead{Lectio I}
\switchcolumn
\EN
\LessonHead{Lesson I}
\switchcolumn*

\LA
\LessonTitle{De Epístola ad Romános}
\switchcolumn
\EN
\LessonTitle{Lesson from the letter of St. Paul the Apostle to the Romans}
\switchcolumn*

\LA
\Cite{Rom 7:1-3}
\switchcolumn
\EN
\Cite{Rom 7:1-3}
\switchcolumn*

\LA
\noindent An ignorátis, fratres, (sciéntibus enim legem loquor) quia lex in hómine dominátur, quanto témpore vivit? Nam quæ sub viro est múlier, vivénte viro, alligáta est legi: si autem mórtuus fúerit vir ejus, solúta est a lege viri. Igitur, vivénte viro, vocábitur adúltera si fúerit cum álio viro: si autem mórtuus fúerit vir ejus, liberáta est a lege viri, ut non sit adúltera si fúerit cum álio viro.\par
\switchcolumn
\EN
\noindent Know you not, brethren, for I speak to them that know the law, that the law hath dominion over a man, as long as it liveth? For the woman that hath an husband, whilst her husband liveth is bound to the law. But if her husband be dead, she is loosed from the law of her husband. Therefore, whilst her husband liveth, she shall be called an adulteress, if she be with another man: but if her husband be dead, she is delivered from the law of her husband; so that she is not an adulteress, if she be with another man.\par
\switchcolumn*

\LA
\begin{Resp}\Rsign Sancta et immaculáta virgínitas, quibus te láudibus éfferam, néscio: \Star{} Quia quem cæli cápere non póterant, tuo grémio contulísti.\end{Resp}
\begin{Resp}\Vsign Benedícta tu in muliéribus, et benedíctus fructus ventris tui.\end{Resp}
\begin{Resp}\Rsign Quia quem cæli cápere non póterant, tuo grémio contulísti.\end{Resp}
\switchcolumn
\EN
\begin{Resp}\Rsign O Mary, how holy and how spotless is thy virginity! I am too dull to praise thee! \Star{} For thou hast borne in thy breast Him Whom the heavens cannot contain.\end{Resp}
\begin{Resp}\Vsign Blessed art thou among women, and blessed is the fruit of thy womb.\end{Resp}
\begin{Resp}\Rsign For thou hast borne in thy breast Him Whom the heavens cannot contain.\end{Resp}
\switchcolumn*

\LA
\LessonHead{Lectio II}
\switchcolumn
\EN
\LessonHead{Lesson II}
\switchcolumn*

\LA
\Cite{Rom 7:4-6}
\switchcolumn
\EN
\Cite{Rom 7:4-6}
\switchcolumn*

\LA
\noindent Itaque, fratres mei, et vos mortificáti estis legi per corpus Christi: ut sitis altérius, qui ex mórtuis resurréxit, ut fructificémus Deo. Cum enim essémus in carne, passiónes peccatórum, quæ per legem erant, operabántur in membris nostris, ut fructificárent morti. Nunc autem solúti sumus a lege mortis, in qua detinebámur, ita ut serviámus in novitáte spíritus, et non in vetustáte lítteræ.\par
\switchcolumn
\EN
\noindent Therefore, my brethren, you also are become dead to the law, by the body of Christ; that you may belong to another, who is risen again from the dead, that we may bring forth fruit to God. For when we were in the flesh, the passions of sins, which were by the law, did work in our members, to bring forth fruit unto death. But now we are loosed from the law of death, wherein we were detained; so that we should serve in newness of spirit, and not in the oldness of the letter.\par
\switchcolumn*

\LA
\begin{Resp}\Rsign Nésciens Mater Virgo virum, péperit sine dolóre: \Star{} Salvatórem sæculórum, ipsum Regem Angelórum, sola Virgo lactábat úbere de cælo pleno.\end{Resp}
\begin{Resp}\Vsign Domus pudíci péctoris templum repénte fit Dei: intácta nésciens virum, verbo concépit Fílium.\end{Resp}
\begin{Resp}\Rsign Salvatórem sæculórum, ipsum Regem Angelórum, sola Virgo lactábat úbere de cælo pleno.\end{Resp}
\begin{Resp}\Vsign Glória Patri, et Fílio, \Star{} et Spirítui Sancto.\end{Resp}
\begin{Resp}\Rsign Salvatórem sæculórum, ipsum Regem Angelórum, sola Virgo lactábat úbere de cælo pleno.\end{Resp}
\switchcolumn
\EN
\begin{Resp}\Rsign The Virgin-Mother that knew not a man, bore, but travailed not. \Star{} She fed the Saviour of the world, The King of Angel hosts above, Jesus, our Redeemer blest, From the fountain of her breast.\end{Resp}
\begin{Resp}\Vsign Soon rises in that modest shrine, The Temple of the Lord Divine; The stainless and unwedded one, Within her womb conceived the Son.\end{Resp}
\begin{Resp}\Rsign She fed the Saviour of the world, The King of Angel hosts above, Jesus, our Redeemer blest, From the fountain of her breast.\end{Resp}
\begin{Resp}\Vsign Glory be to the Father, and to the Son, \Star{} and to the Holy Ghost.\end{Resp}
\begin{Resp}\Rsign She fed the Saviour of the world, The King of Angel hosts above, Jesus, our Redeemer blest, From the fountain of her breast.\end{Resp}
\switchcolumn*

\end{paracol}

\Office{Die 5 Januarii: Ad Romános, cap. 8}{5 January: Epistle to the Romans, chapter 8}{Feria}
\addcontentsline{toc}{subsection}{Die 5 Januarii: Ad Romános, cap. 8\enspace\textit{5 January: Epistle to the Romans, chapter 8}}
\begin{paracol}{2}
\LA
\LessonHead{Lectio I}
\switchcolumn
\EN
\LessonHead{Lesson I}
\switchcolumn*

\LA
\LessonTitle{De Epístola ad Romános}
\switchcolumn
\EN
\LessonTitle{Lesson from the letter of St. Paul the Apostle to the Romans}
\switchcolumn*

\LA
\Cite{Rom 8:1-4}
\switchcolumn
\EN
\Cite{Rom 8:1-4}
\switchcolumn*

\LA
\noindent Nihil ergo nunc damnatiónis est iis, qui sunt in Christo Jesu, qui non secúndum carnem ámbulant. Lex enim spíritus vitæ in Christo Jesu liberávit me a lege peccáti et mortis. Nam quod impossíbile erat legi, in quo infirmabátur per carnem: Deus Fílium suum mittens in similitúdinem carnis peccáti, et de peccáto damnávit peccátum in carne, ut justificátio legis implerétur in nobis, qui non secúndum carnem ambulámus, sed secúndum spíritum.\par
\switchcolumn
\EN
\noindent There is now therefore no condemnation to them that are in Christ Jesus, who walk not according to the flesh. For the law of the spirit of life, in Christ Jesus, hath delivered me from the law of sin and of death. For what the law could not do, in that it was weak through the flesh; God sending his own Son, in the likeness of sinful flesh and of sin, hath condemned sin in the flesh; That the justification of the law might be fulfilled in us, who walk not according to the flesh, but according to the spirit.\par
\switchcolumn*

\LA
\begin{Resp}\Rsign Ecce Agnus Dei, ecce qui tollit peccáta mundi: ecce de quo dicébam vobis: Qui post me venit, ante me factus est: \Star{} Cujus non sum dignus corrígiam calceaménti sólvere.\end{Resp}
\begin{Resp}\Vsign Qui de terra est, de terra lóquitur: qui de cælo venit, super omnes est.\end{Resp}
\begin{Resp}\Rsign Cujus non sum dignus corrígiam calceaménti sólvere.\end{Resp}
\switchcolumn
\EN
\begin{Resp}\Rsign Behold the Lamb of God, behold Him Which taketh away the sins of the world; behold Him of Whom I said unto you: He That cometh after me is preferred before me \Star{} Whose shoe's latchet I am not worthy to unloose.\end{Resp}
\begin{Resp}\Vsign He that is of the earth speaketh of the earth; He That cometh from heaven is above all.\end{Resp}
\begin{Resp}\Rsign Whose shoe's latchet I am not worthy to unloose.\end{Resp}
\switchcolumn*

\LA
\LessonHead{Lectio II}
\switchcolumn
\EN
\LessonHead{Lesson II}
\switchcolumn*

\LA
\Cite{Rom 8:5-9}
\switchcolumn
\EN
\Cite{Rom 8:5-9}
\switchcolumn*

\LA
\noindent Qui enim secúndum carnem sunt: quæ carnis sunt, sápiunt. Qui vero secúndum spíritum sunt, quæ sunt spíritus, séntiunt. Nam prudéntia carnis, mors est: prudéntia autem spíritus, vita et pax. Quóniam sapiéntia carnis inimíca est Deo: legi enim Dei non est subjécta: nec enim potest. Qui autem in carne sunt, Deo placére non possunt. Vos autem in carne non estis, sed in spíritu: si tamen Spíritus Dei hábitat in vobis.\par
\switchcolumn
\EN
\noindent For they that are according to the flesh, mind the things that are of the flesh; but they that are according to the spirit, mind the things that are of the spirit. For the wisdom of the flesh is death; but the wisdom of the spirit is life and peace. Because the wisdom of the flesh is an enemy to God; for it is not subject to the law of God, neither can it be. And they who are in the flesh, cannot please God. But you are not in the flesh, but in the spirit, if so be that the Spirit of God dwell in you.\par
\switchcolumn*

\LA
\begin{Resp}\Rsign Dies sanctificátus illúxit nobis: veníte, gentes, et adoráte Dóminum: \Star{} Quia hódie descéndit lux magna in terris.\end{Resp}
\begin{Resp}\Vsign Hæc dies quam fecit Dóminus, exsultémus et lætémur in ea.\end{Resp}
\begin{Resp}\Rsign Quia hódie descéndit lux magna in terris.\end{Resp}
\switchcolumn
\EN
\begin{Resp}\Rsign This day which is breaking is holy; O come, ye Gentiles, and worship the Lord. \Star{} For this day is much light come down unto us from heaven.\end{Resp}
\begin{Resp}\Vsign This is the day which the Lord hath made, let us rejoice and be glad in it.\end{Resp}
\begin{Resp}\Rsign For this day is much light come down unto us from heaven.\end{Resp}
\switchcolumn*

\LA
\LessonHead{Lectio III}
\switchcolumn
\EN
\LessonHead{Lesson III}
\switchcolumn*

\LA
\Cite{Rom 8:9-11}
\switchcolumn
\EN
\Cite{Rom 8:9-11}
\switchcolumn*

\LA
\noindent Si quis autem Spíritum Christi non habet: hic non est ejus. Si autem Christus in vobis est: corpus quidem mórtuum est propter peccátum, spíritus vero vivit propter justificatiónem. Quod si Spíritus ejus, qui suscitávit Jesum a mórtuis, hábitat in vobis: qui suscitávit Jesum Christum a mórtuis, vivificábit, et mortália córpora vestra, propter inhabitántem Spíritum ejus in vobis.\par
\switchcolumn
\EN
\noindent Now if any man have not the Spirit of Christ, he is none of his. And if Christ be in you, the body indeed is dead, because of sin; but the spirit liveth, because of justification. And if the Spirit of him that raised up Jesus from the dead, dwell in you; he that raised up Jesus Christ from the dead, shall quicken also your mortal bodies, because of his Spirit that dwelleth in you.\par
\switchcolumn*

\LA
\begin{Resp}\Rsign Benedíctus qui venit in nómine Dómini, Deus Dóminus, et illúxit nobis: \Star{} Allelúja, allelúja.\end{Resp}
\begin{Resp}\Vsign Hæc dies quam fecit Dóminus, exsultémus et lætémur in ea.\end{Resp}
\begin{Resp}\Rsign Allelúja, allelúja.\end{Resp}
\begin{Resp}\Vsign Glória Patri, et Fílio, \Star{} et Spirítui Sancto.\end{Resp}
\begin{Resp}\Rsign Allelúja, allelúja.\end{Resp}
\switchcolumn
\EN
\begin{Resp}\Rsign Blessed be He That cometh in the name of the Lord! God is the Lord Who hath showed us light. \Star{} Alleluia, Alleluia.\end{Resp}
\begin{Resp}\Vsign This is the day which the Lord hath made, let us rejoice and be glad in it.\end{Resp}
\begin{Resp}\Rsign Alleluia, Alleluia.\end{Resp}
\begin{Resp}\Vsign Glory be to the Father, and to the Son, \Star{} and to the Holy Ghost.\end{Resp}
\begin{Resp}\Rsign Alleluia, Alleluia.\end{Resp}
\switchcolumn*

\end{paracol}

\Office{Die 7 Januarii: Ad Romános, cap. 9}{7 January: Epistle to the Romans, chapter 9}{Feria}
\addcontentsline{toc}{subsection}{Die 7 Januarii: Ad Romános, cap. 9\enspace\textit{7 January: Epistle to the Romans, chapter 9}}
\begin{paracol}{2}
\LA
\LessonHead{Lectio I}
\switchcolumn
\EN
\LessonHead{Lesson I}
\switchcolumn*

\LA
\LessonTitle{De Epístola ad Romános}
\switchcolumn
\EN
\LessonTitle{Lesson from the letter of St. Paul the Apostle to the Romans}
\switchcolumn*

\LA
\Cite{Rom 9:1-5}
\switchcolumn
\EN
\Cite{Rom 9:1-5}
\switchcolumn*

\LA
\noindent Veritátem dico in Christo, non méntior: testimónium mihi perhibénte consciéntia mea in Spíritu Sancto: quóniam tristítia mihi magna est, et contínuus dolor cordi meo. Optábam enim ego ipse anáthema esse a Christo pro frátribus meis, qui sunt cognáti mei secúndum carnem, qui sunt Israëlítæ, quorum adóptio est filiórum, et glória, et testaméntum, et legislátio, et obséquium, et promíssa: quorum patres, et ex quibus est Christus secúndum carnem, qui est super ómnia Deus benedíctus in sǽcula. Amen.\par
\switchcolumn
\EN
\noindent I speak the truth in Christ, I lie not, my conscience bearing me witness in the Holy Ghost: That I have great sadness, and continual sorrow in my heart. For I wished myself to be an anathema from Christ, for my brethren, who are my kinsmen according to the flesh, Who are Israelites, to whom belongeth the adoption as of children, and the glory, and the testament, and the giving of the law, and the service of God, and the promises: Whose are the fathers, and of whom is Christ, according to the flesh, who is over all things, God blessed for ever. Amen.\par
\switchcolumn*

\LA
\begin{Resp}\Rsign Tria sunt múnera pretiósa, quæ obtulérunt Magi Dómino in die ista, et habent in se divína mystéria: \Star{} In auro, ut ostendátur Regis poténtia: in thure, Sacerdótem magnum consídera: et in myrrha, Domínicam sepultúram.\end{Resp}
\begin{Resp}\Vsign Salútis nostræ auctórem Magi veneráti sunt in cunábulis, et de thesáuris suis mýsticas ei múnerum spécies obtulérunt.\end{Resp}
\begin{Resp}\Rsign In auro, ut ostendátur Regis poténtia: in thure, Sacerdótem magnum consídera: et in myrrha, Domínicam sepultúram.\end{Resp}
\switchcolumn
\EN
\begin{Resp}\Rsign There are three precious gifts which the wise men offered unto the Lord on this day, and they speak a mystery of the things of God: \Star{} Gold, to show His kingly power; frankincense, for our Great High Priest; and myrrh, against the Lord's burying.\end{Resp}
\begin{Resp}\Vsign The wise men worshipped the Captain of our Salvation, as He lay in the manger, and when they had opened their treasures, they presented unto Him mystic gifts.\end{Resp}
\begin{Resp}\Rsign Gold, to show His kingly power; frankincense, for our Great High Priest; and myrrh, against the Lord's burying.\end{Resp}
\switchcolumn*

\LA
\LessonHead{Lectio II}
\switchcolumn
\EN
\LessonHead{Lesson II}
\switchcolumn*

\LA
\Cite{Rom 9:6-10}
\switchcolumn
\EN
\Cite{Rom 9:6-10}
\switchcolumn*

\LA
\noindent Non autem quod excíderit verbum Dei. Non enim omnes qui ex Israël sunt, ii sunt Israëlítæ: neque qui semen sunt Abrahæ, omnes fílii: sed in Isaac vocábitur tibi semen: id est, non qui fílii carnis, hi fílii Dei: sed qui fílii sunt promissiónis, æstimántur in sémine. Promissiónis enim verbum hoc est: Secúndum hoc tempus véniam: et erit Saræ fílius. Non solum autem illa: sed et Rebécca ex uno concúbitu habens, Isaac patris nostri.\par
\switchcolumn
\EN
\noindent Not as though the word of God hath miscarried. For all are not Israelites that are of Israel: Neither are all they that are the seed of Abraham, children; but in Isaac shall thy seed be called: That is to say, not they that are the children of the flesh, are the children of God; but they, that are the children of the promise, are accounted for the seed. For this is the word of promise: According to this time will I come; and Sara shall have a son. And not only she. But when Rebecca also had conceived at once, of Isaac our father.\par
\switchcolumn*

\LA
\begin{Resp}\Rsign In colúmbæ spécie Spíritus Sanctus visus est, Patérna vox audíta est: \Star{} Hic est Fílius meus diléctus, in quo mihi bene complácui.\end{Resp}
\begin{Resp}\Vsign Cæli apérti sunt super eum, et vox Patris intónuit.\end{Resp}
\begin{Resp}\Rsign Hic est Fílius meus diléctus, in quo mihi bene complácui.\end{Resp}
\switchcolumn
\EN
\begin{Resp}\Rsign The Holy Ghost appeared in a bodily shape like a dove, and the voice of the Father was heard: \Star{} This is My beloved Son, in Whom I am well pleased.\end{Resp}
\begin{Resp}\Vsign The heavens were opened unto him, and, lo, the voice of the Father was heard, like unto thunder, saying:\end{Resp}
\begin{Resp}\Rsign This is My beloved Son, in Whom I am well pleased.\end{Resp}
\switchcolumn*

\LA
\LessonHead{Lectio III}
\switchcolumn
\EN
\LessonHead{Lesson III}
\switchcolumn*

\LA
\Cite{Rom 9:11-16}
\switchcolumn
\EN
\Cite{Rom 9:11-16}
\switchcolumn*

\LA
\noindent Cum enim nondum nati fuíssent, aut áliquid boni egíssent, aut mali, (ut secúndum electiónem propósitum Dei manéret), non ex opéribus, sed ex vocánte dictum est ei: Quia major sérviet minóri, sicut scriptum est: Jacob diléxi, Esau autem ódio hábui. Quid ergo dicémus? numquid iníquitas apud Deum? Absit. Móysi enim dicit: Miserébor cujus miséreor: et misericórdiam præstábo cujus miserébor. Igitur non voléntis, neque curréntis, sed miseréntis est Dei.\par
\switchcolumn
\EN
\noindent For when the children were not yet born, nor had done any good or evil that the purpose of God, according to election, might stand, Not of works, but of him that calleth, it was said to her: The elder shall serve the younger. As it is written: Jacob I have loved, but Esau I have hated. What shall we say then? Is there injustice with God? God forbid. For he saith to Moses: I will have mercy on whom I will have mercy; and I will shew mercy to whom I will shew mercy. So then it is not of him that willeth, nor of him that runneth, but of God that sheweth mercy.\par
\switchcolumn*

\LA
\begin{Resp}\Rsign Reges Tharsis et ínsulæ múnera ófferent: \Star{} Reges Arabum et Saba dona Dómino Deo addúcent.\end{Resp}
\begin{Resp}\Vsign Omnes de Saba vénient, aurum et thus deferéntes.\end{Resp}
\begin{Resp}\Rsign Reges Arabum et Saba dona Dómino Deo addúcent.\end{Resp}
\begin{Resp}\Vsign Glória Patri, et Fílio, \Star{} et Spirítui Sancto.\end{Resp}
\begin{Resp}\Rsign Reges Arabum et Saba dona Dómino Deo addúcent.\end{Resp}
\switchcolumn
\EN
\begin{Resp}\Rsign The kings of Tarshish and of the isles shall bring presents. \Star{} The kings of Arabia and Saba shall offer gifts unto the Lord God.\end{Resp}
\begin{Resp}\Vsign All they from Saba shall come, they shall bring gold and incense.\end{Resp}
\begin{Resp}\Rsign The kings of Arabia and Saba shall offer gifts unto the Lord God.\end{Resp}
\begin{Resp}\Vsign Glory be to the Father, and to the Son, \Star{} and to the Holy Ghost.\end{Resp}
\begin{Resp}\Rsign The kings of Arabia and Saba shall offer gifts unto the Lord God.\end{Resp}
\switchcolumn*

\end{paracol}

\Office{Die 8 Januarii: Ad Romános, cap. 12}{8 January: Epistle to the Romans, chapter 12}{Feria}
\addcontentsline{toc}{subsection}{Die 8 Januarii: Ad Romános, cap. 12\enspace\textit{8 January: Epistle to the Romans, chapter 12}}
\begin{paracol}{2}
\LA
\LessonHead{Lectio I}
\switchcolumn
\EN
\LessonHead{Lesson I}
\switchcolumn*

\LA
\LessonTitle{De Epístola ad Romános}
\switchcolumn
\EN
\LessonTitle{Lesson from the letter of St. Paul the Apostle to the Romans}
\switchcolumn*

\LA
\Cite{Rom 12:1-3}
\switchcolumn
\EN
\Cite{Rom 12:1-3}
\switchcolumn*

\LA
\noindent Obsecro ítaque vos, fratres, per misericórdiam Dei, ut exhibeátis córpora vestra hóstiam vivéntem, sanctam, Deo placéntem, rationábile obséquium vestrum. Et nolíte conformári huic sǽculo, sed reformámini in novitáte sensus vestri: ut probétis quæ sit volúntas Dei bona, et benéplacens, et perfécta. Dico enim per grátiam quæ data est mihi, ómnibus qui sunt inter vos: Non plus sápere quam opórtet sápere, sed sápere ad sobrietátem: et unicuíque sicut Deus divísit mensúram fídei.\par
\switchcolumn
\EN
\noindent I beseech you therefore, brethren, by the mercy of God, that you present your bodies a living sacrifice, holy, pleasing unto God, your reasonable service. And be not conformed to this world; but be reformed in the newness of your mind, that you may prove what is the good, and the acceptable, and the perfect will of God. For I say, by the grace that is given me, to all that are among you, not to be more wise than it behoveth to be wise, but to be wise unto sobriety, and according as God hath divided to every one the measure of faith.\par
\switchcolumn*

\LA
\begin{Resp}\Rsign Tria sunt múnera pretiósa, quæ obtulérunt Magi Dómino in die ista, et habent in se divína mystéria: \Star{} In auro, ut ostendátur Regis poténtia: in thure, Sacerdótem magnum consídera: et in myrrha, Domínicam sepultúram.\end{Resp}
\begin{Resp}\Vsign Salútis nostræ auctórem Magi veneráti sunt in cunábulis, et de thesáuris suis mýsticas ei múnerum spécies obtulérunt.\end{Resp}
\begin{Resp}\Rsign In auro, ut ostendátur Regis poténtia: in thure, Sacerdótem magnum consídera: et in myrrha, Domínicam sepultúram.\end{Resp}
\switchcolumn
\EN
\begin{Resp}\Rsign There are three precious gifts which the wise men offered unto the Lord on this day, and they speak a mystery of the things of God: \Star{} Gold, to show His kingly power; frankincense, for our Great High Priest; and myrrh, against the Lord's burying.\end{Resp}
\begin{Resp}\Vsign The wise men worshipped the Captain of our Salvation, as He lay in the manger, and when they had opened their treasures, they presented unto Him mystic gifts.\end{Resp}
\begin{Resp}\Rsign Gold, to show His kingly power; frankincense, for our Great High Priest; and myrrh, against the Lord's burying.\end{Resp}
\switchcolumn*

\LA
\LessonHead{Lectio II}
\switchcolumn
\EN
\LessonHead{Lesson II}
\switchcolumn*

\LA
\Cite{Rom 12:4-8}
\switchcolumn
\EN
\Cite{Rom 12:4-8}
\switchcolumn*

\LA
\noindent Sicut enim in uno córpore multa membra habémus, ómnia autem membra non eúndem actum habent: ita multi unum corpus sumus in Christo, sínguli autem alter altérius membra. Habéntes autem donatiónes secúndum grátiam, quæ data est nobis, differéntes: sive prophetíam secúndum ratiónem fídei, sive ministérium in ministrándo, sive qui docet in doctrína, qui exhortátur in exhortándo, qui tríbuit in simplicitáte, qui præest in sollicitúdine, qui miserétur in hilaritáte.\par
\switchcolumn
\EN
\noindent For as in one body we have many members, but all the members have not the same office: So we being many, are one body in Christ, and every one members one of another. And having different gifts, according to the grace that is given us, either prophecy, to be used according to the rule of faith; Or ministry, in ministering; or he that teacheth, in doctrine; He that exhorteth, in exhorting; he that giveth, with simplicity; he that ruleth, with carefulness; he that sheweth mercy, with cheerfulness.\par
\switchcolumn*

\LA
\begin{Resp}\Rsign In colúmbæ spécie Spíritus Sanctus visus est, Patérna vox audíta est: \Star{} Hic est Fílius meus diléctus, in quo mihi bene complácui.\end{Resp}
\begin{Resp}\Vsign Cæli apérti sunt super eum, et vox Patris intónuit.\end{Resp}
\begin{Resp}\Rsign Hic est Fílius meus diléctus, in quo mihi bene complácui.\end{Resp}
\switchcolumn
\EN
\begin{Resp}\Rsign The Holy Ghost appeared in a bodily shape like a dove, and the voice of the Father was heard: \Star{} This is My beloved Son, in Whom I am well pleased.\end{Resp}
\begin{Resp}\Vsign The heavens were opened unto him, and, lo, the voice of the Father was heard, like unto thunder, saying:\end{Resp}
\begin{Resp}\Rsign This is My beloved Son, in Whom I am well pleased.\end{Resp}
\switchcolumn*

\LA
\LessonHead{Lectio III}
\switchcolumn
\EN
\LessonHead{Lesson III}
\switchcolumn*

\LA
\Cite{Rom 12:9-16}
\switchcolumn
\EN
\Cite{Rom 12:9-16}
\switchcolumn*

\LA
\noindent Diléctio sine simulatióne. Odiéntes malum, adhæréntes bono: Caritáte fraternitátis ínvicem diligéntes: Honóre ínvicem præveniéntes: Sollicitúdine non pigri: Spíritu fervéntes: Dómino serviéntes: Spe gaudéntes: In tribulatióne patiéntes: Oratióni instántes: Necessitátibus sanctórum communicántes: Hospitalitátem sectántes. Benedícite persequéntibus vos: benedícite, et nolíte maledícere. Gaudére cum gaudéntibus, flere cum fléntibus: Idípsum ínvicem sentiéntes: Non alta sapiéntes, sed humílibus consentiéntes.\par
\switchcolumn
\EN
\noindent Let love be without dissimulation. Hating that which is evil, cleaving to that which is good. Loving one another with the charity of brotherhood, with honour preventing one another. In carefulness not slothful. In spirit fervent. Serving the Lord. Rejoicing in hope. Patient in tribulation. Instant in prayer. Communicating to the necessities of the saints. Pursuing hospitality. Bless them that persecute you: bless, and curse not. Rejoice with them that rejoice; weep with them that weep. Being of one mind one towards another. Not minding high things, but consenting to the humble. Be not wise in your own conceits.\par
\switchcolumn*

\LA
\begin{Resp}\Rsign Reges Tharsis et ínsulæ múnera ófferent: \Star{} Reges Arabum et Saba dona Dómino Deo addúcent.\end{Resp}
\begin{Resp}\Vsign Omnes de Saba vénient, aurum et thus deferéntes.\end{Resp}
\begin{Resp}\Rsign Reges Arabum et Saba dona Dómino Deo addúcent.\end{Resp}
\begin{Resp}\Vsign Glória Patri, et Fílio, \Star{} et Spirítui Sancto.\end{Resp}
\begin{Resp}\Rsign Reges Arabum et Saba dona Dómino Deo addúcent.\end{Resp}
\switchcolumn
\EN
\begin{Resp}\Rsign The kings of Tarshish and of the isles shall bring presents. \Star{} The kings of Arabia and Saba shall offer gifts unto the Lord God.\end{Resp}
\begin{Resp}\Vsign All they from Saba shall come, they shall bring gold and incense.\end{Resp}
\begin{Resp}\Rsign The kings of Arabia and Saba shall offer gifts unto the Lord God.\end{Resp}
\begin{Resp}\Vsign Glory be to the Father, and to the Son, \Star{} and to the Holy Ghost.\end{Resp}
\begin{Resp}\Rsign The kings of Arabia and Saba shall offer gifts unto the Lord God.\end{Resp}
\switchcolumn*

\end{paracol}

\Office{Die 9 Januarii: Ad Romános, cap. 13}{9 January: Epistle to the Romans, chapter 13}{Feria}
\addcontentsline{toc}{subsection}{Die 9 Januarii: Ad Romános, cap. 13\enspace\textit{9 January: Epistle to the Romans, chapter 13}}
\begin{paracol}{2}
\LA
\LessonHead{Lectio I}
\switchcolumn
\EN
\LessonHead{Lesson I}
\switchcolumn*

\LA
\LessonTitle{De Epístola ad Romános}
\switchcolumn
\EN
\LessonTitle{Lesson from the letter of St. Paul the Apostle to the Romans}
\switchcolumn*

\LA
\Cite{Rom 13:1-4}
\switchcolumn
\EN
\Cite{Rom 13:1-4}
\switchcolumn*

\LA
\noindent Omnis ánima potestátibus sublimióribus súbdita sit: non est enim potéstas nisi a Deo: quæ autem sunt, a Deo ordinátæ sunt. Itaque qui resístit potestáti, Dei ordinatióni resístit. Qui autem resístunt, ipsi sibi damnatiónem acquírunt: nam príncipes non sunt timóri boni óperis, sed mali. Vis autem non timére potestátem? Bonum fac: et habébis laudem ex illa; Dei enim miníster est tibi in bonum.\par
\switchcolumn
\EN
\noindent Let every soul be subject to higher powers: for there is no power but from God: and those that are, are ordained of God. Therefore he that resisteth the power, resisteth the ordinance of God. And they that resist, purchase to themselves damnation. For princes are not a terror to the good work, but to the evil. Wilt thou then not be afraid of the power? Do that which is good: and thou shalt have praise from the same. For he is God's minister to thee, for good. But if thou do that which is evil, fear: for he beareth not the sword in vain.\par
\switchcolumn*

\LA
\begin{Resp}\Rsign Tria sunt múnera pretiósa, quæ obtulérunt Magi Dómino in die ista, et habent in se divína mystéria: \Star{} In auro, ut ostendátur Regis poténtia: in thure, Sacerdótem magnum consídera: et in myrrha, Domínicam sepultúram.\end{Resp}
\begin{Resp}\Vsign Salútis nostræ auctórem Magi veneráti sunt in cunábulis, et de thesáuris suis mýsticas ei múnerum spécies obtulérunt.\end{Resp}
\begin{Resp}\Rsign In auro, ut ostendátur Regis poténtia: in thure, Sacerdótem magnum consídera: et in myrrha, Domínicam sepultúram.\end{Resp}
\switchcolumn
\EN
\begin{Resp}\Rsign There are three precious gifts which the wise men offered unto the Lord on this day, and they speak a mystery of the things of God: \Star{} Gold, to show His kingly power; frankincense, for our Great High Priest; and myrrh, against the Lord's burying.\end{Resp}
\begin{Resp}\Vsign The wise men worshipped the Captain of our Salvation, as He lay in the manger, and when they had opened their treasures, they presented unto Him mystic gifts.\end{Resp}
\begin{Resp}\Rsign Gold, to show His kingly power; frankincense, for our Great High Priest; and myrrh, against the Lord's burying.\end{Resp}
\switchcolumn*

\LA
\LessonHead{Lectio II}
\switchcolumn
\EN
\LessonHead{Lesson II}
\switchcolumn*

\LA
\Cite{Rom 13:4-7}
\switchcolumn
\EN
\Cite{Rom 13:4-7}
\switchcolumn*

\LA
\noindent Si autem malum féceris, time: non enim sine causa gládium portat. Dei enim miníster est: vindex in iram ei, qui malum agit. Ideo necessitáte súbditi estóte non solum propter iram, sed étiam propter consciéntiam. Ideo enim et tribúta præstátis: minístri enim Dei sunt, in hoc ipsum serviéntes. Réddite ómnibus débita: cui tribútum, tribútum: cui vectígal, vectígal: cui timórem, timórem: cui honórem, honórem.\par
\switchcolumn
\EN
\noindent For he is God's minister: an avenger to execute wrath upon him that doth evil. Wherefore be subject of necessity, not only for wrath, but also for conscience' sake. For therefore also you pay tribute. For they are the ministers of God, serving unto this purpose. Render therefore to all men their dues. Tribute, to whom tribute is due: custom, to whom custom: fear, to whom fear: honour, to whom honour.\par
\switchcolumn*

\LA
\begin{Resp}\Rsign In colúmbæ spécie Spíritus Sanctus visus est, Patérna vox audíta est: \Star{} Hic est Fílius meus diléctus, in quo mihi bene complácui.\end{Resp}
\begin{Resp}\Vsign Cæli apérti sunt super eum, et vox Patris intónuit.\end{Resp}
\begin{Resp}\Rsign Hic est Fílius meus diléctus, in quo mihi bene complácui.\end{Resp}
\switchcolumn
\EN
\begin{Resp}\Rsign The Holy Ghost appeared in a bodily shape like a dove, and the voice of the Father was heard: \Star{} This is My beloved Son, in Whom I am well pleased.\end{Resp}
\begin{Resp}\Vsign The heavens were opened unto him, and, lo, the voice of the Father was heard, like unto thunder, saying:\end{Resp}
\begin{Resp}\Rsign This is My beloved Son, in Whom I am well pleased.\end{Resp}
\switchcolumn*

\LA
\LessonHead{Lectio III}
\switchcolumn
\EN
\LessonHead{Lesson III}
\switchcolumn*

\LA
\Cite{Rom 13:8-10}
\switchcolumn
\EN
\Cite{Rom 13:8-10}
\switchcolumn*

\LA
\noindent Némini quidquam debeátis, nisi ut ínvicem diligátis: qui enim díligit próximum, legem implévit. Nam, Non adulterábis, Non occídes, Non furáberis, Non falsum testimónium dices, Non concupísces: et si quod est áliud mandátum, in hoc verbo instaurátur: Díliges próximum tuum sicut teípsum. Diléctio próximi malum non operátur. Plenitúdo ergo legis est diléctio.\par
\switchcolumn
\EN
\noindent Owe no man any thing, but to love one another. For he that loveth his neighbour hath fulfilled the law. For: Thou shalt not commit adultery: Thou shalt not kill: Thou shalt not steal: Thou shalt not bear false witness: Thou shalt not covet. And if there be any other commandment, it is comprised in this word: Thou shalt love thy neighbour as thyself. The love of our neighbour worketh no evil. Love therefore is the fulfilling of the law.\par
\switchcolumn*

\LA
\begin{Resp}\Rsign Reges Tharsis et ínsulæ múnera ófferent: \Star{} Reges Arabum et Saba dona Dómino Deo addúcent.\end{Resp}
\begin{Resp}\Vsign Omnes de Saba vénient, aurum et thus deferéntes.\end{Resp}
\begin{Resp}\Rsign Reges Arabum et Saba dona Dómino Deo addúcent.\end{Resp}
\begin{Resp}\Vsign Glória Patri, et Fílio, \Star{} et Spirítui Sancto.\end{Resp}
\begin{Resp}\Rsign Reges Arabum et Saba dona Dómino Deo addúcent.\end{Resp}
\switchcolumn
\EN
\begin{Resp}\Rsign The kings of Tarshish and of the isles shall bring presents. \Star{} The kings of Arabia and Saba shall offer gifts unto the Lord God.\end{Resp}
\begin{Resp}\Vsign All they from Saba shall come, they shall bring gold and incense.\end{Resp}
\begin{Resp}\Rsign The kings of Arabia and Saba shall offer gifts unto the Lord God.\end{Resp}
\begin{Resp}\Vsign Glory be to the Father, and to the Son, \Star{} and to the Holy Ghost.\end{Resp}
\begin{Resp}\Rsign The kings of Arabia and Saba shall offer gifts unto the Lord God.\end{Resp}
\switchcolumn*

\end{paracol}

\Office{Die 10 Januarii: Ad Romános, cap. 14}{10 January: Epistle to the Romans, chapter 14}{Feria}
\addcontentsline{toc}{subsection}{Die 10 Januarii: Ad Romános, cap. 14\enspace\textit{10 January: Epistle to the Romans, chapter 14}}
\begin{paracol}{2}
\LA
\LessonHead{Lectio I}
\switchcolumn
\EN
\LessonHead{Lesson I}
\switchcolumn*

\LA
\LessonTitle{De Epístola ad Romános}
\switchcolumn
\EN
\LessonTitle{Lesson from the letter of St. Paul the Apostle to the Romans}
\switchcolumn*

\LA
\Cite{Rom 14:1-4}
\switchcolumn
\EN
\Cite{Rom 14:1-4}
\switchcolumn*

\LA
\noindent Infírmum autem in fide assúmite, non in disceptatiónibus cogitatiónum. Alius enim credit se manducáre ómnia: qui autem infírmus est, olus mandúcet. Is qui mandúcat, non manducántem non spernat: et qui non mandúcat, manducántem non júdicet: Deus enim illum assúmpsit. Tu quis es, qui júdicas aliénum servum? Dómino suo stat, aut cadit: stabit autem: potens est enim Deus statúere illum.\par
\switchcolumn
\EN
\noindent Now him that is weak in faith, take unto you: not in disputes about thoughts. For one believeth that he may eat all things: but he that is weak, let him eat herbs. Let not him that eateth, despise him that eateth not: and he that eateth not, let him not judge him that eateth. For God hath taken him to him. Who art thou that judgest another man's servant? To his own lord he standeth or falleth. And he shall stand: for God is able to make him stand.\par
\switchcolumn*

\LA
\begin{Resp}\Rsign Tria sunt múnera pretiósa, quæ obtulérunt Magi Dómino in die ista, et habent in se divína mystéria: \Star{} In auro, ut ostendátur Regis poténtia: in thure, Sacerdótem magnum consídera: et in myrrha, Domínicam sepultúram.\end{Resp}
\begin{Resp}\Vsign Salútis nostræ auctórem Magi veneráti sunt in cunábulis, et de thesáuris suis mýsticas ei múnerum spécies obtulérunt.\end{Resp}
\begin{Resp}\Rsign In auro, ut ostendátur Regis poténtia: in thure, Sacerdótem magnum consídera: et in myrrha, Domínicam sepultúram.\end{Resp}
\switchcolumn
\EN
\begin{Resp}\Rsign There are three precious gifts which the wise men offered unto the Lord on this day, and they speak a mystery of the things of God: \Star{} Gold, to show His kingly power; frankincense, for our Great High Priest; and myrrh, against the Lord's burying.\end{Resp}
\begin{Resp}\Vsign The wise men worshipped the Captain of our Salvation, as He lay in the manger, and when they had opened their treasures, they presented unto Him mystic gifts.\end{Resp}
\begin{Resp}\Rsign Gold, to show His kingly power; frankincense, for our Great High Priest; and myrrh, against the Lord's burying.\end{Resp}
\switchcolumn*

\LA
\LessonHead{Lectio II}
\switchcolumn
\EN
\LessonHead{Lesson II}
\switchcolumn*

\LA
\Cite{Rom 14:5-8}
\switchcolumn
\EN
\Cite{Rom 14:5-8}
\switchcolumn*

\LA
\noindent Nam álius júdicat diem inter diem, álius autem júdicat omnem diem: unusquísque in suo sensu abúndet. Qui sapit diem, Dómino sapit. Et qui mandúcat, Dómino mandúcat: grátias enim agit Deo. Et qui non mandúcat, Dómino non mandúcat, et grátias agit Deo. Nemo enim nostrum sibi vivit, et nemo sibi móritur. Sive enim vívimus, Dómino vívimus: sive mórimur, Dómino mórimur. Sive ergo vívimus, sive mórimur, Dómini sumus.\par
\switchcolumn
\EN
\noindent For one judgeth between day and day: and another judgeth every day: let every man abound in his own sense. He that regardeth the day, regardeth it unto the Lord. And he that eateth, eateth to the Lord: for he giveth thanks to God. And he that eateth not, to the Lord he eateth not, and giveth thanks to God. For none of us liveth to himself; and no man dieth to himself. For whether we live, we live unto the Lord; or whether we die, we die unto the Lord. Therefore, whether we live, or whether we die, we are the Lord's.\par
\switchcolumn*

\LA
\begin{Resp}\Rsign In colúmbæ spécie Spíritus Sanctus visus est, Patérna vox audíta est: \Star{} Hic est Fílius meus diléctus, in quo mihi bene complácui.\end{Resp}
\begin{Resp}\Vsign Cæli apérti sunt super eum, et vox Patris intónuit.\end{Resp}
\begin{Resp}\Rsign Hic est Fílius meus diléctus, in quo mihi bene complácui.\end{Resp}
\switchcolumn
\EN
\begin{Resp}\Rsign The Holy Ghost appeared in a bodily shape like a dove, and the voice of the Father was heard: \Star{} This is My beloved Son, in Whom I am well pleased.\end{Resp}
\begin{Resp}\Vsign The heavens were opened unto him, and, lo, the voice of the Father was heard, like unto thunder, saying:\end{Resp}
\begin{Resp}\Rsign This is My beloved Son, in Whom I am well pleased.\end{Resp}
\switchcolumn*

\LA
\LessonHead{Lectio III}
\switchcolumn
\EN
\LessonHead{Lesson III}
\switchcolumn*

\LA
\Cite{Rom 14:9-13}
\switchcolumn
\EN
\Cite{Rom 14:9-13}
\switchcolumn*

\LA
\noindent In hoc enim Christus mórtuus est, et resurréxit: ut et mortuórum et vivórum dominétur. Tu autem, quid júdicas fratrem tuum? aut tu quare spernis fratrem tuum? Omnes enim stábimus ante tribúnal Christi. Scriptum est enim: Vivo ego, dicit Dóminus, quóniam mihi flectétur omne genu: et omnis lingua confitébitur Deo. Itaque unusquísque nostrum pro se ratiónem reddet Deo. Non ergo ámplius ínvicem judicémus.\par
\switchcolumn
\EN
\noindent For to this end Christ died and rose again; that he might be Lord both of the dead and of the living. But thou, why judgest thou thy brother? or thou, why dost thou despise thy brother? For we shall all stand before the judgment seat of Christ. For it is written: As I live, saith the Lord, every knee shall bow to me, and every tongue shall confess to God. Therefore every one of us shall render account to God for himself. Let us not therefore judge one another any more. But judge this rather, that you put not a stumblingblock or a scandal in your brother's way.\par
\switchcolumn*

\LA
\begin{Resp}\Rsign Reges Tharsis et ínsulæ múnera ófferent: \Star{} Reges Arabum et Saba dona Dómino Deo addúcent.\end{Resp}
\begin{Resp}\Vsign Omnes de Saba vénient, aurum et thus deferéntes.\end{Resp}
\begin{Resp}\Rsign Reges Arabum et Saba dona Dómino Deo addúcent.\end{Resp}
\begin{Resp}\Vsign Glória Patri, et Fílio, \Star{} et Spirítui Sancto.\end{Resp}
\begin{Resp}\Rsign Reges Arabum et Saba dona Dómino Deo addúcent.\end{Resp}
\switchcolumn
\EN
\begin{Resp}\Rsign The kings of Tarshish and of the isles shall bring presents. \Star{} The kings of Arabia and Saba shall offer gifts unto the Lord God.\end{Resp}
\begin{Resp}\Vsign All they from Saba shall come, they shall bring gold and incense.\end{Resp}
\begin{Resp}\Rsign The kings of Arabia and Saba shall offer gifts unto the Lord God.\end{Resp}
\begin{Resp}\Vsign Glory be to the Father, and to the Son, \Star{} and to the Holy Ghost.\end{Resp}
\begin{Resp}\Rsign The kings of Arabia and Saba shall offer gifts unto the Lord God.\end{Resp}
\switchcolumn*

\end{paracol}

\Office{Die 11 Januarii: Ad Romános, cap. 15}{11 January: Epistle to the Romans, chapter 15}{Feria}
\addcontentsline{toc}{subsection}{Die 11 Januarii: Ad Romános, cap. 15\enspace\textit{11 January: Epistle to the Romans, chapter 15}}
\begin{paracol}{2}
\LA
\LessonHead{Lectio I}
\switchcolumn
\EN
\LessonHead{Lesson I}
\switchcolumn*

\LA
\LessonTitle{De Epístola ad Romános}
\switchcolumn
\EN
\LessonTitle{Lesson from the letter of St. Paul the Apostle to the Romans}
\switchcolumn*

\LA
\Cite{Rom 15:1-4}
\switchcolumn
\EN
\Cite{Rom 15:1-4}
\switchcolumn*

\LA
\noindent Debémus autem nos firmióres imbecillitátes infirmórum sustinére, et non nobis placére. Unusquísque vestrum próximo suo pláceat in bonum, ad ædificatiónem. Etenim Christus non sibi plácuit, sed sicut scriptum est: Impropéria improperántium tibi cecidérunt super me. Quæcúmque enim scripta sunt, ad nostram doctrínam scripta sunt: ut per patiéntiam et consolatiónem Scripturárum, spem habeámus.\par
\switchcolumn
\EN
\noindent Now we that are stronger, ought to bear the infirmities of the weak, and not to please ourselves. Let every one of you please his neighbour unto good, to edification. For Christ did not please himself, but as it is written: The reproaches of them that reproached thee, fell upon me. For what things soever were written, were written for our learning: that through patience and the comfort of the scriptures, we might have hope.\par
\switchcolumn*

\LA
\begin{Resp}\Rsign Tria sunt múnera pretiósa, quæ obtulérunt Magi Dómino in die ista, et habent in se divína mystéria: \Star{} In auro, ut ostendátur Regis poténtia: in thure, Sacerdótem magnum consídera: et in myrrha, Domínicam sepultúram.\end{Resp}
\begin{Resp}\Vsign Salútis nostræ auctórem Magi veneráti sunt in cunábulis, et de thesáuris suis mýsticas ei múnerum spécies obtulérunt.\end{Resp}
\begin{Resp}\Rsign In auro, ut ostendátur Regis poténtia: in thure, Sacerdótem magnum consídera: et in myrrha, Domínicam sepultúram.\end{Resp}
\switchcolumn
\EN
\begin{Resp}\Rsign There are three precious gifts which the wise men offered unto the Lord on this day, and they speak a mystery of the things of God: \Star{} Gold, to show His kingly power; frankincense, for our Great High Priest; and myrrh, against the Lord's burying.\end{Resp}
\begin{Resp}\Vsign The wise men worshipped the Captain of our Salvation, as He lay in the manger, and when they had opened their treasures, they presented unto Him mystic gifts.\end{Resp}
\begin{Resp}\Rsign Gold, to show His kingly power; frankincense, for our Great High Priest; and myrrh, against the Lord's burying.\end{Resp}
\switchcolumn*

\LA
\LessonHead{Lectio II}
\switchcolumn
\EN
\LessonHead{Lesson II}
\switchcolumn*

\LA
\Cite{Rom 15:5-11}
\switchcolumn
\EN
\Cite{Rom 15:5-11}
\switchcolumn*

\LA
\noindent Deus autem patiéntiæ et solátii, det vobis idípsum sápere in altérutrum secúndum Jesum Christum: ut unánimes, uno ore honorificétis Deum, et Patrem Dómini nostri Jesu Christi. Propter quod suscípite ínvicem, sicut et Christus suscépit vos in honórem Dei. Dico enim Christum Jesum minístrum fuísse circumcisiónis propter veritátem Dei, ad confirmándas promissiónes patrum: Gentes autem super misericórdiam honoráre Deum, sicut scriptum est: Proptérea confitébor tibi in géntibus, Dómine, et nómini tuo cantábo. Et íterum dicit: Lætámini, gentes, cum plebe ejus. Et íterum: Laudáte, omnes gentes, Dóminum, et magnificáte eum, omnes pópuli.\par
\switchcolumn
\EN
\noindent Now the God of patience and of comfort grant you to be of one mind one towards another, according to Jesus Christ: That with one mind, and with one mouth, you may glorify God and the Father of our Lord Jesus Christ. Wherefore receive one another, as Christ also hath received you unto the honour of God. For I say that Christ Jesus was minister of the circumcision for the truth of God, to confirm the promises made unto the fathers. But that the Gentiles are to glorify God for his mercy, as it is written: Therefore will I confess to thee, O Lord, among the Gentiles, and will sing to thy name. And again he saith: Rejoice, ye Gentiles, with his people. And again: Praise the Lord, all ye Gentiles; and magnify him, all ye people.\par
\switchcolumn*

\LA
\begin{Resp}\Rsign In colúmbæ spécie Spíritus Sanctus visus est, Patérna vox audíta est: \Star{} Hic est Fílius meus diléctus, in quo mihi bene complácui.\end{Resp}
\begin{Resp}\Vsign Cæli apérti sunt super eum, et vox Patris intónuit.\end{Resp}
\begin{Resp}\Rsign Hic est Fílius meus diléctus, in quo mihi bene complácui.\end{Resp}
\switchcolumn
\EN
\begin{Resp}\Rsign The Holy Ghost appeared in a bodily shape like a dove, and the voice of the Father was heard: \Star{} This is My beloved Son, in Whom I am well pleased.\end{Resp}
\begin{Resp}\Vsign The heavens were opened unto him, and, lo, the voice of the Father was heard, like unto thunder, saying:\end{Resp}
\begin{Resp}\Rsign This is My beloved Son, in Whom I am well pleased.\end{Resp}
\switchcolumn*

\LA
\LessonHead{Lectio III}
\switchcolumn
\EN
\LessonHead{Lesson III}
\switchcolumn*

\LA
\Cite{Rom 15:12-16}
\switchcolumn
\EN
\Cite{Rom 15:12-16}
\switchcolumn*

\LA
\noindent Et rursus Isaías ait: Erit radix Jesse, et qui exsúrget régere gentes, in eum gentes sperábunt. Deus autem spei répleat vos omni gáudio et pace in credéndo: ut abundétis in spe, et virtúte Spíritus Sancti. Certus sum autem, fratres mei, et ego ipse de vobis, quóniam et ipsi pleni estis dilectióne, repléti omni sciéntia, ita ut possítis altérutrum monére. Audácius autem scripsi vobis, fratres, ex parte, tamquam in memóriam vos redúcens: propter grátiam, quæ data est mihi a Deo, ut sim miníster Christi Jesu in géntibus: sanctíficans Evangélium Dei, ut fiat oblátio géntium accépta, et sanctificáta in Spíritu Sancto.\par
\switchcolumn
\EN
\noindent And again Isaias saith: There shall be a root of Jesse; and he that shall rise up to rule the Gentiles, in him the Gentiles shall hope. Now the God of hope fill you with all joy and peace in believing; that you may abound in hope, and in the power of the Holy Ghost. And I myself also, my brethren, am assured of you, that you also are full of love, replenished with all knowledge, so that you are able to admonish one another. But I have written to you, brethren, more boldly in some sort, as it were putting you in mind: because of the grace which is given me from God. That I should be the minister of Christ Jesus among the Gentiles; sanctifying the gospel of God, that the oblation of the Gentiles may be made acceptable and sanctified in the Holy Ghost.\par
\switchcolumn*

\LA
\begin{Resp}\Rsign Reges Tharsis et ínsulæ múnera ófferent: \Star{} Reges Arabum et Saba dona Dómino Deo addúcent.\end{Resp}
\begin{Resp}\Vsign Omnes de Saba vénient, aurum et thus deferéntes.\end{Resp}
\begin{Resp}\Rsign Reges Arabum et Saba dona Dómino Deo addúcent.\end{Resp}
\begin{Resp}\Vsign Glória Patri, et Fílio, \Star{} et Spirítui Sancto.\end{Resp}
\begin{Resp}\Rsign Reges Arabum et Saba dona Dómino Deo addúcent.\end{Resp}
\switchcolumn
\EN
\begin{Resp}\Rsign The kings of Tarshish and of the isles shall bring presents. \Star{} The kings of Arabia and Saba shall offer gifts unto the Lord God.\end{Resp}
\begin{Resp}\Vsign All they from Saba shall come, they shall bring gold and incense.\end{Resp}
\begin{Resp}\Rsign The kings of Arabia and Saba shall offer gifts unto the Lord God.\end{Resp}
\begin{Resp}\Vsign Glory be to the Father, and to the Son, \Star{} and to the Holy Ghost.\end{Resp}
\begin{Resp}\Rsign The kings of Arabia and Saba shall offer gifts unto the Lord God.\end{Resp}
\switchcolumn*

\end{paracol}

\SeasonTitle{Tempus post Epiphaniam}{Time after Epiphany}

\addcontentsline{toc}{section}{Tempus post Epiphaniam\enspace\textit{Time after Epiphany}}

\Office{Sanctæ Familiæ Jesu Mariæ Joseph}{The Holy Family of Jesus, Mary and Joseph}{Duplex majus}
\addcontentsline{toc}{subsection}{Sanctæ Familiæ Jesu Mariæ Joseph\enspace\textit{The Holy Family of Jesus, Mary and Joseph}}
\begin{paracol}{2}
\LA
\LessonHead{Lectio I}
\switchcolumn
\EN
\LessonHead{Lesson I}
\switchcolumn*

\LA
\LessonTitle{De Epístola beáti Pauli Apóstoli ad Colossénses}
\switchcolumn
\EN
\LessonTitle{Lesson from the letter of St. Paul the Apostle to the Colossians}
\switchcolumn*

\LA
\Cite{Col 3:12-16}
\switchcolumn
\EN
\Cite{Col 3:12-16}
\switchcolumn*

\LA
\noindent Indúite vos ergo sicut elécti Dei, sancti, et dilécti, víscera misericórdiæ, benignitátem, humilitátem, modéstiam, patiéntiam: supportántes ínvicem, et donántes vobismetípsis si quis advérsus áliquem habet querélam: sicut et Dóminus donávit vobis, ita et vos. Super ómnia autem hæc, caritátem habéte, quod est vínculum perfectiónis: et pax Christi exsúltet in córdibus vestris, in qua et vocáti estis in uno córpore: et grati estóte. Verbum Christi hábitet in vobis abundánter, in omni sapiéntia, docéntes, et commonéntes vosmetípsos, psalmis, hymnis et cánticis spirituálibus, in grátia cantántes in córdibus vestris Deo.\par
\switchcolumn
\EN
\noindent Put ye on therefore, as the elect of God, holy, and beloved, the bowels of mercy, benignity, humility, modesty, patience: Bearing with one another, and forgiving one another, if any have a complaint against another: even as the Lord hath forgiven you, so do you also. But above all these things have charity, which is the bond of perfection: And let the peace of Christ rejoice in your hearts, wherein also you are called in one body: and be ye thankful. Let the word of Christ dwell in you abundantly, in all wisdom: teaching and admonishing one another in psalms, hymns, and spiritual canticles, singing in grace in your hearts to God.\par
\switchcolumn*

\LA
\begin{Resp}\Rsign Deus noster in terris visus est, \Star{} Et cum homínibus conversátus est.\end{Resp}
\begin{Resp}\Vsign Hic adinvénit omnem viam disciplínæ, et trádidit illam Jacob púero suo.\end{Resp}
\begin{Resp}\Rsign Et cum homínibus conversátus est.\end{Resp}
\switchcolumn
\EN
\begin{Resp}\Rsign Our God did shew himself upon earth, \Star{} And held converse with men.\end{Resp}
\begin{Resp}\Vsign He hath found out all the way of knowledge, and hath given it unto Jacob his servant.\end{Resp}
\begin{Resp}\Rsign And held converse with men.\end{Resp}
\switchcolumn*

\LA
\LessonHead{Lectio II}
\switchcolumn
\EN
\LessonHead{Lesson II}
\switchcolumn*

\LA
\Cite{Col 3:17-21}
\switchcolumn
\EN
\Cite{Col 3:17-21}
\switchcolumn*

\LA
\noindent Omne, quodcúmque fácitis in verbo aut in ópere, ómnia in nómine Dómini Jesu Christi, grátias agéntes Deo et Patri per ipsum. Mulíeres, súbditæ estóte viris, sicut opórtet in Dómino. Viri, dilígite uxóres vestras, et nolíte amári esse ad illas. Fílii, obedíte paréntibus per ómnia: hoc enim plácitum est in Dómino. Patres, nolíte ad indignatiónem provocáre fílios vestros, ut non pusíllo ánimo fiant.\par
\switchcolumn
\EN
\noindent All whatsoever you do in word or in work, do all in the name of the Lord Jesus Christ, giving thanks to God and the Father by him. Wives, be subject to your husbands, as it behoveth in the Lord. Husbands, love your wives, and be not bitter towards them. Children, obey your parents in all things: for this is well pleasing to the Lord. Fathers, provoke not your children to indignation, lest they be discouraged.\par
\switchcolumn*

\LA
\begin{Resp}\Rsign Beáti qui hábitant \Star{} In domo tua, Dómine.\end{Resp}
\begin{Resp}\Vsign In sǽcula sæculórum laudábunt te.\end{Resp}
\begin{Resp}\Rsign In domo tua, Dómine.\end{Resp}
\switchcolumn
\EN
\begin{Resp}\Rsign Blessed are they that live \Star{} In thy house, O Lord.\end{Resp}
\begin{Resp}\Vsign They shall praise thee for ever and ever.\end{Resp}
\begin{Resp}\Rsign In thy house, O Lord.\end{Resp}
\switchcolumn*

\LA
\LessonHead{Lectio III}
\switchcolumn
\EN
\LessonHead{Lesson III}
\switchcolumn*

\LA
\Cite{Col 3:22-25; 4:1-2}
\switchcolumn
\EN
\Cite{Col 3:22-25; 4:1-2}
\switchcolumn*

\LA
\noindent Servi, obedíte per ómnia dóminis carnálibus, non ad óculum serviéntes, quasi homínibus placéntes, sed in simplicitáte cordis, timéntes Deum. Quodcúmque fácitis, ex ánimo operámini sicut Dómino, et non homínibus: sciéntes quod a Dómino accipiétis retributiónem hereditátis. Dómino Christo servíte. Qui enim injúriam facit, recípiet id, quod iníque gessit: et non est personárum accéptio apud Deum. Dómini, quod justum est et æquum, servis præstáte: sciéntes quod et vos Dóminum habétis in cælo. Oratióni instáte, vigilántes in ea in gratiárum actióne.\par
\switchcolumn
\EN
\noindent Servants, obey in all things your masters according to the flesh, not serving to the eye, as pleasing men, but in simplicity of heart, fearing God. Whatsoever you do, do it from the heart, as to the Lord, and not to men: Knowing that you shall receive of the Lord the reward of inheritance. Serve ye the Lord Christ. For he that doth wrong, shall receive for that which he hath done wrongfully: and there is no respect of persons with God. Masters, do to your servants that which is just and equal: knowing that you also have a master in heaven. Be instant in prayer; watching in it with thanksgiving:\par
\switchcolumn*

\LA
\begin{Resp}\Rsign Débuit per ómnia frátribus assimilári, \Star{} Ut miséricors fíeret.\end{Resp}
\begin{Resp}\Vsign Cum esset Fílius Dei, dídicit ex iis quæ passus est, obediéntiam.\end{Resp}
\begin{Resp}\Rsign Ut miséricors fíeret.\end{Resp}
\begin{Resp}\Vsign Glória Patri, et Fílio, \Star{} et Spirítui Sancto.\end{Resp}
\begin{Resp}\Rsign Ut miséricors fíeret.\end{Resp}
\switchcolumn
\EN
\begin{Resp}\Rsign In all things it behoved him to be made like unto his brethren, \Star{} That he might be merciful.\end{Resp}
\begin{Resp}\Vsign Though he was the Son of God, yet learned he obedience by the things which he suffered.\end{Resp}
\begin{Resp}\Rsign That he might be merciful.\end{Resp}
\begin{Resp}\Vsign Glory be to the Father, and to the Son, \Star{} and to the Holy Ghost.\end{Resp}
\begin{Resp}\Rsign That he might be merciful.\end{Resp}
\switchcolumn*

\LA
\LessonHead{Lectio IV}
\switchcolumn
\EN
\LessonHead{Lesson IV}
\switchcolumn*

\LA
\LessonTitle{Ex Lítteris Apostólicis Leónis Papæ décimi tértii}
\switchcolumn
\EN
\LessonTitle{From the Apostolic Letters of Pope Leo XIII}
\switchcolumn*

\LA
\Source{Breve « Neminem fugit » 14 junii 1892}
\switchcolumn
\EN
\Source{Breve « Neminem fugit » 14 junii 1892}
\switchcolumn*

\LA
\noindent Miséricors Deus cum humánæ reparatiónis opus, quod diu sǽcula exspectábant, perfícere decrevísset, ita ejúsdem óperis ratiónem ordinémque dispósuit, ut prima ipsa ejúsdem inítia augústam mundo exhibérent spéciem Famíliæ divínitus constitútæ, in qua omnes hómines absolutíssimum domésticæ societátis, omnísque virtútis ac sanctitátis intueréntur exémplar. Talis quidem Família éxstitit Nazarethána illa, in qua, ántequam géntibus univérsis pleno lúmine emicuísset, Sol justítiæ erat abscónditus: nimírum Christus Deus, Servátor noster, cum Vírgine Matre et Joseph, viro sanctíssimo, qui erga Jesum patérno fungebátur múnere. Mínime dúbium est quin ex iis láudibus, quæ in societáte et consuetúdine doméstica ex mútuis caritátis offíciis, ex sanctitáte morum, ex pietátis exercitatióne proficiscúntur, máxima quæque enitúerit in sacra illa Família, quæ síquidem eárum futúra erat céteris documénto. Ac proptérea benígno Providéntiæ consílio sic illa cónstitit, ut sínguli christiáni, qualicúmque conditióne vel loco, si ad eam ánimum advértant, fácile possint cujuscúmque virtútis exercéndæ habére causam et invitaméntum.\par
\switchcolumn
\EN
\noindent When God in his mercy determined to accomplish the work of man's renewal, which same had so many long ages awaited, he appointed and ordained this work on such wise that its very beginning might shew to the world the august spectacle of a Family which was known to be divinely constituted; that therein all men might behold a perfect model, as well of domestic life as of every virtue and pattern of holiness: for such indeed was the Holy Family of Nazareth. There in secret dwelt the Sun of Righteousness, until the time when he should shine out in full splendour in the sight of all nations. There Christ, our God and Saviour, lived with his Virgin Mother, and with that most holy man Joseph, who held to him the place of father. No one can doubt that in this Holy Family was displayed every virtue which can be called forth by an ordinary home life, with its mutual services of charity, its holy intercourse, and its practices of godly piety, since the Holy Family was destined to be a pattern to all others. For that very reason was it established by the merciful designs of Providence, namely, that every Christian, in every walk of life and in every place, might easily, if he would but give heed to it, have before him a motive and a pattern for the good life.\par
\switchcolumn*

\LA
\begin{Resp}\Rsign Ego autem mendícus sum et pauper: \Star{} Dóminus sollícitus est mei.\end{Resp}
\begin{Resp}\Vsign Labóres mánuum tuárum quia manducábis, beátus es, et bene tibi erit.\end{Resp}
\begin{Resp}\Rsign Dóminus sollícitus est mei.\end{Resp}
\switchcolumn
\EN
\begin{Resp}\Rsign As for me, I am poor and needy: \Star{} But the Lord careth for me.\end{Resp}
\begin{Resp}\Vsign Thou shalt eat the labours of thine hands; blessed art thou, and happy shalt thou be.\end{Resp}
\begin{Resp}\Rsign But the Lord careth for me.\end{Resp}
\switchcolumn*

\LA
\LessonHead{Lectio V}
\switchcolumn
\EN
\LessonHead{Lesson V}
\switchcolumn*

\LA
\noindent Habent revéra patresfamílias in Joseph vigilántiæ providentiǽque patérnæ præclaríssimam normam: habent matres in sanctíssima Vírgine Deípara amóris, verecúndiæ, submissiónis ánimi perfectǽque fídei insígne spécimen: fílii vero famílias in Jesu, qui erat súbditus illis, habent divínum obediéntiæ exémplar quod admiréntur, colant, imiténtur. Qui nóbiles nati sunt, discent a Família régii sánguinis quómodo et in édita fortúna se témperent, et in afflícta retíneant dignitátem: qui dites, noscent ab ea quantum sint virtútibus posthabéndæ divítiæ. Operárii autem et ii omnes, qui familiárum rerum angústiis ac tenuióre conditióne tam ácriter irritántur, si ad sanctíssimos illíus domésticæ societátis consórtes respéctent, non déerit eis causa cur loco, qui sibi óbtigit, delecténtur pótius quam dóleant. Commúnes enim cum Sacra Família sunt illis labóres: commúnes curæ de vita cotidiána: débuit et Joseph de mercéde sua vitæ ratiónibus consúlere: imo ipsæ divínæ manus se fabrílibus exercuérunt. Nec mirum sane est si sapientíssimi hómines divítiis affluéntes, eas abícere volúerint, sociámque cum Jesu, María et Joseph sibi elígere paupertátem.\par
\switchcolumn
\EN
\noindent To all fathers of families, Joseph is verily the best model of paternal vigilance and care. In the most holy Virgin Mother of God, mothers may find an excellent example of love, modesty, resignation of spirit, and the perfecting of faith. And in Jesus, who was subject to his parents, the children of the family have a divine pattern of obedience which they can admire, reverence, and imitate. Those who are of noble birth may learn, from this Family of royal blood, how to live simply in times of prosperity, and how to retain their dignity in times of distress. The rich may learn that moral worth is to be more highly esteemed than wealth. Artisans, and all such as are bitterly grieved by the narrow and slender means of their families, if they would but consider the sublime holiness of the members of this domestic fellowship, cannot fail to find some cause for rejoicing in their lot, rather than for being merely dissatisfied with it. In common with the Holy Family, they have to work, and to provide for the daily wants of life. Joseph had to engage in trade, in order to live; even the divine hands laboured at an artisan's calling. It is not to be wondered at, that the wealthiest men, if truly wise, have been willing to cast away their riches, and to embrace a life of poverty with Jesus, Mary, and Joseph.\par
\switchcolumn*

\LA
\begin{Resp}\Rsign Vulpes fóveas habent, et vólucres cæli nidos, \Star{} Fílius autem hóminis non habet ubi caput reclínet.\end{Resp}
\begin{Resp}\Vsign Pauper sum ego, et in labóribus a juventúte mea.\end{Resp}
\begin{Resp}\Rsign Fílius autem hóminis non habet ubi caput reclínet.\end{Resp}
\switchcolumn
\EN
\begin{Resp}\Rsign The foxes have holes, and the birds of the air have nests, \Star{} But the Son of Man hath not where to lay his head.\end{Resp}
\begin{Resp}\Vsign I am poor and needy, and in agony from my youth up.\end{Resp}
\begin{Resp}\Rsign But the Son of Man hath not where to lay his head.\end{Resp}
\switchcolumn*

\LA
\LessonHead{Lectio VI}
\switchcolumn
\EN
\LessonHead{Lesson VI}
\switchcolumn*

\LA
\noindent Quibus e rebus jure meritóque apud cathólicos Sacræ Famíliæ cultus matúre invéctus, majus in dies síngulos increméntum capit. Id quidem probant tum christianórum sodalitátes sub invocatióne Sacræ Famíliæ institútæ, tum singuláres honóres ei rédditi, tum potíssimum a decessóribus Nostris ad excitándum erga eam pietátis stúdium impertíta privilégia et grátiæ. Hujúsmodi cultus magno in honóre hábitus est jam inde a sǽculo décimo séptimo, latéque per Itáliam, Gálliam et Bélgium propagátus, totam fere Európam pervásit: deínde prætergréssus vastos Océani tractus, in América per Canadénsem regiónem, sese exténdit, faustísque efflóruit auspíciis. Síquidem nihil magis salutáre aut éfficax famíliis christiánis cogitári potest exémplo Sacræ Famíliæ, quæ perfectiónem absolutionémque compléctitur ómnium virtútum domesticárum. Sic imploráti inter domésticos paríetes adsint propítii Jesus, María et Joseph, caritátem alant, mores regant, ad virtútem próvocent imitatióne sui, et quæ úndique instant mortáles ærúmnæ, eas leniéndo fáciant tolerabilióres.\par
\switchcolumn
\EN
\noindent From the foregoing it is evident how natural and fitting it was that devotion to the Holy Family should in due time have grown up amongst Catholics; and once begun, that it should spread far and wide. Proof of this lieth first in the sodalities instituted under the ínvocation of the Holy Family; then in the unique honours bestowed upon it; and above all, by the privileges and favours granted to this devotion by our predecessors to stimulate fervour and piety in its regard. This devotion was already held in great esteem in the seventeenth century. Widely propagated in Italy, France, and Belgium, it spread over almost the whole of Europe; thence, crossing the wide ocean, through Canada it made is way in the Americas, and finding favour there, became very flourishing. Indeed, among Christian families, nothing more salutary nor efficacious can be imagined than the example of the Holy Family, where are to be found all domestic virtues in perfection and completeness. When Jesus, Mary, and Joseph are invoked in the home, charity is likely to be maintained in the family through their example and heavenly entreaty; a good influence is thus exerted over conduct; the practice of virtue is thus incited; and thus the hardships which are everywhere wont to harass mankind, are both mitigated and made easier to bear. To increase devotion to the Holy Family, Pope Leo XIII prescribed that Christian families should be consecrated thereto. Benedict XV extended the Mass and Office to the whole Church.\par
\switchcolumn*

\LA
\begin{Resp}\Rsign Cum in forma Dei esset, semetípsum exinanívit, \Star{} Formam servi accípiens.\end{Resp}
\begin{Resp}\Vsign Humiliávit semetípsum, factus obédiens usque ad mortem.\end{Resp}
\begin{Resp}\Rsign Formam servi accípiens.\end{Resp}
\begin{Resp}\Vsign Glória Patri, et Fílio, \Star{} et Spirítui Sancto.\end{Resp}
\begin{Resp}\Rsign Formam servi accípiens.\end{Resp}
\switchcolumn
\EN
\begin{Resp}\Rsign Being in the form of God, he made himself of no reputation, \Star{} And he took upon him the form of a servant.\end{Resp}
\begin{Resp}\Vsign He humbled himself, and became obedient, even unto death.\end{Resp}
\begin{Resp}\Rsign And he took upon him the form of a servant.\end{Resp}
\begin{Resp}\Vsign Glory be to the Father, and to the Son, \Star{} and to the Holy Ghost.\end{Resp}
\begin{Resp}\Rsign And he took upon him the form of a servant.\end{Resp}
\switchcolumn*

\LA
\LessonHead{Lectio VII}
\switchcolumn
\EN
\LessonHead{Lesson VII}
\switchcolumn*

\LA
\LessonTitle{Léctio sancti Evangélii secúndum Lucam}
\switchcolumn
\EN
\LessonTitle{From the Holy Gospel according to Luke}
\switchcolumn*

\LA
\Cite{Luc 2:42-52}
\switchcolumn
\EN
\Cite{Luke 2:42-52}
\switchcolumn*

\LA
\noindent Cum factus esset Jesus annórum duódecim, ascendéntibus illis Jerosólymam secúndum consuetúdinem diéi festi, consummatísque diébus, cum redirent, remánsit puer Jesus in Jerúsalem, et non cognovérunt paréntes ejus. Et réliqua.\par
\switchcolumn
\EN
\noindent And when he was twelve years old, they going up into Jerusalem, according to the custom of the feast, And having fulfilled the days, when they returned, the child Jesus remained in Jerusalem; and his parents knew it not. And so on.\par
\switchcolumn*

\LA
\LessonTitle{Homilía sancti Bernárdi Abbátis}
\switchcolumn
\EN
\LessonTitle{Homily of St. Bernard Abbot}
\switchcolumn*

\LA
\Source{Homilia 1 super Missus est, n. 7-8}
\switchcolumn
\EN
\Source{Homilia super Missus est, n. 7-8}
\switchcolumn*

\LA
\noindent Et erat súbditus illis. Quis, quibus? Deus homínibus; Deus, inquam, cui Angeli súbditi sunt, cui Principátus et Potestátes obédiunt, súbditus erat Maríæ, nec tantum Maríæ, sed étiam Joseph propter Maríam. Miráre ergo utrúmlibet, et élige quid ámplius miréris, sive Fílii benigníssimam dignatiónem, sive Matris excellentíssimam dignitátem. Utrínque stupor, utrínque miráculum. Et quod Deus féminæ obtémperet, humílitas absque exémplo; et quod Deo fémina principétur, sublímitas sine sócio. In láudibus vírginum singuláriter cánitur quod sequúntur Agnum quocúmque íerit. Quibus ergo láudibus júdicas dignam, quæ étiam præit?\par
\switchcolumn
\EN
\noindent And he was subject unto them. Who was subject? And to whom? God to man! God, I repeat, to whom the Angels are subject, whom the Principalities and Powers do obey, was subject to Mary; and not only to Mary, but to Joseph also for Mary's sake. Marvel, therefore, both at God and man, and choose that which giveth greater wonder, whether it be the most loving condescension of the Son, or the exceeding great dignity of his Mother. Both amaze us, both are marvellous. That God should obey a woman is lowliness without parallel, that woman should rule over God, an elevation beyond comparison. In praise of virgins it is sung of them alone, that they follow the Lamb whithersoever he goeth. Of what praise do ye judge that woman to be worthy who is thus placed before the Lamb of God.\par
\switchcolumn*

\LA
\begin{Resp}\Rsign Vere tu es Rex abscónditus, \Star{} Deus Israël Salvátor.\end{Resp}
\begin{Resp}\Vsign Tu doces hóminem sciéntiam.\end{Resp}
\begin{Resp}\Rsign Deus Israël Salvátor.\end{Resp}
\switchcolumn
\EN
\begin{Resp}\Rsign Verily, thou art a King that hidest thyself, \Star{} O God of Israel, the Saviour.\end{Resp}
\begin{Resp}\Vsign Thou teachest man knowledge.\end{Resp}
\begin{Resp}\Rsign O God of Israel, the Saviour.\end{Resp}
\switchcolumn*

\LA
\LessonHead{Lectio VIII}
\switchcolumn
\EN
\LessonHead{Lesson VIII}
\switchcolumn*

\LA
\noindent Disce, homo, obedíre; disce, terra, subdi; disce, pulvis, obtemperáre. De auctóre tuo loquens Evangelísta. Et erat, inquit, súbditus illis; haud dúbium, quin Maríæ et Joseph. Erubésce, supérbe cinis! Deus se humíliat, et tu te exáltas? Deus se homínibus subdit, et tu, dominári géstiens homínibus, tuo te præpónis auctóri? Utinam mihi, aliquándo tale áliquid cogitánti, Deus respondére dignétur quod et suo increpándo respóndit Apóstolo: Vade, inquit, post me, sátana, quia non sapis ea quæ Dei sunt. Quóties enim homínibus præésse desídero, tóties Deum præíre conténdo; et tunc vere non sápio ea quæ Dei sunt. De ipso namque dictum est: Et erat súbditus illis. Si hóminis, o homo, imitári dedignáris exémplum, certe non erit tibi indígnum sequi auctórem tuum. Si non potes fórsitan sequi eum quocúmque íerit, dignáre vel sequi quo tibi condescéndit.\par
\switchcolumn
\EN
\noindent Learn, O man, to obey! Learn, O earth, to be subject! Learn, O dust, to submit! The Evangelist speaking of thy Creator saith: And he was subject unto them. And there is no doubt that this sheweth us that God was subject to Mary and Joseph. Shame on you, ye proud entities of dust and ashes! God abaseth himself, and dost thou, O creature sprung from the earth, exalt thyself? God maketh himself subject to man, and dost thou, who art always so eager to lord it over men, set up thyself to lord it over thy Creator? For as often soever as I desire pre-eminency over men, so often do I strive to excel God. For of him it was said: And he was subject unto them. If thou disdainest, O man, to follow the example of man, at least thou canst follow thy Creator without dishonour. If thou canst not, perchance, follow him whithersoever he goeth, deign at least to follow him in this thing wherein he hath emptied himself, and made himself of no reputation, for the sake of such as thou.\par
\switchcolumn*

\LA
\begin{Resp}\Rsign Sicut per inobediéntiam uníus hóminis peccatóres constitúti sunt multi: \Star{} Ita et per uníus obeditiónem justi constituéntur multi.\end{Resp}
\begin{Resp}\Vsign Venit Názareth et erat súbditus illis.\end{Resp}
\begin{Resp}\Rsign Ita et per uníus obeditiónem justi constituéntur multi.\end{Resp}
\begin{Resp}\Vsign Glória Patri, et Fílio, \Star{} et Spirítui Sancto.\end{Resp}
\begin{Resp}\Rsign Ita et per uníus obeditiónem justi constituéntur multi.\end{Resp}
\switchcolumn
\EN
\begin{Resp}\Rsign As by one man's disobedience many were made sinners, \Star{} So by the obedience of one Man shall many be made righteous.\end{Resp}
\begin{Resp}\Vsign He came to Nazareth, and was subject unto them.\end{Resp}
\begin{Resp}\Rsign So by the obedience of one Man shall many be made righteous.\end{Resp}
\begin{Resp}\Vsign Glory be to the Father, and to the Son, \Star{} and to the Holy Ghost.\end{Resp}
\begin{Resp}\Rsign So by the obedience of one Man shall many be made righteous.\end{Resp}
\switchcolumn*

\LA
\LessonHead{Lectio IX}
\switchcolumn
\EN
\LessonHead{Lesson IX}
\switchcolumn*

\LA
\noindent Si non potes sublímen incédere sémitam virginitátis, séquere vel Deum per tutíssimam viam humilitátis; a cujus rectitúdine si qui étiam de virgínibus deviáverint, ut verum fátear, nec ipsi sequúntur Agnum quocúmque íerit. Séquitur quidem Agnum coinquinátus húmilis, séquitur et virgo supérbus, sed néuter quocúmque íerit: quia nec ille ascéndere potest ad mundítiam Agni, qui sine mácula est; nec is ad ejúsdem mansuetúdinem descéndere dignátur, qua scílicet non coram tondénte, sed coram occidénte se obmútuit. Attamen salubriórem elégit sequéndi partem in humilitáte peccátor, quam in virginitáte supérbus; cum et illíus immundítiam sua húmilis satisfáctio purget, et hujus pudicítiam supérbia ínquinet.\par
\switchcolumn
\EN
\noindent If thou canst not enter upon the lofty paths of virginity, at least follow God by the most safe road of humility. If any turn aside from this straight way, though they be virgins, they do not follow the Lamb, if the truth be told, whithersoever he goeth. The humble man, though stained with sin, followeth the Lamb; the virgin, though proud, also followeth; but neither of these twain followeth whithersoever he goeth. The former cannot attain unto the purity of the Lamb, for he is without spot; the latter deigneth not to descend to his meekness, who was dumb, not before the shearer, but before the murderer. Yet the sinner who followeth in humility hath chosen a more saving way than the proud man who followeth in virginity; for the humble one maketh satisfaction, and is cleansed of his impurity, but the proud one's chastity is stained by his pride.\par
\switchcolumn*

\LA
\TeDeum{Te Deum}
\switchcolumn
\EN
\TeDeum{Te Deum}
\switchcolumn*

\end{paracol}

\Office{Dominica infra Octavam Epiphaniæ}{Sunday within the Octave of the Epiphany}{Semiduplex II classis}
\addcontentsline{toc}{subsection}{Dominica infra Octavam Epiphaniæ\enspace\textit{Sunday within the Octave of the Epiphany}}
\begin{paracol}{2}
\LA
\LessonHead{Lectio I}
\switchcolumn
\EN
\LessonHead{Lesson I}
\switchcolumn*

\LA
\LessonTitle{Incipit Epístola prima beáti Pauli Apóstoli ad Corínthios}
\switchcolumn
\EN
\LessonTitle{Lesson from the first letter of St. Paul the Apostle to the Corinthians}
\switchcolumn*

\LA
\Cite{1 Cor 1:1-3}
\switchcolumn
\EN
\Cite{1 Cor 1:1-3}
\switchcolumn*

\LA
\noindent Paulus vocátus Apóstolus Jesu Christi per voluntátem Dei, et Sósthenes frater, ecclésiæ Dei, quæ est Corínthi, sanctificátis in Christo Jesu, vocátis sanctis, cum ómnibus qui ínvocant nomen Dómini nostri Jesu Christi, in omni loco ipsórum et nostro. Grátia vobis, et pax a Deo Patre nostro, et Dómino Jesu Christo.\par
\switchcolumn
\EN
\noindent Paul, called to be an apostle of Jesus Christ by the will of God, and Sosthenes a brother, To the church of God that is at Corinth, to them that are sanctified in Christ Jesus, called to be saints, with all that invoke the name of our Lord Jesus Christ, in every place of theirs and ours. Grace to you, and peace from God our Father, and from the Lord Jesus Christ.\par
\switchcolumn*

\LA
\begin{Resp}\Rsign Tria sunt múnera pretiósa, quæ obtulérunt Magi Dómino in die ista, et habent in se divína mystéria: \Star{} In auro, ut ostendátur Regis poténtia: in thure, Sacerdótem magnum consídera: et in myrrha, Domínicam sepultúram.\end{Resp}
\begin{Resp}\Vsign Salútis nostræ auctórem Magi veneráti sunt in cunábulis, et de thesáuris suis mýsticas ei múnerum spécies obtulérunt.\end{Resp}
\begin{Resp}\Rsign In auro, ut ostendátur Regis poténtia: in thure, Sacerdótem magnum consídera: et in myrrha, Domínicam sepultúram.\end{Resp}
\switchcolumn
\EN
\begin{Resp}\Rsign There are three precious gifts which the wise men offered unto the Lord on this day, and they speak a mystery of the things of God: \Star{} Gold, to show His kingly power; frankincense, for our Great High Priest; and myrrh, against the Lord's burying.\end{Resp}
\begin{Resp}\Vsign The wise men worshipped the Captain of our Salvation, as He lay in the manger, and when they had opened their treasures, they presented unto Him mystic gifts.\end{Resp}
\begin{Resp}\Rsign Gold, to show His kingly power; frankincense, for our Great High Priest; and myrrh, against the Lord's burying.\end{Resp}
\switchcolumn*

\LA
\LessonHead{Lectio II}
\switchcolumn
\EN
\LessonHead{Lesson II}
\switchcolumn*

\LA
\Cite{1 Cor 1:4-9}
\switchcolumn
\EN
\Cite{1 Cor 1:4-9}
\switchcolumn*

\LA
\noindent Grátias ago Deo meo semper pro vobis in grátia Dei, quæ data est vobis in Christo Jesu: quod in ómnibus dívites facti estis in illo, in omni verbo, et in omni sciéntia. Sicut testimónium Christi confirmátum est in vobis: ita ut nihil vobis desit in ulla grátia, exspectántibus revelatiónem Dómini nostri Jesu Christi, qui et confirmábit vos usque in finem sine crímine, in die advéntus Dómini nostri Jesu Christi. Fidélis Deus: per quem vocáti estis in societátem fílii ejus Jesu Christi Dómini nostri.\par
\switchcolumn
\EN
\noindent I give thanks to my God always for you, for the grace of God that is given you in Christ Jesus, That in all things you are made rich in him, in all utterance, and in all knowledge; As the testimony of Christ was confirmed in you, So that nothing is wanting to you in any grace, waiting for the manifestation of our Lord Jesus Christ. Who also will confirm you unto the end without crime, in the day of the coming of our Lord Jesus Christ. God is faithful: by whom you are called unto the fellowship of his Son Jesus Christ our Lord.\par
\switchcolumn*

\LA
\begin{Resp}\Rsign In colúmbæ spécie Spíritus Sanctus visus est, Patérna vox audíta est: \Star{} Hic est Fílius meus diléctus, in quo mihi bene complácui.\end{Resp}
\begin{Resp}\Vsign Cæli apérti sunt super eum, et vox Patris intónuit.\end{Resp}
\begin{Resp}\Rsign Hic est Fílius meus diléctus, in quo mihi bene complácui.\end{Resp}
\switchcolumn
\EN
\begin{Resp}\Rsign The Holy Ghost appeared in a bodily shape like a dove, and the voice of the Father was heard: \Star{} This is My beloved Son, in Whom I am well pleased.\end{Resp}
\begin{Resp}\Vsign The heavens were opened unto him, and, lo, the voice of the Father was heard, like unto thunder, saying:\end{Resp}
\begin{Resp}\Rsign This is My beloved Son, in Whom I am well pleased.\end{Resp}
\switchcolumn*

\LA
\LessonHead{Lectio III}
\switchcolumn
\EN
\LessonHead{Lesson III}
\switchcolumn*

\LA
\Cite{1 Cor 1:10-13}
\switchcolumn
\EN
\Cite{1 Cor 1:10-13}
\switchcolumn*

\LA
\noindent Obsécro autem vos fratres per nomen Dómini nostri Jesu Christi: ut idípsum dicátis omnes, et non sint in vobis schísmata: sitis autem perfécti in eódem sensu, et in eádem senténtia. Significátum est enim mihi de vobis fratres mei ab iis, qui sunt Chloes, quia contentiónes sunt inter vos. Hoc autem dico, quod unusquísque vestrum dicit: Ego quidem sum Pauli: ego autem Apóllo: ego vero Cephæ: ego autem Christi. Divísus est Christus? numquid Paulus crucifíxus est pro vobis? aut in nómine Pauli baptizáti estis?\par
\switchcolumn
\EN
\noindent Now I beseech you, brethren, by the name of our Lord Jesus Christ, that you all speak the same thing, and that there be no schisms among you; but that you be perfect in the same mind, and in the same judgment. For it hath been signified unto me, my brethren, of you, by them that are of the house of Chloe, that there are contentions among you. Now this I say, that every one of you saith: I indeed am of Paul; and I am of Apollo; and I am of Cephas; and I of Christ. Is Christ divided? Was Paul then crucified for you? or were you baptized in the name of Paul?\par
\switchcolumn*

\LA
\begin{Resp}\Rsign Reges Tharsis et ínsulæ múnera ófferent: \Star{} Reges Arabum et Saba dona Dómino Deo addúcent.\end{Resp}
\begin{Resp}\Vsign Omnes de Saba vénient, aurum et thus deferéntes.\end{Resp}
\begin{Resp}\Rsign Reges Arabum et Saba dona Dómino Deo addúcent.\end{Resp}
\begin{Resp}\Vsign Glória Patri, et Fílio, \Star{} et Spirítui Sancto.\end{Resp}
\begin{Resp}\Rsign Reges Arabum et Saba dona Dómino Deo addúcent.\end{Resp}
\switchcolumn
\EN
\begin{Resp}\Rsign The kings of Tarshish and of the isles shall bring presents. \Star{} The kings of Arabia and Saba shall offer gifts unto the Lord God.\end{Resp}
\begin{Resp}\Vsign All they from Saba shall come, they shall bring gold and incense.\end{Resp}
\begin{Resp}\Rsign The kings of Arabia and Saba shall offer gifts unto the Lord God.\end{Resp}
\begin{Resp}\Vsign Glory be to the Father, and to the Son, \Star{} and to the Holy Ghost.\end{Resp}
\begin{Resp}\Rsign The kings of Arabia and Saba shall offer gifts unto the Lord God.\end{Resp}
\switchcolumn*

\end{paracol}

\Office{Feria II infra Hebdomadam I post Epiphaniam}{Monday of the First Week after Epiphany}{Feria}
\addcontentsline{toc}{subsection}{Feria II infra Hebdomadam I post Epiphaniam\enspace\textit{Monday of the First Week after Epiphany}}
\begin{paracol}{2}
\LA
\LessonHead{Lectio I}
\switchcolumn
\EN
\LessonHead{Lesson I}
\switchcolumn*

\LA
\LessonTitle{De Epístola prima ad Corínthios}
\switchcolumn
\EN
\LessonTitle{Lesson from the first letter of St. Paul the Apostle to the Corinthians}
\switchcolumn*

\LA
\Cite{1 Cor 2:1-5}
\switchcolumn
\EN
\Cite{1 Cor 2:1-5}
\switchcolumn*

\LA
\noindent Et ego, cum veníssem ad vos, fratres, veni non in sublimitáte sermónis, aut sapiéntiæ, annúntians vobis testimónium Christi. Non enim judicávi me scire áliquid inter vos, nisi Jesum Christum, et hunc crucifíxum. Et ego in infirmitáte, et timóre, et tremóre multo fui apud vos: et sermo meus, et prædicátio mea non in persuasibílibus humánæ sapiéntiæ verbis, sed in ostensióne spíritus et virtútis: ut fides vestra non sit in sapiéntia hóminum, sed in virtúte Dei.\par
\switchcolumn
\EN
\noindent And I, brethren, when I came to you, came not in loftiness of speech or of wisdom, declaring unto you the testimony of Christ. For I judged not myself to know anything among you, but Jesus Christ, and him crucified. And I was with you in weakness, and in fear, and in much trembling. And my speech and my preaching was not in the persuasive words of human wisdom, but in shewing of the Spirit and power; That your faith might not stand on the wisdom of men, but on the power of God.\par
\switchcolumn*

\LA
\begin{Resp}\Rsign Dómine, ne in ira tua árguas me, neque in furóre tuo corrípias me: \Star{} Miserére mei, Dómine, quóniam infírmus sum.\end{Resp}
\begin{Resp}\Vsign Timor et tremor venérunt super me, et contexérunt me ténebræ.\end{Resp}
\begin{Resp}\Rsign Miserére mei, Dómine, quóniam infírmus sum.\end{Resp}
\switchcolumn
\EN
\begin{Resp}\Rsign O Lord, rebuke me not in thine anger, neither chasten me in thine hot displeasure. \Star{} Have mercy upon me, O Lord, for I am weak.\end{Resp}
\begin{Resp}\Vsign Fearfulness and trembling are come upon me, and darkness hath overwhelmed me.\end{Resp}
\begin{Resp}\Rsign Have mercy upon me, O Lord, for I am weak.\end{Resp}
\switchcolumn*

\LA
\LessonHead{Lectio II}
\switchcolumn
\EN
\LessonHead{Lesson II}
\switchcolumn*

\LA
\Cite{1 Cor 2:6-9}
\switchcolumn
\EN
\Cite{1 Cor 2:6-9}
\switchcolumn*

\LA
\noindent Sapiéntiam autem lóquimur inter perféctos: sapiéntiam vero non hujus sǽculi, neque príncipum hujus sǽculi, qui destruúntur: sed lóquimur Dei sapiéntiam in mystério, quæ abscóndita est, quam prædestinávit Deus ante sǽcula in glóriam nostram, quam nemo príncipum hujus sǽculi cognóvit: si enim cognovíssent, nunquam Dóminum glóriæ crucifixíssent. Sed sicut scriptum est: Quod óculus non vidit, nec auris audívit, nec in cor hóminis ascéndit, quæ præparávit Deus iis, qui díligunt illum.\par
\switchcolumn
\EN
\noindent Howbeit we speak wisdom among the perfect: yet not the wisdom of this world, neither of the princes of this world that come to nought; But we speak the wisdom of God in a mystery, a wisdom which is hidden, which God ordained before the world, unto our glory Which none of the princes of this world knew; for if they had known it, they would never have crucified the Lord of glory. But, as it is written: That eye hath not seen, nor ear heard, neither hath it entered into the heart of man, what things God hath prepared for them that love him.\par
\switchcolumn*

\LA
\begin{Resp}\Rsign Deus, qui sedes super thronum, et júdicas æquitátem, esto refúgium páuperum in tribulatióne: \Star{} Quia tu solus labórem et dolórem consíderas.\end{Resp}
\begin{Resp}\Vsign Tibi enim derelíctus est pauper, pupíllo tu eris adjútor.\end{Resp}
\begin{Resp}\Rsign Quia tu solus labórem et dolórem consíderas.\end{Resp}
\switchcolumn
\EN
\begin{Resp}\Rsign O God, Which satest in the throne judging right, be Thou a refuge for the poor, a refuge in times of trouble. \Star{} For Thou alone beholdest mischief and spite.\end{Resp}
\begin{Resp}\Vsign The poor leaveth himself unto thee; Thou wilt be the helper of the fatherless.\end{Resp}
\begin{Resp}\Rsign For Thou alone beholdest mischief and spite.\end{Resp}
\switchcolumn*

\LA
\LessonHead{Lectio III}
\switchcolumn
\EN
\LessonHead{Lesson III}
\switchcolumn*

\LA
\Cite{1 Cor 2:10-13}
\switchcolumn
\EN
\Cite{1 Cor 2:10-13}
\switchcolumn*

\LA
\noindent Nobis autem revelávit Deus per Spíritum suum: Spíritus enim ómnia scrutátur, étiam profúnda Dei. Quis enim hóminum scit quæ sunt hóminis, nisi spíritus hóminis, qui in ipso est? ita et quæ Dei sunt, nemo cognóvit, nisi Spíritus Dei. Nos autem non spíritum hujus mundi accépimus, sed Spíritum, qui ex Deo est, ut sciámus, quæ a Deo donáta sunt nobis: quæ et lóquimur non in doctis humánæ sapiéntiæ verbis, sed in doctrína Spíritus, spirituálibus spirituália comparántes.\par
\switchcolumn
\EN
\noindent But to us God hath revealed them, by this Spirit. For the Spirit searcheth all things, yea, the deep things of God. For what man knoweth the things of a man, but the spirit of a man that is in him? So the things also that are of God no man knoweth, but the Spirit of God. Now we have received not the spirit of this world, but the Spirit that is of God; that we may know the things that are given us from God. Which things also we speak, not in the learned words of human wisdom; but in the doctrine of the Spirit, comparing spiritual things with spiritual.\par
\switchcolumn*

\LA
\begin{Resp}\Rsign A dextris est mihi Dóminus, ne commóvear: \Star{} Propter hoc dilatátum est cor meum, et exsultávit lingua mea.\end{Resp}
\begin{Resp}\Vsign Dóminus pars hereditátis meæ, et cálicis mei.\end{Resp}
\begin{Resp}\Rsign Propter hoc dilatátum est cor meum, et exsultávit lingua mea.\end{Resp}
\begin{Resp}\Vsign Glória Patri, et Fílio, \Star{} et Spirítui Sancto.\end{Resp}
\begin{Resp}\Rsign Propter hoc dilatátum est cor meum, et exsultávit lingua mea.\end{Resp}
\switchcolumn
\EN
\begin{Resp}\Rsign The Lord is at my right hand, I shall never be moved. \Star{} Therefore my heart is glad, and my tongue rejoiceth.\end{Resp}
\begin{Resp}\Vsign The Lord is the portion of mine inheritance, and of my cup.\end{Resp}
\begin{Resp}\Rsign Therefore my heart is glad, and my tongue rejoiceth.\end{Resp}
\begin{Resp}\Vsign Glory be to the Father, and to the Son, \Star{} and to the Holy Ghost.\end{Resp}
\begin{Resp}\Rsign Therefore my heart is glad, and my tongue rejoiceth.\end{Resp}
\switchcolumn*

\end{paracol}

\Office{Feria III infra Hebdomadam I post Epiphaniam}{Tuesday of the First Week after Epiphany}{Feria}
\addcontentsline{toc}{subsection}{Feria III infra Hebdomadam I post Epiphaniam\enspace\textit{Tuesday of the First Week after Epiphany}}
\begin{paracol}{2}
\LA
\LessonHead{Lectio I}
\switchcolumn
\EN
\LessonHead{Lesson I}
\switchcolumn*

\LA
\LessonTitle{De Epístola prima ad Corínthios}
\switchcolumn
\EN
\LessonTitle{Lesson from the first letter of St. Paul the Apostle to the Corinthians}
\switchcolumn*

\LA
\Cite{1 Cor 5:1-5}
\switchcolumn
\EN
\Cite{1 Cor 5:1-5}
\switchcolumn*

\LA
\noindent Omníno audítur inter vos fornicátio, et talis fornicátio qualis nec inter gentes, ita ut uxórem patris sui áliquis hábeat. Et vos infláti estis: et non magis luctum habuístis, ut tollátur de médio vestrum qui hoc opus fecit. Ego quidem absens córpore, præsens autem spíritu, jam judicávi ut præsens, eum, qui sic operátus est, in nómine Dómini nostri Jesu Christi, congregátis vobis et meo spíritu, cum virtúte Dómini nostri Jesu, trádere hujúsmodi sátanæ in intéritum carnis, ut spíritus salvus sit in die Dómini nostri Jesu Christi.\par
\switchcolumn
\EN
\noindent It is absolutely heard, that there is fornication among you, and such fornication as the like is not among the heathens; that one should have his father's wife. And you are puffed up; and have not rather mourned, that he might be taken away from among you, that hath done this deed. I indeed, absent in body, but present in spirit, have already judged, as though I were present, him that hath so done, In the name of our Lord Jesus Christ, you being gathered together, and my spirit, with the power of our Lord Jesus; To deliver such a one to Satan for the destruction of the flesh, that the spirit may be saved in the day of our Lord Jesus Christ.\par
\switchcolumn*

\LA
\begin{Resp}\Rsign Tria sunt múnera pretiósa, quæ obtulérunt Magi Dómino in die ista, et habent in se divína mystéria: \Star{} In auro, ut ostendátur Regis poténtia: in thure, Sacerdótem magnum consídera: et in myrrha, Domínicam sepultúram.\end{Resp}
\begin{Resp}\Vsign Salútis nostræ auctórem Magi veneráti sunt in cunábulis, et de thesáuris suis mýsticas ei múnerum spécies obtulérunt.\end{Resp}
\begin{Resp}\Rsign In auro, ut ostendátur Regis poténtia: in thure, Sacerdótem magnum consídera: et in myrrha, Domínicam sepultúram.\end{Resp}
\switchcolumn
\EN
\begin{Resp}\Rsign There are three precious gifts which the wise men offered unto the Lord on this day, and they speak a mystery of the things of God: \Star{} Gold, to show His kingly power; frankincense, for our Great High Priest; and myrrh, against the Lord's burying.\end{Resp}
\begin{Resp}\Vsign The wise men worshipped the Captain of our Salvation, as He lay in the manger, and when they had opened their treasures, they presented unto Him mystic gifts.\end{Resp}
\begin{Resp}\Rsign Gold, to show His kingly power; frankincense, for our Great High Priest; and myrrh, against the Lord's burying.\end{Resp}
\switchcolumn*

\LA
\LessonHead{Lectio II}
\switchcolumn
\EN
\LessonHead{Lesson II}
\switchcolumn*

\LA
\Cite{1 Cor 5:6-8}
\switchcolumn
\EN
\Cite{1 Cor 5:6-8}
\switchcolumn*

\LA
\noindent Non est bona gloriátio vestra. Nescítis, quia módicum ferméntum totam massam corrúmpit? Expurgáte vetus ferméntum, ut sitis nova conspérsio, sicut estis ázymi. Etenim Pascha nostrum immolátus est Christus. Itaque epulémur: non in ferménto véteri, neque in ferménto malítiæ, et nequítiæ: sed in ázymis sinceritátis et veritátis.\par
\switchcolumn
\EN
\noindent Your glorying is not good. Know you not that a little leaven corrupteth the whole lump? Purge out the old leaven, that you may be a new paste, as you are unleavened. For Christ our pasch is sacrificed. Therefore let us feast, not with the old leaven, nor with the leaven of malice and wickedness; but with the unleavened bread of sincerity and truth.\par
\switchcolumn*

\LA
\begin{Resp}\Rsign In colúmbæ spécie Spíritus Sanctus visus est, Patérna vox audíta est: \Star{} Hic est Fílius meus diléctus, in quo mihi bene complácui.\end{Resp}
\begin{Resp}\Vsign Cæli apérti sunt super eum, et vox Patris intónuit.\end{Resp}
\begin{Resp}\Rsign Hic est Fílius meus diléctus, in quo mihi bene complácui.\end{Resp}
\switchcolumn
\EN
\begin{Resp}\Rsign The Holy Ghost appeared in a bodily shape like a dove, and the voice of the Father was heard: \Star{} This is My beloved Son, in Whom I am well pleased.\end{Resp}
\begin{Resp}\Vsign The heavens were opened unto him, and, lo, the voice of the Father was heard, like unto thunder, saying:\end{Resp}
\begin{Resp}\Rsign This is My beloved Son, in Whom I am well pleased.\end{Resp}
\switchcolumn*

\LA
\LessonHead{Lectio III}
\switchcolumn
\EN
\LessonHead{Lesson III}
\switchcolumn*

\LA
\Cite{1 Cor 5:9-11}
\switchcolumn
\EN
\Cite{1 Cor 5:9-11}
\switchcolumn*

\LA
\noindent Scripsi vobis in epístola: Ne commisceámini fornicáriis. Non útique fornicáriis hujus mundi, aut aváris, aut rapácibus, aut idólis serviéntibus: alióquin debuerátis de hoc mundo exiísse. Nunc autem scripsi vobis non commiscéri: si is, qui frater nominátur, est fornicátor, aut avárus, aut idólis sérviens, aut malédicus, aut ebriósus, aut rapax: cum hujúsmodi nec cibum súmere.\par
\switchcolumn
\EN
\noindent I wrote to you in an epistle, not to keep company with fornicators. I mean not with the fornicators of this world, or with the covetous, or the extortioners, or the servers of idols; otherwise you must needs go out of this world. But now I have written to you, not to keep company, if any man that is named a brother, be a fornicator, or covetous, or a server of idols, or a railer, or a drunkard, or an extortioner: with such a one, not so much as to eat.\par
\switchcolumn*

\LA
\begin{Resp}\Rsign Reges Tharsis et ínsulæ múnera ófferent: \Star{} Reges Arabum et Saba dona Dómino Deo addúcent.\end{Resp}
\begin{Resp}\Vsign Omnes de Saba vénient, aurum et thus deferéntes.\end{Resp}
\begin{Resp}\Rsign Reges Arabum et Saba dona Dómino Deo addúcent.\end{Resp}
\begin{Resp}\Vsign Glória Patri, et Fílio, \Star{} et Spirítui Sancto.\end{Resp}
\begin{Resp}\Rsign Reges Arabum et Saba dona Dómino Deo addúcent.\end{Resp}
\switchcolumn
\EN
\begin{Resp}\Rsign The kings of Tarshish and of the isles shall bring presents. \Star{} The kings of Arabia and Saba shall offer gifts unto the Lord God.\end{Resp}
\begin{Resp}\Vsign All they from Saba shall come, they shall bring gold and incense.\end{Resp}
\begin{Resp}\Rsign The kings of Arabia and Saba shall offer gifts unto the Lord God.\end{Resp}
\begin{Resp}\Vsign Glory be to the Father, and to the Son, \Star{} and to the Holy Ghost.\end{Resp}
\begin{Resp}\Rsign The kings of Arabia and Saba shall offer gifts unto the Lord God.\end{Resp}
\switchcolumn*

\end{paracol}

\Office{Feria IV infra Hebdomadam I post Epiphaniam}{Wednesday of the First Week after Epiphany}{Feria}
\addcontentsline{toc}{subsection}{Feria IV infra Hebdomadam I post Epiphaniam\enspace\textit{Wednesday of the First Week after Epiphany}}
\begin{paracol}{2}
\LA
\LessonHead{Lectio I}
\switchcolumn
\EN
\LessonHead{Lesson I}
\switchcolumn*

\LA
\LessonTitle{De Epístola prima ad Corínthios}
\switchcolumn
\EN
\LessonTitle{Lesson from the first letter of St. Paul the Apostle to the Corinthians}
\switchcolumn*

\LA
\Cite{1 Cor 6:1-6}
\switchcolumn
\EN
\Cite{1 Cor 6:1-6}
\switchcolumn*

\LA
\noindent Audet áliquis vestrum, habens negótium advérsus álterum, judicári apud iníquos, et non apud sanctos? An nescítis quóniam sancti de hoc mundo judicábunt? Et si in vobis judicábitur mundus, indígni estis qui de mínimis judicétis? Nescítis quóniam ángelos judicábimus? quanto magis sæculária? Sæculária ígitur judícia, si habuéritis: contemptíbiles, qui sunt in Ecclésia, illos constitúite ad judicándum. Ad verecúndiam vestram dico. Sic non est inter vos sapiens quisquam, qui possit judicáre inter fratrem suum? Sed frater cum fratre judício conténdit: et hoc apud infidéles?\par
\switchcolumn
\EN
\noindent Dare any of you, having a matter against another, go to be judged before the unjust, and not before the saints Know you not that the saints shall judge this world And if the world shall be judged by you, are you unworthy to judge the smallest matters Know you not that we shall judge angels how much more things of this world If therefore you have judgments of things pertaining to this world, set them to judge, who are the most despised in the church. I speak to your shame. Is it so that there is not among you any one wise man, that is able to judge between his brethren But brother goeth to law with brother, and that before unbelievers.\par
\switchcolumn*

\LA
\begin{Resp}\Rsign Tria sunt múnera pretiósa, quæ obtulérunt Magi Dómino in die ista, et habent in se divína mystéria: \Star{} In auro, ut ostendátur Regis poténtia: in thure, Sacerdótem magnum consídera: et in myrrha, Domínicam sepultúram.\end{Resp}
\begin{Resp}\Vsign Salútis nostræ auctórem Magi veneráti sunt in cunábulis, et de thesáuris suis mýsticas ei múnerum spécies obtulérunt.\end{Resp}
\begin{Resp}\Rsign In auro, ut ostendátur Regis poténtia: in thure, Sacerdótem magnum consídera: et in myrrha, Domínicam sepultúram.\end{Resp}
\switchcolumn
\EN
\begin{Resp}\Rsign There are three precious gifts which the wise men offered unto the Lord on this day, and they speak a mystery of the things of God: \Star{} Gold, to show His kingly power; frankincense, for our Great High Priest; and myrrh, against the Lord's burying.\end{Resp}
\begin{Resp}\Vsign The wise men worshipped the Captain of our Salvation, as He lay in the manger, and when they had opened their treasures, they presented unto Him mystic gifts.\end{Resp}
\begin{Resp}\Rsign Gold, to show His kingly power; frankincense, for our Great High Priest; and myrrh, against the Lord's burying.\end{Resp}
\switchcolumn*

\LA
\LessonHead{Lectio II}
\switchcolumn
\EN
\LessonHead{Lesson II}
\switchcolumn*

\LA
\Cite{1 Cor 6:7-11}
\switchcolumn
\EN
\Cite{1 Cor 6:7-11}
\switchcolumn*

\LA
\noindent Jam quidem omníno delíctum est in vobis, quod judícia habétis inter vos. Quare non magis injúriam accípitis? quare non magis fraudem patímini? Sed vos injúriam fácitis, et fraudátis: et hoc frátribus. An nescítis quia iníqui regnum Dei non possidébunt? Nolíte erráre: neque fornicárii, neque idólis serviéntes, neque adúlteri, neque molles, neque masculórum concubitóres, neque fures, neque avári, neque ebriósi, neque malédici, neque rapáces, regnum Dei possidébunt. Et hæc quidam fuístis: sed ablúti estis, sed sanctificáti estis, sed justificáti estis in nómine Dómini nostri Jesu Christi, et in Spíritu Dei nostri.\par
\switchcolumn
\EN
\noindent Already indeed there is plainly a fault among you, that you have lawsuits one with another. Why do you not rather take wrong? Why do you not rather suffer yourselves to be defrauded? But you do wrong and defraud, and that to your brethren. Know you not that the unjust shall not possess the kingdom of God? Do not err: neither fornicators, nor idolaters, nor adulterers, Nor the effeminate, nor liers with mankind, nor thieves, nor covetous, nor drunkards, nor railers, nor extortioners, shall possess the kingdom of God. And such some of you were; but you are washed, but you are sanctified, but you are justified in the name of our Lord Jesus Christ, and the Spirit of our God.\par
\switchcolumn*

\LA
\begin{Resp}\Rsign In colúmbæ spécie Spíritus Sanctus visus est, Patérna vox audíta est: \Star{} Hic est Fílius meus diléctus, in quo mihi bene complácui.\end{Resp}
\begin{Resp}\Vsign Cæli apérti sunt super eum, et vox Patris intónuit.\end{Resp}
\begin{Resp}\Rsign Hic est Fílius meus diléctus, in quo mihi bene complácui.\end{Resp}
\switchcolumn
\EN
\begin{Resp}\Rsign The Holy Ghost appeared in a bodily shape like a dove, and the voice of the Father was heard: \Star{} This is My beloved Son, in Whom I am well pleased.\end{Resp}
\begin{Resp}\Vsign The heavens were opened unto him, and, lo, the voice of the Father was heard, like unto thunder, saying:\end{Resp}
\begin{Resp}\Rsign This is My beloved Son, in Whom I am well pleased.\end{Resp}
\switchcolumn*

\LA
\LessonHead{Lectio III}
\switchcolumn
\EN
\LessonHead{Lesson III}
\switchcolumn*

\LA
\Cite{1 Cor 6:12-18}
\switchcolumn
\EN
\Cite{1 Cor 6:12-18}
\switchcolumn*

\LA
\noindent Omnia mihi licent, sed non ómnia expédiunt. Omnia mihi licent, sed ego sub nullíus rédigar potestáte. Esca ventri, et venter escis: Deus autem et hunc, et has déstruet: corpus autem non fornicatióni, sed Dómino: et Dóminus córpori. Deus vero et Dóminum suscitávit: et nos suscitábit per virtútem suam. Nescítis quóniam córpora vestra membra sunt Christi? Tollens ergo membra Christi, fáciam membra meretrícis? Absit. An nescítis quóniam qui adhǽret meretríci, unum corpus effícitur? Erunt enim (inquit) duo in carne una. Qui autem adhǽret Dómino, unus spíritus est. Fúgite fornicatiónem.\par
\switchcolumn
\EN
\noindent All things are lawful to me, but all things are not expedient. All things are lawful to me, but I will not be brought under the power of any. Meat for the belly, and the belly for the meats; but God shall destroy both it and them: but the body is not for fornication, but for the Lord, and the Lord for the body. Now God hath both raised up the Lord, and will raise us up also by his power. Know you not that your bodies are the members of Christ Shall I then take the members of Christ, and make them the members of an harlot God forbid. Or know you not, that he who is joined to a harlot, is made one body For they shall be, saith he, two in one flesh. But he who is joined to the Lord, is one spirit. Fly fornication.\par
\switchcolumn*

\LA
\begin{Resp}\Rsign Reges Tharsis et ínsulæ múnera ófferent: \Star{} Reges Arabum et Saba dona Dómino Deo addúcent.\end{Resp}
\begin{Resp}\Vsign Omnes de Saba vénient, aurum et thus deferéntes.\end{Resp}
\begin{Resp}\Rsign Reges Arabum et Saba dona Dómino Deo addúcent.\end{Resp}
\begin{Resp}\Vsign Glória Patri, et Fílio, \Star{} et Spirítui Sancto.\end{Resp}
\begin{Resp}\Rsign Reges Arabum et Saba dona Dómino Deo addúcent.\end{Resp}
\switchcolumn
\EN
\begin{Resp}\Rsign The kings of Tarshish and of the isles shall bring presents. \Star{} The kings of Arabia and Saba shall offer gifts unto the Lord God.\end{Resp}
\begin{Resp}\Vsign All they from Saba shall come, they shall bring gold and incense.\end{Resp}
\begin{Resp}\Rsign The kings of Arabia and Saba shall offer gifts unto the Lord God.\end{Resp}
\begin{Resp}\Vsign Glory be to the Father, and to the Son, \Star{} and to the Holy Ghost.\end{Resp}
\begin{Resp}\Rsign The kings of Arabia and Saba shall offer gifts unto the Lord God.\end{Resp}
\switchcolumn*

\end{paracol}

\Office{Feria V infra Hebdomadam I post Epiphaniam}{Thursday of the First Week after Epiphany}{Feria}
\addcontentsline{toc}{subsection}{Feria V infra Hebdomadam I post Epiphaniam\enspace\textit{Thursday of the First Week after Epiphany}}
\begin{paracol}{2}
\LA
\LessonHead{Lectio I}
\switchcolumn
\EN
\LessonHead{Lesson I}
\switchcolumn*

\LA
\LessonTitle{De Epístola prima ad Corínthios}
\switchcolumn
\EN
\LessonTitle{Lesson from the first letter of St. Paul the Apostle to the Corinthians}
\switchcolumn*

\LA
\Cite{1 Cor 7:1-4}
\switchcolumn
\EN
\Cite{1 Cor 7:1-4}
\switchcolumn*

\LA
\noindent De quibus autem scripsístis mihi: Bonum est hómini mulíerem non tángere: propter fornicatiónem autem, unusquísque suam uxórem hábeat, et unaquǽque suum virum hábeat. Uxóri vir débitum reddat: simíliter autem et uxor viro. Múlier sui córporis potestátem non habet, sed vir. Simíliter autem et vir sui córporis potestátem non habet, sed múlier.\par
\switchcolumn
\EN
\noindent Now concerning the thing whereof you wrote to me: It is good for a man not to touch a woman. But for fear of fornication, let every man have his own wife, and let every woman have her own husband. Let the husband render the debt to his wife, and the wife also in like manner to the husband. The wife hath not power of her own body, but the husband. And in like manner the husband also hath not power of his own body, but the wife.\par
\switchcolumn*

\LA
\begin{Resp}\Rsign Deus, in te sperávi, Dómine, non confúndar in ætérnum: in justítia tua líbera me, \Star{} Et éripe me.\end{Resp}
\begin{Resp}\Vsign Inclína ad me aurem tuam, et salva me.\end{Resp}
\begin{Resp}\Rsign Et éripe me.\end{Resp}
\switchcolumn
\EN
\begin{Resp}\Rsign In thee, O God, do I put my trust; let me never be put to confusion, O Lord deliver me in thy righteousness, \Star{} And cause me to escape.\end{Resp}
\begin{Resp}\Vsign Incline thine ear unto me, deliver me speedily.\end{Resp}
\begin{Resp}\Rsign And cause me to escape.\end{Resp}
\switchcolumn*

\LA
\LessonHead{Lectio II}
\switchcolumn
\EN
\LessonHead{Lesson II}
\switchcolumn*

\LA
\Cite{1 Cor 7:5-9}
\switchcolumn
\EN
\Cite{1 Cor 7:5-9}
\switchcolumn*

\LA
\noindent Nolíte fraudáre ínvicem, nisi forte ex consénsu ad tempus, ut vacétis oratióni: et íterum revertímini in idípsum, ne tentet vos sátanas propter incontinéntiam vestram. Hoc autem dico secúndum indulgéntiam, non secúndum impérium. Volo enim omnes vos esse sicut meípsum: sed unusquísque próprium donum habet ex Deo: álius quidem sic, álius vero sic. Dico autem non nuptis, et víduis: bonum est illis si sic permáneant, sicut et ego. Quod si non se cóntinent, nubant. Mélius est enim núbere, quam uri.\par
\switchcolumn
\EN
\noindent Defraud not one another, except, perhaps, by consent, for a time, that you may give yourselves to prayer; and return together again, lest Satan tempt you for your incontinency. But I speak this by indulgence, not by commandment. For I would that all men were even as myself: but every one hath his proper gift from God; one after this manner, and another after that. But I say to the unmarried, and to the widows: It is good for them if they so continue, even as I. But if they do not contain themselves, let them marry. For it is better to marry than to be burnt.\par
\switchcolumn*

\LA
\begin{Resp}\Rsign Repleátur os meum laude tua, ut hymnum dicam glóriæ tuæ, tota die magnitúdinem tuam: noli me proícere in témpore senectútis: \Star{} Dum defécerit in me virtus mea, ne derelínquas me.\end{Resp}
\begin{Resp}\Vsign Gaudébunt lábia mea cum cantávero tibi.\end{Resp}
\begin{Resp}\Rsign Dum defécerit in me virtus mea, ne derelínquas me.\end{Resp}
\switchcolumn
\EN
\begin{Resp}\Rsign Let my mouth be filled with thy praise, that I may sing of thy glory, all the day long of thy greatness. Cast me not off in the time of old age; \Star{} Forsake me not when my strength faileth.\end{Resp}
\begin{Resp}\Vsign My lips shall be fain when I sing unto thee.\end{Resp}
\begin{Resp}\Rsign Forsake me not when my strength faileth.\end{Resp}
\switchcolumn*

\LA
\LessonHead{Lectio III}
\switchcolumn
\EN
\LessonHead{Lesson III}
\switchcolumn*

\LA
\Cite{1 Cor 7:10-14}
\switchcolumn
\EN
\Cite{1 Cor 7:10-14}
\switchcolumn*

\LA
\noindent Iis autem, qui matrimónio juncti sunt, præcípio non ego, sed Dóminus, uxórem a viro non discédere: quod si discésserit, manére innúptam, aut viro suo reconciliári. Et vir uxórem non dimíttat. Nam céteris ego dico, non Dóminus. Si quis frater uxórem habet infidélem et hæc conséntit habitáre cum illo, non dimíttat illam. Et si qua múlier fidélis habet virum infidélem, et hic conséntit habitáre cum illa, non dimíttat virum: sanctificátus est enim vir infidélis per mulíerem fidélem, et sanctificáta est múlier infidélis per virum fidélem: alióquin fílii vestri immúndi essent, nunc autem sancti sunt.\par
\switchcolumn
\EN
\noindent But to them that are married, not I but the Lord commandeth, that the wife depart not from her husband. And if she depart, that she remain unmarried, or be reconciled to her husband. And let not the husband put away his wife. For to the rest I speak, not the Lord. If any brother hath a wife that believeth not, and she consent to dwell with him, let him not put her away. And if any woman hath a husband that believeth not, and he consent to dwell with her, let her not put away her husband. For the unbelieving husband is sanctified by the believing wife; and the unbelieving wife is sanctified by the believing husband: otherwise your children should be unclean; but now they are holy.\par
\switchcolumn*

\LA
\begin{Resp}\Rsign Gaudébunt lábia mea cum cantávero tibi: \Star{} Et ánima mea, quam redemísti, Dómine.\end{Resp}
\begin{Resp}\Vsign Sed et lingua mea meditábitur justítiam tuam, tota die laudem tuam.\end{Resp}
\begin{Resp}\Rsign Et ánima mea, quam redemísti, Dómine.\end{Resp}
\begin{Resp}\Vsign Glória Patri, et Fílio, \Star{} et Spirítui Sancto.\end{Resp}
\begin{Resp}\Rsign Et ánima mea, quam redemísti, Dómine.\end{Resp}
\switchcolumn
\EN
\begin{Resp}\Rsign My lips shall be fain when I sing unto thee; \Star{} And my soul, which Thou, O Lord, hast redeemed.\end{Resp}
\begin{Resp}\Vsign My tongue shall also talk of thy righteousness, all the day long of thy praise.\end{Resp}
\begin{Resp}\Rsign And my soul, which Thou, O Lord, hast redeemed.\end{Resp}
\begin{Resp}\Vsign Glory be to the Father, and to the Son, \Star{} and to the Holy Ghost.\end{Resp}
\begin{Resp}\Rsign And my soul, which Thou, O Lord, hast redeemed.\end{Resp}
\switchcolumn*

\end{paracol}

\Office{Feria VI infra Hebdomadam I post Epiphaniam}{Friday of the First Week after Epiphany}{Feria}
\addcontentsline{toc}{subsection}{Feria VI infra Hebdomadam I post Epiphaniam\enspace\textit{Friday of the First Week after Epiphany}}
\begin{paracol}{2}
\LA
\LessonHead{Lectio I}
\switchcolumn
\EN
\LessonHead{Lesson I}
\switchcolumn*

\LA
\LessonTitle{De Epístola prima ad Corínthios}
\switchcolumn
\EN
\LessonTitle{Lesson from the first letter of St. Paul the Apostle to the Corinthians}
\switchcolumn*

\LA
\Cite{1 Cor 13:1-3}
\switchcolumn
\EN
\Cite{1 Cor 13:1-3}
\switchcolumn*

\LA
\noindent Si linguis hóminum loquar et Angelórum, caritátem autem non hábeam, factus sum velut æs sonans, aut cýmbalum tínniens. Et si habúero prophetíam, et nóverim mystéria ómnia, et omnem sciéntiam: et si habúero omnem fidem, ita ut montes tránsferam, caritátem autem non habúero, nihil sum. Et si distribúero in cibos páuperum omnes facultátes meas, et si tradídero corpus meum, ita ut árdeam, caritátem autem non habúero, nihil mihi prodest.\par
\switchcolumn
\EN
\noindent If I speak with the tongues of men, and of angels, and have not charity, I am become as sounding brass, or a tinkling cymbal. And if I should have prophecy and should know all mysteries, and all knowledge, and if I should have all faith, so that I could remove mountains, and have not charity, I am nothing. And if I should distribute all my goods to feed the poor, and if I should deliver my body to be burned, and have not charity, it profiteth me nothing.\par
\switchcolumn*

\LA
\begin{Resp}\Rsign Confitébor tibi, Dómine Deus, in toto corde meo, et honorificábo nomen tuum in ætérnum: \Star{} Quia misericórdia tua, Dómine, magna est super me.\end{Resp}
\begin{Resp}\Vsign Deus meus es tu, et confitébor tibi: Deus meus es tu, et exaltábo te.\end{Resp}
\begin{Resp}\Rsign Quia misericórdia tua, Dómine, magna est super me.\end{Resp}
\switchcolumn
\EN
\begin{Resp}\Rsign I will praise thee, O Lord my God, with all my heart, and I will glorify thy Name for evermore. \Star{} For great is thy mercy, O Lord, toward me.\end{Resp}
\begin{Resp}\Vsign Thou art my God, and I will praise thee; Thou art my God, and I will exalt thee.\end{Resp}
\begin{Resp}\Rsign For great is thy mercy, O Lord, toward me.\end{Resp}
\switchcolumn*

\LA
\LessonHead{Lectio II}
\switchcolumn
\EN
\LessonHead{Lesson II}
\switchcolumn*

\LA
\Cite{1 Cor 13:4-10}
\switchcolumn
\EN
\Cite{1 Cor 13:4-10}
\switchcolumn*

\LA
\noindent Cáritas pátiens est, benígna est: cáritas non æmulátur, non agit pérperam, non inflátur, non est ambitiósa, non quærit quæ sua sunt, non irritátur, non cógitat malum, non gaudet super iniquitáte, congáudet autem veritáti: ómnia suffert, ómnia credit, ómnia sperat, ómnia sústinet. Cáritas nunquam éxcidit: sive prophetíæ evacuabúntur, sive linguæ cessábunt, sive sciéntia destruétur. Ex parte enim cognóscimus, et ex parte prophetámus. Cum autem vénerit quod perféctum est, evacuábitur quod ex parte est.\par
\switchcolumn
\EN
\noindent Charity is patient, is kind: charity envieth not, dealeth not perversely; is not puffed up; Is not ambitious, seeketh not her own, is not provoked to anger, thinketh no evil; Rejoiceth not in iniquity, but rejoiceth with the truth; Beareth all things, believeth all things, hopeth all things, endureth all things. Charity never falleth away: whether prophecies shall be made void, or tongues shall cease, or knowledge shall be destroyed. For we know in part, and we prophesy in part. But when that which is perfect is come, that which is in part shall be done away.\par
\switchcolumn*

\LA
\begin{Resp}\Rsign Misericórdia tua, Dómine, magna est super me: \Star{} Et liberásti ánimam meam ex inférno inferióri.\end{Resp}
\begin{Resp}\Vsign In die tribulatiónis meæ clamávi ad te, quia exaudísti me.\end{Resp}
\begin{Resp}\Rsign Et liberásti ánimam meam ex inférno inferióri.\end{Resp}
\switchcolumn
\EN
\begin{Resp}\Rsign Great, O Lord, is thy mercy toward me. \Star{} And Thou hast delivered my soul from the lowest hell.\end{Resp}
\begin{Resp}\Vsign In the day of my trouble I called upon thee, for Thou hast heard me.\end{Resp}
\begin{Resp}\Rsign And Thou hast delivered my soul from the lowest hell.\end{Resp}
\switchcolumn*

\LA
\LessonHead{Lectio III}
\switchcolumn
\EN
\LessonHead{Lesson III}
\switchcolumn*

\LA
\Cite{1 Cor 13:11-13}
\switchcolumn
\EN
\Cite{1 Cor 13:11-13}
\switchcolumn*

\LA
\noindent Cum essem párvulus, loquébar ut párvulus, sapiébam ut párvulus, cogitábam ut párvulus. Quando autem factus sum vir, evacuávi quæ erant párvuli. Vidémus nunc per spéculum in ænígmate: tunc autem fácie ad fáciem. Nunc cognósco ex parte: tunc autem cognóscam sicut et cógnitus sum. Nunc autem manent fides, spes, cáritas: tria hæc. Major autem horum est cáritas.\par
\switchcolumn
\EN
\noindent When I was a child, I spoke as a child, I understood as a child, I thought as a child. But, when I became a man, I put away the things of a child. We see now through a glass in a dark manner; but then face to face. Now I know in part; but then I shall know even as I am known. And now there remain faith, hope, and charity, these three: but the greatest of these is charity.\par
\switchcolumn*

\LA
\begin{Resp}\Rsign Factus est mihi Dóminus in refúgium: \Star{} Et Deus meus in auxílium spei meæ.\end{Resp}
\begin{Resp}\Vsign Erípuit me de inimícis meis fortíssimis, et factus est Dóminus protéctor meus.\end{Resp}
\begin{Resp}\Rsign Et Deus meus in auxílium spei meæ.\end{Resp}
\begin{Resp}\Vsign Glória Patri, et Fílio, \Star{} et Spirítui Sancto.\end{Resp}
\begin{Resp}\Rsign Et Deus meus in auxílium spei meæ.\end{Resp}
\switchcolumn
\EN
\begin{Resp}\Rsign The Lord is my refuge. \Star{} And my God is the stay of my trust.\end{Resp}
\begin{Resp}\Vsign He delivered me from the strongest of mine enemies, and the Lord was my stay.\end{Resp}
\begin{Resp}\Rsign And my God is the stay of my trust.\end{Resp}
\begin{Resp}\Vsign Glory be to the Father, and to the Son, \Star{} and to the Holy Ghost.\end{Resp}
\begin{Resp}\Rsign And my God is the stay of my trust.\end{Resp}
\switchcolumn*

\end{paracol}

\Office{Sabbato infra Hebdomadam I post Epiphaniam}{Saturday of the First Week after Epiphany}{Feria}
\addcontentsline{toc}{subsection}{Sabbato infra Hebdomadam I post Epiphaniam\enspace\textit{Saturday of the First Week after Epiphany}}
\begin{paracol}{2}
\LA
\LessonHead{Lectio I}
\switchcolumn
\EN
\LessonHead{Lesson I}
\switchcolumn*

\LA
\LessonTitle{De Epístola prima ad Corínthios}
\switchcolumn
\EN
\LessonTitle{Lesson from the first letter of St. Paul the Apostle to the Corinthians}
\switchcolumn*

\LA
\Cite{1 Cor 16:1-4}
\switchcolumn
\EN
\Cite{1 Cor 16:1-4}
\switchcolumn*

\LA
\noindent De colléctis autem, quæ fiunt in sanctos, sicut ordinávi ecclésiis Galátiæ, ita et vos fácite. Per unam sábbati, unusquísque vestrum apud se sepónat, recóndens quod ei bene placúerit: ut non, cum vénero, tunc colléctæ fiant. Cum autem præsens fúero: quos probavéritis per epístolas, hos mittam perférre grátiam vestram in Jerúsalem. Quod si dignum fúerit, ut et ego eam, mecum ibunt.\par
\switchcolumn
\EN
\noindent Now concerning the collections that are made for the saints, as I have given order to the churches of Galatia, so do ye also. On the first day of the week let every one of you put apart with himself, laying up what it shall well please him; that when I come, the collections be not then to be made. And when I shall be with you, whomsoever you shall approve by letters, them will I send to carry your grace to Jerusalem. And if it be meet that I also go, they shall go with me.\par
\switchcolumn*

\LA
\begin{Resp}\Rsign Ne perdíderis me cum iniquitátibus meis: \Star{} Neque in finem irátus resérves mala mea.\end{Resp}
\begin{Resp}\Vsign Non intres in judícium cum servo tuo, Dómine.\end{Resp}
\begin{Resp}\Rsign Neque in finem irátus resérves mala mea.\end{Resp}
\switchcolumn
\EN
\begin{Resp}\Rsign Cut me not off in the midst of my sins. \Star{} Nor keep thy wrath against me for my latter end.\end{Resp}
\begin{Resp}\Vsign Enter not into judgment with thy servant, O Lord.\end{Resp}
\begin{Resp}\Rsign Nor keep thy wrath against me for my latter end.\end{Resp}
\switchcolumn*

\LA
\LessonHead{Lectio II}
\switchcolumn
\EN
\LessonHead{Lesson II}
\switchcolumn*

\LA
\Cite{1 Cor 16:5-9}
\switchcolumn
\EN
\Cite{1 Cor 16:5-9}
\switchcolumn*

\LA
\noindent Véniam autem ad vos, cum Macedóniam pertransíero: nam Macedóniam pertransíbo. Apud vos autem fórsitan manébo, vel étiam hiemábo: ut vos me deducátis, quocúmque íero. Nolo enim vos modo in tránsitu vidére: spero enim me aliquántulum témporis manére apud vos, si Dóminus permíserit. Permanébo autem Ephesi usque ad Pentecósten. Ostium enim mihi apértum est magnum, et évidens: et adversárii multi.\par
\switchcolumn
\EN
\noindent Now I will come to you, when I shall have passed through Macedonia. For I shall pass through Macedonia. And with you perhaps I shall abide, or even spend the winter: that you may bring me on my way whithersoever I shall go. For I will not see you now by the way, for I trust that I shall abide with you some time, if the Lord permit. But I will tarry at Ephesus until Pentecost. For a great door and evident is opened unto me: and many adversaries.\par
\switchcolumn*

\LA
\begin{Resp}\Rsign Parátum cor meum, Deus, parátum cor meum: \Star{} Cantábo et psalmum dicam Dómino.\end{Resp}
\begin{Resp}\Vsign Exsúrge, glória mea, exsúrge, psaltérium et cíthara, exsúrgam dilúculo.\end{Resp}
\begin{Resp}\Rsign Cantábo et psalmum dicam Dómino.\end{Resp}
\switchcolumn
\EN
\begin{Resp}\Rsign My heart is ready, O God, my heart is ready. \Star{} I will sing and give praise to the Lord.\end{Resp}
\begin{Resp}\Vsign Awake up, my glory, awake, psaltery and harp! I will awake early.\end{Resp}
\begin{Resp}\Rsign I will sing and give praise to the Lord.\end{Resp}
\switchcolumn*

\LA
\LessonHead{Lectio III}
\switchcolumn
\EN
\LessonHead{Lesson III}
\switchcolumn*

\LA
\Cite{1 Cor 16:10-14}
\switchcolumn
\EN
\Cite{1 Cor 16:10-14}
\switchcolumn*

\LA
\noindent Si autem vénerit Timótheus, vidéte ut sine timóre sit apud vos: opus enim Dómini operátur, sicut et ego. Ne quis ergo illum spernat: dedúcite autem illum in pace, ut véniat ad me: exspécto enim illum cum frátribus. De Apóllo autem fratre vobis notum fácio, quóniam multum rogávi eum, ut veníret ad vos cum frátribus: et útique non fuit volúntas, ut nunc veníret: véniet autem, cum ei vácuum fúerit. Vigiláte, state in fide, viríliter ágite, et confortámini. Omnia vestra in caritáte fiant.\par
\switchcolumn
\EN
\noindent Now if Timothy come, see that he be with you without fear, for he worketh the work of the Lord, as I also do. Let no man therefore despise him, but conduct ye him on his way in peace: that he may come to me. For I look for him with the brethren. And touching our brother Apollo, I give you to understand, that I much entreated him to come unto you with the brethren: and indeed it was not his will at all to come at this time. But he will come when he shall have leisure. Watch ye, stand fast in the faith, do manfully, and be strengthened. Let all your things be done in charity.\par
\switchcolumn*

\LA
\begin{Resp}\Rsign Adjútor meus, tibi psallam, quia, Deus, suscéptor meus es: \Star{} Deus meus, misericórdia mea.\end{Resp}
\begin{Resp}\Vsign Lætábor, et exsultábo in te, psallam nómini tuo, Altíssime.\end{Resp}
\begin{Resp}\Rsign Deus meus, misericórdia mea.\end{Resp}
\begin{Resp}\Vsign Glória Patri, et Fílio, \Star{} et Spirítui Sancto.\end{Resp}
\begin{Resp}\Rsign Deus meus, misericórdia mea.\end{Resp}
\switchcolumn
\EN
\begin{Resp}\Rsign Unto thee, O my Strength, will I sing, for God is my defence, \Star{} The God of my mercy.\end{Resp}
\begin{Resp}\Vsign I will be glad and rejoice in thee, I will sing praise to thy Name, O Thou Most High.\end{Resp}
\begin{Resp}\Rsign The God of my mercy.\end{Resp}
\begin{Resp}\Vsign Glory be to the Father, and to the Son, \Star{} and to the Holy Ghost.\end{Resp}
\begin{Resp}\Rsign The God of my mercy.\end{Resp}
\switchcolumn*

\end{paracol}

\Office{Dominica II post Epiphaniam}{Second Sunday after Epiphany}{Semiduplex}
\addcontentsline{toc}{subsection}{Dominica II post Epiphaniam\enspace\textit{Second Sunday after Epiphany}}
\begin{paracol}{2}
\LA
\LessonHead{Lectio I}
\switchcolumn
\EN
\LessonHead{Lesson I}
\switchcolumn*

\LA
\LessonTitle{Incipit Epístola secúnda beáti Pauli Apóstoli ad Corínthios}
\switchcolumn
\EN
\LessonTitle{Lesson from the second letter of St. Paul the Apostle to the Corinthians}
\switchcolumn*

\LA
\Cite{2 Cor 1:1-5}
\switchcolumn
\EN
\Cite{2 Cor 1:1-5}
\switchcolumn*

\LA
\noindent Paulus, Apóstolus Jesu Christi per voluntátem Dei, et Timótheus frater, Ecclésiæ Dei, quæ est Corínthi cum ómnibus sanctis, qui sunt in univérsa Achája. Grátia vobis, et pax a Deo Patre nostro, et Dómino Jesu Christo. Benedíctus Deus et Pater Dómini nostri Jesu Christi, Pater misericordiárum, et Deus totíus consolatiónis, qui consolátur nos in omni tribulatióne nostra: ut possímus et ipsi consolári eos qui in omni pressúra sunt, per exhortatiónem, qua exhortámur et ipsi a Deo. Quóniam sicut abúndant passiónes Christi in nobis: ita et per Christum abúndat consolátio nostra.\par
\switchcolumn
\EN
\noindent Paul, an apostle of Jesus Christ by the will of God, and Timothy our brother: to the church of God that is at Corinth, with all the saints that are in all Achaia: Grace unto you and peace from God our Father, and from the Lord Jesus Christ. Blessed be the God and Father of our Lord Jesus Christ, the Father of mercies, and the God of all comfort. Who comforteth us in all our tribulation; that we also may be able to comfort them who are in all distress, by the exhortation wherewith we also are exhorted by God. For as the sufferings of Christ abound in us: so also by Christ doth our comfort abound.\par
\switchcolumn*

\LA
\begin{Resp}\Rsign Dómine, ne in ira tua árguas me, neque in furóre tuo corrípias me: \Star{} Miserére mei, Dómine, quóniam infírmus sum.\end{Resp}
\begin{Resp}\Vsign Timor et tremor venérunt super me, et contexérunt me ténebræ.\end{Resp}
\begin{Resp}\Rsign Miserére mei, Dómine, quóniam infírmus sum.\end{Resp}
\switchcolumn
\EN
\begin{Resp}\Rsign O Lord, rebuke me not in thine anger, neither chasten me in thine hot displeasure. \Star{} Have mercy upon me, O Lord, for I am weak.\end{Resp}
\begin{Resp}\Vsign Fearfulness and trembling are come upon me, and darkness hath overwhelmed me.\end{Resp}
\begin{Resp}\Rsign Have mercy upon me, O Lord, for I am weak.\end{Resp}
\switchcolumn*

\LA
\LessonHead{Lectio II}
\switchcolumn
\EN
\LessonHead{Lesson II}
\switchcolumn*

\LA
\Cite{2 Cor 1:6-7}
\switchcolumn
\EN
\Cite{2 Cor 1:6-7}
\switchcolumn*

\LA
\noindent Sive autem tribulámur pro vestra exhortatióne et salúte, sive consolámur pro vestra consolatióne, sive exhortámur pro vestra exhortatióne et salúte, quæ operátur tolerántiam earúndem passiónum, quas et nos pátimur: ut spes nostra firma sit pro vobis: sciéntes quod sicut sócii passiónum estis, sic éritis et consolatiónis.\par
\switchcolumn
\EN
\noindent Now whether we be in tribulation, it is for your exhortation and salvation: or whether we be comforted, it is for your consolation: or whether we be exhorted, it is for your exhortation and salvation, which worketh the enduring of the same sufferings which we also suffer. That our hope for you may be steadfast: knowing that as you are partakers of the sufferings, so shall you be also of the consolation.\par
\switchcolumn*

\LA
\begin{Resp}\Rsign Deus, qui sedes super thronum, et júdicas æquitátem, esto refúgium páuperum in tribulatióne: \Star{} Quia tu solus labórem et dolórem consíderas.\end{Resp}
\begin{Resp}\Vsign Tibi enim derelíctus est pauper, pupíllo tu eris adjútor.\end{Resp}
\begin{Resp}\Rsign Quia tu solus labórem et dolórem consíderas.\end{Resp}
\switchcolumn
\EN
\begin{Resp}\Rsign O God, Which satest in the throne judging right, be Thou a refuge for the poor, a refuge in times of trouble. \Star{} For Thou alone beholdest mischief and spite.\end{Resp}
\begin{Resp}\Vsign The poor leaveth himself unto thee; Thou wilt be the helper of the fatherless.\end{Resp}
\begin{Resp}\Rsign For Thou alone beholdest mischief and spite.\end{Resp}
\switchcolumn*

\LA
\LessonHead{Lectio III}
\switchcolumn
\EN
\LessonHead{Lesson III}
\switchcolumn*

\LA
\Cite{2 Cor 1:8-11}
\switchcolumn
\EN
\Cite{2 Cor 1:8-11}
\switchcolumn*

\LA
\noindent Non enim vólumus ignoráre vos, fratres, de tribulatióne nostra, quæ facta est in Asia, quóniam supra modum graváti sumus supra virtútem, ita ut tædéret nos étiam vívere. Sed ipsi in nobismetípsis respónsum mortis habúimus, ut non simus fidéntes in nobis, sed in Deo, qui súscitat mórtuos: qui de tantis perículis nos erípuit, et éruit: in quem sperámus quóniam et adhuc erípiet, adjuvántibus et vobis in oratióne pro nobis: ut ex multórum persónis, ejus quæ in nobis est donatiónis, per multos grátiæ agántur pro nobis.\par
\switchcolumn
\EN
\noindent For we would not have you ignorant, brethren, of our tribulation, which came to us in Asia, that we were pressed out of measure above our strength, so that we were weary even of life. But we had in ourselves the answer of death, that we should not trust in ourselves, but in God who raiseth the dead. Who hath delivered and doth deliver us out of so great dangers: in whom we trust that he will yet also deliver us. You helping withal in prayer for us: that for this gift obtained for us, by the means of many persons, thanks may be given by many in our behalf.\par
\switchcolumn*

\LA
\begin{Resp}\Rsign A dextris est mihi Dóminus, ne commóvear: \Star{} Propter hoc dilatátum est cor meum, et exsultávit lingua mea.\end{Resp}
\begin{Resp}\Vsign Dóminus pars hereditátis meæ, et cálicis mei.\end{Resp}
\begin{Resp}\Rsign Propter hoc dilatátum est cor meum, et exsultávit lingua mea.\end{Resp}
\begin{Resp}\Vsign Glória Patri, et Fílio, \Star{} et Spirítui Sancto.\end{Resp}
\begin{Resp}\Rsign Propter hoc dilatátum est cor meum, et exsultávit lingua mea.\end{Resp}
\switchcolumn
\EN
\begin{Resp}\Rsign The Lord is at my right hand, I shall never be moved. \Star{} Therefore my heart is glad, and my tongue rejoiceth.\end{Resp}
\begin{Resp}\Vsign The Lord is the portion of mine inheritance, and of my cup.\end{Resp}
\begin{Resp}\Rsign Therefore my heart is glad, and my tongue rejoiceth.\end{Resp}
\begin{Resp}\Vsign Glory be to the Father, and to the Son, \Star{} and to the Holy Ghost.\end{Resp}
\begin{Resp}\Rsign Therefore my heart is glad, and my tongue rejoiceth.\end{Resp}
\switchcolumn*

\LA
\LessonHead{Lectio IV}
\switchcolumn
\EN
\LessonHead{Lesson IV}
\switchcolumn*

\LA
\LessonTitle{Sermo sancti Joánnis Chrysóstomi}
\switchcolumn
\EN
\LessonTitle{Sermon of St. John Chrysostom}
\switchcolumn*

\LA
\Source{Præfatio in Epistolas Beati Pauli}
\switchcolumn
\EN
\Source{Preface for the letters of St. Paul}
\switchcolumn*

\LA
\noindent Beáti Pauli Epistolárum lectiónem dum assídue auscúlto, perque hebdómadas síngulas bis sæpe, et ter et quater, quotiescúmque sanctórum Mártyrum memórias celebrámus, gáudio exsúlto, tuba illa spiritáli pérfruens, et éxcitor, ac desidério incalésco, vocem mihi amícam agnóscens, et fere præséntem ipsum intuéri, et disseréntem audíre vídeor. Sed tamen dóleo et moléste fero, quod virum hunc non omnes, sicut par est, cognóscunt: verum ita illum nonnúlli ignórant, ut ne Epistolárum quidem ejus númerum plane sciant. Hoc vero non imperítia facit: sed quod nolint beáti hujus viri scripta in mánibus habére.\par
\switchcolumn
\EN
\noindent As I listen intently to the reading of St. Paul's Epistles, often two or three times a week, whenever we commemorate the holy martyrs, I am filled with joy, delighting in the sound of that spiritual trumpet. And as I recognize the voice of a friend, I am roused, and enkindled with love so that I almost seem to see him present, and to hear him speaking. But nevertheless I am grieved, and am troubled, that all do not know this great man as he deserves to be known. Indeed, many are so ignorant that they do not even know how many epistles he wrote. But this ignorance is not due to a want of intelligence on their part, but because they will not carefully study the writings of this great man.\par
\switchcolumn*

\LA
\begin{Resp}\Rsign Notas mihi fecísti, Dómine, vias vitæ: \Star{} Adimplébis me lætítia cum vultu tuo: delectatiónes in déxtera tua usque in finem.\end{Resp}
\begin{Resp}\Vsign Tu es qui restítues hereditátem meam mihi.\end{Resp}
\begin{Resp}\Rsign Adimplébis me lætítia cum vultu tuo: delectatiónes in déxtera tua usque in finem.\end{Resp}
\switchcolumn
\EN
\begin{Resp}\Rsign O Lord, Thou hast shown me the path of life. \Star{} Thou shalt fill me with joy in thy presence, at thy right hand there are pleasures for evermore.\end{Resp}
\begin{Resp}\Vsign Thou art He That shalt restore mine inheritance unto me.\end{Resp}
\begin{Resp}\Rsign Thou shalt fill me with joy in thy presence, at thy right hand there are pleasures for evermore.\end{Resp}
\switchcolumn*

\LA
\LessonHead{Lectio V}
\switchcolumn
\EN
\LessonHead{Lesson V}
\switchcolumn*

\LA
\noindent Neque enim nos, quæ scimus, si quid scimus, ab ingénii bonitáte atque acúmine scimus: sed quod erga hunc virum impénse affécti, ab illíus lectióne numquam discédimus: síquidem qui amant, ii plus quam céteri omnes eórum facta norunt, quos amant, ut qui de iis ipsis sint sollíciti. Id quod beátus hic véluti osténdens, ad Philippénses ait: Sicut est mihi justum, ut hoc de vobis ómnibus séntiam, eo quod hábeam vos in corde, et in vínculis meis, et in defensióne et confirmatióne Evangélii.\par
\switchcolumn
\EN
\noindent For what we know, if we know anything, we do not know it owing to any superlative talent or penetration, but, being strongly drawn towards this great man, we never cease from reading his works. For so it is that those who love any one usually know better than others what he has done, because they take the trouble to learn all about him. The blessed Paul himself shows that this is so, when he says to the Philippians: As it is meet for me to think this for you all: for that I have you in my heart; and in my bands, and in the defence and confirmation of the gospel.\par
\switchcolumn*

\LA
\begin{Resp}\Rsign Díligam te, Dómine, virtus mea: Dóminus firmaméntum meum, \Star{} Et refúgium meum.\end{Resp}
\begin{Resp}\Vsign Liberátor meus, Deus meus, adjútor meus.\end{Resp}
\begin{Resp}\Rsign Et refúgium meum.\end{Resp}
\switchcolumn
\EN
\begin{Resp}\Rsign I will love thee, O Lord, my strength; the Lord is my rock \Star{} And my fortress.\end{Resp}
\begin{Resp}\Vsign My Deliverer, my God, mine Helper.\end{Resp}
\begin{Resp}\Rsign And my fortress.\end{Resp}
\switchcolumn*

\LA
\LessonHead{Lectio VI}
\switchcolumn
\EN
\LessonHead{Lesson VI}
\switchcolumn*

\LA
\noindent Quaprópter si et vos quoque lectióni diligénter atténdere voluéritis, nihil áliud vobis erit requiréndum. Verax est enim Christi sermo dicéntis: Quǽrite, et inveniétis: pulsáte, et aperiétur vobis. Céterum, quandóquidem complúres ex iis, qui huc nobíscum convéniunt, et liberórum educatiónem, et uxóris curam, et famíliæ providéntiam suscepére, ob idque totos sese huic labóri dare non sústinent: at certe ipsi vos excitáte ad ea saltem capiénda, quæ álii collégerint; stúdii tantúmdem iis, quæ dicta fúerint, auscultándis, quantum pecúniis colligéndis impertiéntes. Nam etsi turpe sit, non nisi tantum stúdii a vobis exígere: optábile tamen erit, si tantum saltem tribuátis.\par
\switchcolumn
\EN
\noindent And if you also will diligently attend to the reading, you will have no need of other instruction. Most true are those words of Christ: Seek and you shall find: knock and it shall be opened unto you. For the rest, since many of those who are assembled here are charged with the care of a wife, and with providing for a family, and with the bringing-up of children, and therefore cannot devote themselves wholly to this study; let them at least bestir themselves to receive what others have gathered; showing as much eagerness in listening to what is said about him as in acquiring wealth. For though it is unseemly to demand from you no more than this, yet it is to be wished that you do this at least.\par
\switchcolumn*

\LA
\begin{Resp}\Rsign Dómini est terra, et plenitúdo ejus: \Star{} Orbis terrárum, et univérsi qui hábitant in eo.\end{Resp}
\begin{Resp}\Vsign Ipse super mária fundávit eam, et super flúmina præparávit illam.\end{Resp}
\begin{Resp}\Rsign Orbis terrárum, et univérsi qui hábitant in eo.\end{Resp}
\begin{Resp}\Vsign Glória Patri, et Fílio, \Star{} et Spirítui Sancto.\end{Resp}
\begin{Resp}\Rsign Orbis terrárum, et univérsi qui hábitant in eo.\end{Resp}
\switchcolumn
\EN
\begin{Resp}\Rsign The earth is the Lord's, and the fulness thereof \Star{} The world, and they that dwell therein.\end{Resp}
\begin{Resp}\Vsign For He hath founded it upon the seas, and established it upon the floods.\end{Resp}
\begin{Resp}\Rsign The world, and they that dwell therein.\end{Resp}
\begin{Resp}\Vsign Glory be to the Father, and to the Son, \Star{} and to the Holy Ghost.\end{Resp}
\begin{Resp}\Rsign The world, and they that dwell therein.\end{Resp}
\switchcolumn*

\LA
\LessonHead{Lectio VII}
\switchcolumn
\EN
\LessonHead{Lesson VII}
\switchcolumn*

\LA
\LessonTitle{Léctio sancti Evangélii secúndum Joánnem}
\switchcolumn
\EN
\LessonTitle{From the Holy Gospel according to John}
\switchcolumn*

\LA
\Cite{Joannes 2:1-11}
\switchcolumn
\EN
\Cite{John 2:1-11}
\switchcolumn*

\LA
\noindent In illo témpore: Núptiæ factæ sunt in Cana Galilǽæ, et erat mater Jesu ibi. Vocátus est autem et Jesus, et discípuli ejus ad núptias. Et réliqua.\par
\switchcolumn
\EN
\noindent At that time, there was a marriage in Cana of Galilee: and the mother of Jesus was there. And Jesus also was invited, and his disciples, to the marriage. And so on.\par
\switchcolumn*

\LA
\LessonTitle{Homilía sancti Augustíni Epíscopi}
\switchcolumn
\EN
\LessonTitle{Homily by St. Augustine, Bishop (of Hippo.)}
\switchcolumn*

\LA
\Source{Tract. 9 in Joannem, post initium}
\switchcolumn
\EN
\Source{Tract. 9 in John}
\switchcolumn*

\LA
\noindent Quod Dóminus invitátus venit ad núptias, étiam excépta mýstica significatióne, confirmáre vóluit, quod ipse fecit núptias. Futúri enim erant, de quibus dixit Apóstolus, prohibéntes núbere, et dicéntes quod malum essent núptiæ, et quod diábolus eas fecísset: cum idem Dóminus dicat in Evangélio interrogátus, utrum líceat hómini dimíttere uxórem suam ex quálibet causa, non licére, excépta causa fornicatiónis. In qua responsióne, si meminístis, hoc ait: Quod Deus conjúnxit, homo non séparet.\par
\switchcolumn
\EN
\noindent Even setting aside any mystical interpretation, the fact that the Lord was pleased to be asked, and to go to a marriage, showeth plainly enough that He is the Author and Blesser of marriage. There were yet to be those of whom the Apostle hath warned us as forbidding to marry; who say that marriage is a bad thing in itself, and a work of the devil. Yet we read in the Gospel that when the Lord was asked, Is it lawful for a man to put away his wife for every cause? He answered that it was not lawful, except it were for fornication. In which answer ye will remember that He used these words: What God hath joined together, let not man put asunder.\par
\switchcolumn*

\LA
\begin{Resp}\Rsign Ad te, Dómine, levávi ánimam meam: \Star{} Deus meus, in te confído, non erubéscam.\end{Resp}
\begin{Resp}\Vsign Custódi ánimam meam, et éripe me.\end{Resp}
\begin{Resp}\Rsign Deus meus, in te confído, non erubéscam.\end{Resp}
\switchcolumn
\EN
\begin{Resp}\Rsign Unto thee, O Lord, do I lift up my soul. \Star{} O my God, I trust in thee, let me not be ashamed.\end{Resp}
\begin{Resp}\Vsign O keep my soul and deliver me.\end{Resp}
\begin{Resp}\Rsign O my God, I trust in thee, let me not be ashamed.\end{Resp}
\switchcolumn*

\LA
\LessonHead{Lectio VIII}
\switchcolumn
\EN
\LessonHead{Lesson VIII}
\switchcolumn*

\LA
\noindent Et qui bene erudíti sunt in fide catholica novérunt, quod Deus fécerit nuptias: et sicut conjúnctio a Deo, ita divortium a diabolo sit. Sed proptérea in causa fornicatiónis licet uxórem dimíttere: quia ipsa esse uxor prior nóluit, quæ fidem conjugalem maríto non servávit. Nec illæ, quæ virginitátem Deo vovent, quamquam ampliórem gradum honoris et sanctitátis in Ecclésia teneant, sine nuptiis sunt: nam et ipsæ pertinent ad nuptias cum tota Ecclésia, in quibus nuptiis sponsus est Christus.\par
\switchcolumn
\EN
\noindent They who are well instructed in the Catholic religion know that God is the Author and Blesser of marriage; and that, whereas joining together in marriage is of God, divorce is of the devil. But it is lawful for a man to put away his wife in case of fornication, For by not keeping a wife's faith to her husband she herself hath first willed not to be wife. They also who have made a vow of their virginity to God and have thereby attained to an higher degree of honour and holiness in the Church, are not unmarried, for they are a special part of the marriage of the whole Church, which is the Bride of Christ.\par
\switchcolumn*

\LA
\begin{Resp}\Rsign Duo Séraphim clamábant alter ad álterum: \Star{} Sanctus, sanctus, sanctus Dóminus Deus Sábaoth: * Plena est omnis terra glória ejus.\end{Resp}
\begin{Resp}\Vsign Tres sunt qui testimónium dant in cælo: Pater, Verbum, et Spíritus Sanctus: et hi tres unum sunt.\end{Resp}
\begin{Resp}\Rsign Sanctus, sanctus, sanctus Dóminus Deus Sábaoth.\end{Resp}
\begin{Resp}\Vsign Glória Patri, et Fílio, \Star{} et Spirítui Sancto.\end{Resp}
\begin{Resp}\Rsign Plena est omnis terra glória ejus.\end{Resp}
\switchcolumn
\EN
\begin{Resp}\Rsign One Seraph cried unto another: \Star{} Holy, Holy, Holy is the Lord God of hosts; the whole earth is full of His glory.\end{Resp}
\begin{Resp}\Vsign There are Three That bear record in heaven: the Father, the Word, and the Holy Ghost, and these Three are One.\end{Resp}
\begin{Resp}\Rsign Holy, Holy, Holy is the Lord God of hosts.\end{Resp}
\begin{Resp}\Vsign Glory be to the Father, and to the Son, \Star{} and to the Holy Ghost.\end{Resp}
\begin{Resp}\Rsign The whole earth is full of His glory.\end{Resp}
\switchcolumn*

\LA
\LessonHead{Lectio IX}
\switchcolumn
\EN
\LessonHead{Lesson IX}
\switchcolumn*

\LA
\noindent Ac per hoc ergo Dóminus invitátus venit ad nuptias, ut conjugális cástitas firmarétur, et ostenderétur sacraméntum nuptiárum: quia et illárum nuptiárum sponsus personam Dómini figurábat cui dictum est: Servásti vinum bonum usque adhuc. Bonum enim vinum Christus servávit usque adhuc, id est, Evangélium suum.\par
\switchcolumn
\EN
\noindent Lord, being asked, went to the marriage, to strengthen the marriage tie, and to shed light on the hidden meaning of matrimony. In that marriage feast the Bridegroom to whom it was said, “Thou hast kept the good wine until now,” was a figure of the Lord Christ, Who hath kept until now the good wine, namely the Gospel.\par
\switchcolumn*

\LA
\TeDeum{Te Deum}
\switchcolumn
\EN
\TeDeum{Te Deum}
\switchcolumn*

\end{paracol}

\Office{Feria II infra Hebdomadam II post Epiphaniam}{Monday of the Second Week after Epiphany}{Feria}
\addcontentsline{toc}{subsection}{Feria II infra Hebdomadam II post Epiphaniam\enspace\textit{Monday of the Second Week after Epiphany}}
\begin{paracol}{2}
\LA
\LessonHead{Lectio I}
\switchcolumn
\EN
\LessonHead{Lesson I}
\switchcolumn*

\LA
\LessonTitle{De Epístola secúnda ad Corínthios}
\switchcolumn
\EN
\LessonTitle{Lesson from the second letter of St. Paul the Apostle to the Corinthians}
\switchcolumn*

\LA
\Cite{2 Cor 3:1-3}
\switchcolumn
\EN
\Cite{2 Cor 3:1-3}
\switchcolumn*

\LA
\noindent Incípimus íterum nosmetípsos commendáre? aut numquid egémus (sicut quidam) commendatítiis epístolis ad vos, aut ex vobis? Epístola nostra vos estis, scripta in córdibus nostris, quæ scitur, et legitur ab ómnibus homínibus: manifestáti quod epístola estis Christi, ministráta a nobis, et scripta non atraménto, sed Spíritu Dei vivi: non in tábulis lapídeis, sed in tábulis cordis carnálibus.\par
\switchcolumn
\EN
\noindent Do we begin again to commend ourselves? Or do we need (as some do) epistles of commendation to you, or from you? You are our epistle, written in our hearts, which is known and read by all men: Being manifested, that you are the epistle of Christ, ministered by us, and written not with ink, but with the Spirit of the living God; not in tables of stone, but in the fleshly tables of the heart.\par
\switchcolumn*

\LA
\begin{Resp}\Rsign Quam magna multitúdo dulcédinis tuæ, Dómine, \Star{} Quam abscondísti timéntibus te!\end{Resp}
\begin{Resp}\Vsign Et perfecísti eis qui sperant in te, Dómine, in conspéctu filiórum hóminum.\end{Resp}
\begin{Resp}\Rsign Quam abscondísti timéntibus te!\end{Resp}
\switchcolumn
\EN
\begin{Resp}\Rsign O how great is thy goodness, O Lord \Star{} Which Thou hast laid up for them that fear thee!\end{Resp}
\begin{Resp}\Vsign Which Thou, O Lord, hast wrought for them that trust in thee before the sons of men!\end{Resp}
\begin{Resp}\Rsign Which Thou hast laid up for them that fear thee!\end{Resp}
\switchcolumn*

\LA
\LessonHead{Lectio II}
\switchcolumn
\EN
\LessonHead{Lesson II}
\switchcolumn*

\LA
\Cite{2 Cor 3:4-8}
\switchcolumn
\EN
\Cite{2 Cor 3:4-8}
\switchcolumn*

\LA
\noindent Fidúciam autem talem habémus per Christum ad Deum: non quod sufficiéntes simus cogitáre áliquid a nobis, quasi ex nobis: sed sufficiéntia nostra ex Deo est: qui et idóneos nos fecit minístros novi Testaménti, non líttera, sed Spíritu: líttera enim occídit, Spíritus autem vivíficat. Quod si ministrátio mortis lítteris deformáta in lapídibus, fuit in glória, ita ut non possent inténdere fílii Israël in fáciem Móysi propter glóriam vultus ejus, quæ evacuátur: quómodo non magis ministrátio Spíritus erit in glória?\par
\switchcolumn
\EN
\noindent And such confidence we have, through Christ, towards God. Not that we are sufficient to think any thing of ourselves, as of ourselves: but our sufficiency is from God. Who also hath made us fit ministers of the new testament, not in the letter, but in the spirit. For the letter killeth, but the spirit quickeneth. Now if the ministration of death, engraven with letters upon stones, was glorious; so that the children of Israel could not steadfastly behold the face of Moses, for the glory of his countenance, which is made void: How shall not the ministration of the spirit be rather in glory?\par
\switchcolumn*

\LA
\begin{Resp}\Rsign Adjútor meus esto, Deus: \Star{} Ne derelínquas me.\end{Resp}
\begin{Resp}\Vsign Neque despícias me, Deus, salutáris meus.\end{Resp}
\begin{Resp}\Rsign Ne derelínquas me.\end{Resp}
\switchcolumn
\EN
\begin{Resp}\Rsign O God, be Thou my helper. \Star{} Neither leave me.\end{Resp}
\begin{Resp}\Vsign Nor forsake me O God of my salvation.\end{Resp}
\begin{Resp}\Rsign Neither leave me.\end{Resp}
\switchcolumn*

\LA
\LessonHead{Lectio III}
\switchcolumn
\EN
\LessonHead{Lesson III}
\switchcolumn*

\LA
\Cite{2 Cor 3:9-14}
\switchcolumn
\EN
\Cite{2 Cor 3:9-14}
\switchcolumn*

\LA
\noindent Nam si ministrátio damnatiónis glória est: multo magis abúndat ministérium justítiæ in glória. Nam nec glorificátum est, quod cláruit in hac parte, propter excelléntem glóriam. Si enim quod evacuátur, per glóriam est: multo magis quod manet, in glória est. Habéntes ígitur talem spem, multa fidúcia útimur: et non sicut Móyses ponébat velámen super fáciem suam, ut non inténderent fílii Israël in fáciem ejus, quod evacuátur, sed obtúsi sunt sensus eórum. Usque in hodiérnum enim diem idípsum velámen in lectióne véteris Testaménti manet non revelátum, quóniam in Christo evacuátur.\par
\switchcolumn
\EN
\noindent For if the ministration of condemnation be glory, much more the ministration of justice aboundeth in glory. For even that which was glorious in this part was not glorified, by reason of the glory that excelleth. For if that which is done away was glorious, much more that which remaineth is in glory. Having therefore such hope, we use much confidence: And not as Moses put a veil upon his face, that the children of Israel might not steadfastly look on the face of that which is made void. But their senses were made dull. For, until this present day, the selfsame veil, in the reading of the old testament, remaineth not taken away (because in Christ it is made void).\par
\switchcolumn*

\LA
\begin{Resp}\Rsign Benedícam Dóminum in omni témpore: \Star{} Semper laus ejus in ore meo.\end{Resp}
\begin{Resp}\Vsign In Dómino laudábitur ánima mea, áudiant mansuéti, et læténtur.\end{Resp}
\begin{Resp}\Rsign Semper laus ejus in ore meo.\end{Resp}
\begin{Resp}\Vsign Glória Patri, et Fílio, \Star{} et Spirítui Sancto.\end{Resp}
\begin{Resp}\Rsign Semper laus ejus in ore meo.\end{Resp}
\switchcolumn
\EN
\begin{Resp}\Rsign I will bless the Lord at all times. \Star{} His praise shall continually be in my mouth.\end{Resp}
\begin{Resp}\Vsign My soul shall make her boast in the Lord; the humble shall hear thereof and be glad.\end{Resp}
\begin{Resp}\Rsign His praise shall continually be in my mouth.\end{Resp}
\begin{Resp}\Vsign Glory be to the Father, and to the Son, \Star{} and to the Holy Ghost.\end{Resp}
\begin{Resp}\Rsign His praise shall continually be in my mouth.\end{Resp}
\switchcolumn*

\end{paracol}

\Office{Feria III infra Hebdomadam II post Epiphaniam}{Tuesday of the Second Week after Epiphany}{Feria}
\addcontentsline{toc}{subsection}{Feria III infra Hebdomadam II post Epiphaniam\enspace\textit{Tuesday of the Second Week after Epiphany}}
\begin{paracol}{2}
\LA
\LessonHead{Lectio I}
\switchcolumn
\EN
\LessonHead{Lesson I}
\switchcolumn*

\LA
\LessonTitle{De Epístola secúnda ad Corínthios}
\switchcolumn
\EN
\LessonTitle{Lesson from the second letter of St. Paul the Apostle to the Corinthians}
\switchcolumn*

\LA
\Cite{2 Cor 5:1-4}
\switchcolumn
\EN
\Cite{2 Cor 5:1-4}
\switchcolumn*

\LA
\noindent Scimus enim quóniam, si terréstris domus nostra hujus habitatiónis dissolvátur, quod ædificatiónem ex Deo habémus, domum non manufáctam, ætérnam in cælis. Nam et in hoc ingemíscimus, habitatiónem nostram, quæ de cælo est, superíndui cupiéntes: si tamen vestíti, non nudi inveniámur. Nam et qui sumus in hoc tabernáculo, ingemíscimus graváti: eo quod nólumus exspoliári, sed supervestíri, ut absorbeátur quod mortále est, a vita.\par
\switchcolumn
\EN
\noindent For we know, if our earthly house of this habitation be dissolved, that we have a building of God, a house not made with hands, eternal in heaven. For in this also we groan, desiring to be clothed upon with our habitation that is from heaven. Yet so that we be found clothed, not naked. For we also, who are in this tabernacle, do groan, being burthened; because we would not be unclothed, but clothed upon, that that which is mortal may be swallowed up by life.\par
\switchcolumn*

\LA
\begin{Resp}\Rsign Auribus pércipe, Deus, lácrimas meas: ne síleas a me, remítte mihi: \Star{} Quóniam íncola ego sum apud te, et peregrínus.\end{Resp}
\begin{Resp}\Vsign Compláceat tibi, ut erípias me: Dómine, ad adjuvándum me festína.\end{Resp}
\begin{Resp}\Rsign Quóniam íncola ego sum apud te, et peregrínus.\end{Resp}
\switchcolumn
\EN
\begin{Resp}\Rsign O God, give ear unto my tears; hold not thy peace, but, O, spare me! \Star{} For I am a stranger with thee, and a sojourner.\end{Resp}
\begin{Resp}\Vsign Be pleased, O Lord, to deliver me; O Lord, look upon me to help me.\end{Resp}
\begin{Resp}\Rsign For I am a stranger with thee, and a sojourner.\end{Resp}
\switchcolumn*

\LA
\LessonHead{Lectio II}
\switchcolumn
\EN
\LessonHead{Lesson II}
\switchcolumn*

\LA
\Cite{2 Cor 5:6-10}
\switchcolumn
\EN
\Cite{2 Cor 5:6-10}
\switchcolumn*

\LA
\noindent Audéntes ígitur semper, sciéntes, quóniam dum sumus in córpore, peregrinámur a Dómino: (per fidem enim ambulámus, et non per spéciem) audémus autem, et bonam voluntátem habémus magis peregrinári a córpore, et præséntes esse ad Dóminum. Et ídeo conténdimus sive abséntes, sive præséntes, placére illi. Omnes enim nos manifestári opórtet ante tribúnal Christi, ut réferat unusquísque própria córporis, prout gessit, sive bonum, sive malum.\par
\switchcolumn
\EN
\noindent Therefore having always confidence, knowing that, while we are in the body, we are absent from the Lord. (For we walk by faith, and not by sight.) But we are confident, and have a good will to be absent rather from the body, and to be present with the Lord. And therefore we labour, whether absent or present, to please him. For we must all be manifested before the judgement seat of Christ, that every one may receive the proper things of the body, according as he hath done, whether it be good or evil.\par
\switchcolumn*

\LA
\begin{Resp}\Rsign Státuit Dóminus supra petram pedes meos, et diréxit gressus meos Deus: \Star{} Et misit in os meum cánticum novum.\end{Resp}
\begin{Resp}\Vsign Exaudívit preces meas: et edúxit me de lacu misériæ.\end{Resp}
\begin{Resp}\Rsign Et misit in os meum cánticum novum.\end{Resp}
\switchcolumn
\EN
\begin{Resp}\Rsign The Lord hath set my feet upon a rock, and ordered my goings. \Star{} And He hath put a new song in my mouth.\end{Resp}
\begin{Resp}\Vsign He heard my cry He brought me up also out of an horrible pit.\end{Resp}
\begin{Resp}\Rsign And He hath put a new song in my mouth.\end{Resp}
\switchcolumn*

\LA
\LessonHead{Lectio III}
\switchcolumn
\EN
\LessonHead{Lesson III}
\switchcolumn*

\LA
\Cite{2 Cor 5:11-15}
\switchcolumn
\EN
\Cite{2 Cor 5:11-15}
\switchcolumn*

\LA
\noindent Sciéntes ergo timórem Dómini homínibus suadémus, Deo autem manifésti sumus. Spero autem et in consciéntiis vestris maniféstos nos esse. Non íterum commendámus nos vobis, sed occasiónem damus vobis gloriándi pro nobis: ut habeátis ad eos qui in fácie gloriántur, et non in corde. Sive enim mente excédimus, Deo: sive sóbrii sumus, vobis. Cáritas enim Christi urget nos: æstimántes hoc, quóniam si unus pro ómnibus mórtuus est, ergo omnes mórtui sunt: et pro ómnibus mórtuus est Christus: ut et qui vivunt, jam non sibi vivant, sed ei qui pro ipsis mórtuus est et resurréxit.\par
\switchcolumn
\EN
\noindent Knowing therefore the fear of the Lord, we use persuasion to men; but to God we are manifest. And I trust also that in your consciences we are manifest. We commend not ourselves again to you, but give you occasion to glory in our behalf; that you may have somewhat to answer them who glory in face, and not in heart. For whether we be transported in mind, it is to God; or whether we be sober, it is for you. For the charity of Christ presseth us: judging this, that if one died for all, then all were dead. And Christ died for all; that they also who live, may not now live to themselves, but unto him who died for them, and rose again.\par
\switchcolumn*

\LA
\begin{Resp}\Rsign Ego dixi, Dómine, miserére mei: \Star{} Sana ánimam meam, quia peccávi tibi.\end{Resp}
\begin{Resp}\Vsign Ab ómnibus iniquitátibus meis éripe me, Dómine.\end{Resp}
\begin{Resp}\Rsign Sana ánimam meam, quia peccávi tibi.\end{Resp}
\begin{Resp}\Vsign Glória Patri, et Fílio, \Star{} et Spirítui Sancto.\end{Resp}
\begin{Resp}\Rsign Sana ánimam meam, quia peccávi tibi.\end{Resp}
\switchcolumn
\EN
\begin{Resp}\Rsign I said: Lord, be merciful unto me. \Star{} Heal my soul, for I have sinned against thee.\end{Resp}
\begin{Resp}\Vsign Deliver me from all mine iniquities, O Lord.\end{Resp}
\begin{Resp}\Rsign Heal my soul, for I have sinned against thee.\end{Resp}
\begin{Resp}\Vsign Glory be to the Father, and to the Son, \Star{} and to the Holy Ghost.\end{Resp}
\begin{Resp}\Rsign Heal my soul, for I have sinned against thee.\end{Resp}
\switchcolumn*

\end{paracol}

\Office{Feria IV infra Hebdomadam II post Epiphaniam}{Wednesday of the Second Week after Epiphany}{Feria}
\addcontentsline{toc}{subsection}{Feria IV infra Hebdomadam II post Epiphaniam\enspace\textit{Wednesday of the Second Week after Epiphany}}
\begin{paracol}{2}
\LA
\LessonHead{Lectio I}
\switchcolumn
\EN
\LessonHead{Lesson I}
\switchcolumn*

\LA
\LessonTitle{De Epístola secúnda ad Corínthios}
\switchcolumn
\EN
\LessonTitle{Lesson from the second letter of St. Paul the Apostle to the Corinthians}
\switchcolumn*

\LA
\Cite{2 Cor 7:1-3}
\switchcolumn
\EN
\Cite{2 Cor 7:1-3}
\switchcolumn*

\LA
\noindent Has ergo habéntes promissiónes, caríssimi, mundémus nos ab omni inquinaménto carnis et spíritus, perficiéntes sanctificatiónem in timóre Dei. Cápite nos. Néminem lǽsimus, néminem corrúpimus, néminem circumvénimus. Non ad condemnatiónem vestram dico. Prædíximus enim, quod in córdibus nostris estis ad commoriéndum, et ad convivéndum.\par
\switchcolumn
\EN
\noindent Having therefore these promises, dearly beloved, let us cleanse ourselves from all defilement of the flesh and of the spirit, perfecting sanctification in the fear of God. Receive us. We have injured no man, we have corrupted no man, we have overreached no man. I speak not this to your condemnation. For we have said before, that you are in our hearts, to die together, and to live together.\par
\switchcolumn*

\LA
\begin{Resp}\Rsign Ne perdíderis me cum iniquitátibus meis: \Star{} Neque in finem irátus resérves mala mea.\end{Resp}
\begin{Resp}\Vsign Non intres in judícium cum servo tuo, Dómine.\end{Resp}
\begin{Resp}\Rsign Neque in finem irátus resérves mala mea.\end{Resp}
\switchcolumn
\EN
\begin{Resp}\Rsign Cut me not off in the midst of my sins. \Star{} Nor keep thy wrath against me for my latter end.\end{Resp}
\begin{Resp}\Vsign Enter not into judgment with thy servant, O Lord.\end{Resp}
\begin{Resp}\Rsign Nor keep thy wrath against me for my latter end.\end{Resp}
\switchcolumn*

\LA
\LessonHead{Lectio II}
\switchcolumn
\EN
\LessonHead{Lesson II}
\switchcolumn*

\LA
\Cite{2 Cor 7:4-7}
\switchcolumn
\EN
\Cite{2 Cor 7:4-7}
\switchcolumn*

\LA
\noindent Multa mihi fidúcia est apud vos, multa mihi gloriátio pro vobis, replétus sum consolatióne, superabúndo gáudio in omni tribulatióne nostra. Nam et cum venissémus in Macedóniam, nullam réquiem hábuit caro nostra, sed omnem tribulatiónem passi sumus: foris pugnæ, intus timóres. Sed qui consolátur húmiles, consolátus est nos Deus in advéntu Titi. Non solum autem in advéntu ejus, sed étiam in consolatióne, qua consolátus est in vobis, réferens nobis vestrum desidérium, vestrum fletum, vestram æmulatiónem pro me, ita ut magis gaudérem.\par
\switchcolumn
\EN
\noindent Great is my confidence for you, great is my glorying for you. I am filled with comfort: I exceedingly abound with joy in all our tribulation. For also when we were come into Macedonia, our flesh had no rest, but we suffered all tribulation; combats without, fears within. But God, who comforteth the humble, comforted us by the coming of Titus. And not by his coming only, but also by the consolation, wherewith he was comforted in you, relating to us your desire, your mourning, your zeal for me, so that I rejoiced the more.\par
\switchcolumn*

\LA
\begin{Resp}\Rsign Parátum cor meum, Deus, parátum cor meum: \Star{} Cantábo et psalmum dicam Dómino.\end{Resp}
\begin{Resp}\Vsign Exsúrge, glória mea, exsúrge, psaltérium et cíthara, exsúrgam dilúculo.\end{Resp}
\begin{Resp}\Rsign Cantábo et psalmum dicam Dómino.\end{Resp}
\switchcolumn
\EN
\begin{Resp}\Rsign My heart is ready, O God, my heart is ready. \Star{} I will sing and give praise to the Lord.\end{Resp}
\begin{Resp}\Vsign Awake up, my glory, awake, psaltery and harp! I will awake early.\end{Resp}
\begin{Resp}\Rsign I will sing and give praise to the Lord.\end{Resp}
\switchcolumn*

\LA
\LessonHead{Lectio III}
\switchcolumn
\EN
\LessonHead{Lesson III}
\switchcolumn*

\LA
\Cite{2 Cor 7:8-10}
\switchcolumn
\EN
\Cite{2 Cor 7:8-10}
\switchcolumn*

\LA
\noindent Quóniam etsi contristávi vos in epístola, non me pœ́nitet: et si pœnitéret, videns quod epístola illa (etsi ad horam) vos contristávit; nunc gáudeo: non quia contristáti estis, sed quia contristáti estis ad pœniténtiam. Contristáti enim estis secúndum Deum, ut in nullo detriméntum patiámini ex nobis. Quæ enim secúndum Deum tristítia est, pœniténtiam in salútem stábilem operátur: sǽculi autem tristítia mortem operátur.\par
\switchcolumn
\EN
\noindent For although I made you sorrowful by my epistle, I do not repent; and if I did repent, seeing that the same epistle (although but for a time) did make you sorrowful; Now I am glad: not because you were made sorrowful; but because you were made sorrowful unto penance. For you were made sorrowful according to God, that you might suffer damage by us in nothing. For the sorrow that is according to God worketh penance, steadfast unto salvation; but the sorrow of the world worketh death.\par
\switchcolumn*

\LA
\begin{Resp}\Rsign Adjútor meus, tibi psallam, quia, Deus, suscéptor meus es: \Star{} Deus meus, misericórdia mea.\end{Resp}
\begin{Resp}\Vsign Lætábor, et exsultábo in te, psallam nómini tuo, Altíssime.\end{Resp}
\begin{Resp}\Rsign Deus meus, misericórdia mea.\end{Resp}
\begin{Resp}\Vsign Glória Patri, et Fílio, \Star{} et Spirítui Sancto.\end{Resp}
\begin{Resp}\Rsign Deus meus, misericórdia mea.\end{Resp}
\switchcolumn
\EN
\begin{Resp}\Rsign Unto thee, O my Strength, will I sing, for God is my defence, \Star{} The God of my mercy.\end{Resp}
\begin{Resp}\Vsign I will be glad and rejoice in thee, I will sing praise to thy Name, O Thou Most High.\end{Resp}
\begin{Resp}\Rsign The God of my mercy.\end{Resp}
\begin{Resp}\Vsign Glory be to the Father, and to the Son, \Star{} and to the Holy Ghost.\end{Resp}
\begin{Resp}\Rsign The God of my mercy.\end{Resp}
\switchcolumn*

\end{paracol}

\Office{Feria V infra Hebdomadam II post Epiphaniam}{Thursday of the Second Week after Epiphany}{Feria}
\addcontentsline{toc}{subsection}{Feria V infra Hebdomadam II post Epiphaniam\enspace\textit{Thursday of the Second Week after Epiphany}}
\begin{paracol}{2}
\LA
\LessonHead{Lectio I}
\switchcolumn
\EN
\LessonHead{Lesson I}
\switchcolumn*

\LA
\LessonTitle{De Epístola secúnda ad Corínthios}
\switchcolumn
\EN
\LessonTitle{Lesson from the second letter of St. Paul the Apostle to the Corinthians}
\switchcolumn*

\LA
\Cite{2 Cor 10:1-3}
\switchcolumn
\EN
\Cite{2 Cor 10:1-3}
\switchcolumn*

\LA
\noindent Ipse autem ego Paulus óbsecro vos per mansuetúdinem et modéstiam Christi, qui in fácie quidem húmilis sum inter vos, absens autem confído in vobis. Rogo autem vos, ne præsens áudeam per eam confidéntiam, qua exístimor audére in quosdam, qui arbitrántur nos tamquam secúndum carnem ambulémus. In carne enim ambulántes, non secúndum carnem militámus.\par
\switchcolumn
\EN
\noindent Now I Paul myself beseech you, by the mildness and modesty of Christ, who in presence indeed am lowly among you, but being absent, am bold toward you. But I beseech you, that I may not be bold when I am present, with that confidence wherewith I am thought to be bold, against some, who reckon us as if we walked according to the flesh. For though we walk in the flesh, we do not war according to the flesh.\par
\switchcolumn*

\LA
\begin{Resp}\Rsign Deus, in te sperávi, Dómine, non confúndar in ætérnum: in justítia tua líbera me, \Star{} Et éripe me.\end{Resp}
\begin{Resp}\Vsign Inclína ad me aurem tuam, et salva me.\end{Resp}
\begin{Resp}\Rsign Et éripe me.\end{Resp}
\switchcolumn
\EN
\begin{Resp}\Rsign In thee, O God, do I put my trust; let me never be put to confusion, O Lord deliver me in thy righteousness, \Star{} And cause me to escape.\end{Resp}
\begin{Resp}\Vsign Incline thine ear unto me, deliver me speedily.\end{Resp}
\begin{Resp}\Rsign And cause me to escape.\end{Resp}
\switchcolumn*

\LA
\LessonHead{Lectio II}
\switchcolumn
\EN
\LessonHead{Lesson II}
\switchcolumn*

\LA
\Cite{2 Cor 10:4-7}
\switchcolumn
\EN
\Cite{2 Cor 10:4-7}
\switchcolumn*

\LA
\noindent Nam arma milítiæ nostræ non carnália sunt, sed poténtia Deo ad destructiónem munitiónum, consília destruéntes, et omnem altitúdinem extolléntem se advérsus sciéntiam Dei, et in captivitátem redigéntes omnem intelléctum in obséquium Christi, et in promptu habéntes ulcísci omnem inobediéntiam, cum impléta fúerit vestra obediéntia. Quæ secúndum fáciem sunt, vidéte. Si quis confídit sibi Christi se esse, hoc cógitet íterum apud se: quia sicut ipse Christi est, ita et nos.\par
\switchcolumn
\EN
\noindent For the weapons of our warfare are not carnal, but mighty to God unto the pulling down of fortifications, destroying counsels, And every height that exhalteth itself against the knowledge of God, and bringing into captivity every understanding unto the obedience of Christ; And having in readiness to revenge all disobedience, when your obedience shall be fulfilled. See the things that are according to outward appearance. If any man trust to himself, that he is Christ's, let him think this again with himself, that as he is Christ's, so are we also.\par
\switchcolumn*

\LA
\begin{Resp}\Rsign Repleátur os meum laude tua, ut hymnum dicam glóriæ tuæ, tota die magnitúdinem tuam: noli me proícere in témpore senectútis: \Star{} Dum defécerit in me virtus mea, ne derelínquas me.\end{Resp}
\begin{Resp}\Vsign Gaudébunt lábia mea cum cantávero tibi.\end{Resp}
\begin{Resp}\Rsign Dum defécerit in me virtus mea, ne derelínquas me.\end{Resp}
\switchcolumn
\EN
\begin{Resp}\Rsign Let my mouth be filled with thy praise, that I may sing of thy glory, all the day long of thy greatness. Cast me not off in the time of old age; \Star{} Forsake me not when my strength faileth.\end{Resp}
\begin{Resp}\Vsign My lips shall be fain when I sing unto thee.\end{Resp}
\begin{Resp}\Rsign Forsake me not when my strength faileth.\end{Resp}
\switchcolumn*

\LA
\LessonHead{Lectio III}
\switchcolumn
\EN
\LessonHead{Lesson III}
\switchcolumn*

\LA
\Cite{2 Cor 10:8-12}
\switchcolumn
\EN
\Cite{2 Cor 10:8-12}
\switchcolumn*

\LA
\noindent Nam, et si ámplius áliquid gloriátus fúero de potestáte nostra, quam dedit nobis Dóminus in ædificatiónem, et non in destructiónem vestram: non erubéscam. Ut autem non exístimer tamquam terrére vos per epístolas: quóniam quidem epístolæ, ínquiunt, graves sunt et fortes: præséntia autem córporis infírma, et sermo contemptíbilis: hoc cógitet qui ejúsmodi est, quia quales sumus verbo per epístolas abséntes, tales et præséntes in facto. Non enim audémus insérere, aut comparáre nos quibúsdam, qui seípsos comméndant: sed ipsi in nobis nosmetípsos metiéntes, et comparántes nosmetípsos nobis.\par
\switchcolumn
\EN
\noindent For if also I should boast somewhat more of our power, which the Lord hath given us unto edification, and not for your destruction, I should not be ashamed. But that I may not be thought as it were to terrify you by epistles, (For his epistles indeed, say they, are weighty and strong; but his bodily presence is weak, and his speech contemptible,) Let such a one think this, that such as we are in word by epistles, when absent, such also we will be indeed when present. For we dare not match, or compare ourselves with some, that commend themselves; but we measure ourselves by ourselves, and compare ourselves with ourselves.\par
\switchcolumn*

\LA
\begin{Resp}\Rsign Gaudébunt lábia mea cum cantávero tibi: \Star{} Et ánima mea, quam redemísti, Dómine.\end{Resp}
\begin{Resp}\Vsign Sed et lingua mea meditábitur justítiam tuam, tota die laudem tuam.\end{Resp}
\begin{Resp}\Rsign Et ánima mea, quam redemísti, Dómine.\end{Resp}
\begin{Resp}\Vsign Glória Patri, et Fílio, \Star{} et Spirítui Sancto.\end{Resp}
\begin{Resp}\Rsign Et ánima mea, quam redemísti, Dómine.\end{Resp}
\switchcolumn
\EN
\begin{Resp}\Rsign My lips shall be fain when I sing unto thee; \Star{} And my soul, which Thou, O Lord, hast redeemed.\end{Resp}
\begin{Resp}\Vsign My tongue shall also talk of thy righteousness, all the day long of thy praise.\end{Resp}
\begin{Resp}\Rsign And my soul, which Thou, O Lord, hast redeemed.\end{Resp}
\begin{Resp}\Vsign Glory be to the Father, and to the Son, \Star{} and to the Holy Ghost.\end{Resp}
\begin{Resp}\Rsign And my soul, which Thou, O Lord, hast redeemed.\end{Resp}
\switchcolumn*

\end{paracol}

\Office{Feria VI infra Hebdomadam II post Epiphaniam}{Friday of the Second Week after Epiphany}{Feria}
\addcontentsline{toc}{subsection}{Feria VI infra Hebdomadam II post Epiphaniam\enspace\textit{Friday of the Second Week after Epiphany}}
\begin{paracol}{2}
\LA
\LessonHead{Lectio I}
\switchcolumn
\EN
\LessonHead{Lesson I}
\switchcolumn*

\LA
\LessonTitle{De Epístola secúnda ad Corínthios}
\switchcolumn
\EN
\LessonTitle{Lesson from the second letter of St. Paul the Apostle to the Corinthians}
\switchcolumn*

\LA
\Cite{2 Cor 12:1-4}
\switchcolumn
\EN
\Cite{2 Cor 12:1-4}
\switchcolumn*

\LA
\noindent Si gloriári opórtet (non éxpedit quidem) véniam autem ad visiónes, et revelatiónes Dómini. Scio hóminem in Christo ante annos quatuórdecim (sive in córpore néscio, sive extra corpus, néscio, Deus scit:) raptum hujúsmodi usque ad tértium cælum. Et scio hujúsmodi hóminem, (sive in córpore, sive extra corpus, néscio, Deus scit:) quóniam raptus est in paradísum: et audívit arcána verba, quæ non licet hómini loqui.\par
\switchcolumn
\EN
\noindent If I must glory (it is not expedient indeed): but I will come to visions and revelations of the Lord. I know a man in Christ above fourteen years ago (whether in the body, I know not, or out of the body, I know not; God knoweth), such a one caught up to the third heaven. And I know such a man (whether in the body, or out of the body, I know not: God knoweth): That he was caught up into paradise, and heard secret words, which it is not granted to man to utter.\par
\switchcolumn*

\LA
\begin{Resp}\Rsign Confitébor tibi, Dómine Deus, in toto corde meo, et honorificábo nomen tuum in ætérnum: \Star{} Quia misericórdia tua, Dómine, magna est super me.\end{Resp}
\begin{Resp}\Vsign Deus meus es tu, et confitébor tibi: Deus meus es tu, et exaltábo te.\end{Resp}
\begin{Resp}\Rsign Quia misericórdia tua, Dómine, magna est super me.\end{Resp}
\switchcolumn
\EN
\begin{Resp}\Rsign I will praise thee, O Lord my God, with all my heart, and I will glorify thy Name for evermore. \Star{} For great is thy mercy, O Lord, toward me.\end{Resp}
\begin{Resp}\Vsign Thou art my God, and I will praise thee; Thou art my God, and I will exalt thee.\end{Resp}
\begin{Resp}\Rsign For great is thy mercy, O Lord, toward me.\end{Resp}
\switchcolumn*

\LA
\LessonHead{Lectio II}
\switchcolumn
\EN
\LessonHead{Lesson II}
\switchcolumn*

\LA
\Cite{2 Cor 12:5-9}
\switchcolumn
\EN
\Cite{2 Cor 12:5-9}
\switchcolumn*

\LA
\noindent Pro hujúsmodi gloriábor: pro me autem nihil gloriábor, nisi in infirmitátibus meis. Nam, et si volúero gloriári, non ero insípiens: veritátem enim dicam: parco autem, ne quis me exístimet supra id, quod videt in me, aut áliquid audit ex me. Et ne magnitúdo revelatiónum extóllat me, datus est mihi stímulus carnis meæ ángelus sátanæ, qui me colaphízet. Propter quod ter Dóminum rogávi, ut discéderet a me: et dixit mihi: Súfficit tibi grátia mea: nam virtus in infirmitáte perfícitur.\par
\switchcolumn
\EN
\noindent For such an one I will glory; but for myself I will glory nothing, but in my infirmities. For though I should have a mind to glory, I shall not be foolish; for I will say the truth. But I forbear, lest any man should think of me above that which he seeth in me, or any thing he heareth from me. And lest the greatness of the revelations should exalt me, there was given me a sting of my flesh, an angel of Satan, to buffet me. For which thing thrice I besought the Lord, that it might depart from me. And he said to me: My grace is sufficient for thee; for power is made perfect in infirmity.\par
\switchcolumn*

\LA
\begin{Resp}\Rsign Misericórdia tua, Dómine, magna est super me: \Star{} Et liberásti ánimam meam ex inférno inferióri.\end{Resp}
\begin{Resp}\Vsign In die tribulatiónis meæ clamávi ad te, quia exaudísti me.\end{Resp}
\begin{Resp}\Rsign Et liberásti ánimam meam ex inférno inferióri.\end{Resp}
\switchcolumn
\EN
\begin{Resp}\Rsign Great, O Lord, is thy mercy toward me. \Star{} And Thou hast delivered my soul from the lowest hell.\end{Resp}
\begin{Resp}\Vsign In the day of my trouble I called upon thee, for Thou hast heard me.\end{Resp}
\begin{Resp}\Rsign And Thou hast delivered my soul from the lowest hell.\end{Resp}
\switchcolumn*

\LA
\LessonHead{Lectio III}
\switchcolumn
\EN
\LessonHead{Lesson III}
\switchcolumn*

\LA
\Cite{2 Cor 12:9-11}
\switchcolumn
\EN
\Cite{2 Cor 12:9-11}
\switchcolumn*

\LA
\noindent Libénter ígitur gloriábor in infirmitátibus meis, ut inhábitet in me virtus Christi. Propter quod pláceo mihi in infirmitátibus meis, in contuméliis, in necessitátibus, in persecutiónibus, in angústiis pro Christo: Cum enim infírmor, tunc potens sum. Factus sum insípiens, vos me coëgístis. Ego enim a vobis débui commendári: nihil enim minus fui ab iis, qui sunt supra modum, Apóstoli: tamétsi nihil sum.\par
\switchcolumn
\EN
\noindent Gladly therefore will I glory in my infirmities, that the power of Christ may dwell in me. For which cause I please myself in my infirmities, in reproaches, in necessities, in persecutions, in distresses, for Christ. For when I am weak, then am I powerful. I am become foolish: you have compelled me. For I ought to have been commended by you: for I have no way come short of them that are above measure apostles, although I be nothing.\par
\switchcolumn*

\LA
\begin{Resp}\Rsign Factus est mihi Dóminus in refúgium: \Star{} Et Deus meus in auxílium spei meæ.\end{Resp}
\begin{Resp}\Vsign Erípuit me de inimícis meis fortíssimis, et factus est Dóminus protéctor meus.\end{Resp}
\begin{Resp}\Rsign Et Deus meus in auxílium spei meæ.\end{Resp}
\begin{Resp}\Vsign Glória Patri, et Fílio, \Star{} et Spirítui Sancto.\end{Resp}
\begin{Resp}\Rsign Et Deus meus in auxílium spei meæ.\end{Resp}
\switchcolumn
\EN
\begin{Resp}\Rsign The Lord is my refuge. \Star{} And my God is the stay of my trust.\end{Resp}
\begin{Resp}\Vsign He delivered me from the strongest of mine enemies, and the Lord was my stay.\end{Resp}
\begin{Resp}\Rsign And my God is the stay of my trust.\end{Resp}
\begin{Resp}\Vsign Glory be to the Father, and to the Son, \Star{} and to the Holy Ghost.\end{Resp}
\begin{Resp}\Rsign And my God is the stay of my trust.\end{Resp}
\switchcolumn*

\end{paracol}

\Office{Sabbato infra Hebdomadam II post Epiphaniam}{Saturday of the Second Week after Epiphany}{Feria}
\addcontentsline{toc}{subsection}{Sabbato infra Hebdomadam II post Epiphaniam\enspace\textit{Saturday of the Second Week after Epiphany}}
\begin{paracol}{2}
\LA
\LessonHead{Lectio I}
\switchcolumn
\EN
\LessonHead{Lesson I}
\switchcolumn*

\LA
\LessonTitle{De Epístola secúnda ad Corínthios}
\switchcolumn
\EN
\LessonTitle{Lesson from the second letter of St. Paul the Apostle to the Corinthians}
\switchcolumn*

\LA
\Cite{2 Cor 13:1-4}
\switchcolumn
\EN
\Cite{2 Cor 13:1-4}
\switchcolumn*

\LA
\noindent Ecce tértio hoc vénio ad vos: In ore duórum vel trium téstium stabit omne verbum. Prædíxi, et prædíco, ut præsens, et nunc absens iis, qui ante peccavérunt, et céteris ómnibus, quóniam si vénero íterum, non parcam. An experiméntum quǽritis, ejus, qui in me lóquitur Christus, qui in vobis non infirmátur, sed potens est in vobis? Nam etsi crucifíxus est ex infirmitáte: sed vivit ex virtúte Dei. Nam et nos infírmi sumus in illo: sed vivémus cum eo ex virtúte Dei in vobis.\par
\switchcolumn
\EN
\noindent Behold, this is the third time I am coming to you: In the mouth of two or three witnesses shall every word stand. I have told before, and foretell, as present, and now absent, to them that sinned before, and to all the rest, that if I come again, I will not spare. Do you seek a proof of Christ that speaketh in me, who towards you is not weak, but is mighty in you? For although he was crucified through weakness, yet he liveth by the power of God. For we also are weak in him: but we shall live with him by the power of God towards you.\par
\switchcolumn*

\LA
\begin{Resp}\Rsign Misericórdiam et judícium cantábo tibi, Dómine: \Star{} Psallam et intéllegam in via immaculáta, quando vénies ad me.\end{Resp}
\begin{Resp}\Vsign Perambulábam in innocéntia cordis mei, in médio domus meæ.\end{Resp}
\begin{Resp}\Rsign Psallam et intéllegam in via immaculáta, quando vénies ad me.\end{Resp}
\switchcolumn
\EN
\begin{Resp}\Rsign Mercy and judgment I will sing to thee, O Lord: \Star{} I will sing, and I will understand in the unspotted way, when thou shalt come to me.\end{Resp}
\begin{Resp}\Vsign I walked in the innocence of my heart, in the midst of my house.\end{Resp}
\begin{Resp}\Rsign Mercy and judgment I will sing to thee, O Lord\end{Resp}
\switchcolumn*

\LA
\LessonHead{Lectio II}
\switchcolumn
\EN
\LessonHead{Lesson II}
\switchcolumn*

\LA
\Cite{2 Cor 13:5-9}
\switchcolumn
\EN
\Cite{2 Cor 13:5-9}
\switchcolumn*

\LA
\noindent Vosmetípsos tentáte si estis in fide: ipsi vos probáte. An non cognóscitis vosmetípsos, quia Christus Jesus in vobis est? nisi forte réprobi estis. Spero autem quod cognoscétis, quia nos non sumus réprobi. Orámus autem Deum, ut nihil mali faciátis, non ut nos probáti appareámus, sed ut vos quod bonum est faciátis: nos autem ut réprobi simus. Non enim póssumus áliquid advérsus veritátem, sed pro veritáte. Gaudémus enim, quóniam nos infírmi sumus, vos autem poténtes estis. Hoc et orámus vestram consummatiónem.\par
\switchcolumn
\EN
\noindent Try your own selves if you be in the faith; prove ye yourselves. Know you not your own selves, that Christ Jesus is in you, unless perhaps you be reprobates? But I trust that you shall know that we are not reprobates. Now we pray God, that you may do no evil, not that we may appear approved, but that you may do that which is good, and that we may be as reprobates. For we can do nothing against the truth; but for the truth. For we rejoice that we are weak, and you are strong. This also we pray for, your perfection.\par
\switchcolumn*

\LA
\begin{Resp}\Rsign Dómine, exáudi oratiónem meam, et clamor meus ad te pervéniat: \Star{} Quia non spernis, Deus, preces páuperum.\end{Resp}
\begin{Resp}\Vsign Fiant aures tuæ intendéntes in oratiónem servi tui.\end{Resp}
\begin{Resp}\Rsign Quia non spernis, Deus, preces páuperum.\end{Resp}
\switchcolumn
\EN
\begin{Resp}\Rsign Hear my prayer, O Lord, and let my crying come unto thee. \Star{} For thou wilt not despise, O God, the prayer of the destitute\end{Resp}
\begin{Resp}\Vsign O let thine ears consider well the voice of my complaint\end{Resp}
\begin{Resp}\Rsign For thou wilt not despise, O God, the prayer of the destitute.\end{Resp}
\switchcolumn*

\LA
\LessonHead{Lectio III}
\switchcolumn
\EN
\LessonHead{Lesson III}
\switchcolumn*

\LA
\Cite{2 Cor 13:10-13}
\switchcolumn
\EN
\Cite{2 Cor 13:10-13}
\switchcolumn*

\LA
\noindent Ideo hæc absens scribo, ut non præsens dúrius agam secúndum potestátem, quam Dóminus dedit mihi in ædificatiónem, et non in destructiónem. De cétero, fratres, gaudéte, perfécti estóte, exhortámini, idem sápite, pacem habéte, et Deus pacis et dilectiónis erit vobíscum. Salutáte ínvicem in ósculo sancto. Salútant vos omnes sancti. Grátia Dómini nostri Jesu Christi, et cáritas Dei, et communicátio Sancti Spíritus sit cum ómnibus vobis. Amen.\par
\switchcolumn
\EN
\noindent Therefore I write these things, being absent, that, being present, I may not deal more severely, according to the power which the Lord hath given me unto edification, and not unto destruction. For the rest, brethren, rejoice, be perfect, take exhortation, be of one mind, have peace; and the God of peace and of love shall be with you. Salute one another with a holy kiss. All the saints salute you. The grace of our Lord Jesus Christ, and the charity of God, and the communication of the Holy Ghost be with you all. Amen.\par
\switchcolumn*

\LA
\begin{Resp}\Rsign Velóciter exáudi me, Deus, \Star{} Tu autem idem ipse es, et anni tui non defícient.\end{Resp}
\begin{Resp}\Vsign Dies mei sicut umbra declinavérunt, et ego sicut fœnum árui.\end{Resp}
\begin{Resp}\Rsign Quia defecérunt sicut fumus dies mei.\end{Resp}
\begin{Resp}\Vsign Glória Patri, et Fílio, \Star{} et Spirítui Sancto.\end{Resp}
\begin{Resp}\Rsign Tu autem idem ipse es, et anni tui non defícient.\end{Resp}
\switchcolumn
\EN
\begin{Resp}\Rsign Haste thee to hear me, O Lord. \Star{} But thou art the same, and thy years shall not fail.\end{Resp}
\begin{Resp}\Vsign My days have declined like a shadow, and I am withered like grass.\end{Resp}
\begin{Resp}\Rsign For my days are consumed away like smoke.\end{Resp}
\begin{Resp}\Vsign Glory be to the Father, and to the Son, \Star{} and to the Holy Ghost.\end{Resp}
\begin{Resp}\Rsign But thou art the same, and thy years shall not fail.\end{Resp}
\switchcolumn*

\end{paracol}

\Office{Dominica III Post Epiphaniam}{Third Sunday after Epiphany}{Semiduplex}
\addcontentsline{toc}{subsection}{Dominica III Post Epiphaniam\enspace\textit{Third Sunday after Epiphany}}
\begin{paracol}{2}
\LA
\LessonHead{Lectio I}
\switchcolumn
\EN
\LessonHead{Lesson I}
\switchcolumn*

\LA
\LessonTitle{Incipit Epístola beáti Pauli Apóstoli ad Gálatas}
\switchcolumn
\EN
\LessonTitle{Lesson from the letter of St. Paul the Apostle to the Galatians}
\switchcolumn*

\LA
\Cite{Gal 1:1-5}
\switchcolumn
\EN
\Cite{Gal 1:1-5}
\switchcolumn*

\LA
\noindent Paulus Apóstolus non ab homínibus, neque per hóminem, sed per Jesum Christum, et Deum Patrem, qui suscitávit eum a mórtuis: et qui mecum sunt omnes fratres, ecclésiis Galátiæ. Grátia vobis, et pax a Deo Patre, et Dómino nostro Jesu Christo, qui dedit semetípsum pro peccátis nostris, ut eríperet nos de præsénti sǽculo nequam, secúndum voluntátem Dei et Patris nostri, cui est glória in sǽcula sæculórum. Amen.\par
\switchcolumn
\EN
\noindent Paul, an apostle, not of men, neither by man, but by Jesus Christ, and God the Father, who raised him from the dead, And all the brethren who are with me, to the churches of Galatia. Grace be to you, and peace from God the Father, and from our Lord Jesus Christ, Who gave himself for our sins, that he might deliver us from this present wicked world, according to the will of God and our Father: To whom is glory for ever and ever. Amen.\par
\switchcolumn*

\LA
\begin{Resp}\Rsign Dómine, ne in ira tua árguas me, neque in furóre tuo corrípias me: \Star{} Miserére mei, Dómine, quóniam infírmus sum.\end{Resp}
\begin{Resp}\Vsign Timor et tremor venérunt super me, et contexérunt me ténebræ.\end{Resp}
\begin{Resp}\Rsign Miserére mei, Dómine, quóniam infírmus sum.\end{Resp}
\switchcolumn
\EN
\begin{Resp}\Rsign O Lord, rebuke me not in thine anger, neither chasten me in thine hot displeasure. \Star{} Have mercy upon me, O Lord, for I am weak.\end{Resp}
\begin{Resp}\Vsign Fearfulness and trembling are come upon me, and darkness hath overwhelmed me.\end{Resp}
\begin{Resp}\Rsign Have mercy upon me, O Lord, for I am weak.\end{Resp}
\switchcolumn*

\LA
\LessonHead{Lectio II}
\switchcolumn
\EN
\LessonHead{Lesson II}
\switchcolumn*

\LA
\Cite{Gal 1:6-10}
\switchcolumn
\EN
\Cite{Gal 1:6-10}
\switchcolumn*

\LA
\noindent Miror quod sic tam cito transferímini ab eo, qui vos vocávit in grátiam Christi in áliud Evangélium: quod non est áliud, nisi sunt áliqui, qui vos contúrbant, et volunt convértere Evangélium Christi. Sed licet nos, aut Angelus de cælo evangelízet vobis prætérquam quod evangelizávimus vobis, anáthema sit. Sicut prædíximus, et nunc íterum dico: Si quis vobis evangelizáverit præter id, quod accepístis, anáthema sit. Modo enim homínibus suádeo, an Deo? An quæro homínibus placére? Si adhuc homínibus placérem, Christi servus non essem.\par
\switchcolumn
\EN
\noindent I wonder that you are so soon removed from him that called you into the grace of Christ, unto another gospel. Which is not another, only there are some that trouble you, and would pervert the gospel of Christ. But though we, or an angel from heaven, preach a gospel to you besides that which we have preached to you, let him be anathema. As we said before, so now I say again: If any one preach to you a gospel, besides that which you have received, let him be anathema. For do I now persuade men, or God? Or do I seek to please men? If I yet pleased men, I should not be the servant of Christ.\par
\switchcolumn*

\LA
\begin{Resp}\Rsign Deus, qui sedes super thronum, et júdicas æquitátem, esto refúgium páuperum in tribulatióne: \Star{} Quia tu solus labórem et dolórem consíderas.\end{Resp}
\begin{Resp}\Vsign Tibi enim derelíctus est pauper, pupíllo tu eris adjútor.\end{Resp}
\begin{Resp}\Rsign Quia tu solus labórem et dolórem consíderas.\end{Resp}
\switchcolumn
\EN
\begin{Resp}\Rsign O God, Which satest in the throne judging right, be Thou a refuge for the poor, a refuge in times of trouble. \Star{} For Thou alone beholdest mischief and spite.\end{Resp}
\begin{Resp}\Vsign The poor leaveth himself unto thee; Thou wilt be the helper of the fatherless.\end{Resp}
\begin{Resp}\Rsign For Thou alone beholdest mischief and spite.\end{Resp}
\switchcolumn*

\LA
\LessonHead{Lectio III}
\switchcolumn
\EN
\LessonHead{Lesson III}
\switchcolumn*

\LA
\Cite{Gal 1:11-14}
\switchcolumn
\EN
\Cite{Gal 1:11-14}
\switchcolumn*

\LA
\noindent Notum enim vobis fácio, fratres, Evangélium, quod evangelizátum est a me, quia non est secúndum hóminem: neque enim ego ab hómine accépi illud, neque dídici, sed per revelatiónem Jesu Christi. Audístis enim conversatiónem meam aliquándo in Judaísmo: quóniam supra modum persequébar Ecclésiam Dei, et expugnábam illam, et proficiébam in Judaísmo supra multos coætáneos meos in génere meo, abundántius æmulátor exístens paternárum meárum traditiónum.\par
\switchcolumn
\EN
\noindent For I give you to understand, brethren, that the gospel which was preached by me is not according to man. For neither did I receive it of man, nor did I learn it; but by the revelation of Jesus Christ. For you have heard of my conversation in time past in the Jews' religion: how that, beyond measure, I persecuted the church of God, and wasted it. And I made progress in the Jews' religion above many of my equals in my own nation, being more abundantly zealous for the traditions of my fathers.\par
\switchcolumn*

\LA
\begin{Resp}\Rsign A dextris est mihi Dóminus, ne commóvear: \Star{} Propter hoc dilatátum est cor meum, et exsultávit lingua mea.\end{Resp}
\begin{Resp}\Vsign Dóminus pars hereditátis meæ, et cálicis mei.\end{Resp}
\begin{Resp}\Rsign Propter hoc dilatátum est cor meum, et exsultávit lingua mea.\end{Resp}
\begin{Resp}\Vsign Glória Patri, et Fílio, \Star{} et Spirítui Sancto.\end{Resp}
\begin{Resp}\Rsign Propter hoc dilatátum est cor meum, et exsultávit lingua mea.\end{Resp}
\switchcolumn
\EN
\begin{Resp}\Rsign The Lord is at my right hand, I shall never be moved. \Star{} Therefore my heart is glad, and my tongue rejoiceth.\end{Resp}
\begin{Resp}\Vsign The Lord is the portion of mine inheritance, and of my cup.\end{Resp}
\begin{Resp}\Rsign Therefore my heart is glad, and my tongue rejoiceth.\end{Resp}
\begin{Resp}\Vsign Glory be to the Father, and to the Son, \Star{} and to the Holy Ghost.\end{Resp}
\begin{Resp}\Rsign Therefore my heart is glad, and my tongue rejoiceth.\end{Resp}
\switchcolumn*

\LA
\LessonHead{Lectio IV}
\switchcolumn
\EN
\LessonHead{Lesson IV}
\switchcolumn*

\LA
\LessonTitle{De Expositióne sancti Augustíni Epíscopi in Epístolam ad Gálatas}
\switchcolumn
\EN
\LessonTitle{From the Exposition of the Epistle to the Galatians by St. Augustine, Bishop (of Hippo.)}
\switchcolumn*

\LA
\Source{Præfátio tom. 4}
\switchcolumn
\EN
\Source{Preface, Bk. iv.}
\switchcolumn*

\LA
\noindent Causa, propter quam scribit Apóstolus ad Gálatas, hæc est: ut intélligant grátiam Dei id secum ágere, ut sub lege jam non sint. Cum enim prædicáta eis esset Evangélii grátia, non defuérunt quidam ex circumcisióne, quamvis Christiáni nómine, nondum tamen tenéntes ipsum gratiæ benefícium, et adhuc voléntes esse sub onéribus legis, quæ Dóminus Deus imposúerat non justítiæ serviéntibus, sed peccáto, justam scílicet legem injústis homínibus dando ad demonstranda peccáta eórum, non auferenda. Non enim aufert peccáta, nisi grátia fidei, quæ per dilectiónem operátur.\par
\switchcolumn
\EN
\noindent The reason of the Apostle's writing to the Galatians was this that they might understand that the grace of God had worked in them that they were no longer under the law. For when the grace of the Gospel was preached to them, there had not been wanting to them some of them of the circumcision, Christians indeed in name, but who had not yet apprehended that great benefit of grace, and desiring still to be bound with burdens of the law burdens which the Lord God had laid, not upon such as serve righteousness, but upon such as serve sin, laying, that is to say, upon the unrighteous a righteous law, whereby their unrighteousness was made manifest, not taken away. For there is not anything which taketh away sin, save only the grace of faith which worketh by love.\par
\switchcolumn*

\LA
\begin{Resp}\Rsign Notas mihi fecísti, Dómine, vias vitæ: \Star{} Adimplébis me lætítia cum vultu tuo: delectatiónes in déxtera tua usque in finem.\end{Resp}
\begin{Resp}\Vsign Tu es qui restítues hereditátem meam mihi.\end{Resp}
\begin{Resp}\Rsign Adimplébis me lætítia cum vultu tuo: delectatiónes in déxtera tua usque in finem.\end{Resp}
\switchcolumn
\EN
\begin{Resp}\Rsign O Lord, Thou hast shown me the path of life. \Star{} Thou shalt fill me with joy in thy presence, at thy right hand there are pleasures for evermore.\end{Resp}
\begin{Resp}\Vsign Thou art He That shalt restore mine inheritance unto me.\end{Resp}
\begin{Resp}\Rsign Thou shalt fill me with joy in thy presence, at thy right hand there are pleasures for evermore.\end{Resp}
\switchcolumn*

\LA
\LessonHead{Lectio V}
\switchcolumn
\EN
\LessonHead{Lesson V}
\switchcolumn*

\LA
\noindent Sub hac ergo grátia jam Gálatas constitútos illi volébant constitúere sub onéribus legis, asseverántes nihil eis prodésse Evangélium, nisi circumcideréntur, et céteras carnáles Judáici ritus observatiónes subírent. Et ídeo Paulum Apóstolum suspéctum habére cœ́perant, a quo illis Evangélium prædicátum erat, tamquam non tenéntem disciplínam ceterórum Apostolórum, qui gentes cogébant judaizáre.\par
\switchcolumn
\EN
\noindent The men of the circumcision would have the Galatians, who were under grace, to be under the burdens of the law, persuading them that the Gospel profited them nothing, unless they should be circumcised, and take on them the other outward observances of the Jews' religion. Whence the Galatians began to have doubts of the Apostle Paul, by whom the Gospel had been preached to them, as one that held not the doctrine of the other Apostles, who compelled the Gentiles to come under the law.\par
\switchcolumn*

\LA
\begin{Resp}\Rsign Díligam te, Dómine, virtus mea: Dóminus firmaméntum meum, \Star{} Et refúgium meum.\end{Resp}
\begin{Resp}\Vsign Liberátor meus, Deus meus, adjútor meus.\end{Resp}
\begin{Resp}\Rsign Et refúgium meum.\end{Resp}
\switchcolumn
\EN
\begin{Resp}\Rsign I will love thee, O Lord, my strength; the Lord is my rock \Star{} And my fortress.\end{Resp}
\begin{Resp}\Vsign My Deliverer, my God, mine Helper.\end{Resp}
\begin{Resp}\Rsign And my fortress.\end{Resp}
\switchcolumn*

\LA
\LessonHead{Lectio VI}
\switchcolumn
\EN
\LessonHead{Lesson VI}
\switchcolumn*

\LA
\noindent Talis quidem quǽstio est et in Epístola ad Romános: verúmtamen vidétur áliquid interésse, quod ibi contentiónem ipsam dírimit, litémque compónit, quæ inter eos, qui ex Judǽis, et eos, qui ex Géntibus credíderant, orta erat: cum illi tamquam ex méritis óperum legis, sibi rédditum Evangélii prǽmium arbitraréntur, quod prǽmium incircumcísis tamquam imméritis nolébant dari: illi contra Judǽis se præférre gestírent, tamquam interfectóribus Dómini. In hac vero epístola ad eos scribit, qui jam commóti erant auctoritáte illórum, qui ex Judǽis erant, et ad observatiónes legis cogébant.\par
\switchcolumn
\EN
\noindent The same question is discussed in the Epistle to the Romans, but with this difference in that case the Apostle putteth an end to the discussion, and stilleth the strife which had arisen between the Jewish and the Gentile converts, in consequence of the Jews holding that they had earned the knowledge of the Gospel as a reward for their observance of the law, and grudging the same knowledge to the uncircumcised, as to men who had done nothing to deserve it; and the Gentiles, on the contrary, maintaining that they were superior to the Jews, in that they were not the murderers of the Lord. Now, in this Epistle to the Galatians, the Apostle addresseth himself to those who were troubled by the authority claimed by them who were of the circumcision, and sought to bring into subjection to the law them who were of the uncircumcision.\par
\switchcolumn*

\LA
\begin{Resp}\Rsign Dómini est terra, et plenitúdo ejus: \Star{} Orbis terrárum, et univérsi qui hábitant in eo.\end{Resp}
\begin{Resp}\Vsign Ipse super mária fundávit eam, et super flúmina præparávit illam.\end{Resp}
\begin{Resp}\Rsign Orbis terrárum, et univérsi qui hábitant in eo.\end{Resp}
\begin{Resp}\Vsign Glória Patri, et Fílio, \Star{} et Spirítui Sancto.\end{Resp}
\begin{Resp}\Rsign Orbis terrárum, et univérsi qui hábitant in eo.\end{Resp}
\switchcolumn
\EN
\begin{Resp}\Rsign The earth is the Lord's, and the fulness thereof \Star{} The world, and they that dwell therein.\end{Resp}
\begin{Resp}\Vsign For He hath founded it upon the seas, and established it upon the floods.\end{Resp}
\begin{Resp}\Rsign The world, and they that dwell therein.\end{Resp}
\begin{Resp}\Vsign Glory be to the Father, and to the Son, \Star{} and to the Holy Ghost.\end{Resp}
\begin{Resp}\Rsign The world, and they that dwell therein.\end{Resp}
\switchcolumn*

\LA
\LessonHead{Lectio VII}
\switchcolumn
\EN
\LessonHead{Lesson VII}
\switchcolumn*

\LA
\LessonTitle{Léctio sancti Evangélii secúndum Matthǽum}
\switchcolumn
\EN
\LessonTitle{From the Holy Gospel according to Matthew}
\switchcolumn*

\LA
\Cite{Matt 8:1-13}
\switchcolumn
\EN
\Cite{Matt 8:1-13}
\switchcolumn*

\LA
\noindent In illo témpore: Cum descendísset Jesus de monte, secútæ sunt eum turbæ multæ: et ecce leprósus véniens, adorábat eum. Et réliqua.\par
\switchcolumn
\EN
\noindent At that time, when Jesus was come down from the mountain, great multitudes followed him: And behold a leper came and adored him. And so on.\par
\switchcolumn*

\LA
\LessonTitle{Homilía sancti Hierónymi Presbýteri}
\switchcolumn
\EN
\LessonTitle{Homily by St. Jerome, Priest (at Bethlehem.)}
\switchcolumn*

\LA
\Source{Liber 1 Comment. in cap. 8 Matth.}
\switchcolumn
\EN
\Source{Bk. i, Comm. on Matth. viii}
\switchcolumn*

\LA
\noindent De monte Dómino descendénte, occúrrunt turbæ, quia ad altióra ascéndere non valuérunt. Et primus ei occúrrit leprósus: necdum enim póterat cum lepra tam multíplicem in monte Salvatóris audíre sermónem. Et notándum, quod hic primus speciáliter curátus sit: secúndo, puer centuriónis: tértio, socrus Petri fébriens in Caphárnaum: quarto loco, qui obláti sunt ei a dæmónio vexáti: quorum spíritus verbo eiciébat, quando et omnes male habéntes curávit.\par
\switchcolumn
\EN
\noindent When the Lord was come down from the mountain, great multitudes followed Him. They were not able to follow Him when He went up. And first there came a leper. This poor creature's disease had prevented him from hearing the Saviour's long sermon on the Mount. Let it be noted that he is the first person specially named as being healed. The second was the Centurion's servant; the third was Peter's wife's mother, who was sick of a fever at Capernaum; the fourth were they who were brought unto Christ as being troubled with evil spirits, from whom He by His word cast out the evil spirits, at the same time that He healed all that were sick.\par
\switchcolumn*

\LA
\begin{Resp}\Rsign Ad te, Dómine, levávi ánimam meam: \Star{} Deus meus, in te confído, non erubéscam.\end{Resp}
\begin{Resp}\Vsign Custódi ánimam meam, et éripe me.\end{Resp}
\begin{Resp}\Rsign Deus meus, in te confído, non erubéscam.\end{Resp}
\switchcolumn
\EN
\begin{Resp}\Rsign Unto thee, O Lord, do I lift up my soul. \Star{} O my God, I trust in thee, let me not be ashamed.\end{Resp}
\begin{Resp}\Vsign O keep my soul and deliver me.\end{Resp}
\begin{Resp}\Rsign O my God, I trust in thee, let me not be ashamed.\end{Resp}
\switchcolumn*

\LA
\LessonHead{Lectio VIII}
\switchcolumn
\EN
\LessonHead{Lesson VIII}
\switchcolumn*

\LA
\noindent Et ecce leprósus véniens adorábat eum, dicens. Recte post prædicatiónem atque doctrínam, signi offértur occásio, ut per virtútem miráculi, prætéritus apud audiéntes sermo firmétur. Dómine, si vis, potes me mundáre. Qui voluntátem rogat, de virtúte non dúbitat. Et exténdens Jesus manum tétigit eum, dicens: Volo, mundáre. Extendénte manum Dómino, statim lepra fugit. Simúlque consídera, quam húmilis, et sine jactántia respónsio. Ille díxerat, Si vis: Dóminus respóndit, Volo. Ille præmíserat, Potes me mundáre: Dóminus jungit, et dicit, Mundáre. Non ergo, ut pleríque Latinórum putant, jungéndum est, et legéndum, Volo mundáre: sed separátim, ut primum dicat, Volo; deínde ímperet, Mundáre.\par
\switchcolumn
\EN
\noindent And, behold, there came a leper, and worshipped Him, saying: Properly after preaching and doctrine cometh occasion for a sign, that the power of the miracle might confirm in the hearers the truth of the teaching that had gone before. Lord, if Thou wilt, Thou canst make me clean. He that prayeth the Lord to have the will, doubteth not but that He hath the power. And Jesus put forth His hand, and touched him, saying: I will; be thou clean. As soon as the Lord put forth His Hand, the leprosy departed. Let us remark how lowly and unbragging is the Lord's language. The leper had said, If Thou wilt; the Lord answereth, I will. The leper, Thou canst make me clean; the Lord, Be thou clean. Most Latin readers, misled by the identity of form in that language between the Present Infinitive Active and the Second Person Singular Present Imperative Passive of the Verb, read Christ's answer as if it were, I will to make thee clean. This is wrong. The sentences are separate. First cometh the expression of volition, I will, then the command, Be thou clean.\par
\switchcolumn*

\LA
\begin{Resp}\Rsign Duo Séraphim clamábant alter ad álterum: \Star{} Sanctus, sanctus, sanctus Dóminus Deus Sábaoth: * Plena est omnis terra glória ejus.\end{Resp}
\begin{Resp}\Vsign Tres sunt qui testimónium dant in cælo: Pater, Verbum, et Spíritus Sanctus: et hi tres unum sunt.\end{Resp}
\begin{Resp}\Rsign Sanctus, sanctus, sanctus Dóminus Deus Sábaoth.\end{Resp}
\begin{Resp}\Vsign Glória Patri, et Fílio, \Star{} et Spirítui Sancto.\end{Resp}
\begin{Resp}\Rsign Plena est omnis terra glória ejus.\end{Resp}
\switchcolumn
\EN
\begin{Resp}\Rsign One Seraph cried unto another: \Star{} Holy, Holy, Holy is the Lord God of hosts; the whole earth is full of His glory.\end{Resp}
\begin{Resp}\Vsign There are Three That bear record in heaven: the Father, the Word, and the Holy Ghost, and these Three are One.\end{Resp}
\begin{Resp}\Rsign Holy, Holy, Holy is the Lord God of hosts.\end{Resp}
\begin{Resp}\Vsign Glory be to the Father, and to the Son, \Star{} and to the Holy Ghost.\end{Resp}
\begin{Resp}\Rsign The whole earth is full of His glory.\end{Resp}
\switchcolumn*

\LA
\LessonHead{Lectio IX}
\switchcolumn
\EN
\LessonHead{Lesson IX}
\switchcolumn*

\LA
\noindent Et ait illi Jesus: Vide, némini díxeris. Et revéra quid erat necésse ut sermóne jactáret, quod córpore præferébat? Sed vade, osténde te sacerdóti. Várias ob causas mittit eum ad sacerdótem: primum propter humilitátem, ut sacerdótibus deférre honórem videátur. Erat enim lege præcéptum, ut, qui mundáti fúerant a lepra, offérrent múnera sacerdótibus. Deínde, ut mundátum vidéntes leprósum, aut créderent Salvatóri, aut non créderent: si créderent, salvaréntur; si non créderent, inexcusábiles forent. Et simul, ne, quod in eo sæpíssime criminabántur, legem viderétur infríngere.\par
\switchcolumn
\EN
\noindent And Jesus saith unto him: See thou tell no man. What need was there to tell what his body showed? But go thy way, show thyself to the Priest. There were diverse reasons why Christ should send him to the Priest. First, for humility's sake, that He might show reverence to God's Priest. Then there was a command in the law that they that were cleansed of leprosy should make an offering to the Priests. Moreover, that, when the Priests saw the leper cleansed, they might either believe in the Saviour, or refuse to believe; if they believed, that they might be saved, and, if they believed not, that they might have no excuse. Lastly, that He might give no ground for the accusation that was so often brought against Him, that He was unobservant of the law.\par
\switchcolumn*

\LA
\TeDeum{Te Deum}
\switchcolumn
\EN
\TeDeum{Te Deum}
\switchcolumn*

\end{paracol}

\Office{Feria II infra Hebdomadam III post Epiphaniam}{Monday of the Third Week after Epiphany}{Feria}
\addcontentsline{toc}{subsection}{Feria II infra Hebdomadam III post Epiphaniam\enspace\textit{Monday of the Third Week after Epiphany}}
\begin{paracol}{2}
\LA
\RefLine{Lectiones I–III: ex Commúni Evangelistarum (C1a).}
\switchcolumn
\EN
\RefLine{Lessons I–III: from the Commune Evangelistarum (C1a).}
\switchcolumn*

\LA
\LessonHead{Lectio I}
\switchcolumn
\EN
\LessonHead{Lesson I}
\switchcolumn*

\LA
\LessonTitle{De Epístola ad Gálatas}
\switchcolumn
\EN
\LessonTitle{Lesson from the letter of St. Paul the Apostle to the Galatians}
\switchcolumn*

\LA
\Cite{Gal 3:1-6}
\switchcolumn
\EN
\Cite{Gal 3:1-6}
\switchcolumn*

\LA
\noindent O insensáti Gálatæ, quis vos fascinávit non obedíre veritáti, ante quorum óculos Jesus Christus præscríptus est, in vobis crucifíxus? Hoc solum a vobis volo díscere: ex opéribus legis Spíritum accepístis, an ex audítu fídei? sic stulti estis, ut cum Spíritu cœpéritis, nunc carne consummémini? tanta passi estis sine causa? si tamen sine causa. Qui ergo tríbuit vobis Spíritum, et operátur virtútes in vobis: ex opéribus legis, an ex audítu fídei? Sicut scriptum est: Abraham crédidit Deo, et reputátum est illi ad justítiam:\par
\switchcolumn
\EN
\noindent O senseless Galatians, who hath bewitched you that you should not obey the truth, before whose eyes Jesus Christ hath been set forth, crucified among you? This only would I learn of you: Did you receive the Spirit by the works of the law, or by the hearing of faith? Are you so foolish, that, whereas you began in the Spirit, you would now be made perfect by the flesh? Have you suffered so great things in vain? If it be yet in vain. He therefore who giveth to you the Spirit, and worketh miracles among you; doth he do it by the works of the law, or by the hearing of the faith? As it is written: Abraham believed God, and it was reputed to him unto justice.\par
\switchcolumn*

\LA
\begin{Resp}\Rsign Quam magna multitúdo dulcédinis tuæ, Dómine, \Star{} Quam abscondísti timéntibus te!\end{Resp}
\begin{Resp}\Vsign Et perfecísti eis qui sperant in te, Dómine, in conspéctu filiórum hóminum.\end{Resp}
\begin{Resp}\Rsign Quam abscondísti timéntibus te!\end{Resp}
\switchcolumn
\EN
\begin{Resp}\Rsign O how great is thy goodness, O Lord \Star{} Which Thou hast laid up for them that fear thee!\end{Resp}
\begin{Resp}\Vsign Which Thou, O Lord, hast wrought for them that trust in thee before the sons of men!\end{Resp}
\begin{Resp}\Rsign Which Thou hast laid up for them that fear thee!\end{Resp}
\switchcolumn*

\LA
\LessonHead{Lectio II}
\switchcolumn
\EN
\LessonHead{Lesson II}
\switchcolumn*

\LA
\Cite{Gal 3:7-10}
\switchcolumn
\EN
\Cite{Gal 3:7-10}
\switchcolumn*

\LA
\noindent Cognóscite ergo quia qui ex fide sunt, ii sunt fílii Abrahæ. Próvidens autem Scriptúra quia ex fide justíficat gentes Deus, prænuntiávit Abrahæ: Quia benedicéntur in te omnes gentes. Igitur qui ex fide sunt, benedicéntur cum fidéli Abraham. Quicúmque enim ex opéribus legis sunt, sub maledícto sunt. Scriptum est enim: Maledíctus omnis qui non permánserit in ómnibus quæ scripta sunt in libro legis ut fáciat ea.\par
\switchcolumn
\EN
\noindent Know ye therefore, that they who are of faith, the same are the children of Abraham. And the scripture, foreseeing, that God justifieth the Gentiles by faith, told unto Abraham before: In thee shall all nations be blessed. Therefore they that are of faith, shall be blessed with faithful Abraham. For as many as are of the works of the law, are under a curse. For it is written: Cursed is every one, that abideth not in all things, which are written in the book of the law to do them.\par
\switchcolumn*

\LA
\begin{Resp}\Rsign Adjútor meus esto, Deus: \Star{} Ne derelínquas me.\end{Resp}
\begin{Resp}\Vsign Neque despícias me, Deus, salutáris meus.\end{Resp}
\begin{Resp}\Rsign Ne derelínquas me.\end{Resp}
\switchcolumn
\EN
\begin{Resp}\Rsign O God, be Thou my helper. \Star{} Neither leave me.\end{Resp}
\begin{Resp}\Vsign Nor forsake me O God of my salvation.\end{Resp}
\begin{Resp}\Rsign Neither leave me.\end{Resp}
\switchcolumn*

\LA
\LessonHead{Lectio III}
\switchcolumn
\EN
\LessonHead{Lesson III}
\switchcolumn*

\LA
\Cite{Gal 3:11-14}
\switchcolumn
\EN
\Cite{Gal 3:11-14}
\switchcolumn*

\LA
\noindent Quóniam autem in lege nemo justificátur apud Deum, maniféstum est: quia justus ex fide vivit. Lex autem non est ex fide, sed: Qui fécerit ea, vivet in illis. Christus nos redémit de maledícto legis, factus pro nobis maledíctum: quia scriptum est: Maledíctus omnis qui pendet in ligno: ut in géntibus benedíctio Abrahæ fíeret in Christo Jesu, ut pollicitatiónem Spíritus accipiámus per fidem.\par
\switchcolumn
\EN
\noindent But that in the law no man is justified with God, it is manifest: because the just man liveth by faith. But the law is not of faith: but, He that doth those things, shall live in them. Christ hath redeemed us from the curse of the law, being made a curse for us: for it is written: Cursed is every one that hangeth on a tree: That the blessing of Abraham might come on the Gentiles through Christ Jesus: that we may receive the promise of the Spirit by faith.\par
\switchcolumn*

\LA
\begin{Resp}\Rsign Benedícam Dóminum in omni témpore: \Star{} Semper laus ejus in ore meo.\end{Resp}
\begin{Resp}\Vsign In Dómino laudábitur ánima mea, áudiant mansuéti, et læténtur.\end{Resp}
\begin{Resp}\Rsign Semper laus ejus in ore meo.\end{Resp}
\begin{Resp}\Vsign Glória Patri, et Fílio, \Star{} et Spirítui Sancto.\end{Resp}
\begin{Resp}\Rsign Semper laus ejus in ore meo.\end{Resp}
\switchcolumn
\EN
\begin{Resp}\Rsign I will bless the Lord at all times. \Star{} His praise shall continually be in my mouth.\end{Resp}
\begin{Resp}\Vsign My soul shall make her boast in the Lord; the humble shall hear thereof and be glad.\end{Resp}
\begin{Resp}\Rsign His praise shall continually be in my mouth.\end{Resp}
\begin{Resp}\Vsign Glory be to the Father, and to the Son, \Star{} and to the Holy Ghost.\end{Resp}
\begin{Resp}\Rsign His praise shall continually be in my mouth.\end{Resp}
\switchcolumn*

\end{paracol}

\Office{Feria III infra Hebdomadam III post Epiphaniam}{Tuesday of the Third Week after Epiphany}{Feria}
\addcontentsline{toc}{subsection}{Feria III infra Hebdomadam III post Epiphaniam\enspace\textit{Tuesday of the Third Week after Epiphany}}
\begin{paracol}{2}
\LA
\RefLine{Lectiones I–III: de Scriptura occurrente.}
\switchcolumn
\EN
\RefLine{Lessons I–III: from the Scripture of the day.}
\switchcolumn*

\LA
\LessonHead{Lectio I}
\switchcolumn
\EN
\LessonHead{Lesson I}
\switchcolumn*

\LA
\LessonTitle{De Epístola ad Gálatas}
\switchcolumn
\EN
\LessonTitle{Lesson from the letter of St. Paul the Apostle to the Galatians}
\switchcolumn*

\LA
\Cite{Gal 5:1-5}
\switchcolumn
\EN
\Cite{Gal 5:1-5}
\switchcolumn*

\LA
\noindent State, et nolíte íterum jugo servitútis continéri. Ecce ego Paulus dico vobis: quóniam si circumcidámini, Christus vobis nihil próderit. Testíficor autem rursus omni hómini circumcidénti se, quóniam débitor est univérsæ legis faciéndæ. Evacuáti estis a Christo, qui in lege justificámini: a grátia excidístis. Nos enim spíritu ex fide, spem justítiæ exspectámus.\par
\switchcolumn
\EN
\noindent Stand fast, and be not held again under the yoke of bondage. Behold, I Paul tell you, that if you be circumcised, Christ shall profit you nothing. And I testify again to every man circumcising himself, that he is a debtor to the whole law. You are made void of Christ, you who are justified in the law: you are fallen from grace. For we in spirit, by faith, wait for the hope of justice.\par
\switchcolumn*

\LA
\begin{Resp}\Rsign Auribus pércipe, Deus, lácrimas meas: ne síleas a me, remítte mihi: \Star{} Quóniam íncola ego sum apud te, et peregrínus.\end{Resp}
\begin{Resp}\Vsign Compláceat tibi, ut erípias me: Dómine, ad adjuvándum me festína.\end{Resp}
\begin{Resp}\Rsign Quóniam íncola ego sum apud te, et peregrínus.\end{Resp}
\switchcolumn
\EN
\begin{Resp}\Rsign O God, give ear unto my tears; hold not thy peace, but, O, spare me! \Star{} For I am a stranger with thee, and a sojourner.\end{Resp}
\begin{Resp}\Vsign Be pleased, O Lord, to deliver me; O Lord, look upon me to help me.\end{Resp}
\begin{Resp}\Rsign For I am a stranger with thee, and a sojourner.\end{Resp}
\switchcolumn*

\LA
\LessonHead{Lectio II}
\switchcolumn
\EN
\LessonHead{Lesson II}
\switchcolumn*

\LA
\Cite{Gal 5:6-10}
\switchcolumn
\EN
\Cite{Gal 5:6-10}
\switchcolumn*

\LA
\noindent Nam in Christo Jesu neque circumcísio áliquid valet, neque præpútium: sed fides, quæ per caritátem operátur. Currebátis bene: quis vos impedívit veritáti non obedíre? persuásio hæc non est ex eo, qui vocat vos. Módicum ferméntum totam massam corrúmpit. Ego confído in vobis in Dómino, quod nihil áliud sapiétis: qui autem contúrbat vos, portábit judícium, quicúmque est ille.\par
\switchcolumn
\EN
\noindent For in Christ Jesus neither circumcision availeth any thing, nor uncircumcision: but faith that worketh by charity. You did run well, who hath hindered you, that you should not obey the truth? This persuasion is not from him that calleth you. A little leaven corrupteth the whole lump. I have confidence in you in the Lord: that you will not be of another mind: but he that troubleth you, shall bear the judgment, whosoever he be.\par
\switchcolumn*

\LA
\begin{Resp}\Rsign Státuit Dóminus supra petram pedes meos, et diréxit gressus meos Deus: \Star{} Et misit in os meum cánticum novum.\end{Resp}
\begin{Resp}\Vsign Exaudívit preces meas: et edúxit me de lacu misériæ.\end{Resp}
\begin{Resp}\Rsign Et misit in os meum cánticum novum.\end{Resp}
\switchcolumn
\EN
\begin{Resp}\Rsign The Lord hath set my feet upon a rock, and ordered my goings. \Star{} And He hath put a new song in my mouth.\end{Resp}
\begin{Resp}\Vsign He heard my cry He brought me up also out of an horrible pit.\end{Resp}
\begin{Resp}\Rsign And He hath put a new song in my mouth.\end{Resp}
\switchcolumn*

\LA
\LessonHead{Lectio III}
\switchcolumn
\EN
\LessonHead{Lesson III}
\switchcolumn*

\LA
\Cite{Gal 5:11-17}
\switchcolumn
\EN
\Cite{Gal 5:11-17}
\switchcolumn*

\LA
\noindent Ego autem, fratres, si circumcisiónem adhuc prǽdico: quid adhuc persecutiónem pátior? ergo evacuátum est scándalum crucis. Utinam et abscindántur qui vos contúrbant. Vos enim in libertátem vocáti estis, fratres: tantum ne libertátem in occasiónem detis carnis, sed per caritátem Spíritus servíte ínvicem. Omnis enim lex in uno sermóne implétur: Díliges próximum tuum sicut teípsum. Quod si ínvicem mordétis, et coméditis: vidéte ne ab ínvicem consumámini. Dico autem: Spíritu ambuláte, et desidéria carnis non perficiétis. Caro enim concupíscit advérsus spíritum, spíritus autem advérsus carnem: hæc enim sibi ínvicem adversántur, ut non quæcúmque vultis, illa faciátis.\par
\switchcolumn
\EN
\noindent And I, brethren, if I yet preach circumcision, why do I yet suffer persecution? Then is the scandal of the cross made void. I would they were even cut off, who trouble you. For you, brethren, have been called unto liberty: only make not liberty an occasion to the flesh, but by charity of the spirit serve one another. For all the law is fulfilled in one word: Thou shalt love thy neighbour as thyself. But if you bite and devour one another; take heed you be not consumed one of another. I say then, walk in the spirit, and you shall not fulfil the lusts of the flesh. For the flesh lusteth against the spirit: and the spirit against the flesh; for these are contrary one to another: so that you do not the things that you would.\par
\switchcolumn*

\LA
\begin{Resp}\Rsign Ego dixi, Dómine, miserére mei: \Star{} Sana ánimam meam, quia peccávi tibi.\end{Resp}
\begin{Resp}\Vsign Ab ómnibus iniquitátibus meis éripe me, Dómine.\end{Resp}
\begin{Resp}\Rsign Sana ánimam meam, quia peccávi tibi.\end{Resp}
\begin{Resp}\Vsign Glória Patri, et Fílio, \Star{} et Spirítui Sancto.\end{Resp}
\begin{Resp}\Rsign Sana ánimam meam, quia peccávi tibi.\end{Resp}
\switchcolumn
\EN
\begin{Resp}\Rsign I said: Lord, be merciful unto me. \Star{} Heal my soul, for I have sinned against thee.\end{Resp}
\begin{Resp}\Vsign Deliver me from all mine iniquities, O Lord.\end{Resp}
\begin{Resp}\Rsign Heal my soul, for I have sinned against thee.\end{Resp}
\begin{Resp}\Vsign Glory be to the Father, and to the Son, \Star{} and to the Holy Ghost.\end{Resp}
\begin{Resp}\Rsign Heal my soul, for I have sinned against thee.\end{Resp}
\switchcolumn*

\end{paracol}

\Office{Feria IV infra Hebdomadam III post Epiphaniam}{Wednesday of the Third Week after Epiphany}{Feria}
\addcontentsline{toc}{subsection}{Feria IV infra Hebdomadam III post Epiphaniam\enspace\textit{Wednesday of the Third Week after Epiphany}}
\begin{paracol}{2}
\LA
\RefLine{Lectiones I–III: de Scriptura occurrente.}
\switchcolumn
\EN
\RefLine{Lessons I–III: from the Scripture of the day.}
\switchcolumn*

\LA
\LessonHead{Lectio I}
\switchcolumn
\EN
\LessonHead{Lesson I}
\switchcolumn*

\LA
\LessonTitle{Incipit Epístola beáti Pauli Apóstoli ad Ephésios}
\switchcolumn
\EN
\LessonTitle{Lesson from the letter of St. Paul the Apostle to the Ephesians}
\switchcolumn*

\LA
\Cite{Eph 1:1-4}
\switchcolumn
\EN
\Cite{Eph 1:1-4}
\switchcolumn*

\LA
\noindent Paulus Apóstolus Jesu Christi per voluntátem Dei, ómnibus sanctis qui sunt Ephesi, et fidélibus in Christo Jesu. Grátia vobis, et pax a Deo Patre nostro, et Dómino Jesu Christo. Benedíctus Deus et Pater Dómini nostri Jesu Christi, qui benedíxit nos in omni benedictióne spirituáli in cæléstibus in Christo, sicut elégit nos in ipso ante mundi constitutiónem, ut essémus sancti et immaculáti in conspéctu ejus in caritáte.\par
\switchcolumn
\EN
\noindent Paul, an apostle of Jesus Christ, by the will of God, to all the saints who are at Ephesus, and to the faithful in Christ Jesus. Grace be to you, and peace from God the Father, and from the Lord Jesus Christ. Blessed by the God and Father of our Lord Jesus Christ, who hath blessed us with spiritual blessings in heavenly places, in Christ: As he chose us in him before the foundation of the world, that we should be holy and unspotted in his sight in charity.\par
\switchcolumn*

\LA
\begin{Resp}\Rsign Ne perdíderis me cum iniquitátibus meis: \Star{} Neque in finem irátus resérves mala mea.\end{Resp}
\begin{Resp}\Vsign Non intres in judícium cum servo tuo, Dómine.\end{Resp}
\begin{Resp}\Rsign Neque in finem irátus resérves mala mea.\end{Resp}
\switchcolumn
\EN
\begin{Resp}\Rsign Cut me not off in the midst of my sins. \Star{} Nor keep thy wrath against me for my latter end.\end{Resp}
\begin{Resp}\Vsign Enter not into judgment with thy servant, O Lord.\end{Resp}
\begin{Resp}\Rsign Nor keep thy wrath against me for my latter end.\end{Resp}
\switchcolumn*

\LA
\LessonHead{Lectio II}
\switchcolumn
\EN
\LessonHead{Lesson II}
\switchcolumn*

\LA
\Cite{Eph 1:5-10}
\switchcolumn
\EN
\Cite{Eph 1:5-10}
\switchcolumn*

\LA
\noindent Qui prædestinávit nos in adoptiónem filiórum per Jesum Christum in ipsum: secúndum propósitum voluntátis suæ, in laudem glóriæ grátiæ suæ, in qua gratificávit nos in dilécto Fílio suo. In quo habémus redemptiónem per sánguinem ejus, remissiónem peccatórum secúndum divítias grátiæ ejus, quæ superabundávit in nobis in omni sapiéntia et prudéntia: ut notum fáceret nobis sacraméntum voluntátis suæ, secúndum beneplácitum ejus, quod propósuit in eo, in dispensatióne plenitúdinis témporum, instauráre ómnia in Christo, quæ in cælis et quæ in terra sunt, in ipso;\par
\switchcolumn
\EN
\noindent Who hath predestinated us unto the adoption of children through Jesus Christ unto himself: according to the purpose of his will: Unto the praise of the glory of his grace, in which he hath graced us in his beloved son. In whom we have redemption through his blood, the remission of sins, according to the riches of his grace, Which hath superabounded in us in all wisdom and prudence, That he might make known unto us the mystery of his will, according to his good pleasure, which he hath purposed in him, In the dispensation of the fulness of times, to re-establish all things in Christ, that are in heaven and on earth, in him.\par
\switchcolumn*

\LA
\begin{Resp}\Rsign Parátum cor meum, Deus, parátum cor meum: \Star{} Cantábo et psalmum dicam Dómino.\end{Resp}
\begin{Resp}\Vsign Exsúrge, glória mea, exsúrge, psaltérium et cíthara, exsúrgam dilúculo.\end{Resp}
\begin{Resp}\Rsign Cantábo et psalmum dicam Dómino.\end{Resp}
\switchcolumn
\EN
\begin{Resp}\Rsign My heart is ready, O God, my heart is ready. \Star{} I will sing and give praise to the Lord.\end{Resp}
\begin{Resp}\Vsign Awake up, my glory, awake, psaltery and harp! I will awake early.\end{Resp}
\begin{Resp}\Rsign I will sing and give praise to the Lord.\end{Resp}
\switchcolumn*

\LA
\LessonHead{Lectio III}
\switchcolumn
\EN
\LessonHead{Lesson III}
\switchcolumn*

\LA
\Cite{Eph 1:11-14}
\switchcolumn
\EN
\Cite{Eph 1:11-14}
\switchcolumn*

\LA
\noindent In quo étiam et nos sorte vocáti sumus prædestináti secúndum propósitum ejus qui operátur ómnia secúndum consílium voluntátis suæ: ut simus in laudem glóriæ ejus nos, qui ante sperávimus in Christo; in quo et vos, cum audissétis verbum veritátis, Evangélium salútis vestræ, in quo et credéntes signáti estis Spíritu promissiónis Sancto, qui est pignus hereditátis nostræ, in redemptiónem acquisitiónis, in laudem glóriæ ipsíus.\par
\switchcolumn
\EN
\noindent In whom we also are called by lot, being predestinated according to the purpose of him who worketh all things according to the counsel of his will. That we may be unto the praise of his glory, we who before hoped Christ: In whom you also, after you had heard the word of truth, (the gospel of your salvation;) in whom also believing, you were signed with the holy Spirit of promise, Who is the pledge of our inheritance, unto the redemption of acquisition, unto the praise of his glory.\par
\switchcolumn*

\LA
\begin{Resp}\Rsign Adjútor meus, tibi psallam, quia, Deus, suscéptor meus es: \Star{} Deus meus, misericórdia mea.\end{Resp}
\begin{Resp}\Vsign Lætábor, et exsultábo in te, psallam nómini tuo, Altíssime.\end{Resp}
\begin{Resp}\Rsign Deus meus, misericórdia mea.\end{Resp}
\begin{Resp}\Vsign Glória Patri, et Fílio, \Star{} et Spirítui Sancto.\end{Resp}
\begin{Resp}\Rsign Deus meus, misericórdia mea.\end{Resp}
\switchcolumn
\EN
\begin{Resp}\Rsign Unto thee, O my Strength, will I sing, for God is my defence, \Star{} The God of my mercy.\end{Resp}
\begin{Resp}\Vsign I will be glad and rejoice in thee, I will sing praise to thy Name, O Thou Most High.\end{Resp}
\begin{Resp}\Rsign The God of my mercy.\end{Resp}
\begin{Resp}\Vsign Glory be to the Father, and to the Son, \Star{} and to the Holy Ghost.\end{Resp}
\begin{Resp}\Rsign The God of my mercy.\end{Resp}
\switchcolumn*

\end{paracol}

\Office{Feria V infra Hebdomadam III post Epiphaniam}{Thursday of the Third Week after Epiphany}{Feria}
\addcontentsline{toc}{subsection}{Feria V infra Hebdomadam III post Epiphaniam\enspace\textit{Thursday of the Third Week after Epiphany}}
\begin{paracol}{2}
\LA
\RefLine{Lectiones I–III: de Scriptura occurrente.}
\switchcolumn
\EN
\RefLine{Lessons I–III: from the Scripture of the day.}
\switchcolumn*

\LA
\LessonHead{Lectio I}
\switchcolumn
\EN
\LessonHead{Lesson I}
\switchcolumn*

\LA
\LessonTitle{De Epístola ad Ephésios}
\switchcolumn
\EN
\LessonTitle{Lesson from the letter of St. Paul the Apostle to the Ephesians}
\switchcolumn*

\LA
\Cite{Eph 4:1-6}
\switchcolumn
\EN
\Cite{Eph 4:1-6}
\switchcolumn*

\LA
\noindent Obsecro ítaque vos ego vinctus in Dómino, ut digne ambulétis vocatióne, qua vocáti estis, cum omni humilitáte, et mansuetúdine, cum patiéntia, supportántes ínvicem in caritáte, sollíciti serváre unitátem Spíritus in vínculo pacis. Unum corpus, et unus Spíritus, sicut vocáti estis in una spe vocatiónis vestræ. Unus Dóminus, una fides, unum baptísma. Unus Deus et Pater ómnium, qui est super omnes, et per ómnia, et in ómnibus nobis.\par
\switchcolumn
\EN
\noindent I therefore, a prisoner in the Lord, beseech you that you walk worthy of the vocation in which you are called, With all humility and mildness, with patience, supporting one another in charity. Careful to keep the unity of the Spirit in the bond of peace. One body and one Spirit; as you are called in one hope of your calling. One Lord, one faith, one baptism. One God and Father of all, who is above all, and through all, and in us all.\par
\switchcolumn*

\LA
\begin{Resp}\Rsign Deus, in te sperávi, Dómine, non confúndar in ætérnum: in justítia tua líbera me, \Star{} Et éripe me.\end{Resp}
\begin{Resp}\Vsign Inclína ad me aurem tuam, et salva me.\end{Resp}
\begin{Resp}\Rsign Et éripe me.\end{Resp}
\switchcolumn
\EN
\begin{Resp}\Rsign In thee, O God, do I put my trust; let me never be put to confusion, O Lord deliver me in thy righteousness, \Star{} And cause me to escape.\end{Resp}
\begin{Resp}\Vsign Incline thine ear unto me, deliver me speedily.\end{Resp}
\begin{Resp}\Rsign And cause me to escape.\end{Resp}
\switchcolumn*

\LA
\LessonHead{Lectio II}
\switchcolumn
\EN
\LessonHead{Lesson II}
\switchcolumn*

\LA
\Cite{Eph 4:7-10}
\switchcolumn
\EN
\Cite{Eph 4:7-10}
\switchcolumn*

\LA
\noindent Unicuíque autem nostrum data est grátia secúndum mensúram donatiónis Christi. Propter quod dicit: Ascéndens in altum, captívam duxit captivitátem: dedit dona homínibus. Quod autem ascéndit, quid est, nisi quia et descéndit primum in inferióres partes terræ? Qui descéndit, ipse est et qui ascéndit super omnes cælos, ut impléret ómnia.\par
\switchcolumn
\EN
\noindent But to every one of us is given grace, according to the measure of the giving of Christ. Wherefore he saith: Ascending on high, he led captivity captive; he gave gifts to men. Now that he ascended, what is it, but because he also descended first into the lower parts of the earth? He that descended is the same also that ascended above all the heavens, that he might fill all things.\par
\switchcolumn*

\LA
\begin{Resp}\Rsign Repleátur os meum laude tua, ut hymnum dicam glóriæ tuæ, tota die magnitúdinem tuam: noli me proícere in témpore senectútis: \Star{} Dum defécerit in me virtus mea, ne derelínquas me.\end{Resp}
\begin{Resp}\Vsign Gaudébunt lábia mea cum cantávero tibi.\end{Resp}
\begin{Resp}\Rsign Dum defécerit in me virtus mea, ne derelínquas me.\end{Resp}
\switchcolumn
\EN
\begin{Resp}\Rsign Let my mouth be filled with thy praise, that I may sing of thy glory, all the day long of thy greatness. Cast me not off in the time of old age; \Star{} Forsake me not when my strength faileth.\end{Resp}
\begin{Resp}\Vsign My lips shall be fain when I sing unto thee.\end{Resp}
\begin{Resp}\Rsign Forsake me not when my strength faileth.\end{Resp}
\switchcolumn*

\LA
\LessonHead{Lectio III}
\switchcolumn
\EN
\LessonHead{Lesson III}
\switchcolumn*

\LA
\Cite{Eph 4:11-15}
\switchcolumn
\EN
\Cite{Eph 4:11-15}
\switchcolumn*

\LA
\noindent Et ipse dedit quosdam quidem apóstolos, quosdam autem prophétas, álios vero evangelístas, álios autem pastóres et doctóres, ad consummatiónem sanctórum in opus ministérii, in ædificatiónem córporis Christi: donec occurrámus omnes in unitátem fídei, et agnitiónis Fílii Dei, in virum perféctum, in mensúram ætátis plenitúdinis Christi: ut jam non simus párvuli fluctuántes, et circumferámur omni vento doctrínæ in nequítia hóminum, in astútia ad circumventiónem erróris. Veritátem autem faciéntes in caritáte, crescámus in illo per ómnia, qui est caput Christus:\par
\switchcolumn
\EN
\noindent And he gave some apostles, and some prophets, and other some evangelists, and other some pastors and doctors, For the perfecting of the saints, for the work of the ministry, for the edifying of the body of Christ: Until we all meet into the unity of faith, and of the knowledge of the Son of God, unto a perfect man, unto the measure of the age of the fulness of Christ; That henceforth we be no more children tossed to and fro, and carried about with every wind of doctrine by the wickedness of men, by cunning craftiness, by which they lie in wait to deceive But doing the truth in charity, we may in all things grow up in him who is the head, even Christ:\par
\switchcolumn*

\LA
\begin{Resp}\Rsign Gaudébunt lábia mea cum cantávero tibi: \Star{} Et ánima mea, quam redemísti, Dómine.\end{Resp}
\begin{Resp}\Vsign Sed et lingua mea meditábitur justítiam tuam, tota die laudem tuam.\end{Resp}
\begin{Resp}\Rsign Et ánima mea, quam redemísti, Dómine.\end{Resp}
\begin{Resp}\Vsign Glória Patri, et Fílio, \Star{} et Spirítui Sancto.\end{Resp}
\begin{Resp}\Rsign Et ánima mea, quam redemísti, Dómine.\end{Resp}
\switchcolumn
\EN
\begin{Resp}\Rsign My lips shall be fain when I sing unto thee; \Star{} And my soul, which Thou, O Lord, hast redeemed.\end{Resp}
\begin{Resp}\Vsign My tongue shall also talk of thy righteousness, all the day long of thy praise.\end{Resp}
\begin{Resp}\Rsign And my soul, which Thou, O Lord, hast redeemed.\end{Resp}
\begin{Resp}\Vsign Glory be to the Father, and to the Son, \Star{} and to the Holy Ghost.\end{Resp}
\begin{Resp}\Rsign And my soul, which Thou, O Lord, hast redeemed.\end{Resp}
\switchcolumn*

\end{paracol}

\Office{Feria VI infra Hebdomadam III post Epiphaniam}{Friday of the Third Week after Epiphany}{Feria}
\addcontentsline{toc}{subsection}{Feria VI infra Hebdomadam III post Epiphaniam\enspace\textit{Friday of the Third Week after Epiphany}}
\begin{paracol}{2}
\LA
\LessonHead{Lectio I}
\switchcolumn
\EN
\LessonHead{Lesson I}
\switchcolumn*

\LA
\LessonTitle{De Epístola ad Ephésios}
\switchcolumn
\EN
\LessonTitle{Lesson from the letter of St. Paul the Apostle to the Ephesians}
\switchcolumn*

\LA
\Cite{Eph 5:1-4}
\switchcolumn
\EN
\Cite{Eph 5:1-4}
\switchcolumn*

\LA
\noindent Estóte ergo imitatóres Dei, sicut fílii caríssimi, et ambuláte in dilectióne, sicut et Christus diléxit nos, et trádidit semetípsum pro nobis, oblatiónem et hóstiam Deo in odórem suavitátis. Fornicátio autem, et omnis immundítia, aut avarítia, nec nominétur in vobis, sicut decet sanctos: aut turpitúdo, aut stultilóquium, aut scurrílitas, quæ ad rem non pértinet: sed magis gratiárum áctio.\par
\switchcolumn
\EN
\noindent Be ye therefore followers of God, as most dear children; And walk in love, as Christ also hath loved us, and hath delivered himself for us, an oblation and a sacrifice to God for an odour of sweetness. But fornication, and all uncleanness, or covetousness, let it not so much as be named among you, as becometh saints: Or obscenity, or foolish talking, or scurrility, which is to no purpose; but rather giving of thanks.\par
\switchcolumn*

\LA
\begin{Resp}\Rsign Confitébor tibi, Dómine Deus, in toto corde meo, et honorificábo nomen tuum in ætérnum: \Star{} Quia misericórdia tua, Dómine, magna est super me.\end{Resp}
\begin{Resp}\Vsign Deus meus es tu, et confitébor tibi: Deus meus es tu, et exaltábo te.\end{Resp}
\begin{Resp}\Rsign Quia misericórdia tua, Dómine, magna est super me.\end{Resp}
\switchcolumn
\EN
\begin{Resp}\Rsign I will praise thee, O Lord my God, with all my heart, and I will glorify thy Name for evermore. \Star{} For great is thy mercy, O Lord, toward me.\end{Resp}
\begin{Resp}\Vsign Thou art my God, and I will praise thee; Thou art my God, and I will exalt thee.\end{Resp}
\begin{Resp}\Rsign For great is thy mercy, O Lord, toward me.\end{Resp}
\switchcolumn*

\LA
\LessonHead{Lectio II}
\switchcolumn
\EN
\LessonHead{Lesson II}
\switchcolumn*

\LA
\Cite{Eph 5:5-8}
\switchcolumn
\EN
\Cite{Eph 5:5-8}
\switchcolumn*

\LA
\noindent Hoc enim scitóte intellegéntes: quod omnis fornicátor, aut immúndus, aut avárus, quod est idolórum sérvitus, non habet hereditátem in regno Christi et Dei. Nemo vos sedúcat inánibus verbis: propter hæc enim venit ira Dei in fílios diffidéntiæ. Nolíte ergo éffici partícipes eórum. Erátis enim aliquándo ténebræ: nunc autem lux in Dómino. Ut fílii lucis ambuláte:\par
\switchcolumn
\EN
\noindent For know you this and understand, that no fornicator, or unclean, or covetous person (which is a serving of idols), hath inheritance in the kingdom of Christ and of God. Let no man deceive you with vain words. For because of these things cometh the anger of God upon the children of unbelief. Be ye not therefore partakers with them. For you were heretofore darkness, but now light in the Lord. Walk then as children of the light.\par
\switchcolumn*

\LA
\begin{Resp}\Rsign Misericórdia tua, Dómine, magna est super me: \Star{} Et liberásti ánimam meam ex inférno inferióri.\end{Resp}
\begin{Resp}\Vsign In die tribulatiónis meæ clamávi ad te, quia exaudísti me.\end{Resp}
\begin{Resp}\Rsign Et liberásti ánimam meam ex inférno inferióri.\end{Resp}
\switchcolumn
\EN
\begin{Resp}\Rsign Great, O Lord, is thy mercy toward me. \Star{} And Thou hast delivered my soul from the lowest hell.\end{Resp}
\begin{Resp}\Vsign In the day of my trouble I called upon thee, for Thou hast heard me.\end{Resp}
\begin{Resp}\Rsign And Thou hast delivered my soul from the lowest hell.\end{Resp}
\switchcolumn*

\LA
\LessonHead{Lectio III}
\switchcolumn
\EN
\LessonHead{Lesson III}
\switchcolumn*

\LA
\Cite{Eph 5:9-14}
\switchcolumn
\EN
\Cite{Eph 5:9-14}
\switchcolumn*

\LA
\noindent fructus enim lucis est in omni bonitáte, et justítia, et veritáte: probántes quid sit beneplácitum Deo: et nolíte communicáre opéribus infructuósis tenebrárum, magis autem redargúite. Quæ enim in occúlto fiunt ab ipsis, turpe est et dícere. Omnia autem, quæ arguúntur, a lúmine manifestántur: omne enim, quod manifestátur, lumen est. Propter quod dicit: Surge qui dormis, et exsúrge a mórtuis, et illuminábit te Christus.\par
\switchcolumn
\EN
\noindent For the fruit of the light is in all goodness, and justice, and truth; Proving what is well pleasing to God: And have no fellowship with the unfruitful works of darkness, but rather reprove them. For the things that are done by them in secret, it is a shame even to speak of. But all things that are reproved, are made manifest by the light; for all that is made manifest is light. Wherefore he saith: Rise thou that sleepest, and arise from the dead: and Christ shall enlighten thee.\par
\switchcolumn*

\LA
\begin{Resp}\Rsign Factus est mihi Dóminus in refúgium: \Star{} Et Deus meus in auxílium spei meæ.\end{Resp}
\begin{Resp}\Vsign Erípuit me de inimícis meis fortíssimis, et factus est Dóminus protéctor meus.\end{Resp}
\begin{Resp}\Rsign Et Deus meus in auxílium spei meæ.\end{Resp}
\begin{Resp}\Vsign Glória Patri, et Fílio, \Star{} et Spirítui Sancto.\end{Resp}
\begin{Resp}\Rsign Et Deus meus in auxílium spei meæ.\end{Resp}
\switchcolumn
\EN
\begin{Resp}\Rsign The Lord is my refuge. \Star{} And my God is the stay of my trust.\end{Resp}
\begin{Resp}\Vsign He delivered me from the strongest of mine enemies, and the Lord was my stay.\end{Resp}
\begin{Resp}\Rsign And my God is the stay of my trust.\end{Resp}
\begin{Resp}\Vsign Glory be to the Father, and to the Son, \Star{} and to the Holy Ghost.\end{Resp}
\begin{Resp}\Rsign And my God is the stay of my trust.\end{Resp}
\switchcolumn*

\end{paracol}

\Office{Sabbato infra Hebdomadam III post Epiphaniam}{Saturday of the Third Week after Epiphany}{Feria}
\addcontentsline{toc}{subsection}{Sabbato infra Hebdomadam III post Epiphaniam\enspace\textit{Saturday of the Third Week after Epiphany}}
\begin{paracol}{2}
\LA
\LessonHead{Lectio I}
\switchcolumn
\EN
\LessonHead{Lesson I}
\switchcolumn*

\LA
\LessonTitle{De Epístola ad Ephésios}
\switchcolumn
\EN
\LessonTitle{Lesson from the letter of St. Paul the Apostle to the Ephesians}
\switchcolumn*

\LA
\Cite{Eph 6:1-4}
\switchcolumn
\EN
\Cite{Eph 6:1-4}
\switchcolumn*

\LA
\noindent Fílii, obedíte paréntibus vestris in Dómino: hoc enim justum est. Honóra patrem tuum, et matrem tuam, quod est mandátum primum in promissióne: ut bene sit tibi, et sis longǽvus super terram. Et vos patres, nolíte ad iracúndiam provocáre fílios vestros: sed educáte illos in disciplína et correptióne Dómini.\par
\switchcolumn
\EN
\noindent Children, obey your parents in the Lord, for this is just. Honour thy father and thy mother, which is the first commandment with a promise: That it may be well with thee, and thou mayest be long lived upon earth. And you, fathers, provoke not your children to anger; but bring them up in the discipline and correction of the Lord.\par
\switchcolumn*

\LA
\begin{Resp}\Rsign Misericórdiam et judícium cantábo tibi, Dómine: \Star{} Psallam et intéllegam in via immaculáta, quando vénies ad me.\end{Resp}
\begin{Resp}\Vsign Perambulábam in innocéntia cordis mei, in médio domus meæ.\end{Resp}
\begin{Resp}\Rsign Psallam et intéllegam in via immaculáta, quando vénies ad me.\end{Resp}
\switchcolumn
\EN
\begin{Resp}\Rsign Mercy and judgment I will sing to thee, O Lord: \Star{} I will sing, and I will understand in the unspotted way, when thou shalt come to me.\end{Resp}
\begin{Resp}\Vsign I walked in the innocence of my heart, in the midst of my house.\end{Resp}
\begin{Resp}\Rsign Mercy and judgment I will sing to thee, O Lord\end{Resp}
\switchcolumn*

\LA
\LessonHead{Lectio II}
\switchcolumn
\EN
\LessonHead{Lesson II}
\switchcolumn*

\LA
\Cite{Eph 6:5-9}
\switchcolumn
\EN
\Cite{Eph 6:5-9}
\switchcolumn*

\LA
\noindent Servi, obedíte dóminis carnálibus cum timóre et tremóre, in simplicitáte cordis vestri, sicut Christo: non ad óculum serviéntes, quasi homínibus placéntes, sed ut servi Christi, faciéntes voluntátem Dei ex ánimo, cum bona voluntáte serviéntes, sicut Dómino, et non homínibus: sciéntes quóniam unusquísque quodcúmque fécerit bonum, hoc recípiet a Dómino, sive servus, sive liber. Et vos dómini, éadem fácite illis, remitténtes minas: sciéntes quia et illórum et vester Dóminus est in cælis: et personárum accéptio non est apud eum.\par
\switchcolumn
\EN
\noindent Servants, be obedient to them that are your lords according to the flesh, with fear and trembling, in the simplicity of your heart, as to Christ: Not serving to the eye, as it were pleasing men, but, as the servants of Christ doing the will of God from the heart, With a good will serving, as to the Lord, and not to men. Knowing that whatsoever good thing any man shall do, the same shall he receive from the Lord, whether he be bond, or free. And you, masters, do the same things to them, forbearing threatenings, knowing that the Lord both of them and you is in heaven; and there is no respect of persons with him.\par
\switchcolumn*

\LA
\begin{Resp}\Rsign Dómine, exáudi oratiónem meam, et clamor meus ad te pervéniat: \Star{} Quia non spernis, Deus, preces páuperum.\end{Resp}
\begin{Resp}\Vsign Fiant aures tuæ intendéntes in oratiónem servi tui.\end{Resp}
\begin{Resp}\Rsign Quia non spernis, Deus, preces páuperum.\end{Resp}
\switchcolumn
\EN
\begin{Resp}\Rsign Hear my prayer, O Lord, and let my crying come unto thee. \Star{} For thou wilt not despise, O God, the prayer of the destitute\end{Resp}
\begin{Resp}\Vsign O let thine ears consider well the voice of my complaint\end{Resp}
\begin{Resp}\Rsign For thou wilt not despise, O God, the prayer of the destitute.\end{Resp}
\switchcolumn*

\LA
\LessonHead{Lectio III}
\switchcolumn
\EN
\LessonHead{Lesson III}
\switchcolumn*

\LA
\Cite{Eph 6:10-13}
\switchcolumn
\EN
\Cite{Eph 6:10-13}
\switchcolumn*

\LA
\noindent De cétero, fratres, confortámini in Dómino, et in poténtia virtútis ejus. Indúite vos armatúram Dei, ut possítis stare advérsus insídias diáboli: quóniam non est nobis colluctátio advérsus carnem et sánguinem, sed advérsus príncipes, et potestátes, advérsus mundi rectóres tenebrárum harum, contra spirituália nequítiæ, in cæléstibus. Proptérea accípite armatúram Dei, ut possítis resístere in die malo, et in ómnibus perfécti stare.\par
\switchcolumn
\EN
\noindent Finally, brethren, be strengthened in the Lord, and in the might of his power. Put you on the armour of God, that you may be able to stand against the deceits of the devil. For our wrestling is not against flesh and blood; but against principalities and power, against the rulers of the world of this darkness, against the spirits of wickedness in the high places. Therefore take unto you the armour of God, that you may be able to resist in the evil day, and to stand in all things perfect.\par
\switchcolumn*

\LA
\begin{Resp}\Rsign Velóciter exáudi me, Deus, \Star{} Tu autem idem ipse es, et anni tui non defícient.\end{Resp}
\begin{Resp}\Vsign Dies mei sicut umbra declinavérunt, et ego sicut fœnum árui.\end{Resp}
\begin{Resp}\Rsign Quia defecérunt sicut fumus dies mei.\end{Resp}
\begin{Resp}\Vsign Glória Patri, et Fílio, \Star{} et Spirítui Sancto.\end{Resp}
\begin{Resp}\Rsign Tu autem idem ipse es, et anni tui non defícient.\end{Resp}
\switchcolumn
\EN
\begin{Resp}\Rsign Haste thee to hear me, O Lord. \Star{} But thou art the same, and thy years shall not fail.\end{Resp}
\begin{Resp}\Vsign My days have declined like a shadow, and I am withered like grass.\end{Resp}
\begin{Resp}\Rsign For my days are consumed away like smoke.\end{Resp}
\begin{Resp}\Vsign Glory be to the Father, and to the Son, \Star{} and to the Holy Ghost.\end{Resp}
\begin{Resp}\Rsign But thou art the same, and thy years shall not fail.\end{Resp}
\switchcolumn*

\end{paracol}

\Office{Dominica IV Post Epiphaniam}{Fourth Sunday after Epiphany}{Semiduplex}
\addcontentsline{toc}{subsection}{Dominica IV Post Epiphaniam\enspace\textit{Fourth Sunday after Epiphany}}
\begin{paracol}{2}
\LA
\LessonHead{Lectio I}
\switchcolumn
\EN
\LessonHead{Lesson I}
\switchcolumn*

\LA
\LessonTitle{Incipit Epístola beáti Pauli Apóstoli ad Philippénses}
\switchcolumn
\EN
\LessonTitle{Lesson from the letter of St. Paul the Apostle to the Philippians}
\switchcolumn*

\LA
\Cite{Phil 1:1-7}
\switchcolumn
\EN
\Cite{Phil 1:1-7}
\switchcolumn*

\LA
\noindent Paulus et Timótheus, servi Jesu Christi, ómnibus sanctis in Christo Jesu, qui sunt Philíppis, cum epíscopis et diacónibus. Grátia vobis, et pax a Deo Patre nostro, et Dómino Jesu Christo. Grátias ago Deo meo in omni memória vestri, semper in cunctis oratiónibus meis pro ómnibus vobis, cum gáudio deprecatiónem fáciens, super communicatióne vestra in Evangélio Christi a prima die usque nunc. Confídens hoc ipsum, quia qui cœpit in vobis opus bonum, perfíciet usque in diem Christi Jesu: sicut est mihi justum hoc sentíre pro ómnibus vobis: eo quod hábeam vos in corde, et in vínculis meis, et in defensióne, et confirmatióne Evangélii, sócios gáudii mei omnes vos esse.\par
\switchcolumn
\EN
\noindent Paul and Timothy, the servants of Jesus Christ; to all the saints in Christ Jesus, who are at Philippi, with the bishops and deacons. Grace be unto you, and peace from God our Father, and from the Lord Jesus Christ. I give thanks to my God in every remembrance of you, Always in all my prayers making supplication for you all, with joy; For your communication in the gospel of Christ from the first day until now. Being confident of this very thing, that he, who hath begun a good work in you, will perfect it unto the day of Christ Jesus. As it is meet for me to think this for you all, for that I have you in my heart; and that in my bands, and in the defence and confirmation of the gospel, you all are partakers of my joy.\par
\switchcolumn*

\LA
\begin{Resp}\Rsign Dómine, ne in ira tua árguas me, neque in furóre tuo corrípias me: \Star{} Miserére mei, Dómine, quóniam infírmus sum.\end{Resp}
\begin{Resp}\Vsign Timor et tremor venérunt super me, et contexérunt me ténebræ.\end{Resp}
\begin{Resp}\Rsign Miserére mei, Dómine, quóniam infírmus sum.\end{Resp}
\switchcolumn
\EN
\begin{Resp}\Rsign O Lord, rebuke me not in thine anger, neither chasten me in thine hot displeasure. \Star{} Have mercy upon me, O Lord, for I am weak.\end{Resp}
\begin{Resp}\Vsign Fearfulness and trembling are come upon me, and darkness hath overwhelmed me.\end{Resp}
\begin{Resp}\Rsign Have mercy upon me, O Lord, for I am weak.\end{Resp}
\switchcolumn*

\LA
\LessonHead{Lectio II}
\switchcolumn
\EN
\LessonHead{Lesson II}
\switchcolumn*

\LA
\Cite{Phil 1:8-14}
\switchcolumn
\EN
\Cite{Phil 1:8-14}
\switchcolumn*

\LA
\noindent Testis enim mihi est Deus, quómodo cúpiam omnes vos in viscéribus Jesu Christi. Et hoc oro, ut cáritas vestra magis ac magis abúndet in sciéntia, et in omni sensu: ut probétis potióra, ut sitis sincéri, et sine offénsa in diem Christi, repléti fructu justítiæ per Jesum Christum, in glóriam et laudem Dei. Scire autem vos volo fratres, quia quæ circa me sunt, magis ad proféctum venérunt Evangélii: ita ut víncula mea manifésta fíerent in Christo in omni prætório, et in céteris ómnibus, et plures e frátribus in Dómino confidéntes vínculis meis, abundántius audérent sine timóre verbum Dei loqui.\par
\switchcolumn
\EN
\noindent For God is my witness, how I long after you all in the bowels of Jesus Christ. And this I pray, that your charity may more and more abound in knowledge, and in all understanding: That you may approve the better things, that you may be sincere and without offence unto the day of Christ, Filled with the fruit of justice, through Jesus Christ, unto the glory and praise of God. Now, brethren, I desire you should know, that the things which have happened to me, have fallen out rather to the furtherance of the gospel: So that my bands are made manifest in Christ, in all the court, and in all other places; And many of the brethren in the Lord, growing confident by my bands, are much more bold to speak the word of God without fear.\par
\switchcolumn*

\LA
\begin{Resp}\Rsign Deus, qui sedes super thronum, et júdicas æquitátem, esto refúgium páuperum in tribulatióne: \Star{} Quia tu solus labórem et dolórem consíderas.\end{Resp}
\begin{Resp}\Vsign Tibi enim derelíctus est pauper, pupíllo tu eris adjútor.\end{Resp}
\begin{Resp}\Rsign Quia tu solus labórem et dolórem consíderas.\end{Resp}
\switchcolumn
\EN
\begin{Resp}\Rsign O God, Which satest in the throne judging right, be Thou a refuge for the poor, a refuge in times of trouble. \Star{} For Thou alone beholdest mischief and spite.\end{Resp}
\begin{Resp}\Vsign The poor leaveth himself unto thee; Thou wilt be the helper of the fatherless.\end{Resp}
\begin{Resp}\Rsign For Thou alone beholdest mischief and spite.\end{Resp}
\switchcolumn*

\LA
\LessonHead{Lectio III}
\switchcolumn
\EN
\LessonHead{Lesson III}
\switchcolumn*

\LA
\Cite{Phil 1:15-18}
\switchcolumn
\EN
\Cite{Phil 1:15-18}
\switchcolumn*

\LA
\noindent Quidam quidem et propter invídiam et contentiónem: quidam autem et propter bonam voluntátem Christum prǽdicant: quidam ex caritáte, sciéntes quóniam in defensiónem Evangélii pósitus sum. Quidam autem ex contentióne Christum annúntiant non sincére, existimántes pressúram se suscitáre vínculis meis. Quid enim? Dum omni modo sive per occasiónem, sive per veritátem, Christus annuntiétur: et in hoc gáudeo, sed et gaudébo.\par
\switchcolumn
\EN
\noindent Some indeed, even out of envy and contention; but some also for good will preach Christ. Some out of charity, knowing that I am set for the defence of the gospel. And some out of contention preach Christ not sincerely: supposing that they raise affliction to my bands. But what then? So that by all means, whether by occasion, or by truth, Christ be preached: in this also I rejoice, yea, and will rejoice.\par
\switchcolumn*

\LA
\begin{Resp}\Rsign A dextris est mihi Dóminus, ne commóvear: \Star{} Propter hoc dilatátum est cor meum, et exsultávit lingua mea.\end{Resp}
\begin{Resp}\Vsign Dóminus pars hereditátis meæ, et cálicis mei.\end{Resp}
\begin{Resp}\Rsign Propter hoc dilatátum est cor meum, et exsultávit lingua mea.\end{Resp}
\begin{Resp}\Vsign Glória Patri, et Fílio, \Star{} et Spirítui Sancto.\end{Resp}
\begin{Resp}\Rsign Propter hoc dilatátum est cor meum, et exsultávit lingua mea.\end{Resp}
\switchcolumn
\EN
\begin{Resp}\Rsign The Lord is at my right hand, I shall never be moved. \Star{} Therefore my heart is glad, and my tongue rejoiceth.\end{Resp}
\begin{Resp}\Vsign The Lord is the portion of mine inheritance, and of my cup.\end{Resp}
\begin{Resp}\Rsign Therefore my heart is glad, and my tongue rejoiceth.\end{Resp}
\begin{Resp}\Vsign Glory be to the Father, and to the Son, \Star{} and to the Holy Ghost.\end{Resp}
\begin{Resp}\Rsign Therefore my heart is glad, and my tongue rejoiceth.\end{Resp}
\switchcolumn*

\LA
\LessonHead{Lectio IV}
\switchcolumn
\EN
\LessonHead{Lesson IV}
\switchcolumn*

\LA
\LessonTitle{Ex libro Morálium sancti Gregórii Papæ}
\switchcolumn
\EN
\LessonTitle{From the Book of Moral (Reflections on Job) written by Pope St. Gregory (the Great.)}
\switchcolumn*

\LA
\Source{Lib. 4. Cap. 30.}
\switchcolumn
\EN
\Source{iv, 30}
\switchcolumn*

\LA
\noindent Replémus refectiónibus corpus, ne extenuátum defíciat; extenuámus abstinéntia, ne nos replétum premat: vegetámus hoc mótibus, ne situ immobilitátis intéreat; sed cítius hoc collocándo sístimus, ne ipsa sua vegetatióne succúmbat: adjuméntis hoc véstium tégimus, ne frigus intérimat; et quæsíta adjuménta projícimus, ne calor exúrat. Tot ígitur diversitátibus occurréntes, quid ágimus, nisi corruptibilitáti sérvimus, ut saltem multiplícitas impénsi obséquii corpus sustíneat, quod anxíetas infírmæ mutabilitátis gravat?\par
\switchcolumn
\EN
\noindent We refresh the body lest it should grow too weak and fail us; we chasten it by abstinence, lest it should wax gross, and become lord over us; we strengthen it with exercise, lest it perish by the not using; and straightway we give it rest, lest it faint through weariness; we succour it with raiment, lest the cold should blight it; and we strip it of the raiment wherewith we have clothed it, lest the heat should afflict it. In all these so many offices what do we but serve the corruptible? Upon what is all this care spent but upon that wherover hangeth the doom of weakness and change?\par
\switchcolumn*

\LA
\begin{Resp}\Rsign Notas mihi fecísti, Dómine, vias vitæ: \Star{} Adimplébis me lætítia cum vultu tuo: delectatiónes in déxtera tua usque in finem.\end{Resp}
\begin{Resp}\Vsign Tu es qui restítues hereditátem meam mihi.\end{Resp}
\begin{Resp}\Rsign Adimplébis me lætítia cum vultu tuo: delectatiónes in déxtera tua usque in finem.\end{Resp}
\switchcolumn
\EN
\begin{Resp}\Rsign O Lord, Thou hast shown me the path of life. \Star{} Thou shalt fill me with joy in thy presence, at thy right hand there are pleasures for evermore.\end{Resp}
\begin{Resp}\Vsign Thou art He That shalt restore mine inheritance unto me.\end{Resp}
\begin{Resp}\Rsign Thou shalt fill me with joy in thy presence, at thy right hand there are pleasures for evermore.\end{Resp}
\switchcolumn*

\LA
\LessonHead{Lectio V}
\switchcolumn
\EN
\LessonHead{Lesson V}
\switchcolumn*

\LA
\noindent Unde bene per Paulum dícitur: Vanitáti enim subjécta est creatúra non volens, sed propter eum qui subjécit eam in spe: quia et ipsa creatúra liberábitur a servitúte corruptiónis, in libertátem glóriæ filiórum Dei. Vanitáti quippe creatúra non volens súbditur: quia homo, qui ingénitæ constántiæ statum volens deséruit, pressus justæ mortalitátis póndere, nolens mutabilitátis suæ corruptióni servit. Sed creatúra hæc tunc a servitúte corruptiónis erípitur, cum ad filiórum Dei glóriam incorrúpta resurgéndo sublevátur.\par
\switchcolumn
\EN
\noindent Therefore saith Paul tells: For the creature was made subject to vanity, not willingly, but by reason of Him Who hath subjected the same in hope because the creature itself also shall be delivered from the bondage of corruption into the glorious liberty of the children of God. (Rom. viii. 20.) The creature was made subject to vanity, not willingly for when man had of his own free will abdicated his state of unchangeable blessedness, the just sentence of death was passed upon him, and whether he willed or not, he became subject to the state of change and corruption. But the creature itself also shall be delivered from the bondage of corruption when it shall rise again incorruptible and be made partaker of the glory of the children of God.\par
\switchcolumn*

\LA
\begin{Resp}\Rsign Díligam te, Dómine, virtus mea: Dóminus firmaméntum meum, \Star{} Et refúgium meum.\end{Resp}
\begin{Resp}\Vsign Liberátor meus, Deus meus, adjútor meus.\end{Resp}
\begin{Resp}\Rsign Et refúgium meum.\end{Resp}
\switchcolumn
\EN
\begin{Resp}\Rsign I will love thee, O Lord, my strength; the Lord is my rock \Star{} And my fortress.\end{Resp}
\begin{Resp}\Vsign My Deliverer, my God, mine Helper.\end{Resp}
\begin{Resp}\Rsign And my fortress.\end{Resp}
\switchcolumn*

\LA
\LessonHead{Lectio VI}
\switchcolumn
\EN
\LessonHead{Lesson VI}
\switchcolumn*

\LA
\noindent Hic ítaque elécti moléstia vincti sunt, quia adhuc corruptiónis suæ pœna deprimúntur: sed cum corruptíbili carne exúimur, quasi ab his, quibus nunc astríngimur, moléstiæ vínculis relaxámur. Præsentári namque jam Deo cúpimus, sed adhuc mortális córporis obligatióne præpedímur. Jure ergo vincti dícimur, quia adhuc incéssum nostri desidérii ad Deum líberum non habémus. Unde bene Paulus, ætérna desíderans, sed tamen adhuc corruptiónis suæ sárcinam portans, vinctus clamat: Cúpio dissólvi, et esse cum Christo. Dissólvi enim non quǽreret, nisi se proculdúbio vinctum vidéret.\par
\switchcolumn
\EN
\noindent Where, then, the elect are still subject to sorrow, being yet bound by the sentence of corruption; but when we shall have put off this corruptible we shall be loosed from that sentence, and shall sorrow no more. For though we earnestly desire to appear before God, we are still hindered by the burden of this dying body. Rightly then are we called prisoners, since we are not free to go whither we will, that is to say, to God; and rightly did the prisoner Paul, yearning after the things which are eternal, and still weighed down with the burden of this corruptible, rightly did he cry out I have a desire to depart and to be with Christ. (Phil. i. 23.) He would not have felt this keenness if he had not felt himself bound down.\par
\switchcolumn*

\LA
\begin{Resp}\Rsign Dómini est terra, et plenitúdo ejus: \Star{} Orbis terrárum, et univérsi qui hábitant in eo.\end{Resp}
\begin{Resp}\Vsign Ipse super mária fundávit eam, et super flúmina præparávit illam.\end{Resp}
\begin{Resp}\Rsign Orbis terrárum, et univérsi qui hábitant in eo.\end{Resp}
\begin{Resp}\Vsign Glória Patri, et Fílio, \Star{} et Spirítui Sancto.\end{Resp}
\begin{Resp}\Rsign Orbis terrárum, et univérsi qui hábitant in eo.\end{Resp}
\switchcolumn
\EN
\begin{Resp}\Rsign The earth is the Lord's, and the fulness thereof \Star{} The world, and they that dwell therein.\end{Resp}
\begin{Resp}\Vsign For He hath founded it upon the seas, and established it upon the floods.\end{Resp}
\begin{Resp}\Rsign The world, and they that dwell therein.\end{Resp}
\begin{Resp}\Vsign Glory be to the Father, and to the Son, \Star{} and to the Holy Ghost.\end{Resp}
\begin{Resp}\Rsign The world, and they that dwell therein.\end{Resp}
\switchcolumn*

\LA
\LessonHead{Lectio VII}
\switchcolumn
\EN
\LessonHead{Lesson VII}
\switchcolumn*

\LA
\LessonTitle{Léctio sancti Evangélii secúndum Matthǽum}
\switchcolumn
\EN
\LessonTitle{From the Holy Gospel according to Matthew}
\switchcolumn*

\LA
\Cite{Matt 8:23-27}
\switchcolumn
\EN
\Cite{Matt 8:23-27}
\switchcolumn*

\LA
\noindent In illo témpore: Ascendénte Jesu in navículam, secúti sunt eum discípuli ejus: et ecce motus magnus factus est in mari, ita ut navícula operirétur flúctibus: ipse vero dormiébat. Et réliqua.\par
\switchcolumn
\EN
\noindent At that time, when Jesus entered into the boat, his disciples followed him: And behold a great tempest arose in the sea, so that the boat was covered with waves, but he was asleep. And so on.\par
\switchcolumn*

\LA
\LessonTitle{Homilía sancti Hierónymi Presbýteri}
\switchcolumn
\EN
\LessonTitle{Homily by St. Jerome, Priest (at Bethlehem.)}
\switchcolumn*

\LA
\Source{Liber 1 Comment. in cap. 8 Matth.}
\switchcolumn
\EN
\Source{Bk. i Comm. on Matth. viii}
\switchcolumn*

\LA
\noindent Quintum signum fecit, quando ascéndens navem de Caphárnaum, ventis imperávit et mari. Sextum, quando in regióne Gerasenórum dedit potestátem dæmónibus in porcos. Séptimum, quando ingrédiens civitátem suam, paralýticum secúndum curávit in léctulo. Primus enim paralýticus est puer centuriónis.\par
\switchcolumn
\EN
\noindent The fifth sign that He did was when He took ship at Capernaum, and commanded the winds and the sea the sixth, when, in the country of the Gergesenes, He suffered the devils to enter into the swine the seventh, when, as He came into His own city, He cured the man sick of the palsy lying on a bed. The first man sick of the palsy that He cured was the centurion's servant.\par
\switchcolumn*

\LA
\begin{Resp}\Rsign Ad te, Dómine, levávi ánimam meam: \Star{} Deus meus, in te confído, non erubéscam.\end{Resp}
\begin{Resp}\Vsign Custódi ánimam meam, et éripe me.\end{Resp}
\begin{Resp}\Rsign Deus meus, in te confído, non erubéscam.\end{Resp}
\switchcolumn
\EN
\begin{Resp}\Rsign Unto thee, O Lord, do I lift up my soul. \Star{} O my God, I trust in thee, let me not be ashamed.\end{Resp}
\begin{Resp}\Vsign O keep my soul and deliver me.\end{Resp}
\begin{Resp}\Rsign O my God, I trust in thee, let me not be ashamed.\end{Resp}
\switchcolumn*

\LA
\LessonHead{Lectio VIII}
\switchcolumn
\EN
\LessonHead{Lesson VIII}
\switchcolumn*

\LA
\noindent Ipse vero dormiébat: et accessérunt ad eum, et suscitavérunt eum, dicéntes: Dómine, salva nos. Hujus signi typum in Jona légimus, quando céteris periclitántibus, ipse secúrus est, et dormit, et suscitátur; et império ac sacraménto passiónis suæ líberat suscitántes. Tunc surgens imperávit ventis et mari. Ex hoc loco intellígimus, quod omnes creatúræ séntiant Creatórem. Quas enim increpávit, et quibus imperávit, sentiunt imperántem: non erróre hæreticórum qui ómnia putant animántia, sed majestáte Conditóris, quæ apud nos insensibília, illi sensibília sunt.\par
\switchcolumn
\EN
\noindent But He was asleep; and His disciples came to Him, and awoke Him, saying Lord, save us. There is a type of this in the history of Jonah, who, when the storm arose, was lying fast asleep, and whom the sailors woke to help them; who also saved the sailors by commanding them to throw him into the sea, the said casting of him into the sea, being, as we know, a figure of Christ's Passion. Then He arose and rebuked the winds and the sea. From these words we understand that all things, which have been made, are sentient to their Maker. All things which He rebuketh or commandeth, hear His voice. This is not the error of the heretics who will have it that everything is quick, but part of the majesty of the Creator, Who maketh to feel Him things which we cannot make to feel us.\par
\switchcolumn*

\LA
\begin{Resp}\Rsign Duo Séraphim clamábant alter ad álterum: \Star{} Sanctus, sanctus, sanctus Dóminus Deus Sábaoth: * Plena est omnis terra glória ejus.\end{Resp}
\begin{Resp}\Vsign Tres sunt qui testimónium dant in cælo: Pater, Verbum, et Spíritus Sanctus: et hi tres unum sunt.\end{Resp}
\begin{Resp}\Rsign Sanctus, sanctus, sanctus Dóminus Deus Sábaoth.\end{Resp}
\begin{Resp}\Vsign Glória Patri, et Fílio, \Star{} et Spirítui Sancto.\end{Resp}
\begin{Resp}\Rsign Plena est omnis terra glória ejus.\end{Resp}
\switchcolumn
\EN
\begin{Resp}\Rsign One Seraph cried unto another: \Star{} Holy, Holy, Holy is the Lord God of hosts; the whole earth is full of His glory.\end{Resp}
\begin{Resp}\Vsign There are Three That bear record in heaven: the Father, the Word, and the Holy Ghost, and these Three are One.\end{Resp}
\begin{Resp}\Rsign Holy, Holy, Holy is the Lord God of hosts.\end{Resp}
\begin{Resp}\Vsign Glory be to the Father, and to the Son, \Star{} and to the Holy Ghost.\end{Resp}
\begin{Resp}\Rsign The whole earth is full of His glory.\end{Resp}
\switchcolumn*

\LA
\LessonHead{Lectio IX}
\switchcolumn
\EN
\LessonHead{Lesson IX}
\switchcolumn*

\LA
\noindent Porro hómines miráti sunt, dicéntes: Qualis est hic, quia venti et mare obédiunt ei? Non discípuli, sed nautæ, et céteri, qui in navi erant, mirabántur. Sin autem quis contentióse volúerit, eos, qui mirabántur, fuísse discípulos: respondébimus, recte hómines appellátos, qui necdum nóverant poténtiam Salvatóris.\par
\switchcolumn
\EN
\noindent But the men marvelled, saying What manner of man is this, that even the winds and the sea obey Him? It was not His disciples that marvelled, but the sailors, and the others that were in the ship. If however, any one willeth to withstand this our interpretation and to maintain that it was the disciples who marvelled, we are ready to answer them that they who knew not before the power of the Saviour deserve to be stripped of the title of disciples, and to be called simply the men.\par
\switchcolumn*

\LA
\TeDeum{Te Deum}
\switchcolumn
\EN
\TeDeum{Te Deum}
\switchcolumn*

\end{paracol}

\Office{Feria Secunda infra Hebdomadam IV post Epiphaniam}{Monday of the Fourth Week after Epiphany}{Feria}
\addcontentsline{toc}{subsection}{Feria Secunda infra Hebdomadam IV post Epiphaniam\enspace\textit{Monday of the Fourth Week after Epiphany}}
\begin{paracol}{2}
\LA
\LessonHead{Lectio I}
\switchcolumn
\EN
\LessonHead{Lesson I}
\switchcolumn*

\LA
\LessonTitle{De Epístola ad Philippénses}
\switchcolumn
\EN
\LessonTitle{Lesson from the letter of St. Paul the Apostle to the Philippians}
\switchcolumn*

\LA
\Cite{Phil 4:1-3}
\switchcolumn
\EN
\Cite{Phil 4:1-3}
\switchcolumn*

\LA
\noindent Itaque fratres mei caríssimi, et desideratíssimi, gáudium meum, et coróna mea: sic state in Dómino, caríssimi. Evódiam rogo, et Sýntychen déprecor, idípsum sápere in Dómino. Etiam rogo et te, germáne compar, ádjuva illas, quæ mecum laboravérunt in Evangélio cum Cleménte, et céteris adjutóribus meis, quorum nómina sunt in libro vitæ.\par
\switchcolumn
\EN
\noindent Therefore, my dearly beloved brethren, and most desired, my joy and my crown; so stand fast in the Lord, my dearly beloved. I beg of Evodia, and I beseech Syntyche, to be of one mind in the Lord. And I entreat thee also, my sincere companion, help those women who have laboured with me in the gospel, with Clement and the rest of my fellow labourers, whose names are in the book of life.\par
\switchcolumn*

\LA
\begin{Resp}\Rsign Quam magna multitúdo dulcédinis tuæ, Dómine, \Star{} Quam abscondísti timéntibus te!\end{Resp}
\begin{Resp}\Vsign Et perfecísti eis qui sperant in te, Dómine, in conspéctu filiórum hóminum.\end{Resp}
\begin{Resp}\Rsign Quam abscondísti timéntibus te!\end{Resp}
\switchcolumn
\EN
\begin{Resp}\Rsign O how great is thy goodness, O Lord \Star{} Which Thou hast laid up for them that fear thee!\end{Resp}
\begin{Resp}\Vsign Which Thou, O Lord, hast wrought for them that trust in thee before the sons of men!\end{Resp}
\begin{Resp}\Rsign Which Thou hast laid up for them that fear thee!\end{Resp}
\switchcolumn*

\LA
\LessonHead{Lectio II}
\switchcolumn
\EN
\LessonHead{Lesson II}
\switchcolumn*

\LA
\Cite{Phil 4:4-7}
\switchcolumn
\EN
\Cite{Phil 4:4-7}
\switchcolumn*

\LA
\noindent Gaudéte in Dómino semper: íterum dico gaudéte. Modéstia vestra nota sit ómnibus homínibus: Dóminus prope est. Nihil sollíciti sitis: sed in omni oratióne, et obsecratióne, cum gratiárum actióne petitiónes vestræ innotéscant apud Deum. Et pax Dei, quæ exúperat omnem sensum, custódiat corda vestra, et intellegéntias vestras in Christo Jesu.\par
\switchcolumn
\EN
\noindent Rejoice in the Lord always; again, I say, rejoice. Let your modesty be known to all men. The Lord is nigh. Be nothing solicitous; but in every thing, by prayer and supplication, with thanksgiving, let your petitions be made known to God. And the peace of God, which surpasseth all understanding, keep your hearts and minds in Christ Jesus.\par
\switchcolumn*

\LA
\begin{Resp}\Rsign Adjútor meus esto, Deus: \Star{} Ne derelínquas me.\end{Resp}
\begin{Resp}\Vsign Neque despícias me, Deus, salutáris meus.\end{Resp}
\begin{Resp}\Rsign Ne derelínquas me.\end{Resp}
\switchcolumn
\EN
\begin{Resp}\Rsign O God, be Thou my helper. \Star{} Neither leave me.\end{Resp}
\begin{Resp}\Vsign Nor forsake me O God of my salvation.\end{Resp}
\begin{Resp}\Rsign Neither leave me.\end{Resp}
\switchcolumn*

\LA
\LessonHead{Lectio III}
\switchcolumn
\EN
\LessonHead{Lesson III}
\switchcolumn*

\LA
\Cite{Phil 4:8-10}
\switchcolumn
\EN
\Cite{Phil 4:8-10}
\switchcolumn*

\LA
\noindent De cétero fratres, quæcúmque sunt vera, quæcúmque pudíca, quæcúmque justa, quæcúmque sancta, quæcúmque amabília, quæcúmque bonæ famæ, siqua virtus, siqua laus disciplínæ, hæc cogitáte. Quæ et didicístis, et accepístis, et audístis, et vidístis in me, hæc ágite: et Deus pacis erit vobíscum. Gavísus sum autem in Dómino veheménter, quóniam tandem aliquándo refloruístis pro me sentíre, sicut et sentiebátis: occupáti autem erátis.\par
\switchcolumn
\EN
\noindent For the rest, brethren, whatsoever things are true, whatsoever modest, whatsoever just, whatsoever holy, whatsoever lovely, whatsoever of good fame, if there be any virtue, if any praise of discipline, think on these things. The things which you have both learned, and received, and heard, and seen in me, these do ye, and the God of peace shall be with you. Now I rejoice in the Lord exceedingly, that now at length your thought for me hath flourished again, as you did also think; but you were busied.\par
\switchcolumn*

\LA
\begin{Resp}\Rsign Benedícam Dóminum in omni témpore: \Star{} Semper laus ejus in ore meo.\end{Resp}
\begin{Resp}\Vsign In Dómino laudábitur ánima mea, áudiant mansuéti, et læténtur.\end{Resp}
\begin{Resp}\Rsign Semper laus ejus in ore meo.\end{Resp}
\begin{Resp}\Vsign Glória Patri, et Fílio, \Star{} et Spirítui Sancto.\end{Resp}
\begin{Resp}\Rsign Semper laus ejus in ore meo.\end{Resp}
\switchcolumn
\EN
\begin{Resp}\Rsign I will bless the Lord at all times. \Star{} His praise shall continually be in my mouth.\end{Resp}
\begin{Resp}\Vsign My soul shall make her boast in the Lord; the humble shall hear thereof and be glad.\end{Resp}
\begin{Resp}\Rsign His praise shall continually be in my mouth.\end{Resp}
\begin{Resp}\Vsign Glory be to the Father, and to the Son, \Star{} and to the Holy Ghost.\end{Resp}
\begin{Resp}\Rsign His praise shall continually be in my mouth.\end{Resp}
\switchcolumn*

\end{paracol}

\Office{Feria Tertia infra Hebdomadam IV post Epiphaniam}{Tuesday of the Fourth Week after Epiphany}{Feria}
\addcontentsline{toc}{subsection}{Feria Tertia infra Hebdomadam IV post Epiphaniam\enspace\textit{Tuesday of the Fourth Week after Epiphany}}
\begin{paracol}{2}
\LA
\LessonHead{Lectio I}
\switchcolumn
\EN
\LessonHead{Lesson I}
\switchcolumn*

\LA
\LessonTitle{Incipit Epístola beáti Pauli Apóstoli ad Colossénses}
\switchcolumn
\EN
\LessonTitle{Lesson from the letter of St. Paul the Apostle to the Colossians}
\switchcolumn*

\LA
\Cite{Col 1:1-8}
\switchcolumn
\EN
\Cite{Col 1:1-8}
\switchcolumn*

\LA
\noindent Paulus Apóstolus Jesu Christi per voluntátem Dei, et Timótheus frater: eis, qui sunt Colóssis, sanctis, et fidélibus frátribus in Christo Jesu. Grátia vobis, et pax a Deo Patre nostro, et Dómino Jesu Christo. Grátias ágimus Deo, et Patri Dómini nostri Jesu Christi semper pro vobis orántes: audiéntes fidem vestram in Christo Jesu, et dilectiónem quam habétis in sanctos omnes propter spem, quæ repósita est vobis in cælis: quam audístis in verbo veritátis Evangélii: quod pervénit ad vos, sicut et in univérso mundo est, et fructíficat, et crescit sicut in vobis, ex ea die, qua audístis, et cognovístis grátiam Dei in veritáte, sicut didicístis ab Epaphra caríssimo consérvo nostro, qui est fidélis pro vobis miníster Christi Jesu, qui étiam manifestávit nobis dilectiónem vestram in spíritu.\par
\switchcolumn
\EN
\noindent Paul, an apostle of Jesus Christ, by the will of God, and Timothy, a brother, To the saints and faithful brethren in Christ Jesus, who are at Colossa. Grace be to you and peace from God our Father, and from the Lord Jesus Christ. We give thanks to God, and the Father of our Lord Jesus Christ, praying always for you. Hearing your faith in Christ Jesus, and the love which you have towards all the saints. For the hope that is laid up for you in heaven, which you have heard in the word of the truth of the gospel, Which is come unto you, as also it is in the whole world, and bringeth forth fruit and groweth, even as it doth in you, since the day you heard and knew the grace of God in truth. As you learned of Epaphras, our most beloved fellow servant, who is for you a faithful minister of Christ Jesus; Who also hath manifested to us your love in the spirit.\par
\switchcolumn*

\LA
\begin{Resp}\Rsign Auribus pércipe, Deus, lácrimas meas: ne síleas a me, remítte mihi: \Star{} Quóniam íncola ego sum apud te, et peregrínus.\end{Resp}
\begin{Resp}\Vsign Compláceat tibi, ut erípias me: Dómine, ad adjuvándum me festína.\end{Resp}
\begin{Resp}\Rsign Quóniam íncola ego sum apud te, et peregrínus.\end{Resp}
\switchcolumn
\EN
\begin{Resp}\Rsign O God, give ear unto my tears; hold not thy peace, but, O, spare me! \Star{} For I am a stranger with thee, and a sojourner.\end{Resp}
\begin{Resp}\Vsign Be pleased, O Lord, to deliver me; O Lord, look upon me to help me.\end{Resp}
\begin{Resp}\Rsign For I am a stranger with thee, and a sojourner.\end{Resp}
\switchcolumn*

\LA
\LessonHead{Lectio II}
\switchcolumn
\EN
\LessonHead{Lesson II}
\switchcolumn*

\LA
\Cite{Col 1:9-12}
\switchcolumn
\EN
\Cite{Col 1:9-12}
\switchcolumn*

\LA
\noindent Ideo et nos ex qua die audívimus, non cessámus pro vobis orántes, et postulántes ut impleámini agnitióne voluntátis ejus, in omni sapiéntia et intelléctu spiritáli: ut ambulétis digne Deo per ómnia placéntes: in omni ópere bono fructificántes, et crescéntes in sciéntia Dei: in omni virtúte confortáti secúndum poténtiam claritátis ejus, in omni patiéntia et longanimitáte cum gáudio, grátias agéntes Deo Patri, qui dignos nos fecit in partem sortis sanctórum in lúmine:\par
\switchcolumn
\EN
\noindent Therefore we also, from the day that we heard it, cease not to pray for you, and to beg that you may be filled with the knowledge of his will, in all wisdom, and spiritual understanding: That you may walk worthy of God, in all things pleasing; being fruitful in every good work, and increasing in the knowledge of God: Strengthened with all might, according to the power of his glory, in all patience and longsuffering with joy, Giving thanks to God the Father, who hath made us worthy to be partakers of the lot of the saints in light:\par
\switchcolumn*

\LA
\begin{Resp}\Rsign Státuit Dóminus supra petram pedes meos, et diréxit gressus meos Deus: \Star{} Et misit in os meum cánticum novum.\end{Resp}
\begin{Resp}\Vsign Exaudívit preces meas: et edúxit me de lacu misériæ.\end{Resp}
\begin{Resp}\Rsign Et misit in os meum cánticum novum.\end{Resp}
\switchcolumn
\EN
\begin{Resp}\Rsign The Lord hath set my feet upon a rock, and ordered my goings. \Star{} And He hath put a new song in my mouth.\end{Resp}
\begin{Resp}\Vsign He heard my cry He brought me up also out of an horrible pit.\end{Resp}
\begin{Resp}\Rsign And He hath put a new song in my mouth.\end{Resp}
\switchcolumn*

\LA
\LessonHead{Lectio III}
\switchcolumn
\EN
\LessonHead{Lesson III}
\switchcolumn*

\LA
\Cite{Col 1:13-18}
\switchcolumn
\EN
\Cite{Col 1:13-18}
\switchcolumn*

\LA
\noindent qui erípuit nos de potestáte tenebrárum, et tránstulit in regnum fílii dilectiónis suæ, in quo habémus redemptiónem per sánguinem ejus, remissiónem peccatórum: qui est imágo Dei invisíbilis, primogénitus omnis creatúræ: quóniam in ipso cóndita sunt univérsa in cælis, et in terra, visibília, et invisibília, sive throni, sive dominatiónes, sive principátus, sive potestátes: ómnia per ipsum et in ipso creáta sunt: et ipse est ante omnes, et ómnia in ipso constant. Et ipse est caput córporis Ecclésiæ, qui est princípium, primogénitus ex mórtuis: ut sit in ómnibus ipse primátum tenens:\par
\switchcolumn
\EN
\noindent Who hath delivered us from the power of darkness, and hath translated us into the kingdom of the Son of his love, In whom we have redemption through his blood, the remission of sins; Who is the image of the invisible God, the firstborn of every creature: For in him were all things created in heaven and on earth, visible and invisible, whether thrones, or dominations, or principalities, or powers: all things were created by him and in him. And he is before all, and by him all things consist. And he is the head of the body, the church, who is the beginning, the firstborn from the dead.\par
\switchcolumn*

\LA
\begin{Resp}\Rsign Ego dixi, Dómine, miserére mei: \Star{} Sana ánimam meam, quia peccávi tibi.\end{Resp}
\begin{Resp}\Vsign Ab ómnibus iniquitátibus meis éripe me, Dómine.\end{Resp}
\begin{Resp}\Rsign Sana ánimam meam, quia peccávi tibi.\end{Resp}
\begin{Resp}\Vsign Glória Patri, et Fílio, \Star{} et Spirítui Sancto.\end{Resp}
\begin{Resp}\Rsign Sana ánimam meam, quia peccávi tibi.\end{Resp}
\switchcolumn
\EN
\begin{Resp}\Rsign I said: Lord, be merciful unto me. \Star{} Heal my soul, for I have sinned against thee.\end{Resp}
\begin{Resp}\Vsign Deliver me from all mine iniquities, O Lord.\end{Resp}
\begin{Resp}\Rsign Heal my soul, for I have sinned against thee.\end{Resp}
\begin{Resp}\Vsign Glory be to the Father, and to the Son, \Star{} and to the Holy Ghost.\end{Resp}
\begin{Resp}\Rsign Heal my soul, for I have sinned against thee.\end{Resp}
\switchcolumn*

\end{paracol}

\Office{Feria IV infra Hebdomadam IV post Epiphaniam}{Wednesday of the Fourth Week after Epiphany}{Feria}
\addcontentsline{toc}{subsection}{Feria IV infra Hebdomadam IV post Epiphaniam\enspace\textit{Wednesday of the Fourth Week after Epiphany}}
\begin{paracol}{2}
\LA
\LessonHead{Lectio I}
\switchcolumn
\EN
\LessonHead{Lesson I}
\switchcolumn*

\LA
\LessonTitle{De Epístola ad Colossénses}
\switchcolumn
\EN
\LessonTitle{Lesson from the letter of St. Paul the Apostle to the Colossians}
\switchcolumn*

\LA
\Cite{Col 3:12-15}
\switchcolumn
\EN
\Cite{Col 3:12-15}
\switchcolumn*

\LA
\noindent Indúite vos ergo, sicut elécti Dei, sancti, et dilécti, víscera misericórdiæ, benignitátem, humilitátem, modéstiam, patiéntiam: supportántes ínvicem, et donántes vobismetípsis si quis advérsus áliquem habet querélam: sicut et Dóminus donávit vobis, ita et vos. Super ómnia autem hæc, caritátem habéte, quod est vínculum perfectiónis: et pax Christi exsúltet in córdibus vestris, in qua et vocáti estis in uno córpore: et grati estóte.\par
\switchcolumn
\EN
\noindent Put ye on therefore, as the elect of God, holy, and beloved, the bowels of mercy, benignity, humility, modesty, patience: Bearing with one another, and forgiving one another, if any have a complaint against another: even as the Lord hath forgiven you, so do you also. But above all these things have charity, which is the bond of perfection: And let the peace of Christ rejoice in your hearts, wherein also you are called in one body: and be ye thankful.\par
\switchcolumn*

\LA
\begin{Resp}\Rsign Ne perdíderis me cum iniquitátibus meis: \Star{} Neque in finem irátus resérves mala mea.\end{Resp}
\begin{Resp}\Vsign Non intres in judícium cum servo tuo, Dómine.\end{Resp}
\begin{Resp}\Rsign Neque in finem irátus resérves mala mea.\end{Resp}
\switchcolumn
\EN
\begin{Resp}\Rsign Cut me not off in the midst of my sins. \Star{} Nor keep thy wrath against me for my latter end.\end{Resp}
\begin{Resp}\Vsign Enter not into judgment with thy servant, O Lord.\end{Resp}
\begin{Resp}\Rsign Nor keep thy wrath against me for my latter end.\end{Resp}
\switchcolumn*

\LA
\LessonHead{Lectio II}
\switchcolumn
\EN
\LessonHead{Lesson II}
\switchcolumn*

\LA
\Cite{Col 3:16-21}
\switchcolumn
\EN
\Cite{Col 3:16-21}
\switchcolumn*

\LA
\noindent Verbum Christi hábitet in vobis abundánter, in omni sapiéntia, docéntes, et commonéntes vosmetípsos, psalmis, hymnis, et cánticis spirituálibus, in grátia cantántes in córdibus vestris Deo. Omne, quodcúmque fácitis in verbo aut in ópere, ómnia in nómine Dómini Jesu Christi, grátias agéntes Deo et Patri per ipsum. Mulíeres, súbditæ estóte viris, sicut opórtet, in Dómino. Viri, dilígite uxóres vestras, et nolíte amári esse ad illas. Fílii, obedíte paréntibus per ómnia: hoc enim plácitum est in Dómino. Patres, nolíte ad indignatiónem provocáre fílios vestros, ut non pusíllo ánimo fiant.\par
\switchcolumn
\EN
\noindent Let the word of Christ dwell in you abundantly, in all wisdom: teaching and admonishing one another in psalms, hymns, and spiritual canticles, singing in grace in your hearts to God. All whatsoever you do in word or in work, do all in the name of the Lord Jesus Christ, giving thanks to God and the Father by him. Wives, be subject to your husbands, as it behoveth in the Lord. Husbands, love your wives, and be not bitter towards them. Children, obey your parents in all things: for this is well pleasing to the Lord. Fathers, provoke not your children to indignation, lest they be discouraged.\par
\switchcolumn*

\LA
\begin{Resp}\Rsign Parátum cor meum, Deus, parátum cor meum: \Star{} Cantábo et psalmum dicam Dómino.\end{Resp}
\begin{Resp}\Vsign Exsúrge, glória mea, exsúrge, psaltérium et cíthara, exsúrgam dilúculo.\end{Resp}
\begin{Resp}\Rsign Cantábo et psalmum dicam Dómino.\end{Resp}
\switchcolumn
\EN
\begin{Resp}\Rsign My heart is ready, O God, my heart is ready. \Star{} I will sing and give praise to the Lord.\end{Resp}
\begin{Resp}\Vsign Awake up, my glory, awake, psaltery and harp! I will awake early.\end{Resp}
\begin{Resp}\Rsign I will sing and give praise to the Lord.\end{Resp}
\switchcolumn*

\LA
\LessonHead{Lectio III}
\switchcolumn
\EN
\LessonHead{Lesson III}
\switchcolumn*

\LA
\Cite{Col 3:22-25;4:1-2}
\switchcolumn
\EN
\Cite{Col 3:22-25; 4:1-2}
\switchcolumn*

\LA
\noindent Servi, obedíte per ómnia dóminis carnálibus, non ad óculum serviéntes, quasi homínibus placéntes, sed in simplicitáte cordis, timéntes Deum. Quodcúmque fácitis, ex ánimo operámini sicut Dómino, et non homínibus: sciéntes quod a Dómino accipiétis retributiónem hæreditátis. Dómino Christo servíte. Qui enim injúriam facit, recípiet id quod iníque gessit: et non est personárum accéptio apud Deum. Dómini, quod justum est et æquum, servis præstáte: sciéntes quod et vos Dóminum habétis in cælo. Oratióni instáte, vigilántes in ea in gratiárum actióne:\par
\switchcolumn
\EN
\noindent Servants, obey in all things your masters according to the flesh, not serving to the eye, as pleasing men, but in simplicity of heart, fearing God. Whatsoever you do, do it from the heart, as to the Lord, and not to men: Knowing that you shall receive of the Lord the reward of inheritance. Serve ye the Lord Christ. For he that doth wrong, shall receive for that which he hath done wrongfully: and there is no respect of persons with God. Masters, do to your servants that which is just and equal: knowing that you also have a master in heaven. Be instant in prayer; watching in it with thanksgiving:\par
\switchcolumn*

\LA
\begin{Resp}\Rsign Adjútor meus, tibi psallam, quia, Deus, suscéptor meus es: \Star{} Deus meus, misericórdia mea.\end{Resp}
\begin{Resp}\Vsign Lætábor, et exsultábo in te, psallam nómini tuo, Altíssime.\end{Resp}
\begin{Resp}\Rsign Deus meus, misericórdia mea.\end{Resp}
\begin{Resp}\Vsign Glória Patri, et Fílio, \Star{} et Spirítui Sancto.\end{Resp}
\begin{Resp}\Rsign Deus meus, misericórdia mea.\end{Resp}
\switchcolumn
\EN
\begin{Resp}\Rsign Unto thee, O my Strength, will I sing, for God is my defence, \Star{} The God of my mercy.\end{Resp}
\begin{Resp}\Vsign I will be glad and rejoice in thee, I will sing praise to thy Name, O Thou Most High.\end{Resp}
\begin{Resp}\Rsign The God of my mercy.\end{Resp}
\begin{Resp}\Vsign Glory be to the Father, and to the Son, \Star{} and to the Holy Ghost.\end{Resp}
\begin{Resp}\Rsign The God of my mercy.\end{Resp}
\switchcolumn*

\end{paracol}

\Office{Feria V infra Hebdomadam IV post Epiphaniam}{Thursday of the Fourth Week after Epiphany}{Feria}
\addcontentsline{toc}{subsection}{Feria V infra Hebdomadam IV post Epiphaniam\enspace\textit{Thursday of the Fourth Week after Epiphany}}
\begin{paracol}{2}
\LA
\LessonHead{Lectio I}
\switchcolumn
\EN
\LessonHead{Lesson I}
\switchcolumn*

\LA
\LessonTitle{Incipit Epístola prima beáti Pauli Apóstoli ad Thessalonicénses}
\switchcolumn
\EN
\LessonTitle{Lesson from the first letter of St. Paul the Apostle to the Thessalonians}
\switchcolumn*

\LA
\Cite{1 Thess 1:1-5}
\switchcolumn
\EN
\Cite{1 Thess 1:1-5}
\switchcolumn*

\LA
\noindent Paulus, et Silvánus, et Timótheus ecclésiæ Thessalonicénsium in Deo Patre, et Dómino Jesu Christo. Grátia vobis, et pax. Grátias ágimus Deo semper pro ómnibus vobis, memóriam vestri faciéntes in oratiónibus nostris sine intermissióne, mémores óperis fídei vestræ, et labóris, et caritátis, et sustinéntiæ spei Dómini nostri Jesu Christi, ante Deum et Patrem nostrum: sciéntes, fratres dilécti a Deo, electiónem vestram: quia Evangélium nostrum non fuit ad vos in sermóne tantum, sed et in virtúte, et in Spíritu Sancto, et in plenitúdine multa, sicut scitis quales fuérimus in vobis propter vos.\par
\switchcolumn
\EN
\noindent Paul and Sylvanus and Timothy: to the church of the Thessalonians, in God the Father, and in the Lord Jesus Christ. Grace be to you and peace. We give thanks to God always for you all; making a remembrance of you in our prayers without ceasing, Being mindful of the work of your faith, and labour, and charity, and of the enduring of the hope of our Lord Jesus Christ before God and our Father: Knowing, brethren beloved of God, your election: For our gospel hath not been unto you in word only, but in power also, and in the Holy Ghost, and in much fulness, as you know what manner of men we have been among you for your sakes.\par
\switchcolumn*

\LA
\begin{Resp}\Rsign Deus, in te sperávi, Dómine, non confúndar in ætérnum: in justítia tua líbera me, \Star{} Et éripe me.\end{Resp}
\begin{Resp}\Vsign Inclína ad me aurem tuam, et salva me.\end{Resp}
\begin{Resp}\Rsign Et éripe me.\end{Resp}
\switchcolumn
\EN
\begin{Resp}\Rsign In thee, O God, do I put my trust; let me never be put to confusion, O Lord deliver me in thy righteousness, \Star{} And cause me to escape.\end{Resp}
\begin{Resp}\Vsign Incline thine ear unto me, deliver me speedily.\end{Resp}
\begin{Resp}\Rsign And cause me to escape.\end{Resp}
\switchcolumn*

\LA
\LessonHead{Lectio II}
\switchcolumn
\EN
\LessonHead{Lesson II}
\switchcolumn*

\LA
\Cite{1 Thess 1:6-10}
\switchcolumn
\EN
\Cite{1 Thess 1:6-10}
\switchcolumn*

\LA
\noindent Et vos imitatóres nostri facti estis, et Dómini, excipiéntes verbum in tribulatióne multa, cum gáudio Spíritus Sancti: ita ut facti sitis forma ómnibus credéntibus in Macedónia, et in Achája. A vobis enim diffamátus est sermo Dómini, non solum in Macedónia, et in Achája, sed et in omni loco fides vestra, quæ est ad Deum, profécta est ita ut non sit nobis necésse quidquam loqui. Ipsi enim de nobis annúntiant qualem intróitum habuérimus ad vos: et quómodo convérsi estis ad Deum a simulácris, servíre Deo vivo, et vero, et exspectáre Fílium ejus de cælis (quem suscitávit a mórtuis) Jesum, qui erípuit nos ab ira ventúra.\par
\switchcolumn
\EN
\noindent And you became followers of us, and of the Lord; receiving the word in much tribulation, with joy of the Holy Ghost: So that you were made a pattern to all that believe in Macedonia and in Achaia. For from you was spread abroad the word of the Lord, not only in Macedonia, and in Achaia, but also in every place, your faith which is towards God, is gone forth, so that we need not to speak any thing. For they themselves relate of us, what manner of entering in we had unto you; and how you turned to God from idols, to serve the living and true God. And to wait for his Son from heaven (whom he raised up from the dead,) Jesus, who hath delivered us from the wrath to come.\par
\switchcolumn*

\LA
\begin{Resp}\Rsign Repleátur os meum laude tua, ut hymnum dicam glóriæ tuæ, tota die magnitúdinem tuam: noli me proícere in témpore senectútis: \Star{} Dum defécerit in me virtus mea, ne derelínquas me.\end{Resp}
\begin{Resp}\Vsign Gaudébunt lábia mea cum cantávero tibi.\end{Resp}
\begin{Resp}\Rsign Dum defécerit in me virtus mea, ne derelínquas me.\end{Resp}
\switchcolumn
\EN
\begin{Resp}\Rsign Let my mouth be filled with thy praise, that I may sing of thy glory, all the day long of thy greatness. Cast me not off in the time of old age; \Star{} Forsake me not when my strength faileth.\end{Resp}
\begin{Resp}\Vsign My lips shall be fain when I sing unto thee.\end{Resp}
\begin{Resp}\Rsign Forsake me not when my strength faileth.\end{Resp}
\switchcolumn*

\LA
\LessonHead{Lectio III}
\switchcolumn
\EN
\LessonHead{Lesson III}
\switchcolumn*

\LA
\Cite{1 Thess 2:1-6}
\switchcolumn
\EN
\Cite{1 Thess 2:1-6}
\switchcolumn*

\LA
\noindent Nam ipsi scitis, fratres, intróitum nostrum ad vos, quia non inánis fuit: sed ante passi, et contuméliis affécti (sicut scitis) in Philíppis, fidúciam habúimus in Deo nostro, loqui ad vos Evangélium Dei in multa sollicitúdine. Exhortátio enim nostra non de erróre, neque de immundítia, neque in dolo, sed sicut probáti sumus a Deo ut crederétur nobis Evangélium: ita lóquimur non quasi homínibus placéntes, sed Deo, qui probat corda nostra. Neque enim aliquándo fúimus in sermóne adulatiónis, sicut scitis: neque in occasióne avarítiæ: Deus testis est: nec quæréntes ab homínibus glóriam, neque a vobis, neque ab áliis.\par
\switchcolumn
\EN
\noindent For yourselves know, brethren, our entrance in unto you, that it was not in vain: But having suffered many things before, and been shamefully treated (as you know) at Philippi, we had confidence in our God, to speak unto you the gospel of God in much carefulness. For our exhortation was not of error, nor of uncleanness, nor in deceit: But as we were approved by God that the gospel should be committed to us: even so we speak, not as pleasing men, but God, who proveth our hearts. For neither have we used, at any time, the speech of flattery, as you know; nor taken an occasion of covetousness, God is witness: Nor sought we glory of men, neither of you, nor of others. Whereas we might have been burdensome to you, as the apostles of Christ: but we became little ones in the midst of you, as if a nurse should cherish her children:\par
\switchcolumn*

\LA
\begin{Resp}\Rsign Gaudébunt lábia mea cum cantávero tibi: \Star{} Et ánima mea, quam redemísti, Dómine.\end{Resp}
\begin{Resp}\Vsign Sed et lingua mea meditábitur justítiam tuam, tota die laudem tuam.\end{Resp}
\begin{Resp}\Rsign Et ánima mea, quam redemísti, Dómine.\end{Resp}
\begin{Resp}\Vsign Glória Patri, et Fílio, \Star{} et Spirítui Sancto.\end{Resp}
\begin{Resp}\Rsign Et ánima mea, quam redemísti, Dómine.\end{Resp}
\switchcolumn
\EN
\begin{Resp}\Rsign My lips shall be fain when I sing unto thee; \Star{} And my soul, which Thou, O Lord, hast redeemed.\end{Resp}
\begin{Resp}\Vsign My tongue shall also talk of thy righteousness, all the day long of thy praise.\end{Resp}
\begin{Resp}\Rsign And my soul, which Thou, O Lord, hast redeemed.\end{Resp}
\begin{Resp}\Vsign Glory be to the Father, and to the Son, \Star{} and to the Holy Ghost.\end{Resp}
\begin{Resp}\Rsign And my soul, which Thou, O Lord, hast redeemed.\end{Resp}
\switchcolumn*

\end{paracol}

\Office{Feria VI infra Hebdomadam IV post Epiphaniam}{Friday of the Fourth Week after Epiphany}{Feria}
\addcontentsline{toc}{subsection}{Feria VI infra Hebdomadam IV post Epiphaniam\enspace\textit{Friday of the Fourth Week after Epiphany}}
\begin{paracol}{2}
\LA
\LessonHead{Lectio I}
\switchcolumn
\EN
\LessonHead{Lesson I}
\switchcolumn*

\LA
\LessonTitle{De Epístola prima ad Thessalonicénses}
\switchcolumn
\EN
\LessonTitle{Lesson from the first letter of St. Paul the Apostle to the Thessalonians}
\switchcolumn*

\LA
\Cite{1 Thess 4:1-5}
\switchcolumn
\EN
\Cite{1 Thess 4:1-5}
\switchcolumn*

\LA
\noindent De cétero ergo, fratres, rogámus vos et obsecrámus in Dómino Jesu, ut quemádmodum accepístis a nobis quómodo opórteat vos ambuláre, et placére Deo, sic et ambulétis ut abundétis magis. Scitis enim quæ præcépta déderim vobis per Dóminum Jesum. Hæc est enim volúntas Dei, sanctificátio vestra: ut abstineátis vos a fornicatióne, ut sciat unusquísque vestrum vas suum possidére in sanctificatióne, et honóre: non in passióne desidérii, sicut et gentes, quæ ignórant Deum:\par
\switchcolumn
\EN
\noindent For the rest therefore, brethren, we pray and beseech you in the Lord Jesus, that as you have received from us, how you ought to walk, and to please God, so also you would walk, that you may abound the more. For you know what precepts I have given to you by the Lord Jesus. For this is the will of God, your sanctification; that you should abstain from fornication; That every one of you should know how to possess his vessel in sanctification and honour: Not in the passion of lust, like the Gentiles that know not God:\par
\switchcolumn*

\LA
\begin{Resp}\Rsign Confitébor tibi, Dómine Deus, in toto corde meo, et honorificábo nomen tuum in ætérnum: \Star{} Quia misericórdia tua, Dómine, magna est super me.\end{Resp}
\begin{Resp}\Vsign Deus meus es tu, et confitébor tibi: Deus meus es tu, et exaltábo te.\end{Resp}
\begin{Resp}\Rsign Quia misericórdia tua, Dómine, magna est super me.\end{Resp}
\switchcolumn
\EN
\begin{Resp}\Rsign I will praise thee, O Lord my God, with all my heart, and I will glorify thy Name for evermore. \Star{} For great is thy mercy, O Lord, toward me.\end{Resp}
\begin{Resp}\Vsign Thou art my God, and I will praise thee; Thou art my God, and I will exalt thee.\end{Resp}
\begin{Resp}\Rsign For great is thy mercy, O Lord, toward me.\end{Resp}
\switchcolumn*

\LA
\LessonHead{Lectio II}
\switchcolumn
\EN
\LessonHead{Lesson II}
\switchcolumn*

\LA
\Cite{1 Thess 4:6-8}
\switchcolumn
\EN
\Cite{1 Thess 4:6-8}
\switchcolumn*

\LA
\noindent et ne quis supergrediátur, neque circumvéniat in negótio fratrem suum: quóniam vindex est Dóminus de his ómnibus, sicut prædíximus vobis, et testificáti sumus. Non enim vocávit nos Deus in immundítiam, sed in sanctificatiónem. Itaque qui hæc spernit, non hóminem spernit, sed Deum: qui étiam dedit Spíritum suum Sanctum in nobis.\par
\switchcolumn
\EN
\noindent And that no man overreach, nor circumvent his brother in business: because the Lord is the avenger of all these things, as we have told you before, and have testified. For God hath not called us unto uncleanness, but unto sanctification. Therefore, he that despiseth these things, despiseth not man, but God, who also hath given his holy Spirit in us.\par
\switchcolumn*

\LA
\begin{Resp}\Rsign Misericórdia tua, Dómine, magna est super me: \Star{} Et liberásti ánimam meam ex inférno inferióri.\end{Resp}
\begin{Resp}\Vsign In die tribulatiónis meæ clamávi ad te, quia exaudísti me.\end{Resp}
\begin{Resp}\Rsign Et liberásti ánimam meam ex inférno inferióri.\end{Resp}
\switchcolumn
\EN
\begin{Resp}\Rsign Great, O Lord, is thy mercy toward me. \Star{} And Thou hast delivered my soul from the lowest hell.\end{Resp}
\begin{Resp}\Vsign In the day of my trouble I called upon thee, for Thou hast heard me.\end{Resp}
\begin{Resp}\Rsign And Thou hast delivered my soul from the lowest hell.\end{Resp}
\switchcolumn*

\LA
\LessonHead{Lectio III}
\switchcolumn
\EN
\LessonHead{Lesson III}
\switchcolumn*

\LA
\Cite{1 Thess 4:9-12}
\switchcolumn
\EN
\Cite{1 Thess 4:9-11}
\switchcolumn*

\LA
\noindent De caritáte autem fraternitátis non necésse habémus scríbere vobis: ipsi enim vos a Deo didicístis ut diligátis ínvicem. Etenim illud fácitis in omnes fratres in univérsa Macedónia. Rogámus autem vos, fratres, ut abundétis magis, et óperam detis ut quiéti sitis, et ut vestrum negótium agátis, et operémini mánibus vestris, sicut præcépimus vobis: et ut honéste ambulétis ad eos qui foris sunt: et nullíus áliquid desiderétis.\par
\switchcolumn
\EN
\noindent But as touching the charity of brotherhood, we have no need to write to you: for yourselves have learned of God to love one another. For indeed you do it towards all the brethren in all Macedonia. But we entreat you, brethren, that you abound more: And that you use your endeavour to be quiet, and that you do your own business, and work with your own hands, as we commanded you: and that you walk honestly towards them that are without; and that you want nothing of any man's.\par
\switchcolumn*

\LA
\begin{Resp}\Rsign Factus est mihi Dóminus in refúgium: \Star{} Et Deus meus in auxílium spei meæ.\end{Resp}
\begin{Resp}\Vsign Erípuit me de inimícis meis fortíssimis, et factus est Dóminus protéctor meus.\end{Resp}
\begin{Resp}\Rsign Et Deus meus in auxílium spei meæ.\end{Resp}
\begin{Resp}\Vsign Glória Patri, et Fílio, \Star{} et Spirítui Sancto.\end{Resp}
\begin{Resp}\Rsign Et Deus meus in auxílium spei meæ.\end{Resp}
\switchcolumn
\EN
\begin{Resp}\Rsign The Lord is my refuge. \Star{} And my God is the stay of my trust.\end{Resp}
\begin{Resp}\Vsign He delivered me from the strongest of mine enemies, and the Lord was my stay.\end{Resp}
\begin{Resp}\Rsign And my God is the stay of my trust.\end{Resp}
\begin{Resp}\Vsign Glory be to the Father, and to the Son, \Star{} and to the Holy Ghost.\end{Resp}
\begin{Resp}\Rsign And my God is the stay of my trust.\end{Resp}
\switchcolumn*

\end{paracol}

\Office{Sabbato infra Hebdomadam IV post Epiphaniam}{Saturday of the Fourth Week after Epiphany}{Feria}
\addcontentsline{toc}{subsection}{Sabbato infra Hebdomadam IV post Epiphaniam\enspace\textit{Saturday of the Fourth Week after Epiphany}}
\begin{paracol}{2}
\LA
\LessonHead{Lectio I}
\switchcolumn
\EN
\LessonHead{Lesson I}
\switchcolumn*

\LA
\LessonTitle{Incipit Epístola secúnda beáti Pauli Apóstoli ad Thessalonicénses}
\switchcolumn
\EN
\LessonTitle{Lesson from the second letter of St. Paul the Apostle to the Thessalonians}
\switchcolumn*

\LA
\Cite{2 Thess 1:1-5}
\switchcolumn
\EN
\Cite{2 Thess 1:1-5}
\switchcolumn*

\LA
\noindent Paulus, et Sylvánus, et Timótheus, ecclésiæ Thessalonicénsium in Deo Patre nostro, et Dómino Jesu Christo. Grátia vobis, et pax a Deo Patre nostro, et Dómino Jesu Christo. Grátias ágere debémus semper Deo pro vobis, fratres, ita ut dignum est, quóniam supercréscit fides vestra, et abúndat cáritas uniuscujúsque vestrum in ínvicem: ita ut et nos ipsi in vobis gloriémur in ecclésiis Dei, pro patiéntia vestra, et fide, et in ómnibus persecutiónibus vestris, et tribulatiónibus, quas sustinétis in exémplum justi judícii Dei, ut digni habeámini in regno Dei, pro quo et patímini.\par
\switchcolumn
\EN
\noindent Paul, and Sylvanus, and Timothy, to the church of the Thessalonians in God our Father, and the Lord Jesus Christ. Grace unto you, and peace from God our Father, and from the Lord Jesus Christ. We are bound to give thanks always to God for you, brethren, as it is fitting, because your faith groweth exceedingly, and the charity of every one of you towards each other, aboundeth: So that we ourselves also glory in you in the churches of God, for your patience and faith, and in all your persecutions and tribulations, which you endure, For an example of the just judgment of God, that you may be counted worthy of the kingdom of God, for which also you suffer.\par
\switchcolumn*

\LA
\begin{Resp}\Rsign Misericórdiam et judícium cantábo tibi, Dómine: \Star{} Psallam et intéllegam in via immaculáta, quando vénies ad me.\end{Resp}
\begin{Resp}\Vsign Perambulábam in innocéntia cordis mei, in médio domus meæ.\end{Resp}
\begin{Resp}\Rsign Psallam et intéllegam in via immaculáta, quando vénies ad me.\end{Resp}
\switchcolumn
\EN
\begin{Resp}\Rsign Mercy and judgment I will sing to thee, O Lord: \Star{} I will sing, and I will understand in the unspotted way, when thou shalt come to me.\end{Resp}
\begin{Resp}\Vsign I walked in the innocence of my heart, in the midst of my house.\end{Resp}
\begin{Resp}\Rsign Mercy and judgment I will sing to thee, O Lord\end{Resp}
\switchcolumn*

\LA
\LessonHead{Lectio II}
\switchcolumn
\EN
\LessonHead{Lesson II}
\switchcolumn*

\LA
\Cite{2 Thess 1:6-12}
\switchcolumn
\EN
\Cite{2 Thess 1:6-12}
\switchcolumn*

\LA
\noindent Si tamen justum est apud Deum retribúere tribulatiónem iis qui vos tríbulant: et vobis, qui tribulámini, réquiem nobíscum in revelatióne Dómini Jesu de cælo cum ángelis virtútis ejus, in flamma ignis dantis vindíctam iis qui non novérunt Deum, et qui non obédiunt Evangélio Dómini nostri Jesu Christi, qui pœnas dabunt in intéritu ætérnas a fácie Dómini, et a glória virtútis ejus: cum vénerit glorificári in sanctis suis, et admirábilis fíeri in ómnibus, qui credidérunt, quia créditum est testimónium nostrum super vos in die illo. In quo étiam orámus semper pro vobis: ut dignétur vos vocatióne sua Deus noster, et ímpleat omnem voluntátem bonitátis, et opus fídei in virtúte, ut clarificétur nomen Dómini nostri Jesu Christi in vobis, et vos in illo secúndum grátiam Dei nostri, et Dómini Jesu Christi.\par
\switchcolumn
\EN
\noindent Seeing it is a just thing with God to repay tribulation to them that trouble you: And to you who are troubled, rest with us when the Lord Jesus shall be revealed from heaven, with the angels of his power: In a flame of fire, giving vengeance to them who know not God, and who obey not the gospel of our Lord Jesus Christ. Who shall suffer eternal punishment in destruction, from the face of the Lord, and from the glory of his power: When he shall come to be glorified in his saints, and to be made wonderful in all them who have believed; because our testimony was believed upon you in that day. Wherefore also we pray always for you; that our God would make you worthy of his vocation, and fulfill all the good pleasure of his goodness and the work of faith in power; That the name of our Lord Jesus may be glorified in you, and you in him, according to the grace of our God, and of the Lord Jesus Christ.\par
\switchcolumn*

\LA
\begin{Resp}\Rsign Dómine, exáudi oratiónem meam, et clamor meus ad te pervéniat: \Star{} Quia non spernis, Deus, preces páuperum.\end{Resp}
\begin{Resp}\Vsign Fiant aures tuæ intendéntes in oratiónem servi tui.\end{Resp}
\begin{Resp}\Rsign Quia non spernis, Deus, preces páuperum.\end{Resp}
\switchcolumn
\EN
\begin{Resp}\Rsign Hear my prayer, O Lord, and let my crying come unto thee. \Star{} For thou wilt not despise, O God, the prayer of the destitute\end{Resp}
\begin{Resp}\Vsign O let thine ears consider well the voice of my complaint\end{Resp}
\begin{Resp}\Rsign For thou wilt not despise, O God, the prayer of the destitute.\end{Resp}
\switchcolumn*

\LA
\LessonHead{Lectio III}
\switchcolumn
\EN
\LessonHead{Lesson III}
\switchcolumn*

\LA
\Cite{2 Thess 2:1-4}
\switchcolumn
\EN
\Cite{2 Thess 2:1-4}
\switchcolumn*

\LA
\noindent Rogámus autem vos, fratres, per advéntum Dómini nostri Jesu Christi, et nostræ congregatiónis in ipsum: ut non cito moveámini a vestro sensu, neque terreámini, neque per spíritum, neque per sermónem, neque per epístolam tamquam per nos missam, quasi instet dies Dómini. Ne quis vos sedúcat ullo modo: quóniam nisi vénerit discéssio primum, et revelátus fúerit homo peccáti fílius perditiónis, qui adversátur, et extóllitur supra omne, quod dícitur Deus, aut quod cólitur, ita ut in templo Dei sédeat osténdens se tamquam sit Deus.\par
\switchcolumn
\EN
\noindent And we beseech you, brethren, by the coming of our Lord Jesus Christ, and of our gathering together unto him: That you be not easily moved from your sense, nor be terrified, neither by spirit, nor by word, nor by epistle, as sent from us, as if the day of the Lord were at hand. Let no man deceive you by any means, for unless there come a revolt first, and the man of sin be revealed, the son of perdition, Who opposeth, and is lifted up above all that is called God, or that is worshipped, so that he sitteth in the temple of God, shewing himself as if he were God.\par
\switchcolumn*

\LA
\begin{Resp}\Rsign Velóciter exáudi me, Deus, \Star{} Tu autem idem ipse es, et anni tui non defícient.\end{Resp}
\begin{Resp}\Vsign Dies mei sicut umbra declinavérunt, et ego sicut fœnum árui.\end{Resp}
\begin{Resp}\Rsign Quia defecérunt sicut fumus dies mei.\end{Resp}
\begin{Resp}\Vsign Glória Patri, et Fílio, \Star{} et Spirítui Sancto.\end{Resp}
\begin{Resp}\Rsign Tu autem idem ipse es, et anni tui non defícient.\end{Resp}
\switchcolumn
\EN
\begin{Resp}\Rsign Haste thee to hear me, O Lord. \Star{} But thou art the same, and thy years shall not fail.\end{Resp}
\begin{Resp}\Vsign My days have declined like a shadow, and I am withered like grass.\end{Resp}
\begin{Resp}\Rsign For my days are consumed away like smoke.\end{Resp}
\begin{Resp}\Vsign Glory be to the Father, and to the Son, \Star{} and to the Holy Ghost.\end{Resp}
\begin{Resp}\Rsign But thou art the same, and thy years shall not fail.\end{Resp}
\switchcolumn*

\end{paracol}

\Office{Dominica V Post Epiphaniam}{Fifth Sunday after Epiphany}{Semiduplex}
\addcontentsline{toc}{subsection}{Dominica V Post Epiphaniam\enspace\textit{Fifth Sunday after Epiphany}}
\begin{paracol}{2}
\LA
\RefLine{Lectiones I–VI: de Scriptura occurrente.}
\switchcolumn
\EN
\RefLine{Lessons I–VI: from the Scripture of the day.}
\switchcolumn*

\LA
\LessonHead{Lectio I}
\switchcolumn
\EN
\LessonHead{Lesson I}
\switchcolumn*

\LA
\LessonTitle{Incipit Epístola prima beáti Pauli Apóstoli ad Timótheum}
\switchcolumn
\EN
\LessonTitle{Lesson from the first letter of St. Paul the Apostle to Timothy}
\switchcolumn*

\LA
\Cite{1 Tim 1:1-4}
\switchcolumn
\EN
\Cite{1 Tim 1:1-4}
\switchcolumn*

\LA
\noindent Paulus Apóstolus Jesu Christi secúndum impérium Dei Salvatóris nostri, et Christi Jesu spei nostræ: Timótheo dilécto fílio in fide. Grátia, misericórdia, et pax a Deo Patre, et Christo Jesu Dómino nostro. Sicut rogávi te ut remanéres Ephesi, cum irem in Macedóniam, ut denuntiáres quibúsdam ne áliter docérent, neque inténderent fábulis, et genealogíis interminátis: quæ quæstiónes præstant magis quam ædificatiónem Dei, quæ est in fide.\par
\switchcolumn
\EN
\noindent Paul, an apostle of Jesus Christ, according to the commandment of God our Saviour, and of Christ Jesus our hope: To Timothy, his beloved son in faith. Grace, mercy, and peace from God the Father, and from Christ Jesus our Lord. As I desired thee to remain at Ephesus when I went into Macedonia, that thou mightest charge some not to teach otherwise, Not to give heed to fables and endless genealogies: which furnish questions rather than the edification of God, which is in faith.\par
\switchcolumn*

\LA
\begin{Resp}\Rsign Dómine, ne in ira tua árguas me, neque in furóre tuo corrípias me: \Star{} Miserére mei, Dómine, quóniam infírmus sum.\end{Resp}
\begin{Resp}\Vsign Timor et tremor venérunt super me, et contexérunt me ténebræ.\end{Resp}
\begin{Resp}\Rsign Miserére mei, Dómine, quóniam infírmus sum.\end{Resp}
\switchcolumn
\EN
\begin{Resp}\Rsign O Lord, rebuke me not in thine anger, neither chasten me in thine hot displeasure. \Star{} Have mercy upon me, O Lord, for I am weak.\end{Resp}
\begin{Resp}\Vsign Fearfulness and trembling are come upon me, and darkness hath overwhelmed me.\end{Resp}
\begin{Resp}\Rsign Have mercy upon me, O Lord, for I am weak.\end{Resp}
\switchcolumn*

\LA
\LessonHead{Lectio II}
\switchcolumn
\EN
\LessonHead{Lesson II}
\switchcolumn*

\LA
\Cite{1 Tim 1:5-11}
\switchcolumn
\EN
\Cite{1 Tim 1:5-11}
\switchcolumn*

\LA
\noindent Finis autem præcépti est cáritas de corde puro, et consciéntia bona, et fide non ficta. A quibus quidam aberrántes, convérsi sunt in vanilóquium, voléntes esse legis doctóres, non intellegéntes neque quæ loquúntur, neque de quibus affírmant. Scimus autem quia bona est lex, si quis ea legítime utátur: sciens hoc quia lex justo non est pósita, sed injústis, et non súbditis, ímpiis et peccatóribus, scelerátis, et contaminátis, parricídis et matricídis, homicídis, fornicáriis, masculórum concubitóribus, plagiáriis, mendácibus et perjúris, et si quid áliud sanæ doctrínæ adversátur, quæ est secúndum Evangélium glóriæ beáti Dei, quod créditum est mihi.\par
\switchcolumn
\EN
\noindent Now the end of the commandment is charity, from a pure heart, and a good conscience, and an unfeigned faith. From which things some going astray, are turned aside unto vain babbling: Desiring to be teachers of the law, understanding neither the things they say, nor whereof they affirm. But we know that the law is good, if a man use it lawfully: Knowing this, that the law is not made for the just man, but for the unjust and disobedient, for the ungodly, and for sinners, for the wicked and defiled, for murderers of fathers, and murderers of mothers, for manslayers, For fornicators, for them who defile themselves with mankind, for menstealers, for liars, for perjured persons, and whatever other thing is contrary to sound doctrine, Which is according to the gospel of the glory of the blessed God, which hath been committed to my trust.\par
\switchcolumn*

\LA
\begin{Resp}\Rsign Deus, qui sedes super thronum, et júdicas æquitátem, esto refúgium páuperum in tribulatióne: \Star{} Quia tu solus labórem et dolórem consíderas.\end{Resp}
\begin{Resp}\Vsign Tibi enim derelíctus est pauper, pupíllo tu eris adjútor.\end{Resp}
\begin{Resp}\Rsign Quia tu solus labórem et dolórem consíderas.\end{Resp}
\switchcolumn
\EN
\begin{Resp}\Rsign O God, Which satest in the throne judging right, be Thou a refuge for the poor, a refuge in times of trouble. \Star{} For Thou alone beholdest mischief and spite.\end{Resp}
\begin{Resp}\Vsign The poor leaveth himself unto thee; Thou wilt be the helper of the fatherless.\end{Resp}
\begin{Resp}\Rsign For Thou alone beholdest mischief and spite.\end{Resp}
\switchcolumn*

\LA
\LessonHead{Lectio III}
\switchcolumn
\EN
\LessonHead{Lesson III}
\switchcolumn*

\LA
\Cite{1 Tim 1:12-16}
\switchcolumn
\EN
\Cite{1 Tim 1:12-16}
\switchcolumn*

\LA
\noindent Grátias ago ei, qui me confortávit Christo Jesu Dómino nostro, quia fidélem me existimávit, ponens in ministério: qui prius blasphémus fui, et persecútor, et contumeliósus: sed misericórdiam Dei consecútus sum, quia ignórans feci in incredulitáte. Superabundávit autem grátia Dómini nostri cum fide et dilectióne, quæ est in Christo Jesu. Fidélis sermo, et omni acceptióne dignus: quod Christus Jesus venit in hunc mundum peccatóres salvos fácere, quorum primus ego sum. Sed ídeo misericórdiam consecútus sum: ut in me primo osténderet Christus Jesus omnem patiéntiam ad informatiónem eórum, qui creditúri sunt illi, in vitam ætérnam.\par
\switchcolumn
\EN
\noindent I give thanks to him who hath strengthened me, even to Christ Jesus our Lord, for that he hath counted me faithful, putting me in the ministry; Who before was a blasphemer, and a persecutor, and contumelious. But I obtained the mercy of God, because I did it ignorantly in unbelief. Now the grace of our Lord hath abounded exceedingly with faith and love, which is in Christ Jesus. A faithful saying, and worthy of all acceptation, that Christ Jesus came into this world to save sinners, of whom I am the chief. But for this cause have I obtained mercy: that in me first Christ Jesus might shew forth all patience, for the information of them that shall believe in him unto life everlasting.\par
\switchcolumn*

\LA
\begin{Resp}\Rsign A dextris est mihi Dóminus, ne commóvear: \Star{} Propter hoc dilatátum est cor meum, et exsultávit lingua mea.\end{Resp}
\begin{Resp}\Vsign Dóminus pars hereditátis meæ, et cálicis mei.\end{Resp}
\begin{Resp}\Rsign Propter hoc dilatátum est cor meum, et exsultávit lingua mea.\end{Resp}
\begin{Resp}\Vsign Glória Patri, et Fílio, \Star{} et Spirítui Sancto.\end{Resp}
\begin{Resp}\Rsign Propter hoc dilatátum est cor meum, et exsultávit lingua mea.\end{Resp}
\switchcolumn
\EN
\begin{Resp}\Rsign The Lord is at my right hand, I shall never be moved. \Star{} Therefore my heart is glad, and my tongue rejoiceth.\end{Resp}
\begin{Resp}\Vsign The Lord is the portion of mine inheritance, and of my cup.\end{Resp}
\begin{Resp}\Rsign Therefore my heart is glad, and my tongue rejoiceth.\end{Resp}
\begin{Resp}\Vsign Glory be to the Father, and to the Son, \Star{} and to the Holy Ghost.\end{Resp}
\begin{Resp}\Rsign Therefore my heart is glad, and my tongue rejoiceth.\end{Resp}
\switchcolumn*

\LA
\LessonHead{Lectio IV}
\switchcolumn
\EN
\LessonHead{Lesson IV}
\switchcolumn*

\LA
\LessonTitle{Sermo sancti Augustíni Epíscopi}
\switchcolumn
\EN
\LessonTitle{From the Sermons of St. Augustine, Bishop (of Hippo.)}
\switchcolumn*

\LA
\Source{De Verbis Apóstoli, Sermo 8 sub initio}
\switchcolumn
\EN
\Source{On the words of the Apostles, 8}
\switchcolumn*

\LA
\noindent Humánus sermo, et omni acceptióne dignus, quia Christus Jesus venit in hunc mundum peccatóres salvos fácere. Atténde Evangélium: Venit enim Fílius hóminis quærére, et salváre, quod períerat. Si homo non periísset, Fílius hóminis non venísset. Ergo períerat homo: venit Deus homo, et invéntus est homo. Períerat homo per líberam voluntátem: venit Deus homo per grátiam liberatrícem.\par
\switchcolumn
\EN
\noindent This is a saying made for man, and worthy of all acceptation, that Christ Jesus came into this world to save sinners. Listen to the words of the Gospel The Son of man is come to seek, and to save that which was lost. If man had not been lost, the Son of man would not have come. Wherefore, man had been lost; God came made Man, and man was found; man had perished by his own free will God made Man came by grace which setteth free.\par
\switchcolumn*

\LA
\begin{Resp}\Rsign Notas mihi fecísti, Dómine, vias vitæ: \Star{} Adimplébis me lætítia cum vultu tuo: delectatiónes in déxtera tua usque in finem.\end{Resp}
\begin{Resp}\Vsign Tu es qui restítues hereditátem meam mihi.\end{Resp}
\begin{Resp}\Rsign Adimplébis me lætítia cum vultu tuo: delectatiónes in déxtera tua usque in finem.\end{Resp}
\switchcolumn
\EN
\begin{Resp}\Rsign O Lord, Thou hast shown me the path of life. \Star{} Thou shalt fill me with joy in thy presence, at thy right hand there are pleasures for evermore.\end{Resp}
\begin{Resp}\Vsign Thou art He That shalt restore mine inheritance unto me.\end{Resp}
\begin{Resp}\Rsign Thou shalt fill me with joy in thy presence, at thy right hand there are pleasures for evermore.\end{Resp}
\switchcolumn*

\LA
\LessonHead{Lectio V}
\switchcolumn
\EN
\LessonHead{Lesson V}
\switchcolumn*

\LA
\noindent Quæris, quid váleat ad malum líberum arbítrium? Recóle hóminem peccántem. Quæris, quid váleat ad auxílium Deus et homo? Atténde in eo grátiam liberántem. Nusquam pótuit sic osténdi, quantum váleat volúntas hóminis usurpáta per supérbiam, ad uténdum sine adjutório Dei: malum non pótuit plus, et maniféstius éxprimi, quam in hómine primo. Ecce perit primus homo, et ubi esset, nisi venísset secúndus homo? quia et ille homo, ídeo et iste homo; et ídeo humánus sermo.\par
\switchcolumn
\EN
\noindent Dost thou ask how free-will availeth to evil? Call to mind a sinner: Dost thou ask what God made Man availeth to help? Consider in Him the grace which setteth free. There is no example which so showeth what availeth the free will of man, when it is taken possession of by pride, to use it without God's help, of evil is there no greater and plainer example, than the first man. The first man fell and where had he been if the second Man had not come? As the first was man, so was the second Man, and therefore is this saying a saying. made for man.\par
\switchcolumn*

\LA
\begin{Resp}\Rsign Díligam te, Dómine, virtus mea: Dóminus firmaméntum meum, \Star{} Et refúgium meum.\end{Resp}
\begin{Resp}\Vsign Liberátor meus, Deus meus, adjútor meus.\end{Resp}
\begin{Resp}\Rsign Et refúgium meum.\end{Resp}
\switchcolumn
\EN
\begin{Resp}\Rsign I will love thee, O Lord, my strength; the Lord is my rock \Star{} And my fortress.\end{Resp}
\begin{Resp}\Vsign My Deliverer, my God, mine Helper.\end{Resp}
\begin{Resp}\Rsign And my fortress.\end{Resp}
\switchcolumn*

\LA
\LessonHead{Lectio VI}
\switchcolumn
\EN
\LessonHead{Lesson VI}
\switchcolumn*

\LA
\noindent Prorsus nusquam sic appáret benígnitas grátiæ, et liberálitas omnipoténtiæ Dei, quam in hómine mediatóre Dei et hóminum, hómine Christo Jesu. Quid enim dícimus, fratres mei? In fide cathólica nutrítis loquor, vel in pacem cathólicam lucrátis. Nóvimus et tenémus, mediatórem Dei et hóminum, hóminem Christum Jesum, in quantum homo erat, ejus esse natúræ, cujus et nos sumus. Non enim altérius natúræ caro nostra, et caro illíus: nec altérius natúræ ánima nostra, et ánima illíus. Hanc suscépit natúram, quam salvándam esse judicávit.\par
\switchcolumn
\EN
\noindent Neither is there any example which so showeth what availeth the tenderness of the grace and the abundance of the All-might of God, as the Man That is the Mediator between God and men, the Man Christ Jesus. For what do we say, my brethren? I speak to them that have been bred up in the Catholic Church, or who have been reconciled to that Church. We know and hold that the Mediator between God and men, the Man Christ Jesus, as touching upon His Manhood, is of the same nature as we. For our flesh is not of one nature, and His Flesh of another nature, neither our soul of one nature and His Soul of another nature. He took upon Himself the same nature which He had freely ordained to save.\par
\switchcolumn*

\LA
\begin{Resp}\Rsign Dómini est terra, et plenitúdo ejus: \Star{} Orbis terrárum, et univérsi qui hábitant in eo.\end{Resp}
\begin{Resp}\Vsign Ipse super mária fundávit eam, et super flúmina præparávit illam.\end{Resp}
\begin{Resp}\Rsign Orbis terrárum, et univérsi qui hábitant in eo.\end{Resp}
\begin{Resp}\Vsign Glória Patri, et Fílio, \Star{} et Spirítui Sancto.\end{Resp}
\begin{Resp}\Rsign Orbis terrárum, et univérsi qui hábitant in eo.\end{Resp}
\switchcolumn
\EN
\begin{Resp}\Rsign The earth is the Lord's, and the fulness thereof \Star{} The world, and they that dwell therein.\end{Resp}
\begin{Resp}\Vsign For He hath founded it upon the seas, and established it upon the floods.\end{Resp}
\begin{Resp}\Rsign The world, and they that dwell therein.\end{Resp}
\begin{Resp}\Vsign Glory be to the Father, and to the Son, \Star{} and to the Holy Ghost.\end{Resp}
\begin{Resp}\Rsign The world, and they that dwell therein.\end{Resp}
\switchcolumn*

\LA
\LessonHead{Lectio VII}
\switchcolumn
\EN
\LessonHead{Lesson VII}
\switchcolumn*

\LA
\LessonTitle{Léctio sancti Evangélii secúndum Matthǽum}
\switchcolumn
\EN
\LessonTitle{From the Holy Gospel according to Matthew}
\switchcolumn*

\LA
\Cite{Matt 13:24-30}
\switchcolumn
\EN
\Cite{Matt 13:24-30}
\switchcolumn*

\LA
\noindent In illo témpore: Dixit Jesus turbis parábolam hanc: Símile factum est regnum cælórum hómini, qui seminávit bonum semen in agro suo. Et réliqua.\par
\switchcolumn
\EN
\noindent At that time, Jesus said to the people a parable: The kingdom of heaven is likened to a man that sowed good seeds in his field. And so on.\par
\switchcolumn*

\LA
\LessonTitle{Homilía S. Augustíni Epíscopi}
\switchcolumn
\EN
\LessonTitle{Homily by St. Augustine, Bishop (of Hippo.)}
\switchcolumn*

\LA
\Source{Liber Quæst. Evang. in Matth. cap. 11, tom. 4}
\switchcolumn
\EN
\Source{Quest. Evang. Matth. xi, Bk. 4}
\switchcolumn*

\LA
\noindent Cum negligéntius ágerent præpósiti Ecclésiæ, aut cum dormitiónem mortis accíperent Apóstoli, venit diábolus, et superseminávit eos, quos malos fílios Dóminus interpretátur. Sed quǽritur: utrum hærétici sint, an male vivéntes cathólici? Possunt enim dici fílii mali étiam hærétici, quia ex eódem Evangélii sémine, et Christi nómine procreáti, pravis opiniónibus ad falsa dógmata convertúntur.\par
\switchcolumn
\EN
\noindent When the Shepherds of the Church wax careless, and since the Apostles sleep the sleep of death, cometh the devil, and soweth them whom the Lord calleth a seed of evil-doers. Now, are these seed of evil-doers the heretics, or Catholics of bad lives? It is possible to call even the heretics a seed of evil-doers because they have sprung up from the seed of the Gospel, and been begotten in the Name of Christ, though afterwards they have turned after crooked ways and lying doctrines.\par
\switchcolumn*

\LA
\begin{Resp}\Rsign Ad te, Dómine, levávi ánimam meam: \Star{} Deus meus, in te confído, non erubéscam.\end{Resp}
\begin{Resp}\Vsign Custódi ánimam meam, et éripe me.\end{Resp}
\begin{Resp}\Rsign Deus meus, in te confído, non erubéscam.\end{Resp}
\switchcolumn
\EN
\begin{Resp}\Rsign Unto thee, O Lord, do I lift up my soul. \Star{} O my God, I trust in thee, let me not be ashamed.\end{Resp}
\begin{Resp}\Vsign O keep my soul and deliver me.\end{Resp}
\begin{Resp}\Rsign O my God, I trust in thee, let me not be ashamed.\end{Resp}
\switchcolumn*

\LA
\LessonHead{Lectio VIII}
\switchcolumn
\EN
\LessonHead{Lesson VIII}
\switchcolumn*

\LA
\noindent Sed quod dicit eos in médio trítici seminátos, quasi vidéntur illi significári, qui uníus communiónis sunt. Verúmtamen quóniam Dóminus agrum ipsum, non Ecclésiam, sed hunc mundum interpretátus est: bene intelligúntur hærétici, quia non societáte uníus Ecclésiæ, vel uníus fídei, sed societáte solíus nóminis christiáni in hoc mundo permiscéntur bonis. At illi, qui in eádem fide mali sunt, pálea pótius quam zizánia reputántur: quia pálea étiam fundaméntum ipsum habet cum fruménto, radicémque commúnem.\par
\switchcolumn
\EN
\noindent But whereas it is written that they were sown in the midst of the wheat, we ought haply to understand that they are of one communion with the righteous. Nevertheless, forasmuch as the Lord saith, The field is the world, (and not, the Church,) we may well understand that the seed of evil doers are the heretics, since in this world they are mingled together with the good, not in one common Communion, but only under one common name of Christian. But they which are of one faith with the good seed, and yet are themselves worthless, may more fitly be likened to straw than to tares, since the straw springeth from one soil and one root with the good ear.\par
\switchcolumn*

\LA
\begin{Resp}\Rsign Duo Séraphim clamábant alter ad álterum: \Star{} Sanctus, sanctus, sanctus Dóminus Deus Sábaoth: * Plena est omnis terra glória ejus.\end{Resp}
\begin{Resp}\Vsign Tres sunt qui testimónium dant in cælo: Pater, Verbum, et Spíritus Sanctus: et hi tres unum sunt.\end{Resp}
\begin{Resp}\Rsign Sanctus, sanctus, sanctus Dóminus Deus Sábaoth.\end{Resp}
\begin{Resp}\Vsign Glória Patri, et Fílio, \Star{} et Spirítui Sancto.\end{Resp}
\begin{Resp}\Rsign Plena est omnis terra glória ejus.\end{Resp}
\switchcolumn
\EN
\begin{Resp}\Rsign One Seraph cried unto another: \Star{} Holy, Holy, Holy is the Lord God of hosts; the whole earth is full of His glory.\end{Resp}
\begin{Resp}\Vsign There are Three That bear record in heaven: the Father, the Word, and the Holy Ghost, and these Three are One.\end{Resp}
\begin{Resp}\Rsign Holy, Holy, Holy is the Lord God of hosts.\end{Resp}
\begin{Resp}\Vsign Glory be to the Father, and to the Son, \Star{} and to the Holy Ghost.\end{Resp}
\begin{Resp}\Rsign The whole earth is full of His glory.\end{Resp}
\switchcolumn*

\LA
\LessonHead{Lectio IX}
\switchcolumn
\EN
\LessonHead{Lesson IX}
\switchcolumn*

\LA
\noindent In illa plane sagéna, qua concludúntur et mali et boni pisces, non absúrde mali cathólici intelligúntur. Aliud est enim mare, quod magis mundum istum signíficat: áliud sagéna, quæ uníus fídei, vel uníus Ecclésiæ communiónem vidétur osténdere. Inter hæréticos et malos cathólicos hoc ínterest, quod hærétici falsa credunt: illi autem, vera credéntes, non vivunt ita ut credunt.\par
\switchcolumn
\EN
\noindent However, as touching the net cast into the sea, and enclosing a great multitude of fishes, both bad and good, we may well understand that by the bad are meant Catholics of bad lives. For the sea is one thing whereby we may understand to be signified the world; and the net another, which seemeth to signify our faith, or the Communion of one Church. Between heretics and sinful Catholics there is this difference, that heretics believe a lie, and sinful Catholics believe the truth, but live not as they believe.\par
\switchcolumn*

\LA
\TeDeum{Te Deum}
\switchcolumn
\EN
\TeDeum{Te Deum}
\switchcolumn*

\end{paracol}

\Office{Feria II infra Hebdomadam V post Epiphaniam}{Monday of the Fifth Week after Epiphany}{Feria}
\addcontentsline{toc}{subsection}{Feria II infra Hebdomadam V post Epiphaniam\enspace\textit{Monday of the Fifth Week after Epiphany}}
\begin{paracol}{2}
\LA
\LessonHead{Lectio I}
\switchcolumn
\EN
\LessonHead{Lesson I}
\switchcolumn*

\LA
\LessonTitle{De Epístola prima ad Timótheum}
\switchcolumn
\EN
\LessonTitle{Lesson from the first letter of St. Paul the Apostle to Timothy}
\switchcolumn*

\LA
\Cite{1 Tim 3:1-7}
\switchcolumn
\EN
\Cite{1 Tim 3:1-7}
\switchcolumn*

\LA
\noindent Fidélis sermo: si quis episcopátum desíderat, bonum opus desíderat. Opórtet ergo epíscopum irreprehensíbilem esse, uníus uxóris virum, sóbrium, prudéntem, ornátum, pudícum, hospitálem, doctórem, non vinoléntum, non percussórem, sed modéstum: non litigiósum, non cúpidum, sed suæ dómui bene præpósitum: fílios habéntem súbditos cum omni castitáte. Si quis autem dómui suæ præésse nescit, quómodo ecclésiæ Dei diligéntiam habébit? Non neóphytum: ne in supérbiam elátus, in judícium íncidat diáboli. Opórtet autem illum et testimónium habére bonum ab iis qui foris sunt, ut non in oppróbrium íncidat, et in láqueum diáboli.\par
\switchcolumn
\EN
\noindent A faithful saying: if a man desire the office of a bishop, he desireth a good work. It behoveth therefore a bishop to be blameless, the husband of one wife, sober, prudent, of good behaviour, chaste, given to hospitality, a teacher, Not given to wine, no striker, but modest, not quarrelsome, not covetous, but One that ruleth well his own house, having his children in subjection with all chastity. But if a man know not how to rule his own house, how shall he take care of the church of God? Not a neophyte: lest being puffed up with pride, he fall into the judgment of the devil. Moreover he must have a good testimony of them who are without: lest he fall into reproach and the snare of the devil.\par
\switchcolumn*

\LA
\begin{Resp}\Rsign Quam magna multitúdo dulcédinis tuæ, Dómine, \Star{} Quam abscondísti timéntibus te!\end{Resp}
\begin{Resp}\Vsign Et perfecísti eis qui sperant in te, Dómine, in conspéctu filiórum hóminum.\end{Resp}
\begin{Resp}\Rsign Quam abscondísti timéntibus te!\end{Resp}
\switchcolumn
\EN
\begin{Resp}\Rsign O how great is thy goodness, O Lord \Star{} Which Thou hast laid up for them that fear thee!\end{Resp}
\begin{Resp}\Vsign Which Thou, O Lord, hast wrought for them that trust in thee before the sons of men!\end{Resp}
\begin{Resp}\Rsign Which Thou hast laid up for them that fear thee!\end{Resp}
\switchcolumn*

\LA
\LessonHead{Lectio II}
\switchcolumn
\EN
\LessonHead{Lesson II}
\switchcolumn*

\LA
\Cite{1 Tim 3:8-13}
\switchcolumn
\EN
\Cite{1 Tim 3:8-13}
\switchcolumn*

\LA
\noindent Diáconos simíliter pudícos, non bilíngues, non multo vino déditos, non turpe lucrum sectántes: habéntes mystérium fídei in consciéntia pura. Et hi autem probéntur primum: et sic minístrent, nullum crimen habéntes. Mulíeres simíliter pudícas, non detrahéntes, sóbrias, fidéles in ómnibus. Diáconi sint uníus uxóris viri, qui fíliis suis bene præsint, et suis dómibus. Qui enim bene ministráverint, gradum bonum sibi acquírent, et multam fidúciam in fide, quæ est in Christo Jesu.\par
\switchcolumn
\EN
\noindent Deacons in like manner chaste, not double tongued, not given to much wine, not greedy of filthy lucre: Holding the mystery of faith in a pure conscience. And let these also first be proved: and so let them minister, having no crime. The women in like manner chaste, not slanderers, but sober, faithful in all things. Let deacons be the husbands of one wife: who rule well their children, and their own houses. For they that have ministered well, shall purchase to themselves a good degree, and much confidence in the faith which is in Christ Jesus.\par
\switchcolumn*

\LA
\begin{Resp}\Rsign Adjútor meus esto, Deus: \Star{} Ne derelínquas me.\end{Resp}
\begin{Resp}\Vsign Neque despícias me, Deus, salutáris meus.\end{Resp}
\begin{Resp}\Rsign Ne derelínquas me.\end{Resp}
\switchcolumn
\EN
\begin{Resp}\Rsign O God, be Thou my helper. \Star{} Neither leave me.\end{Resp}
\begin{Resp}\Vsign Nor forsake me O God of my salvation.\end{Resp}
\begin{Resp}\Rsign Neither leave me.\end{Resp}
\switchcolumn*

\LA
\LessonHead{Lectio III}
\switchcolumn
\EN
\LessonHead{Lesson III}
\switchcolumn*

\LA
\Cite{1 Tim 3:14-16; 4:1}
\switchcolumn
\EN
\Cite{1 Tim 3:14-16;4:1}
\switchcolumn*

\LA
\noindent Hæc tibi scribo, sperans me ad te veníre cito: si autem tardávero, ut scias quómodo opórteat te in domo Dei conversári, quæ est ecclésia Dei vivi, colúmna et firmaméntum veritátis. Et maniféste magnum est pietátis sacraméntum, quod manifestátum est in carne, justificátum est in spíritu, appáruit ángelis, prædicátum est géntibus, créditum est in mundo, assúmptum est in glória. Spíritus autem maniféste dicit, quia in novíssimis tempóribus discédent quidam a fide, attendéntes spirítibus erróris, et doctrínis dæmoniórum.\par
\switchcolumn
\EN
\noindent These things I write to thee, hoping that I shall come to thee shortly. But if I tarry long, that thou mayest know how thou oughtest to behave thyself in the house of God, which is the church of the living God, the pillar and ground of the truth. And evidently great is the mystery of godliness, which was manifested in the flesh, was justified in the spirit, appeared unto angels, hath been preached unto the Gentiles, is believed in the world, is taken up in glory. Now the Spirit manifestly saith, that in the last times some shall depart from the faith, giving heed to spirits of error, and doctrines of devils,\par
\switchcolumn*

\LA
\begin{Resp}\Rsign Benedícam Dóminum in omni témpore: \Star{} Semper laus ejus in ore meo.\end{Resp}
\begin{Resp}\Vsign In Dómino laudábitur ánima mea, áudiant mansuéti, et læténtur.\end{Resp}
\begin{Resp}\Rsign Semper laus ejus in ore meo.\end{Resp}
\begin{Resp}\Vsign Glória Patri, et Fílio, \Star{} et Spirítui Sancto.\end{Resp}
\begin{Resp}\Rsign Semper laus ejus in ore meo.\end{Resp}
\switchcolumn
\EN
\begin{Resp}\Rsign I will bless the Lord at all times. \Star{} His praise shall continually be in my mouth.\end{Resp}
\begin{Resp}\Vsign My soul shall make her boast in the Lord; the humble shall hear thereof and be glad.\end{Resp}
\begin{Resp}\Rsign His praise shall continually be in my mouth.\end{Resp}
\begin{Resp}\Vsign Glory be to the Father, and to the Son, \Star{} and to the Holy Ghost.\end{Resp}
\begin{Resp}\Rsign His praise shall continually be in my mouth.\end{Resp}
\switchcolumn*

\end{paracol}

\Office{Feria III infra Hebdomadam V post Epiphaniam}{Tuesday of the Fifth Week after Epiphany}{Feria}
\addcontentsline{toc}{subsection}{Feria III infra Hebdomadam V post Epiphaniam\enspace\textit{Tuesday of the Fifth Week after Epiphany}}
\begin{paracol}{2}
\LA
\LessonHead{Lectio I}
\switchcolumn
\EN
\LessonHead{Lesson I}
\switchcolumn*

\LA
\LessonTitle{Incipit Epístola secúnda beáti Pauli Apóstoli ad Timótheum}
\switchcolumn
\EN
\LessonTitle{Lesson from the secons letter of St. Paul the Apostle to Timothy}
\switchcolumn*

\LA
\Cite{2 Tim 1:1-5}
\switchcolumn
\EN
\Cite{2 Tim 1:1-5}
\switchcolumn*

\LA
\noindent Paulus Apóstolus Jesu Christi per voluntátem Dei, secúndum promissiónem vitæ, quæ est in Christo Jesu, Timótheo caríssimo fílio: grátia, misericórdia, pax a Deo Patre, et Christo Jesu Dómino nostro. Grátias ago Deo, cui sérvio a progenitóribus in consciéntia pura, quod sine intermissióne hábeam tui memóriam in oratiónibus meis, nocte ac die desíderans te vidére, memor lacrimárum tuárum, ut gáudio ímplear, recordatiónem accípiens ejus fídei, quæ est in te non ficta, quæ et habitávit primum in ávia tua Loíde, et matre tua Euníce, certus sum autem quod et in te.\par
\switchcolumn
\EN
\noindent Paul, an apostle of Jesus Christ, by the will of God, according to the promise of life, which is in Christ Jesus. To Timothy my dearly beloved son, grace, mercy, and peace, from God the Father, and from Christ Jesus our Lord. I give thanks to God, whom I serve from my forefathers with a pure conscience, that without ceasing, I have a remembrance of thee in my prayers, night and day. Desiring to see thee, being mindful of thy tears, that I may be filled with joy, Calling to mind that faith which is in thee unfeigned, which also dwelt first in thy grandmother Lois, and in thy mother Eunice, and I am certain that in thee also.\par
\switchcolumn*

\LA
\begin{Resp}\Rsign Auribus pércipe, Deus, lácrimas meas: ne síleas a me, remítte mihi: \Star{} Quóniam íncola ego sum apud te, et peregrínus.\end{Resp}
\begin{Resp}\Vsign Compláceat tibi, ut erípias me: Dómine, ad adjuvándum me festína.\end{Resp}
\begin{Resp}\Rsign Quóniam íncola ego sum apud te, et peregrínus.\end{Resp}
\switchcolumn
\EN
\begin{Resp}\Rsign O God, give ear unto my tears; hold not thy peace, but, O, spare me! \Star{} For I am a stranger with thee, and a sojourner.\end{Resp}
\begin{Resp}\Vsign Be pleased, O Lord, to deliver me; O Lord, look upon me to help me.\end{Resp}
\begin{Resp}\Rsign For I am a stranger with thee, and a sojourner.\end{Resp}
\switchcolumn*

\LA
\LessonHead{Lectio II}
\switchcolumn
\EN
\LessonHead{Lesson II}
\switchcolumn*

\LA
\Cite{2 Tim 1:6-9}
\switchcolumn
\EN
\Cite{2 Tim 1:6-9}
\switchcolumn*

\LA
\noindent Propter quam causam admóneo te ut resúscites grátiam Dei, quæ est in te per impositiónem mánuum meárum. Non enim dedit nobis Deus spíritum timóris: sed virtútis, et dilectiónis, et sobrietátis. Noli ítaque erubéscere testimónium Dómini nostri, neque me vinctum ejus: sed collabóra Evangélio secúndum virtútem Dei: qui nos liberávit, et vocávit vocatióne sua sancta, non secúndum ópera nostra, sed secúndum propósitum suum, et grátiam, quæ data est nobis in Christo Jesu ante témpora sæculária.\par
\switchcolumn
\EN
\noindent For which cause I admonish thee, that thou stir up the grace of God which is in thee, by the imposition of my hands. For God hath not given us the spirit of fear: but of power, and of love, and of sobriety. Be not thou therefore ashamed of the testimony of our Lord, nor of me his prisoner: but labour with the gospel, according to the power of God, Who hath delivered us and called us by his holy calling, not according to our works, but according to his own purpose and grace, which was given us in Christ Jesus before the times of the world.\par
\switchcolumn*

\LA
\begin{Resp}\Rsign Státuit Dóminus supra petram pedes meos, et diréxit gressus meos Deus: \Star{} Et misit in os meum cánticum novum.\end{Resp}
\begin{Resp}\Vsign Exaudívit preces meas: et edúxit me de lacu misériæ.\end{Resp}
\begin{Resp}\Rsign Et misit in os meum cánticum novum.\end{Resp}
\switchcolumn
\EN
\begin{Resp}\Rsign The Lord hath set my feet upon a rock, and ordered my goings. \Star{} And He hath put a new song in my mouth.\end{Resp}
\begin{Resp}\Vsign He heard my cry He brought me up also out of an horrible pit.\end{Resp}
\begin{Resp}\Rsign And He hath put a new song in my mouth.\end{Resp}
\switchcolumn*

\LA
\LessonHead{Lectio III}
\switchcolumn
\EN
\LessonHead{Lesson III}
\switchcolumn*

\LA
\Cite{2 Tim 1:10-13}
\switchcolumn
\EN
\Cite{2 Tim 1:10-13}
\switchcolumn*

\LA
\noindent Manifestáta est autem nunc per illuminatiónem Salvatóris nostri Jesu Christi, qui destrúxit quidem mortem, illuminávit autem vitam, et incorruptiónem per Evangélium: in quo pósitus sum ego prædicátor, et Apóstolus, et magíster géntium. Ob quam causam étiam hæc pátior, sed non confúndor. Scio enim cui crédidi, et certus sum quia potens est depósitum meum serváre in illum diem. Formam habe sanórum verbórum, quæ a me audísti in fide, et in dilectióne in Christo Jesu.\par
\switchcolumn
\EN
\noindent But is now made manifest by the illumination of our Saviour Jesus Christ, who hath destroyed death, and hath brought to light life and incorruption by the gospel: Wherein I am appointed a preacher, and an apostle, and teacher of the Gentiles. For which cause I also suffer these things: but I am not ashamed. For I know whom I have believed, and I am certain that he is able to keep that which I have committed unto him, against that day. Hold the form of sound words, which thou hast heard of me in faith, and in the love which is in Christ Jesus.\par
\switchcolumn*

\LA
\begin{Resp}\Rsign Ego dixi, Dómine, miserére mei: \Star{} Sana ánimam meam, quia peccávi tibi.\end{Resp}
\begin{Resp}\Vsign Ab ómnibus iniquitátibus meis éripe me, Dómine.\end{Resp}
\begin{Resp}\Rsign Sana ánimam meam, quia peccávi tibi.\end{Resp}
\begin{Resp}\Vsign Glória Patri, et Fílio, \Star{} et Spirítui Sancto.\end{Resp}
\begin{Resp}\Rsign Sana ánimam meam, quia peccávi tibi.\end{Resp}
\switchcolumn
\EN
\begin{Resp}\Rsign I said: Lord, be merciful unto me. \Star{} Heal my soul, for I have sinned against thee.\end{Resp}
\begin{Resp}\Vsign Deliver me from all mine iniquities, O Lord.\end{Resp}
\begin{Resp}\Rsign Heal my soul, for I have sinned against thee.\end{Resp}
\begin{Resp}\Vsign Glory be to the Father, and to the Son, \Star{} and to the Holy Ghost.\end{Resp}
\begin{Resp}\Rsign Heal my soul, for I have sinned against thee.\end{Resp}
\switchcolumn*

\end{paracol}

\Office{Feria IV infra Hebdomadam V post Epiphaniam}{Wednesday of the Fifth Week after Epiphany}{Feria}
\addcontentsline{toc}{subsection}{Feria IV infra Hebdomadam V post Epiphaniam\enspace\textit{Wednesday of the Fifth Week after Epiphany}}
\begin{paracol}{2}
\LA
\LessonHead{Lectio I}
\switchcolumn
\EN
\LessonHead{Lesson I}
\switchcolumn*

\LA
\LessonTitle{De Epístola secúnda ad Timótheum}
\switchcolumn
\EN
\LessonTitle{Lesson from the second letter of St. Paul the Apostle to Timothy}
\switchcolumn*

\LA
\Cite{2 Tim 3:1-5}
\switchcolumn
\EN
\Cite{2 Tim 3:1-5}
\switchcolumn*

\LA
\noindent Hoc autem scito, quod in novíssimis diébus instábunt témpora periculósa: erunt hómines seípsos amántes, cúpidi, eláti, supérbi, blasphémi, paréntibus non obediéntes, ingráti, scelésti, sine affectióne, sine pace, criminatóres, incontinéntes, immítes, sine benignitáte, proditóres, protérvi, túmidi, et voluptátum amatóres magis quam Dei: habéntes spéciem quidem pietátis, virtútem autem ejus abnegántes. Et hos devíta:\par
\switchcolumn
\EN
\noindent Know also this, that, in the last days, shall come dangerous times. Men shall be lovers of themselves, covetous, haughty, proud, blasphemers, disobedient to parents, ungrateful, wicked, Without affection, without peace, slanderers, incontinent, unmerciful, without kindness, Traitors, stubborn, puffed up, and lovers of pleasures more than of God: Having an appearance indeed of godliness, but denying the power thereof. Now these avoid.\par
\switchcolumn*

\LA
\begin{Resp}\Rsign Ne perdíderis me cum iniquitátibus meis: \Star{} Neque in finem irátus resérves mala mea.\end{Resp}
\begin{Resp}\Vsign Non intres in judícium cum servo tuo, Dómine.\end{Resp}
\begin{Resp}\Rsign Neque in finem irátus resérves mala mea.\end{Resp}
\switchcolumn
\EN
\begin{Resp}\Rsign Cut me not off in the midst of my sins. \Star{} Nor keep thy wrath against me for my latter end.\end{Resp}
\begin{Resp}\Vsign Enter not into judgment with thy servant, O Lord.\end{Resp}
\begin{Resp}\Rsign Nor keep thy wrath against me for my latter end.\end{Resp}
\switchcolumn*

\LA
\LessonHead{Lectio II}
\switchcolumn
\EN
\LessonHead{Lesson II}
\switchcolumn*

\LA
\Cite{2 Tim 3:6-9}
\switchcolumn
\EN
\Cite{2 Tim 3:6-9}
\switchcolumn*

\LA
\noindent ex his enim sunt qui pénetrant domos, et captívas ducunt muliérculas onerátas peccátis, quæ ducúntur váriis desidériis: semper discéntes, et nunquam ad sciéntiam veritátis perveniéntes. Quemádmodum autem Jannes et Mambres restitérunt Móysi: ita et hi resístunt veritáti, hómines corrúpti mente, réprobi circa fidem; sed ultra non profícient: insipiéntia enim eórum manifésta erit ómnibus, sicut et illórum fuit.\par
\switchcolumn
\EN
\noindent For of these sort are they who creep into houses, and lead captive silly women laden with sins, who are led away with diverse desires: Ever learning, and never attaining to the knowledge of the truth. Now as Jannes and Mambres resisted Moses, so these also resist the truth, men corrupted in mind, reprobate concerning the faith. But they shall proceed no farther; for their folly shall be manifest to all men, as theirs also was.\par
\switchcolumn*

\LA
\begin{Resp}\Rsign Parátum cor meum, Deus, parátum cor meum: \Star{} Cantábo et psalmum dicam Dómino.\end{Resp}
\begin{Resp}\Vsign Exsúrge, glória mea, exsúrge, psaltérium et cíthara, exsúrgam dilúculo.\end{Resp}
\begin{Resp}\Rsign Cantábo et psalmum dicam Dómino.\end{Resp}
\switchcolumn
\EN
\begin{Resp}\Rsign My heart is ready, O God, my heart is ready. \Star{} I will sing and give praise to the Lord.\end{Resp}
\begin{Resp}\Vsign Awake up, my glory, awake, psaltery and harp! I will awake early.\end{Resp}
\begin{Resp}\Rsign I will sing and give praise to the Lord.\end{Resp}
\switchcolumn*

\LA
\LessonHead{Lectio III}
\switchcolumn
\EN
\LessonHead{Lesson III}
\switchcolumn*

\LA
\Cite{2 Tim 3:10-13}
\switchcolumn
\EN
\Cite{2 Tim 3:10-13}
\switchcolumn*

\LA
\noindent Tu autem assecútus es meam doctrínam, institutiónem, propósitum, fidem, longanimitátem, dilectiónem, patiéntiam, persecutiónes, passiónes: quália mihi facta sunt Antiochíæ, Icónii, et Lystris: quales persecutiónes sustínui, et ex ómnibus erípuit me Dóminus. Et omnes, qui pie volunt vívere in Christo Jesu, persecutiónem patiéntur. Mali autem hómines et seductóres profícient in pejus, errántes, et in errórem mitténtes.\par
\switchcolumn
\EN
\noindent But thou hast fully known my doctrine, manner of life, purpose, faith, longsuffering, love, patience, persecutions, afflictions: such as came upon me at Antioch, at Iconium, and at Lystra: what persecutions I endured, and out of them all the Lord delivered me. And all that will live godly in Christ Jesus, shall suffer persecution. But evil men and seducers shall grow worse and worse: erring, and driving into error.\par
\switchcolumn*

\LA
\begin{Resp}\Rsign Adjútor meus, tibi psallam, quia, Deus, suscéptor meus es: \Star{} Deus meus, misericórdia mea.\end{Resp}
\begin{Resp}\Vsign Lætábor, et exsultábo in te, psallam nómini tuo, Altíssime.\end{Resp}
\begin{Resp}\Rsign Deus meus, misericórdia mea.\end{Resp}
\begin{Resp}\Vsign Glória Patri, et Fílio, \Star{} et Spirítui Sancto.\end{Resp}
\begin{Resp}\Rsign Deus meus, misericórdia mea.\end{Resp}
\switchcolumn
\EN
\begin{Resp}\Rsign Unto thee, O my Strength, will I sing, for God is my defence, \Star{} The God of my mercy.\end{Resp}
\begin{Resp}\Vsign I will be glad and rejoice in thee, I will sing praise to thy Name, O Thou Most High.\end{Resp}
\begin{Resp}\Rsign The God of my mercy.\end{Resp}
\begin{Resp}\Vsign Glory be to the Father, and to the Son, \Star{} and to the Holy Ghost.\end{Resp}
\begin{Resp}\Rsign The God of my mercy.\end{Resp}
\switchcolumn*

\end{paracol}

\Office{Feria V infra Hebdomadam V post Epiphaniam}{Thursday of the Fifth Week after Epiphany}{Feria}
\addcontentsline{toc}{subsection}{Feria V infra Hebdomadam V post Epiphaniam\enspace\textit{Thursday of the Fifth Week after Epiphany}}
\begin{paracol}{2}
\LA
\LessonHead{Lectio I}
\switchcolumn
\EN
\LessonHead{Lesson I}
\switchcolumn*

\LA
\LessonTitle{Incipit Epístola beáti Pauli Apóstoli ad Titum}
\switchcolumn
\EN
\LessonTitle{Lesson from the letter of St. Paul the Apostle to Titus}
\switchcolumn*

\LA
\Cite{Titus 1:1-4}
\switchcolumn
\EN
\Cite{Titus 1:1-4}
\switchcolumn*

\LA
\noindent Paulus servus Dei, Apóstolus autem Jesu Christi secúndum fidem electórum Dei, et agnitiónem veritátis, quæ secúndum pietátem est in spem vitæ ætérnæ, quam promísit qui non mentítur, Deus, ante témpora sæculária: manifestávit autem tempóribus suis verbum suum in prædicatióne, quæ crédita est mihi secúndum præcéptum Salvatóris nostri Dei: Tito dilécto fílio secúndum commúnem fidem, grátia, et pax a Deo Patre, et Christo Jesu Salvatóre nostro.\par
\switchcolumn
\EN
\noindent Paul, a servant of God, and an apostle of Jesus Christ, according to the faith of the elect of God and the acknowledging of the truth, which is according to godliness: Unto the hope of life everlasting, which God, who lieth not, hath promised before the times of the world: But hath in due times manifested his word in preaching, which is committed to me according to the commandment of God our Saviour: To Titus my beloved son, according to the common faith, grace and peace from God the Father, and from Christ Jesus our Saviour.\par
\switchcolumn*

\LA
\begin{Resp}\Rsign Deus, in te sperávi, Dómine, non confúndar in ætérnum: in justítia tua líbera me, \Star{} Et éripe me.\end{Resp}
\begin{Resp}\Vsign Inclína ad me aurem tuam, et salva me.\end{Resp}
\begin{Resp}\Rsign Et éripe me.\end{Resp}
\switchcolumn
\EN
\begin{Resp}\Rsign In thee, O God, do I put my trust; let me never be put to confusion, O Lord deliver me in thy righteousness, \Star{} And cause me to escape.\end{Resp}
\begin{Resp}\Vsign Incline thine ear unto me, deliver me speedily.\end{Resp}
\begin{Resp}\Rsign And cause me to escape.\end{Resp}
\switchcolumn*

\LA
\LessonHead{Lectio II}
\switchcolumn
\EN
\LessonHead{Lesson II}
\switchcolumn*

\LA
\Cite{Titus 1:5-9}
\switchcolumn
\EN
\Cite{Titus 1:5-9}
\switchcolumn*

\LA
\noindent Hujus rei grátia relíqui te Cretæ, ut ea quæ desunt, córrigas, et constítuas per civitátes presbýteros, sicut et ego dispósui tibi, si quis sine crímine est, uníus uxóris vir, fílios habens fidéles, non in accusatióne luxúriæ, aut non súbditos. Opórtet enim epíscopum sine crímine esse, sicut Dei dispensatórem: non supérbum, non iracúndum, non vinoléntum, non percussórem, non turpis lucri cúpidum: sed hospitálem, benígnum, sóbrium, justum, sanctum, continéntem, amplecténtem eum, qui secúndum doctrínam est, fidélem sermónem: ut potens sit exhortári in doctrína sana, et eos qui contradícunt, argúere.\par
\switchcolumn
\EN
\noindent For this cause I left thee in Crete, that thou shouldest set in order the things that are wanting, and shouldest ordain priests in every city, as I also appointed thee: If any be without crime, the husband of one wife, having faithful children, not accused of riot, or unruly. For a bishop must be without crime, as the steward of God: not proud, not subject to anger, not given to wine, no striker, not greedy of filthy lucre: But given to hospitality, gentle, sober, just, holy, continent: Embracing that faithful word which is according to doctrine, that he may be able to exhort in sound doctrine, and to convince the gainsayers.\par
\switchcolumn*

\LA
\begin{Resp}\Rsign Repleátur os meum laude tua, ut hymnum dicam glóriæ tuæ, tota die magnitúdinem tuam: noli me proícere in témpore senectútis: \Star{} Dum defécerit in me virtus mea, ne derelínquas me.\end{Resp}
\begin{Resp}\Vsign Gaudébunt lábia mea cum cantávero tibi.\end{Resp}
\begin{Resp}\Rsign Dum defécerit in me virtus mea, ne derelínquas me.\end{Resp}
\switchcolumn
\EN
\begin{Resp}\Rsign Let my mouth be filled with thy praise, that I may sing of thy glory, all the day long of thy greatness. Cast me not off in the time of old age; \Star{} Forsake me not when my strength faileth.\end{Resp}
\begin{Resp}\Vsign My lips shall be fain when I sing unto thee.\end{Resp}
\begin{Resp}\Rsign Forsake me not when my strength faileth.\end{Resp}
\switchcolumn*

\LA
\LessonHead{Lectio III}
\switchcolumn
\EN
\LessonHead{Lesson III}
\switchcolumn*

\LA
\Cite{Titus 1:10-15}
\switchcolumn
\EN
\Cite{Titus 1:10-15}
\switchcolumn*

\LA
\noindent Sunt enim multi étiam inobediéntes, vaníloqui, et seductóres: máxime qui de circumcisióne sunt: quos opórtet redárgui: qui univérsas domos subvértunt, docéntes quæ non opórtet, turpis lucri grátia. Dixit quidam ex illis, próprius ipsórum prophéta: Creténses semper mendáces, malæ béstiæ, ventres pigri. Testimónium hoc verum est. Quam ob causam íncrepa illos dure, ut sani sint in fide, non intendéntes Judáicis fábulis, et mandátis hóminum, aversántium se a veritáte. Omnia munda mundis: coinquinátis autem et infidélibus, nihil est mundum.\par
\switchcolumn
\EN
\noindent For there are also many disobedient, vain talkers, and seducers: especially they who are of the circumcision: Who must be reproved, who subvert whole houses, teaching things which they ought not, for filthy lucre's sake. One of them a prophet of their own, said, The Cretians are always liars, evil beasts, slothful bellies. This testimony is true. Wherefore rebuke them sharply, that they may be sound in the faith; Not giving heed to Jewish fables and commandments of men, who turn themselves away from the truth. All things are clean to the clean: but to them that are defiled, and to unbelievers, nothing is clean: but both their mind and their conscience are defiled.\par
\switchcolumn*

\LA
\begin{Resp}\Rsign Gaudébunt lábia mea cum cantávero tibi: \Star{} Et ánima mea, quam redemísti, Dómine.\end{Resp}
\begin{Resp}\Vsign Sed et lingua mea meditábitur justítiam tuam, tota die laudem tuam.\end{Resp}
\begin{Resp}\Rsign Et ánima mea, quam redemísti, Dómine.\end{Resp}
\begin{Resp}\Vsign Glória Patri, et Fílio, \Star{} et Spirítui Sancto.\end{Resp}
\begin{Resp}\Rsign Et ánima mea, quam redemísti, Dómine.\end{Resp}
\switchcolumn
\EN
\begin{Resp}\Rsign My lips shall be fain when I sing unto thee; \Star{} And my soul, which Thou, O Lord, hast redeemed.\end{Resp}
\begin{Resp}\Vsign My tongue shall also talk of thy righteousness, all the day long of thy praise.\end{Resp}
\begin{Resp}\Rsign And my soul, which Thou, O Lord, hast redeemed.\end{Resp}
\begin{Resp}\Vsign Glory be to the Father, and to the Son, \Star{} and to the Holy Ghost.\end{Resp}
\begin{Resp}\Rsign And my soul, which Thou, O Lord, hast redeemed.\end{Resp}
\switchcolumn*

\end{paracol}

\Office{Feria VI infra Hebdomadam V post Epiphaniam}{Friday of the Fifth Week after Epiphany}{Feria}
\addcontentsline{toc}{subsection}{Feria VI infra Hebdomadam V post Epiphaniam\enspace\textit{Friday of the Fifth Week after Epiphany}}
\begin{paracol}{2}
\LA
\RefLine{Lectiones I–III: de Scriptura occurrente.}
\switchcolumn
\EN
\RefLine{Lessons I–III: from the Scripture of the day.}
\switchcolumn*

\LA
\LessonHead{Lectio I}
\switchcolumn
\EN
\LessonHead{Lesson I}
\switchcolumn*

\LA
\LessonTitle{De epístola ad Titum}
\switchcolumn
\EN
\LessonTitle{Lesson from the letter of St. Paul the Apostle to Titus}
\switchcolumn*

\LA
\Cite{Titus 2:15; 3:1-2}
\switchcolumn
\EN
\Cite{Titus 2:15;3:1-2}
\switchcolumn*

\LA
\noindent Hæc lóquere, et exhortáre, et árgue cum omni império. Nemo te contémnat. Admone illos princípibus, et potestátibus súbditos esse, dicto obedíre, ad omne opus bonum parátos esse: néminem blasphemáre, non litigiósos esse, sed modéstos, omnem ostendéntes mansuetúdinem ad omnes hómines.\par
\switchcolumn
\EN
\noindent These things speak, and exhort and rebuke with all authority. Let no man despise thee. Admonish them to be subject to princes and powers, to obey at a word, to be ready to every good work. To speak evil of no man, not to be litigious, but gentle: shewing all mildness towards all men.\par
\switchcolumn*

\LA
\begin{Resp}\Rsign Confitébor tibi, Dómine Deus, in toto corde meo, et honorificábo nomen tuum in ætérnum: \Star{} Quia misericórdia tua, Dómine, magna est super me.\end{Resp}
\begin{Resp}\Vsign Deus meus es tu, et confitébor tibi: Deus meus es tu, et exaltábo te.\end{Resp}
\begin{Resp}\Rsign Quia misericórdia tua, Dómine, magna est super me.\end{Resp}
\switchcolumn
\EN
\begin{Resp}\Rsign I will praise thee, O Lord my God, with all my heart, and I will glorify thy Name for evermore. \Star{} For great is thy mercy, O Lord, toward me.\end{Resp}
\begin{Resp}\Vsign Thou art my God, and I will praise thee; Thou art my God, and I will exalt thee.\end{Resp}
\begin{Resp}\Rsign For great is thy mercy, O Lord, toward me.\end{Resp}
\switchcolumn*

\LA
\LessonHead{Lectio II}
\switchcolumn
\EN
\LessonHead{Lesson II}
\switchcolumn*

\LA
\Cite{Titus 3:3-7}
\switchcolumn
\EN
\Cite{Titus 3:3-7}
\switchcolumn*

\LA
\noindent Erámus enim aliquándo et nos insipiéntes, incréduli, errántes, serviéntes desidériis, et voluptátibus váriis, in malítia et invídia agéntes, odíbiles, odiéntes ínvicem. Cum autem benígnitas et humánitas appáruit Salvatóris nostri Dei, non ex opéribus justítiæ, quæ fécimus nos, sed secúndum suam misericórdiam salvos nos fecit per lavácrum regeneratiónis et renovatiónis Spíritus Sancti, quem effúdit in nos abúnde per Jesum Christum Salvatórem nostrum: ut justificáti grátia ipsíus, hærédes simus secúndum spem vitæ ætérnæ.\par
\switchcolumn
\EN
\noindent For we ourselves also were some time unwise, incredulous, erring, slaves to diverse desires and pleasures, living in malice and envy, hateful, and hating one another. But when the goodness and kindness of God our Saviour appeared: Not by the works of justice, which we have done, but according to his mercy, he saved us, by the laver of regeneration, and renovation of the Holy Ghost; Whom he hath poured forth upon us abundantly, through Jesus Christ our Saviour: That, being justified by his grace, we may be heirs, according to hope of life everlasting.\par
\switchcolumn*

\LA
\begin{Resp}\Rsign Misericórdia tua, Dómine, magna est super me: \Star{} Et liberásti ánimam meam ex inférno inferióri.\end{Resp}
\begin{Resp}\Vsign In die tribulatiónis meæ clamávi ad te, quia exaudísti me.\end{Resp}
\begin{Resp}\Rsign Et liberásti ánimam meam ex inférno inferióri.\end{Resp}
\switchcolumn
\EN
\begin{Resp}\Rsign Great, O Lord, is thy mercy toward me. \Star{} And Thou hast delivered my soul from the lowest hell.\end{Resp}
\begin{Resp}\Vsign In the day of my trouble I called upon thee, for Thou hast heard me.\end{Resp}
\begin{Resp}\Rsign And Thou hast delivered my soul from the lowest hell.\end{Resp}
\switchcolumn*

\LA
\LessonHead{Lectio III}
\switchcolumn
\EN
\LessonHead{Lesson III}
\switchcolumn*

\LA
\Cite{Titus 3:8-11}
\switchcolumn
\EN
\Cite{Titus 3:8-11}
\switchcolumn*

\LA
\noindent Fidélis sermo est: et de his volo te confirmáre: ut curent bonis opéribus præésse qui credunt Deo. Hæc sunt bona, et utília homínibus. Stultas autem quæstiónes, et genealógias, et contentiónes, et pugnas legis devíta: sunt enim inútiles, et vanæ. Hæréticum hóminem post unam et secúndam correptiónem devíta: sciens quia subvérsus est, qui ejúsmodi est, et delínquit, cum sit próprio judício condemnátus.\par
\switchcolumn
\EN
\noindent It is a faithful saying: and these things I will have thee affirm constantly: that they, who believe in God, may be careful to excel in good works. These things are good and profitable unto men. But avoid foolish questions, and genealogies, and contentions, and strivings about the law. For they are unprofitable and vain. A man that is a heretic, after the first and second admonition, avoid: Knowing that he, that is such an one, is subverted, and sinneth, being condemned by his own judgment.\par
\switchcolumn*

\LA
\begin{Resp}\Rsign Factus est mihi Dóminus in refúgium: \Star{} Et Deus meus in auxílium spei meæ.\end{Resp}
\begin{Resp}\Vsign Erípuit me de inimícis meis fortíssimis, et factus est Dóminus protéctor meus.\end{Resp}
\begin{Resp}\Rsign Et Deus meus in auxílium spei meæ.\end{Resp}
\begin{Resp}\Vsign Glória Patri, et Fílio, \Star{} et Spirítui Sancto.\end{Resp}
\begin{Resp}\Rsign Et Deus meus in auxílium spei meæ.\end{Resp}
\switchcolumn
\EN
\begin{Resp}\Rsign The Lord is my refuge. \Star{} And my God is the stay of my trust.\end{Resp}
\begin{Resp}\Vsign He delivered me from the strongest of mine enemies, and the Lord was my stay.\end{Resp}
\begin{Resp}\Rsign And my God is the stay of my trust.\end{Resp}
\begin{Resp}\Vsign Glory be to the Father, and to the Son, \Star{} and to the Holy Ghost.\end{Resp}
\begin{Resp}\Rsign And my God is the stay of my trust.\end{Resp}
\switchcolumn*

\end{paracol}

\Office{Sabbato infra Hebdomadam V post Epiphaniam}{Saturday of the Fifth Week after Epiphany}{Feria}
\addcontentsline{toc}{subsection}{Sabbato infra Hebdomadam V post Epiphaniam\enspace\textit{Saturday of the Fifth Week after Epiphany}}
\begin{paracol}{2}
\LA
\LessonHead{Lectio I}
\switchcolumn
\EN
\LessonHead{Lesson I}
\switchcolumn*

\LA
\LessonTitle{Incipit epístola beáti Pauli Apóstoli ad Philémonem}
\switchcolumn
\EN
\LessonTitle{Lesson from the letter of St. Paul the Apostle to Philemon}
\switchcolumn*

\LA
\Cite{Phlm 1:1-6}
\switchcolumn
\EN
\Cite{Phlm 1:1-6}
\switchcolumn*

\LA
\noindent Paulus vinctus Christi Jesu, et Timótheus frater, Philémoni dilécto, et adjutóri nostro, et Appíæ soróri caríssimæ, et Archíppo commilitóni nostro, et ecclésiæ, quæ in domo tua est. Grátia vobis, et pax a Deo Patre nostro, et Dómino Jesu Christo. Grátias ago Deo meo, semper memóriam tui fáciens in oratiónibus meis, áudiens caritátem tuam, et fidem, quam habes in Dómino Jesu, et in omnes sanctos: ut communicátio fídei tuæ évidens fiat in agnitióne omnis óperis boni, quod est in vobis in Christo Jesu.\par
\switchcolumn
\EN
\noindent Paul, a prisoner of Christ Jesus, and Timothy, a brother: to Philemon, our beloved and fellow labourer; And to Appia, our dearest sister, and to Archippus, our fellow soldier, and to the church which is in thy house: Grace to you and peace from God our Father, and from the Lord Jesus Christ. I give thanks to my God, always making a remembrance of thee in my prayers. Hearing of thy charity and faith, which thou hast in the Lord Jesus, and towards all the saints: That the communication of thy faith may be made evident in the acknowledgment of every good work, that is in you in Christ Jesus.\par
\switchcolumn*

\LA
\begin{Resp}\Rsign Confitébor tibi, Dómine Deus, in toto corde meo, et honorificábo nomen tuum in ætérnum: \Star{} Quia misericórdia tua, Dómine, magna est super me.\end{Resp}
\begin{Resp}\Vsign Deus meus es tu, et confitébor tibi: Deus meus es tu, et exaltábo te.\end{Resp}
\begin{Resp}\Rsign Quia misericórdia tua, Dómine, magna est super me.\end{Resp}
\switchcolumn
\EN
\begin{Resp}\Rsign I will praise thee, O Lord my God, with all my heart, and I will glorify thy Name for evermore. \Star{} For great is thy mercy, O Lord, toward me.\end{Resp}
\begin{Resp}\Vsign Thou art my God, and I will praise thee; Thou art my God, and I will exalt thee.\end{Resp}
\begin{Resp}\Rsign For great is thy mercy, O Lord, toward me.\end{Resp}
\switchcolumn*

\LA
\LessonHead{Lectio II}
\switchcolumn
\EN
\LessonHead{Lesson II}
\switchcolumn*

\LA
\Cite{Phlm 1:7-11}
\switchcolumn
\EN
\Cite{Phlm 1:7-11}
\switchcolumn*

\LA
\noindent Gáudium enim magnum hábui, et consolatiónem in caritáte tua: quia víscera sanctórum requievérunt per te, frater. Propter quod multam fidúciam habens in Christo Jesu imperándi tibi quod ad rem pértinet: propter caritátem magis óbsecro, cum sis talis, ut Paulus senex, nunc autem et vinctus Jesu Christi: óbsecro te pro meo fílio, quem génui in vínculis, Onésimo, qui tibi aliquándo inútilis fuit, nunc autem et mihi et tibi útilis,\par
\switchcolumn
\EN
\noindent For I have had great joy and consolation in thy charity, because the bowels of the saints have been refreshed by thee, brother. Wherefore though I have much confidence in Christ Jesus, to command thee that which is to the purpose: For charity sake I rather beseech, whereas thou art such a one, as Paul an old man, and now a prisoner also of Jesus Christ. I beseech thee for my son, whom I have begotten in my bands, Onesimus, Who hath been heretofore unprofitable to thee, but now is profitable both to me and thee, whom I have sent back to thee.\par
\switchcolumn*

\LA
\begin{Resp}\Rsign Misericórdia tua, Dómine, magna est super me: \Star{} Et liberásti ánimam meam ex inférno inferióri.\end{Resp}
\begin{Resp}\Vsign In die tribulatiónis meæ clamávi ad te, quia exaudísti me.\end{Resp}
\begin{Resp}\Rsign Et liberásti ánimam meam ex inférno inferióri.\end{Resp}
\switchcolumn
\EN
\begin{Resp}\Rsign Great, O Lord, is thy mercy toward me. \Star{} And Thou hast delivered my soul from the lowest hell.\end{Resp}
\begin{Resp}\Vsign In the day of my trouble I called upon thee, for Thou hast heard me.\end{Resp}
\begin{Resp}\Rsign And Thou hast delivered my soul from the lowest hell.\end{Resp}
\switchcolumn*

\LA
\LessonHead{Lectio III}
\switchcolumn
\EN
\LessonHead{Lesson III}
\switchcolumn*

\LA
\Cite{Phlm 1:12-19}
\switchcolumn
\EN
\Cite{Phlm 1:12-19}
\switchcolumn*

\LA
\noindent quem remísi tibi. Tu autem illum, ut mea víscera, súscipe: quem ego volúeram mecum detinére, ut pro te mihi ministráret in vínculis Evangélii: sine consílio autem tuo nihil vólui fácere, uti ne velut ex necessitáte bonum tuum esset, sed voluntárium. Fórsitan enim ídeo discéssit ad horam a te, ut ætérnum illum recíperes: jam non ut servum, sed pro servo caríssimum fratrem, máxime mihi: quanto autem magis tibi et in carne, et in Dómino? Si ergo habes me sócium, súscipe illum sicut me: si autem áliquid nócuit tibi, aut debet, hoc mihi ímputa. Ego Paulus scripsi mea manu.\par
\switchcolumn
\EN
\noindent And do thou, therefore, receive him as my own bowels. Whom I would have retained with me, that in thy stead he might have ministered to me in the bands of the gospel: But without thy counsel I would do nothing: that thy good deed might not be as it were of necessity, but voluntary. For perhaps he therefore departed for a season from thee, that thou mightest receive him again for ever: Not now as a servant, but instead of a servant, a most dear brother, especially to me: but how much more to thee both in the flesh and in the Lord? If therefore thou count me a partner, receive him as myself. And if he hath wronged thee in any thing, or is in thy debt, put that to my account. I Paul have written it with my own hand: I will repay it: not to say to thee, that thou owest me thy own self also.\par
\switchcolumn*

\LA
\begin{Resp}\Rsign Factus est mihi Dóminus in refúgium: \Star{} Et Deus meus in auxílium spei meæ.\end{Resp}
\begin{Resp}\Vsign Erípuit me de inimícis meis fortíssimis, et factus est Dóminus protéctor meus.\end{Resp}
\begin{Resp}\Rsign Et Deus meus in auxílium spei meæ.\end{Resp}
\begin{Resp}\Vsign Glória Patri, et Fílio, \Star{} et Spirítui Sancto.\end{Resp}
\begin{Resp}\Rsign Et Deus meus in auxílium spei meæ.\end{Resp}
\switchcolumn
\EN
\begin{Resp}\Rsign The Lord is my refuge. \Star{} And my God is the stay of my trust.\end{Resp}
\begin{Resp}\Vsign He delivered me from the strongest of mine enemies, and the Lord was my stay.\end{Resp}
\begin{Resp}\Rsign And my God is the stay of my trust.\end{Resp}
\begin{Resp}\Vsign Glory be to the Father, and to the Son, \Star{} and to the Holy Ghost.\end{Resp}
\begin{Resp}\Rsign And my God is the stay of my trust.\end{Resp}
\switchcolumn*

\end{paracol}

\Office{Dominica VI Post Epiphaniam}{Sixth Sunday after Epiphany}{Semiduplex}
\addcontentsline{toc}{subsection}{Dominica VI Post Epiphaniam\enspace\textit{Sixth Sunday after Epiphany}}
\begin{paracol}{2}
\LA
\RefLine{Lectiones I–VI: de Scriptura occurrente.}
\switchcolumn
\EN
\RefLine{Lessons I–VI: from the Scripture of the day.}
\switchcolumn*

\LA
\LessonHead{Lectio I}
\switchcolumn
\EN
\LessonHead{Lesson I}
\switchcolumn*

\LA
\LessonTitle{Incipit Epístola beáti Pauli Apóstoli ad Hebrǽos}
\switchcolumn
\EN
\LessonTitle{Lesson from the letter of St. Paul the Apostle to the Hebrews}
\switchcolumn*

\LA
\Cite{Heb 1:1-4}
\switchcolumn
\EN
\Cite{Heb 1:1-4}
\switchcolumn*

\LA
\noindent Multifáriam, multísque modis olim Deus loquens pátribus in prophétis: novíssime, diébus istis locútus est nobis in Fílio, quem constítuit herédem universórum, per quem fecit et sǽcula: qui cum sit splendor glóriæ, et figúra substántiæ ejus, portánsque ómnia verbo virtútis suæ, purgatiónem peccatórum fáciens, sedet ad déxteram majestátis in excélsis: tanto mélior ángelis efféctus, quanto differéntius præ illis nomen hereditávit.\par
\switchcolumn
\EN
\noindent God, who, at sundry times and in diverse manners, spoke in times past to the fathers by the prophets, last of all, In these days hath spoken to us by his Son, whom he hath appointed heir of all things, by whom also he made the world. Who being the brightness of his glory, and the figure of his substance, and upholding all things by the word of his power, making purgation of sins, sitteth on the right hand of the majesty on high. Being made so much better than the angels, as he hath inherited a more excellent name than they.\par
\switchcolumn*

\LA
\begin{Resp}\Rsign Dómine, ne in ira tua árguas me, neque in furóre tuo corrípias me: \Star{} Miserére mei, Dómine, quóniam infírmus sum.\end{Resp}
\begin{Resp}\Vsign Timor et tremor venérunt super me, et contexérunt me ténebræ.\end{Resp}
\begin{Resp}\Rsign Miserére mei, Dómine, quóniam infírmus sum.\end{Resp}
\switchcolumn
\EN
\begin{Resp}\Rsign O Lord, rebuke me not in thine anger, neither chasten me in thine hot displeasure. \Star{} Have mercy upon me, O Lord, for I am weak.\end{Resp}
\begin{Resp}\Vsign Fearfulness and trembling are come upon me, and darkness hath overwhelmed me.\end{Resp}
\begin{Resp}\Rsign Have mercy upon me, O Lord, for I am weak.\end{Resp}
\switchcolumn*

\LA
\LessonHead{Lectio II}
\switchcolumn
\EN
\LessonHead{Lesson II}
\switchcolumn*

\LA
\Cite{Heb 1:5-9}
\switchcolumn
\EN
\Cite{Heb 1:5-9}
\switchcolumn*

\LA
\noindent Cui enim dixit aliquándo angelórum: Fílius meus es tu, ego hódie génui te? Et rursum: Ego ero illi in patrem, et ipse erit mihi in fílium? Et cum íterum introdúcit primogénitum in orbem terræ, dicit: Et adórent eum omnes ángeli Dei. Et ad ángelos quidem dicit: Qui facit ángelos suos spíritus, et minístros suos flammam ignis. Ad Fílium autem: Thronus tuus Deus in sǽculum sǽculi: virga æquitátis, virga regni tui. Dilexísti justítiam, et odísti iniquitátem: proptérea unxit te Deus, Deus tuus, óleo exultatiónis præ particípibus tuis.\par
\switchcolumn
\EN
\noindent For to which of the angels hath he said at any time, Thou art my Son, to day have I begotten thee? And again, I will be to him a Father, and he shall be to me a Son? And again, when he bringeth in the first begotten into the world, he saith: And let all the angels of God adore him. And to the angels indeed he saith: He that maketh his angels spirits, and his ministers a flame of fire. But to the Son: thy throne, O God, is for ever and ever: a sceptre of justice is the sceptre of thy kingdom. Thou hast loved justice, and hated iniquity: therefore God, thy God, hath anointed thee with the oil of gladness above thy fellows.\par
\switchcolumn*

\LA
\begin{Resp}\Rsign Deus, qui sedes super thronum, et júdicas æquitátem, esto refúgium páuperum in tribulatióne: \Star{} Quia tu solus labórem et dolórem consíderas.\end{Resp}
\begin{Resp}\Vsign Tibi enim derelíctus est pauper, pupíllo tu eris adjútor.\end{Resp}
\begin{Resp}\Rsign Quia tu solus labórem et dolórem consíderas.\end{Resp}
\switchcolumn
\EN
\begin{Resp}\Rsign O God, Which satest in the throne judging right, be Thou a refuge for the poor, a refuge in times of trouble. \Star{} For Thou alone beholdest mischief and spite.\end{Resp}
\begin{Resp}\Vsign The poor leaveth himself unto thee; Thou wilt be the helper of the fatherless.\end{Resp}
\begin{Resp}\Rsign For Thou alone beholdest mischief and spite.\end{Resp}
\switchcolumn*

\LA
\LessonHead{Lectio III}
\switchcolumn
\EN
\LessonHead{Lesson III}
\switchcolumn*

\LA
\Cite{Heb 1:10-14}
\switchcolumn
\EN
\Cite{Heb 1:10-14}
\switchcolumn*

\LA
\noindent Et: Tu in princípio, Dómine, terram fundásti: et ópera mánuum tuárum sunt cæli. Ipsi períbunt, tu autem permanébis, et omnes ut vestiméntum veteráscent: et velut amíctum mutábis eos, et mutabúntur: tu autem idem ipse es, et anni tui non defícient. Ad quem autem angelórum dixit aliquándo: Sede a dextris meis, quoadúsque ponam inimícos tuos scabéllum pedum tuórum? Nonne omnes sunt administratórii spíritus, in ministérium missi propter eos, qui hæreditátem cápient salútis?\par
\switchcolumn
\EN
\noindent And: Thou in the beginning, O Lord, didst found the earth: and the works of thy hands are the heavens. They shall perish, but thou shalt continue: and they shall all grow old as a garment. And as a vesture shalt thou change them, and they shall be changed: but thou art the selfsame, and thy years shall not fail. But to which of the angels said he at any time: Sit on my right hand, until I make thy enemies thy footstool? Are they not all ministering spirits, sent to minister for them, who shall receive the inheritance of salvation?\par
\switchcolumn*

\LA
\begin{Resp}\Rsign A dextris est mihi Dóminus, ne commóvear: \Star{} Propter hoc dilatátum est cor meum, et exsultávit lingua mea.\end{Resp}
\begin{Resp}\Vsign Dóminus pars hereditátis meæ, et cálicis mei.\end{Resp}
\begin{Resp}\Rsign Propter hoc dilatátum est cor meum, et exsultávit lingua mea.\end{Resp}
\begin{Resp}\Vsign Glória Patri, et Fílio, \Star{} et Spirítui Sancto.\end{Resp}
\begin{Resp}\Rsign Propter hoc dilatátum est cor meum, et exsultávit lingua mea.\end{Resp}
\switchcolumn
\EN
\begin{Resp}\Rsign The Lord is at my right hand, I shall never be moved. \Star{} Therefore my heart is glad, and my tongue rejoiceth.\end{Resp}
\begin{Resp}\Vsign The Lord is the portion of mine inheritance, and of my cup.\end{Resp}
\begin{Resp}\Rsign Therefore my heart is glad, and my tongue rejoiceth.\end{Resp}
\begin{Resp}\Vsign Glory be to the Father, and to the Son, \Star{} and to the Holy Ghost.\end{Resp}
\begin{Resp}\Rsign Therefore my heart is glad, and my tongue rejoiceth.\end{Resp}
\switchcolumn*

\LA
\LessonHead{Lectio IV}
\switchcolumn
\EN
\LessonHead{Lesson IV}
\switchcolumn*

\LA
\LessonTitle{Sermo sancti Athanásii Epíscopi}
\switchcolumn
\EN
\LessonTitle{From the Sermons of St. Athanasius, Pope (of Alexandria.)}
\switchcolumn*

\LA
\Source{Orat. 2 contra Ariános, post medium.}
\switchcolumn
\EN
\Source{2nd against the Arians.}
\switchcolumn*

\LA
\noindent Si persónam, rem, tempus apostólici dicti cognóscerent hærétici, numquam humána in Deitátem transferéntes, tam ímpie et stulte advérsus Christum sese habuíssent. Id intuéri licébit, si inítium lectiónis dénuo repetítum probe excípias. Dicit enim Apóstolus: Multifáriam, multísque modis olim Deus locútus est pátribus nostris per Prophétas: últimis autem diébus locútus est nobis in Fílio. Atque ita paulo post dicit: Perfécta ab eo nostrórum peccatórum purificatióne, ipsum sedére ad déxteram majestátis in excélsis, tanto meliórem Angelis factum, quanto præstántius nomen præ illis sortítus est. De eo ígitur témpore, quo nobis per Fílium locútus est, cum peccatórum purgátio fíeret, apostólicum dictum mentiónem facit. Quando autem nobis locútus est in Fílio, aut quando purgátio peccatórum facta, aut quando natus est homo, nisi post Prophétas, idque in últimis diébus?\par
\switchcolumn
\EN
\noindent If the heretics had but known the person, the matter, and the times of the Apostle who spoke, they would never have spoken of Godhead as if It were human, nor borne themselves so wickedly, and withal so foolishly against Christ. It will be permitted to us to return, and to take again the first words of the Lesson. The Apostle then saith God, Who at sundry times and in diverse manners, spake in time past unto the fathers by the Prophets, hath in these last days spoken unto us by His Son and again, a little farther on: When the Son had purged our sins, He sat down on the right hand of the Majesty on high being made so much better than the angels as He hath by inheritance obtained a more excellent name than they. The Apostle here expressly nameth the times wherein God hath spoken unto us by His Son, and wherein the Same His Son hath purged our sins; for when hath He spoken unto us by His Son, when did the Son purge our sins, or when was He born a Man, but since God spake unto the Fathers by the Prophets, namely, in these last days?\par
\switchcolumn*

\LA
\begin{Resp}\Rsign Notas mihi fecísti, Dómine, vias vitæ: \Star{} Adimplébis me lætítia cum vultu tuo: delectatiónes in déxtera tua usque in finem.\end{Resp}
\begin{Resp}\Vsign Tu es qui restítues hereditátem meam mihi.\end{Resp}
\begin{Resp}\Rsign Adimplébis me lætítia cum vultu tuo: delectatiónes in déxtera tua usque in finem.\end{Resp}
\switchcolumn
\EN
\begin{Resp}\Rsign O Lord, Thou hast shown me the path of life. \Star{} Thou shalt fill me with joy in thy presence, at thy right hand there are pleasures for evermore.\end{Resp}
\begin{Resp}\Vsign Thou art He That shalt restore mine inheritance unto me.\end{Resp}
\begin{Resp}\Rsign Thou shalt fill me with joy in thy presence, at thy right hand there are pleasures for evermore.\end{Resp}
\switchcolumn*

\LA
\LessonHead{Lectio V}
\switchcolumn
\EN
\LessonHead{Lesson V}
\switchcolumn*

\LA
\noindent Deínde cum narrátio institúta esset de humána Verbi dispensatióne, deque últimis tempóribus: consequénter commemorávit, Deum neque superióribus ætátibus tacuísse, sed locútum esse per Prophétas: et postquam Prophétæ suo offício perfúncti sunt, et lex per Angelos pronuntiáta est, et Fílius étiam ad nos descéndit, et ad ministrándum accéssit; tunc demum necessário subíntulit: Tanto mélior Angelis factus: osténdere volens, quanto Fílius præ servo excéllit, tanto functióne officióque servórum, Fílii administratiónem meliórem fuísse.\par
\switchcolumn
\EN
\noindent The Apostle, about to enter on the subject of the Word's human dispensation and the last days, naturally mentioneth first that God had not up to those days been silent, but had spoken unto the fathers by the Prophets and, after the Prophets had discharged their office, and the law had been given by the ministry of angels, that the Son also came down unto us to minister and then he addeth, being made so much better than the angels, to show that as the Son differeth from a servant, so is the ministry of the Son better than the duty and office of servants.\par
\switchcolumn*

\LA
\begin{Resp}\Rsign Díligam te, Dómine, virtus mea: Dóminus firmaméntum meum, \Star{} Et refúgium meum.\end{Resp}
\begin{Resp}\Vsign Liberátor meus, Deus meus, adjútor meus.\end{Resp}
\begin{Resp}\Rsign Et refúgium meum.\end{Resp}
\switchcolumn
\EN
\begin{Resp}\Rsign I will love thee, O Lord, my strength; the Lord is my rock \Star{} And my fortress.\end{Resp}
\begin{Resp}\Vsign My Deliverer, my God, mine Helper.\end{Resp}
\begin{Resp}\Rsign And my fortress.\end{Resp}
\switchcolumn*

\LA
\LessonHead{Lectio VI}
\switchcolumn
\EN
\LessonHead{Lesson VI}
\switchcolumn*

\LA
\noindent Functiónem ígitur discérnens Apóstolus, tum véterem, tum novam, magna dicéndi libertáte útitur, ad Judǽos scribens et loquens. Propter hoc ígitur non in univérsum ex própria comparatiónis ratióne dixit, quod major aut honorátior esset: ne quis quasi de ejúsdem géneris, et cum eo commúnibus rebus hæc verba intellígeret: sed ídeo meliórem illum dixit, ut discrímen natúræ Fílii ad res creátas indicáret.\par
\switchcolumn
\EN
\noindent The Apostle, therefore, seeing the difference between the new ministry and the old, maketh very bold in writing and speaking to the Jews. For this cause, therefore, he doth not compare the details of the two ministries, and then come to the general conclusion that the new was greater or more honourable than the old, (lest any should understand that the two ministries were of the same kind, and that the conclusion that the new is better is arrived at by comparing the degrees in each of things which they had in common,) but he saith that the Son was made better, to distinguish at once and completely the nature of the Son from the nature of things created.\par
\switchcolumn*

\LA
\begin{Resp}\Rsign Dómini est terra, et plenitúdo ejus: \Star{} Orbis terrárum, et univérsi qui hábitant in eo.\end{Resp}
\begin{Resp}\Vsign Ipse super mária fundávit eam, et super flúmina præparávit illam.\end{Resp}
\begin{Resp}\Rsign Orbis terrárum, et univérsi qui hábitant in eo.\end{Resp}
\begin{Resp}\Vsign Glória Patri, et Fílio, \Star{} et Spirítui Sancto.\end{Resp}
\begin{Resp}\Rsign Orbis terrárum, et univérsi qui hábitant in eo.\end{Resp}
\switchcolumn
\EN
\begin{Resp}\Rsign The earth is the Lord's, and the fulness thereof \Star{} The world, and they that dwell therein.\end{Resp}
\begin{Resp}\Vsign For He hath founded it upon the seas, and established it upon the floods.\end{Resp}
\begin{Resp}\Rsign The world, and they that dwell therein.\end{Resp}
\begin{Resp}\Vsign Glory be to the Father, and to the Son, \Star{} and to the Holy Ghost.\end{Resp}
\begin{Resp}\Rsign The world, and they that dwell therein.\end{Resp}
\switchcolumn*

\LA
\LessonHead{Lectio VII}
\switchcolumn
\EN
\LessonHead{Lesson VII}
\switchcolumn*

\LA
\LessonTitle{Léctio sancti Evangélii secúndum Matthǽum}
\switchcolumn
\EN
\LessonTitle{From the Holy Gospel according to Matthew}
\switchcolumn*

\LA
\Cite{Matt 13:31-35}
\switchcolumn
\EN
\Cite{Matt 13:31-32}
\switchcolumn*

\LA
\noindent In illo témpore: Dixit Jesus turbis parábolam hanc: Símile est regnum cælórum grano sinápis, quod accípiens homo seminávit in agro suo. Et réliqua.\par
\switchcolumn
\EN
\noindent At that time, Jesus said to the people a parable: The kingdom of heaven is like to a grain of mustard seed, which a man took and sowed in his field. And so on.\par
\switchcolumn*

\LA
\LessonTitle{Homilía S. Hierónymi Presbýteri}
\switchcolumn
\EN
\LessonTitle{Homily by St. Jerome, Priest (at Bethlehem.)}
\switchcolumn*

\LA
\Source{Lib. 2. Comment. in cap. 13. Matth.}
\switchcolumn
\EN
\Source{Book II Comment. on Matth. xiii.}
\switchcolumn*

\LA
\noindent Regnum cælórum prædicátio Evangélii est, et notítia Scripturárum, quæ ducit ad vitam, et de qua dícitur ad Judǽos: Auferétur a vobis regnum Dei, et dábitur genti faciénti fructus ejus. Símile est ergo hujuscémodi regnum grano sinápis, quod accípiens homo seminávit in agro suo. Homo qui séminat in agro suo, a plerísque Salvátor intellégitur, quod in ánimis credéntium séminet: ab áliis ipse homo séminans in agro suo, hoc est in semetípso, et in corde suo.\par
\switchcolumn
\EN
\noindent The kingdom of heaven is the proclamation of the Gospel, and that knowledge of the Scriptures, which leadeth unto life, and whereof it is said to the Jews: The kingdom of God shall be taken from you, and given to a nation bringing forth the fruits thereof. (Matth. xxi. 43) Therefore is this kingdom like to a grain of mustard-seed, which a man took and sowed in his field. By the man that sowed it in his field, many understand to be meant the Saviour, because He is the Sower That soweth in the souls of believers; others understand every man that soweth good seed in his own field, that is, in himself and in his own heart.\par
\switchcolumn*

\LA
\begin{Resp}\Rsign Laudábilis pópulus, \Star{} Quem Dóminus exercítuum benedíxit dicens: Opus mánuum meárum tu es, heréditas mea Israël.\end{Resp}
\begin{Resp}\Vsign Beáta gens, cujus est Dóminus Deus, pópulus eléctus in hereditátem.\end{Resp}
\begin{Resp}\Rsign Quem Dóminus exercítuum benedíxit dicens: Opus mánuum meárum tu es, heréditas mea Israël.\end{Resp}
\switchcolumn
\EN
\begin{Resp}\Rsign Blessed is the people \Star{} Whom the Lord of hosts hath blessed, saying O Israel thou art the work of Mine own hands, thou art Mine own inheritance.\end{Resp}
\begin{Resp}\Vsign Blessed is the nation whose God is the Lord, and the people whom He hath chosen for His own inheritance.\end{Resp}
\begin{Resp}\Rsign Whom the Lord of hosts hath blessed, saying O Israel thou art the work of Mine own hands, thou art Mine own inheritance.\end{Resp}
\switchcolumn*

\LA
\LessonHead{Lectio VIII}
\switchcolumn
\EN
\LessonHead{Lesson VIII}
\switchcolumn*

\LA
\noindent Quis est iste, qui séminat, nisi sensus noster et ánimus; qui suscípiens granum prædicatiónis, et fovens seméntem, humóre fídei facit in agro sui péctoris pulluláre? Prædicátio Evangélii mínima est ómnibus disciplínis. Ad primam quippe doctrínam, fidem non habet veritátis, hóminem Deum, Christum mórtuum, et scándalum crucis prǽdicans. Confer hujuscémodi doctrínam dogmátibus philosophórum, et libris eórum, et splendóri eloquéntiæ, et compositióni sermónum: et vidébis quanto minor sit céteris semínibus seméntis Evangélii.\par
\switchcolumn
\EN
\noindent Who is he that soweth, but our own mind and soul, which take the grain from preaching, and by nourishing it in the soil, cause it to sprout in the field of our own breast? The preaching of the Gospel is the least of all doctrines. He that preacheth, for his first lesson, God made man, Christ dead, and the stumbling-block of the Cross, receiveth at first but little credit. Compare such teaching as this with the doctrines of the Philosophers, with their books, their magnificent eloquence, and their rounded sentences, and thou shalt see how the grain of the Gospel, when it is sown, is the humblest of all seeds.\par
\switchcolumn*

\LA
\begin{Resp}\Rsign Duo Séraphim clamábant alter ad álterum: \Star{} Sanctus, sanctus, sanctus Dóminus Deus Sábaoth: * Plena est omnis terra glória ejus.\end{Resp}
\begin{Resp}\Vsign Tres sunt qui testimónium dant in cælo: Pater, Verbum, et Spíritus Sanctus: et hi tres unum sunt.\end{Resp}
\begin{Resp}\Rsign Sanctus, sanctus, sanctus Dóminus Deus Sábaoth.\end{Resp}
\begin{Resp}\Vsign Glória Patri, et Fílio, \Star{} et Spirítui Sancto.\end{Resp}
\begin{Resp}\Rsign Plena est omnis terra glória ejus.\end{Resp}
\switchcolumn
\EN
\begin{Resp}\Rsign One Seraph cried unto another: \Star{} Holy, Holy, Holy is the Lord God of hosts; the whole earth is full of His glory.\end{Resp}
\begin{Resp}\Vsign There are Three That bear record in heaven: the Father, the Word, and the Holy Ghost, and these Three are One.\end{Resp}
\begin{Resp}\Rsign Holy, Holy, Holy is the Lord God of hosts.\end{Resp}
\begin{Resp}\Vsign Glory be to the Father, and to the Son, \Star{} and to the Holy Ghost.\end{Resp}
\begin{Resp}\Rsign The whole earth is full of His glory.\end{Resp}
\switchcolumn*

\LA
\LessonHead{Lectio IX}
\switchcolumn
\EN
\LessonHead{Lesson IX}
\switchcolumn*

\LA
\noindent Sed illa cum créverint, nihil mordax, nihil vívidum, nihil vitále demónstrant: sed totum fláccidum marcidúmque et mollítum ebúllit in ólera et in herbas, quæ cito aréscunt et córruunt. Hæc autem prædicátio, quæ parva videbátur in princípio, cum vel in ánima credéntis, vel in tot mundo sata fúerit, non exsúrgit in ólera, sed crescit in árborem: ita ut vólucres cæli (quas vel ánimas credéntium, vel fortitúdines, Dei servítio mancipátas, sentíre debémus) véniant et hábitent in ramis ejus. Ramos puto evangélicæ árboris, quæ de grano sinápis créverit, dógmatum esse diversitátes, in quibus supradictárum vólucrum unaquǽque requiéscit.\par
\switchcolumn
\EN
\noindent But when the doctrines of men grow up, there is therein nothing piercing, nothing healthy, nothing life-giving. The plant is drooping, and delicate, and soft. There are herbs and grass whereof it may truly be said that the grass withereth and the flower fadeth. (Isa. xl. 8.) But the grain of Gospel seed, though, when it was sown, it seemed to be the least of all seeds, when once it is rooted in the soul of man, or in the whole world, groweth not into an herb, but becometh a tree so that the birds of the air (whereby we may understand, either the souls of believers, or the (angelic) powers bound to the service of God,) come and lodge in the branches thereof. I consider that the branches of the Gospel tree, which groweth from the grain of mustard-seed, are the diverse developments of doctrine, on which the birds above mentioned find resting-places.\par
\switchcolumn*

\LA
\TeDeum{Te Deum}
\switchcolumn
\EN
\TeDeum{Te Deum}
\switchcolumn*

\end{paracol}

\Office{Feria II infra Hebdomadam VI post Epiphaniam}{Monday of the Sixth Week after Epiphany}{Feria}
\addcontentsline{toc}{subsection}{Feria II infra Hebdomadam VI post Epiphaniam\enspace\textit{Monday of the Sixth Week after Epiphany}}
\begin{paracol}{2}
\LA
\LessonHead{Lectio I}
\switchcolumn
\EN
\LessonHead{Lesson I}
\switchcolumn*

\LA
\LessonTitle{De Epístola ad Hebrǽos}
\switchcolumn
\EN
\LessonTitle{Lesson from the letter of St. Paul the Apostle to the Hebrews}
\switchcolumn*

\LA
\Cite{Heb 3:1-4}
\switchcolumn
\EN
\Cite{Heb 3:1-4}
\switchcolumn*

\LA
\noindent Unde, fratres sancti, vocatiónis cæléstis partícipes, consideráte Apóstolum, et pontíficem confessiónis nostræ Jesum: qui fidélis est ei, qui fecit illum, sicut et Móyses in omni domo ejus. Amplióris enim glóriæ iste præ Móyse dignus est hábitus, quanto ampliórem honórem habet domus, qui fabricávit illam. Omnis namque domus fabricátur ab áliquo: qui autem ómnia creávit, Deus est.\par
\switchcolumn
\EN
\noindent Wherefore, holy brethren, partakers of the heavenly vocation, consider the apostle and high priest of our confession, Jesus: Who is faithful to him that made him, as was also Moses in all his house. For this man was counted worthy of greater glory than Moses, by so much as he that hath built the house, hath greater honour than the house. For every house is built by some man: but he that created all things, is God.\par
\switchcolumn*

\LA
\begin{Resp}\Rsign Quam magna multitúdo dulcédinis tuæ, Dómine, \Star{} Quam abscondísti timéntibus te!\end{Resp}
\begin{Resp}\Vsign Et perfecísti eis qui sperant in te, Dómine, in conspéctu filiórum hóminum.\end{Resp}
\begin{Resp}\Rsign Quam abscondísti timéntibus te!\end{Resp}
\switchcolumn
\EN
\begin{Resp}\Rsign O how great is thy goodness, O Lord \Star{} Which Thou hast laid up for them that fear thee!\end{Resp}
\begin{Resp}\Vsign Which Thou, O Lord, hast wrought for them that trust in thee before the sons of men!\end{Resp}
\begin{Resp}\Rsign Which Thou hast laid up for them that fear thee!\end{Resp}
\switchcolumn*

\LA
\LessonHead{Lectio II}
\switchcolumn
\EN
\LessonHead{Lesson II}
\switchcolumn*

\LA
\Cite{Heb 3:5-8}
\switchcolumn
\EN
\Cite{Heb 3:5-8}
\switchcolumn*

\LA
\noindent Et Móyses quidem fidélis erat in tota domo ejus tamquam fámulus, in testimónium eórum, quæ dicénda erant: Christus vero tamquam fílius in domo sua: quæ domus sumus nos, si fidúciam, et glóriam spei usque ad finem, firmam retineámus. Quaprópter sicut dicit Spíritus Sanctus: Hódie si vocem ejus audiéritis, nolíte obduráre corda vestra, sicut in exacerbatióne secúndum diem tentatiónis in desérto.\par
\switchcolumn
\EN
\noindent And Moses indeed was faithful in all his house as a servant, for a testimony of those things which were to be said: But Christ as the Son in his own house: which house are we, if we hold fast the confidence and glory of hope unto the end. Wherefore, as the Holy Ghost saith: To day if you shall hear his voice, Harden not your hearts, as in the provocation; in the day of temptation in the desert,\par
\switchcolumn*

\LA
\begin{Resp}\Rsign Adjútor meus esto, Deus: \Star{} Ne derelínquas me.\end{Resp}
\begin{Resp}\Vsign Neque despícias me, Deus, salutáris meus.\end{Resp}
\begin{Resp}\Rsign Ne derelínquas me.\end{Resp}
\switchcolumn
\EN
\begin{Resp}\Rsign O God, be Thou my helper. \Star{} Neither leave me.\end{Resp}
\begin{Resp}\Vsign Nor forsake me O God of my salvation.\end{Resp}
\begin{Resp}\Rsign Neither leave me.\end{Resp}
\switchcolumn*

\LA
\LessonHead{Lectio III}
\switchcolumn
\EN
\LessonHead{Lesson III}
\switchcolumn*

\LA
\Cite{Heb 3:12-16}
\switchcolumn
\EN
\Cite{Heb 3:12-16}
\switchcolumn*

\LA
\noindent Vidéte fratres, ne forte sit in áliquo vestrum cor malum incredulitátis, discedéndi a Deo vivo: sed adhortámini vosmetípsos per síngulos dies, donec hódie cognominátur, ut non obdurétur quis ex vobis fallácia peccáti. Partícipes enim Christi effécti sumus, si tamen inítium substántiæ ejus usque ad finem firmum retineámus. Dum dícitur: Hódie si vocem ejus audiéritis, nolíte obduráre corda vestra, quemádmodum in illa exacerbatióne. Quidam enim audiéntes exacerbavérunt: sed non univérsi qui profécti sunt ex Ægýpto per Móysen.\par
\switchcolumn
\EN
\noindent Take heed, brethren, lest perhaps there be in any of you an evil heart of unbelief, to depart from the living God. But exhort one another every day, whilst it is called to day, that none of you be hardened through the deceitfulness of sin. For we are made partakers of Christ: yet so, if we hold the beginning of his substance firm unto the end. While it is said, To day if you shall hear his voice, harden not your hearts, as in that provocation. For some who heard did provoke: but not all that came out of Egypt by Moses.\par
\switchcolumn*

\LA
\begin{Resp}\Rsign Benedícam Dóminum in omni témpore: \Star{} Semper laus ejus in ore meo.\end{Resp}
\begin{Resp}\Vsign In Dómino laudábitur ánima mea, áudiant mansuéti, et læténtur.\end{Resp}
\begin{Resp}\Rsign Semper laus ejus in ore meo.\end{Resp}
\begin{Resp}\Vsign Glória Patri, et Fílio, \Star{} et Spirítui Sancto.\end{Resp}
\begin{Resp}\Rsign Semper laus ejus in ore meo.\end{Resp}
\switchcolumn
\EN
\begin{Resp}\Rsign I will bless the Lord at all times. \Star{} His praise shall continually be in my mouth.\end{Resp}
\begin{Resp}\Vsign My soul shall make her boast in the Lord; the humble shall hear thereof and be glad.\end{Resp}
\begin{Resp}\Rsign His praise shall continually be in my mouth.\end{Resp}
\begin{Resp}\Vsign Glory be to the Father, and to the Son, \Star{} and to the Holy Ghost.\end{Resp}
\begin{Resp}\Rsign His praise shall continually be in my mouth.\end{Resp}
\switchcolumn*

\end{paracol}

\Office{Feria III infra Hebdomadam VI post Epiphaniam}{Tuesday of the Sixth Week after Epiphany}{Feria}
\addcontentsline{toc}{subsection}{Feria III infra Hebdomadam VI post Epiphaniam\enspace\textit{Tuesday of the Sixth Week after Epiphany}}
\begin{paracol}{2}
\LA
\LessonHead{Lectio I}
\switchcolumn
\EN
\LessonHead{Lesson I}
\switchcolumn*

\LA
\LessonTitle{De Epístola ad Hebrǽos}
\switchcolumn
\EN
\LessonTitle{Lesson from the letter of St. Paul the Apostle to the Hebrews}
\switchcolumn*

\LA
\Cite{Heb 4:1-3}
\switchcolumn
\EN
\Cite{Heb 4:1-3}
\switchcolumn*

\LA
\noindent Timeámus ergo ne forte relícta pollicitatióne introeúndi in réquiem ejus, existimétur áliquis ex vobis deésse. Etenim et nobis nuntiátum est, quemádmodum et illis: sed non prófuit illis sermo audítus, non admístus fídei ex iis quæ audiérunt. Ingrediémur enim in réquiem, qui credídimus: quemádmodum dixit: Sicut jurávi in ira mea: Si introíbunt in réquiem meam: et quidem opéribus ab institutióne mundi perféctis.\par
\switchcolumn
\EN
\noindent Let us fear therefore lest the promise being left of entering into his rest, any of you should be thought to be wanting. For unto us also it hath been declared, in like manner as unto them. But the word of hearing did not profit them, not being mixed with faith of those things they heard. For we, who have believed, shall enter into rest; as he said: As I have sworn in my wrath; If they shall enter into my rest; and this indeed when the works from the foundation of the world were finished.\par
\switchcolumn*

\LA
\begin{Resp}\Rsign Auribus pércipe, Deus, lácrimas meas: ne síleas a me, remítte mihi: \Star{} Quóniam íncola ego sum apud te, et peregrínus.\end{Resp}
\begin{Resp}\Vsign Compláceat tibi, ut erípias me: Dómine, ad adjuvándum me festína.\end{Resp}
\begin{Resp}\Rsign Quóniam íncola ego sum apud te, et peregrínus.\end{Resp}
\switchcolumn
\EN
\begin{Resp}\Rsign O God, give ear unto my tears; hold not thy peace, but, O, spare me! \Star{} For I am a stranger with thee, and a sojourner.\end{Resp}
\begin{Resp}\Vsign Be pleased, O Lord, to deliver me; O Lord, look upon me to help me.\end{Resp}
\begin{Resp}\Rsign For I am a stranger with thee, and a sojourner.\end{Resp}
\switchcolumn*

\LA
\LessonHead{Lectio II}
\switchcolumn
\EN
\LessonHead{Lesson II}
\switchcolumn*

\LA
\Cite{Heb 4:4-7}
\switchcolumn
\EN
\Cite{Heb 4:4-7}
\switchcolumn*

\LA
\noindent Dixit enim in quodam loco de die séptima sic: Et requiévit Deus die séptima ab ómnibus opéribus suis. Et in isto rursum: Si introíbunt in réquiem meam. Quóniam ergo súperest introíre quosdam in illam, et ii, quibus prióribus annuntiátum est, non introiérunt propter incredulitátem: íterum términat diem quemdam, Hódie, in David dicéndo, post tantum témporis, sicut supra dictum est: Hódie si vocem ejus audiéritis, nolíte obduráre corda vestra.\par
\switchcolumn
\EN
\noindent For in a certain place he spoke of the seventh day thus: And God rested the seventh day from all his works. And in this place again: If they shall enter into my rest. Seeing then it remaineth that some are to enter into it, and they, to whom it was first preached, did not enter because of unbelief: Again he limiteth a certain day, saying in David, To day, after so long a time, as it is above said: To day if you shall hear his voice, harden not your hearts.\par
\switchcolumn*

\LA
\begin{Resp}\Rsign Adjútor meus esto, Deus: \Star{} Ne derelínquas me.\end{Resp}
\begin{Resp}\Vsign Neque despícias me, Deus, salutáris meus.\end{Resp}
\begin{Resp}\Rsign Ne derelínquas me.\end{Resp}
\switchcolumn
\EN
\begin{Resp}\Rsign O God, be Thou my helper. \Star{} Neither leave me.\end{Resp}
\begin{Resp}\Vsign Nor forsake me O God of my salvation.\end{Resp}
\begin{Resp}\Rsign Neither leave me.\end{Resp}
\switchcolumn*

\LA
\LessonHead{Lectio III}
\switchcolumn
\EN
\LessonHead{Lesson III}
\switchcolumn*

\LA
\Cite{Heb 4:8-12}
\switchcolumn
\EN
\Cite{Heb 4:8-12}
\switchcolumn*

\LA
\noindent Nam si eis Jesus réquiem præstitísset, numquam de ália loquerétur, posthac, die. Itaque relínquitur sabbatísmus pópulo Dei. Qui enim ingréssus est in réquiem ejus, étiam ipse requiévit ab opéribus suis, sicut a suis Deus. Festinémus ergo íngredi in illam réquiem: ut ne in idípsum quis íncidat incredulitátis exémplum. Vivus est enim sermo Dei, et éfficax et penetrabílior omni gládio ancípiti: et pertíngens usque ad divisiónem ánimæ ac spíritus: compágum quoque ac medullárum, et discrétor cogitatiónum et intentiónum cordis.\par
\switchcolumn
\EN
\noindent For if Jesus had given them rest, he would never have afterwards spoken of another day. There remaineth therefore a day of rest for the people of God. For he that is entered into his rest, the same also hath rested from his works, as God did from his. Let us hasten therefore to enter into that rest; lest any man fall into the same example of unbelief. For the word of God is living and effectual, and more piercing than any two edged sword; and reaching unto the division of the soul and the spirit, of the joints also and the marrow, and is a discerner of the thoughts and intents of the heart.\par
\switchcolumn*

\LA
\begin{Resp}\Rsign Benedícam Dóminum in omni témpore: \Star{} Semper laus ejus in ore meo.\end{Resp}
\begin{Resp}\Vsign In Dómino laudábitur ánima mea, áudiant mansuéti, et læténtur.\end{Resp}
\begin{Resp}\Rsign Semper laus ejus in ore meo.\end{Resp}
\begin{Resp}\Vsign Glória Patri, et Fílio, \Star{} et Spirítui Sancto.\end{Resp}
\begin{Resp}\Rsign Semper laus ejus in ore meo.\end{Resp}
\switchcolumn
\EN
\begin{Resp}\Rsign I will bless the Lord at all times. \Star{} His praise shall continually be in my mouth.\end{Resp}
\begin{Resp}\Vsign My soul shall make her boast in the Lord; the humble shall hear thereof and be glad.\end{Resp}
\begin{Resp}\Rsign His praise shall continually be in my mouth.\end{Resp}
\begin{Resp}\Vsign Glory be to the Father, and to the Son, \Star{} and to the Holy Ghost.\end{Resp}
\begin{Resp}\Rsign His praise shall continually be in my mouth.\end{Resp}
\switchcolumn*

\end{paracol}

\Office{Feria IV infra Hebdomadam VI post Epiphaniam}{Wednesday of the Sixth Week after Epiphany}{Feria}
\addcontentsline{toc}{subsection}{Feria IV infra Hebdomadam VI post Epiphaniam\enspace\textit{Wednesday of the Sixth Week after Epiphany}}
\begin{paracol}{2}
\LA
\LessonHead{Lectio I}
\switchcolumn
\EN
\LessonHead{Lesson I}
\switchcolumn*

\LA
\LessonTitle{De Epístola ad Hebrǽos}
\switchcolumn
\EN
\LessonTitle{Lesson from the letter of St. Paul the Apostle to the Hebrews}
\switchcolumn*

\LA
\Cite{Heb 6:1-3}
\switchcolumn
\EN
\Cite{Heb 6:1-3}
\switchcolumn*

\LA
\noindent Quaprópter intermitténtes inchoatiónis Christi sermónem, ad perfectióra ferámur, non rursum jaciéntes fundaméntum pœniténtiæ ab opéribus mórtuis, et fídei ad Deum, baptísmatum doctrínæ, impositiónis quoque mánuum, ac resurrectiónis mortuórum, et judícii ætérni. Et hoc faciémus, si quidem permíserit Deus.\par
\switchcolumn
\EN
\noindent Wherefore leaving the word of the beginning of Christ, let us go on to things more perfect, not laying again the foundation of penance from dead works, and of faith towards God, Of the doctrine of baptisms, and imposition of hands, and of the resurrection of the dead, and of eternal judgment. And this will we do, if God permit.\par
\switchcolumn*

\LA
\begin{Resp}\Rsign Ne perdíderis me cum iniquitátibus meis: \Star{} Neque in finem irátus resérves mala mea.\end{Resp}
\begin{Resp}\Vsign Non intres in judícium cum servo tuo, Dómine.\end{Resp}
\begin{Resp}\Rsign Neque in finem irátus resérves mala mea.\end{Resp}
\switchcolumn
\EN
\begin{Resp}\Rsign Cut me not off in the midst of my sins. \Star{} Nor keep thy wrath against me for my latter end.\end{Resp}
\begin{Resp}\Vsign Enter not into judgment with thy servant, O Lord.\end{Resp}
\begin{Resp}\Rsign Nor keep thy wrath against me for my latter end.\end{Resp}
\switchcolumn*

\LA
\LessonHead{Lectio II}
\switchcolumn
\EN
\LessonHead{Lesson II}
\switchcolumn*

\LA
\Cite{Heb 6:4-6}
\switchcolumn
\EN
\Cite{Heb 6:4-6}
\switchcolumn*

\LA
\noindent Impossíbile est enim eos qui semel sunt illumináti, gustavérunt étiam donum cæléste, et partícipes facti sunt Spíritus Sancti, gustavérunt nihilóminus bonum Dei verbum, virtutésque sǽculi ventúri, et prolápsi sunt; rursus renovári ad pœniténtiam, rursum crucifigéntes sibimetípsis Fílium Dei, et osténtui habéntes.\par
\switchcolumn
\EN
\noindent For it is impossible for those who were once illuminated, have tasted also the heavenly gift, and were made partakers of the Holy Ghost, Have moreover tasted the good word of God, and the powers of the world to come, And are fallen away: to be renewed again to penance, crucifying again to themselves the Son of God, and making him a mockery.\par
\switchcolumn*

\LA
\begin{Resp}\Rsign Parátum cor meum, Deus, parátum cor meum: \Star{} Cantábo et psalmum dicam Dómino.\end{Resp}
\begin{Resp}\Vsign Exsúrge, glória mea, exsúrge, psaltérium et cíthara, exsúrgam dilúculo.\end{Resp}
\begin{Resp}\Rsign Cantábo et psalmum dicam Dómino.\end{Resp}
\switchcolumn
\EN
\begin{Resp}\Rsign My heart is ready, O God, my heart is ready. \Star{} I will sing and give praise to the Lord.\end{Resp}
\begin{Resp}\Vsign Awake up, my glory, awake, psaltery and harp! I will awake early.\end{Resp}
\begin{Resp}\Rsign I will sing and give praise to the Lord.\end{Resp}
\switchcolumn*

\LA
\LessonHead{Lectio III}
\switchcolumn
\EN
\LessonHead{Lesson III}
\switchcolumn*

\LA
\Cite{Heb 6:7-10}
\switchcolumn
\EN
\Cite{Heb 6:7-10}
\switchcolumn*

\LA
\noindent Terra enim sæpe veniéntem super se bibens imbrem, et génerans herbam opportúnam illis, a quibus cólitur, áccipit benedictiónem a Deo: próferens autem spinas ac tríbulos, réproba est, et maledícto próxima: cujus consummátio in combustiónem. Confídimus autem de vobis dilectíssimi melióra, et vicinióra salúti: tamétsi ita lóquimur. Non enim injústus Deus, ut obliviscátur óperis vestri, et dilectiónis, quam ostendístis in nómine ipsíus, qui ministrástis sanctis, et ministrátis.\par
\switchcolumn
\EN
\noindent For the earth that drinketh in the rain which cometh often upon it, and bringeth forth herbs meet for them by whom it is tilled, receiveth blessing from God. But that which bringeth forth thorns and briers, is reprobate, and very near unto a curse, whose end is to be burnt. But, my dearly beloved, we trust better things of you, and nearer to salvation; though we speak thus. For God is not unjust, that he should forget your work, and the love which you have shewn in his name, you who have ministered, and do minister to the saints.\par
\switchcolumn*

\LA
\begin{Resp}\Rsign Adjútor meus, tibi psallam, quia, Deus, suscéptor meus es: \Star{} Deus meus, misericórdia mea.\end{Resp}
\begin{Resp}\Vsign Lætábor, et exsultábo in te, psallam nómini tuo, Altíssime.\end{Resp}
\begin{Resp}\Rsign Deus meus, misericórdia mea.\end{Resp}
\begin{Resp}\Vsign Glória Patri, et Fílio, \Star{} et Spirítui Sancto.\end{Resp}
\begin{Resp}\Rsign Deus meus, misericórdia mea.\end{Resp}
\switchcolumn
\EN
\begin{Resp}\Rsign Unto thee, O my Strength, will I sing, for God is my defence, \Star{} The God of my mercy.\end{Resp}
\begin{Resp}\Vsign I will be glad and rejoice in thee, I will sing praise to thy Name, O Thou Most High.\end{Resp}
\begin{Resp}\Rsign The God of my mercy.\end{Resp}
\begin{Resp}\Vsign Glory be to the Father, and to the Son, \Star{} and to the Holy Ghost.\end{Resp}
\begin{Resp}\Rsign The God of my mercy.\end{Resp}
\switchcolumn*

\end{paracol}

\Office{Feria V infra Hebdomadam VI post Epiphaniam}{Thursday of the Sixth Week after Epiphany}{Feria}
\addcontentsline{toc}{subsection}{Feria V infra Hebdomadam VI post Epiphaniam\enspace\textit{Thursday of the Sixth Week after Epiphany}}
\begin{paracol}{2}
\LA
\LessonHead{Lectio I}
\switchcolumn
\EN
\LessonHead{Lesson I}
\switchcolumn*

\LA
\LessonTitle{De Epístola ad Hebrǽos}
\switchcolumn
\EN
\LessonTitle{Lesson from the letter of St. Paul the Apostle to the Hebrews}
\switchcolumn*

\LA
\Cite{Heb 7:1-3}
\switchcolumn
\EN
\Cite{Heb 7:1-3}
\switchcolumn*

\LA
\noindent Hic enim Melchísedech, rex Salem, sacérdos Dei summi, qui obviávit Abrahæ regrésso a cæde regum, et benedíxit ei: cui et décimas ómnium divísit Abraham: primum quidem qui interpretátur rex justítiæ: deínde autem et rex Salem, quod est, rex pacis, sine patre, sine matre, sine genealógia, neque inítium diérum, neque finem vitæ habens, assimilátus autem Fílio Dei, manet sacérdos in perpétuum.\par
\switchcolumn
\EN
\noindent For this Melchisedech was king of Salem, priest of the most high God, who met Abraham returning from the slaughter of the kings, and blessed him: To whom also Abraham divided the tithes of all: who first indeed by interpretation, is king of justice: and then also king of Salem, that is, king of peace: Without father, without mother, without genealogy, having neither beginning of days nor end of life, but likened unto the Son of God, continueth a priest for ever.\par
\switchcolumn*

\LA
\begin{Resp}\Rsign Deus, in te sperávi, Dómine, non confúndar in ætérnum: in justítia tua líbera me, \Star{} Et éripe me.\end{Resp}
\begin{Resp}\Vsign Inclína ad me aurem tuam, et salva me.\end{Resp}
\begin{Resp}\Rsign Et éripe me.\end{Resp}
\switchcolumn
\EN
\begin{Resp}\Rsign In thee, O God, do I put my trust; let me never be put to confusion, O Lord deliver me in thy righteousness, \Star{} And cause me to escape.\end{Resp}
\begin{Resp}\Vsign Incline thine ear unto me, deliver me speedily.\end{Resp}
\begin{Resp}\Rsign And cause me to escape.\end{Resp}
\switchcolumn*

\LA
\LessonHead{Lectio II}
\switchcolumn
\EN
\LessonHead{Lesson II}
\switchcolumn*

\LA
\Cite{Heb 7:4-6}
\switchcolumn
\EN
\Cite{Heb 7:4-6}
\switchcolumn*

\LA
\noindent Intuémini autem quantus sit hic, cui et décimas dedit de præcípuis Abraham patriárcha. Et quidem de fíliis Levi sacerdótium accipiéntes, mandátum habent décimas súmere a pópulo secúndum legem, id est, a frátribus suis: quamquam et ipsi exíerint de lumbis Abrahæ. Cujus autem generátio non annumerátur in eis, décimas sumpsit ab Abraham, et hunc, qui habébat repromissiónes, benedíxit.\par
\switchcolumn
\EN
\noindent Now consider how great this man is, to whom also Abraham the patriarch gave tithes out of the principal things. And indeed they that are of the sons of Levi, who receive the priesthood, have a commandment to take tithes of the people according to the law, that is to say, of their brethren: though they themselves also came out of the loins of Abraham. But he, whose pedigree is not numbered among them, received tithes of Abraham, and blessed him that had the promises.\par
\switchcolumn*

\LA
\begin{Resp}\Rsign Repleátur os meum laude tua, ut hymnum dicam glóriæ tuæ, tota die magnitúdinem tuam: noli me proícere in témpore senectútis: \Star{} Dum defécerit in me virtus mea, ne derelínquas me.\end{Resp}
\begin{Resp}\Vsign Gaudébunt lábia mea cum cantávero tibi.\end{Resp}
\begin{Resp}\Rsign Dum defécerit in me virtus mea, ne derelínquas me.\end{Resp}
\switchcolumn
\EN
\begin{Resp}\Rsign Let my mouth be filled with thy praise, that I may sing of thy glory, all the day long of thy greatness. Cast me not off in the time of old age; \Star{} Forsake me not when my strength faileth.\end{Resp}
\begin{Resp}\Vsign My lips shall be fain when I sing unto thee.\end{Resp}
\begin{Resp}\Rsign Forsake me not when my strength faileth.\end{Resp}
\switchcolumn*

\LA
\LessonHead{Lectio III}
\switchcolumn
\EN
\LessonHead{Lesson III}
\switchcolumn*

\LA
\Cite{Heb 7:7-12}
\switchcolumn
\EN
\Cite{Heb 7:7-12}
\switchcolumn*

\LA
\noindent Sine ulla autem contradictióne, quod minus est, a melióre benedícitur. Et hic quidem, décimas moriéntes hómines accípiunt: ibi autem contestátur, quia vivit. Et (ut ita dictum sit) per Abraham, et Levi, qui décimas accépit, decimátus est: adhuc enim in lumbis patris erat, quando obviávit ei Melchísedech. Si ergo consummátio per sacerdótium Levíticum erat (pópulus enim sub ipso legem accépit) quid adhuc necessárium fuit secúndum órdinem Melchísedech, álium súrgere sacerdótem, et non secúndum órdinem Aaron dici? Transláto enim sacerdótio, necésse est ut et legis translátio fiat.\par
\switchcolumn
\EN
\noindent And without all contradiction, that which is less, is blessed by the better. And here indeed, men that die, receive thithes: but there he hath witness, that he liveth. And (as it may be said) even Levi who received tithes, paid tithes in Abraham: For he was yet in the loins of his father, when Melchisedech met him. If then perfection was by the Levitical priesthood, (for under it the people received the law,) what further need was there that another priest should rise according to the order of Melchisedech, and not be called according to the order of Aaron? For the priesthood being translated, it is necessary that a translation also be made of the law.\par
\switchcolumn*

\LA
\begin{Resp}\Rsign Gaudébunt lábia mea cum cantávero tibi: \Star{} Et ánima mea, quam redemísti, Dómine.\end{Resp}
\begin{Resp}\Vsign Sed et lingua mea meditábitur justítiam tuam, tota die laudem tuam.\end{Resp}
\begin{Resp}\Rsign Et ánima mea, quam redemísti, Dómine.\end{Resp}
\begin{Resp}\Vsign Glória Patri, et Fílio, \Star{} et Spirítui Sancto.\end{Resp}
\begin{Resp}\Rsign Et ánima mea, quam redemísti, Dómine.\end{Resp}
\switchcolumn
\EN
\begin{Resp}\Rsign My lips shall be fain when I sing unto thee; \Star{} And my soul, which Thou, O Lord, hast redeemed.\end{Resp}
\begin{Resp}\Vsign My tongue shall also talk of thy righteousness, all the day long of thy praise.\end{Resp}
\begin{Resp}\Rsign And my soul, which Thou, O Lord, hast redeemed.\end{Resp}
\begin{Resp}\Vsign Glory be to the Father, and to the Son, \Star{} and to the Holy Ghost.\end{Resp}
\begin{Resp}\Rsign And my soul, which Thou, O Lord, hast redeemed.\end{Resp}
\switchcolumn*

\end{paracol}

\Office{Feria VI infra Hebdomadam VI post Epiphaniam}{Friday of the Sixth Week after Epiphany}{Feria}
\addcontentsline{toc}{subsection}{Feria VI infra Hebdomadam VI post Epiphaniam\enspace\textit{Friday of the Sixth Week after Epiphany}}
\begin{paracol}{2}
\LA
\LessonHead{Lectio I}
\switchcolumn
\EN
\LessonHead{Lesson I}
\switchcolumn*

\LA
\LessonTitle{De Epístola ad Hebrǽos}
\switchcolumn
\EN
\LessonTitle{Lesson from the letter of St. Paul the Apostle to the Hebrews}
\switchcolumn*

\LA
\Cite{Heb 11:1-4}
\switchcolumn
\EN
\Cite{Heb 11:1-4}
\switchcolumn*

\LA
\noindent Est autem fides sperandárum substántia rerum, arguméntum non apparéntium. In hac enim testimónium consecúti sunt senes. Fide intellígimus aptáta esse sǽcula verbo Dei: ut ex invisibílibus visibília fíerent. Fide plúrimam hóstiam Abel, quam Cain, óbtulit Deo, per quam testimónium consecútus est esse justus, testimónium perhibénte munéribus ejus Deo, et per illam defúnctus adhuc lóquitur.\par
\switchcolumn
\EN
\noindent Now faith is the substance of things to be hoped for, the evidence of things that appear not. For by this the ancients obtained a testimony. By faith we understand that the world was framed by the word of God; that from invisible things visible things might be made. By faith Abel offered to God a sacrifice exceeding that of Cain, by which he obtained a testimony that he was just, God giving testimony to his gifts; and by it he being dead yet speaketh.\par
\switchcolumn*

\LA
\begin{Resp}\Rsign Confitébor tibi, Dómine Deus, in toto corde meo, et honorificábo nomen tuum in ætérnum: \Star{} Quia misericórdia tua, Dómine, magna est super me.\end{Resp}
\begin{Resp}\Vsign Deus meus es tu, et confitébor tibi: Deus meus es tu, et exaltábo te.\end{Resp}
\begin{Resp}\Rsign Quia misericórdia tua, Dómine, magna est super me.\end{Resp}
\switchcolumn
\EN
\begin{Resp}\Rsign I will praise thee, O Lord my God, with all my heart, and I will glorify thy Name for evermore. \Star{} For great is thy mercy, O Lord, toward me.\end{Resp}
\begin{Resp}\Vsign Thou art my God, and I will praise thee; Thou art my God, and I will exalt thee.\end{Resp}
\begin{Resp}\Rsign For great is thy mercy, O Lord, toward me.\end{Resp}
\switchcolumn*

\LA
\LessonHead{Lectio II}
\switchcolumn
\EN
\LessonHead{Lesson II}
\switchcolumn*

\LA
\Cite{Heb 11:5-7}
\switchcolumn
\EN
\Cite{Heb 11:5-7}
\switchcolumn*

\LA
\noindent Fide Henoch translátus est ne vidéret mortem, et non inveniebátur, quia tránstulit illum Deus: ante translatiónem enim testimónium hábuit placuísse Deo. Sine fide autem impossíbile est placére Deo. Crédere enim opórtet accedéntem ad Deum quia est, et inquiréntibus se remunerátor sit. Fide Noe respónso accépto de iis quæ adhuc non videbántur, métuens aptávit arcam in salútem domus suæ, per quam damnávit mundum: et justítiæ, quæ per fidem est, hæres est institútus.\par
\switchcolumn
\EN
\noindent By faith Henoch was translated, that he should not see death; and he was not found, because God had translated him: for before his translation he had testimony that he pleased God. But without faith it is impossible to please God. For he that cometh to God, must believe that he is, and is a rewarder to them that seek him. By faith Noe, having received an answer concerning those things which as yet were not seen, moved with fear, framed the ark for the saving of his house, by the which he condemned the world; and was instituted heir of the justice which is by faith.\par
\switchcolumn*

\LA
\begin{Resp}\Rsign Misericórdia tua, Dómine, magna est super me: \Star{} Et liberásti ánimam meam ex inférno inferióri.\end{Resp}
\begin{Resp}\Vsign In die tribulatiónis meæ clamávi ad te, quia exaudísti me.\end{Resp}
\begin{Resp}\Rsign Et liberásti ánimam meam ex inférno inferióri.\end{Resp}
\switchcolumn
\EN
\begin{Resp}\Rsign Great, O Lord, is thy mercy toward me. \Star{} And Thou hast delivered my soul from the lowest hell.\end{Resp}
\begin{Resp}\Vsign In the day of my trouble I called upon thee, for Thou hast heard me.\end{Resp}
\begin{Resp}\Rsign And Thou hast delivered my soul from the lowest hell.\end{Resp}
\switchcolumn*

\LA
\LessonHead{Lectio III}
\switchcolumn
\EN
\LessonHead{Lesson III}
\switchcolumn*

\LA
\Cite{Heb 11:8-10}
\switchcolumn
\EN
\Cite{Heb 1:8-10}
\switchcolumn*

\LA
\noindent Fide qui vocátur Abraham obedívit in locum exíre, quem acceptúrus erat in hæreditátem: et éxiit, nésciens quo iret. Fide demorátus est in terra repromissiónis, tamquam in aliéna, in cásulis habitándo cum Isaac et Jacob cohærédibus repromissiónis ejúsdem. Exspectábat enim fundaménta habéntem civitátem: cujus ártifex et cónditor Deus.\par
\switchcolumn
\EN
\noindent But to the Son: thy throne, O God, is for ever and ever: a sceptre of justice is the sceptre of thy kingdom. Thou hast loved justice, and hated iniquity: therefore God, thy God, hath anointed thee with the oil of gladness above thy fellows. And: Thou in the beginning, O Lord, didst found the earth: and the works of thy hands are the heavens.\par
\switchcolumn*

\LA
\begin{Resp}\Rsign Factus est mihi Dóminus in refúgium: \Star{} Et Deus meus in auxílium spei meæ.\end{Resp}
\begin{Resp}\Vsign Erípuit me de inimícis meis fortíssimis, et factus est Dóminus protéctor meus.\end{Resp}
\begin{Resp}\Rsign Et Deus meus in auxílium spei meæ.\end{Resp}
\begin{Resp}\Vsign Glória Patri, et Fílio, \Star{} et Spirítui Sancto.\end{Resp}
\begin{Resp}\Rsign Et Deus meus in auxílium spei meæ.\end{Resp}
\switchcolumn
\EN
\begin{Resp}\Rsign The Lord is my refuge. \Star{} And my God is the stay of my trust.\end{Resp}
\begin{Resp}\Vsign He delivered me from the strongest of mine enemies, and the Lord was my stay.\end{Resp}
\begin{Resp}\Rsign And my God is the stay of my trust.\end{Resp}
\begin{Resp}\Vsign Glory be to the Father, and to the Son, \Star{} and to the Holy Ghost.\end{Resp}
\begin{Resp}\Rsign And my God is the stay of my trust.\end{Resp}
\switchcolumn*

\end{paracol}

\Office{Sabbato infra Hebdomadam VI post Epiphaniam}{Saturday of the Sixth Week after Epiphany}{Feria}
\addcontentsline{toc}{subsection}{Sabbato infra Hebdomadam VI post Epiphaniam\enspace\textit{Saturday of the Sixth Week after Epiphany}}
\begin{paracol}{2}
\LA
\LessonHead{Lectio I}
\switchcolumn
\EN
\LessonHead{Lesson I}
\switchcolumn*

\LA
\LessonTitle{De Epístola ad Hebrǽos}
\switchcolumn
\EN
\LessonTitle{Lesson from the letter of St. Paul the Apostle to the Hebrews}
\switchcolumn*

\LA
\Cite{Heb 13:1-4}
\switchcolumn
\EN
\Cite{Heb 13:1-4}
\switchcolumn*

\LA
\noindent Cáritas fraternitátis máneat in vobis, et hospitalitátem nolíte oblivísci: per hanc enim latuérunt quidam, ángelis hospítio recéptis. Mementóte vinctórum, tamquam simul vincti: et laborántium, tamquam et ipsi in córpore morántes. Honorábile connúbium in ómnibus, et thorus immaculátus. Fornicatóres enim, et adúlteros judicábit Deus.\par
\switchcolumn
\EN
\noindent Let the charity of the brotherhood abide in you. And hospitality do not forget; for by this some, being not aware of it, have entertained angels. Remember them that are in bands, as if you were bound with them; and them that labour, as being yourselves also in the body. Marriage honourable in all, and the bed undefiled. For fornicators and adulterers God will judge.\par
\switchcolumn*

\LA
\begin{Resp}\Rsign Misericórdiam et judícium cantábo tibi, Dómine: \Star{} Psallam et intéllegam in via immaculáta, quando vénies ad me.\end{Resp}
\begin{Resp}\Vsign Perambulábam in innocéntia cordis mei, in médio domus meæ.\end{Resp}
\begin{Resp}\Rsign Psallam et intéllegam in via immaculáta, quando vénies ad me.\end{Resp}
\switchcolumn
\EN
\begin{Resp}\Rsign Mercy and judgment I will sing to thee, O Lord: \Star{} I will sing, and I will understand in the unspotted way, when thou shalt come to me.\end{Resp}
\begin{Resp}\Vsign I walked in the innocence of my heart, in the midst of my house.\end{Resp}
\begin{Resp}\Rsign Mercy and judgment I will sing to thee, O Lord\end{Resp}
\switchcolumn*

\LA
\LessonHead{Lectio II}
\switchcolumn
\EN
\LessonHead{Lesson II}
\switchcolumn*

\LA
\Cite{Heb 13:5-8}
\switchcolumn
\EN
\Cite{Heb 13:5-8}
\switchcolumn*

\LA
\noindent Sint mores sine avarítia, conténti præséntibus: ipse enim dixit: Non te déseram, neque derelínquam: ita ut confidénter dicámus: Dóminus mihi adjútor: non timébo quid fáciat mihi homo. Mementóte præpositórum vestrórum, qui vobis locúti sunt verbum Dei: quorum intuéntes éxitum conversatiónis, imitámini fidem. Jesus Christus heri, et hódie: ipse et in sǽcula.\par
\switchcolumn
\EN
\noindent Let your manners be without covetousness, contented with such things as you have; for he hath said: I will not leave thee, neither will I forsake thee. So that we may confidently say: The Lord is my helper: I will not fear what man shall do to me. Remember your prelates who have spoken the word of God to you; whose faith follow, considering the end of their conversation, Jesus Christ, yesterday, and to day; and the same for ever.\par
\switchcolumn*

\LA
\begin{Resp}\Rsign Dómine, exáudi oratiónem meam, et clamor meus ad te pervéniat: \Star{} Quia non spernis, Deus, preces páuperum.\end{Resp}
\begin{Resp}\Vsign Fiant aures tuæ intendéntes in oratiónem servi tui.\end{Resp}
\begin{Resp}\Rsign Quia non spernis, Deus, preces páuperum.\end{Resp}
\begin{Resp}\Vsign Glória Patri, et Fílio, \Star{} et Spirítui Sancto.\end{Resp}
\begin{Resp}\Rsign Quia non spernis, Deus, preces páuperum.\end{Resp}
\switchcolumn
\EN
\begin{Resp}\Rsign Hear my prayer, O Lord, and let my crying come unto thee. \Star{} For thou wilt not despise, O God, the prayer of the destitute\end{Resp}
\begin{Resp}\Vsign O let thine ears consider well the voice of my complaint\end{Resp}
\begin{Resp}\Rsign For thou wilt not despise, O God, the prayer of the destitute.\end{Resp}
\begin{Resp}\Vsign Glory be to the Father, and to the Son, \Star{} and to the Holy Ghost.\end{Resp}
\begin{Resp}\Rsign For thou wilt not despise, O God, the prayer of the destitute.\end{Resp}
\switchcolumn*

\end{paracol}

\SeasonTitle{Tempus Septuagesimæ}{Septuagesima}

\addcontentsline{toc}{section}{Tempus Septuagesimæ\enspace\textit{Septuagesima}}

\Office{Dominica in Septuagesima}{Septuagesima Sunday}{Semiduplex II classis}
\addcontentsline{toc}{subsection}{Dominica in Septuagesima\enspace\textit{Septuagesima Sunday}}
\begin{paracol}{2}
\LA
\LessonHead{Lectio I}
\switchcolumn
\EN
\LessonHead{Lesson I}
\switchcolumn*

\LA
\LessonTitle{Incipit liber Génesis}
\switchcolumn
\EN
\LessonTitle{Lesson from the book of Genesis}
\switchcolumn*

\LA
\Cite{Gen 1:1-8}
\switchcolumn
\EN
\Cite{Gen 1:1-8}
\switchcolumn*

\LA
\noindent In princípio creávit Deus cælum et terram. Terra autem erat inánis et vácua, et ténebræ erant super fáciem abýssi: et spíritus Dei ferebátur super aquas. Dixítque Deus: Fiat lux. Et facta est lux. Et vidit Deus lucem quod esset bona: et divísit lucem a ténebris. Appellavítque lucem Diem, et ténebras Noctem: factúmque est véspere et mane, dies unus. Dixit quoque Deus: Fiat firmaméntum in médio aquárum: et dívidat aquas ab aquis. Et fecit Deus firmaméntum, divisítque aquas, quæ erant sub firmaménto, ab his, quæ erant super firmaméntum. Et factum est ita. Vocavítque Deus firmaméntum, Cælum: et factum est véspere et mane, dies secúndus.\par
\switchcolumn
\EN
\noindent In the beginning God created heaven, and earth. And the earth was void and empty, and darkness was upon the face of the deep; and the spirit of God moved over the waters. And God said: Be light made. And light was made. And God saw the light that it was good; and he divided the light from the darkness. And he called the light Day, and the darkness Night; and there was evening and morning one day. And God said: Let there be a firmament made amidst the waters: and let it divide the waters from the waters. And God made a firmament, and divided the waters that were under the firmament, from those that were above the firmament, and it was so. And God called the firmament, Heaven; and the evening and morning were the second day.\par
\switchcolumn*

\LA
\begin{Resp}\Rsign In princípio creávit Deus cælum, et terram, et fecit in ea hóminem, \Star{} Ad imáginem et similitúdinem suam.\end{Resp}
\begin{Resp}\Vsign Formávit ígitur Deus hóminem de limo terræ, et inspirávit in fáciem ejus spiráculum vitæ.\end{Resp}
\begin{Resp}\Rsign Ad imáginem et similitúdinem suam.\end{Resp}
\switchcolumn
\EN
\begin{Resp}\Rsign In the beginning God created the heavens and the earth, wherein He made man also \Star{} After His own image and likeness.\end{Resp}
\begin{Resp}\Vsign So God formed man of the dust of the ground, and breathed into his face the breath of life.\end{Resp}
\begin{Resp}\Rsign After His own image and likeness.\end{Resp}
\switchcolumn*

\LA
\LessonHead{Lectio II}
\switchcolumn
\EN
\LessonHead{Lesson II}
\switchcolumn*

\LA
\Cite{Gen 1:9-19}
\switchcolumn
\EN
\Cite{Gen 1:9-19}
\switchcolumn*

\LA
\noindent Dixit vero Deus: Congregéntur aquæ, quæ sub cælo sunt, in locum unum: et appáreat árida. Et factum est ita. Et vocávit Deus áridam Terram, congregationésque aquárum appellávit Mária. Et vidit Deus quod esset bonum. Et ait: Gérminet terra herbam viréntem, et faciéntem semen, et lignum pomíferum fáciens fructum juxta genus suum, cujus semen in semetípso sit super terram. Et factum est ita. Et prótulit terra herbam viréntem, et faciéntem semen juxta genus suum, lignúmque fáciens fructum, et habens unumquódque seméntem secúndum spéciem suam. Et vidit Deus quod esset bonum. Et factum est véspere et mane, dies tértius. Dixit autem Deus: Fiant luminária in firmaménto cæli, et dívidant diem ac noctem, et sint in signa et témpora, et dies et annos: ut lúceant in firmaménto cæli, et illúminent terram. Et factum est ita. Fecítque Deus duo luminária magna: lumináre majus, ut præésset diéi: et lumináre minus, ut præésset nocti: et stellas. Et pósuit eas in firmaménto cæli, ut lucérent super terram, et præéssent diéi ac nocti, et divíderent lucem ac ténebras. Et vidit Deus quod esset bonum. Et factum est véspere et mane, dies quartus.\par
\switchcolumn
\EN
\noindent God also said: Let the waters that are under the heaven, be gathered together into one place: and let the dry land appear. And it was so done. And God called the dry land, Earth; and the gathering together of the waters, he called Seas. And God saw that it was good. And he said: Let the earth bring forth the green herb, and such as may seed, and the fruit tree yielding fruit after its kind, which may have seed in itself upon the earth. And it was so done. And the earth brought forth the green herb, and such as yieldeth seed according to its kind, and the tree that beareth fruit having seed each one according to its kind. And God saw that it was good. And the evening and the morning were the third day. And God said: Let there be lights made in the firmament of heaven, to divide the day and the night, and let them be for signs, and for seasons, and for days and years: To shine in the firmament of heaven, and to give light upon the earth. And it was so done. And God made two great lights: a greater light to rule the day; and a lesser light to rule the night: and the stars. And he set them in the firmament of heaven to shine upon the earth. And to rule the day and the night, and to divide the light and the darkness. And God saw that it was good. And the evening and morning were the fourth day.\par
\switchcolumn*

\LA
\begin{Resp}\Rsign In princípio creávit Deus cælum et terram, et Spíritus Dei ferebátur super aquas: \Star{} Et vidit Deus cuncta quæ fécerat, et erant valde bona.\end{Resp}
\begin{Resp}\Vsign Igitur perfécti sunt cæli et terra, et omnis ornátus eórum.\end{Resp}
\begin{Resp}\Rsign Et vidit Deus cuncta quæ fécerat, et erant valde bona.\end{Resp}
\switchcolumn
\EN
\begin{Resp}\Rsign In the beginning God created the heavens and the earth, and the Spirit of God moved upon the face of the waters. \Star{} And God saw everything that He had made, and it was very good.\end{Resp}
\begin{Resp}\Vsign Thus the heavens and the earth were finished, and all the hosts of them.\end{Resp}
\begin{Resp}\Rsign And God saw everything that He had made, and it was very good.\end{Resp}
\switchcolumn*

\LA
\LessonHead{Lectio III}
\switchcolumn
\EN
\LessonHead{Lesson III}
\switchcolumn*

\LA
\Cite{Gen 1:20-26}
\switchcolumn
\EN
\Cite{Gen 1:20-26}
\switchcolumn*

\LA
\noindent Dixit étiam Deus: Prodúcant aquæ réptile ánimæ vivéntis, et volátile super terram sub firmaménto cæli. Creavítque Deus cete grándia, et omnem ánimam vivéntem atque motábilem, quam prodúxerant aquæ in spécies suas, et omne volátile secúndum genus suum. Et vidit Deus quod esset bonum. Benedixítque eis, dicens: Créscite, et multiplicámini, et repléte aquas maris: avésque multiplicéntur super terram. Et factum est véspere et mane, dies quintus. Dixit quoque Deus: Prodúcat terra ánimam vivéntem in génere suo, juménta, et reptília, et béstias terræ secúndum spécies suas. Factúmque est ita. Et fecit Deus béstias terræ juxta spécies suas, et juménta, et omne réptile terræ in génere suo. Et vidit Deus quod esset bonum, et ait: Faciámus hóminem ad imáginem et similitúdinem nostram: et præsit píscibus maris, et volatílibus cæli, et béstiis, universǽque terræ, omníque réptili, quod movétur in terra.\par
\switchcolumn
\EN
\noindent God also said: Let the waters bring forth the creeping creature having life, and the fowl that may fly over the earth under the firmament of heaven. And God created the great whales, and every living and moving creature, which the waters brought forth, according to their kinds, and every winged fowl according to its kind. And God saw that it was good. And he blessed them, saying: Increase and multiply, and fill the waters of the sea: and let the birds be multiplied upon the earth. And the evening and morning were the fifth day. And God said: Let the earth bring forth the living creature in its kind, cattle and creeping things, and beasts of the earth, according to their kinds. And it was so done. And God made the beasts of the earth according to their kinds, and cattle, and every thing that creepeth on the earth after its kind. And God saw that it was good. And he said: Let us make man to our image and likeness: and let him have dominion over the fishes of the sea, and the fowls of the air, and the beasts, and the whole earth, and every creeping creature that moveth upon the earth.\par
\switchcolumn*

\LA
\begin{Resp}\Rsign Formávit Dóminus hóminem de limo terræ, \Star{} Et inspirávit in fáciem ejus spiráculum vitæ, et factus est homo in ánimam vivéntem.\end{Resp}
\begin{Resp}\Vsign In princípio fecit Deus cælum et terram, et plasmávit in ea hóminem.\end{Resp}
\begin{Resp}\Rsign Et inspirávit in fáciem ejus spiráculum vitæ, et factus est homo in ánimam vivéntem.\end{Resp}
\begin{Resp}\Vsign Glória Patri, et Fílio, \Star{} et Spirítui Sancto.\end{Resp}
\begin{Resp}\Rsign Et inspirávit in fáciem ejus spiráculum vitæ, et factus est homo in ánimam vivéntem.\end{Resp}
\switchcolumn
\EN
\begin{Resp}\Rsign The Lord formed man of the dust of the ground \Star{} And breathed into his face the breath of life, and man became a living soul.\end{Resp}
\begin{Resp}\Vsign In the beginning God created the heavens and the earth, wherein He made man also.\end{Resp}
\begin{Resp}\Rsign And breathed into his face the breath of life, and man became a living soul.\end{Resp}
\begin{Resp}\Vsign Glory be to the Father, and to the Son, \Star{} and to the Holy Ghost.\end{Resp}
\begin{Resp}\Rsign And breathed into his face the breath of life, and man became a living soul.\end{Resp}
\switchcolumn*

\LA
\LessonHead{Lectio IV}
\switchcolumn
\EN
\LessonHead{Lesson IV}
\switchcolumn*

\LA
\LessonTitle{Ex libro Enchirídii sancti Augustíni Epíscopi}
\switchcolumn
\EN
\LessonTitle{From the Handbook, written by St. Augustine, Bishop (of Hippo.)}
\switchcolumn*

\LA
\Source{Cap 25, 26 et 27 tomi 3}
\switchcolumn
\EN
\Source{Chaps. xxv., xxvi., xxvii., tom. 3}
\switchcolumn*

\LA
\noindent Mortis supplícium Dóminus hómini comminátus fúerat, si peccáret: sic eum múnerans líbero arbítrio, ut tamen régeret império, terréret exítio: atque in paradísi felicitáte, tamquam in umbra vitæ, unde justítia custodíta in melióra conscénderet, collocávit. Hinc post peccátum exsul efféctus, stirpem quoque suam, quam peccándo in se tamquam in radíce vitiáverat, pœna mortis et damnatióne obstrínxit: ut quidquid prolis ex illo, et simul damnáta, per quam peccáverat, cónjuge, per carnálem concupiscéntiam, in qua inobediéntiæ pœna símilis retribúta est, nascerétur, tráheret originále peccátum, quo traherétur per erróres dolorésque divérsos ad illud extrémum cum desertóribus ángelis, vitiatóribus et possessóribus et consórtibus suis, sine fine supplícium.\par
\switchcolumn
\EN
\noindent The Lord threatened man with the punishment of death, in case he sinned. Thus did He gift him with free will, while He yet kept His lordship over him, and helped him with the dread of destruction. And so He put him in that happy garden, under the very shadow of the tree of life, in that good place from whence, had he kept his righteousness, he might have passed to a better. But the first man sinned, and was banished from Eden, and infected all his descendants with the disease of sin, poisoning their very root, and bringing upon all that sentence of death and damnation, which he had earned for himself. So that all that descend by fleshly generation from Adam, and from the guilty woman, who was the cause of his sin and the partaker of his punishment, derive from them original sin; whereby they are drawn through a way of diverse sins and sorrows, towards that final ruin which they shall share with the rebel angels who are at once their corrupters, their lords, and their comrades.\par
\switchcolumn*

\LA
\begin{Resp}\Rsign Tulit Dóminus hóminem, et pósuit eum in paradíso voluptátis: \Star{} Ut operarétur et custodíret illum.\end{Resp}
\begin{Resp}\Vsign Plantáverat autem Dóminus Deus paradísum voluptátis a princípio, in quo pósuit hóminem quem formáverat.\end{Resp}
\begin{Resp}\Rsign Ut operarétur et custodíret illum.\end{Resp}
\switchcolumn
\EN
\begin{Resp}\Rsign God took the man and put him into the garden of Eden \Star{} To dress it and to keep it.\end{Resp}
\begin{Resp}\Vsign And the Lord God had planted a garden aforetime in Eden, and there He put the man whom He had formed.\end{Resp}
\begin{Resp}\Rsign To dress it and to keep it.\end{Resp}
\switchcolumn*

\LA
\LessonHead{Lectio V}
\switchcolumn
\EN
\LessonHead{Lesson V}
\switchcolumn*

\LA
\noindent Sic per unum hóminem peccátum intrávit in mundum, et per peccátum mors: et ita in omnes hómines pertránsiit, in quo omnes peccavérunt. Mundum quippe appellávit eo loco Apóstolus univérsum genus humánum. Ita ergo res se habébant. Jacébat in malis, vel étiam volvebátur, et de malis in mala præcipitabátur totíus humáni géneris massa damnáta: et adjúncta parti eórum, qui peccáverant, angelórum, luébat ímpiæ desertiónis digníssimas pœnas.\par
\switchcolumn
\EN
\noindent So by one man sin entered into the world, and death by sin, (and so death passed upon all men,) in whom all have sinned. (Rom. v. 12.) By the world the Apostle signified! in this place all mankind. Thus then hath the matter stood. The damned mass of humanity lay in misery, or rather wallowed in it, and fell from bad to worse, till it joined the company of the sinning angels, and both together suffered the deserved punishment of their vile treason.\par
\switchcolumn*

\LA
\begin{Resp}\Rsign Dixit Dóminus Deus: Non est bonum hóminem esse solum: \Star{} Faciámus ei adjutórium símile sibi.\end{Resp}
\begin{Resp}\Vsign Adæ vero non inveniebátur adjútor símilis sibi: dixit vero Deus.\end{Resp}
\begin{Resp}\Rsign Faciámus ei adjutórium símile sibi.\end{Resp}
\switchcolumn
\EN
\begin{Resp}\Rsign The Lord God said: It is not good that the man should be alone. \Star{} Let Us make an help meet for him.\end{Resp}
\begin{Resp}\Vsign But for Adam there was not found an help meet for him; and God said:\end{Resp}
\begin{Resp}\Rsign Let Us make an help meet for him.\end{Resp}
\switchcolumn*

\LA
\LessonHead{Lectio VI}
\switchcolumn
\EN
\LessonHead{Lesson VI}
\switchcolumn*

\LA
\noindent Ad iram quippe Dei pértinet justam, quidquid cæca et indómita concupiscéntia fáciunt libénter mali, et quidquid maniféstis opertísque pœnis patiúntur invíti: non sane Creatóris desisténte bonitáte, et malis ángelis subministráre vitam, vivacémque poténtiam, (quæ subministrátio si auferátur, interíbunt) et hóminum, quamvis de propágine vitiáta damnatáque nascéntium, formáre sémina, et animáre, et ordináre membra per témporum ætátes, per locórum spátia vegetáre sensus, aliménta donáre. Mélius enim judicáre de malis bene fácere, quam mala nulla esse permíttere.\par
\switchcolumn
\EN
\noindent So the wrath of God appertained whatever sin man, through the blind and untamed sting of his flesh, willingly committeth, and whatever punishment, declared and open, he unwillingly suffereth. There is, indeed, no pause in that goodness of the Creator whereby He giveth even to the traitor angels life and strength, (which if He gave not, they would be annihilated,) and whereby He formeth the seed of men, though they come of a corrupt and condemned stock, quickeneth them, strengtheneth and fitteth their limbs for the changing seasons of their life, extendeth their knowledge in diverse places, and giveth them whereon to live. It hath been His will rather to draw good out of evil, than to suffer that there should be no evil.\par
\switchcolumn*

\LA
\begin{Resp}\Rsign Immísit Dóminus sopórem in Adam, et tulit unam de costis ejus: † Et ædificávit costam, quam túlerat Dóminus de Adam, in mulíerem, et addúxit eam ad Adam, ut vidéret quid vocáret eam: \Star{} Et vocávit nomen ejus Virágo, quia de viro sumpta est.\end{Resp}
\begin{Resp}\Vsign Cumque obdormísset, tulit unam de costis ejus, et replévit carnem pro ea.\end{Resp}
\begin{Resp}\Rsign Et ædificávit costam, quam túlerat Dóminus de Adam, in mulíerem, et addúxit eam ad Adam, ut vidéret quid vocáret eam.\end{Resp}
\begin{Resp}\Vsign Glória Patri, et Fílio, \Star{} et Spirítui Sancto.\end{Resp}
\begin{Resp}\Rsign Et vocávit nomen ejus Virágo, quia de viro sumpta est.\end{Resp}
\switchcolumn
\EN
\begin{Resp}\Rsign The Lord caused a deep sleep to fall upon Adam, and He took one of his ribs. † And the rib which the Lord had taken from Adam made He a woman, and brought her unto Adam, to see what he would call her. \Star{} And he called her name Woman, because she was taken out of Man.\end{Resp}
\begin{Resp}\Vsign And while he slept He took one of his ribs, and closed up the flesh instead thereof.\end{Resp}
\begin{Resp}\Rsign And the rib which the Lord had taken from Adam made He a woman, and brought her unto Adam, to see what he would call her.\end{Resp}
\begin{Resp}\Vsign Glory be to the Father, and to the Son, \Star{} and to the Holy Ghost.\end{Resp}
\begin{Resp}\Rsign And he called her name Woman, because she was taken out of Man.\end{Resp}
\switchcolumn*

\LA
\LessonHead{Lectio VII}
\switchcolumn
\EN
\LessonHead{Lesson VII}
\switchcolumn*

\LA
\LessonTitle{Léctio sancti Evangélii secúndum Matthǽum}
\switchcolumn
\EN
\LessonTitle{From the Holy Gospel according to Matthew}
\switchcolumn*

\LA
\Cite{Matt 20:1-16}
\switchcolumn
\EN
\Cite{Matt 20:1-16}
\switchcolumn*

\LA
\noindent In illo témpore: Dixit Jesus discípulis suis parábolam hanc: Símile est regnum cælórum hómini patrifamílias, qui éxiit primo mane condúcere operários in víneam suam. Et réliqua.\par
\switchcolumn
\EN
\noindent In that time, Jesus said to his disciples: The kingdom of heaven is like to an householder, who went out early in the morning to hire labourers into his vineyard. And so on.\par
\switchcolumn*

\LA
\LessonTitle{Homilía sancti Gregórii Papæ}
\switchcolumn
\EN
\LessonTitle{Homily by Pope St. Gregory (the Great.)}
\switchcolumn*

\LA
\Source{Homil. 19 in Evangelium post princ.}
\switchcolumn
\EN
\Source{19th on the Gospels}
\switchcolumn*

\LA
\noindent Regnum cælórum hómini patrifamílias símile dícitur, qui ad excoléndam víneam suam operários condúcit. Quis vero patrisfamílias similitúdinem réctius tenet, quam Cónditor noster, qui regit quos cóndidit, et eléctos suos sic in hoc mundo póssidet, quasi subjéctos dóminus in domo? Qui habet víneam, universálem scílicet Ecclésiam, quæ ab Abel justo usque ad últimum eléctum, qui in fine mundi nascitúrus est, quot Sanctos prótulit, quasi tot pálmites misit.\par
\switchcolumn
\EN
\noindent We hear that the kingdom of heaven is like unto a man that is an householder, which went out early in the morning, to hire labourers into his vineyard. Who indeed is more justly to be likened to an householder than our Maker, Who is the Head of the household of faith, bearing rule over them whom He hath made, and being Master of His chosen ones in the world, as a Master over those that are in his house? He it is That hath the Church for a vineyard, a vineyard that ceaseth not to bring forth branches of the True Vine, from righteous Abel to the last of the elect that shall be born in the world.\par
\switchcolumn*

\LA
\begin{Resp}\Rsign Plantáverat autem Dóminus Deus paradísum voluptátis a princípio: \Star{} In quo pósuit hóminem, quem formáverat.\end{Resp}
\begin{Resp}\Vsign Produxítque Dóminus Deus de humo omne lignum pulchrum visu, et ad vescéndum suáve; lignum étiam vitæ in médio paradísi.\end{Resp}
\begin{Resp}\Rsign In quo pósuit hóminem, quem formáverat.\end{Resp}
\switchcolumn
\EN
\begin{Resp}\Rsign And the Lord God had planted a garden aforetime in Eden, \Star{} And there He put the man whom He had formed.\end{Resp}
\begin{Resp}\Vsign And out of the ground made the Lord God to grow every tree that is pleasant to the sight, and good for food, the tree of life also in the midst of the garden.\end{Resp}
\begin{Resp}\Rsign And there He put the man whom He had formed.\end{Resp}
\switchcolumn*

\LA
\LessonHead{Lectio VIII}
\switchcolumn
\EN
\LessonHead{Lesson VIII}
\switchcolumn*

\LA
\noindent Hic ítaque paterfamílias ad excoléndam víneam suam, mane, hora tértia, sexta, nona et undécima operários condúcit: quia a mundi hujus inítio usque in finem ad erudiéndam plebem fidélium, prædicatóres congregáre non desístit. Mane étenim mundi fuit ab Adam usque ad Noë: hora vero tértia a Noë usque ad Abraham: sexta quoque ab Abraham usque ad Móysen: nona autem a Móyse usque ad advéntum Dómini: undécima vero ab advéntu Dómini usque ad finem mundi. In qua prædicatóres sancti Apóstoli missi sunt, qui mercédem plenam et tarde veniéntes accepérunt.\par
\switchcolumn
\EN
\noindent This householder, then, for the cultivation of his vineyard, goeth out early in the morning, and at the third hour, and the sixth hour, and the ninth hour, and the eleventh hour, to hire labourers into his vineyard. Thus the Lord, from the beginning to the end of the world, ceaseth not to gather together preachers for the instruction of His faithful people. The early morning of the world was from Adam until Noah; the third hour from Noah until Abraham; the sixth hour from Abraham until Moses; the ninth hour from Moses until the coming of the Lord; the eleventh hour from the coming of the Lord until the end of the world. At this eleventh hour are sent forth as preachers the Holy Apostles, who have received full wages, albeit they be come in late.\par
\switchcolumn*

\LA
\begin{Resp}\Rsign Ecce Adam quasi unus ex nobis factus est sciens bonum et malum: \Star{} Vidéte, ne forte sumat de ligno vitæ, et vivat in ætérnum.\end{Resp}
\begin{Resp}\Vsign Fecit quoque Dóminus Deus Adæ túnicam pellíceam, et índuit eum, et dixit.\end{Resp}
\begin{Resp}\Rsign Vidéte, ne forte sumat de ligno vitæ, et vivat in ætérnum.\end{Resp}
\switchcolumn
\EN
\begin{Resp}\Rsign Behold, Adam is become as One of Us, to know good and evil. \Star{} See lest he take of the tree of life and live for ever.\end{Resp}
\begin{Resp}\Vsign Unto Adam also did the Lord God make a coat of skins, and clothed him, and said:\end{Resp}
\begin{Resp}\Rsign See lest he take of the tree of life and live for ever.\end{Resp}
\switchcolumn*

\LA
\LessonHead{Lectio IX}
\switchcolumn
\EN
\LessonHead{Lesson IX}
\switchcolumn*

\LA
\noindent Ad erudiéndam ergo Dóminus plebem suam, quasi ad excoléndam víneam suam, nullo témpore déstitit operários míttere: quia et prius per Patres, et póstmodum per legis Doctóres et Prophétas, ad extrémum vero per Apóstolos, dum plebis suæ mores excóluit, quasi per operários in víneæ cultúra laborávit: quamvis in quólibet módulo vel mensúra, quisquis cum fide recta bonæ prædicátor actiónis éxstitit, hujus víneæ operárius fuit. Operátor ergo mane, hora tértia, sexta, et nona, antíquus ille et Hebráicus pópulus designátur: qui in eléctis suis ab ipso mundi exórdio, dum recta fide Deum stúduit cólere, quasi non déstitit in víneæ cultúra laboráre. Ad undécimam vero Gentíles vocántur, quibus et dícitur: Quid hic statis tota die otiósi?\par
\switchcolumn
\EN
\noindent For the cultivation of His vineyard, (that is, the instruction of His people,) the Lord hath never ceased to send into it labourers. First, by the Fathers, then, by the Prophets and Teachers of the Law, and lastly, by the Apostles He hath dressed and tended the lives of His people, as the owner of a vineyard dresseth and tendeth it by means of workmen. Whoever in whatever degree joined to a right faith the teaching of righteousness, was so far one of God's labourers in God's vineyard. By the labourers at early morning, and at the third hour, the sixth hour, and the ninth hour, may be understood God's ancient people, the Hebrews, who strove to worship Him with a right faith in company with His chosen ones from the very beginning of the world, and thus continually laboured in His vineyard. And now, at the eleventh hour, it is said unto the Gentiles also Why stand ye here all the day idle?\par
\switchcolumn*

\LA
\begin{Resp}\Rsign Ubi est Abel frater tuus? dixit Dóminus ad Cain. Néscio, Dómine, numquid custos fratris mei sum ego? Et dixit ad eum: Quid fecísti? \Star{} Ecce vox sánguinis fratris tui Abel clamat ad me de terra.\end{Resp}
\begin{Resp}\Vsign Maledíctus eris super terram, quæ apéruit os suum, et suscépit sánguinem fratris tui de manu tua.\end{Resp}
\begin{Resp}\Rsign Ecce vox sánguinis fratris tui Abel clamat ad me de terra.\end{Resp}
\begin{Resp}\Vsign Glória Patri, et Fílio, \Star{} et Spirítui Sancto.\end{Resp}
\begin{Resp}\Rsign Ecce vox sánguinis fratris tui Abel clamat ad me de terra.\end{Resp}
\switchcolumn
\EN
\begin{Resp}\Rsign The Lord said unto Cain: Where is Abel thy brother? Lord, I know not am I my brother's keeper? And He said unto him What hast thou done? \Star{} Behold, the voice of thy brother Abel's blood crieth unto Me from the ground.\end{Resp}
\begin{Resp}\Vsign Cursed shalt thou be upon the earth, which hath opened her mouth to receive thy brother's blood from thy hand.\end{Resp}
\begin{Resp}\Rsign Behold, the voice of thy brother Abel's blood crieth unto Me from the ground.\end{Resp}
\begin{Resp}\Vsign Glory be to the Father, and to the Son, \Star{} and to the Holy Ghost.\end{Resp}
\begin{Resp}\Rsign Behold, the voice of thy brother Abel's blood crieth unto Me from the ground.\end{Resp}
\switchcolumn*

\end{paracol}

\Office{Feria II infra Hebdomadam Septuagesimæ}{Monday after Septuagesima Sunday}{Feria}
\addcontentsline{toc}{subsection}{Feria II infra Hebdomadam Septuagesimæ\enspace\textit{Monday after Septuagesima Sunday}}
\begin{paracol}{2}
\LA
\LessonHead{Lectio I}
\switchcolumn
\EN
\LessonHead{Lesson I}
\switchcolumn*

\LA
\LessonTitle{De libro Génesis}
\switchcolumn
\EN
\LessonTitle{Lesson from the book of Genesis}
\switchcolumn*

\LA
\Cite{Gen 1:27-31}
\switchcolumn
\EN
\Cite{Gen 1:27-31}
\switchcolumn*

\LA
\noindent Et creávit Deus hóminem ad imáginem suam: ad imáginem Dei creávit illum, másculum et féminam creávit eos. Benedixítque illis Deus, et ait: Créscite et multiplicámini, et repléte terram, et subícite eam, et dominámini píscibus maris, et volatílibus cæli, et univérsis animántibus, quæ movéntur super terram. Dixítque Deus: Ecce dedi vobis omnem herbam afferéntem semen super terram, et univérsa ligna quæ habent in semetípsis seméntem géneris sui, ut sint vobis in escam: et cunctis animántibus terræ, omníque vólucri cæli, et univérsis quæ movéntur in terra, et in quibus est ánima vivens, ut hábeant ad vescéndum. Et factum est ita. Vidítque Deus cuncta quæ fécerat, et erant valde bona. Et factum est véspere et mane, dies sextus.\par
\switchcolumn
\EN
\noindent And God created man to his own image: to the image of God he created him: male and female he created them. And God blessed them, saying: Increase and multiply, and fill the earth, and subdue it, and rule over the fishes of the sea, and the fowls of the air, and all living creatures that move upon the earth. And God said: Behold I have given you every herb bearing seed upon the earth, and all trees that have in themselves seed of their own kind, to be your meat: And to all the beasts of the earth, and to every fowl of the air, and to all that move upon the earth, and wherein there is life, that they may have to feed upon. And it was so done. And God saw all the things that he had made, and they were very good. And the evening and morning were the sixth day.\par
\switchcolumn*

\LA
\begin{Resp}\Rsign Dum deambuláret Dóminus in paradíso ad auram post merídiem, clamávit, et dixit: Adam, ubi es? Audívi, Dómine, vocem tuam, \Star{} Et abscóndi me.\end{Resp}
\begin{Resp}\Vsign Vocem tuam audívi in paradiso, et tímui, eo quod nudus essem.\end{Resp}
\begin{Resp}\Rsign Et abscóndi me.\end{Resp}
\switchcolumn
\EN
\begin{Resp}\Rsign When the Lord walked in the garden in the cool of the day, He called, and said Adam, where art thou? Lord, I heard thy voice. \Star{} And I hid myself.\end{Resp}
\begin{Resp}\Vsign I heard thy voice in the garden, and I was afraid, because I was naked.\end{Resp}
\begin{Resp}\Rsign And I hid myself.\end{Resp}
\switchcolumn*

\LA
\LessonHead{Lectio II}
\switchcolumn
\EN
\LessonHead{Lesson II}
\switchcolumn*

\LA
\Cite{Gen 2:1-6}
\switchcolumn
\EN
\Cite{Gen 2:1-6}
\switchcolumn*

\LA
\noindent Igitur perfécti sunt cæli et terra, et omnis ornátus eórum. Complevítque Deus die séptimo opus suum quod fécerat: et requiévit die séptimo ab univérso ópere quod patrárat. Et benedíxit diéi séptimo, et sanctificávit illum, quia in ipso cessáverat ab omni ópere suo quod creávit Deus ut fáceret. Istæ sunt generatiónes cæli et terræ, quando creáta sunt, in die quo fecit Dóminus Deus cælum et terram, et omne virgúltum agri ántequam orirétur in terra, omnémque herbam regiónis priúsquam germináret: non enim plúerat Dóminus Deus super terram, et homo non erat qui operarétur terram: sed fons ascendébat e terra, írrigans univérsam superfíciem terræ.\par
\switchcolumn
\EN
\noindent So the heavens and the earth were finished, and all the furniture of them. And on the seventh day God ended his work which he had made: and he rested on the seventh day from all his work which he had done. And he blessed the seventh day, and sanctified it: because in it he had rested from all his work which God created and made. These are the generations of the heaven and the earth, when they were created, in the day that the Lord God made the heaven and the earth: And every plant of the field before it spring up in the earth, and every herb of the ground before it grew: for the Lord God had not rained upon the earth; and there was not a man to till the earth. But a spring rose out the earth, watering all the surface of the earth.\par
\switchcolumn*

\LA
\begin{Resp}\Rsign In sudóre vultus tui vescéris pane tuo, dixit Dóminus ad Adam: cum operátus fúeris terram, non dabit fructus suos: \Star{} Sed spinas et tríbulos germinábit tibi.\end{Resp}
\begin{Resp}\Vsign Quia audísti vocem uxóris tuæ, et comedísti de ligno, ex quo præcéperam tibi ne coméderes, maledícta terra in opere tuo.\end{Resp}
\begin{Resp}\Rsign Sed spinas et tríbulos germinábit tibi.\end{Resp}
\switchcolumn
\EN
\begin{Resp}\Rsign The Lord said unto Adam: In the sweat of thy face shalt thou eat bread; when thou tillest the ground, it shall not henceforth yield unto thee her fruits. \Star{} Thorns also and thistles shall it bring forth to thee.\end{Resp}
\begin{Resp}\Vsign Because thou hast hearkened unto the voice of thy wife, and hast eaten of the tree which I commanded thee, saying: Thou shalt not eat of it, cursed is the ground whereon thou shalt labour.\end{Resp}
\begin{Resp}\Rsign Thorns also and thistles shall it bring forth to thee.\end{Resp}
\switchcolumn*

\LA
\LessonHead{Lectio III}
\switchcolumn
\EN
\LessonHead{Lesson III}
\switchcolumn*

\LA
\Cite{Gen 2:7-10}
\switchcolumn
\EN
\Cite{Gen 2:7-10}
\switchcolumn*

\LA
\noindent Formávit ígitur Dóminus Deus hóminem de limo terræ, et inspirávit in fáciem ejus spiráculum vitæ, et factus est homo in ánimam vivéntem. Plantáverat autem Dóminus Deus paradísum voluptátis a princípio, in quo pósuit hóminem quem formáverat. Produxítque Dóminus Deus de humo omne lignum pulchrum visu, et ad vescéndum suave lignum étiam vitæ in médio paradísi, lignúmque sciéntiæ boni et mali. Et flúvius egrediebátur de loco voluptátis ad irrigándum paradísum, qui inde divíditur in quátuor cápita.\par
\switchcolumn
\EN
\noindent And the Lord God formed man of the slime of the earth: and breathed into his face the breath of life, and man became a living soul. And the Lord God had planted a paradise of pleasure from the beginning: wherein he placed man whom he had formed. And the Lord God brought forth of the ground all manner of trees, fair to behold, and pleasant to eat of: the tree of life also in the midst of paradise: and the tree of knowledge of good and evil. And a river went out the place of pleasure to water paradise, which from thence is divided into four heads.\par
\switchcolumn*

\LA
\begin{Resp}\Rsign Formávit Dóminus hóminem de limo terræ, \Star{} Et inspirávit in fáciem ejus spiráculum vitæ, et factus est homo in ánimam vivéntem.\end{Resp}
\begin{Resp}\Vsign In princípio fecit Deus cælum et terram, et plasmávit in ea hóminem.\end{Resp}
\begin{Resp}\Rsign Et inspirávit in fáciem ejus spiráculum vitæ, et factus est homo in ánimam vivéntem.\end{Resp}
\begin{Resp}\Vsign Glória Patri, et Fílio, \Star{} et Spirítui Sancto.\end{Resp}
\begin{Resp}\Rsign Et inspirávit in fáciem ejus spiráculum vitæ, et factus est homo in ánimam vivéntem.\end{Resp}
\switchcolumn
\EN
\begin{Resp}\Rsign The Lord formed man of the dust of the ground \Star{} And breathed into his face the breath of life, and man became a living soul.\end{Resp}
\begin{Resp}\Vsign In the beginning God created the heavens and the earth, wherein He made man also.\end{Resp}
\begin{Resp}\Rsign And breathed into his face the breath of life, and man became a living soul.\end{Resp}
\begin{Resp}\Vsign Glory be to the Father, and to the Son, \Star{} and to the Holy Ghost.\end{Resp}
\begin{Resp}\Rsign And breathed into his face the breath of life, and man became a living soul.\end{Resp}
\switchcolumn*

\end{paracol}

\Office{Feria III infra Hebdomadam Septuagesimæ}{Tuesday after Septuagesima Sunday}{Feria}
\addcontentsline{toc}{subsection}{Feria III infra Hebdomadam Septuagesimæ\enspace\textit{Tuesday after Septuagesima Sunday}}
\begin{paracol}{2}
\LA
\LessonHead{Lectio I}
\switchcolumn
\EN
\LessonHead{Lesson I}
\switchcolumn*

\LA
\LessonTitle{De libro Génesis}
\switchcolumn
\EN
\LessonTitle{Lesson from the book of Genesis}
\switchcolumn*

\LA
\Cite{Gen 2:15-18}
\switchcolumn
\EN
\Cite{Gen 2:15-18}
\switchcolumn*

\LA
\noindent Tulit ergo Dóminus Deus hóminem, et pósuit eum in paradíso voluptátis, ut operarétur, et custodíret illum: præcepítque ei, dicens: Ex omni ligno paradísi cómede; de ligno autem sciéntiæ boni et mali ne cómedas: in quocúmque enim die coméderis ex eo, morte moriéris. Dixit quoque Dóminus Deus: Non est bonum esse hóminem solum: faciámus ei adjutórium símile sibi.\par
\switchcolumn
\EN
\noindent And the Lord God took man, and put him into the paradise for pleasure, to dress it, and keep it. And he commanded him, saying: Of every tree of paradise thou shalt eat: But of the tree of knowledge of good and evil, thou shalt not eat. for in what day soever thou shalt eat of it, thou shalt die the death. And the Lord God said: It is not good for man to be alone: let us make him a help like unto himself.\par
\switchcolumn*

\LA
\begin{Resp}\Rsign Tulit Dóminus hóminem, et pósuit eum in paradíso voluptátis: \Star{} Ut operarétur et custodíret illum.\end{Resp}
\begin{Resp}\Vsign Plantáverat autem Dóminus Deus paradísum voluptátis a princípio, in quo pósuit hóminem quem formáverat.\end{Resp}
\begin{Resp}\Rsign Ut operarétur et custodíret illum.\end{Resp}
\switchcolumn
\EN
\begin{Resp}\Rsign God took the man and put him into the garden of Eden \Star{} To dress it and to keep it.\end{Resp}
\begin{Resp}\Vsign And the Lord God had planted a garden aforetime in Eden, and there He put the man whom He had formed.\end{Resp}
\begin{Resp}\Rsign To dress it and to keep it.\end{Resp}
\switchcolumn*

\LA
\LessonHead{Lectio II}
\switchcolumn
\EN
\LessonHead{Lesson II}
\switchcolumn*

\LA
\Cite{Gen 2:19-20}
\switchcolumn
\EN
\Cite{Gen 2:19-20}
\switchcolumn*

\LA
\noindent Formátis ígitur Dóminus Deus de humo cunctis animántibus terræ, et univérsis volatílibus cæli, addúxit ea ad Adam, ut vidéret quid vocáret ea: omne enim quod vocávit Adam ánimæ vivéntis, ipsum est nomen ejus. Appellavítque Adam nomínibus suis cuncta animántia, et univérsa volatília cæli, et omnes béstias terræ: Adæ vero non inveniebátur adjútor símilis ejus.\par
\switchcolumn
\EN
\noindent And the Lord God having formed out of the ground all the beasts of the earth, and all the fowls of the air, brought them to Adam to see what he would call them: for whatsoever Adam called any living creature the same is its name. And Adam called all the beasts by their names, and all the fowls of the air, and all the cattle of the field: but for Adam there was not found a helper like himself.\par
\switchcolumn*

\LA
\begin{Resp}\Rsign Dixit Dóminus Deus: Non est bonum hóminem esse solum: \Star{} Faciámus ei adjutórium símile sibi.\end{Resp}
\begin{Resp}\Vsign Adæ vero non inveniebátur adjútor símilis sibi: dixit vero Deus.\end{Resp}
\begin{Resp}\Rsign Faciámus ei adjutórium símile sibi.\end{Resp}
\switchcolumn
\EN
\begin{Resp}\Rsign The Lord God said: It is not good that the man should be alone. \Star{} Let Us make an help meet for him.\end{Resp}
\begin{Resp}\Vsign But for Adam there was not found an help meet for him; and God said:\end{Resp}
\begin{Resp}\Rsign Let Us make an help meet for him.\end{Resp}
\switchcolumn*

\LA
\LessonHead{Lectio III}
\switchcolumn
\EN
\LessonHead{Lesson III}
\switchcolumn*

\LA
\Cite{Gen 2:21-24}
\switchcolumn
\EN
\Cite{Gen 2:21-24}
\switchcolumn*

\LA
\noindent Immísit ergo Dóminus Deus sopórem in Adam: cumque obdormísset, tulit unam de costis ejus, et replévit carnem pro ea. Et ædificávit Dóminus Deus costam, quam túlerat de Adam, in mulíerem: et addúxit eam ad Adam. Dixítque Adam: Hoc nunc os ex óssibus meis, et caro de carne mea: hæc vocábitur Virágo, quóniam de viro sumpta est. Quam ob rem relínquet homo patrem suum, et matrem, et adhærébit uxóri suæ: et erunt duo in carne una.\par
\switchcolumn
\EN
\noindent Then the Lord God cast a deep sleep upon Adam: and when he was fast asleep, he took one of his ribs, and filled up flesh for it. And the Lord God built the rib which he took from Adam into a woman: and brought her to Adam. And Adam said: This now is bone of my bones, and flesh of my flesh; she shall be called woman, because she was taken out of man. Wherefore a man shall leave father and mother, and shall cleave to his wife: and they shall be two in one flesh.\par
\switchcolumn*

\LA
\begin{Resp}\Rsign Immísit Dóminus sopórem in Adam, et tulit unam de costis ejus: † Et ædificávit costam, quam túlerat Dóminus de Adam, in mulíerem, et addúxit eam ad Adam, ut vidéret quid vocáret eam: \Star{} Et vocávit nomen ejus Virágo, quia de viro sumpta est.\end{Resp}
\begin{Resp}\Vsign Cumque obdormísset, tulit unam de costis ejus, et replévit carnem pro ea.\end{Resp}
\begin{Resp}\Rsign Et ædificávit costam, quam túlerat Dóminus de Adam, in mulíerem, et addúxit eam ad Adam, ut vidéret quid vocáret eam.\end{Resp}
\begin{Resp}\Vsign Glória Patri, et Fílio, \Star{} et Spirítui Sancto.\end{Resp}
\begin{Resp}\Rsign Et vocávit nomen ejus Virágo, quia de viro sumpta est.\end{Resp}
\switchcolumn
\EN
\begin{Resp}\Rsign The Lord caused a deep sleep to fall upon Adam, and He took one of his ribs. † And the rib which the Lord had taken from Adam made He a woman, and brought her unto Adam, to see what he would call her. \Star{} And he called her name Woman, because she was taken out of Man.\end{Resp}
\begin{Resp}\Vsign And while he slept He took one of his ribs, and closed up the flesh instead thereof.\end{Resp}
\begin{Resp}\Rsign And the rib which the Lord had taken from Adam made He a woman, and brought her unto Adam, to see what he would call her.\end{Resp}
\begin{Resp}\Vsign Glory be to the Father, and to the Son, \Star{} and to the Holy Ghost.\end{Resp}
\begin{Resp}\Rsign And he called her name Woman, because she was taken out of Man.\end{Resp}
\switchcolumn*

\end{paracol}

\Office{Feria IV infra Hebdomadam Septuagesimæ}{Wednesday after Septuagesima Sunday}{Feria}
\addcontentsline{toc}{subsection}{Feria IV infra Hebdomadam Septuagesimæ\enspace\textit{Wednesday after Septuagesima Sunday}}
\begin{paracol}{2}
\LA
\LessonHead{Lectio I}
\switchcolumn
\EN
\LessonHead{Lesson I}
\switchcolumn*

\LA
\LessonTitle{De libro Génesis}
\switchcolumn
\EN
\LessonTitle{Lesson from the book of Genesis}
\switchcolumn*

\LA
\Cite{Gen 3:1-7}
\switchcolumn
\EN
\Cite{Gen 3:1-7}
\switchcolumn*

\LA
\noindent Sed et serpens erat callídior cunctis animántibus terræ quæ fécerat Dóminus Deus. Qui dixit ad mulíerem: Cur præcépit vobis Deus ut non comederétis de omni ligno paradísi? Cui respóndit múlier: De fructu lignórum, quæ sunt in paradíso, véscimur: de fructu vero ligni quod est in médio paradísi, præcépit nobis Deus ne comederémus, et ne tangerémus illud, ne forte moriámur. Dixit autem serpens ad mulíerem: Nequáquam morte moriémini. Scit enim Deus quod in quocúmque die comedéritis ex eo, aperiéntur óculi vestri, et éritis sicut dii, sciéntes bonum et malum. Vidit ígitur múlier quod bonum esset lignum ad vescéndum, et pulchrum óculis, aspectúque delectábile: et tulit de fructu illíus, et comédit: dedítque viro suo, qui comédit. Et apérti sunt óculi ambórum;\par
\switchcolumn
\EN
\noindent Now the serpent was more subtle than any of the beasts of the earth which the Lord God made. And he said to the woman: Why hath God commanded you, that you should not eat of every tree of paradise? And the woman answered him, saying: Of the fruit of the trees that are in paradise we do eat: But of the fruit of the tree which is in the midst of paradise, God hath commanded us that we should not eat; and that we should not touch it, lest perhaps we die. And the serpent said to the woman: No, you shall not die the death. For God doth know that in what day soever you shall eat thereof, your eyes shall be opened: and you shall be as Gods, knowing good and evil. And the woman saw that the tree was good to eat, and fair to the eyes, and delightful to behold: and she took of the fruit thereof, and did eat, and gave to her husband who did eat. And the eyes of them both were opened.\par
\switchcolumn*

\LA
\begin{Resp}\Rsign Plantáverat autem Dóminus Deus paradísum voluptátis a princípio: \Star{} In quo pósuit hóminem, quem formáverat.\end{Resp}
\begin{Resp}\Vsign Produxítque Dóminus Deus de humo omne lignum pulchrum visu, et ad vescéndum suáve; lignum étiam vitæ in médio paradísi.\end{Resp}
\begin{Resp}\Rsign In quo pósuit hóminem, quem formáverat.\end{Resp}
\switchcolumn
\EN
\begin{Resp}\Rsign And the Lord God had planted a garden aforetime in Eden, \Star{} And there He put the man whom He had formed.\end{Resp}
\begin{Resp}\Vsign And out of the ground made the Lord God to grow every tree that is pleasant to the sight, and good for food, the tree of life also in the midst of the garden.\end{Resp}
\begin{Resp}\Rsign And there He put the man whom He had formed.\end{Resp}
\switchcolumn*

\LA
\LessonHead{Lectio II}
\switchcolumn
\EN
\LessonHead{Lesson II}
\switchcolumn*

\LA
\Cite{Gen 3:7-13}
\switchcolumn
\EN
\Cite{Gen 3:7-13}
\switchcolumn*

\LA
\noindent Cumque cognovíssent se esse nudos, consuérunt fólia ficus, et fecérunt sibi perizómata. Et cum audíssent vocem Dómini Dei deambulántis in paradíso ad auram post merídiem, abscóndit se Adam et uxor ejus a fácie Dómini Dei in médio ligni paradísi. Vocavítque Dóminus Deus Adam, et dixit ei: Ubi es? Qui ait: Vocem tuam audívi in paradíso, et tímui, eo quod nudus essem, et abscóndi me. Cui dixit: Quis enim indicávit tibi quod nudus esses, nisi quod ex ligno de quo præcéperam tibi ne coméderes, comedísti? Dixítque Adam: Múlier, quam dedísti mihi sóciam, dedit mihi de ligno, et comédi. Et dixit Dóminus Deus ad mulíerem: Quare hoc fecísti? Quæ respóndit: Serpens decépit me, et comédi.\par
\switchcolumn
\EN
\noindent When they perceived themselves to be naked, they sewed together fig leaves, and made themselves aprons. And when they heard the voice of the Lord God walking in paradise at the afternoon air, Adam and his wife hid themselves from the face of the Lord God, amidst the trees of paradise. And the Lord God called Adam, and said to him: Where art thou? And he said: I heard thy voice in paradise; and I was afraid, because I was naked, and I hid myself. And he said to him: And who hath told thee that thou wast naked, but that thou hast eaten of the tree whereof I commanded thee that thou shouldst not eat? And Adam said: The woman, whom thou gavest me to be my companion, gave me of the tree, and I did eat. And the Lord God said to the woman: Why hast thou done this? And she answered: The serpent deceived me, and I did eat.\par
\switchcolumn*

\LA
\begin{Resp}\Rsign Ecce Adam quasi unus ex nobis factus est sciens bonum et malum: \Star{} Vidéte, ne forte sumat de ligno vitæ, et vivat in ætérnum.\end{Resp}
\begin{Resp}\Vsign Fecit quoque Dóminus Deus Adæ túnicam pellíceam, et índuit eum, et dixit.\end{Resp}
\begin{Resp}\Rsign Vidéte, ne forte sumat de ligno vitæ, et vivat in ætérnum.\end{Resp}
\switchcolumn
\EN
\begin{Resp}\Rsign Behold, Adam is become as One of Us, to know good and evil. \Star{} See lest he take of the tree of life and live for ever.\end{Resp}
\begin{Resp}\Vsign Unto Adam also did the Lord God make a coat of skins, and clothed him, and said:\end{Resp}
\begin{Resp}\Rsign See lest he take of the tree of life and live for ever.\end{Resp}
\switchcolumn*

\LA
\LessonHead{Lectio III}
\switchcolumn
\EN
\LessonHead{Lesson III}
\switchcolumn*

\LA
\Cite{Gen 3:14-20}
\switchcolumn
\EN
\Cite{Gen 3:14-20}
\switchcolumn*

\LA
\noindent Et ait Dóminus Deus ad serpéntem: Quia fecísti hoc, maledíctus es inter ómnia animántia, et béstias terræ: super pectus tuum gradiéris, et terram cómedes cunctis diébus vitæ tuæ. Inimicítias ponam inter te et mulíerem, et semen tuum et semen illíus: ipsa cónteret caput tuum, et tu insidiáberis calcáneo ejus. Mulíeri quoque dixit: Multiplicábo ærúmnas tuas, et concéptus tuos: in dolóre páries fílios, et sub viri potestáte eris, et ipse dominábitur tui. Adæ vero dixit: Quia audísti vocem uxóris tuæ, et comedísti de ligno, ex quo præcéperam tibi ne coméderes, maledícta terra in ópere tuo: in labóribus cómedes ex ea cunctis diébus vitæ tuæ. Spinas et tríbulos germinábit tibi, et cómedes herbam terræ. In sudóre vultus tui vescéris pane, donec revertáris in terram de qua sumptus es: quia pulvis es et in púlverem revertéris. Et vocávit Adam nomen uxóris suæ, Heva: eo quod mater esset cunctórum vivéntium.\par
\switchcolumn
\EN
\noindent And the Lord God said to the serpent: Because thou hast done this thing, thou art cursed among all cattle, and the beasts of the earth: upon thy breast shalt thou go, and earth shalt thou eat all the days of thy life. I will put enmities between thee and the woman, and thy seed and her seed: she shall crush thy head, and thou shalt lie in wait for her heel. To the woman also he said: I will multiply thy sorrows, and thy conceptions: in sorrow shalt thou bring forth children, and thou shalt be under thy husband's power, and he shall have dominion over thee. And to Adam he said: Because thou hast hearkened to the voice of thy wife, and hast eaten of the tree, whereof I commanded thee that thou shouldst not eat, cursed is the earth in thy work; with labour and toil shalt thou eat thereof all the days of thy life. Thorns and thistles shall it bring forth to thee; and thou eat the herbs of the earth. In the sweat of thy face shalt thou eat bread till thou return to the earth, out of which thou wast taken: for dust thou art, and into dust thou shalt return. And Adam called the name of his wife Eve: because she was the mother of all the living.\par
\switchcolumn*

\LA
\begin{Resp}\Rsign Ubi est Abel frater tuus? dixit Dóminus ad Cain. Néscio, Dómine, numquid custos fratris mei sum ego? Et dixit ad eum: Quid fecísti? \Star{} Ecce vox sánguinis fratris tui Abel clamat ad me de terra.\end{Resp}
\begin{Resp}\Vsign Maledíctus eris super terram, quæ apéruit os suum, et suscépit sánguinem fratris tui de manu tua.\end{Resp}
\begin{Resp}\Rsign Ecce vox sánguinis fratris tui Abel clamat ad me de terra.\end{Resp}
\begin{Resp}\Vsign Glória Patri, et Fílio, \Star{} et Spirítui Sancto.\end{Resp}
\begin{Resp}\Rsign Ecce vox sánguinis fratris tui Abel clamat ad me de terra.\end{Resp}
\switchcolumn
\EN
\begin{Resp}\Rsign The Lord said unto Cain: Where is Abel thy brother? Lord, I know not am I my brother's keeper? And He said unto him What hast thou done? \Star{} Behold, the voice of thy brother Abel's blood crieth unto Me from the ground.\end{Resp}
\begin{Resp}\Vsign Cursed shalt thou be upon the earth, which hath opened her mouth to receive thy brother's blood from thy hand.\end{Resp}
\begin{Resp}\Rsign Behold, the voice of thy brother Abel's blood crieth unto Me from the ground.\end{Resp}
\begin{Resp}\Vsign Glory be to the Father, and to the Son, \Star{} and to the Holy Ghost.\end{Resp}
\begin{Resp}\Rsign Behold, the voice of thy brother Abel's blood crieth unto Me from the ground.\end{Resp}
\switchcolumn*

\end{paracol}

\Office{Feria V infra Hebdomadam Septuagesimæ}{Thursday after Septuagesima Sunday}{Feria}
\addcontentsline{toc}{subsection}{Feria V infra Hebdomadam Septuagesimæ\enspace\textit{Thursday after Septuagesima Sunday}}
\begin{paracol}{2}
\LA
\LessonHead{Lectio I}
\switchcolumn
\EN
\LessonHead{Lesson I}
\switchcolumn*

\LA
\LessonTitle{De libro Génesis}
\switchcolumn
\EN
\LessonTitle{Lesson from the book of Genesis}
\switchcolumn*

\LA
\Cite{Gen 4:1-7}
\switchcolumn
\EN
\Cite{Gen 4:1-7}
\switchcolumn*

\LA
\noindent Adam vero cognóvit uxórem suam Hevam, quæ concépit et péperit Cain, dicens: Possédi hóminem per Deum. Rursúmque péperit fratrem ejus Abel. Fuit autem Abel pastor óvium, et Cain agrícola. Factum est autem post multos dies ut offérret Cain de frúctibus terræ múnera Dómino. Abel quoque óbtulit de primogénitis gregis sui, et de adípibus eórum: et respéxit Dóminus ad Abel, et ad múnera ejus. Ad Cain vero, et ad múnera illíus non respéxit: iratúsque est Cain veheménter, et cóncidit vultus ejus. Dixítque Dóminus ad eum: Quare irátus es? et cur cóncidit fácies tua? nonne si bene égeris, recípies: sin autem male, statim in fóribus peccátum áderit? sed sub te erit appetítus ejus, et tu domináberis illíus.\par
\switchcolumn
\EN
\noindent And Adam knew Eve his wife: who conceived and brought forth Cain, saying: I have gotten a man through God. And again she brought forth his brother Abel. And Abel was a shepherd, and Cain a husbandman. And it came to pass after many days, that Cain offered, of the fruits of the earth, gifts to the Lord. Abel also offered of the firstlings of his flock, and of their fat: and the Lord had respect to Abel, and to his offerings. But to Cain and his offerings he had no respect: and Cain was exceedingly angry, and his countenance fell. And the Lord said to him: Why art thou angry? and why is thy countenance fallen? If thou do well, shalt thou not receive? but if ill, shall not sin forthwith be present at the door? but the lust thereof shall be under thee, and thou shalt have dominion over it.\par
\switchcolumn*

\LA
\begin{Resp}\Rsign In princípio creávit Deus cælum, et terram, et fecit in ea hóminem, \Star{} Ad imáginem et similitúdinem suam.\end{Resp}
\begin{Resp}\Vsign Formávit ígitur Deus hóminem de limo terræ, et inspirávit in fáciem ejus spiráculum vitæ.\end{Resp}
\begin{Resp}\Rsign Ad imáginem et similitúdinem suam.\end{Resp}
\switchcolumn
\EN
\begin{Resp}\Rsign In the beginning God created the heavens and the earth, wherein He made man also \Star{} After His own image and likeness.\end{Resp}
\begin{Resp}\Vsign So God formed man of the dust of the ground, and breathed into his face the breath of life.\end{Resp}
\begin{Resp}\Rsign After His own image and likeness.\end{Resp}
\switchcolumn*

\LA
\LessonHead{Lectio II}
\switchcolumn
\EN
\LessonHead{Lesson II}
\switchcolumn*

\LA
\Cite{Gen 4:8-12}
\switchcolumn
\EN
\Cite{Gen 4:8-12}
\switchcolumn*

\LA
\noindent Dixítque Cain ad Abel fratrem suum: Egrediámur foras. Cumque essent in agro, consurréxit Cain advérsus fratrem suum Abel, et interfécit eum. Et ait Dóminus ad Cain: Ubi est Abel frater tuus? Qui respóndit: Néscio: num custos fratris mei sum ego? Dixítque ad eum: Quid fecísti? vox sánguinis fratris tui clamat ad me de terra. Nunc ígitur maledíctus eris super terram, quæ apéruit os suum, et suscépit sánguinem fratris tui de manu tua. Cum operátus fúeris eam, non dabit tibi fructus suos: vagus et prófugus eris super terram.\par
\switchcolumn
\EN
\noindent And Cain said to Abel his brother: Let us go forth abroad. And when they were in the field, Cain rose up against his brother Abel, and slew him. And the Lord said to Cain: Where is thy brother Abel? And he answered, I know not: am I my brother's keeper? And he said to him: What hast thou done? the voice of thy brother's blood crieth to me from the earth. Now, therefore, cursed shalt thou be upon the earth, which hath opened her mouth and received the blood of thy brother at thy hand, When thou shalt till it, it shall not yield to thee its fruit: a fugitive and vagabond shalt thou be upon the earth.\par
\switchcolumn*

\LA
\begin{Resp}\Rsign In princípio creávit Deus cælum et terram, et Spíritus Dei ferebátur super aquas: \Star{} Et vidit Deus cuncta quæ fécerat, et erant valde bona.\end{Resp}
\begin{Resp}\Vsign Igitur perfécti sunt cæli et terra, et omnis ornátus eórum.\end{Resp}
\begin{Resp}\Rsign Et vidit Deus cuncta quæ fécerat, et erant valde bona.\end{Resp}
\switchcolumn
\EN
\begin{Resp}\Rsign In the beginning God created the heavens and the earth, and the Spirit of God moved upon the face of the waters. \Star{} And God saw everything that He had made, and it was very good.\end{Resp}
\begin{Resp}\Vsign Thus the heavens and the earth were finished, and all the hosts of them.\end{Resp}
\begin{Resp}\Rsign And God saw everything that He had made, and it was very good.\end{Resp}
\switchcolumn*

\LA
\LessonHead{Lectio III}
\switchcolumn
\EN
\LessonHead{Lesson III}
\switchcolumn*

\LA
\Cite{Gen 4:13-16}
\switchcolumn
\EN
\Cite{Gen 4:13-16}
\switchcolumn*

\LA
\noindent Dixítque Cain ad Dóminum: Major est iníquitas mea, quam ut véniam mérear. Ecce éicis me hódie a fácie terræ, et a fácie tua abscóndar, et ero vagus et prófugus in terra: omnis ígitur qui invénerit me, occídet me. Dixítque ei Dóminus: Nequáquam ita fiet: sed omnis qui occíderit Cain, séptuplum puniétur. Posuítque Dóminus Cain signum, ut non interfíceret eum omnis qui invenísset eum. Egressúsque Cain a fácie Dómini, habitávit prófugus in terra ad orientálem plagam Eden.\par
\switchcolumn
\EN
\noindent And Cain said to the Lord: My iniquity is greater than that I may deserve pardon. Behold thou dost cast me out this day from the face of the earth, and I shall be hidden from thy face, and I shall be a vagabond and a fugitive on the earth: everyone, therefore, that findeth me, shall kill me. And the Lord said to him: No, it shall not be so: but whosoever shall kill Cain, shall be punished sevenfold. And the Lord set a mark upon Cain, that whosoever found him should not kill him. And Cain went out from the face of the Lord, and dwelt as a fugitive on the earth, at the east side of Eden.\par
\switchcolumn*

\LA
\begin{Resp}\Rsign Formávit Dóminus hóminem de limo terræ, \Star{} Et inspirávit in fáciem ejus spiráculum vitæ, et factus est homo in ánimam vivéntem.\end{Resp}
\begin{Resp}\Vsign In princípio fecit Deus cælum et terram, et plasmávit in ea hóminem.\end{Resp}
\begin{Resp}\Rsign Et inspirávit in fáciem ejus spiráculum vitæ, et factus est homo in ánimam vivéntem.\end{Resp}
\begin{Resp}\Vsign Glória Patri, et Fílio, \Star{} et Spirítui Sancto.\end{Resp}
\begin{Resp}\Rsign Et inspirávit in fáciem ejus spiráculum vitæ, et factus est homo in ánimam vivéntem.\end{Resp}
\switchcolumn
\EN
\begin{Resp}\Rsign The Lord formed man of the dust of the ground \Star{} And breathed into his face the breath of life, and man became a living soul.\end{Resp}
\begin{Resp}\Vsign In the beginning God created the heavens and the earth, wherein He made man also.\end{Resp}
\begin{Resp}\Rsign And breathed into his face the breath of life, and man became a living soul.\end{Resp}
\begin{Resp}\Vsign Glory be to the Father, and to the Son, \Star{} and to the Holy Ghost.\end{Resp}
\begin{Resp}\Rsign And breathed into his face the breath of life, and man became a living soul.\end{Resp}
\switchcolumn*

\end{paracol}

\Office{Feria VI infra Hebdomadam Septuagesimæ}{Friday after Septuagesima Sunday}{Feria}
\addcontentsline{toc}{subsection}{Feria VI infra Hebdomadam Septuagesimæ\enspace\textit{Friday after Septuagesima Sunday}}
\begin{paracol}{2}
\LA
\LessonHead{Lectio I}
\switchcolumn
\EN
\LessonHead{Lesson I}
\switchcolumn*

\LA
\LessonTitle{De libro Génesis}
\switchcolumn
\EN
\LessonTitle{Lesson from the book of Genesis}
\switchcolumn*

\LA
\Cite{Gen 4:17-22}
\switchcolumn
\EN
\Cite{Gen 4:17-22}
\switchcolumn*

\LA
\noindent Cognóvit autem Cain uxórem suam, quæ concépit, et péperit Henoch: et ædificávit civitátem, vocavítque nomen ejus ex nómine fílii sui, Henoch. Porro Henoch génuit Irad, et Irad génuit Mavíaël, et Mavíaël génuit Mathúsaël, et Mathúsaël génuit Lamech. Qui accépit duas uxóres, nomen uni Ada, et nomen álteri Sella. Genuítque Ada Jabel, qui fuit pater habitántium in tentóriis, atque pastórum. Et nomen fratris ejus Jubal: ipse fuit pater canéntium cíthara et órgano. Sella quoque génuit Tubálcain, qui fuit malleátor et faber in cuncta ópera æris et ferri. Soror vero Tubálcain, Noéma.\par
\switchcolumn
\EN
\noindent And Cain knew his wife, and she conceived, and brought forth Henoch: and he built a city, and called the name thereof by the name of his son Henoch. And Henoch begot Irad, and Irad begot Maviael, and Maviael begot Mathusael, and Mathusael begot Lamech: Who took two wives: the name of the one was Ada, and the name of the other was Sella. And Ada brought forth Jabel: who was the father of such as dwell in tents, and of herdsmen. And his brother's name was Jubal; he was the father of them that play upon the harp and the organs. Sella also brought forth Tubalcain, who was a hammerer and artificer in every work of brass and iron. And the sister of Tubalcain was Noema.\par
\switchcolumn*

\LA
\begin{Resp}\Rsign Tulit Dóminus hóminem, et pósuit eum in paradíso voluptátis: \Star{} Ut operarétur et custodíret illum.\end{Resp}
\begin{Resp}\Vsign Plantáverat autem Dóminus Deus paradísum voluptátis a princípio, in quo pósuit hóminem quem formáverat.\end{Resp}
\begin{Resp}\Rsign Ut operarétur et custodíret illum.\end{Resp}
\switchcolumn
\EN
\begin{Resp}\Rsign God took the man and put him into the garden of Eden \Star{} To dress it and to keep it.\end{Resp}
\begin{Resp}\Vsign And the Lord God had planted a garden aforetime in Eden, and there He put the man whom He had formed.\end{Resp}
\begin{Resp}\Rsign To dress it and to keep it.\end{Resp}
\switchcolumn*

\LA
\LessonHead{Lectio II}
\switchcolumn
\EN
\LessonHead{Lesson II}
\switchcolumn*

\LA
\Cite{Gen 4:23-26}
\switchcolumn
\EN
\Cite{Gen 4:23-26}
\switchcolumn*

\LA
\noindent Dixítque Lamech uxóribus suis Adæ et Sellæ: Audíte vocem meam, uxóres Lamech; auscultáte sermónem meum: quóniam occídi virum in vulnus meum, et adolescéntulum in livórem meum. Séptuplum últio dábitur de Cain: de Lamech vero septuágies sépties. Cognóvit quoque adhuc Adam uxórem suam: et péperit fílium, vocavítque nomen ejus Seth, dicens: Pósuit mihi Deus semen áliud pro Abel, quem occídit Cain. Sed et Seth natus est fílius, quem vocávit Enos: iste cœpit invocáre nomen Dómini.\par
\switchcolumn
\EN
\noindent And Lamech said to his wives Ada and Sella: Hear my voice, ye wives of Lamech, hearken to my speech: for I have slain a man to the wounding of myself, and a stripling to my own bruising. Sevenfold vengeance shall be taken for Cain: but for Lamech seventy times sevenfold. Adam also knew his wife again: and she brought forth a son, and called his name Seth, saying: God hath given me another seed, for Abel whom Cain slew. But to Seth also was born a son, whom he called Enos; this man began to call upon the name of the Lord.\par
\switchcolumn*

\LA
\begin{Resp}\Rsign Dixit Dóminus Deus: Non est bonum hóminem esse solum: \Star{} Faciámus ei adjutórium símile sibi.\end{Resp}
\begin{Resp}\Vsign Adæ vero non inveniebátur adjútor símilis sibi: dixit vero Deus.\end{Resp}
\begin{Resp}\Rsign Faciámus ei adjutórium símile sibi.\end{Resp}
\switchcolumn
\EN
\begin{Resp}\Rsign The Lord God said: It is not good that the man should be alone. \Star{} Let Us make an help meet for him.\end{Resp}
\begin{Resp}\Vsign But for Adam there was not found an help meet for him; and God said:\end{Resp}
\begin{Resp}\Rsign Let Us make an help meet for him.\end{Resp}
\switchcolumn*

\LA
\LessonHead{Lectio III}
\switchcolumn
\EN
\LessonHead{Lesson III}
\switchcolumn*

\LA
\Cite{Gen 5:1-5}
\switchcolumn
\EN
\Cite{Gen 5:1-5}
\switchcolumn*

\LA
\noindent Hic est liber generatiónis Adam. In die qua creávit Deus hóminem, ad similitúdinem Dei fecit illum. Másculum et féminam creávit eos, et benedíxit illis: et vocávit nomen eórum Adam, in die quo creáti sunt. Vixit autem Adam centum trigínta annis: et génuit ad imáginem et similitúdinem suam, vocavítque nomen ejus Seth. Et facti sunt dies Adam, postquam génuit Seth, octingénti anni: genuítque fílios et fílias. Et factum est omne tempus quod vixit Adam, anni nongénti trigínta, et mórtuus est.\par
\switchcolumn
\EN
\noindent This is the book of the generation of Adam. In the day that God created man, he made him to the likeness of God. He created them male and female; and blessed them: and called their name Adam, in the day when they were created. And Adam lived a hundred and thirty years, and begot a son to his own image and likeness, and called his name Seth. And the days of Adam, after he begot Seth, were eight hundred years: and he begot sons and daughters. And all the time that Adam lived came to nine hundred and thirty years, and he died.\par
\switchcolumn*

\LA
\begin{Resp}\Rsign Immísit Dóminus sopórem in Adam, et tulit unam de costis ejus: † Et ædificávit costam, quam túlerat Dóminus de Adam, in mulíerem, et addúxit eam ad Adam, ut vidéret quid vocáret eam: \Star{} Et vocávit nomen ejus Virágo, quia de viro sumpta est.\end{Resp}
\begin{Resp}\Vsign Cumque obdormísset, tulit unam de costis ejus, et replévit carnem pro ea.\end{Resp}
\begin{Resp}\Rsign Et ædificávit costam, quam túlerat Dóminus de Adam, in mulíerem, et addúxit eam ad Adam, ut vidéret quid vocáret eam.\end{Resp}
\begin{Resp}\Vsign Glória Patri, et Fílio, \Star{} et Spirítui Sancto.\end{Resp}
\begin{Resp}\Rsign Et vocávit nomen ejus Virágo, quia de viro sumpta est.\end{Resp}
\switchcolumn
\EN
\begin{Resp}\Rsign The Lord caused a deep sleep to fall upon Adam, and He took one of his ribs. † And the rib which the Lord had taken from Adam made He a woman, and brought her unto Adam, to see what he would call her. \Star{} And he called her name Woman, because she was taken out of Man.\end{Resp}
\begin{Resp}\Vsign And while he slept He took one of his ribs, and closed up the flesh instead thereof.\end{Resp}
\begin{Resp}\Rsign And the rib which the Lord had taken from Adam made He a woman, and brought her unto Adam, to see what he would call her.\end{Resp}
\begin{Resp}\Vsign Glory be to the Father, and to the Son, \Star{} and to the Holy Ghost.\end{Resp}
\begin{Resp}\Rsign And he called her name Woman, because she was taken out of Man.\end{Resp}
\switchcolumn*

\end{paracol}

\Office{Sabbato infra Hebdomadam Septuagesimæ}{Saturday after Septuagesima Sunday}{Feria}
\addcontentsline{toc}{subsection}{Sabbato infra Hebdomadam Septuagesimæ\enspace\textit{Saturday after Septuagesima Sunday}}
\begin{paracol}{2}
\LA
\LessonHead{Lectio I}
\switchcolumn
\EN
\LessonHead{Lesson I}
\switchcolumn*

\LA
\LessonTitle{De libro Génesis}
\switchcolumn
\EN
\LessonTitle{Lesson from the book of Genesis}
\switchcolumn*

\LA
\Cite{Gen 5:15-21}
\switchcolumn
\EN
\Cite{Gen 5:15-21}
\switchcolumn*

\LA
\noindent Vixit autem Maláleel sexagínta quinque annis, et génuit Jared. Et vixit Maláleel, postquam génuit Jared, octingéntis trigínta annis, et génuit fílios et fílias. Et facti sunt omnes dies Maláleel octingénti nonagínta quinque anni, et mórtuus est. Vixítque Jared centum sexagínta duóbus annis, et génuit Henoch. Et vixit Jared, postquam génuit Henoch, octingéntis annis, et génuit fílios et fílias. Et facti sunt omnes dies Jared nongénti sexagínta duo anni, et mórtuus est. Porro Henoch vixit sexagínta quinque annis, et génuit Mathúsalam.\par
\switchcolumn
\EN
\noindent And Malaleel lived sixty-five years, and begot Jared. And Malaleel lived after he begot Jared, eight hundred and thirty years, and begot sons and daughters. And all the days of Malaleel were eight hundred and ninety-five years, and he died. And Jared lived a hundred and sixty-two years, and begot Henoch. And Jared lived after he begot Henoch, eight hundred years, and begot sons and daughters. And all the days of Jared were nine hundred and sixty-two years, and he died. And Henoch lived sixty-five years, and begot Mathusala.\par
\switchcolumn*

\LA
\begin{Resp}\Rsign Plantáverat autem Dóminus Deus paradísum voluptátis a princípio: \Star{} In quo pósuit hóminem, quem formáverat.\end{Resp}
\begin{Resp}\Vsign Produxítque Dóminus Deus de humo omne lignum pulchrum visu, et ad vescéndum suáve; lignum étiam vitæ in médio paradísi.\end{Resp}
\begin{Resp}\Rsign In quo pósuit hóminem, quem formáverat.\end{Resp}
\switchcolumn
\EN
\begin{Resp}\Rsign And the Lord God had planted a garden aforetime in Eden, \Star{} And there He put the man whom He had formed.\end{Resp}
\begin{Resp}\Vsign And out of the ground made the Lord God to grow every tree that is pleasant to the sight, and good for food, the tree of life also in the midst of the garden.\end{Resp}
\begin{Resp}\Rsign And there He put the man whom He had formed.\end{Resp}
\switchcolumn*

\LA
\LessonHead{Lectio II}
\switchcolumn
\EN
\LessonHead{Lesson II}
\switchcolumn*

\LA
\Cite{Gen 5:22-27}
\switchcolumn
\EN
\Cite{Gen 5:22-27}
\switchcolumn*

\LA
\noindent Et ambulávit Henoch cum Deo: et vixit, postquam génuit Mathúsalam, trecéntis annis, et génuit fílios et fílias. Et facti sunt omnes dies Henoch trecénti sexagínta quinque anni. Ambulavítque cum Deo, et non appáruit: quia tulit eum Deus. Vixit quoque Mathúsala centum octogínta septem annis, et génuit Lamech. Et vixit Mathúsala, postquam génuit Lamech, septingéntis octogínta duóbus annis, et génuit fílios et fílias. Et facti sunt omnes dies Mathúsala nongénti sexagínta novem anni, et mórtuus est.\par
\switchcolumn
\EN
\noindent And Henoch walked with God: and lived after he begot Mathusala, three hundred years, and begot sons and daughters. And all the days of Henoch were three hundred and sixty-five years. And he walked with God, and was seen no more: because God took him. And Mathusala lived a hundred and eighty-seven years, and begot Lamech. And Mathusala lived after he begot Lamech, seven hundred and eighty-two years, and begot sons and daughters. And all the days of Mathusala were nine hundred and sixty-nine years, and he died.\par
\switchcolumn*

\LA
\begin{Resp}\Rsign Ecce Adam quasi unus ex nobis factus est sciens bonum et malum: \Star{} Vidéte, ne forte sumat de ligno vitæ, et vivat in ætérnum.\end{Resp}
\begin{Resp}\Vsign Fecit quoque Dóminus Deus Adæ túnicam pellíceam, et índuit eum, et dixit.\end{Resp}
\begin{Resp}\Rsign Vidéte, ne forte sumat de ligno vitæ, et vivat in ætérnum.\end{Resp}
\switchcolumn
\EN
\begin{Resp}\Rsign Behold, Adam is become as One of Us, to know good and evil. \Star{} See lest he take of the tree of life and live for ever.\end{Resp}
\begin{Resp}\Vsign Unto Adam also did the Lord God make a coat of skins, and clothed him, and said:\end{Resp}
\begin{Resp}\Rsign See lest he take of the tree of life and live for ever.\end{Resp}
\switchcolumn*

\LA
\LessonHead{Lectio III}
\switchcolumn
\EN
\LessonHead{Lesson III}
\switchcolumn*

\LA
\Cite{Gen 5:28-31}
\switchcolumn
\EN
\Cite{Gen 5:28-31}
\switchcolumn*

\LA
\noindent Vixit autem Lamech centum octogínta duóbus annis, et génuit fílium: vocavítque nomen ejus Noë, dicens: Iste consolábitur nos ab opéribus et labóribus mánuum nostrárum in terra, cui maledíxit Dóminus. Vixítque Lamech, postquam génuit Noë, quingéntis nonagínta quinque annis, et génuit fílios et fílias. Et facti sunt omnes dies Lamech septingénti septuagínta septem anni, et mórtuus est.\par
\switchcolumn
\EN
\noindent And Lamech lived a hundred and eighty-two years, and begot a son. And he called his name Noe, saying: This same shall comfort us from the works and labours of our hands on the earth which the Lord hath cursed. And Lamech lived after he begot Noe, five hundred and ninety-five years, and he begot sons and daughters. And all the days of Lamech came to seven hundred and seventy-seven years, and he died.\par
\switchcolumn*

\LA
\begin{Resp}\Rsign Ubi est Abel frater tuus? dixit Dóminus ad Cain. Néscio, Dómine, numquid custos fratris mei sum ego? Et dixit ad eum: Quid fecísti? \Star{} Ecce vox sánguinis fratris tui Abel clamat ad me de terra.\end{Resp}
\begin{Resp}\Vsign Maledíctus eris super terram, quæ apéruit os suum, et suscépit sánguinem fratris tui de manu tua.\end{Resp}
\begin{Resp}\Rsign Ecce vox sánguinis fratris tui Abel clamat ad me de terra.\end{Resp}
\begin{Resp}\Vsign Glória Patri, et Fílio, \Star{} et Spirítui Sancto.\end{Resp}
\begin{Resp}\Rsign Ecce vox sánguinis fratris tui Abel clamat ad me de terra.\end{Resp}
\switchcolumn
\EN
\begin{Resp}\Rsign The Lord said unto Cain: Where is Abel thy brother? Lord, I know not am I my brother's keeper? And He said unto him What hast thou done? \Star{} Behold, the voice of thy brother Abel's blood crieth unto Me from the ground.\end{Resp}
\begin{Resp}\Vsign Cursed shalt thou be upon the earth, which hath opened her mouth to receive thy brother's blood from thy hand.\end{Resp}
\begin{Resp}\Rsign Behold, the voice of thy brother Abel's blood crieth unto Me from the ground.\end{Resp}
\begin{Resp}\Vsign Glory be to the Father, and to the Son, \Star{} and to the Holy Ghost.\end{Resp}
\begin{Resp}\Rsign Behold, the voice of thy brother Abel's blood crieth unto Me from the ground.\end{Resp}
\switchcolumn*

\end{paracol}

\Office{Dominica in Sexagesima}{Sexagesima Sunday}{Semiduplex II classis}
\addcontentsline{toc}{subsection}{Dominica in Sexagesima\enspace\textit{Sexagesima Sunday}}
\begin{paracol}{2}
\LA
\LessonHead{Lectio I}
\switchcolumn
\EN
\LessonHead{Lesson I}
\switchcolumn*

\LA
\LessonTitle{De libro Génesis}
\switchcolumn
\EN
\LessonTitle{Lesson from the book of Genesis}
\switchcolumn*

\LA
\Cite{Gen 5:31; 6:1-4}
\switchcolumn
\EN
\Cite{Gen 5:31; 6:1-4}
\switchcolumn*

\LA
\noindent Noë vero cum quingentórum esset annórum, génuit Sem, Cham et Japheth. Cumque cœpíssent hómines multiplicári super terram, et fílias procreássent, vidéntes fílii Dei fílias hóminum quod essent pulchræ, accepérunt sibi uxóres ex ómnibus, quas elégerant. Dixítque Deus: Non permanébit spíritus meus in hómine in ætérnum, quia caro est: erúntque dies illíus centum vigínti annórum. Gigántes autem erant super terram in diébus illis: postquam enim ingréssi sunt fílii Dei ad fílias hóminum, illǽque genuérunt, isti sunt poténtes a sǽculo viri famósi.\par
\switchcolumn
\EN
\noindent Noe, when he was five hundred years old, begot Sem, Cham, and Japheth. And after that men began to be multiplied upon the earth, and daughters were born to them, The sons of God seeing the daughters of men, that they were fair, took to themselves wives of all which they chose. And God said: My spirit shall not remain in man for ever, because he is flesh, and his days shall be a hundred and twenty years. Now giants were upon the earth in those days. For after the sons of God went in to the daughters of men, and they brought forth children, these are the mighty men of old, men of renown.\par
\switchcolumn*

\LA
\begin{Resp}\Rsign Dixit Dóminus ad Noë: Finis univérsæ carnis venit coram me: repléta est terra iniquitáte eórum, \Star{} Et ego dispérdam eos cum terra.\end{Resp}
\begin{Resp}\Vsign Fac tibi arcam de lignis lævigátis, mansiúnculas in ea fácies.\end{Resp}
\begin{Resp}\Rsign Et ego dispérdam eos cum terra.\end{Resp}
\switchcolumn
\EN
\begin{Resp}\Rsign The Lord said unto Noah: The end of all flesh is come before Me; for the earth is filled with violence through them. \Star{} And I will destroy them with the earth.\end{Resp}
\begin{Resp}\Vsign Make thee an ark of planed timber, rooms shalt thou make in it.\end{Resp}
\begin{Resp}\Rsign And I will destroy them with the earth.\end{Resp}
\switchcolumn*

\LA
\LessonHead{Lectio II}
\switchcolumn
\EN
\LessonHead{Lesson II}
\switchcolumn*

\LA
\Cite{Gen 6:5-8}
\switchcolumn
\EN
\Cite{Gen 6:5-8}
\switchcolumn*

\LA
\noindent Videns autem Deus quod multa malítia hóminum esset in terra, et cuncta cogitátio cordis inténta esset ad malum omni témpore, pœnítuit eum quod hóminem fecísset in terra. Et tactus dolóre cordis intrínsecus, Delébo, inquit, hóminem, quem creávi, a fácie terræ, ab hómine usque ad animántia, a réptili usque ad vólucres cæli: pœ́nitet enim me fecísse eos. Noë vero invénit grátiam coram Dómino.\par
\switchcolumn
\EN
\noindent And God seeing that the wickedness of men was great on the earth, and that all the thought of their heart was bent upon evil at all times, It repented him that he had made man on the earth. And being touched inwardly with sorrow of heart, He said: I will destroy man, whom I have created, from the face of the earth, from man even to beasts, from the creeping thing even to the fowls of the air, for it repenteth me that I have made them. But Noe found grace before the Lord.\par
\switchcolumn*

\LA
\begin{Resp}\Rsign Noë, vir justus atque perféctus, cum Deo ambulávit: \Star{} Et fecit ómnia quæcúmque præcépit ei Deus.\end{Resp}
\begin{Resp}\Vsign Fecit sibi arcam, ut salvarétur univérsum semen.\end{Resp}
\begin{Resp}\Rsign Et fecit ómnia quæcúmque præcépit ei Deus.\end{Resp}
\switchcolumn
\EN
\begin{Resp}\Rsign Noah was a just man and perfect; he walked with God. \Star{} According to all that God commanded him, so did he.\end{Resp}
\begin{Resp}\Vsign He made him an ark, that a seed of every sort might be saved alive.\end{Resp}
\begin{Resp}\Rsign According to all that God commanded him, so did he.\end{Resp}
\switchcolumn*

\LA
\LessonHead{Lectio III}
\switchcolumn
\EN
\LessonHead{Lesson III}
\switchcolumn*

\LA
\Cite{Gen 6:9-15}
\switchcolumn
\EN
\Cite{Gen 6:9-15}
\switchcolumn*

\LA
\noindent Hæ sunt generatiónes Noë: Noë vir justus atque perféctus fuit in generatiónibus suis; cum Deo ambulávit. Et génuit tres fílios, Sem, Cham et Japheth. Corrúpta est autem terra coram Deo, et repléta est iniquitáte. Cumque vidísset Deus terram esse corrúptam (omnis quippe caro corrúperat viam suam super terram), dixit ad Noë: Finis univérsæ carnis venit coram me: repléta est terra iniquitáte a fácie eórum, et ego dispérdam eos cum terra. Fac tibi arcam de lignis lævigátis; mansiúnculas in arca fácies, et bitúmine línies intrínsecus et extrínsecus. Et sic fácies eam: trecentórum cubitórum erit longitúdo arcæ, quinquagínta cubitórum latitúdo, et trigínta cubitórum altitúdo illíus.\par
\switchcolumn
\EN
\noindent These are the generations of Noe: Noe was a just and perfect man in his generations, he walked with God. And he begot three sons, Sem, Cham, and Japheth. And the earth was corrupted before God, and was filled with iniquity. And when God had seen that the earth was corrupted (for all flesh had corrupted its way upon the earth,) He said to Noe: The end of all flesh is come before me, the earth is filled with iniquity through them, and I will destroy them with the earth. Make thee an ark of timber planks: thou shalt make little rooms in the ark, and thou shalt pitch it within and without. And thus shalt thou make it: The length of the ark shall be three hundred cubits: the breadth of it fifty cubits, and the height of it thirty cubits.\par
\switchcolumn*

\LA
\begin{Resp}\Rsign Quadragínta dies et noctes apérti sunt cæli, et ex omni carne habénte spíritum vitæ ingréssa sunt in arcam: \Star{} Et clausit a foris óstium Dóminus.\end{Resp}
\begin{Resp}\Vsign In artículo diéi illíus ingréssus est Noë in arcam et fílii ejus, et uxor illíus et uxóres filiórum ejus.\end{Resp}
\begin{Resp}\Rsign Et clausit a foris óstium Dóminus.\end{Resp}
\begin{Resp}\Vsign Glória Patri, et Fílio, \Star{} et Spirítui Sancto.\end{Resp}
\begin{Resp}\Rsign Et clausit a foris óstium Dóminus.\end{Resp}
\switchcolumn
\EN
\begin{Resp}\Rsign Forty days and forty nights were the heavens opened; and there went into the ark two and two of all flesh wherein is the breath of life. \Star{} And the Lord shut them in.\end{Resp}
\begin{Resp}\Vsign In the self-same day entered Noah into the ark, and his sons, and his wife, and the wives of his sons.\end{Resp}
\begin{Resp}\Rsign And the Lord shut them in.\end{Resp}
\begin{Resp}\Vsign Glory be to the Father, and to the Son, \Star{} and to the Holy Ghost.\end{Resp}
\begin{Resp}\Rsign And the Lord shut them in.\end{Resp}
\switchcolumn*

\LA
\LessonHead{Lectio IV}
\switchcolumn
\EN
\LessonHead{Lesson IV}
\switchcolumn*

\LA
\LessonTitle{Ex libro sancti Ambrósii Epíscopi de Noë et arca}
\switchcolumn
\EN
\LessonTitle{From the Book upon Noah's Ark by St. Ambrose, Bishop (of Milan.)}
\switchcolumn*

\LA
\Source{Cap. 4. circa med.}
\switchcolumn
\EN
\Source{Chap, iv.}
\switchcolumn*

\LA
\noindent Habes, quia irátus Dóminus est: quóniam quamvis cogitáret, hoc est sciret, quia homo pósitus in terræ regióne, carnem portans, sine peccáto esse non possit, (terra enim velut quidam tentatiónum locus est, caróque corruptélæ illécebra) tamen cum habérent mentem ratiónis capácem, virtutémque ánimæ infusam córpori, sine consideratióne áliqua in lapsum ruérunt, ex quo revocáre se nollent. Neque enim Deus cógitat sicut hómines, ut áliqua ei nova succédat senténtia, neque iráscitur quasi mutábilis: sed ídeo hæc legúntur, ut exprimátur peccatórum nostrórum acérbitas, quæ divínam merúerit offénsam: tamquam eoúsque incréverit culpa, ut étiam Deus, qui naturáliter non movétur aut ira, aut ódio, aut passióne ulla, provocátus videátur ad iracúndiam.\par
\switchcolumn
\EN
\noindent We read that the Lord was angry. It is in the thoughts, that is to say, in the knowledge of God, that man being put on earth and weighted with the body cannot be without sin, for earth is the home of temptations, and the flesh is a bait for corruption. Yet man had a reasonable soul, and his soul had power to control his body; and, being so made, he made no struggle to keep himself from falling into that from whence he would not return. God's thoughts are not as man's thoughts; in Him there is no such thing as change of mind, no such thing as to be angry and then cool down again. These things are written that we may know the bitterness of our sins, whereby we have earned the Divine wrath. To such a degree had iniquity grown that God, Who by His nature cannot be moved by anger, or hatred, or any passion whatsoever, is represented as provoked to anger.\par
\switchcolumn*

\LA
\begin{Resp}\Rsign Ædificávit Noë altáre Dómino, ófferens super illud holocáustum: odoratúsque est Dóminus odórem suavitátis, et benedíxit ei, dicens: \Star{} Créscite, et multiplicámini, et repléte terram.\end{Resp}
\begin{Resp}\Vsign Ecce ego státuam pactum meum vobíscum, et cum sémine vestro post vos.\end{Resp}
\begin{Resp}\Rsign Créscite, et multiplicámini, et repléte terram.\end{Resp}
\switchcolumn
\EN
\begin{Resp}\Rsign Noah builded an Altar unto the Lord, and offered burnt offerings on the Altar; and the Lord smelled a sweet savour, and blessed Noah, and said: \Star{} Be fruitful, and multiply, and replenish the earth.\end{Resp}
\begin{Resp}\Vsign Behold, I establish My covenant with you, and with your seed after you.\end{Resp}
\begin{Resp}\Rsign Be fruitful, and multiply, and replenish the earth.\end{Resp}
\switchcolumn*

\LA
\LessonHead{Lectio V}
\switchcolumn
\EN
\LessonHead{Lesson V}
\switchcolumn*

\LA
\noindent Minitátus est prætérea, quod deléret hóminem. Ab hómine, inquit, usque ad pecus, et a reptílibus usque ad volatília delébo. Quid læserant irrationabília? Sed quia propter hóminem illa facta erant, eo útique deléto, propter quem facta sunt, cónsequens erat, ut étiam illa deleréntur, quia non erat qui his uterétur. Sensu autem altióre illud manifestátur: quia homo mens est, quæ est ratiónis capax. Homo enim definítur ánimal vivum, mortále, rationábile. Principáli ígitur exstíncto, étiam sensus omnis exstínguitur: eo quod nihil réliqui ad salútem supérsit, cum salútis fundaméntum virtus defécerit.\par
\switchcolumn
\EN
\noindent And God threatened that He would destroy man. He said: I will destroy man, whom I have created, from the face of the earth; both man and beast, and the creeping thing, and the fowls of the air. What harm had the animals done? For man's use had they been created, and, when man was wiped away, they were of use no longer. And there is an higher reason. Man is a living soul, capable of reason, who may be described as a living animal, subject to death, and endowed with reason. When then the highest animal is gone, why should the lower branches remain? Why should anything be saved alive, when righteousness, the basis of salvation, is to be no more?\par
\switchcolumn*

\LA
\begin{Resp}\Rsign Ponam arcum meum in núbibus cæli, dixit Dóminus ad Noë: \Star{} Et recordábor fœ́deris mei, quod pépigi tecum.\end{Resp}
\begin{Resp}\Vsign Cumque obdúxero núbibus cælum, apparébit arcus meus in núbibus.\end{Resp}
\begin{Resp}\Rsign Et recordábor fœ́deris mei, quod pépigi tecum.\end{Resp}
\switchcolumn
\EN
\begin{Resp}\Rsign The Lord said unto Noah: I do set My bow in the clouds of heaven, \Star{} And I will remember My covenant which is between Me and you.\end{Resp}
\begin{Resp}\Vsign And it shall come to pass, when I bring a cloud over the heaven, that My bow shall be seen in the cloud.\end{Resp}
\begin{Resp}\Rsign And I will remember My covenant which is between Me and you.\end{Resp}
\switchcolumn*

\LA
\LessonHead{Lectio VI}
\switchcolumn
\EN
\LessonHead{Lesson VI}
\switchcolumn*

\LA
\noindent Ad condemnatiónem autem ceterórum, et ad expressiónem pietátis divínæ, dícitur Noë apud Deum grátiam invenísse. Simul osténditur, quod hóminem justum non obúmbret aliórum offénsio, quando ipse ad totíus géneris reservátur seminárium. Qui non generatiónis nobilitáte, sed justítiæ et perfectiónis mérito laudátur. Probáti enim viri genus, virtútis prosápia est: quia sicut hóminum genus hómines, ita animárum genus virtútes sunt. Etenim famíliæ hóminum splendóre géneris nobilitántur, animárum autem clarificátur grátia splendóre virtútis.\par
\switchcolumn
\EN
\noindent But more effectually to condemn the rest of men, and to manifest the goodness of God, it is written that Noah found grace in the eyes of the Lord. Here we learn also that the sin of his neighbour casteth no shadow on the righteous, when he is kept as a stock from whence the whole race are to spring. He is praised, not because he was of a noble race, but because he was a just man and perfect. The stock of a just man yieldeth men of just souls; for virtues, like blood, are hereditary. Among men are some families illustrious for honourable pedigrees, and so there are also races of souls whose comeliness is the lustre of virtues.\par
\switchcolumn*

\LA
\begin{Resp}\Rsign Per memetípsum jurávi, dicit Dóminus, non adíciam ultra aquas dilúvii super terram: pacti mei recordábor, \Star{} Ut non perdam aquis dilúvii omnem carnem.\end{Resp}
\begin{Resp}\Vsign Arcum meum ponam in núbibus, et erit signum fœ́deris inter me et inter terram.\end{Resp}
\begin{Resp}\Rsign Ut non perdam aquis dilúvii omnem carnem.\end{Resp}
\begin{Resp}\Vsign Glória Patri, et Fílio, \Star{} et Spirítui Sancto.\end{Resp}
\begin{Resp}\Rsign Ut non perdam aquis dilúvii omnem carnem.\end{Resp}
\switchcolumn
\EN
\begin{Resp}\Rsign By Myself have I sworn, saith the Lord. I will not again bring the waters of the flood upon the earth, I will remember My covenant. \Star{} And the waters shall become no more a flood to destroy all flesh.\end{Resp}
\begin{Resp}\Vsign I do set My bow in the clouds, and it shall be for a token of a covenant between Me and the earth.\end{Resp}
\begin{Resp}\Rsign And the waters shall no more become a flood to destroy all flesh.\end{Resp}
\begin{Resp}\Vsign Glory be to the Father, and to the Son, \Star{} and to the Holy Ghost.\end{Resp}
\begin{Resp}\Rsign And the waters shall no more become a flood to destroy all flesh.\end{Resp}
\switchcolumn*

\LA
\LessonHead{Lectio VII}
\switchcolumn
\EN
\LessonHead{Lesson VII}
\switchcolumn*

\LA
\LessonTitle{Léctio sancti Evangélii secúndum Lucam}
\switchcolumn
\EN
\LessonTitle{From the Holy Gospel according to Luke}
\switchcolumn*

\LA
\Cite{Luc 8:4-15}
\switchcolumn
\EN
\Cite{Luke 8:4-15}
\switchcolumn*

\LA
\noindent In illo témpore: Cum autem turba plúrima convenírent, et de civitátibus properárent ad Jesum, dixit per similitúdinem: Exiit qui séminat, semináre semen suum. Et réliqua.\par
\switchcolumn
\EN
\noindent At that time: When much people were gathered together, and were come to Jesus out of every city, he spake by a parable. A sower went out to sow his seed. And so on, and that which followeth.\par
\switchcolumn*

\LA
\LessonTitle{Homilía sancti Gregórii Papæ}
\switchcolumn
\EN
\LessonTitle{Homily by Pope St. Gregory (the Great.)}
\switchcolumn*

\LA
\Source{Homilia 15 in Evangelia}
\switchcolumn
\EN
\Source{15th on the Gospels}
\switchcolumn*

\LA
\noindent Léctio sancti Evangélii, quam modo, fratres caríssimi, audístis, expositióne non índiget, sed admonitióne. Quam enim per semetípsam Véritas expósuit, hanc discútere humána fragílitas non præsúmat. Sed est quod sollícite in hac ipsa expositióne Domínica pensáre debeámus: quia si nos vobis semen verbum, agrum mundum, vólucres dæmónia, spinas divítias significáre dicerémus ad credéndum nobis mens fórsitan vestra dubitáret. Unde et idem Dóminus per semetípsum dignátus est expónere quod dicébat, ut sciátis rerum significatiónes quǽrere in iis étiam, quæ per semetípsum nóluit explanáre.\par
\switchcolumn
\EN
\noindent Dearly beloved brethren, the passage from the Holy Gospel which ye have just heard, needeth not so much that I should explain it, as that I should seek to enforce its lesson. The Truth Himself hath explained it, and, after that, it beseemeth not man's frailty to fritter away His exposition by any further comment. But there is, in that very explanation by the Lord, somewhat, which it behoveth us well to weigh. If it were but we who bade you believe that by the seed is signified the word; by the field, the world; by the birds, the devils; and by the thorns, riches ye would perchance doubt of the truth of our explanation. Therefore the Lord Himself hath vouchsafed to give this explanation, and that, not for this parable only, but that ye may know in what manner to interpret others, whereof He hath not given the meaning.\par
\switchcolumn*

\LA
\begin{Resp}\Rsign Benedíxit Deus Noë, et fíliis ejus, et dixit ad eos: \Star{} Créscite, et multiplicámini, et repléte terram.\end{Resp}
\begin{Resp}\Vsign Ecce ego státuam pactum meum vobíscum, et cum sémine vestro post vos.\end{Resp}
\begin{Resp}\Rsign Créscite, et multiplicámini, et repléte terram.\end{Resp}
\switchcolumn
\EN
\begin{Resp}\Rsign God blessed Noah and his sons, and said unto them: \Star{} Be fruitful, and multiply, and replenish the earth.\end{Resp}
\begin{Resp}\Vsign Behold, I establish My covenant with you, and with your seed after you.\end{Resp}
\begin{Resp}\Rsign Be fruitful, and multiply, and replenish the earth.\end{Resp}
\switchcolumn*

\LA
\LessonHead{Lectio VIII}
\switchcolumn
\EN
\LessonHead{Lesson VIII}
\switchcolumn*

\LA
\noindent Exponéndo ergo quod dixit, figuráte se loqui innótuit: quátenus certos nos rédderet, cum vobis nostra fragílitas verbórum illíus figúras aperíret. Quis enim mihi umquam créderet, si spinas divítias interpretári voluíssem? máxime cum illæ pugnant, istæ deléctent. Et tamen spinæ sunt, quia cogitatiónum suárum punctiónibus mentem lácerant: et cum usque ad peccátum pértrahunt, quasi inflícto vúlnere cruéntant. Quas bene hoc in loco, álio Evangelísta testánte, nequáquam Dóminus divítias, sed falláces divítias appéllat.\par
\switchcolumn
\EN
\noindent Beginning His explanation, the Lord saith that He speaketh in parables. Hereby He doth certify us, when our weakness would unveil to you the hidden meaning of His words. If I spake of myself, who would believe me when I say that riches are thorns? Thorns prick, but riches lull to rest. And yet riches are indeed thorns, for the anxiety they bring is a ceaseless pricking to the minds of their owners, and, if they lead into sin, they are thorns which bloodily tear the soul. But we understand from another Evangelist (Matth. xiii. 22) that in this place the Lord speaketh, not of riches themselves, but of the deceitfulness of riches.\par
\switchcolumn*

\LA
\begin{Resp}\Rsign Ecce ego státuam pactum meum vobíscum, et cum sémine vestro post vos: \Star{} Neque erit deínceps dilúvium díssipans terram.\end{Resp}
\begin{Resp}\Vsign Arcum meum ponam in núbibus, et erit signum fœ́deris inter me et inter terram.\end{Resp}
\begin{Resp}\Rsign Neque erit deínceps dilúvium díssipans terram.\end{Resp}
\switchcolumn
\EN
\begin{Resp}\Rsign Behold, I establish My covenant with you, and with your seed after you. \Star{} Neither shall there any more be a flood to destroy the earth.\end{Resp}
\begin{Resp}\Vsign I do set My bow in the clouds, and it shall be for a token of a covenant between Me and the earth.\end{Resp}
\begin{Resp}\Rsign Neither shall there any more be a flood to destroy the earth.\end{Resp}
\switchcolumn*

\LA
\LessonHead{Lectio IX}
\switchcolumn
\EN
\LessonHead{Lesson IX}
\switchcolumn*

\LA
\noindent Falláces enim sunt, quæ nobíscum diu permanére non possunt: falláces sunt, quæ mentis nostræ inópiam non expéllunt. Solæ autem divítiæ veræ sunt, quæ nos dívites virtútibus fáciunt. Si ergo, fratres caríssimi, esse dívites cúpitis, veras divítias amáte. Si culmen veri honóris quæritis, ad cæléste regnum téndite. Si glóriam dignitátum dilígitis, in illa supérna Angelórum cúria adscríbi festináte. Verba Dómini, quæ aure percípitis, mente retinéte. Cibus enim mentis est sermo Dei: et quasi accéptus cibus stómacho languénte rejícitur, quando audítus sermo in ventre memóriæ non tenétur. Sed quisquis aliménta non rétinet, hujus profécto vita desperátur.\par
\switchcolumn
\EN
\noindent Those riches are deceitful riches, which can be ours only for a little while; those riches are deceitful riches, which cannot relieve the poverty of our souls. They are the only true riches, which make us rich in virtues. If then, dearly beloved brethren, ye seek to be rich, earnestly desire the true riches. If ye would be truly honourable, strive after the kingdom of heaven. If ye love the bravery of titles, hasten to have your names written down at Court above, where Angels are. Take to heart the Lord's words which your ear heareth. The food of the soul is the word of God when the stomach is sick it throweth up again the food which is put into it, and so is the soul sick when a man heareth and digesteth not in his memory the Word of God. And if any man cannot keep his food, that man's life is in desperate case.\par
\switchcolumn*

\LA
\begin{Resp}\Rsign Cum turba plúrima convenírent ad Jesum, et de civitátibus properárent ad eum, dixit per similitúdinem: \Star{} Exiit qui séminat, semináre semen suum.\end{Resp}
\begin{Resp}\Vsign Et dum séminat, áliud cécidit in terram bonam, et ortum fecit fructum céntuplum.\end{Resp}
\begin{Resp}\Rsign Exiit qui séminat, semináre semen suum.\end{Resp}
\begin{Resp}\Vsign Glória Patri, et Fílio, \Star{} et Spirítui Sancto.\end{Resp}
\begin{Resp}\Rsign Exiit qui séminat, semináre semen suum.\end{Resp}
\switchcolumn
\EN
\begin{Resp}\Rsign When much people were gathered together to Jesus, and were come to Him out of every city, He spake by a parable: \Star{} A sower went out to sow his seed.\end{Resp}
\begin{Resp}\Vsign And, as he sowed, some fell on good ground, and sprang up, and bare fruit an hundred-fold.\end{Resp}
\begin{Resp}\Rsign A sower went out to sow his seed.\end{Resp}
\begin{Resp}\Vsign Glory be to the Father, and to the Son, \Star{} and to the Holy Ghost.\end{Resp}
\begin{Resp}\Rsign A sower went out to sow his seed.\end{Resp}
\switchcolumn*

\end{paracol}

\Office{Feria II infra Hebdomadam Sexagesimæ}{Monday after Sexagesima Sunday}{Feria}
\addcontentsline{toc}{subsection}{Feria II infra Hebdomadam Sexagesimæ\enspace\textit{Monday after Sexagesima Sunday}}
\begin{paracol}{2}
\LA
\LessonHead{Lectio I}
\switchcolumn
\EN
\LessonHead{Lesson I}
\switchcolumn*

\LA
\LessonTitle{De libro Génesis}
\switchcolumn
\EN
\LessonTitle{Lesson from the book of Genesis}
\switchcolumn*

\LA
\Cite{Gen 7:1-4}
\switchcolumn
\EN
\Cite{Gen 7:1-4}
\switchcolumn*

\LA
\noindent Dixítque Dóminus ad eum: Ingrédere tu et omnis domus tua in arcam: te enim vidi justum coram me in generatióne hac. Ex ómnibus animántibus mundis tolle septéna et septéna, másculum et féminam: de animántibus vero immúndis duo et duo, másculum et féminam. Sed et de volatílibus cæli septéna et septéna, másculum et féminam: ut salvétur semen super fáciem univérsæ terræ. Adhuc enim, et post dies septem ego pluam super terram quadragínta diébus et quadragínta nóctibus: et delébo omnem substántiam, quam feci, de superfície terræ.\par
\switchcolumn
\EN
\noindent And the Lord said to him: Go in thou and all thy house into the ark: for thee I have seen just before me in this generation. Of all clean beasts take seven and seven, the male and female. But of the beasts that are unclean two and two, the male and female. Of the fowls also of the air seven and seven, the male and the female: that seed may be saved upon the face of the whole earth. For yet a while, and after seven days, I will rain upon the earth forty days and forty nights; and I will destroy every substance that I have made, from the face of the earth.\par
\switchcolumn*

\LA
\begin{Resp}\Rsign In artículo diéi illíus ingréssus est Noë in arcam et fílii ejus, \Star{} Uxor illíus et uxóres filiórum ejus.\end{Resp}
\begin{Resp}\Vsign Deléta sunt univérsa de terra, remánsit autem solus Noë, et qui cum eo erant in arca.\end{Resp}
\begin{Resp}\Rsign Uxor illíus et uxóres filiórum ejus.\end{Resp}
\switchcolumn
\EN
\begin{Resp}\Rsign In the self-same day entered Noah into the ark, and his sons, \Star{} And his wife and the wives of his sons.\end{Resp}
\begin{Resp}\Vsign Every living substance was destroyed from the earth, and Noah only remained alive, and they that were with him in the ark.\end{Resp}
\begin{Resp}\Rsign And his wife, and the wives of his sons.\end{Resp}
\switchcolumn*

\LA
\LessonHead{Lectio II}
\switchcolumn
\EN
\LessonHead{Lesson II}
\switchcolumn*

\LA
\Cite{Gen 7:5; 7:10-12}
\switchcolumn
\EN
\Cite{Gen 7:5; 7:10-12}
\switchcolumn*

\LA
\noindent Fecit ergo Noë ómnia quæ mandáverat ei Dóminus. Cumque transíssent septem dies, aquæ dilúvii inundavérunt super terram. Anno sexcentésimo vitæ Noë, mense secúndo, séptimo décimo die mensis, rupti sunt omnes fontes abýssi magnæ, et cataráctæ cæli apértæ sunt: et facta est plúvia super terram quadragínta diébus et quadragínta nóctibus.\par
\switchcolumn
\EN
\noindent And Noe did all things which the Lord had commanded him. And after seven days were passed, the waters of the flood overflowed the earth. In the six hundreth year of the life of Noe in the second month, in the seventeenth day of the month, all the fountains of the great deep were broken up, and the flood gates of heaven were open: And the rain fell upon the earth forty days and forty nights.\par
\switchcolumn*

\LA
\begin{Resp}\Rsign Recordátus Dóminus Noë, addúxit spíritum super terram, et imminútæ sunt aquæ: \Star{} Et prohíbitæ sunt plúviæ de cælis.\end{Resp}
\begin{Resp}\Vsign Reversǽque sunt aquæ de terra eúntes et redeúntes, et cœpérunt mínui post centum quinquagínta dies.\end{Resp}
\begin{Resp}\Rsign Et prohíbitæ sunt plúviæ de cælis.\end{Resp}
\switchcolumn
\EN
\begin{Resp}\Rsign The Lord remembered Noah, and made a wind to pass over the earth, and the waters assuaged, \Star{} And the rain from heaven was restrained.\end{Resp}
\begin{Resp}\Vsign And the waters returned from off the earth continually, and after the end of the hundred and fifty days the waters were abated.\end{Resp}
\begin{Resp}\Rsign And the rain from heaven was restrained.\end{Resp}
\switchcolumn*

\LA
\LessonHead{Lectio III}
\switchcolumn
\EN
\LessonHead{Lesson III}
\switchcolumn*

\LA
\Cite{Gen 7:13-14; 7:17}
\switchcolumn
\EN
\Cite{Gen 7:13-14; 7:17}
\switchcolumn*

\LA
\noindent In artículo diéi illíus ingréssus est Noë, et Sem, et Cham, et Japheth fílii ejus; uxor illíus, et tres uxóres filiórum ejus cum eis in arcam: ipsi et omne ánimal secúndum genus suum, universáque juménta in génere suo, et omne quod movétur super terram in génere suo, cunctúmque volátile secúndum genus suum. Factúmque est dilúvium quadragínta diébus super terram: et multiplicátæ sunt aquæ, et elevavérunt arcam in sublíme a terra.\par
\switchcolumn
\EN
\noindent In the selfsame day Noe, and Sem, and Cham, and Japheth his sons: his wife, and the three wives of his sons with them, went into the ark: They and every beast according to its kind, and all the cattle in their kind, and every thing that moveth upon the earth according to its kind, and every fowl according to its kind. And the flood was forty days upon the earth, and the waters increased, and lifted up the ark on high from earth.\par
\switchcolumn*

\LA
\begin{Resp}\Rsign Quadragínta dies et noctes apérti sunt cæli, et ex omni carne habénte spíritum vitæ ingréssa sunt in arcam: \Star{} Et clausit a foris óstium Dóminus.\end{Resp}
\begin{Resp}\Vsign In artículo diéi illíus ingréssus est Noë in arcam et fílii ejus, et uxor illíus et uxóres filiórum ejus.\end{Resp}
\begin{Resp}\Rsign Et clausit a foris óstium Dóminus.\end{Resp}
\begin{Resp}\Vsign Glória Patri, et Fílio, \Star{} et Spirítui Sancto.\end{Resp}
\begin{Resp}\Rsign Et clausit a foris óstium Dóminus.\end{Resp}
\switchcolumn
\EN
\begin{Resp}\Rsign Forty days and forty nights were the heavens opened; and there went into the ark two and two of all flesh wherein is the breath of life. \Star{} And the Lord shut them in.\end{Resp}
\begin{Resp}\Vsign In the self-same day entered Noah into the ark, and his sons, and his wife, and the wives of his sons.\end{Resp}
\begin{Resp}\Rsign And the Lord shut them in.\end{Resp}
\begin{Resp}\Vsign Glory be to the Father, and to the Son, \Star{} and to the Holy Ghost.\end{Resp}
\begin{Resp}\Rsign And the Lord shut them in.\end{Resp}
\switchcolumn*

\end{paracol}

\Office{Feria III infra Hebdomadam Sexagesimæ}{Tuesday after Sexagesima Sunday}{Feria}
\addcontentsline{toc}{subsection}{Feria III infra Hebdomadam Sexagesimæ\enspace\textit{Tuesday after Sexagesima Sunday}}
\begin{paracol}{2}
\LA
\LessonHead{Lectio I}
\switchcolumn
\EN
\LessonHead{Lesson I}
\switchcolumn*

\LA
\LessonTitle{De libro Génesis}
\switchcolumn
\EN
\LessonTitle{Lesson from the book of Genesis}
\switchcolumn*

\LA
\Cite{Gen 8:1-4}
\switchcolumn
\EN
\Cite{Gen 8:1-4}
\switchcolumn*

\LA
\noindent Recordátus autem Deus Noë, cunctorúmque animántium, et ómnium jumentórum, quæ erant cum eo in arca, addúxit spíritum super terram, et imminútæ sunt aquæ. Et clausi sunt fontes abýssi, et cataráctæ cæli: et prohíbitæ sunt plúviæ de cælo. Reversǽque sunt aquæ de terra eúntes et redeúntes: et cœpérunt mínui post centum quinquagínta dies. Requievítque arca mense séptimo, vigésimo séptimo die mensis, super montes Arméniæ.\par
\switchcolumn
\EN
\noindent And God remembered Noe, and all the living creatures, and all the cattle which were with him in the ark, and brought a wind upon the earth, and the waters were abated. The fountains also of the deep, and the flood gates of heaven were shut up, and the rain from heaven was restrained. And the waters returned from off the earth going and coming: and they began to be abated after a hundred and fifty days. And the ark rested in the seventh month, the seven and twentieth day of the month, upon the mountains of Armenia.\par
\switchcolumn*

\LA
\begin{Resp}\Rsign Ædificávit Noë altáre Dómino, ófferens super illud holocáustum: odoratúsque est Dóminus odórem suavitátis, et benedíxit ei, dicens: \Star{} Créscite, et multiplicámini, et repléte terram.\end{Resp}
\begin{Resp}\Vsign Ecce ego státuam pactum meum vobíscum, et cum sémine vestro post vos.\end{Resp}
\begin{Resp}\Rsign Créscite, et multiplicámini, et repléte terram.\end{Resp}
\switchcolumn
\EN
\begin{Resp}\Rsign Noah builded an Altar unto the Lord, and offered burnt offerings on the Altar; and the Lord smelled a sweet savour, and blessed Noah, and said: \Star{} Be fruitful, and multiply, and replenish the earth.\end{Resp}
\begin{Resp}\Vsign Behold, I establish My covenant with you, and with your seed after you.\end{Resp}
\begin{Resp}\Rsign Be fruitful, and multiply, and replenish the earth.\end{Resp}
\switchcolumn*

\LA
\LessonHead{Lectio II}
\switchcolumn
\EN
\LessonHead{Lesson II}
\switchcolumn*

\LA
\Cite{Gen 8:5-9}
\switchcolumn
\EN
\Cite{Gen 8:5-9}
\switchcolumn*

\LA
\noindent At vero aquæ ibant et decrescébant usque ad décimum mensem: décimo enim mense, primo die mensis, apparuérunt cacúmina móntium. Cumque transíssent quadragínta dies, apériens Noë fenéstram arcæ, quam fécerat, dimísit corvum, qui egrediebátur, et non revertebátur, donec siccaréntur aquæ super terram. Emísit quoque colúmbam post eum, ut vidéret si jam cessássent aquæ super fáciem terræ. Quæ cum non invenísset ubi requiésceret pes ejus, revérsa est ad eum in arcam.\par
\switchcolumn
\EN
\noindent And the waters were going and decreasing until the tenth month: for in the tenth month, the first day of the month, the tops of the mountains appeared. And after that forty days were passed, Noe, opening the window of the ark which he had made, sent forth a raven: Which went forth and did not return, till the waters were dried up upon the earth. He sent forth also a dove after him, to see if the waters had now ceased upon the face of the earth. But she, not finding where her foot might rest, returned to him into the ark.\par
\switchcolumn*

\LA
\begin{Resp}\Rsign Ponam arcum meum in núbibus cæli, dixit Dóminus ad Noë: \Star{} Et recordábor fœ́deris mei, quod pépigi tecum.\end{Resp}
\begin{Resp}\Vsign Cumque obdúxero núbibus cælum, apparébit arcus meus in núbibus.\end{Resp}
\begin{Resp}\Rsign Et recordábor fœ́deris mei, quod pépigi tecum.\end{Resp}
\switchcolumn
\EN
\begin{Resp}\Rsign The Lord said unto Noah: I do set My bow in the clouds of heaven, \Star{} And I will remember My covenant which is between Me and you.\end{Resp}
\begin{Resp}\Vsign And it shall come to pass, when I bring a cloud over the heaven, that My bow shall be seen in the cloud.\end{Resp}
\begin{Resp}\Rsign And I will remember My covenant which is between Me and you.\end{Resp}
\switchcolumn*

\LA
\LessonHead{Lectio III}
\switchcolumn
\EN
\LessonHead{Lesson III}
\switchcolumn*

\LA
\Cite{Gen 8:10-13}
\switchcolumn
\EN
\Cite{Gen 8:10-13}
\switchcolumn*

\LA
\noindent Expectátis autem ultra septem diébus áliis, rursum dimísit colúmbam ex arca. At illa venit ad eum ad vésperam, portans ramum olívæ viréntibus fóliis in ore suo: intelléxit ergo Noë quod cessássent aquæ super terram. Expectavítque nihilóminus septem álios dies: et emísit colúmbam, quæ non est revérsa ultra ad eum. Igitur sexcentésimo primo anno, primo mense, prima die mensis, imminútæ sunt aquæ super terram.\par
\switchcolumn
\EN
\noindent And having waited yet seven other days, he again sent forth the dove out of the ark. And she came to him in the evening, carrying a bough of an olive tree, with green leaves, in her mouth. Noe therefore understood that the waters were ceased upon the earth. And he stayed yet other seven days: and he sent forth the dove, which returned not any more unto him. Therefore in the six hundreth and first year, the first month, the first day of the month, the waters were lessened upon the earth.\par
\switchcolumn*

\LA
\begin{Resp}\Rsign Per memetípsum jurávi, dicit Dóminus, non adíciam ultra aquas dilúvii super terram: pacti mei recordábor, \Star{} Ut non perdam aquis dilúvii omnem carnem.\end{Resp}
\begin{Resp}\Vsign Arcum meum ponam in núbibus, et erit signum fœ́deris inter me et inter terram.\end{Resp}
\begin{Resp}\Rsign Ut non perdam aquis dilúvii omnem carnem.\end{Resp}
\begin{Resp}\Vsign Glória Patri, et Fílio, \Star{} et Spirítui Sancto.\end{Resp}
\begin{Resp}\Rsign Ut non perdam aquis dilúvii omnem carnem.\end{Resp}
\switchcolumn
\EN
\begin{Resp}\Rsign By Myself have I sworn, saith the Lord. I will not again bring the waters of the flood upon the earth, I will remember My covenant. \Star{} And the waters shall become no more a flood to destroy all flesh.\end{Resp}
\begin{Resp}\Vsign I do set My bow in the clouds, and it shall be for a token of a covenant between Me and the earth.\end{Resp}
\begin{Resp}\Rsign And the waters shall no more become a flood to destroy all flesh.\end{Resp}
\begin{Resp}\Vsign Glory be to the Father, and to the Son, \Star{} and to the Holy Ghost.\end{Resp}
\begin{Resp}\Rsign And the waters shall no more become a flood to destroy all flesh.\end{Resp}
\switchcolumn*

\end{paracol}

\Office{Feria IV infra Hebdomadam Sexagesimæ}{Wednesday after Sexagesima Sunday}{Feria}
\addcontentsline{toc}{subsection}{Feria IV infra Hebdomadam Sexagesimæ\enspace\textit{Wednesday after Sexagesima Sunday}}
\begin{paracol}{2}
\LA
\LessonHead{Lectio I}
\switchcolumn
\EN
\LessonHead{Lesson I}
\switchcolumn*

\LA
\LessonTitle{De libro Génesis}
\switchcolumn
\EN
\LessonTitle{Lesson from the book of Genesis}
\switchcolumn*

\LA
\Cite{Gen 8:15-19}
\switchcolumn
\EN
\Cite{Gen 8:15-19}
\switchcolumn*

\LA
\noindent Locútus est autem Deus ad Noë, dicens: Egrédere de arca, tu et uxor tua, fílii tui et uxóres filiórum tuórum tecum. Cuncta animántia, quæ sunt apud te, ex omni carne, tam in volatílibus quam in béstiis et univérsis reptílibus, quæ reptant super terram, educ tecum, et ingredímini super terram: créscite et multiplicámini super eam. Egréssus est ergo Noë, et fílii ejus: uxor illíus, et uxóres filiórum ejus cum eo. Sed et ómnia animántia, juménta, et reptília quæ reptant super terram, secúndum genus suum, egréssa sunt de arca.\par
\switchcolumn
\EN
\noindent And God spoke to Noe, saying: Go out of the ark, thou and thy wife, thy sons, and the wives of thy sons with thee. All livings things that are with thee of all flesh, as well in fowls as in beasts, and all creeping things that creep upon the earth, bring out with thee, and go ye upon the earth: increase and multiply upon it. So Noe went out, he and his sons: his wife, and the wives of his sons with him. And all living things, and cattle, and creeping things that creep upon the earth, according to their kinds, went out of the ark.\par
\switchcolumn*

\LA
\begin{Resp}\Rsign Benedíxit Deus Noë, et fíliis ejus, et dixit ad eos: \Star{} Créscite, et multiplicámini, et repléte terram.\end{Resp}
\begin{Resp}\Vsign Ecce ego státuam pactum meum vobíscum, et cum sémine vestro post vos.\end{Resp}
\begin{Resp}\Rsign Créscite, et multiplicámini, et repléte terram.\end{Resp}
\switchcolumn
\EN
\begin{Resp}\Rsign God blessed Noah and his sons, and said unto them: \Star{} Be fruitful, and multiply, and replenish the earth.\end{Resp}
\begin{Resp}\Vsign Behold, I establish My covenant with you, and with your seed after you.\end{Resp}
\begin{Resp}\Rsign Be fruitful, and multiply, and replenish the earth.\end{Resp}
\switchcolumn*

\LA
\LessonHead{Lectio II}
\switchcolumn
\EN
\LessonHead{Lesson II}
\switchcolumn*

\LA
\Cite{Gen 8:20-22}
\switchcolumn
\EN
\Cite{Gen 8:20-22}
\switchcolumn*

\LA
\noindent Ædificávit autem Noë altáre Dómino: et tollens de cunctis pecóribus et volúcribus mundis, óbtulit holocáusta super altáre. Odoratúsque est Dóminus odórem suavitátis, et ait: Nequáquam ultra maledícam terræ propter hómines: sensus enim et cogitátio humáni cordis in malum prona sunt ab adolescéntia sua: non ígitur ultra percútiam omnem ánimam vivéntem sicut feci. Cunctis diébus terræ, seméntis et messis, frigus et æstus, æstas et hiems, nox et dies non requiéscent.\par
\switchcolumn
\EN
\noindent And Noe built an altar unto the Lord: and taking of all cattle and fowls that were clean, offered holocausts upon the altar. And the Lord smelled a sweet savour, and said: I will no more curse the earth for the sake of man: for the imagination and thought of man's heart are prone to evil from his youth: therefore I will no more destroy every living soul as I have done. All the days of the earth, seedtime and harvest, cold and heat, summer and winter, night and day, shall not cease.\par
\switchcolumn*

\LA
\begin{Resp}\Rsign Ecce ego státuam pactum meum vobíscum, et cum sémine vestro post vos: \Star{} Neque erit deínceps dilúvium díssipans terram.\end{Resp}
\begin{Resp}\Vsign Arcum meum ponam in núbibus, et erit signum fœ́deris inter me et inter terram.\end{Resp}
\begin{Resp}\Rsign Neque erit deínceps dilúvium díssipans terram.\end{Resp}
\switchcolumn
\EN
\begin{Resp}\Rsign Behold, I establish My covenant with you, and with your seed after you. \Star{} Neither shall there any more be a flood to destroy the earth.\end{Resp}
\begin{Resp}\Vsign I do set My bow in the clouds, and it shall be for a token of a covenant between Me and the earth.\end{Resp}
\begin{Resp}\Rsign Neither shall there any more be a flood to destroy the earth.\end{Resp}
\switchcolumn*

\LA
\LessonHead{Lectio III}
\switchcolumn
\EN
\LessonHead{Lesson III}
\switchcolumn*

\LA
\Cite{Gen 9:1-6}
\switchcolumn
\EN
\Cite{Gen 9:1-6}
\switchcolumn*

\LA
\noindent Benedixítque Deus Noë et fíliis ejus. Et dixit ad eos: Créscite, et multiplicámini, et repléte terram. Et terror vester ac tremor sit super cuncta animália terræ, et super omnes vólucres cæli, cum univérsis quæ movéntur super terram: omnes pisces maris mánui vestræ tráditi sunt. Et omne, quod movétur et vivit, erit vobis in cibum: quasi ólera viréntia trádidi vobis ómnia. Excépto, quod carnem cum sánguine non comedétis. Sánguinem enim animárum vestrárum requíram de manu cunctárum bestiárum: et de manu hóminis, de manu viri, et fratris ejus requíram ánimam hóminis. Quicúmque effúderit humánum sánguinem, fundétur sanguis illíus: ad imáginem quippe Dei factus est homo.\par
\switchcolumn
\EN
\noindent And God blessed Noe and his sons. And he said to them: Increase and multiply, and fill the earth. And let the fear and dread of you be upon all the beasts of the earth, and upon all the fowls of the air, and all that move upon the earth: all the fishes of the sea are delivered into your hand. And every thing that moveth and liveth shall be meat for you: even as the green herbs have I delivered them all to you: Saving that flesh with blood you shall not eat. For I will require the blood of your lives at the hand of every beast, and at the hand of man, at the hand of every man, and of his brother, will I require the life of man. Whosoever shall shed man's blood, his blood shall be shed: for man was made to the image of God.\par
\switchcolumn*

\LA
\begin{Resp}\Rsign In artículo diéi illíus ingréssus est Noë in arcam et fílii ejus, \Star{} Uxor illíus et uxóres filiórum ejus.\end{Resp}
\begin{Resp}\Vsign Deléta sunt univérsa de terra, remánsit autem solus Noë, et qui cum eo erant in arca.\end{Resp}
\begin{Resp}\Rsign Uxor illíus et uxóres filiórum ejus.\end{Resp}
\begin{Resp}\Vsign Glória Patri, et Fílio, \Star{} et Spirítui Sancto.\end{Resp}
\begin{Resp}\Rsign Uxor illíus et uxóres filiórum ejus.\end{Resp}
\switchcolumn
\EN
\begin{Resp}\Rsign In the self-same day entered Noah into the ark, and his sons, \Star{} And his wife and the wives of his sons.\end{Resp}
\begin{Resp}\Vsign Every living substance was destroyed from the earth, and Noah only remained alive, and they that were with him in the ark.\end{Resp}
\begin{Resp}\Rsign And his wife, and the wives of his sons.\end{Resp}
\begin{Resp}\Vsign Glory be to the Father, and to the Son, \Star{} and to the Holy Ghost.\end{Resp}
\begin{Resp}\Rsign And his wife, and the wives of his sons.\end{Resp}
\switchcolumn*

\end{paracol}

\Office{Feria V infra Hebdomadam Sexagesimæ}{Thursday after Sexagesima Sunday}{Feria}
\addcontentsline{toc}{subsection}{Feria V infra Hebdomadam Sexagesimæ\enspace\textit{Thursday after Sexagesima Sunday}}
\begin{paracol}{2}
\LA
\LessonHead{Lectio I}
\switchcolumn
\EN
\LessonHead{Lesson I}
\switchcolumn*

\LA
\LessonTitle{De libro Génesis}
\switchcolumn
\EN
\LessonTitle{Lesson from the book of Genesis}
\switchcolumn*

\LA
\Cite{Gen 9:12-15}
\switchcolumn
\EN
\Cite{Gen 9:12-15}
\switchcolumn*

\LA
\noindent Dixítque Deus: Hoc signum fœ́deris quod do inter me et vos, et ad omnem ánimam vivéntem, quæ est vobíscum in generatiónes sempitérnas: arcum meum ponam in núbibus, et erit signum fœ́deris inter me et inter terram. Cumque obdúxero núbibus cælum, apparébit arcus meus in núbibus: et recordábor fœ́deris mei vobíscum, et cum omni ánima vivénte quæ carnem végetat: et non erunt ultra aquæ dilúvii ad deléndum univérsam carnem.\par
\switchcolumn
\EN
\noindent And God said: This is the sign of the covenant which I will give between me and you, and to every living soul that is with you, for perpetual generations. I will set my bow in the clouds, and it shall be the sign of a covenant between me, and between the earth. And when I shall cover the sky with clouds, my bow shall appear in the clouds: And I will remember my covenant with you, and with every living soul that beareth flesh: and there shall no more be waters of a flood to destroy all flesh.\par
\switchcolumn*

\LA
\begin{Resp}\Rsign Dixit Dóminus ad Noë: Finis univérsæ carnis venit coram me: repléta est terra iniquitáte eórum, \Star{} Et ego dispérdam eos cum terra.\end{Resp}
\begin{Resp}\Vsign Fac tibi arcam de lignis lævigátis, mansiúnculas in ea fácies.\end{Resp}
\begin{Resp}\Rsign Et ego dispérdam eos cum terra.\end{Resp}
\switchcolumn
\EN
\begin{Resp}\Rsign The Lord said unto Noah: The end of all flesh is come before Me; for the earth is filled with violence through them. \Star{} And I will destroy them with the earth.\end{Resp}
\begin{Resp}\Vsign Make thee an ark of planed timber, rooms shalt thou make in it.\end{Resp}
\begin{Resp}\Rsign And I will destroy them with the earth.\end{Resp}
\switchcolumn*

\LA
\LessonHead{Lectio II}
\switchcolumn
\EN
\LessonHead{Lesson II}
\switchcolumn*

\LA
\Cite{Gen 9:20-23}
\switchcolumn
\EN
\Cite{Gen 9:20-23}
\switchcolumn*

\LA
\noindent Cœpítque Noë vir agrícola exercére terram, et plantávit víneam. Bibénsque vinum inebriátus est, et nudátus in tabernáculo suo. Quod cum vidísset Cham, pater Chánaan, verénda scílicet patris sui esse nudáta, nuntiávit duóbus frátribus suis foras. At vero Sem et Japheth pállium imposuérunt húmeris suis, et incedéntes retrórsum, operuérunt verénda patris sui: faciésque eórum avérsæ erant, et patris virília non vidérunt.\par
\switchcolumn
\EN
\noindent And Noe, a husbandman, began to till the ground, and planted a vineyard. And drinking of the wine was made drunk, and was uncovered in his tent. Which when Cham the father of Chaanan had seen, to wit, that his father's nakedness was uncovered, he told it to his two brethren without. But Sem and Japheth put a cloak upon their shoulders, and going backward, covered the nakedness of their father: and their faces were turned away, and they saw not their father's nakedness.\par
\switchcolumn*

\LA
\begin{Resp}\Rsign Noë, vir justus atque perféctus, cum Deo ambulávit: \Star{} Et fecit ómnia quæcúmque præcépit ei Deus.\end{Resp}
\begin{Resp}\Vsign Fecit sibi arcam, ut salvarétur univérsum semen.\end{Resp}
\begin{Resp}\Rsign Et fecit ómnia quæcúmque præcépit ei Deus.\end{Resp}
\switchcolumn
\EN
\begin{Resp}\Rsign Noah was a just man and perfect; he walked with God. \Star{} According to all that God commanded him, so did he.\end{Resp}
\begin{Resp}\Vsign He made him an ark, that a seed of every sort might be saved alive.\end{Resp}
\begin{Resp}\Rsign According to all that God commanded him, so did he.\end{Resp}
\switchcolumn*

\LA
\LessonHead{Lectio III}
\switchcolumn
\EN
\LessonHead{Lesson III}
\switchcolumn*

\LA
\Cite{Gen 9:24-29}
\switchcolumn
\EN
\Cite{Gen 9:24-29}
\switchcolumn*

\LA
\noindent Evígilans autem Noë ex vino, cum didicísset quæ fécerat ei fílius suus minor, ait: Maledíctus Chánaan, servus servórum erit frátribus suis. Dixítque: Benedíctus Dóminus Deus Sem, sit Chánaan servus ejus. Dilátet Deus Japheth, et hábitet in tabernáculis Sem, sitque Chánaan servus ejus. Vixit autem Noë post dilúvium trecéntis quinquagínta annis. Et impléti sunt omnes dies ejus nongentórum quinquagínta annórum: et mórtuus est.\par
\switchcolumn
\EN
\noindent And Noe awaking from the wine, when he had learned what his younger son had done to him, He said: Cursed be Chaanan, a servant of servants, shall he be unto his brethren. And he said: Blessed be the Lord God of Sem, be Chanaan his servant. May God enlarge Japheth, and may he dwell in the tents of Sem, and Chanaan be his servant. And Noe lived after the flood three hundred and fifty years: And all his days were in the whole nine hundred and fifty years: and he died.\par
\switchcolumn*

\LA
\begin{Resp}\Rsign Quadragínta dies et noctes apérti sunt cæli, et ex omni carne habénte spíritum vitæ ingréssa sunt in arcam: \Star{} Et clausit a foris óstium Dóminus.\end{Resp}
\begin{Resp}\Vsign In artículo diéi illíus ingréssus est Noë in arcam et fílii ejus, et uxor illíus et uxóres filiórum ejus.\end{Resp}
\begin{Resp}\Rsign Et clausit a foris óstium Dóminus.\end{Resp}
\begin{Resp}\Vsign Glória Patri, et Fílio, \Star{} et Spirítui Sancto.\end{Resp}
\begin{Resp}\Rsign Et clausit a foris óstium Dóminus.\end{Resp}
\switchcolumn
\EN
\begin{Resp}\Rsign Forty days and forty nights were the heavens opened; and there went into the ark two and two of all flesh wherein is the breath of life. \Star{} And the Lord shut them in.\end{Resp}
\begin{Resp}\Vsign In the self-same day entered Noah into the ark, and his sons, and his wife, and the wives of his sons.\end{Resp}
\begin{Resp}\Rsign And the Lord shut them in.\end{Resp}
\begin{Resp}\Vsign Glory be to the Father, and to the Son, \Star{} and to the Holy Ghost.\end{Resp}
\begin{Resp}\Rsign And the Lord shut them in.\end{Resp}
\switchcolumn*

\end{paracol}

\Office{Feria VI infra Hebdomadam Sexagesimæ}{Friday after Sexagesima Sunday}{Feria}
\addcontentsline{toc}{subsection}{Feria VI infra Hebdomadam Sexagesimæ\enspace\textit{Friday after Sexagesima Sunday}}
\begin{paracol}{2}
\LA
\LessonHead{Lectio I}
\switchcolumn
\EN
\LessonHead{Lesson I}
\switchcolumn*

\LA
\LessonTitle{De libro Génesis}
\switchcolumn
\EN
\LessonTitle{Lesson from the book of Genesis}
\switchcolumn*

\LA
\Cite{Gen 10:1-6}
\switchcolumn
\EN
\Cite{Gen 10:1-6}
\switchcolumn*

\LA
\noindent Hæ sunt generatiónes filiórum Noë, Sem, Cham et Japheth: natíque sunt eis fílii post dilúvium. Fílii Japheth: Gomer, et Magog, et Mádai, et Javan, et Thubal, et Mosoch, et Thiras. Porro fílii Gomer: Ascenez et Riphath et Thogórma. Fílii autem Javan: Elísa et Tharsis, Cetthim et Dódanim. Ab his divísæ sunt ínsulæ géntium in regiónibus suis, unusquísque secúndum linguam suam et famílias suas in natiónibus suis. Fílii autem Cham: Chus, et Mésraim et Phuth, et Chánaan.\par
\switchcolumn
\EN
\noindent These are the generations of the sons of Noe: Sem, Cham, and Japheth: and unto them sons were born after the flood. The sons of Japheth: Gomer, and Magog, and Madai, and Javan, and Thubal, and Mosoch, and Thiras. And the sons of Gomer: Ascenez and Riphath and Thogorma. And the sons of Javan: Elisa and Tharsis, Cetthim, and Dodanim. By these were divided the islands of the Gentiles in their lands, every one according to his tongue and their families in their nations. And the sons of Cham: Chus, and Mesram, and Phuth, and Chanaan.\par
\switchcolumn*

\LA
\begin{Resp}\Rsign Ædificávit Noë altáre Dómino, ófferens super illud holocáustum: odoratúsque est Dóminus odórem suavitátis, et benedíxit ei, dicens: \Star{} Créscite, et multiplicámini, et repléte terram.\end{Resp}
\begin{Resp}\Vsign Ecce ego státuam pactum meum vobíscum, et cum sémine vestro post vos.\end{Resp}
\begin{Resp}\Rsign Créscite, et multiplicámini, et repléte terram.\end{Resp}
\switchcolumn
\EN
\begin{Resp}\Rsign Noah builded an Altar unto the Lord, and offered burnt offerings on the Altar; and the Lord smelled a sweet savour, and blessed Noah, and said: \Star{} Be fruitful, and multiply, and replenish the earth.\end{Resp}
\begin{Resp}\Vsign Behold, I establish My covenant with you, and with your seed after you.\end{Resp}
\begin{Resp}\Rsign Be fruitful, and multiply, and replenish the earth.\end{Resp}
\switchcolumn*

\LA
\LessonHead{Lectio II}
\switchcolumn
\EN
\LessonHead{Lesson II}
\switchcolumn*

\LA
\Cite{Gen 11:1-4}
\switchcolumn
\EN
\Cite{Gen 11:1-4}
\switchcolumn*

\LA
\noindent Erat autem terra lábii uníus, et sermónum eorúndem. Cumque proficisceréntur de oriénte, invenérunt campum in terra Sénnaar, et habitavérunt in eo. Dixítque alter ad próximum suum: Veníte, faciámus láteres, et coquámus eos igni. Habuerúntque láteres pro saxis, et bitúmen pro cæménto: et dixérunt: Veníte, faciámus nobis civitátem et turrim, cujus culmen pertíngat ad cælum: et celebrémus nomen nostrum ántequam dividámur in univérsas terras.\par
\switchcolumn
\EN
\noindent And the earth was of one tongue, and of the same speech. And when they removed from the east, they found a plain in the land of Sennaar, and dwelt in it. And each one said to his neighbour: Come, let us make brick, and bake them of stones, and slime instead of mortar. And they said: Come, let us make a city and a tower, the top whereof may reach to heaven: and let us make our name famous before we be scattered abroad into all lands.\par
\switchcolumn*

\LA
\begin{Resp}\Rsign Ponam arcum meum in núbibus cæli, dixit Dóminus ad Noë: \Star{} Et recordábor fœ́deris mei, quod pépigi tecum.\end{Resp}
\begin{Resp}\Vsign Cumque obdúxero núbibus cælum, apparébit arcus meus in núbibus.\end{Resp}
\begin{Resp}\Rsign Et recordábor fœ́deris mei, quod pépigi tecum.\end{Resp}
\switchcolumn
\EN
\begin{Resp}\Rsign The Lord said unto Noah: I do set My bow in the clouds of heaven, \Star{} And I will remember My covenant which is between Me and you.\end{Resp}
\begin{Resp}\Vsign And it shall come to pass, when I bring a cloud over the heaven, that My bow shall be seen in the cloud.\end{Resp}
\begin{Resp}\Rsign And I will remember My covenant which is between Me and you.\end{Resp}
\switchcolumn*

\LA
\LessonHead{Lectio III}
\switchcolumn
\EN
\LessonHead{Lesson III}
\switchcolumn*

\LA
\Cite{Gen 11:5-8}
\switchcolumn
\EN
\Cite{Gen 11:5-8}
\switchcolumn*

\LA
\noindent Descéndit autem Dóminus ut vidéret civitátem et turrim, quam ædificábant fílii Adam, et dixit: Ecce, unus est pópulus, et unum lábium ómnibus: cœperúntque hoc fácere, nec desístent a cogitatiónibus suis, donec eas ópere cómpleant. Veníte ígitur, descendámus, et confundámus ibi linguam eórum, ut non áudiat unusquísque vocem próximi sui. Atque ita divísit eos Dóminus ex illo loco in univérsas terras, et cessavérunt ædificáre civitátem.\par
\switchcolumn
\EN
\noindent And the Lord came down to see the city and the tower, which the children of Adam were building. And he said: Behold, it is one people, and all have one tongue: and they have begun to do this, neither will they leave off from their designs, till they accomplish them in deed. Come ye, therefore, let us go down, and there may not understand one another's speech. And so the Lord scattered them from that place into all lands, and they ceased to build the city.\par
\switchcolumn*

\LA
\begin{Resp}\Rsign Per memetípsum jurávi, dicit Dóminus, non adíciam ultra aquas dilúvii super terram: pacti mei recordábor, \Star{} Ut non perdam aquis dilúvii omnem carnem.\end{Resp}
\begin{Resp}\Vsign Arcum meum ponam in núbibus, et erit signum fœ́deris inter me et inter terram.\end{Resp}
\begin{Resp}\Rsign Ut non perdam aquis dilúvii omnem carnem.\end{Resp}
\begin{Resp}\Vsign Glória Patri, et Fílio, \Star{} et Spirítui Sancto.\end{Resp}
\begin{Resp}\Rsign Ut non perdam aquis dilúvii omnem carnem.\end{Resp}
\switchcolumn
\EN
\begin{Resp}\Rsign By Myself have I sworn, saith the Lord. I will not again bring the waters of the flood upon the earth, I will remember My covenant. \Star{} And the waters shall become no more a flood to destroy all flesh.\end{Resp}
\begin{Resp}\Vsign I do set My bow in the clouds, and it shall be for a token of a covenant between Me and the earth.\end{Resp}
\begin{Resp}\Rsign And the waters shall no more become a flood to destroy all flesh.\end{Resp}
\begin{Resp}\Vsign Glory be to the Father, and to the Son, \Star{} and to the Holy Ghost.\end{Resp}
\begin{Resp}\Rsign And the waters shall no more become a flood to destroy all flesh.\end{Resp}
\switchcolumn*

\end{paracol}

\Office{Sabbato infra Hebdomadam Sexagesimæ}{Saturday after Sexagesima Sunday}{Feria}
\addcontentsline{toc}{subsection}{Sabbato infra Hebdomadam Sexagesimæ\enspace\textit{Saturday after Sexagesima Sunday}}
\begin{paracol}{2}
\LA
\LessonHead{Lectio I}
\switchcolumn
\EN
\LessonHead{Lesson I}
\switchcolumn*

\LA
\LessonTitle{De libro Génesis}
\switchcolumn
\EN
\LessonTitle{Lesson from the book of Genesis}
\switchcolumn*

\LA
\Cite{Gen 11:10-15}
\switchcolumn
\EN
\Cite{Gen 11:10-15}
\switchcolumn*

\LA
\noindent Hæ sunt generatiónes Sem: Sem erat centum annórum quando génuit Arpháxad, biénnio post dilúvium. Vixítque Sem, postquam génuit Arpháxad, quingéntis annis: et génuit fílios et fílias. Porro Arpháxad vixit trigínta quinque annis, et génuit Sale. Vixítque Arpháxad, postquam génuit Sale, trecéntis tribus annis: et génuit fílios et fílias. Sale quoque vixit trigínta annis, et génuit Heber. Vixítque Sale, postquam génuit Heber, quadringéntis tribus annis: et génuit fílios et fílias.\par
\switchcolumn
\EN
\noindent These are the generations of Sem: Sem was a hundred years old when he begot Arphaxad, two years after the flood. And Sem lived after he begot Arphaxad, five hundred years, and begot sons and daughters. And Arphaxad lived thirty-five years, and begot Sale. And Arphaxad lived after he begot Sale, three hundred and three years; and begot sons and daughters. Sale also lived thirty years, and begot Heber. And Sale lived after he begot Heber, four hundred and three years; and begot sons and daughters.\par
\switchcolumn*

\LA
\begin{Resp}\Rsign Benedíxit Deus Noë, et fíliis ejus, et dixit ad eos: \Star{} Créscite, et multiplicámini, et repléte terram.\end{Resp}
\begin{Resp}\Vsign Ecce ego státuam pactum meum vobíscum, et cum sémine vestro post vos.\end{Resp}
\begin{Resp}\Rsign Créscite, et multiplicámini, et repléte terram.\end{Resp}
\switchcolumn
\EN
\begin{Resp}\Rsign God blessed Noah and his sons, and said unto them: \Star{} Be fruitful, and multiply, and replenish the earth.\end{Resp}
\begin{Resp}\Vsign Behold, I establish My covenant with you, and with your seed after you.\end{Resp}
\begin{Resp}\Rsign Be fruitful, and multiply, and replenish the earth.\end{Resp}
\switchcolumn*

\LA
\LessonHead{Lectio II}
\switchcolumn
\EN
\LessonHead{Lesson II}
\switchcolumn*

\LA
\Cite{Gen 11:16-23}
\switchcolumn
\EN
\Cite{Gen 11:16-23}
\switchcolumn*

\LA
\noindent Vixit autem Heber trigínta quátuor annis, et génuit Phaleg. Et vixit Heber postquam génuit Phaleg, quadringéntis trigínta annis: et génuit fílios et fílias. Vixit quoque Phaleg trigínta annis, et génuit Reu. Vixítque Phaleg, postquam génuit Reu, ducéntis novem annis: et génuit fílios et fílias. Vixit autem Reu trigínta duóbus annis, et génuit Sarug. Vixit quoque Reu, postquam génuit Sarug, ducéntis septem annis: et génuit fílios et fílias. Vixit vero Sarug trigínta annis, et génuit Nachor. Vixítque Sarug, postquam génuit Nachor, ducéntis annis: et génuit fílios et fílias.\par
\switchcolumn
\EN
\noindent And Heber lived thirty-four years, and begot Phaleg. And Heber lived after he begot Phaleg, four hundred and thirty years: and begot sons and daughters. Phaleg also lived thirty years, and begot Reu. And Phaleg lived after he begot Reu, two hundred and nine years, and begot sons and daughters. And Reu lived thirty-two years, and begot Sarug. And Reu lived after he begot Sarug, two hundred and seven years, and begot sons and daughters. And Sarug lived thirty years, and begot Nachor. And Sarug lived after he begot Nachor, two hundred years: and begot sons and daughters.\par
\switchcolumn*

\LA
\begin{Resp}\Rsign Ecce ego státuam pactum meum vobíscum, et cum sémine vestro post vos: \Star{} Neque erit deínceps dilúvium díssipans terram.\end{Resp}
\begin{Resp}\Vsign Arcum meum ponam in núbibus, et erit signum fœ́deris inter me et inter terram.\end{Resp}
\begin{Resp}\Rsign Neque erit deínceps dilúvium díssipans terram.\end{Resp}
\begin{Resp}\Vsign Glória Patri, et Fílio, \Star{} et Spirítui Sancto.\end{Resp}
\begin{Resp}\Rsign Neque erit deínceps dilúvium díssipans terram.\end{Resp}
\switchcolumn
\EN
\begin{Resp}\Rsign Behold, I establish My covenant with you, and with your seed after you. \Star{} Neither shall there any more be a flood to destroy the earth.\end{Resp}
\begin{Resp}\Vsign I do set My bow in the clouds, and it shall be for a token of a covenant between Me and the earth.\end{Resp}
\begin{Resp}\Rsign Neither shall there any more be a flood to destroy the earth.\end{Resp}
\begin{Resp}\Vsign Glory be to the Father, and to the Son, \Star{} and to the Holy Ghost.\end{Resp}
\begin{Resp}\Rsign Neither shall there any more be a flood to destroy the earth.\end{Resp}
\switchcolumn*

\LA
\LessonHead{Lectio III}
\switchcolumn
\EN
\LessonHead{Lesson III}
\switchcolumn*

\LA
\Cite{Gen 11:24-30}
\switchcolumn
\EN
\Cite{Gen 11:24-30}
\switchcolumn*

\LA
\noindent Vixit autem Nachor vigínti novem annis, et génuit Thare. Vixítque Nachor, postquam génuit Thare, centum decem et novem annis: et génuit fílios et fílias. Vixítque Thare septuagínta annis, et génuit Abram, et Nachor, et Aran. Hæ sunt autem generatiónes Thare: Thare génuit Abram, Nachor et Aran. Porro Aran génuit Lot. Mortuúsque est Aran ante Thare patrem suum, in terra nativitátis suæ, in Ur Chaldæórum. Duxérunt autem Abram et Nachor uxóres: nomen uxóris Abram, Sárai: et nomen uxóris Nachor, Melcha fília Aran, patris Melchæ, et patris Jeschæ. Erat autem Sárai stérilis, nec habébat líberos.\par
\switchcolumn
\EN
\noindent And Nachor lived nine and twenty years, and begot Thare. And Nachor lived after he begot Thare, a hundred and nineteen years: and begot sons and daughters. And Thare lived seventy years, and begot Abram, and Nachor, and Aran. And these are the generations of Thare: Thare begot Abram, Nachor, and Aran. And Aran begot Lot. And Aran died before Thare his father, in the land of his nativity in Ur of the Chaldees. And Abram and Nachor married wives: the name of Abram's wife was Sarai: and the name of Nachor's wife, Melcha, the daughter of Aran, father of Melcha, and father of Jescha. And Sarai was barren, and had no children.\par
\switchcolumn*

\LA
\begin{Resp}\Rsign In artículo diéi illíus ingréssus est Noë in arcam et fílii ejus, \Star{} Uxor illíus et uxóres filiórum ejus.\end{Resp}
\begin{Resp}\Vsign Deléta sunt univérsa de terra, remánsit autem solus Noë, et qui cum eo erant in arca.\end{Resp}
\begin{Resp}\Rsign Uxor illíus et uxóres filiórum ejus.\end{Resp}
\begin{Resp}\Vsign Glória Patri, et Fílio, \Star{} et Spirítui Sancto.\end{Resp}
\begin{Resp}\Rsign Uxor illíus et uxóres filiórum ejus.\end{Resp}
\switchcolumn
\EN
\begin{Resp}\Rsign In the self-same day entered Noah into the ark, and his sons, \Star{} And his wife and the wives of his sons.\end{Resp}
\begin{Resp}\Vsign Every living substance was destroyed from the earth, and Noah only remained alive, and they that were with him in the ark.\end{Resp}
\begin{Resp}\Rsign And his wife, and the wives of his sons.\end{Resp}
\begin{Resp}\Vsign Glory be to the Father, and to the Son, \Star{} and to the Holy Ghost.\end{Resp}
\begin{Resp}\Rsign And his wife, and the wives of his sons.\end{Resp}
\switchcolumn*

\end{paracol}

\Office{Dominica in Quinquagesima}{Quinquagesima Sunday}{Semiduplex II classis}
\addcontentsline{toc}{subsection}{Dominica in Quinquagesima\enspace\textit{Quinquagesima Sunday}}
\begin{paracol}{2}
\LA
\LessonHead{Lectio I}
\switchcolumn
\EN
\LessonHead{Lesson I}
\switchcolumn*

\LA
\LessonTitle{De libro Génesis}
\switchcolumn
\EN
\LessonTitle{Lesson from the book of Genesis}
\switchcolumn*

\LA
\Cite{Gen 12:1-6}
\switchcolumn
\EN
\Cite{Gen 12:1-6}
\switchcolumn*

\LA
\noindent Dixit autem Dóminus ad Abram: Egrédere de terra tua, et de cognatióne tua, et de domo patris tui, et veni in terram quam monstrábo tibi. Faciámque te in gentem magnam, et benedícam tibi, et magnificábo nomen tuum, erísque benedíctus. Benedícam benedicéntibus tibi, et maledícam maledicéntibus tibi, atque in te benedicéntur univérsæ cognatiónes terræ. Egréssus est ítaque Abram sicut præcéperat ei Dóminus, et ivit cum eo Lot: septuagínta quinque annórum erat Abram cum egrederétur de Haran. Tulítque Sárai uxórem suam, et Lot fílium fratris sui, universámque substántiam quam posséderant, et ánimas quas fécerant in Haran: et egréssi sunt ut irent in terram Chánaan. Cumque veníssent in eam, pertransívit Abram terram usque ad locum Sichem, usque ad convállem illústrem: Chananǽus autem tunc erat in terra.\par
\switchcolumn
\EN
\noindent And the Lord said to Abram: Go forth out of thy country, and from thy kindred, and out of thy father's house, and come into the land which I shall show thee. And I will make of thee a great nation, and I will bless thee, and magnify thy name, and thou shalt be blessed. I will bless them that bless thee, and curse them that curse thee, and in thee shall all the kindred of the earth be blessed: So Abram went out as the Lord had commanded him, and Lot went with him: Abram was seventy-five years old when he went forth from Haran. And he took Sarai his wife, and Lot his brother's son, and all the substance which they had gathered, and the souls which they had gotten in Haran: and they went out to go into the land of Chanaan. And when they were come into it, Abram passed through the country into the place of Sichem, as far as the noble vale: now the Chanaanite was at that time in the land.\par
\switchcolumn*

\LA
\begin{Resp}\Rsign Locútus est Dóminus ad Abram, dicens: Egrédere de terra tua, et de cognatióne tua, et veni in terram quam monstrávero tibi: \Star{} Et fáciam te in gentem magnam.\end{Resp}
\begin{Resp}\Vsign Benedícens benedícam tibi, et magnificábo nomen tuum, erísque benedíctus.\end{Resp}
\begin{Resp}\Rsign Et fáciam te in gentem magnam.\end{Resp}
\switchcolumn
\EN
\begin{Resp}\Rsign The Lord spake unto Abram, saying: Get thee out of thy country, and from thy kindred, and go unto the land that I will show thee, \Star{} And I will make of thee a great nation.\end{Resp}
\begin{Resp}\Vsign I will surely bless thee and make thy name great, and thou shalt be blessed.\end{Resp}
\begin{Resp}\Rsign And I will make of thee a great nation.\end{Resp}
\switchcolumn*

\LA
\LessonHead{Lectio II}
\switchcolumn
\EN
\LessonHead{Lesson II}
\switchcolumn*

\LA
\Cite{Gen 12:7-13}
\switchcolumn
\EN
\Cite{Gen 12:7-13}
\switchcolumn*

\LA
\noindent Appáruit autem Dóminus Abram, et dixit ei: Sémini tuo dabo terram hanc. Qui ædificávit ibi altáre Dómino, qui apparúerat ei. Et inde transgrédiens ad montem, qui erat contra oriéntem Bethel, teténdit ibi tabernáculum suum, ab occidénte habens Bethel, et ab oriénte Hai: ædificávit quoque ibi altáre Dómino, et invocávit nomen ejus. Perrexítque Abram vadens, et ultra progrédiens ad merídiem. Facta est autem fames in terra: descendítque Abram in Ægýptum, ut peregrinarétur ibi: prævalúerat enim fames in terra. Cumque prope esset ut ingrederétur Ægýptum, dixit Sárai uxóri suæ: Novi quod pulchra sis múlier: et quod cum víderint te Ægýptii, dictúri sunt: Uxor ipsíus est: et interfícient me, et te reservábunt. Dic ergo, óbsecro te, quod soror mea sis: ut bene sit mihi propter te, et vivat ánima mea ob grátiam tui.\par
\switchcolumn
\EN
\noindent And the Lord appeared to Abram, and said to him: To thy seed will I give this land. And he built there an altar to the Lord, who had appeared to him. And passing on from thence to a mountain, that was on the east side of Bethel, he there pitched his tent, having Bethel on the west, and Hai on the east; he built there also an altar to the Lord, and called upon his name. And Abram went forward, going, and proceeding on to the south. And there came a famine in the country; and Abram went down into Egypt, to sojourn there: for the famine was very grievous in the land. And when he was near to enter into Egypt, he said to Sarai his wife: I know that thou art a beautiful woman: And that when the Egyptians shall see thee, they will say: She is his wife: and they will kill me, and keep thee. Say, therefore, I pray thee, that thou art my sister: that I may be well used for thee, and that my soul may live for thy sake.\par
\switchcolumn*

\LA
\begin{Resp}\Rsign Dum staret Abraham ad ílicem Mambre, vidit tres viros ascendéntes per viam: \Star{} Tres vidit, et unum adorávit.\end{Resp}
\begin{Resp}\Vsign Ecce Sara uxor tua páriet tibi fílium, et vocábis nomen ejus Isaac.\end{Resp}
\begin{Resp}\Rsign Tres vidit, et unum adorávit.\end{Resp}
\switchcolumn
\EN
\begin{Resp}\Rsign Abraham stood by the oak of Mamre, and he saw three men coming up by the path. \Star{} He saw three, and worshipped One.\end{Resp}
\begin{Resp}\Vsign Behold, Sarah thy wife shall bear thee a son, and thou shalt call his name Isaac.\end{Resp}
\begin{Resp}\Rsign He saw three, and worshipped One.\end{Resp}
\switchcolumn*

\LA
\LessonHead{Lectio III}
\switchcolumn
\EN
\LessonHead{Lesson III}
\switchcolumn*

\LA
\Cite{Gen 12:14-19}
\switchcolumn
\EN
\Cite{Gen 12:14-19}
\switchcolumn*

\LA
\noindent Cum ítaque ingréssus esset Abram Ægýptum, vidérunt Ægýptii mulíerem quod esset pulchra nimis. Et nuntiavérunt príncipes Pharaóni, et laudavérunt eam apud illum: et subláta est múlier in domum Pharaónis. Abram vero bene usi sunt propter illam: fuerúntque ei oves et boves et ásini, et servi et fámulæ, et ásinæ et caméli. Flagellávit autem Dóminus Pharaónem plagis máximis, et domum ejus, propter Sárai uxórem Abram. Vocavítque Phárao Abram, et dixit ei: Quidnam est hoc quod fecísti mihi? quare non indicásti quod uxor tua esset? quam ob causam dixísti esse sorórem tuam, ut tóllerem eam mihi in uxórem? Nunc ígitur ecce conjux tua, áccipe eam, et vade.\par
\switchcolumn
\EN
\noindent And when Abram was come into Egypt, the Egyptians saw the woman that she was very beautiful. And the princes told Pharao, and praised her before him: and the woman was taken into the house of Pharao. And they used Abram well for her sake. And he had sheep and oxen, and he asses, and men servants and maid servants, and she asses, and camels. But the Lord scourged Pharao and his house with most grievous stripes for Sarai, Abram's wife. And Pharao called Abram, and said to him: What is this that thou hast done to me? Why didst thou not tell me that she was thy wife? For what cause didst thou say, she was thy sister, that I might take her to my wife? Now, therefore, there is thy wife, take her, and go thy way.\par
\switchcolumn*

\LA
\begin{Resp}\Rsign Tentávit Dóminus Abraham, et dixit ad eum: \Star{} Tolle fílium tuum, quem díligis, Isaac, et offer illum ibi in holocáustum super unum móntium, quem díxero tibi.\end{Resp}
\begin{Resp}\Vsign Vocátus quoque a Dómino, respóndit, Adsum: et ait ei Dóminus.\end{Resp}
\begin{Resp}\Rsign Tolle fílium tuum, quem díligis, Isaac, et offer illum ibi in holocáustum super unum móntium, quem díxero tibi.\end{Resp}
\begin{Resp}\Vsign Glória Patri, et Fílio, \Star{} et Spirítui Sancto.\end{Resp}
\begin{Resp}\Rsign Tolle fílium tuum, quem díligis, Isaac, et offer illum ibi in holocáustum super unum móntium, quem díxero tibi.\end{Resp}
\switchcolumn
\EN
\begin{Resp}\Rsign The Lord did tempt Abraham, and said unto him: \Star{} Take thy son Isaac whom thou lovest, and offer him there for a burnt-offering upon one of the mountains which I will tell thee of.\end{Resp}
\begin{Resp}\Vsign And when the Lord called him, he answered: Behold, here I am. And the Lord said unto him:\end{Resp}
\begin{Resp}\Rsign Take thy son Isaac whom thou lovest, and offer him there for a burnt-offering upon one of the mountains which I will tell thee of.\end{Resp}
\begin{Resp}\Vsign Glory be to the Father, and to the Son, \Star{} and to the Holy Ghost.\end{Resp}
\begin{Resp}\Rsign Take thy son Isaac whom thou lovest, and offer him there for a burnt-offering upon one of the mountains which I will tell thee of.\end{Resp}
\switchcolumn*

\LA
\LessonHead{Lectio IV}
\switchcolumn
\EN
\LessonHead{Lesson IV}
\switchcolumn*

\LA
\LessonTitle{Ex libro sancti Ambrósii Epíscopi de Abraham Patriárcha}
\switchcolumn
\EN
\LessonTitle{From the Book upon the Patriarch Abraham written by St. Ambrose, Bishop (of Milan.)}
\switchcolumn*

\LA
\Source{Liber 1, cap 2}
\switchcolumn
\EN
\Source{Bk. I, 2}
\switchcolumn*

\LA
\noindent Magnus plane vir Abraham, et multárum virtútum clarus insígnibus, quem votis suis philosophía non pótuit æquare. Dénique minus est quod illa finxit, quam quod iste gessit: majórque ambitióso eloquéntiæ mendácio simplex veritátis fides. Itaque cujúsmodi fúerit in eo viro devótio, considerémus. Ea enim virtus órdine prima est, quæ est fundaméntum ceterárum: meritóque hanc ab eo primam exégit Deus, dicens: Exi de terra tua, et de cognatióne tua, et de domo patris tui. Satis fúerat dixísse: De terra tua. Ibi enim erat exíre de cognatióne, exíre de patérna domo.\par
\switchcolumn
\EN
\noindent Abraham was truly a great man, illustrious as an example of many virtues; one the like of whom the day-dreams of Philosophy have not been able to produce. That which she imagineth is less than that which he did; his simple truth and faith were something grander than her lying rounded periods. Let us then consider what this man's loyalty was. For that virtue is first to be taken which was the source of all the others, and thus this was the first which God called for from him, when He said: Get thee out of thy country, and from thy kindred, and from thy father's house. It would have been enough to have said, Get thee out of thy country, for there were his kindred, and there his father's house.\par
\switchcolumn*

\LA
\begin{Resp}\Rsign Angelus Dómini vocávit Abraham, dicens: \Star{} Ne exténdas manum tuam super púerum, eo quod tímeas Dóminum.\end{Resp}
\begin{Resp}\Vsign Cumque extendísset manum ut immoláret fílium, ecce Angelus Dómini de cælo clamávit, dicens.\end{Resp}
\begin{Resp}\Rsign Ne exténdas manum tuam super púerum, eo quod tímeas Dóminum.\end{Resp}
\switchcolumn
\EN
\begin{Resp}\Rsign The Angel of the Lord called unto Abraham and said: \Star{} Lay not thine hand upon the lad; for now I know that thou fearest God.\end{Resp}
\begin{Resp}\Vsign And Abraham stretched forth his hand to slay his son; and, behold, the Angel of the Lord called unto him out of heaven, and said:\end{Resp}
\begin{Resp}\Rsign Lay not thy hand upon the lad; for now I know that thou fearest God.\end{Resp}
\switchcolumn*

\LA
\LessonHead{Lectio V}
\switchcolumn
\EN
\LessonHead{Lesson V}
\switchcolumn*

\LA
\noindent Sed ideo áddidit síngula, ut ejus afféctum probáret: ne forte aut imprudénter cœpísse viderétur, aut fraus áliqua mandátis cæléstibus pararétur. Sed sicut coacervánda fuérunt præcépta, ne quid latéret: ita étiam proponénda præmia, ne forte desperáret. Tentátur ut fortis, incitátur ut fidélis, provocátur ut justus: meritóque exívit, quemádmodum locútus est illi Dóminus. Et exívit cum eo Lot. Hoc autem, quod pro magno inter septem sapiéntum dicta celebrátur: Séquere Deum; perfécit Abraham, factóque sapiéntum dicta prævénit, et secútus Deum exívit de terra sua.\par
\switchcolumn
\EN
\noindent But He gave the details of his sacrifice one by one, that He might see whether he loved Him, lest also he should begin rashly, or should seek to evade the heavenly commandment. But as the whole of the precept was plainly set forth, lest anything should be unconsidered, so also were the rewards set forth, lest the burden should seem hopeless. He was tried as one that is strong, he was roused as one that is true, he was called as one that is righteous; and he departed loyally as the Lord had spoken unto him. And Lot went forth with him. That saying of the Seven Wise Men of Greece is much spoken of Follow God. But this did Abraham before the Seven Wise Men were thought of; he followed God, and went out of his own land.\par
\switchcolumn*

\LA
\begin{Resp}\Rsign Vocávit Angelus Dómini Abraham de cælo, secúndo, dicens: Benedícam tibi, \Star{} Et multiplicábo te sicut stellas cæli.\end{Resp}
\begin{Resp}\Vsign Possidébit semen tuum portas inimicórum tuórum, et benedicéntur in sémine tuo omnes tribus terræ.\end{Resp}
\begin{Resp}\Rsign Et multiplicábo te sicut stellas cæli.\end{Resp}
\switchcolumn
\EN
\begin{Resp}\Rsign The Angel of the Lord called unto Abraham out of heaven the second time, \Star{} And said: I will bless thee, and I will multiply thy seed as the stars of the heaven.\end{Resp}
\begin{Resp}\Vsign Thy seed shall possess the gate of his enemies, and in thy seed shall all the nations of the earth be blessed.\end{Resp}
\begin{Resp}\Rsign And I will multiply thy seed as the stars of the heaven.\end{Resp}
\switchcolumn*

\LA
\LessonHead{Lectio VI}
\switchcolumn
\EN
\LessonHead{Lesson VI}
\switchcolumn*

\LA
\noindent Sed quia ante terra ei fúerat ália, hoc est, régio Chaldæórum, de qua exívit Thare pater Abrahæ, et in Charran demigrávit: et quia secum edúxit nepótem suum, cui dictum fúerat, Exi de cognatióne tua: considerémus, ne forte hoc sit exíre de terra sua, de hujus terræ, hoc est, de córporis nostri quadam commemoratióne égredi, de qua exívit Paulus, qui dixit: Nostra autem conversátio in cælis est.\par
\switchcolumn
\EN
\noindent But, forasmuch as Abraham had before had another country, namely, the land of the Chaldees, whence went forth Terah the father of Abraham, and came unto Haran, and forasmuch as he to whom it had been said, Get thee out from thy kindred, took Lot, his brother's son, with him, let us consider whether this Get thee out of thy country signifieth not get thee out of this earthly dwelling, namely, our body, from which Paul came forth, who said, Our conversation is in heaven. (Phil. iii. 20.)\par
\switchcolumn*

\LA
\begin{Resp}\Rsign Deus dómini mei Abraham, dírige viam meam: \Star{} Ut cum salúte revértar in domum dómini mei.\end{Resp}
\begin{Resp}\Vsign Obsecro, Dómine, fac misericórdiam cum servo tuo.\end{Resp}
\begin{Resp}\Rsign Ut cum salúte revértar in domum dómini mei.\end{Resp}
\begin{Resp}\Vsign Glória Patri, et Fílio, \Star{} et Spirítui Sancto.\end{Resp}
\begin{Resp}\Rsign Ut cum salúte revértar in domum dómini mei.\end{Resp}
\switchcolumn
\EN
\begin{Resp}\Rsign O God of my master Abraham, prosper my way which I go; \Star{} That I may return again in safety unto the house of my master.\end{Resp}
\begin{Resp}\Vsign O Lord, I pray thee, be merciful unto thy servant.\end{Resp}
\begin{Resp}\Rsign That I may return again in safety unto the house of my master.\end{Resp}
\begin{Resp}\Vsign Glory be to the Father, and to the Son, \Star{} and to the Holy Ghost.\end{Resp}
\begin{Resp}\Rsign That I may return again in safety unto the house of my master.\end{Resp}
\switchcolumn*

\LA
\LessonHead{Lectio VII}
\switchcolumn
\EN
\LessonHead{Lesson VII}
\switchcolumn*

\LA
\LessonTitle{Léctio sancti Evangélii secúndum Lucam}
\switchcolumn
\EN
\LessonTitle{From the Holy Gospel according to Luke}
\switchcolumn*

\LA
\Cite{Luc 18:31-43}
\switchcolumn
\EN
\Cite{Luke 18:31-43}
\switchcolumn*

\LA
\noindent In illo témpore: Assúmpsit autem Jesus duódecim, et ait illis: Ecce ascéndimus Jerosólymam, et consummabúntur ómnia quæ scripta sunt per Prophétas de Fílio hóminis. Et réliqua.\par
\switchcolumn
\EN
\noindent In that time, Jesus took unto him the twelve, and said to them: Behold, we go up to Jerusalem, and all things shall be accomplished which were written by the prophets concerning the Son of man. And so on.\par
\switchcolumn*

\LA
\LessonTitle{Homilía sancti Gregórii Papæ}
\switchcolumn
\EN
\LessonTitle{Homily by Pope St. Gregory (the Great.)}
\switchcolumn*

\LA
\Source{Homilia 2 in Evangelia}
\switchcolumn
\EN
\Source{2nd on the Gospels}
\switchcolumn*

\LA
\noindent Redémptor noster, prǽvidens ex passióne sua discipulórum ánimos perturbándos, eis longe ante et ejúsdem passiónis pœnam, et resurrectiónis suæ glóriam prædíxit: ut cum eum moriéntem, sicut prædíctum est, cérnerent; étiam resurrectúrum non dubitárent. Sed quia carnáles adhuc discípuli nullo modo valébant cápere verba mystérii, venítur ad miráculum. Ante eórum óculos cæcus lumen recépit: ut qui cæléstis mystérii verba non cáperent, eos ad fidem cæléstia facta solidárent.\par
\switchcolumn
\EN
\noindent Our Redeemer, foreseeing that the minds of His disciples would be troubled by His suffering, told them long before both of the pains of that suffering, and of the glory of His rising again, to the end that, when they should see Him die as He had prophesied, they might not doubt that He was likewise to rise again. But, since His disciples were yet carnal, and could not receive the words telling of this mystery, He wrought a miracle before them. A blind man received his sight before their eyes, that if they could not receive heavenly things by words, they might be persuaded of heavenly things by deeds.\par
\switchcolumn*

\LA
\begin{Resp}\Rsign Veni hódie ad fontem aquæ, et orávi Dóminum, dicens: \Star{} Dómine, Deus Abraham, tu prósperum fecísti desidérium meum.\end{Resp}
\begin{Resp}\Vsign Igitur puélla, cui díxero, Da mihi aquam de hýdria tua, ut bibam: et illa díxerit, Bibe, dómine, et camélis tuis potum tríbuam: ipsa est, quam præparávit Dóminus fílio dómini mei.\end{Resp}
\begin{Resp}\Rsign Dómine, Deus Abraham, tu prósperum fecísti desidérium meum.\end{Resp}
\switchcolumn
\EN
\begin{Resp}\Rsign I came this day unto the well, and I besought the Lord, and said: \Star{} O Lord God of Abraham, Thou hast prospered my way.\end{Resp}
\begin{Resp}\Vsign Therefore the virgin to whom I shall say: Give me water of thy pitcher to drink; and she shall say to me: Drink, my lord, and I will give thy camels drink also; let the same be the woman whom the Lord hath appointed out for my master's son.\end{Resp}
\begin{Resp}\Rsign O Lord God of Abraham, Thou hast prospered my way.\end{Resp}
\switchcolumn*

\LA
\LessonHead{Lectio VIII}
\switchcolumn
\EN
\LessonHead{Lesson VIII}
\switchcolumn*

\LA
\noindent Sed mirácula Dómini et Salvatóris nostri sic accipiénda sunt, fratres caríssimi, ut et in veritáte credántur facta, et tamen per significatiónem nobis aliquid ínnuant. Opera quippe ejus et per poténtiam áliud osténdunt, et per mystérium áliud loquúntur. Ecce enim, quis juxta históriam cæcus iste fúerit, ignorámus: sed tamen quid per mystérium signíficet, nóvimus. Cæcum quippe est genus humánum, quod in parénte primo a paradísi gáudiis expúlsum, claritátem supérnæ lucis ignórans, damnatiónis suæ ténebras pátitur. Sed tamen per Redemptóris sui præséntiam illuminátur: ut intérnæ lucis gáudia jam per desidérium vídeat, atque in via vitæ boni óperis gressus ponat.\par
\switchcolumn
\EN
\noindent But, dearly beloved brethren, we must so take the miracles of our Lord and Saviour, as believing both that they were actually wrought, and that they have some mystic interpretation for our instruction. For in His works, power speaketh one thing and mystery again another. Behold here, for instance. We know not historically who this blind man was, but we do know of what he was mystically the figure. Mankind is blind, driven out from Eden in the persons of his first parents, knowing not the light of heaven, and suffering the darkness of condemnation. But, nevertheless, through the coming of his Redeemer, he is enlightened, so that now he seeth by hope already the gladness of inward light, and walketh by good works in the path of life.\par
\switchcolumn*

\LA
\begin{Resp}\Rsign Factus est sermo Dómini ad Abram, dicens: \Star{} Noli timére, Abram: ego protéctor tuus sum, et merces tua magna nimis.\end{Resp}
\begin{Resp}\Vsign Ego enim sum Dóminus Deus tuus, qui edúxi te de Ur Chaldæórum.\end{Resp}
\begin{Resp}\Rsign Noli timére, Abram: ego protéctor tuus sum, et merces tua magna nimis.\end{Resp}
\switchcolumn
\EN
\begin{Resp}\Rsign The word of the Lord came unto Abram, saying: \Star{} Fear not, Abram I am thy shield, and thy exceeding great reward.\end{Resp}
\begin{Resp}\Vsign For I am the Lord thy God That brought thee out of Ur of the Chaldees.\end{Resp}
\begin{Resp}\Rsign Fear not, Abram I am thy shield, and thy exceeding great reward.\end{Resp}
\switchcolumn*

\LA
\LessonHead{Lectio IX}
\switchcolumn
\EN
\LessonHead{Lesson IX}
\switchcolumn*

\LA
\noindent Notándum vero est, quod cum Jesus Jéricho appropinquáre dícitur, cæcus illuminátur. Jéricho quippe luna interpretátur: luna autem in sacro elóquio pro deféctu carnis pónitur: quia dum ménstruis moméntis decréscit, deféctum nostræ mortalitátis desígnat. Dum ígitur Cónditor noster appropínquat Jéricho, cæcus ad lumen redit: quia dum divínitas deféctum nostræ carnis suscépit, humánum genus lumen, quod amíserat, recépit. Unde enim Deus humána pátitur, inde homo ad divína sublevátur. Qui vidélicet cæcus recte et juxta viam sedére, et mendícans esse descríbitur. Ipsa enim Véritas dicit: Ego sum via.\par
\switchcolumn
\EN
\noindent One must note that as Jesus drew to Jericho a blind man received his sight. Now, this name Jericho, being interpreted, signifieth the city of the moon and in Holy Scripture the moon is used as a figure of our imperfect flesh, of whose gradual corruption her monthly waning is a type. As, therefore, our Maker draweth nigh to Jericho, a blind man receiveth his sight. While the Godhead taketh. into itself our weak manhood, man receiveth again the light which he had lost. By God's suffering in the Manhood, man is raised up toward God. This blind man is also well described as sitting by the wayside begging, for the Truth saith: “I am the Way.” (John XIV, 6.)\par
\switchcolumn*

\LA
\begin{Resp}\Rsign Cæcus sedébat secus viam, transeúnte Dómino, et clamávit ad eum: et ait illi Dóminus: \Star{} Dómine, ut vídeam lumen.\end{Resp}
\begin{Resp}\Vsign Stans autem Jesus, jussit illum duci ad se, et cum appropinquásset, interrogávit eum, dicens.\end{Resp}
\begin{Resp}\Rsign Quid vis ut fáciam tibi?\end{Resp}
\begin{Resp}\Vsign Glória Patri, et Fílio, \Star{} et Spirítui Sancto.\end{Resp}
\begin{Resp}\Rsign Dómine, ut vídeam lumen.\end{Resp}
\switchcolumn
\EN
\begin{Resp}\Rsign As the Lord passed by, a certain blind man sat by the way-side, and cried unto Him. And the Lord asked him, saying: \Star{} What wilt thou that I shall do unto thee? * Lord, that I may receive my sight.\end{Resp}
\begin{Resp}\Vsign And Jesus stood, and commanded him to be brought unto Him; and when he was come near, He asked him, saying:\end{Resp}
\begin{Resp}\Rsign What wilt thou that I shall do unto thee?\end{Resp}
\begin{Resp}\Vsign Glory be to the Father, and to the Son, \Star{} and to the Holy Ghost.\end{Resp}
\begin{Resp}\Rsign Lord, that I may receive my sight.\end{Resp}
\switchcolumn*

\end{paracol}

\Office{Feria II infra Hebdomadam Quinquagesimæ}{Monday after Quinquagesima Sunday}{Feria}
\addcontentsline{toc}{subsection}{Feria II infra Hebdomadam Quinquagesimæ\enspace\textit{Monday after Quinquagesima Sunday}}
\begin{paracol}{2}
\LA
\LessonHead{Lectio I}
\switchcolumn
\EN
\LessonHead{Lesson I}
\switchcolumn*

\LA
\LessonTitle{De libro Génesis}
\switchcolumn
\EN
\LessonTitle{Lesson from the book of Genesis}
\switchcolumn*

\LA
\Cite{Gen 13:1-6}
\switchcolumn
\EN
\Cite{Gen 13:1-6}
\switchcolumn*

\LA
\noindent Ascéndit ergo Abram de Ægýpto, ipse et uxor ejus, et ómnia quæ habébat, et Lot cum eo, ad austrálem plagam. Erat autem dives valde in possessióne auri et argénti. Reversúsque est per iter, quo vénerat, a merídie in Bethel, usque ad locum ubi prius fíxerat tabernáculum inter Bethel et Hai, in loco altáris quod fécerat prius: et invocávit ibi nomen Dómini. Sed et Lot qui erat cum Abram, fuérunt greges óvium, et arménta, et tabernácula. Nec póterat eos cápere terra, ut habitárent simul: erat quippe substántia eórum multa, et nequíbant habitáre commúniter.\par
\switchcolumn
\EN
\noindent And Abram went up out of Egypt, he and his wife, and all that he had, and Lot with him, into the south. And he was very rich in possession of gold and silver. And he returned by the way that he came, from the south to Bethel, to the place where before he had pitched his tent between Bethel and Hai: In the place of the altar which he had made before; and there he called upon the name of the Lord. But Lot also, who was with Abram, had flocks of sheep, and herds of beasts, and tents. Neither was the land able to bear them, that they might dwell together: for their substance was great, and they could not dwell together.\par
\switchcolumn*

\LA
\begin{Resp}\Rsign Movens Abram tabernáculum suum, venit et habitávit juxta convállem Mambre: \Star{} Ædificavítque ibi altáre Dómino.\end{Resp}
\begin{Resp}\Vsign Dixit autem Dóminus ad eum: Leva óculos tuos, et vide: omnem terram, quam cónspicis tibi dabo, et sémini tuo in sempitérnum.\end{Resp}
\begin{Resp}\Rsign Ædificavítque ibi altáre Dómino.\end{Resp}
\switchcolumn
\EN
\begin{Resp}\Rsign Abram removed his tent, and came, and dwelt by the vale of Mamre \Star{} And built there an altar unto the Lord.\end{Resp}
\begin{Resp}\Vsign And the Lord said unto him: Lift up thine eyes, and look; all the land which thou seest, to thee will I give it, and to thy seed for ever.\end{Resp}
\begin{Resp}\Rsign And built there an altar unto the Lord.\end{Resp}
\switchcolumn*

\LA
\LessonHead{Lectio II}
\switchcolumn
\EN
\LessonHead{Lesson II}
\switchcolumn*

\LA
\Cite{Gen 13:7-11}
\switchcolumn
\EN
\Cite{Gen 13:7-11}
\switchcolumn*

\LA
\noindent Unde et facta est rixa inter pastóres gregum Abram et Lot. Eo autem témpore Chananǽus et Pherezǽus habitábant in terra illa. Dixit ergo Abram ad Lot: Ne quæso sit iúrgium inter me et te, et inter pastóres meos et pastóres tuos: fratres enim sumus. Ecce univérsa terra coram te est: recéde a me, óbsecro: si ad sinístram íeris, ego déxteram tenébo: si tu déxteram elégeris, ego ad sinístram pergam. Elevátis ítaque Lot óculis, vidit omnem circa regiónem Jordánis, quæ univérsa irrigabátur ántequam subvérteret Dóminus Sódomam et Gomórrham, sicut paradísus Dómini, et sicut Ægýptus veniéntibus in Segor. Elegítque sibi Lot regiónem circa Jordánem, et recéssit ab oriénte.\par
\switchcolumn
\EN
\noindent Whereupon also there arose a strife between the herdsmen of Abram and of Lot. And at that time the Chanaanite and the Pherezite dwelled in that country. Abram therefore said to Lot: Let there be no quarrel, I beseech thee, between me and thee, and between my herdsmen and thy herdsmen: for we are brethren. Behold the whole land is before thee: depart from me I pray thee: if thou wilt go to the left hand, I will take the right: if thou choose the right hand, I will pass to the left. And Lot, lifting up his eyes, saw all the country about the Jordan, which was watered throughout, before the Lord destroyed Sodom and Gomorrha, as the paradise of the Lord, and like Egypt as one comes to Segor. And Lot chose to himself the country about the Jordan, and he departed from the east.\par
\switchcolumn*

\LA
\begin{Resp}\Rsign Crédidit Abram Deo, et reputátum est ei ad justítiam: \Star{} Et ídeo amícus Dei factus est.\end{Resp}
\begin{Resp}\Vsign Fuit autem justus coram Dómino, et ambulávit in viis ejus.\end{Resp}
\begin{Resp}\Rsign Et ídeo amícus Dei factus est.\end{Resp}
\switchcolumn
\EN
\begin{Resp}\Rsign Abraham believed God, and it was counted unto him for righteousness. \Star{} And therefore he became the friend of God.\end{Resp}
\begin{Resp}\Vsign For he was righteous in the sight of the Lord, and walked in His ways.\end{Resp}
\begin{Resp}\Rsign And therefore he became the friend of God.\end{Resp}
\switchcolumn*

\LA
\LessonHead{Lectio III}
\switchcolumn
\EN
\LessonHead{Lesson III}
\switchcolumn*

\LA
\Cite{Gen 13:11-16}
\switchcolumn
\EN
\Cite{Gen 13:11-16}
\switchcolumn*

\LA
\noindent Divisíque sunt altérutrum a fratre suo. Abram habitávit in terra Chánaan; Lot vero morátus est in óppidis, quæ erant circa Jordánem, et habitávit in Sódomis. Hómines autem Sodomítæ péssimi erant, et peccatóres coram Dómino nimis. Dixítque Dóminus ad Abram, postquam divísus est ab eo Lot: Leva óculos tuos et vide a loco, in quo nunc es, ad aquilónem et merídiem, ad oriéntem et occidéntem. Omnem terram, quam cónspicis, tibi dabo, et sémini tuo usque in sempitérnum. Faciámque semen tuum sicut púlverem terræ.\par
\switchcolumn
\EN
\noindent And they were separated one brother from the other. Abram dwelt in the land of Chanaan; and Lot abode in the towns that were about the Jordan, and dwelt in Sodom. And the men of Sodom were very wicked, and sinners before the face of the Lord, beyond measure. And the Lord said to Abram, after Lot was separated from him: Lift up thy eyes, and look from the place wherein thou now art, to the north and to the south, to the east and to the west. All the land which thou seest, I will give to thee, and to thy seed for ever. And I will make thy seed as the dust of the earth.\par
\switchcolumn*

\LA
\begin{Resp}\Rsign Tentávit Dóminus Abraham, et dixit ad eum: \Star{} Tolle fílium tuum, quem díligis, Isaac, et offer illum ibi in holocáustum super unum móntium, quem díxero tibi.\end{Resp}
\begin{Resp}\Vsign Vocátus quoque a Dómino, respóndit, Adsum: et ait ei Dóminus.\end{Resp}
\begin{Resp}\Rsign Tolle fílium tuum, quem díligis, Isaac, et offer illum ibi in holocáustum super unum móntium, quem díxero tibi.\end{Resp}
\begin{Resp}\Vsign Glória Patri, et Fílio, \Star{} et Spirítui Sancto.\end{Resp}
\begin{Resp}\Rsign Tolle fílium tuum, quem díligis, Isaac, et offer illum ibi in holocáustum super unum móntium, quem díxero tibi.\end{Resp}
\switchcolumn
\EN
\begin{Resp}\Rsign The Lord did tempt Abraham, and said unto him: \Star{} Take thy son Isaac whom thou lovest, and offer him there for a burnt-offering upon one of the mountains which I will tell thee of.\end{Resp}
\begin{Resp}\Vsign And when the Lord called him, he answered: Behold, here I am. And the Lord said unto him:\end{Resp}
\begin{Resp}\Rsign Take thy son Isaac whom thou lovest, and offer him there for a burnt-offering upon one of the mountains which I will tell thee of.\end{Resp}
\begin{Resp}\Vsign Glory be to the Father, and to the Son, \Star{} and to the Holy Ghost.\end{Resp}
\begin{Resp}\Rsign Take thy son Isaac whom thou lovest, and offer him there for a burnt-offering upon one of the mountains which I will tell thee of.\end{Resp}
\switchcolumn*

\end{paracol}

\Office{Feria III infra Hebdomadam Quinquagesimæ}{Tuesday after Quinquagesima Sunday}{Feria}
\addcontentsline{toc}{subsection}{Feria III infra Hebdomadam Quinquagesimæ\enspace\textit{Tuesday after Quinquagesima Sunday}}
\begin{paracol}{2}
\LA
\LessonHead{Lectio I}
\switchcolumn
\EN
\LessonHead{Lesson I}
\switchcolumn*

\LA
\LessonTitle{De libro Génesis}
\switchcolumn
\EN
\LessonTitle{Lesson from the book of Genesis}
\switchcolumn*

\LA
\Cite{Gen 14:8-12}
\switchcolumn
\EN
\Cite{Gen 14:8-12}
\switchcolumn*

\LA
\noindent Et egréssi sunt rex Sodomórum, et rex Gomórrhæ, rexque Adamæ, et rex Séboim, necnon et rex Balæ, quæ est Segor: et direxérunt áciem contra eos in valle Silvéstri: scílicet advérsus Chodorláhomor regem Elamitárum, et Thadal regem Géntium, et Amraphel regem Sénnaar, et Arioch regem Ponti: quátuor reges advérsus quinque. Vallis autem Silvéstris habébat púteos multos bitúminis. Itaque rex Sodomórum, et Gomórrhæ, terga vertérunt, ceciderúntque ibi: et qui remánserant, fugérunt ad montem. Tulérunt autem omnem substántiam Sodomórum et Gomórrhæ, et univérsa quæ ad cibum pértinent, et abiérunt: necnon et Lot, et substántiam ejus, fílium fratris Abram, qui habitábat in Sódomis.\par
\switchcolumn
\EN
\noindent And the king of Sodom, and the king of Gomorrha, and the king of Adama, and the king of Seboim, and the king of Bala, which is Segor, went out: and they set themselves against them in battle array in the woodland vale: To wit, against Chodorlahomor king of the Elamites, and Thadal king of nations, and Amraphel king of Sennaar, and Arioch king of Pontus: four kings against five. Now the woodland vale had many pits of slime. And the king of Sodom, and the king of Gomorrha turned their backs and were overthrown there: and they that remained fled to the mountain. And they took all the substance of the Sodomites, and Gomorrhites, and all their victuals, and went their way: And Lot also, the son of Abram's brother, who dwelt in Sodom, and his substance.\par
\switchcolumn*

\LA
\begin{Resp}\Rsign Angelus Dómini vocávit Abraham, dicens: \Star{} Ne exténdas manum tuam super púerum, eo quod tímeas Dóminum.\end{Resp}
\begin{Resp}\Vsign Cumque extendísset manum ut immoláret fílium, ecce Angelus Dómini de cælo clamávit, dicens.\end{Resp}
\begin{Resp}\Rsign Ne exténdas manum tuam super púerum, eo quod tímeas Dóminum.\end{Resp}
\switchcolumn
\EN
\begin{Resp}\Rsign The Angel of the Lord called unto Abraham and said: \Star{} Lay not thine hand upon the lad; for now I know that thou fearest God.\end{Resp}
\begin{Resp}\Vsign And Abraham stretched forth his hand to slay his son; and, behold, the Angel of the Lord called unto him out of heaven, and said:\end{Resp}
\begin{Resp}\Rsign Lay not thy hand upon the lad; for now I know that thou fearest God.\end{Resp}
\switchcolumn*

\LA
\LessonHead{Lectio II}
\switchcolumn
\EN
\LessonHead{Lesson II}
\switchcolumn*

\LA
\Cite{Gen 14:13-16}
\switchcolumn
\EN
\Cite{Gen 14:13-16}
\switchcolumn*

\LA
\noindent Et ecce unus, qui eváserat, nuntiávit Abram Hebrǽo, qui habitábat in conválle Mambre Amorrhǽi, fratris Escol, et fratris Aner: hi enim pepígerant fœdus cum Abram. Quod cum audísset Abram, captum vidélicet Lot fratrem suum, numerávit expedítos vernáculos suos trecéntos decem et octo: et persecútus est usque Dan. Et divísis sóciis, írruit super eos nocte: percussítque eos, et persecútus est eos usque Hoba, quæ est ad lævam Damásci. Reduxítque omnem substántiam, et Lot fratrem suum cum substántia illíus, mulíeres quoque et pópulum.\par
\switchcolumn
\EN
\noindent And behold one that had escaped told Abram the Hebrew, who dwelt in the vale of Mambre the Amorrhite, the brother of Escol, and the brother of Aner: for these had made league with Abram. Which when Abram had heard, to wit, that his brother Lot was taken, he numbered of the servants born in his house, three hundred and eighteen well appointed: and pursued them to Dan. And dividing his company, he rushed upon them in the night: and defeated them, and pursued them as far as Hoba, which is on the left hand of Damascus. And he brought back all the substance, and Lot his brother, with his substance, the women also the people.\par
\switchcolumn*

\LA
\begin{Resp}\Rsign Vocávit Angelus Dómini Abraham de cælo, secúndo, dicens: Benedícam tibi, \Star{} Et multiplicábo te sicut stellas cæli.\end{Resp}
\begin{Resp}\Vsign Possidébit semen tuum portas inimicórum tuórum, et benedicéntur in sémine tuo omnes tribus terræ.\end{Resp}
\begin{Resp}\Rsign Et multiplicábo te sicut stellas cæli.\end{Resp}
\switchcolumn
\EN
\begin{Resp}\Rsign The Angel of the Lord called unto Abraham out of heaven the second time, \Star{} And said: I will bless thee, and I will multiply thy seed as the stars of the heaven.\end{Resp}
\begin{Resp}\Vsign Thy seed shall possess the gate of his enemies, and in thy seed shall all the nations of the earth be blessed.\end{Resp}
\begin{Resp}\Rsign And I will multiply thy seed as the stars of the heaven.\end{Resp}
\switchcolumn*

\LA
\LessonHead{Lectio III}
\switchcolumn
\EN
\LessonHead{Lesson III}
\switchcolumn*

\LA
\Cite{Gen 14:17-20}
\switchcolumn
\EN
\Cite{Gen 14:17-20}
\switchcolumn*

\LA
\noindent Egréssus est autem rex Sodomórum in occúrsum ejus postquam revérsus est a cæde Chodorláhomor, et regum qui cum eo erant in valle Save, quæ est vallis regis. At vero Melchísedech rex Salem, próferens panem et vinum, erat enim sacérdos Dei altíssimi, benedíxit ei, et ait: Benedíctus Abram Deo excélso, qui creávit cælum et terram: et benedíctus Deus excélsus, quo protegénte, hostes in mánibus tuis sunt. Et dedit ei décimas ex ómnibus.\par
\switchcolumn
\EN
\noindent And the king of Sodom went out to meet him, after he returned from the slaughter of Chodorlahomor, and of the kings that were with him in the vale of Save, which is the king's vale. But Melchisedech the king of Salem, bringing forth bread and wine, for he was the priest of the most high God, Blessed him, and said: Blessed be Abram by the most high God, who created heaven and earth. And blessed be the most high God, by whose protection the enemies are in thy hands. And he gave him the tithes of all.\par
\switchcolumn*

\LA
\begin{Resp}\Rsign Deus dómini mei Abraham, dírige viam meam: \Star{} Ut cum salúte revértar in domum dómini mei.\end{Resp}
\begin{Resp}\Vsign Obsecro, Dómine, fac misericórdiam cum servo tuo.\end{Resp}
\begin{Resp}\Rsign Ut cum salúte revértar in domum dómini mei.\end{Resp}
\begin{Resp}\Vsign Glória Patri, et Fílio, \Star{} et Spirítui Sancto.\end{Resp}
\begin{Resp}\Rsign Ut cum salúte revértar in domum dómini mei.\end{Resp}
\switchcolumn
\EN
\begin{Resp}\Rsign O God of my master Abraham, prosper my way which I go; \Star{} That I may return again in safety unto the house of my master.\end{Resp}
\begin{Resp}\Vsign O Lord, I pray thee, be merciful unto thy servant.\end{Resp}
\begin{Resp}\Rsign That I may return again in safety unto the house of my master.\end{Resp}
\begin{Resp}\Vsign Glory be to the Father, and to the Son, \Star{} and to the Holy Ghost.\end{Resp}
\begin{Resp}\Rsign That I may return again in safety unto the house of my master.\end{Resp}
\switchcolumn*

\end{paracol}

\Office{Feria IV Cinerum}{Ash Wednesday}{Feria privilegiata}
\addcontentsline{toc}{subsection}{Feria IV Cinerum\enspace\textit{Ash Wednesday}}
\begin{paracol}{2}
\LA
\LessonHead{Lectio I}
\switchcolumn
\EN
\LessonHead{Lesson I}
\switchcolumn*

\LA
\LessonTitle{Léctio sancti Evangélii secúndum Matthǽum}
\switchcolumn
\EN
\LessonTitle{Continuation of the Holy Gospel according to Matthew}
\switchcolumn*

\LA
\Cite{Matt 6:16-21}
\switchcolumn
\EN
\Cite{Matt 6:16-21}
\switchcolumn*

\LA
\noindent In illo témpore: Dixit Jesus discípulis suis: Cum jejunátis, nolíte fíeri sicut hypócritæ, tristes. Et réliqua.\par
\switchcolumn
\EN
\noindent In that time Jesus said to his disciples: And when you fast, be not as the hypocrites, sad. For they disfigure their faces, that they may appear unto men to fast. And so on.\par
\switchcolumn*

\LA
\LessonTitle{Homilía sancti Augustíni Epíscopi}
\switchcolumn
\EN
\LessonTitle{Homily by St. Augustine, Bishop (of Hippo.)}
\switchcolumn*

\LA
\Source{Lib. 2 de Serm. Domini in monte, cap. 12 tom. 4}
\switchcolumn
\EN
\Source{Bk. ii on the Lord's Sermon on the Mount, ch. xii, tom. 4}
\switchcolumn*

\LA
\noindent Maniféstum est, his præcéptis omnem nostram intentiónem in interióra gáudia dírigi: ne foris quæréntes mercédem, huic sǽculo conformémur, et amittámus promissiónem tanto solidióris atque firmióris, quanto interióris beatitúdinis, qua nos elégit Deus confórmes fíeri imáginis Fílii sui. In hoc autem capítulo máxime adverténdum est, non in solo rerum corporeárum nitóre atque pompa, sed étiam in ipsis sórdibus luctuósis esse posse jactántiam et eo periculosiórem, quo sub nómine servitútis Dei décipit.\par
\switchcolumn
\EN
\noindent It is evident that by these precepts we are bidden to seek for inner gladness, lest, by running after that reward which is without, we should become conformed to the fashion of this world, and should so lose the promise of that blessing which is all the truer and more stable that it is inward, that blessing wherein God hath chosen us to be conformed to the likeness of His Son. In this chapter we will principally consider the fact that vain-glory findeth a ground for its exercise in struggling poverty as much as in worldly distinction and display; and this development is the most dangerous, because it entices under pretence of being the serving of God.\par
\switchcolumn*

\LA
\begin{Resp}\Rsign Veni hódie ad fontem aquæ, et orávi Dóminum, dicens: \Star{} Dómine, Deus Abraham, tu prósperum fecísti desidérium meum.\end{Resp}
\begin{Resp}\Vsign Igitur puélla, cui díxero, Da mihi aquam de hýdria tua, ut bibam: et illa díxerit, Bibe, dómine, et camélis tuis potum tríbuam: ipsa est, quam præparávit Dóminus fílio dómini mei.\end{Resp}
\begin{Resp}\Rsign Dómine, Deus Abraham, tu prósperum fecísti desidérium meum.\end{Resp}
\switchcolumn
\EN
\begin{Resp}\Rsign I came this day unto the well, and I besought the Lord, and said: \Star{} O Lord God of Abraham, Thou hast prospered my way.\end{Resp}
\begin{Resp}\Vsign Therefore the virgin to whom I shall say: Give me water of thy pitcher to drink; and she shall say to me: Drink, my lord, and I will give thy camels drink also; let the same be the woman whom the Lord hath appointed out for my master's son.\end{Resp}
\begin{Resp}\Rsign O Lord God of Abraham, Thou hast prospered my way.\end{Resp}
\switchcolumn*

\LA
\LessonHead{Lectio II}
\switchcolumn
\EN
\LessonHead{Lesson II}
\switchcolumn*

\LA
\noindent Qui ergo immoderáto cultu córporis atque vestítu, vel ceterárum rerum nitóre præfúlget, fácile convíncitur rebus ipsis, pompárum sǽculi esse sectátor, nec quemquam fallit dolósa imágine sanctitátis. Qui autem in professióne christianitátis, inusitáto squalóre ac sórdibus inténtos in se óculos hóminum facit, cum id voluntáte fáciat, non necessitáte patiátur: ex céteris ejus opéribus potest cónici, utrum hoc contémptu supérflui cultus, an ambitióne áliqua fáciat: quia et sub ovína pelle cavéndos lupos Dóminus præcépit: Sed ex frúctibus, inquit, eórum cognoscétis eos.\par
\switchcolumn
\EN
\noindent He that is characterised by unbridled indulgence in luxury or in dress, or any other display, is by these very things easily shown to be a follower of worldly vanities, and deceiveth no one by putting on an hypocritical mask of godliness. But those professors of Christianity, who turn all eyes on themselves by an eccentric show of grovelling and dirtiness, not suffered by necessity, but by their own choice, of them we must judge by their other works whether their conduct really proceedeth from the desire of mortification by giving up unnecessary comfort, or is only the mean of some ambition the Lord biddeth us beware of wolves in sheep's clothing, but by their fruits, saith He, ye shall know them.\par
\switchcolumn*

\LA
\begin{Resp}\Rsign Factus est sermo Dómini ad Abram, dicens: \Star{} Noli timére, Abram: ego protéctor tuus sum, et merces tua magna nimis.\end{Resp}
\begin{Resp}\Vsign Ego enim sum Dóminus Deus tuus, qui edúxi te de Ur Chaldæórum.\end{Resp}
\begin{Resp}\Rsign Noli timére, Abram: ego protéctor tuus sum, et merces tua magna nimis.\end{Resp}
\switchcolumn
\EN
\begin{Resp}\Rsign The word of the Lord came unto Abram, saying: \Star{} Fear not, Abram I am thy shield, and thy exceeding great reward.\end{Resp}
\begin{Resp}\Vsign For I am the Lord thy God That brought thee out of Ur of the Chaldees.\end{Resp}
\begin{Resp}\Rsign Fear not, Abram I am thy shield, and thy exceeding great reward.\end{Resp}
\switchcolumn*

\LA
\LessonHead{Lectio III}
\switchcolumn
\EN
\LessonHead{Lesson III}
\switchcolumn*

\LA
\noindent Cum enim cœ́perint alíquibus tentatiónibus ea ipsa scílicet illis súbtrahi, vel negári, quæ isto velámine vel consecúti sunt, vel cónsequi cúpiunt: tunc necésse est ut appáreat, utrum lupus in ovína pelle sit, an ovis in sua. Non tamen proptérea ornátu supérfluo debet aspéctus hóminum mulcére Christiánus, quia illum parcum hábitum ac necessárium étiam simulatóres sǽpius usúrpant, ut incáutos decípiant: quia et illæ oves non debent pelles suas depónere, si aliquándo eis lupi se cóntegant.\par
\switchcolumn
\EN
\noindent The test is when, by diverse trials, such persons lose those things which under the cover of seeming unworldliness they have either gained or sought to gain. Then must it needs appear whether they be wolves in sheep's clothing, or indeed sheep in their own. But that hypocrites do the contrary maketh it no duty of a Christian to shine before the eyes of men with a display of needless luxury the sheep need not to lay aside their own clothing because wolves sometimes falsely assume it.\par
\switchcolumn*

\LA
\begin{Resp}\Rsign Movens Abram tabernáculum suum, venit et habitávit juxta convállem Mambre: \Star{} Ædificavítque ibi altáre Dómino.\end{Resp}
\begin{Resp}\Vsign Dixit autem Dóminus ad eum: Leva óculos tuos, et vide: omnem terram, quam cónspicis tibi dabo, et sémini tuo in sempitérnum.\end{Resp}
\begin{Resp}\Rsign Ædificavítque ibi altáre Dómino.\end{Resp}
\begin{Resp}\Vsign Glória Patri, et Fílio, \Star{} et Spirítui Sancto.\end{Resp}
\begin{Resp}\Rsign Ædificavítque ibi altáre Dómino.\end{Resp}
\switchcolumn
\EN
\begin{Resp}\Rsign Abram removed his tent, and came, and dwelt by the vale of Mamre \Star{} And built there an altar unto the Lord.\end{Resp}
\begin{Resp}\Vsign And the Lord said unto him: Lift up thine eyes, and look; all the land which thou seest, to thee will I give it, and to thy seed for ever.\end{Resp}
\begin{Resp}\Rsign And built there an altar unto the Lord.\end{Resp}
\begin{Resp}\Vsign Glory be to the Father, and to the Son, \Star{} and to the Holy Ghost.\end{Resp}
\begin{Resp}\Rsign And built there an altar unto the Lord.\end{Resp}
\switchcolumn*

\end{paracol}

\Office{Feria V post Cineres}{Thursday after Ash Wednesday}{Feria major}
\addcontentsline{toc}{subsection}{Feria V post Cineres\enspace\textit{Thursday after Ash Wednesday}}
\begin{paracol}{2}
\LA
\LessonHead{Lectio I}
\switchcolumn
\EN
\LessonHead{Lesson I}
\switchcolumn*

\LA
\LessonTitle{Léctio sancti Evangélii secúndum Matthǽum}
\switchcolumn
\EN
\LessonTitle{Continuation of the Holy Gospel according to Matthew}
\switchcolumn*

\LA
\Cite{Matt 8:5-13}
\switchcolumn
\EN
\Cite{Matt 8:5-13}
\switchcolumn*

\LA
\noindent In illo témpore: Cum introísset Jesus Caphárnaum, accéssit ad eum centúrio, rogans eum, et dicens: Dómine, puer meus jacet in domo paralýticus, et male torquétur. Et réliqua.\par
\switchcolumn
\EN
\noindent At that time, when Jesus had entered into Capharnaum, there came to him a centurion, beseeching him and saying: Lord, my servant lieth at home sick of the palsy, and is grieviously tormented. And so on.\par
\switchcolumn*

\LA
\LessonTitle{Homilía sancti Augustíni Epíscopi}
\switchcolumn
\EN
\LessonTitle{Homily by St. Augustine, Bishop (of Hippo.)}
\switchcolumn*

\LA
\Source{Lib. 2 de Consensu Evangelist. cap. 20, tom. 4}
\switchcolumn
\EN
\Source{Book ii on the Agreement of the Evangelists, ch. xx, tom. 4}
\switchcolumn*

\LA
\noindent Videámus, utrum sibi de hoc servo centuriónis Matthǽus Lucásque conséntiant. Matthǽus enim dicit: Accéssit ad eum centúrio, rogans eum, et dicens: Puer meus jacet in domo paralýticus. Cui vidétur repugnáre quod ait Lucas: Et cum audísset de Jesu, misit ad eum senióres Judæórum, rogans eum ut veníret, et sanáret servum ejus. At illi cum veníssent ad Jesum, rogábant eum sollícite, dicéntes ei: Quia dignus est ut hoc illi præstes: díligit enim gentem nostram, et synagógam ipse ædificávit nobis. Jesus autem ibat cum illis: et cum jam non longe esset a domo, misit ad eum centúrio amícos, dicens: Dómine, noli vexári: non enim dignus sum ut sub tectum meum intres.\par
\switchcolumn
\EN
\noindent Let us consider whether Matthew and Luke are at one as touching this centurion's servant. Matthew saith: “There came unto Him a centurion, beseeching Him, and saying: ‘Lord, my servant lieth at home sick, of the palsy.’” This seemeth to differ from what Luke saith, namely: “And when he heard of Jesus, he sent unto Him the elders of the Jews, beseeching Him that He would come and heal his servant. And when they came to Jesus, they besought him instantly, saying: “That he was worthy for whom He should do this; for he loveth our nation, and he hath built us a synagogue.” Then Jesus went with them; and when He was now not far from the house, the centurion sent friends to Him, saying unto Him: “Lord, trouble not thyself; for I am not worthy that Thou shouldest enter under my roof.” (vii. 6, seq.)\par
\switchcolumn*

\LA
\begin{Resp}\Rsign Dómine, puer meus jacet paralýticus in domo, et male torquétur: \Star{} Amen dico tibi, ego véniam, et curábo eum.\end{Resp}
\begin{Resp}\Vsign Dómine, non sum dignus ut intres sub tectum meum: sed tantum dic verbo, et sanábitur puer meus.\end{Resp}
\begin{Resp}\Rsign Amen dico tibi, ego véniam, et curábo eum.\end{Resp}
\switchcolumn
\EN
\begin{Resp}\Rsign Lord, my servant lieth at home sick of the palsy, and grievously tormented. \Star{} Amen, I say unto thee, I will come and heal him:\end{Resp}
\begin{Resp}\Vsign Lord, I am not worthy that Thou shouldest enter under my roof, but speak the word only, and my servant shall be healed:\end{Resp}
\begin{Resp}\Rsign Amen, I say unto thee, I will come and heal him.\end{Resp}
\switchcolumn*

\LA
\LessonHead{Lectio II}
\switchcolumn
\EN
\LessonHead{Lesson II}
\switchcolumn*

\LA
\noindent Si enim hoc ita gestum est, quómodo erit verum, quod Matthǽus narrat: Accéssit ad eum quidam centúrio, cum ipse non accésserit, sed amícos míserit: nisi diligénter adverténtes intellegámus Matthǽum non omnímodo deseruísse usitátum morem loquéndi? Non solum enim dícere solémus, accessísse áliquem étiam antequam pervéniat illuc, quo dícitur accessísse: unde étiam dícimus: Parum accéssit, vel multum accéssit eo, quo áppetit perveníre: verum étiam ipsam perventiónem cujus adipiscéndi causa accéditur, dícimus plerúmque factam, etsi eum, ad quem pérvenit, non vídeat ille qui pérvenit, cum per amícum pérvenit ad áliquem, cujus ei favor est necessárius. Quod ita ténuit consuetúdo, ut jam étiam vulgo perventóres appelléntur, qui poténtium quorúmlibet tamquam inaccessíbiles ánimos, per conveniéntium personárum interpositiónem, ambitiónis arte pertíngunt.\par
\switchcolumn
\EN
\noindent If it were done thus, how is Matthew truthful, when he saith that the centurion came unto Him, seeing that, in fact, he sent his friends? We must then look well into this, and we shall see that Matthew only made use of a common form of speech. Now, we use to say of a man that he cometh to a place even though he be not already come: whence also we say, He arrived close; or He arrived long way off, that is, to that place to which he would come; yea, we speak of that coming, toward which he tendeth, as though it had already taken place, when he that should be come at, seeth not yet him that cometh, but is come at for him by friends, to obtain his favour, which is needful for him that would come to him. And so much doth this manner of speaking hold, that they are commonly said to come at a great man, (who is beyond their personal reach,) who, by means of suitable persons, succeed in laying before him such things as they desire.\par
\switchcolumn*

\LA
\begin{Resp}\Rsign Dum staret Abraham ad ílicem Mambre, vidit tres viros ascendéntes per viam: \Star{} Tres vidit, et unum adorávit.\end{Resp}
\begin{Resp}\Vsign Ecce Sara uxor tua páriet tibi fílium, et vocábis nomen ejus Isaac.\end{Resp}
\begin{Resp}\Rsign Tres vidit, et unum adorávit.\end{Resp}
\switchcolumn
\EN
\begin{Resp}\Rsign Abraham stood by the oak of Mamre, and he saw three men coming up by the path. \Star{} He saw three, and worshipped One.\end{Resp}
\begin{Resp}\Vsign Behold, Sarah thy wife shall bear thee a son, and thou shalt call his name Isaac.\end{Resp}
\begin{Resp}\Rsign He saw three, and worshipped One.\end{Resp}
\switchcolumn*

\LA
\LessonHead{Lectio III}
\switchcolumn
\EN
\LessonHead{Lesson III}
\switchcolumn*

\LA
\noindent Non ergo absúrde Matthǽus, étiam quod vulgo possit intéllegi, per álios facto accéssu centuriónis ad Dóminum, compéndio dícere vóluit: Accéssit ad eum centúrio. Verúmtamen non negligénter intuénda est étiam sancti Evangelístæ altitúdo mýsticæ locutiónis, secúndum quam scriptum est in Psalmo: Accédite ad eum, et illuminámini. Proínde quia fidem centuriónis, qua vere accéditur ad Jesum, ipse ita laudávit, ut díceret: Non invéni tantam fidem in Israël: ipsum pótius accessísse ad Christum dícere vóluit prudens Evangelísta, quam illos, per quos verba sua míserat.\par
\switchcolumn
\EN
\noindent Therefore it is not strange that Matthew should make use of the common short phrase, and say of the centurion, who reached the Lord's sympathies, by mean of friends, that he came unto Him. Also we must needs not pass by lightly the mystic depth which underlieth the words of this holy Evangelist. It is written in the Psalms (xxxiii. 6): Draw near unto Him and be enlightened. Thus did the centurion in faith draw near unto Jesus, and the Lord so praised him that He said: I have not found so great faith, no, not in Israel. Of him of whom these words were spoken, the Evangelist deemeth it wiser to say that he had found his way to Jesus; that he had got to Christ, than that they came, through whom he sent his message unto Him.\par
\switchcolumn*

\LA
\begin{Resp}\Rsign Tentávit Dóminus Abraham, et dixit ad eum: \Star{} Tolle fílium tuum, quem díligis, Isaac, et offer illum ibi in holocáustum super unum móntium, quem díxero tibi.\end{Resp}
\begin{Resp}\Vsign Vocátus quoque a Dómino, respóndit, Adsum: et ait ei Dóminus.\end{Resp}
\begin{Resp}\Rsign Tolle fílium tuum, quem díligis, Isaac, et offer illum ibi in holocáustum super unum móntium, quem díxero tibi.\end{Resp}
\begin{Resp}\Vsign Glória Patri, et Fílio, \Star{} et Spirítui Sancto.\end{Resp}
\begin{Resp}\Rsign Tolle fílium tuum, quem díligis, Isaac, et offer illum ibi in holocáustum super unum móntium, quem díxero tibi.\end{Resp}
\switchcolumn
\EN
\begin{Resp}\Rsign The Lord did tempt Abraham, and said unto him: \Star{} Take thy son Isaac whom thou lovest, and offer him there for a burnt-offering upon one of the mountains which I will tell thee of.\end{Resp}
\begin{Resp}\Vsign And when the Lord called him, he answered: Behold, here I am. And the Lord said unto him:\end{Resp}
\begin{Resp}\Rsign Take thy son Isaac whom thou lovest, and offer him there for a burnt-offering upon one of the mountains which I will tell thee of.\end{Resp}
\begin{Resp}\Vsign Glory be to the Father, and to the Son, \Star{} and to the Holy Ghost.\end{Resp}
\begin{Resp}\Rsign Take thy son Isaac whom thou lovest, and offer him there for a burnt-offering upon one of the mountains which I will tell thee of.\end{Resp}
\switchcolumn*

\end{paracol}

\Office{Feria VI post Cineres}{Friday after Ash Wednesday}{Feria major}
\addcontentsline{toc}{subsection}{Feria VI post Cineres\enspace\textit{Friday after Ash Wednesday}}
\begin{paracol}{2}
\LA
\LessonHead{Lectio I}
\switchcolumn
\EN
\LessonHead{Lesson I}
\switchcolumn*

\LA
\LessonTitle{Léctio sancti Evangélii secúndum Matthǽum}
\switchcolumn
\EN
\LessonTitle{Continuation of the Holy Gospel according to Matthew}
\switchcolumn*

\LA
\Cite{Matt 5:43-48; 6:1-4}
\switchcolumn
\EN
\Cite{Matt 5:43-48; 6:1-4}
\switchcolumn*

\LA
\noindent In illo témpore: Dixit Jesus discípulis suis: Audístis quia dictum est: Díliges próximum tuum, et ódio habébis inimícum tuum. Et réliqua.\par
\switchcolumn
\EN
\noindent In that time, Jesus said to his disciples: You have heard that it hath been said, Thou shalt love thy neighbour, and hate thy enemy. And so on.\par
\switchcolumn*

\LA
\LessonTitle{Homilía sancti Hierónymi Presbýteri}
\switchcolumn
\EN
\LessonTitle{Homily by St. Jerome, Priest (at Bethlehem.)}
\switchcolumn*

\LA
\Source{Lib. 1 Comm. in cap. 5 et 6 Matth.}
\switchcolumn
\EN
\Source{Bk. i, Comm. on Matth. v and vi}
\switchcolumn*

\LA
\noindent Ego autem dico vobis: Dilígite inimícos vestros; benefácite his qui odérunt vos. Multi præcépta Dei, imbecillitáte sua, non Sanctórum víribus æstimántes, putant esse impossibília quæ præcépta sunt: et dicunt suffícere virtútibus, non odísse inimícos: céterum dilígere, plus prǽcipi, quam humána natúra patiátur. Sciéndum est ergo, Christum non impossibília præcípere, sed perfécta. Quæ fecit David in Saul, et in Absalom: Stéphanus quoque Martyr pro inimícis lapidántibus deprecátus est: et Paulus anáthema cupit esse pro persecutóribus suis. Hæc autem Jesus et dócuit et fecit, dicens: Pater, ignósce illis: quod enim fáciunt, nésciunt.\par
\switchcolumn
\EN
\noindent But I say unto you: “Love your enemies, do good to them that hate you.” There are many who judge of the commandments of the Lord by their own weakness, and not by the strength of His Saints; and so deem Him to have commanded things impossible. These are they who think that not to hate their enemies is all that they are able to do; and that to command us to love them, is to command more than man's nature can bear. It behoveth them to know, that this which Christ commandeth is not impossible, albeit perfect. This is what David did in respect of Saul and Absalom; the martyr Stephen also prayed for his enemies, even while they were stoning him; and Paul could wish that himself were accursed from Christ for his persecutors. (Rom. ix. 3.) And this, Jesus Himself did, as well as taught, when He said: “Father, forgive them for they know not what they do.” (Luke xxiii. 34.)\par
\switchcolumn*

\LA
\begin{Resp}\Rsign Angelus Dómini vocávit Abraham, dicens: \Star{} Ne exténdas manum tuam super púerum, eo quod tímeas Dóminum.\end{Resp}
\begin{Resp}\Vsign Cumque extendísset manum ut immoláret fílium, ecce Angelus Dómini de cælo clamávit, dicens.\end{Resp}
\begin{Resp}\Rsign Ne exténdas manum tuam super púerum, eo quod tímeas Dóminum.\end{Resp}
\switchcolumn
\EN
\begin{Resp}\Rsign The Angel of the Lord called unto Abraham and said: \Star{} Lay not thine hand upon the lad; for now I know that thou fearest God.\end{Resp}
\begin{Resp}\Vsign And Abraham stretched forth his hand to slay his son; and, behold, the Angel of the Lord called unto him out of heaven, and said:\end{Resp}
\begin{Resp}\Rsign Lay not thy hand upon the lad; for now I know that thou fearest God.\end{Resp}
\switchcolumn*

\LA
\LessonHead{Lectio II}
\switchcolumn
\EN
\LessonHead{Lesson II}
\switchcolumn*

\LA
\noindent Ut sitis fílii Patris vestri, qui in cælis est. Si Dei præcépta custódiens, fílius quis effícitur Dei: ergo non est natúra fílius, sed arbítrio suo. Cum ergo facis eleemósynam, noli tuba cánere ante te, sicut hypócritæ fáciunt in synagógis et in vicis, ut honorificéntur ab homínibus. Qui tuba canit, eleemósynam fáciens, hypócrita est. Qui jejúnans demolítur fáciem suam, ut ventris inanitátem monstret in vultu, et hic hypócrita est. Qui in synagógis et in ángulis plateárum orat, ut videátur ab homínibus, hypócrita est.\par
\switchcolumn
\EN
\noindent That ye may be the children of your Father Which is in heaven. If he that doeth the commandments of God becometh a son of God, then is he not a son by nature, but by his own choice. Therefore when thou doest thine alms, do not sound a trumpet before thee, as the hypocrites do in the synagogues, and in the streets, that they may have glory of men. He that soundeth a trumpet before him, when he doeth alms, is an hypocrite. He that disfigureth his face, when he fasteth, to the end that he may show the emptiness of his belly in his looks, he also is an hypocrite.\par
\switchcolumn*

\LA
\begin{Resp}\Rsign Vocávit Angelus Dómini Abraham de cælo, secúndo, dicens: Benedícam tibi, \Star{} Et multiplicábo te sicut stellas cæli.\end{Resp}
\begin{Resp}\Vsign Possidébit semen tuum portas inimicórum tuórum, et benedicéntur in sémine tuo omnes tribus terræ.\end{Resp}
\begin{Resp}\Rsign Et multiplicábo te sicut stellas cæli.\end{Resp}
\switchcolumn
\EN
\begin{Resp}\Rsign The Angel of the Lord called unto Abraham out of heaven the second time, \Star{} And said: I will bless thee, and I will multiply thy seed as the stars of the heaven.\end{Resp}
\begin{Resp}\Vsign Thy seed shall possess the gate of his enemies, and in thy seed shall all the nations of the earth be blessed.\end{Resp}
\begin{Resp}\Rsign And I will multiply thy seed as the stars of the heaven.\end{Resp}
\switchcolumn*

\LA
\LessonHead{Lectio III}
\switchcolumn
\EN
\LessonHead{Lesson III}
\switchcolumn*

\LA
\noindent Ex quibus ómnibus collígitur hypócritas esse, qui quódlibet fáciunt ut ab homínibus glorificéntur. Mihi vidétur et ille, qui dicit fratri suo: Dimítte ut tollam festúcam de óculo tuo: nam propter glóriam hoc fácere vidétur, ut ipse justus esse videátur. Unde dícitur ei a Dómino: Hypócrita, éice primum trabem de óculo tuo. Non ítaque virtus, sed causa virtútis apud Deum mercédem habet. Et si a recta via páululum declináveris, non ínterest, utrum ad déxteram vadas, an ad sinístram, cum verum iter amíseris.\par
\switchcolumn
\EN
\noindent He that prayeth in the synagogues and in the corners of the streets, that he may be seen of men, is an hypocrite. From all which, we gather that an hypocrite is one which doeth anything that he may have glory of men. To me also it seemeth that he which saith unto his brother: “Let me pull out the mote out of thine eye”, (vii. 4) that he also is an hypocrite; for he proposeth to take upon him that office for vainglory's sake, that he himself may appear righteous. Therefore the Lord saith unto him: “Thou hypocrite, first cast out the beam out of thine own eye.” Thus we see that it is, not the doing good, but the motive which moveth us to do good, which will meet with reward from God; and, if thou stray but a little from the right way, it is of small moment whether thou wander to the right hand or to the left, when once thou hast lost the straight path.\par
\switchcolumn*

\LA
\begin{Resp}\Rsign Deus dómini mei Abraham, dírige viam meam: \Star{} Ut cum salúte revértar in domum dómini mei.\end{Resp}
\begin{Resp}\Vsign Obsecro, Dómine, fac misericórdiam cum servo tuo.\end{Resp}
\begin{Resp}\Rsign Ut cum salúte revértar in domum dómini mei.\end{Resp}
\begin{Resp}\Vsign Glória Patri, et Fílio, \Star{} et Spirítui Sancto.\end{Resp}
\begin{Resp}\Rsign Ut cum salúte revértar in domum dómini mei.\end{Resp}
\switchcolumn
\EN
\begin{Resp}\Rsign O God of my master Abraham, prosper my way which I go; \Star{} That I may return again in safety unto the house of my master.\end{Resp}
\begin{Resp}\Vsign O Lord, I pray thee, be merciful unto thy servant.\end{Resp}
\begin{Resp}\Rsign That I may return again in safety unto the house of my master.\end{Resp}
\begin{Resp}\Vsign Glory be to the Father, and to the Son, \Star{} and to the Holy Ghost.\end{Resp}
\begin{Resp}\Rsign That I may return again in safety unto the house of my master.\end{Resp}
\switchcolumn*

\end{paracol}

\Office{Sabbato post Cineres}{Saturday after Ash Wednesday}{Feria major}
\addcontentsline{toc}{subsection}{Sabbato post Cineres\enspace\textit{Saturday after Ash Wednesday}}
\begin{paracol}{2}
\LA
\LessonHead{Lectio I}
\switchcolumn
\EN
\LessonHead{Lesson I}
\switchcolumn*

\LA
\LessonTitle{Léctio sancti Evangélii secúndum Marcum}
\switchcolumn
\EN
\LessonTitle{Continuation of the Holy Gospel according to Mark}
\switchcolumn*

\LA
\Cite{Marc 6:47-56}
\switchcolumn
\EN
\Cite{Mark 6:47-56}
\switchcolumn*

\LA
\noindent In illo témpore: Cum sero esset, erat navis in médio mari, et Jesus solus in terra. Et réliqua.\par
\switchcolumn
\EN
\noindent In that time when it was late, the ship was in the midst of the sea, and Jesus himself alone on the land. And so on.\par
\switchcolumn*

\LA
\LessonTitle{Homilía sancti venerábilis Bedæ Presbýteri}
\switchcolumn
\EN
\LessonTitle{Homily by the Venerable Bede, Priest (at Jarrow) and Doctor of the Church}
\switchcolumn*

\LA
\Source{Lib. 2, cap. 28 in cap. 6 Marci, tom. 4}
\switchcolumn
\EN
\Source{Bk. ii, cap. 6, on Mark vi, 45}
\switchcolumn*

\LA
\noindent Labor discipulórum in remigándo, et contrárius eis ventus, labóres sanctæ Ecclésiæ varios desígnat: quæ inter undas sǽculi adversántis, et immundórum flatus spirítuum, ad quiétem pátriæ cæléstis, quasi ad fidam líttoris statiónem, perveníre conátur. Ubi bene dícitur, quia navis erat in médio mari, et ipse solus in terra: quia nonnúnquam Ecclésia tantis gentílium pressúris non solum afflícta, sed et fœdáta est, ut, si fíeri posset, Redémptor ipsíus eam prorsus deseruísse ad tempus viderétur.\par
\switchcolumn
\EN
\noindent The toil of the disciples in rowing, and the wind contrary to them, is a figure of the diverse toils of the Holy Church, as, amid the waves of a world that fighteth against her, and the stormy blasts of unclean spirits, she laboureth to reach the rest of her Fatherland above, as a shore safe for her anchor. Here also it is well said that the ship was in the midst of the sea, and He alone on the land; for sometimes it cometh to pass that the Church is, by the great pressure of the Gentiles, not only so afflicted, but also befouled, that it seemeth as though, if it were possible, her Redeemer had for the time forsaken her.\par
\switchcolumn*

\LA
\begin{Resp}\Rsign Veni hódie ad fontem aquæ, et orávi Dóminum, dicens: \Star{} Dómine, Deus Abraham, tu prósperum fecísti desidérium meum.\end{Resp}
\begin{Resp}\Vsign Igitur puélla, cui díxero, Da mihi aquam de hýdria tua, ut bibam: et illa díxerit, Bibe, dómine, et camélis tuis potum tríbuam: ipsa est, quam præparávit Dóminus fílio dómini mei.\end{Resp}
\begin{Resp}\Rsign Dómine, Deus Abraham, tu prósperum fecísti desidérium meum.\end{Resp}
\switchcolumn
\EN
\begin{Resp}\Rsign I came this day unto the well, and I besought the Lord, and said: \Star{} O Lord God of Abraham, Thou hast prospered my way.\end{Resp}
\begin{Resp}\Vsign Therefore the virgin to whom I shall say: Give me water of thy pitcher to drink; and she shall say to me: Drink, my lord, and I will give thy camels drink also; let the same be the woman whom the Lord hath appointed out for my master's son.\end{Resp}
\begin{Resp}\Rsign O Lord God of Abraham, Thou hast prospered my way.\end{Resp}
\switchcolumn*

\LA
\LessonHead{Lectio II}
\switchcolumn
\EN
\LessonHead{Lesson II}
\switchcolumn*

\LA
\noindent Unde est illa vox ejus inter undas procellásque tentatiónum irruéntium deprehénsæ, atque auxílium protectiónis illíus gemebúndo clamóre quæréntis: Ut quid, Dómine, recessísti longe, déspicis in opportunitátibus, in tribulatióne? Quæ páriter vocem inimíci persequéntibus expónit, in sequéntibus Psalmi subíciens: Dixit enim in corde suo: Oblítus est Deus, avértit fáciem suam, ne vídeat usque in finem.\par
\switchcolumn
\EN
\noindent Whence it is that there cometh that cry of hers, when she is taken amid the waves, and the winds of temptations that break upon her, and with piteous entreaty she calleth on Him to protect her: Why standest Thou afar off, O Lord, why hidest Thou thyself in times of trouble? (Ps. ix. 22.) And then, in the verses that follow, she telleth Him what saith the enemy that persecuted her, saying: For he hath said in his heart God hath forgotten; He hideth His face, He will never see it. (32.)\par
\switchcolumn*

\LA
\begin{Resp}\Rsign Factus est sermo Dómini ad Abram, dicens: \Star{} Noli timére, Abram: ego protéctor tuus sum, et merces tua magna nimis.\end{Resp}
\begin{Resp}\Vsign Ego enim sum Dóminus Deus tuus, qui edúxi te de Ur Chaldæórum.\end{Resp}
\begin{Resp}\Rsign Noli timére, Abram: ego protéctor tuus sum, et merces tua magna nimis.\end{Resp}
\switchcolumn
\EN
\begin{Resp}\Rsign The word of the Lord came unto Abram, saying: \Star{} Fear not, Abram I am thy shield, and thy exceeding great reward.\end{Resp}
\begin{Resp}\Vsign For I am the Lord thy God That brought thee out of Ur of the Chaldees.\end{Resp}
\begin{Resp}\Rsign Fear not, Abram I am thy shield, and thy exceeding great reward.\end{Resp}
\switchcolumn*

\LA
\LessonHead{Lectio III}
\switchcolumn
\EN
\LessonHead{Lesson III}
\switchcolumn*

\LA
\noindent Verum ille non oblivíscitur oratiónem páuperum, neque avértit fáciem suam a sperántibus in se: quin pótius et certántes cum hóstibus, ut vincant, ádjuvat, et victóres in ætérnum corónat. Unde hic quoque apérte dícitur, quia vidit eos laborántes in remigándo. Videt quippe Dóminus laborántes in mari, quamvis ipse pósitus in terra: quia etsi ad horam différre videátur auxílium tribulátis impéndere, nihilóminus eos, ne in tribulatiónibus defíciant, suæ respéctu pietátis corróborat: et aliquándo étiam manifésto adjutório, victis adversitátibus, quasi calcátis sedatísque flúctuum volumínibus líberat.\par
\switchcolumn
\EN
\noindent Verily, He forgetteth not the prayer of the poor, neither turneth He His face away from any that putteth his trust in Him; yea, rather, to him whosoever is striving with the enemy, He giveth help to conquer, and, whosoever conquereth, to him He giveth an everlasting crown. For the which reason also it is here said plainly He saw them toiling in rowing. The Lord seeth them that are toiling in the sea, albeit He be Himself on the land. Although He seem for a moment to tarry in succouring the distressed, nevertheless the look of His love is strengthening them, all the while, lest they should faint and sometimes He setteth them free, even by an open deliverance, conquering all their adversaries for them, as when He walked upon the swelling of the waves, and stilled them.\par
\switchcolumn*

\LA
\begin{Resp}\Rsign Movens Abram tabernáculum suum, venit et habitávit juxta convállem Mambre: \Star{} Ædificavítque ibi altáre Dómino.\end{Resp}
\begin{Resp}\Vsign Dixit autem Dóminus ad eum: Leva óculos tuos, et vide: omnem terram, quam cónspicis tibi dabo, et sémini tuo in sempitérnum.\end{Resp}
\begin{Resp}\Rsign Ædificavítque ibi altáre Dómino.\end{Resp}
\begin{Resp}\Vsign Glória Patri, et Fílio, \Star{} et Spirítui Sancto.\end{Resp}
\begin{Resp}\Rsign Ædificavítque ibi altáre Dómino.\end{Resp}
\switchcolumn
\EN
\begin{Resp}\Rsign Abram removed his tent, and came, and dwelt by the vale of Mamre \Star{} And built there an altar unto the Lord.\end{Resp}
\begin{Resp}\Vsign And the Lord said unto him: Lift up thine eyes, and look; all the land which thou seest, to thee will I give it, and to thy seed for ever.\end{Resp}
\begin{Resp}\Rsign And built there an altar unto the Lord.\end{Resp}
\begin{Resp}\Vsign Glory be to the Father, and to the Son, \Star{} and to the Holy Ghost.\end{Resp}
\begin{Resp}\Rsign And built there an altar unto the Lord.\end{Resp}
\switchcolumn*

\end{paracol}

\PartTitle{Proprium de Sanctis}{Proper of Saints}

\addcontentsline{toc}{chapter}{Proprium de Sanctis\enspace\textit{Proper of Saints}}

\SeasonTitle{Mensis November}{November}

\addcontentsline{toc}{section}{Mensis November\enspace\textit{November}}

\par\needspace{10\baselineskip}\addvspace{1.2em}{\centering\small\textsc{Die 26 Novembris} · \textsc{26 November}\par}
\Office{S. Silvestri Abbatis}{St. Sylvester, Abbot}{Duplex}
\addcontentsline{toc}{subsection}{Die 26 Novembris · S. Silvestri Abbatis\enspace\textit{St. Sylvester, Abbot}}
\begin{paracol}{2}
\LA
\RefLine{Lectiones I–III: de Scriptura occurrente.}
\switchcolumn
\EN
\RefLine{Lessons I–III: from the Scripture of the day.}
\switchcolumn*

\LA
\RefLine{Lectio IX: de S. Petri Alexandrini Martyris (11-26t).}
\switchcolumn
\EN
\RefLine{Lesson IX: of St. Peter of Alexandria, Martyr (11-26t).}
\switchcolumn*

\LA
\LessonHead{Lectio IV}
\switchcolumn
\EN
\LessonHead{Lesson IV}
\switchcolumn*

\LA
\noindent Silvester, Auximi in Piceno nobili genere ortus, statim puerílem ætátem litteris ac bonis móribus mirifice exornavit. Adoléscens, Bonóniam ad stúdia jurisprudéntiæ missus a patre, cum sacris litteris, a Deo mónitus, dedísset operam, parentis incurrit indignatiónem; quam æquo animo toto decennio pértulit. Ob egregiam ejus virtútem a canonicis cathedralis Auximanæ ecclésiæ socius honoris eléctus est; in quo munere pópulo oratiónibus, exemplo et conciónibus opem tulit.\par
\switchcolumn
\EN
\noindent This Silvester was born of a noble family at Osimo, in Picenum, and in his childhood was a wonderful example both in regard to letters and good living. When he grew older his father sent him to Bologna to study the law, but God warned him to give himself to divinity, and he thereby incurred the wrath of his father, which he bore with complacency for ten full years. On account of his eminent graces he was elected an honorary canon of the Cathedral of Osimo, in the which dignity he ministered to the people by his prayers, his example, and his sermons.\par
\switchcolumn*

\LA
\begin{Resp}\Rsign Honéstum fecit illum Dóminus, et custodívit eum ab inimícis, et a seductóribus tutávit illum: \Star{} Et dedit illi claritátem ætérnam.\end{Resp}
\begin{Resp}\Vsign Justum dedúxit Dóminus per vias rectas, et osténdit illi regnum Dei.\end{Resp}
\begin{Resp}\Rsign Et dedit illi claritátem ætérnam.\end{Resp}
\switchcolumn
\EN
\begin{Resp}\Rsign The Lord made him honourable, and defended him from his enemies, and kept him safe from those that lay in wait for him; \Star{} And gave him perpetual glory.\end{Resp}
\begin{Resp}\Vsign He went down with him into the pit, and left him not in bonds.\end{Resp}
\begin{Resp}\Rsign And gave him perpetual glory.\end{Resp}
\switchcolumn*

\LA
\LessonHead{Lectio V}
\switchcolumn
\EN
\LessonHead{Lesson V}
\switchcolumn*

\LA
\noindent Inter funus nobilis cujusdam defuncti, in aperto tumulo formosi viri suique propinqui deforme cadáver conspíciens: Ego, inquit, sum, quod hic fuit; quod hic est, ego ero. Et mox, peracto funere, illa sibi Dómini occurrénte senténtia: Qui vult veníre post me, ábneget semetípsum, et tollat crucem suam, et sequátur me; in solitúdinem, majoris perfectiónis studio, secessit, ibique vigiliis, oratiónibus jejuniisque deditus, crudas tantum herbas in cibum sæpius adhibuit. Ut autem magis latéret hómines, varias mutavit sedes; ac demum pervenit ad montem Fanum, locum, quamvis prope Fabrianum, eo tamen témpore desértum, ibique in honórem sanctíssimi patris Benedicti templum eréxit, congregatiónisque Silvestrinórum fundaménta jecit, sub regula et habitu in visióne sibi ab eodem Sancto osténsis.\par
\switchcolumn
\EN
\noindent At the funeral of a certain nobleman he perceived in an open grave the disfigured corpse of a kinsman of his own who had been very comely in his lifetime, and he said to himself, I am what he was, and what he is I shall be. Straightway after the funeral he read the words of the Lord, If any man will come after Me let him deny himself and take up his cross and follow Me (Matth. xvi. 24.) Thereupon he withdrew into the desert to seek after greater perfection, and then gave himself up to watching, praying, and fasting, very often taking no food but uncooked herbs. In order, however, to cut himself off the more from men, he moved from one place to another, and at length came to Mount Fano, which is hard by Fabriano, but was itself then absolutely uninhabited. Then he built a church in honour of the holy Father Benedict, and founded the congregation of Silvestrians, with a rule and dress which were revealed to him in a vision by the holy Patriarch himself.\par
\switchcolumn*

\LA
\begin{Resp}\Rsign Amávit eum Dóminus, et ornávit eum: stolam glóriæ índuit eum, \Star{} Et ad portas paradísi coronávit eum.\end{Resp}
\begin{Resp}\Vsign Índuit eum Dóminus lorícam fídei, et ornávit eum.\end{Resp}
\begin{Resp}\Rsign Et ad portas paradísi coronávit eum.\end{Resp}
\switchcolumn
\EN
\begin{Resp}\Rsign The Lord loved him and beautified him; He clothed him with a robe of glory, \Star{} And crowned him at the gates of Paradise.\end{Resp}
\begin{Resp}\Vsign The Lord hath put on him the breast-plate of faith, and hath adorned him.\end{Resp}
\begin{Resp}\Rsign And crowned him at the gates of Paradise.\end{Resp}
\switchcolumn*

\LA
\LessonHead{Lectio VI}
\switchcolumn
\EN
\LessonHead{Lesson VI}
\switchcolumn*

\LA
\noindent At ínvidens sátanas variis terróribus illíus monachos turbare nitebátur, noctu monasterii jánuas hostíliter invadens. Sed vir Dei, hostis impetum ita repressit, ut monachi in sancto instituto magis confirmaréntur ac patris sanctitátem agnoscerent. Spíritu prophetíæ aliisque donis enituit. Quæ ut semper profúnda humilitate conservavit, ita contra se dæmonis invidiam concitavit; a quo præceps actus per scalas oratórii, et prope interimendus, præsentíssimo Vírginis beneficio incolumitáti redditus est. Quod benefícium perpetua et singulari in illam pietáte commendavit ad ultimum usque vitæ spíritum, quem, fere nonagenarius, sanctitáte et miraculis clarus, Deo réddidit anno salútis millesimo ducentésimo sexagesimo septimo, sexto Kalendas Decembris. Ejus Offícium ac Missam Leo décimus tertius Pontifex maximus ad universam exténdit Ecclésiam.\par
\switchcolumn
\EN
\noindent Satan envied him, strove to trouble his monks by diverse terrors, and made an hostile attack by night upon the gates of his monastery, but the man of God so overcame the assault of the enemy that his monks were the more confirmed in their Institute and recognised the holiness of their father. He shone with the spirit of prophecy and other gifts. These things he always preserved by the deepest lowliness, whereby he so stirred up against him the ill-will of the devil that that evil spirit cast him headlong down the stairs of his oratory, and went near to slay him, but he was restored to soundness by the helpful gift of the Virgin. This help he remembered with an unceasing and singular love toward her until the last breath of his life, the which breath he resigned to God, famous for holiness and miracles, aged almost ninety years, upon the 26th day of November, in the year of salvation 1267. The Supreme Pontiff Leo XIII extended his Office and Mass to the whole Church.\par
\switchcolumn*

\LA
\begin{Resp}\Rsign Iste homo perfécit ómnia quæ locútus est ei Deus, et dixit ad eum: Ingrédere in réquiem meam: \Star{} Quia te vidi justum coram me ex ómnibus géntibus.\end{Resp}
\begin{Resp}\Vsign Iste est, qui contémpsit vitam mundi, et pervénit ad cæléstia regna.\end{Resp}
\begin{Resp}\Rsign Quia te vidi justum coram me ex ómnibus géntibus.\end{Resp}
\begin{Resp}\Vsign Glória Patri, et Fílio, \Star{} et Spirítui Sancto.\end{Resp}
\begin{Resp}\Rsign Quia te vidi justum coram me ex ómnibus géntibus.\end{Resp}
\switchcolumn
\EN
\begin{Resp}\Rsign This is he which did according unto all that God commanded him; and God said unto him: Enter thou into My rest, \Star{} For thee have I seen righteous before Me among all people.\end{Resp}
\begin{Resp}\Vsign This is he which loved not his life in this world, and is come unto an everlasting kingdom.\end{Resp}
\begin{Resp}\Rsign For thee have I seen righteous before Me among all people.\end{Resp}
\begin{Resp}\Vsign Glory be to the Father, and to the Son, \Star{} and to the Holy Ghost.\end{Resp}
\begin{Resp}\Rsign For thee have I seen righteous before Me among all people.\end{Resp}
\switchcolumn*

\LA
\LessonHead{Lectio VII}
\switchcolumn
\EN
\LessonHead{Lesson VII}
\switchcolumn*

\LA
\LessonTitle{Léctio sancti Evangélii secúndum Matthǽum.}
\switchcolumn
\EN
\LessonTitle{From the Holy Gospel according to Matthew}
\switchcolumn*

\LA
\Cite{Matt 19:27-29}
\switchcolumn
\EN
\Cite{Matt 19:27-29}
\switchcolumn*

\LA
\noindent In illo témpore: Dixit Petrus ad Jesum: Ecce nos relíquimus ómnia, et secúti sumus te: quid ergo erit nobis? Et réliqua.\par
\switchcolumn
\EN
\noindent At that time, Peter said unto Jesus: Behold, we have forsaken all, and followed thee what shall we have therefore? And so on.\par
\switchcolumn*

\LA
\LessonTitle{Homilía sancti Hierónymi Presbýteri.}
\switchcolumn
\EN
\LessonTitle{Homily by St. Jerome, Priest (at Bethlehem.)}
\switchcolumn*

\LA
\Source{Lib. 3. in Matth. cap. 19.}
\switchcolumn
\EN
\Source{Bk. iii. on Matth xix.}
\switchcolumn*

\LA
\noindent Grandis fidúcia! Petrus piscátor erat, dives non fúerat, cibos manu et arte quærébat: et tamen lóquitur confidénter: Relíquimus ómnia. Et quia non súfficit tantum relínquere, jungit quod perféctum est: Et secúti sumus te. Fécimus quod jussísti: quid ígitur nobis dabis prǽmii? Jesus autem dixit illis: Amen dico vobis, quod vos qui secúti estis me, in regeneratióne, cum séderit Fílius hóminis in sede majestátis suæ, sedébitis et vos super sedes duódecim, judicántes duódecim tribus Ísraël. Non dixit: Qui reliquístis ómnia; hoc enim et Crates fecit philósophus, et multi álii divítias contempsérunt: sed, Qui secúti estis me: quod próprie Apostolórum est, atque credéntium.\par
\switchcolumn
\EN
\noindent Peter was a fisherman, he was not rich, he earned his bread by his hand and skill, and nevertheless he is thus bold, and saith confidently: We have forsaken all. And because it sufficeth not to forsake only, he addeth that which to do is to be perfect: and followed thee. We have done that which Thou hast commanded us, what reward therefore wilt Thou give us? And Jesus said unto them Amen I say unto you, that ye which have followed Me, in the regeneration, when the Son of Man shall sit in the throne of His glory, ye also shall sit upon twelve thrones, judging the twelve tribes of Israel. He said not, Ye which have forsaken all, for this did even Crates the philosopher, and they which have set nothing by riches are many, but, Ye which have followed Me. This did the Apostles, and this do believers do.\par
\switchcolumn*

\LA
\begin{Resp}\Rsign Iste est qui ante Deum magnas virtútes operátus est, et de omni corde suo laudávit Dóminum: \Star{} Ipse intercédat pro peccátis ómnium populórum.\end{Resp}
\begin{Resp}\Vsign Ecce homo sine queréla, verus Dei cultor, ábstinens se ab omni ópere malo, et pérmanens in innocéntia sua.\end{Resp}
\begin{Resp}\Rsign Ipse intercédat pro peccátis ómnium populórum.\end{Resp}
\switchcolumn
\EN
\begin{Resp}\Rsign This is he which wrought great wonders before God, and praised the Lord with all his heart. \Star{} May he pray for all people, that their sins may be forgiven unto them.\end{Resp}
\begin{Resp}\Vsign Behold a man without blame, a worshipper of God in truth, keeping himself clean from every evil work, and abiding still in his innocency.\end{Resp}
\begin{Resp}\Rsign May he pray for all people, that their sins may be forgiven unto them.\end{Resp}
\switchcolumn*

\LA
\LessonHead{Lectio VIII}
\switchcolumn
\EN
\LessonHead{Lesson VIII}
\switchcolumn*

\LA
\noindent In regeneratióne, cum séderit Fílius hóminis in sede majestátis suæ (quando et mórtui de corruptióne resúrgent incorrúpti) sedébitis et vos in sóliis judicántium, condemnántes duódecim tribus Ísraël: quia vobis credéntibus, illi crédere noluérunt. Et omnis qui relíquerit domum, vel fratres, aut soróres, aut patrem, aut matrem, aut uxórem, aut fílios, aut agros propter nomen meum, céntuplum accípiet, et vitam ætérnam possidébit. Locus iste cum illa senténtia cóngruit, in qua Salvátor lóquitur: Non veni pacem míttere, sed gládium. Veni enim separáre hóminem a patre suo, et matrem a fília, et nurum a socru: et inimíci hóminis doméstici ejus. Qui ergo propter fidem Christi, et prædicatiónem Evangélii, omnes afféctus contémpserint, atque divítias et sǽculi voluptátes: isti céntuplum recípient, et vitam ætérnam possidébunt.\par
\switchcolumn
\EN
\noindent In the regeneration, when the Son of Man shall sit in the throne of His glory, and when the dead shall rise again from corruption incorruptible, (i Cor. xv. 53,) ye also shall sit upon twelve thrones of judgment, condemning the twelve tribes of Israel, because, when ye believed in Me, they would not. (John iii. 18.) And every one that hath forsaken houses, or brethren, or sisters, or father, or mother, or wife, or children, or lands, for My Name's sake, shall receive an hundredfold, and shall inherit everlasting life. This place agreeth well with that other where the Saviour saith I came not to send peace, but a sword. For I am come to set a man at variance against his father, and the daughter against her mother, and the daughter-in-law against her mother-in-law; and a man's foes shall be they of his own household. (Matth. x. 34.) Every one, therefore, that hath set no store by affection, and riches, and the pleasures of the world, for Christ's faith's sake, and the preaching of the Gospel, shall receive an hundred-fold, and shall inherit everlasting life.\par
\switchcolumn*

\LA
\begin{Resp}\Rsign Sint lumbi vestri præcíncti, et lucérnæ ardéntes in mánibus vestris: \Star{} Et vos símiles homínibus exspectántibus dóminum suum, quando revertátur a núptiis.\end{Resp}
\begin{Resp}\Vsign Vigiláte ergo, quia nescítis qua hora Dóminus vester ventúrus sit.\end{Resp}
\begin{Resp}\Rsign Et vos símiles homínibus exspectántibus dóminum suum, quando revertátur a núptiis.\end{Resp}
\begin{Resp}\Vsign Glória Patri, et Fílio, \Star{} et Spirítui Sancto.\end{Resp}
\begin{Resp}\Rsign Et vos símiles homínibus exspectántibus dóminum suum, quando revertátur a núptiis.\end{Resp}
\switchcolumn
\EN
\begin{Resp}\Rsign Let your loins be girded about, and your lights burning; \Star{} And ye yourselves like unto men that wait for their lord, when he will return from the wedding.\end{Resp}
\begin{Resp}\Vsign Watch therefore, for ye know not what hour your Lord doth come.\end{Resp}
\begin{Resp}\Rsign And ye yourselves like unto men that wait for their lord, when he will return from the wedding.\end{Resp}
\begin{Resp}\Vsign Glory be to the Father, and to the Son, \Star{} and to the Holy Ghost.\end{Resp}
\begin{Resp}\Rsign And ye yourselves like unto men that wait for their lord, when he will return from the wedding.\end{Resp}
\switchcolumn*

\end{paracol}

\par\needspace{10\baselineskip}\addvspace{1.2em}{\centering\small\textsc{Die 26 Novembris} · \textsc{26 November}\par}
\Office{S. Petri Alexandrini Martyris}{St. Peter of Alexandria, Martyr}{Simplex}
\addcontentsline{toc}{subsection}{Die 26 Novembris · S. Petri Alexandrini Martyris\enspace\textit{St. Peter of Alexandria, Martyr}}
\begin{paracol}{2}
\LA
\LessonHead{Lectio IX (pro commemoratione)}
\switchcolumn
\EN
\LessonHead{Lesson IX (when commemorated)}
\switchcolumn*

\LA
\LessonTitle{Commemoratio: S. Petri Alexandrini Martyris}
\switchcolumn
\EN
\LessonTitle{Commemoration of: St. Peter of Alexandria, Martyr}
\switchcolumn*

\LA
\noindent Petrus, epíscopus Alexandríæ, post Theonam virum sanctíssimum, sanctitátis et doctrinæ splendore non solum illustravit Ægyptum, sed toti luxit Ecclésiæ Dei. Qui in persecutióne Maximiáni Galerii illam témporum acerbitátem ita pértulit, ut multi, admirábilem ejus patiéntiam intuéntes, plurimum in christiana virtúte proficerent. Is primus Aríum diaconum Alexandrinum, propter schisma Meletianum cui favebat, a fidelium communióne sejunxit. Ad eum, cápitis ab eodem Maximiáno damnátum, in carcere cum Achillas et Alexander presbyteri deprecatores Aríi venissent, respóndit, noctu apparuísse sibi Jesum veste discissa, causamque rei sciscitanti dixisse: Aríus vestem meam, quæ est Ecclésia, dilaceravit. Quibus étiam prædícens fore, ut sibi in episcopatu succéderent, præcepit ne umquam Aríum in communiónem recíperent, quem Deo mortuum esse sciret. Et hanc divinam prænotiónem veram fuísse, non diu post rei probavit eventus. Denique, duodecimo sui episcopátus anno, sexto Kalendas Decembris, abscisso cápite, ad martyrii corónam evolávit.\par
\switchcolumn
\EN
\noindent This Peter succeeded that eminent Saint, Theonas, as Pope of Alexandria, (in the year of our Lord 300,) and the glory of his holiness and teaching hath enlightened not Egypt only, but the whole Church of God. The wondrous patience wherewith he bore the roughness of the times in the persecution under Maximian Galerius caused many greatly to increase in Christian graces. He was the first who cut off Arius, then a Deacon of Alexandria, from the Communion of the faithful, on account of his leaning to the Meletian schism. He was condemned to death by Maximian, and was in prison when there came to him the two Priests Achilles and Alexander to plead for Arius, but Peter told them that Jesus had appeared to him in the night clad in a rent garment, and when he asked what was thereby signified, had said unto him Arius hath torn My vesture, which is the Church. Also, he foretold to them that they should be Popes of Alexandria after him, and strictly commanded them never to receive Arius into Communion, because he knew him to be dead in the sight of God. That this was a true prophecy the event did shortly prove. At length, in the twelfth year of his papacy, upon the 26th day of November, (in the year of salvation 311,) his head was cut off, and he went hence to receive the crown of his testimony.\par
\switchcolumn*

\LA
\TeDeum{Te Deum}
\switchcolumn
\EN
\TeDeum{Te Deum}
\switchcolumn*

\end{paracol}

\par\needspace{10\baselineskip}\addvspace{1.2em}{\centering\small\textsc{Die 29 Novembris} · \textsc{29 November}\par}
\Office{In Vigilia S. Andreæ Apostoli}{Vigil of St. Andrew the Apostle}{Vigilia}
\addcontentsline{toc}{subsection}{Die 29 Novembris · In Vigilia S. Andreæ Apostoli\enspace\textit{Vigil of St. Andrew the Apostle}}
\begin{paracol}{2}
\LA
\LessonHead{Lectio I}
\switchcolumn
\EN
\LessonHead{Lesson I}
\switchcolumn*

\LA
\LessonTitle{Léctio sancti Evangélii secúndum Joánnem}
\switchcolumn
\EN
\LessonTitle{Continuation of the Holy Gospel according to John}
\switchcolumn*

\LA
\Cite{Joannes 1:35-51}
\switchcolumn
\EN
\Cite{John 1:35-51}
\switchcolumn*

\LA
\noindent In illo témpore: Stabat Joannes, et ex discipulis ejus duo. Et respiciens Jesum ambulantem, dicit: Ecce Agnus Dei. Et réliqua.\par
\switchcolumn
\EN
\noindent At that time again John stood, and two of his disciples. And beholding Jesus walking, he saith: Behold the Lamb of God. And so on.\par
\switchcolumn*

\LA
\LessonTitle{Homilia sancti Augustíni Epíscopi}
\switchcolumn
\EN
\LessonTitle{Homily of St. Augustine, bishop of Hippo}
\switchcolumn*

\LA
\Source{Tract. 7 in Joann., post init.}
\switchcolumn
\EN
\Source{7th Tract on John}
\switchcolumn*

\LA
\noindent Quia talis erat Joannes amicus sponsi, non quærebat gloriam suam, sed testimonium perhibebat veritati: numquid voluit apud se remanere discipulos suos, ut non sequerentur Dominum? Magis ipse ostendit discipulis suis quem sequerentur: habebant enim illum tamquam agnum. Et ille: Quid me attenditis: Ego non sum agnus. Ecce Agnus Dei. De quo et superius dixerat: Ecce Agnus Dei. Et quid nobis prodest Agnus Dei? Ecce, ait, qui tollit peccatum mundi. Secuti sunt illum, hoc audito, duo qui erant cum Joanne.\par
\switchcolumn
\EN
\noindent Since John was the friend of the Bridegroom, (iii. 29,) he sought not his own glory, but bare witness to the truth. Would he that his disciples should remain with him rather than that they should follow the Lord? Nay, he showed his disciples Whom they should follow. They thought that he himself was the Lamb; but he saith Why wait ye on me? I am not the Lamb. Behold the Lamb of God! This was He of Whom he had already said above (29) Behold the Lamb of God! And what use to us is the Lamb of God? Behold the Lamb of God, saith John, Which taketh away the sin of the world. And the two disciples heard him speak, and they followed Jesus.\par
\switchcolumn*

\LA
\begin{Resp}\Rsign Laudábilis pópulus, \Star{} Quem Dóminus exercítuum benedíxit dicens: Opus mánuum meárum tu es, heréditas mea Israël.\end{Resp}
\begin{Resp}\Vsign Beáta gens, cujus est Dóminus Deus, pópulus eléctus in hereditátem.\end{Resp}
\begin{Resp}\Rsign Quem Dóminus exercítuum benedíxit dicens: Opus mánuum meárum tu es, heréditas mea Israël.\end{Resp}
\switchcolumn
\EN
\begin{Resp}\Rsign Blessed is the people \Star{} Whom the Lord of hosts hath blessed, saying O Israel thou art the work of Mine own hands, thou art Mine own inheritance.\end{Resp}
\begin{Resp}\Vsign Blessed is the nation whose God is the Lord, and the people whom He hath chosen for His own inheritance.\end{Resp}
\begin{Resp}\Rsign Whom the Lord of hosts hath blessed, saying O Israel thou art the work of Mine own hands, thou art Mine own inheritance.\end{Resp}
\switchcolumn*

\LA
\LessonHead{Lectio II}
\switchcolumn
\EN
\LessonHead{Lesson II}
\switchcolumn*

\LA
\noindent Videamus sequentia. Ecce Agnus Dei. Hoc Joannes. Et audierunt eum duo discipuli loquentem, et secuti sunt Jesum. Non sic illum sequebantur, quasi jam ut inhærerent illi: nam manifestum est, quando illi inhæserunt, quia de navi eos vocavit. In his enim duobus erat Andreas, sicut modo audistis. Andreas autem frater Petri erat. Et novimus in Evangelio quod Petrum et Andream Dominus de navi vocavit, dicens: Venite post me, et faciam vos piscatores hominum. Et ex illo jam inhæserunt illi, ut non recederent.\par
\switchcolumn
\EN
\noindent Let us see what followed. John stood, and two of his disciples; and looking on Jesus as He walked, he saith Behold the Lamb of God! And the two disciples heard him speak, and they followed Jesus. They followed Him, not yet to cleave unto Him, for it is manifest that they clave unto Him only after that He had called them out of the ship. (Matth. iv. 18.) One of the two which heard John speak, and followed Him, was Andrew, Simon Peter's brother. And we know how it is written in the Gospel of Matthew: Jesus walking by the sea of Galilee, saw two brethren, Simon called Peter, and Andrew his brother, casting a net into the sea, for they were fishers. And He saith unto them Follow Me, and I will make you fishers of men. And they straightway left their nets and followed Him.\par
\switchcolumn*

\LA
\begin{Resp}\Rsign Angústiæ mihi sunt úndique, et quid éligam ignóro; \Star{} Mélius est mihi incídere in manus hóminum, quam derelínquere legem Dei mei.\end{Resp}
\begin{Resp}\Vsign Si enim hoc égero, mors mihi est; si autem non égero, non effúgiam manus vestras.\end{Resp}
\begin{Resp}\Rsign Mélius est mihi incídere in manus hóminum, quam derelínquere legem Dei mei.\end{Resp}
\switchcolumn
\EN
\begin{Resp}\Rsign I am straitened on every side, and know not what to choose. \Star{} It is better for me to fall into the hands of men, than to sin against the law of my God.\end{Resp}
\begin{Resp}\Vsign For if I do this thing, it is death unto me and if I do it not, I cannot escape your hands.\end{Resp}
\begin{Resp}\Rsign It is better for me to fall into the hands of men, than to sin against the law of my God.\end{Resp}
\switchcolumn*

\LA
\LessonHead{Lectio III}
\switchcolumn
\EN
\LessonHead{Lesson III}
\switchcolumn*

\LA
\noindent Modo ergo quod illum sequuntur isti duo, non quasi non recessuri sequuntur, sed videre voluerunt ubi habitaret, et facere quod scriptum est: Limen ostiorum ejus exterat pes tuus: surge ad illum venire assidue, et erudire præceptis ejus. Ostendit eis ille ubi maneret: venerunt, et fuerunt cum illo. Quam beatum diem duxerunt, quam beatam noctem! Quis est, qui nobis dicat quæ audierint illi a Domino? Aedificemus et nosmetipsi in corde nostro, et faciamus domum, quo veniat ille, et doceat nos, et colloquatur nobis.\par
\switchcolumn
\EN
\noindent From that time, therefore, was it that they clave unto Him continuously. Then Jesus turned, and saw them following, and saith unto them What seek ye? They said unto Him Rabbi, (which is to say, being interpreted, Master,) where dwellest Thou? So they follow Him now, not as to cleave unto Him for ever, but as to know where He dwelt, and to obey that which is written If thou see a man of understanding, go to him early in the morning, and let thy foot wear the steps of his doors. (Ecclus. vi. 36.) He saith unto them Come and see. They came and saw where He dwelt, and abode with Him that day. O what a blessed day! O what a blessed night! for it was about the tenth hour. Who shall tell what they heard from the Lord? O let us also build an house in our hearts, where He may come, and teach us, and talk with us!\par
\switchcolumn*

\LA
\begin{Resp}\Rsign Misit Dóminus Angelum suum et conclúsit ora leónum, \Star{} Et non contaminavérunt, quia coram eo injustítia invénta non est in me.\end{Resp}
\begin{Resp}\Vsign Misit Deus misericórdiam suam et veritátem suam: ánimam meam erípuit de médio catulórum leónum.\end{Resp}
\begin{Resp}\Rsign Et non contaminavérunt, quia coram eo injustítia invénta non est in me.\end{Resp}
\begin{Resp}\Vsign Glória Patri, et Fílio, \Star{} et Spirítui Sancto.\end{Resp}
\begin{Resp}\Rsign Et non contaminavérunt: quia coram eo injustítia invénta non est in me.\end{Resp}
\switchcolumn
\EN
\begin{Resp}\Rsign The Lord hath sent His angel, and hath shut the lions' mouths, that \Star{} They have not hurt me forasmuch as before Him innocency was found in me.\end{Resp}
\begin{Resp}\Vsign God hath sent forth His mercy and His truth, (and delivered) my soul from among the lions' whelps.\end{Resp}
\begin{Resp}\Rsign They have not hurt me forasmuch as before Him innocency was found in me.\end{Resp}
\begin{Resp}\Vsign Glory be to the Father, and to the Son, \Star{} and to the Holy Ghost.\end{Resp}
\begin{Resp}\Rsign They have not hurt me forasmuch as before Him innocency was found in me.\end{Resp}
\switchcolumn*

\end{paracol}

\par\needspace{10\baselineskip}\addvspace{1.2em}{\centering\small\textsc{Die 30 Novembris} · \textsc{30 November}\par}
\Office{S. Andreæ Apostoli}{St. Andrew the Apostle}{Duplex II classis}
\addcontentsline{toc}{subsection}{Die 30 Novembris · S. Andreæ Apostoli\enspace\textit{St. Andrew the Apostle}}
\begin{paracol}{2}
\LA
\LessonHead{Lectio I}
\switchcolumn
\EN
\LessonHead{Lesson I}
\switchcolumn*

\LA
\LessonTitle{De Epístola beáti Pauli Apóstoli ad Romános}
\switchcolumn
\EN
\LessonTitle{Lesson from the letter of St. Paul the Apostle to the Romans}
\switchcolumn*

\LA
\Cite{Rom 10:4-9}
\switchcolumn
\EN
\Cite{Rom 10:4-9}
\switchcolumn*

\LA
\noindent Finis legis, Christus, ad justítiam omni credénti. Móyses enim scripsit, quóniam justítiam, quæ ex lege est, qui fécerit homo, vivet in ea. Quæ autem ex fide est justítia, sic dicit: Ne díxeris in corde tuo: Quis ascéndet in cælum? id est, Christum dedúcere: aut, Quis descéndet in abýssum? hoc est, Christum a mórtuis revocáre. Sed quid dicit Scriptúra? Prope est verbum in ore tuo, et in corde tuo: hoc est verbum fídei, quod prædicámus. Quia si confiteáris in ore tuo Dóminum Jesum, et in corde tuo credíderis quod Deus illum suscitávit a mórtuis, salvus eris.\par
\switchcolumn
\EN
\noindent For the end of the law is Christ, unto justice to every one that believeth. For Moses wrote, that the justice which is of the law, the man that shall do it, shall live by it. But the justice which is of faith, speaketh thus: Say not in thy heart, Who shall ascend into heaven? that is, to bring Christ down; Or who shall descend into the deep? that is, to bring up Christ again from the dead. But what saith the scripture? The word is nigh thee, even in thy mouth, and in thy heart. This is the word of faith, which we preach. For if thou confess with thy mouth the Lord Jesus, and believe in thy heart that God hath raised him up from the dead, thou shalt be saved.\par
\switchcolumn*

\LA
\begin{Resp}\Rsign Cum perambuláret Dóminus juxta mare Galilǽæ, vidit Petrum et Andréam rétia mitténtes in mare, et vocávit eos, dicens: \Star{} Veníte post me, fáciam vos fíeri piscatóres hóminum.\end{Resp}
\begin{Resp}\Vsign Erant enim piscatóres, et ait illis.\end{Resp}
\begin{Resp}\Rsign Veníte post me, fáciam vos fíeri piscatóres hóminum.\end{Resp}
\switchcolumn
\EN
\begin{Resp}\Rsign The Lord, walking by the Sea of Galilee, saw Peter and Andrew casting their nets into the sea, and He called them saying: \Star{} Follow Me, and I will make you fishers of men.\end{Resp}
\begin{Resp}\Vsign For they were fishers, and He saith unto them:\end{Resp}
\begin{Resp}\Rsign Follow Me, and I will make you fishers of men.\end{Resp}
\switchcolumn*

\LA
\LessonHead{Lectio II}
\switchcolumn
\EN
\LessonHead{Lesson II}
\switchcolumn*

\LA
\Cite{Rom 10:10-15}
\switchcolumn
\EN
\Cite{Rom 10:10-15}
\switchcolumn*

\LA
\noindent Corde enim créditur ad justítiam: ore autem conféssio fit ad salútem. Dicit enim Scriptúra: Omnis qui credit in illum, non confundétur. Non enim est distínctio Judǽi et Græci: nam idem Dóminus ómnium, dives in omnes qui ínvocant illum. Omnis enim quicúmque invocáverit nomen Dómini, salvus erit. Quómodo ergo invocábunt, in quem non credidérunt? aut quómodo credent ei, quem non audiérunt? quómodo autem áudient sine prædicánte? Quómodo vero prædicábunt nisi mittántur? sicut scriptum est: Quam speciósi pedes evangelizántium pacem, evangelizántium bona!\par
\switchcolumn
\EN
\noindent For, with the heart, we believe unto justice; but, with the mouth, confession is made unto salvation. For the scripture saith: Whosoever believeth in him, shall not be confounded. For there is no distinction of the Jew and the Greek: for the same is Lord over all, rich unto all that call upon him. For whosoever shall call upon the name of the Lord, shall be saved. How then shall they call on him, in whom they have not believed? Or how shall they believe him, of whom they have not heard? And how shall they hear, without a preacher? And how shall they preach unless they be sent, as it is written: How beautiful are the feet of them that preach the gospel of peace, of them that bring glad tidings of good things!\par
\switchcolumn*

\LA
\begin{Resp}\Rsign Mox ut vocem Dómini prædicántis audívit beátus Andréas, relíctis rétibus, quorum usu actúque vivébat, \Star{} Ætérnæ vitæ secútus est prǽmia largiéntem.\end{Resp}
\begin{Resp}\Vsign Hic est qui pro amóre Christi pepéndit in cruce, et pro lege ejus sustínuit passiónem.\end{Resp}
\begin{Resp}\Rsign Ætérnæ vitæ secútus est prǽmia largiéntem.\end{Resp}
\switchcolumn
\EN
\begin{Resp}\Rsign As soon as the blessed Andrew heard the voice of the Lord calling him, he left his nets, by the exercise and use whereof he lived; \Star{} And followed Him Who giveth life everlasting.\end{Resp}
\begin{Resp}\Vsign This is that disciple who for the love of Christ hung upon the cross, and suffered for the law of his God.\end{Resp}
\begin{Resp}\Rsign And followed Him Who giveth life everlasting.\end{Resp}
\switchcolumn*

\LA
\LessonHead{Lectio III}
\switchcolumn
\EN
\LessonHead{Lesson III}
\switchcolumn*

\LA
\Cite{Rom 10:16-21}
\switchcolumn
\EN
\Cite{Rom 10:16-21}
\switchcolumn*

\LA
\noindent Sed non omnes obédiunt Evangélio. Isaías enim dicit: Dómine, quis crédidit audítui nostro? Ergo fides ex audítu, audítus autem per verbum Christi. Sed dico: Numquid non audiérunt? Et quidem in omnem terram exívit sonus eórum, et in fines orbis terræ verba eórum. Sed dico: Numquid Israël non cognóvit? Primus Móyses dicit: Ego ad æmulatiónem vos addúcam in non gentem: in gentem insipiéntem, in iram vos mittam. Isaías autem audet, et dicit: Invéntus sum a non quæréntibus me: palam appárui iis, qui me non interrogábant. Ad Israël autem dicit: Tota die expándi manus meas ad pópulum non credéntem, et contradicéntem.\par
\switchcolumn
\EN
\noindent But all do not obey the gospel. For Isaias saith: Lord, who hath believed our report? Faith then cometh by hearing; and hearing by the word of Christ. But I say: Have they not heard? Yes, verily, their sound hath gone forth into all the earth, and their words unto the ends of the whole world. But I say: Hath not Israel known? First, Moses saith: I will provoke you to jealousy by that which is not a nation; by a foolish nation I will anger you. But Isaias is bold, and saith: I was found by them that did not seek me: I appeared openly to them that asked not after me. But to Israel he saith: All the day long have I spread my hands to a people that believeth not, and contradicteth me.\par
\switchcolumn*

\LA
\begin{Resp}\Rsign Doctor bonus et amícus Dei Andréas dúcitur ad crucem, quam a longe aspíciens dixit: Salve, crux, \Star{} Súscipe discípulum ejus, qui pepéndit in te magíster meus Christus.\end{Resp}
\begin{Resp}\Vsign Salve, crux, quæ in córpore Christi dedicáta es, et ex membris ejus tamquam margarítis ornáta.\end{Resp}
\begin{Resp}\Rsign Súscipe discípulum ejus, qui pepéndit in te magíster meus Christus.\end{Resp}
\begin{Resp}\Vsign Glória Patri, et Fílio, \Star{} et Spirítui Sancto.\end{Resp}
\begin{Resp}\Rsign Súscipe discípulum ejus, qui pepéndit in te magíster meus Christus.\end{Resp}
\switchcolumn
\EN
\begin{Resp}\Rsign Andrew the good teacher, the friend of God, was led to the cross, and when he saw it afar off, he said: God bless thee, O cross; \Star{} Welcome to the follower of Him That hung on thee, even my Master Christ.\end{Resp}
\begin{Resp}\Vsign God bless thee, O cross, thou art hallowed by the Body of Christ; His Members make thee goodly as with pearls.\end{Resp}
\begin{Resp}\Rsign Be welcome to the follower of Him That hung on thee, even my Master Christ.\end{Resp}
\begin{Resp}\Vsign Glory be to the Father, and to the Son, \Star{} and to the Holy Ghost.\end{Resp}
\begin{Resp}\Rsign Be welcome to the follower of Him That hung on thee, even my Master Christ.\end{Resp}
\switchcolumn*

\LA
\LessonHead{Lectio IV}
\switchcolumn
\EN
\LessonHead{Lesson IV}
\switchcolumn*

\LA
\noindent Andréas Apóstolus, Bethsáidæ natus, qui est Galilǽæ vicus, frater Petri, discípulus Joánnis Baptistæ, cum eum de Christo dicéntem audísset: Ecce Agnus Dei; secútus Jesum, fratrem quoque suum ad eúndem perdúxit. Cum póstea una cum fratre piscarétur in mari Galilǽæ, ambo a prætereúnte Christo Dómino ante álios Apóstolos vocáti illis verbis: Veníte post me, fáciam vos fíeri piscatóres hóminum; nullam interponéntes moram, et relíctis rétibus, secúti sunt eum. Post cujus passiónem et resurrectiónem, Andréas, cum in Scýthiam Európæ, quæ ei província ad Christi fidem disseminándam obtígerat, venísset, deínde Epírum ac Thráciam peragrásset; doctrína et miráculis innumerábiles hómines ad Christum convértit. Post, Patras Achájæ proféctus, et in ea urbe plúrimis ad veritátem evangélicam perdúctis, Ægéam procónsulem, prædicatióni evangélicæ resisténtem, libérrime increpávit, quod, qui judex hóminum habéri vellet, Christum Deum ómnium júdicem, a dæmónibus elúsus, non agnósceret.\par
\switchcolumn
\EN
\noindent The Apostle Andrew was born at Bethsaida, a town of Galilee, and was the brother of Peter. He was a disciple of John the Baptist, and heard him say of Christ, Behold the Lamb of God, (John i. 35-37, 40,) whereupon he immediately followed Jesus, bringing his brother also with him. Some while after, they were both fishing in the Sea of Galilee, and the Lord Christ, going by, called them both, before any other of the Apostles, in the words, Follow Me, and I will make you fishers of men. They made no delay, but left their nets, and followed Him. \textasciitilde{}(Matth. iv. 18-20.) After the death and Resurrection of Christ, Andrew was allotted Scythia as the province of his preaching, and, after labouring there, he went through Epirus and Thrace, where he turned vast multitudes to Christ by his teaching and miracles. Finally he went to Patras in Achaia, and there also he brought many to the knowledge of Gospel truth. Aegeas the Pro-consul resisted the preaching of the Gospel, and the Apostle freely rebuked him, bidding him know that while he held himself a judge of his fellow men, he was himself hindered by devils from knowing Christ our God, the Judge of all.\par
\switchcolumn*

\LA
\begin{Resp}\Rsign Homo Dei ducebátur ut crucifígerent eum: pópulus autem clamábat voce magna, dicens: \Star{} Innocens ejus sanguis sine causa damnátur.\end{Resp}
\begin{Resp}\Vsign Cumque dúcerent eum ut crucifigerétur, factus est concúrsus populórum clamántium et dicéntium.\end{Resp}
\begin{Resp}\Rsign Innocens ejus sanguis sine causa damnátur.\end{Resp}
\switchcolumn
\EN
\begin{Resp}\Rsign The man of God was led to be crucified, and the people cried with a loud voice, saying: \Star{} The innocent blood of this just person is condemned without a cause.\end{Resp}
\begin{Resp}\Vsign And when they led him out to crucify him, all the people ran together and cried, saying\end{Resp}
\begin{Resp}\Rsign The innocent blood of this just person is condemned without a cause.\end{Resp}
\switchcolumn*

\LA
\LessonHead{Lectio V}
\switchcolumn
\EN
\LessonHead{Lesson V}
\switchcolumn*

\LA
\noindent Tum Ægéas irátus, Désine, inquit, Christum jactáre, cui simília verba nihil profuérunt, quóminus a Judǽis crucifigerétur. Andréam vero de Christo nihilóminus líbere prædicántem quod pro salúte humáni géneris se crucifigéndum obtulísset, ímpia oratióne interpéllat, ac demum hortátur, ut sibi cónsulens, diis velit immoláre. Cui Andréas: Ego omnipoténti Deo, qui unus et verus est, ímmolo cotídie, non taurórum carnes, nec hircórum sánguinem, sed immaculátum Agnum in altári; cujus carnem posteáquam omnis pópulus credéntium manducáverit, Agnus, qui sacrificátus est, ínteger persevérat et vivus. Quam ob rem ira accénsus Ægéas, jubet eum in cárcerem detrúdi unde pópulus Andréam fácile liberásset, nisi ipse sedásset multitúdinem, veheméntius rogans, ne se ad optatíssimam martýrii corónam properántem impedírent.\par
\switchcolumn
\EN
\noindent Then Egeas, being angry, answered him, Boast no more of this thy Christ. He spake words even such as thine, but they availed Him not, and He was crucified by the Jews. Whereto Andrew boldly answered that Christ had given Himself up to die for man's salvation; but the Pro-consul blasphemously interrupted him, and bade him look to himself, and sacrifice to the gods. Then said Andrew, We have an altar, whereon day by day I offer up to God, the Almighty, the One, and the True, not the flesh of bulls nor the blood of goats, but a Lamb without spot and when all they that believe have eaten of the Flesh Thereof, the Lamb That was slain abideth whole and liveth. Then Aegeas being filled with wrath, bound the Apostle in prison. Now, the people would have delivered him, but he himself calmed the multitude, and earnestly besought them not to take away from him the crown of martyrdom, for which he longed and which was now drawing near.\par
\switchcolumn*

\LA
\begin{Resp}\Rsign O bona crux, quæ decórem et pulchritúdinem de membris Dómini suscepísti; áccipe me ab homínibus, et redde me magístro meo: \Star{} Ut per te me recípiat, qui per te me redémit.\end{Resp}
\begin{Resp}\Vsign Beátus Andréas expánsis mánibus ad cælum orábat, dicens: Salva me, bona crux.\end{Resp}
\begin{Resp}\Rsign Ut per te me recípiat, qui per te me redémit.\end{Resp}
\switchcolumn
\EN
\begin{Resp}\Rsign O precious cross, which the Members of my Lord have made so fair and goodly, welcome me from among men, and join me again to my Master, \Star{} That, as by thee He redeemed me, so by thee also He may take me unto Himself.\end{Resp}
\begin{Resp}\Vsign The blessed Andrew stretched forth his hands to heaven and prayed, saying: Precious cross, be my salvation,\end{Resp}
\begin{Resp}\Rsign That, as by thee He redeemed me, so by thee also He may take me unto Himself.\end{Resp}
\switchcolumn*

\LA
\LessonHead{Lectio VI}
\switchcolumn
\EN
\LessonHead{Lesson VI}
\switchcolumn*

\LA
\noindent Igitur paulo post in tribúnal prodúctum, cum Ægéas crucis extolléntem mystéria sibíque suam impietátem exprobrántem diútius ferre non posset, in crucem tolli, et Christi mortem imitári jussit. Addúctus Andréas ad locum martýrii, cum crucem vidísset, longe exclamáre cœpit: O bona crux, quæ decórem ex membris Dómini suscepísti, diu desideráta, sollícite amáta, sine intermissióne quæsíta, et aliquándo cupiénti ánimo præparáta: áccipe me ab homínibus, et redde me magístro meo; ut per te me recípiat, qui per te me redémit. Itaque cruci affíxus est: in qua bíduum vivus pendens, et Christi fidem prædicáre nunquam intermíttens, ad eum migrávit, cujus mortis similitúdinem concupíerat. Quæ ómnia presbýteri et diáconi Achájæ, qui ejus passiónem scripsérunt, se ita ut commemoráta sunt, audísse et vidísse testántur. Ejus ossa primum Constantíno imperatóre Constantinópolim, deínde Amálphim transláta sunt. Caput, Pio secúndo Pontífice, Romam allátum, in basílica sancti Petri collocátum est.\par
\switchcolumn
\EN
\noindent Come a short while after, he was brought before the judgment-seat, where he extolled the mystery of the cross, and rebuked Aegeas for his ungodliness. Then Aegeas could bear with him no longer, but commanded him to be crucified, in imitation of Christ. Andrew, then, was led to the place of martyrdom, and, as soon as he came in sight of the cross, he cried out, O precious cross, which the Members of my Lord have made so goodly, how long have I desired thee! how warmly have I loved thee! how constantly have I sought thee! And, now that thou art come to me, how is my soul drawn to thee! Welcome me from among men, and join me again to my Master, that as by thee He redeemed me, so by thee also He may take me unto Himself. So he was fastened to the cross, whereon he hung living for two days, during which time he ceased not to preach the faith of Christ, and, finally, passed into the Presence of Him the likeness of Whose death he had loved so well. All the above particulars of his last sufferings were written by the Priests and Deacons of Achaia, who bear witness to them of their own knowledge. Under the Emperor Constantine the bones of the Apostle were first taken to Constantinople, whence they were afterwards brought to Amalfi. In the Pontificate of Pope Pius II his head was carried to Rome, where it is kept in the Basilica of St. Peter.\par
\switchcolumn*

\LA
\begin{Resp}\Rsign Expándi manus meas tota die in cruce ad pópulum non credéntem, sed contradicéntem mihi: \Star{} Qui ámbulant vias non bonas, sed post peccáta sua.\end{Resp}
\begin{Resp}\Vsign Deus ultiónum Dóminus, Deus ultiónum líbere egit: exaltáre, qui júdicas terram; redde retributiónem supérbis.\end{Resp}
\begin{Resp}\Rsign Qui ámbulant vias non bonas, sed post peccáta sua.\end{Resp}
\begin{Resp}\Vsign Glória Patri, et Fílio, \Star{} et Spirítui Sancto.\end{Resp}
\begin{Resp}\Rsign Qui ámbulant vias non bonas, sed post peccáta sua.\end{Resp}
\switchcolumn
\EN
\begin{Resp}\Rsign All day long I have stretched forth my hands upon the cross unto a disobedient and gainsaying people; \Star{} Which walketh in a way that is not good, but after their own sins.\end{Resp}
\begin{Resp}\Vsign The Lord God to Whom vengeance belongeth, the God to Whom vengeance belongeth, hath shown Himself: lift up thyself, Thou Judge of the earth, render a reward to the proud.\end{Resp}
\begin{Resp}\Rsign Which walketh in a way that is not good, but after their own sins.\end{Resp}
\begin{Resp}\Vsign Glory be to the Father, and to the Son, \Star{} and to the Holy Ghost.\end{Resp}
\begin{Resp}\Rsign Which walketh in a way that is not good, but after their own sins.\end{Resp}
\switchcolumn*

\LA
\LessonHead{Lectio VII}
\switchcolumn
\EN
\LessonHead{Lesson VII}
\switchcolumn*

\LA
\LessonTitle{Léctio sancti Evangélii secúndum Matthǽum}
\switchcolumn
\EN
\LessonTitle{From the Holy Gospel according to Matthew}
\switchcolumn*

\LA
\Cite{Matt 4:18-22}
\switchcolumn
\EN
\Cite{Matt 4:18-22}
\switchcolumn*

\LA
\noindent In illo témpore: Ambulans Jesus juxta mare Galilǽæ, vidit duos fratres, Simónem, qui vocátur Petrus, et Andréam fratrem ejus, mitténtes rete in mare. Et réliqua.\par
\switchcolumn
\EN
\noindent At that time, Jesus, walking by the sea of Galilee, saw two brethren, Simon who is called Peter, and Andrew his brother, casting a net into the sea (for they were fishers). And so on.\par
\switchcolumn*

\LA
\LessonTitle{Homilía sancti Gregórii Papæ}
\switchcolumn
\EN
\LessonTitle{Homily by Pope St. Gregory (the Great.)}
\switchcolumn*

\LA
\Source{Homilia 5 in Evangelia}
\switchcolumn
\EN
\Source{15th on the Gospels}
\switchcolumn*

\LA
\noindent Audístis, fratres caríssimi, quia ad uníus jussiónis vocem Petrus et Andréas relíctis rétibus secúti sunt Redemptórem. Nulla vero hunc fácere adhuc mirácula víderant, nihil ab eo de prǽmio ætérnæ retributiónis audíerant: et tamen ad unum Dómini præcéptum hoc quod possidére videbántur, oblíti sunt. Quanta nos ejus mirácula vidémus, quot flagéllis afflígimur, quantis minárum asperitátibus deterrémur, et tamen vocántem sequi contémnimus.\par
\switchcolumn
\EN
\noindent Dearly beloved brethren, ye hear how that Peter and Andrew, having once heard the Lord call them, left their nets, and followed their Saviour. As yet they had seen none of His miracles, as yet they had received no promise of their exceeding and eternal reward; nevertheless, at one word of the Lord they forgot all those things which they seemed to have. We have seen many of His miracles; we have received many of His gracious chastenings; many times hath He warned us of the wrath to come and yet Christ calleth and we do not follow.\par
\switchcolumn*

\LA
\begin{Resp}\Rsign Orávit sanctus Andréas, dum respíceret in cælum, et voce magna clamávit et dixit: Tu es Deus meus, quem vidi: ne me patiáris ab ímpio júdice depóni: \Star{} Quia virtútem sanctæ crucis agnóvi.\end{Resp}
\begin{Resp}\Vsign Tu es magíster meus Christus, quem diléxi, quem cognóvi, quem conféssus sum: tantúmmodo in ista voce exáudi me.\end{Resp}
\begin{Resp}\Rsign Quia virtútem sanctæ crucis agnóvi.\end{Resp}
\switchcolumn
\EN
\begin{Resp}\Rsign The holy Andrew lifted up his eyes to heaven, and prayed, and cried with a loud voice, and said: Thou art my God, Whom I have seen; suffer not the unjust judge to take me down from the cross \Star{} For now I know what the power of thy holy Cross is.\end{Resp}
\begin{Resp}\Vsign Thou art Christ my Master, Whom I have loved, Whom I have known, Whom I have confessed in this thing hear me.\end{Resp}
\begin{Resp}\Rsign For now I know what the power of thy holy Cross is.\end{Resp}
\switchcolumn*

\LA
\LessonHead{Lectio VIII}
\switchcolumn
\EN
\LessonHead{Lesson VIII}
\switchcolumn*

\LA
\noindent In cælo jam sedet, qui de conversióne nos ádmonet; jam jugo fídei colla géntium súbdidit, jam mundi glóriam stravit, jam ruínis ejus crebrescéntibus, distrícti sui judícii diem propinquántem denúntiat: et tamen supérba mens nostra adhuc non vult hoc sponte desérere, quod cotídie perdit invíta. Quid ergo, fratres caríssimi, quid in ejus judício dictúri sumus, qui ab amóre præséntis sǽculi nec præcéptis fléctimur, nec verbéribus emendámur?\par
\switchcolumn
\EN
\noindent He who calleth us to be converted is now enthroned in heaven; He hath broken the necks of the Gentiles to the yoke of the faith, He hath laid low the glory of the world, and the wrecks thereof, falling ever more and more to decay, do preach unto us that the coming of that day when He is to be revealed as our Judge is drawing nigh and yet, so stubborn is our mind, that we will not yet freely abandon that which, will we, nill we, we lose day by day. Dearly beloved brethren, what shall we answer at His judgment-seat, we whom no lessons can persuade, and no stripes can break of the love of this present world?\par
\switchcolumn*

\LA
\begin{Resp}\Rsign Videns crucem Andréas exclamávit, dicens: O crux admirábilis, o crux desiderábilis, o crux quæ per totum mundum rútilas: \Star{} Súscipe discípulum Christi, ac per te me recípiat, qui per te móriens me redémit.\end{Resp}
\begin{Resp}\Vsign O bona crux, quæ decórem et pulchritúdinem de membris Dómini suscepísti.\end{Resp}
\begin{Resp}\Rsign Súscipe discípulum Christi, ac per te me recípiat, qui per te móriens me redémit.\end{Resp}
\begin{Resp}\Vsign Glória Patri, et Fílio, \Star{} et Spirítui Sancto.\end{Resp}
\begin{Resp}\Rsign Súscipe discípulum Christi, ac per te me recípiat, qui per te móriens me redémit.\end{Resp}
\switchcolumn
\EN
\begin{Resp}\Rsign When Andrew saw the cross he cried, saying: How wonderful art thou, O cross! O cross, how loveable art thou! O cross, thy bright beams enlighten the darkness of the whole world! \Star{} Welcome a follower of Jesus, that, as by thee He died to redeem me, so by thee also He may take me unto Himself.\end{Resp}
\begin{Resp}\Vsign O precious cross, which the Members of my Lord have made so fair and goodly,\end{Resp}
\begin{Resp}\Rsign Welcome a follower of Jesus, that, as by thee He died to redeem me, so by thee also He may take me unto Himself.\end{Resp}
\begin{Resp}\Vsign Glory be to the Father, and to the Son, \Star{} and to the Holy Ghost.\end{Resp}
\begin{Resp}\Rsign Welcome a follower of Jesus, that, as by thee He died to redeem me, so by thee also He may take me unto Himself.\end{Resp}
\switchcolumn*

\LA
\LessonHead{Lectio IX}
\switchcolumn
\EN
\LessonHead{Lesson IX}
\switchcolumn*

\LA
\noindent Sed fortásse áliquis tácitis sibi cogitatiónibus dicat: Ad vocem Domínicam utérque iste piscátor quid, aut quantum dimísit, qui pene nihil hábuit? Sed hac in re, fratres caríssimi, afféctum debémus pótius pensáre quam censum. Multum relíquit, qui sibi nihil retínuit: multum relíquit, qui quantúmlibet parum, totum deséruit. Certe nos et hábita cum amóre possidémus, et ea, quæ mínime habémus, ex desidério quǽrimus. Multum ergo Petrus et Andréas dimísit, quando utérque étiam desidéria habéndi derelíquit.\par
\switchcolumn
\EN
\noindent Perhaps one perchance will ask in his heart, what Peter or Andrew had to lose by obeying the call of the Lord? Dearly beloved brethren, we must consider here rather the intention than the loss incurred by this obedience. He that keepeth nothing for himself, giveth up much; he that sacrificeth his all, sacrificed what is to him a great deal. Beyond doubt, we cling to whatever we have, and what we have least, that we desire most. Peter and Andrew therefore gave up much when they gave up even the desire of possessing anything.\par
\switchcolumn*

\LA
\TeDeum{Te Deum}
\switchcolumn
\EN
\TeDeum{Te Deum}
\switchcolumn*

\end{paracol}

\SeasonTitle{Mensis December}{December}

\addcontentsline{toc}{section}{Mensis December\enspace\textit{December}}

\par\needspace{10\baselineskip}\addvspace{1.2em}{\centering\small\textsc{Die 2 Decembris} · \textsc{2 December}\par}
\Office{S. Bibianæ Virginis et Martyris}{St. Bibiana, Virgin and Martyr}{Semiduplex}
\addcontentsline{toc}{subsection}{Die 2 Decembris · S. Bibianæ Virginis et Martyris\enspace\textit{St. Bibiana, Virgin and Martyr}}
\begin{paracol}{2}
\LA
\RefLine{Lectiones I–III: de Scriptura occurrente.}
\switchcolumn
\EN
\RefLine{Lessons I–III: from the Scripture of the day.}
\switchcolumn*

\LA
\RefLine{Lectiones VII–IX: ex Commúni Unius Virginis et Martyris (C6), in 2 loco.}
\switchcolumn
\EN
\RefLine{Lessons VII–IX: from the Common of a Virgin Martyr (C6), set 2.}
\switchcolumn*

\LA
\LessonHead{Lectio IV}
\switchcolumn
\EN
\LessonHead{Lesson IV}
\switchcolumn*

\LA
\noindent Bibiana virgo Romana, nobili genere nata, christiana fide nobilior fuit. Ejus enim pater Flavianus sub Juliano Apostata impiissimo tyranno expræfectus, servilibusque notis compunctus, ad aquas Taurinas deportatus, martyr occubuit. Mater Dafrosa, et filiæ primum conclusæ domi, ut inedia conficerentur; mox relegata mater extra Urbem capite plexa est. Mortuis autem piis parentibus, Bibiana cum sorore sua Demetria bonis omnibus exspoliatur; Apronianus Urbis prætor, pecuniis inhians, sorores persequitur; quas humana prorsus ope destitutas, Deo mirabiliter, qui dat escam esurientibus, enutriente, cum vivaciores vegetioresque conspexisset, vehementer est admiratus.\par
\switchcolumn
\EN
\noindent Bibiana was a Roman maiden, distinguished on account of the nobility of her family, but now far more distinguished for her confession of Christ. In the reign of the foul tyrant, Julian the Apostate, her father Flavian, although he was an ex-Praefect, was branded as a slave and banished to Acquapendente, not far from Rome, where he soon died a martyr for his faith. His wife, Dafrosa, and his two daughters, Bibiana and Demetria, were first imprisoned in their own house, with the idea of starving them to death; but the mother was afterwards taken outside the city and beheaded. Bibiana and her sister Demetria, after the death of their holy parents, were stripped of all they had in the world. Apronianus, Praetor of the city, who hankered after their property, continued to persecute them, but although they were destitute of all human support, God, Who giveth bread to the hungry, fed them, and kept them in health, life, and strength, to the wonder of their enemies.\par
\switchcolumn*

\LA
\begin{Resp}\Rsign Propter veritátem, et mansuetúdinem, et justítiam: \Star{} Et dedúcet te mirabíliter déxtera tua.\end{Resp}
\begin{Resp}\Vsign Spécie tua et pulchritúdine tua inténde, próspere procéde, et regna.\end{Resp}
\begin{Resp}\Rsign Et dedúcet te mirabíliter déxtera tua.\end{Resp}
\switchcolumn
\EN
\begin{Resp}\Rsign Because of truth, and meekness, and righteousness; \Star{} And thy right hand shall lead thee wonderfully.\end{Resp}
\begin{Resp}\Vsign In thy comeliness, and thy beauty, go forward, fare prosperously, and reign.\end{Resp}
\begin{Resp}\Rsign And thy right hand shall lead thee wonderfully.\end{Resp}
\switchcolumn*

\LA
\LessonHead{Lectio V}
\switchcolumn
\EN
\LessonHead{Lesson V}
\switchcolumn*

\LA
\noindent Suadet nihilominus Apronianus, ut venerentur deos Gentium, amissas ideo opes, imperatoris gratiam, præclarissimas nuptias consecuturæ. Si secus fecerint, minatur carceres, virgas, secures. At illæ neque blanditiis, neque minis a recta fide declinantes, paratæ potius mori, quam fœdari moribus ethnicorum, prætoris impietatem constantissime detestantur. Quare Demetria ob oculos Bibianæ repente corruens, obiit in Domino: Bibiana Rufinæ mulieri vaferrimæ seducenda traditur, quæ ab incunabulis edocta, christianas leges, et illibatum servare virginitatis florem, se ipsa fortior feminæ superavit insidias, et prætoris astus delusit.\par
\switchcolumn
\EN
\noindent Apronianus then attacked them, to make them worship the gods of the Gentiles, and promised them the restoration of their property, the favour of the Emperor, and a great marriage for each of them, if they would give way, and, on the other hand, imprisonment, stripes, and death. But neither promises nor threats availed, for they remained firm in the faith, being resolved rather to die than to pollute themselves by doing according to the deeds of the heathen; and, as for the iniquity of the Praetor, they loathed it continually. At length the strength of Demetria gave way, and she fell down suddenly, and died in the Lord, before the eyes of her sister Bibiana. Then Bibiana was put into the hands of an artful woman named Rufina, to seduce her if possible; but she had known the law of Christ from her childhood, and kept the lily of her purity undefiled, triumphing over the efforts of that vile person, and disappointing the lust of the Praetor.\par
\switchcolumn*

\LA
\begin{Resp}\Rsign Dilexísti justítiam, et odísti iniquitátem: \Star{} Proptérea unxit te Deus, Deus tuus, óleo lætítiæ.\end{Resp}
\begin{Resp}\Vsign Propter veritátem, et mansuetúdinem, et justítiam.\end{Resp}
\begin{Resp}\Rsign Proptérea unxit te Deus, Deus tuus, óleo lætítiæ.\end{Resp}
\switchcolumn
\EN
\begin{Resp}\Rsign Thou hast loved righteousness, and hated iniquity; \Star{} Therefore God, thy God, hath anointed thee with the oil of gladness.\end{Resp}
\begin{Resp}\Vsign Because of truth, and meekness, and righteousness.\end{Resp}
\begin{Resp}\Rsign Therefore God, thy God, hath anointed thee with the oil of gladness.\end{Resp}
\switchcolumn*

\LA
\LessonHead{Lectio VI}
\switchcolumn
\EN
\LessonHead{Lesson VI}
\switchcolumn*

\LA
\noindent Nihil autem proficiente Rufina, quæ, præter dolosa verba, illam quotidie verberibus affligebat, ut de sancto proposito dimoveret, spe sua frustratus prætor, accensus ira, quod in Bibiana perdidisset operam, a lictoribus eam denudari, vinctisque manibus columnæ alligari, eamque plumbatis cædi jubet, donec efflaret animam. Cujus sacrum corpus objectum canibus biduo jacuit in foro Tauri, illæsum tamen, et divinitus servatum; quod deínde Joannes presbyter sepelivit noctu juxta sepulcrum sororis et matris ad palatium Licinianum, ubi usque in præsens exstat ecclesia Deo, sanctæ Bibianæ nomine, dicata, quam Urbanus octavus instauravit, sanctarum Bibianæ, Demetriæ, et Dafrosæ corporibus in ea repertis, et sub ara maxima collocatis.\par
\switchcolumn
\EN
\noindent Then, when Rufina saw that her false words availed not, she took to blows, and scourged Bibiana daily, but the saint was not staggered in her holy resolution. At last the Praetor, mad with baffled lust, when he found his labour was thrown away, ordered his lictors to strip her naked, hang her up by the hands to a pillar, and flog her to death with whips weighted with lead. When all was over, her sacred body was thrown out for the dogs to eat. It lay two days in the Forum Tauri, but the animals would not touch it; and, at last, a Priest, named John, took it, and buried it by night beside the graves of her mother and sister, near the Licinian Palace. This is the place where there is still a church, dedicated in the name of St. Bibiana. When this church was being restored by Urban VIII., the bodies of these three holy women, Bibiana, Demetria, and Dafrosa, were found, and were re-buried under the High Altar.\par
\switchcolumn*

\LA
\begin{Resp}\Rsign Afferéntur Regi vírgines post eam, próximæ ejus; \Star{} Afferéntur tibi in lætítia et exsultatióne.\end{Resp}
\begin{Resp}\Vsign Spécie tua et pulchritúdine tua; inténde, próspere procéde, et regna.\end{Resp}
\begin{Resp}\Rsign Afferéntur tibi in lætítia et exsultatióne.\end{Resp}
\begin{Resp}\Vsign Glória Patri, et Fílio, \Star{} et Spirítui Sancto.\end{Resp}
\begin{Resp}\Rsign Afferéntur tibi in lætítia et exsultatióne.\end{Resp}
\switchcolumn
\EN
\begin{Resp}\Rsign After her shall virgins be brought unto the King, her fellows \Star{} They shall be brought unto thee with gladness and rejoicing.\end{Resp}
\begin{Resp}\Vsign In thy comeliness and thy beauty, go forward, fare prosperously, and reign.\end{Resp}
\begin{Resp}\Rsign They shall be brought unto thee with gladness and rejoicing.\end{Resp}
\begin{Resp}\Vsign Glory be to the Father, and to the Son, \Star{} and to the Holy Ghost.\end{Resp}
\begin{Resp}\Rsign They shall be brought unto thee with gladness and rejoicing.\end{Resp}
\switchcolumn*

\end{paracol}

\par\needspace{10\baselineskip}\addvspace{1.2em}{\centering\small\textsc{Die 3 Decembris} · \textsc{3 December}\par}
\Office{S. Francisci Xaverii Confessoris}{St. Francis Xavier, Confessor}{Duplex majus}
\addcontentsline{toc}{subsection}{Die 3 Decembris · S. Francisci Xaverii Confessoris\enspace\textit{St. Francis Xavier, Confessor}}
\begin{paracol}{2}
\LA
\RefLine{Lectiones I–III: de Scriptura occurrente.}
\switchcolumn
\EN
\RefLine{Lessons I–III: from the Scripture of the day.}
\switchcolumn*

\LA
\LessonHead{Lectio IV}
\switchcolumn
\EN
\LessonHead{Lesson IV}
\switchcolumn*

\LA
\noindent Francíscus, in Xavério diœcésis Pampelonénsis nobílibus paréntibus natus, Parísiis sancto Ignátio sese cómitem et discípulum junxit. Ipso magistro eo brevi devénit, ut in rerum divinárum contemplatióne defíxus, a terra aliquándo sublímis elevarétur: quod illi sacrificánti coram pópuli multitúdine aliquóties evénit. Has ánimi delícias magnis sui córporis cruciátibus merebátur. Nam interdícto sibi, non carnis solum et vini, sed panis quoque tritícei usu, vílibus cibis vesci sólitus, per bíduum subínde triduúmque omni prorsus aliménto abstínuit. Férreis in se flagéllis ita sǽviit, ut sæpe copióso cruóre difflúeret; somnum brevíssimum humi jacens carpébat.\par
\switchcolumn
\EN
\noindent Francis was of noble family, and was born in the castle of Xavier, in the diocese of Pampeluna, (in the year of our Lord 1506.) He was a companion of St. Ignatius at Paris, and one of his earliest disciples. Under his teaching, he learnt to become so wrapt in the contemplation of divine things, that he was sometimes lifted in ecstasy off the ground, which happened to him several times when he was saying Mass in public before large congregations. He earned these refreshments of the soul by the sharpest punishment of the body. He gave up the use not only of meat and wine, but also of wheaten bread; he lived on the vilest food, and ate only once every two or three days. He used an iron scourge till his blood ran freely; he shortened the hours of his rest, and lay only on the ground.\par
\switchcolumn*

\LA
\begin{Resp}\Rsign Honéstum fecit illum Dóminus, et custodívit eum ab inimícis, et a seductóribus tutávit illum: \Star{} Et dedit illi claritátem ætérnam.\end{Resp}
\begin{Resp}\Vsign Justum dedúxit Dóminus per vias rectas, et osténdit illi regnum Dei.\end{Resp}
\begin{Resp}\Rsign Et dedit illi claritátem ætérnam.\end{Resp}
\switchcolumn
\EN
\begin{Resp}\Rsign The Lord made him honourable, and defended him from his enemies, and kept him safe from those that lay in wait for him; \Star{} And gave him perpetual glory.\end{Resp}
\begin{Resp}\Vsign He went down with him into the pit, and left him not in bonds.\end{Resp}
\begin{Resp}\Rsign And gave him perpetual glory.\end{Resp}
\switchcolumn*

\LA
\LessonHead{Lectio V}
\switchcolumn
\EN
\LessonHead{Lesson V}
\switchcolumn*

\LA
\noindent Vitæ austeritáte ac sanctitáte apostólico múneri jam matúrus, cum Joánnes tértius Lusitániæ rex áliquot nascéntis societátis viros a Paulo tértio pro Indiis postulásset, sancti Ignátii hortátu, ab eódem Pontífice ad tantum opus cum apostólici núntii potestáte delígitur. Eo appúlsus, íllico variárum géntium difficíllimis et váriis linguis divínitus instrúctus appáruit. Quin eum quandóque único idiómate ad divérsas gentes concionántem, unaquǽque sua lingua loquéntem audívit. Províncias innúmeras pédibus semper, et sæpe nudis, peragrávit. Fidem Japóniæ et sex áliis regiónibus invéxit. Multa centéna hóminum mília ad Christum in Indiis convértit; magnósque príncipes, regésque complúres sacro fonte expiávit. Et cum tam magna pro Deo ágeret, ea erat humilitáte, ut sancto Ignátio, tunc præpósito suo, flexis génibus scríberet.\par
\switchcolumn
\EN
\noindent The hardness and holiness of his life had made him meet to be called to be an Apostle, and when John III, King of Portugal, asked Pope Paul III to send to the Indies some members of the then new Society of Jesus, the Pontiff, by the advice of St. Ignatius, sent Francis to enter on that vast field of labour with the powers of Apostolic Nuncio. He arrived (in India on the 6th day of May, in the year 1542.) When he began his work, it seemed as though God Himself taught him the many and difficult languages of the natives. It even happened that when he preached in one language to a mixed congregation of different nationalities, each one heard him in his own tongue wherein he was born. He travelled over countless districts, always walking, and often bare-footed. He introduced the faith into Japan, and six other countries. In India he turned many hundred thousands to Christ, and regenerated many chiefs and kings in the holy font. And notwithstanding that he was doing all these great things for God's service, so deep was his lowliness that when he wrote to St. Ignatius, the General of the Society, he did so on his knees.\par
\switchcolumn*

\LA
\begin{Resp}\Rsign Amávit eum Dóminus, et ornávit eum: stolam glóriæ índuit eum, \Star{} Et ad portas paradísi coronávit eum.\end{Resp}
\begin{Resp}\Vsign Índuit eum Dóminus lorícam fídei, et ornávit eum.\end{Resp}
\begin{Resp}\Rsign Et ad portas paradísi coronávit eum.\end{Resp}
\switchcolumn
\EN
\begin{Resp}\Rsign The Lord loved him and beautified him; He clothed him with a robe of glory, \Star{} And crowned him at the gates of Paradise.\end{Resp}
\begin{Resp}\Vsign The Lord hath put on him the breast-plate of faith, and hath adorned him.\end{Resp}
\begin{Resp}\Rsign And crowned him at the gates of Paradise.\end{Resp}
\switchcolumn*

\LA
\LessonHead{Lectio VI}
\switchcolumn
\EN
\LessonHead{Lesson VI}
\switchcolumn*

\LA
\noindent Hunc dilatándi Evangélii ardórem multitúdine et excelléntia miraculórum Dóminus roborávit. Cæco visum réddidit. Tantum marínæ aquæ signo crucis convértit in dulcem, quantum quingéntis vectóribus, qui siti adigebántur ad mortem, diu suffécit: qua in várias quoque regiónes asportáta, ægri plúrimi súbito curáti sunt. Plures mórtuos revocávit ad vitam, inter quos prídie sepúltum érui jussum e túmulo suscitávit, duósque álios, dum efferréntur, apprehénsa eórum manu, paréntibus e féretro vivos restituit Prophetíæ spíritu passim afflátus, plúrima et loco et témpore remotíssima enuntiávit. Demum in Sanciáno, Sinárum ínsula, die secúnda Decémbris, óbiit plenus méritis, laboribúsque conféctus. Demórtui cadáver viva calce per multos menses bis óbrutum, sed pénitus incorrúptum, odóre et sánguine manávit; et ubi Málacam delátum est, pestem sævíssimam conféstim exstínxit. Dénique ubíque terrárum novis maximísque fulgéntem miráculis Gregórius décimus quintus Sanctis adscrípsit. Pius autem décimus ipsum sodalitáti et óperi Propagándæ Fídei cæléstem patrónum elégit atque constítuit.\par
\switchcolumn
\EN
\noindent God was pleased to support his zeal for spreading the Gospel with many and great miracles. He gave sight to a blind man. On one occasion the supply of fresh water failed when he was at sea, and five hundred sailors were in danger of perishing by thirst, but the servant of God, by the sign of the Cross, turned salt water into fresh, and they used it for a considerable time. Some of this water was also carried into different countries, and a great number of sick persons were instantaneously cured by it. He called several dead men to life, among whom was one who had been buried the day before, and who was disinterred by command of the saint; and likewise two others who were being carried to the grave, and whom he took by the hand and restored living to their parents. He had the spirit of prophecy, and foretold many things, remote both in place and time. Utterly worn out with his labours, he died full of good works in the island of San-Chan in the Canton River, (upon the 2nd day of December, in the year of our Lord 1552.) His body was buried in quick lime, and, being again taken up, was again buried in the same, but at the end of many months it was found entirely incorrupt, and sweet, and, when cut, blood flowed freely from it. From China it was carried to Malacca, and, as soon as it reached that place, a plague, which was raging there, ceased. At length, when he had become famous throughout the whole world for new and wonderful miracles, Gregory XV added his name to the list of the Saints.\par
\switchcolumn*

\LA
\begin{Resp}\Rsign Iste homo perfécit ómnia quæ locútus est ei Deus, et dixit ad eum: Ingrédere in réquiem meam: \Star{} Quia te vidi justum coram me ex ómnibus géntibus.\end{Resp}
\begin{Resp}\Vsign Iste est, qui contémpsit vitam mundi, et pervénit ad cæléstia regna.\end{Resp}
\begin{Resp}\Rsign Quia te vidi justum coram me ex ómnibus géntibus.\end{Resp}
\begin{Resp}\Vsign Glória Patri, et Fílio, \Star{} et Spirítui Sancto.\end{Resp}
\begin{Resp}\Rsign Quia te vidi justum coram me ex ómnibus géntibus.\end{Resp}
\switchcolumn
\EN
\begin{Resp}\Rsign This is he which did according unto all that God commanded him; and God said unto him: Enter thou into My rest, \Star{} For thee have I seen righteous before Me among all people.\end{Resp}
\begin{Resp}\Vsign This is he which loved not his life in this world, and is come unto an everlasting kingdom.\end{Resp}
\begin{Resp}\Rsign For thee have I seen righteous before Me among all people.\end{Resp}
\begin{Resp}\Vsign Glory be to the Father, and to the Son, \Star{} and to the Holy Ghost.\end{Resp}
\begin{Resp}\Rsign For thee have I seen righteous before Me among all people.\end{Resp}
\switchcolumn*

\LA
\LessonHead{Lectio VII}
\switchcolumn
\EN
\LessonHead{Lesson VII}
\switchcolumn*

\LA
\LessonTitle{Léctio sancti Evangélii secúndum Marcum}
\switchcolumn
\EN
\LessonTitle{From the Holy Gospel according to Mark}
\switchcolumn*

\LA
\Cite{Marc 16:15-18}
\switchcolumn
\EN
\Cite{Mark 16:15-18}
\switchcolumn*

\LA
\noindent In illo témpore: Dixit Jesus discípulis suis: Euntes in mundum univérsum prædicate Evangélium omni creaturæ. Et réliqua.\par
\switchcolumn
\EN
\noindent In that time, Jesus said to his disciples: Go ye into the whole world, and preach the gospel to every creature. And so on.\par
\switchcolumn*

\LA
\LessonTitle{Homilía sancti Gregórii Papæ}
\switchcolumn
\EN
\LessonTitle{Homily by Pope St. Gregory (the Great.)}
\switchcolumn*

\LA
\Source{Homilia 29 in Evangelia, post initium}
\switchcolumn
\EN
\Source{29th on the Gospels.}
\switchcolumn*

\LA
\noindent Potest omnis creaturæ nómine omnis natio Géntium designari. Ante enim dictum fuerat: In viam Géntium ne abieritis; nunc autem dícitur: Prædicate omni creaturæ: ut scilicet prius a Judǽa Apostolórum repulsa prædicátio tunc nobis in adjutórium fieret, cum hanc illa ad damnatiónis suæ testimónium superba repulísset. Sed cum discipulos ad prædicándum Véritas mittit, quid aliud in mundo facit, nisi grana seminis spargit? Et pauca grana mittit in semine, ut multárum messium fruges recipiat ex nostra fide.\par
\switchcolumn
\EN
\noindent By the words every creature we may understand every tribe of the Gentiles. Of aforetime it had been said, Go not into the way of the Gentiles, (Matth. x. 5,) but now, Preach the Gospel to every creature, that, since the Jews had proudly rejected the preaching of the Apostles, that might become our gain which was the seal of their condemnation. But when the Eternal Truth sendeth forth His disciples to preach, what doth He but scatter seed over the field of the world? He scattereth abroad a few grains for seed, that He may afterward reap an abundant harvest in our faith.\par
\switchcolumn*

\LA
\begin{Resp}\Rsign Iste est qui ante Deum magnas virtútes operátus est, et de omni corde suo laudávit Dóminum: \Star{} Ipse intercédat pro peccátis ómnium populórum.\end{Resp}
\begin{Resp}\Vsign Ecce homo sine queréla, verus Dei cultor, ábstinens se ab omni ópere malo, et pérmanens in innocéntia sua.\end{Resp}
\begin{Resp}\Rsign Ipse intercédat pro peccátis ómnium populórum.\end{Resp}
\switchcolumn
\EN
\begin{Resp}\Rsign This is he which wrought great wonders before God, and praised the Lord with all his heart. \Star{} May he pray for all people, that their sins may be forgiven unto them.\end{Resp}
\begin{Resp}\Vsign Behold a man without blame, a worshipper of God in truth, keeping himself clean from every evil work, and abiding still in his innocency.\end{Resp}
\begin{Resp}\Rsign May he pray for all people, that their sins may be forgiven unto them.\end{Resp}
\switchcolumn*

\LA
\LessonHead{Lectio VIII}
\switchcolumn
\EN
\LessonHead{Lesson VIII}
\switchcolumn*

\LA
\noindent Neque étenim in universo mundo tanta fidelium messis exsúrgeret, si de manu Dómini super rationalem terram illa electa grana prædicántium non venissent. Sequitur: Qui crediderit et baptizátus fúerit, salvus erit: qui vero non crediderit, condemnábitur. Fortasse unusquísque apud semetípsum dicat: Ego jam credidi, salvus ero. Verum dicit, si fidem opéribus tenet. Vera étenim fides est, quæ in hoc, quod verbis dicit, móribus non contradicit. Hinc est enim quod de quibusdam falsis fidelibus Paulus dicit: Qui confiténtur se nosse Deum, factis autem negant.\par
\switchcolumn
\EN
\noindent The great harvest of faithful souls throughout the whole world would never have sprung up, if the hand of the Lord had not first scattered those chosen grains of preachers over the reasonable soil of men's minds. Then is written, He that believeth and is baptized shall be saved but he that believeth not, shall be damned. Now, perchance, thou sayest in thine heart I believe, and therefore I shall be saved. True, if to thy faith thou dost add works. He only hath a living faith whose life doth not give the lie to his profession. It is of this that Paul speaketh, where he saith of certain vain believers, They profess that they know God; but in works they deny Him. (Tit. i. 16.)\par
\switchcolumn*

\LA
\begin{Resp}\Rsign Sint lumbi vestri præcíncti, et lucérnæ ardéntes in mánibus vestris: \Star{} Et vos símiles homínibus exspectántibus dóminum suum, quando revertátur a núptiis.\end{Resp}
\begin{Resp}\Vsign Vigiláte ergo, quia nescítis qua hora Dóminus vester ventúrus sit.\end{Resp}
\begin{Resp}\Rsign Et vos símiles homínibus exspectántibus dóminum suum, quando revertátur a núptiis.\end{Resp}
\begin{Resp}\Vsign Glória Patri, et Fílio, \Star{} et Spirítui Sancto.\end{Resp}
\begin{Resp}\Rsign Et vos símiles homínibus exspectántibus dóminum suum, quando revertátur a núptiis.\end{Resp}
\switchcolumn
\EN
\begin{Resp}\Rsign Let your loins be girded about, and your lights burning; \Star{} And ye yourselves like unto men that wait for their lord, when he will return from the wedding.\end{Resp}
\begin{Resp}\Vsign Watch therefore, for ye know not what hour your Lord doth come.\end{Resp}
\begin{Resp}\Rsign And ye yourselves like unto men that wait for their lord, when he will return from the wedding.\end{Resp}
\begin{Resp}\Vsign Glory be to the Father, and to the Son, \Star{} and to the Holy Ghost.\end{Resp}
\begin{Resp}\Rsign And ye yourselves like unto men that wait for their lord, when he will return from the wedding.\end{Resp}
\switchcolumn*

\LA
\LessonHead{Lectio IX}
\switchcolumn
\EN
\LessonHead{Lesson IX}
\switchcolumn*

\LA
\noindent Signa autem eos qui credituri sunt, hæc sequéntur: In nómine meo dæmónia ejícient, linguis loquéntur novis, serpéntes tollent: et si mortiferum quid biberint, non eis nocébit: super ægros manus impónent, et bene habébunt. Numquidnam, fratres mei, quia ista signa non fácitis, minime creditis? Sed hæc necessaria in exordio Ecclésiæ fuérunt. Ut enim ad fidem cresceret multitúdo credéntium, miraculis fuerat nutriénda: quia et nos, cum arbústa plantamus, tamdiu eis aquam infúndimus, quoúsque ea in terra jam coaluisse videámus: et si semel radícem fixerint, irrigátio cessabit. Hinc est enim quod Paulus dicit: Linguæ in signum sunt non fidelibus, sed infidelibus.\par
\switchcolumn
\EN
\noindent And these signs shall follow them that believe In My name they shall cast out devils, they shall speak with new tongues, they shall take up serpents; and if they drink any deadly thing, it shall not hurt them; they shall lay hands on the sick, and they shall recover. My brethren, these signs do not follow us. Do we, then, not believe? Nay. The truth is, these things were needful when the Church was young. That she might grow by the increase of the faithful, she needed to be nourished with miracles. So we, when we plant a young tree, continually water and tend it, till we see that it hath taken firm root in the earth but when once it hath taken firm root, it can grow of itself. Hence Paul saith of tongues Tongues are for a sign, not to them that believe, but to them that believe not. \textasciitilde{}(i Cor. xiv. 22.)\par
\switchcolumn*

\LA
\TeDeum{Te Deum}
\switchcolumn
\EN
\TeDeum{Te Deum}
\switchcolumn*

\end{paracol}

\par\needspace{10\baselineskip}\addvspace{1.2em}{\centering\small\textsc{Die 4 Decembris} · \textsc{4 December}\par}
\Office{S. Petri Chrysologi Episcopi Confessoris et Ecclesiæ Doctoris}{St. Peter Chrysologus, Bishop, Confessor and Doctor of the Church}{Duplex}
\addcontentsline{toc}{subsection}{Die 4 Decembris · S. Petri Chrysologi Episcopi Confessoris et Ecclesiæ Doctoris\enspace\textit{St. Peter Chrysologus, Bishop, Confessor and Doctor of the Church}}
\begin{paracol}{2}
\LA
\RefLine{Lectiones I–III: de Scriptura occurrente.}
\switchcolumn
\EN
\RefLine{Lessons I–III: from the Scripture of the day.}
\switchcolumn*

\LA
\RefLine{Lectiones VII–IX: ex Commúni Doctoris Pontificis (C4a).}
\switchcolumn
\EN
\RefLine{Lessons VII–IX: from the Common of a Doctor Bishop (C4a).}
\switchcolumn*

\LA
\LessonHead{Lectio IV}
\switchcolumn
\EN
\LessonHead{Lesson IV}
\switchcolumn*

\LA
\noindent Petrus, qui ob auream ejus eloquéntiam Chrysólogi cognomen adeptus est, Foro Cornelii in Æmília honestis paréntibus natus, a prima ætate ánimum ad religiónem adjíciens, Cornelio Romano, tunc ejusdem urbis Corneliénsis episcopo, operam dedit: a quo etiam, sciéntia et vitæ sanctitáte cum brevi profecísset, diaconus creátus est. Postmodum contigit, ut Ravennates ob mortem archipræsulis sui alium (ut moris erat) ab eis eléctum Romam ad sanctum Xystum Papam tertium pro confirmatióne miserint una cum legátis suis et cum prædicto Cornelio, qui eumdem levitam secum perdúxit. Interim sanctus Petrus Apóstolus, et Martyr Apollináris summo Pontifici in somnis apparuérunt, mediúmque habéntes hunc juvenem, jussérunt, ut illum, et non alium, in archiepiscopum Ravennæ crearet. Hinc Pontifex, mox ut vidit Petrum, cognóvit eum a Dómino eo præeléctum: proptérea, rejecto illo quem ipsi offerébant, hunc solum, anno Christi quadringentésimo trigesimo tertio, illi metropolitanæ præfecit ecclésiæ. Quod cum legáti Ravennatenses ægre ferrent, audíta visióne, divinæ voluntáti libenter acquiescéntes, novum archiepiscopum maxima cum reveréntia suscepérunt.\par
\switchcolumn
\EN
\noindent Peter, called in Greek Chrysologos, or, of the golden words, on account of his wonderful eloquence, was born of respectable parents at Imola, near Ravenna. He displayed a very early leaning to godliness, and became a disciple of Cornelius of Rome, Bishop of Imola. This Prelate, having experience of his learning and holiness of life, soon ordained him Deacon. On the death of the Archbishop of Ravenna, the people of that place elected a successor, and sent him, according to custom, to Rome, to be confirmed in his appointment by Pope Sixtus III. The Archbishop elect accordingly set forth, along with the ambassadors of the people of Ravenna and Cornelius, Bishop of Imola, attended by Peter the Deacon. While they were yet on the way, the holy Apostle Peter and Apollinaris the Martyr appeared to the Supreme Pontiff in a dream, leading a young man between them, whom they commanded him to make Archbishop. As soon as the embassy arrived at Rome the Pope knew in Peter the young man of his dream, chosen of God to the Archbishopric. Wherefore he set aside him that the people of Ravenna had presented, and preferred Peter to that Metropolitan Church, in the year of our Lord, 433. The ambassadors of the people of Ravenna took it ill, till they heard the vision then they gave themselves up to the will of God, and received the new Archbishop with great reverence.\par
\switchcolumn*

\LA
\begin{Resp}\Rsign Invéni David servum meum, óleo sancto meo unxi eum: \Star{} Manus enim mea auxiliábitur ei.\end{Resp}
\begin{Resp}\Vsign Nihil profíciet inimícus in eo, et fílius iniquitátis non nocébit ei.\end{Resp}
\begin{Resp}\Rsign Manus enim mea auxiliábitur ei.\end{Resp}
\switchcolumn
\EN
\begin{Resp}\Rsign I have found David My servant, with My holy oil have I anointed him \Star{} For My hand shall help him.\end{Resp}
\begin{Resp}\Vsign The enemy shall prevail nothing against him, nor the son of wickedness afflict him.\end{Resp}
\begin{Resp}\Rsign For My hand shall help him.\end{Resp}
\switchcolumn*

\LA
\LessonHead{Lectio V}
\switchcolumn
\EN
\LessonHead{Lesson V}
\switchcolumn*

\LA
\noindent Petrus ígitur, licet invitus, in archipræsulem consecratus, Ravennam dedúcitur; ubi a Valentiniáno imperatóre, et a Galla Placidia ejus matre, ac ab universo pópulo maxima lætítia exceptus est. Et ille ab eis id unum pétere dixit, ut quando tantum oneris pro ipsórum salúte subire non recusávit, studérent ipsi mónitis suis obtemperare, divinisque præceptis non obsistere. Duórum Sanctórum tunc ibi defunctórum córpora optimis unguentis condíta sepelívit, Barbatiáni videlicet presbyteri, et Germani Antisiodorénsis epíscopi; cujus étiam cucúllam et cilícium sibi vindicávit in hereditátem. Projectum et Marcellinum in episcopos ordinávit. In Classe fontem exstruxit magnitúdinis vere admirábilis, et templa quædam magnifica ædificávit tum beato Andreæ Apóstolo, tum aliis Sanctis. Ludos ab homínibus personátis cum variis saltatiónibus Kaléndis Januarii fíeri sólitos concióne cohibuit acerrima; ubi inter alia illud præcláre dixit: Qui jocari volúerit cum diabolo, non póterit gaudere cum Christo. Jussu sancti Leónis Papæ primi scripsit póstea ad Chalcedonénse Concílium advérsus hæresim Eutychétis. Respóndit præterea ad Eutychen ipsum et alia epistola, quæ eidem concílio in novis editiónibus præfixa, et in annales ecclesiásticos reláta fuit.\par
\switchcolumn
\EN
\noindent Peter being against his will consecrated Archbishop, arrived at Ravenna, where he was received with great joy by the Emperor Valentinian, the Empress-Mother Galla Placidia, and all the people. And this one thing he asked of them, that, as he, for the saving of their souls, had not refused to bear the heavy weight of the Archbishopric, so they would strive to follow his warnings, and live in submission to the law of God. He took the bodies of the two Saints, namely, Barbatian the Priest, and German, Bishop of Auxerre, and caused them to be embalmed with rich ointments and honourably buried, and he kept the cowl and haircloth shirt of German for a legacy for himself. At Classis, three miles from Ravenna, he built a Baptistery of extraordinary size, and several splendid churches, in honour of the blessed Apostle Andrew and other Saints. He preached a most severe sermon against the acting and dancing of guisards about New Year time, in which discourse he said among other things, He that jesteth with the devil will never rejoice with Christ. By command of Pope Leo I. he addressed an Epistle to the Council of Chalcedon against the heretic Eutyches. He also confuted Eutyches himself in another letter, which is likewise published in the new editions of the Acts of the Council, and is matter of Church History.\par
\switchcolumn*

\LA
\begin{Resp}\Rsign Pósui adjutórium super poténtem, et exaltávi eléctum de plebe mea: \Star{} Manus enim mea auxiliábitur ei.\end{Resp}
\begin{Resp}\Vsign Invéni David servum meum, óleo sancto meo unxi eum.\end{Resp}
\begin{Resp}\Rsign Manus enim mea auxiliábitur ei.\end{Resp}
\switchcolumn
\EN
\begin{Resp}\Rsign I have laid help upon one that is mighty, and have exalted one chosen out of My people \Star{} For My hand shall help him.\end{Resp}
\begin{Resp}\Vsign I have found David My servant, with My holy oil have I anointed him.\end{Resp}
\begin{Resp}\Rsign For My hand shall help him.\end{Resp}
\switchcolumn*

\LA
\LessonHead{Lectio VI}
\switchcolumn
\EN
\LessonHead{Lesson VI}
\switchcolumn*

\LA
\noindent Dum publice sermónes haberet ad pópulum, adeo vehemens erat in dicendo, ut præ nimio ardore vox illi intérdum defécerit, sicut cóntigit in concióne mulíeris hæmorrhoíssæ. Unde Ravennates commóti, tot lácrimis, clamóribus et oratiónibus locum replevérunt, ut póstea ipse grátias ágeret Deo, quod in lucrum amoris verterit damnum ejusdem sermónis. Cum tandem annos circiter decem et octo eam ecclésiam sanctíssime rexísset, labórum suórum finem adesse divinitus prænoscens, in pátriam se cóntulit; ubi sancti Cassiáni templum ingréssus, magnum diadéma aureum gemmis distinctum pretiosíssimis offerens, super altáre majus pósuit, necnon aureum cratérem, et patenam argenteam, quam tum rábidi canis morsus, tum febres sanare sæpius expertum est, aqua inde demíssa. Ex tunc Ravennates, qui eumdem secúti fuerant, dimísit, ádmonens, ut in eligéndo optimo pastore invigilárent attente. Mox Deum humíliter precátus et sanctum Cassianum patronum, ut benígne ánimam ejus exciperet, quarto Nonas Decembris plácide ex hac vita migrávit, anno Dómini circiter quadringentésimo quinquagesimo. Sacrum illíus corpus, communi totius civitátis fletu ac pietáte prope corpus ejusdem sancti Cassiáni honorifice cónditum, nostris étiam tempóribus religiose cólitur: cujus tamen brácchium, auro et gemmis ornátum Ravennam delátum, in Ursiana æde venerátur.\par
\switchcolumn
\EN
\noindent When he preached in public his vehemence was such that he sometimes became speechless from excitement. This happened to him once when he was preaching on the subject of the woman who had an issue of blood. (Matth. ix. 20-22.) The congregation on this occasion were so wrought up, that they filled the whole place with tears, cries, and prayers, and Peter afterwards thanked God, Who had turned his failure to the profit of their souls. When he had ruled the Church of Ravenna in holiness for about eighteen years, God gave him knowledge that the end of his labours was at hand, and he returned to his home at Imola, to die. When he arrived at Imola, he entered the church of St. Cassian, and offered upon the High Altar a great circlet of gold, set with stones of great price, a golden chalice and a silver paten. Water poured out of these vessels hath often healed hydrophobia and fevers. Some of the people of Ravenna had followed the Archbishop, but he now dismissed them, with a charge to use great prudence in their choice of his successor. Then he fell to prayer, that God would mercifully receive his spirit, asking the same likewise for the sake of his patron St Cassian, and so he passed in peace to a better life, on a 2nd of December, about the year of our Lord 450. His holy body was buried, amid the sorrow and veneration of the whole city, hard by the remains of St. Cassian, where it lieth even to this day, guarded with great reverence. One arm was cut off and sent to Ravenna, where it is preserved in the Ursian Church, in a reliquary of gold and precious stones.\par
\switchcolumn*

\LA
\begin{Resp}\Rsign Iste est, qui ante Deum magnas virtútes operátus est, et omnis terra doctrína ejus repléta est: \Star{} Ipse intercédat pro peccátis ómnium populórum.\end{Resp}
\begin{Resp}\Vsign Iste est, qui contémpsit vitam mundi, et pervénit ad cæléstia regna.\end{Resp}
\begin{Resp}\Rsign Ipse intercédat pro peccátis ómnium populórum.\end{Resp}
\begin{Resp}\Vsign Glória Patri, et Fílio, \Star{} et Spirítui Sancto.\end{Resp}
\begin{Resp}\Rsign Ipse intercédat pro peccátis ómnium populórum.\end{Resp}
\switchcolumn
\EN
\begin{Resp}\Rsign This is he which wrought great wonders before God, and the whole earth is full of his teaching \Star{} May he pray for all people, that their sins may be forgiven unto them\end{Resp}
\begin{Resp}\Vsign This is he which loved not his life in this world, and hath attained unto the kingdom of heaven.\end{Resp}
\begin{Resp}\Rsign May he pray for all people, that their sins may be forgiven unto them\end{Resp}
\begin{Resp}\Vsign Glory be to the Father, and to the Son, \Star{} and to the Holy Ghost.\end{Resp}
\begin{Resp}\Rsign May he pray for all people, that their sins may be forgiven unto them\end{Resp}
\switchcolumn*

\end{paracol}

\par\needspace{10\baselineskip}\addvspace{1.2em}{\centering\small\textsc{Die 6 Decembris} · \textsc{6 December}\par}
\Office{S. Nicolai Episcopi et Confessoris}{St. Nicholas, Bishop and Confessor}{Duplex}
\addcontentsline{toc}{subsection}{Die 6 Decembris · S. Nicolai Episcopi et Confessoris\enspace\textit{St. Nicholas, Bishop and Confessor}}
\begin{paracol}{2}
\LA
\RefLine{Lectiones I–III: de Scriptura occurrente.}
\switchcolumn
\EN
\RefLine{Lessons I–III: from the Scripture of the day.}
\switchcolumn*

\LA
\RefLine{Lectiones VII–IX: ex Commúni unius Confessoris Pontificis (C4).}
\switchcolumn
\EN
\RefLine{Lessons VII–IX: from the Common of a Confessor Bishop (C4).}
\switchcolumn*

\LA
\LessonHead{Lectio IV}
\switchcolumn
\EN
\LessonHead{Lesson IV}
\switchcolumn*

\LA
\noindent Nicoláum, illustri loco Pátaræ in Lycia natum, paréntes a Deo precibus impetrarunt. Cujus viri sanctitas quanta futura esset, jam ab incunábulis appáruit. Nam infans, cum réliquos dies lac nutrícis frequens súgeret, quarta et sexta feria semel dumtaxat, idque vesperi sugebat: quam jejunii consuetúdine in réliqua vita semper ténuit. Adoléscens paréntibus orbatus, facultates suas paupéribus distríbuit. Cujus illud insigne est christianæ benignitátis exemplum, quod cum ejus civis egens tres filias jam nubiles in matrimonio collocare non posset, earumque pudicítiam prostitúere cogitaret; re cognita, Nicolaus noctu per fenestram tantum pecuniæ in ejus domum injecit, quantum uníus vírginis doti satis esset: quod cum íterum et tertio fecísset, tres illæ vírgines honestis viris in matrimónium datæ sunt.\par
\switchcolumn
\EN
\noindent Nicolas was born at the famous city of Patara in Lycia. His parents obtained him from God by prayer, and the holiness of his life was marked even from the cradle. When he was at the breast he never would suck more than once on Wednesdays and Fridays, and that always after sunset, though he sucked freely on other days. This custom of fasting he never broke through during his whole life. While he was still a young man he lost both his father and mother, after which he gave his whole property away to the poor. One particular example is given of his Christian charity. There was a certain needy man in the city who had three marriageable daughters, for whom he could not get husbands, and so thought to make them harlots. When Nicolas heard of it, he went to the house by night and threw in by the window such a sum of money as made a dowry for one of them. This he did a second and a third time, and thus by his charity they were honourably given in marriage.\par
\switchcolumn*

\LA
\begin{Resp}\Rsign Invéni David servum meum, óleo sancto meo unxi eum: \Star{} Manus enim mea auxiliábitur ei.\end{Resp}
\begin{Resp}\Vsign Nihil profíciet inimícus in eo, et fílius iniquitátis non nocébit ei.\end{Resp}
\begin{Resp}\Rsign Manus enim mea auxiliábitur ei.\end{Resp}
\switchcolumn
\EN
\begin{Resp}\Rsign I have found David My servant, with My holy oil have I anointed him \Star{} For My hand shall help him.\end{Resp}
\begin{Resp}\Vsign The enemy shall prevail nothing against him, nor the son of wickedness afflict him.\end{Resp}
\begin{Resp}\Rsign For My hand shall help him.\end{Resp}
\switchcolumn*

\LA
\LessonHead{Lectio V}
\switchcolumn
\EN
\LessonHead{Lesson V}
\switchcolumn*

\LA
\noindent Cum vero se totum Deo dedísset, in Palæstinam profectus est, ut loca sancta viseret, et præsens venerarétur. Qua in peregrinatióne navem conscéndens sereno cælo et tranquillo mari, horribilem nautis tempestátem prædixit; moxque ortam, cum essent omnes in summo periculo, orans mirabíliter sedávit. Unde cum domum reversus singuláris sanctitátis ómnibus documénta præberet, Dei admónitu Myram, quæ Lyciæ metropolis erat, venit; quo témpore ejus urbis episcopo mortuo, provinciales epíscopi de successore deligéndo consultábant. Itaque in ea deliberatióne divinitus admoniti sunt, ut eam eligerent, qui postridie mane primus in ecclésiam ingrederétur, Nicolaus nómine. Qua observatióne adhibita, in ecclésiæ janua deprehénsus est Nicoláus, et summo ómnium consénsu Myræ epíscopus creatur. In episcopatu castitátem, quam semper colúerat, gravitátem, oratiónis assiduitátem, vigílias abstinéntiam, liberalitátem et hospitalitátem, in adhortando mansuetúdinem, in reprehendéndo severitátem perpetuo adhibuit.\par
\switchcolumn
\EN
\noindent When he had given himself entirely to God he set forth for Palestine, that he might see the Holy Places, and worship therein. During this pilgrimage he embarked once on board a ship when the sky was clear and the sea calm, but he foretold a great storm, which afterwards arose and raged until the sailors were afraid; and then the saint by prayer stilled the tempest. After he had returned home, and his holy life was known to all men, God bade him go to Myra, which is the chief city of Lycia, at a time when the Bishop had just died and the Bishops of the Province were called together to choose a successor. While they deliberated, they received a warning from heaven to choose that Nicolas who should first come into the church in the morning. In obedience to that warning, Nicolas was seized at the door of the church, and with universal consent consecrated Archbishop. In his great office he was an unceasing model of purity, as he had always been, of gravity, of regularity in prayer, of watching, of abstinence, of charity, of hospitality, of meekness in exhortation, and of sternness in rebuke.\par
\switchcolumn*

\LA
\begin{Resp}\Rsign Pósui adjutórium super poténtem, et exaltávi eléctum de plebe mea: \Star{} Manus enim mea auxiliábitur ei.\end{Resp}
\begin{Resp}\Vsign Invéni David servum meum, óleo sancto meo unxi eum.\end{Resp}
\begin{Resp}\Rsign Manus enim mea auxiliábitur ei.\end{Resp}
\switchcolumn
\EN
\begin{Resp}\Rsign I have laid help upon one that is mighty, and have exalted one chosen out of My people \Star{} For My hand shall help him.\end{Resp}
\begin{Resp}\Vsign I have found David My servant, with My holy oil have I anointed him.\end{Resp}
\begin{Resp}\Rsign For My hand shall help him.\end{Resp}
\switchcolumn*

\LA
\LessonHead{Lectio VI}
\switchcolumn
\EN
\LessonHead{Lesson VI}
\switchcolumn*

\LA
\noindent Viduis et orphanis pecúnia, consílio, ópera non defuit; oppressos adeo sublevávit, ut étiam tres tribunos, per calumniam a Constantino Augusto condemnatos, qui se propter famam ejus miraculórum oratiónibus longíssime absenti commendarant, adhuc vivens, cum imperatóri, mináciter eum terrens, apparuísset, liberaverit. Cum vero contra edictum Diocletiáni et Maximiáni christianæ fidei veritátem Myræ prædicaret, ab imperatórum satellítibus comprehénsus, et longíssime abductus, in carcerem conjéctus est; ubi fuit usque ad Constantinum imperatórem, cujus jussu ex custódia ereptus, Myram rediit. Mox ad Nicænum concílium se cóntulit; ubi cum trecentis illis decem et octo Pátribus Arianam hæresim condemnávit. Inde reversus ad episcopátum, non ita multo post, instante morte, suspíciens in cælum, cum Angelos sibi occurréntes intuerétur, illo Psalmo pronuntiato, In te Dómine, sperávi; usque ad eum locum, In manus tuas comméndo spíritum meum; in cælestem pátriam migrávit. Ejus corpus Bárium in Apulia translátum, ibidem summa celebritate ac veneratióne cólitur.\par
\switchcolumn
\EN
\noindent He was the comforter of widows and orphans by money, by advice, and by labour. He was the deliverer of the oppressed, so mightily, that it is related that the Emperor Constantine once unjustly condemned three Tribunes to death, and these unhappy men called upon Nicolas, though living and absent, to save them, who yet appeared in a vision to the Emperor, and forced him by threats to set them free. When the Emperors Diocletian and Maximian published their edict against Christianity, Nicolas did not cease to preach the truth at Myra, wherefore he was seized by the soldiers of the Emperors, carried away from his See, and thrown into prison, where he remained until the accession of Constantine. This Prince set him free, and he returned to Myra. He betook himself to the first Council of Nicea, where he was one of the 318 Bishops who condemned the heresy of Arius. He returned thence to his Bishopric, and, not long after, became aware of the approach of death. When his last moment was come, he lifted up his eyes to heaven, and, when he saw the Angels coming to meet him, he began to recite the thirtieth Psalm, In thee, O Lord, do I put my trust, and when he had said, Into thy hands I commend my spirit, he passed to the heavenly Fatherland. His body was finally removed to Bari in Apulia, where it is kept with great fame and honour.\par
\switchcolumn*

\LA
\begin{Resp}\Rsign Iste est, qui ante Deum magnas virtútes operátus est, et omnis terra doctrína ejus repléta est: \Star{} Ipse intercédat pro peccátis ómnium populórum.\end{Resp}
\begin{Resp}\Vsign Iste est, qui contémpsit vitam mundi, et pervénit ad cæléstia regna.\end{Resp}
\begin{Resp}\Rsign Ipse intercédat pro peccátis ómnium populórum.\end{Resp}
\begin{Resp}\Vsign Glória Patri, et Fílio, \Star{} et Spirítui Sancto.\end{Resp}
\begin{Resp}\Rsign Ipse intercédat pro peccátis ómnium populórum.\end{Resp}
\switchcolumn
\EN
\begin{Resp}\Rsign This is he which wrought great wonders before God, and the whole earth is full of his teaching \Star{} May he pray for all people, that their sins may be forgiven unto them\end{Resp}
\begin{Resp}\Vsign This is he which loved not his life in this world, and hath attained unto the kingdom of heaven.\end{Resp}
\begin{Resp}\Rsign May he pray for all people, that their sins may be forgiven unto them\end{Resp}
\begin{Resp}\Vsign Glory be to the Father, and to the Son, \Star{} and to the Holy Ghost.\end{Resp}
\begin{Resp}\Rsign May he pray for all people, that their sins may be forgiven unto them\end{Resp}
\switchcolumn*

\end{paracol}

\par\needspace{10\baselineskip}\addvspace{1.2em}{\centering\small\textsc{Die 7 Decembris} · \textsc{7 December}\par}
\Office{S. Ambrosii Episcopi Confessoris et Ecclesiæ Doctoris}{St. Ambrose, Bishop, Confessor and Doctor of the Church}{Duplex}
\addcontentsline{toc}{subsection}{Die 7 Decembris · S. Ambrosii Episcopi Confessoris et Ecclesiæ Doctoris\enspace\textit{St. Ambrose, Bishop, Confessor and Doctor of the Church}}
\begin{paracol}{2}
\LA
\RefLine{Lectiones I–III: de Scriptura occurrente.}
\switchcolumn
\EN
\RefLine{Lessons I–III: from the Scripture of the day.}
\switchcolumn*

\LA
\RefLine{Lectiones VII–IX: ex Commúni Doctoris Pontificis (C4a).}
\switchcolumn
\EN
\RefLine{Lessons VII–IX: from the Common of a Doctor Bishop (C4a).}
\switchcolumn*

\LA
\LessonHead{Lectio IV}
\switchcolumn
\EN
\LessonHead{Lesson IV}
\switchcolumn*

\LA
\noindent Ambrósius epíscopus Mediolanénsis, Ambrósii civis Románi fílius, patre Gálliæ præfécto natus est. In hujus infántis ore exámen apum consedísse dícitur: quæ res divínam viri eloquéntiam præmonstrábat. Romæ liberálibus disciplínis erudítus est. Post a Probo præfécto Ligúriæ et Æmíliæ præpósitus: unde póstea, ejúsdem Probi jussu, cum potestáte Mediolánum venit; ubi, mórtuo Auxéntio Ariáno epíscopo, pópulus de successóre deligéndo dissidébat. Quare Ambrósius, pro offícii sui múnere ecclésiam ingréssus, ut commótam seditiónem sedáret, cum multa de quiéte et tranquillitáte reipúblicæ præcláre dixísset, derepénte púero Ambrósium epíscopum exclamánte, univérsi pópuli vox erúpit, Ambrósium epíscopum deposcéntis.\par
\switchcolumn
\EN
\noindent Ambrose, Bishop of Milan, was the son of another Ambrose, a Roman citizen, and was born when his father was Prefect of Gaul, (about the year of our Lord 340.) A swarm of bees settled upon his face when he was in his cradle, which was considered an omen of his future eloquence. He received a liberal education at Rome. He was afterwards, under the Prefect Probus, made governor of Liguria and Emilia, and so came with authority to Milan. Auxentius, an Arian, who had been intruded into the Bishopric of Milan, happening to die, the most violent disputes arose about the choice of a successor. Ambrose came to the church in his official capacity, and urged upon the contending factions, in a long and powerful speech, the necessity of keeping the public peace; whereupon a child suddenly cried out, Ambrose, Bishop, and the whole assembly took it up, and unanimously called for his election.\par
\switchcolumn*

\LA
\begin{Resp}\Rsign Invéni David servum meum, óleo sancto meo unxi eum: \Star{} Manus enim mea auxiliábitur ei.\end{Resp}
\begin{Resp}\Vsign Nihil profíciet inimícus in eo, et fílius iniquitátis non nocébit ei.\end{Resp}
\begin{Resp}\Rsign Manus enim mea auxiliábitur ei.\end{Resp}
\switchcolumn
\EN
\begin{Resp}\Rsign I have found David My servant, with My holy oil have I anointed him \Star{} For My hand shall help him.\end{Resp}
\begin{Resp}\Vsign The enemy shall prevail nothing against him, nor the son of wickedness afflict him.\end{Resp}
\begin{Resp}\Rsign For My hand shall help him.\end{Resp}
\switchcolumn*

\LA
\LessonHead{Lectio V}
\switchcolumn
\EN
\LessonHead{Lesson V}
\switchcolumn*

\LA
\noindent Recusánte illo et eórum précibus resisténte, ardens pópuli stúdium ad Valentiniánum imperatórem delátum est; cui gratíssimum fuit, a se deléctos júdices ad sacerdótium postulári. Fuit id étiam Probo præfécto jucúndum, qui Ambrósio proficiscénti quasi divínans díxerat: Vade, age, non ut judex, sed ut epíscopus. Itaque cum ad pópuli desidérium imperatóris volúntas accéderet, Ambrósius baptizátus (erat enim catechúmenus) sacrísque initiátus, ac servátis ómnibus ex institúto Ecclésiæ órdinum grádibus, octávo die, qui fuit séptimo Idus Decémbris, episcopále onus suscépit. Factus epíscopus, cathólicam fidem et disciplínam ecclesiásticam acérrime deféndit; multósque Ariános, et álios hæréticos ad fídei veritátem convértit, in quibus claríssimum Ecclésiæ lumen sanctum Augustínum Jesu Christo péperit.\par
\switchcolumn
\EN
\noindent Ambrose refused, and would not yield to their prayers, whereupon they carried their petition to the Emperor Valentinian. It was very pleasing to this Prince that those he had appointed as judges should be chosen Bishops, as also to the Prefect Probus, who had, as it were prophetically, said to him when he appointed him, Go and govern them more like a Bishop than a Judge. When the will of the Emperor was added to the desire of the people, Ambrose yielded, and received Baptism, (for hitherto he was only a Catechumen,) Confirmation, and Communion, and then the several Orders on successive days, till on the eighth day, which was the 7th of December, (in the year 374,) the weight of the Episcopate was laid upon his shoulders. Being made Bishop, he showed himself a stout upholder of the Catholic faith, and the discipline of the Church, and turned to the truth great numbers of Arians and other heretics, and, among them, he begat in Christ Jesus that burning and shining light of the Church, Augustine.\par
\switchcolumn*

\LA
\begin{Resp}\Rsign Pósui adjutórium super poténtem, et exaltávi eléctum de plebe mea: \Star{} Manus enim mea auxiliábitur ei.\end{Resp}
\begin{Resp}\Vsign Invéni David servum meum, óleo sancto meo unxi eum.\end{Resp}
\begin{Resp}\Rsign Manus enim mea auxiliábitur ei.\end{Resp}
\switchcolumn
\EN
\begin{Resp}\Rsign I have laid help upon one that is mighty, and have exalted one chosen out of My people \Star{} For My hand shall help him.\end{Resp}
\begin{Resp}\Vsign I have found David My servant, with My holy oil have I anointed him.\end{Resp}
\begin{Resp}\Rsign For My hand shall help him.\end{Resp}
\switchcolumn*

\LA
\LessonHead{Lectio VI}
\switchcolumn
\EN
\LessonHead{Lesson VI}
\switchcolumn*

\LA
\noindent Gratiáno imperatóre occíso, ad Máximum ejus interfectórem legátus íterum proféctus est; eóque pœniténtiam ágere recusánte, se ab ejus communióne semóvit. Theodósium imperatórem, propter cædem Thessalonícæ factam, ingréssu ecclésiæ prohíbuit. Cui, cum ille David quoque regem adúlterum et homicídam fuísse dixísset, respóndit Ambrósius: Qui secútus es errántem, séquere pœniténtem. Quare Theodósius sibi ab eo impósitam públicam pœniténtiam humíliter egit. Ergo sanctus epíscopus pro Ecclésia Dei máximis labóribus curísque perfúnctus, multis libris étiam egrégie conscrípsit, ántequam in morbum incíderet, mortis suæ diem prædíxit. Ad quem ægrótum Honorátus Vercellénsis epíscopus, Dei voce ter admónitus, accúrrit, eíque sanctum Dómini corpus prǽbuit: quo ille sumpto, conformátis in crucis similitúdinem mánibus, orans, ánimam Deo réddidit prídie Nonas Aprílis, anno post Christum natum trecentésimo nonagésimo séptimo.\par
\switchcolumn
\EN
\noindent After the murder of the Emperor Gratian, (in 383,) Ambrose was sent as an ambassador to Maximus, by whom he had been slain, and, as he refused to repent, the Bishop renounced his communion. After the massacre which the Emperor Theodosius had commanded at Thessalonica, (in 390,) he refused to permit that Prince to enter a church. The Emperor pleaded that he was no worse than David, who had been guilty of adultery and murder, to which Ambrose answered him, As thou hast followed him in his sin, follow him also in his repentance. Then Theodosius humbly did public penance laid upon him by the Bishop. At length the Saint was worn out with his continual labour and care for the Church, (for the which also he composed many excellent books,) and foretold that the day of his death was at hand, though he had not then fallen into his last sickness. As he lay dying, Honoratus, Bishop of Vercelli, heard a voice from God three times crying to him that the hour of Ambrose's departure was come, whereupon he went to him quickly, and gave him the sacred Body of our Lord. When he had received It, the Saint, still praying, with his hands stretched out in the form of a cross, gave his spirit to God, upon the 4th day of April, in the year of Christ, 397.\par
\switchcolumn*

\LA
\begin{Resp}\Rsign Iste est, qui ante Deum magnas virtútes operátus est, et omnis terra doctrína ejus repléta est: \Star{} Ipse intercédat pro peccátis ómnium populórum.\end{Resp}
\begin{Resp}\Vsign Iste est, qui contémpsit vitam mundi, et pervénit ad cæléstia regna.\end{Resp}
\begin{Resp}\Rsign Ipse intercédat pro peccátis ómnium populórum.\end{Resp}
\begin{Resp}\Vsign Glória Patri, et Fílio, \Star{} et Spirítui Sancto.\end{Resp}
\begin{Resp}\Rsign Ipse intercédat pro peccátis ómnium populórum.\end{Resp}
\switchcolumn
\EN
\begin{Resp}\Rsign This is he which wrought great wonders before God, and the whole earth is full of his teaching \Star{} May he pray for all people, that their sins may be forgiven unto them\end{Resp}
\begin{Resp}\Vsign This is he which loved not his life in this world, and hath attained unto the kingdom of heaven.\end{Resp}
\begin{Resp}\Rsign May he pray for all people, that their sins may be forgiven unto them\end{Resp}
\begin{Resp}\Vsign Glory be to the Father, and to the Son, \Star{} and to the Holy Ghost.\end{Resp}
\begin{Resp}\Rsign May he pray for all people, that their sins may be forgiven unto them\end{Resp}
\switchcolumn*

\end{paracol}

\par\needspace{10\baselineskip}\addvspace{1.2em}{\centering\small\textsc{Die 8 Decembris} · \textsc{8 December}\par}
\Office{In Conceptione Immaculata Beatæ Mariæ Virginis}{In Conceptione Immaculata Beatæ Mariæ Virginis}{Duplex I classis cum Octava communi}
\addcontentsline{toc}{subsection}{Die 8 Decembris · In Conceptione Immaculata Beatæ Mariæ Virginis\enspace\textit{In Conceptione Immaculata Beatæ Mariæ Virginis}}
\begin{paracol}{2}
\LA
\LessonHead{Lectio I}
\switchcolumn
\EN
\LessonHead{Lesson I}
\switchcolumn*

\LA
\LessonTitle{De libro Génesis}
\switchcolumn
\EN
\LessonTitle{From the book of Genesis}
\switchcolumn*

\LA
\Cite{Gen 3:1-5}
\switchcolumn
\EN
\Cite{Gen 3:1-5}
\switchcolumn*

\LA
\noindent Serpens erat callídior cunctis animántibus terræ quæ fécerat Dóminus Deus. Qui dixit ad mulíerem: Cur præcépit vobis Deus ut non comederétis de omni ligno paradísi? Cui respóndit múlier: De fructu lignórum, quæ sunt in paradíso, véscimur: de fructu vero ligni quod est in médio paradísi, præcépit nobis Deus ne comederémus, et ne tangerémus illud, ne forte moriámur. Dixit autem serpens ad mulíerem: Nequáquam morte moriémini. Scit enim Deus quod in quocúmque die comedéritis ex eo, aperiéntur óculi vestri, et éritis sicut dii, sciéntes bonum et malum.\par
\switchcolumn
\EN
\noindent Now the serpent was more subtle than any of the beasts of the earth which the Lord God made. And he said to the woman: Why hath God commanded you, that you should not eat of every tree of paradise? And the woman answered him, saying: Of the fruit of the trees that are in paradise we do eat: But of the fruit of the tree which is in the midst of paradise, God hath commanded us that we should not eat; and that we should not touch it, lest perhaps we die. And the serpent said to the woman: No, you shall not die the death. For God doth know that in what day soever you shall eat thereof, your eyes shall be opened: and you shall be as Gods, knowing good and evil.\par
\switchcolumn*

\LA
\begin{Resp}\Rsign Per unum hóminem peccátum in hunc mundum intrávit, in quo omnes peccavérunt. \Star{} Ne tímeas, María, invenísti grátiam apud Deum.\end{Resp}
\begin{Resp}\Vsign Erípuit Dóminus ánimam tuam de morte, et contra inimícum factus est protéctor tuus.\end{Resp}
\begin{Resp}\Rsign Ne tímeas, María, invenísti grátiam apud Deum.\end{Resp}
\switchcolumn
\EN
\begin{Resp}\Rsign By one man sin entered into the world, in whom all have sinned. \Star{} Fear not, Mary, for thou hast found grace with God.\end{Resp}
\begin{Resp}\Vsign The Lord hath delivered thy soul from death, yea, the Lord was thy stay.\end{Resp}
\begin{Resp}\Rsign Fear not, Mary, for thou hast found grace with God.\end{Resp}
\switchcolumn*

\LA
\LessonHead{Lectio II}
\switchcolumn
\EN
\LessonHead{Lesson II}
\switchcolumn*

\LA
\Cite{Gen 3:6-8}
\switchcolumn
\EN
\Cite{Gen 3:6-8}
\switchcolumn*

\LA
\noindent Vidit ígitur múlier quod bonum esset lignum ad vescéndum, et pulchrum óculis, aspectúque delectábile: et tulit de fructu illíus, et comédit: dedítque viro suo, qui comédit. Et apérti sunt óculi ambórum; cumque cognovíssent se esse nudos, consuérunt fólia ficus, et fecérunt sibi perizómata. Et cum audíssent vocem Dómini Dei deambulántis in paradíso ad auram post merídiem, abscóndit se Adam et uxor ejus a fácie Dómini Dei in médio ligni paradísi.\par
\switchcolumn
\EN
\noindent And the woman saw that the tree was good to eat, and fair to the eyes, and delightful to behold: and she took of the fruit thereof, and did eat, and gave to her husband who did eat. And the eyes of them both were opened: and when they perceived themselves to be naked, they sewed together fig leaves, and made themselves aprons. And when they heard the voice of the Lord God walking in paradise at the afternoon air, Adam and his wife hid themselves from the face of the Lord God, amidst the trees of paradise.\par
\switchcolumn*

\LA
\begin{Resp}\Rsign Transíte ad me, omnes qui concupíscitis me: \Star{} Et narrábo vobis quanta fecit Deus ánimæ meæ.\end{Resp}
\begin{Resp}\Vsign Vivit Dóminus, quóniam adimplévit in me misericórdiam suam.\end{Resp}
\begin{Resp}\Rsign Et narrábo vobis quanta fecit Deus ánimæ meæ.\end{Resp}
\switchcolumn
\EN
\begin{Resp}\Rsign Come unto me all ye that be desirous of me \Star{} And I will declare what God hath done for my soul.\end{Resp}
\begin{Resp}\Vsign As the Lord liveth, by me He hath fulfilled His mercy.\end{Resp}
\begin{Resp}\Rsign And I will declare what God hath done for my soul.\end{Resp}
\switchcolumn*

\LA
\LessonHead{Lectio III}
\switchcolumn
\EN
\LessonHead{Lesson III}
\switchcolumn*

\LA
\Cite{Gen 3:9-15}
\switchcolumn
\EN
\Cite{Gen 3:9-15}
\switchcolumn*

\LA
\noindent Vocavítque Dóminus Deus Adam, et dixit ei: Ubi es? Qui ait: Vocem tuam audívi in paradíso, et tímui, eo quod nudus essem, et abscóndi me. Cui dixit: Quis enim indicávit tibi quod nudus esses, nisi quod ex ligno de quo præcéperam tibi ne coméderes, comedísti? Dixítque Adam: Múlier, quam dedísti mihi sóciam, dedit mihi de ligno, et comédi. Et dixit Dóminus Deus ad mulíerem: Quare hoc fecísti? Quæ respóndit: Serpens decépit me, et comédi. Et ait Dóminus Deus ad serpéntem: Quia fecísti hoc, maledíctus es inter ómnia animántia, et béstias terræ: super pectus tuum gradiéris, et terram cómedes cunctis diébus vitæ tuæ. Inimicítias ponam inter te et mulíerem, et semen tuum et semen illíus: ipsa cónteret caput tuum, et tu insidiáberis calcáneo ejus.\par
\switchcolumn
\EN
\noindent And the Lord God called Adam, and said to him: Where art thou? And he said: I heard thy voice in paradise; and I was afraid, because I was naked, and I hid myself. And he said to him: And who hath told thee that thou wast naked, but that thou hast eaten of the tree whereof I commanded thee that thou shouldst not eat? And Adam said: The woman, whom thou gavest me to be my companion, gave me of the tree, and I did eat. And the Lord God said to the woman: Why hast thou done this? And she answered: The serpent deceived me, and I did eat. And the Lord God said to the serpent: Because thou hast done this thing, thou art cursed among all cattle, and the beasts of the earth: upon thy breast shalt thou go, and earth shalt thou eat all the days of thy life. I will put enmities between thee and the woman, and thy seed and her seed: she shall crush thy head, and thou shalt lie in wait for her heel.\par
\switchcolumn*

\LA
\begin{Resp}\Rsign Elécta mea cándida sicut nix in Líbano; sicut favus distíllans lábia ejus: \Star{} Mel et lac sub lingua illíus.\end{Resp}
\begin{Resp}\Vsign Veni de Líbano, sponsa mea, veni, coronáberis coróna gratiárum.\end{Resp}
\begin{Resp}\Rsign Mel et lac sub lingua illíus.\end{Resp}
\begin{Resp}\Vsign Glória Patri, et Fílio, \Star{} et Spirítui Sancto.\end{Resp}
\begin{Resp}\Rsign Mel et lac sub lingua illíus.\end{Resp}
\switchcolumn
\EN
\begin{Resp}\Rsign My beloved is white like snow in Lebanon, her lips drop as the honeycomb. \Star{} Honey and milk are under her tongue.\end{Resp}
\begin{Resp}\Vsign Come from Lebanon, My Spouse, thou shalt be crowned with a crown of grace.\end{Resp}
\begin{Resp}\Rsign Honey and milk are under her tongue.\end{Resp}
\begin{Resp}\Vsign Glory be to the Father, and to the Son, \Star{} and to the Holy Ghost.\end{Resp}
\begin{Resp}\Rsign Honey and milk are under her tongue.\end{Resp}
\switchcolumn*

\LA
\LessonHead{Lectio IV}
\switchcolumn
\EN
\LessonHead{Lesson IV}
\switchcolumn*

\LA
\LessonTitle{Sermo sancti Hierónymi Presbýteri}
\switchcolumn
\EN
\LessonTitle{From the Sermons of St. Jerome, Priest (at Bethlehem.)}
\switchcolumn*

\LA
\Source{De Assumptione B. M. V.}
\switchcolumn
\EN
\Source{On the Assumption}
\switchcolumn*

\LA
\noindent Qualis et quanta esset beáta et gloriósa semper Virgo María, ab Angelo divínitus declarátur, cum dícitur: Ave, grátia plena; Dóminus tecum: benedícta tu in muliéribus. Tálibus namque decébat Vírginem oppignorári munéribus, ut esset grátia plena, quæ dedit cælis glóriam, terris Dóminum, pacémque refúdit, fidem géntibus, finem vítiis, vitæ órdinem, móribus disciplínam. Et bene plena, quia céteris per partes præstátur; Maríæ vero simul se tota infúdit plenitúdo grátiæ. Vere plena, quia etsi in sanctis Pátribus et Prophétis grátia fuísse créditur, non tamen eátenus plena; in Maríam vero totíus grátiæ, quæ in Christo est, plenitúdo venit, quamquam áliter. Et ídeo inquit: Benedícta tu in muliéribus; id est plus benedícta quam omnes mulíeres. Ac per hoc quidquid maledictiónis infúsum est per Hevam, totum ábstulit benedíctio Maríæ. De ipsa Sálomon in Cánticis, quasi in laudem ejus, Veni, inquit, colúmba mea, immaculáta mea. Jam enim hiems tránsiit, imber ábiit et recéssit. Ac deínde inquit: Veni de Líbano, veni, coronáberis.\par
\switchcolumn
\EN
\noindent Who and what was the blessed and glorious Mary, always a Virgin, hath been revealed by God by the message of an Angel, in these words: Hail, thou that art full of grace, the Lord is with thee blessed art thou among women. It was fitting that a fullness of grace should be poured into that Virgin who hath given to God glory and to man a Saviour, who hath brought peace to earth, who hath given faith to the Gentiles, who hath killed sin, who hath given law to life, who hath made the crooked ways straight. Verily, she is full of grace. To others grace cometh measure by measure; in Mary grace dwelleth at once in all fullness. Verily, she is full of grace. We believe that the holy Fathers and Prophets had grace; but they were not full of grace. But into Mary came a fullness of all the grace which is in Christ, albeit otherwise than as it is in Him. Therefore is it said: Blessed art thou among women, that is, Blessed art thou above all women. The fulness of blessing in Mary utterly neutralized in her any effects of the curse of Eve. In her praise Solomon writeth in the Song of Songs, (ii. 10,): Rise up, my dove, my fair one, for the winter is past, the rain is over and gone. And again: Come from Lebanon, my Spouse, come, thou shalt be crowned. \textasciitilde{}(iv. 8.)\par
\switchcolumn*

\LA
\begin{Resp}\Rsign Ego ex ore Altíssimi prodívi, primogénita ante omnem creatúram: ego feci in cælis, ut orirétur lumen indefíciens. \Star{} Nondum erant abýssi, et ego jam concépta eram.\end{Resp}
\begin{Resp}\Vsign Deus enim creávit me in justítia, et apprehéndit manum meam, et servávit me.\end{Resp}
\begin{Resp}\Rsign Nondum erant abýssi, et ego jam concépta eram.\end{Resp}
\switchcolumn
\EN
\begin{Resp}\Rsign I came out of the mouth of the Most High, the first-begotten before every creature. I made the unfading light to arise in the heavens. \Star{} When there were no depths I was conceived.\end{Resp}
\begin{Resp}\Vsign For the Lord hath created me in righteousness, and hath held mine hand, and hath kept me.\end{Resp}
\begin{Resp}\Rsign When there were no depths I was conceived.\end{Resp}
\switchcolumn*

\LA
\LessonHead{Lectio V}
\switchcolumn
\EN
\LessonHead{Lesson V}
\switchcolumn*

\LA
\noindent Non immérito ígitur veníre de Líbano jubétur, quia Líbanus candidátio interpretátur. Erat enim candidáta multis meritórum virtútibus, et dealbáta nive candídior, Spíritus Sancti munéribus, simplicitátem colúmbæ in ómnibus repræséntans: quóniam, quidquid in ea gestum est, totum púritas et simplícitas, totum véritas et grátia fuit; totum misericórdia et justítia, quæ de cælo prospéxit: et ídeo immaculáta, quia in nullo corrúpta. Circúmdedit enim virum in útero, sicut Jeremías sanctus testátur, et non aliúnde accépit. Fáciet, inquit, Dóminus novum super terram, et múlier circúmdabit virum. Vere novum, et ómnium novitátum superéminens nóvitas virtútum, quando Deus (quem ferre non potest mundus, neque vidére áliquis, ut vívere possit) sic ingréssus est hospítium ventris, ut córporis claustrum nescíret; sicque gestátus, ut totus Deus in eo esset; et sic exívit inde, ut esset (sicut Ezéchiel fatétur) porta omníno clausa. Unde cánitur in eísdem Cánticis de ea: Hortus conclúsus, fons signátus, emissiónes tuæ paradísus. Vere hortus deliciárum, in quo cónsita sunt univérsa florum génera, et odoraménta virtútum; sicque conclúsus, ut nésciat violári, neque corrúmpi ullis insidiárum fráudibus. Fons ítaque signátus sigíllo totíus Trinitátis.\par
\switchcolumn
\EN
\noindent Not unjustly then is she bidden to come from Lebanon, for Lebanon is so named on account of its stainless and glistening whiteness. The earthly Lebanon is white with snow, but the lonely heights of Mary's holiness are white with purity and grace, brilliantly fair, whiter far than snow, sparkling with the gifts of the Holy Ghost she is undefiled like a dove, all clean, all upright, full of grace and truth. She is full of mercy, and of the righteousness that hath looked down from heaven, and therefore is she without stain because in her hath never been any corruption. She hath compassed a man in her womb, saith holy Jeremiah, but she conceived not by the will of fallen man. The Lord, saith the Prophet, hath created a new thing in the earth; a woman shall compass a man. (xxxi. 22.) Verily, it is a new thing. Verily, it was a new work of power, greater than all other works, when God, Whom the world cannot bear, and Whom no man shall see and live, entered the lodging of her womb, breaking not the blissful cloister of her virgin flesh. And in her body He was borne, the Infinite inclosed within her womb. And from her womb He came forth, so that it was fulfilled which was spoken of the Prophet Ezekiel, saying: This gate shall be shut, it shall not be opened, and no man shall enter in by it; because the Lord, the God of Israel, hath entered in by it, therefore it shall be shut. (xliv. 2.) Hence also in the Song of Songs it is said of her (iv. 12,): A garden enclosed is my sister, my spouse, a garden enclosed, a fountain sealed, thy perfumes are a garden of delights. Verily a garden of delights, filled with the perfumes of all flowers, rich with the sweet savour of grace. And the most holy Virgin herself is a garden enclosed, whereinto sin and Satan have never entered to sully the blossoms, a fountain sealed, sealed with the seal of the Trinity.\par
\switchcolumn*

\LA
\begin{Resp}\Rsign Nihil inquinátum in eam incúrrit: \Star{} Candor est lucis ætérnæ et spéculum sine mácula.\end{Resp}
\begin{Resp}\Vsign Est enim hæc speciósior sole, et luci comparáta invenítur púrior.\end{Resp}
\begin{Resp}\Rsign Candor est lucis ætérnæ et spéculum sine mácula.\end{Resp}
\switchcolumn
\EN
\begin{Resp}\Rsign No defiled thing can fall into her; \Star{} She is the brightness of the everlasting light, and the unspotted mirror of the power of God.\end{Resp}
\begin{Resp}\Vsign For she is more beautiful than the sun, and being compared with the light, she is found before it.\end{Resp}
\begin{Resp}\Rsign She is the brightness of the everlasting light, and the unspotted mirror of the power of God.\end{Resp}
\switchcolumn*

\LA
\LessonHead{Lectio VI}
\switchcolumn
\EN
\LessonHead{Lesson VI}
\switchcolumn*

\LA
\noindent Ex Actis Pii Papæ noni\par
\switchcolumn
\EN
\noindent From the acts of Pope Pius IX\par
\switchcolumn*

\LA
\LessonTitle{Deíparæ autem Vírginis in sua Conceptióne de tetérrimo humáni géneris hoste victóriam, quam divína elóquia, veneránda tradítio, perpétuus Ecclésiæ sensus, singuláris episcopórum ac fidélium conspirátio, insígnia quoque summórum Pontíficum acta atque constitutiónes mirífice jam illustrábant, Pius nonus Póntifex máximus totíus Ecclésiæ votis ánnuens státuit suprémo suo atque infallíbili oráculo solémniter proclamáre. Itaque sexto idus decémbris anni millésimi octingentésimi quinquagésimi quarti in basílica Vaticána, ingénti sanctæ Románæ Ecclésiæ patrum cardinálium et Episcopórum ex díssitis étiam regiónibus astánte cœtu, universóque plaudénte orbe, solémniter pronuntiávit et definívit: Doctrínam quæ tenet beatíssimam Vírginem Maríam in primo instánti suæ Conceptiónis fuísse, singulári Dei privilégio, ab omni originális culpæ labe præservátam immúnem, esse a Deo revelátam, ac proínde ab ómnibus fidélibus fírmiter constantérque credéndam.}
\switchcolumn
\EN
\LessonTitle{The fact that the Virgin Mother of God had at the moment of her conception triumphed over the foul enemy of man, hath ever been borne out by the Holy Scriptures, by the venerable tradition of the Church, and by her unceasing belief, as well as by the common conviction of all Bishops and faithful Catholics, and by marked acts and constitutions of the Holy See. At length the Supreme Pontiff Pius IX, in compliance with the wishes of the Universal Church, determined to publish it as a truth of faith, on his own absolute and unerring authority, and accordingly, on the 8th day of December, 1854, in the Vatican Basilica, in presence of a great multitude composed of the Fathers Cardinals of the Holy Roman Church, and Bishops from all parts of the earth, he, with the consent and jubilation of the whole world, declared and defined as follows That doctrine which declareth that the most Blessed Virgin Mary was in the first instant of her Conception preserved, by a special privilege granted unto her by God, from any stain of original sin, is a doctrine taught and revealed by God, and therefore is to be held by all faithful Christians firmly and constantly.}
\switchcolumn*

\LA
\begin{Resp}\Rsign Signum magnum appáruit in cælo: Múlier amícta sole, et luna sub pédibus ejus, \Star{} Et in cápite ejus coróna stellárum duódecim.\end{Resp}
\begin{Resp}\Vsign Induit eam Dóminus vestiméntis salútis, induménto justítiæ, et quasi sponsam ornávit eam monílibus suis.\end{Resp}
\begin{Resp}\Rsign Et in cápite ejus coróna stellárum duódecim.\end{Resp}
\begin{Resp}\Vsign Glória Patri, et Fílio, \Star{} et Spirítui Sancto.\end{Resp}
\begin{Resp}\Rsign Et in cápite ejus coróna stellárum duódecim.\end{Resp}
\switchcolumn
\EN
\begin{Resp}\Rsign There appeared a great wonder in heaven; a Woman clothed with the sun, and the moon under her feet \Star{} And upon her head a crown of twelve stars.\end{Resp}
\begin{Resp}\Vsign The Lord hath clothed her with the garments of salvation, and hath covered her with the robe of righteousness, yea, as a bride He hath adorned her with jewels.\end{Resp}
\begin{Resp}\Rsign And upon her head a crown of twelve stars.\end{Resp}
\begin{Resp}\Vsign Glory be to the Father, and to the Son, \Star{} and to the Holy Ghost.\end{Resp}
\begin{Resp}\Rsign And upon her head a crown of twelve stars.\end{Resp}
\switchcolumn*

\LA
\LessonHead{Lectio VII}
\switchcolumn
\EN
\LessonHead{Lesson VII}
\switchcolumn*

\LA
\LessonTitle{Léctio sancti Evangélii secúndum Lucam}
\switchcolumn
\EN
\LessonTitle{From the Holy Gospel according to Luke}
\switchcolumn*

\LA
\Cite{Luc 1:26-28}
\switchcolumn
\EN
\Cite{Luke 1:26-28}
\switchcolumn*

\LA
\noindent In illo témpore: Missus est Angelus Gábriel a Deo in civitátem Galilǽæ, cui nomen Názareth, ad Vírginem desponsátam viro, cui nomen erat Joseph, de domo David, et nomen Vírginis María. Et réliqua.\par
\switchcolumn
\EN
\noindent In that time, the angel Gabriel was sent from God into a city of Galilee, called Nazareth, to a virgin espoused to a man whose name was Joseph, of the house of David; and the virgin's name was Mary. And so on.\par
\switchcolumn*

\LA
\LessonTitle{Homilía sancti Germáni Epíscopi}
\switchcolumn
\EN
\LessonTitle{Homily by St. Germain, Patriarch (of Constantinople.)}
\switchcolumn*

\LA
\Source{In Præsentatione Deiparæ}
\switchcolumn
\EN
\Source{On the Presentation of the Blessed Virgin}
\switchcolumn*

\LA
\noindent Ave, María, grátia plena, Sanctis sánctior, et cælis excélsior, et Chérubim gloriósior, et Séraphim honorabílior, et super omnem creatúram venerabílior. Ave, colúmba, quæ nobis et fructum fers olívæ, et servatórem a spiritáli dilúvio ac portum salútis annúntias; cujus pennæ deargentátæ, et posterióra dorsi in pallóre auri sanctíssimi et illuminántis Spíritus fulgóre irradiántur. Ave, amœníssimus et rationális Dei paradísus, benevolentíssima et omnipoténti ejúsdem dextra hódie ad Oriéntem plantátus, et ipsi suáve olens lílium, et rosam immarcescíbilem gérminans in eórum medélam, qui pestíferam animǽque exitiálem amaritúdinem mortis ad Occidéntem ebíberant; paradísus, in quo ad veritátis agnitiónem lignum vivíficum effloréscit, e quo qui gustáverint, immortalitátem consequúntur. Ave, sacrosáncte ædificátum, immaculátum, purissimúmque Dei summi Regis palátium, ejúsdem Dei Regis magnificéntia circumornátum, omnésque hospítio recípiens ac mýsticis refíciens delíciis; in quo non manufáctus et vário decóre nitens situs est spirituális Sponsi thálamus; in quo Verbum errántem humánam stirpem revocáre volens, carnem sibi desponsávit, ut eos, qui voluntáte própria extórres facti fúerant, Patri reconciliáret.\par
\switchcolumn
\EN
\noindent Hail, Mary, full of grace, holier than the Saints, higher than the heavens, more glorious than the Cherubim, more honourable than the Seraphim, and the most worshipful thing that the hands of God have made. Hail, O dove, bearing in thy beak the olive-branch of peace that telleth us of salvation from the spiritual flood, (Gen. viii. 10, n,) dove, blessed omen of a safe harbour, whose wings are of silver, and thy feathers of gold, shining in the bright beams of the Most Holy and Light-giving Spirit. (Ps. lxvii. 14.) Hail, thou living garden of Eden, planted towards the East by the right hand of the Most Merciful and Mighty God, wherein do grow to His glory rich lilies and unfading roses, for the healing of them that have drunk in death from the blighting and pestilential breezes of the bitter West, (Gen. ii. 8, 9); Eden, wherein hath sprung that Tree of life, Whereof if any man eat he shall live for ever. (Gen. ii. 9; iii. 22. John vi. 52.) Hail, stately Palace of the King, most holy, stainless, purest, House of the Most High God, adorned with His Royal splendour, open to all, filled with Kingly dainties; Palace wherein is that spiritual bridal chamber, not made with hands, nor hung with diverse colours, in the which the Eternal Word, when He would raise up fallen man, wedded flesh unto Himself, that He might reconcile unto the Father them who had cast themselves away.\par
\switchcolumn*

\LA
\begin{Resp}\Rsign Hortus conclúsus soror mea sponsa, hortus conclúsus, fons signátus: \Star{} Emissiónes tuæ paradísus, o María.\end{Resp}
\begin{Resp}\Vsign Aperi mihi, soror mea, amíca mea, colúmba mea, immaculáta mea.\end{Resp}
\begin{Resp}\Rsign Emissiónes tuæ paradísus, o María.\end{Resp}
\switchcolumn
\EN
\begin{Resp}\Rsign A garden enclosed is my sister, my spouse, a garden enclosed, a fountain sealed. \Star{} O Mary, thy perfumes are a garden of delights.\end{Resp}
\begin{Resp}\Vsign Open to me, my sister, my love, my dove, my undefiled.\end{Resp}
\begin{Resp}\Rsign O Mary, thy perfumes are a garden of delights.\end{Resp}
\switchcolumn*

\LA
\LessonHead{Lectio VIII}
\switchcolumn
\EN
\LessonHead{Lesson VIII}
\switchcolumn*

\LA
\noindent Ave, Dei mons præpínguis et umbrósus, in quo enutrítus Agnus rationális peccáta atque infirmitátes nostras portávit: mons, e quo devolútus ille nulla manu præcísus lapis, contrívit aras idolórum, et factus est in caput ánguli, mirábilis in óculis nostris. Ave, sanctus Dei thronus, divínum donárium, domus glóriæ, perpúlchrum ornaméntum, cimélium eléctum, et totíus orbis propitiatórium, cælúmque Dei glóriam enárrans. Ave, urna ex puro auro confláta, et suavíssimam animárum nostrárum dulcédinem, Christum scílicet qui manna est, cóntinens. O puríssima et omni laude et obséquio digníssima Virgo, Deo dicátum donárium omni creaturárum conditióni præcéllens, terra non secta, inarátus ager, vitis floridíssima, fons aquas effúndens, virgo génerans, et mater viri néscia, innocéntiæ thesáurus abscónditus, et sanctimóniæ decus; acceptíssimis tuis ac matérna auctoritáte válidis précibus ad Dóminum ac Deum ómnium Conditórem, Fílium tuum ex te sine patre génitum, ecclesiástici órdinis gubernácula fac dírigas, et ad portum tranquíllum perdúcas.\par
\switchcolumn
\EN
\noindent Hail, O rich and shady Mountain of God, whereon pastured the True Lamb, Who hath taken away our sins and infirmities, (Hab. iii. 3; Isa. liii. 4; John i. 29,) mountain, whereout hath been cut without hands that Stone which hath smitten the altars of the idols, and become the head-stone of the corner, marvellous in our eyes. (Dan. ii. 34; Ps. cxvii. 22, 23.) Hail, thou holy Throne of God, thou divinest store-house, thou temple of glory, thou bright crown, thou chosen treasure, thou mercy-seat for the whole world, thou heaven declaring the glory of God. (Ps. xviii. 2.) Hail, thou vessel of pure gold, made to hold the manna that came down from heaven, the sweet food of our souls, even Christ. (Ex. xvi.33; Heb. ix. 4; John vi. 49-51.) Hail, O purest Virgin, most praiseworthy and most worshipful, hallowed treasury for the wants of all creatures; thou art the untitled earth, the unploughed field; thou art the vine full of flowers, the well overflowing with waters, Maiden and Mother; thou art the Mother that knew not a man, the hidden treasure of guilelessness, and the clear, bright star of holiness; by thy most acceptable prayers, strong from thy motherly mouth, obtain for all estates of men in the Church that they may continually tend unto Him Who is the Lord, and God, and Maker of thee, and of them, and of all, but of thee the Son also, conceived without man's intervention; obtain this, O Mother, pilot them to the harbour of peace.\par
\switchcolumn*

\LA
\begin{Resp}\Rsign Magníficat ánima mea Dóminum: \Star{} Quia fecit mihi magna qui potens est, et sanctum nomen ejus.\end{Resp}
\begin{Resp}\Vsign Ecce enim ex hoc beátam me dicent omnes generatiónes.\end{Resp}
\begin{Resp}\Rsign Quia fecit mihi magna qui potens est, et sanctum nomen ejus.\end{Resp}
\begin{Resp}\Vsign Glória Patri, et Fílio, \Star{} et Spirítui Sancto.\end{Resp}
\begin{Resp}\Rsign Quia fecit mihi magna qui potens est, et sanctum nomen ejus.\end{Resp}
\switchcolumn
\EN
\begin{Resp}\Rsign My soul doth magnify the Lord; \Star{} For He That is mighty hath done to me great things, and holy is His name.\end{Resp}
\begin{Resp}\Vsign For, behold, from henceforth all generations shall call me blessed.\end{Resp}
\begin{Resp}\Rsign For He That is mighty hath done to me great things, and holy is His name.\end{Resp}
\begin{Resp}\Vsign Glory be to the Father, and to the Son, \Star{} and to the Holy Ghost.\end{Resp}
\begin{Resp}\Rsign For He That is mighty hath done to me great things, and holy is His name.\end{Resp}
\switchcolumn*

\LA
\LessonHead{Lectio IX}
\switchcolumn
\EN
\LessonHead{Lesson IX}
\switchcolumn*

\LA
\noindent Sacerdótes justítia, et probátæ, immaculátæ ac sincéræ fídei exsultatióne splendidíssime indúito. Orthodóxis princípibus, qui præ omni púrpuræ aut auri splendóre, et præ margarítis ac lapídibus pretiósis, te nacti sunt diadéma et induméntum ac firmíssimum regni sui ornaméntum, in tranquíllo ac próspero statu sceptra dírige. Male fidas natiónes in te ac Deum ex te génitum blasphemántes, eórum pédibus sternens subícito; subjectúmque pópulum, ut secúndum Dei præcéptum in suávi obediéntiæ obséquio persevéret, confirmáto. Tuam hanc civitátem, quæ te tamquam turrim ac fundaméntum habet, victóriæ triúmphis coronáto, et fortitúdine circumcíngens custódito Dei habitatiónem, templi decórem semper conserváto; laudatóres tuos ab omni discrímine et ánimi angóre éxime; captívis redemptiónem tribúito; peregrínis tecto et quovis præsídio destitútis, solámen te éxhibe. Univérso mundo auxiliatrícem manum tuam pórrige, ut in lætítia et exsultatióne solemnitátes tuas simul cum ista, quam modo celebrámus, festivitáte splendidissímo éxitu transigámus, in Christo Jesu universórum Rege ac vero Deo nostro, cui glória et fortitúdo una cum sancto vitǽque princípio Patre, et coætérno et consubstantiáli et conregnánte Spíritu, nunc et semper et in sǽcula sæculórum. Amen.\par
\switchcolumn
\EN
\noindent Be it thine to clothe God's priests with righteousness, and to make them shout aloud for joy (Ps. cxxxi. 9, 16,) in approved and stainless, and upright and glorious faith. thine be it to guide in peace the sceptres of orthodox princes, even of princes who put their trust in thee to be the crown of their Majesty, and the Royal Robe of their greatness, and the firm foundation of their dominion, more than in purple, or fine gold, or pearls, or precious stones; thine be it to put under their feet the unfaithful nations, nations that blaspheme thee, and the God That was born of thee; thine be it to keep in meek obedience the people that are under them, according to the commandment of God. Behold, this is thine own city, which hath thee for her towers and her foundations, crown her with victory, gird the house of God with strength, keep undefiled the loveliness of His tabernacles, as for them that praise thy name, be thou their deliverer from strife and bitterness of spirit. Free thou the prisoner, protect the wanderer, and if there be any that hath no refuge, be thou to him a consolation. Stretch forth thine hand and help the whole earth so shall we year by year keep this and all thy feasts, and at last be found with thee in Christ Jesus, Who is Lord of all, and verily our God. To Him, with the Holy Father, Who is the Fountain of Life, and the coeternal Spirit, Three Persons and One Substance, even as there is one Kingdom, be glory and strength, now and for ever. Amen.\par
\switchcolumn*

\LA
\TeDeum{Te Deum}
\switchcolumn
\EN
\TeDeum{Te Deum}
\switchcolumn*

\end{paracol}

\par\needspace{10\baselineskip}\addvspace{1.2em}{\centering\small\textsc{Die 9 Decembris} · \textsc{9 December}\par}
\Office{De II die infra Octavam Concept. Immac. Beatæ Mariæ Virginis}{De II die infra Octavam Concept. Immac. Beatæ Mariæ Virginis}{Semiduplex}
\addcontentsline{toc}{subsection}{Die 9 Decembris · De II die infra Octavam Concept. Immac. Beatæ Mariæ Virginis\enspace\textit{De II die infra Octavam Concept. Immac. Beatæ Mariæ Virginis}}
\begin{paracol}{2}
\LA
\RefLine{Lectiones I–III: de Scriptura occurrente.}
\switchcolumn
\EN
\RefLine{Lessons I–III: from the Scripture of the day.}
\switchcolumn*

\LA
\LessonHead{Lectio IV}
\switchcolumn
\EN
\LessonHead{Lesson IV}
\switchcolumn*

\LA
\noindent Ex Bulla dogmatica Pii Papæ IX\par
\switchcolumn
\EN
\noindent God is unspeakable. His ways are mercy and truth; His Will is Almighty Power; and His wisdom reacheth mightily from one end to another, and sweetly ordereth all things. (Wisd. viii. 1.) He from all eternity foresaw the sorrowful fall of man by the transgression of Adam, and, in His mysterious purpose, He decreed, before the worlds were, that the Word should be made flesh, to the end that man, who had been seduced by the fraud of the devil, might not perish, but that as in the first Adam all die, in Christ all might be made alive. (1 Cor. xv. 22.) And to this end, the Eternal Creator from the beginning, and before all ages, chose and ordained a woman to be the Mother of His Only - begotten Son, of whom He should take flesh and be born, in the blessed fulness of time. (Gal. iv. 4.) And this woman He loved with so great a love that He allowed His Will to be freely wrought in her. (1 Thess. iv. 3.)\par
\switchcolumn*

\LA
\begin{Resp}\Rsign Ego ex ore Altíssimi prodívi, primogénita ante omnem creatúram: ego feci in cælis, ut orirétur lumen indefíciens. \Star{} Nondum erant abýssi, et ego jam concépta eram.\end{Resp}
\begin{Resp}\Vsign Deus enim creávit me in justítia, et apprehéndit manum meam, et servávit me.\end{Resp}
\begin{Resp}\Rsign Nondum erant abýssi, et ego jam concépta eram.\end{Resp}
\switchcolumn
\EN
\begin{Resp}\Rsign I came out of the mouth of the Most High, the first-begotten before every creature. I made the unfading light to arise in the heavens. \Star{} When there were no depths I was conceived.\end{Resp}
\begin{Resp}\Vsign For the Lord hath created me in righteousness, and hath held mine hand, and hath kept me.\end{Resp}
\begin{Resp}\Rsign When there were no depths I was conceived.\end{Resp}
\switchcolumn*

\LA
\LessonHead{Lectio V}
\switchcolumn
\EN
\LessonHead{Lesson V}
\switchcolumn*

\LA
\noindent Quapropter illam longe ante omnes angelicos spiritus, cunctosque Sanctos cælestium omnium charismatum copia de thesauro divinitatis deprompta ita mirifice cumulavit, ut ipsa ab omni prorsus peccati labe semper libera, ac tota pulchra et perfecta eam innocentiæ et sanctitatis plenitudinem præ se ferret, qua major sub Deo nullatenus intelligitur et quam præter Deum nemo assequi cogitando potest. Et quidem decebat omnino, ut perfectissimæ sanctitatis splendoribus semper ornata fulgeret, ac vel ab ipsa originalis culpæ labe plane immunis amplissimum de antiquo serpente triumphum referret tam venerabilis mater, cui Deus Pater unicum Filium suum, quem de corde suo æqualem sibi genitum, tamquam seipsum diligit, ita dare disposuit, ut naturaliter esset unus idemque communis Dei Patris et Virginis Filius.\par
\switchcolumn
\EN
\noindent Therefore, He bestowed upon her, out of the treasure of the Divinity, such a wealth of gifts of grace as He hath bestowed upon none of the Angels and none of the Saints. He made her always free from any the slightest pollution of sin, so fair and so upright that no other of His works are like to her, and only Himselfcan we understand to excel her. Verily, this was most fitting, that this most worshipful Mother should be made bright with the brightness of uncontaminated holiness, and should conquer the old serpent by escaping altogether the stain of original sin, for she was that Mother to whom the Eternal Father was willing to give the Co-Eternal and Co-Equal Onlybegotten Son of His love, to be her Son also.\par
\switchcolumn*

\LA
\begin{Resp}\Rsign Nihil inquinátum in eam incúrrit: \Star{} Candor est lucis ætérnæ et spéculum sine mácula.\end{Resp}
\begin{Resp}\Vsign Est enim hæc speciósior sole, et luci comparáta invenítur púrior.\end{Resp}
\begin{Resp}\Rsign Candor est lucis ætérnæ et spéculum sine mácula.\end{Resp}
\switchcolumn
\EN
\begin{Resp}\Rsign No defiled thing can fall into her; \Star{} She is the brightness of the everlasting light, and the unspotted mirror of the power of God.\end{Resp}
\begin{Resp}\Vsign For she is more beautiful than the sun, and being compared with the light, she is found before it.\end{Resp}
\begin{Resp}\Rsign She is the brightness of the everlasting light, and the unspotted mirror of the power of God.\end{Resp}
\switchcolumn*

\LA
\LessonHead{Lectio VI}
\switchcolumn
\EN
\LessonHead{Lesson VI}
\switchcolumn*

\LA
\noindent Quam originalem augustæ Virginis innocentiam cum admirabili ejusdem sanctitate præcelsaque Dei Matris dignitate omnino cohærentem catholica Ecclesia, quæ sancto semper edocta Spiritu columna est ac firmamentum veritatis, tamquam doctrinam possidens divinitus acceptam, et cælestis revelationis deposito comprehensam multiplici continenter ratione, splendidisque factis magis in dies explicare, proponere ac fovere numquam destitit Hanc enim doctrinam ab antiquissimis temporibus vigentem, ac fidelium animis penitus insitam, et sacrorum Antistitum curis studiisque per catholicum orbem mirifice propagatam, ipsa Ecclesia luculentissime significavit, cum ejusdem Virginis Conceptionem publico fidelium cultui ac venerationi proponere non dubitavit. Quo illustri quidem facto ipsius Virginis Conceptionem veluti singularem, miram et a reliquorum hominum primordiis longissime secretam, et omnino sanctam colendam exhibuit, cum Ecclesia nonnisi de Sanctis dies festos concelebret.\par
\switchcolumn
\EN
\noindent The Catholic Church, which, through the perpetual teaching of the Holy Ghost, is the pillar and ground of the truth, (1 Tim. iii. 15,) hath always held the original innocence of this most exalted Virgin to be bound up with her wonderful holiness, and her mighty dignity of Mother of God. This doctrine she hath felt herself to hold by the gift of God, and as part of that faith once delivered from heaven unto the Saints, (Jude 3,) and as time hath gone on, she hath continually explained, put forth, and upheld it. This belief is found strong in the earliest times, and rooted as it were in the hearts of Christ's faithful people; by the care and study of holy Bishops it hath been taught in all parts of the Catholic world; and the Church herself pointed to it when she allowed the Conception of the Blessed Virgin Mary to be held as a feast, for exciting the piety and devotion of her children. In the case of the Saints the Church celebrateth only the day of their being made perfect at death, (but of her Divine Lord, of His Blessed Mother, and of St. John the Baptist she venerateth the birth also, as of those sanctified in the womb.) When (therefore) she (goeth further and) maketh the case of the Blessed Virgin an exception to all others besides that of Christ, keeping holiday in honour of her conception (as well as of her birth,) it is manifest that she regardeth that Conception as altogether singular, wonderful, and different to all other conceptions, except only Christ's, namely, as holy.\par
\switchcolumn*

\LA
\begin{Resp}\Rsign Signum magnum appáruit in cælo: Múlier amícta sole, et luna sub pédibus ejus, \Star{} Et in cápite ejus coróna stellárum duódecim.\end{Resp}
\begin{Resp}\Vsign Induit eam Dóminus vestiméntis salútis, induménto justítiæ, et quasi sponsam ornávit eam monílibus suis.\end{Resp}
\begin{Resp}\Rsign Et in cápite ejus coróna stellárum duódecim.\end{Resp}
\begin{Resp}\Vsign Glória Patri, et Fílio, \Star{} et Spirítui Sancto.\end{Resp}
\begin{Resp}\Rsign Et in cápite ejus coróna stellárum duódecim.\end{Resp}
\switchcolumn
\EN
\begin{Resp}\Rsign There appeared a great wonder in heaven; a Woman clothed with the sun, and the moon under her feet \Star{} And upon her head a crown of twelve stars.\end{Resp}
\begin{Resp}\Vsign The Lord hath clothed her with the garments of salvation, and hath covered her with the robe of righteousness, yea, as a bride He hath adorned her with jewels.\end{Resp}
\begin{Resp}\Rsign And upon her head a crown of twelve stars.\end{Resp}
\begin{Resp}\Vsign Glory be to the Father, and to the Son, \Star{} and to the Holy Ghost.\end{Resp}
\begin{Resp}\Rsign And upon her head a crown of twelve stars.\end{Resp}
\switchcolumn*

\LA
\LessonHead{Lectio VII}
\switchcolumn
\EN
\LessonHead{Lesson VII}
\switchcolumn*

\LA
\LessonTitle{Léctio sancti Evangélii secúndum Lucam}
\switchcolumn
\EN
\LessonTitle{From the Holy Gospel according to Luke}
\switchcolumn*

\LA
\Cite{Luc 1:26-28}
\switchcolumn
\EN
\Cite{Luke 1:26-28}
\switchcolumn*

\LA
\noindent In illo témpore: Missus est Angelus Gabriel a Deo in civitatem Galilææ, cui nomen Nazareth, ad Virginem desponsatam viro, cui nomen erat Joseph, de domo David, et nomen Virginis Maria. Et réliqua.\par
\switchcolumn
\EN
\noindent In that time, the angel Gabriel was sent from God into a city of Galilee, called Nazareth, To a virgin espoused to a man whose name was Joseph, of the house of David; and the virgin's name was Mary. And so on.\par
\switchcolumn*

\LA
\LessonTitle{Homilía sancti Sophrónii Epíscopi}
\switchcolumn
\EN
\LessonTitle{Homily by St. Sophronius, Patriarch (of Jerusalem.)}
\switchcolumn*

\LA
\Source{Homilia in Deiparæ Annunt.}
\switchcolumn
\EN
\Source{On the Annunciation.}
\switchcolumn*

\LA
\noindent Quid missus beatus ille Angelus ad Virginem integerrimam dicit? Aut quomodo faustissimum hoc nuntium ipsi defert? Ave gratia plena, Dominus tecum. A gaudio incipit eam alloqui ille gaudii nuntius. Noverat enim et plane sciebat, nuntium illud suum universis hominibus atque omnibus pariter creaturis gaudium parare, et quoslibet a quibuscumque dolores expellere; noverat, ex divina hujus mysterii cognitione mundum lumine collustrari; noverat, erroris disjici caliginem; noverat, retundi mortis aculeum; noverat, vim corruptionis infringi; noverat, inferno victoriam auferri; noverat, salutem perdito affulgere homini, qui horum malorum jugo jamdiu premebatur, ex quo scilicet a paradisi deliciis expulsus, et a beato illo domicilio ejectus fuerat. Propterea legationis suæ exordium a gaudio ducit; propterea sermonibus suis gaudii voces præmittit; propterea faustis hisce nuntiis gaudium antecedit, utpote quæ omnibus credentibus gaudio futura erant.\par
\switchcolumn
\EN
\noindent When this blessed Angel was sent to the most pure virgin what did he say? In what words did he break the happy news of Redemption? Hail, thou that art full of grace, the Lord is with thee. (Now this word Hail is in the original Chaire, which being interpreted signifieth Rejoice.) The messenger of joy in his first word biddeth her rejoice. He knew well that his message was a message of good tidings of great joy to men, (Luke ii. 10,) yea, to all creatures, a message of healing to all sicknesses. He knew well that his message was a message of God's light to a dark world. He knew well that it proclaimed the end of error. He knew well that it blunted the sting of death. He knew well that it broke the power of corruption. He knew well that it brought victory over hell. He knew well that it told of salvation to all the fallen children of Adam, groaning under that yoke of malediction which fell on them when they were thrust out of Eden, and banished from that happy home. Therefore, when he began to speak, he spoke in tones of rejoicing, and opened his message with sounds of gladness. Therefore made he the name of joy to herald the tidings of good, which were to be for a joy unto all people, whosoever should believe.\par
\switchcolumn*

\LA
\begin{Resp}\Rsign Hortus conclúsus soror mea sponsa, hortus conclúsus, fons signátus: \Star{} Emissiónes tuæ paradísus, o María.\end{Resp}
\begin{Resp}\Vsign Aperi mihi, soror mea, amíca mea, colúmba mea, immaculáta mea.\end{Resp}
\begin{Resp}\Rsign Emissiónes tuæ paradísus, o María.\end{Resp}
\switchcolumn
\EN
\begin{Resp}\Rsign A garden enclosed is my sister, my spouse, a garden enclosed, a fountain sealed. \Star{} O Mary, thy perfumes are a garden of delights.\end{Resp}
\begin{Resp}\Vsign Open to me, my sister, my love, my dove, my undefiled.\end{Resp}
\begin{Resp}\Rsign O Mary, thy perfumes are a garden of delights.\end{Resp}
\switchcolumn*

\LA
\LessonHead{Lectio VIII}
\switchcolumn
\EN
\LessonHead{Lesson VIII}
\switchcolumn*

\LA
\noindent Et sane par omnino erat, ut divina gaudii denuntiatio a sermonibus verbisque gaudium elicientibus sumeret initium. Propterea enim et Angelus gaudium ante omnia renuntiat, quia faustæ legationis suæ non ignorat exitum, ac probe novit colloquium quod habebatur, in totius mundi gaudium manifeste esse cessurum. Et profecto quodnam gaudium, aut quænam reperiri potest jucunditas, quam non longe excedat alloquium ad Virginem illam beatam ac gaudii parentem habitum? Gaude, o supercælestis gaudii genitrix. Gaude, o sublimissimi gaudii nutrix. Gaude, o salutaris gaudii sedes princeps. Gaude o immortalis gaudii auctrix. Gaude o ineffabilis gaudii mysticum diversorium. Gaude, o indeficientis gaudii fons beatissime. Gaude, o gaudii æterni deiferum cimelium. Gaude, o vivificantis gaudii arbor virentissima. Gaude o innupta Dei Mater. Gaude, o Virgo post partum integerrima. Gaude, o spectaculum præ mirabilibus omnibus summe admirandum.\par
\switchcolumn
\EN
\noindent And, of a truth, it was fitting that God's proclamation of joy should open with the accents of gladness. And this is the reason why the angel nameth joy first, because he knew the coming fruits of his message, and that his converse with the Virgin was to bring joy to the whole world. Can we find any joy or any brightness like the joy and the brightness of that salutation addressed to the Blessed Mother of gladness? Rejoice, O mother of joy more than heavenly! Rejoice, O thou that nourishest joy in the highest! Rejoice, O Lady, full of.the joy of salvation! Rejoice, O thou that bringest a joy that passeth not away! Rejoice, O mysterious treasury dispensing unspeakable joy! Rejoice, O most blessed fountain, overflowing with unfailing joy! Rejoice, O store-house of God, filled with the everlasting joy of eternity! Rejoice, O fair tree, bearing fruit of life-giving joy! Rejoice, O Maiden Mother of God! Rejoice, O thou that after child-birth remainest a virgin! Rejoice, O wonder, who, after all wonders, art still the most wonderful!\par
\switchcolumn*

\LA
\begin{Resp}\Rsign Magníficat ánima mea Dóminum: \Star{} Quia fecit mihi magna qui potens est, et sanctum nomen ejus.\end{Resp}
\begin{Resp}\Vsign Ecce enim ex hoc beátam me dicent omnes generatiónes.\end{Resp}
\begin{Resp}\Rsign Quia fecit mihi magna qui potens est, et sanctum nomen ejus.\end{Resp}
\begin{Resp}\Vsign Glória Patri, et Fílio, \Star{} et Spirítui Sancto.\end{Resp}
\begin{Resp}\Rsign Quia fecit mihi magna qui potens est, et sanctum nomen ejus.\end{Resp}
\switchcolumn
\EN
\begin{Resp}\Rsign My soul doth magnify the Lord; \Star{} For He That is mighty hath done to me great things, and holy is His name.\end{Resp}
\begin{Resp}\Vsign For, behold, from henceforth all generations shall call me blessed.\end{Resp}
\begin{Resp}\Rsign For He That is mighty hath done to me great things, and holy is His name.\end{Resp}
\begin{Resp}\Vsign Glory be to the Father, and to the Son, \Star{} and to the Holy Ghost.\end{Resp}
\begin{Resp}\Rsign For He That is mighty hath done to me great things, and holy is His name.\end{Resp}
\switchcolumn*

\LA
\LessonHead{Lectio IX}
\switchcolumn
\EN
\LessonHead{Lesson IX}
\switchcolumn*

\LA
\noindent Quisnam tuum eloqui splendorem poterit? Quisnam portentum, quod ipsa es, enarrare verbis audeat? Quisnam magnificentiam tuam effari se posse confidet? Tu hominum exornasti naturam; tu Angelorum ordines superasti; tu fulgores Archangelorum obtenebrasti; tu sublimes Thronorum sedes infra te ostendisti; tu altitudinem Dominationum depressisti; tu Principatuum ducatibus præcucurristi; tu enervasti fortitudinem Potestatum; tu ipsis Virtutibus potentior virtus prodiisti; tu Cherubim oculatissimum visum terrestribus oculis vicisti; tu Seraphim sex alas habentium volatus animæ pennis divinitus agitatis transvolasti tu denique omnem creaturam longe transgressa es; quippe quæ præ omni creatura enituisti puritate; et omnium creaturarum Conditorem in te excepisti; ipsumque et sinu tuo gestasti, et genuisti; et sola ex omnibus creaturis Dei Mater effecta es.\par
\switchcolumn
\EN
\noindent Who shall worthily set forth thy glory? Who shall make bold to say what thou art? Who will hold himself able to tell of all thy splendour? Thou art the exaltation of humanity; thou art made much higher than the Angels; thy brightness hath thrown the brightness of the Archangels into shadow; thou lookest down upon the lofty seats of the Thrones; thou makest the height of the Lordships to seem low; thy rank taketh precedence before the rank of the Principalities; compared with thee the Powers are weakness; thou art a Mighty one mightier than all the Mighty; thine earthly eyes see further than the contemplation of the Cherubim can reach; the Seraphim have six wings, but thy flight is nobler than theirs; in a word, thou hast far excelled every other work of God; thou wast far purer than any other creature; and thou hast conceived the Creator of all creatures, carried Him in thy womb, and brought Him forth; thou hast been chosen, out of all that He has made, to be His mother.\par
\switchcolumn*

\LA
\TeDeum{Te Deum}
\switchcolumn
\EN
\TeDeum{Te Deum}
\switchcolumn*

\end{paracol}

\par\needspace{10\baselineskip}\addvspace{1.2em}{\centering\small\textsc{Die 9 Decembris} · \textsc{9 December}\par}
\Office{De II die infra Octavam Concept. Immac. Beatæ Mariæ Virginis}{De II die infra Octavam Concept. Immac. Beatæ Mariæ Virginis}{Semiduplex}
\addcontentsline{toc}{subsection}{Die 9 Decembris · De II die infra Octavam Concept. Immac. Beatæ Mariæ Virginis\enspace\textit{De II die infra Octavam Concept. Immac. Beatæ Mariæ Virginis}}
\begin{paracol}{2}
\LA
\RefLine{Lectiones I–III: de Scriptura occurrente.}
\switchcolumn
\EN
\RefLine{Lessons I–III: from the Scripture of the day.}
\switchcolumn*

\LA
\LessonHead{Lectio IV}
\switchcolumn
\EN
\LessonHead{Lesson IV}
\switchcolumn*

\LA
\noindent Ex Bulla dogmatica Pii Papæ IX\par
\switchcolumn
\EN
\noindent God is unspeakable. His ways are mercy and truth; His Will is Almighty Power; and His wisdom reacheth mightily from one end to another, and sweetly ordereth all things. (Wisd. viii. 1.) He from all eternity foresaw the sorrowful fall of man by the transgression of Adam, and, in His mysterious purpose, He decreed, before the worlds were, that the Word should be made flesh, to the end that man, who had been seduced by the fraud of the devil, might not perish, but that as in the first Adam all die, in Christ all might be made alive. (1 Cor. xv. 22.) And to this end, the Eternal Creator from the beginning, and before all ages, chose and ordained a woman to be the Mother of His Only - begotten Son, of whom He should take flesh and be born, in the blessed fulness of time. (Gal. iv. 4.) And this woman He loved with so great a love that He allowed His Will to be freely wrought in her. (1 Thess. iv. 3.)\par
\switchcolumn*

\LA
\begin{Resp}\Rsign Ego ex ore Altíssimi prodívi, primogénita ante omnem creatúram: ego feci in cælis, ut orirétur lumen indefíciens. \Star{} Nondum erant abýssi, et ego jam concépta eram.\end{Resp}
\begin{Resp}\Vsign Deus enim creávit me in justítia, et apprehéndit manum meam, et servávit me.\end{Resp}
\begin{Resp}\Rsign Nondum erant abýssi, et ego jam concépta eram.\end{Resp}
\switchcolumn
\EN
\begin{Resp}\Rsign I came out of the mouth of the Most High, the first-begotten before every creature. I made the unfading light to arise in the heavens. \Star{} When there were no depths I was conceived.\end{Resp}
\begin{Resp}\Vsign For the Lord hath created me in righteousness, and hath held mine hand, and hath kept me.\end{Resp}
\begin{Resp}\Rsign When there were no depths I was conceived.\end{Resp}
\switchcolumn*

\LA
\LessonHead{Lectio V}
\switchcolumn
\EN
\LessonHead{Lesson V}
\switchcolumn*

\LA
\noindent Quapropter illam longe ante omnes angelicos spiritus, cunctosque Sanctos cælestium omnium charismatum copia de thesauro divinitatis deprompta ita mirifice cumulavit, ut ipsa ab omni prorsus peccati labe semper libera, ac tota pulchra et perfecta eam innocentiæ et sanctitatis plenitudinem præ se ferret, qua major sub Deo nullatenus intelligitur et quam præter Deum nemo assequi cogitando potest. Et quidem decebat omnino, ut perfectissimæ sanctitatis splendoribus semper ornata fulgeret, ac vel ab ipsa originalis culpæ labe plane immunis amplissimum de antiquo serpente triumphum referret tam venerabilis mater, cui Deus Pater unicum Filium suum, quem de corde suo æqualem sibi genitum, tamquam seipsum diligit, ita dare disposuit, ut naturaliter esset unus idemque communis Dei Patris et Virginis Filius.\par
\switchcolumn
\EN
\noindent Therefore, He bestowed upon her, out of the treasure of the Divinity, such a wealth of gifts of grace as He hath bestowed upon none of the Angels and none of the Saints. He made her always free from any the slightest pollution of sin, so fair and so upright that no other of His works are like to her, and only Himselfcan we understand to excel her. Verily, this was most fitting, that this most worshipful Mother should be made bright with the brightness of uncontaminated holiness, and should conquer the old serpent by escaping altogether the stain of original sin, for she was that Mother to whom the Eternal Father was willing to give the Co-Eternal and Co-Equal Onlybegotten Son of His love, to be her Son also.\par
\switchcolumn*

\LA
\begin{Resp}\Rsign Nihil inquinátum in eam incúrrit: \Star{} Candor est lucis ætérnæ et spéculum sine mácula.\end{Resp}
\begin{Resp}\Vsign Est enim hæc speciósior sole, et luci comparáta invenítur púrior.\end{Resp}
\begin{Resp}\Rsign Candor est lucis ætérnæ et spéculum sine mácula.\end{Resp}
\switchcolumn
\EN
\begin{Resp}\Rsign No defiled thing can fall into her; \Star{} She is the brightness of the everlasting light, and the unspotted mirror of the power of God.\end{Resp}
\begin{Resp}\Vsign For she is more beautiful than the sun, and being compared with the light, she is found before it.\end{Resp}
\begin{Resp}\Rsign She is the brightness of the everlasting light, and the unspotted mirror of the power of God.\end{Resp}
\switchcolumn*

\LA
\LessonHead{Lectio VI}
\switchcolumn
\EN
\LessonHead{Lesson VI}
\switchcolumn*

\LA
\noindent Quam originalem augustæ Virginis innocentiam cum admirabili ejusdem sanctitate præcelsaque Dei Matris dignitate omnino cohærentem catholica Ecclesia, quæ sancto semper edocta Spiritu columna est ac firmamentum veritatis, tamquam doctrinam possidens divinitus acceptam, et cælestis revelationis deposito comprehensam multiplici continenter ratione, splendidisque factis magis in dies explicare, proponere ac fovere numquam destitit Hanc enim doctrinam ab antiquissimis temporibus vigentem, ac fidelium animis penitus insitam, et sacrorum Antistitum curis studiisque per catholicum orbem mirifice propagatam, ipsa Ecclesia luculentissime significavit, cum ejusdem Virginis Conceptionem publico fidelium cultui ac venerationi proponere non dubitavit. Quo illustri quidem facto ipsius Virginis Conceptionem veluti singularem, miram et a reliquorum hominum primordiis longissime secretam, et omnino sanctam colendam exhibuit, cum Ecclesia nonnisi de Sanctis dies festos concelebret.\par
\switchcolumn
\EN
\noindent The Catholic Church, which, through the perpetual teaching of the Holy Ghost, is the pillar and ground of the truth, (1 Tim. iii. 15,) hath always held the original innocence of this most exalted Virgin to be bound up with her wonderful holiness, and her mighty dignity of Mother of God. This doctrine she hath felt herself to hold by the gift of God, and as part of that faith once delivered from heaven unto the Saints, (Jude 3,) and as time hath gone on, she hath continually explained, put forth, and upheld it. This belief is found strong in the earliest times, and rooted as it were in the hearts of Christ's faithful people; by the care and study of holy Bishops it hath been taught in all parts of the Catholic world; and the Church herself pointed to it when she allowed the Conception of the Blessed Virgin Mary to be held as a feast, for exciting the piety and devotion of her children. In the case of the Saints the Church celebrateth only the day of their being made perfect at death, (but of her Divine Lord, of His Blessed Mother, and of St. John the Baptist she venerateth the birth also, as of those sanctified in the womb.) When (therefore) she (goeth further and) maketh the case of the Blessed Virgin an exception to all others besides that of Christ, keeping holiday in honour of her conception (as well as of her birth,) it is manifest that she regardeth that Conception as altogether singular, wonderful, and different to all other conceptions, except only Christ's, namely, as holy.\par
\switchcolumn*

\LA
\begin{Resp}\Rsign Signum magnum appáruit in cælo: Múlier amícta sole, et luna sub pédibus ejus, \Star{} Et in cápite ejus coróna stellárum duódecim.\end{Resp}
\begin{Resp}\Vsign Induit eam Dóminus vestiméntis salútis, induménto justítiæ, et quasi sponsam ornávit eam monílibus suis.\end{Resp}
\begin{Resp}\Rsign Et in cápite ejus coróna stellárum duódecim.\end{Resp}
\begin{Resp}\Vsign Glória Patri, et Fílio, \Star{} et Spirítui Sancto.\end{Resp}
\begin{Resp}\Rsign Et in cápite ejus coróna stellárum duódecim.\end{Resp}
\switchcolumn
\EN
\begin{Resp}\Rsign There appeared a great wonder in heaven; a Woman clothed with the sun, and the moon under her feet \Star{} And upon her head a crown of twelve stars.\end{Resp}
\begin{Resp}\Vsign The Lord hath clothed her with the garments of salvation, and hath covered her with the robe of righteousness, yea, as a bride He hath adorned her with jewels.\end{Resp}
\begin{Resp}\Rsign And upon her head a crown of twelve stars.\end{Resp}
\begin{Resp}\Vsign Glory be to the Father, and to the Son, \Star{} and to the Holy Ghost.\end{Resp}
\begin{Resp}\Rsign And upon her head a crown of twelve stars.\end{Resp}
\switchcolumn*

\LA
\LessonHead{Lectio VII}
\switchcolumn
\EN
\LessonHead{Lesson VII}
\switchcolumn*

\LA
\LessonTitle{Léctio sancti Evangélii secúndum Lucam}
\switchcolumn
\EN
\LessonTitle{From the Holy Gospel according to Luke}
\switchcolumn*

\LA
\Cite{Luc 1:26-28}
\switchcolumn
\EN
\Cite{Luke 1:26-28}
\switchcolumn*

\LA
\noindent In illo témpore: Missus est Angelus Gabriel a Deo in civitatem Galilææ, cui nomen Nazareth, ad Virginem desponsatam viro, cui nomen erat Joseph, de domo David, et nomen Virginis Maria. Et réliqua.\par
\switchcolumn
\EN
\noindent In that time, the angel Gabriel was sent from God into a city of Galilee, called Nazareth, To a virgin espoused to a man whose name was Joseph, of the house of David; and the virgin's name was Mary. And so on.\par
\switchcolumn*

\LA
\LessonTitle{Homilía sancti Sophrónii Epíscopi}
\switchcolumn
\EN
\LessonTitle{Homily by St. Sophronius, Patriarch (of Jerusalem.)}
\switchcolumn*

\LA
\Source{Homilia in Deiparæ Annunt.}
\switchcolumn
\EN
\Source{On the Annunciation.}
\switchcolumn*

\LA
\noindent Quid missus beatus ille Angelus ad Virginem integerrimam dicit? Aut quomodo faustissimum hoc nuntium ipsi defert? Ave gratia plena, Dominus tecum. A gaudio incipit eam alloqui ille gaudii nuntius. Noverat enim et plane sciebat, nuntium illud suum universis hominibus atque omnibus pariter creaturis gaudium parare, et quoslibet a quibuscumque dolores expellere; noverat, ex divina hujus mysterii cognitione mundum lumine collustrari; noverat, erroris disjici caliginem; noverat, retundi mortis aculeum; noverat, vim corruptionis infringi; noverat, inferno victoriam auferri; noverat, salutem perdito affulgere homini, qui horum malorum jugo jamdiu premebatur, ex quo scilicet a paradisi deliciis expulsus, et a beato illo domicilio ejectus fuerat. Propterea legationis suæ exordium a gaudio ducit; propterea sermonibus suis gaudii voces præmittit; propterea faustis hisce nuntiis gaudium antecedit, utpote quæ omnibus credentibus gaudio futura erant.\par
\switchcolumn
\EN
\noindent When this blessed Angel was sent to the most pure virgin what did he say? In what words did he break the happy news of Redemption? Hail, thou that art full of grace, the Lord is with thee. (Now this word Hail is in the original Chaire, which being interpreted signifieth Rejoice.) The messenger of joy in his first word biddeth her rejoice. He knew well that his message was a message of good tidings of great joy to men, (Luke ii. 10,) yea, to all creatures, a message of healing to all sicknesses. He knew well that his message was a message of God's light to a dark world. He knew well that it proclaimed the end of error. He knew well that it blunted the sting of death. He knew well that it broke the power of corruption. He knew well that it brought victory over hell. He knew well that it told of salvation to all the fallen children of Adam, groaning under that yoke of malediction which fell on them when they were thrust out of Eden, and banished from that happy home. Therefore, when he began to speak, he spoke in tones of rejoicing, and opened his message with sounds of gladness. Therefore made he the name of joy to herald the tidings of good, which were to be for a joy unto all people, whosoever should believe.\par
\switchcolumn*

\LA
\begin{Resp}\Rsign Hortus conclúsus soror mea sponsa, hortus conclúsus, fons signátus: \Star{} Emissiónes tuæ paradísus, o María.\end{Resp}
\begin{Resp}\Vsign Aperi mihi, soror mea, amíca mea, colúmba mea, immaculáta mea.\end{Resp}
\begin{Resp}\Rsign Emissiónes tuæ paradísus, o María.\end{Resp}
\switchcolumn
\EN
\begin{Resp}\Rsign A garden enclosed is my sister, my spouse, a garden enclosed, a fountain sealed. \Star{} O Mary, thy perfumes are a garden of delights.\end{Resp}
\begin{Resp}\Vsign Open to me, my sister, my love, my dove, my undefiled.\end{Resp}
\begin{Resp}\Rsign O Mary, thy perfumes are a garden of delights.\end{Resp}
\switchcolumn*

\LA
\LessonHead{Lectio VIII}
\switchcolumn
\EN
\LessonHead{Lesson VIII}
\switchcolumn*

\LA
\noindent Et sane par omnino erat, ut divina gaudii denuntiatio a sermonibus verbisque gaudium elicientibus sumeret initium. Propterea enim et Angelus gaudium ante omnia renuntiat, quia faustæ legationis suæ non ignorat exitum, ac probe novit colloquium quod habebatur, in totius mundi gaudium manifeste esse cessurum. Et profecto quodnam gaudium, aut quænam reperiri potest jucunditas, quam non longe excedat alloquium ad Virginem illam beatam ac gaudii parentem habitum? Gaude, o supercælestis gaudii genitrix. Gaude, o sublimissimi gaudii nutrix. Gaude, o salutaris gaudii sedes princeps. Gaude o immortalis gaudii auctrix. Gaude o ineffabilis gaudii mysticum diversorium. Gaude, o indeficientis gaudii fons beatissime. Gaude, o gaudii æterni deiferum cimelium. Gaude, o vivificantis gaudii arbor virentissima. Gaude o innupta Dei Mater. Gaude, o Virgo post partum integerrima. Gaude, o spectaculum præ mirabilibus omnibus summe admirandum.\par
\switchcolumn
\EN
\noindent And, of a truth, it was fitting that God's proclamation of joy should open with the accents of gladness. And this is the reason why the angel nameth joy first, because he knew the coming fruits of his message, and that his converse with the Virgin was to bring joy to the whole world. Can we find any joy or any brightness like the joy and the brightness of that salutation addressed to the Blessed Mother of gladness? Rejoice, O mother of joy more than heavenly! Rejoice, O thou that nourishest joy in the highest! Rejoice, O Lady, full of.the joy of salvation! Rejoice, O thou that bringest a joy that passeth not away! Rejoice, O mysterious treasury dispensing unspeakable joy! Rejoice, O most blessed fountain, overflowing with unfailing joy! Rejoice, O store-house of God, filled with the everlasting joy of eternity! Rejoice, O fair tree, bearing fruit of life-giving joy! Rejoice, O Maiden Mother of God! Rejoice, O thou that after child-birth remainest a virgin! Rejoice, O wonder, who, after all wonders, art still the most wonderful!\par
\switchcolumn*

\LA
\begin{Resp}\Rsign Magníficat ánima mea Dóminum: \Star{} Quia fecit mihi magna qui potens est, et sanctum nomen ejus.\end{Resp}
\begin{Resp}\Vsign Ecce enim ex hoc beátam me dicent omnes generatiónes.\end{Resp}
\begin{Resp}\Rsign Quia fecit mihi magna qui potens est, et sanctum nomen ejus.\end{Resp}
\begin{Resp}\Vsign Glória Patri, et Fílio, \Star{} et Spirítui Sancto.\end{Resp}
\begin{Resp}\Rsign Quia fecit mihi magna qui potens est, et sanctum nomen ejus.\end{Resp}
\switchcolumn
\EN
\begin{Resp}\Rsign My soul doth magnify the Lord; \Star{} For He That is mighty hath done to me great things, and holy is His name.\end{Resp}
\begin{Resp}\Vsign For, behold, from henceforth all generations shall call me blessed.\end{Resp}
\begin{Resp}\Rsign For He That is mighty hath done to me great things, and holy is His name.\end{Resp}
\begin{Resp}\Vsign Glory be to the Father, and to the Son, \Star{} and to the Holy Ghost.\end{Resp}
\begin{Resp}\Rsign For He That is mighty hath done to me great things, and holy is His name.\end{Resp}
\switchcolumn*

\LA
\LessonHead{Lectio IX}
\switchcolumn
\EN
\LessonHead{Lesson IX}
\switchcolumn*

\LA
\noindent Quisnam tuum eloqui splendorem poterit? Quisnam portentum, quod ipsa es, enarrare verbis audeat? Quisnam magnificentiam tuam effari se posse confidet? Tu hominum exornasti naturam; tu Angelorum ordines superasti; tu fulgores Archangelorum obtenebrasti; tu sublimes Thronorum sedes infra te ostendisti; tu altitudinem Dominationum depressisti; tu Principatuum ducatibus præcucurristi; tu enervasti fortitudinem Potestatum; tu ipsis Virtutibus potentior virtus prodiisti; tu Cherubim oculatissimum visum terrestribus oculis vicisti; tu Seraphim sex alas habentium volatus animæ pennis divinitus agitatis transvolasti tu denique omnem creaturam longe transgressa es; quippe quæ præ omni creatura enituisti puritate; et omnium creaturarum Conditorem in te excepisti; ipsumque et sinu tuo gestasti, et genuisti; et sola ex omnibus creaturis Dei Mater effecta es.\par
\switchcolumn
\EN
\noindent Who shall worthily set forth thy glory? Who shall make bold to say what thou art? Who will hold himself able to tell of all thy splendour? Thou art the exaltation of humanity; thou art made much higher than the Angels; thy brightness hath thrown the brightness of the Archangels into shadow; thou lookest down upon the lofty seats of the Thrones; thou makest the height of the Lordships to seem low; thy rank taketh precedence before the rank of the Principalities; compared with thee the Powers are weakness; thou art a Mighty one mightier than all the Mighty; thine earthly eyes see further than the contemplation of the Cherubim can reach; the Seraphim have six wings, but thy flight is nobler than theirs; in a word, thou hast far excelled every other work of God; thou wast far purer than any other creature; and thou hast conceived the Creator of all creatures, carried Him in thy womb, and brought Him forth; thou hast been chosen, out of all that He has made, to be His mother.\par
\switchcolumn*

\LA
\TeDeum{Te Deum}
\switchcolumn
\EN
\TeDeum{Te Deum}
\switchcolumn*

\end{paracol}

\par\needspace{10\baselineskip}\addvspace{1.2em}{\centering\small\textsc{Die 10 Decembris} · \textsc{10 December}\par}
\Office{De III die infra Octavam Conceptionis Immaculatæ Beatæ Mariæ Virginis}{De III die infra Octavam Conceptionis Immaculatæ Beatæ Mariæ Virginis}{Semiduplex}
\addcontentsline{toc}{subsection}{Die 10 Decembris · De III die infra Octavam Conceptionis Immaculatæ Beatæ Mariæ Virginis\enspace\textit{De III die infra Octavam Conceptionis Immaculatæ Beatæ Mariæ Virginis}}
\begin{paracol}{2}
\LA
\RefLine{Lectiones I–III: de Scriptura occurrente.}
\switchcolumn
\EN
\RefLine{Lessons I–III: from the Scripture of the day.}
\switchcolumn*

\LA
\LessonHead{Lectio IV}
\switchcolumn
\EN
\LessonHead{Lesson IV}
\switchcolumn*

\LA
\noindent Ex Bulla dogmatica Pii Papæ IX\par
\switchcolumn
\EN
\noindent Both in her Offices and in the most holy Liturgy the Church hath been accustomed to apply to the creation of Mary the language in which the Holy Scriptures set forth the Eternal Generation of the Uncreated Wisdom, and that, because Mary was predestined in the decree of the Incarnation of the same Wisdom. This practice hath been received by the faithful in all quarters, and plainly showeth what hath been the mind of the Church of Rome, which is the mother and mistress of all Churches, on the subject of the sinless conception of the Virgin. Nevertheless, it is fitting to set forth in greater detail the celebrated acts of this Church, on account of that pre-eminent rank and power which all other Churches are bound to yield her, because she is the centre of Catholic truth and unity, wherein alone Doctrine is always preserved pure, and from whom all the other Churches must needs receive the tradition of the Faith.\par
\switchcolumn*

\LA
\begin{Resp}\Rsign Ego ex ore Altíssimi prodívi, primogénita ante omnem creatúram: ego feci in cælis, ut orirétur lumen indefíciens. \Star{} Nondum erant abýssi, et ego jam concépta eram.\end{Resp}
\begin{Resp}\Vsign Deus enim creávit me in justítia, et apprehéndit manum meam, et servávit me.\end{Resp}
\begin{Resp}\Rsign Nondum erant abýssi, et ego jam concépta eram.\end{Resp}
\switchcolumn
\EN
\begin{Resp}\Rsign I came out of the mouth of the Most High, the first-begotten before every creature. I made the unfading light to arise in the heavens. \Star{} When there were no depths I was conceived.\end{Resp}
\begin{Resp}\Vsign For the Lord hath created me in righteousness, and hath held mine hand, and hath kept me.\end{Resp}
\begin{Resp}\Rsign When there were no depths I was conceived.\end{Resp}
\switchcolumn*

\LA
\LessonHead{Lectio V}
\switchcolumn
\EN
\LessonHead{Lesson V}
\switchcolumn*

\LA
\noindent Itaque eadem Romana Ecclesia nihil potius habuit, quam eloquentissimis quibusque modis Immaculatam Virginis Conceptionem, ejusque cultum et doctrinam asserere, tueri, promovere et vindicare. Enim vero prædecessores nostri vehementer gloriati sunt apostolica sua auctoritate festum Conceptionis in Romana Ecclesia instituere, ac proprio Officio, propriaque Missa, quibus prærogativa immunitatis ab hereditaria labe manifestissime asserebatur, augere, honestare et cultum jam institutum omni ope promovere, amplificare, sive erogatis indulgentiis, sive facultate tributa civitatibus, provinciis regnisque, ut Deiparam sub titulo Immaculatæ Conceptionis patronam sibi deligerent, sive comprobatis sodalitatibus congregationibus, religiosisque familiis ad Immaculatæ Conceptionis honorem institutis, sive laudibus eorum pietati delatis, qui monasteria, xenodochia, altaria, templa sub Immaculati Conceptus titulo erexerint, aut sacramenti religione interposita Immaculatam Deiparæ Conceptionem strenue propugnare spoponderint.\par
\switchcolumn
\EN
\noindent Thus it hath always been one of the most striking features of the Roman Church that she hath most powerfully asserted, guarded, promoted, and vindicated the doctrine that the Virgin was conceived without sin. It hath been the boast of Our Predecessors that by their authority they instituted in the Roman Church the Feast of the Conception of Mary, and caused it to be observed with an Office and a Mass wherein her privilege of immunity from original sin was openly asserted. Our said Predecessors have done everything in their power to increase the love of the faithful for this doctrine by granting Indulgences in its honour; by giving permission to cities, provinces, and kingdoms to choose for their Patroness the Mother of God, under her title Conceived without sin; by approving of Guilds, Congregations, and Associations of persons under vows, all instituted in honour of the sinless Conception; by praising the piety of those who have founded Convents, Hospitals, Altars and Churches named from this belief; and lastly, by encouraging those who have taken an oath to defend this opinion to the utmost of their power.\par
\switchcolumn*

\LA
\begin{Resp}\Rsign Nihil inquinátum in eam incúrrit: \Star{} Candor est lucis ætérnæ et spéculum sine mácula.\end{Resp}
\begin{Resp}\Vsign Est enim hæc speciósior sole, et luci comparáta invenítur púrior.\end{Resp}
\begin{Resp}\Rsign Candor est lucis ætérnæ et spéculum sine mácula.\end{Resp}
\switchcolumn
\EN
\begin{Resp}\Rsign No defiled thing can fall into her; \Star{} She is the brightness of the everlasting light, and the unspotted mirror of the power of God.\end{Resp}
\begin{Resp}\Vsign For she is more beautiful than the sun, and being compared with the light, she is found before it.\end{Resp}
\begin{Resp}\Rsign She is the brightness of the everlasting light, and the unspotted mirror of the power of God.\end{Resp}
\switchcolumn*

\LA
\LessonHead{Lectio VI}
\switchcolumn
\EN
\LessonHead{Lesson VI}
\switchcolumn*

\LA
\noindent Insuper summopere lætati sunt decernere Conceptionis festum ab omni ecclesia esse habendum eodem censu ac numero quo festum Nativitatis, idemque Conceptionis festum cum octava ab universa Ecclesia celebrandum, et ab omnibus inter ea quæ præcepta sunt, sancte colendum, ac pontificiam capellam in patriarchali nostra Liberiana basilica die Virginis Conceptionis sacro quotannis esse peragendam. Atque exoptantes in fidelium animis quotidie magis fovere hanc de Immaculata Deiparæ Conceptione doctrinam, eorumque pietatem excitare ad ipsam Virginem sine labe originali conceptam colendam et venerandam, gavisi sunt quam libentissime facultatem tribuere, ut in Lauretanis Litaniis, et in ipsa Missæ Præfatione Immaculatus ejusdem Virginis proclamaretur Conceptus, atque adeo lex credendi ipsa supplicandi lege statueretur.\par
\switchcolumn
\EN
\noindent Moreover, Our said Predecessors with great joy ordained that the Feast of the said Conception should be observed as of the same rank as that of the Nativity of the Blessed Virgin, and appointed that it should be kept with an Octave throughout the whole Church. They added this Feast to those which are commanded to be kept with solemnity, and ordered that the ceremony called a Papal Chapel should take place every year on this Feast in our Patriarchal Basilica of our Lady of the Snows. And above all did they rejoice in the hope of strengthening this belief in the minds of the faithful, and stirring them up to love and venerate the Virgin conceived without sin, when they granted permission to add to the Litany of Loretto the invocation, Queen conceived without original sin, and to insert the word stainless into the Preface of the Mass on this Feast, that so the law of prayer might become the law of belief.\par
\switchcolumn*

\LA
\begin{Resp}\Rsign Signum magnum appáruit in cælo: Múlier amícta sole, et luna sub pédibus ejus, \Star{} Et in cápite ejus coróna stellárum duódecim.\end{Resp}
\begin{Resp}\Vsign Induit eam Dóminus vestiméntis salútis, induménto justítiæ, et quasi sponsam ornávit eam monílibus suis.\end{Resp}
\begin{Resp}\Rsign Et in cápite ejus coróna stellárum duódecim.\end{Resp}
\begin{Resp}\Vsign Glória Patri, et Fílio, \Star{} et Spirítui Sancto.\end{Resp}
\begin{Resp}\Rsign Et in cápite ejus coróna stellárum duódecim.\end{Resp}
\switchcolumn
\EN
\begin{Resp}\Rsign There appeared a great wonder in heaven; a Woman clothed with the sun, and the moon under her feet \Star{} And upon her head a crown of twelve stars.\end{Resp}
\begin{Resp}\Vsign The Lord hath clothed her with the garments of salvation, and hath covered her with the robe of righteousness, yea, as a bride He hath adorned her with jewels.\end{Resp}
\begin{Resp}\Rsign And upon her head a crown of twelve stars.\end{Resp}
\begin{Resp}\Vsign Glory be to the Father, and to the Son, \Star{} and to the Holy Ghost.\end{Resp}
\begin{Resp}\Rsign And upon her head a crown of twelve stars.\end{Resp}
\switchcolumn*

\LA
\LessonHead{Lectio VII}
\switchcolumn
\EN
\LessonHead{Lesson VII}
\switchcolumn*

\LA
\LessonTitle{Léctio sancti Evangélii secúndum Lucam}
\switchcolumn
\EN
\LessonTitle{From the Holy Gospel according to Luke}
\switchcolumn*

\LA
\Cite{Luc 1:26-28}
\switchcolumn
\EN
\Cite{Luke 1:26-28}
\switchcolumn*

\LA
\noindent In illo témpore: Missus est Angelus Gábriel a Deo in civitátem Galilææ, cui nomen Názareth, ad Vírginem desponsátam viro, cui nomen erat Joseph, de domo David, et nomen Vírginis María. Et réliqua.\par
\switchcolumn
\EN
\noindent In that time the angel Gabriel was sent from God into a city of Galilee, called Nazareth, To a virgin espoused to a man whose name was Joseph, of the house of David; and the virgin's name was Mary. And so on.\par
\switchcolumn*

\LA
\LessonTitle{Homilía sancti Bernárdi Abbátis}
\switchcolumn
\EN
\LessonTitle{Homily by St. Bernard, Abbot (of Clairvaux.)}
\switchcolumn*

\LA
\Source{Ex Homilía 2 super Missus}
\switchcolumn
\EN
\Source{2nd on this text.}
\switchcolumn*

\LA
\noindent Lætáre, pater Adam, sed magis tu, o Heva mater, exsúlta, qui sicut ómnium paréntes, ita ómnium fuístis peremptóres; et quod infelícius est, prius peremptóres quam paréntes. Ambo, inquam, consolámini super fília, et tali fília, sed illa ámplius de qua malum ortum est prius, cujus oppróbrium in omnes pertransívit mulíeres. Instat namque tempus, quo jam tollátur oppróbrium, nec hábeat vir quid causétur advérsus féminam: qui útique, dum se imprudénter excusáre conarétur, crudéliter illam accusáre non cunctátus est, dicens: Múlier quam dedísti mihi, dedit mihi de ligno, et comédi. Proptérea curre Heva ad Maríam; curre mater ad fíliam; fília pro matre respóndeat: ipsa matris oppróbrium áuferat; ipsa patri pro matre satisfáciat: quia ecce si vir cécidit per féminam, jam non erígitur nisi per féminam.\par
\switchcolumn
\EN
\noindent Rejoice, father Adam, and yet more thou mother Eve, ye that are the source of all, and the ruin of all, and the unhappy cause of their ruin before ye gave them birth. Be comforted both in your daughter, and such a daughter; but chiefly thou, O woman, of whom the first evil came, and who hast cast thy slur upon all women. The time is come for the slur to be taken away, and for the man to have nothing to say against the woman. At the first, when he unwisely began to make excuse, he scrupled not to throw the blame upon her, saying, The woman whom Thou gavest to be with me, she gave me of the tree, and I did eat. Wherefore, O Eve, betake thyself to Mary Mother, betake thyself to thy daughter let the daughter answer for the mother let her take away her mother's reproach; let her make up to her father for her mother's fault for if man be fallen by means of woman, it is by means of woman that he is raised up again.\par
\switchcolumn*

\LA
\begin{Resp}\Rsign Hortus conclúsus soror mea sponsa, hortus conclúsus, fons signátus: \Star{} Emissiónes tuæ paradísus, o María.\end{Resp}
\begin{Resp}\Vsign Aperi mihi, soror mea, amíca mea, colúmba mea, immaculáta mea.\end{Resp}
\begin{Resp}\Rsign Emissiónes tuæ paradísus, o María.\end{Resp}
\switchcolumn
\EN
\begin{Resp}\Rsign A garden enclosed is my sister, my spouse, a garden enclosed, a fountain sealed. \Star{} O Mary, thy perfumes are a garden of delights.\end{Resp}
\begin{Resp}\Vsign Open to me, my sister, my love, my dove, my undefiled.\end{Resp}
\begin{Resp}\Rsign O Mary, thy perfumes are a garden of delights.\end{Resp}
\switchcolumn*

\LA
\LessonHead{Lectio VIII}
\switchcolumn
\EN
\LessonHead{Lesson VIII}
\switchcolumn*

\LA
\noindent Quid dicébas, o Adam? Múlier quam dedísti mihi, dedit mihi de ligno, et comédi. Verba malítiæ sunt hæc quibus magis áugeas quam déleas culpam. Verúmtamen Sapiéntia vicit malítiam, cum occasiónem véniæ, quam a te Deus interrogándo elícere tentávit, sed non pótuit, in thesáuro indeficiéntis suæ pietátis invénit. Rédditur nempe fémina pro fémina, prudens pro fátua, húmilis pro supérba, quæ pro ligno mortis gustum tibi pórrigat vitæ, et pro venenóso cibo illo amaritúdinis, dulcédinem páriat fructus ætérni. Muta ergo iníquæ excusatiónis verbum in vocem gratiárum actiónis, et dic: Dómine, múlier, quam dedísti mihi, dedit mihi de ligno vitæ, et comédi; et dulce factum est super mel ori meo, quia in ipso vivificásti me. Ecce enim ad hoc missus est Angelus ad Vírginem. O admirándam et omni honóre digníssimam Vírginem; o féminam singuláriter venerándam, super omnes féminas admirábilem, paréntum reparatrícem, posterórum vivificatrícem.\par
\switchcolumn
\EN
\noindent What didst thou say, O Adam? The woman whom Thou gavest to be with me, she gave me of the tree, and I did eat. These are wrathful words, by the which thou dost rather magnify than diminish thine offence. Nevertheless, Wisdom hath defeated thy malice. God asked thee that He might find in thee an occasion of pardon, but, in that He found it not, He hath sought and found it in the Treasure of His Own mercy. One woman answereth for another; the wise for the foolish; the lowly for the proud; for her that gave thee of the tree of death, another that giveth thee to taste of the tree of life; for her that brought thee the bitter food of sin, another that giveth thee of the sweet fruits of righteousness. Wherefore accuse the woman no more, but speak in thanksgiving, and say, Lord, the woman whom Thou hast given me, she hath given me of the tree of life, and I have eaten; and it is in my mouth sweeter than honey, for thereby hast Thou quickened me. (Ps. cxviii. 103, 93.) Behold, it was for this that the angel Gabriel was sent to the Virgin, to the most worshipful of women, a woman more wonderful than all women, the restorer of them that went before, and the quickener of them that come after her.\par
\switchcolumn*

\LA
\begin{Resp}\Rsign Magníficat ánima mea Dóminum: \Star{} Quia fecit mihi magna qui potens est, et sanctum nomen ejus.\end{Resp}
\begin{Resp}\Vsign Ecce enim ex hoc beátam me dicent omnes generatiónes.\end{Resp}
\begin{Resp}\Rsign Quia fecit mihi magna qui potens est, et sanctum nomen ejus.\end{Resp}
\begin{Resp}\Vsign Glória Patri, et Fílio, \Star{} et Spirítui Sancto.\end{Resp}
\begin{Resp}\Rsign Quia fecit mihi magna qui potens est, et sanctum nomen ejus.\end{Resp}
\switchcolumn
\EN
\begin{Resp}\Rsign My soul doth magnify the Lord; \Star{} For He That is mighty hath done to me great things, and holy is His name.\end{Resp}
\begin{Resp}\Vsign For, behold, from henceforth all generations shall call me blessed.\end{Resp}
\begin{Resp}\Rsign For He That is mighty hath done to me great things, and holy is His name.\end{Resp}
\begin{Resp}\Vsign Glory be to the Father, and to the Son, \Star{} and to the Holy Ghost.\end{Resp}
\begin{Resp}\Rsign For He That is mighty hath done to me great things, and holy is His name.\end{Resp}
\switchcolumn*

\LA
\LessonHead{Lectio IX}
\switchcolumn
\EN
\LessonHead{Lesson IX}
\switchcolumn*

\LA
\noindent Quam tibi áliam prædixísse Deus vidétur, quando ad serpéntem ait: Inimicítias ponam inter te et mulíerem? Et si adhuc dúbitas, quod de María díxerit, audi quod séquitur: Ipsa cónteret caput tuum. Cui hæc serváta victória est, nisi Maríæ? Ipsa procul dúbio caput contrívit venenátum, quæ omnimódam malígni suggestiónem tam de carnis illécebra, quam de mentis supérbia dedúxit ad níhilum. Quam vero áliam Sálomon requirébat, cum dicébat: Mulíerem fortem quis invéniet? Nóverat quippe vir sápiens hujus sexus infirmitátem, frágile corpus, lúbricam mentem. Quia tamen et Deum légerat promisísse, et ita vidébat congrúere, ut qui vícerat per féminam, vincerétur per ipsam, veheménter admírans ajébat: Mulíerem fortem quis invéniet? Quod est dícere: si ita de manu féminæ pendet et nostra ómnium salus, et innocéntiæ restitútio, et de hoste victória: fortis omníno necésse est ut provideátur, quæ ad tantum opus possit esse idónea.\par
\switchcolumn
\EN
\noindent Has it not of this thy daughter, O Adam, that God spake when He said unto the serpent, I will put enmity between thee and the woman And if thou wilt still doubt that He speaketh of Mary, hear what followeth She shall bruise thy head. Who won this conquest but Mary? She brought to nought the whole wiles of Satan, whether for the pollution of her body or the injury of her soul. Was it not of her that Solomon spake, where he saith, Who shall find a virtuous woman? (Prov. xxxi. 10.) The wise man knew the weaknesses of women, how frail they are in body, and how changeable in mind. But he had read that God had promised that the enemy, who had prevailed by means of a woman, was by a woman to be overthrown, and he believed. But he wondered greatly, and said, Who shall find a virtuous woman? that is to say If our salvation, and the bringing back of that which is lost, and the final triumph over the enemy, is in the hand of a woman, it must needs be that a virtuous woman be found, meet to work in that matter\par
\switchcolumn*

\LA
\TeDeum{Te Deum}
\switchcolumn
\EN
\TeDeum{Te Deum}
\switchcolumn*

\end{paracol}

\par\needspace{10\baselineskip}\addvspace{1.2em}{\centering\small\textsc{Die 11 Decembris} · \textsc{11 December}\par}
\Office{S. Damasi Papæ et Confessoris}{St. Damasus, Pope and Confessor}{Semiduplex}
\addcontentsline{toc}{subsection}{Die 11 Decembris · S. Damasi Papæ et Confessoris\enspace\textit{St. Damasus, Pope and Confessor}}
\begin{paracol}{2}
\LA
\RefLine{Lectiones I–III: de Scriptura occurrente.}
\switchcolumn
\EN
\RefLine{Lessons I–III: from the Scripture of the day.}
\switchcolumn*

\LA
\RefLine{Lectiones VII–IX: ex Commúni Summorum Pontificum Confessorum (C4b).}
\switchcolumn
\EN
\RefLine{Lessons VII–IX: from the Common of Popes (C4b).}
\switchcolumn*

\LA
\LessonHead{Lectio IV}
\switchcolumn
\EN
\LessonHead{Lesson IV}
\switchcolumn*

\LA
\noindent Dámasus Hispánus, vir egrégius et erudítus in Scriptúris, indícto primo Constantinopolitáno concílio nefáriam Eunómii et Macedónii hæresim exstínxit. Idem Ariminénsem convéntum a Libério jam ante rejéctum, íterum condemnávit: in quo, ut scribit sanctus Hierónymus, Valéntis potíssimum et Ursácii fráudibus damnátio Nicænæ fídei conclamáta fuit: et ingemíscens orbis terrárum, se Ariánum esse mirátus est.\par
\switchcolumn
\EN
\noindent Damasus was a Spaniard, a man of eminence and of great learning in the Scriptures, (and was elected to the Chair of Peter in the year of our Lord 381) he convoked the First Council of Constantinople, wherein he crushed the wicked heresy of Eunomius and Macedonius. He confirmed the condemnation of the Assembly, at Rimini, which condemnation had already been pronounced by Liberius. This Assembly of Rimini was that in which, to use the language of St. Jerome, Valens and Ursacius brought it about through trickery that the Faith of Nice was abrogated by mob law, and the world afterwards groaned in amazement to find itself Arian.\par
\switchcolumn*

\LA
\begin{Resp}\Rsign Invéni David servum meum, óleo sancto meo unxi eum: \Star{} Manus enim mea auxiliábitur ei.\end{Resp}
\begin{Resp}\Vsign Nihil profíciet inimícus in eo, et fílius iniquitátis non nocébit ei.\end{Resp}
\begin{Resp}\Rsign Manus enim mea auxiliábitur ei.\end{Resp}
\switchcolumn
\EN
\begin{Resp}\Rsign I have found David My servant, with My holy oil have I anointed him \Star{} For My hand shall help him.\end{Resp}
\begin{Resp}\Vsign The enemy shall prevail nothing against him, nor the son of wickedness afflict him.\end{Resp}
\begin{Resp}\Rsign For My hand shall help him.\end{Resp}
\switchcolumn*

\LA
\LessonHead{Lectio V}
\switchcolumn
\EN
\LessonHead{Lesson V}
\switchcolumn*

\LA
\noindent Basílicas duas ædificávit, álteram sancti Lauréntii nómine ad theátrum Pompéji, quam máximis munéribus auxit, eíque domos et prædia attríbuit: álteram via Ardeatína ad Catacúmbas. Platóniam étiam, ubi córpora sanctórum Petri et Pauli aliquámdiu jacuérunt, dedicávit, et exornávit elegántibus vérsibus. Idémque prosa et versu scripsit de virginitáte, múltaque ália metro édidit.\par
\switchcolumn
\EN
\noindent This Pope built two Basilicas, first, St. Lawrence's, near Pompey's Theatre, which he magnificently enriched, and endowed with houses and farms; and, secondly, another, over the Catacombs on the Road to Ardea. He also consecrated the Platonia, where the bodies of St. Peter and St. Paul lay for some time, and decorated it with elegant inscriptions in poetry composed by himself. He wrote on the subject of virginity both in prose and verse, and likewise many other poems on various subjects.\par
\switchcolumn*

\LA
\begin{Resp}\Rsign Pósui adjutórium super poténtem, et exaltávi eléctum de plebe mea: \Star{} Manus enim mea auxiliábitur ei.\end{Resp}
\begin{Resp}\Vsign Invéni David servum meum, óleo sancto meo unxi eum.\end{Resp}
\begin{Resp}\Rsign Manus enim mea auxiliábitur ei.\end{Resp}
\switchcolumn
\EN
\begin{Resp}\Rsign I have laid help upon one that is mighty, and have exalted one chosen out of My people \Star{} For My hand shall help him.\end{Resp}
\begin{Resp}\Vsign I have found David My servant, with My holy oil have I anointed him.\end{Resp}
\begin{Resp}\Rsign For My hand shall help him.\end{Resp}
\switchcolumn*

\LA
\LessonHead{Lectio VI}
\switchcolumn
\EN
\LessonHead{Lesson VI}
\switchcolumn*

\LA
\noindent Pœnam taliónis constítuit iis, qui álterum falsi críminis accusássent. Státuit, ut, quod plúribus jam locis erat in usu, Psalmi per omnes ecclésias die noctúque ab altérnis caneréntur; et in fine cujúsque Psalmi dicerétur: Glória Patri, et Fílio, et Spirítui Sancto. Ejus jussu sanctus Hierónymus novum testaméntum Græcæ fídei réddidit. Cum Ecclésiam rexísset annos decem et septem, menses duos, dies vigínti sex, et habuísset ordinatiónes quinque mense Decémbri, quibus creávit presbyteros trigínta unum, diáconos úndecim, epíscopos per divérsa loca sexagínta duos; virtúte, doctrína, ac prudéntia clarus, prope octogenárius, Theodósio senióre imperánte, obdormívit in Dómino: et via Ardeatína una cum matre et soróre sepúltus est in basílica, quam ipse ædificáverat. Illíus relíquiæ póstea translátæ sunt in ecclésiam sancti Lauréntii, ab ejus nómine in Dámaso vocátam.\par
\switchcolumn
\EN
\noindent He ordained that false accusers should be punished for the offences which they had falsely laid to the charge of their neighbours. He established the usage, which already prevailed in many churches, of singing the Psalms, both by day and by night, by alternate choirs, and of adding at the end of each Psalm the words, Glory be to the Father, and to the Son, and to the Holy Ghost. It was at his command that St. Jerome revised the translation of the New Testament to accord with the Greek text. He ruled the Church for seventeen years, two months, and twenty-six days. He held five Advent ordinations, wherein he ordained thirty-one Priests, eleven Deacons, and sixty-two Bishops for diverse Sees. At length he fell asleep in the Lord, in the reign of Theodosius the Elder, (upon the 10th day of December, in the year 384, being) aged nearly eighty years, and full of righteousness, truth, and judgment. He was buried beside his mother and sister in the Church which he had himself founded on the Road to Ardea. His reliques were afterwards taken to the Basilica of St. Lawrence, which is thence sometimes called San Lorenzo in Damaso.\par
\switchcolumn*

\LA
\begin{Resp}\Rsign Iste est, qui ante Deum magnas virtútes operátus est, et omnis terra doctrína ejus repléta est: \Star{} Ipse intercédat pro peccátis ómnium populórum.\end{Resp}
\begin{Resp}\Vsign Iste est, qui contémpsit vitam mundi, et pervénit ad cæléstia regna.\end{Resp}
\begin{Resp}\Rsign Ipse intercédat pro peccátis ómnium populórum.\end{Resp}
\begin{Resp}\Vsign Glória Patri, et Fílio, \Star{} et Spirítui Sancto.\end{Resp}
\begin{Resp}\Rsign Ipse intercédat pro peccátis ómnium populórum.\end{Resp}
\switchcolumn
\EN
\begin{Resp}\Rsign This is he which wrought great wonders before God, and the whole earth is full of his teaching \Star{} May he pray for all people, that their sins may be forgiven unto them\end{Resp}
\begin{Resp}\Vsign This is he which loved not his life in this world, and hath attained unto the kingdom of heaven.\end{Resp}
\begin{Resp}\Rsign May he pray for all people, that their sins may be forgiven unto them\end{Resp}
\begin{Resp}\Vsign Glory be to the Father, and to the Son, \Star{} and to the Holy Ghost.\end{Resp}
\begin{Resp}\Rsign May he pray for all people, that their sins may be forgiven unto them\end{Resp}
\switchcolumn*

\end{paracol}

\par\needspace{10\baselineskip}\addvspace{1.2em}{\centering\small\textsc{Die 12 Decembris} · \textsc{12 December}\par}
\Office{De V die infra Octavam Concept. Immac. Beatæ Mariæ Virginis}{De V die infra Octavam Concept. Immac. Beatæ Mariæ Virginis}{Semiduplex}
\addcontentsline{toc}{subsection}{Die 12 Decembris · De V die infra Octavam Concept. Immac. Beatæ Mariæ Virginis\enspace\textit{De V die infra Octavam Concept. Immac. Beatæ Mariæ Virginis}}
\begin{paracol}{2}
\LA
\RefLine{Lectiones I–III: de Scriptura occurrente.}
\switchcolumn
\EN
\RefLine{Lessons I–III: from the Scripture of the day.}
\switchcolumn*

\LA
\LessonHead{Lectio IV}
\switchcolumn
\EN
\LessonHead{Lesson IV}
\switchcolumn*

\LA
\Source{Ex Bulla dogmatica Pii Papæ noni}
\switchcolumn
\EN
\Source{From the Dogmatic Bull of Pope Pius IX.}
\switchcolumn*

\LA
\noindent Quoniam quæ ad cultum pertinent, intimo plane vinculo cum ejusdem objecto conserta sunt, neque rata et fixa manere possunt, si illud anceps sit et in ambiguo versetur, idcirco decessores nostri Romani Pontifices omni cura Conceptionis cultum amplificantes illius etiam objectum ac doctrinam declarare, et inculcare impensissime studuerunt. Etenim clare aperteque docuere, festum agi de Virginis Conceptione, atque uti falsam et ab Ecclesiæ mente alienissimam proscripserunt illorum opinionem qui non Conceptionem ipsam, sed sanctificationem ab Ecclesia coli arbitrarentur et affirmarent.\par
\switchcolumn
\EN
\noindent The language used in public worship is the necessary offspring of the teaching which it expresseth, and the former can have no safety unless the latter be settled. Wherefore Our Predecessors the Roman Pontiffs, while encouraging the pious love of the faithful for the Conception of the Blessed Virgin, have taken care ceaselessly to inculcate the sinlessness of the same. They have always particularly insisted that the Feast should be observed not in honour of Mary's sanctification, a false opinion, most foreign to the mind of the Church (but which hath nevertheless been maintained by some,) but in honour of her Conception itself.\par
\switchcolumn*

\LA
\begin{Resp}\Rsign Ego ex ore Altíssimi prodívi, primogénita ante omnem creatúram: ego feci in cælis, ut orirétur lumen indefíciens. \Star{} Nondum erant abýssi, et ego jam concépta eram.\end{Resp}
\begin{Resp}\Vsign Deus enim creávit me in justítia, et apprehéndit manum meam, et servávit me.\end{Resp}
\begin{Resp}\Rsign Nondum erant abýssi, et ego jam concépta eram.\end{Resp}
\switchcolumn
\EN
\begin{Resp}\Rsign I came out of the mouth of the Most High, the first-begotten before every creature. I made the unfading light to arise in the heavens. \Star{} When there were no depths I was conceived.\end{Resp}
\begin{Resp}\Vsign For the Lord hath created me in righteousness, and hath held mine hand, and hath kept me.\end{Resp}
\begin{Resp}\Rsign When there were no depths I was conceived.\end{Resp}
\switchcolumn*

\LA
\LessonHead{Lectio V}
\switchcolumn
\EN
\LessonHead{Lesson V}
\switchcolumn*

\LA
\noindent Neque mitius cum iis agendum esse existimarunt, qui ad labefactandam de Immaculata Virginis Conceptione doctrinam excogitato inter primum atque alterum Conceptionis instans et momentum discrimine, asserebant, celebrari quidem Conceptionem sed non pro primo instanti atque momento. Ipsi namque prædecessores nostri suarum partium esse duxerunt, et beatissimæ Virginis Conceptionis festum, et Conceptionem pro primo instanti tamquam verum cultus objectum omni studio tueri ac propugnare. Hinc decretoria plane verba, quibus Alexander septimus decessor noster sinceram Ecclesiæ mentem declaravit, inquiens: Sane vetus est Christifidelium erga ejus beatissimam Matrem Virginem Mariam pietas sentientium, ejus animam in primo instanti creationis atque infusionis in corpus fuísse speciali Dei gratia et privilegio, intuitu meritorum Jesu Christi ejus Filii humani generis Redemptoris, a macula peccati originalis præservatam immunem, atque in hoc sensu ejus Conceptionis festivitatem solemni ritu colentium et celebrantium.\par
\switchcolumn
\EN
\noindent The same Our Predecessors have likewise resisted the dreams of those who have imagined that in the sinless Conception there were Two Instants, and that the Church celebrateth the Second and not the First. Indeed, Our said Predecessors have considered the sinlessness of the First Instant to be as much a truth for their assertion, protection, and promulgation, as the sinlessness of the Conception at all. Hence came those words in which Our Predecessor Alexander VII. in a decree declareth the mind of the Church, and saith, Christ's faithful people, drawn by love to His most blessed Mother, the Virgin Mary, have of a long time believed that God, at the very First Instant in which He made her soul and joined it to her body, by a special grace and privilege granted to her, through the merits of His dear Son, Christ Jesus, the Saviour of the world, Whose precious death He foreknew, cleansed her from all sin, original as well as actual; and it is in this belief, and no other, that the said faithful of Christ have always kept with devotion and joy\par
\switchcolumn*

\LA
\begin{Resp}\Rsign Nihil inquinátum in eam incúrrit: \Star{} Candor est lucis ætérnæ et spéculum sine mácula.\end{Resp}
\begin{Resp}\Vsign Est enim hæc speciósior sole, et luci comparáta invenítur púrior.\end{Resp}
\begin{Resp}\Rsign Candor est lucis ætérnæ et spéculum sine mácula.\end{Resp}
\switchcolumn
\EN
\begin{Resp}\Rsign No defiled thing can fall into her; \Star{} She is the brightness of the everlasting light, and the unspotted mirror of the power of God.\end{Resp}
\begin{Resp}\Vsign For she is more beautiful than the sun, and being compared with the light, she is found before it.\end{Resp}
\begin{Resp}\Rsign She is the brightness of the everlasting light, and the unspotted mirror of the power of God.\end{Resp}
\switchcolumn*

\LA
\LessonHead{Lectio VI}
\switchcolumn
\EN
\LessonHead{Lesson VI}
\switchcolumn*

\LA
\noindent Atque illud in primis solemne quoque fuit iisdem decessoribus nostris, doctrinam de Immaculata Dei Matris Conceptione sartam tectamque omni cura, studio et contentione tueri. Etenim non solum nullatenus passi sunt, ipsam doctrinam quovis modo a quopiam notari atque traduci, verum étiam longe ulterius progressi, perspicuis declarationibus iteratisque vicibus edixerunt: Doctrinam qua Immaculatam Virginis Conceptionem profitemur, esse, suoque merito haberi cum ecclesiastico cultu plane consonam, eamque veterem, ac prope universalem et ejusmodi, quam Romana Ecclesia sibi fovendam tuendamque susceperit, atque omnino dignam, quæ in sacra ipsa liturgia, solemnioribusque precibus usurparetur. Neque his contenti, ut ipsa de Immaculato Virginis Conceptu doctrina inviolata persisteret, opinionem huic doctrinæ adversam, sive publice, sive privatim, defendi posse severissime prohibuere, eamque multiplici veluti vulnere confectam esse voluerunt.\par
\switchcolumn
\EN
\noindent It hath always been one of the most weighty cares of Our said Predecessors the Roman Pontiffs to protect the doctrine of the sinlessness of Mary's Conception from any sort of attack or corruption. Not only have they suffered no one to condemn and traduce it, but they have gone much further, and in public and repeated declarations have averred That that doctrine which holdeth that the Virgin was conceived without sin is a doctrine, the arguments in support of which are strong enough to enable the time of public worship, which is ancient, which is almost universal, which is one of those which the Church of Rome encourageth and protecteth, and which is worthy even to be expressed in the Liturgy itself, and in the most solemn prayers of the Church. Our said Predecessors did not stop even here, but in order to preserve the doctrine of the Virgin's sinless Conception from an injury they strictly forbade that the opposite opinion should be maintained either in public or in private, to the end that it might at length die out under disapprobation.\par
\switchcolumn*

\LA
\begin{Resp}\Rsign Signum magnum appáruit in cælo: Múlier amícta sole, et luna sub pédibus ejus, \Star{} Et in cápite ejus coróna stellárum duódecim.\end{Resp}
\begin{Resp}\Vsign Induit eam Dóminus vestiméntis salútis, induménto justítiæ, et quasi sponsam ornávit eam monílibus suis.\end{Resp}
\begin{Resp}\Rsign Et in cápite ejus coróna stellárum duódecim.\end{Resp}
\begin{Resp}\Vsign Glória Patri, et Fílio, \Star{} et Spirítui Sancto.\end{Resp}
\begin{Resp}\Rsign Et in cápite ejus coróna stellárum duódecim.\end{Resp}
\switchcolumn
\EN
\begin{Resp}\Rsign There appeared a great wonder in heaven; a Woman clothed with the sun, and the moon under her feet \Star{} And upon her head a crown of twelve stars.\end{Resp}
\begin{Resp}\Vsign The Lord hath clothed her with the garments of salvation, and hath covered her with the robe of righteousness, yea, as a bride He hath adorned her with jewels.\end{Resp}
\begin{Resp}\Rsign And upon her head a crown of twelve stars.\end{Resp}
\begin{Resp}\Vsign Glory be to the Father, and to the Son, \Star{} and to the Holy Ghost.\end{Resp}
\begin{Resp}\Rsign And upon her head a crown of twelve stars.\end{Resp}
\switchcolumn*

\LA
\LessonHead{Lectio VII}
\switchcolumn
\EN
\LessonHead{Lesson VII}
\switchcolumn*

\LA
\LessonTitle{Léctio sancti Evangélii secúndum Lucam}
\switchcolumn
\EN
\LessonTitle{From the Holy Gospel according to Luke}
\switchcolumn*

\LA
\Cite{Luc 1:26-28}
\switchcolumn
\EN
\Cite{Luke 1:26-28}
\switchcolumn*

\LA
\noindent In illo témpore: Missus est Angelus Gabriel a Deo in civitatem Galilææ, cui nomen Nazareth, ad Virginem desponsatam viro, cui nomen erat Joseph, de domo David, et nomen Virginis Maria. Et réliqua.\par
\switchcolumn
\EN
\noindent At that time the angel Gabriel was sent from God into a city of Galilee, called Nazareth, to a virgin espoused to a man whose name was Joseph, of the house of David; and the virgin's name was Mary. And so on.\par
\switchcolumn*

\LA
\LessonTitle{Homilia sancti Tharásii Epíscopi}
\switchcolumn
\EN
\LessonTitle{Homily by St. Tarasius, Patriarch (of Constantinople.)}
\switchcolumn*

\LA
\Source{De Præsentatióne Deiparæ}
\switchcolumn
\EN
\Source{On the Presentation of the Mother of God.}
\switchcolumn*

\LA
\noindent Quibus te laudibus cumulabimus, Maria? O puella immaculata; o virgo impolluta; mulierum ornamentum filiarum nitor! O mater Virgo sancta, tu benedicta inter mulieres; tu celebrata propter innocentiam; tu obsignata virginitate. Tu Adami maledicti expiatio; tu debiti Hevæ solutio. Tu Abelis purissima oblatio; primogenitorum delectus; immaculatum sacrificium. Tu Enos in Deum spes non pudore suffusa; tu Enoch inita gratia, et in securam vitam migratio. Tu Noë arca, et secundæ regenerationis apud Deum conciliatio. Tu regni et sacerdotii Melchisedech perillustris splendor; tu Abrahami firma fiducia et promissionis futuræ posteritatis obsequens fides. Tu Isaac novum sacrificium et rationale holocaustum; tu Jacob in scalam ascensus causa, et fœcunditatis in duodecim tribus permanentis expressio nobilissima. Tu Judæ apparuisti secundum stirpem filia; tu Josephi pudicitia et veteris AEgypti, nimirum synagogæ Judæorum eversio, o Immaculata! Tu Moysis ejusdemque legislatoris liber divinitus concinnatus, in quo scriptum est sacramentum regenerationis, et divinis digitis insculpta in tabulis lex est tamquam in monte Sina, ubi novus Israël ab intelligibilium Aegyptiorum servitute vindicabitur, quemadmodum antiquus populus in solitudine manna et aqua de petra satiatus est, petra autem erat Christus e tuo gremio proditurus tamquam sponsus de thalamo. Tu Aaronis virga florescens; tu es Davidis filia fimbriis aureis circumvestita vario ornatu nitescens.\par
\switchcolumn
\EN
\noindent O Mary, where shall I find words to praise thee? Maiden undefiled, virgin unstained, exaltation of women, glory of daughters! Holy Maiden Mother, blessed art thou among women, thy glory is in thy guilelessness, and thy name is a name of purity. In thee the curse of Adam is done away, and the debt of Eve paid. Thou art the clean offering of Abel, chosen out of the firstlings profession thereof to be made at of the flock, a pure sacrifice. Thou art the hope of Enoch, that firm hope that he had in God, and was not ashamed. Thou art the grace that was in Enoch in this life, and his transit to a better. Thou art the Holy Ark of Noah, and the bond of reconciliation with God in a new regeneration. Thou art the exceeding glory of the kingdom and Priesthood of Melchisedech. Thou art the unshaken trust of Abraham, and his faith in the promise of children that were to-be. Thou art the renewed oblation and the reasonable burntoffering of Isaac. Thou art the ladder that Jacob saw going up to heaven, and the most noble of all his children throughout the twelve tribes of Israel. According to the flesh thou art the daughter of Judah. Thou art the modesty of Joseph, and the overthrow of the old Egypt, yea, and of the Synagogue of the Jews. O purest! Thou art the book of Moses the Lawgiver, whereon the new covenant is written with the finger of God, for the new Israel, fleeing from the spiritual Egypt, even as the old law was written upon Sinai, for the old Israel, that Israel which was fed in the wilderness upon manna and water from the rock, whereof both were types of Christ, which was yet to come from thy womb, as a bridegroom from his chamber. Thou art Aaron's rod that budded. Thou art David's daughter, all glorious within, clothed in a vesture of gold, wrought about with diverse colours.\par
\switchcolumn*

\LA
\begin{Resp}\Rsign Hortus conclúsus soror mea sponsa, hortus conclúsus, fons signátus: \Star{} Emissiónes tuæ paradísus, o María.\end{Resp}
\begin{Resp}\Vsign Aperi mihi, soror mea, amíca mea, colúmba mea, immaculáta mea.\end{Resp}
\begin{Resp}\Rsign Emissiónes tuæ paradísus, o María.\end{Resp}
\switchcolumn
\EN
\begin{Resp}\Rsign A garden enclosed is my sister, my spouse, a garden enclosed, a fountain sealed. \Star{} O Mary, thy perfumes are a garden of delights.\end{Resp}
\begin{Resp}\Vsign Open to me, my sister, my love, my dove, my undefiled.\end{Resp}
\begin{Resp}\Rsign O Mary, thy perfumes are a garden of delights.\end{Resp}
\switchcolumn*

\LA
\LessonHead{Lectio VIII}
\switchcolumn
\EN
\LessonHead{Lesson VIII}
\switchcolumn*

\LA
\noindent Tu es prophetarum speculum, et rerum ab illis prænuntiatarum exitus. Te enixe Ezechiel vaticinans appellavit portam clausam, per quam nemo hominum umquam transibit nisi Dominus Deus solus, et portam clausam conservabit. Te Isaias ille in primis grandiloquus prænuntiat virgam Jesse, ex qua flos Christus orietur, et fruticibus vitiorum exstirpatis radicitus, plantas divinæ cognitionis in agro inseret. Te Jeremias præmonstravit inquiens: Ecce dies venient, dicit Dominus, et feriam domui Israël et domui Judæ fœdus novum, quod constitui cum patribus eorum, ita significans adventum ortumque filii tui, et populum Gentium vocans ad Deum adorandum inde usque a finibus terræ. Te étiam Daniel vir desideriorum proclamavit montem ingentem, e quo Christus lapis angularis abscindetur, et simulacrum multiformis serpentis ruina atque exitio dissipabit. Te honoro agnam immaculatam, te prædico gratia plenam, te cano Dei habitationem puram et immaculatam. Et sane ubi abundavit delictum, superabundavit gratia. Per mulierem mortem lucrati sumus, per mulierem universa ipse rursus instaurabit. Per serpentem cibum accepimus amari saporis, per ipsum vero rursum vescemur cibo immortalitatis. Prima parens Heva Cainum in lucem edidit invidiæ et nequitiæ principem; unigenitus Filius tuus erit primogenitus vitæ et resurrectionis, O inauditum prodigium! O admirandam novitatem! O sapientiam nullis verbis coæquandam!\par
\switchcolumn
\EN
\noindent Thou art the vision of the Prophets and the fulfilment of those things which they foretold. Thou art the gate whereof Ezekiel spake, when he prophesied, and said, This gate shall be shut, it shall not be opened, and no man shall enter in by it; because the Lord, the God of Israel, hath entered in by it, therefore it shall be shut (xliv. 2.) Thou art the Rod of Jesse, whereof Isaiah spake, \textasciitilde{}(xi. I,) even that Rod whose Flower is Christ, and whose offshoots shall choke out all the seedlings of sin, and fill the earth with plants of grace. Thou art the Covenant foretold by Jeremiah when he said (xxxi. 31) Behold, the days come, saith the Lord, that I will make a new covenant with the house of Israel, and with the house of Judah, not according to the covenant that I made with their fathers thereby signifying the coming of thy Son, and calling upon all nations to worship Him for their God, even to the uttermost parts of the earth. Thou art the great mountain spoken of by Daniel, the man greatly beloved, wherefrom is cut without man's hands the corner-stone, that is, Christ, which hath smitten in pieces the parti-coloured image of the old serpent. I honour thee as the unpolluted fountain, I proclaim that thou art full of grace, I praise thee as the clean and undefiled tabernacle of God. Verily, where sin abounded, grace did much more abound. As by a woman death entered into the world, by a woman came the power to rise again. The serpent gave us to eat deadly fruit, but that fall hath ended in the lifegiving Bread of Immortality. Eve, our first mother, brought forth Cain the first murderer; thou, O Mary, hast brought forth Christ, the first-fruits of life and of the resurrection. Ear hath not heard the like. It hath not entered into the heart of man to conceive this new thing. Blessed be the unspeakable depths of the Wisdom of God.\par
\switchcolumn*

\LA
\begin{Resp}\Rsign Magníficat ánima mea Dóminum: \Star{} Quia fecit mihi magna qui potens est, et sanctum nomen ejus.\end{Resp}
\begin{Resp}\Vsign Ecce enim ex hoc beátam me dicent omnes generatiónes.\end{Resp}
\begin{Resp}\Rsign Quia fecit mihi magna qui potens est, et sanctum nomen ejus.\end{Resp}
\begin{Resp}\Vsign Glória Patri, et Fílio, \Star{} et Spirítui Sancto.\end{Resp}
\begin{Resp}\Rsign Quia fecit mihi magna qui potens est, et sanctum nomen ejus.\end{Resp}
\switchcolumn
\EN
\begin{Resp}\Rsign My soul doth magnify the Lord; \Star{} For He That is mighty hath done to me great things, and holy is His name.\end{Resp}
\begin{Resp}\Vsign For, behold, from henceforth all generations shall call me blessed.\end{Resp}
\begin{Resp}\Rsign For He That is mighty hath done to me great things, and holy is His name.\end{Resp}
\begin{Resp}\Vsign Glory be to the Father, and to the Son, \Star{} and to the Holy Ghost.\end{Resp}
\begin{Resp}\Rsign For He That is mighty hath done to me great things, and holy is His name.\end{Resp}
\switchcolumn*

\LA
\LessonHead{Lectio IX}
\switchcolumn
\EN
\LessonHead{Lesson IX}
\switchcolumn*

\LA
\noindent Nos autem populus Dei, gens sancta, congregatio acceptabilis, filii columbæ, soboles gratiæ in hac Virginis celebritate puris animis, impollutis labiis, multisonis linguis hymnos suavidicos extollamus. Illustre hoc festum, principem solemnitatem Angelis lætam et hominum prædicatione dignissimam, prouti par est, venerantes illud Ave Gabrielis cum reverentia et gaudio sancto conclamemus. Ave delicium Patris, per quam ad ultimos terræ fines Dei cognitio manavit. Ave Filii domicilium, de qua ille carne indutus prodivit. Ave sancti Spiritus habitaculum ineffabile. Ave sanctior Cherubim; ave gloriosior Seraphim; ave cælo altior; ave sole splendidior; ave luna micantior; ave multiplex astrorum nitor; ave levis nubes, quæ cælestem pluviam inspergis. Ave aura sancta, quæ spiritum malitiæ a terra dissipasti. Ave nobile præconium Prophetarum; ave Apostolorum auditus per totum orbem sonus; ave Mártyrum excellens confessio; ave Patriarcharum laudatissima prædicatio; ave Sanctorum summum ornamentum. Ave causa salutis omnium mortalium; ave regina pacis conciliatrix; ave matrum splendor immaculatus. Ave mediatrix omnium, qui sub cælo sunt; ave totius orbis reparatio; ave gratia plena, Dominus tecum, qui ante te, et ex te, et nobiscum. Ipsi laus cum Patre et sanctissimo et vivifico Spiritu, nunc et semper, et in infinita sæcula sæculorum. Amen.\par
\switchcolumn
\EN
\noindent And now we, the people of God, a holy generation, an acceptable congregation, the nestlings of the dove of peace, children of grace, do with purified minds and unpolluted lips, praise God in the tongues of all nations in this joyful solemnity of the Virgin. This is a noble Feast wherein the Angels keep holiday and men do most fitly offer praise, even a feast wherein we echo with reverence and joy that salutation first spoken by Gabriel. Hail Mary! Hail, thou Paradise of God the Father, whence the knowledge of Him floweth in broad rivers to the ends of the earth! Hail, Dwelling-place of God the Son, whence He came forth clothed in flesh! Hail, mysterious Tabernacle of God the Holy Ghost! Hail, thou that art holier than the Cherubim! Hail, thou that art more glorious than the Seraphim! Hail, thou that art nobler than the heavens! Hail, thou that art brighter than the sun! Hail, thou that art fairer than the moon! Hail, manifold splendour of the stars! Hail, light cloud, dropping the dew of heaven! Hail, holy breeze, clearing the air of the vapours of sin! Hail, royal theme of the Prophets! Hail, sound of the Apostles gone out into all the earth! Hail, most excellent confession of the Martyrs! Hail, just hope of the Patriarchs! Hail, peculiar honour of all the Saints! Hail, source of health to dying creatures! Hail, O Queen, ambassadress of peace! Hail, stainless crown of motherhood! Hail, advocate of all under heaven! Hail, restoration of the whole world! Hail, thou that art full of grace, the Lord is with thee, even the Lord that is before thee, and from thee, and that is with us. To Him, with the Father, and the most holy and Life-giving Spirit, be ascribed all praise, now and ever, world without end. Amen.\par
\switchcolumn*

\LA
\TeDeum{Te Deum}
\switchcolumn
\EN
\TeDeum{Te Deum}
\switchcolumn*

\end{paracol}

\par\needspace{10\baselineskip}\addvspace{1.2em}{\centering\small\textsc{Die 13 Decembris} · \textsc{13 December}\par}
\Office{S. Luciæ Virginis et Martyris}{St. Lucy, Virgin and Martyr}{Duplex}
\addcontentsline{toc}{subsection}{Die 13 Decembris · S. Luciæ Virginis et Martyris\enspace\textit{St. Lucy, Virgin and Martyr}}
\begin{paracol}{2}
\LA
\RefLine{Lectiones I–III: de Scriptura occurrente.}
\switchcolumn
\EN
\RefLine{Lessons I–III: from the Scripture of the day.}
\switchcolumn*

\LA
\RefLine{Lectiones VII–IX: ex Commúni Unius Virginis et Martyris (C6), in 2 loco.}
\switchcolumn
\EN
\RefLine{Lessons VII–IX: from the Common of a Virgin Martyr (C6), set 2.}
\switchcolumn*

\LA
\LessonHead{Lectio IV}
\switchcolumn
\EN
\LessonHead{Lesson IV}
\switchcolumn*

\LA
\noindent Lúcia virgo Syracusana, genere et christiana fide ab infántia nobilis, una cum matre Eutychia, quæ sánguinis fluxu laborabat, Cátanam ad venerándum corpus beátæ Agathæ venit: quæ ad ejus sepúlcrum cum suppliciter orasset, Agathæ intercessióne matri sanitátem impetrávit. Statim vero matrem exorávit, ut quam dotem sibi datura esset, Christi paupéribus tribui paterétur. Ut ígitur Syracusas rediit, omnem pecúniam, quam ex facultátibus venditis redegerat, paupéribus distríbuit.\par
\switchcolumn
\EN
\noindent Lucy was a maiden of Syracuse, the daughter of a noble Christian family. Her mother Eutychia, being afflicted with an issue of blood, went with her to Catania, to pray before the body of the blessed Agatha. Lucy, by her earnest prayers at the grave, obtained her mother's cure, through the intercession of Agatha, and then immediately begged her to give to Christ's poor the whole dowry which had been set apart for herself. As soon, therefore, as they returned to Syracuse, they sold the property, and distributed the money among the poor.\par
\switchcolumn*

\LA
\begin{Resp}\Rsign Lúcia virgo, quid a me petis quod ipsa póteris præstáre contínuo matri tuæ? nam et fides tua illi subvénit, et ecce salváta est: \Star{} Quia jucúndum Deo in tua virginitáte habitáculum præparásti.\end{Resp}
\begin{Resp}\Vsign Sicut per me cívitas Catanénsium sublimátur a Christo, ita per te Syracusána cívitas decorábitur.\end{Resp}
\begin{Resp}\Rsign Quia jucúndum Deo in tua virginitáte habitáculum præparásti.\end{Resp}
\switchcolumn
\EN
\begin{Resp}\Rsign Maiden Lucy, why seekest thou of me that which thou thyself canst presently give thy mother? For thy faith hath helped her, and, behold, she is made whole; \Star{} Because thou hast made in thy virginity a pleasant dwelling-place for thy God.\end{Resp}
\begin{Resp}\Vsign Even as Christ hath by me glorified Catania, so by thee shall He glorify Syracuse.\end{Resp}
\begin{Resp}\Rsign Because thou hast made in thy virginity a pleasant dwellingplace for thy God.\end{Resp}
\switchcolumn*

\LA
\LessonHead{Lectio V}
\switchcolumn
\EN
\LessonHead{Lesson V}
\switchcolumn*

\LA
\noindent Quod ubi rescivísset is, cui eam paréntes contra Vírginis voluntátem desponderant, apud Paschasium præfectum Lúciam, quod christiana esset, accusávit. Quam ille cum nec precibus, nec minis ad cultum idolórum posset perducere; immo tanto magis incensam vidéret ad celebrandas christianæ fidei laudes, quanto magis ipse eam a senténtia avértere conabátur: Cessábunt, inquit, verba, cum ventum erit ad verbera. Cui Virgo: Dei servis verba deésse non possunt, quibus a Christo Dómino dictum est: Cum steteritis ante reges et præsides, nolíte cogitare quómodo aut quid loquámini; dábitur enim vobis in illa hora quid loquámini; non enim vos estis qui loquímini, sed Spíritus Sanctus, qui lóquitur in vobis.\par
\switchcolumn
\EN
\noindent When this came to the ears of one to whom her parents had betrothed her against her will, he accused Lucy before Paschasius, the Prefect of being a Christian. The Prefect could not move her to commit idolatry, either by his entreaties or his threats; nay, the more he strove to persuade her, so much the bolder did she become in her confession. Then, seeing that he could prevail nothing, words, saith he, will cease when we come to blows. To whom the virgin answered, God's servants will never want words, for the Lord Christ hath said: When ye shall stand before kings and governors, take no thought how or what ye shall speak, for it shall be given you in that same hour what ye shall speak, for it is not ye that speak, but the Holy Ghost Which speaketh in you. (Matth. x. 18-20; Mark xiii. 9-11.)\par
\switchcolumn*

\LA
\begin{Resp}\Rsign Rogávi Dóminum meum Jesum Christum, ut ignis iste non dominétur mei: \Star{} Et impetrávi a Dómino inducias martyrii mei.\end{Resp}
\begin{Resp}\Vsign Pro eo ut me diligerent, detrahébant mihi: ego autem orábam.\end{Resp}
\begin{Resp}\Rsign Et impetrávi a Dómino inducias martyrii mei.\end{Resp}
\switchcolumn
\EN
\begin{Resp}\Rsign I besought my Lord Jesus Christ that this fire might not take hold upon me; \Star{} And I obtained from the Lord that I should not finish my testimony for yet a while.\end{Resp}
\begin{Resp}\Vsign For so much as they loved me, so bitterly spake they against me: but I gave myself unto prayer.\end{Resp}
\begin{Resp}\Rsign And I obtained from the Lord that I should not finish my testimony for yet a while.\end{Resp}
\switchcolumn*

\LA
\LessonHead{Lectio VI}
\switchcolumn
\EN
\LessonHead{Lesson VI}
\switchcolumn*

\LA
\noindent Quam cum Paschasius interrogasset, Estne in te Spíritus Sanctus? respóndit: Caste et pie vivéntes templum sunt Spíritus Sancti. At ille, Jubébo te ad lupánar duci, ut te Spíritus Sanctus déserat. Cui virgo: Si invitam jusseris violari, castitas mihi duplicábitur ad corónam. Quare Paschasius ira inflammatus, Lúciam eo trahi jussit, ubi ejus virginitas violarétur: sed divinitus factum est, ut firma virgo ita consísteret, ut nulla vi de loco dimoveri posset. Quam ob rem præfectus circum ipsam, pice, resína, ac fervénti oleo perfusam, ignem accéndi imperávit: sed cum ne flamma quidem eam læderet, multis tormentis excruciátæ guttur gládio transfígitur. Quo vulnere accepto, Lúcia prædícens Ecclésiæ tranquillitátem, quæ futura erat Diocletiáno et Maximiáno mórtuis, Idibus Decembris spíritum Deo réddidit. Cujus corpus Syracusis sepúltum, deínde Constantinopolim, postremo Venetias translátum est.\par
\switchcolumn
\EN
\noindent Then Paschasius asked her, saying: Is the Holy Ghost in thee? Whereto she answered: They that live in chastity and piety are the temples of the Holy Ghost. Then, said he, I will send thee to be prostituted in a brothel, and get the Holy Ghost out of thee. To whom she made reply: Thou canst not prostitute my will. If thou cause this poor body to be violated, the crown of my soul's purity will be brighter through suffering. Then he bade them take her to the place of shame, but by the power of God it became impossible to move her. Whereupon, being inflamed with anger, he had pitch, resin, and boiling oil poured upon her, and then set on fire. But the fire did not take hold upon her. Therefore he practised many other cruelties upon her, and at last thrust a sword through her neck. When Lucy had received this wound, she began to speak of the peace of the Church, which it should enjoy after the death of Diocletian and Maximian, and presently returned her soul into the hands of God. She testified on the thirteenth day of December. Her body was buried at Syracuse, but afterwards taken to Constantinople, and lastly to Venice.\par
\switchcolumn*

\LA
\begin{Resp}\Rsign Grata facta est a Dómino in certámine, quia apud Deum et apud hómines glorificáta est: in conspéctu príncipis loquebátur sapiéntiam: \Star{} Et Dóminus ómnium diléxit eam.\end{Resp}
\begin{Resp}\Vsign Adjuvábit eam Deus vultu suo: Deus in médio ejus, non commovébitur.\end{Resp}
\begin{Resp}\Rsign Et Dóminus ómnium diléxit eam.\end{Resp}
\begin{Resp}\Vsign Glória Patri, et Fílio, \Star{} et Spirítui Sancto.\end{Resp}
\begin{Resp}\Rsign Et Dóminus ómnium diléxit eam.\end{Resp}
\switchcolumn
\EN
\begin{Resp}\Rsign The Lord made her to prevail in the battle, and she was glorified in the sight of God and man; she spake wisdom before princes; \Star{} And the Lord of all loved her.\end{Resp}
\begin{Resp}\Vsign God shall help her with His countenance; God is in the midst of her, she shall not be moved.\end{Resp}
\begin{Resp}\Rsign And the Lord of all loved her.\end{Resp}
\begin{Resp}\Vsign Glory be to the Father, and to the Son, \Star{} and to the Holy Ghost.\end{Resp}
\begin{Resp}\Rsign And the Lord of all loved her.\end{Resp}
\switchcolumn*

\end{paracol}

\par\needspace{10\baselineskip}\addvspace{1.2em}{\centering\small\textsc{Die 14 Decembris} · \textsc{14 December}\par}
\Office{De VII die infra Octavam Concept. Immac. Beatæ Mariæ Virginis}{De VII die infra Octavam Concept. Immac. Beatæ Mariæ Virginis}{Semiduplex}
\addcontentsline{toc}{subsection}{Die 14 Decembris · De VII die infra Octavam Concept. Immac. Beatæ Mariæ Virginis\enspace\textit{De VII die infra Octavam Concept. Immac. Beatæ Mariæ Virginis}}
\begin{paracol}{2}
\LA
\RefLine{Lectiones I–III: de Scriptura occurrente.}
\switchcolumn
\EN
\RefLine{Lessons I–III: from the Scripture of the day.}
\switchcolumn*

\LA
\LessonHead{Lectio IV}
\switchcolumn
\EN
\LessonHead{Lesson IV}
\switchcolumn*

\LA
\noindent Ex Bulla dogmatica Pii Papæ IX\par
\switchcolumn
\EN
\noindent It is known to all men, with what care this doctrine of the sinlessness of the conception of the Mother of God hath been handed down, set forth, and defended by the most distinguished Religious Orders, Theologians, Universities, and Doctors skilled in the things of God. All men know likewise how carefully Christian Bishops, even in their public teaching, have professed the doctrine that through the merits of Christ our Lord and Saviour, foreknown by God, the Holy Virgin Mary, Mother of God, was delivered from ever being the victim of original sin, but, on the contrary, had the fruits of redemption applied to her at the very moment of her Conception, and was therefore redeemed in a nobler way than others. But the weightiest fact of all is that the most holy Council of Trent, when, in accordance with the Holy Scriptures, as interpreted by the holy Fathers and the approved Councils, it decreed that all men are conceived in sin, expressly added that it did not mean thereby to say that the blessed and stainless Mary, Mother of God, did not form an exception to the rule. From this declaration of the Fathers of Trent it can clearly be drawn that there is nothing in the Bible, nothing in tradition, and nothing in the Fathers which can rightly be adduced against this prerogative of the most Blessed Virgin; nay, as far as circumstances demanded, they as much as declared her free from the original stain.\par
\switchcolumn*

\LA
\begin{Resp}\Rsign Ego ex ore Altíssimi prodívi, primogénita ante omnem creatúram: ego feci in cælis, ut orirétur lumen indefíciens. \Star{} Nondum erant abýssi, et ego jam concépta eram.\end{Resp}
\begin{Resp}\Vsign Deus enim creávit me in justítia, et apprehéndit manum meam, et servávit me.\end{Resp}
\begin{Resp}\Rsign Nondum erant abýssi, et ego jam concépta eram.\end{Resp}
\switchcolumn
\EN
\begin{Resp}\Rsign I came out of the mouth of the Most High, the first-begotten before every creature. I made the unfading light to arise in the heavens. \Star{} When there were no depths I was conceived.\end{Resp}
\begin{Resp}\Vsign For the Lord hath created me in righteousness, and hath held mine hand, and hath kept me.\end{Resp}
\begin{Resp}\Rsign When there were no depths I was conceived.\end{Resp}
\switchcolumn*

\LA
\LessonHead{Lectio V}
\switchcolumn
\EN
\LessonHead{Lesson V}
\switchcolumn*

\LA
\noindent Et re quidem vera hanc de immaculáta beatíssimæ Vírginis Conceptióne doctrinam quotídie magis gravíssimo Ecclésiæ sensu, magisterio, studio, sciéntia ac sapiéntia tam splendide explicatam, declaratam, confírmatam, et apud omnes catholici orbis pópulos ac natiónes mirándum in modum propagatam, in ipsa Ecclésia semper exstitisse véluti a majóribus acceptam, ac revelátæ doctrinæ charactére insígnitam, illustria venerandæ antiquitátis Ecclésiæ orientalis et occidentalis monuménta validíssime testántur. Equidem Patres Ecclésiæque scriptores, cæléstibus edocti eloquiis, nihil antiquius habuere, quam in libris ad explicandas Scripturas, vindicánda dógmata, erudiendosque fidéles elucubratis, summam Vírginis sanctitátem, dignitátem, atque ab omni peccáti labe integritátem, ejusque præcláram de teterrimo humani generis hoste victoriam, multis mirisque modis certatim prædicáre atque efferre.\par
\switchcolumn
\EN
\noindent In truth, this doctrine upon the Conception of the most Blessed Virgin is day by day more earnestly set forth by the graver thought of the Church, by her teaching, by her care, by her learning, and by her wisdom. It is explained, taught, confirmed, and wonderfully spread among all peoples and nations of the Catholic world. The Church hath received it from the Fathers, as a part of the original faith, attested strongly by the most ancient and venerable monuments of both the Eastern and Western Churches. Indeed, the Fathers and Ecclesiastical writers, learned in Holy Scripture, are marked by no more earnest feature than that in all their books and Scriptural Commentaries, written for the confirmation of doctrine, and the edification of the faithful, they do all in diverse ways preach and teach the excelling holiness of this Virgin, her dignity, her freedom from any stain of sin, and the glory of her victory over the dark enemy of our race.\par
\switchcolumn*

\LA
\begin{Resp}\Rsign Nihil inquinátum in eam incúrrit: \Star{} Candor est lucis ætérnæ et spéculum sine mácula.\end{Resp}
\begin{Resp}\Vsign Est enim hæc speciósior sole, et luci comparáta invenítur púrior.\end{Resp}
\begin{Resp}\Rsign Candor est lucis ætérnæ et spéculum sine mácula.\end{Resp}
\switchcolumn
\EN
\begin{Resp}\Rsign No defiled thing can fall into her; \Star{} She is the brightness of the everlasting light, and the unspotted mirror of the power of God.\end{Resp}
\begin{Resp}\Vsign For she is more beautiful than the sun, and being compared with the light, she is found before it.\end{Resp}
\begin{Resp}\Rsign She is the brightness of the everlasting light, and the unspotted mirror of the power of God.\end{Resp}
\switchcolumn*

\LA
\LessonHead{Lectio VI}
\switchcolumn
\EN
\LessonHead{Lesson VI}
\switchcolumn*

\LA
\noindent Quapropter enarrántes verba, quibus Deus præparáta renovandis mortálibus suæ pietátis remedia inter ipsa mundi primordia prænuntians; et deceptoris serpentis rétudit audaciam et nostri generis spem mirifice eréxit, inquiens, Inimicítias ponam inter te et mulierem, et semen tuum et semen illíus; docuére, divino hoc oraculo clare aperteque demonstrátum fuísse misericordem humani generis Redemptórem, scilicet unigénitum Dei Fílium Christum Jesum, ac designatam beatíssimam ejus Matrem Vírginem Mariam, ac simul ipsíssimas utriusque contra diabolum inimicítias insígniter expressas. Quocirca sicut Christus Dei hominúmque mediator, humana assumpta natúra, delens quod advérsus nos erat chirógraphum decreti, illud cruci triumphator affixit; sic sanctíssima Virgo, arctíssimo et indissolubili vinculo cum eo conjuncta, una cum illo et per illum, sempiternas contra venenósum serpéntem inimicítias exercens ac de ipso pleníssime triumphans, illíus caput immaculáto pede contrívit.\par
\switchcolumn
\EN
\noindent All Commentators on the Book of Genesis remark that passage where God at the very time of the Fall speaketh of the Atonement, to the confusion of the lying serpent, and the comfortable hope of man, and saith I will put enmity between thee and the woman, and between thy seed and her seed, and all the ancients teach that by this passage is meant the most merciful Saviour of mankind, namely, our Lord Jesus Christ, the OnlyBegotten Son of God and His most blessed Mother the Virgin Mary, as if the enmity which both He and she felt against the devil was, in a sense, of a kind common to them Both. Christ took our nature upon Him, and is become the Mediator between God and man, (Tim. ii. 5,) blotting out the handwriting that was against us, nailing it to His Cross, (Col. ii. 14,) and the most Holy Virgin, by that subtle, close, and abiding tie which bindeth mother to Child, feeleth along with Him His truceless enmity to the serpent, and He, through His merits, hath granted to her that moment of victory wherein her stainless foot bruised the serpent's head.\par
\switchcolumn*

\LA
\begin{Resp}\Rsign Signum magnum appáruit in cælo: Múlier amícta sole, et luna sub pédibus ejus, \Star{} Et in cápite ejus coróna stellárum duódecim.\end{Resp}
\begin{Resp}\Vsign Induit eam Dóminus vestiméntis salútis, induménto justítiæ, et quasi sponsam ornávit eam monílibus suis.\end{Resp}
\begin{Resp}\Rsign Et in cápite ejus coróna stellárum duódecim.\end{Resp}
\begin{Resp}\Vsign Glória Patri, et Fílio, \Star{} et Spirítui Sancto.\end{Resp}
\begin{Resp}\Rsign Et in cápite ejus coróna stellárum duódecim.\end{Resp}
\switchcolumn
\EN
\begin{Resp}\Rsign There appeared a great wonder in heaven; a Woman clothed with the sun, and the moon under her feet \Star{} And upon her head a crown of twelve stars.\end{Resp}
\begin{Resp}\Vsign The Lord hath clothed her with the garments of salvation, and hath covered her with the robe of righteousness, yea, as a bride He hath adorned her with jewels.\end{Resp}
\begin{Resp}\Rsign And upon her head a crown of twelve stars.\end{Resp}
\begin{Resp}\Vsign Glory be to the Father, and to the Son, \Star{} and to the Holy Ghost.\end{Resp}
\begin{Resp}\Rsign And upon her head a crown of twelve stars.\end{Resp}
\switchcolumn*

\LA
\LessonHead{Lectio VII}
\switchcolumn
\EN
\LessonHead{Lesson VII}
\switchcolumn*

\LA
\LessonTitle{Léctio sancti Evangélii secúndum Lucam}
\switchcolumn
\EN
\LessonTitle{From the Holy Gospel according to Luke}
\switchcolumn*

\LA
\Cite{Luc 1:26-28}
\switchcolumn
\EN
\Cite{Luke 1:26-28}
\switchcolumn*

\LA
\noindent In illo témpore: Missus est Angelus Gabriel a Deo in civitátem Galilææ, cui nomen Nazareth, ad Vírginem desponsatam viro, cui nomen erat Joseph, de domo David, et nomen Vírginis Maria. Et réliqua.\par
\switchcolumn
\EN
\noindent In that time the angel Gabriel was sent from God into a city of Galilee, called Nazareth, To a virgin espoused to a man whose name was Joseph, of the house of David; and the virgin's name was Mary. And so on.\par
\switchcolumn*

\LA
\LessonTitle{Homilía sancti Sophrónii Epíscopi.}
\switchcolumn
\EN
\LessonTitle{Homily by St. Sophronius, Patriarch (of Jerusalem.)}
\switchcolumn*

\LA
\Source{Homilia in Deiparæ Annunt.}
\switchcolumn
\EN
\Source{On the Annunciation.}
\switchcolumn*

\LA
\noindent Vere benedicta tu in muliéribus, quóniam Hevæ maledictiónem in benedictiónem commutásti; quóniam Adam, qui prius jacebat exsecratióne perculsus, ut per te benedicerétur effecísti. Vere benedicta tu in muliéribus, quóniam benedíctio Patris per te affulsit homínibus, eosque a veteri maledicto liberávit. Vere benedicta tu in muliéribus, quia per te progenitores tui salútem invéniunt; tu síquidem genitúra es Servatórem, qui divinam ipsis salútem comparabit. Vere benedicta tu in muliéribus, quóniam sine semine eum protulísti fructum, qui benedictiónem terrárum orbi elargítur, ipsumque a maledictióne spinas germinante rédimit. Vere benedicta tu in muliéribus, quia mulier naturali conditióne cum sis, Dei tamen Génetrix reipsa fies. Etenim si, qui ex te nascitúrus est, secúndum veritátem Deus est incarnatus, ipsa jure meritoque diceris Deipara, quippe quæ Deum veríssime paris.\par
\switchcolumn
\EN
\noindent Blessed indeed art thou among women, for thou hast turned the curse of Eve into a blessing; thou hast even brought a blessing upon Adam, when he lay smitten by the first sentence of death. Blessed indeed art thou among women, for thou art the mean whereby the Father's blessing hath come upon man, and delivered him from the old curse. Blessed indeed art thou among women, for by thee thy fathers have found salvation; the salutation of the Angel telleth thee that thou art about to bear them a Deliverer. Blessed indeed art thou among women, for thou, not knowing a man, conceivest a Son through Whom the whole earth shall be blessed, and bring forth thorns and thistles no more. Blessed indeed art thou among women, for thou remainest thyself no more than a woman, and yet art made Mother of God. If That holy Thing Which shall be born of thee be truly God made Man, then art thou truly Mother of God, for God is made thine Offspring.\par
\switchcolumn*

\LA
\begin{Resp}\Rsign Hortus conclúsus soror mea sponsa, hortus conclúsus, fons signátus: \Star{} Emissiónes tuæ paradísus, o María.\end{Resp}
\begin{Resp}\Vsign Aperi mihi, soror mea, amíca mea, colúmba mea, immaculáta mea.\end{Resp}
\begin{Resp}\Rsign Emissiónes tuæ paradísus, o María.\end{Resp}
\switchcolumn
\EN
\begin{Resp}\Rsign A garden enclosed is my sister, my spouse, a garden enclosed, a fountain sealed. \Star{} O Mary, thy perfumes are a garden of delights.\end{Resp}
\begin{Resp}\Vsign Open to me, my sister, my love, my dove, my undefiled.\end{Resp}
\begin{Resp}\Rsign O Mary, thy perfumes are a garden of delights.\end{Resp}
\switchcolumn*

\LA
\LessonHead{Lectio VIII}
\switchcolumn
\EN
\LessonHead{Lesson VIII}
\switchcolumn*

\LA
\noindent Ne timeas Maria, invenísti enim apud Deum grátiam perire nesciam; invenísti apud Deum grátiam præ quálibet exímiam; invenísti apud Deum grátiam ómnibus votis expeténdam; invenísti apud Deum grátiam ómnium gratiárum splendidíssimam; invenísti apud Deum grátiam numquam elanguéntem; invenísti apud Deum grátiam, quæ te salvam præstet; invenísti apud Deum grátiam, quæ nullo ímpetu quatiátur; invenísti apud Deum grátiam plene invictam; invenísti apud Deum grátiam perenniter duratúram. Et alii quidem iíque plures, ante te exímia sanctitáte floruérunt. Sed nemini, quemádmodum tibi, plena grátia impertita est: nemo sicut tu, ad tantum magnificéntiæ est evectus; nemo sicut tu, purificante grátia præoccupátus est; nemo sicut tu, cælésti lúmine refulsit; nemo sicut tu, præ omni celsitúdine exaltátus est.\par
\switchcolumn
\EN
\noindent Fear not, Mary, for thou hast found grace with God abiding grace. Thou hast found grace with God exceeding grace. Thou hast found grace with God all desirable grace. Thou hast found grace with God - greater grace than any other. Thou hast found grace with God unfailing grace. Thou hast found grace with God saving grace. Thou hast found grace with God immoveable grace. Thou hast found grace with God invincible grace. Thou hast found grace with God everlasting grace. Before thee there have been others, many others, made wonderful in holiness, but to none hath it been given, as to thee, to be full of grace; to none hath it been given, as to thee, to attain to such divine riches; to none, as to thee, to be prevented by purifying grace; to none, as to thee, to shine from the dayspring with light from heaven; to none, as to thee, to be exalted above all things before created.\par
\switchcolumn*

\LA
\begin{Resp}\Rsign Magníficat ánima mea Dóminum: \Star{} Quia fecit mihi magna qui potens est, et sanctum nomen ejus.\end{Resp}
\begin{Resp}\Vsign Ecce enim ex hoc beátam me dicent omnes generatiónes.\end{Resp}
\begin{Resp}\Rsign Quia fecit mihi magna qui potens est, et sanctum nomen ejus.\end{Resp}
\begin{Resp}\Vsign Glória Patri, et Fílio, \Star{} et Spirítui Sancto.\end{Resp}
\begin{Resp}\Rsign Quia fecit mihi magna qui potens est, et sanctum nomen ejus.\end{Resp}
\switchcolumn
\EN
\begin{Resp}\Rsign My soul doth magnify the Lord; \Star{} For He That is mighty hath done to me great things, and holy is His name.\end{Resp}
\begin{Resp}\Vsign For, behold, from henceforth all generations shall call me blessed.\end{Resp}
\begin{Resp}\Rsign For He That is mighty hath done to me great things, and holy is His name.\end{Resp}
\begin{Resp}\Vsign Glory be to the Father, and to the Son, \Star{} and to the Holy Ghost.\end{Resp}
\begin{Resp}\Rsign For He That is mighty hath done to me great things, and holy is His name.\end{Resp}
\switchcolumn*

\LA
\LessonHead{Lectio IX}
\switchcolumn
\EN
\LessonHead{Lesson IX}
\switchcolumn*

\LA
\noindent Et mérito quidem; nemo enim sicut tu, ad Deum tam prope accessit; nemo sicut tu, Dei donis ditátus est; nemo sicut tu, Dei gratiæ particeps fuit. Omnia vincis, quæcúmque inter hómines emineant; ómnia superas múnera, quæ effusa a Dei largitate in quóslibet dimanáverint. Plus enim ómnibus inhabitántis Dei possessióne ditescis. Nemo sic in se Deum complecti pótuit; nemo pótuit sic frui Dei præséntia; nemo dignus est habitus, qui sic per Deum illuminarétur: ac proptérea non modo Deum ómnium Conditórem ac Dóminum in te excepísti, sed eum habes ex te ineffabíliter incarnátum, atque in útero tuo gestátum, et post hæc génitum, et universos hómines a paterna condemnatióne rediméntem, ac salútem, quæ nullum finem habitúra sit, ipsis elargiéntem. Et proptérea tibi clamávi, atque íterum veheménter clamábo: Ave grátia plena; Dóminus tecum: benedicta tu in muliéribus.\par
\switchcolumn
\EN
\noindent And justly; for none hath ever drawn so near to God as thou hast; none hath ever been gifted by God with good gifts as thou hast; none hath ever received of God's grace as thou hast. Thou art mightier than all things which are called mighty among men; thou hast received more than the goodness of God hath conferred on any other. It is indeed because God made His home in thee that thou thus aboundest. There hath never been any save thee that hath comprehended the Incomprehensible; none save thee that hath enjoyed His presence so much; none that He hath made so ready therefor; none on whom the uncreated light hath shone so clearly; and therefore none who hath, like thee, sheltered the Lord God, the Maker and Lord of all, conceived Him in thy womb, and brought Him into the world, to redeem men lying under the Father's sentence, and to offer to them everlasting salvation. Wherefore, O Lady, I have already cried unto thee with the Angel, and full of grace, the Lord is with thee! Blessed art thou among women!\par
\switchcolumn*

\LA
\TeDeum{Te Deum}
\switchcolumn
\EN
\TeDeum{Te Deum}
\switchcolumn*

\end{paracol}

\par\needspace{10\baselineskip}\addvspace{1.2em}{\centering\small\textsc{Die 15 Decembris} · \textsc{15 December}\par}
\Office{In Octava Concept. Immac. Beatæ Mariæ Virginis}{In Octava Concept. Immac. Beatæ Mariæ Virginis}{Duplex majus}
\addcontentsline{toc}{subsection}{Die 15 Decembris · In Octava Concept. Immac. Beatæ Mariæ Virginis\enspace\textit{In Octava Concept. Immac. Beatæ Mariæ Virginis}}
\begin{paracol}{2}
\LA
\RefLine{Lectiones I–III: de Scriptura occurrente.}
\switchcolumn
\EN
\RefLine{Lessons I–III: from the Scripture of the day.}
\switchcolumn*

\LA
\LessonHead{Lectio IV}
\switchcolumn
\EN
\LessonHead{Lesson IV}
\switchcolumn*

\LA
\noindent Ex Bulla dogmatica Pii Papæ IX\par
\switchcolumn
\EN
\noindent From of old time continual prayers have been offered to this Apostolic See not only by Bishops, Churchmen, and the Regular Orders, but also by Emperors and Kings, beseeching that the sinlessness of the Conception of the Mother of God might be made the subject of a Dogmatic Definition. These prayers have been still more urgently addressed in recent times to Our Predecessor, of happy memory, Gregory XVI., and to Ourselves, by the Bishops, by the Secular Clergy, by the Regular Orders, and by the most eminent Christian Princes and nations. The knowledge of these things hath caused Our heart to rejoice, and they have been the serious occupation of Our thoughts ever since, in spite of our unworthiness, the inscrutable Providence of God was pleased to set Us in this supreme Chair of Peter, and to put His Church into Our hands, for Us to govern. Since then, We have had nothing so much at heart as to yield to the desires of the Church in this matter, to the increasing of the veneration of which the most holy Virgin is already the object, and to the setting in a clearer light the singular graces with which God hath adorned her, being Ourselves especially drawn to the same by the reverence, love, and affection with which We have been from Our childhood animated towards the same most holy Virgin Mary, Mother of God.\par
\switchcolumn*

\LA
\begin{Resp}\Rsign Ego ex ore Altíssimi prodívi, primogénita ante omnem creatúram: ego feci in cælis, ut orirétur lumen indefíciens. \Star{} Nondum erant abýssi, et ego jam concépta eram.\end{Resp}
\begin{Resp}\Vsign Deus enim creávit me in justítia, et apprehéndit manum meam, et servávit me.\end{Resp}
\begin{Resp}\Rsign Nondum erant abýssi, et ego jam concépta eram.\end{Resp}
\switchcolumn
\EN
\begin{Resp}\Rsign I came out of the mouth of the Most High, the first-begotten before every creature. I made the unfading light to arise in the heavens. \Star{} When there were no depths I was conceived.\end{Resp}
\begin{Resp}\Vsign For the Lord hath created me in righteousness, and hath held mine hand, and hath kept me.\end{Resp}
\begin{Resp}\Rsign When there were no depths I was conceived.\end{Resp}
\switchcolumn*

\LA
\LessonHead{Lectio V}
\switchcolumn
\EN
\LessonHead{Lesson V}
\switchcolumn*

\LA
\noindent Itaque plurimum in Dómino confisi advenisse témporum opportunitátem pro immaculáta sanctíssimæ Dei Genitrícis Vírginis Maríæ Conceptióne definienda, quam divina elóquia, veneranda traditio, perpetuus Ecclésiæ sensus, singuláris catholicórum antístitum ac fidelium conspirátio, et insígnia prædecessórum nostrórum acta et constitutiónes mirifice illustrant atque declarant; rebus ómnibus diligentíssime perpénsis et assiduis fervidisque ad Deum precibus effúsis, minime cunctándum nobis esse censúimus, supremo nostro judício immaculátam ipsíus Vírginis Conceptiónem sancire, definire, atque ita pientíssimis catholici orbis desidériis, nostræque in ipsam sanctíssimam Vírginem pietáti satisfácere, ac simul in ipsa unigénitum Fílium suum Dóminum nostrum Jesum Christum magis atque magis honorificare, cum in Fílium redundet quidquid honoris et laudis in Matrem impénditur.\par
\switchcolumn
\EN
\noindent And now We trust in the Lord that the time is come to define as a truth of faith the doctrine of the stainless Conception of the most holy Virgin Mary, Mother of God, that doctrine already set forth by Holy Scripture, by the ancient tradition, by the unbroken belief of the Universal Church, by the one common opinion of the Catholic Episcopate and laity, and by the marked acts and decrees of Our Predecessors. We have weighed everything in Our mind, and We have without ceasing implored the help and light of God's Holy Spirit by earnest prayer, and We are of opinion that it is Our duty no longer to delay, but by Our Supreme Decision to settle and declare that the Virgin was conceived without sin, and thus to satisfy the godly cravings of the whole Christian world, as well as the instinct of Our own love for the said most holy Virgin, and, above all, because she is His Mother, to glorify our Lord Jesus Christ by this act, since whatever we do rightly to honour the Mother must redound to the glory of the Son.\par
\switchcolumn*

\LA
\begin{Resp}\Rsign Nihil inquinátum in eam incúrrit: \Star{} Candor est lucis ætérnæ et spéculum sine mácula.\end{Resp}
\begin{Resp}\Vsign Est enim hæc speciósior sole, et luci comparáta invenítur púrior.\end{Resp}
\begin{Resp}\Rsign Candor est lucis ætérnæ et spéculum sine mácula.\end{Resp}
\switchcolumn
\EN
\begin{Resp}\Rsign No defiled thing can fall into her; \Star{} She is the brightness of the everlasting light, and the unspotted mirror of the power of God.\end{Resp}
\begin{Resp}\Vsign For she is more beautiful than the sun, and being compared with the light, she is found before it.\end{Resp}
\begin{Resp}\Rsign She is the brightness of the everlasting light, and the unspotted mirror of the power of God.\end{Resp}
\switchcolumn*

\LA
\LessonHead{Lectio VI}
\switchcolumn
\EN
\LessonHead{Lesson VI}
\switchcolumn*

\LA
\noindent Quare postquam numquam intermísimus in humilitate et jejúnio privatas nostras et publicas Ecclésiæ preces Deo Patri per Fílium ejus offerre, ut Spíritus Sancti virtúte mentem nostram dirigere et confírmare dignarétur, implorato univérsæ cæléstis curiæ præsidio, et advocato cum gemítibus Paraclito Spiritu, eoque sic aspirante; ad honórem sanctæ et indivíduæ Trinitátis, ad decus et ornaméntum Vírginis Deiparæ, ad exaltatiónem fidei catholicæ et christianæ religiónis augméntum, auctoritate Dómini nostri Jesu Christi, beatórum Apostolórum Petri et Pauli, ac nostra, declaramus, pronuntiámus et definímus: Doctrinam quæ tenet beatíssimam Vírginem Mariam in primo instanti suæ Conceptiónis fuísse singulari omnipoténtis Dei grátia et privilegio, intuitu meritórum Christi Jesu Salvatoris humani generis, ab omni originalis culpæ labe præservatam immunem, esse a Deo revelátum, atque idcirco ab ómnibus fidelibus firmiter constanterque credéndam. Quapropter si qui secus ac a nobis definítum est, quod Deus avértat, præsumpserint corde sentire, ii nóverint ac porro sciant se proprio judício condemnatos, naufragium circa fidem passos esse, et ab unitate Ecclésiæ defecisse.\par
\switchcolumn
\EN
\noindent Therefore, having in all fasting and humbleness of heart continually implored God the Father through His Blessed Son, to hear Our own prayers and those of His whole Church, and to teach and strengthen Our mind by the power of His Holy Spirit, having begged the intercession of the Church Triumphant, and, above all, with groans called on the Holy Ghost the Comforter, We now, moved by Him, for the honour of the Holy and Undivided Trinity, for the greater praise and exaltation of the Virgin Mother of God, the glory of the Catholic Faith, and the good of Christianity, in the name and authority of our Lord Jesus Christ, of the Blessed Apostles Peter and Paul, and in Our own, declare and define that That doctrine which declareth that the most blessed Virgin Mary was, in the first instant of her Conception, preserved, by a special grace and privilege granted to her by Almighty God, through the merits of Christ Jesus, Saviour of mankind, which He foreknew, from any stain of original sin, is a doctrine taught and revealed by God, and therefore from this time forward must be held by all faithful Christians firmly and constantly. Wherefore if any one, which God forbid, shall at any time think in his heart any thing contrary to this Our definition, let him know that he is condemned by his own judgment, that he hath made shipwreck of the faith, and that he hath cut himself off from the body of the Church.\par
\switchcolumn*

\LA
\begin{Resp}\Rsign Signum magnum appáruit in cælo: Múlier amícta sole, et luna sub pédibus ejus, \Star{} Et in cápite ejus coróna stellárum duódecim.\end{Resp}
\begin{Resp}\Vsign Induit eam Dóminus vestiméntis salútis, induménto justítiæ, et quasi sponsam ornávit eam monílibus suis.\end{Resp}
\begin{Resp}\Rsign Et in cápite ejus coróna stellárum duódecim.\end{Resp}
\begin{Resp}\Vsign Glória Patri, et Fílio, \Star{} et Spirítui Sancto.\end{Resp}
\begin{Resp}\Rsign Et in cápite ejus coróna stellárum duódecim.\end{Resp}
\switchcolumn
\EN
\begin{Resp}\Rsign There appeared a great wonder in heaven; a Woman clothed with the sun, and the moon under her feet \Star{} And upon her head a crown of twelve stars.\end{Resp}
\begin{Resp}\Vsign The Lord hath clothed her with the garments of salvation, and hath covered her with the robe of righteousness, yea, as a bride He hath adorned her with jewels.\end{Resp}
\begin{Resp}\Rsign And upon her head a crown of twelve stars.\end{Resp}
\begin{Resp}\Vsign Glory be to the Father, and to the Son, \Star{} and to the Holy Ghost.\end{Resp}
\begin{Resp}\Rsign And upon her head a crown of twelve stars.\end{Resp}
\switchcolumn*

\LA
\LessonHead{Lectio VII}
\switchcolumn
\EN
\LessonHead{Lesson VII}
\switchcolumn*

\LA
\LessonTitle{Léctio sancti Evangélii secúndum Lucam}
\switchcolumn
\EN
\LessonTitle{From the Holy Gospel according to Luke}
\switchcolumn*

\LA
\Cite{Luc 1:26-28}
\switchcolumn
\EN
\Cite{Luke 1:26-28}
\switchcolumn*

\LA
\noindent In illo témpore: Missus est Angelus Gabriel a Deo in civitátem Galilææ, cui nomen Nazareth, ad Vírginem desponsatam viro, cui nomen erat Joseph, de domo David, et nomen Vírginis Maria. Et réliqua.\par
\switchcolumn
\EN
\noindent In that time the angel Gabriel was sent from God into a city of Galilee, called Nazareth, To a virgin espoused to a man whose name was Joseph, of the house of David; and the virgin's name was Mary. And so on.\par
\switchcolumn*

\LA
\LessonTitle{Homilía sancti Epiphánii Epíscopi}
\switchcolumn
\EN
\LessonTitle{Homily by St. Epiphanius, Bishop (of Salamis-in-Cyprus.)}
\switchcolumn*

\LA
\Source{Oratio de laudibus S. Maríæ Deiparæ}
\switchcolumn
\EN
\Source{Discourse in praise of the Blessed Virgin.}
\switchcolumn*

\LA
\noindent Quid dicam, aut quid próloquar de præclára et sancta Vírgine? Solo enim Deo excepto, cunctis superior éxstitit: natúra formosior est ipsis Chérubim et Séraphim, et omni exercitu Angelico: cui prædicandæ cæléstis ac terrena lingua minime sufficit, immo vero nec Angelórum. O beáta Virgo, columba pura et sponsa cæléstis, Maria, cælum, templum, et thronus divinitátis, quæ coruscántem in cælo et in terra solem habes Christum! Nubes lucida, quæ fulgur de cælo lucidíssimum ad illuminándum mundum deduxísti Christum. Ave, grátia plena, porta cælórum, de qua in Canticis prophéta in decursu oratiónis plane et aperte prolóquitur inclámans: Hortus conclusus soror mea sponsa, hortus conclusus, fons signátus.\par
\switchcolumn
\EN
\noindent I am at a loss what words or terms I ought to employ in speaking of this illustrious and holy Virgin. She is raised above all things except God; she was made much higher than the Cherubim and Seraphim, and the whole host of heaven; neither the voices of heaven nor of earth are full enough to set forth her majesty, no, not the voices of Angels. O blessed Virgin! O pure dove and Bride of heaven! O Mary! At once the heaven, the temple, and the throne of God! Mother of the Sun that shineth both on heaven and on earth, even Christ! Bright cloud, through which the Son of Man hath come as the lightning, that lighteneth from the East even unto the West! Hail, gate of heaven, full of grace, of whom the Prophet in the Song of Songs openly speaketh in the course of his prayer, saying, A garden enclosed is My sister, My Spouse, a garden enclosed, a fountain sealed. \textasciitilde{}(iv. 12.)\par
\switchcolumn*

\LA
\begin{Resp}\Rsign Hortus conclúsus soror mea sponsa, hortus conclúsus, fons signátus: \Star{} Emissiónes tuæ paradísus, o María.\end{Resp}
\begin{Resp}\Vsign Aperi mihi, soror mea, amíca mea, colúmba mea, immaculáta mea.\end{Resp}
\begin{Resp}\Rsign Emissiónes tuæ paradísus, o María.\end{Resp}
\switchcolumn
\EN
\begin{Resp}\Rsign A garden enclosed is my sister, my spouse, a garden enclosed, a fountain sealed. \Star{} O Mary, thy perfumes are a garden of delights.\end{Resp}
\begin{Resp}\Vsign Open to me, my sister, my love, my dove, my undefiled.\end{Resp}
\begin{Resp}\Rsign O Mary, thy perfumes are a garden of delights.\end{Resp}
\switchcolumn*

\LA
\LessonHead{Lectio VIII}
\switchcolumn
\EN
\LessonHead{Lesson VIII}
\switchcolumn*

\LA
\noindent Virgo est lílium immaculátum, quæ rosam immarcescíbilem génuit Christum. O sancta Deipara, ovis immaculáta quæ Verbum ex te incarnátum Agnum Christum peperisti! O Virgo sanctíssima, quæ exercitus Angelórum in stupórem dedúxisti! Stupéndum est miraculum in cælis, mulier amícta sole, gestans lucem in ulnis: stupéndum miraculum in cælis, thálamus Vírginis habens Fílium Dei: stupéndum miraculum in cælis, Dóminus Angelórum infans Vírginis effectus est. Angeli accusábant Hevam, nunc vero Mariam glória prosequúntur, quæ lapsam Hevam eréxit, et Adámum e paradiso dejectum in cælos misit. Ipsa enim est cæli et terræ mediatrix, quæ uniónem naturáliter peregit.\par
\switchcolumn
\EN
\noindent The Virgin is that stainless lily whence hath sprung the Rose that fadeth not, even Christ. O Holy Mother of God! Ewe without spot, that hast borne the Lamb That took flesh of thee, even Christ! O Maiden whose holiness hath dazzled the heavenly armies! There hath appeared a great sign in heaven, a woman clothed with the sun, and with the Light in her arms; a great sign in heaven, the Virgin's womb the chamber of the Son of God; a great sign in heaven, the Lord of angels made the Virgin's child. The angels accused Eve, but now they praise Mary, who hath raised Eve fallen, and restored to heaven Adam banished from Paradise. For Mary is the bridge between heaven and earth, the ambassadress who hath reconciled them in her womb.\par
\switchcolumn*

\LA
\begin{Resp}\Rsign Magníficat ánima mea Dóminum: \Star{} Quia fecit mihi magna qui potens est, et sanctum nomen ejus.\end{Resp}
\begin{Resp}\Vsign Ecce enim ex hoc beátam me dicent omnes generatiónes.\end{Resp}
\begin{Resp}\Rsign Quia fecit mihi magna qui potens est, et sanctum nomen ejus.\end{Resp}
\begin{Resp}\Vsign Glória Patri, et Fílio, \Star{} et Spirítui Sancto.\end{Resp}
\begin{Resp}\Rsign Quia fecit mihi magna qui potens est, et sanctum nomen ejus.\end{Resp}
\switchcolumn
\EN
\begin{Resp}\Rsign My soul doth magnify the Lord; \Star{} For He That is mighty hath done to me great things, and holy is His name.\end{Resp}
\begin{Resp}\Vsign For, behold, from henceforth all generations shall call me blessed.\end{Resp}
\begin{Resp}\Rsign For He That is mighty hath done to me great things, and holy is His name.\end{Resp}
\begin{Resp}\Vsign Glory be to the Father, and to the Son, \Star{} and to the Holy Ghost.\end{Resp}
\begin{Resp}\Rsign For He That is mighty hath done to me great things, and holy is His name.\end{Resp}
\switchcolumn*

\LA
\LessonHead{Lectio IX}
\switchcolumn
\EN
\LessonHead{Lesson IX}
\switchcolumn*

\LA
\noindent Gratia sanctæ Vírginis est immensa. Hinc Gabriel in primis salutat Vírginem, dicens: Ave, grátia plena, quæ es splendidum cælum. Ave, grátia plena, multis virtútibus exornáta Virgo. Ave, grátia plena, quæ es urna aurea cóntinens manna cæleste. Ave, grátia plena, quæ sitiéntes perennis fontis dulcédine satias. Ave, sanctíssima Mater immaculáta, quæ genuísti Christum, qui est ante te. Ave, púrpura regia, quæ cæli terræque Regem induísti. Ave, liber incomprehénsus, quæ Verbum et Fílium Patris mundo legéndum exhibuisti.\par
\switchcolumn
\EN
\noindent We cannot measure the grace bestowed upon this holy Virgin. Hence the salutation addressed to her by Gabriel, Hail, thou glorious heaven, full of grace. Hail, Virgin adorned with many graces, yea, full of grace. Hail, thou vessel of gold that holdest the manna that came down from heaven, full of grace. Hail, thou unfailing fountain, that satisfiest the thirsty soul with sweet waters, full of grace. Hail, holy, sinless Mother of Him That was before thee, even Christ. Hail, thou Queenly purple, mantle of the King of heaven and earth. Hail, thou Book that no man can understand, and yet which the Eternal Word, the Son of the Father, hath opened for earth to read.\par
\switchcolumn*

\LA
\TeDeum{Te Deum}
\switchcolumn
\EN
\TeDeum{Te Deum}
\switchcolumn*

\end{paracol}

\par\needspace{10\baselineskip}\addvspace{1.2em}{\centering\small\textsc{Die 16 Decembris} · \textsc{16 December}\par}
\Office{S. Eusebii Episcopi et Martyris}{St. Eusebius, Bishop and Martyr}{Semiduplex}
\addcontentsline{toc}{subsection}{Die 16 Decembris · S. Eusebii Episcopi et Martyris\enspace\textit{St. Eusebius, Bishop and Martyr}}
\begin{paracol}{2}
\LA
\RefLine{Lectiones I–III: de Scriptura occurrente.}
\switchcolumn
\EN
\RefLine{Lessons I–III: from the Scripture of the day.}
\switchcolumn*

\LA
\LessonHead{Lectio IV}
\switchcolumn
\EN
\LessonHead{Lesson IV}
\switchcolumn*

\LA
\noindent Eusebius, natióne Sardus, Romanæ urbis lector, post Vercellénsis epíscopus, ad hanc regéndam ecclésiam mérito est créditus divino eléctus judício: nam quem numquam ante constitúti electóres cognoverant, posthábitis civibus, simul ut vidérunt, et probavérunt; tantumque intérfuit, ut probarétur, quantum ut viderétur. Primus in Occidéntis partibus in eadem ecclésia eosdem monachos instituit esse quos clericos, ut esset in ipsis viris et contemptus rerum, et accurátio levitárum. Ariánis impietátibus ea tempestáte per Occidéntem longe lateque traductis, advérsus eas viríliter sic dimicávit, ut ejus invicta fides Libérium summum Pontificem ad vitæ solatium erigeret. Quare hic sciens in ipso fervére Spíritum Dei, cum ei significasset ut penes imperatórem una cum suis legátis patrocinium fidei susciperet, mox cum illis profectus est ad Constántium; apud quem enixius agens, quidquid legatióne petebátur, obtinuit, ut episcopórum nempe cœtus celebrarétur.\par
\switchcolumn
\EN
\noindent Eusebius was a Sardinian by birth, first a Reader in the Roman Church, and then Bishop of Vercelli. It seemed specially designed by Providence that he should be called to govern that Church, for the electors, who had never before known him, passed over, with a strange unanimity, all their own fellow-citizens, and chose Eusebius, as soon as they had seen him. He was the first Bishop in the Western Church who established an Order of Regular Clergy, to combine the active with the contemplative life. At this time the storm of Arian blasphemy and sin was sweeping far and wide over the West, and Eusebius set himself to fight against it so manfully, that his unshaken faith brought back Liberius again to life and hope. This Pope, knowing that the Spirit of God was in him, sent him with his Legates to the Emperor Constantius, (in the year 354,) to plead the cause of the Catholic Faith. His earnestness prevailed with that Prince, so that he obtained all that was asked for, and, among other things, permission for a Council to be summoned.\par
\switchcolumn*

\LA
\begin{Resp}\Rsign Honéstum fecit illum Dóminus, et custodívit eum ab inimícis, et a seductóribus tutávit illum: \Star{} Et dedit illi claritátem ætérnam.\end{Resp}
\begin{Resp}\Vsign Descendítque cum illo in fóveam, et in vínculis non derelíquit eum.\end{Resp}
\begin{Resp}\Rsign Et dedit illi claritátem ætérnam.\end{Resp}
\switchcolumn
\EN
\begin{Resp}\Rsign The Lord made him honourable, and defended him from his enemies, and kept him safe from those that lay in wait for him, \Star{} And gave him perpetual glory.\end{Resp}
\begin{Resp}\Vsign He went down with him into the pit, and left him not in bonds.\end{Resp}
\begin{Resp}\Rsign And gave him perpetual glory.\end{Resp}
\switchcolumn*

\LA
\LessonHead{Lectio V}
\switchcolumn
\EN
\LessonHead{Lesson V}
\switchcolumn*

\LA
\noindent Colléctum est Mediolani anno sequénti concílium, ad quod a Constantio invitátum Eusebium concupitúmque, ac vocátum a Liberii legatis, tantum abest, ut malignántium synagóga Arianórum contra sanctum Athanasium furéntium in suas partes adduceret, ut potius diserte statim ipse declarans, e præséntibus quosdam sibi compertos hæretica labe pollutos, Nicænam immo fidem proposúerit iis subscribéndam, ántequam cetera tractaréntur. Quod Ariánis acerbe irátis negántibus, nedum in Athanasium recusávit ipse subscribere, quin sancti Dionysii Mártyris, qui decéptus ab ipsis subscripserat, captivatam simplicitátem ingeniosíssime liberávit. Quam ob rem illi gráviter indignántes, post multas illatas injurias, exsílio illum mulctarunt: sed sanctus vir, excusso púlvere, nec cæsaris minas veritus, nec enses obstrictos, exsílium véluti sui ministerii offícium accépit; missusque Scythópolim, famem, sitim, verbera, divérsaque supplícia perpéssus, pro fide strenue vitam contempsit, mortem non metuit, corpus carnificibus trádidit.\par
\switchcolumn
\EN
\noindent The year following, the Council met at Milan, and Eusebius, by the invitation of the Emperor, and the desire and command of the Papal Legates, attended. Here the Arians, assembled in a perfect synagogue of Satan, and all furiously raging together against holy Athanasius, found Eusebius one of the stoutest enemies of their faction. As soon as he entered the Council, he delivered a long harangue, wherein he remarked that, of those there gathered together, some were notoriously defiled with heresy, and therefore he proposed that everyone should first of all subscribe the Nicene Creed, before proceeding to any other business. The Arians, in a violent passion, refused, whereupon he on his part refused to subscribe any proceedings against Athanasius, and even skilfully procured the withdrawal of the signature of the holy martyr Denys, then Bishop of Milan, which they had lyingly procured by practising on his simplicity. The Arians were now entirely enraged, and, after many persecutions, procured a decree of banishment against Eusebius. The Saint shook off the dust of his feet against them, and, defying alike the threats of Caesar and the drawn swords of the soldiery, accepted the sentence as one of the dignities of his office. He was sent to Bethshan in the Holy Land, suffering hunger, thirst, stripes, and all manner of violence, but for the Faith's sake he despised this life, and feared not death, but freely delivered his body to the tormentors.\par
\switchcolumn*

\LA
\begin{Resp}\Rsign Desidérium ánimæ ejus tribuísti ei Dómine, \Star{} Et voluntáte labiórum ejus non fraudásti eum.\end{Resp}
\begin{Resp}\Vsign Quóniam prævenísti eum in benedictiónibus dulcédinis: posuísti in cápite ejus corónam de lápide pretióso.\end{Resp}
\begin{Resp}\Rsign Et voluntáte labiórum ejus non fraudásti eum.\end{Resp}
\switchcolumn
\EN
\begin{Resp}\Rsign Lord, Thou hast given him his heart's desire. \Star{} And hast not withholden the request of his lips.\end{Resp}
\begin{Resp}\Vsign For Thou hast prevented him with the blessings of sweetness: Thou hast set a crown of precious stones upon his head.\end{Resp}
\begin{Resp}\Rsign And hast not withholden the request of his lips.\end{Resp}
\switchcolumn*

\LA
\LessonHead{Lectio VI}
\switchcolumn
\EN
\LessonHead{Lesson VI}
\switchcolumn*

\LA
\noindent Quanta in eum tunc Arianórum crudelitas fúerit ac effrons inverecúndia, ostendunt graves litteræ plenæ róboris, pietátis ac religiónis, quas e Scythópoli scripsit ad Vercellensem clerum et pópulum, aliosque finítimos; e quibus étiam est explorátum, ipsórum nec minis, inhumanáque sævítia potuísse umquam eum deterreri, nec serpentina blanda subtilitate ad eórum societátem perduci. Hinc in Cappadociam, postremoque ad superiores Ægypti Thebáidas pro constántia sua deportatus, exsilii rigores tulit ad mortem usque Constantii: post quam ad gregem suum reverti permissus, non prius redire vóluit, quam reparandis fidei jacturis ad Alexandrinam synodum sese conferret, postque médici præstántis instar, péragrans Oriéntis provincias, in fide infirmos ad integram valetúdinem restitúeret, eos instítuens in Ecclésiæ doctrina. Inde salubritate pari digresso in Illyricum, tandemque in Italiam delato, ad ejus réditum lúgubres vestes Italia mutávit: ubi postquam Psalmórum ómnium expurgatos a se commentarios Orígenis édidit, Eusebiíque Cæsareénsis, quos verterat de Græco in Latinum; demum tot egregie factis illustris ad immarcescíbilem glóriæ corónam tantis ærumnis proméritam, sub Valentiniáno et Valénte Vercellis migrávit.\par
\switchcolumn
\EN
\noindent He wrote a solemn letter from Bethshan, addressed to the clergy and people of Vercelli and that neighbourhood, full of constancy, devotion, and piety, describing the frightful cruelty and brazen impudence of the Arians. From this letter we know how completely they failed to scare him by their threats and their inhuman brutality, or to seduce him by their serpent-like cunning into receiving their communion. In consequence of his unshaken resolution, he was moved from Bethshan into Cappadocia, and then again, to the deserts of Upper Egypt. He suffered exile until the death of Constantius, (in 361,) after which he was allowed to return to his flock. First, however, he took care to attend the Council at Alexandria, called to heal the wounds of the Church, and, afterwards, like a skilful physician, he made a progress through all the provinces of the East, strengthening those that were weak in the Faith, and confirming them in Christian doctrine. Then, with the same healthful results, he passed through Illyricum into Italy, who, at his coming, laid aside her garments of mourning. After his return, he published an expurgated edition of Origen's Commentary on the Psalms, and likewise of the works of Eusebius of Caesarea, both which he translated from Greek into Latin. At length, distinguished by all these great works, he passed to that crown of glory which fadeth not away, promised to them who suffer for the truth. He departed this present life at Vercelli, (in 371,) in the reign of Valentinian and Valens.\par
\switchcolumn*

\LA
\begin{Resp}\Rsign Stola jucunditátis índuit eum Dóminus: \Star{} Et corónam pulchritúdinis pósuit super caput ejus.\end{Resp}
\begin{Resp}\Vsign Cibávit illum Dóminus pane vitæ et intelléctus: et aqua sapiéntiæ salutáris potávit illum.\end{Resp}
\begin{Resp}\Rsign Et corónam pulchritúdinis pósuit super caput ejus.\end{Resp}
\begin{Resp}\Vsign Glória Patri, et Fílio, \Star{} et Spirítui Sancto.\end{Resp}
\begin{Resp}\Rsign Et corónam pulchritúdinis pósuit super caput ejus.\end{Resp}
\switchcolumn
\EN
\begin{Resp}\Rsign The Lord hath put on him a robe of honour, \Star{} And put about his head a crown of joy.\end{Resp}
\begin{Resp}\Vsign With the bread of life and understanding hath the Lord fed him, and given him the water of health and wisdom to drink.\end{Resp}
\begin{Resp}\Rsign And put about his head a crown of joy.\end{Resp}
\begin{Resp}\Vsign Glory be to the Father, and to the Son, \Star{} and to the Holy Ghost.\end{Resp}
\begin{Resp}\Rsign And put about his head a crown of joy.\end{Resp}
\switchcolumn*

\LA
\LessonHead{Lectio VII}
\switchcolumn
\EN
\LessonHead{Lesson VII}
\switchcolumn*

\LA
\LessonTitle{Léctio sancti Evangélii secúndum Matthǽum}
\switchcolumn
\EN
\LessonTitle{From the Holy Gospel according to Matthew}
\switchcolumn*

\LA
\Cite{Matt 16:24-27}
\switchcolumn
\EN
\Cite{Matt 16:24-27}
\switchcolumn*

\LA
\noindent In illo témpore: Dixit Jesus discípulis suis: Si quis vult post me veníre, ábneget semetípsum, et tollat crucem suam, et sequátur me. Et réliqua.\par
\switchcolumn
\EN
\noindent At that time, Jesus said unto His disciples: If any man will come after Me, let him deny himself, and take up his cross, and follow Me. And so on.\par
\switchcolumn*

\LA
\LessonTitle{Homilía sancti Gregórii Papæ}
\switchcolumn
\EN
\LessonTitle{Homily by Pope St. Gregory (the Great.)}
\switchcolumn*

\LA
\Source{Homilia 32 in Evang.}
\switchcolumn
\EN
\Source{32nd on the Gospels}
\switchcolumn*

\LA
\noindent Quia Dóminus ac Redémptor noster novus homo venit in mundum, nova præcépta dedit mundo. Vitæ étenim nostræ véteri in vítiis enutrítæ contrarietátem oppósuit novitátis suæ. Quid enim vetus, quid carnális homo nóverat, nisi sua retinére, aliéna rápere, si posset; concupíscere, si non posset? Sed cæléstis médicus síngulis quibúsque vítiis obviántia ádhibet medicaménta. Nam sicut arte medicínæ cálida frígidis, frígida cálidis curántur: ita Dóminus noster contrária oppósuit medicaménta peccátis, ut lúbricis continéntiam, tenácibus largitátem, iracúndis mansuetúdinem, elátis præcíperet humilitátem.\par
\switchcolumn
\EN
\noindent Our Lord and Redeemer came into the world a new Man, and gave the world new commandments. For against the ways of our old life, brought and bred up in sin, He set the contrast of His new life. It was the old way, according to the knowledge of the carnal man, for every man to keep his own goods, and, if he were able to do it, to take his neighbour's goods also, and, if he were not able to take them, at least to lust after them. But the Heavenly Physician hath medicines wherewith to meet all the diseases of sin. For, even, as by the art of the physician, things hot are healed by things cold, and things cold by things hot, so doth our Lord set against sin holiness, ordaining for the lecherous purity, for the miserly munificence, for the hot-tempered meekness, and for the proud lowliness.\par
\switchcolumn*

\LA
\begin{Resp}\Rsign Coróna áurea super caput ejus, \Star{} Expréssa signo sanctitátis, glória honóris, et opus fortitúdinis.\end{Resp}
\begin{Resp}\Vsign Quóniam prævenísti eum in benedictiónibus dulcédinis, posuísti in cápite ejus corónam de lápide pretióso.\end{Resp}
\begin{Resp}\Rsign Expréssa signo sanctitátis, glória honóris, et opus fortitúdinis.\end{Resp}
\switchcolumn
\EN
\begin{Resp}\Rsign A crown of gold upon his head, \Star{} Wherein is engraved Holiness, an ornament of honour, a costly work.\end{Resp}
\begin{Resp}\Vsign For Thou hast prevented him with the blessings of sweetness, Thou hast set a crown of precious stones upon his head.\end{Resp}
\begin{Resp}\Rsign Wherein is engraved Holiness, an ornament of honour, a costly work.\end{Resp}
\switchcolumn*

\LA
\LessonHead{Lectio VIII}
\switchcolumn
\EN
\LessonHead{Lesson VIII}
\switchcolumn*

\LA
\noindent Certe cum se sequéntibus nova mandáta propóneret, dixit: Nisi quis renuntiáverit ómnibus quæ póssidet, non potest meus esse discípulus. Ac si apérte dicat: Qui per vitam véterem aliéna concupíscitis, per novæ conversatiónis stúdium et vestra largímini. Quid vero in hac lectióne dicat, audiámus: Qui vult post me veníre, ábneget semetípsum. Ibi dícitur, ut abnegémus nostra: hic dícitur, ut abnegémus nos. Et fortásse laboriósum non est hómini relínquere sua; sed valde laboriósum est relínquere semetípsum. Minus quippe est abnegáre quod habet; valde autem multum est abnegáre quod est.\par
\switchcolumn
\EN
\noindent So the Lord, when He would give a new commandment unto them that came to Him, said Whosoever he be of you that forsaketh not all that he hath, he cannot be My disciple, (Luke xiv. 33,) as though He had said openly: All ye that according to the old man lust after your neighbour's goods, must, according to the zeal of the new man, give away even that which is your own. But let us hear again what He saith in this place If any man will come after Me, let him deny himself. First He saith that we must deny to ourselves that which is our own, and now that we must even deny ourselves to ourselves. Perchance it is not hard for a man to give up that which is his own, but it is exceeding hard to give up himself. To deny himself his possessions is little but to deny himself himself is a denial exceeding great.\par
\switchcolumn*

\LA
\begin{Resp}\Rsign Dómine, prævenísti eum in benedictiónibus dulcédinis: \Star{} Posuísti in cápite ejus corónam de lápide pretióso.\end{Resp}
\begin{Resp}\Vsign Vitam pétiit a te, et tribuísti ei longitúdinem diérum in sǽculum sǽculi.\end{Resp}
\begin{Resp}\Rsign Posuísti in cápite ejus corónam de lápide pretióso.\end{Resp}
\begin{Resp}\Vsign Glória Patri, et Fílio, \Star{} et Spirítui Sancto.\end{Resp}
\begin{Resp}\Rsign Posuísti in cápite ejus corónam de lápide pretióso.\end{Resp}
\switchcolumn
\EN
\begin{Resp}\Rsign O Lord, thou hast prevented him with blessings of sweetness: \Star{} Thou hast set on his head a crown of precious stones.\end{Resp}
\begin{Resp}\Vsign He asked life of thee: and thou hast given him length of days for ever and ever.\end{Resp}
\begin{Resp}\Rsign Thou hast set on his head a crown of precious stones.\end{Resp}
\begin{Resp}\Vsign Glory be to the Father, and to the Son, \Star{} and to the Holy Ghost.\end{Resp}
\begin{Resp}\Rsign Thou hast set on his head a crown of precious stones.\end{Resp}
\switchcolumn*

\LA
\LessonHead{Lectio IX}
\switchcolumn
\EN
\LessonHead{Lesson IX}
\switchcolumn*

\LA
\noindent Ad se autem nobis veniéntibus Dóminus præcépit, ut renuntiémus nostris: quia quicúmque ad fídei agónem vénimus, luctámen contra malígnos spíritus súmimus. Nihil autem malígni spíritus in hoc mundo próprium póssident: nudi ergo cum nudis luctári debémus. Nam si vestítus quisque cum nudo luctátur, cítius ad terram dejícitur, quia habet unde teneátur. Quid enim sunt terréna ómnia, nisi quædam córporis induménta? Qui ergo contra diábolum ad certámen próperat, vestiménta abjíciat, ne succúmbat.\par
\switchcolumn
\EN
\noindent Let when we come unto Him the Lord will have us deny to ourselves even ourselves, since as many of us as are entered into the battle of faith, are entered into a contention against evil spirits. But the evil spirits have nothing of their own in this world, and therefore must we wrestle with them, naked with naked. For if he that is clothed, wrestle with him that is naked, he faileth swiftly, because he hath whereon he that is naked taketh hold. And what are all things earthly but things wherewith the soul is clothed upon? whosoever therefore will wrestle with Satan, let him cast away his clothes, lest he be thereby endangered.\par
\switchcolumn*

\LA
\TeDeum{Te Deum}
\switchcolumn
\EN
\TeDeum{Te Deum}
\switchcolumn*

\end{paracol}

\par\needspace{10\baselineskip}\addvspace{1.2em}{\centering\small\textsc{Die 21 Decembris} · \textsc{21 December}\par}
\Office{S. Thomæ Apostoli}{St. Thomas the Apostle}{Duplex II classis}
\addcontentsline{toc}{subsection}{Die 21 Decembris · S. Thomæ Apostoli\enspace\textit{St. Thomas the Apostle}}
\begin{paracol}{2}
\LA
\RefLine{Lectiones I–III: ex Commúni Apostolorum (C1).}
\switchcolumn
\EN
\RefLine{Lessons I–III: from the Commune Apostolorum (C1).}
\switchcolumn*

\LA
\LessonHead{Lectio IV}
\switchcolumn
\EN
\LessonHead{Lesson IV}
\switchcolumn*

\LA
\noindent Thomas Apóstolus, qui et Dídymus, Galilǽus, post accéptum Spíritum Sanctum in multas províncias proféctus est ad prædicándum Christi Evangélium: Parthis, Medis, Persis, Hyrcánis et Bactris christiánæ fídei et vitæ præcépta trádidit. Postrémo ad Indos se cónferens, eos in christiána religióne erudívit. Qui ad extrémum, vitæ doctrinǽque sanctitáte et miraculórum magnitúdine, cum céteris ómnibus sui admiratiónem et Jesu Christi amórem commovísset; illíus gentis regem, idolórum cultórem, magis ad iram accéndit: cujus senténtia condemnátus, telísque confóssus, Calamínæ apostolátus honórem martýrii coróna decorávit.\par
\switchcolumn
\EN
\noindent The Apostle Thomas, called Didymus, or the Twin, was a Galilean. After the descent of the Holy Ghost, he went into many provinces to preach Christ's Gospel. He gave knowledge of the rules of Christian faith and life to the Parthians, Medes, Persians, Hyrcanians, and Bactrians. He went last to the East Indies. Here he provoked the anger of one of the idolatrous kings, because the holiness of his life and teaching, and the number of his miracles, drew many after him, and brought them to the love of Christ Jesus. He was therefore condemned, and slain with lances. He crowned the dignity of the Apostleship with the glory of martyrdom, on the Coromandel coast, not far from Madras.\par
\switchcolumn*

\LA
\begin{Resp}\Rsign Vidi conjúnctos viros, habéntes spléndidas vestes, et Ángelus Dómini locútus est ad me, dicens: \Star{} Isti sunt viri sancti facti amíci Dei.\end{Resp}
\begin{Resp}\Vsign Vidi Ángelum Dei fortem, volántem per médium cælum, voce magna clamántem et dicéntem.\end{Resp}
\begin{Resp}\Rsign Isti sunt viri sancti facti amíci Dei.\end{Resp}
\switchcolumn
\EN
\begin{Resp}\Rsign I saw men standing together, clad in shining raiment, and the Angel of the Lord spoke unto me, saying: \Star{} These men are holy, for they are the friends of God.\end{Resp}
\begin{Resp}\Vsign I saw a strong Angel of God fly into the midst of heaven, saying with a loud voice:\end{Resp}
\begin{Resp}\Rsign These men are holy, for they are the friends of God.\end{Resp}
\switchcolumn*

\LA
\LessonHead{Lectio V}
\switchcolumn
\EN
\LessonHead{Lesson V}
\switchcolumn*

\LA
\LessonTitle{Sermo sancti Gregórii Papæ}
\switchcolumn
\EN
\LessonTitle{Sermon of St. Gregory, Pope}
\switchcolumn*

\LA
\Source{Homilia 30 in Evangelia, post medium}
\switchcolumn
\EN
\Source{Sermon 30 on the Gospels}
\switchcolumn*

\LA
\noindent Scriptum est: Spíritus Dómini ornávit cælos. Ornaménta enim cælórum sunt virtútes prædicántium. Quæ vidélicet ornaménta Paulus enúmerat, dicens: Álii datur per Spíritum sermo sapiéntiæ, álii sermo sciéntiæ secúndum eúndem Spíritum, álteri fides in eódem Spíritu, álii grátia sanitátum in uno Spíritu, álii operátio virtútum, álii prophetía, álii discrétio spirítuum, álii génera linguárum, álii interpretátio sermónum. Hæc autem ómnia operátur unus atque idem Spíritus, dívidens síngulis prout vult.\par
\switchcolumn
\EN
\noindent It is written: By His Spirit the Lord hath adorned the heavens. (Job xxvi. 13.) Now the ornament of the heavens are the godly powers of preachers, and this ornament, what it is, Paul teacheth us thus To one is given by the Spirit the word of wisdom, to another the word of knowledge by the same Spirit; to another faith by the same Spirit; to another the gifts of healing by the same Spirit, to another the working of miracles, to another prophecy, to another discerning of spirits, to another diverse kinds of tongues, to another the interpretation of tongues. But all these worketh that one and the self-same Spirit, dividing to every man severally as He will. (1 Cor. xii. 8.)\par
\switchcolumn*

\LA
\begin{Resp}\Rsign Beáti estis, cum maledíxerint vobis hómines, et persecúti vos fúerint, et díxerint omne malum advérsum vos, mentiéntes, propter me: \Star{} Gaudéte et exsultáte, quóniam merces vestra copiósa est in cælis.\end{Resp}
\begin{Resp}\Vsign Cum vos óderint hómines, et cum separáverint vos, et exprobráverint, et ejécerint nomen vestrum tamquam malum propter Fílium hóminis.\end{Resp}
\begin{Resp}\Rsign Gaudéte et exsultáte, quóniam merces vestra copiósa est in cælis.\end{Resp}
\switchcolumn
\EN
\begin{Resp}\Rsign Blessed are ye when men shall revile you, and persecute you, and shall say all manner of evil against you falsely, for My sake; \Star{} Rejoice, and be exceeding glad, for great is your reward in heaven.\end{Resp}
\begin{Resp}\Vsign When men shall hate you, and when they shall separate you from their company, and shall reproach you, and cast out your name as evil, for the Son of Man's sake.\end{Resp}
\begin{Resp}\Rsign Rejoice, and be exceeding glad, for great is your reward in heaven.\end{Resp}
\switchcolumn*

\LA
\LessonHead{Lectio VI}
\switchcolumn
\EN
\LessonHead{Lesson VI}
\switchcolumn*

\LA
\noindent Quot ergo sunt bona prædicántium, tot sunt ornaménta cælórum. Hinc rursus scriptum est: Verbo Dómini cæli firmáti sunt. Verbum enim Dómini Fílius est Patris. Sed eósdem cælos, vidélicet sanctos Apóstolos, ut tota simul sancta Trínitas ostendátur operáta, repénte de Sancti Spíritus divinitáte adjúngitur: Et Spíritu oris ejus omnis virtus eórum. Cælórum ergo virtus de Spíritu sumpta est: quia mundi hujus potestátibus contraíre non præsúmerent, nisi eos Sancti Spíritus fortitúdo solidásset. Quales namque doctóres sanctæ Ecclésiæ ante advéntum hujus Spíritus fúerint, scimus: et post advéntum illíus, cujus fortitúdinis facti sint, conspícimus.\par
\switchcolumn
\EN
\noindent So much power then as have preachers, so much ornament have the heavens. Wherefore again it is written By the word of the Lord were the heavens made. \textasciitilde{}(Ps. xxxii. 6.) For the Word of the Lord is the Son of the Father. But, to the end that all the Holy Trinity may be made manifest as the Maker of the heavens, that is, of the Apostles, it is straightway added touching God the Holy Ghost: you and all the host of them by the Breath of His mouth. Therefore the might of the same heavens is the might of the Spirit, for they had not braved the powers of this world, unless the strength of the Holy Ghost had comforted them. For we know what manner of men the Teachers of the Holy Church were before the coming of this Spirit and since He came we see in Whose strength they are made strong.\par
\switchcolumn*

\LA
\begin{Resp}\Rsign Isti sunt triumphatóres et amíci Dei, qui contemnéntes jussa príncipum, meruérunt prǽmia ætérna: \Star{} Modo coronántur, et accípiunt palmam.\end{Resp}
\begin{Resp}\Vsign Isti sunt qui venérunt ex magna tribulatióne, et lavérunt stolas suas in sánguine Agni.\end{Resp}
\begin{Resp}\Rsign Modo coronántur, et accípiunt palmam.\end{Resp}
\begin{Resp}\Vsign Glória Patri, et Fílio, \Star{} et Spirítui Sancto.\end{Resp}
\begin{Resp}\Rsign Modo coronántur, et accípiunt palmam.\end{Resp}
\switchcolumn
\EN
\begin{Resp}\Rsign These are they which have conquered, and are become the friends of God, who recked not of the commandments of princes, and earned the everlasting reward. \Star{} And now have they crowns on their heads, and palms in their hands.\end{Resp}
\begin{Resp}\Vsign These are they which came out of great tribulation, and have washed their robes in the blood of the Lamb.\end{Resp}
\begin{Resp}\Rsign And now have they crowns on their heads, and palms in their hands.\end{Resp}
\begin{Resp}\Vsign Glory be to the Father, and to the Son, \Star{} and to the Holy Ghost.\end{Resp}
\begin{Resp}\Rsign And now have they crowns on their heads, and palms in their hands.\end{Resp}
\switchcolumn*

\LA
\LessonHead{Lectio VII}
\switchcolumn
\EN
\LessonHead{Lesson VII}
\switchcolumn*

\LA
\LessonTitle{Léctio sancti Evangélii secúndum Joánnem}
\switchcolumn
\EN
\LessonTitle{From the Holy Gospel according to John}
\switchcolumn*

\LA
\Cite{Joannes 20:24-29}
\switchcolumn
\EN
\Cite{John 20:24-29}
\switchcolumn*

\LA
\noindent In illo témpore: Thomas, unus ex duódecim, qui dícitur Dídymus, non erat cum eis, quando venit Jesus. Et réliqua.\par
\switchcolumn
\EN
\noindent In that time, Thomas, one of the twelve, who is called Didymus, was not with them when Jesus came. And so on.\par
\switchcolumn*

\LA
\LessonTitle{Homilía sancti Gregórii Papæ}
\switchcolumn
\EN
\LessonTitle{Homily by Pope St. Gregory (the Great.)}
\switchcolumn*

\LA
\Source{Homilia 26 in Evang., post med.}
\switchcolumn
\EN
\Source{26th on the Gospels}
\switchcolumn*

\LA
\noindent Quid, fratres caríssimi, quid inter hæc animadvértitis? Numquid casu gestum créditis, ut eléctus ille discípulus tunc deésset, post autem véniens audíret, áudiens dubitáret, dúbitans palpáret, palpans créderet? Non hoc casu, sed divína dispensatióne gestum est. Egit namque miro modo supérna cleméntia, ut discípulus ille dúbitans, dum in magístro suo vúlnera palpáret carnis, in nobis vúlnera sanáret infidelitátis. Plus enim nobis Thomæ infidélitas ad fidem, quam fides credéntium discipulórum prófuit; quia dum ille ad fidem palpándo redúcitur, nostra mens, omni dubitatióne postpósita, in fide solidátur.\par
\switchcolumn
\EN
\noindent Dearly beloved brethren, what is it in this passage which particularly claimeth our attention? Think ye that it was by accident that this chosen Apostle was not with them when Jesus came? or, when he came, heard? or, when he heard, doubted? or, when he doubted, felt? or when he had felt, believed? All these things were not accidental, but Providential. It was a wonderful provision of Divine mercy, that this incredulous disciple, by thrusting his fingers into the bodily Wounds of his Master, should apply a remedy to the spiritual wounds of unbelief in our souls. The doubts of Thomas have done us more good than the faith of all the disciples that believed. While he feeleth his way to faith, our minds are freed from doubt, and settled in faith.\par
\switchcolumn*

\LA
\begin{Resp}\Rsign Isti sunt qui vivéntes in carne, plantavérunt Ecclésiam sánguine suo: \Star{} Cálicem Dómini bibérunt, et amíci Dei facti sunt.\end{Resp}
\begin{Resp}\Vsign In omnem terram exívit sonus eórum, et in fines orbis terræ verba eórum.\end{Resp}
\begin{Resp}\Rsign Cálicem Dómini bibérunt, et amíci Dei facti sunt.\end{Resp}
\switchcolumn
\EN
\begin{Resp}\Rsign These are they who while yet they lived in the flesh, planted the Church in their own blood; \Star{} They drank of the Lord's cup, and became the friends of God.\end{Resp}
\begin{Resp}\Vsign Their sound is gone out through all the earth, and their words to the ends of the world.\end{Resp}
\begin{Resp}\Rsign They drank of the Lord's cup, and became the friends of God.\end{Resp}
\switchcolumn*

\LA
\LessonHead{Lectio VIII}
\switchcolumn
\EN
\LessonHead{Lesson VIII}
\switchcolumn*

\LA
\noindent Sic quippe discípulum Dóminus post resurrectiónem suam dubitáre permísit, nec tamen in dubitatióne deséruit; sicut ante nativitátem suam habére Maríam sponsum vóluit, qui tamen ad ejus núptias non pervénit. Nam ita factus est discípulus dúbitans et palpans, testis veræ resurrectiónis, sicut sponsus matris fúerat custos integérrimæ virginitátis. Palpávit autem, et exclamávit: Dóminus meus, et Deus meus. Dicit ei Jesus: Quia vidísti me, credidísti. Cum Paulus Apóstolus dicat, Est autem fides sperandárum substántia rerum, arguméntum non apparéntium; profécto liquet quia fides illárum rerum arguméntum est, quæ apparére non possunt. Quæ étenim appárent, jam fidem non habent, sed agnitiónem.\par
\switchcolumn
\EN
\noindent Even as the Lord before His birth willed that Mary should be espoused, and yet never lose her virginity, so, after His Resurrection, He willed that His disciple should doubt, and yet not lose his faith. For, even as the espoused husband was the keeper of the virginity of the Mother, so was the disciple who doubted and felt, the witness of the truth of the Resurrection. He felt, and cried out My Lord and my God. Jesus saith unto him Thomas, because thou hast seen Me, thou hast believed. When the Apostle Paul saith (Heb. xi. i): Faith is the substance of things hoped for, the evidence of things not seen, he plainly meaneth that faith is the evidence of things that cannot be seen. When they are seen, there remaineth not faith, but knowledge.\par
\switchcolumn*

\LA
\begin{Resp}\Rsign Isti sunt viri sancti, quos elégit Dóminus in caritáte non ficta, et dedit illis glóriam sempitérnam: \Star{} Quorum doctrína fulget Ecclésia, ut sole luna.\end{Resp}
\begin{Resp}\Vsign Sancti per fidem vicérunt regna: operáti sunt justítiam.\end{Resp}
\begin{Resp}\Rsign Quorum doctrína fulget Ecclésia, ut sole luna.\end{Resp}
\begin{Resp}\Vsign Glória Patri, et Fílio, \Star{} et Spirítui Sancto.\end{Resp}
\begin{Resp}\Rsign Quorum doctrína fulget Ecclésia, ut sole luna.\end{Resp}
\switchcolumn
\EN
\begin{Resp}\Rsign These men are saints, whom the Lord hath chosen in love unfeigned, and hath given them glory everlasting. These are they \Star{} By the light of whose teaching the Church is glorified, even as the moon is glorified by the light of the sun.\end{Resp}
\begin{Resp}\Vsign The saints through faith subdued kingdoms, wrought righteousness.\end{Resp}
\begin{Resp}\Rsign By the light of whose teaching the Church is glorified, even as the moon is glorified by the light of the sun.\end{Resp}
\begin{Resp}\Vsign Glory be to the Father, and to the Son, \Star{} and to the Holy Ghost.\end{Resp}
\begin{Resp}\Rsign By the light of whose teaching the Church is glorified, even as the moon is glorified by the light of the sun.\end{Resp}
\switchcolumn*

\LA
\LessonHead{Lectio IX}
\switchcolumn
\EN
\LessonHead{Lesson IX}
\switchcolumn*

\LA
\noindent Dum ergo vidit Thomas, dum palpávit, cur ei dícitur: Quia vidísti me, credidísti? Sed áliud vidit, áliud crédidit. A mortáli quippe hómine divínitas vidéri non pótuit. Hóminem ergo vidit, et Deum conféssus est, dicens: Dóminus meus, et Deus meus. Vidéndo ergo crédidit, qui considerándo verum hóminem, hunc Deum, quem vidére non póterat, exclamávit. Lætíficat valde quod séquitur: Beáti qui non vidérunt, et credidérunt. In qua nimírum senténtia nos speciáliter signáti sumus, qui eum, quem carne non vídimus, mente retinémus. Nos signáti sumus; sed si fidem nostram opéribus séquimur. Ille étenim vere credit, qui exércet operándo quod credit.\par
\switchcolumn
\EN
\noindent Thomas, then, seeth, and believeth. Why is it said to him Because thou hast seen Me, thou hast believed? The truth is, he saw one thing, and so believed another. To mortal man it is not given to see God. He therefore saw only the Manhood, and yet had faith in the Godhead: My Lord and my God. This he said, seeing and believing, seeing Perfect Man, and yet believing in Perfect God, Whom he could not see. O what a comfort are the words which follow! Blessed are they that have not seen, and yet have believed. These words are specially meant for us, who have not seen even the Flesh, and who yet do believe. They are specially meant for us if we believe and do not, by our lives, give the lie to our belief. He only hath a saving faith, whose faith beareth fruit.\par
\switchcolumn*

\LA
\TeDeum{Te Deum}
\switchcolumn
\EN
\TeDeum{Te Deum}
\switchcolumn*

\end{paracol}

\par\needspace{10\baselineskip}\addvspace{1.2em}{\centering\small\textsc{Die 24 Decembris} · \textsc{24 December}\par}
\Office{In Vigilia Nativitatis Domini}{In Vigilia Nativitatis Domini}{Duplex I classis}
\addcontentsline{toc}{subsection}{Die 24 Decembris · In Vigilia Nativitatis Domini\enspace\textit{In Vigilia Nativitatis Domini}}
\begin{paracol}{2}
\LA
\LessonHead{Lectio I}
\switchcolumn
\EN
\LessonHead{Lesson I}
\switchcolumn*

\LA
\LessonTitle{Léctio sancti Evangélii secúndum Matthǽum}
\switchcolumn
\EN
\LessonTitle{Continuation of the Holy Gospel according to Matthew}
\switchcolumn*

\LA
\Cite{Matt 1:18-21}
\switchcolumn
\EN
\Cite{Matt 1:18-21}
\switchcolumn*

\LA
\noindent Cum esset desponsáta Mater Jesu María Joseph, ántequam convenírent, invénta est in útero habens de Spíritu Sancto. Et réliqua.\par
\switchcolumn
\EN
\noindent When as his mother Mary was espoused to Joseph, before they came together, she was found with child, of the Holy Ghost. And so on.\par
\switchcolumn*

\LA
\LessonTitle{Homilía sancti Hierónymi Presbýteri}
\switchcolumn
\EN
\LessonTitle{Homily by St. Jerome, Priest (at Bethlehem.)}
\switchcolumn*

\LA
\Source{Lib. 1 Comment. in c. 1 Matth.}
\switchcolumn
\EN
\Source{Book of Commentaries, on Matth. i.}
\switchcolumn*

\LA
\noindent Quare non de símplici vírgine, sed de desponsáta concípitur? Primum, ut per generatiónem Joseph, orígo Maríæ monstrarétur: secúndo, ne lapidarétur a Judǽis ut adúltera: tértio, ut in Ægýptum fúgiens habéret solátium. Martyr Ignátius étiam quartam áddidit causam, cur a desponsáta concéptus sit: Ut partus, ínquiens, ejus celarétur diábolo, dum eum putat non de vírgine, sed de uxóre generátum.\par
\switchcolumn
\EN
\noindent Why was the Lord conceived of an espoused virgin rather than of a free? First, for the sake of the genealogy of Mary, which we have obtained by that of Joseph. Secondly, because she was thus saved from being stoned by the Jews as an adulteress. Thirdly, that Himself and His mother might have a guardian on their journey into Egypt. To these, Ignatius, the martyr of Antioch, has added a fourth reason namely, that the birth might take place unknown to the devil, who would naturally suppose that Mary had conceived by Joseph.\par
\switchcolumn*

\LA
\begin{Resp}\Rsign Sanctificámini hódie, et estóte paráti: quia die crástina vidébitis \Star{} Majestátem Dei in vobis.\end{Resp}
\begin{Resp}\Vsign Hódie sciétis quia véniet Dóminus, et mane vidébitis.\end{Resp}
\begin{Resp}\Rsign Majestátem Dei in vobis.\end{Resp}
\switchcolumn
\EN
\begin{Resp}\Rsign Sanctify yourselves today, and be ready; for on the morrow, ye shall see \Star{} The majesty of God upon you.\end{Resp}
\begin{Resp}\Vsign This day ye shall know that the Lord cometh, and in the morning, then ye shall see.\end{Resp}
\begin{Resp}\Rsign The majesty of God upon you.\end{Resp}
\switchcolumn*

\LA
\LessonHead{Lectio II}
\switchcolumn
\EN
\LessonHead{Lesson II}
\switchcolumn*

\LA
\noindent Antequam convenírent, invénta est in útero habens de Spíritu Sancto. Non ab álio invénta est, nisi a Joseph, qui pene licéntia maritáli futúræ uxóris ómnia nóverat. Quod autem dícitur, Antequam convenírent, non séquitur ut póstea convénerint: sed Scriptúra quod factum non sit, osténdit.\par
\switchcolumn
\EN
\noindent Before they came together, she was found with child of the Holy Ghost. She was found, that is, by Joseph, but by no one else. He had already almost an husband's privilege to know all that concerned her. Before they came together. This doth not imply that they ever did come together the Scripture merely showeth the absolute fact that up to this time they had not done so.\par
\switchcolumn*

\LA
\begin{Resp}\Rsign Constántes estóte, vidébitis auxílium Dómini super vos: Judǽa et Jerúsalem, nolíte timére: \Star{} Cras egrediémini, et Dóminus erit vobíscum.\end{Resp}
\begin{Resp}\Vsign Sanctificámini, fílii Israël, et estóte paráti.\end{Resp}
\begin{Resp}\Rsign Cras egrediémini, et Dóminus erit vobíscum.\end{Resp}
\switchcolumn
\EN
\begin{Resp}\Rsign Stand still, and ye shall see the help of the Lord with you; O Judah and Jerusalem, fear not. \Star{} Tomorrow ye shall go out, and the Lord will be with you.\end{Resp}
\begin{Resp}\Vsign Sanctify yourselves, O ye children of Israel, and be ready.\end{Resp}
\begin{Resp}\Rsign Tomorrow ye shall go out, and the Lord will be with you.\end{Resp}
\switchcolumn*

\LA
\LessonHead{Lectio III}
\switchcolumn
\EN
\LessonHead{Lesson III}
\switchcolumn*

\LA
\noindent Joseph autem vir ejus, cum esset justus, et nollet eam tradúcere, vóluit occúlte dimíttere eam. Si quis fornicáriæ conjúngitur, unum corpus effícitur, et in lege præcéptum est, non solum reos, sed et cónscios críminum obnóxios esse peccáti: quómodo Joseph, cum crimen celáret uxóris, justus scríbitur? Sed hoc testimónium Maríæ est, quod Joseph sciens illíus castitátem, et admírans quod evénerat, celat siléntio, cujus mystérium nesciébat.\par
\switchcolumn
\EN
\noindent Then Joseph her husband, being a just man, and not willing to make her a public example, was minded to put her away privily. If any man be joined to a fornicatress they become one body; and according to the law they that are privy to a crime are thereby guilty. How then can it be that Joseph is described as a just man, at the very time he was compounding the criminality of his espoused? It must have been that he knew her to be pure, and yet understood not the mystery of her pregnancy, but, while he wondered at that which had happened, was willing to hold his peace.\par
\switchcolumn*

\LA
\begin{Resp}\Rsign Sanctificámini, fílii Israël, dicit Dóminus: die enim crástina descéndet Dóminus, \Star{} Et áuferet a vobis omnem languórem.\end{Resp}
\begin{Resp}\Vsign Crástina die delébitur iníquitas terræ, et regnábit super nos Salvátor mundi.\end{Resp}
\begin{Resp}\Rsign Et áuferet a vobis omnem languórem.\end{Resp}
\begin{Resp}\Vsign Glória Patri, et Fílio, \Star{} et Spirítui Sancto.\end{Resp}
\begin{Resp}\Rsign Et áuferet a vobis omnem languórem.\end{Resp}
\switchcolumn
\EN
\begin{Resp}\Rsign Sanctify yourselves, O ye children of Israel, saith the Lord, for on the morrow the Lord will come down. \Star{} And will take away from you all sickness.\end{Resp}
\begin{Resp}\Vsign On the morrow the sins of the earth shall be washed away, and the Saviour of the world will be our King.\end{Resp}
\begin{Resp}\Rsign And will take away from you all sickness.\end{Resp}
\begin{Resp}\Vsign Glory be to the Father, and to the Son, \Star{} and to the Holy Ghost.\end{Resp}
\begin{Resp}\Rsign And will take away from you all sickness.\end{Resp}
\switchcolumn*

\end{paracol}

\par\needspace{10\baselineskip}\addvspace{1.2em}{\centering\small\textsc{Die 24 Decembris} · \textsc{24 December}\par}
\Office{In Vigilia Nativitatis Domini}{In Vigilia Nativitatis Domini}{Duplex I classis}
\addcontentsline{toc}{subsection}{Die 24 Decembris · In Vigilia Nativitatis Domini\enspace\textit{In Vigilia Nativitatis Domini}}
\begin{paracol}{2}
\LA
\LessonHead{Lectio I}
\switchcolumn
\EN
\LessonHead{Lesson I}
\switchcolumn*

\LA
\LessonTitle{De Isaía Prophéta}
\switchcolumn
\EN
\LessonTitle{Lesson from the book of Isaias}
\switchcolumn*

\LA
\Cite{Isa 35:1-7}
\switchcolumn
\EN
\Cite{Isa 35:1-7}
\switchcolumn*

\LA
\noindent Lætábitur desérta et ínvia, et exsultábit solitúdo, et florébit quasi lílium. Gérminans germinábit, et exsultábit lætabúnda et laudans: glória Líbani data est ei: decor Carméli et Saron, ipsi vidébunt glóriam Dómini, et decórem Dei nostri. Confortáte manus dissolútas, et génua debília roboráte. Dícite pusillánimis: Confortámini, et nolíte timére: ecce Deus vester ultiónem addúcet retributiónis: Deus ipse véniet, et salvábit vos. Tunc aperiéntur óculi cæcórum, et aures surdórum patébunt. Tunc sáliet sicut cervus claudus, et apérta erit lingua mutórum: quia scissæ sunt in desérto aquæ, et torréntes in solitúdine. Et quæ erat árida, erit in stagnum, et sítiens in fontes aquárum.\par
\switchcolumn
\EN
\noindent The land that was desolate and impassable shall be glad, and the wilderness shall rejoice, and shall flourish like the lily. It shall bud forth and blossom, and shall rejoice with joy and praise: the glory of Libanus is given to it: the beauty of Carmel, and Saron, they shall see the glory of the Lord, and the beauty of our God. Strengthen ye the feeble hands, and confirm the weak knees. Say to the fainthearted: Take courage, and fear not: behold your God will bring the revenge of recompense: God himself will come and will save you. Then shall the eyes of the blind be opened, and the ears of the deaf shall be unstopped. Then shall the lame man leap as a hart, and the tongue of the dumb shall be free: for waters are broken out in the desert, and streams in the wilderness. And that which was dry land, shall become a pool, and the thirsty land springs of water.\par
\switchcolumn*

\LA
\begin{Resp}\Rsign Cánite tuba in Sion, vocáte gentes, annuntiáte pópulis, et dícite: \Star{} Ecce Deus Salvátor noster advéniet.\end{Resp}
\begin{Resp}\Vsign Annuntiáte, et audítum fácite: loquímini, et clamáte.\end{Resp}
\begin{Resp}\Rsign Ecce Deus Salvátor noster advéniet.\end{Resp}
\switchcolumn
\EN
\begin{Resp}\Rsign Blow ye the trumpet in Zion, call together the nations, tell it out among the people, and say: \Star{} Behold, God our Saviour cometh.\end{Resp}
\begin{Resp}\Vsign Tell it out and make it to be heard; speak aloud and cry.\end{Resp}
\begin{Resp}\Rsign Behold, God our Saviour cometh.\end{Resp}
\switchcolumn*

\LA
\LessonHead{Lectio II}
\switchcolumn
\EN
\LessonHead{Lesson II}
\switchcolumn*

\LA
\Cite{Isa 35:7-10}
\switchcolumn
\EN
\Cite{Isa 35:7-10}
\switchcolumn*

\LA
\noindent In cubílibus, in quibus prius dracónes habitábant, oriétur viror cálami et junci. Et erit ibi sémita et via, et via sancta vocábitur: non transíbit per eam pollútus, et hæc erit vobis dirécta via, ita ut stulti non errent per eam. Non erit ibi leo, et mala béstia non ascéndet per eam, nec inveniétur ibi: et ambulábunt, qui liberáti fúerint. Et redémpti a Dómino converténtur, et vénient in Sion cum laude: et lætítia sempitérna super caput eórum: gáudium et lætítiam obtinébunt, et fúgiet dolor et gémitus.\par
\switchcolumn
\EN
\noindent In the dens where dragons dwell before, shall rise up the verdure of the reed and the bulrush. And a path and a way shall be there, and it shall be called the holy way: the unclean shall not pass over it, and this shall be unto you a straight way, so that fools shall not err therein. No lion shall be there, nor shall any mischievous beast go up by it, nor be found there: but they shall walk there that shall be delivered. And the redeemed of the Lord shall return, and shall come into Sion with praise, and everlasting joy shall be upon their heads: they shall obtain joy and gladness, and sorrow and mourning shall flee away.\par
\switchcolumn*

\LA
\begin{Resp}\Rsign Non auferétur sceptrum de Juda, et dux de fémore ejus, donec véniat qui mitténdus est: \Star{} Et ipse erit exspectátio géntium.\end{Resp}
\begin{Resp}\Vsign Pulchrióres sunt óculi ejus vino, et dentes ejus lacte candidióres.\end{Resp}
\begin{Resp}\Rsign Et ipse erit exspectátio géntium.\end{Resp}
\switchcolumn
\EN
\begin{Resp}\Rsign The sceptre shall not depart from Judah, nor the law-giver from his loins, until he that shall be sent cometh. \Star{} And unto him shall the longing of the Gentiles be.\end{Resp}
\begin{Resp}\Vsign His eyes shall be bright with wine, and his teeth white with milk.\end{Resp}
\begin{Resp}\Rsign And unto him shall the desire of the Gentiles be.\end{Resp}
\switchcolumn*

\LA
\LessonHead{Lectio III}
\switchcolumn
\EN
\LessonHead{Lesson III}
\switchcolumn*

\LA
\Cite{Isa 41:1-4}
\switchcolumn
\EN
\Cite{Isa 41:1-4}
\switchcolumn*

\LA
\noindent Táceant ad me ínsulæ, et gentes mutent fortitúdinem: accédant, et tunc loquántur, simul ad judícium propinquémus. Quis suscitávit ab Oriénte justum, vocávit eum ut sequerétur se? dabit in conspéctu ejus gentes et reges obtinébit: dabit quasi púlverem gládio ejus, sicut stípulam vento raptam árcui ejus. Persequétur eos, transíbit in pace, sémita in pédibus ejus non apparébit. Quis hæc operátus est, et fecit, vocans generatiónes ab exórdio? Ego Dóminus, primus et novíssimus ego sum.\par
\switchcolumn
\EN
\noindent Let the islands keep silence before me, and the nations take new strength: let them come near, and then speak, let us come near to judgment together. Who hath raised up the just one from the east, hath called him to follow him? he shall give the nations in his sight, and he shall rule over kings: he shall give them as the dust to his sword, as stubble driven by the wind, to his bow. He shall pursue them, he shall pass in peace, no path shall appear after his feet. Who hath wrought and done these things, calling the generations from the beginning? I the Lord, I am the first and the last.\par
\switchcolumn*

\LA
\begin{Resp}\Rsign Me opórtet mínui, illum autem créscere: qui autem post me venit, ante me factus est: \Star{} Cujus non sum dignus corrígiam calceamentórum sólvere.\end{Resp}
\begin{Resp}\Vsign Ego baptizávi vos aqua: ille autem baptizábit vos Spíritu Sancto.\end{Resp}
\begin{Resp}\Rsign Cujus non sum dignus corrígiam calceamentórum sólvere.\end{Resp}
\begin{Resp}\Vsign Glória Patri, et Fílio, \Star{} et Spirítui Sancto.\end{Resp}
\begin{Resp}\Rsign Cujus non sum dignus corrígiam calceamentórum sólvere.\end{Resp}
\switchcolumn
\EN
\begin{Resp}\Rsign I must decrease, but He must increase; He it is Who, coming after me is preferred before me; \Star{} Whose shoe's latchet I am not worthy to unloose.\end{Resp}
\begin{Resp}\Vsign I baptize you with water; but He shall baptize you with the Holy Ghost.\end{Resp}
\begin{Resp}\Rsign Whose shoe's latchet I am not worthy to unloose.\end{Resp}
\begin{Resp}\Vsign Glory be to the Father, and to the Son, \Star{} and to the Holy Ghost.\end{Resp}
\begin{Resp}\Rsign Whose shoe's latchet I am not worthy to unloose.\end{Resp}
\switchcolumn*

\LA
\LessonHead{Lectio IV}
\switchcolumn
\EN
\LessonHead{Lesson IV}
\switchcolumn*

\LA
\LessonTitle{Sermo sancti Leónis Papæ}
\switchcolumn
\EN
\LessonTitle{From the Sermons of Pope St. Leo (the Great)}
\switchcolumn*

\LA
\Source{Sermo 1 de jejúnio décimi mensis, et collectis}
\switchcolumn
\EN
\Source{1st on the December Fast, and Almsgiving}
\switchcolumn*

\LA
\noindent Si fidéliter, dilectíssimi, atque sapiénter creatiónis nostræ intelligámus exórdium, inveniémus hóminem ideo ad imáginem Dei cónditum, ut imitátor sui esset auctóris: et hanc esse naturálem nostri géneris dignitátem, si in nobis, quasi in quodam spéculo, divínæ benignitátis forma respléndeat. Ad quam quotídie nos útique réparat grátia Salvatóris, dum quod cécidit in Adam primo, erígitur in secúndo.\par
\switchcolumn
\EN
\noindent Dearly beloved brethren, if we study attentively the history of the creation of our race, we shall find that man was made in the image of God, that his ways also might be an imitation of the ways of his Maker. This is the natural, real, and highest dignity to which we are capable of attaining, that the goodness of the Divine nature should have a reflection in us, as in a glass. As a mean of reaching this dignity, we are daily offered the grace of our Saviour, for as in the first Adam all men are fallen, so in the Second Adam can all men be raised up again (i Cor. xv. 22).\par
\switchcolumn*

\LA
\begin{Resp}\Rsign Nascétur nobis párvulus, et vocábitur Deus, Fortis: \Star{} Ipse sedébit super thronum David patris sui, et imperábit: cujus potéstas super húmerum ejus.\end{Resp}
\begin{Resp}\Vsign In ipso benedicéntur omnes tribus terræ, omnes gentes sérvient ei.\end{Resp}
\begin{Resp}\Rsign Ipse sedébit super thronum David patris sui, et imperábit: cujus potéstas super húmerum ejus.\end{Resp}
\switchcolumn
\EN
\begin{Resp}\Rsign Unto us shall a Child be born, and His name shall be called the Mighty God. \Star{} He shall sit upon the throne of His father David, and shall reign, and the government shall be upon His shoulder.\end{Resp}
\begin{Resp}\Vsign In Him shall all the kindreds of the earth be blessed; all nations shall serve Him.\end{Resp}
\begin{Resp}\Rsign He shall sit upon the throne of His father David, and shall reign, and the government shall be upon His shoulder.\end{Resp}
\switchcolumn*

\LA
\LessonHead{Lectio V}
\switchcolumn
\EN
\LessonHead{Lesson V}
\switchcolumn*

\LA
\noindent Causa autem reparatiónis nostræ non est nisi misericórdia Dei: quem non diligerémus, nisi prius nos ipse dilígeret, et ténebras ignorántiæ nostræ, suæ veritátis luce discúteret. Quod per sanctum Isaíam Dóminus denúntians, ait: Addúcam cæcos in viam quam ignorábant, et sémitas quas nesciébant, fáciam illos calcáre: fáciam illis ténebras in lucem, et prava in dirécta. Hæc verba fáciam illis, et non relínquam eos. Et íterum: Invéntus sum, inquit, a non quæréntibus me, et palam appárui iis qui me non interrogábant.\par
\switchcolumn
\EN
\noindent Our restoration from the consequences of Adam's fall is sheer mercy of God, and nothing else; we should not have loved Him unless He had first loved us, (1 John iv. 19,) and scattered the darkness of our ignorance by the light of His truth. This the Lord promised by the mouth of Isaiah, where He saith: (Isa. xlii. 16,) I will bring the blind by a way that they knew not, and I will lead them in paths that they have not known I will make darkness light before them, and crooked things straight. These things will I do unto them and not forsake them. And again: (Isa. lxv. 1, 2; Rom. x. 20,) I was found of them that sought Me not; I was made manifest unto them that asked not after Me.\par
\switchcolumn*

\LA
\begin{Resp}\Rsign Ecce jam venit plenitúdo témporis, in quo misit Deus Fílium suum in terras, natum de Vírgine, factum sub lege: \Star{} Ut eos, qui sub lege erant, redímeret.\end{Resp}
\begin{Resp}\Vsign Propter nímiam caritátem suam, qua diléxit nos Deus, Fílium suum misit in similitúdinem carnis peccáti.\end{Resp}
\begin{Resp}\Rsign Ut eos, qui sub lege erant, redímeret.\end{Resp}
\switchcolumn
\EN
\begin{Resp}\Rsign Behold, the fullness of the time is come, wherein God hath sent forth His Son into the world, born of a Virgin, made under the law; \Star{} To redeem them that were under the law.\end{Resp}
\begin{Resp}\Vsign God, for His great love wherewith He loved us, hath sent His own Son in the likeness of sinful flesh.\end{Resp}
\begin{Resp}\Rsign To redeem them that were under the law.\end{Resp}
\switchcolumn*

\LA
\LessonHead{Lectio VI}
\switchcolumn
\EN
\LessonHead{Lesson VI}
\switchcolumn*

\LA
\noindent Quod quómodo implétum sit, Joánnes Apóstolus docet, dicens: Scimus quóniam Fílius Dei venit, et dedit nobis sensum, ut cognoscámus verum, et simus in vero Fílio ejus. Et íterum: Nos ergo diligámus Deum, quóniam ipse prior diléxit nos. Diligéndo ítaque nos Deus, ad imáginem suam nos réparat: et ut in nobis formam suæ bonitátis invéniat, dat unde ipsi quoque quod operátur operémur, accéndens scílicet méntium nostrárum lucérnas, et igne nos suæ caritátis inflámmans, ut non solum ipsum, sed étiam quidquid díligit, diligámus.\par
\switchcolumn
\EN
\noindent And we know from the Apostle John how God fulfilled His promise, (1 John v. 20.) We know that the Son of God is come, and hath given us an understanding, that we may know Him That is True, and be in Him That is True, even in His Son. And again: (iv. 19,) Let us therefore love God, because He first loved us. For His great love then wherewith he hath loved us, (Eph. ii. 4,) God reneweth His likeness in us. And, moreover, in order that He may find in us the reflection of His goodness, He giveth us that whereby to work along with Himself, (Who worketh all in all,) lighting, as it were, candles in our dark minds, and kindling in us the fire of His love, to make us love not Himself only, but likewise, in Him, whatsoever He loveth.\par
\switchcolumn*

\LA
\begin{Resp}\Rsign Virgo Israël, revértere ad civitátes tuas: \Star{} Usquequo dolens avertéris? generábis Dóminum Salvatórem, oblatiónem novam in terra: * Ambulábunt hómines in salvatiónem.\end{Resp}
\begin{Resp}\Vsign In caritáte perpétua diléxi te: ídeo attráxi te míserans tui.\end{Resp}
\begin{Resp}\Rsign Usquequo dolens avertéris? generábis Dóminum Salvatórem, oblatiónem novam in terra.\end{Resp}
\begin{Resp}\Vsign Glória Patri, et Fílio, \Star{} et Spirítui Sancto.\end{Resp}
\begin{Resp}\Rsign Ambulábunt hómines in salvatiónem.\end{Resp}
\switchcolumn
\EN
\begin{Resp}\Rsign O virgin of Israel, turn again to thy cities. \Star{} How long wilt thou go about sorrowing? Thou shalt bring forth the Lord thy Saviour, a new offering in the earth; men shall walk in paths of salvation.\end{Resp}
\begin{Resp}\Vsign I have loved thee with an everlasting love; therefore with loving-kindness have I drawn thee.\end{Resp}
\begin{Resp}\Rsign How long wilt thou go about sorrowing? Thou shalt bring forth the Lord thy Saviour, a new offering in the earth.\end{Resp}
\begin{Resp}\Vsign Glory be to the Father, and to the Son, \Star{} and to the Holy Ghost.\end{Resp}
\begin{Resp}\Rsign Men shall walk in paths of salvation.\end{Resp}
\switchcolumn*

\LA
\LessonHead{Lectio VII}
\switchcolumn
\EN
\LessonHead{Lesson VII}
\switchcolumn*

\LA
\LessonTitle{Léctio sancti Evangélii secúndum Matthǽum}
\switchcolumn
\EN
\LessonTitle{Continuation of the Holy Gospel according to Matthew}
\switchcolumn*

\LA
\Cite{Matt 1:18-21}
\switchcolumn
\EN
\Cite{Matt 1:18-21}
\switchcolumn*

\LA
\noindent Cum esset desponsáta Mater Jesu María Joseph, ántequam convenírent, invénta est in útero habens de Spíritu Sancto. Et réliqua.\par
\switchcolumn
\EN
\noindent When as his mother Mary was espoused to Joseph, before they came together, she was found with child, of the Holy Ghost. And so on.\par
\switchcolumn*

\LA
\LessonTitle{Homilía sancti Hierónymi Presbýteri}
\switchcolumn
\EN
\LessonTitle{Homily by St. Jerome, Priest (at Bethlehem.)}
\switchcolumn*

\LA
\Source{Lib. 1 Comment. in c. 1 Matth.}
\switchcolumn
\EN
\Source{Book of Commentaries, on Matth. i.}
\switchcolumn*

\LA
\noindent Quare non de símplici vírgine, sed de desponsáta concípitur? Primum, ut per generatiónem Joseph, orígo Maríæ monstrarétur: secúndo, ne lapidarétur a Judǽis ut adúltera: tértio, ut in Ægýptum fúgiens habéret solátium. Martyr Ignátius étiam quartam áddidit causam, cur a desponsáta concéptus sit: Ut partus, ínquiens, ejus celarétur diábolo, dum eum putat non de vírgine, sed de uxóre generátum.\par
\switchcolumn
\EN
\noindent Why was the Lord conceived of an espoused virgin rather than of a free? First, for the sake of the genealogy of Mary, which we have obtained by that of Joseph. Secondly, because she was thus saved from being stoned by the Jews as an adulteress. Thirdly, that Himself and His mother might have a guardian on their journey into Egypt. To these, Ignatius, the martyr of Antioch, has added a fourth reason namely, that the birth might take place unknown to the devil, who would naturally suppose that Mary had conceived by Joseph.\par
\switchcolumn*

\LA
\begin{Resp}\Rsign Sanctificámini hódie, et estóte paráti: quia die crástina vidébitis \Star{} Majestátem Dei in vobis.\end{Resp}
\begin{Resp}\Vsign Hódie sciétis quia véniet Dóminus, et mane vidébitis.\end{Resp}
\begin{Resp}\Rsign Majestátem Dei in vobis.\end{Resp}
\switchcolumn
\EN
\begin{Resp}\Rsign Sanctify yourselves today, and be ready; for on the morrow, ye shall see \Star{} The majesty of God upon you.\end{Resp}
\begin{Resp}\Vsign This day ye shall know that the Lord cometh, and in the morning, then ye shall see.\end{Resp}
\begin{Resp}\Rsign The majesty of God upon you.\end{Resp}
\switchcolumn*

\LA
\LessonHead{Lectio VIII}
\switchcolumn
\EN
\LessonHead{Lesson VIII}
\switchcolumn*

\LA
\noindent Antequam convenírent, invénta est in útero habens de Spíritu Sancto. Non ab álio invénta est, nisi a Joseph, qui pene licéntia maritáli futúræ uxóris ómnia nóverat. Quod autem dícitur, Antequam convenírent, non séquitur ut póstea convénerint: sed Scriptúra quod factum non sit, osténdit.\par
\switchcolumn
\EN
\noindent Before they came together, she was found with child of the Holy Ghost. She was found, that is, by Joseph, but by no one else. He had already almost an husband's privilege to know all that concerned her. Before they came together. This doth not imply that they ever did come together the Scripture merely showeth the absolute fact that up to this time they had not done so.\par
\switchcolumn*

\LA
\begin{Resp}\Rsign Constántes estóte, vidébitis auxílium Dómini super vos: Judǽa et Jerúsalem, nolíte timére: \Star{} Cras egrediémini, et Dóminus erit vobíscum.\end{Resp}
\begin{Resp}\Vsign Sanctificámini, fílii Israël, et estóte paráti.\end{Resp}
\begin{Resp}\Rsign Cras egrediémini, et Dóminus erit vobíscum.\end{Resp}
\switchcolumn
\EN
\begin{Resp}\Rsign Stand still, and ye shall see the help of the Lord with you; O Judah and Jerusalem, fear not. \Star{} Tomorrow ye shall go out, and the Lord will be with you.\end{Resp}
\begin{Resp}\Vsign Sanctify yourselves, O ye children of Israel, and be ready.\end{Resp}
\begin{Resp}\Rsign Tomorrow ye shall go out, and the Lord will be with you.\end{Resp}
\switchcolumn*

\LA
\LessonHead{Lectio IX}
\switchcolumn
\EN
\LessonHead{Lesson IX}
\switchcolumn*

\LA
\noindent Joseph autem vir ejus, cum esset justus, et nollet eam tradúcere, vóluit occúlte dimíttere eam. Si quis fornicáriæ conjúngitur, unum corpus effícitur, et in lege præcéptum est, non solum reos, sed et cónscios críminum obnóxios esse peccáti: quómodo Joseph, cum crimen celáret uxóris, justus scríbitur? Sed hoc testimónium Maríæ est, quod Joseph sciens illíus castitátem, et admírans quod evénerat, celat siléntio, cujus mystérium nesciébat.\par
\switchcolumn
\EN
\noindent Then Joseph her husband, being a just man, and not willing to make her a public example, was minded to put her away privily. If any man be joined to a fornicatress they become one body; and according to the law they that are privy to a crime are thereby guilty. How then can it be that Joseph is described as a just man, at the very time he was compounding the criminality of his espoused? It must have been that he knew her to be pure, and yet understood not the mystery of her pregnancy, but, while he wondered at that which had happened, was willing to hold his peace.\par
\switchcolumn*

\LA
\begin{Resp}\Rsign Sanctificámini, fílii Israël, dicit Dóminus: die enim crástina descéndet Dóminus, \Star{} Et áuferet a vobis omnem languórem.\end{Resp}
\begin{Resp}\Vsign Crástina die delébitur iníquitas terræ, et regnábit super nos Salvátor mundi.\end{Resp}
\begin{Resp}\Rsign Et áuferet a vobis omnem languórem.\end{Resp}
\begin{Resp}\Vsign Glória Patri, et Fílio, \Star{} et Spirítui Sancto.\end{Resp}
\begin{Resp}\Rsign Et áuferet a vobis omnem languórem.\end{Resp}
\switchcolumn
\EN
\begin{Resp}\Rsign Sanctify yourselves, O ye children of Israel, saith the Lord, for on the morrow the Lord will come down. \Star{} And will take away from you all sickness.\end{Resp}
\begin{Resp}\Vsign On the morrow the sins of the earth shall be washed away, and the Saviour of the world will be our King.\end{Resp}
\begin{Resp}\Rsign And will take away from you all sickness.\end{Resp}
\begin{Resp}\Vsign Glory be to the Father, and to the Son, \Star{} and to the Holy Ghost.\end{Resp}
\begin{Resp}\Rsign And will take away from you all sickness.\end{Resp}
\switchcolumn*

\end{paracol}

\par\needspace{10\baselineskip}\addvspace{1.2em}{\centering\small\textsc{Die 25 Decembris} · \textsc{25 December}\par}
\Office{In Nativitate Domini}{In Nativitate Domini}{Duplex I classis}
\addcontentsline{toc}{subsection}{Die 25 Decembris · In Nativitate Domini\enspace\textit{In Nativitate Domini}}
\begin{paracol}{2}
\LA
\LessonHead{Lectio I}
\switchcolumn
\EN
\LessonHead{Lesson I}
\switchcolumn*

\LA
\Cite{Isa 9:1-6}
\switchcolumn
\EN
\Cite{Isa 9:1-6}
\switchcolumn*

\LA
\noindent Primo témpore alleviáta est terra Zábulon, et terra Néphthali: et novíssimo aggraváta est via maris trans Jordánem Galilǽæ géntium. Pópulus qui ambulábat in ténebris, vidit lucem magnam: habitántibus in regióne umbræ mortis, lux orta est eis. Multiplicásti gentem, et non magnificásti lætítiam. Lætabúntur coram te, sicut qui lætántur in messe, sicut exsúltant victóres, capta præda, quando dívidunt spólia. Jugum enim óneris ejus, et virgam húmeri ejus, et sceptrum exactóris ejus superásti sicut in die Mádian. Quia omnis violénta prædátio cum tumúltu, et vestiméntum mistum sánguine, erit in combustiónem, et cibus ignis. Párvulus enim natus est nobis, et fílius datus est nobis, et factus est principátus super húmerum ejus: et vocábitur nomen ejus, Admirábilis, Consiliárius, Deus, Fortis, Pater futúri sǽculi, Princeps pacis.\par
\switchcolumn
\EN
\noindent At the first time the land of Zabulon, and the land of Nephtali was lightly touched: and at the last the way of the sea beyond the Jordan of the Galilee of the Gentiles was heavily loaded. The people that walked in darkness, have seen a great light: to them that dwelt in the region of the shadow of death, light is risen. Thou hast multiplied the nation, and hast not increased the joy. They shall rejoice before thee, as they that rejoice in the harvest, as conquerors rejoice after taking a prey, when they divide the spoils. For the yoke of their burden, and the rod of their shoulder, and the sceptre of their oppressor thou hast overcome, as in the day of Median. For every violent taking of spoils, with tumult, and garment mingled with blood, shall be burnt, and be fuel for the fire. For a child is born to us, and a son is given to us, and the government is upon his shoulder: and his name shall be called, Wonderful, Counsellor, God the Mighty, the Father of the world to come, the Prince of Peace.\par
\switchcolumn*

\LA
\begin{Resp}\Rsign Hódie nobis cælórum Rex de Vírgine nasci dignátus est, ut hóminem pérditum ad cæléstia regna revocáret: \Star{} Gaudet exércitus Angelórum: quia salus ætérna humáno géneri appáruit.\end{Resp}
\begin{Resp}\Vsign Glória in excélsis Deo, et in terra pax homínibus bonæ voluntátis.\end{Resp}
\begin{Resp}\Rsign Gaudet exércitus Angelórum: quia salus ætérna humáno géneri appáruit.\end{Resp}
\begin{Resp}\Vsign Glória Patri, et Fílio, et Spirítui Sancto.\end{Resp}
\begin{Resp}\Rsign Hódie nobis cælórum Rex de Vírgine nasci dignátus est, ut hóminem pérditum ad cæléstia regna revocáret: * Gaudet exércitus Angelórum: quia salus ætérna humáno géneri appáruit.\end{Resp}
\switchcolumn
\EN
\begin{Resp}\Rsign This is the day whereon the King of heaven was pleased to be born of a Virgin, that He might bring back to heaven man who was lost. \Star{} There is joy among the hosts of Angels, because eternal salvation hath appeared unto men.\end{Resp}
\begin{Resp}\Vsign Glory to God in the highest, and on earth peace, to men of goodwill.\end{Resp}
\begin{Resp}\Rsign There is joy among the hosts of Angels, because eternal salvation hath appeared unto men.\end{Resp}
\begin{Resp}\Vsign Glory be to the Father, and to the Son, and to the Holy Ghost.\end{Resp}
\begin{Resp}\Rsign This is the day whereon the King of heaven was pleased to be born of a Virgin, that He might bring back to heaven man who was lost. There is joy among the hosts of Angels, because eternal salvation hath appeared unto men.\end{Resp}
\switchcolumn*

\LA
\LessonHead{Lectio II}
\switchcolumn
\EN
\LessonHead{Lesson II}
\switchcolumn*

\LA
\Cite{Isa 40:1-8}
\switchcolumn
\EN
\Cite{Isa 40:1-8}
\switchcolumn*

\LA
\noindent Consolámini, consolámini, pópule meus, dicit Deus vester. Loquímini ad cor Jerúsalem, et advocáte eam: quóniam compléta est malítia ejus, dimíssa est iníquitas illíus: suscépit de manu Dómini duplícia pro ómnibus peccátis suis. Vox clamántis in desérto: Paráte viam Dómini, rectas fácite in solitúdine sémitas Dei nostri. Omnis vallis exaltábitur, et omnis mons et collis humiliábitur: et erunt prava in dirécta, et áspera in vias planas. Et revelábitur glória Dómini: et vidébit omnis caro páriter quod os Dómini locútum est. Vox dicéntis: Clama. Et dixi: Quid clamábo? Omnis caro fœnum, et omnis glória ejus quasi flos agri. Exsiccátum est fœnum, et cécidit flos: quia spíritus Dómini sufflávit in eo. Vere fœnum est pópulus: exsiccátum est fœnum, et cécidit flos: Verbum autem Dómini nostri manet in ætérnum.\par
\switchcolumn
\EN
\noindent Be comforted, be comforted, my people, saith your God. Speak ye to the heart of Jerusalem, and call to her: for her evil is come to an end, her iniquity is forgiven: she hath received of the hand of the Lord double for all her sins. The voice of one crying in the desert: Prepare ye the way of the Lord, make straight in the wilderness the paths of our God. Every valley shall be exalted, and every mountain and hill shall be made low, and the crooked shall become straight, and the rough ways plain. And the glory of the Lord shall be revealed, and all flesh together shall see, that the mouth of the Lord hath spoken. The voice of one, saying: Cry. And I said: What shall I cry? All flesh is grass, and all the glory thereof as the flower of the field. The grass is withered, and the flower is fallen, because the spirit of the Lord hath blown upon it. Indeed the people is grass: The grass is withered, and the flower is fallen: but the word of our Lord endureth for ever.\par
\switchcolumn*

\LA
\begin{Resp}\Rsign Hódie nobis de cælo pax vera descéndit: \Star{} Hódie per totum mundum mellíflui facti sunt cæli.\end{Resp}
\begin{Resp}\Vsign Hódie illúxit nobis dies redemptiónis novæ, reparatiónis antíquæ, felicitátis ætérnæ.\end{Resp}
\begin{Resp}\Rsign Hódie per totum mundum mellíflui facti sunt cæli.\end{Resp}
\switchcolumn
\EN
\begin{Resp}\Rsign This day is the true peace come down unto us from heaven. \Star{} This day throughout the whole world the skies drop down sweetness.\end{Resp}
\begin{Resp}\Vsign This day is the daybreak of our new redemption, of the restoring of the old, of everlasting joy.\end{Resp}
\begin{Resp}\Rsign This day throughout the whole world the skies drop down sweetness.\end{Resp}
\switchcolumn*

\LA
\LessonHead{Lectio III}
\switchcolumn
\EN
\LessonHead{Lesson III}
\switchcolumn*

\LA
\Cite{Isa 52:1-6}
\switchcolumn
\EN
\Cite{Isa 52:1-6}
\switchcolumn*

\LA
\noindent Consúrge, consúrge, indúere fortitúdine tua, Sion, indúere vestiméntis glóriæ tuæ, Jerúsalem, cívitas sancti: quia non adíciet ultra ut pertránseat per te incircumcísus, et immúndus. Excútere de púlvere, consúrge, sede, Jerúsalem: solve víncula colli tui, captíva fília Sion. Quia hæc dicit Dóminus: Gratis venumdáti estis, et sine argénto redimémini. Quia hæc dicit Dóminus Deus: In Ægýptum descéndit pópulus meus in princípio, ut colónus esset ibi: et Assur absque ulla causa calumniátus est eum. Et nunc quid mihi est hic, dicit Dóminus, quóniam ablátus est pópulus meus gratis? Dominatóres ejus iníque agunt, dicit Dóminus: et júgiter tota die nomen meum blasphemátur. Propter hoc sciet pópulus meus nomen meum, in die illa: quia ego ipse qui loquébar, ecce adsum.\par
\switchcolumn
\EN
\noindent Arise, arise, put on thy strength, O Sion, put on the garments of thy glory, O Jerusalem, the city of the Holy One: for henceforth the uncircumcised, and unclean shall no more pass through thee. Shake thyself from the dust, arise, sit up, O Jerusalem: loose the bonds from off thy neck, O captive daughter of Sion. For thus saith the Lord: You were sold gratis, and you shall be redeemed without money. For thus saith the Lord God: My people went down into Egypt at the beginning to sojourn there: and the Assyrian hath oppressed them without any cause at all. And now what have I here, saith the Lord: for my people is taken away gratis. They that rule over them treat them unjustly, saith the Lord, and my name is continually blasphemed all the day long. Therefore my people shall know my name in that day: for I myself that spoke, behold I am here.\par
\switchcolumn*

\LA
\begin{Resp}\Rsign Quem vidístis, pastóres? dícite, annuntiáte nobis, in terris quis appáruit? \Star{} Natum vídimus, et choros Angelórum collaudántes Dóminum.\end{Resp}
\begin{Resp}\Vsign Dícite, quidnam vidístis? et annuntiáte Christi nativitátem.\end{Resp}
\begin{Resp}\Rsign Natum vídimus, et choros Angelórum collaudántes Dóminum.\end{Resp}
\begin{Resp}\Vsign Glória Patri, et Fílio, \Star{} et Spirítui Sancto.\end{Resp}
\begin{Resp}\Rsign Natum vídimus, et choros Angelórum collaudántes Dóminum.\end{Resp}
\switchcolumn
\EN
\begin{Resp}\Rsign O ye shepherds, speak, and tell us what ye have seen; who is appeared in the earth? \Star{} We saw the new-born Child, and Angels singing praise to the Lord.\end{Resp}
\begin{Resp}\Vsign Speak; what have ye seen? And tell us of the Birth of Christ.\end{Resp}
\begin{Resp}\Rsign We saw the new-born Child, and Angels singing praise to the Lord.\end{Resp}
\begin{Resp}\Vsign Glory be to the Father, and to the Son, \Star{} and to the Holy Ghost.\end{Resp}
\begin{Resp}\Rsign We saw the new-born Child, and Angels singing praise to the Lord.\end{Resp}
\switchcolumn*

\LA
\LessonHead{Lectio IV}
\switchcolumn
\EN
\LessonHead{Lesson IV}
\switchcolumn*

\LA
\LessonTitle{Sermo sancti Leónis Papæ}
\switchcolumn
\EN
\LessonTitle{From the Sermons of Pope St. Leo (the Great)}
\switchcolumn*

\LA
\Source{Sermo 1 de Nativitate Domini}
\switchcolumn
\EN
\Source{1st for Christmas}
\switchcolumn*

\LA
\noindent Salvátor noster, dilectíssimi, hódie natus est: gaudeámus. Neque enim fas est locum esse tristítiæ, ubi natális est vitæ: quæ, consúmpto mortalitátis timóre, nobis íngerit de promíssa æternitáte lætítiam. Nemo ab hujus alacritátis participatióne secérnitur. Una cunctis lætítiæ commúnis est rátio: quia Dóminus noster, peccáti mortísque destrúctor, sicut nullum a reátu líberum réperit, ita liberándis ómnibus venit. Exsúltet sanctus, quia appropínquat ad palmam: gáudeat peccátor, quia invitátur ad véniam: animétur gentílis, quia vocátur ad vitam. Dei namque Fílius secúndum plenitúdinem témporis, quam divíni consílii inscrutábilis altitúdo dispósuit, reconciliándam auctóri suo natúram géneris assúmpsit humáni, ut invéntor mortis diábolus, per ipsam, quam vícerat, vincerétur.\par
\switchcolumn
\EN
\noindent Dearly beloved brethren, Unto us is born this day a Saviour, (Luke ii. 11.) Let us rejoice. It would be unlawful to be sad to day, for today is Life's Birthday; the Birthday of that Life, Which, for us dying creatures, taketh away the sting of death, and bringeth the bright promise of the eternal gladness hereafter. It would be unlawful for any man to refuse to partake in our rejoicing. All men have an equal share in the great cause of our joy, for, since our Lord, Who is the destroyer of sin and of death, findeth that all are bound under the condemnation, He is come to make all free. Rejoice, O thou that art holy, thou drawest nearer to thy crown! Rejoice, O thou that art sinful, thy Saviour offereth thee pardon! Rejoice also, O thou Gentile, God calleth thee to life! For the Son of God, when the fulness of the time was come, which had been fixed by the unsearchable counsel of God, took upon Him the nature of man, that He might reconcile that nature to Him Who made it, and so the devil, the inventor of death, is met and beaten in that very flesh which hath been the field of his victory.\par
\switchcolumn*

\LA
\begin{Resp}\Rsign O magnum mystérium, et admirábile sacraméntum, ut animália vidérent Dóminum natum, jacéntem in præsépio: \Star{} Beáta Virgo, cujus víscera meruérunt portáre Dóminum Christum.\end{Resp}
\begin{Resp}\Vsign Ave, María, grátia plena: Dóminus tecum.\end{Resp}
\begin{Resp}\Rsign Beáta Virgo, cujus víscera meruérunt portáre Dóminum Christum.\end{Resp}
\switchcolumn
\EN
\begin{Resp}\Rsign How great is this mystery, how wonderful is the teaching of the faith! The beasts saw the new-born Lord lying in a manger. \Star{} Blessed is that Virgin whose womb was made meet to bear our Lord Christ.\end{Resp}
\begin{Resp}\Vsign Hail, Mary, full of grace the Lord is with thee.\end{Resp}
\begin{Resp}\Rsign Blessed is that Virgin whose womb was made meet to bear our Lord Christ.\end{Resp}
\switchcolumn*

\LA
\LessonHead{Lectio V}
\switchcolumn
\EN
\LessonHead{Lesson V}
\switchcolumn*

\LA
\noindent In quo conflíctu pro nobis ínito, magno et mirábili æquitátis jure certátum est, dum omnípotens Dóminus cum sævíssimo hoste non in sua majestáte, sed in nostra congréditur humilitáte: obíciens ei eándem formam, eandémque natúram, mortalitátis quidem nostræ partícipem, sed peccáti totíus expértem. Aliénum quippe ab hac nativitáte est, quod de ómnibus légitur: Nemo mundus a sorde, nec infans, cujus est uníus diéi vita super terram. Nihil ergo in istam singulárem nativitátem de carnis concupiscéntia transívit, nihil de peccáti lege manávit. Virgo régia Davídicæ stirpis elígitur, quæ sacro gravidánda fœtu, divínam humanámque prolem prius concíperet mente, quam córpore. Et ne supérni ignára consílii ad inusitátos pavéret affátus, quod in ea operándum erat a Spíritu Sancto, collóquio discit angélico: nec damnum credit pudóris, Dei Génetrix mox futúra.\par
\switchcolumn
\EN
\noindent When our Lord entered the field of battle against the devil, He did so with a great and wonderful fairness. Being Himself the Almighty, He laid aside His uncreated Majesty to fight with our cruel enemy in our weak flesh. He brought against him the very shape, the very nature of our mortality, yet without sin. \textasciitilde{}(Heb. iv. 15.) His birth however was not a birth like other births for no other is born pure, nay, not the little child whose life endureth but a day on the earth. To His birth alone the throes of human passion had not contributed, in His alone no consequence of sin had had part. For His Mother was chosen a Virgin of the kingly lineage of David, and when she was to grow heavy with the sacred Child, her soul had already conceived Him before her body. She knew the counsel of God announced to her by the Angel, lest the unwonted events should alarm her. The future Mother of God knew what was to be wrought in her by the Holy Ghost, and that her modesty was absolutely safe.\par
\switchcolumn*

\LA
\begin{Resp}\Rsign Beáta Dei Génetrix María, cujus víscera intácta pérmanent: \Star{} Hódie génuit Salvatórem sǽculi.\end{Resp}
\begin{Resp}\Vsign Beáta, quæ crédidit: quóniam perfécta sunt ómnia, quæ dicta sunt ei a Dómino.\end{Resp}
\begin{Resp}\Rsign Hódie génuit Salvatórem sǽculi.\end{Resp}
\switchcolumn
\EN
\begin{Resp}\Rsign Blessed is God's holy Mother, Mary, maiden undefiled. \Star{} This day hath she brought forth the Saviour of the world.\end{Resp}
\begin{Resp}\Vsign Blessed is she that believed; for there is a performance of all those things which were told her from the Lord.\end{Resp}
\begin{Resp}\Rsign This day hath she brought forth the Saviour of the world.\end{Resp}
\switchcolumn*

\LA
\LessonHead{Lectio VI}
\switchcolumn
\EN
\LessonHead{Lesson VI}
\switchcolumn*

\LA
\noindent Agámus ergo, dilectíssimi, grátias Deo Patri, per Fílium ejus in Spíritu Sancto: qui propter multam caritátem suam, qua diléxit nos, misértus est nostri: et cum essémus mórtui peccátis, convivificávit nos Christo, ut essémus in ipso nova creatúra, novúmque figméntum. Deponámus ergo véterem hóminem cum áctibus suis: et adépti participatiónem generatiónis Christi, carnis renuntiémus opéribus. Agnósce, o Christiáne, dignitátem tuam: et divínæ consors factus natúræ, noli in véterem vilitátem degéneri conversatióne redíre. Meménto, cujus cápitis et cujus córporis sis membrum. Reminíscere, quia érutus de potestáte tenebrárum, translátus es in Dei lumen et regnum.\par
\switchcolumn
\EN
\noindent Therefore, dearly beloved brethren, let us give thanks to God the Father, through His Son, in the Holy Ghost: Who, for His great love wherewith He loved us, hath had mercy on us and, even when we were dead in sins, hath quickened us together with Christ, (Eph. ii. 4, 5,) that in Him we might be a new creature, and a new workmanship. Let us then put off the old man with his deeds (Col. iii. 9); and, having obtained a share in the Sonship of Christ, let us renounce the deeds of the flesh. Learn, O Christian, how great thou art, who hast been made partaker of the Divine nature, (2 Pet. i. 4,) and fall not again by corrupt conversation into the beggarly elements above which thou art lifted. Remember Whose Body it is Whereof thou art made a member, and Who is its Head, (1 Cor. vi. 15.) Remember that it is He That hath delivered thee from the power of darkness and hath translated thee into God's light, and God's kingdom, (Col. i. 13.)\par
\switchcolumn*

\LA
\begin{Resp}\Rsign Sancta et immaculáta virgínitas, quibus te láudibus éfferam, néscio: \Star{} Quia quem cæli cápere non póterant, tuo grémio contulísti.\end{Resp}
\begin{Resp}\Vsign Benedícta tu in muliéribus, et benedíctus fructus ventris tui.\end{Resp}
\begin{Resp}\Rsign Quia quem cæli cápere non póterant, tuo grémio contulísti.\end{Resp}
\begin{Resp}\Vsign Glória Patri, et Fílio, \Star{} et Spirítui Sancto.\end{Resp}
\begin{Resp}\Rsign Quia quem cæli cápere non póterant, tuo grémio contulísti.\end{Resp}
\switchcolumn
\EN
\begin{Resp}\Rsign O Mary, how holy and how spotless is thy virginity! I am too dull to praise thee! \Star{} For thou hast borne in thy breast Him Whom the heavens cannot contain.\end{Resp}
\begin{Resp}\Vsign Blessed art thou among women, and blessed is the fruit of thy womb.\end{Resp}
\begin{Resp}\Rsign For thou hast borne in thy breast Him Whom the heavens cannot contain.\end{Resp}
\begin{Resp}\Vsign Glory be to the Father, and to the Son, \Star{} and to the Holy Ghost.\end{Resp}
\begin{Resp}\Rsign For thou hast borne in thy breast Him Whom the heavens cannot contain.\end{Resp}
\switchcolumn*

\LA
\LessonHead{Lectio VII}
\switchcolumn
\EN
\LessonHead{Lesson VII}
\switchcolumn*

\LA
\LessonTitle{Léctio sancti Evangélii secúndum Lucam}
\switchcolumn
\EN
\LessonTitle{From the Holy Gospel according to Luke}
\switchcolumn*

\LA
\Cite{Luc 2:1-14}
\switchcolumn
\EN
\Cite{Luke 2:1-14}
\switchcolumn*

\LA
\noindent In illo témpore: Exiit edíctum a Cǽsare Augústo, ut describerétur univérsus orbis. Et réliqua.\par
\switchcolumn
\EN
\noindent At that time, it came to pass that there went out a decree from Caesar Augustus, that the whole world should be enrolled. And so on.\par
\switchcolumn*

\LA
\LessonTitle{Homilía sancti Gregórii Papæ}
\switchcolumn
\EN
\LessonTitle{Homily by Pope St. Gregory (the Great)}
\switchcolumn*

\LA
\Source{Homilia 8 in Evangelia}
\switchcolumn
\EN
\Source{8th on the Gospels}
\switchcolumn*

\LA
\noindent Quia, largiénte Dómino, Missárum solémnia ter hódie celebratúri sumus, loqui diu de evangélica lectióne non póssumus: sed nos áliquid vel bréviter dícere, Redemptóris nostri Natívitas ipsa compéllit. Quid est enim, quod nascitúro Dómino mundus descríbitur, nisi hoc, quod apérte monstrátur, quia ille apparébat in carne, qui eléctos suos adscríberet in æternitáte? Quo contra de réprobis per prophétam dícitur: Deleántur de libro vivéntium, et cum justis non scribántur. Qui bene étiam in Béthlehem náscitur: Béthlehem quippe domus panis interpretátur. Ipse namque est, qui ait: Ego sum panis vivus, qui de cælo descéndi. Locus ergo, in quo Dóminus náscitur, domus panis ántea vocátus est; quia futúrum profécto erat, ut ille ibi per matériam carnis apparéret, qui electórum mentes intérna satietáte refíceret. Qui non in paréntum domo, sed in via náscitur: ut profécto osténderet, quia per humanitátem suam, quam assúmpserat, quasi in aliéno nascebátur.\par
\switchcolumn
\EN
\noindent By God's mercy we are to say three Masses today, so that there is not much time left for preaching; but at the same time the occasion of the Lord's Birthday itself obliges me to speak a few words. I will first ask why, when the Lord was to be born, the world was enrolled? Was it not to herald the appearing of Him by Whom the elect are enrolled in the book of life? Whereas the Prophet saith of\par
the reprobate: Let them be blotted out of the book of the living, and not be written with the righteous. (Ps. lxviii. 29.) Then, the Lord is born in Bethlehem. Now the name Bethlehem signifieth the House of Bread, and thus it is the birth-place of Him Who hath said: I am the Living Bread, Which came down from heaven. (John vi. 51.) We see then that this name of Bethlehem was prophetically given to the place where Christ was born, because it was there that He was to appear in the flesh by Whom the souls of the faithful are fed unto life eternal. He was born, not in His Mother's house, but away from home. And this is a mystery, showing that this our mortality into which He was born was not the home of Him Who is begotten of the Father before the worlds.\par
\switchcolumn*

\LA
\begin{Resp}\Rsign Beáta víscera Maríæ Vírginis, quæ portavérunt ætérni Patris Fílium: et beáta úbera, quæ lactavérunt Christum Dóminum: \Star{} Qui hódie pro salúte mundi de Vírgine nasci dignátus est.\end{Resp}
\begin{Resp}\Vsign Dies sanctificátus illúxit nobis: veníte, gentes, et adoráte Dóminum.\end{Resp}
\begin{Resp}\Rsign Qui hódie pro salúte mundi de Vírgine nasci dignátus est.\end{Resp}
\switchcolumn
\EN
\begin{Resp}\Rsign Blessed be the womb of the Virgin Mary, which bore the Son of the Eternal Father, and blessed be the paps which give suck to Christ our Lord. \Star{} This day hath He been pleased for the salvation of the world to be born of a Virgin.\end{Resp}
\begin{Resp}\Vsign This day which is breaking is holy; O come, ye Gentiles, and worship the Lord.\end{Resp}
\begin{Resp}\Rsign This day hath He been pleased for the salvation of the world to be born of a Virgin.\end{Resp}
\switchcolumn*

\LA
\LessonHead{Lectio VIII}
\switchcolumn
\EN
\LessonHead{Lesson VIII}
\switchcolumn*

\LA
\LessonTitle{Léctio sancti Evangélii secúndum Lucam}
\switchcolumn
\EN
\LessonTitle{From the holy Gospel according to Luke}
\switchcolumn*

\LA
\Cite{Luc 2:15-20}
\switchcolumn
\EN
\Cite{Luke 2:15-20}
\switchcolumn*

\LA
\noindent In illo témpore: Pastóres loquebántur ad ínvicem: Transeámus usque Béthlehem, et videámus hoc verbum, quod factum est, quod Dóminus osténdit nobis. Et réliqua.\par
\switchcolumn
\EN
\noindent At that time, the shepherds said one to another: Let us go over to Bethlehem, and let us see this word that is come to pass, which the Lord hath showed to us. And so on.\par
\switchcolumn*

\LA
\LessonTitle{Homilía sancti Ambrósii Epíscopi}
\switchcolumn
\EN
\LessonTitle{Homily by St. Ambrose, Bishop (of Milan)}
\switchcolumn*

\LA
\Source{Lib. 2 in cap. 2. Lucæ, circa medium}
\switchcolumn
\EN
\Source{Book II on Luke II}
\switchcolumn*

\LA
\noindent Vidéte Ecclésiæ surgéntis exórdium: Christus náscitur, et pastóres vigiláre cœpérunt; qui géntium greges, pécudum modo ante vivéntes, in caulam Dómini congregárent, ne quos spiritálium bestiárum per offúsas nóctium ténebras pateréntur incúrsus. Et bene pastóres vígilant, quos bonus pastor infórmat. Grex ígitur pópulus, nox sǽculum, pastóres sunt sacerdótes. Aut fortásse étiam ille sit pastor, cui dícitur: Esto vígilans, et confírma. Quia non solum epíscopos ad tuéndum gregem Dóminus ordinávit, sed étiam Angelos destinávit.\par
\switchcolumn
\EN
\noindent Behold the beginning of the Church. Christ is born, and the shepherds watch; shepherds, to gather together the scattered sheep of the Gentiles, and to lead them into the fold of Christ, that they might no longer be a prey to the ravages of spiritual wolves in the night of this world's darkness. And that shepherd is wide awake, whom the Good Shepherd stirreth up. The flock then is the people, the night is the world, and the shepherds are the Priests. And perhaps he is a shepherd to whom it is said: Be watchful and strengthen, (Apoc. iii. 2,) for God hath ordained as the shepherds of His flock not Bishops only, but also Angels.\par
\switchcolumn*

\LA
\begin{Resp}\Rsign Verbum caro factum est, et habitávit in nobis: \Star{} Et vídimus glóriam ejus, glóriam quasi Unigéniti a Patre, plenum grátiæ et veritátis.\end{Resp}
\begin{Resp}\Vsign Omnia per ipsum facta sunt, et sine ipso factum est nihil.\end{Resp}
\begin{Resp}\Rsign Et vídimus glóriam ejus, glóriam quasi Unigéniti a Patre, plenum grátiæ et veritátis.\end{Resp}
\begin{Resp}\Vsign Glória Patri, et Fílio, \Star{} et Spirítui Sancto.\end{Resp}
\begin{Resp}\Rsign Et vídimus glóriam ejus, glóriam quasi Unigéniti a Patre, plenum grátiæ et veritátis.\end{Resp}
\switchcolumn
\EN
\begin{Resp}\Rsign The Word was made flesh, and dwelt among us. \Star{} And we beheld His glory, the glory as of the Only-begotten of the Father, full of grace and truth.\end{Resp}
\begin{Resp}\Vsign All things were made by Him, and without Him was not anything made.\end{Resp}
\begin{Resp}\Rsign And we beheld His glory, the glory as of the Only-begotten of the Father, full of grace and truth.\end{Resp}
\begin{Resp}\Vsign Glory be to the Father, and to the Son, \Star{} and to the Holy Ghost.\end{Resp}
\begin{Resp}\Rsign And we beheld His glory, the glory as of the Only-begotten of the Father, full of grace and truth.\end{Resp}
\switchcolumn*

\LA
\LessonHead{Lectio IX}
\switchcolumn
\EN
\LessonHead{Lesson IX}
\switchcolumn*

\LA
\LessonTitle{Léctio sancti Evangélii secúndum Joánnem}
\switchcolumn
\EN
\LessonTitle{From the holy Gospel according to John}
\switchcolumn*

\LA
\Cite{Joannes 1:1-14}
\switchcolumn
\EN
\Cite{John 1:1-14}
\switchcolumn*

\LA
\noindent In princípio erat Verbum, et Verbum erat apud Deum, et Deus erat Verbum. Et réliqua.\par
\switchcolumn
\EN
\noindent In the beginning was the Word, and the Word was with God, and the Word was God. The same was in the beginning with God. And so on.\par
\switchcolumn*

\LA
\LessonTitle{Homilía sancti Augustíni Epíscopi}
\switchcolumn
\EN
\LessonTitle{Homily by St. Augustine, Bishop (of Hippo)}
\switchcolumn*

\LA
\Source{Tract. 1 in Joann., circa médium}
\switchcolumn
\EN
\Source{1st Tract on John}
\switchcolumn*

\LA
\noindent Ne vile áliquid putáres quale consuevísti cogitáre, cum verba humána soléres audíre, audi quid cógites: Deus erat Verbum. Exeat nunc néscio quis infidélis Ariánus, et dicat quia Verbum Dei factum est. Quómodo potest fíeri, ut Verbum Dei factum sit, quando Deus per Verbum fecit ómnia? Si et Verbum Dei ipsum factum est: per quod áliud verbum factum est? Si hoc dicis, quia hoc est verbum Verbi, per quod factum est illud; ipsum dico ego únicum Fílium Dei. Si autem non dicis verbum Verbi, concéde non factum, per quod facta sunt ómnia. Non enim per seípsum fíeri pótuit, per quod facta sunt ómnia. Crede ergo Evangelístæ.\par
\switchcolumn
\EN
\noindent Lest thou shouldest think all things mean, as thou art accustomed to think of things human, hear and digest this The Word was God. Now perhaps there will come forward some Arian unbeliever, and say that the Word of God was a creature. How can the Word of God be a creature, when it was by the Word that all creatures were made? If He be a creature, then there must have been some other Word, not a creature, by which He was made. And what Word is that? If thou sayest that it was by the word of the Word Himself that He was made, I tell thee that God had no other, but One Only-begotten Son. But if thou say not that it was by the word of the Word Himself that He was made, thou art forced to confess that. He by Whom all things were made was not Himself made at all. Believe the Gospel.\par
\switchcolumn*

\LA
\TeDeum{Te Deum}
\switchcolumn
\EN
\TeDeum{Te Deum}
\switchcolumn*

\end{paracol}

\par\needspace{10\baselineskip}\addvspace{1.2em}{\centering\small\textsc{Die 26 Decembris} · \textsc{26 December}\par}
\Office{S. Stephani Protomartyris}{St. Stephen the First Martyr}{Duplex II classis cum Octava simplici}
\addcontentsline{toc}{subsection}{Die 26 Decembris · S. Stephani Protomartyris\enspace\textit{St. Stephen the First Martyr}}
\begin{paracol}{2}
\LA
\LessonHead{Lectio I}
\switchcolumn
\EN
\LessonHead{Lesson I}
\switchcolumn*

\LA
\LessonTitle{De Actibus Apostolórum}
\switchcolumn
\EN
\LessonTitle{Lesson from the Acts of Apostles}
\switchcolumn*

\LA
\Cite{Act. 6:1-6}
\switchcolumn
\EN
\Cite{Acts 6:1-6}
\switchcolumn*

\LA
\noindent In diébus illis, crescénte número discipulórum, factum est murmur Græcórum advérsus Hebrǽos, eo quod despiceréntur in ministério cotidiáno víduæ eórum. Convocántes autem duódecim multitúdinem discipulórum, dixérunt: Non est æquum nos derelínquere verbum Dei et ministráre mensis. Consideráte ergo, fratres, viros ex vobis boni testimónii, septem, plenos Spíritu Sancto et sapiéntia, quos constituámus super hoc opus. Nos vero oratióni et ministério verbi instántes érimus. Et plácuit sermo coram omni multitúdine. Et elegérunt Stéphanum virum plenum fide et Spíritu Sancto, et Philíppum, et Próchorum, et Nicánorem, et Timónem, et Pármenam, et Nicoláum, ádvenam Antiochénum. Hos statuérunt ante conspéctum Apostolórum: et orántes imposuérunt eis manus.\par
\switchcolumn
\EN
\noindent And in those days, the number of the disciples increasing, there arose a murmuring of the Greeks against the Hebrews, for that their widows were neglected in the daily ministration. Then the twelve calling together the multitude of the disciples, said: It is not reason that we should leave the word of God, and serve tables. Wherefore, brethren, look ye out among you seven men of good reputation, full of the Holy Ghost and wisdom, whom we may appoint over this business. But we will give ourselves continually to prayer, and to the ministry of the word. And the saying was liked by all the multitude. And they chose Stephen, a man full of faith, and of the Holy Ghost, and Philip, and Prochorus, and Nicanor, and Timon, and Parmenas, and Nicolas, a proselyte of Antioch. These they set before the apostles; and they praying, imposed hands upon them.\par
\switchcolumn*

\LA
\begin{Resp}\Rsign Stéphanus autem plenus grátia et fortitúdine, \Star{} Faciébat prodígia et signa magna in pópulo.\end{Resp}
\begin{Resp}\Vsign Surrexérunt quidam de synagóga disputántes cum Stéphano: et non póterant resístere sapiéntiæ, et Spirítui qui loquebátur.\end{Resp}
\begin{Resp}\Rsign Faciébat prodígia et signa magna in pópulo.\end{Resp}
\switchcolumn
\EN
\begin{Resp}\Rsign And Stephen, full of grace and power \Star{} Did great wonders and miracles among the people.\end{Resp}
\begin{Resp}\Vsign There arose certain of the synagogue, disputing with Stephen; and they were not able to resist the wisdom, and the Spirit which spoke.\end{Resp}
\begin{Resp}\Rsign Did great wonders and miracles among the people.\end{Resp}
\switchcolumn*

\LA
\LessonHead{Lectio II}
\switchcolumn
\EN
\LessonHead{Lesson II}
\switchcolumn*

\LA
\Cite{Act. 6:7-10;7:54}
\switchcolumn
\EN
\Cite{Acts 6:7-10;7:54}
\switchcolumn*

\LA
\noindent Et verbum Dómini crescébat, et multiplicabátur númerus discipulórum in Jerúsalem valde: multa étiam turba sacerdótum obediébat fídei. Stéphanus autem plenus grátia et fortitúdine, faciébat prodígia et signa magna in pópulo. Surrexérunt autem quidam de synagóga, quæ appellátur Libertinórum, et Cyrenénsium et Alexandrinórum, et eórum qui erant a Cilícia, et Asia, disputántes cum Stéphano: et non póterant resístere sapiéntiæ, et Spirítui qui loquebátur. Audiéntes autem hæc, dissecabántur córdibus suis, et stridébant déntibus in eum.\par
\switchcolumn
\EN
\noindent And the word of the Lord increased; and the number of the disciples was multiplied in Jerusalem exceedingly: a great multitude also of the priests obeyed the faith. And Stephen, full of grace and fortitude, did great wonders and signs among the people. Now there arose some of that which is called the synagogue of the Libertines, and of the Cyrenians, and of the Alexandrians, and of them that were of Cilicia and Asia, disputing with Stephen. And they were not able to resist the wisdom and the spirit that spoke. Now hearing these things, they were cut to the heart, and they gnashed with their teeth at him.\par
\switchcolumn*

\LA
\begin{Resp}\Rsign Vidébant omnes Stéphanum, qui erant in concílio: \Star{} Et intuebántur vultum ejus tamquam vultum Angeli stantis inter illos.\end{Resp}
\begin{Resp}\Vsign Plenus grátia et fortitúdine, faciébat prodígia et signa magna in pópulo.\end{Resp}
\begin{Resp}\Rsign Et intuebántur vultum ejus tamquam vultum Angeli stantis inter illos.\end{Resp}
\switchcolumn
\EN
\begin{Resp}\Rsign All that sat in the council, looking steadfastly on Stephen, saw \Star{} His face as it had been the face of an angel standing among them.\end{Resp}
\begin{Resp}\Vsign Full of grace and power, he did great wonders and miracles among the people.\end{Resp}
\begin{Resp}\Rsign His face as it had been the face of an angel standing among them.\end{Resp}
\switchcolumn*

\LA
\LessonHead{Lectio III}
\switchcolumn
\EN
\LessonHead{Lesson III}
\switchcolumn*

\LA
\Cite{Act. 7:55-59}
\switchcolumn
\EN
\Cite{Acts 7:55-59}
\switchcolumn*

\LA
\noindent Cum autem esset plenus Spíritu Sancto, inténdens in cælum, vidit glóriam Dei, et Jesum stantem a dextris Dei. Et ait: Ecce vídeo cælos apértos, et Fílium hóminis stantem a dextris Dei. Exclamántes autem voce magna, continuérunt aures suas, et ímpetum fecérunt unanímiter in eum. Et eiciéntes eum extra civitátem lapidábant: et testes deposuérunt vestiménta sua secus pedes adolescéntis, qui vocabátur Saulus. Et lapidábant Stéphanum, invocántem et dicéntem: Dómine Jesu, súscipe spíritum meum. Pósitis autem génibus, clamávit voce magna, dicens: Dómine, ne státuas illis hoc peccátum. Et cum hoc dixísset, obdormívit in Dómino.\par
\switchcolumn
\EN
\noindent But he, being full of the Holy Ghost, looking up steadfastly to heaven, saw the glory of God, and Jesus standing on the right hand of God. And he said: Behold, I see the heavens opened, and the Son of man standing on the right hand of God. And they crying out with a loud voice, stopped their ears, and with one accord ran violently upon him. And casting him forth without the city, they stoned him; and the witnesses laid down their garments at the feet of a young man, whose name was Saul. And they stoned Stephen, invoking, and saying: Lord Jesus, receive my spirit. And falling on his knees, he cried with a loud voice, saying: Lord, lay not this sin to their charge. And when he had said this, he fell asleep in the Lord.\par
\switchcolumn*

\LA
\begin{Resp}\Rsign Intuens in cælum beátus Stéphanus, vidit glóriam Dei, et ait: \Star{} Ecce vídeo cælos apértos, et Fílium hóminis stantem a dextris virtútis Dei.\end{Resp}
\begin{Resp}\Vsign Cum autem esset Stéphanus plenus Spíritu Sancto, inténdens in cælum, vidit glóriam Dei, et ait.\end{Resp}
\begin{Resp}\Rsign Ecce vídeo cælos apértos, et Fílium hóminis stantem a dextris virtútis Dei.\end{Resp}
\begin{Resp}\Vsign Glória Patri, et Fílio, \Star{} et Spirítui Sancto.\end{Resp}
\begin{Resp}\Rsign Ecce vídeo cælos apértos, et Fílium hóminis stantem a dextris virtútis Dei.\end{Resp}
\switchcolumn
\EN
\begin{Resp}\Rsign The blessed Stephen looked up steadfastly into heaven, and saw the glory of God, and said \Star{} Behold, I see the heavens opened, and the Son of Man standing at the right hand of the power of God.\end{Resp}
\begin{Resp}\Vsign But Stephen, being full of the Holy Ghost, looked up steadfastly into heaven, and saw the glory of God, and said\end{Resp}
\begin{Resp}\Rsign Behold, I see the heavens opened, and the Son of man standing at the right hand of the power of God.\end{Resp}
\begin{Resp}\Vsign Glory be to the Father, and to the Son, \Star{} and to the Holy Ghost.\end{Resp}
\begin{Resp}\Rsign Behold, I see the heavens opened, and the Son of Man standing at the right hand of the power of God.\end{Resp}
\switchcolumn*

\LA
\LessonHead{Lectio IV}
\switchcolumn
\EN
\LessonHead{Lesson IV}
\switchcolumn*

\LA
\LessonTitle{Sermo sancti Fulgéntii Epíscopi}
\switchcolumn
\EN
\LessonTitle{From the Sermons of St. Fulgentius, Bishop (of Ruspe.)}
\switchcolumn*

\LA
\Source{Sermo 3, de S. Stéphano, circa init.}
\switchcolumn
\EN
\Source{On St. Stephen}
\switchcolumn*

\LA
\noindent Heri celebrávimus temporálem sempitérni Regis nostri natálem: hódie celebrámus triumphálem mílitis passiónem. Heri enim Rex noster trábea carnis indútus, de aula úteri virginális egrédiens, visitáre dignátus est mundum: hódie miles de tabernáculo córporis éxiens, triumphátor migrávit ad cælum. Ille, sempitérnæ Deitátis majestáte serváta, servíle cinctórium carnis assúmens, in hujus sǽculi campum pugnatúrus intrávit; iste, depósitis corruptibílibus córporis induméntis, ad cæli palátium perénniter regnatúrus ascéndit. Ille descéndit carne velátus, iste ascéndit sánguine laureátus.\par
\switchcolumn
\EN
\noindent Yesterday we were celebrating the birth in time of our Eternal King; today we celebrate the victory, through suffering, of one of His soldiers. Yesterday our King was pleased to come forth from His royal palace of the Virgin's womb, clothed in a robe of flesh, to visit the world; today His soldier, laying aside the tabernacle of the body, entereth in triumph into the heavenly palaces. The One, preserving unchanged that glory of the Godhead which He had before the world was, girded Himself with the form of a servant, and entered the arena of this world to fight sin; the other taketh off the garments of this corruptible body, and entereth into the heavenly mansions, where he will reign for ever. The One cometh down, veiled in flesh; the other goeth up, clothed in a robe of glory, red with blood.\par
\switchcolumn*

\LA
\begin{Resp}\Rsign Lapidábant Stéphanum invocántem et dicéntem: \Star{} Dómine Jesu Christe, áccipe spíritum meum: et ne státuas illis hoc peccátum.\end{Resp}
\begin{Resp}\Vsign Pósitis autem génibus, clamávit voce magna, dicens.\end{Resp}
\begin{Resp}\Rsign Dómine Jesu Christe, áccipe spíritum meum: et ne státuas illis hoc peccátum.\end{Resp}
\switchcolumn
\EN
\begin{Resp}\Rsign They stoned Stephen, calling upon God and saying: \Star{} Lord Jesus Christ, receive my spirit; and lay not this sin to their charge.\end{Resp}
\begin{Resp}\Vsign And he kneeled down, and cried with a loud voice, saying\end{Resp}
\begin{Resp}\Rsign Lord Jesus Christ, receive my spirit; and lay not this sin to their charge.\end{Resp}
\switchcolumn*

\LA
\LessonHead{Lectio V}
\switchcolumn
\EN
\LessonHead{Lesson V}
\switchcolumn*

\LA
\noindent Ascéndit iste lapidántibus Judǽis, quia ille descéndit lætántibus Angelis. Glória in excélsis Deo, heri sancti Angeli exsultántes cantavérunt: hódie Stéphanum lætántes in suum consórtium suscepérunt. Heri Dóminus exívit de útero Vírginis: hódie miles egréssus est de ergástulo carnis. Heri Christus pro nobis pannis est involútus: hódie Stéphanus stola est ab eo immortalitátis indútus. Heri præsépis angústia Christum portávit infántem: hódie imménsitas cæli suscépit Stéphanum triumphántem. Solus Dóminus descéndit, ut multos eleváret: humiliávit se Rex noster, ut suos mílites exaltáret.\par
\switchcolumn
\EN
\noindent The One cometh down amid the jubilation of angels; the other goeth up amid the stoning of the Jews. Yesterday the holy angels were singing, Glory to God in the highest; today there is joy among them, for they receive Stephen into their company. Yesterday the Lord came forth from the Virgin's womb; today His soldier is delivered from the prison of the body. Yesterday Christ was for our sakes wrapped in swaddling bands; today He girdeth Stephen with a robe of immortality. Yesterday the new-born Christ lay in a narrow manger; today Stephen entereth victorious into the boundless heavens. The Lord came down alone that He might raise many up; our King humbled Himself that He might set His soldiers in high places.\par
\switchcolumn*

\LA
\begin{Resp}\Rsign Impetum fecérunt unanímiter in eum, et ejecérunt eum extra civitátem, invocántem et dicéntem: \Star{} Dómine Jesu, áccipe spíritum meum.\end{Resp}
\begin{Resp}\Vsign Et testes deposuérunt vestiménta sua secus pedes adolescéntis, qui vocabátur Saulus: et lapidábant Stéphanum invocántem et dicéntem.\end{Resp}
\begin{Resp}\Rsign Dómine Jesu, áccipe spíritum meum.\end{Resp}
\switchcolumn
\EN
\begin{Resp}\Rsign They ran upon him with one accord, and cast him out of the city, calling upon God, and saying: \Star{} Lord Jesus, receive my spirit.\end{Resp}
\begin{Resp}\Vsign And the witnesses laid down their clothes at a young man's feet, whose name was Saul; and they stoned Stephen, calling upon God, and saying\end{Resp}
\begin{Resp}\Rsign Lord Jesus, receive my spirit.\end{Resp}
\switchcolumn*

\LA
\LessonHead{Lectio VI}
\switchcolumn
\EN
\LessonHead{Lesson VI}
\switchcolumn*

\LA
\noindent Necessárium tamen nobis est, fratres, agnóscere, quibus armis præcínctus Stéphanus sævítiam Judæórum pótuit superáre, ut ita meruísset felíciter triumpháre. Stéphanus ergo, ut nóminis sui corónam meruísset accípere, caritátem pro armis habébat, et per ipsam ubíque vincébat. Per caritátem Dei sæviéntibus Judǽis non cessit: per caritátem próximi pro lapidántibus intercéssit. Per caritátem arguébat errántes, ut corrigeréntur: per caritátem pro lapidántibus orábat, ne puniréntur. Caritátis virtúte subníxus, vicit Saulum crudéliter sæviéntem; et quem hábuit in terra persecutórem, in cælo méruit habére consórtem.\par
\switchcolumn
\EN
\noindent Why brethren, it behoveth us to consider with what arms Stephen was able, amid all the cruelty of the Jews, to remain more than conqueror, and worthily to attain to so blessed a triumph. Stephen, in that struggle which brought him to the crown whereof his name is a prophecy, had for armour the love of God and man, and by it he remained victorious on all hands. The love of God strengthened him against the cruelty of the Jews; and the love of his neighbour made him pray even for his murderers. Through love he rebuked the wandering, that they might be corrected; through love he prayed for them that stoned him, that they might not be punished. By the might of his love he overcame Saul his cruel persecutor; and earned for a comrade in heaven, the very man who had done him to death upon earth.\par
\switchcolumn*

\LA
\begin{Resp}\Rsign Impii super justum jactúram fecérunt, ut eum morti tráderent: \Star{} At ille gaudens suscépit lápides, ut mererétur accípere corónam glóriæ.\end{Resp}
\begin{Resp}\Vsign Continuérunt aures suas, et ímpetum fecérunt unanímiter in eum.\end{Resp}
\begin{Resp}\Rsign At ille gaudens suscépit lápides, ut mererétur accípere corónam glóriæ.\end{Resp}
\begin{Resp}\Vsign Glória Patri, et Fílio, \Star{} et Spirítui Sancto.\end{Resp}
\begin{Resp}\Rsign At ille gaudens suscépit lápides, ut mererétur accípere corónam glóriæ.\end{Resp}
\switchcolumn
\EN
\begin{Resp}\Rsign The ungodly fell upon the righteous, to put him to death. \Star{} But he received the stones with joy, that he might earn a crown of glory.\end{Resp}
\begin{Resp}\Vsign They stopped their ears, and ran upon him with one accord.\end{Resp}
\begin{Resp}\Rsign But he received the stones with joy, that he might earn a crown of glory.\end{Resp}
\begin{Resp}\Vsign Glory be to the Father, and to the Son, \Star{} and to the Holy Ghost.\end{Resp}
\begin{Resp}\Rsign But he received the stones with joy, that he might earn a crown of glory.\end{Resp}
\switchcolumn*

\LA
\LessonHead{Lectio VII}
\switchcolumn
\EN
\LessonHead{Lesson VII}
\switchcolumn*

\LA
\LessonTitle{Léctio sancti Evangélii secúndum Matthǽum}
\switchcolumn
\EN
\LessonTitle{From the Holy Gospel according to Matthew}
\switchcolumn*

\LA
\Cite{Matt 23:34-39}
\switchcolumn
\EN
\Cite{Matt 23:34-39}
\switchcolumn*

\LA
\noindent In illo témpore: Dicébat Jesus scribis et pharisǽis: Ecce ego mitto ad vos prophétas, et sapiéntes, et scribas: et ex illis occidétis, et crucifigétis. Et réliqua.\par
\switchcolumn
\EN
\noindent At that time, Jesus said to the scribes and pharisee: I send to you prophets, and wise men, and scribes: and some of them you will put to death and crucify. And so on,\par
\switchcolumn*

\LA
\LessonTitle{Homilía sancti Hierónymi Presbýteri}
\switchcolumn
\EN
\LessonTitle{Homily by St. Jerome, Priest (at Bethlehem.)}
\switchcolumn*

\LA
\Source{Liber 4 Comment. in cap. 23 Matthæi}
\switchcolumn
\EN
\Source{Book IV, Commentary on Matth. XXIII}
\switchcolumn*

\LA
\noindent Hoc quod ántea dixerámus, Impléte mensúram patrum vestrórum; ad persónam Dómini pertinére, eo quod occidéndus esset ab eis, potest et ad discípulos ejus reférri, de quibus nunc dicit: Ecce ego mitto ad vos prophétas, et sapiéntes, et scribas. Simúlque obsérva, juxta Apóstolum scribéntem ad Corínthios, vária esse dona discipulórum Christi; álios prophétas, qui ventúra prædícent; álios sapiéntes, qui novérunt quando débeant proférre sermónem; álios scribas in lege doctíssimos, ex quibus lapidátus est Stéphanus, Paulus occísus, crucifíxus Petrus, flagelláti in Actibus Apostolórum discípuli.\par
\switchcolumn
\EN
\noindent We have already remarked that the Lord's words, Fill ye up the measure of your fathers, (32,) refer in the first place to Himself, Whom the Jews afterwards put to death. In a secondary sense it may likewise be applied to His disciples, of whom He saith, Behold, I send unto you Prophets, and wise men, and Scribes. Here observe that, according to the Apostle writing to the Corinthians, (1 Cor. xii. 4,) there are diversities of gifts among Christ's followers. Some are Prophets of that which is to come; some are wise men, who know the due season for rebuke and exhortation; some are Scribes learned in the law. And of these they stoned Stephen, slew Paul with the sword, crucified Peter, and scourged the Disciples mentioned in the Acts of the Apostles. (v, 40; xvi, 23.)\par
\switchcolumn*

\LA
\begin{Resp}\Rsign Stéphanus servus Dei, quem lapidábant Judǽi, vidit cælos apértos: vidit, et introívit: \Star{} Beátus homo, cui cæli patébant.\end{Resp}
\begin{Resp}\Vsign Cum ígitur saxórum crepitántium túrbine quaterétur, inter æthéreos aulæ cæléstis sinus divína ei cláritas fulsit.\end{Resp}
\begin{Resp}\Rsign Beátus homo, cui cæli patébant.\end{Resp}
\switchcolumn
\EN
\begin{Resp}\Rsign Stephen, the servant of God, who was stoned by the Jews, saw the heavens opened he saw and entered in. \Star{} Blessed is he, unto whom the heavens were opened.\end{Resp}
\begin{Resp}\Vsign While his poor body was crushed by the hurtling shower of stones, God's brightness broke upon him out of the heavenly palaces.\end{Resp}
\begin{Resp}\Rsign Blessed is he unto whom the heavens were opened.\end{Resp}
\switchcolumn*

\LA
\LessonHead{Lectio VIII}
\switchcolumn
\EN
\LessonHead{Lesson VIII}
\switchcolumn*

\LA
\noindent Quǽrimus, quis iste sit Zacharías fílius Barachíæ: quia multos légimus Zacharías. Et ne líbera nobis tribuerétur erróris facúltas, ádditum est: Quem occidístis inter templum et altáre. In divérsis divérsa legi: et débeo singulórum opiniónes pónere. Alii Zacharíam fílium Barachíæ dicunt, qui in duódecim prophétis undécimus est, patrísque in eo nomen conséntit: sed ubi occísus sit inter templum et altáre, Scriptúra non lóquitur; máxime cum tempóribus ejus vix ruínæ templi fúerint. Alii Zacharíam patrem Joánnis intellégi volunt, ex quibúsdam apocryphórum sómniis approbántes, quod proptérea occísus sit, quia Salvatóris prædicáverit advéntum.\par
\switchcolumn
\EN
\noindent It is a subject of dispute among commentators who is meant by Zacharias the son of Barachias. We read of several persons of the name of Zacharias. But here, in order to prevent any mistake, it is particularly said, Whom ye slew between the temple and the altar. I have read various opinions in various places upon this question, and I will give each. First, some hold that Zacharias the son of Barachias is the eleventh of the twelve Minor Prophets; and this opinion is supported by the father's name. But the Bible nowhere telleth us that this Prophet was slain between the temple and the altar; and it is hardly possible that he can have been, for in his time it could scarcely be said that even the ruins of the temple were in existence. Secondly, others maintain that this Zacharias was Zacharias, the father of John the Baptist. This interpretation is derived from the dreams of the Apocryphal Gospels, wherein it is asserted that he was martyred for preaching Christ's coming.\par
\switchcolumn*

\LA
\begin{Resp}\Rsign Patefáctæ sunt jánuæ cæli Christi Mártyri beáto Stéphano, qui in número Mártyrum invéntus est primus: \Star{} Et ídeo triúmphat in cælis coronátus.\end{Resp}
\begin{Resp}\Vsign Mortem enim, quam Salvátor noster dignátus est pro nobis pati, hanc ille primus réddidit Salvatóri.\end{Resp}
\begin{Resp}\Rsign Et ídeo triúmphat in cælis coronátus.\end{Resp}
\begin{Resp}\Vsign Glória Patri, et Fílio, \Star{} et Spirítui Sancto.\end{Resp}
\begin{Resp}\Rsign Et ídeo triúmphat in cælis coronátus.\end{Resp}
\switchcolumn
\EN
\begin{Resp}\Rsign The gates of heaven were opened to Christ's blessed martyr Stephen, and he is the first of all the martyrs. \Star{} Wherefore he reigneth crowned in heaven.\end{Resp}
\begin{Resp}\Vsign For he was the first to make an offering of his death to that Saviour Who vouchsafed to suffer death for us.\end{Resp}
\begin{Resp}\Rsign Wherefore he reigneth crowned in heaven.\end{Resp}
\begin{Resp}\Vsign Glory be to the Father, and to the Son, \Star{} and to the Holy Ghost.\end{Resp}
\begin{Resp}\Rsign Wherefore he reigneth crowned in heaven.\end{Resp}
\switchcolumn*

\LA
\LessonHead{Lectio IX}
\switchcolumn
\EN
\LessonHead{Lesson IX}
\switchcolumn*

\LA
\noindent Alii istum volunt esse Zacharíam, qui occísus est a Joas rege Judæ inter templum et altáre, sicut Regum narrat história. Sed observándum, quod ille Zacharías non sit fílius Barachíæ, sed fílius Jójadæ sacerdótis: unde et Scriptúra refert: Non fuit recordátus Joas patris ejus Jójadæ, quia sibi fecísset bona. Cum ergo et Zacharíam teneámus et occisiónis conséntiat locus: quǽrimus, quare Barachíæ dicátur fílius, et non Jójadæ? Barachía lingua nostra Benedíctus Dómini dícitur, et sacerdótis Jójadæ justítia Hebrǽo nómine demonstrátur. In Evangélio, quo utúntur Nazaréni, pro fílio Barachíæ, fílium Jójadæ reperímus scriptum.\par
\switchcolumn
\EN
\noindent A third school will have it that this Zacharias, the son of Barachias, was that Zacharias of whom we read, in Chron. xxiv. 22, that he was slain by Joash, king of Judah, between the temple and the altar. Against this it is to be remarked, that that Zacharias was not the son of Barachias, but of Jehoiada the priest; whence it is written, Joash remembered not the kindness which Jehoiada his father had done to him. The question therefore ariseth, if this opinion be true, why, the name and manner of death both agreeing with this explanation, Zacharias is called the son, not of Jehoiada, but of Barachias. In Hebrew, Barachias signifieth the Blessed of the Lord, and Jehoiada proves his Righteousness. In the Gospel used by the Nazarenes the name of Jehoiada is used instead of Barachias.\par
\switchcolumn*

\LA
\TeDeum{Te Deum}
\switchcolumn
\EN
\TeDeum{Te Deum}
\switchcolumn*

\end{paracol}

\par\needspace{10\baselineskip}\addvspace{1.2em}{\centering\small\textsc{Die 27 Decembris} · \textsc{27 December}\par}
\Office{S. Joannis Apostoli et Evangelistæ}{St. John the Apostle, Evangelist}{Duplex II classis cum Octava simplici}
\addcontentsline{toc}{subsection}{Die 27 Decembris · S. Joannis Apostoli et Evangelistæ\enspace\textit{St. John the Apostle, Evangelist}}
\begin{paracol}{2}
\LA
\LessonHead{Lectio I}
\switchcolumn
\EN
\LessonHead{Lesson I}
\switchcolumn*

\LA
\LessonTitle{Incipit Epístola prima beáti Joánnis Apóstoli}
\switchcolumn
\EN
\LessonTitle{Lesson from the first letter of St. John the Apostle}
\switchcolumn*

\LA
\Cite{1 Joannes 1:1-5}
\switchcolumn
\EN
\Cite{1 John 1:1-5}
\switchcolumn*

\LA
\noindent Quod fuit ab inítio, quod audívimus, quod vídimus óculis nostris, quod perspéximus, et manus nostræ contrectavérunt de verbo vitæ: et vita manifestáta est, et vídimus, et testámur, et annuntiámus vobis vitam ætérnam, quæ erat apud Patrem, et appáruit nobis: quod vídimus, et audívimus, annuntiámus vobis, ut et vos societátem habeátis nobíscum, et socíetas nostra sit cum Patre et cum Fílio ejus Jesu Christo. Et hæc scríbimus vobis ut gaudeátis, et gáudium vestrum sit plenum. Et hæc est annuntiátio, quam audívimus ab eo, et annuntiámus vobis: Quóniam Deus lux est, et ténebræ in eo non sunt ullæ.\par
\switchcolumn
\EN
\noindent That which was from the beginning, which we have heard, which we have seen with our eyes, which we have looked upon, and our hands have handled, of the word of life: For the life was manifested; and we have seen and do bear witness, and declare unto you the life eternal, which was with the Father, and hath appeared to us: That which we have seen and have heard, we declare unto you, that you also may have fellowship with us, and our fellowship may be with the Father, and with his Son Jesus Christ. And these things we write to you, that you may rejoice, and your joy may be full. And this is the declaration which we have heard from him, and declare unto you: That God is light, and in him there is no darkness.\par
\switchcolumn*

\LA
\begin{Resp}\Rsign Valde honorándus est beátus Joánnes, qui supra pectus Dómini in cena recúbuit: \Star{} Cui Christus in cruce Matrem vírginem vírgini commendávit.\end{Resp}
\begin{Resp}\Vsign Virgo est eléctus a Dómino, atque inter céteros magis diléctus.\end{Resp}
\begin{Resp}\Rsign Cui Christus in cruce Matrem vírginem vírgini commendávit.\end{Resp}
\switchcolumn
\EN
\begin{Resp}\Rsign Right worthy of honour is the blessed Apostle John, who leaned on the Lord's bosom at the Last Supper. \Star{} To Him did Christ upon the Cross commit His mother, virgin to virgin.\end{Resp}
\begin{Resp}\Vsign The Lord chose him for his clean maidenhood, and loved him more than all the rest.\end{Resp}
\begin{Resp}\Rsign To him did Christ upon the Cross commit His mother, virgin to virgin.\end{Resp}
\switchcolumn*

\LA
\LessonHead{Lectio II}
\switchcolumn
\EN
\LessonHead{Lesson II}
\switchcolumn*

\LA
\Cite{1 Joannes 1:6-10}
\switchcolumn
\EN
\Cite{1 John 1:6-10}
\switchcolumn*

\LA
\noindent Si dixérimus quóniam societátem habémus cum eo, et in ténebris ambulámus, mentímur, et veritátem non fácimus. Si autem in luce ambulámus sicut et ipse est in luce, societátem habémus ad ínvicem, et sanguis Jesu Christi, Fílii ejus, emúndat nos ab omni peccáto. Si dixérimus quóniam peccátum non habémus, ipsi nos sedúcimus, et véritas in nobis non est. Si confiteámur peccáta nostra, fidélis est et justus, ut remíttat nobis peccáta nostra, et emúndet nos ab omni iniquitáte. Si dixérimus quóniam non peccávimus, mendácem fácimus eum, et verbum ejus non est in nobis.\par
\switchcolumn
\EN
\noindent If we say that we have fellowship with him, and walk in darkness, we lie, and do not the truth. But if we walk in the light, as he also is in the light, we have fellowship one with another, and the blood of Jesus Christ his Son cleanseth us from all sin. If we say that we have no sin, we deceive ourselves, and the truth is not in us. If we confess our sins, he is faithful and just, to forgive us our sins, and to cleanse us from all iniquity. If we say that we have not sinned, we make him a liar, and his word is not in us.\par
\switchcolumn*

\LA
\begin{Resp}\Rsign Hic est discípulus ille, qui testimónium pérhibet de his, et scripsit hæc: \Star{} Et scimus quia verum est testimónium ejus.\end{Resp}
\begin{Resp}\Vsign Fluénta Evangélii de ipso sacro Domínici péctoris fonte potávit.\end{Resp}
\begin{Resp}\Rsign Et scimus quia verum est testimónium ejus.\end{Resp}
\switchcolumn
\EN
\begin{Resp}\Rsign This is the disciple which testifieth of these things, and wrote these things. \Star{} And we know that his testimony is true.\end{Resp}
\begin{Resp}\Vsign He drank in the rivers of the Gospel from the Lord's Breast as from an holy fountain.\end{Resp}
\begin{Resp}\Rsign And we know that his testimony is true.\end{Resp}
\switchcolumn*

\LA
\LessonHead{Lectio III}
\switchcolumn
\EN
\LessonHead{Lesson III}
\switchcolumn*

\LA
\Cite{1 Joannes 2:1-5}
\switchcolumn
\EN
\Cite{1 John 2:1-5}
\switchcolumn*

\LA
\noindent Filíoli mei, hæc scribo vobis, ut non peccétis. Sed et si quis peccáverit, advocátum habémus apud Patrem, Jesum Christum justum: et ipse est propitiátio pro peccátis nostris: non pro nostris autem tantum, sed étiam pro totíus mundi. Et in hoc scimus quóniam cognóvimus eum, si mandáta ejus observémus. Qui dicit se nosse eum, et mandáta ejus non custódit, mendax est, et in hoc véritas non est. Qui autem servat verbum ejus, vere in hoc cáritas Dei perfécta est.\par
\switchcolumn
\EN
\noindent My little children, these things I write to you, that you may not sin. But if any man sin, we have an advocate with the Father, Jesus Christ the just: And he is the propitiation for our sins: and not for ours only, but also for those of the whole world. And by this we know that we have known him, if we keep his commandments. He who saith that he knoweth him, and keepeth not his commandments, is a liar, and the truth is not in him. But he that keepeth his word, in him in very deed the charity of God is perfected.\par
\switchcolumn*

\LA
\begin{Resp}\Rsign Hic est beatíssimus Evangelísta et Apóstolus Joánnes, \Star{} Qui privilégio amóris præcípui, céteris áltius a Dómino méruit honorári.\end{Resp}
\begin{Resp}\Vsign Hic est discípulus ille, quem diligébat Jesus, qui supra pectus Dómini in cena recúbuit.\end{Resp}
\begin{Resp}\Rsign Qui privilégio amóris præcípui, céteris áltius a Dómino méruit honorári.\end{Resp}
\begin{Resp}\Vsign Glória Patri, et Fílio, \Star{} et Spirítui Sancto.\end{Resp}
\begin{Resp}\Rsign Qui privilégio amóris præcípui, céteris áltius a Dómino méruit honorári.\end{Resp}
\switchcolumn
\EN
\begin{Resp}\Rsign This is that most blessed Evangelist and Apostle John. \Star{} Who was found worthy that the Lord should honour him more than all the rest, by a special privilege of love.\end{Resp}
\begin{Resp}\Vsign This is the disciple whom Jesus loved, which also leaned on the Lord's Breast at supper.\end{Resp}
\begin{Resp}\Rsign Who was found worthy that the Lord should honour him more than all the rest, by a special privilege of love.\end{Resp}
\begin{Resp}\Vsign Glory be to the Father, and to the Son, \Star{} and to the Holy Ghost.\end{Resp}
\begin{Resp}\Rsign Who was found worthy that the Lord should honour him more than all the rest, by a special privilege of love.\end{Resp}
\switchcolumn*

\LA
\LessonHead{Lectio IV}
\switchcolumn
\EN
\LessonHead{Lesson IV}
\switchcolumn*

\LA
\LessonTitle{Ex libro sancti Hierónymi Presbýteri de Scriptóribus ecclesiásticis}
\switchcolumn
\EN
\LessonTitle{From the Book on Ecclesiastical writers, written by St. Jerome, Priest (at Bethlehem.)}
\switchcolumn*

\LA
\Source{Cap. 9}
\switchcolumn
\EN
\Source{Ch. 9}
\switchcolumn*

\LA
\noindent Joánnes Apóstolus, quem Jesus amávit plúrimum, fílius Zebedǽi, frater Jacóbi Apóstoli, quem Heródes post passiónem Dómini decollávit, novíssimus ómnium scripsit Evangélium, rogátus ab Asiæ epíscopis, advérsus Cerínthum, aliósque hæréticos, et máxime tunc Ebionitárum dogma consúrgens, qui ásserunt Christum ante Maríam non fuísse: unde et compúlsus est divínam ejus nativitátem edícere.\par
\switchcolumn
\EN
\noindent The Apostle John whom Jesus loved was a son of Zebedee, and brother of the Apostle James, who was beheaded by Herod soon after our Lord suffered. He was the last of the Evangelists to write his Gospel, which he published at the request of the Bishops of Asia, against Cerinthus and other heretics, and particularly against the then spreading doctrine of the Ebionites, who asserted that Christ had had no existence before Mary. It was therefore needful for the Evangelist to declare His Eternal and Divine Generation.\par
\switchcolumn*

\LA
\begin{Resp}\Rsign Qui vícerit, fáciam illum colúmnam in templo meo, dicit Dóminus: \Star{} Et scribam super eum nomen meum, et nomen civitátis novæ Jerúsalem.\end{Resp}
\begin{Resp}\Vsign Vincénti dabo édere de ligno vitæ, quod est in paradíso Dei mei.\end{Resp}
\begin{Resp}\Rsign Et scribam super eum nomen meum, et nomen civitátis novæ Jerúsalem.\end{Resp}
\switchcolumn
\EN
\begin{Resp}\Rsign Him that overcometh will I make a pillar in My temple, saith the Lord, \Star{} And I will write My name upon him, and the name of the city, which is New Jerusalem.\end{Resp}
\begin{Resp}\Vsign To him that overcometh will I give to eat of the tree of life, which is in the midst of the Paradise of My God.\end{Resp}
\begin{Resp}\Rsign And I will write My name upon him, and the name of the city, which is New Jerusalem.\end{Resp}
\switchcolumn*

\LA
\LessonHead{Lectio V}
\switchcolumn
\EN
\LessonHead{Lesson V}
\switchcolumn*

\LA
\noindent Quarto décimo ígitur anno, secúndam post Nerónem persecutiónem movénte Domitiáno, in Patmos ínsulam relegátus, scripsit Apocalýpsim, quam interpretátur Justínus Martyr et Irenǽus. Interfécto autem Domitiáno, et actis ejus ob nímiam crudelitátem a senátu rescíssis, sub Nerva príncipe rédiit Ephesum: ibíque usque ad Trajánum príncipem persevérans, totas Asiæ fundávit rexítque ecclésias: et conféctus sénio, sexagésimo octávo post passiónem Dómini anno mórtuus, juxta eándem urbem sepúltus est.\par
\switchcolumn
\EN
\noindent In the fourteenth year after Nero, Domitian stirred up the second persecution, and John was exiled to the island of Patmos, where he wrote his Apocalypse, which hath been explained by Justin the Martyr and Irenaeus. When Domitian was killed, the Senate annulled all his acts, on account of his savage cruelty, and the Apostle returned to Ephesus, during the reign of Nerva. He remained at Ephesus until the time of Trajan, and founded and governed all the Churches of Asia. There, in an extreme old age, he died, in the sixty-eighth year after the Lord's passion, and was buried near the city.\par
\switchcolumn*

\LA
\begin{Resp}\Rsign Diligébat autem eum Jesus, quóniam speciális prærogatíva castitátis amplióri dilectióne fécerat dignum: \Star{} Quia virgo eléctus ab ipso, virgo in ævum permánsit.\end{Resp}
\begin{Resp}\Vsign In cruce dénique moritúrus, huic Matrem suam vírginem vírgini commendávit.\end{Resp}
\begin{Resp}\Rsign Quia virgo eléctus ab ipso, virgo in ævum permánsit.\end{Resp}
\switchcolumn
\EN
\begin{Resp}\Rsign Jesus loved him, because his singular gift of purity made him more worthy of love. \Star{} He chose him for a virgin unto Himself, and he remaineth a virgin for ever.\end{Resp}
\begin{Resp}\Vsign At the end, when He was dying upon the Cross, to him did He commit His mother, maiden to maiden.\end{Resp}
\begin{Resp}\Rsign He chose him for a virgin unto Himself, and he remaineth a virgin for ever.\end{Resp}
\switchcolumn*

\LA
\LessonHead{Lectio VI}
\switchcolumn
\EN
\LessonHead{Lesson VI}
\switchcolumn*

\LA
\LessonTitle{Ex Commentáriis ejúsdem in Epístolam ad Gálatas}
\switchcolumn
\EN
\LessonTitle{From the Commentary upon the Epistle to the Galatians, by the same author}
\switchcolumn*

\LA
\Source{Lib. 3, Cap. 6}
\switchcolumn
\EN
\Source{Book 3, Ch. 6}
\switchcolumn*

\LA
\noindent Beátus Joánnes Evangelísta, cum Ephesi morarétur usque ad últimam senectútem, et vix inter discipulórum manus ad ecclésiam deferrétur, nec posset in plura vocem verba contéxere; nihil áliud per síngulas solébat proférre colléctas, nisi hoc: Filíoli, dilígite altérutrum. Tandem discípuli et fratres, qui áderant, tǽdio affécti quod éadem semper audírent, dixérunt: Magíster, quare semper hoc lóqueris? Qui respóndit dignam Joánne senténtiam: Quia præcéptum Dómini est; et, si solum fiat, súfficit.\par
\switchcolumn
\EN
\noindent The Blessed Evangelist John lived at Ephesus down to an extreme old age, and, at length, when he was with difficulty carried to the Church, and was not able to exhort the congregation at length, he was used simply to say at each meeting, My little children, love one another. At last the disciples and brethren were weary with hearing these words continually, and asked him, Master, wherefore ever sayest thou this only? Whereto he replied to them, (worthy of John,) It is the commandment of the Lord, and if this only be done, it is enough.\par
\switchcolumn*

\LA
\begin{Resp}\Rsign In médio Ecclésiæ apéruit os ejus, \Star{} Et implévit eum Dóminus spíritu sapiéntiæ et intelléctus.\end{Resp}
\begin{Resp}\Vsign Jucunditátem et exsultatiónem thesaurizávit super eum.\end{Resp}
\begin{Resp}\Rsign Et implévit eum Dóminus spíritu sapiéntiæ et intelléctus.\end{Resp}
\begin{Resp}\Vsign Glória Patri, et Fílio, \Star{} et Spirítui Sancto.\end{Resp}
\begin{Resp}\Rsign Et implévit eum Dóminus spíritu sapiéntiæ et intelléctus.\end{Resp}
\switchcolumn
\EN
\begin{Resp}\Rsign In the midst of the congregation did the Lord open his mouth. \Star{} And filled him with the spirit of wisdom and understanding.\end{Resp}
\begin{Resp}\Vsign He made him rich with joy and gladness.\end{Resp}
\begin{Resp}\Rsign And filled him with the spirit of wisdom and understanding.\end{Resp}
\begin{Resp}\Vsign Glory be to the Father, and to the Son, \Star{} and to the Holy Ghost.\end{Resp}
\begin{Resp}\Rsign And filled him with the spirit of wisdom and understanding.\end{Resp}
\switchcolumn*

\LA
\LessonHead{Lectio VII}
\switchcolumn
\EN
\LessonHead{Lesson VII}
\switchcolumn*

\LA
\LessonTitle{Léctio sancti Evangélii secúndum Joánnem}
\switchcolumn
\EN
\LessonTitle{From the Holy Gospel according to John}
\switchcolumn*

\LA
\Cite{Joannes 21:19-24}
\switchcolumn
\EN
\Cite{John 21:19-24}
\switchcolumn*

\LA
\noindent In illo témpore: Dixit Jesus Petro: Séquere me. Convérsus Petrus, vidit illum discípulum, quem diligébat Jesus, sequéntem. Et réliqua.\par
\switchcolumn
\EN
\noindent In that time Jesus said to Peter: Follow me. Peter turning about, saw that disciple whom Jesus loved following. And so on.\par
\switchcolumn*

\LA
\LessonTitle{Homilía sancti Augustíni Epíscopi}
\switchcolumn
\EN
\LessonTitle{Homily on this passage by St. Augustine, Bishop (of Hippo)}
\switchcolumn*

\LA
\Source{Tract. 124 in Joann., post med.}
\switchcolumn
\EN
\Source{124 on the Tract on John}
\switchcolumn*

\LA
\noindent Duas vitas sibi divínitus prædicátas, et commendátas novit Ecclésia: quarum est una in fide, áltera in spécie; una in témpore peregrinatiónis, áltera in æternitáte mansiónis; una in labóre, áltera in réquie; una in via, áltera in pátria; una in ópere actiónis, áltera in mercéde contemplatiónis. Una declínat a malo, et facit bonum: áltera nullum habet, a quo declínet, malum; et magnum habet, quo fruátur, bonum. Una cum hoste pugnat; áltera sine hoste regnat.\par
\switchcolumn
\EN
\noindent The Church knoweth of two different lives, which God hath revealed and blessed: one is the life of faith, the other the life of knowledge; one the life of this pilgrimage, the other the life of the eternal mansions; one the life of work, the other the life of rest; one the life of the journey, the other the life of home; one the life of action, the other the life of contemplation. The one escheweth evil and doeth good; the other hath no evil to eschew, and only an exceeding good to enjoy. The one striveth with the enemy, the other hath no enemies, and reigneth.\par
\switchcolumn*

\LA
\begin{Resp}\Rsign In illum diem suscípiam te servum meum, et ponam te sicut signáculum in conspéctu meo: \Star{} Quóniam ego elégi te, dicit Dóminus.\end{Resp}
\begin{Resp}\Vsign Esto fidélis usque ad mortem, et dabo tibi corónam vitæ.\end{Resp}
\begin{Resp}\Rsign Quóniam ego elégi te, dicit Dóminus.\end{Resp}
\switchcolumn
\EN
\begin{Resp}\Rsign In that day will I take thee, O My servant, and will make thee as a signet before Me. \Star{} For I have chosen thee, saith the Lord.\end{Resp}
\begin{Resp}\Vsign Be thou faithful unto death, and I will give thee a crown of life.\end{Resp}
\begin{Resp}\Rsign For I have chosen thee, saith the Lord.\end{Resp}
\switchcolumn*

\LA
\LessonHead{Lectio VIII}
\switchcolumn
\EN
\LessonHead{Lesson VIII}
\switchcolumn*

\LA
\noindent Una súbvenit indigénti; áltera ibi est, ubi nullum ínvenit indigéntem. Una aliéna peccáta, ut sua sibi ignoscántur, ignóscit; áltera nec pátitur quod ignóscat, nec facit quod sibi poscat ignósci. Una flagellátur malis, ne extollátur in bonis; áltera tanta plenitúdine grátiæ caret omni malo, ut sine ulla tentatióne supérbiæ cohǽreat summo bono.\par
\switchcolumn
\EN
\noindent The one succoureth the needy; the other is where there are no needy to succour. The one forgiveth them that trespass against it, that its own trespasses may be forgiven; the other neither hath trespasses to forgive nor to be forgiven. The one is chastened with evil, lest it be exalted above measure by good; the other enjoyeth such a fulness of grace that it feeleth no evil, and cleaveth so firmly unto the Highest Good, that it hath no temptation to pride.\par
\switchcolumn*

\LA
\begin{Resp}\Rsign Iste est Joánnes, qui supra pectus Dómini in cena recúbuit: \Star{} Beátus Apóstolus, cui reveláta sunt secréta cæléstia.\end{Resp}
\begin{Resp}\Vsign Fluénta Evangélii de ipso sacro Domínici péctoris fonte potávit.\end{Resp}
\begin{Resp}\Rsign Beátus Apóstolus, cui reveláta sunt secréta cæléstia.\end{Resp}
\begin{Resp}\Vsign Glória Patri, et Fílio, \Star{} et Spirítui Sancto.\end{Resp}
\begin{Resp}\Rsign Beátus Apóstolus, cui reveláta sunt secréta cæléstia.\end{Resp}
\switchcolumn
\EN
\begin{Resp}\Rsign This is that John which leaned on the Lord's Breast at supper \Star{} Even that blessed Apostle unto whom were made known the secret things of heaven.\end{Resp}
\begin{Resp}\Vsign He drank in the rivers of the Gospel from the Lord's Breast, as from an holy fountain.\end{Resp}
\begin{Resp}\Rsign Even that blessed Apostle unto whom were made known the secret things of heaven.\end{Resp}
\begin{Resp}\Vsign Glory be to the Father, and to the Son, \Star{} and to the Holy Ghost.\end{Resp}
\begin{Resp}\Rsign Even that blessed Apostle unto whom were made known the secret things of heaven.\end{Resp}
\switchcolumn*

\LA
\LessonHead{Lectio IX}
\switchcolumn
\EN
\LessonHead{Lesson IX}
\switchcolumn*

\LA
\noindent Ergo una bona est, sed adhuc mísera; áltera mélior, et beáta. Ista significáta est per Apóstolum Petrum, illa per Joánnem. Tota hic ágitur ista usque in hujus sǽculi finem, et illic ínvenit finem; différtur illa complénda post hujus sǽculi finem, sed in futúro sǽculo non habet finem. Ideo dícitur huic: Séquere me. De illo autem: Sic eum volo manére, donec véniam: quid ad te? tu me séquere. Quid enim est hoc? Quantum sápio, quantum cápio, quid est hoc? nisi tu me séquere per imitatiónem perferéndi temporália mala: ille máneat, donec sempitérna vénio redditúrus bona.\par
\switchcolumn
\EN
\noindent Therefore the one is good, but still sorrowful; the other is better and perfectly blessed. And of these two lives there are types, of the one in the Apostle Peter, of the other in John. The one laboureth here even unto the end, and findeth its end hereafter; the other stretcheth out into the hereafter, and in eternity findeth no end. Therefore is it said unto the one, Follow Me; but of the other, If I will that he tarry till I come, what is that to thee? Follow thou Me. What is the meaning of these words? who can know? who can understand? what is it? is it Follow thou Me, imitating Me in the bearing of earthly sorrow; let him tarry till I come again, bringing the everlasting reward?\par
\switchcolumn*

\LA
\TeDeum{Te Deum}
\switchcolumn
\EN
\TeDeum{Te Deum}
\switchcolumn*

\end{paracol}

\par\needspace{10\baselineskip}\addvspace{1.2em}{\centering\small\textsc{Die 28 Decembris} · \textsc{28 December}\par}
\Office{Ss. Innocentium}{Holy Innocents}{Duplex II classis cum Octava simplici}
\addcontentsline{toc}{subsection}{Die 28 Decembris · Ss. Innocentium\enspace\textit{Holy Innocents}}
\begin{paracol}{2}
\LA
\LessonHead{Lectio I}
\switchcolumn
\EN
\LessonHead{Lesson I}
\switchcolumn*

\LA
\LessonTitle{De Jeremía Prophéta}
\switchcolumn
\EN
\LessonTitle{Lesson from the book of Jeremias}
\switchcolumn*

\LA
\Cite{Jer 31:15-17}
\switchcolumn
\EN
\Cite{Jer 31:15-17}
\switchcolumn*

\LA
\noindent Hæc dicit Dóminus: Vox in excélso audíta est lamentatiónis, luctus, et fletus Rachel plorántis fílios suos, et noléntis consolári super eis, quia non sunt. Hæc dicit Dóminus: Quiéscat vox tua a plorátu, et óculi tui a lácrimis: quia est merces óperi tuo, ait Dóminus: et reverténtur de terra inimíci. Et est spes novíssimis tuis, ait Dóminus: et reverténtur fílii ad términos suos.\par
\switchcolumn
\EN
\noindent Thus saith the Lord: A voice was heard on high of lamentation, of mourning, and weeping, of Rachel weeping for her children, and refusing to be comforted for them, because they are not. Thus saith the Lord: Let thy voice cease from weeping, and thy eyes from tears: for there is a reward for thy work, saith the Lord: and they shall return out of the land of the enemy. And here is hope for thy last end, saith the Lord: and the children shall return to their own borders.\par
\switchcolumn*

\LA
\begin{Resp}\Rsign Centum quadragínta quátuor míllia, qui empti sunt de terra, hi sunt qui cum muliéribus non sunt coinquináti: \Star{} Vírgines enim permansérunt, ídeo regnant cum Deo, et Agnus Dei cum illis.\end{Resp}
\begin{Resp}\Vsign Isti sunt qui venérunt ex magna tribulatióne, et lavérunt stolas suas in sánguine Agni.\end{Resp}
\begin{Resp}\Rsign Vírgines enim permansérunt, ídeo regnant cum Deo, et Agnus Dei cum illis.\end{Resp}
\switchcolumn
\EN
\begin{Resp}\Rsign A hundred forty four thousand, which were redeemed from the earth; these are they which were not defiled with women. \Star{} For they remained virgins; therefore are they kings before God, and the Lamb of God is with them.\end{Resp}
\begin{Resp}\Vsign These are they which came out of great tribulation, and have washed their robes in the Blood of the Lamb.\end{Resp}
\begin{Resp}\Rsign For they remained virgins; therefore are they kings before God, and the Lamb of God is with them.\end{Resp}
\switchcolumn*

\LA
\LessonHead{Lectio II}
\switchcolumn
\EN
\LessonHead{Lesson II}
\switchcolumn*

\LA
\Cite{Jer 31:18-20}
\switchcolumn
\EN
\Cite{Jer 31:18-20}
\switchcolumn*

\LA
\noindent Audiens audívi Ephraim transmigrántem: Castigásti me, et erudítus sum, quasi juvénculus indómitus: convérte me, et convértar: quia tu Dóminus Deus meus. Postquam enim convertísti me, egi pœniténtiam: et postquam ostendísti mihi, percússi femur meum. Confúsus sum, et erúbui: quóniam sustínui oppróbrium adolescéntiæ meæ. Si fílius honorábilis mihi Ephraim, si puer delicátus: quia ex quo locútus sum de eo, adhuc recordábor ejus.\par
\switchcolumn
\EN
\noindent Hearing I heard Ephraim when he went into captivity: thou hast chastised me, and I was instructed, as a young bullock unaccustomed to the yoke. Convert me, and I shall be converted, for thou art the Lord my God. For after thou didst convert me, I did penance: and after thou didst shew unto me, I struck my thigh: I am confounded and ashamed, because I have borne the reproach of my youth. Surely Ephraim is an honourable son to me, surely he is a tender child: for since I spoke of him, I will still remember him.\par
\switchcolumn*

\LA
\begin{Resp}\Rsign Sub altáre Dei audívi voces occisórum dicéntium: \Star{} Quare non deféndis sánguinem nostrum? Et accepérunt divínum respónsum: Adhuc sustinéte módicum tempus, donec impleátur númerus fratrum vestrórum.\end{Resp}
\begin{Resp}\Vsign Vidi sub altáre Dei ánimas interfectórum propter verbum Dei, et propter testimónium quod habébant, et clamábant voce magna, dicéntes.\end{Resp}
\begin{Resp}\Rsign Quare non deféndis sánguinem nostrum? Et accepérunt divínum respónsum: Adhuc sustinéte módicum tempus, donec impleátur númerus fratrum vestrórum.\end{Resp}
\switchcolumn
\EN
\begin{Resp}\Rsign I heard under the altar the voices of them that were slain, saying: \Star{} How long dost Thou not avenge our blood? And it was said unto them from God: Rest yet for a little season, until the number of your brethren be fulfilled.\end{Resp}
\begin{Resp}\Vsign I saw under the altar of God the souls of them that were slain for the Word of God, and for the testimony which they held, and they cried with a loud voice, saying:\end{Resp}
\begin{Resp}\Rsign How long dost Thou not avenge our blood? And it was said unto them from God: Rest yet for a little season, until the number of your brethren be fulfilled.\end{Resp}
\switchcolumn*

\LA
\LessonHead{Lectio III}
\switchcolumn
\EN
\LessonHead{Lesson III}
\switchcolumn*

\LA
\Cite{Jer 31:21-23}
\switchcolumn
\EN
\Cite{Jer 31:21-23}
\switchcolumn*

\LA
\noindent Státue tibi spéculam, pone tibi amaritúdines: dírige cor tuum in viam rectam, in qua ambulásti: revértere, virgo Israël, revértere ad civitátes tuas istas. Usquequo delíciis dissólveris, fília vaga? quia creávit Dóminus novum super terram: Fémina circúmdabit virum. Hæc dicit Dóminus exercítuum Deus Israël: Adhuc dicent verbum istud in terra Juda, et in úrbibus ejus, cum convértero captivitátem eórum: Benedícat tibi Dóminus, pulchritúdo justítiæ, mons sanctus.\par
\switchcolumn
\EN
\noindent Set thee up a watchtower, make to thee bitterness: direct thy heart into the right way, wherein thou hast walked: return, O virgin of Israel, return to these thy cities. How long wilt thou be dissolute in deliciousness, O wandering daughter? for the Lord hath created a new thing upon the earth: A woman shall compass a man. Thus saith the Lord of hosts, the God of Israel: As yet shall they say this word in the land of Juda, and in the cities thereof, when I shall bring back their captivity: The Lord bless thee, the beauty of justice, the holy mountain.\par
\switchcolumn*

\LA
\begin{Resp}\Rsign Adoravérunt vivéntem in sǽcula sæculórum, \Star{} Mitténtes corónas suas ante thronum Dómini Dei sui.\end{Resp}
\begin{Resp}\Vsign Et cecidérunt in conspéctu throni in fácies suas: et benedixérunt vivéntem in sǽcula sæculórum.\end{Resp}
\begin{Resp}\Rsign Mitténtes corónas suas ante thronum Dómini Dei sui.\end{Resp}
\begin{Resp}\Vsign Glória Patri, et Fílio, \Star{} et Spirítui Sancto.\end{Resp}
\begin{Resp}\Rsign Mitténtes corónas suas ante thronum Dómini Dei sui.\end{Resp}
\switchcolumn
\EN
\begin{Resp}\Rsign They worshipped Him That liveth for ever and ever. \Star{} And cast their crowns before the throne of the Lord their God.\end{Resp}
\begin{Resp}\Vsign And they fell down upon their faces before the throne, and blessed Him That liveth for ever and ever.\end{Resp}
\begin{Resp}\Rsign And cast their crowns before the throne of the Lord their God.\end{Resp}
\begin{Resp}\Vsign Glory be to the Father, and to the Son, \Star{} and to the Holy Ghost.\end{Resp}
\begin{Resp}\Rsign And cast their crowns before the throne of the Lord their God.\end{Resp}
\switchcolumn*

\LA
\LessonHead{Lectio IV}
\switchcolumn
\EN
\LessonHead{Lesson IV}
\switchcolumn*

\LA
\LessonTitle{Sermo sancti Augustíni Epíscopi}
\switchcolumn
\EN
\LessonTitle{From the Sermons of St. Augustine, Bishop (of Hippo)}
\switchcolumn*

\LA
\Source{Sermo 10 de Sanctis}
\switchcolumn
\EN
\Source{10th on the Saints}
\switchcolumn*

\LA
\noindent Hódie, fratres caríssimi, natálem illórum infántium cólimus, quos ab Heróde crudelíssimo rege interféctos esse, Evangélii textus elóquitur. Et ídeo cum summa exsultatióne gáudeat terra, cæléstium mílitum et tantárum parens fecúnda virtútum. Ecce profánus hostis nunquam beátis párvulis tantum prodésse potuísset obséquio, quantum prófuit ódio. Nam, sicut sacratíssimum præséntis diéi festum maniféstat, quantum in beátos párvulos iníquitas abundávit, tantum in eis grátia benedictiónis refúdit.\par
\switchcolumn
\EN
\noindent Dearly beloved brethren, today we keep the birthday of those children, who, as we are informed by the Gospel, were massacred by the savage King Herod. Therefore let earth rejoice with exceeding joy, for she is the mother of these heavenly soldiers, and of this numerous host. The love of the vile Herod could never have crowned these blessed ones as hath his hatred. For the Church testifieth by this holy solemnity, that whereas iniquity did specially abound against these little saints, so much the more were heavenly blessings poured out upon them.\par
\switchcolumn*

\LA
\begin{Resp}\Rsign Effudérunt sánguinem Sanctórum velut aquam in circúitu Jerúsalem: \Star{} Et non erat qui sepelíret.\end{Resp}
\begin{Resp}\Vsign Posuérunt mortália servórum tuórum escas volatílibus cæli, carnes sanctórum tuórum béstiis terræ.\end{Resp}
\begin{Resp}\Rsign Et non erat qui sepelíret.\end{Resp}
\switchcolumn
\EN
\begin{Resp}\Rsign The blood of thy saints have they shed like water round about Jerusalem. \Star{} And there was none to bury them.\end{Resp}
\begin{Resp}\Vsign The dead bodies of thy servants have they given to be meat unto the fowls of the air, the flesh of thy saints unto the beasts of the earth.\end{Resp}
\begin{Resp}\Rsign And there was none to bury them.\end{Resp}
\switchcolumn*

\LA
\LessonHead{Lectio V}
\switchcolumn
\EN
\LessonHead{Lesson V}
\switchcolumn*

\LA
\noindent Beáta es, o Béthlehem terra Juda, quæ Heródis regis immanitátem in puerórum exstinctióne perpéssa es: quæ sub uno témpore candidátam plebem imbéllis infántiæ Deo offérre meruísti. Digne tamen natálem illórum cólimus, quos beátius ætérnæ vitæ mundus édidit, quam quos maternórum víscerum partus effúdit. Síquidem ante vitæ perpétuæ adépti sunt dignitátem, quam usúram præséntis accéperint.\par
\switchcolumn
\EN
\noindent Blessed art thou, O Bethlehem in the land of Judah, which hast suffered the cruelty of King Herod in the slaughter of thy children; who art found worthy to offer at once to God a whole white-robed army of guileless martyrs! Surely, it is well to keep their birthday, even that blessed birthday which gave them from earth to heaven, more blessed than the day that brought them out of their mother's womb. Scarcely had they entered on the life that now is, when they obtained that glorious life which is to come.\par
\switchcolumn*

\LA
\begin{Resp}\Rsign Isti sunt Sancti, qui passi sunt propter te, Dómine: víndica eos, \Star{} Quia clamant ad te cotídie.\end{Resp}
\begin{Resp}\Vsign Víndica, Dómine, sánguinem Sanctórum tuórum, qui effúsus est.\end{Resp}
\begin{Resp}\Rsign Quia clamant ad te cotídie.\end{Resp}
\switchcolumn
\EN
\begin{Resp}\Rsign These holy ones suffered for thy sake, O Lord take vengeance for them. \Star{} For day by day they cry unto thee.\end{Resp}
\begin{Resp}\Vsign Avenge, O Lord, the blood of thy saints which is shed.\end{Resp}
\begin{Resp}\Rsign For day by day they cry unto thee.\end{Resp}
\switchcolumn*

\LA
\LessonHead{Lectio VI}
\switchcolumn
\EN
\LessonHead{Lesson VI}
\switchcolumn*

\LA
\noindent Aliórum quidem pretiósa mors Mártyrum laudem in confessióne proméruit, horum in consummatióne complácuit; quia incipiéntis vitæ primórdiis, ipse eis occásus inítium glóriæ dedit, qui præséntis términum impósuit. Quos Heródis impíetas lactántes matrum ubéribus abstráxit; qui jure dicúntur Mártyrum flores, quos in médio frígore infidelitátis exórtos, velut primas erumpéntes Ecclésiæ gemmas, quædam persecutiónis pruína decóxit.\par
\switchcolumn
\EN
\noindent We praise the death of other martyrs because it was the crowning act of an undaunted and persistent testimony; but these were crowned at once. He That maketh an end to this present life, gave to them at its very gates that eternal blessedness which we hope for at its close. They whom the wickedness of Herod tore from their mothers' breasts are rightfully called the flowers of martyrdom; hardly had these buds of the Church shown their heads above the soil, in the winter of unbelief, when the frost of persecution nipped them.\par
\switchcolumn*

\LA
\begin{Resp}\Rsign Isti sunt, qui non inquinavérunt vestiménta sua: \Star{} Ambulábunt mecum in albis, quia digni sunt.\end{Resp}
\begin{Resp}\Vsign Hi sunt, qui cum muliéribus non sunt coinquináti: vírgines enim sunt.\end{Resp}
\begin{Resp}\Rsign Ambulábunt mecum in albis, quia digni sunt.\end{Resp}
\begin{Resp}\Vsign Glória Patri, et Fílio, \Star{} et Spirítui Sancto.\end{Resp}
\begin{Resp}\Rsign Ambulábunt mecum in albis, quia digni sunt.\end{Resp}
\switchcolumn
\EN
\begin{Resp}\Rsign These are they which have not defiled their garments. \Star{} They shall walk with Me in white, for they are worthy.\end{Resp}
\begin{Resp}\Vsign These are they which were not defiled with women for they are virgins.\end{Resp}
\begin{Resp}\Rsign They shall walk with Me in white, for they are worthy.\end{Resp}
\begin{Resp}\Vsign Glory be to the Father, and to the Son, \Star{} and to the Holy Ghost.\end{Resp}
\begin{Resp}\Rsign They shall walk with Me in white, for they are worthy.\end{Resp}
\switchcolumn*

\LA
\LessonHead{Lectio VII}
\switchcolumn
\EN
\LessonHead{Lesson VII}
\switchcolumn*

\LA
\LessonTitle{Léctio sancti Evangélii secúndum Matthǽum}
\switchcolumn
\EN
\LessonTitle{From the Holy Gospel according to Matthew}
\switchcolumn*

\LA
\Cite{Matt 2:13-18}
\switchcolumn
\EN
\Cite{Matt 2:13-18}
\switchcolumn*

\LA
\noindent In illo témpore: Angelus Dómini appáruit in somnis Joseph, dicens: Surge, et áccipe Púerum et Matrem ejus, et fuge in Ægýptum, et esto ibi usque dum dicam tibi. Et réliqua.\par
\switchcolumn
\EN
\noindent At that time, an angel of the Lord appeared in sleep to Joseph, saying: Arise, and take the child and his mother, and fly into Egypt: and be there until I shall tell thee. And so on.\par
\switchcolumn*

\LA
\LessonTitle{Homilía sancti Hierónymi Presbýteri}
\switchcolumn
\EN
\LessonTitle{Homily by St. Jerome, Priest (at Bethlehem)}
\switchcolumn*

\LA
\Source{Lib. 1 Comm. in c. 2 Matth. et apud Gloss. ord.}
\switchcolumn
\EN
\Source{Book I Comment, on Matth. II}
\switchcolumn*

\LA
\noindent Quando tulit Púerum et Matrem ejus, ut in Ægýptum tránseat, nocte tulit, et in ténebris: quia noctem ignorántiæ his, a quibus ipse recéssit, relíquit incrédulis. Quando vero revértitur in Judǽam, nec nox, nec ténebræ ponúntur in Evangélio: quia in fine mundi Judǽi fidem, tamquam Christum ab Ægýpto reverténtem suscipiéntes, illuminabúntur.\par
\switchcolumn
\EN
\noindent He took the young Child, and His mother, and fled into Egypt, by night and in darkness; and the darkness of that night was a figure of the darkness of ignorance in which they left the unbelievers from whom they fled. But when they returned into Judaea, we learn not from the Gospel that it was by night, or in darkness; which is an image of that light which will lighten the Jews, when, at the end of the world, they shall receive back the faith, which now lighteneth the Gentiles, even as Judaea received Christ returning from Egypt.\par
\switchcolumn*

\LA
\begin{Resp}\Rsign Cantábant Sancti cánticum novum ante sedem Dei et Agni: \Star{} Et resonábat terra in voces eórum.\end{Resp}
\begin{Resp}\Vsign Hi empti sunt ex homínibus primítiæ Deo et Agno, et in ore ipsórum non est invéntum mendácium.\end{Resp}
\begin{Resp}\Rsign Et resonábat terra in voces eórum.\end{Resp}
\switchcolumn
\EN
\begin{Resp}\Rsign The saints sung a new song before the throne of God and the Lamb. \Star{} And their voices were echoed on earth.\end{Resp}
\begin{Resp}\Vsign These were redeemed from among men, being the first-fruits unto God, and to the Lamb, and in their mouth was found no guile.\end{Resp}
\begin{Resp}\Rsign And their voices were echoed on earth.\end{Resp}
\switchcolumn*

\LA
\LessonHead{Lectio VIII}
\switchcolumn
\EN
\LessonHead{Lesson VIII}
\switchcolumn*

\LA
\noindent Ut adimplerétur quod dictum est a Dómino per prophétam dicéntem: Ex Ægýpto vocávi Fílium meum. Respóndeant, qui Hebræórum volúminum dénegant veritátem, ubi hoc in Septuagínta legátur interprétibus. Quod cum non invénerint, nos eis dicémus, in Osée prophéta scriptum, sicut et exemplária probáre possunt, quæ nuper edídimus.\par
\switchcolumn
\EN
\noindent That it might be fulfilled which was spoken of the Lord by the Prophet, saying Out of Egypt have I called My Son. Those who go about to deny the authority of the Hebrew Scriptures, ask where any such passage is to be found in the Septuagint. But, although they find it not there, I tell them that the fact of its being written in the Prophet Hosea (xi. i) can be proved by the texts which I have lately published.\par
\switchcolumn*

\LA
\begin{Resp}\Rsign Vidi sub altáre Dei ánimas interfectórum propter verbum Dei quod habébant, et clara voce dicébant: \Star{} Víndica, Dómine, sánguinem Sanctórum tuórum, qui effúsus est.\end{Resp}
\begin{Resp}\Vsign Sub throno Dei omnes Sancti clamant.\end{Resp}
\begin{Resp}\Rsign Víndica, Dómine, sánguinem Sanctórum tuórum, qui effúsus est.\end{Resp}
\switchcolumn
\EN
\begin{Resp}\Rsign I saw under the Altar of God the souls of them that were slain for the word of God, which they held, and they cried with a loud voice \Star{} Avenge, O Lord, the blood of thy saints, which is shed.\end{Resp}
\begin{Resp}\Vsign Under the throne of God all the saints cry aloud\end{Resp}
\begin{Resp}\Rsign Avenge, O Lord, the blood of thy saints, which is shed.\end{Resp}
\switchcolumn*

\LA
\LessonHead{Lectio IX}
\switchcolumn
\EN
\LessonHead{Lesson IX}
\switchcolumn*

\LA
\noindent Tunc adimplétum est, quod dictum est per Jeremíam prophétam dicéntem: Vox in Rama audíta est, plorátus et ululátus multus, Rachel plorans fílios suos. De Rachel natus est Bénjamin, in cujus tribu non est Béthlehem. Quǽritur ergo, quómodo Rachel fílios Judæ, id est Béthlehem, quasi suos ploret. Respondébimus bréviter, quia sepúlta sit juxta Béthlehem in Ephrata; et ex matérno corpúsculi hospítio matris nomen accéperit; sive quóniam Juda et Bénjamin duæ tribus junctæ erant, et Heródes præcéperat non solum in Béthlehem intérfici púeros, sed et in ómnibus fínibus ejus.\par
\switchcolumn
\EN
\noindent Then was fulfilled that which was spoken by Jeremias the Prophet, saying; In Rama was there a voice heard, weeping and great mourning; Rachel weeping for her children. The child of Rachel was Benjamin, and Bethlehem is not a town belonging to his tribe. We must therefore seek another reason why Rachel should weep for the children of Judah, to whom Bethlehem belongeth, as for her own. The plain answer is that she is buried at Ephrath close to Bethlehem, and she is called Mother on account of the resting-place of her earthly tabernacle being there. It is possible also that she is called Mother because the tribes of Judah and Benjamin were joined together, and Herod slew not only all the children that were in Bethlehem, but also in all the coasts thereof.\par
\switchcolumn*

\LA
\begin{Resp}\Rsign Isti, qui amícti sunt stolis albis, qui sunt, et unde venérunt? Et dixit mihi: \Star{} Hi sunt, qui venérunt de tribulatióne magna, et lavérunt stolas suas, et dealbavérunt eas in sánguine Agni.\end{Resp}
\begin{Resp}\Vsign Vidi sub altáre Dei ánimas interfectórum propter verbum Dei, et propter testimónium quod habébant.\end{Resp}
\begin{Resp}\Rsign Hi sunt, qui venérunt de tribulatióne magna, et lavérunt stolas suas, et dealbavérunt eas in sánguine Agni.\end{Resp}
\begin{Resp}\Vsign Glória Patri, et Fílio, \Star{} et Spirítui Sancto.\end{Resp}
\begin{Resp}\Rsign Hi sunt, qui venérunt de tribulatióne magna, et lavérunt stolas suas, et dealbavérunt eas in sánguine Agni.\end{Resp}
\switchcolumn
\EN
\begin{Resp}\Rsign What are these which are arrayed in white robes? and whence came they? And he said to me: \Star{} These are they which came out of great tribulation, and have washed their robes, and made them white in the Blood of the Lamb.\end{Resp}
\begin{Resp}\Vsign I saw under the Altar of God the souls of them which were slain for the word of God, and for the testimony which they held.\end{Resp}
\begin{Resp}\Rsign These are they which came out of great tribulation, and have washed their robes, and made them white in the Blood of the Lamb.\end{Resp}
\begin{Resp}\Vsign Glory be to the Father, and to the Son, \Star{} and to the Holy Ghost.\end{Resp}
\begin{Resp}\Rsign These are they which came out of great tribulation, and have washed their robes, and made them white in the Blood of the Lamb.\end{Resp}
\switchcolumn*

\end{paracol}

\par\needspace{10\baselineskip}\addvspace{1.2em}{\centering\small\textsc{Die 29 Decembris} · \textsc{29 December}\par}
\Office{S. Thomæ Cantuariensis Episcopi et Martyris}{St. Thomas of Canterbury, Bishop and Martyr}{Duplex}
\addcontentsline{toc}{subsection}{Die 29 Decembris · S. Thomæ Cantuariensis Episcopi et Martyris\enspace\textit{St. Thomas of Canterbury, Bishop and Martyr}}
\begin{paracol}{2}
\LA
\RefLine{Lectiones I–III: de Scriptura occurrente.}
\switchcolumn
\EN
\RefLine{Lessons I–III: from the Scripture of the day.}
\switchcolumn*

\LA
\LessonHead{Lectio IV}
\switchcolumn
\EN
\LessonHead{Lesson IV}
\switchcolumn*

\LA
\noindent Thomas, Londíni in Anglia natus, Theobáldo succéssit Cantuariénsi epíscopo: et qui ántea in administrándo cancelláriæ múnere præcláre se gésserat, in episcopáli offício fortis et invíctus fuit. Cum enim Henrícus secúndus Angliæ rex, convocátis ad se epíscopis et procéribus regni, leges ferret utilitáti ac dignitáti ecclesiásticæ repugnántes, ádeo constánter óbstitit régiæ cupiditáti, ut neque pollicitatiónibus, neque terróribus de senténtia decédens, próxime conjiciéndus in cárcerem clam recésserit. Inde propínqui ejus omnis ætátis ejécti, amíci, fautóres omnes, iis, quibus per ætátem licéret, jurejurándo adstríctis, univérsos Thomam aditúros, si fortásse miserábili suórum calamitátis aspéctu moverétur, qui a sancto propósito privátis incómmodis deterréri mínime potuísset. Non respéxit carnem aut sánguinem, neque ullus in eo humanitátis sensus offícii constántiam labefactávit.\par
\switchcolumn
\EN
\noindent Thomas was born in London, (in the year of our Lord 11 17,) and succeeded Theobald in the Archbishopric of Canterbury (in 1162). He had previously filled with great distinction the office of Lord Chancellor, and showed an indomitable firmness in his duty as Primate. When Henry II., King of England, in an assembly of the Bishops and great men of his realm, endeavoured to pass laws detrimental to the advantage and dignity of the Church, he opposed himself so steadily to the king's wishes, that, neither promises, nor threats availing to shake him, he was about to be cast into prison, had he not made good his escape in time. The whole of his kinsfolk without regard to age or sex, his friends, and his advisers, were then banished the kingdom, and those who were able, were bound by an oath to make their way to the presence of Thomas, in the hope that though careless of his own sufferings, he might yield at the sight of their misery. But neither flesh and blood, nor the pleadings of natural affection could make him swerve from the line of his pastoral duty.\par
\switchcolumn*

\LA
\begin{Resp}\Rsign Honéstum fecit illum Dóminus, et custodívit eum ab inimícis, et a seductóribus tutávit illum: \Star{} Et dedit illi claritátem ætérnam.\end{Resp}
\begin{Resp}\Vsign Descendítque cum illo in fóveam, et in vínculis non derelíquit eum.\end{Resp}
\begin{Resp}\Rsign Et dedit illi claritátem ætérnam.\end{Resp}
\switchcolumn
\EN
\begin{Resp}\Rsign The Lord made him honourable, and defended him from his enemies, and kept him safe from those that lay in wait for him, \Star{} And gave him perpetual glory.\end{Resp}
\begin{Resp}\Vsign He went down with him into the pit, and left him not in bonds.\end{Resp}
\begin{Resp}\Rsign And gave him perpetual glory.\end{Resp}
\switchcolumn*

\LA
\LessonHead{Lectio V}
\switchcolumn
\EN
\LessonHead{Lesson V}
\switchcolumn*

\LA
\noindent Cóntulit ígitur se ad Alexándrum tértium Pontíficem, a quo benígne accéptus est: et inde proféctus, mónachis Pontiniacénsis monastérii, Cisterciénsis órdinis, ab eódem commendátus. Quod ut cognóvit Henrícus, missis ad convéntum fratrum Cisterciénsium minácibus lítteris, Thomam e Pontiníaco monastério exturbáre conátur. Quare vir sanctus, véritus ne sua causa mali áliquid Cisterciénsis família paterétur, sponte discéssit, et Ludovícum Gálliæ regem ejus invitátu convénit; ubi támdiu fuit, quoad, Pontífice máximo et ipso rege agéntibus, ab exsílio, summa totíus regnis gratulatióne revocátur. Quid dum boni pastóris offícium secúrus exséquitur, ecce calumniatóres ad regem déferunt, eum multa contra regnum et públicam quiétem molíri, ut proptérea sǽpius conquererétur rex, se in suo regno cum uno sacerdóte pacem habére non posse.\par
\switchcolumn
\EN
\noindent Lie betook himself to Pope Alexander III., by whom he was graciously received, and who committed him to the care of the Cistercians at Pontigni. As soon as this came to the knowledge of King Henry, he sent threatening letters to the monks, in order to drive Thomas from this shelter. The saint was unwilling that the Cistercian Order should suffer on his account, and therefore voluntarily withdrew from Pontigni, and accepted the invitation of Lewis VII., King of France, to go to his court. He remained here, until his banishment was recalled at the intercession of the Pope and of the King of France, and he returned to England amid great public joy. He was quietly continuing the work of a faithful shepherd of souls, when certain calumniators denounced him to the king as a plotter against the crown and the public peace. Henry, deceived by these libels, cried out that it was hard that one priest should never let him have quiet in his kingdom.\par
\switchcolumn*

\LA
\begin{Resp}\Rsign Desidérium ánimæ ejus tribuísti ei Dómine, \Star{} Et voluntáte labiórum ejus non fraudásti eum.\end{Resp}
\begin{Resp}\Vsign Quóniam prævenísti eum in benedictiónibus dulcédinis: posuísti in cápite ejus corónam de lápide pretióso.\end{Resp}
\begin{Resp}\Rsign Et voluntáte labiórum ejus non fraudásti eum.\end{Resp}
\switchcolumn
\EN
\begin{Resp}\Rsign Lord, Thou hast given him his heart's desire. \Star{} And hast not withholden the request of his lips.\end{Resp}
\begin{Resp}\Vsign For Thou hast prevented him with the blessings of sweetness: Thou hast set a crown of precious stones upon his head.\end{Resp}
\begin{Resp}\Rsign And hast not withholden the request of his lips.\end{Resp}
\switchcolumn*

\LA
\LessonHead{Lectio VI}
\switchcolumn
\EN
\LessonHead{Lesson VI}
\switchcolumn*

\LA
\noindent Ex qua regis voce nefárii satéllites sperántes gratum se regi factúros, si Thomam e médio tóllerent, clam conveniéntes Cantuáriam, epíscopum in templo vespertínis horis óperam dantem aggrediúntur. Qui cléricis templi áditus præclúdere conántibus accúrrens, óstium apéruit, illis usus verbis ad suos: Non est Dei Ecclésia custodiénda more castrórum; et ego pro Ecclésia Dei libénter mortem subíbo. Tum ad mílites: Vos Dei jussu cavéte, ne cuípiam meórum noceátis. Deínde flexis génibus, Deo, beátæ Maríæ, sancto Dionýsio et réliquis Sanctis ejus ecclésiæ patrónis, Ecclésiam et seípsum comméndans, sacrum caput eádem constántia, qua iniquíssimi regis légibus restíterat, ímpio ferro præcidéndum óbtulit, quarto Kaléndas Januárii, anno Dómini millésimo centésimo septuagésimo primo: cujus cérebro respérsum est totíus templi paviméntum. Quem multis póstea illústrem miráculis idem Alexánder Póntifex rétulit in Sanctórum númerum.\par
\switchcolumn
\EN
\noindent Come wicked servants of the king, hearing his words, and thinking to do him pleasure, betook themselves to Canterbury to rid him of the Archbishop. They entered the cathedral in the evening as Thomas was proceeding to assist at Evensong. The clergy in attendance on him, conscious of the attempt about to be made, wished to bolt the doors. But the saint caused them to be again opened, saying, The Church of God is not to be made a castle of, and for the cause of God's Church I am willing to die. He then said to his murderers, I charge you in the name of the Almighty God to hurt none of my people. With these words he fell on his knees, and commended himself to God, to the Blessed Virgin Mary, to St Denis, and to the other holy Patrons of the Church of Canterbury. He presently offered his sacred head for the stroke of death, and received it from the swords of those wicked men with the same constancy with which he had withstood the commands of the unrighteous king. The murderers pulled out his brains and strewed them all about the floor of the Church. He testified on the 29th day of December, in the (53rd) year of (his own age and of) our Lord 11 70, and, being afterwards honoured with many miracles, was canonised by Pope Alexander IIL (in 1173).\par
\switchcolumn*

\LA
\begin{Resp}\Rsign Stola jucunditátis índuit eum Dóminus: \Star{} Et corónam pulchritúdinis pósuit super caput ejus.\end{Resp}
\begin{Resp}\Vsign Cibávit illum Dóminus pane vitæ et intelléctus: et aqua sapiéntiæ salutáris potávit illum.\end{Resp}
\begin{Resp}\Rsign Et corónam pulchritúdinis pósuit super caput ejus.\end{Resp}
\begin{Resp}\Vsign Glória Patri, et Fílio, \Star{} et Spirítui Sancto.\end{Resp}
\begin{Resp}\Rsign Et corónam pulchritúdinis pósuit super caput ejus.\end{Resp}
\switchcolumn
\EN
\begin{Resp}\Rsign The Lord hath put on him a robe of honour, \Star{} And put about his head a crown of joy.\end{Resp}
\begin{Resp}\Vsign With the bread of life and understanding hath the Lord fed him, and given him the water of health and wisdom to drink.\end{Resp}
\begin{Resp}\Rsign And put about his head a crown of joy.\end{Resp}
\begin{Resp}\Vsign Glory be to the Father, and to the Son, \Star{} and to the Holy Ghost.\end{Resp}
\begin{Resp}\Rsign And put about his head a crown of joy.\end{Resp}
\switchcolumn*

\LA
\LessonHead{Lectio VII}
\switchcolumn
\EN
\LessonHead{Lesson VII}
\switchcolumn*

\LA
\LessonTitle{Léctio sancti Evangélii secúndum Joánnem}
\switchcolumn
\EN
\LessonTitle{From the Holy Gospel according to John}
\switchcolumn*

\LA
\Cite{Joannes 10:11-16}
\switchcolumn
\EN
\Cite{John 10:11-16}
\switchcolumn*

\LA
\noindent In illo témpore: Dixit Jesus pharisǽis: Ego sum pastor bonus. Bonus pastor ánimam suam dat pro óvibus suis. Et réliqua.\par
\switchcolumn
\EN
\noindent In that time Jesus said to the pharisees: I am the good shepherd. The good shepherd giveth his life for his sheep. And so on.\par
\switchcolumn*

\LA
\LessonTitle{Homilía sancti Joánnis Chrysóstomi}
\switchcolumn
\EN
\LessonTitle{Homily by St. John Chrysostom, Patriarch (of Constantinople.)}
\switchcolumn*

\LA
\Source{Homilia 59 in Joánnem}
\switchcolumn
\EN
\Source{59th on John.}
\switchcolumn*

\LA
\noindent Magnum quiddam, dilectíssimi, magnum, inquam, est Ecclésiæ prælátio, et quæ multa índiget sapiéntia et fortitúdine, qualem Christus propósuit: ut ánimam pro óvibus ponámus, et numquam illas deserámus; ut lupo generóse resistámus. Hæc enim inter pastórem et mercenárium est differéntia: alter própriæ, contémptis óvibus; alter, sua contémpta, óvium semper salúti invígilat. Pastóris ergo exémplo demonstráto, deceptóres duos méminit: furem mactántem et rapiéntem oves; et mercenárium permitténtem, neque defendéntem commíssas.\par
\switchcolumn
\EN
\noindent Dearly beloved brethren, the Bishops of the Church hold a great office, an office that needeth much that wisdom and strength whereof Christ hath given us an example. We must learn of Him to lay down our lives for the sheep and never to leave them; and to fight bravely against the wolf. This is the difference between the true shepherd and the hireling. The one leaveth the sheep and seeketh his own safety, but the other recketh not of his own safety, so as he may watch over the sheep. Christ then having given us the pattern of a good shepherd, warneth us against two enemies; first, the thief that cometh not but to kill and to steal, and, secondly, the hireling that standeth by, and defendeth not them that are committed to his charge.\par
\switchcolumn*

\LA
\begin{Resp}\Rsign Coróna áurea super caput ejus, \Star{} Expréssa signo sanctitátis, glória honóris, et opus fortitúdinis.\end{Resp}
\begin{Resp}\Vsign Quóniam prævenísti eum in benedictiónibus dulcédinis, posuísti in cápite ejus corónam de lápide pretióso.\end{Resp}
\begin{Resp}\Rsign Expréssa signo sanctitátis, glória honóris, et opus fortitúdinis.\end{Resp}
\switchcolumn
\EN
\begin{Resp}\Rsign A crown of gold upon his head, \Star{} Wherein is engraved Holiness, an ornament of honour, a costly work.\end{Resp}
\begin{Resp}\Vsign For Thou hast prevented him with the blessings of sweetness, Thou hast set a crown of precious stones upon his head.\end{Resp}
\begin{Resp}\Rsign Wherein is engraved Holiness, an ornament of honour, a costly work.\end{Resp}
\switchcolumn*

\LA
\LessonHead{Lectio VIII}
\switchcolumn
\EN
\LessonHead{Lesson VIII}
\switchcolumn*

\LA
\noindent Quod superióri témpore Ezéchiel his verbis insectátus est: Væ pastóribus Israël: nonne pascébant semetípsos? Nonne greges pascúntur a pastóribus? Sed illi contrárium faciébant, quod máximæ malítiæ genus est, et plurimórum causa malórum. Idcírco inquit: Neque quod abjéctum erat, reducébant: neque quod períerat, quærébant: neque confráctum alligábant: neque infírmum consolidábant: quóniam se, non gregem pascébant. Idem et Paulus áliis verbis signíficat: Omnes quæ sua sunt, quærunt, non quæ Jesu Christi.\par
\switchcolumn
\EN
\noindent Ezechiel hath said of old time, (xxxiv. 2): Woe be to the shepherds of Israel! do they not feed themselves? Should not the shepherds feed the flocks? But they did the contrary, a great wickedness and the root of many evils. Therefore, he saith, they brought not back that which was gone astray neither did they search for that which was lost neither did they bind up that which was broken, nor strengthen that which was sick; for they fed themselves, and not the flock. And Paul hath the same in other words, where he saith, (Phil. ii. 21): All seek their own, not the things which are Jesus Christ's.\par
\switchcolumn*

\LA
\begin{Resp}\Rsign Hic est vere Martyr, qui pro Christi nómine sánguinem suum fudit: \Star{} Qui minas júdicum non tímuit, nec terrénæ dignitátis glóriam quæsívit, sed ad cæléstia regna pervénit.\end{Resp}
\begin{Resp}\Vsign Justum dedúxit Dóminus per vias rectas, et osténdit illi regnum Dei.\end{Resp}
\begin{Resp}\Rsign Qui minas júdicum non tímuit, nec terrénæ dignitátis glóriam quæsívit, sed ad cæléstia regna pervénit.\end{Resp}
\begin{Resp}\Vsign Glória Patri, et Fílio, \Star{} et Spirítui Sancto.\end{Resp}
\begin{Resp}\Rsign Qui minas júdicum non tímuit, nec terrénæ dignitátis glóriam quæsívit, sed ad cæléstia regna pervénit.\end{Resp}
\switchcolumn
\EN
\begin{Resp}\Rsign This is a Martyr indeed, who shed his blood for Christ's Name's sake \Star{} Who feared not for the threats of judges, nor sought to be great with the glory of this world, but pressed on unto the kingdom of heaven.\end{Resp}
\begin{Resp}\Vsign The Lord guided the righteous in right paths, and showed him the kingdom of God.\end{Resp}
\begin{Resp}\Rsign Who feared not for the threats of judges, nor sought to be great with the glory of this world, but pressed on unto the kingdom of heaven.\end{Resp}
\begin{Resp}\Vsign Glory be to the Father, and to the Son, \Star{} and to the Holy Ghost.\end{Resp}
\begin{Resp}\Rsign Who feared not for the threats of judges, nor sought to be great with the glory of this world, but pressed on unto the kingdom of heaven.\end{Resp}
\switchcolumn*

\LA
\LessonHead{Lectio IX}
\switchcolumn
\EN
\LessonHead{Lesson IX}
\switchcolumn*

\LA
\noindent Verum ab utróque se dissidére osténdit Christus: ab illis quidem, qui in aliórum perníciem véniunt, cum dicat, se proptérea venísse, ut vitam habérent, et abundántius habérent: ab his autem, qui oves a lupis rapi negligébant, dicéndo se propter eas ánimam pónere, ne oves perírent. Nam cum Judǽi ipsum interímere vellent, non proptérea déstitit a doctrína, neque credéntes trádidit, sed pérstitit, et pértulit mortem: ídeo frequénter inquit: Ego sum pastor bonus. Quæ cum nullo niti testimónio videréntur, (quod enim póneret ánimam suam, non multo post re monstrátum est; quod autem vitam habérent, et abundántius habérent, eventúrum erat in futúro sǽculo) álterum ab áltero confírmat.\par
\switchcolumn
\EN
\noindent Christ showeth Himself very different from either the thief or the hireling; whereas the thief cometh to destroy, He came that they might have life, and that they might have it more abundantly. The hireling fleeth, but He layeth down His life for the sheep, that the sheep perish not. When then the Jews went about to kill Him, He ceased not to teach He gave not up them that believed in Him, but stood steadfast and died. Wherefore He hath good title often to say, I am the Good Shepherd. It was but a little while, and He showed us how He could lay down His life for the sheep. And if it appeareth not as yet how they have life, and have it more abundantly, (but it shall appear, in the world which is to come,) we may well be persuaded of the truth of the second promise, who have seen the fulfilment of the first.\par
\switchcolumn*

\LA
\TeDeum{Te Deum}
\switchcolumn
\EN
\TeDeum{Te Deum}
\switchcolumn*

\end{paracol}

\par\needspace{10\baselineskip}\addvspace{1.2em}{\centering\small\textsc{Die 31 Decembris} · \textsc{31 December}\par}
\Office{S. Silvestri Papæ et Confessoris}{St. Sylvester, Pope and Confessor}{Duplex}
\addcontentsline{toc}{subsection}{Die 31 Decembris · S. Silvestri Papæ et Confessoris\enspace\textit{St. Sylvester, Pope and Confessor}}
\begin{paracol}{2}
\LA
\RefLine{Lectiones I–III: de Scriptura occurrente.}
\switchcolumn
\EN
\RefLine{Lessons I–III: from the Scripture of the day.}
\switchcolumn*

\LA
\RefLine{Lectiones VII–IX: ex Commúni Summorum Pontificum Confessorum (C4b).}
\switchcolumn
\EN
\RefLine{Lessons VII–IX: from the Common of Popes (C4b).}
\switchcolumn*

\LA
\LessonHead{Lectio IV}
\switchcolumn
\EN
\LessonHead{Lesson IV}
\switchcolumn*

\LA
\noindent Silvéster Románus, patre Rufíno, a prima ætáte óperam dedit Cyríno presbýtero; cujus doctrínam et mores egrégie imitátus est. Hic, sæviénte persecutióne, in Sorácte monte latitáverat; et trigésimum annum agens, présbyter sanctæ Románæ Ecclésiæ a Marcellíno Pontífice creátur. Quo in múnere cum omni laude cléricis áliis antecélleret, in Melchíadis póstea locum succéssit, imperatóre Constantíno, qui ántea pacem Ecclésiæ Christi lege pública déderat. Vix Ecclésiæ gubernácula tractánda suscéperat, Constantínum, jam Crucis signo cǽlitus illustrátum, et de hoste Maxéntio victórem, ad tuéndam propagandámque christiánam religiónem impénse fovit. Quem étiam, uti vetus Ecclésiæ Románæ refert tradítio, Apostolórum imágines recognóscere fecit, sacro baptísmate tinxit, et ab infidelitátis lepra mundávit.\par
\switchcolumn
\EN
\noindent Sylvester was a Roman by birth, and his father's name was Rufinus. He was brought up from a very early age under a Priest named Cyrinus, of whose teaching and example he was a diligent learner. In his thirtieth year he was ordained Priest of the Holy Roman Church by Pope Marcellinus. In the discharge of his duties he became a model for all the clergy, and, after the death of Melchiades, he succeeded him on the Papal throne, (in the year of our Lord 314,) during the reign of Constantine, who had already by public decree proclaimed peace to the Church of Christ. Hardly had he undertaken the government of the Church when he betook himself to stir up the Emperor to protect and propagate the religion of Christ. Constantine was fresh from his victory over his enemy Maxentius, on the Eve whereof the sign of the Cross had been revealed to him limned in light upon the sky; and there was an old story in the Church of Rome that it was Sylvester who caused him to recognise the images of the Apostles, administered to him holy Baptism, and cleansed him from the leprosy of misbelief.\par
\switchcolumn*

\LA
\begin{Resp}\Rsign Invéni David servum meum, óleo sancto meo unxi eum: \Star{} Manus enim mea auxiliábitur ei.\end{Resp}
\begin{Resp}\Vsign Nihil profíciet inimícus in eo, et fílius iniquitátis non nocébit ei.\end{Resp}
\begin{Resp}\Rsign Manus enim mea auxiliábitur ei.\end{Resp}
\switchcolumn
\EN
\begin{Resp}\Rsign I have found David My servant, with My holy oil have I anointed him \Star{} For My hand shall help him.\end{Resp}
\begin{Resp}\Vsign The enemy shall prevail nothing against him, nor the son of wickedness afflict him.\end{Resp}
\begin{Resp}\Rsign For My hand shall help him.\end{Resp}
\switchcolumn*

\LA
\LessonHead{Lectio V}
\switchcolumn
\EN
\LessonHead{Lesson V}
\switchcolumn*

\LA
\noindent Itaque, auctóre Silvéstro, pius imperátor quam christifidélibus públice templa exstruéndi tribúerat facultátem, eam exémplo suo confirmávit. Multas enim eréxit Basílicas, scílicet Lateranénsem Christo Servatóri, Vaticánam sancto Petro, Ostiénsem sancto Paulo, sancto Lauréntio in Agro Veráno, sanctæ Cruci in Atrio Sessoriáno, sancto Petro et Marcellíno ac sanctæ Agnéti in viis Lavicána et Nomentána, et álias; quas sacris imagínibus spléndide exornávit, et, munéribus prædiísque attribútis, magnificentíssime locupletávit. Hoc Pontífice hábitum est Concílium Nicǽnum primum, ubi præsidéntibus ejus legátis, præsénteque Constantíno et trecéntis decem et octo epíscopis, sancta et cathólica fides explicáta est, Arío ejúsdem sectatóribus condemnátis; quod étiam Concílium, peténtibus Pátribus, confirmávit in sýnodo Romæ hábita, ubi íterum Aríus condemnátus est. Multa item decréta fecit Ecclésiæ Dei utília, quæ sub ejus nómine recenséntur: id est, ut a solo epíscopo chrisma conficerétur; ut présbyter chrísmate baptizáti summum liníret vérticem; ut diáconi dalmáticis in ecclésia, et palla linóstima ad lævam uteréntur; ut in líneo tantum velo sacrifícium altáris conficerétur.\par
\switchcolumn
\EN
\noindent The godly Emperor had already granted to Christ's faithful people permission to build public churches, and by the advice of Sylvester he himself set them the example. He built many Basilicas, and magnificently adorned them with holy images, and gifted them with gifts and endowments. Among these there were, besides others, the Church of Christ the Saviour, hard by the Lateran Palace; that of St. Peter, upon the Vatican Mount; that of St. Paul, upon the road to Ostia; that of St. Lawrence, in Verus' field; that of the Holy Cross at the Sessorian hall; that of St. Peter and St. Marcellinus, upon the Lavican Way; and that of St. Agnes, upon the road to Mentana. Under this Pope was held the first Council of Nicea, presided over by the Papal Legates, and in the Presence of Constantine, and three hundred and eighteen Bishops, where the holy and Catholic Faith was declared, and Arius and his followers condemned; which Council was finally confirmed by the Pope, at the request of all the assembled Fathers, in a synod held at Rome, where Arius was again condemned. This Pope issued many useful ordinances for the Church of God. He reserved to Bishops the right of consecrating the Holy Chrism; ordered Priests to anoint with Chrism the heads of the newly baptised; settled the officiating dress of Deacons as a dalmatic and a linen maniple; and forbade the consecration of the Sacrament of the Altar on anything but a linen corporal.\par
\switchcolumn*

\LA
\begin{Resp}\Rsign Pósui adjutórium super poténtem, et exaltávi eléctum de plebe mea: \Star{} Manus enim mea auxiliábitur ei.\end{Resp}
\begin{Resp}\Vsign Invéni David servum meum, óleo sancto meo unxi eum.\end{Resp}
\begin{Resp}\Rsign Manus enim mea auxiliábitur ei.\end{Resp}
\switchcolumn
\EN
\begin{Resp}\Rsign I have laid help upon one that is mighty, and have exalted one chosen out of My people \Star{} For My hand shall help him.\end{Resp}
\begin{Resp}\Vsign I have found David My servant, with My holy oil have I anointed him.\end{Resp}
\begin{Resp}\Rsign For My hand shall help him.\end{Resp}
\switchcolumn*

\LA
\LessonHead{Lectio VI}
\switchcolumn
\EN
\LessonHead{Lesson VI}
\switchcolumn*

\LA
\noindent Idem præscripsísse tráditur tempus ómnibus, qui ordínibus initiáti essent, exercéndi síngulos órdines in Ecclésia, ántequam quisque ad altiórem gradum ascénderet; ut láicus clérico non inférret crimen, ne cléricus apud profánum júdicem causam díceret. Sábbati et Domínici diéi nómine reténto, réliquos hebdómadæ dies feriárum nómine distínctos, ut jam ante in Ecclésia vocári cœperant, appellári vóluit, quo significarétur, quotídie cléricos, abjécta ceterárum rerum cura, uni Deo prorsus vacáre debére. Huic cælésti prudéntiæ, qua ecclésiam administrábat, insígnis vitæ sánctitas, et benígnitas in páuperes perpétuo respóndit. Quo in génere provídit, ut cléricis copiósis egéntes conjúngeret, et sacris virgínibus, quæ ad victum necessária essent, suppeditaréntur. Vixit in pontificátu annos vigínti unum, menses decem, diem unum. Sepúltus est in cœmeterio Priscíllæ, via Salária. Fecit ordinatiónes septem mense Decémbri, quibus creávit presbýteros quadragínta duos, diáconos vigínti quinque, epíscopos per divérsa loca sexagínta quinque.\par
\switchcolumn
\EN
\noindent This Sylvester is likewise said to have ordained that all persons taking Holy Orders should remain awhile in each grade before being promoted to a higher; that laymen should not go to law against the clergy; and that the clergy themselves were not to plead before civil tribunals. He decreed that the first and seventh days of the week should be called respectively the Lord's Day and the Sabbath, and the others, Second Day, Third Day, and so on. In this he confirmed the use of the word Feria for the weekdays, the which use had already begun in the Church. This word signifieth an holiday, and pointeth to the duty of the clergy ever to lay aside all worldly labour, and leave themselves free to do continually the work of the Lord. The heavenly wisdom with which he ruled the Church of God, was joined in him to a singular holiness of life, and an inexhaustible tenderness towards the poor; in which matter he ordained that the wealthy clergy should each relieve a certain number of needy persons; and he also made arrangements for supplying the consecrated virgins with the necessaries of life. He lived as Pope twenty-one years, ten months and one day, and was buried in the cemetery of Priscilla on the Salarian Way, \textasciitilde{}(in the year 335.) He held seven Advent ordinations, and made forty-two Priests, twenty-five Deacons, and sixty-five Bishops of various sees.\par
\switchcolumn*

\LA
\begin{Resp}\Rsign Iste est, qui ante Deum magnas virtútes operátus est, et omnis terra doctrína ejus repléta est: \Star{} Ipse intercédat pro peccátis ómnium populórum.\end{Resp}
\begin{Resp}\Vsign Iste est, qui contémpsit vitam mundi, et pervénit ad cæléstia regna.\end{Resp}
\begin{Resp}\Rsign Ipse intercédat pro peccátis ómnium populórum.\end{Resp}
\begin{Resp}\Vsign Glória Patri, et Fílio, \Star{} et Spirítui Sancto.\end{Resp}
\begin{Resp}\Rsign Ipse intercédat pro peccátis ómnium populórum.\end{Resp}
\switchcolumn
\EN
\begin{Resp}\Rsign This is he which wrought great wonders before God, and the whole earth is full of his teaching \Star{} May he pray for all people, that their sins may be forgiven unto them\end{Resp}
\begin{Resp}\Vsign This is he which loved not his life in this world, and hath attained unto the kingdom of heaven.\end{Resp}
\begin{Resp}\Rsign May he pray for all people, that their sins may be forgiven unto them\end{Resp}
\begin{Resp}\Vsign Glory be to the Father, and to the Son, \Star{} and to the Holy Ghost.\end{Resp}
\begin{Resp}\Rsign May he pray for all people, that their sins may be forgiven unto them\end{Resp}
\switchcolumn*

\end{paracol}

\SeasonTitle{Mensis Januarius}{January}

\addcontentsline{toc}{section}{Mensis Januarius\enspace\textit{January}}

\par\needspace{10\baselineskip}\addvspace{1.2em}{\centering\small\textsc{Die 1 Januarii} · \textsc{1 January}\par}
\Office{In Circumcisione Domini}{In Circumcisione Domini}{Duplex II classis}
\addcontentsline{toc}{subsection}{Die 1 Januarii · In Circumcisione Domini\enspace\textit{In Circumcisione Domini}}
\begin{paracol}{2}
\LA
\RefLine{Lectiones I–III: de Scriptura occurrente.}
\switchcolumn
\EN
\RefLine{Lessons I–III: from the Scripture of the day.}
\switchcolumn*

\LA
\LessonHead{Lectio IV}
\switchcolumn
\EN
\LessonHead{Lesson IV}
\switchcolumn*

\LA
\LessonTitle{Sermo sancti Leónis Papæ}
\switchcolumn
\EN
\LessonTitle{From the Sermons of Pope St. Leo (the Great.)}
\switchcolumn*

\LA
\Source{Sermo 7 de Nativitate Domini}
\switchcolumn
\EN
\Source{7th for Christmas}
\switchcolumn*

\LA
\noindent Festivitátis hodiérnæ, dilectíssimi, verus venerátor est, et pius cultor, qui nec de Incarnatióne Dómini áliquid falsum, nec de Deitáte áliquid sentit indígnum. Paris enim perículi malum est, si illi aut natúræ nostræ véritas, aut Patérnæ glóriæ negátur æquálitas. Cum ergo ad intellegéndum sacraméntum nativitátis Christi, qua de Matre vírgine est ortus, accédimus, abigátur procul terrenárum calígo ratiónum, et ab illuminátæ fídei óculo mundánæ sapiéntiæ fumus abscédat.\par
\switchcolumn
\EN
\noindent Dearly beloved brethren, whosoever will keep truly and honour piously this day's festival, it is necessary for him neither to think falsely of the Lord's Incarnation, nor meanly of the Lord's Godhead. For as there is danger, on the one hand, of denying the truth of Christ's participation of our nature, so is there no less danger, on the other, of doing despite to the equality of His glory with the glory of the Father. Wherefore, when we draw near to understand the mystery of Christ's Birth, wherein He was born of the Virgin Mary, we must leave the clouds of earthly imagination behind and pierce the fog of human wisdom with the eye of enlightened faith.\par
\switchcolumn*

\LA
\begin{Resp}\Rsign Congratulámini mihi, omnes qui dilígitis Dóminum: \Star{} Quia, cum essem párvula, plácui Altíssimo, et de meis viscéribus génui Deum et hóminem.\end{Resp}
\begin{Resp}\Vsign Beátam me dicent omnes generatiónes, quia ancíllam húmilem respéxit Deus.\end{Resp}
\begin{Resp}\Rsign Quia, cum essem párvula, plácui Altíssimo, et de meis viscéribus génui Deum et hóminem.\end{Resp}
\switchcolumn
\EN
\begin{Resp}\Rsign Rejoice with me, all ye that love the Lord. \Star{} For while I was yet little I pleased the Most High, and from my womb have I brought forth God and man.\end{Resp}
\begin{Resp}\Vsign All generations shall call me blessed, for God hath regarded the lowliness of His hand-maiden.\end{Resp}
\begin{Resp}\Rsign For while I was yet little I pleased the Most High, and from my womb have I brought forth God and man.\end{Resp}
\switchcolumn*

\LA
\LessonHead{Lectio V}
\switchcolumn
\EN
\LessonHead{Lesson V}
\switchcolumn*

\LA
\noindent Divína est enim auctóritas, cui crédimus: divína est doctrína, quam séquimur. Quóniam sive legis testificatióni, sive oráculis prophetárum, sive evangélicæ tubæ interiórem admoveámus audítum; verum est, quod Joánnes plenus Spíritu Sancto intónuit: In princípio erat Verbum, et Verbum erat apud Deum, et Deus erat Verbum. Hoc erat in princípio apud Deum. Omnia per ipsum facta sunt, et sine ipso factum est nihil. Et simíliter verum est, quod idem prædicátor adjécit: Verbum caro factum est et habitávit in nobis: et vídimus glóriam ejus, glóriam quasi Unigéniti a Patre.\par
\switchcolumn
\EN
\noindent The authority on which we believe is the authority of God Himself; the teaching which we follow is the teaching of God Himself. Wherefore whether we lend the ear of our mind to the testimony of the Law, or to the revelations of the Prophets, or to the full pealing of the Gospel trumpet, that is true, which John the Son of Thunder, uttered, when he was filled with the Holy Ghost, and said: In the beginning was the Word, and the Word was with God, and the Word was God. The Same was in the beginning with God. All things were made by Him, and without Him was not anything made. True also is his witness when he saith: The Word was made Flesh, and dwelt among us, and we beheld His glory, the glory as of the Only-begotten of the Father. (John i. 1-3, 14.)\par
\switchcolumn*

\LA
\begin{Resp}\Rsign Confirmátum est cor Vírginis, in quo divína mystéria, Angelo nuntiánte, concépit: tunc speciósum forma præ fíliis hóminum castis suscépit viscéribus: \Star{} Et benedícta in ætérnum, Deum nobis prótulit et hóminem.\end{Resp}
\begin{Resp}\Vsign Domus pudíci péctoris templum repénte fit Dei: intácta nésciens virum, verbo concépit Fílium.\end{Resp}
\begin{Resp}\Rsign Et benedícta in ætérnum, Deum nobis prótulit et hóminem.\end{Resp}
\switchcolumn
\EN
\begin{Resp}\Rsign The heart of the Virgin was fixed, when the Angel declared unto her the mystery of God, and she conceived, then did she receive in her pure womb Him That is fairer than the children of men. \Star{} And, she that is blessed for ever, brought forth for us God and man.\end{Resp}
\begin{Resp}\Vsign Soon rises, in that modest shrine, The Temple of the Lord Divine The stainless and unwedded one, Within her womb conceived the Son.\end{Resp}
\begin{Resp}\Rsign And, she that is blessed for ever, brought forth for us God and man.\end{Resp}
\switchcolumn*

\LA
\LessonHead{Lectio VI}
\switchcolumn
\EN
\LessonHead{Lesson VI}
\switchcolumn*

\LA
\noindent In utráque ergo natúra idem est Dei Fílius, nostra suscípiens, et própria non amíttens: in hómine hóminem rénovans, in se incommutábilis persevérans. Déitas enim, quæ illi cum Patre commúnis est, nullum detriméntum omnipoténtiæ súbiit, nec Dei formam servi forma violávit: quia summa et sempitérna esséntia, quæ se ad humáni géneris inclinávit salútem, nos quidem in suam glóriam tránstulit, sed quod erat, esse non déstitit. Unde cum Unigénitus Dei minórem se Patre confitétur, cui se dicit æquálem, veritátem in se formæ utriúsque demónstrat: ut et humánam probet imparílitas, et divínam decláret æquálitas.\par
\switchcolumn
\EN
\noindent The Person of the Son of God therefore remaineth unchanged and one, though He have two natures, keeping His own, and taking ours. He appeareth as man to be the restorer of men, but abideth all the while in His immutable Godhead. That Godhead which He shareth with the Father was not a whit the less Almighty, nor did the form of a servant touch the form of God to derogate from it. The Most High and Everlasting Being, bending down for man's salvation, took the Manhood into His glory; He ceased not to be That which He is from everlasting. Hence we see the Only-begotten Son of God in one place confessing that the Father is greater than He, (John xiv. 28,) and in another declaring that He and the Father are One, (x. 30.) This is an evident proof of the distinction of His two natures, and the unity of His Person; for He is inferior to the Father as touching His Manhood, and yet equal to the Father as touching His Godhead, and yet, though He be God and Man, He is not two, but One Christ.\par
\switchcolumn*

\LA
\begin{Resp}\Rsign Benedícta et venerábilis es, Virgo María, quæ sine tactu pudóris invénta es mater Salvatóris: \Star{} Jacébat in præsépio, et fulgébat in cælo.\end{Resp}
\begin{Resp}\Vsign Dómine, audívi audítum tuum, et tímui: considerávi ópera tua, et expávi: in médio duórum animálium.\end{Resp}
\begin{Resp}\Rsign Jacébat in præsépio, et fulgébat in cælo.\end{Resp}
\begin{Resp}\Vsign Glória Patri, et Fílio, \Star{} et Spirítui Sancto.\end{Resp}
\begin{Resp}\Rsign Jacébat in præsépio, et fulgébat in cælo.\end{Resp}
\switchcolumn
\EN
\begin{Resp}\Rsign Blessed and worshipful art thou, O Virgin Mary; from thee, still maiden undefiled, the Saviour came a little Child. \Star{} He Whose glory filled the heavens lay in a manger.\end{Resp}
\begin{Resp}\Vsign O Lord, I have heard thy speech and was afraid; I considered thy works and trembled. O Thou That dwellest between the two living creatures!\end{Resp}
\begin{Resp}\Rsign He Whose glory filled the heavens lay in a manger.\end{Resp}
\begin{Resp}\Vsign Glory be to the Father, and to the Son, \Star{} and to the Holy Ghost.\end{Resp}
\begin{Resp}\Rsign He Whose glory filled the heavens lay in a manger.\end{Resp}
\switchcolumn*

\LA
\LessonHead{Lectio VII}
\switchcolumn
\EN
\LessonHead{Lesson VII}
\switchcolumn*

\LA
\LessonTitle{Léctio sancti Evangélii secúndum Lucam}
\switchcolumn
\EN
\LessonTitle{From the Holy Gospel according to Luke}
\switchcolumn*

\LA
\Cite{Luc 2:21}
\switchcolumn
\EN
\Cite{Luke 2:21}
\switchcolumn*

\LA
\noindent In illo témpore: Postquam consummáti sunt dies octo, ut circumciderétur Puer: vocátum est nomen ejus Jesus. Et réliqua.\par
\switchcolumn
\EN
\noindent In that time, after eight days were accomplished, that the child should be circumcised, his name was called Jesus. And so on.\par
\switchcolumn*

\LA
\LessonTitle{Homilía sancti Ambrósii Epíscopi}
\switchcolumn
\EN
\LessonTitle{Homily by St. Ambrose, Bishop (of Milan.)}
\switchcolumn*

\LA
\Source{Liber 2 in cap. 2 Lucæ, circa medium}
\switchcolumn
\EN
\Source{Book ii on Luke ii}
\switchcolumn*

\LA
\noindent Circumcíditur ítaque Puer. Quis est iste puer, nisi ille, de quo dictum est: Puer natus est nobis, Fílius datus est nobis? Factus est enim sub lege, ut eos, qui sub lege essent, lucrifáceret. Ut sísterent eum Dómino. Quid sit autem in Jerúsalem sisti Dómino, dícerem, nisi in Isaíæ comméntis ante dixíssem. Circumcísus enim vítiis, Domínico dignus judicátur obtútu; quia, Oculi Dómini super justos. Vides omnem legis véteris sériem fuísse typum futúri: nam et circumcísio purgatiónem signíficat delictórum.\par
\switchcolumn
\EN
\noindent So the Child is circumcised. This is the Child of Whom it is said: Unto us a Child is born, unto us a Son is given, (Isa. ix. 6.) Made under the law to redeem them that were under the law. (Gal. iv. 4.) To present Him to the Lord, \textasciitilde{}(22.) In my Commentary on Isaiah I have already explained what is meant by being presented to the Lord in Jerusalem, and therefore I will not enter into the subject again. He that is circumcised in heart gaineth the protection of God, for the eyes of the Lord are upon the righteous. (Ps. xxxiii. 16.) Ye will see that as all the ceremonies of the old law were types of realities in the new, so the circumcision of the body signified the cleansing of the heart from the guilt of sin.\par
\switchcolumn*

\LA
\begin{Resp}\Rsign Sancta et immaculáta virgínitas, quibus te láudibus éfferam, néscio: \Star{} Quia quem cæli cápere non póterant, tuo grémio contulísti.\end{Resp}
\begin{Resp}\Vsign Benedícta tu in muliéribus, et benedíctus fructus ventris tui.\end{Resp}
\begin{Resp}\Rsign Quia quem cæli cápere non póterant, tuo grémio contulísti.\end{Resp}
\switchcolumn
\EN
\begin{Resp}\Rsign O Mary, how holy and how spotless is thy virginity! I am too dull to praise thee! \Star{} For thou hast borne in thy breast Him Whom the heavens cannot contain.\end{Resp}
\begin{Resp}\Vsign Blessed art thou among women, and blessed is the fruit of thy womb.\end{Resp}
\begin{Resp}\Rsign For thou hast borne in thy breast Him Whom the heavens cannot contain.\end{Resp}
\switchcolumn*

\LA
\LessonHead{Lectio VIII}
\switchcolumn
\EN
\LessonHead{Lesson VIII}
\switchcolumn*

\LA
\noindent Sed quóniam prona quadam cupiditáte peccándi, humánæ carnis et mentis fragílitas inextricabílibus vítiis implicátur, eo per octávum circumcisiónis diem culpæ totíus futúra purgátio, resurrectiónis præfigurabátur ætáte. Hoc est enim illud: Quia omne masculínum adapériens vulvam, sanctum Dómino vocábitur. Verbis enim legis promittebátur Vírginis partus. Et vere sanctus, quia immaculátus. Dénique ipsum esse, qui lege signétur, in eúndem modum ab Angelo repetíta verba declárant: Quod nascétur, inquit, Sanctum, vocábitur Fílius Dei.\par
\switchcolumn
\EN
\noindent But since the body and mind of man remain yet infected with a proneness' to sin, the circumcision of the eighth day is also a type of that complete cleansing from sin which we shall have at the resurrection. This ceremony was also performed in obedience to the commandment of God: Every male that openeth the womb shall be called holy unto the Lord. These words were written with especial reference to the delivery of the Blessed Virgin. Truly He That opened her womb was holy, for He was altogether without spot, and we may gather that the law was written specially for Him from the words of the Angel: That Holy Thing Which shall be born of thee, shall be called the Son of God.\par
\switchcolumn*

\LA
\begin{Resp}\Rsign Nésciens Mater Virgo virum, péperit sine dolóre: \Star{} Salvatórem sæculórum, ipsum Regem Angelórum, sola Virgo lactábat úbere de cælo pleno.\end{Resp}
\begin{Resp}\Vsign Domus pudíci péctoris templum repénte fit Dei: intácta nésciens virum, verbo concépit Fílium.\end{Resp}
\begin{Resp}\Rsign Salvatórem sæculórum, ipsum Regem Angelórum, sola Virgo lactábat úbere de cælo pleno.\end{Resp}
\begin{Resp}\Vsign Glória Patri, et Fílio, \Star{} et Spirítui Sancto.\end{Resp}
\begin{Resp}\Rsign Salvatórem sæculórum, ipsum Regem Angelórum, sola Virgo lactábat úbere de cælo pleno.\end{Resp}
\switchcolumn
\EN
\begin{Resp}\Rsign The Virgin-Mother that knew not a man, bore, but travailed not. \Star{} She fed the Saviour of the world, The King of Angel hosts above, Jesus, our Redeemer blest, From the fountain of her breast.\end{Resp}
\begin{Resp}\Vsign Soon rises in that modest shrine, The Temple of the Lord Divine; The stainless and unwedded one, Within her womb conceived the Son.\end{Resp}
\begin{Resp}\Rsign She fed the Saviour of the world, The King of Angel hosts above, Jesus, our Redeemer blest, From the fountain of her breast.\end{Resp}
\begin{Resp}\Vsign Glory be to the Father, and to the Son, \Star{} and to the Holy Ghost.\end{Resp}
\begin{Resp}\Rsign She fed the Saviour of the world, The King of Angel hosts above, Jesus, our Redeemer blest, From the fountain of her breast.\end{Resp}
\switchcolumn*

\LA
\LessonHead{Lectio IX}
\switchcolumn
\EN
\LessonHead{Lesson IX}
\switchcolumn*

\LA
\noindent Solus enim per ómnia ex natis de fémina sanctus Dóminus Jesus, qui terrénæ contágia corruptélæ, immaculáti partus novitáte non sénserit, et cælésti majestáte depúlerit. Nam si lítteram sequámur: quómodo sanctus omnis másculus, cum multos sceleratíssimos fuísse non láteat? Numquid sanctus Achab? numquid sancti pseudoprophétæ, quos ad Elíæ preces ultor cæléstis injúriæ ignis absúmpsit? Sed ille sanctus, quem in figúra futúri mystérii pia legis divínæ præscrípta signábant; eo quod solus sanctæ Ecclésiæ vírginis ad generándos pópulos Dei, immaculátæ fecunditátis aperíret genitále secrétum.\par
\switchcolumn
\EN
\noindent Among all that are born of women the Lord Jesus Christ stood alone in holiness. Fresh from His immaculate Birth, He felt no contagion from human corruption, and His heavenly Majesty drove it away. If we are to follow the letter and say that every male that openeth the womb is holy, how shall we explain that so many have been unrighteous? Was Ahab holy? Were the false prophets holy? Were they holy on whom Elijah justly called down fire from heaven? But He to Whom the sacred commandment of the law of God is mystically directed is the Holy One of Israel; Who also alone hath opened the secret womb of His holy Virgin-bride the Church, filling her with a sinless fruitfulness to give birth to Christian souls.\par
\switchcolumn*

\LA
\TeDeum{Te Deum}
\switchcolumn
\EN
\TeDeum{Te Deum}
\switchcolumn*

\end{paracol}

\par\needspace{10\baselineskip}\addvspace{1.2em}{\centering\small\textsc{Die 2 Januarii} · \textsc{2 January}\par}
\Office{In Octava S. Stephani Protomartyris}{Octave of St. Stephen}{Simplex}
\addcontentsline{toc}{subsection}{Die 2 Januarii · In Octava S. Stephani Protomartyris\enspace\textit{Octave of St. Stephen}}
\begin{paracol}{2}
\LA
\RefLine{Lectiones I–II: de Scriptura occurrente.}
\switchcolumn
\EN
\RefLine{Lessons I–II: from the Scripture of the day.}
\switchcolumn*

\LA
\LessonHead{Lectio III (pro commemoratione)}
\switchcolumn
\EN
\LessonHead{Lesson III (when commemorated)}
\switchcolumn*

\LA
\LessonTitle{Sermo sancti Augustíni Epíscopi}
\switchcolumn
\EN
\LessonTitle{From the Sermons of St. Augustine, Bishop (of Hippo.)}
\switchcolumn*

\LA
\Source{Sermo 2 de S. Stephano}
\switchcolumn
\EN
\Source{2nd on St. Stephen.}
\switchcolumn*

\LA
\noindent Christus caput Mártyrum prior passus est pro nobis, relínquens vobis exémplum, ut sequámini vestígia ejus. Cujus passiónis vestígia prosecútus beatíssimus Stéphanus, confiténdo Christum lapidátus a Judǽis, corónam méruit tamquam suo sibi nómine pósitam. Stéphanus enim Græce, Latíne coróna appellátur. Jam coróna nomen habébat, et ídeo palmam martýrii suo nómine præferébat. Qui cum lapidarétur, non solum non exspectábat de persecutóribus reportáre vindíctam, sed eis pótius a Deo véniam postulábat.\par
\switchcolumn
\EN
\noindent Christ, the Captain of the Martyrs, hath first suffered for us, leaving us an example that we should follow His steps, (i Pet. ii. 21.) And truly, Blessed Stephen followed them, when, having confessed Christ, he was stoned to death by the Jews, and obtained the crown which his name had foreshown. For the meaning of the Greek name Stephanos is a crown. Already he had a crown for his name, a foreshadowing of the martyr's palm which he beareth in heaven. Then they stoned him he did not rejoice at the thought that God would take vengeance on his persecutors. On the contrary, he prayed that they might be forgiven.\par
\switchcolumn*

\LA
\TeDeum{Te Deum}
\switchcolumn
\EN
\TeDeum{Te Deum}
\switchcolumn*

\LA
\LessonHead{Lectio IV}
\switchcolumn
\EN
\LessonHead{Lesson IV}
\switchcolumn*

\LA
\LessonTitle{Commemoratio: In Octava S. Stephani Protomartyris}
\switchcolumn
\EN
\LessonTitle{Commemoration of: Octave of St. Stephen}
\switchcolumn*

\LA
\Source{Serm. 2. de S. Stephano.}
\switchcolumn
\EN
\Source{2nd on St. Stephen.}
\switchcolumn*

\LA
\noindent Sermo sancti Augustíni Epíscopi.\par
\switchcolumn
\EN
\noindent From the Sermons of St. Augustine, Bishop (of Hippo.)\par
\switchcolumn*

\LA
Post hestérnum festivíssimum diem, quo Salvatóris nostri Christi nobis Natívitas illúxit, étiam hodiérnus dies beáti Mártyris Stéphani coróna illustrátur. Mártyris illíus mérita nulla pars orbis ignórat: passus enim est in ipso princípio Ecclésiæ, id est, in ipsa urbe Jerosólyma. Ibi enim Diáconus ministrávit: et in ipso juventútis flore decórem ætátis suæ sánguine purpurávit. Pássio ejus insígnis est, multúmque mirábilis. Hanc modo de libro Actuum Apostolórum, cum legerétur, non solum audívimus, sed étiam óculis spectávimus.\par
\switchcolumn
\EN
Even after the glory of yesterday, bright with the splendour of Christ our Saviour's Birth, this day findeth itself an illumination of its own from the crown of the blessed Martyr Stephen. The whole earth knoweth how manfully he fought and conquered for he suffered at the very fountain-head of the Church, that is to say, in Jerusalem. It was in the Church there that he ministered as a Deacon and in the youthful springtime of life dyed with his blood the lily of his purity. His Passion is very glorious, and many ways wonderful, and when we read it in the Acts of the Apostles, we seem rather to see than to hear.\par
\switchcolumn*

\LA
Christus ergo caput Mártyrum prior passus est pro nobis, relínquens vobis exémplum, ut sequámini vestígia ejus. Cujus passiónis vestígia prosecútus beatíssimus Stéphanus, confiténdo Christum lapidátus a Judǽis, corónam méruit tamquam suo sibi nómine pósitam. Stéphanus enim Græce, Latíne coróna appellátur. Jam coróna nomen habébat, et ídeo palmam martýrii suo nómine præferébat.\par
\switchcolumn
\EN
Christ, the Captain of the Martyrs, hath first suffered for us, leaving us an example that we should follow His steps, (i Pet. ii. 21.) And truly, Blessed Stephen followed them, when, having confessed Christ, he was stoned to death by the Jews, and obtained the crown which his name had foreshown. For the meaning of the Greek name Stephanos is a crown. Already he had a crown for his name, a foreshadowing of the martyr's palm which he beareth in heaven.\par
\switchcolumn*

\LA
Qui cum lapidarétur, non solum non exspectábat de persecutóribus reportáre vindíctam, sed eis pótius a Deo véniam postulábat. Memínerat enim dixísse Dóminum: Mihi vindíctam, ego retríbuam, dicit Dóminus. Et íterum: Ne díxeris, Ulcíscar me de inimícis meis; sed exspécta Dóminum, ut tibi auxílio sit. Exspectáre nos jubet Dóminus Deus, ut in die futúræ retributiónis cum sanctis Martýribus vindicémur.\par
\switchcolumn
\EN
Then they stoned him he did not rejoice at the thought that God would take vengeance on his persecutors. On the contrary, he prayed that they might be forgiven. For he remembered the word of the Lord, that saith Vengeance belongeth unto Me, I will recompense, saith the Lord, (Heb. x. 30,) and again: Say not thou; I will recompense evil to mine enemies, but wait on the Lord, and He shall save thee. (Prov. xx. 22.) The Lord God biddeth us also be patient, knowing that in the great day of retribution, we, as well as His holy martyrs, shall be righted.\par
\switchcolumn*

\LA
\TeDeum{Te Deum}
\switchcolumn
\EN
\TeDeum{Te Deum}
\switchcolumn*

\end{paracol}

\par\needspace{10\baselineskip}\addvspace{1.2em}{\centering\small\textsc{Die 3 Januarii} · \textsc{3 January}\par}
\Office{In Octava S. Joannis Apostoli et Evangelistæ}{Octave Day of St. John the Evangelist}{Simplex}
\addcontentsline{toc}{subsection}{Die 3 Januarii · In Octava S. Joannis Apostoli et Evangelistæ\enspace\textit{Octave Day of St. John the Evangelist}}
\begin{paracol}{2}
\LA
\RefLine{Lectiones I–II: de Scriptura occurrente.}
\switchcolumn
\EN
\RefLine{Lessons I–II: from the Scripture of the day.}
\switchcolumn*

\LA
\LessonHead{Lectio IV}
\switchcolumn
\EN
\LessonHead{Lesson IV}
\switchcolumn*

\LA
\LessonTitle{Ex Tractatu sancti Augustíni Epíscopi in Joánnem}
\switchcolumn
\EN
\LessonTitle{The Lesson is taken from a treatise of St. Augustine, Bishop (of Hippo.)}
\switchcolumn*

\LA
\Source{Tractatus 36}
\switchcolumn
\EN
\Source{on John. (36.)}
\switchcolumn*

\LA
\noindent In quátuor Evangéliis, vel pótius quátuor libris uníus Evangélii, sanctus Joánnes Apóstolus non immérito secúndum intelligéntiam spiritálem áquilæ comparátus, áltius multóque sublímius áliis tribus eréxit prædicatiónem suam: et in ejus erectióne étiam corda nostra érigi vóluit. Nam céteri tres Evangelístæ tamquam cum hómine Dómino in terra ámbulant, et de divinitáte ejus pauca dixérunt: istum autem quasi pigúerit in terra ambuláre, sicut ipso exórdio sui sermónis intónuit, eréxit se non solum super terram, et super omnem ámbitum áëris et cæli, sed super omnem étiam exércitum Angelórum, omnémque constitutiónem invisibílium Potestátum: et pervénit ad eum, per quem facta sunt ómnia, dicéndo: In princípio erat Verbum, et Verbum erat apud Deum, et Deus erat Verbum.\par
\switchcolumn
\EN
\noindent Of the Four Evangelists, or, rather, the Four Writers of the 'one Evangel, the holy Apostle John hath not unworthily been compared by spiritual writers to an eagle, because of the lofty and glorious flight of his teaching, soaring above the other three; a flight that raiseth not himself alone, but also the hearts of all, whosoever will hear him. The other three writers walk with the Lord upon earth, as with a man, and enlarge little upon His Godhead; but John, as though it had wearied him to walk upon earth, in the very first words of his writing, riseth not above the earth only, or above the firmament, and the heavens, but above every angel, and above every power of things unseen, and flieth directly to Him by Whom all things were made, saying In the beginning was the Word, and the Word was with God, and the Word was God.\par
\switchcolumn*

\LA
\TeDeum{Te Deum}
\switchcolumn
\EN
\TeDeum{Te Deum}
\switchcolumn*

\LA
\LessonHead{Lectio IV}
\switchcolumn
\EN
\LessonHead{Lesson IV}
\switchcolumn*

\LA
\LessonTitle{Commemoratio: In Octava S. Joannis Apostoli et Evangelistæ}
\switchcolumn
\EN
\LessonTitle{Commemoration of: Octave Day of St. John the Evangelist}
\switchcolumn*

\LA
\Source{Tractatus 36}
\switchcolumn
\EN
\Source{on John. (36.)}
\switchcolumn*

\LA
\noindent Ex Tractatu sancti Augustíni Epíscopi in Joánnem\par
\switchcolumn
\EN
\noindent The Lesson is taken from a treatise of St. Augustine, Bishop (of Hippo.)\par
\switchcolumn*

\LA
In quátuor Evangéliis, vel pótius quátuor libris uníus Evangélii, sanctus Joánnes Apóstolus non immérito secúndum intelligéntiam spiritálem áquilæ comparátus, áltius multóque sublímius áliis tribus eréxit prædicatiónem suam: et in ejus erectióne étiam corda nostra érigi vóluit. Nam céteri tres Evangelístæ tamquam cum hómine Dómino in terra ámbulant, et de divinitáte ejus pauca dixérunt: istum autem quasi pigúerit in terra ambuláre, sicut ipso exórdio sui sermónis intónuit, eréxit se non solum super terram, et super omnem ámbitum áëris et cæli, sed super omnem étiam exércitum Angelórum, omnémque constitutiónem invisibílium Potestátum: et pervénit ad eum, per quem facta sunt ómnia, dicéndo: In princípio erat Verbum, et Verbum erat apud Deum, et Deus erat Verbum.\par
\switchcolumn
\EN
Of the Four Evangelists, or, rather, the Four Writers of the 'one Evangel, the holy Apostle John hath not unworthily been compared by spiritual writers to an eagle, because of the lofty and glorious flight of his teaching, soaring above the other three; a flight that raiseth not himself alone, but also the hearts of all, whosoever will hear him. The other three writers walk with the Lord upon earth, as with a man, and enlarge little upon His Godhead; but John, as though it had wearied him to walk upon earth, in the very first words of his writing, riseth not above the earth only, or above the firmament, and the heavens, but above every angel, and above every power of things unseen, and flieth directly to Him by Whom all things were made, saying In the beginning was the Word, and the Word was with God, and the Word was God.\par
\switchcolumn*

\LA
Huic tantæ sublimitáti princípii étiam cétera cóngrua prædicávit, et de Dómini divinitáte, quómodo nullus álius, est locútus: hoc ructábat, quod bíberat. Non enim sine causa de illo in isto ipso Evangélio narrátur, quia et in convívio super pectus Dómini discumbébat. De illo ergo péctore in secréto bibébat, sed quod in secréto bibit, in manifésto eructávit: ut pervéniat ad omnes gentes non solum incarnátio Fílii Dei, pássio, et resurréctio, sed étiam quod erat ante incarnatiónem únicus Patri, Verbum Patris, coætérnus generánti, æquális ei a quo missus est.\par
\switchcolumn
\EN
Then he goeth on worthily of such a beginning, and speaketh of the Lord's Godhead as none hath ever spoken, uttering freely the things which he had heard. It is not without cause that it is told of him in that Gospel how he lay on Jesus' Breast at supper. Truly he drank secretly from that Breast, and what he drank secretly he hath uttered openly, that all men may know not only how the Son of God became man, suffered, and rose again for us, but likewise how He was with the Father before He took flesh, the Onlybegotten Son, the Word of the Father, co-eternal with Him that begat Him, equal to Him That sent Him.\par
\switchcolumn*

\LA
Aquila ipse est Joánnes, sublímium prædicátor, et lucis intérnæ atque ætérnæ fixis óculis contemplátor. Dicúntur enim et pulli aquilárum a paréntibus sic probári, patris scílicet ungue suspéndi, et rádiis solis oppóni. Qui firme contemplátus fúerit, fílius agnóscitur: si ácie palpitáverit, tamquam adulterínus ab ungue dimíttitur. Jam ergo vidéte, quam sublímia loqui débuit, qui est áquilæ comparátus: et tamen étiam nos humi repéntes, infírmi, et vix ullíus moménti inter hómines, audémus tractáre ista, et ista expónere: et putámus nos aut cápere posse cum cogitámus, aut capi dum dícimus.\par
\switchcolumn
\EN
John is an eagle that soareth up to the things that are not made, and fixeth his eye unquailing upon the secret and eternal Light. It is said that the father-eagle taketh his young in his talons and flieth with them, that they may look at the sun. Them that look at it boldly, he acknowledgeth for his own offspring, but them whose eyes shrink, he letteth fall from his claws as bastards. Let us consider, then, how mightily he must speak who is likened to an eagle; and yet we who are reptiles crawling on earth, weak, and of small consideration even among men, dare take in hand these utterances, to treat of them, and to explain them, and think that our intelligence understandeth them, or our speech commendeth them.\par
\switchcolumn*

\LA
\TeDeum{Te Deum}
\switchcolumn
\EN
\TeDeum{Te Deum}
\switchcolumn*

\end{paracol}

\par\needspace{10\baselineskip}\addvspace{1.2em}{\centering\small\textsc{Die 4 Januarii} · \textsc{4 January}\par}
\Office{In Octava Ss. Innocentium}{Octave Day of the Holy Innocents}{Simplex}
\addcontentsline{toc}{subsection}{Die 4 Januarii · In Octava Ss. Innocentium\enspace\textit{Octave Day of the Holy Innocents}}
\begin{paracol}{2}
\LA
\RefLine{Lectiones I–II: de Scriptura occurrente.}
\switchcolumn
\EN
\RefLine{Lessons I–II: from the Scripture of the day.}
\switchcolumn*

\LA
\LessonHead{Lectio IV}
\switchcolumn
\EN
\LessonHead{Lesson IV}
\switchcolumn*

\LA
\LessonTitle{Sermo sancti Augustíni Epíscopi}
\switchcolumn
\EN
\LessonTitle{From the Sermons of St. Augustine, Bishop (of Hippo.)}
\switchcolumn*

\LA
\Source{Sermo 1 de Innocentibus}
\switchcolumn
\EN
\Source{1st for Childermas.}
\switchcolumn*

\LA
\noindent Nascénte Dómino, luctus cœpit, non cælo, sed mundo: indícitur mátribus lamentátio, Angelis exsultátio, infántibus transmigrátio. Deus est, qui natus est: Innocéntes illi debéntur víctima, qui venit damnáre mundi malítiam. Agnélli debent immolári, quia Agnus futúrus est crucifígi, qui tollit peccáta mundi. Sed oves úlulant matres, quia agnos perdunt sine voce balántes. Grande martýrium, crudéle spectáculum! Exímitur machæra, et nulla intérvenit causa: sola stridet invídia, cum qui natus est, nulli fáciat violéntiam. Sed oves cérnimus matres: quæ super agnos lugent: Vox in Rama audíta est, plorátus et ululátus magnus. Pígnora sunt, non crédita, sed creáta; non depósita, sed expósita.\par
\switchcolumn
\EN
\noindent The Lord is born, and sorrow breaketh out, not in heaven but on earth; to mothers is proclaimed lamentation, to angels joy, to children translation. God is born, and innocence must be offered up to Him Who cometh to condemn the malice of the world. The Lamb that taketh away the sins of the world is come to be crucified, and the tender flock is brought to the sacrifice. But the mothers will lament over them whose inarticulate bleating is silenced for ever. Let us turn a look on this great martyrdom, this heart-rending sorrow. The sword is drawn, though there is no offence to punish, only jealousy shrieking for Him Who is born, and doth no violence. And here are mothers weeping over the lambs of the flock. In Ramah was there a voice heard, weeping and great mourning. which shall be returned hereafter, but they are pledges taken without being given, impounded without being entrusted.\par
\switchcolumn*

\LA
\TeDeum{Te Deum}
\switchcolumn
\EN
\TeDeum{Te Deum}
\switchcolumn*

\end{paracol}

\par\needspace{10\baselineskip}\addvspace{1.2em}{\centering\small\textsc{Die 5 Januarii} · \textsc{5 January}\par}
\Office{In Vigilia Epiphaniæ}{In Vigilia Epiphaniæ}{Semiduplex}
\addcontentsline{toc}{subsection}{Die 5 Januarii · In Vigilia Epiphaniæ\enspace\textit{In Vigilia Epiphaniæ}}
\begin{paracol}{2}
\LA
\RefLine{Lectiones I–III: de Scriptura occurrente.}
\switchcolumn
\EN
\RefLine{Lessons I–III: from the Scripture of the day.}
\switchcolumn*

\LA
\LessonHead{Lectio IV}
\switchcolumn
\EN
\LessonHead{Lesson IV}
\switchcolumn*

\LA
\LessonTitle{Sermo sancti Augustíni Epíscopi}
\switchcolumn
\EN
\LessonTitle{A sermon of St. Augustine, Bishop.}
\switchcolumn*

\LA
\Source{Sermo 13 de Tempore}
\switchcolumn
\EN
\Source{Serm. 13. de Tempore.}
\switchcolumn*

\LA
\noindent Dóminus noster Jesus Christus, fratres caríssimi, qui in ætérnum est cunctórum Creátor, hódie de matre nascéndo factus est nobis Salvátor. Natus est nobis hódie in témpore per voluntátem, ut nos perdúcat ad Patris æternitátem. Factus est Deus homo, ut homo fíeret Deus: ut panem Angelórum manducáret homo, Dóminus Angelórum hódie factus est homo.\par
\switchcolumn
\EN
\noindent Our Lord Jesus Christ, dearest brethren, who in eternity is the Creator of all things, was as at this time born of a mother and became our Saviour. It was as at this time that he willed to be born for us in earthly time, so as to lead us to the Father's eternity. God is made man, that man may be made as God. That man may eat Angels' food, the Lord of Angels was as on this day made man.\par
\switchcolumn*

\LA
\begin{Resp}\Rsign Congratulámini mihi, omnes qui dilígitis Dóminum: \Star{} Quia, cum essem párvula, plácui Altíssimo, et de meis viscéribus génui Deum et hóminem.\end{Resp}
\begin{Resp}\Vsign Beátam me dicent omnes generatiónes, quia ancíllam húmilem respéxit Deus.\end{Resp}
\begin{Resp}\Rsign Quia, cum essem párvula, plácui Altíssimo, et de meis viscéribus génui Deum et hóminem.\end{Resp}
\switchcolumn
\EN
\begin{Resp}\Rsign Rejoice with me, all ye that love the Lord. \Star{} For while I was yet little I pleased the Most High, and from my womb have I brought forth God and man.\end{Resp}
\begin{Resp}\Vsign All generations shall call me blessed, for God hath regarded the lowliness of His hand-maiden.\end{Resp}
\begin{Resp}\Rsign For while I was yet little I pleased the Most High, and from my womb have I brought forth God and man.\end{Resp}
\switchcolumn*

\LA
\LessonHead{Lectio V}
\switchcolumn
\EN
\LessonHead{Lesson V}
\switchcolumn*

\LA
\noindent Hódie impléta est prophetía illa, quæ dicit: Roráte, cæli, désuper, et nubes pluant justum: aperiátur terra, et gérminet Salvatórem. Factus est ígitur qui fécerat, ut invenirétur qui períerat. Sic enim in Psalmis homo confitétur: Priúsquam humiliárer, ego peccávi. Peccávit homo, et factus est reus: natus est homo Deus, ut liberarétur reus. Homo ígitur cécidit, sed Deus descéndit. Cécidit homo miserabíliter, descéndit Deus misericórditer; cécidit homo per supérbiam, descéndit Deus cum grátia.\par
\switchcolumn
\EN
\noindent Now is fulfilled that prophecy: Drop down, ye heavens, from above, and let the skies pour down righteousness: let the earth open, and bring forth a Saviour. He who made all things is therefore himself made, that those who are lost may be found. It is even as man is made to testify of himself in the Psalms: Before I was humbled, I went wrong. Man sinned and became guilty. God is born man, that the guilty may be delivered. Man fell, but God descended. Man fell miserably, God descended mercifully. Man fell by pride, God descended with grace.\par
\switchcolumn*

\LA
\begin{Resp}\Rsign Confirmátum est cor Vírginis, in quo divína mystéria, Angelo nuntiánte, concépit: tunc speciósum forma præ fíliis hóminum castis suscépit viscéribus: \Star{} Et benedícta in ætérnum, Deum nobis prótulit et hóminem.\end{Resp}
\begin{Resp}\Vsign Domus pudíci péctoris templum repénte fit Dei: intácta nésciens virum, verbo concépit Fílium.\end{Resp}
\begin{Resp}\Rsign Et benedícta in ætérnum, Deum nobis prótulit et hóminem.\end{Resp}
\switchcolumn
\EN
\begin{Resp}\Rsign The heart of the Virgin was fixed, when the Angel declared unto her the mystery of God, and she conceived, then did she receive in her pure womb Him That is fairer than the children of men. \Star{} And, she that is blessed for ever, brought forth for us God and man.\end{Resp}
\begin{Resp}\Vsign Soon rises, in that modest shrine, The Temple of the Lord Divine The stainless and unwedded one, Within her womb conceived the Son.\end{Resp}
\begin{Resp}\Rsign And, she that is blessed for ever, brought forth for us God and man.\end{Resp}
\switchcolumn*

\LA
\LessonHead{Lectio VI}
\switchcolumn
\EN
\LessonHead{Lesson VI}
\switchcolumn*

\LA
\noindent O mirácula, o prodígia, fratres mei! Natúræ jura mutántur in hómine; Deus náscitur, virgo sine viro gravidátur, viri nésciam sermo Dei marítat: simul facta est mater et virgo; mater facta, sed incorrúpta; virgo habens fílium, nésciens virum; semper clausa, sed non infœcúnda. Solus enim sine peccáto est natus, quem sine viríli compléxu non concupiscéntia carnis, sed obediéntia génuit mentis.\par
\switchcolumn
\EN
\noindent O my brethren, what a miracle! what a wonder! The laws of nature are changed concerning man: God is born, a Virgin conceiveth without an husband; the Word of God is wedded to one who knoweth no man; she is at once Mother and Virgin. A Mother, yet inviolate: a Virgin having a Son; knowing no man, ever sealed, yet not unfruitful. For he alone was born without sin. He alone was born without human embrace, begotten not of the will of the flesh, but of the obedience of the mind.\par
\switchcolumn*

\LA
\begin{Resp}\Rsign Benedícta et venerábilis es, Virgo María, quæ sine tactu pudóris invénta es mater Salvatóris: \Star{} Jacébat in præsépio, et fulgébat in cælo.\end{Resp}
\begin{Resp}\Vsign Dómine, audívi audítum tuum, et tímui: considerávi ópera tua, et expávi: in médio duórum animálium.\end{Resp}
\begin{Resp}\Rsign Jacébat in præsépio, et fulgébat in cælo.\end{Resp}
\begin{Resp}\Vsign Glória Patri, et Fílio, \Star{} et Spirítui Sancto.\end{Resp}
\begin{Resp}\Rsign Jacébat in præsépio, et fulgébat in cælo.\end{Resp}
\switchcolumn
\EN
\begin{Resp}\Rsign Blessed and worshipful art thou, O Virgin Mary; from thee, still maiden undefiled, the Saviour came a little Child. \Star{} He Whose glory filled the heavens lay in a manger.\end{Resp}
\begin{Resp}\Vsign O Lord, I have heard thy speech and was afraid; I considered thy works and trembled. O Thou That dwellest between the two living creatures!\end{Resp}
\begin{Resp}\Rsign He Whose glory filled the heavens lay in a manger.\end{Resp}
\begin{Resp}\Vsign Glory be to the Father, and to the Son, \Star{} and to the Holy Ghost.\end{Resp}
\begin{Resp}\Rsign He Whose glory filled the heavens lay in a manger.\end{Resp}
\switchcolumn*

\LA
\LessonHead{Lectio VII}
\switchcolumn
\EN
\LessonHead{Lesson VII}
\switchcolumn*

\LA
\LessonTitle{Léctio sancti Evangélii secúndum Matthǽum}
\switchcolumn
\EN
\LessonTitle{From the Holy Gospel according to Matthew}
\switchcolumn*

\LA
\Cite{Matt 2:19-23}
\switchcolumn
\EN
\Cite{Matt 2:19-23}
\switchcolumn*

\LA
\noindent In illo témpore: Defúncto Heróde, ecce Angelus Dómini appáruit in somnis Joseph in Ægýpto, dicens: Surge, et áccipe púerum et matrem ejus, et vade in terram Israël. Et réliqua.\par
\switchcolumn
\EN
\noindent In that time: When Herod was dead, behold an angel of the Lord appeared in sleep to Joseph in Egypt, Saying: Arise, and take the child and his mother, and go into the land of Israel. And so on.\par
\switchcolumn*

\LA
\LessonTitle{Homilía sancti Hierónymi Presbýteri}
\switchcolumn
\EN
\LessonTitle{Homily by St. Jerome, Priest (at Bethlehem.)}
\switchcolumn*

\LA
\Source{Liber 1 in Comment. in cap. 2 Matthæi}
\switchcolumn
\EN
\Source{Bk. i., Comm. on Matth. ii.}
\switchcolumn*

\LA
\noindent Ex hoc loco intellígimus non solum Heródem, sed et sacerdótes et scribas eódem témpore necem Dómini fuísse meditátos. Qui surgens accépit púerum, et matrem ejus. Non dixit: Accépit fílium suum et uxórem suam, sed Púerum et matrem ejus; quasi nutrítius, non marítus.\par
\switchcolumn
\EN
\noindent From the words, they are dead, (in the Plural), which are used in this passage of the Gospel, we may understand that there were others beside Herod which sought the young Child's life probably the Priests and Scribes. And he Joseph arose, and took the young Child and His Mother. It is not written, He took his wife and child, but he took the young Child and His Mother; whence it is clear that the holy Evangelist willeth to imply that Joseph was not the father, but the Guardian of Jesus, not the husband, but the Betrothed of Mary.\par
\switchcolumn*

\LA
\begin{Resp}\Rsign Sancta et immaculáta virgínitas, quibus te láudibus éfferam, néscio: \Star{} Quia quem cæli cápere non póterant, tuo grémio contulísti.\end{Resp}
\begin{Resp}\Vsign Benedícta tu in muliéribus, et benedíctus fructus ventris tui.\end{Resp}
\begin{Resp}\Rsign Quia quem cæli cápere non póterant, tuo grémio contulísti.\end{Resp}
\switchcolumn
\EN
\begin{Resp}\Rsign O Mary, how holy and how spotless is thy virginity! I am too dull to praise thee! \Star{} For thou hast borne in thy breast Him Whom the heavens cannot contain.\end{Resp}
\begin{Resp}\Vsign Blessed art thou among women, and blessed is the fruit of thy womb.\end{Resp}
\begin{Resp}\Rsign For thou hast borne in thy breast Him Whom the heavens cannot contain.\end{Resp}
\switchcolumn*

\LA
\LessonHead{Lectio VIII}
\switchcolumn
\EN
\LessonHead{Lesson VIII}
\switchcolumn*

\LA
\noindent Audiens autem quod Archeláus regnáret in Judǽa pro Heróde patre suo, tímuit illo ire. Multi labúntur erróre propter ignorántiam históriæ, putántes eumdem esse Heródem, a quo in passióne sua Dóminus irridétur, et qui nunc mórtuus esse referátur. Ergo Heródes ille qui cum Piláto póstea amicítias fecit, hujus Heródis fílius est, frater Archelái.\par
\switchcolumn
\EN
\noindent But when he heard that Archelaus did reign in Judea, in the room of his father Herod, he was afraid to go thither. There are some persons so grossly ignorant of history that they confuse themselves over the two Herods, as if the one mentioned here were the same who afterwards set our Lord at nought during His Passion, and they cannot understand how he should now be said to be dead. The Herod who was made friends with Pilate over Christ's death, was the son of the Herod who massacred the infants of Bethlehem, and the brother of Archelaus.\par
\switchcolumn*

\LA
\begin{Resp}\Rsign Nésciens Mater Virgo virum, péperit sine dolóre: \Star{} Salvatórem sæculórum, ipsum Regem Angelórum, sola Virgo lactábat úbere de cælo pleno.\end{Resp}
\begin{Resp}\Vsign Domus pudíci péctoris templum repénte fit Dei: intácta nésciens virum, verbo concépit Fílium.\end{Resp}
\begin{Resp}\Rsign Salvatórem sæculórum, ipsum Regem Angelórum, sola Virgo lactábat úbere de cælo pleno.\end{Resp}
\begin{Resp}\Vsign Glória Patri, et Fílio, \Star{} et Spirítui Sancto.\end{Resp}
\begin{Resp}\Rsign Salvatórem sæculórum, ipsum Regem Angelórum, sola Virgo lactábat úbere de cælo pleno.\end{Resp}
\switchcolumn
\EN
\begin{Resp}\Rsign The Virgin-Mother that knew not a man, bore, but travailed not. \Star{} She fed the Saviour of the world, The King of Angel hosts above, Jesus, our Redeemer blest, From the fountain of her breast.\end{Resp}
\begin{Resp}\Vsign Soon rises in that modest shrine, The Temple of the Lord Divine; The stainless and unwedded one, Within her womb conceived the Son.\end{Resp}
\begin{Resp}\Rsign She fed the Saviour of the world, The King of Angel hosts above, Jesus, our Redeemer blest, From the fountain of her breast.\end{Resp}
\begin{Resp}\Vsign Glory be to the Father, and to the Son, \Star{} and to the Holy Ghost.\end{Resp}
\begin{Resp}\Rsign She fed the Saviour of the world, The King of Angel hosts above, Jesus, our Redeemer blest, From the fountain of her breast.\end{Resp}
\switchcolumn*

\LA
\LessonHead{Lectio IX}
\switchcolumn
\EN
\LessonHead{Lesson IX}
\switchcolumn*

\LA
\noindent Quóniam Nazarǽus vocábitur. Si fixum de Scriptúris posuísset exémplum, numquam díceret: Quod dictum est per prophétas: sed simplíciter: Quod dictum est per prophétam. Nunc autem pluráliter prophétas vocans, osténdit se non verba de Scriptúris sumpsísse, sed sensum. Nazarǽus sanctus interpretátur; sanctum autem Dóminum futúrum omnis Scriptúra commémorat.\par
\switchcolumn
\EN
\noindent He shall be called a Nazarene. The Evangelist, in quoting these words, saith that they were spoken by the Prophets, (Plural). If he had been citing any one precise passage he would have said by the Prophet, in the Singular. But he is citing the sense of the Prophets, and not any individual passage in any of their writings. He seemeth to refer to the fact that in Hebrew the word Nazarene signifieth holy, and that Christ is the Holy One of God is the common declaration of all the Scriptures.\par
\switchcolumn*

\LA
\TeDeum{Te Deum}
\switchcolumn
\EN
\TeDeum{Te Deum}
\switchcolumn*

\end{paracol}

\par\needspace{10\baselineskip}\addvspace{1.2em}{\centering\small\textsc{Die 6 Januarii} · \textsc{6 January}\par}
\Office{In Epiphania Domini}{In Epiphania Domini}{Duplex I classis}
\addcontentsline{toc}{subsection}{Die 6 Januarii · In Epiphania Domini\enspace\textit{In Epiphania Domini}}
\begin{paracol}{2}
\LA
\LessonHead{Lectio I}
\switchcolumn
\EN
\LessonHead{Lesson I}
\switchcolumn*

\LA
\LessonTitle{De Isaía Prophéta}
\switchcolumn
\EN
\LessonTitle{Lesson from the book of Isaias}
\switchcolumn*

\LA
\Cite{Isa 55:1-4}
\switchcolumn
\EN
\Cite{Isa 55:1-4}
\switchcolumn*

\LA
\noindent Omnes sitiéntes, veníte ad aquas, et qui non habétis argéntum, properáte, émite, et comédite: veníte, émite absque argénto et absque ulla commutatióne vinum et lac. Quare appénditis argéntum non in pánibus, et labórem vestrum non in saturitáte? Audíte, audiéntes me, et comédite bonum, et delectábitur in crassitúdine ánima vestra. Inclináte aurem vestram, et veníte ad me; audíte, et vivet ánima vestra, et fériam vobíscum pactum sempitérnum, misericórdias David fidéles. Ecce testem pópulis dedi eum, ducem ac præceptórem géntibus.\par
\switchcolumn
\EN
\noindent All you that thirst, come to the waters: and you that have no money make haste, buy, and eat: come ye, buy wine and milk without money, and without any price. Why do you spend money for that which is not bread, and your labour for that which doth not satisfy you? Hearken diligently to me, and eat that which is good, and your soul shall be delighted in fatness. Incline your ear and come to me: hear and your soul shall live, and I will make an everlasting covenant with you, the faithful mercies of David. Behold I have given him for a witness to the people, for a leader and a master to the Gentiles.\par
\switchcolumn*

\LA
\begin{Resp}\Rsign Hódie in Jordáne baptizáto Dómino apérti sunt cæli, et sicut colúmba super eum Spíritus mansit, et vox Patris intónuit: \Star{} Hic est Fílius meus diléctus, in quo mihi bene complácui.\end{Resp}
\begin{Resp}\Vsign Descéndit Spíritus Sanctus corporáli spécie sicut colúmba in ipsum, et vox de cælo facta est.\end{Resp}
\begin{Resp}\Rsign Hic est Fílius meus diléctus, in quo mihi bene complácui.\end{Resp}
\switchcolumn
\EN
\begin{Resp}\Rsign This day, when the Lord was baptized in Jordan, the heavens were opened, and the Spirit descended like a dove, and abode upon Him, and, lo, the voice of the Father was heard, like unto thunder, saying: \Star{} This is My beloved Son, in Whom I am well pleased.\end{Resp}
\begin{Resp}\Vsign The Holy Ghost descended in a bodily shape like a dove upon Him, and a voice came from heaven\end{Resp}
\begin{Resp}\Rsign This is My beloved Son, in Whom I am well pleased.\end{Resp}
\switchcolumn*

\LA
\LessonHead{Lectio II}
\switchcolumn
\EN
\LessonHead{Lesson II}
\switchcolumn*

\LA
\Cite{Isa 60:1-6}
\switchcolumn
\EN
\Cite{Isa 60:1-6}
\switchcolumn*

\LA
\noindent Surge, illumináre, Jerúsalem, quia venit lumen tuum, et glória Dómini super te orta est. Quia ecce ténebræ opérient terram, et calígo pópulos; super te autem oriétur Dóminus, et glória ejus in te vidébitur. Et ambulábunt gentes in lúmine tuo, et reges in splendóre ortus tui. Leva in circúitu óculos tuos, et vide: omnes isti congregáti sunt, venérunt tibi; fílii tui de longe vénient et fíliæ tuæ de látere surgent. Tunc vidébis, et áfflues; mirábitur et dilatábitur cor tuum: quando convérsa fúerit ad te multitúdo maris; fortitúdo géntium vénerit tibi. Inundátio camelórum opériet te, dromedárii Mádian et Epha; omnes de Saba vénient, aurum et thus deferéntes, et laudem Dómino annuntiántes.\par
\switchcolumn
\EN
\noindent Arise, be enlightened, O Jerusalem: for thy light is come, and the glory of the Lord is risen upon thee. For behold darkness shall cover the earth, and a mist the people: but the Lord shall arise upon thee, and his glory shall be seen upon thee. And the Gentiles shall walk in thy light, and kings in the brightness of thy rising. Lift up thy eyes round about, and see: all these are gathered together, they are come to thee: thy sons shall come from afar, and thy daughters shall rise up at thy side. Then shalt thou see, and abound, and thy heart shall wonder and be enlarged, when the multitude of the sea shall be converted to thee, the strength of the Gentiles shall come to thee. The multitude of camels shall cover thee, the dromedaries of Madian and Epha: all they from Saba shall come, bringing gold and frankincense: and shewing forth praise to the Lord.\par
\switchcolumn*

\LA
\begin{Resp}\Rsign In colúmbæ spécie Spíritus Sanctus visus est, Patérna vox audíta est: \Star{} Hic est Fílius meus diléctus, in quo mihi bene complácui.\end{Resp}
\begin{Resp}\Vsign Cæli apérti sunt super eum, et vox Patris intónuit.\end{Resp}
\begin{Resp}\Rsign Hic est Fílius meus diléctus, in quo mihi bene complácui.\end{Resp}
\switchcolumn
\EN
\begin{Resp}\Rsign The Holy Ghost appeared in a bodily shape like a dove, and the voice of the Father was heard: \Star{} This is My beloved Son, in Whom I am well pleased.\end{Resp}
\begin{Resp}\Vsign The heavens were opened unto him, and, lo, the voice of the Father was heard, like unto thunder, saying:\end{Resp}
\begin{Resp}\Rsign This is My beloved Son, in Whom I am well pleased.\end{Resp}
\switchcolumn*

\LA
\LessonHead{Lectio III}
\switchcolumn
\EN
\LessonHead{Lesson III}
\switchcolumn*

\LA
\Cite{Isa 61:10-11; 62:1}
\switchcolumn
\EN
\Cite{Isa 61:10-11; 62:1}
\switchcolumn*

\LA
\noindent Gaudens gaudébo in Dómino, et exsultábit ánima mea in Deo meo, quia índuit me vestiméntis salútis, et induménto justítiæ circúmdedit me, quasi sponsum decorátum coróna, et quasi sponsam ornátam monílibus suis. Sicut enim terra profert germen suum, et sicut hortus semen suum gérminat, sic Dóminus Deus germinábit justítiam et laudem coram univérsis géntibus. Propter Sion non tacébo, et propter Jerúsalem non quiéscam, donec egrediátur ut splendor justus ejus, et salvátor ejus ut lampas accendátur.\par
\switchcolumn
\EN
\noindent I will greatly rejoice in the Lord, and my soul shall be joyful in my God: for he hath clothed me with the garments of salvation: and with the robe of justice he hath covered me, as a bridegroom decked with a crown, and as a bride adorned with her jewels. For as the earth bringeth forth her bud, and as the garden causeth her seed to shoot forth: so shall the Lord God make justice to spring forth, and praise before all the nations. For Sion's sake I will not hold my peace, and for the sake of Jerusalem, I will not rest till her just one come forth as brightness, and her saviour be lighted as a lamp.\par
\switchcolumn*

\LA
\begin{Resp}\Rsign Reges Tharsis et ínsulæ múnera ófferent: \Star{} Reges Arabum et Saba dona Dómino Deo addúcent.\end{Resp}
\begin{Resp}\Vsign Omnes de Saba vénient, aurum et thus deferéntes.\end{Resp}
\begin{Resp}\Rsign Reges Arabum et Saba dona Dómino Deo addúcent.\end{Resp}
\begin{Resp}\Vsign Glória Patri, et Fílio, \Star{} et Spirítui Sancto.\end{Resp}
\begin{Resp}\Rsign Reges Arabum et Saba dona Dómino Deo addúcent.\end{Resp}
\switchcolumn
\EN
\begin{Resp}\Rsign The kings of Tarshish and of the isles shall bring presents. \Star{} The kings of Arabia and Saba shall offer gifts unto the Lord God.\end{Resp}
\begin{Resp}\Vsign All they from Saba shall come, they shall bring gold and incense.\end{Resp}
\begin{Resp}\Rsign The kings of Arabia and Saba shall offer gifts unto the Lord God.\end{Resp}
\begin{Resp}\Vsign Glory be to the Father, and to the Son, \Star{} and to the Holy Ghost.\end{Resp}
\begin{Resp}\Rsign The kings of Arabia and Saba shall offer gifts unto the Lord God.\end{Resp}
\switchcolumn*

\LA
\LessonHead{Lectio IV}
\switchcolumn
\EN
\LessonHead{Lesson IV}
\switchcolumn*

\LA
\LessonTitle{Sermo sancti Leónis Papæ}
\switchcolumn
\EN
\LessonTitle{From the Sermons of Pope St. Leo (the Great)}
\switchcolumn*

\LA
\Source{Sermo 2 de Epiphania}
\switchcolumn
\EN
\Source{2nd for Twelfth-Day.}
\switchcolumn*

\LA
\noindent Gaudéte in Dómino, dilectíssimi, íterum dico, gaudéte: quóniam brevi intervállo témporis, post solemnitátem Nativitátis Christi, festívitas declaratiónis ejus illúxit: et quem in illo die Virgo péperit, in hoc mundus agnóvit. Verbum enim caro factum, sic susceptiónis nostræ temperávit exórdia, ut natus Jesus et credéntibus maniféstus, et persequéntibus esset occúltus. Jam tunc ergo cæli enarravérunt glóriam Dei, et in omnem terram sonus veritátis exívit, quando et pastóribus exércitus Angelórum Salvatóris éditi annuntiátor appáruit, et Magos ad eum adorándum prǽvia stella perdúxit: ut a solis ortu usque ad occásum veri Regis generátio coruscáret, cum rerum fidem et regna Oriéntis per Magos díscerent, et Románum impérium non latéret.\par
\switchcolumn
\EN
\noindent Dearly beloved brethren, rejoice in the Lord; again I say, rejoice. But a few days are past since the solemnity of Christ's Birth, and now the glorious light of His Manifestation is breaking upon us. On that day the Virgin brought Him forth, and on this the world knew Him. The Word made Flesh was pleased to reveal Himself by degrees to those for whom He had come. When Jesus was born He was manifested indeed to the believing, but hidden from His enemies. Already indeed the heavens declared the glory of God, and their sound went out into all lands, when the Herald Angels appeared to tell to the shepherds the glad tidings of a Saviour's Birth; and now the guiding star leadeth the wise men to worship Him, that from the rising of the sun to the going down thereof, the Birth of the true King may be known abroad; that through those wise men the kingdoms of the east might learn the great truth, and the Roman empire remain no more in darkness.\par
\switchcolumn*

\LA
\begin{Resp}\Rsign Illumináre, illumináre Jerúsalem, quia venit lux tua: \Star{} Et glória Dómini super te orta est.\end{Resp}
\begin{Resp}\Vsign Et ambulábunt gentes in lúmine tuo, et reges in splendóre ortus tui.\end{Resp}
\begin{Resp}\Rsign Et glória Dómini super te orta est.\end{Resp}
\switchcolumn
\EN
\begin{Resp}\Rsign Shine, shine, O Jerusalem, for thy light is come; \Star{} And the glory of the Lord is risen upon thee.\end{Resp}
\begin{Resp}\Vsign And the Gentiles shall walk in thy light, and kings in the brightness of thy rising.\end{Resp}
\begin{Resp}\Rsign And the glory of the Lord is risen upon thee.\end{Resp}
\switchcolumn*

\LA
\LessonHead{Lectio V}
\switchcolumn
\EN
\LessonHead{Lesson V}
\switchcolumn*

\LA
\noindent Nam et sævítia Heródis volens primórdia suspécti sibi Regis exstínguere, huic dispensatíoni nésciens serviébat: ut dum atróci inténtus facínori, ignótum sibi púerum indiscréta infántium cæde perséquitur, annuntiátum cǽlitus dominatóris ortum insígnior ubíque fama loquerétur: quam promptiórem ad narrándum, diligentiorémque faciébat et supérnæ significatiónis nóvitas, et cruentíssimi persecutóris impíetas. Tunc autem étiam Ægýpto Salvátor illátus est, ut gens antíquis erróribus dédita, jam ad vicínam salútem per occúltam grátiam signarétur: et quæ nondum ejécerat ab ánimo superstitiónem, jam hospítio recíperet veritátem.\par
\switchcolumn
\EN
\noindent The very cruelty of Herod, when he strove to crush at His birth this King Whom he alone feared, was made a blind means to carry out this dispensation of mercy. While the tyrant with horrid guilt sought to slay the little Child he did not know, amid an indiscriminate slaughter of innocents, his infamous act served to spread wider abroad the heaven-told news of the Birth of the Lord. Thus were these glad tidings loudly proclaimed, both by the novelty of their story, and the iniquity of their enemies. Then was the Saviour borne into Egypt, that nation, of a long time hardened in idolatry, might by the mysterious virtue which went out of Him, even when His presence was unknown, be prepared for the saving light so soon to dawn on them, and might receive the Truth as a wanderer even before they had banished falsehood.\par
\switchcolumn*

\LA
\begin{Resp}\Rsign Omnes de Saba vénient, aurum et thus deferéntes, et laudem Dómino annuntiántes, \Star{} Allelúja, allelúja, allelúja.\end{Resp}
\begin{Resp}\Vsign Reges Tharsis et ínsulæ múnera ófferent, reges Arabum et Saba dona addúcent.\end{Resp}
\begin{Resp}\Rsign Allelúja, allelúja, allelúja.\end{Resp}
\switchcolumn
\EN
\begin{Resp}\Rsign All they from Saba shall come, they shall bring gold and incense, and they shall show forth the praises of the Lord. \Star{} Alleluia, Alleluia, Alleluia.\end{Resp}
\begin{Resp}\Vsign The kings of Tarshish and of the isles shall bring presents, the kings of Arabia and Saba shall offer gifts.\end{Resp}
\begin{Resp}\Rsign Alleluia, Alleluia, Alleluia.\end{Resp}
\switchcolumn*

\LA
\LessonHead{Lectio VI}
\switchcolumn
\EN
\LessonHead{Lesson VI}
\switchcolumn*

\LA
\noindent Agnoscámus ergo, dilectíssimi, in Magis adoratóribus Christi, vocatiónis nostræ fideíque primítias: et exsultántibus ánimis beátæ spei inítia celebrémus. Exínde enim in ætérnam hereditátem cœ́pimus introíre: exínde nobis Christum loquéntia Scripturárum arcána patuérunt, et véritas, quam Judæórum obcæcátio non récipit, ómnibus natiónibus lumen suum invéxit. Honorétur ítaque a nobis sacratíssimus dies, in quo salútis nostræ Auctor appáruit: et quem Magi infántem veneráti sunt in cunábulis, nos omnipoténtem adorémus in cælis. Ac sicut illi de thesáuris suis mýsticas Dómino múnerum spécies obtulérunt, ita et nos de córdibus nostris, quæ Deo sunt digna, promámus.\par
\switchcolumn
\EN
\noindent Dearly beloved brethren, we recognize in the wise men who came to worship Christ, the first-fruits of that dispensation to the Gentiles wherein we also are called and enlightened. Let us then keep this Feast with grateful hearts, in thanksgiving for our blessed hope, whereof it doth commemorate the dawn. From that worship paid to the new-born Christ is to be dated the entry of us Gentiles upon our heirship of God and co-heirship with Christ. Since that joyful day the Scriptures which testify of Christ have lain open for us as well as for the Jews. Yea, their blindness rejected that Truth, Which, since that day, hath shed Its bright beams upon all nations. Let all observance, then, be paid to this most sacred day, whereon the Author of our salvation was made manifest, and as the wise men fell down and worshipped Him in the manger, so let us fall down and worship Him enthroned Almighty in heaven. As they also opened their treasures and presented unto Him mystic and symbolic gifts, so let us strive to open our hearts to Him, and offer Him from thence some worthy offering.\par
\switchcolumn*

\LA
\begin{Resp}\Rsign Magi véniunt ab Oriénte Jerosólymam, quæréntes, et dicéntes: Ubi est qui natus est, cujus stellam vídimus? \Star{} Et vénimus adoráre Dóminum.\end{Resp}
\begin{Resp}\Vsign Vídimus stellam ejus in Oriénte.\end{Resp}
\begin{Resp}\Rsign Et vénimus adoráre Dóminum.\end{Resp}
\begin{Resp}\Vsign Glória Patri, et Fílio, \Star{} et Spirítui Sancto.\end{Resp}
\begin{Resp}\Rsign Et vénimus adoráre Dóminum.\end{Resp}
\switchcolumn
\EN
\begin{Resp}\Rsign There came wise men from the east to Jerusalem, asking and saying: Where is He That is born King of the Jews? for we have seen His star in the east; \Star{} And are come to worship the Lord.\end{Resp}
\begin{Resp}\Vsign We have seen His star in the east.\end{Resp}
\begin{Resp}\Rsign And are come to worship the Lord.\end{Resp}
\begin{Resp}\Vsign Glory be to the Father, and to the Son, \Star{} and to the Holy Ghost.\end{Resp}
\begin{Resp}\Rsign And are come to worship the Lord.\end{Resp}
\switchcolumn*

\LA
\LessonHead{Lectio VII}
\switchcolumn
\EN
\LessonHead{Lesson VII}
\switchcolumn*

\LA
\LessonTitle{Léctio sancti Evangélii secúndum Matthǽum}
\switchcolumn
\EN
\LessonTitle{From the Holy Gospel according to Matthew}
\switchcolumn*

\LA
\Cite{Matt 2:1-12}
\switchcolumn
\EN
\Cite{Matt 2:1-12}
\switchcolumn*

\LA
\noindent Cum natus esset Jesus in Béthlehem Juda in diébus Heródis regis, ecce Magi ab Oriénte venérunt Jerosólymam, dicéntes: Ubi est qui natus est Rex Judæórum? Et réliqua.\par
\switchcolumn
\EN
\noindent When Jesus therefore was born in Bethlehem of Juda, in the days of king Herod, behold, there came wise men from the east to Jerusalem, saying: Where is he that is born king of the Jews? And so on.\par
\switchcolumn*

\LA
\LessonTitle{Homilía sancti Gregórii Papæ}
\switchcolumn
\EN
\LessonTitle{Homily by Pope St. Gregory (the Great.)}
\switchcolumn*

\LA
\Source{Homilia 10 in Evang.}
\switchcolumn
\EN
\Source{10th on the Gospels}
\switchcolumn*

\LA
\noindent Sicut in lectióne evangélica, fratres caríssimi, audístis, cæli Rege nato, rex terræ turbátus est: quia nimírum terréna altitúdo confúnditur, cum celsitúdo cæléstis aperítur. Sed quæréndum nobis est, quidnam sit, quod, Redemptóre nato, pastóribus in Judǽa Angelus appáruit, atque ad adorándum hunc ab Oriénte Magos non Angelus, sed stella perdúxit? Quia vidélicet Judǽis, tamquam ratióne uténtibus, rationále ánimal, id est, Angelus prædicáre débuit: gentíles vero, quia uti ratióne nesciébant, ad cognoscéndum Dóminum non per vocem, sed per signa perducúntur. Unde étiam per Paulum dícitur: Prophetíæ fidélibus datæ sunt, non infidélibus: signa autem infidélibus, non fidélibus. Quia et illis prophetíæ tamquam fidélibus, non infidélibus: et istis signa tamquam infidélibus, non fidélibus data sunt.\par
\switchcolumn
\EN
\noindent Dearly beloved brethren, hear ye from the Gospel lesson how, when the King of heaven was born, the king of earth was troubled? The heights of heaven are opened and the depths of earth are stirred. Let us now consider the question, why, when the Redeemer was born, an angel brought the news to the shepherds of Judea, but a star led the wise men of the East to worship Him. It seemeth as if the Jews as reasonable creatures received a revelation from a reasonable being, that is, an angel, but the Gentiles without, being as brutes, are roused not by a voice, but by a sign, that is, a star. Hence Paul hath it: a sign, not to them that believe, but to them that believe not but prophesying serveth not for them that believe not, but for them which believe. (1 Cor. xiv. 22.) So the prophesying, that is, of the angel was given to them that believed, and the sign to them that believed not.\par
\switchcolumn*

\LA
\begin{Resp}\Rsign Stella, quam víderant Magi in Oriénte, antecedébat eos, donec venírent ad locum, ubi puer erat. \Star{} Vidéntes autem eam, gavísi sunt gáudio magno.\end{Resp}
\begin{Resp}\Vsign Et intrántes domum, invenérunt púerum cum María matre ejus, et procidéntes adoravérunt eum.\end{Resp}
\begin{Resp}\Rsign Vidéntes autem eam, gavísi sunt gáudio magno.\end{Resp}
\switchcolumn
\EN
\begin{Resp}\Rsign The star which the wise men had seen in the East, went before them, till they came where the young Child was. \Star{} And when they saw the star, they rejoiced with exceeding great joy.\end{Resp}
\begin{Resp}\Vsign And when they were come into the house, they found the young Child with Mary His Mother, and fell down and worshipped Him.\end{Resp}
\begin{Resp}\Rsign And when they saw the star, they rejoiced with exceeding great joy.\end{Resp}
\switchcolumn*

\LA
\LessonHead{Lectio VIII}
\switchcolumn
\EN
\LessonHead{Lesson VIII}
\switchcolumn*

\LA
\noindent Et notándum, quod Redemptórem nostrum, cum jam perféctæ esset ætátis, eísdem gentílibus Apóstoli prǽdicant, eúmque párvulum, et necdum per humáni córporis offícium loquéntem, stella géntibus denúntiat: quia nimírum ratiónis ordo poscébat, ut et loquéntem jam Dóminum loquéntes nobis prædicatóres innotéscerent, et necdum loquéntem eleménta muta prædicárent. Sed in ómnibus signis, quæ vel nascénte Dómino, vel moriénte eo monstráta sunt, considerándum nobis est, quanta fúerit in quorúmdam Judæórum corde durítia, qui hunc nec per prophetíæ donum, nec per mirácula agnovérunt.\par
\switchcolumn
\EN
\noindent Thus also we remark that afterwards the Redeemer was preached among the Gentiles not by Himself, but by His Apostles, even as, when a little Child, He is shown to them, not by the voice of angels, but merely by the vision of a star. When He Himself had begun to speak He was made known to us by speakers, but when He lay silent in the manger, by that silent testimony in heaven. But whether we consider the signs which accompanied His birth or His death, this thing is wonderful, namely, the hardness of heart of the Jews, who would not believe in Him either for prophesying or for miracles.\par
\switchcolumn*

\LA
\begin{Resp}\Rsign Vidéntes stellam Magi, gavísi sunt gáudio magno: \Star{} Et intrántes domum invenérunt púerum cum María matre ejus, et procidéntes adoravérunt eum: * Et apértis thesáuris suis obtulérunt ei múnera, aurum, thus, et myrrham.\end{Resp}
\begin{Resp}\Vsign Stella, quam víderant Magi in Oriénte, antecedébat eos, usque dum véniens staret supra ubi erat puer.\end{Resp}
\begin{Resp}\Rsign Et intrántes domum invenérunt púerum cum María matre ejus, et procidéntes adoravérunt eum:\end{Resp}
\begin{Resp}\Vsign Glória Patri, et Fílio, \Star{} et Spirítui Sancto.\end{Resp}
\begin{Resp}\Rsign Et apértis thesáuris suis obtulérunt ei múnera, aurum, thus, et myrrham.\end{Resp}
\switchcolumn
\EN
\begin{Resp}\Rsign When the wise men saw the star, they rejoiced with exceeding great joy. and when they had opened their treasures, they presented unto Him gifts; gold, and frankincense, and myrrh. \Star{} And when they were come into the house, they found the young Child with Mary His Mother, and fell down and worshipped Him; *\end{Resp}
\begin{Resp}\Vsign The star which the wise men had seen in the East, went before them, till it came and stood over where the young Child was.\end{Resp}
\begin{Resp}\Rsign And when they were come into the house, they found the young Child with Mary His Mother, and fell down and worshipped Him.\end{Resp}
\begin{Resp}\Vsign Glory be to the Father, and to the Son, \Star{} and to the Holy Ghost.\end{Resp}
\begin{Resp}\Rsign And when they had opened their treasures, they presented unto Him gifts; gold, and frankincense, and myrrh.\end{Resp}
\switchcolumn*

\LA
\LessonHead{Lectio IX}
\switchcolumn
\EN
\LessonHead{Lesson IX}
\switchcolumn*

\LA
\noindent Omnia quippe eleménta auctórem suum venísse testáta sunt. Ut enim de eis quiddam usu humáno loquar: Deum hunc cæli esse cognovérunt, quia prótinus stellam misérunt. Mare cognóvit, quia sub plantis ejus se calcábile prǽbuit. Terra cognóvit, quia eo moriénte contrémuit. Sol cognóvit, quia lucis suæ rádios abscóndit. Saxa et paríetes cognovérunt, quia témpore mortis ejus scissa sunt. Inférnus agnóvit, quia hos, quos tenébat mórtuos, réddidit. Et tamen hunc, quem Dóminum ómnia insensibília eleménta sensérunt, adhuc infidélium Judæórum corda Deum esse mínime cognóscunt, et durióra saxis, scindi ad pœniténdum nolunt.\par
\switchcolumn
\EN
\noindent All things which He had made, bore witness that their Maker was come. Let me reckon them after the manner of men. The heavens knew that He was God, and sent a star to shine over where He lay. The sea knew it, and bore Him up when He walked upon it. The earth knew it, and quaked when He died. The sun knew it, and was darkened. The rocks and walls knew it, and rent at the hour of His death. Hell knew it, and gave up the dead that were in it. And yet up to this very hour the hearts of the unbelieving Jews will not acknowledge that He to Whom all nature testified is their God, and, being more hardened than the rocks, refuse to be rent by repentance.\par
\switchcolumn*

\LA
\TeDeum{Te Deum}
\switchcolumn
\EN
\TeDeum{Te Deum}
\switchcolumn*

\end{paracol}

\par\needspace{10\baselineskip}\addvspace{1.2em}{\centering\small\textsc{Die 7 Januarii} · \textsc{7 January}\par}
\Office{Secunda die infra Octavam Epiphaniæ}{Secunda die infra Octavam Epiphaniæ}{Semiduplex II classis}
\addcontentsline{toc}{subsection}{Die 7 Januarii · Secunda die infra Octavam Epiphaniæ\enspace\textit{Secunda die infra Octavam Epiphaniæ}}
\begin{paracol}{2}
\LA
\RefLine{Lectiones I–III: de Scriptura occurrente.}
\switchcolumn
\EN
\RefLine{Lessons I–III: from the Scripture of the day.}
\switchcolumn*

\LA
\LessonHead{Lectio IV}
\switchcolumn
\EN
\LessonHead{Lesson IV}
\switchcolumn*

\LA
\LessonTitle{Sermo sancti Augustíni Epíscopi}
\switchcolumn
\EN
\LessonTitle{From the Sermons of St. Augustine, Bishop (of Hippo.)}
\switchcolumn*

\LA
\Source{Sermo 2 de Epiphania, qui est 30 de Tempore}
\switchcolumn
\EN
\Source{2nd on the Epiphany, 30th on the Season.}
\switchcolumn*

\LA
\noindent Ad partum Vírginis adorándum Magi ab Oriénte venérunt. Hunc diem hódie celebrámus: huic debitum solemnitáti sermónem persolvimus. Illis dies iste primus illuxit, anniversaria nobis festivitate rediit. Illi erant primítiæ Géntium, nos pópulus Géntium. Nobis hoc lingua nuntiávit Apostolórum, stella illis tamquam lingua cælórum; et nobis iídem Apóstoli, tamquam alii cæli, enarravérunt glóriam Dei.\par
\switchcolumn
\EN
\noindent Wise men came from the East to worship the Virgin's Son. This is the event which we this day commemorate, the occasion in honour of which this sermon is preached. On them that day first broke in gladness, which year by year, now cometh round to us for celebration. They were the first-fruits of that Gentile Church whereof we are the in-gathering. To us the voice of Apostles, to them a star, as a voice from heaven, proclaimed the advent of a Saviour; and to us the voice of the Apostolic preachers is also as a voice from heaven, a heaven declaring the glory of God.\par
\switchcolumn*

\LA
\begin{Resp}\Rsign Illumináre, illumináre Jerúsalem, quia venit lux tua: \Star{} Et glória Dómini super te orta est.\end{Resp}
\begin{Resp}\Vsign Et ambulábunt gentes in lúmine tuo, et reges in splendóre ortus tui.\end{Resp}
\begin{Resp}\Rsign Et glória Dómini super te orta est.\end{Resp}
\switchcolumn
\EN
\begin{Resp}\Rsign Shine, shine, O Jerusalem, for thy light is come; \Star{} And the glory of the Lord is risen upon thee.\end{Resp}
\begin{Resp}\Vsign And the Gentiles shall walk in thy light, and kings in the brightness of thy rising.\end{Resp}
\begin{Resp}\Rsign And the glory of the Lord is risen upon thee.\end{Resp}
\switchcolumn*

\LA
\LessonHead{Lectio V}
\switchcolumn
\EN
\LessonHead{Lesson V}
\switchcolumn*

\LA
\noindent Magnum sacraméntum: in præsepi jacebat, et Magos ab Oriénte ducebat. Abscondebátur in stábulo, et agnoscebátur in cælo: ut ágnitus in cælo manifestarétur in stábulo, et appellarétur Epiphanía dies iste, quod Latine manifestátio dici potest: simul ejus celsitúdinem humilitatémque commendans, ut qui in aperto cælo signis sidéreis monstrabátur magnus, in angusto diversorio quæsítus invenirétur inválidus, infantílibus in membris natus, infantilibúsque pannis involutus, adorarétur a Magis, timerétur a malis.\par
\switchcolumn
\EN
\noindent Great is the mystery. While He lay in the manger, He drew to Himself wise men from the East; while He was unknown in the stable, He was recognised in the heavens; and, being recognised in the heavens, made Himself known in the stable. So this day is called in the Greek Epiphaneia, which is, being interpreted, Manifestation. Wherein is manifested both the greatness and the lowliness of Him, Whose greatness was attested in the stars of heaven, and Who, being sought on earth, is found so lowly that there is no room for Him in the inn. And yet, though to be found in fashion as a little Child wrapped in swaddlingclothes, He is the object of worship to the wise men and of terror to the godless.\par
\switchcolumn*

\LA
\begin{Resp}\Rsign Omnes de Saba vénient, aurum et thus deferéntes, et laudem Dómino annuntiántes, \Star{} Allelúja, allelúja, allelúja.\end{Resp}
\begin{Resp}\Vsign Reges Tharsis et ínsulæ múnera ófferent, reges Arabum et Saba dona addúcent.\end{Resp}
\begin{Resp}\Rsign Allelúja, allelúja, allelúja.\end{Resp}
\switchcolumn
\EN
\begin{Resp}\Rsign All they from Saba shall come, they shall bring gold and incense, and they shall show forth the praises of the Lord. \Star{} Alleluia, Alleluia, Alleluia.\end{Resp}
\begin{Resp}\Vsign The kings of Tarshish and of the isles shall bring presents, the kings of Arabia and Saba shall offer gifts.\end{Resp}
\begin{Resp}\Rsign Alleluia, Alleluia, Alleluia.\end{Resp}
\switchcolumn*

\LA
\LessonHead{Lectio VI}
\switchcolumn
\EN
\LessonHead{Lesson VI}
\switchcolumn*

\LA
\noindent Tímuit enim eum rex Herodes, eisdem sibi Magis nuntiántibus, cum adhuc quærerent párvulum, quem cognóverant cælo teste jam natum. Quid erit tribúnal judicántis, quando supérbos reges cunábula terrébant infantis? Quanto consultius reges, non sicut Heródes interfícere quærant, sed sicut Magi potius adorare delecténtur; jam præsertim eum, qui et ipsam mortem, quam cupiebat inimícus inférre, étiam pro ipsis inimícis ab inimícis sustinuit, eamque in suo corpore occisus occídit. Pie timeant reges ad Patris déxteram jam sedéntem, quem rex ímpius tímuit adhuc matris úbera lambéntem.\par
\switchcolumn
\EN
\noindent Nor Herod feared when he heard from the wise men of Him Whom they sought, and of Whose birth they knew by the witness of a star. What will be the fearful judgment-seat of Him, Who, even as a Suckling, struck terror into haughty kings? How much wiser is the thought of those kings who seek Christ like the wise men, to worship Him, than of those who seek Him, like Herod, to slay Him! who seek to put Him to that same death, which He came to suffer from His.enemies for their own salvation, and which, by His death, He hath trodden down! Kings will do well to fear Him Who now sitteth at the right hand of the Father, and Whom Herod feared when He hung upon His mother's breast.\par
\switchcolumn*

\LA
\begin{Resp}\Rsign Magi véniunt ab Oriénte Jerosólymam, quæréntes, et dicéntes: Ubi est qui natus est, cujus stellam vídimus? \Star{} Et vénimus adoráre Dóminum.\end{Resp}
\begin{Resp}\Vsign Vídimus stellam ejus in Oriénte.\end{Resp}
\begin{Resp}\Rsign Et vénimus adoráre Dóminum.\end{Resp}
\begin{Resp}\Vsign Glória Patri, et Fílio, \Star{} et Spirítui Sancto.\end{Resp}
\begin{Resp}\Rsign Et vénimus adoráre Dóminum.\end{Resp}
\switchcolumn
\EN
\begin{Resp}\Rsign There came wise men from the east to Jerusalem, asking and saying: Where is He That is born King of the Jews? for we have seen His star in the east; \Star{} And are come to worship the Lord.\end{Resp}
\begin{Resp}\Vsign We have seen His star in the east.\end{Resp}
\begin{Resp}\Rsign And are come to worship the Lord.\end{Resp}
\begin{Resp}\Vsign Glory be to the Father, and to the Son, \Star{} and to the Holy Ghost.\end{Resp}
\begin{Resp}\Rsign And are come to worship the Lord.\end{Resp}
\switchcolumn*

\LA
\LessonHead{Lectio VII}
\switchcolumn
\EN
\LessonHead{Lesson VII}
\switchcolumn*

\LA
\LessonTitle{Léctio sancti Evangélii secúndum Matthǽum}
\switchcolumn
\EN
\LessonTitle{From the Holy Gospel according to Matthew}
\switchcolumn*

\LA
\Cite{Matt 2:1-12}
\switchcolumn
\EN
\Cite{Matt 2:1-12}
\switchcolumn*

\LA
\noindent Cum natus esset Jesus in Bethlehem Juda in diébus Herodis regis, ecce Magi ab Oriénte venérunt Jerosolymam, dicéntes: Ubi est qui natus est Rex Judæórum? Et réliqua.\par
\switchcolumn
\EN
\noindent When Jesus therefore was born in Bethlehem of Juda, in the days of king Herod, behold, there came wise men from the east to Jerusalem. Saying, Where is he that is born king of the Jews? And so on.\par
\switchcolumn*

\LA
\LessonTitle{De Homilía sancti Gregórii Papæ}
\switchcolumn
\EN
\LessonTitle{Homily by Pope St. Gregory (the Great)}
\switchcolumn*

\LA
\Source{Homilia 10 in Evangelia}
\switchcolumn
\EN
\Source{10th on the Gospels.}
\switchcolumn*

\LA
\noindent Nativitáte Regis nostri cógnita, Heródes ad callida arguménta convértitur; et ne terreno regno privarétur, renuntiari sibi ubi puer inveniátur, postulat. Adorare eum velle se símulat, ut hunc, si inveníri possit, exstinguat. Sed quanta est humana malítia contra consílium divinitátis? Scriptum quippe est: Non est sapiéntia, non est prudéntia, non est consílium contra Dóminum. Nam ea, quæ appáruit stella, Magos perducit: natum Regem reperiunt, múnera déferunt: et ne redire ad Heródem debeant, in somnis admonéntur. Sicque fit, ut Jesum, quem quærit Heródes, inveníre non possit. Cujus persona qui alii, quam hypócritæ designántur, qui dum ficte quærunt, inveníre Dóminum numquam merentur?\par
\switchcolumn
\EN
\noindent When Herod knew of the birth of our King, he betook himself to his cunning wiles, and lest he should be deprived of an earthly kingdom he desired the wise men to search diligently for the young Child, and when they had found Him, to bring him word again. He said, that he also might come and worship Him, but, in reality, that, when he had found Him, he might put Him to death. But, behold, of how light weight is the malice of man, when it is tried against the counsel of God. It is written There is no wisdom, nor understanding, nor counsel, against the Lord, (Prov. xxi. 30.) So the star which the wise men saw in the East, still led them on; they found the new-born King, and presented unto Him gifts; then they were warned in a dream that they should not return to Herod. And as it came to pass that, when Herod sought Jesus, he could not find Him even so is it with hypocrites, who, while they make pretence to seek the Lord to worship Him, find Him not.\par
\switchcolumn*

\LA
\begin{Resp}\Rsign Stella, quam víderant Magi in Oriénte, antecedébat eos, donec venírent ad locum, ubi puer erat. \Star{} Vidéntes autem eam, gavísi sunt gáudio magno.\end{Resp}
\begin{Resp}\Vsign Et intrántes domum, invenérunt púerum cum María matre ejus, et procidéntes adoravérunt eum.\end{Resp}
\begin{Resp}\Rsign Vidéntes autem eam, gavísi sunt gáudio magno.\end{Resp}
\switchcolumn
\EN
\begin{Resp}\Rsign The star which the wise men had seen in the East, went before them, till they came where the young Child was. \Star{} And when they saw the star, they rejoiced with exceeding great joy.\end{Resp}
\begin{Resp}\Vsign And when they were come into the house, they found the young Child with Mary His Mother, and fell down and worshipped Him.\end{Resp}
\begin{Resp}\Rsign And when they saw the star, they rejoiced with exceeding great joy.\end{Resp}
\switchcolumn*

\LA
\LessonHead{Lectio VIII}
\switchcolumn
\EN
\LessonHead{Lesson VIII}
\switchcolumn*

\LA
\noindent Sed inter hæc sciéndum, quod Priscilliánistæ hæretici nasci unumquemque hóminem sub constitutiónibus stellárum putant: et hoc in adjutórium sui erroris assumunt, quod nova stella éxiit cum Dóminus in carne appáruit; cujus fuísse fatum eamdem, quæ appáruit, stellam putant. Sed si Evangélii verba pensamus, quibus de eádem stella dícitur: Usque dum veniens staret supra, ubi erat puer: dum non puer ad stellam, sed stella ad púerum cucúrrit, si dici liceat, non stella fatum púeri, sed fatum stellæ is, qui appáruit puer fuit.\par
\switchcolumn
\EN
\noindent It is as well to know that it is one of the opinions of the Priscillianist heretics l that every man is born under the influence of a star; and, to confirm this notion, they bring forward the instance of the star of Bethlehem, which appeared when the Lord was born; and which they call His star, that is, the star ruling over His fate or destiny. But if we consider the words of the Gospel concerning this star, they are It went before, till it came and stood over where the young Child was. Whence we see that it was not the young Child Who followed the star, but the star which followed the young Child, as if to show that the young Child ruled the star, instead of the star ruling Him.\par
\switchcolumn*

\LA
\begin{Resp}\Rsign Vidéntes stellam Magi, gavísi sunt gáudio magno: \Star{} Et intrántes domum invenérunt púerum cum María matre ejus, et procidéntes adoravérunt eum: * Et apértis thesáuris suis obtulérunt ei múnera, aurum, thus, et myrrham.\end{Resp}
\begin{Resp}\Vsign Stella, quam víderant Magi in Oriénte, antecedébat eos, usque dum véniens staret supra ubi erat puer.\end{Resp}
\begin{Resp}\Rsign Et intrántes domum invenérunt púerum cum María matre ejus, et procidéntes adoravérunt eum:\end{Resp}
\begin{Resp}\Vsign Glória Patri, et Fílio, \Star{} et Spirítui Sancto.\end{Resp}
\begin{Resp}\Rsign Et apértis thesáuris suis obtulérunt ei múnera, aurum, thus, et myrrham.\end{Resp}
\switchcolumn
\EN
\begin{Resp}\Rsign When the wise men saw the star, they rejoiced with exceeding great joy. and when they had opened their treasures, they presented unto Him gifts; gold, and frankincense, and myrrh. \Star{} And when they were come into the house, they found the young Child with Mary His Mother, and fell down and worshipped Him; *\end{Resp}
\begin{Resp}\Vsign The star which the wise men had seen in the East, went before them, till it came and stood over where the young Child was.\end{Resp}
\begin{Resp}\Rsign And when they were come into the house, they found the young Child with Mary His Mother, and fell down and worshipped Him.\end{Resp}
\begin{Resp}\Vsign Glory be to the Father, and to the Son, \Star{} and to the Holy Ghost.\end{Resp}
\begin{Resp}\Rsign And when they had opened their treasures, they presented unto Him gifts; gold, and frankincense, and myrrh.\end{Resp}
\switchcolumn*

\LA
\LessonHead{Lectio IX}
\switchcolumn
\EN
\LessonHead{Lesson IX}
\switchcolumn*

\LA
\noindent Sed a fidelium córdibus absit, ut aliquid esse fatum dicant. Vitam quippe hóminum solus hic Cónditor, qui creávit, adminístrat. Neque enim propter stellas homo, sed stellæ propter hóminem factæ sunt: et si stella fatum hóminis dícitur, ipsis suis ministériis subesse homo perhibétur. Certe cum Jacob de útero egrédiens, prioris fratris plantam tenéret manu, prior perfecte nequáquam égredi pótuit, nisi súbsequens inchoasset: et tamen cum uno témpore, eodémque momento utrumque mater fúderit, non una utriusque vitæ qualitas fuit.\par
\switchcolumn
\EN
\noindent But I pray that the hearts of the faithful may ever be free from the thought that anything ruleth their destiny. In this world there is but One Who ruleth the destiny of man, even He Who made man; neither was man made for the stars, but the stars for man; and if we say that they rule his destiny, we set them above him for whose service they were made. When Jacob came out of his mother's womb, and his hand took hold on his elder brother Esau's heel, he could not have done so unless this his first movement had been behind his brother, and, nevertheless, such was not in after life the position of those two brethren whom their mother brought forth at one birth.\par
\switchcolumn*

\LA
\TeDeum{Te Deum}
\switchcolumn
\EN
\TeDeum{Te Deum}
\switchcolumn*

\end{paracol}

\par\needspace{10\baselineskip}\addvspace{1.2em}{\centering\small\textsc{Die 8 Januarii} · \textsc{8 January}\par}
\Office{Tertia die infra Octavam Epiphaniæ}{Tertia die infra Octavam Epiphaniæ}{Semiduplex II classis}
\addcontentsline{toc}{subsection}{Die 8 Januarii · Tertia die infra Octavam Epiphaniæ\enspace\textit{Tertia die infra Octavam Epiphaniæ}}
\begin{paracol}{2}
\LA
\RefLine{Lectiones I–III: de Scriptura occurrente.}
\switchcolumn
\EN
\RefLine{Lessons I–III: from the Scripture of the day.}
\switchcolumn*

\LA
\LessonHead{Lectio IV}
\switchcolumn
\EN
\LessonHead{Lesson IV}
\switchcolumn*

\LA
\LessonTitle{De Sermóne sancti Augustíni Epíscopi}
\switchcolumn
\EN
\LessonTitle{From the Sermons of St. Augustine, Bishop (of Hippo.)}
\switchcolumn*

\LA
\Source{Sermo 2 de Epiphania}
\switchcolumn
\EN
\Source{2nd on the Epiphany.}
\switchcolumn*

\LA
\noindent Cum tam multi jam nati, atque defuncti essent reges Judæórum, numquid quemquam eórum adorándum Magi quæsiérunt? Non, quia nec quemquam eórum cælo loquénte didicérunt. Verúmtamen, quod prætereundum non est, hæc Magórum illuminátio magnum testimónium cæcitátis exstitit Judæórum. In terra eórum isti requirébant, quem illi in sua non agnoscebant.\par
\switchcolumn
\EN
\noindent Many kings of the Jews had been born, and died, but which of them was sought after by wise men to worship him? Not one for not one had been proclaimed by the voice of heaven. Let us not also pass by the fact that the enlightenment of the wise men standeth in strong contrast to the blindness of the Jews. The first came from far to find Him, Whom, being born in their midst, the second knew not.\par
\switchcolumn*

\LA
\begin{Resp}\Rsign Illumináre, illumináre Jerúsalem, quia venit lux tua: \Star{} Et glória Dómini super te orta est.\end{Resp}
\begin{Resp}\Vsign Et ambulábunt gentes in lúmine tuo, et reges in splendóre ortus tui.\end{Resp}
\begin{Resp}\Rsign Et glória Dómini super te orta est.\end{Resp}
\switchcolumn
\EN
\begin{Resp}\Rsign Shine, shine, O Jerusalem, for thy light is come; \Star{} And the glory of the Lord is risen upon thee.\end{Resp}
\begin{Resp}\Vsign And the Gentiles shall walk in thy light, and kings in the brightness of thy rising.\end{Resp}
\begin{Resp}\Rsign And the glory of the Lord is risen upon thee.\end{Resp}
\switchcolumn*

\LA
\LessonHead{Lectio V}
\switchcolumn
\EN
\LessonHead{Lesson V}
\switchcolumn*

\LA
\noindent Apud eos isti infántem invenérunt, quem illi apud se negavérunt. In his terris de longinquo isti et peregrini púerum Christum nondum verba proméntem adoravérunt, ubi cives illi juvenem miracula faciéntem crucifixérunt. Isti in membris parvis Deum adoravérunt, illi in magnis factis nec tamquam hómini pepercérunt: quasi plus fúerit vidére novam stellam in ejus nativitáte fulgéntem, quam solem ejus in morte lugéntem.\par
\switchcolumn
\EN
\noindent The wise men found the young Child among those who denied Him. These holy pilgrims came and worshipped the yet silent Christ in the land whose inhabitants, after He grew up and worked miracles, crucified Him. They worshipped in that tiny Body the God Whom, amid great signs and wonders, the Jews would not spare even as a man. They who saw the Star which shone at His birth, put it to more profit than they who saw the sun veiled at His death.\par
\switchcolumn*

\LA
\begin{Resp}\Rsign Omnes de Saba vénient, aurum et thus deferéntes, et laudem Dómino annuntiántes, \Star{} Allelúja, allelúja, allelúja.\end{Resp}
\begin{Resp}\Vsign Reges Tharsis et ínsulæ múnera ófferent, reges Arabum et Saba dona addúcent.\end{Resp}
\begin{Resp}\Rsign Allelúja, allelúja, allelúja.\end{Resp}
\switchcolumn
\EN
\begin{Resp}\Rsign All they from Saba shall come, they shall bring gold and incense, and they shall show forth the praises of the Lord. \Star{} Alleluia, Alleluia, Alleluia.\end{Resp}
\begin{Resp}\Vsign The kings of Tarshish and of the isles shall bring presents, the kings of Arabia and Saba shall offer gifts.\end{Resp}
\begin{Resp}\Rsign Alleluia, Alleluia, Alleluia.\end{Resp}
\switchcolumn*

\LA
\LessonHead{Lectio VI}
\switchcolumn
\EN
\LessonHead{Lesson VI}
\switchcolumn*

\LA
\noindent Jam vero quod éadem stella, quæ Magos perdúxit ad locum, ubi erat cum matre Vírgine Deus infans, quæ útique póterat eos ad ipsam perducere civitátem, se tamen subtraxit, nec eis prorsus appáruit, donec de civitáte, in qua Christus nascerétur, iídem ipsi interrogaréntur Judæi, et ipsi eam secúndum divinæ Scripturæ testimónium nominárent, ipsi dicerent: In Bethlehem Judæ. Sic enim scriptum est: Et tu Bethlehem terra Juda, non es minima in princípibus Juda: ex te enim éxiet Dux, qui regat pópulum meum Israël. Quid aliud hic significávit divina providéntia, nisi apud Judæos solas divinas litteras remansuras, quibus gentes instrueréntur, illi excæcarentur?\par
\switchcolumn
\EN
\noindent The star which led the wise men towards the place where the new-born God dwelt with His Virgin Mother, ceased to shine when it came to the city of Jerusalem, while they were inquiring of the Jews where Christ should be born. The Jews answered them according to the testimony of the Divine Scriptures In Bethlehem of Judah for thus it is written And thou Bethlehem in the land of Judah, art not the least among the princes of Judah, for out of thee shall come a Governor, that shall rule My people Israel. What else are we to understand that God's Providence would here signify, than that there should remain among the Jews those Divine Writings only, whereby the Gentiles are enlightened, while they themselves remain dark?\par
\switchcolumn*

\LA
\begin{Resp}\Rsign Magi véniunt ab Oriénte Jerosólymam, quæréntes, et dicéntes: Ubi est qui natus est, cujus stellam vídimus? \Star{} Et vénimus adoráre Dóminum.\end{Resp}
\begin{Resp}\Vsign Vídimus stellam ejus in Oriénte.\end{Resp}
\begin{Resp}\Rsign Et vénimus adoráre Dóminum.\end{Resp}
\begin{Resp}\Vsign Glória Patri, et Fílio, \Star{} et Spirítui Sancto.\end{Resp}
\begin{Resp}\Rsign Et vénimus adoráre Dóminum.\end{Resp}
\switchcolumn
\EN
\begin{Resp}\Rsign There came wise men from the east to Jerusalem, asking and saying: Where is He That is born King of the Jews? for we have seen His star in the east; \Star{} And are come to worship the Lord.\end{Resp}
\begin{Resp}\Vsign We have seen His star in the east.\end{Resp}
\begin{Resp}\Rsign And are come to worship the Lord.\end{Resp}
\begin{Resp}\Vsign Glory be to the Father, and to the Son, \Star{} and to the Holy Ghost.\end{Resp}
\begin{Resp}\Rsign And are come to worship the Lord.\end{Resp}
\switchcolumn*

\LA
\LessonHead{Lectio VII}
\switchcolumn
\EN
\LessonHead{Lesson VII}
\switchcolumn*

\LA
\LessonTitle{Léctio sancti Evangélii secúndum Matthǽum}
\switchcolumn
\EN
\LessonTitle{From the Holy Gospel according to Matthew}
\switchcolumn*

\LA
\Cite{Matt 2:1-12}
\switchcolumn
\EN
\Cite{Matt 2:1-12}
\switchcolumn*

\LA
\noindent Cum natus esset Jesus in Béthlehem Juda in diébus Herodis regis, ecce Magi ab Oriénte venérunt Jerosolymam, dicéntes: Ubi est qui natus est Rex Judæórum? Et réliqua.\par
\switchcolumn
\EN
\noindent When Jesus therefore was born in Bethlehem of Juda, in the days of king Herod, behold, there came wise men from the east to Jerusalem. Saying, Where is he that is born king of the Jews? And so on.\par
\switchcolumn*

\LA
\LessonTitle{De Homilía sancti Gregórii Papæ}
\switchcolumn
\EN
\LessonTitle{Homily by Pope St. Gregory (the Great.)}
\switchcolumn*

\LA
\Source{Ex Homilía 10 in Evangelia}
\switchcolumn
\EN
\Source{10th on the Gospels.}
\switchcolumn*

\LA
\noindent Magi vero aurum, thus et myrrham déferunt. Aurum quippe Regi congruit, thus vero in Dei sacrifícium pónitur, myrrha autem mortuórum córpora condiúntur. Eum ergo Magi, quem adorant, étiam mysticis munéribus prædicant: auro Regem, thure Deum, myrrha mortalem. Sunt vero nonnulli hæretici, qui hunc Deum credunt, sed ubíque regnare nequáquam credunt. Hi profecto ei thus ófferunt, sed offerre étiam aurum nolunt. Et sunt nonnulli, qui hunc Regem existimant, sed Deum negant. Hi videlicet ei aurum ófferunt, sed offerre thus nolunt.\par
\switchcolumn
\EN
\noindent The wise men brought gold, frankincense, and myrrh. Gold is the fitting gift to a King, frankincense is offered in sacrifice to God, and with myrrh are embalmed the bodies of the dead. By the gifts, therefore, which they presented unto Him, the wise men set forth three things concerning Him unto Whom they offered them; by the gold, that He was King; by the frankincense, that He was God; and by the myrrh, that He was to die. There are some heretics who believe Him to be God, but confess not His Kingly dominion over all things; these offer unto Him frankincense, but refuse Him gold. There are some others who admit that He is King, but deny that He is God; these present unto Him gold, but will not give Him frankincense.\par
\switchcolumn*

\LA
\begin{Resp}\Rsign Stella, quam víderant Magi in Oriénte, antecedébat eos, donec venírent ad locum, ubi puer erat. \Star{} Vidéntes autem eam, gavísi sunt gáudio magno.\end{Resp}
\begin{Resp}\Vsign Et intrántes domum, invenérunt púerum cum María matre ejus, et procidéntes adoravérunt eum.\end{Resp}
\begin{Resp}\Rsign Vidéntes autem eam, gavísi sunt gáudio magno.\end{Resp}
\switchcolumn
\EN
\begin{Resp}\Rsign The star which the wise men had seen in the East, went before them, till they came where the young Child was. \Star{} And when they saw the star, they rejoiced with exceeding great joy.\end{Resp}
\begin{Resp}\Vsign And when they were come into the house, they found the young Child with Mary His Mother, and fell down and worshipped Him.\end{Resp}
\begin{Resp}\Rsign And when they saw the star, they rejoiced with exceeding great joy.\end{Resp}
\switchcolumn*

\LA
\LessonHead{Lectio VIII}
\switchcolumn
\EN
\LessonHead{Lesson VIII}
\switchcolumn*

\LA
\noindent Et sunt nonnulli, qui hunc et Deum et Regem faténtur, sed assumpsísse carnem mortalem negant. Hi nimirum ei aurum et thus ófferunt, sed offerre myrrham assumptæ mortalitátis nolunt. Nos ítaque nato Dómino offerámus aurum, ut hunc ubíque regnare fateámur: offerámus thus, ut credamus, quod is, qui in témpore appáruit, Deus ante témpora éxstitit: offerámus myrrham, ut eum, quem crédimus in sua divinitate impassibilem, credámus étiam in nostra fuísse carne mortalem.\par
\switchcolumn
\EN
\noindent There are some other heretics who profess that Christ is both God and King, but not that He took a dying nature; these offer Him gold and frankincense, but not myrrh for the Manhood. Let us, however, present gold unto the new-born Lord, acknowledging His universal Kingship; let us offer unto Him frankincense, confessing that He Who hath been made manifest unto us in time, is God before time was; let us give unto Him myrrh, believing that He Who cannot suffer as touching His Godhead, was made capable of death as touching the manhood which He shareth with us.\par
\switchcolumn*

\LA
\begin{Resp}\Rsign Vidéntes stellam Magi, gavísi sunt gáudio magno: \Star{} Et intrántes domum invenérunt púerum cum María matre ejus, et procidéntes adoravérunt eum: * Et apértis thesáuris suis obtulérunt ei múnera, aurum, thus, et myrrham.\end{Resp}
\begin{Resp}\Vsign Stella, quam víderant Magi in Oriénte, antecedébat eos, usque dum véniens staret supra ubi erat puer.\end{Resp}
\begin{Resp}\Rsign Et intrántes domum invenérunt púerum cum María matre ejus, et procidéntes adoravérunt eum:\end{Resp}
\begin{Resp}\Vsign Glória Patri, et Fílio, \Star{} et Spirítui Sancto.\end{Resp}
\begin{Resp}\Rsign Et apértis thesáuris suis obtulérunt ei múnera, aurum, thus, et myrrham.\end{Resp}
\switchcolumn
\EN
\begin{Resp}\Rsign When the wise men saw the star, they rejoiced with exceeding great joy. and when they had opened their treasures, they presented unto Him gifts; gold, and frankincense, and myrrh. \Star{} And when they were come into the house, they found the young Child with Mary His Mother, and fell down and worshipped Him; *\end{Resp}
\begin{Resp}\Vsign The star which the wise men had seen in the East, went before them, till it came and stood over where the young Child was.\end{Resp}
\begin{Resp}\Rsign And when they were come into the house, they found the young Child with Mary His Mother, and fell down and worshipped Him.\end{Resp}
\begin{Resp}\Vsign Glory be to the Father, and to the Son, \Star{} and to the Holy Ghost.\end{Resp}
\begin{Resp}\Rsign And when they had opened their treasures, they presented unto Him gifts; gold, and frankincense, and myrrh.\end{Resp}
\switchcolumn*

\LA
\LessonHead{Lectio IX}
\switchcolumn
\EN
\LessonHead{Lesson IX}
\switchcolumn*

\LA
\noindent Quamvis in auro, thure, et myrrha intelligi et aliud potest. Auro namque sapiéntia designátur, Salomone attestante, qui ait: Thesaurus desiderábilis requiescit in ore sapiéntis. Thure autem, quod Deo incénditur, virtus oratiónis exprímitur, Psalmista testante, qui dicit: Dirigátur orátio mea sicut incensum in conspéctu tuo. Per myrrham vero carnis nostræ mortificátio figurátur. Unde sancta Ecclésia de suis operariis usque ad mortem pro Deo certántibus dicit: Manus meæ distillavérunt myrrham.\par
\switchcolumn
\EN
\noindent There is also another signification in this gold, frankincense, and myrrh. Gold is a type of wisdom; as saith Solomon x In the mouth of the wise abideth a treasure to be desired. Frankincense, which is burnt in honour of God, is a figure of prayer; witness the words of the Psalmist, (Ps. cxl. 2:) Let my prayer be set forth as incense before thee. By myrrh is represented the putting to death of the body; as where the holy Church saith of her labourers who strive for God even unto death My hands dropped with myrrh. (Cant. v. 5.)\par
\switchcolumn*

\LA
\TeDeum{Te Deum}
\switchcolumn
\EN
\TeDeum{Te Deum}
\switchcolumn*

\end{paracol}

\par\needspace{10\baselineskip}\addvspace{1.2em}{\centering\small\textsc{Die 9 Januarii} · \textsc{9 January}\par}
\Office{Quarta die infra Octavam Epiphaniæ}{Quarta die infra Octavam Epiphaniæ}{Semiduplex II classis}
\addcontentsline{toc}{subsection}{Die 9 Januarii · Quarta die infra Octavam Epiphaniæ\enspace\textit{Quarta die infra Octavam Epiphaniæ}}
\begin{paracol}{2}
\LA
\RefLine{Lectiones I–III: de Scriptura occurrente.}
\switchcolumn
\EN
\RefLine{Lessons I–III: from the Scripture of the day.}
\switchcolumn*

\LA
\LessonHead{Lectio IV}
\switchcolumn
\EN
\LessonHead{Lesson IV}
\switchcolumn*

\LA
\LessonTitle{Sermo sancti Leónis Papæ}
\switchcolumn
\EN
\LessonTitle{From the Sermons of Pope St. Leo (the Great)}
\switchcolumn*

\LA
\Source{Sermo 1 de Epiphania}
\switchcolumn
\EN
\Source{1 st for Twelfth-Day.}
\switchcolumn*

\LA
\noindent Celebrato próximo die, quo intemeráta virginitas humani generis edidit Salvatórem, Epiphaníæ nobis, dilectíssimi, veneranda festívitas dat perseverantiam gaudiórum: ut inter cognatárum solemnitátum vicina sacramenta, exsultatiónis vigor, et fervor fidei non tepescat. Ad ómnium enim hóminum spectat salútem, quod infántia mediatoris Dei et hóminum, jam universo declarabátur mundo, cum adhuc exíguo detinerétur oppídulo.\par
\switchcolumn
\EN
\noindent Dearly beloved brethren, we have but lately celebrated that day whereon the inviolate virginity of Blessed Mary gave to man a Saviour. And now the venerable solemnity of the Epiphany giveth us a continuance of joy. So that by the nearness of these two holy Feasts, the freshness of our gladness and the quickening of our faith hath no time wherein to die away. And truly it concerneth the salvation of all men, that the Mediator between God and men is already made manifest before leaving the humble city of His birth.\par
\switchcolumn*

\LA
\begin{Resp}\Rsign Illumináre, illumináre Jerúsalem, quia venit lux tua: \Star{} Et glória Dómini super te orta est.\end{Resp}
\begin{Resp}\Vsign Et ambulábunt gentes in lúmine tuo, et reges in splendóre ortus tui.\end{Resp}
\begin{Resp}\Rsign Et glória Dómini super te orta est.\end{Resp}
\switchcolumn
\EN
\begin{Resp}\Rsign Shine, shine, O Jerusalem, for thy light is come; \Star{} And the glory of the Lord is risen upon thee.\end{Resp}
\begin{Resp}\Vsign And the Gentiles shall walk in thy light, and kings in the brightness of thy rising.\end{Resp}
\begin{Resp}\Rsign And the glory of the Lord is risen upon thee.\end{Resp}
\switchcolumn*

\LA
\LessonHead{Lectio V}
\switchcolumn
\EN
\LessonHead{Lesson V}
\switchcolumn*

\LA
\noindent Quamvis enim Israëliticam gentem, et ipsíus gentis unam familiam delegísset, de qua natúram univérsæ humanitátis assumeret: nóluit tamen intra maternæ habitatiónis angustias ortus sui latére primordia: sed mox ab ómnibus vóluit agnosci, qui dignátus est pro ómnibus nasci. Tribus ígitur Magis in regióne Oriéntis stella novæ claritátis appáruit, quæ illustrior ceteris pulchriorque sidéribus, fácile in se intuéntium óculos animosque convérteret: ut conféstim adverterétur, non esse otiosum, quod tam insólitum videbátur.\par
\switchcolumn
\EN
\noindent It is true that the Lord chose the nation of Israel, and in that nation one family, whence to take upon Him that nature which He shareth with all mankind, but, at the same time, He would not that the narrow walls of His Mother's house should imprison within them all the brightness of His appearing, and, as He was pleased to be born for all, so willed He to be forthwith made manifest to all. Three wise men in the East, therefore, saw a new and brilliant star, which, by excelling all others in brightness and beauty, attracted the eyes and thoughts of all beholders and thereby it became at once evident that some new and great event had befallen.\par
\switchcolumn*

\LA
\begin{Resp}\Rsign Omnes de Saba vénient, aurum et thus deferéntes, et laudem Dómino annuntiántes, \Star{} Allelúja, allelúja, allelúja.\end{Resp}
\begin{Resp}\Vsign Reges Tharsis et ínsulæ múnera ófferent, reges Arabum et Saba dona addúcent.\end{Resp}
\begin{Resp}\Rsign Allelúja, allelúja, allelúja.\end{Resp}
\switchcolumn
\EN
\begin{Resp}\Rsign All they from Saba shall come, they shall bring gold and incense, and they shall show forth the praises of the Lord. \Star{} Alleluia, Alleluia, Alleluia.\end{Resp}
\begin{Resp}\Vsign The kings of Tarshish and of the isles shall bring presents, the kings of Arabia and Saba shall offer gifts.\end{Resp}
\begin{Resp}\Rsign Alleluia, Alleluia, Alleluia.\end{Resp}
\switchcolumn*

\LA
\LessonHead{Lectio VI}
\switchcolumn
\EN
\LessonHead{Lesson VI}
\switchcolumn*

\LA
\noindent Dedit ergo aspiciéntibus intelléctum, qui præstitit signum: et quod fecit intelligi, fecit inquiri, et se inveniéndum óbtulit requisítus. Sequúntur tres viri superni lúminis ductum, et prævii fulgoris indícium inténta contemplatióne comitántes, ad agnitiónem veritátis, gratiæ splendore ducúntur: qui humano sensu significátum sibi Regis ortum æstimavérunt in civitáte regia esse quæréndum. Sed qui servi susceperat formam, et non judicáre vénerat, sed judicári, Bethlehem præelégit nativitáti, Jerosolymam passióni.\par
\switchcolumn
\EN
\noindent Then He Who had given the sign, gave understanding to those that saw it; and having given to them to understand that He was born, He gave them the grace to seek Him; and, being sought by them, was pleased to be found. The three wise men followed the guiding of the heavenly light, and, with their eyes firmly fixed upon the glory that went before them, were so led by the light of grace as to obtain the knowledge of truth. They, knowing that He was born a King, sought Him in the Royal City; but He Who had taken upon Him the form of a servant, and came not to judge but to be judged, had chosen Bethlehem for His birth, and Jerusalem for His Suffering.\par
\switchcolumn*

\LA
\begin{Resp}\Rsign Magi véniunt ab Oriénte Jerosólymam, quæréntes, et dicéntes: Ubi est qui natus est, cujus stellam vídimus? \Star{} Et vénimus adoráre Dóminum.\end{Resp}
\begin{Resp}\Vsign Vídimus stellam ejus in Oriénte.\end{Resp}
\begin{Resp}\Rsign Et vénimus adoráre Dóminum.\end{Resp}
\begin{Resp}\Vsign Glória Patri, et Fílio, \Star{} et Spirítui Sancto.\end{Resp}
\begin{Resp}\Rsign Et vénimus adoráre Dóminum.\end{Resp}
\switchcolumn
\EN
\begin{Resp}\Rsign There came wise men from the east to Jerusalem, asking and saying: Where is He That is born King of the Jews? for we have seen His star in the east; \Star{} And are come to worship the Lord.\end{Resp}
\begin{Resp}\Vsign We have seen His star in the east.\end{Resp}
\begin{Resp}\Rsign And are come to worship the Lord.\end{Resp}
\begin{Resp}\Vsign Glory be to the Father, and to the Son, \Star{} and to the Holy Ghost.\end{Resp}
\begin{Resp}\Rsign And are come to worship the Lord.\end{Resp}
\switchcolumn*

\LA
\LessonHead{Lectio VII}
\switchcolumn
\EN
\LessonHead{Lesson VII}
\switchcolumn*

\LA
\LessonTitle{Léctio sancti Evangélii secúndum Matthǽum}
\switchcolumn
\EN
\LessonTitle{From the Holy Gospel according to Matthew}
\switchcolumn*

\LA
\Cite{Matt 2:1-12}
\switchcolumn
\EN
\Cite{Matt 2:1-12}
\switchcolumn*

\LA
\noindent Cum natus esset Jesus in Bethlehem Juda in diébus Herodis regis, ecce Magi ab Oriénte venérunt Jerosolymam, dicéntes: Ubi est qui natus est Rex Judæórum? Et réliqua.\par
\switchcolumn
\EN
\noindent When Jesus therefore was born in Bethlehem of Juda, in the days of king Herod, behold, there came wise men from the east to Jerusalem. Saying, Where is he that is born king of the Jews? And so on.\par
\switchcolumn*

\LA
\LessonTitle{De Homilía sancti Gregórii Papæ}
\switchcolumn
\EN
\LessonTitle{Homily by Pope St. Gregory (the Great.)}
\switchcolumn*

\LA
\Source{Ex Homilía 10 in Evangelia}
\switchcolumn
\EN
\Source{10th on the Gospels.}
\switchcolumn*

\LA
\noindent Magnum vero nobis aliquid Magi ínnuunt, quod in regiónem suam per aliam viam revertúntur. In eo namque quod admóniti faciunt, nobis profecto insinuant quid faciámus. Regio quippe nostra paradisus est: ad quam, Jesu cógnito, redire per viam, qua vénimus, prohibemur. A regióne étenim nostra superbiéndo, inobediendo, visibília sequéndo, cibum vétitum gustando discessimus: sed ad eam necesse est, ut flendo, obediendo, visibília contemnendo, atque appetítum carnis refrenando redeamus.\par
\switchcolumn
\EN
\noindent The wise men teach us a great lesson in that they departed into their own country another way. That which they did, being warned of God in a dream, we ought to do. Our country is heaven; and, when we have once known Jesus, we can never get there by returning on the way wherein we walked before we knew Him. We have left our country far, by the way of pride, and disobedience, and worldliness, and forbidden indulgence we must seek that heavenly Fatherland by tears, by subjection, by contempt of the things which are seen, and by curbing the fleshly appetites.\par
\switchcolumn*

\LA
\begin{Resp}\Rsign Stella, quam víderant Magi in Oriénte, antecedébat eos, donec venírent ad locum, ubi puer erat. \Star{} Vidéntes autem eam, gavísi sunt gáudio magno.\end{Resp}
\begin{Resp}\Vsign Et intrántes domum, invenérunt púerum cum María matre ejus, et procidéntes adoravérunt eum.\end{Resp}
\begin{Resp}\Rsign Vidéntes autem eam, gavísi sunt gáudio magno.\end{Resp}
\switchcolumn
\EN
\begin{Resp}\Rsign The star which the wise men had seen in the East, went before them, till they came where the young Child was. \Star{} And when they saw the star, they rejoiced with exceeding great joy.\end{Resp}
\begin{Resp}\Vsign And when they were come into the house, they found the young Child with Mary His Mother, and fell down and worshipped Him.\end{Resp}
\begin{Resp}\Rsign And when they saw the star, they rejoiced with exceeding great joy.\end{Resp}
\switchcolumn*

\LA
\LessonHead{Lectio VIII}
\switchcolumn
\EN
\LessonHead{Lesson VIII}
\switchcolumn*

\LA
\noindent Per aliam ergo viam ad regiónem nostram regrédimur: quóniam qui a paradisi gaudiis per delectaménta discessimus, ad hæc per laménta revocamur. Unde necesse est, fratres caríssimi, ut semper pávidi, semperque suspecti, ponámus ante óculos cordis, hinc culpas operis, illinc judícium extremæ districtiónis. Pensemus, quam districtus judex véniat, qui judícium minátur, et latet: terróres peccatóribus intentat, et tamen adhuc sústinet: et idcirco veníre citius differt, ut minus invéniat quod condemnet.\par
\switchcolumn
\EN
\noindent Let us then depart into our own country another way. They that have by enjoyment put themselves away from it, must seek it again by sorrow. Therefore, my dearly beloved brethren, it behoveth us to be ever fearful and watch, having continually before the eyes of our heart, on the one hand, the guilt of our doings, and, on the other, the judgment at the latter day. It behoveth us to think how that awful Judge will surely come, Whose judgment is hanging over us, and hath not yet fallen the wrath to come is before sinners, and hath not yet smitten them and the Judge yet tarrieth in order that, when He cometh, there may haply be less to condemn.\par
\switchcolumn*

\LA
\begin{Resp}\Rsign Vidéntes stellam Magi, gavísi sunt gáudio magno: \Star{} Et intrántes domum invenérunt púerum cum María matre ejus, et procidéntes adoravérunt eum: * Et apértis thesáuris suis obtulérunt ei múnera, aurum, thus, et myrrham.\end{Resp}
\begin{Resp}\Vsign Stella, quam víderant Magi in Oriénte, antecedébat eos, usque dum véniens staret supra ubi erat puer.\end{Resp}
\begin{Resp}\Rsign Et intrántes domum invenérunt púerum cum María matre ejus, et procidéntes adoravérunt eum:\end{Resp}
\begin{Resp}\Vsign Glória Patri, et Fílio, \Star{} et Spirítui Sancto.\end{Resp}
\begin{Resp}\Rsign Et apértis thesáuris suis obtulérunt ei múnera, aurum, thus, et myrrham.\end{Resp}
\switchcolumn
\EN
\begin{Resp}\Rsign When the wise men saw the star, they rejoiced with exceeding great joy. and when they had opened their treasures, they presented unto Him gifts; gold, and frankincense, and myrrh. \Star{} And when they were come into the house, they found the young Child with Mary His Mother, and fell down and worshipped Him; *\end{Resp}
\begin{Resp}\Vsign The star which the wise men had seen in the East, went before them, till it came and stood over where the young Child was.\end{Resp}
\begin{Resp}\Rsign And when they were come into the house, they found the young Child with Mary His Mother, and fell down and worshipped Him.\end{Resp}
\begin{Resp}\Vsign Glory be to the Father, and to the Son, \Star{} and to the Holy Ghost.\end{Resp}
\begin{Resp}\Rsign And when they had opened their treasures, they presented unto Him gifts; gold, and frankincense, and myrrh.\end{Resp}
\switchcolumn*

\LA
\LessonHead{Lectio IX}
\switchcolumn
\EN
\LessonHead{Lesson IX}
\switchcolumn*

\LA
\noindent Puniámus fletibus culpas, et cum Psalmistæ voce, præveniámus fáciem ejus in confessione. Voluptátum nos ergo fallacia nulla decipiat, nulla vana lætítia seducat. In próximo namque est judex, qui dixit: Væ vobis, qui ridetis nunc; quia lugébitis et flébitis. Hinc enim Sálomon ait: Risus dolori miscébitur: et, Extrema gaudii luctus óccupat. Hinc íterum dicit: Risum deputávi errórem, et gáudio dixi, Quid frustra decíperis? Hinc rursus ait: Cor sapiéntium ubi tristítia est: et cor stultórum ubi lætítia.\par
\switchcolumn
\EN
\noindent Let us afflict ourselves for our faults with weeping, and, with the Psalmist, let us come before His Presence with thanksgiving. Let us take heed that we be not fooled by the appearance of earthly happiness, or seduced by the vanity of earthly pleasure. For the Judge is at hand, even He That saith Woe unto you that laugh now, for ye shall mourn and weep, (Luke vi. 25.) Hence also Solomon saith Even in laughter the heart is sorrowful; and the end of that mirth is heaviness, (Prov. xiv. 13.) And again: I said of laughter, It is mad; and of mirth, What doeth it? (Eccles. ii. 2.) And yet again The heart of the wise is in the house of mourning, but the heart of fools is in the house of mirth, (vii. 5.)\par
\switchcolumn*

\LA
\TeDeum{Te Deum}
\switchcolumn
\EN
\TeDeum{Te Deum}
\switchcolumn*

\end{paracol}

\par\needspace{10\baselineskip}\addvspace{1.2em}{\centering\small\textsc{Die 10 Januarii} · \textsc{10 January}\par}
\Office{Quinta die infra Octavam Epiphaniæ}{Quinta die infra Octavam Epiphaniæ}{Semiduplex II classis}
\addcontentsline{toc}{subsection}{Die 10 Januarii · Quinta die infra Octavam Epiphaniæ\enspace\textit{Quinta die infra Octavam Epiphaniæ}}
\begin{paracol}{2}
\LA
\RefLine{Lectiones I–III: de Scriptura occurrente.}
\switchcolumn
\EN
\RefLine{Lessons I–III: from the Scripture of the day.}
\switchcolumn*

\LA
\LessonHead{Lectio IV}
\switchcolumn
\EN
\LessonHead{Lesson IV}
\switchcolumn*

\LA
\LessonTitle{Sermo sancti Máximi Epíscopi}
\switchcolumn
\EN
\LessonTitle{From the Sermons of St. Maximus, Bishop (of Turin.)}
\switchcolumn*

\LA
\Source{Homilia 1 de Epiphania}
\switchcolumn
\EN
\Source{1st on the Epiphany.}
\switchcolumn*

\LA
\noindent In hac, dilectíssimi, celebritate, sicut relatu paternæ traditiónis instrúimur, multíplici nobis est festivitate lætándum. Ferunt enim, hódie Christum Dóminum nostrum vel stella duce a Géntibus adorátum: vel invitátum ad nuptias, aquas in vina vertisse; vel suscepto a Joánne baptismate, consecrasse fluénta Jordánis, suumque simul purificasse Baptistam.\par
\switchcolumn
\EN
\noindent Dearly beloved brethren, we are instructed by the tradition of the Fathers, that we have to keep holiday on this solemnity in honour of several joyful events. We are taught that on this day, our Lord Christ was, first, manifested to the Gentiles by the leading of a star; secondly, that being bidden to a marriage, He turned water into wine; and, thirdly, that He received baptism from John, whereby He hallowed the waters of the Jordan, and cleansed him that baptised Him.\par
\switchcolumn*

\LA
\begin{Resp}\Rsign Illumináre, illumináre Jerúsalem, quia venit lux tua: \Star{} Et glória Dómini super te orta est.\end{Resp}
\begin{Resp}\Vsign Et ambulábunt gentes in lúmine tuo, et reges in splendóre ortus tui.\end{Resp}
\begin{Resp}\Rsign Et glória Dómini super te orta est.\end{Resp}
\switchcolumn
\EN
\begin{Resp}\Rsign Shine, shine, O Jerusalem, for thy light is come; \Star{} And the glory of the Lord is risen upon thee.\end{Resp}
\begin{Resp}\Vsign And the Gentiles shall walk in thy light, and kings in the brightness of thy rising.\end{Resp}
\begin{Resp}\Rsign And the glory of the Lord is risen upon thee.\end{Resp}
\switchcolumn*

\LA
\LessonHead{Lectio V}
\switchcolumn
\EN
\LessonHead{Lesson V}
\switchcolumn*

\LA
\noindent Sed quid potíssimum præsénti hoc factum sit die, nóverit ipse qui fecit: nos tamen crédere, nec dubitare debemus, quidquid illud est, factum esse pro nobis. Nam quod eum fulgentioris stellæ radiis incitáti adoravére Chaldæi, Deum verum Géntibus spes data est adorandi. Quod aquæ novo sunt ordine in vina mutátæ, novi nobis póculi prælibátum est sacraméntum. Quod autem baptizátus est Agnus Dei, regenerántis baptismi salutáre nobis munus est dedicátum.\par
\switchcolumn
\EN
\noindent Which of these events was the greatest, He knoweth by Whose Will they came to pass; for us it is needful to believe and doubt not that whatever was wrought was wrought for us. For to the Gentiles is given a hope of worshipping that Very God of Very God, to adore Whom the Chaldaeans were led by the rays of a glorious star. So also He That by His Will changed water into wine, hath given us to drink of the cup of His Blood of the New Testament; and the Lamb of God baptized in the Jordan hath hallowed for us that saving Fountain wherein we are born again.\par
\switchcolumn*

\LA
\begin{Resp}\Rsign Omnes de Saba vénient, aurum et thus deferéntes, et laudem Dómino annuntiántes, \Star{} Allelúja, allelúja, allelúja.\end{Resp}
\begin{Resp}\Vsign Reges Tharsis et ínsulæ múnera ófferent, reges Arabum et Saba dona addúcent.\end{Resp}
\begin{Resp}\Rsign Allelúja, allelúja, allelúja.\end{Resp}
\switchcolumn
\EN
\begin{Resp}\Rsign All they from Saba shall come, they shall bring gold and incense, and they shall show forth the praises of the Lord. \Star{} Alleluia, Alleluia, Alleluia.\end{Resp}
\begin{Resp}\Vsign The kings of Tarshish and of the isles shall bring presents, the kings of Arabia and Saba shall offer gifts.\end{Resp}
\begin{Resp}\Rsign Alleluia, Alleluia, Alleluia.\end{Resp}
\switchcolumn*

\LA
\LessonHead{Lectio VI}
\switchcolumn
\EN
\LessonHead{Lesson VI}
\switchcolumn*

\LA
\noindent Oportet ítaque nos, fratres, ad honórem Salvatoris nostri, cujus nativitátem débita nuper cum exsultatióne transégimus, omni cum devotióne étiam hunc virtútum ejus celebrare natalem. Et quam recte tria hæc nobis uno acta in die mysteria prædicántur, qui ineffábilis Trinitátis arcanum uno Dei sub nómine confitemur! Per hæc ergo miracula Christus Dóminus Redémptor noster óculis se vóluit revelare mortalium, quátenus invisibilis ejus divinitas, quæ latebat in hómine, in opere non latéret.\par
\switchcolumn
\EN
\noindent Therefore, my brethren, as we have lately celebrated with gladness the Festival of our Saviour's birth, so now it behoveth us with all earnestness to keep holy in His honour, this the birth-day of His wonderworking. And, verily, these three anniversaries are rightly on one day preached to us, who acknowledge the unspeakable mystery of the Trinity under the name of one God. By these miracles the Lord Christ our Redeemer willed to manifest to men some of the power of that Godhead, Which in Him lay hidden under the Manhood.\par
\switchcolumn*

\LA
\begin{Resp}\Rsign Magi véniunt ab Oriénte Jerosólymam, quæréntes, et dicéntes: Ubi est qui natus est, cujus stellam vídimus? \Star{} Et vénimus adoráre Dóminum.\end{Resp}
\begin{Resp}\Vsign Vídimus stellam ejus in Oriénte.\end{Resp}
\begin{Resp}\Rsign Et vénimus adoráre Dóminum.\end{Resp}
\begin{Resp}\Vsign Glória Patri, et Fílio, \Star{} et Spirítui Sancto.\end{Resp}
\begin{Resp}\Rsign Et vénimus adoráre Dóminum.\end{Resp}
\switchcolumn
\EN
\begin{Resp}\Rsign There came wise men from the east to Jerusalem, asking and saying: Where is He That is born King of the Jews? for we have seen His star in the east; \Star{} And are come to worship the Lord.\end{Resp}
\begin{Resp}\Vsign We have seen His star in the east.\end{Resp}
\begin{Resp}\Rsign And are come to worship the Lord.\end{Resp}
\begin{Resp}\Vsign Glory be to the Father, and to the Son, \Star{} and to the Holy Ghost.\end{Resp}
\begin{Resp}\Rsign And are come to worship the Lord.\end{Resp}
\switchcolumn*

\LA
\LessonHead{Lectio VII}
\switchcolumn
\EN
\LessonHead{Lesson VII}
\switchcolumn*

\LA
\LessonTitle{Léctio sancti Evangélii secúndum Matthǽum}
\switchcolumn
\EN
\LessonTitle{From the Holy Gospel according to Matthew}
\switchcolumn*

\LA
\Cite{Matt 2:1-12}
\switchcolumn
\EN
\Cite{Matt 2:1-12}
\switchcolumn*

\LA
\noindent Cum natus esset Jesus in Bethlehem Juda in diébus Herodis regis, ecce Magi ab Oriénte venérunt Jerosolymam, dicéntes: Ubi est qui natus est Rex Judæórum? Et réliqua.\par
\switchcolumn
\EN
\noindent When Jesus therefore was born in Bethlehem of Juda, in the days of king Herod, behold, there came wise men from the east to Jerusalem. Saying, Where is he that is born king of the Jews? And so on.\par
\switchcolumn*

\LA
\LessonTitle{Homilía sancti Hierónymi Presbýteri}
\switchcolumn
\EN
\LessonTitle{Homily by St. Jerome, Priest (at Bethlehem.)}
\switchcolumn*

\LA
\Source{Liber 1 Comment. in cap. 2. Matth.}
\switchcolumn
\EN
\Source{Bk. i. Comm. on Matth. ii.}
\switchcolumn*

\LA
\noindent Vídimus enim stellam ejus in Oriénte. Ad confusiónem Judæórum, ut nativitátem Christi a Géntibus díscerent, óritur in Oriénte stella; quam futuram, Bálaam, cujus successores erant, vaticinio nóverant. Lege Numerórum librum. Deferúntur autem Magi stellæ indicio in Judǽam, ut sacerdótes a Magis interrogati, ubi Christus nascerétur, inexcusábiles fíerent de adventu ejus.\par
\switchcolumn
\EN
\noindent We have seen His star in the East. In order that the Jews might be confounded by hearing from the Gentiles of the birth of Christ, the star rose in the East. They knew that it would come, by the prophecy of Balaam, whose successors they were. See the Book of Numbers, (xxiv. 17.) The star led the wise men to Judea, that the Priests, having it demanded of them where Christ should be born, might have no power to plead that they knew not of His coming.\par
\switchcolumn*

\LA
\begin{Resp}\Rsign Stella, quam víderant Magi in Oriénte, antecedébat eos, donec venírent ad locum, ubi puer erat. \Star{} Vidéntes autem eam, gavísi sunt gáudio magno.\end{Resp}
\begin{Resp}\Vsign Et intrántes domum, invenérunt púerum cum María matre ejus, et procidéntes adoravérunt eum.\end{Resp}
\begin{Resp}\Rsign Vidéntes autem eam, gavísi sunt gáudio magno.\end{Resp}
\switchcolumn
\EN
\begin{Resp}\Rsign The star which the wise men had seen in the East, went before them, till they came where the young Child was. \Star{} And when they saw the star, they rejoiced with exceeding great joy.\end{Resp}
\begin{Resp}\Vsign And when they were come into the house, they found the young Child with Mary His Mother, and fell down and worshipped Him.\end{Resp}
\begin{Resp}\Rsign And when they saw the star, they rejoiced with exceeding great joy.\end{Resp}
\switchcolumn*

\LA
\LessonHead{Lectio VIII}
\switchcolumn
\EN
\LessonHead{Lesson VIII}
\switchcolumn*

\LA
\noindent At illi dixérunt ei: In Bethlehem Judææ. Librariórum hic error est. Putámus enim ab Evangelista primum éditum, sicut in ipso Hebráico légimus, Judæ, non Judææ. Quæ est enim aliárum géntium Bethlehem, ut ad distinctiónem ejus, hic Judææ ponerétur? Judæ autem idcirco scribitur, quia est et alia Bethlehem in Galilæa. Lege librum Jesu fílii Nave. Dénique et in ipso testimonio, quod de Michææ prophetía sumptum est, ita habétur: Et tu Bethlehem terra Juda.\par
\switchcolumn
\EN
\noindent And they said unto him, In Bethlehem of Judea, this is a mistake of copyists. 1 In our opinion, what the Evangelist wrote must have been, not of Judea, but of Judah. Thus it is in the Hebrew text. Nor is there any town called Bethlehem among any other people, that this should be called of Judea to distinguish it. But it is fitly distinguished as of Judah, because there is in Judea another Bethlehem, namely, the one in Galilee. See the Book of Joshua the son of Nun. \textasciitilde{}(xix. 15.) Finally, the passage cited, which is in the prophet Micah, (v. 2,) hath But thou, Bethlehem Ephratah, though thou be little among the thousands of Judah(, yet out of thee shall He come forth unto me that is to be ruler in Israel.)\par
\switchcolumn*

\LA
\begin{Resp}\Rsign Vidéntes stellam Magi, gavísi sunt gáudio magno: \Star{} Et intrántes domum invenérunt púerum cum María matre ejus, et procidéntes adoravérunt eum: * Et apértis thesáuris suis obtulérunt ei múnera, aurum, thus, et myrrham.\end{Resp}
\begin{Resp}\Vsign Stella, quam víderant Magi in Oriénte, antecedébat eos, usque dum véniens staret supra ubi erat puer.\end{Resp}
\begin{Resp}\Rsign Et intrántes domum invenérunt púerum cum María matre ejus, et procidéntes adoravérunt eum:\end{Resp}
\begin{Resp}\Vsign Glória Patri, et Fílio, \Star{} et Spirítui Sancto.\end{Resp}
\begin{Resp}\Rsign Et apértis thesáuris suis obtulérunt ei múnera, aurum, thus, et myrrham.\end{Resp}
\switchcolumn
\EN
\begin{Resp}\Rsign When the wise men saw the star, they rejoiced with exceeding great joy. and when they had opened their treasures, they presented unto Him gifts; gold, and frankincense, and myrrh. \Star{} And when they were come into the house, they found the young Child with Mary His Mother, and fell down and worshipped Him; *\end{Resp}
\begin{Resp}\Vsign The star which the wise men had seen in the East, went before them, till it came and stood over where the young Child was.\end{Resp}
\begin{Resp}\Rsign And when they were come into the house, they found the young Child with Mary His Mother, and fell down and worshipped Him.\end{Resp}
\begin{Resp}\Vsign Glory be to the Father, and to the Son, \Star{} and to the Holy Ghost.\end{Resp}
\begin{Resp}\Rsign And when they had opened their treasures, they presented unto Him gifts; gold, and frankincense, and myrrh.\end{Resp}
\switchcolumn*

\LA
\LessonHead{Lectio IX}
\switchcolumn
\EN
\LessonHead{Lesson IX}
\switchcolumn*

\LA
\noindent Et, apertis thesáuris suis, obtulérunt ei múnera, aurum, thus et myrrham. Pulcherrime munerum sacraménta Juvencus Presbyter uno versículo comprehéndit: Thus, aurum, myrrham, Regíque, hominíque, Deóque dona ferunt. Et responso accepto in somnis, ne redirent ad Herodem, per aliam viam reversi sunt in regiónem suam. Qui múnera obtulerant Dómino, consequenter responsum accipiunt, non per Angelum, sed per ipsum Dóminum: ut meritórum Joseph privilegium demonstrarétur. Revertúntur autem per aliam viam: quia infidelitáti miscéndi non erant Judæórum.\par
\switchcolumn
\EN
\noindent And treasures they presented unto Him gifts, gold, and frankincense, and myrrh. The mystic meaning of these gifts is thus neatly expressed by Juvencus the Priest, To God made man, born Israel's King, Frankincense, myrrh, and gold they bring. And being warned of God in a dream that they should not return to Herod, they departed into their own country another way. They who had presented unto the Lord gifts, were honoured by receiving a warning, not from an Angel, but from God Himself; whereas even Joseph was warned only by an Angel. They departed into their own country another way, that they might not be brought into contact with the unbelief of the Jews.\par
\switchcolumn*

\LA
\TeDeum{Te Deum}
\switchcolumn
\EN
\TeDeum{Te Deum}
\switchcolumn*

\end{paracol}

\par\needspace{10\baselineskip}\addvspace{1.2em}{\centering\small\textsc{Die 11 Januarii} · \textsc{11 January}\par}
\Office{Sexta die infra Octavam Epiphaniæ}{Sexta die infra Octavam Epiphaniæ}{Semiduplex II classis}
\addcontentsline{toc}{subsection}{Die 11 Januarii · Sexta die infra Octavam Epiphaniæ\enspace\textit{Sexta die infra Octavam Epiphaniæ}}
\begin{paracol}{2}
\LA
\RefLine{Lectiones I–III: de Scriptura occurrente.}
\switchcolumn
\EN
\RefLine{Lessons I–III: from the Scripture of the day.}
\switchcolumn*

\LA
\LessonHead{Lectio IV}
\switchcolumn
\EN
\LessonHead{Lesson IV}
\switchcolumn*

\LA
\LessonTitle{Sermo sancti Fulgéntii Epíscopi}
\switchcolumn
\EN
\LessonTitle{From the Sermons of St. Fulgentius, Bishop (of Ruspe.)}
\switchcolumn*

\LA
\Source{Ex Serm. 5 qui est de Epiphania, sub initium}
\switchcolumn
\EN
\Source{5th on Twelfth-Day.}
\switchcolumn*

\LA
\noindent Ipse Deus, qui sibi in véteri testaménto primítias offérri mandávit, homo natus Géntium primítias suo cultui dedicávit. Pastóres fuérunt primítiæ Judæórum: Magi facti sunt primítiæ Géntium. Illi de próximo addúcti, isti de longínquo redúcti. Ubi est, ínquiunt, qui natus est Rex Judæórum? De Heróde rege Judæórum fílii jam fúerant nati. Archeláus natus est in palátio, Christus in diversório. Archeláus natus, est in lecto argénteo pósitus, Christus autem natus, in præsépio est brevíssimo collocátus: et tamen ille natus in palátio contémnitur, iste natus in diversório quæritur: ille a Magis nullátenus nominátur, iste invéntus supplíciter adorátur.\par
\switchcolumn
\EN
\noindent The same God Who in the Old Testament had commanded the first-fruits to be offered to Himself, being born as a man, Himself consecrated to His own worship the firstfruits of the nations. The Shepherds were the first-fruits of the Jews, and the wise men of the Gentiles. The first came from near at hand, the second from afar. Where is He, say they, that is born King of the Jews? Herod, the king of the Jews, had already had children. Archelaus was born in a palace, Christ at an inn; Archelaus was laid in a silver cradle, Christ in a manger. And yet the wise men sought, not Archelaus, but Christ; they did not even name him that was born in a palace, but when they found Him That lay in a manger, they fell down and worshipped Him.\par
\switchcolumn*

\LA
\begin{Resp}\Rsign Illumináre, illumináre Jerúsalem, quia venit lux tua: \Star{} Et glória Dómini super te orta est.\end{Resp}
\begin{Resp}\Vsign Et ambulábunt gentes in lúmine tuo, et reges in splendóre ortus tui.\end{Resp}
\begin{Resp}\Rsign Et glória Dómini super te orta est.\end{Resp}
\switchcolumn
\EN
\begin{Resp}\Rsign Shine, shine, O Jerusalem, for thy light is come; \Star{} And the glory of the Lord is risen upon thee.\end{Resp}
\begin{Resp}\Vsign And the Gentiles shall walk in thy light, and kings in the brightness of thy rising.\end{Resp}
\begin{Resp}\Rsign And the glory of the Lord is risen upon thee.\end{Resp}
\switchcolumn*

\LA
\LessonHead{Lectio V}
\switchcolumn
\EN
\LessonHead{Lesson V}
\switchcolumn*

\LA
\noindent Quis est iste Rex Judæórum? Páuperet dives, húmilis et sublímis. Quis est iste Rex Judæórum, qui portátur ut párvulus, adorátur ut Deus? Parvus in præsépio, imménsus in cælo: vilis in pannis, pretiósus in stellis. Quid est quod sic turbáris, Heródes? Rex iste, qui natus est, non venit reges pugnándo superáre, sed moriéndo mirabíliter subjugáre. Nec ídeo natus est, ut tibi succédat, sed ut in eum mundus fidéliter credat. Venit ergo, non ut pugnet vivus, sed ut triúmphet occísus.\par
\switchcolumn
\EN
\noindent Who is the King of the Jews? The Poor and the Rich, the Lowly and the Exalted One. Who is the King of the Jews? He Who, being carried at the breast, is adored as the Eternal; He Who lieth tiny in the manger, and is He Whom the heavens cannot contain; He Who is meanly wrapped in swaddling clothes, and is more glorious than all the stars. Why art thou troubled, O Herod? He that is born King of the Jews cometh not by carnal warfare to conquer other kings, but by a marvellous working, by dying, to subdue them to Himself. He is not born to be thy successor, but that the world may faithfully believe in Him. He cometh, not that He may fight in the flesh, but that He may conquer through the suffering of death.\par
\switchcolumn*

\LA
\begin{Resp}\Rsign Omnes de Saba vénient, aurum et thus deferéntes, et laudem Dómino annuntiántes, \Star{} Allelúja, allelúja, allelúja.\end{Resp}
\begin{Resp}\Vsign Reges Tharsis et ínsulæ múnera ófferent, reges Arabum et Saba dona addúcent.\end{Resp}
\begin{Resp}\Rsign Allelúja, allelúja, allelúja.\end{Resp}
\switchcolumn
\EN
\begin{Resp}\Rsign All they from Saba shall come, they shall bring gold and incense, and they shall show forth the praises of the Lord. \Star{} Alleluia, Alleluia, Alleluia.\end{Resp}
\begin{Resp}\Vsign The kings of Tarshish and of the isles shall bring presents, the kings of Arabia and Saba shall offer gifts.\end{Resp}
\begin{Resp}\Rsign Alleluia, Alleluia, Alleluia.\end{Resp}
\switchcolumn*

\LA
\LessonHead{Lectio VI}
\switchcolumn
\EN
\LessonHead{Lesson VI}
\switchcolumn*

\LA
\noindent Puer iste, qui nunc a Magis dícitur Rex Judæórum, idem Creátor est, et Dóminus Angelórum. Quaprópter cujus times infántiam nascéntis, magis timére debes omnipoténtiam judicántis. Noli eum timére regni tui successórem, sed time infidelitátis tuæ justíssimum damnatórem. Ite, inquit, et renuntiáte mihi, ut et ego véniens adórem eum. O callíditas ficta, o incredúlitas ímpia, o nequítia fraudulénta! Sanguis Innocéntium, quem crudéliter effudísti, attestátur quid de hoc púero voluísti.\par
\switchcolumn
\EN
\noindent The little Child, Whom the wise men call the King of the Jews, is the Maker and Lord of Angels. If thou fearest Him at His birth, thou hast more reason to fear Him as the Almighty Judge. Fear Him, not as a pretender to thy kingdom, but fear Him as Him Who will pass a most just sentence of condemnation on thee because thou hast not believed in Him. Go, said Herod, and bring me word again, that I may come and worship Him also. We know thy cunning lying, thy godless unbelief, thine iniquitous treachery. The blood of the innocents which thou didst cruelly shed, is witness to us of what thou wouldst have done to Him.\par
\switchcolumn*

\LA
\begin{Resp}\Rsign Magi véniunt ab Oriénte Jerosólymam, quæréntes, et dicéntes: Ubi est qui natus est, cujus stellam vídimus? \Star{} Et vénimus adoráre Dóminum.\end{Resp}
\begin{Resp}\Vsign Vídimus stellam ejus in Oriénte.\end{Resp}
\begin{Resp}\Rsign Et vénimus adoráre Dóminum.\end{Resp}
\begin{Resp}\Vsign Glória Patri, et Fílio, \Star{} et Spirítui Sancto.\end{Resp}
\begin{Resp}\Rsign Et vénimus adoráre Dóminum.\end{Resp}
\switchcolumn
\EN
\begin{Resp}\Rsign There came wise men from the east to Jerusalem, asking and saying: Where is He That is born King of the Jews? for we have seen His star in the east; \Star{} And are come to worship the Lord.\end{Resp}
\begin{Resp}\Vsign We have seen His star in the east.\end{Resp}
\begin{Resp}\Rsign And are come to worship the Lord.\end{Resp}
\begin{Resp}\Vsign Glory be to the Father, and to the Son, \Star{} and to the Holy Ghost.\end{Resp}
\begin{Resp}\Rsign And are come to worship the Lord.\end{Resp}
\switchcolumn*

\LA
\LessonHead{Lectio VII}
\switchcolumn
\EN
\LessonHead{Lesson VII}
\switchcolumn*

\LA
\LessonTitle{Léctio sancti Evangélii secúndum Matthǽum}
\switchcolumn
\EN
\LessonTitle{From the Holy Gospel according to Matthew}
\switchcolumn*

\LA
\Cite{Matt 2:1-12}
\switchcolumn
\EN
\Cite{Matt 2:1-12}
\switchcolumn*

\LA
\noindent Cum natus esset Jesus in Béthlehem Juda in diébus Heródis regis, ecce Magi ab Oriénte venérunt Jerosólymam, dicéntes: Ubi est qui natus est Rex Judæórum? Et réliqua.\par
\switchcolumn
\EN
\noindent When Jesus therefore was born in Bethlehem of Juda, in the days of king Herod, behold, there came wise men from the east to Jerusalem. Saying, Where is he that is born king of the Jews? And so on.\par
\switchcolumn*

\LA
\LessonTitle{Homilía sancti Ambrósii Epíscopi}
\switchcolumn
\EN
\LessonTitle{Homily by St. Ambrose, Bishop (of Milan.)}
\switchcolumn*

\LA
\Source{Liber 2 in Lucam cap. 2, post init.}
\switchcolumn
\EN
\Source{Book ii. on Luke ii.}
\switchcolumn*

\LA
\noindent Quæ sunt ista veræ fídei múnera? Aurum Regi, thus Deo, myrrha defúncto. Aliud enim Regis insígne, áliud divínæ sacraméntum potestátis, áliud honor est sepultúræ, quæ non corrúmpat corpus mórtui, sed resérvet. Nos quoque, qui hæc audímus et légimus, de thesáuris nostris tália, fratres, múnera proferámus. Habémus enim thesáurum in vasis fictílibus. Si ígitur in teípso quod es, non ex te debes æstimáre, sed ex Christo: quanto magis in Christo non tua debes æstimáre, sed Christi?\par
\switchcolumn
\EN
\noindent What are the gifts of the faithful and true? Gold to our King, frankincense to our God, and myrrh to Him Who died for us. The first is that whereof are made the royal honours of kings, the second is that mystic offering which is used in the worship of the Divine Power, and the third is that wherewith we pay respect to the dead, whose bodies it keepeth from corruption. My brethren, let us who hear and read these things, make offering out of what treasures we have albeit we have it in earthen vessels. (2 Cor. iv. 7.) If we confess that all that we have, we have, not from ourselves, but from Christ, how much more should we confess that whatever we have is not our own, but Christ's?\par
\switchcolumn*

\LA
\begin{Resp}\Rsign Stella, quam víderant Magi in Oriénte, antecedébat eos, donec venírent ad locum, ubi puer erat. \Star{} Vidéntes autem eam, gavísi sunt gáudio magno.\end{Resp}
\begin{Resp}\Vsign Et intrántes domum, invenérunt púerum cum María matre ejus, et procidéntes adoravérunt eum.\end{Resp}
\begin{Resp}\Rsign Vidéntes autem eam, gavísi sunt gáudio magno.\end{Resp}
\switchcolumn
\EN
\begin{Resp}\Rsign The star which the wise men had seen in the East, went before them, till they came where the young Child was. \Star{} And when they saw the star, they rejoiced with exceeding great joy.\end{Resp}
\begin{Resp}\Vsign And when they were come into the house, they found the young Child with Mary His Mother, and fell down and worshipped Him.\end{Resp}
\begin{Resp}\Rsign And when they saw the star, they rejoiced with exceeding great joy.\end{Resp}
\switchcolumn*

\LA
\LessonHead{Lectio VIII}
\switchcolumn
\EN
\LessonHead{Lesson VIII}
\switchcolumn*

\LA
\noindent Ergo Magi de thesáuris suis ófferunt múnera. Vultis scire quam bonum méritum hábeant? Stella ab his vidétur: ubi Heródes est, non vidétur; ubi Christus est, rursus vidétur, et viam demónstrat. Ergo stella hæc via est, et via Christus: quia secúndum incarnatiónis mystérium Christus est stella. Oriétur enim stella ex Jacob, et exsúrget homo ex Israël. Dénique ubi Christus, et stella est. Ipse enim est stella spléndida et matutína. Sua ígitur ipse luce se signat.\par
\switchcolumn
\EN
\noindent The wise men out of their treasures presented unto Him gifts. Wilt thou know how pleasing to Him they were? The star appeared to them, but disappeared when it came near Herod. Then it appeareth to them again, leading them on the way that led to Christ. This star then was the way, and we know that Christ calleth Himself the Way. (John xiv. 6.) And truly also in the mystery of His Incarnation He is called a Star; as it is written There shall come forth a Star out of Jacob, and a Man shall rise out of Israel. Where Christ is, there is a Star; yea, He is Himself the bright and morning Star. (Apoc. xxii. 16.) And the light that leadeth to Jesus is His own.\par
\switchcolumn*

\LA
\begin{Resp}\Rsign Vidéntes stellam Magi, gavísi sunt gáudio magno: \Star{} Et intrántes domum invenérunt púerum cum María matre ejus, et procidéntes adoravérunt eum: * Et apértis thesáuris suis obtulérunt ei múnera, aurum, thus, et myrrham.\end{Resp}
\begin{Resp}\Vsign Stella, quam víderant Magi in Oriénte, antecedébat eos, usque dum véniens staret supra ubi erat puer.\end{Resp}
\begin{Resp}\Rsign Et intrántes domum invenérunt púerum cum María matre ejus, et procidéntes adoravérunt eum:\end{Resp}
\begin{Resp}\Vsign Glória Patri, et Fílio, \Star{} et Spirítui Sancto.\end{Resp}
\begin{Resp}\Rsign Et apértis thesáuris suis obtulérunt ei múnera, aurum, thus, et myrrham.\end{Resp}
\switchcolumn
\EN
\begin{Resp}\Rsign When the wise men saw the star, they rejoiced with exceeding great joy. and when they had opened their treasures, they presented unto Him gifts; gold, and frankincense, and myrrh. \Star{} And when they were come into the house, they found the young Child with Mary His Mother, and fell down and worshipped Him; *\end{Resp}
\begin{Resp}\Vsign The star which the wise men had seen in the East, went before them, till it came and stood over where the young Child was.\end{Resp}
\begin{Resp}\Rsign And when they were come into the house, they found the young Child with Mary His Mother, and fell down and worshipped Him.\end{Resp}
\begin{Resp}\Vsign Glory be to the Father, and to the Son, \Star{} and to the Holy Ghost.\end{Resp}
\begin{Resp}\Rsign And when they had opened their treasures, they presented unto Him gifts; gold, and frankincense, and myrrh.\end{Resp}
\switchcolumn*

\LA
\LessonHead{Lectio IX}
\switchcolumn
\EN
\LessonHead{Lesson IX}
\switchcolumn*

\LA
\noindent Accipe áliud documéntum. Alia venérunt via Magi, ália rédeunt. Qui enim Christum víderant, Christum intelléxerant; melióres útique, quam vénerant, revertúntur. Duæ quippe sunt viæ: una, quæ ducit ad intéritum; alia, quæ ducit ad regnum. Illa peccatórum est, quæ ducit ad Heródem: hæc Christus est, qua redítur ad pátriam. Hic enim temporális est incolátus, sicut scriptum est: Multum íncola facta est ánima mea.\par
\switchcolumn
\EN
\noindent Remark another point. The wise men came by one way and departed by another. They that had seen Christ, knew Christ, and they departed better than they came. There are two ways, the one which leadeth to destruction, the other which leadeth to the kingdom; the one is the way of sin, which leadeth to Herod; the other is Christ, the true Way, Who leadeth us home to the fatherland, from that journeying here, whereof it is said My soul hath long dwelt as an exile. (Ps. cxix. 5.)\par
\switchcolumn*

\LA
\TeDeum{Te Deum}
\switchcolumn
\EN
\TeDeum{Te Deum}
\switchcolumn*

\end{paracol}

\par\needspace{10\baselineskip}\addvspace{1.2em}{\centering\small\textsc{Die 12 Januarii} · \textsc{12 January}\par}
\Office{Septima die infra Octavam Epiphaniæ}{Septima die infra Octavam Epiphaniæ}{Semiduplex II classis}
\addcontentsline{toc}{subsection}{Die 12 Januarii · Septima die infra Octavam Epiphaniæ\enspace\textit{Septima die infra Octavam Epiphaniæ}}
\begin{paracol}{2}
\LA
\RefLine{Lectiones I–III: de Scriptura occurrente.}
\switchcolumn
\EN
\RefLine{Lessons I–III: from the Scripture of the day.}
\switchcolumn*

\LA
\LessonHead{Lectio IV}
\switchcolumn
\EN
\LessonHead{Lesson IV}
\switchcolumn*

\LA
\LessonTitle{Sermo sancti Leónis Papæ}
\switchcolumn
\EN
\LessonTitle{From the Sermons of Pope St. Leo (the Great.)}
\switchcolumn*

\LA
\Source{Sermo 4 de Epiphania}
\switchcolumn
\EN
\Source{5th on the Twelfth day}
\switchcolumn*

\LA
\noindent Justum et rationábile, dilectíssimi, veræ pietátis obséquium est in diébus, qui divínæ misericórdiæ ópera protestántur, toto corde gaudére, et honorífice ea, quæ ad salútem nostram gesta sunt, celebráre, vocánte nos ad hanc devotiónem ipsa recurréntium témporum lege, quæ nobis post diem, in quo coætérnus Patri Fílius Dei natus ex Vírgine est, brevi intervállo Epiphaníæ íntulit festum, ex apparitióne Dómini consecrátum.\par
\switchcolumn
\EN
\noindent This meet and right, dearly beloved brethren, yea, it is our bounden duty and godly service, to rejoice with full hearts upon those days which more especially set forth before us the workings of God's mercy; and to have in honourable memory those things that were done for our salvation. Hereto are we called by the seasons of the year which continually return, and notably by this present, which, after but a short time hath passed since that day whereon the Coeternal Son of God was born of a Virgin, bringeth now the Feast of the Epiphany, hallowed by the Manifestation of the Lord.\par
\switchcolumn*

\LA
\begin{Resp}\Rsign Illumináre, illumináre Jerúsalem, quia venit lux tua: \Star{} Et glória Dómini super te orta est.\end{Resp}
\begin{Resp}\Vsign Et ambulábunt gentes in lúmine tuo, et reges in splendóre ortus tui.\end{Resp}
\begin{Resp}\Rsign Et glória Dómini super te orta est.\end{Resp}
\switchcolumn
\EN
\begin{Resp}\Rsign Shine, shine, O Jerusalem, for thy light is come; \Star{} And the glory of the Lord is risen upon thee.\end{Resp}
\begin{Resp}\Vsign And the Gentiles shall walk in thy light, and kings in the brightness of thy rising.\end{Resp}
\begin{Resp}\Rsign And the glory of the Lord is risen upon thee.\end{Resp}
\switchcolumn*

\LA
\LessonHead{Lectio V}
\switchcolumn
\EN
\LessonHead{Lesson V}
\switchcolumn*

\LA
\noindent In quo fídei nostræ magnum præsídium providéntia divína constítuit: ut, dum solémni veneratióne recólitur adoráta in exórdiis suis Salvatóris infántia, per ipsa originália documénta probarétur veri hóminis in ipso orta natúra. Hoc enim est quod justíficat ímpios, hoc est quod ex peccatóribus facit sanctos, si in uno eodémque Dómino nostro Jesu Christo et vera Déitas, et vera credátur humánitas: Déitas, qua ante ómnia sæcula in forma Dei æquális est Patri: humánitas, qua novíssimis diébus in forma servi unítus est hómini.\par
\switchcolumn
\EN
\noindent In this said Manifestation the good Providence of God hath appointed a strong bulwark to our faith. For now, while in solemn worship we call to mind how the childhood of the Saviour was adored in its first infancy, we receive from the original Scriptures the doctrine that Christ was born with the very nature of man. For this is that which maketh of sinners saints, even to believe that in one and the same our Lord Jesus Christ there is very Godhead and very Manhood very Godhead, as touching Which, He, being in the form of God, is equal to the Father from everlasting to everlasting and very Manhood, wherein He, taking upon Him the form of a servant, hath in these latter days been born Man.\par
\switchcolumn*

\LA
\begin{Resp}\Rsign Omnes de Saba vénient, aurum et thus deferéntes, et laudem Dómino annuntiántes, \Star{} Allelúja, allelúja, allelúja.\end{Resp}
\begin{Resp}\Vsign Reges Tharsis et ínsulæ múnera ófferent, reges Arabum et Saba dona addúcent.\end{Resp}
\begin{Resp}\Rsign Allelúja, allelúja, allelúja.\end{Resp}
\switchcolumn
\EN
\begin{Resp}\Rsign All they from Saba shall come, they shall bring gold and incense, and they shall show forth the praises of the Lord. \Star{} Alleluia, Alleluia, Alleluia.\end{Resp}
\begin{Resp}\Vsign The kings of Tarshish and of the isles shall bring presents, the kings of Arabia and Saba shall offer gifts.\end{Resp}
\begin{Resp}\Rsign Alleluia, Alleluia, Alleluia.\end{Resp}
\switchcolumn*

\LA
\LessonHead{Lectio VI}
\switchcolumn
\EN
\LessonHead{Lesson VI}
\switchcolumn*

\LA
\noindent Ad roborándam ergo hanc fidem, quæ contra omnes præmuniebátur erróres, ex magna gestum est divíni pietáte consílii ut gens in longínqua Orientális plagæ regióne consístens, quæ spectandórum síderum arte pollébat, signum nati púeri, qui supra omnem Israël esset regnatúrus, accíperet. Nova étenim cláritas apud Magos stellæ illustrióris appáruit, et intuéntium ánimos ita admiratióne sui splendóris implévit, ut nequáquam sibi créderent negligéndum, quod tanto nuntiabátur indício.\par
\switchcolumn
\EN
\noindent For the strengthening of this our faith, which we profess in the face of every false doctrine, the mercy of God hath made it to come to pass that one of those peoples who dwell in the uttermost parts of the East, and excel in the skill of reading the stars, should see the sign of the birth of that Child Who was to reign over all Israel. There appeared to the eyes of wise men a new star of such passing beauty, as wrought in the minds of all that saw it the persuasion that the event, which is announced, was of an importance not to be neglected.\par
\switchcolumn*

\LA
\begin{Resp}\Rsign Magi véniunt ab Oriénte Jerosólymam, quæréntes, et dicéntes: Ubi est qui natus est, cujus stellam vídimus? \Star{} Et vénimus adoráre Dóminum.\end{Resp}
\begin{Resp}\Vsign Vídimus stellam ejus in Oriénte.\end{Resp}
\begin{Resp}\Rsign Et vénimus adoráre Dóminum.\end{Resp}
\begin{Resp}\Vsign Glória Patri, et Fílio, \Star{} et Spirítui Sancto.\end{Resp}
\begin{Resp}\Rsign Et vénimus adoráre Dóminum.\end{Resp}
\switchcolumn
\EN
\begin{Resp}\Rsign There came wise men from the east to Jerusalem, asking and saying: Where is He That is born King of the Jews? for we have seen His star in the east; \Star{} And are come to worship the Lord.\end{Resp}
\begin{Resp}\Vsign We have seen His star in the east.\end{Resp}
\begin{Resp}\Rsign And are come to worship the Lord.\end{Resp}
\begin{Resp}\Vsign Glory be to the Father, and to the Son, \Star{} and to the Holy Ghost.\end{Resp}
\begin{Resp}\Rsign And are come to worship the Lord.\end{Resp}
\switchcolumn*

\LA
\LessonHead{Lectio VII}
\switchcolumn
\EN
\LessonHead{Lesson VII}
\switchcolumn*

\LA
\LessonTitle{Léctio sancti Evangélii secúndum Matthǽum}
\switchcolumn
\EN
\LessonTitle{From the Holy Gospel according to Matthew}
\switchcolumn*

\LA
\Cite{Matt 2:1-12}
\switchcolumn
\EN
\Cite{Matt 2:1-12}
\switchcolumn*

\LA
\noindent Cum natus esset Jesus in Béthlehem Juda in diébus Heródis regis, ecce Magi ab Oriénte venérunt Jerosólymam, dicéntes: Ubi est qui natus est Rex Judæórum? Et réliqua.\par
\switchcolumn
\EN
\noindent When Jesus therefore was born in Bethlehem of Juda, in the days of king Herod, behold, there came wise men from the east to Jerusalem, saying: Where is he that is born king of the Jews? And so on.\par
\switchcolumn*

\LA
\LessonTitle{De Homilía sancti Joánnis Chrysóstomi}
\switchcolumn
\EN
\LessonTitle{Homily by St. John Chrysostom}
\switchcolumn*

\LA
\Source{Homilia 8, in Matth. n. 1}
\switchcolumn
\EN
\Source{Homily 8, on Matthew, n. 1}
\switchcolumn*

\LA
\noindent Magi, intrántes domum, vidérunt púerum cum María, matre ejus. Et procidéntes adoravérunt eum; et, apértis thesáuris, obtulérunt ei múnera, aurum, thus et myrrham. Verum quid illos ad púerum adorándum indúxit? Neque enim Virgo insígne quídpiam præ se ferébat, neque domus magnífica erat, neque istic áliquid áliud erat quod posset illos vel percéllere vel allícere. Illi vero non modo adórant, sed, apértis thesáuris suis, múnera ófferunt, múnera, inquam, non quasi hómini, sed quasi Deo. Thus enim et myrrha máxime Dei symbolum erat. Quid ígitur hoc illis suásit? Idípsum quod excitávit illos, ut, relícta domo, tantum iter suscíperent: nimírum stella et illustrátio mentis a Deo ipsis índita, quæ paulátim illos ad perfectiórem notítiam dedúxit.\par
\switchcolumn
\EN
\noindent “The wise men entering into the house, they found the child with Mary his mother, and falling down they adored him: and opening their treasures, they offered him gifts; gold, frankincense, and myrrh.” (Matthew 2:11) But what was it that induced them to worship? For neither was the Virgin conspicuous, nor the house distinguished, nor was any other of the things which they saw apt to amaze or attract them. Yet they not only worship, but also open their treasures, and offer gifts; and gifts, not as to a man, but as to God. For the frankincense and the myrrh were a symbol of this. What then was their inducement? That which wrought upon them to set out from home and to come so long a journey; and this was both the star, and the illumination wrought of God in their mind, guiding them little by little to the more perfect knowledge.\par
\switchcolumn*

\LA
\begin{Resp}\Rsign Stella, quam víderant Magi in Oriénte, antecedébat eos, donec venírent ad locum, ubi puer erat. \Star{} Vidéntes autem eam, gavísi sunt gáudio magno.\end{Resp}
\begin{Resp}\Vsign Et intrántes domum, invenérunt púerum cum María matre ejus, et procidéntes adoravérunt eum.\end{Resp}
\begin{Resp}\Rsign Vidéntes autem eam, gavísi sunt gáudio magno.\end{Resp}
\switchcolumn
\EN
\begin{Resp}\Rsign The star which the wise men had seen in the East, went before them, till they came where the young Child was. \Star{} And when they saw the star, they rejoiced with exceeding great joy.\end{Resp}
\begin{Resp}\Vsign And when they were come into the house, they found the young Child with Mary His Mother, and fell down and worshipped Him.\end{Resp}
\begin{Resp}\Rsign And when they saw the star, they rejoiced with exceeding great joy.\end{Resp}
\switchcolumn*

\LA
\LessonHead{Lectio VIII}
\switchcolumn
\EN
\LessonHead{Lesson VIII}
\switchcolumn*

\LA
\noindent Nisi enim res ita se habéret, cum ómnia quæ istic videbántur, vília essent, non ei tantum exhibuíssent honórem. Ideo autem nihil eórum, quæ sub sensum cadunt, ibi magnum erat, sed præsépe, tugúrium, mater inops; ut nudam Magórum philosophíam perspícias, atque discas, eos non ut hóminem purum, sed ut Deum, et ipsum benéficum, adivísse. Quaprópter nullo eórum quæ extrínsecus videbántur, offénsi sunt, sed adorárunt; et dona obtulérunt, quæ dona multum a Judáica crassítie differébant. Neque enim oves et vítulos immolárunt, sed quæ ecclesiásticæ philosophíæ vicína erant; síquidem sciéntiam, obediéntiam et dilectiónem ipsi offerébant.\par
\switchcolumn
\EN
\noindent For, surely, had it not been so, all that was in sight being ordinary, they would not have shown so great honor. Therefore none of the outward circumstances was great in that instance, but it was a manger, and a shed, and a mother in poor estate; to set before your eyes, naked and bare, those wise men's love of wisdom, and to prove to you, that not as mere man they approached Him, but as a God, and Benefactor. Wherefore neither were they offended by ought of what they saw outwardly, but even worshipped, and brought gifts; gifts not only free from Judaical grossness, in that they sacrificed not sheep and calves, but also coming near to the self-devotion of the Church, for it was knowledge and obedience and love that they offered unto Him.\par
\switchcolumn*

\LA
\begin{Resp}\Rsign Vidéntes stellam Magi, gavísi sunt gáudio magno: \Star{} Et intrántes domum invenérunt púerum cum María matre ejus, et procidéntes adoravérunt eum: * Et apértis thesáuris suis obtulérunt ei múnera, aurum, thus, et myrrham.\end{Resp}
\begin{Resp}\Vsign Stella, quam víderant Magi in Oriénte, antecedébat eos, usque dum véniens staret supra ubi erat puer.\end{Resp}
\begin{Resp}\Rsign Et intrántes domum invenérunt púerum cum María matre ejus, et procidéntes adoravérunt eum:\end{Resp}
\begin{Resp}\Vsign Glória Patri, et Fílio, \Star{} et Spirítui Sancto.\end{Resp}
\begin{Resp}\Rsign Et apértis thesáuris suis obtulérunt ei múnera, aurum, thus, et myrrham.\end{Resp}
\switchcolumn
\EN
\begin{Resp}\Rsign When the wise men saw the star, they rejoiced with exceeding great joy. and when they had opened their treasures, they presented unto Him gifts; gold, and frankincense, and myrrh. \Star{} And when they were come into the house, they found the young Child with Mary His Mother, and fell down and worshipped Him; *\end{Resp}
\begin{Resp}\Vsign The star which the wise men had seen in the East, went before them, till it came and stood over where the young Child was.\end{Resp}
\begin{Resp}\Rsign And when they were come into the house, they found the young Child with Mary His Mother, and fell down and worshipped Him.\end{Resp}
\begin{Resp}\Vsign Glory be to the Father, and to the Son, \Star{} and to the Holy Ghost.\end{Resp}
\begin{Resp}\Rsign And when they had opened their treasures, they presented unto Him gifts; gold, and frankincense, and myrrh.\end{Resp}
\switchcolumn*

\LA
\LessonHead{Lectio IX}
\switchcolumn
\EN
\LessonHead{Lesson IX}
\switchcolumn*

\LA

\switchcolumn
\EN
\LessonTitle{“And having received an answer in sleep that they should not return to Herod, they went back another way into their country.” (Matthew 2:12) See from this also their faith, how they were not offended, but are docile, and considerate; neither are they troubled, nor reason with themselves, saying, And yet, if this Child be great, and has any might, what need of flight, and of a clandestine retreat? And wherefore can it be, that when we have come openly and with boldness, and have stood against so great a people, and against a king's madness, the angel sends us out of the city as runaways and fugitives? But none of these things did they either say or think. For this most especially belongs to faith, not to seek an account of what is enjoined, but merely to obey the commandments laid upon us.}
\switchcolumn*

\LA
\noindent Et, respónso accépto in somnis, ne redírent ad Heródem, per áliam viam revérsi sunt in regiónem suam. Hic mihi vide fidem ipsórum, quómodo non offendántur, sed sint quiéti et obtemperántes, neque turbéntur, neque tália mútuo loquántur: Sane, si magnus hic puer est, et si quam habet poténtiam, quid opus fuga et occúlto discéssu? cur vero nos, qui palam et cum fidúcia accéssimus ad pópulum tantum, et coram furénte rege stétimus, quasi fugitívos ex civitáte dimíttit Angelus? At nihil tale vel dixérunt vel cogitárunt. Illud enim máxime ad fidem pértinet, ut nulla mandatórum exquirátur rátio, sed jussis tantum obediátur.\par
\switchcolumn
\EN

\switchcolumn*

\LA
\TeDeum{Te Deum}
\switchcolumn
\EN
\TeDeum{Te Deum}
\switchcolumn*

\end{paracol}

\par\needspace{10\baselineskip}\addvspace{1.2em}{\centering\small\textsc{Die 13 Januarii} · \textsc{13 January}\par}
\Office{In Octava Epiphaniæ}{In Octava Epiphaniæ}{Duplex majus}
\addcontentsline{toc}{subsection}{Die 13 Januarii · In Octava Epiphaniæ\enspace\textit{In Octava Epiphaniæ}}
\begin{paracol}{2}
\LA
\RefLine{Lectiones I–III: de Scriptura occurrente.}
\switchcolumn
\EN
\RefLine{Lessons I–III: from the Scripture of the day.}
\switchcolumn*

\LA
\LessonHead{Lectio IV}
\switchcolumn
\EN
\LessonHead{Lesson IV}
\switchcolumn*

\LA
\LessonTitle{Sermo sancti Gregórii Nazianzéni}
\switchcolumn
\EN
\LessonTitle{From the Sermons of St. Gregory of Nazianzus, Patriarch (of Constantinople.)}
\switchcolumn*

\LA
\Source{Oratio in sancta Lumina}
\switchcolumn
\EN
\Source{Discourse on the Epiphany, 2}
\switchcolumn*

\LA
\noindent Non possum cohibére lætítiæ voluptátem, sed mente extóllor, et affícior: et própriæ pusillitátis oblítus, offícium magni Joánnis, immo pótius famulátum subíre conténdo, ac géstio: et licet non sim præcúrsor, de erémo tamen vénio. Christus ergo illuminátur, immo pótius fulgóre suo nos illúminat: Christus baptizátur, simul et nos descendámus, ut cum ipso páriter ascendámus.\par
\switchcolumn
\EN
\noindent I am not able to restrain the outbursts of my happiness. I feel changed and elated. I forget my own meanness while I undertake and try to discharge the office of the great John. It is true that I am not the Forerunner, but at least I come from the desert. Christ is enlightened, or rather, He enlighteneth us with His own light. Christ is baptized; let us go down with Him into the water, that we may come up with Him.\par
\switchcolumn*

\LA
\begin{Resp}\Rsign Illumináre, illumináre Jerúsalem, quia venit lux tua: \Star{} Et glória Dómini super te orta est.\end{Resp}
\begin{Resp}\Vsign Et ambulábunt gentes in lúmine tuo, et reges in splendóre ortus tui.\end{Resp}
\begin{Resp}\Rsign Et glória Dómini super te orta est.\end{Resp}
\switchcolumn
\EN
\begin{Resp}\Rsign Shine, shine, O Jerusalem, for thy light is come; \Star{} And the glory of the Lord is risen upon thee.\end{Resp}
\begin{Resp}\Vsign And the Gentiles shall walk in thy light, and kings in the brightness of thy rising.\end{Resp}
\begin{Resp}\Rsign And the glory of the Lord is risen upon thee.\end{Resp}
\switchcolumn*

\LA
\LessonHead{Lectio V}
\switchcolumn
\EN
\LessonHead{Lesson V}
\switchcolumn*

\LA
\noindent Joánnes baptízat, et accédit Jesus, sanctíficans quidem et ipsum, qui baptízat, præcípue tamen ut et véterem Adam sepéliat in aquis, et ante ómnia, ut per hæc sanctificéntur aquæ Jordánis: ut sicut erat spíritus et caro, ita et his, qui baptizándi erant, in spíritu et aqua sanctificatiónis succéssio traderétur. Non súscipit Baptísta, conténdit Jesus. Ego, inquit, opus hábeo a te baptizári. Lucérna dicit ad Solem, et vox lóquitur ad Verbum.\par
\switchcolumn
\EN
\noindent John is baptizing. Jesus cometh. He cometh that He may make holy him who baptizeth Him; He cometh to bury the old Adam in the waters; He cometh to hallow the blessed flood of Jordan. He Who is Flesh and Spirit cometh to open for all that should ever be baptized that power of generation whereby new peoples are constantly begotten of water and the Holy Ghost. The Baptist will not receive Him. Jesus striveth with him. I, saith John, have need to be baptized of thee. Thus speaketh the candle to the Sun, the voice to the Word.\par
\switchcolumn*

\LA
\begin{Resp}\Rsign Omnes de Saba vénient, aurum et thus deferéntes, et laudem Dómino annuntiántes, \Star{} Allelúja, allelúja, allelúja.\end{Resp}
\begin{Resp}\Vsign Reges Tharsis et ínsulæ múnera ófferent, reges Arabum et Saba dona addúcent.\end{Resp}
\begin{Resp}\Rsign Allelúja, allelúja, allelúja.\end{Resp}
\switchcolumn
\EN
\begin{Resp}\Rsign All they from Saba shall come, they shall bring gold and incense, and they shall show forth the praises of the Lord. \Star{} Alleluia, Alleluia, Alleluia.\end{Resp}
\begin{Resp}\Vsign The kings of Tarshish and of the isles shall bring presents, the kings of Arabia and Saba shall offer gifts.\end{Resp}
\begin{Resp}\Rsign Alleluia, Alleluia, Alleluia.\end{Resp}
\switchcolumn*

\LA
\LessonHead{Lectio VI}
\switchcolumn
\EN
\LessonHead{Lesson VI}
\switchcolumn*

\LA
\noindent Ascéndit Jesus de aqua, secum quodámmodo demérsum edúcens et élevans mundum: et vidit non dívidi cælum, sed aperíri, quod sibi ac nobis post se aliquándo Adam ille conclúserat; sicut et ígneo gládio paradísus fúerat conclúsus. Spíritus Sanctus testimónium pérhibet: simília namque sibi ipsa concúrrunt. De cælo testimónium defértur: inde enim erat ille, cui testimónium perhibétur.\par
\switchcolumn
\EN
\noindent Jesus came up out of the water, having, in a manner, washed the whole world, and brought it up with Him. And He saw the heavens opened not divided, even those heavens which Adam had once shut upon himself and us his descendants, when the cherub's fiery sword barred the gates of Paradise. And the Holy Spirit bare witness, witness unto Him Who is of One Substance with Himself. And witness was given from Heaven, unto Him that came down from heaven.\par
\switchcolumn*

\LA
\begin{Resp}\Rsign Magi véniunt ab Oriénte Jerosólymam, quæréntes, et dicéntes: Ubi est qui natus est, cujus stellam vídimus? \Star{} Et vénimus adoráre Dóminum.\end{Resp}
\begin{Resp}\Vsign Vídimus stellam ejus in Oriénte.\end{Resp}
\begin{Resp}\Rsign Et vénimus adoráre Dóminum.\end{Resp}
\begin{Resp}\Vsign Glória Patri, et Fílio, \Star{} et Spirítui Sancto.\end{Resp}
\begin{Resp}\Rsign Et vénimus adoráre Dóminum.\end{Resp}
\switchcolumn
\EN
\begin{Resp}\Rsign There came wise men from the east to Jerusalem, asking and saying: Where is He That is born King of the Jews? for we have seen His star in the east; \Star{} And are come to worship the Lord.\end{Resp}
\begin{Resp}\Vsign We have seen His star in the east.\end{Resp}
\begin{Resp}\Rsign And are come to worship the Lord.\end{Resp}
\begin{Resp}\Vsign Glory be to the Father, and to the Son, \Star{} and to the Holy Ghost.\end{Resp}
\begin{Resp}\Rsign And are come to worship the Lord.\end{Resp}
\switchcolumn*

\LA
\LessonHead{Lectio VII}
\switchcolumn
\EN
\LessonHead{Lesson VII}
\switchcolumn*

\LA
\LessonTitle{Léctio sancti Evangélii secúndum Joánnem}
\switchcolumn
\EN
\LessonTitle{From the Holy Gospel according to John}
\switchcolumn*

\LA
\Cite{Joannes 1:29-34}
\switchcolumn
\EN
\Cite{John 1:29-34}
\switchcolumn*

\LA
\noindent In illo témpore: Vidit Joánnes Jesum veniéntem ad se, et ait: Ecce Agnus Dei, ecce qui tollit peccátum mundi. Et réliqua.\par
\switchcolumn
\EN
\noindent At that time, John saw Jesus coming to him, and he saith: Behold the Lamb of God, behold him who taketh away the sin of the world. And so on.\par
\switchcolumn*

\LA
\LessonTitle{Homilía sancti Augustíni Epíscopi}
\switchcolumn
\EN
\LessonTitle{Homily by St. Augustine, Bishop (of Hippo.)}
\switchcolumn*

\LA
\Source{Tract. 6 in Joann., ante med.}
\switchcolumn
\EN
\Source{6th Tract on John}
\switchcolumn*

\LA
\noindent Antequam veníret Dóminus ut baptizarétur a Joánne in Jordáne, nóverat eum illis vócibus, ubi ait: Tu ad me venis baptizári? ego a te débeo baptizári. Sed ecce Dóminum nóverat, nóverat Fílium Dei. Unde probámus quod jam nóverat, quia ipse baptizáret in Spíritu Sancto? Antequam veníret ad flúvium, cum multi ad Joánnem cúrrerent baptizári, ait illis: Ego quidem baptízo vos in aqua: qui autem post me venit, major me est: cujus non sum dignus corrígiam calceaménti sólvere: ipse vos baptizábit in Spíritu Sancto, et igni. Jam et hoc nóverat.\par
\switchcolumn
\EN
\noindent John knew Jesus even before He came to be baptized of him in Jordan, as we perceive by the words I have need to be baptized of thee, and comest Thou to me? Behold, how he knew that He was the Lord, how he knew that He was the Son of God! How do we prove that he knew that He it was Who should baptize with the Holy Ghost? Before the Lord came to the river, when many betook themselves to John to be baptized of him, the Baptist said I indeed baptize you with water but One Mightier than I cometh; the latchet of Whose shoes I am not worthy to unloose He shall baptize you with the Holy Ghost and with fire. (Luke iii. 16.) Behold, John knew this also.\par
\switchcolumn*

\LA
\begin{Resp}\Rsign Stella, quam víderant Magi in Oriénte, antecedébat eos, donec venírent ad locum, ubi puer erat. \Star{} Vidéntes autem eam, gavísi sunt gáudio magno.\end{Resp}
\begin{Resp}\Vsign Et intrántes domum, invenérunt púerum cum María matre ejus, et procidéntes adoravérunt eum.\end{Resp}
\begin{Resp}\Rsign Vidéntes autem eam, gavísi sunt gáudio magno.\end{Resp}
\switchcolumn
\EN
\begin{Resp}\Rsign The star which the wise men had seen in the East, went before them, till they came where the young Child was. \Star{} And when they saw the star, they rejoiced with exceeding great joy.\end{Resp}
\begin{Resp}\Vsign And when they were come into the house, they found the young Child with Mary His Mother, and fell down and worshipped Him.\end{Resp}
\begin{Resp}\Rsign And when they saw the star, they rejoiced with exceeding great joy.\end{Resp}
\switchcolumn*

\LA
\LessonHead{Lectio VIII}
\switchcolumn
\EN
\LessonHead{Lesson VIII}
\switchcolumn*

\LA
\noindent Quid ergo per colúmbam dídicit, ne mendax póstea inveniátur, (quod avértat a nobis Deus opinári) nisi quamdam proprietátem in Christo talem futúram, ut quamvis multi minístri baptizatúri essent, sive justi, sive injústi, non tribuerétur sánctitas baptísmi nisi illi, super quem descéndit colúmba, de quo dictum est: Hic est qui baptízat in Spíritu Sancto? Petrus baptízet, hic est qui baptízat: Paulus baptízet, hic est qui baptízat: Judas baptízet, hic est qui baptízat. Nam si pro diversitáte meritórum baptísma sanctum est: quia divérsa sunt mérita, divérsa erunt baptísmata: et tanto quisque áliquid mélius putátur accípere, quanto a melióre vidétur accepísse.\par
\switchcolumn
\EN
\noindent Yet John saith: I knew Him Now, how are we to explain this without calling John a liar? and God forbid that we should ever even think anything of the kind. Was it not that when the Dove descended on Christ, John then for the first time knew Him to have that peculiar attribute, that, whosoever should baptize with His Baptism, whether they were themselves just or unjust, the virtue of the Sacrament should proceed, not from them, but from Him on Whom abode the Dove; so that He is the real Baptizer in every Christian Baptism until the end of time, and it is in this sense that it is said of Him the Same is He Which baptizeth with the Holy Ghost Whether it be Peter, or Paul, or Judas, that performeth the ceremony, the real Baptizer and effectual Worker is Christ. For if the holiness of the baptism depended on the holiness of the particular officiator, no two baptisms would be exactly alike, and every one would be supposed to be more or less regenerated according as the minister who baptized him was more or less of a saint.\par
\switchcolumn*

\LA
\begin{Resp}\Rsign Vidéntes stellam Magi, gavísi sunt gáudio magno: \Star{} Et intrántes domum invenérunt púerum cum María matre ejus, et procidéntes adoravérunt eum: * Et apértis thesáuris suis obtulérunt ei múnera, aurum, thus, et myrrham.\end{Resp}
\begin{Resp}\Vsign Stella, quam víderant Magi in Oriénte, antecedébat eos, usque dum véniens staret supra ubi erat puer.\end{Resp}
\begin{Resp}\Rsign Et intrántes domum invenérunt púerum cum María matre ejus, et procidéntes adoravérunt eum:\end{Resp}
\begin{Resp}\Vsign Glória Patri, et Fílio, \Star{} et Spirítui Sancto.\end{Resp}
\begin{Resp}\Rsign Et apértis thesáuris suis obtulérunt ei múnera, aurum, thus, et myrrham.\end{Resp}
\switchcolumn
\EN
\begin{Resp}\Rsign When the wise men saw the star, they rejoiced with exceeding great joy. and when they had opened their treasures, they presented unto Him gifts; gold, and frankincense, and myrrh. \Star{} And when they were come into the house, they found the young Child with Mary His Mother, and fell down and worshipped Him; *\end{Resp}
\begin{Resp}\Vsign The star which the wise men had seen in the East, went before them, till it came and stood over where the young Child was.\end{Resp}
\begin{Resp}\Rsign And when they were come into the house, they found the young Child with Mary His Mother, and fell down and worshipped Him.\end{Resp}
\begin{Resp}\Vsign Glory be to the Father, and to the Son, \Star{} and to the Holy Ghost.\end{Resp}
\begin{Resp}\Rsign And when they had opened their treasures, they presented unto Him gifts; gold, and frankincense, and myrrh.\end{Resp}
\switchcolumn*

\LA
\LessonHead{Lectio IX}
\switchcolumn
\EN
\LessonHead{Lesson IX}
\switchcolumn*

\LA
\noindent Ipsi sancti, (intellégite, fratres) boni pertinéntes ad colúmbam, pertinéntes ad sortem civitátis illíus Jerúsalem, ipsi boni in Ecclésia, de quibus dicit Apóstolus: Novit Dóminus qui sunt ejus; diversárum gratiárum sunt, non omnes pária mérita habent. Sunt álii áliis sanctióres, sunt álii áliis melióres. Quare ergo si unus ab illo, verbi grátia, justo, sancto baptizétur, álius ab álio inferióris mériti apud Deum, inferióris gradus, inferióris continéntiæ, inferióris vitæ; unum tamen et par et æquále est quod accepérunt, nisi quia hic est, qui baptízat?\par
\switchcolumn
\EN
\noindent Now, my brethren, understand me. The saints themselves, those good men who appertain to the Dove, those good men whose portion is in Jerusalem, those good men in the Church, of whom the Apostle saith: the Lord knoweth them that are His, (2 Tim. ii. 19,) these good men differ one from another by diversities of graces, and are not all of the same worthiness. Some are holier than others, and some are better than others. Supposing then for the sake of argument that A is baptized by B, a righteous saint, and C is baptized by D who is less worthy in the sight of God, who hath attained only a lower degree in godliness, who is not so chaste, and whose life is not so good as B's, yet A and C receive just the same thing. And how is this, unless it be that it is Christ Himself Who is the effectual Baptizer?\par
\switchcolumn*

\LA
\TeDeum{Te Deum}
\switchcolumn
\EN
\TeDeum{Te Deum}
\switchcolumn*

\end{paracol}

\par\needspace{10\baselineskip}\addvspace{1.2em}{\centering\small\textsc{Die 14 Januarii} · \textsc{14 January}\par}
\Office{S. Hilarii Episcopi Confessoris Ecclesiæ Doctoris}{St. Hilary, Bishop of Poitiers, Confessor and Doctor of the Church}{Duplex}
\addcontentsline{toc}{subsection}{Die 14 Januarii · S. Hilarii Episcopi Confessoris Ecclesiæ Doctoris\enspace\textit{St. Hilary, Bishop of Poitiers, Confessor and Doctor of the Church}}
\begin{paracol}{2}
\LA
\RefLine{Lectiones I–III: de Scriptura occurrente.}
\switchcolumn
\EN
\RefLine{Lessons I–III: from the Scripture of the day.}
\switchcolumn*

\LA
\RefLine{Lectio IX: de S. Felicis Presbyteri et Martyris (01-14cc).}
\switchcolumn
\EN
\RefLine{Lesson IX: of St. Felix, Priest and Martyr (01-14cc).}
\switchcolumn*

\LA
\LessonHead{Lectio IV}
\switchcolumn
\EN
\LessonHead{Lesson IV}
\switchcolumn*

\LA
\noindent Hilárius, in Aquitánia nóbili génere natus, doctrína et eloquéntia excélluit. Qui primum in matrimónio quasi mónachi vitam egit: deínde propter singuláres virtútes Pictavórum epíscopus creátur: quod munus episcopále sic gessit, ut a fidélibus summam laudem consequerétur. Quo témpore, cum terróribus, bonórum spoliatióne, exsílio, et omni crudelitáte Constántius imperátor Cathólicos vexáret, nisi ad Ariánas partes transírent; Hilárius tamquam firmíssimum murum se Ariánis oppónens, illórum furórem in se concitávit. Itaque multis petítus insídiis, tandem dolo Saturníni Arelaténsis epíscopi, de sýnodo Biterrénsi in Phrýgiam relegátus est: ubi et mórtuum suscitávit, et libros duódecim scripsit de Trinitáte contra Ariános.\par
\switchcolumn
\EN
\noindent This Hilary was born of a noble family in Aquitaine, and is distinguished as a divine and an orator. He was married in his earlier life, but even then lived the life of a monk and on account of his remarkable holiness, was ultimately made Bishop of Poitiers, (about the year 353,) in which office he did his duty so as to gain the universal praise of the faithful. At that time the Emperor Constantius was persecuting the Catholics by threats, by the plundering of their goods, by exile, and at length, by every species of cruelty, in order to force them to yield to the Arian heresy. Against the Arians Hilary set himself up as a brazen wall, and turned upon himself the fierceness of their anger. They assailed him by many artifices, and at last Saturnine, Bishop of Arles, at the Council of Beziers, (in 356,) procured his banishment to Phrygia. During this exile he raised a dead man to life, and wrote a work in twelve books on the Trinity, against the Arians.\par
\switchcolumn*

\LA
\begin{Resp}\Rsign Invéni David servum meum, óleo sancto meo unxi eum: \Star{} Manus enim mea auxiliábitur ei.\end{Resp}
\begin{Resp}\Vsign Nihil profíciet inimícus in eo, et fílius iniquitátis non nocébit ei.\end{Resp}
\begin{Resp}\Rsign Manus enim mea auxiliábitur ei.\end{Resp}
\switchcolumn
\EN
\begin{Resp}\Rsign I have found David My servant, with My holy oil have I anointed him \Star{} For My hand shall help him.\end{Resp}
\begin{Resp}\Vsign The enemy shall prevail nothing against him, nor the son of wickedness afflict him.\end{Resp}
\begin{Resp}\Rsign For My hand shall help him.\end{Resp}
\switchcolumn*

\LA
\LessonHead{Lectio V}
\switchcolumn
\EN
\LessonHead{Lesson V}
\switchcolumn*

\LA
\noindent Quadriénnio post coácto concílio ad Seléuciam, Isáuriæ urbem, Hilárius adésse compúlsus est: ac deínde Constantinópolim proféctus, ubi extrémum fídei perículum animadvértit, tribus libéllis públice datis, audiéntiam imperatóris popóscit, ut de fide cum adversáriis coram disputáret. Verum cum Ursácius et Valens Ariáni epíscopi, quos Hilárius scriptis confutárat, præséntis eruditiónem pertiméscerent, Constántio persuasérunt, ut spécie honóris eum in suum episcopátum restitúeret. Tunc Hilárium e prǽlio hæreticórum reverténtem, ut inquit sanctus Hierónymus, Galliárum ecclésia cómplexa est: quem ad episcopátum secútus est Martínus, qui póstea Turonénsi prǽfuit ecclésiæ; tantúmque illo doctóre profécit, quantum ejus póstea sánctitas declarávit.\par
\switchcolumn
\EN
\noindent In (the year 359,) the fourth of his exile, took place the Council of Seleucia in Isauria, at which Hilary was obliged to be present, but afterwards withdrew to Constantinople. Here he realised more sharply the awful nature of this crisis in the history of Christianity, published three pamphlets in the form of letters to the Emperor, and demanded from that Prince leave to hold a public disputation in his presence. The Arian Bishops Ursacius and Valens, whom Hilary had already confuted in writing, were afraid to meet him in debate, and therefore induced Constantius, under pretence of pardon, to send him back to his Bishopric, (in 360.) His mother, the Church of Gaul, to use the language of Jerome, received him with open arms on his return from the battle with the heretics. He was followed to Poitiers by Martin, afterwards Bishop of Tours, whose after holiness was a fruit of his teaching.\par
\switchcolumn*

\LA
\begin{Resp}\Rsign Pósui adjutórium super poténtem, et exaltávi eléctum de plebe mea: \Star{} Manus enim mea auxiliábitur ei.\end{Resp}
\begin{Resp}\Vsign Invéni David servum meum, óleo sancto meo unxi eum.\end{Resp}
\begin{Resp}\Rsign Manus enim mea auxiliábitur ei.\end{Resp}
\switchcolumn
\EN
\begin{Resp}\Rsign I have laid help upon one that is mighty, and have exalted one chosen out of My people \Star{} For My hand shall help him.\end{Resp}
\begin{Resp}\Vsign I have found David My servant, with My holy oil have I anointed him.\end{Resp}
\begin{Resp}\Rsign For My hand shall help him.\end{Resp}
\switchcolumn*

\LA
\LessonHead{Lectio VI}
\switchcolumn
\EN
\LessonHead{Lesson VI}
\switchcolumn*

\LA
\noindent Magna deínceps tranquillitáte Pictavórum ecclésiam administrávit: Galliámque univérsam addúxit, ut Arianórum impietátem condemnáret. Multos libros scripsit mira eruditióne; quos omnes sanctus Hierónymus ad Lætam, sine ulla erróris suspicióne legi posse testátur illis verbis: Hilárii libros inoffénso decúrrat pede. Migrávit in cælum Idibus Januárii, Valentiniáno et Valénte imperatóribus, anno post Christum natum trecentésimo sexagésimo nono.\par
\switchcolumn
\EN
\noindent Henceforth he ruled the Church of Poitiers in great peace. By his exertions the Church of Gaul was led to denounce the Arian blasphemy. His wonderful learning is seen in his numerous works, of which Jerome write to Laeta that he deemeth them quite faultless. One can follow Hilary in his books, saith he, without tripping once. He passed from earth to heaven upon the thirteenth day of January, in the year of our Lord 369, in the reign of the Emperors Valentinian and Valens.\par
\switchcolumn*

\LA
Eum a multis Pátribus et concíliis insígnem Ecclésiæ Doctórem nuncupátum, atque uti talem in áliquot diœcésibus cultum, tandem instánte sýnodo Burdigalénsi Pius nonus Póntifex máximus, ex sacrórum Rítuum Congregatiónis consúlto, universális Ecclésiæ Doctórem declarávit et confirmávit, ac ipsíus festo die Missam et Offícium de Doctóribus ab ómnibus recitári jussit.\par
\switchcolumn
\EN
He had already been called an illustrious Doctor of the Church, by many Fathers and Councils, and was so styled in the Liturgy in some Dioceses, when at length, (in the year 1850,) Pope Pius IX., at the prayer of the Synod of Bordeaux, and in accordance with a resolution of the Sacred Congregation of Rites, proclaimed and confirmed the title, and commanded that the Mass and Office of his Feast should be everywhere said as those of a Doctor.\par
\switchcolumn*

\LA
\begin{Resp}\Rsign Iste est, qui ante Deum magnas virtútes operátus est, et omnis terra doctrína ejus repléta est: \Star{} Ipse intercédat pro peccátis ómnium populórum.\end{Resp}
\begin{Resp}\Vsign Iste est, qui contémpsit vitam mundi, et pervénit ad cæléstia regna.\end{Resp}
\begin{Resp}\Rsign Ipse intercédat pro peccátis ómnium populórum.\end{Resp}
\begin{Resp}\Vsign Glória Patri, et Fílio, \Star{} et Spirítui Sancto.\end{Resp}
\begin{Resp}\Rsign Ipse intercédat pro peccátis ómnium populórum.\end{Resp}
\switchcolumn
\EN
\begin{Resp}\Rsign This is he which wrought great wonders before God, and the whole earth is full of his teaching \Star{} May he pray for all people, that their sins may be forgiven unto them\end{Resp}
\begin{Resp}\Vsign This is he which loved not his life in this world, and hath attained unto the kingdom of heaven.\end{Resp}
\begin{Resp}\Rsign May he pray for all people, that their sins may be forgiven unto them\end{Resp}
\begin{Resp}\Vsign Glory be to the Father, and to the Son, \Star{} and to the Holy Ghost.\end{Resp}
\begin{Resp}\Rsign May he pray for all people, that their sins may be forgiven unto them\end{Resp}
\switchcolumn*

\LA
\LessonHead{Lectio VII}
\switchcolumn
\EN
\LessonHead{Lesson VII}
\switchcolumn*

\LA
\LessonTitle{Léctio sancti Evangélii secúndum Matthǽum}
\switchcolumn
\EN
\LessonTitle{From the Holy Gospel according to Matthew}
\switchcolumn*

\LA
\Cite{Matt 5:13-19}
\switchcolumn
\EN
\Cite{Matt 5:13-19}
\switchcolumn*

\LA
\noindent In illo témpore: Dixit Jesus discípulis suis: Vos estis sal terræ. Quod si sal evanúerit, in quo saliétur? Et réliqua.\par
\switchcolumn
\EN
\noindent At that time, Jesus said unto His disciples: Ye are the salt of the earth. But if the salt have lost his savour, wherewith shall it be salted. And so on.\par
\switchcolumn*

\LA
\LessonTitle{Homilía sancti Hilárii Epíscopi}
\switchcolumn
\EN
\LessonTitle{Homily by St. Hilary, Bishop of Poitiers.}
\switchcolumn*

\LA
\Source{Comm. in Matthæum, cap. 4, n. 10}
\switchcolumn
\EN
\Source{Comment, on Matth. V}
\switchcolumn*

\LA
\noindent Vos estis sal terræ: Quod si sal infatuátum fúerit, ad níhilum valet id quod saliétur. Sal, ut árbitror, terræ nullum est. Quómodo ergo Apóstolos sal terræ nuncupávit? Sed propríetas est quærénda dictórum, quam et Apostolórum offícium, et ipsíus salis natúra monstrábit. Sal est in se uno cóntinens aquæ et ignis eleméntum, et hoc ex duóbus est unum.\par
\switchcolumn
\EN
\noindent “Ye are the salt of the earth. But if the salt have lost his savour, wherewith shall it be salted It is thenceforth good for nothing, but to be cast out, and to be trodden under foot of men.” There is, I take it, no such thing as salt of the earth. How, then, can the Apostles be called the salt of the earth? But the true meaning of these words will be made plain, when we consider the duty of Apostles, and the nature of salt itself. Now, salt is a compound of the elements of water and fire, out of the which two things in salt there is made one.\par
\switchcolumn*

\LA
\begin{Resp}\Rsign Amávit eum Dóminus, et ornávit eum: stolam glóriæ índuit eum, \Star{} Et ad portas paradísi coronávit eum.\end{Resp}
\begin{Resp}\Vsign Induit eum Dóminus lorícam fídei, et ornávit eum.\end{Resp}
\begin{Resp}\Rsign Et ad portas paradísi coronávit eum.\end{Resp}
\switchcolumn
\EN
\begin{Resp}\Rsign The Lord loved him and beautified him He clothed him with a robe of glory \Star{} And crowned him at the gates of Paradise.\end{Resp}
\begin{Resp}\Vsign The Lord hath put on him the breast-plate of faith, and hath adorned him.\end{Resp}
\begin{Resp}\Rsign And crowned him at the gates of Paradise.\end{Resp}
\switchcolumn*

\LA
\LessonHead{Lectio VIII}
\switchcolumn
\EN
\LessonHead{Lesson VIII}
\switchcolumn*

\LA
\noindent Hic ígitur in omnem usum humáni géneris efféctus, incorruptiónem corpóribus, quibus fúerit aspérsus, impértit, et ad omnem sensum cónditi sapóris aptíssimus est. Apóstoli autem sunt rerum cæléstium prædicatóres, et æternitátis velut satóres, immortalitátem ómnibus corpóribus, quibus eórum sermo aspérsus fúerit, conferéntes. Mérito ígitur sal terræ nuncupáti sunt, per doctrínæ virtútem saliéndi modo æternitáti córpora reservántes.\par
\switchcolumn
\EN
\noindent This thing, therefore, thus made to serve in diverse ways the use of men, doth keep from corruption bodies whereon it is sprinkled, and doth readily yield to all the senses the perception of its inborn savour. And thus are the Apostles, seeing that they are the preachers of the kingdom of heaven, and in a certain sense the sowers of the seed of life everlasting, since that Word of God which they scatter hath power to make this mortal put on immortality. Meetly then are they called salt, the savour of whose teaching doth keep sweet the receiver thereof even unto life everlasting.\par
\switchcolumn*

\LA
\begin{Resp}\Rsign In médio Ecclésiæ apéruit os ejus, \Star{} Et implévit eum Dóminus spíritu sapiéntiæ et intelléctus.\end{Resp}
\begin{Resp}\Vsign Jucunditátem et exsultatiónem thesaurizávit super eum.\end{Resp}
\begin{Resp}\Rsign Et implévit eum Dóminus spíritu sapiéntiæ et intelléctus.\end{Resp}
\begin{Resp}\Vsign Glória Patri, et Fílio, \Star{} et Spirítui Sancto.\end{Resp}
\begin{Resp}\Rsign Et implévit eum Dóminus spíritu sapiéntiæ et intelléctus.\end{Resp}
\switchcolumn
\EN
\begin{Resp}\Rsign In the midst of the congregation did the Lord open his mouth. \Star{} And filled him with the spirit of wisdom and understanding.\end{Resp}
\begin{Resp}\Vsign He made him rich with joy and gladness.\end{Resp}
\begin{Resp}\Rsign And filled him with the spirit of wisdom and understanding.\end{Resp}
\begin{Resp}\Vsign Glory be to the Father, and to the Son, \Star{} and to the Holy Ghost.\end{Resp}
\begin{Resp}\Rsign And filled him with the spirit of wisdom and understanding.\end{Resp}
\switchcolumn*

\end{paracol}

\par\needspace{10\baselineskip}\addvspace{1.2em}{\centering\small\textsc{Die 14 Januarii} · \textsc{14 January}\par}
\Office{S. Felicis Presbyteri et Martyris}{St. Felix, Priest and Martyr}{Simplex}
\addcontentsline{toc}{subsection}{Die 14 Januarii · S. Felicis Presbyteri et Martyris\enspace\textit{St. Felix, Priest and Martyr}}
\begin{paracol}{2}
\LA
\LessonHead{Lectio IX (pro commemoratione)}
\switchcolumn
\EN
\LessonHead{Lesson IX (when commemorated)}
\switchcolumn*

\LA
\LessonTitle{Commemoratio: S. Felicis Presbyteri et Martyris}
\switchcolumn
\EN
\LessonTitle{Commemoration of: St. Felix, Priest and Martyr}
\switchcolumn*

\LA
\noindent Felix Nolanus presbyter, cum in idola vehementius inveherétur, ab infidelibus varie vexatus, in carcerem conjícitur. Unde ab Angelo nocte edúctus, quærere jussus est Maximianum Nolæ episcopum: qui cum senio confectus desperáret se ferre posse supplícia persequéntium, se abdíderat in silvam. Quo cum Felix Deo duce pervenísset, sanctum episcopum humi jacéntem pene mortuum videt; quem recreátum ac sublátum in humeros, apud fidélem víduam reficiéndum curávit. Sed cum is íterum idolórum cultores impietátis argúeret, facto in ipsum ímpetu, fúgiens in angusto duórum paríetum intervallo se occultávit; qui áditus cum repénte araneárum telis pertextus visus esset, némini recentis látebræ suspiciónem relíquit. Inde ígitur evádens Felix in ædibus piæ mulíeris tres menses látuit. Cum vero Dei Ecclésia requiéscere cœpísset, Nolam rédiens, multisque ibi vitæ exemplis, et doctrinæ præceptis, miraculisque ad Christi fidem conversis, constanter étiam recusato ejus urbis episcopatu, obdormívit in Dómino, sepultúsque est prope Nolam in loco, quem in Pincis appellábant.\par
\switchcolumn
\EN
\noindent Felix was a Priest of Nola, who on account of his fiery zeal against idolatry, suffered much persecution from the heathens, and was cast into prison. From thence an angel delivered him by night, and bade him go to Maximian, Bishop of Nola. This Bishop, enfeebled by old age, had at length despaired of power to withstand the torments of the persecutors, and had hidden, himself in a wood. Thither came Felix, by the will of God, and found the holy bishop lying half-dead upon the ground. He succoured him, and carried him upon his shoulders to the house of an holy widow. On another occasion, Felix, having again provoked the anger of the devil-worshippers, became an object of their pursuit, from which he hid himself in a narrow place between two walls. Hardly had he entered, when some spiders wove their webs across the entrance, which the enemy perceiving, concluded that no man had entered, and passed by. After leaving this hiding-place, Felix lay for three months in the house of an holy woman. After the Lord gave peace to His Church, the Saint returned to Nola, where he turned many to Christ by his life, his preaching, and his miracles. He steadily refused to accept the Bishopric, fell asleep in the Lord, and was buried at Nola in the place called The Pines.\par
\switchcolumn*

\LA
\TeDeum{Te Deum}
\switchcolumn
\EN
\TeDeum{Te Deum}
\switchcolumn*

\end{paracol}

\par\needspace{10\baselineskip}\addvspace{1.2em}{\centering\small\textsc{Die 15 Januarii} · \textsc{15 January}\par}
\Office{S. Pauli Primi Eremitæ et Confessoris}{St. Paul the First Hermit, Confessor}{Duplex}
\addcontentsline{toc}{subsection}{Die 15 Januarii · S. Pauli Primi Eremitæ et Confessoris\enspace\textit{St. Paul the First Hermit, Confessor}}
\begin{paracol}{2}
\LA
\RefLine{Lectiones I–III: de Scriptura occurrente.}
\switchcolumn
\EN
\RefLine{Lessons I–III: from the Scripture of the day.}
\switchcolumn*

\LA
\RefLine{Lectio IX: de S. Mauri Abbatis (01-15cc).}
\switchcolumn
\EN
\RefLine{Lesson IX: of S. Maur, Abbot (01-15cc).}
\switchcolumn*

\LA
\LessonHead{Lectio IV}
\switchcolumn
\EN
\LessonHead{Lesson IV}
\switchcolumn*

\LA
\noindent Paulus, eremitárum auctor et magister, apud inferiórem Thebáidem natus, cum quindecim esset annórum, orbátus paréntibus est. Qui postea declinandæ causa persecutiónis Decii et Valeriáni, et Deo liberius inserviéndi, in eremi spelúncam se cóntulit: ubi, palma ei victum et vestítum præbente, vixit ad centésimum et décimum tertium annum, quo témpore ab Antonio nonagenario Dei admónitu invísitur. Quibus inter se, cum ántea non nossent, proprio nómine consalutántibus, et multa de regno Dei colloquéntibus, corvus, qui ántea semper Paulo dimidiátum panem attulerat, integrum détulit.\par
\switchcolumn
\EN
\noindent Paul the first hermit, and model of all others, was a native of the lower Theba'id in Egypt. At the age of fifteen years he lost both his parents. In order to escape from the persecution of Decius and Valerian, and to serve God in quietness, he betook himself to a cave in the desert. Here there was a palm-tree, on the fruit of which he lived, and of whose leaves he made his raiment until he attained the age of 113 years. At that time * Anthony, being now himself aged\par
90 Years, received a command from God to go and see him. They met without knowing one another's names, and saluted one another, after which they fell straightway into a long discourse concerning the kingdom of God. Now it so happened that a raven had of a long time brought Paul every day half a loaf, but on this day while they spake together he brought a whole one.\par
\switchcolumn*

\LA
\begin{Resp}\Rsign Honéstum fecit illum Dóminus, et custodívit eum ab inimícis, et a seductóribus tutávit illum: \Star{} Et dedit illi claritátem ætérnam.\end{Resp}
\begin{Resp}\Vsign Justum dedúxit Dóminus per vias rectas, et osténdit illi regnum Dei.\end{Resp}
\begin{Resp}\Rsign Et dedit illi claritátem ætérnam.\end{Resp}
\switchcolumn
\EN
\begin{Resp}\Rsign The Lord made him honourable, and defended him from his enemies, and kept him safe from those that lay in wait for him; \Star{} And gave him perpetual glory.\end{Resp}
\begin{Resp}\Vsign He went down with him into the pit, and left him not in bonds.\end{Resp}
\begin{Resp}\Rsign And gave him perpetual glory.\end{Resp}
\switchcolumn*

\LA
\LessonHead{Lectio V}
\switchcolumn
\EN
\LessonHead{Lesson V}
\switchcolumn*

\LA
\noindent Post corvi discessum, Eja, inquit Paulus, Dóminus nobis prandium misit, vere pius, vere miséricors. Sexagínta jam anni sunt, cum accipio quotídie dimidii panis fragméntum, nunc ad advéntum tuum milítibus suis Christus duplicávit annonam. Quare cum gratiárum actióne ad fontem capiéntes cibum, ubi tantisper recreáti sunt, íterum gratiis de more Deo actis, noctem in divinis laudibus consumpsérunt. Dilúculo Paulus de morte, quæ sibi instaret, ádmonens Antónium, hortátur, ut pallium, quod ab Athanasio acceperat, ad involvéndum suum corpus afferret. Quo ex itinere rédiens ille, vidit inter Angelórum choros, inter Prophetárum et Apostolórum cœtus Pauli ánimam in cælum ascendere.\par
\switchcolumn
\EN
\noindent When the raven had flown away Well, quoth Paul, the Lord hath sent us our dinner. Truly He is gracious; truly He is merciful. It is now sixty years that I have had half a loaf of bread every day, but now that thou art come, Christ giveth His soldiers double rations. Then they asked a blessing, and ate together, sitting by a spring. When they were refreshed, they returned thanks, as is the custom, and afterwards spent the whole night praising God. At break of day Paul felt the approach of death, and desired Anthony to bring the cloak which Athanasius had given him, to use for his winding-sheet. While Anthony was on his way back from this journey, he saw in a vision the soul of Paul ascending to heaven, surrounded by choirs of angels, and accompanied by the Prophets and Apostles.\par
\switchcolumn*

\LA
\begin{Resp}\Rsign Amávit eum Dóminus, et ornávit eum: stolam glóriæ índuit eum, \Star{} Et ad portas paradísi coronávit eum.\end{Resp}
\begin{Resp}\Vsign Índuit eum Dóminus lorícam fídei, et ornávit eum.\end{Resp}
\begin{Resp}\Rsign Et ad portas paradísi coronávit eum.\end{Resp}
\switchcolumn
\EN
\begin{Resp}\Rsign The Lord loved him and beautified him; He clothed him with a robe of glory, \Star{} And crowned him at the gates of Paradise.\end{Resp}
\begin{Resp}\Vsign The Lord hath put on him the breast-plate of faith, and hath adorned him.\end{Resp}
\begin{Resp}\Rsign And crowned him at the gates of Paradise.\end{Resp}
\switchcolumn*

\LA
\LessonHead{Lectio VI}
\switchcolumn
\EN
\LessonHead{Lesson VI}
\switchcolumn*

\LA
\noindent Cumque ad ejus cellam pervenísset, invénit genibus complicatis, erecta cervice, exténsisque in altum mánibus corpus exanime: quod pallio obvolvens, hymnosque et psalmos ex christiana traditióne decantans, cum sárculum, quo terram fóderet, non haberet, duo leones ex interiore erémo rápido cursu ad beáti senis corpus ferúntur; ut facile intelligerétur, eos, quo modo póterant, plorátum édere: qui certatim terram pédibus effodiéntes, fóveam, quæ hóminem commode cáperet, effecérunt. Qui cum abiísent, Antonius sanctum corpus in eum locum íntulit; et injecta humo, tumulum ex christiáno more compósuit: túnicam vero Pauli, quam in sportæ modum ex palmæ foliis ille sibi contexúerat, secum áuferens, eo vestítu diébus solemnibus Paschæ et Pentecostes, quoad vixit, usus est.\par
\switchcolumn
\EN
\noindent When Anthony reached the cell of Paul, he found the dead body of the Saint in a kneeling posture, with the head thrown up and the hands stretched out towards heaven. He immediately began to chant the psalms and hymns ordained by Christian tradition, while he wrapped the body in the cloak of Athanasius. He had no spade to dig a grave, but two lions came roaring from the desert, as though to attend the burying, and scratched a hole big enough to hold a man's body, with their paws, showing meanwhile such signs of grief as their nature alloweth. When they were gone away, Anthony put the holy body in this hole, covered it with earth, and arranged it like a Christian's grave. He took away for himself Paul's tunic, which he had woven out of the palm-leaves somewhat after the manner of basket-work, and this tunic Anthony was in the habit of wearing on the great days of the Passover and Pentecost as long as he lived.\par
\switchcolumn*

\LA
\begin{Resp}\Rsign Iste homo perfécit ómnia quæ locútus est ei Deus, et dixit ad eum: Ingrédere in réquiem meam: \Star{} Quia te vidi justum coram me ex ómnibus géntibus.\end{Resp}
\begin{Resp}\Vsign Iste est, qui contémpsit vitam mundi, et pervénit ad cæléstia regna.\end{Resp}
\begin{Resp}\Rsign Quia te vidi justum coram me ex ómnibus géntibus.\end{Resp}
\begin{Resp}\Vsign Glória Patri, et Fílio, \Star{} et Spirítui Sancto.\end{Resp}
\begin{Resp}\Rsign Quia te vidi justum coram me ex ómnibus géntibus.\end{Resp}
\switchcolumn
\EN
\begin{Resp}\Rsign This is he which did according unto all that God commanded him; and God said unto him: Enter thou into My rest, \Star{} For thee have I seen righteous before Me among all people.\end{Resp}
\begin{Resp}\Vsign This is he which loved not his life in this world, and is come unto an everlasting kingdom.\end{Resp}
\begin{Resp}\Rsign For thee have I seen righteous before Me among all people.\end{Resp}
\begin{Resp}\Vsign Glory be to the Father, and to the Son, \Star{} and to the Holy Ghost.\end{Resp}
\begin{Resp}\Rsign For thee have I seen righteous before Me among all people.\end{Resp}
\switchcolumn*

\LA
\LessonHead{Lectio VII}
\switchcolumn
\EN
\LessonHead{Lesson VII}
\switchcolumn*

\LA
\LessonTitle{Léctio sancti Evangélii secúndum Matthǽum}
\switchcolumn
\EN
\LessonTitle{From the holy Gospel according to Matthew}
\switchcolumn*

\LA
\Cite{Matt 11:25-30}
\switchcolumn
\EN
\Cite{Matt 11:25-30}
\switchcolumn*

\LA
\noindent In illo témpore: Respóndens Jesus dixit: Confíteor tibi, Pater, Dómine cæli et terræ, quia abscondísti hæc a sapiéntibus et prudéntibus, et revelásti ea párvulis. Et réliqua.\par
\switchcolumn
\EN
\noindent At that time, Jesus answered and said: I confess to thee, O Father, Lord of heaven and earth, because thou hast hid these things from the wise and prudent, and hast revealed them to the little ones. And so on.\par
\switchcolumn*

\LA
\LessonTitle{Homilía sancti Augustíni Epíscopi}
\switchcolumn
\EN
\LessonTitle{Homily by St. Augustine, Bishop (of Hippo.)}
\switchcolumn*

\LA
\Source{Sermo 10 de verbis Domini}
\switchcolumn
\EN
\Source{10th Sermon on the Words of the Lord}
\switchcolumn*

\LA
\noindent Veníte ad me, omnes qui laborátis. Quare enim omnes laborámus, nisi quia sumus hómines mortáles, frágiles, infírmi, lútea vasa portántes, quæ fáciunt ínvicem angústias? Sed, si angustiántur vasa carnis, dilaténtur spátia caritátis. Quid ergo dicit, Veníte ad me, omnes qui laborátis, nisi ut non laborétis? Dénique promíssio ejus in promptu est; quóniam laborántes vocávit, quærent forte, qua mercéde vocáti sunt. Et ego vos, inquit, refíciam. Tóllite jugum meum super vos, et díscite a me, non mundum fabricáre, non cuncta visibília et invisibília creáre, non in ipso mundo mirabília fácere et mórtuos suscitáre; sed, Quóniam mitis sum et húmilis corde.\par
\switchcolumn
\EN
\noindent Come unto Me, all ye that labour! And wherefore labour we all, but because we are frail, sickly, dying creatures, burdened with earthen vessels which distress us? But if these fleshly vessels be distressful, let the open expanse of love be free and wide. Come unto Me, all ye that labour! and why? That we may labour no more. His promise is an instant promise, for He calleth such as are labouring. Perchance they will ask Him what shall be their reward? And I, saith He, will give you rest. Take My yoke upon you, and learn of me not how to make the world, not how to create all things visible and invisible, not to work wonders in the earth, nor to raise the dead but for I am meek and lowly in heart.\par
\switchcolumn*

\LA
\begin{Resp}\Rsign Iste est qui ante Deum magnas virtútes operátus est, et de omni corde suo laudávit Dóminum: \Star{} Ipse intercédat pro peccátis ómnium populórum.\end{Resp}
\begin{Resp}\Vsign Ecce homo sine queréla, verus Dei cultor, ábstinens se ab omni ópere malo, et pérmanens in innocéntia sua.\end{Resp}
\begin{Resp}\Rsign Ipse intercédat pro peccátis ómnium populórum.\end{Resp}
\switchcolumn
\EN
\begin{Resp}\Rsign This is he which wrought great wonders before God, and praised the Lord with all his heart. \Star{} May he pray for all people, that their sins may be forgiven unto them.\end{Resp}
\begin{Resp}\Vsign Behold a man without blame, a worshipper of God in truth, keeping himself clean from every evil work, and abiding still in his innocency.\end{Resp}
\begin{Resp}\Rsign May he pray for all people, that their sins may be forgiven unto them.\end{Resp}
\switchcolumn*

\LA
\LessonHead{Lectio VIII}
\switchcolumn
\EN
\LessonHead{Lesson VIII}
\switchcolumn*

\LA
\noindent Magnus esse vis? a mínimo íncipe. Cógitas magnam fábricam constrúere celsitúdinis? de fundaménto prius cógita humilitátis. Et quantam quisque vult et dispónit superimpónere molem ædifícii, quanto erit majus ædifícium, tanto áltius fodit fundaméntum. Et fábrica quidem cum constrúitur, in supérna consúrgit; qui autem fodit fundaméntum, ad ima deprímitur. Ergo et fábrica ante celsitúdinem humiliátur, et fastígium post humiliatiónem erígitur.\par
\switchcolumn
\EN
\noindent Wilt thou be great? Begin by being little. Dost thou think to raise up a lofty building? Then lay the foundations thereof in lowliness. The greater soever, and the more massy, be that which any man thinketh to build, so much the deeper doth he dig his foundation. And when the house is built, it towereth heavenward; but he which layeth the foundation goeth down into the earth. The building, therefore, is low before it is high, and, after it is low, it riseth high to the roof.\par
\switchcolumn*

\LA
\begin{Resp}\Rsign Sint lumbi vestri præcíncti, et lucérnæ ardéntes in mánibus vestris: \Star{} Et vos símiles homínibus exspectántibus dóminum suum, quando revertátur a núptiis.\end{Resp}
\begin{Resp}\Vsign Vigiláte ergo, quia nescítis qua hora Dóminus vester ventúrus sit.\end{Resp}
\begin{Resp}\Rsign Et vos símiles homínibus exspectántibus dóminum suum, quando revertátur a núptiis.\end{Resp}
\begin{Resp}\Vsign Glória Patri, et Fílio, \Star{} et Spirítui Sancto.\end{Resp}
\begin{Resp}\Rsign Et vos símiles homínibus exspectántibus dóminum suum, quando revertátur a núptiis.\end{Resp}
\switchcolumn
\EN
\begin{Resp}\Rsign Let your loins be girded about, and your lights burning; \Star{} And ye yourselves like unto men that wait for their lord, when he will return from the wedding.\end{Resp}
\begin{Resp}\Vsign Watch therefore, for ye know not what hour your Lord doth come.\end{Resp}
\begin{Resp}\Rsign And ye yourselves like unto men that wait for their lord, when he will return from the wedding.\end{Resp}
\begin{Resp}\Vsign Glory be to the Father, and to the Son, \Star{} and to the Holy Ghost.\end{Resp}
\begin{Resp}\Rsign And ye yourselves like unto men that wait for their lord, when he will return from the wedding.\end{Resp}
\switchcolumn*

\end{paracol}

\par\needspace{10\baselineskip}\addvspace{1.2em}{\centering\small\textsc{Die 15 Januarii} · \textsc{15 January}\par}
\Office{S. Mauri Abbatis}{S. Maur, Abbot}{Simplex}
\addcontentsline{toc}{subsection}{Die 15 Januarii · S. Mauri Abbatis\enspace\textit{S. Maur, Abbot}}
\begin{paracol}{2}
\LA
\LessonHead{Lectio IX (pro commemoratione)}
\switchcolumn
\EN
\LessonHead{Lesson IX (when commemorated)}
\switchcolumn*

\LA
\LessonTitle{Commemoratio: S. Mauri Abbatis}
\switchcolumn
\EN
\LessonTitle{Commemoration of: S. Maur, Abbot}
\switchcolumn*

\LA
\noindent Maurus nobilis Romanus, puer a patre Eutychio Deo sub sancti Benedicti disciplína oblatus, brevi tantum divina grátia profecit, ut ipsi magistro admiratióni esset; qui illum sæpe, véluti regularis observantiæ et virtútum ómnium specimen, ceteris discípulis ad imitándum proponebat. Cujus adhuc adolescentis illud admirábilis obediéntiæ exemplum a sancto Gregório Papa commemorátur. Nam, cum Placidus mónachus in lacum prolapsus, aquárum ímpetu raperétur, sancti Patris jussu accurrens Maurus, et super aquas incédens, socium capíllis apprehensum, ad terram attraxit. Missus in Gálliam ab eodem sancto Benedicto, célebri monasterio exstructo, cui annos quadragínta præfuit, monasticam disciplínam mirifice propagávit. Denique sanctitáte et miraculis clarus, septuagenario major migrávit in cælum, anno salútis quingentésimo sexagésimo quinto.\par
\switchcolumn
\EN
\noindent Maurus was born of a noble Roman family, and while he was yet a child was offered to God by his father Eutychius, in the order and under the personal teaching of St. Benedict. In a short while he made such progress in the life of grace that he became a wonder to his master, who often held him up to his other disciples as a pattern of regular observance and all virtues. While he was yet very young, Pope St. Gregory telleth a wonderful instance of his obedience. Placid the monk having fallen into a lake where he was being swept away by the current, the holy Patriarch called Maurus and bade him run to the rescue, which he did, walking on the water till he reached Placid, whom he took by the hair of the head and dragged to the shore. He was sent by St. Benedict into France, where he founded the celebrated monastery (of Glanfeuil, now called St. Maur-sur-Loire,) which he governed for forty years. He was a zealous and successful propagator of monastic discipline. He passed to heaven, famous for holiness and miracles, when he was more than seventy years of age, in the year of Salvation 565.\par
\switchcolumn*

\LA
\TeDeum{Te Deum}
\switchcolumn
\EN
\TeDeum{Te Deum}
\switchcolumn*

\end{paracol}

\par\needspace{10\baselineskip}\addvspace{1.2em}{\centering\small\textsc{Die 16 Januarii} · \textsc{16 January}\par}
\Office{S. Marcelli Papæ et Martyris}{St. Marcellus, Pope and Martyr}{Semiduplex}
\addcontentsline{toc}{subsection}{Die 16 Januarii · S. Marcelli Papæ et Martyris\enspace\textit{St. Marcellus, Pope and Martyr}}
\begin{paracol}{2}
\LA
\RefLine{Lectiones I–III: de Scriptura occurrente.}
\switchcolumn
\EN
\RefLine{Lessons I–III: from the Scripture of the day.}
\switchcolumn*

\LA
\RefLine{Lectiones VII–IX: ex Commúni unius Summi Pontificis et Martyris (C2b).}
\switchcolumn
\EN
\RefLine{Lessons VII–IX: from the Common of a Single Martyr Pope (C2b).}
\switchcolumn*

\LA
\LessonHead{Lectio IV}
\switchcolumn
\EN
\LessonHead{Lesson IV}
\switchcolumn*

\LA
\noindent Marcellus Romanus, a Constantio et Galerio usque ad Maxéntium pontificátum gessit. Cujus hortatu Lucina matrona Romana, bonórum suórum Dei Ecclésiam fecit heredem. Aucto in Urbe fidelium numero, ad eórum utilitátem, ad baptismum pœnitentiámque dandam eis, qui christianam religiónem suscíperent, et ad Mártyrum sepultúram, novos Titulos instituit, et quasi álteras diœcéses distríbuit. Quibus rebus ira incénsus Maxentius, Marcello gravia supplícia minátur, nisi, depósito pontificatu, idolis immoláret.\par
\switchcolumn
\EN
\noindent This Marcellus was a Roman, and held the supreme Pontificate from (the year of our Lord 304, in) the reign of Constantius and Galerius, till (310, in) that of Maxentius. It was through his persuasion that the Roman lady Lucina left the whole of her property to the Church of God. As the believers increased, he instituted new titles in the City, which he divided after the manner of dioceses for their convenience, and for the baptism and penance of heathens converted to Christianity, and for the burial of the martyrs. These proceedings excited the wrath of Maxentius, who threatened Marcellus with the heaviest punishment, unless he would lay down the Popedom and sacrifice to idols.\par
\switchcolumn*

\LA
\begin{Resp}\Rsign Honéstum fecit illum Dóminus, et custodívit eum ab inimícis, et a seductóribus tutávit illum: \Star{} Et dedit illi claritátem ætérnam.\end{Resp}
\begin{Resp}\Vsign Descendítque cum illo in fóveam, et in vínculis non derelíquit eum.\end{Resp}
\begin{Resp}\Rsign Et dedit illi claritátem ætérnam.\end{Resp}
\switchcolumn
\EN
\begin{Resp}\Rsign The Lord made him honourable, and defended him from his enemies, and kept him safe from those that lay in wait for him, \Star{} And gave him perpetual glory.\end{Resp}
\begin{Resp}\Vsign He went down with him into the pit, and left him not in bonds.\end{Resp}
\begin{Resp}\Rsign And gave him perpetual glory.\end{Resp}
\switchcolumn*

\LA
\LessonHead{Lectio V}
\switchcolumn
\EN
\LessonHead{Lesson V}
\switchcolumn*

\LA
\noindent Qui cum insanas hóminis voces neglígeret, misit eum in catábulum, ut bestiárum, quæ publice alebántur, curam sustinéret: ubi Marcéllus assiduis jejuniis et precibus novem menses vitam duxit, parochias, quas præsens non póterat, visitans per epistolas. Inde ereptus a clericis, hospítio recípitur a beáta Lucina, in cujus ædibus ecclésiam dedicávit, quæ hódie titulo sancti Marcelli nominátur: in qua et Christiáni orábant, et ipse beátus Marcellus prædicabat.\par
\switchcolumn
\EN
\noindent The servant of God treated with contempt the mad cries of this man, who accordingly took him and sent him to a menagerie, to take care of the beasts which were fed at the public cost. Marcellus remained at this place for nine months, which he spent in continual fasting and prayer, and, as he could not visit the parishes in person, he wrote letters to them. Some clerks rescued him, and the blessed Lucina hospitably received him into her house, in which he dedicated a Church, which is now called St. Marcellus'. Here the Christians met to pray, and the blessed Marcellus himself preached.\par
\switchcolumn*

\LA
\begin{Resp}\Rsign Desidérium ánimæ ejus tribuísti ei Dómine, \Star{} Et voluntáte labiórum ejus non fraudásti eum.\end{Resp}
\begin{Resp}\Vsign Quóniam prævenísti eum in benedictiónibus dulcédinis: posuísti in cápite ejus corónam de lápide pretióso.\end{Resp}
\begin{Resp}\Rsign Et voluntáte labiórum ejus non fraudásti eum.\end{Resp}
\switchcolumn
\EN
\begin{Resp}\Rsign Lord, Thou hast given him his heart's desire. \Star{} And hast not withholden the request of his lips.\end{Resp}
\begin{Resp}\Vsign For Thou hast prevented him with the blessings of sweetness: Thou hast set a crown of precious stones upon his head.\end{Resp}
\begin{Resp}\Rsign And hast not withholden the request of his lips.\end{Resp}
\switchcolumn*

\LA
\LessonHead{Lectio VI}
\switchcolumn
\EN
\LessonHead{Lesson VI}
\switchcolumn*

\LA
\noindent Quibus cógnitis, Maxéntius in eam ecclésiam catábuli bestias transférri, et a Marcello custodiri jubet: ubi loci fœditate, multisque ærumnis afflíctus obdormívit in Dómino. Cujus corpus in cœmeterio Priscillæ via Salaria a beáta Lucina sepúltum est décimo septimo Kalendas Februarii. Sedit annos quinque, mensem unum, dies viginti quinque. Scripsit epistolam ad episcopos Antiochenæ provinciæ de primatu Romanæ Ecclésiæ, quam caput ecclesiárum appellandam demonstrat; ubi étiam illud scriptum est, nullum concílium jure celebrari, nisi ex auctoritáte Romani Pontificis. Ordinávit mense Decembri Romæ presbyteros viginti quinque, diaconos duos, episcopos per diversa loca viginti unum.\par
\switchcolumn
\EN
\noindent These proceedings came to the knowledge of Maxentius, who thereupon had the wild beasts brought from the menagerie and located in the church, where Marcellus was made to feed them. The noisomeness of the place and the filthiness of his occupation broke down a constitution already enfeebled by many ailments, and he fell asleep in the Lord. The blessed Lucina buried his body in the cemetery of Priscilla, on the Salarian Way, on the 16th of January. He sat on the throne of Peter for five years, one month, and twenty-five days. He wrote an epistle to the Bishops of the Patriarchate of Antioch on the primacy of the Roman Church, wherein he proveth the right of the same Church to be called the head of all the Churches. In this letter he likewise saith that no Council can be lawfully gathered together except by the authority of the Roman Pontiff. He ordained at Rome in the month of December twenty - five Priests, two Deacons, and twenty - one Bishops for diverse Sees.\par
\switchcolumn*

\LA
\begin{Resp}\Rsign Stola jucunditátis índuit eum Dóminus: \Star{} Et corónam pulchritúdinis pósuit super caput ejus.\end{Resp}
\begin{Resp}\Vsign Cibávit illum Dóminus pane vitæ et intelléctus: et aqua sapiéntiæ salutáris potávit illum.\end{Resp}
\begin{Resp}\Rsign Et corónam pulchritúdinis pósuit super caput ejus.\end{Resp}
\begin{Resp}\Vsign Glória Patri, et Fílio, \Star{} et Spirítui Sancto.\end{Resp}
\begin{Resp}\Rsign Et corónam pulchritúdinis pósuit super caput ejus.\end{Resp}
\switchcolumn
\EN
\begin{Resp}\Rsign The Lord hath put on him a robe of honour, \Star{} And put about his head a crown of joy.\end{Resp}
\begin{Resp}\Vsign With the bread of life and understanding hath the Lord fed him, and given him the water of health and wisdom to drink.\end{Resp}
\begin{Resp}\Rsign And put about his head a crown of joy.\end{Resp}
\begin{Resp}\Vsign Glory be to the Father, and to the Son, \Star{} and to the Holy Ghost.\end{Resp}
\begin{Resp}\Rsign And put about his head a crown of joy.\end{Resp}
\switchcolumn*

\end{paracol}

\par\needspace{10\baselineskip}\addvspace{1.2em}{\centering\small\textsc{Die 17 Januarii} · \textsc{17 January}\par}
\Office{S. Antonii Abbatis}{S. Anthony, Abbot}{Duplex}
\addcontentsline{toc}{subsection}{Die 17 Januarii · S. Antonii Abbatis\enspace\textit{S. Anthony, Abbot}}
\begin{paracol}{2}
\LA
\RefLine{Lectiones I–III: de Scriptura occurrente.}
\switchcolumn
\EN
\RefLine{Lessons I–III: from the Scripture of the day.}
\switchcolumn*

\LA
\RefLine{Lectiones VII–IX: ex Commúni Abbatis (C5b).}
\switchcolumn
\EN
\RefLine{Lessons VII–IX: from the Common of an Abbot (C5b).}
\switchcolumn*

\LA
\LessonHead{Lectio IV}
\switchcolumn
\EN
\LessonHead{Lesson IV}
\switchcolumn*

\LA
\noindent Antonius Ægyptius, nobílibus et christiánis paréntibus natus, quibus adoléscens orbátus est, cum ingréssus ecclésiam ex Evangelio audivísset: Si vis perfectus esse, vade, et vende ómnia quæ habes, et da paupéribus; tamquam ea sibi dicta essent, sic Christo Dómino obtemperándum existimávit. Itaque, véndita re familiari, pecúniam omnem paupéribus distríbuit. Quibus solutus impedimentis, cæléstis vitæ genus in terris cólere instituit. Sed cum in periculósum illud certamen descenderet, ad fidei præsidium, quo erat armatus, adhibéndum sibi putávit subsidium reliquárum virtútum; quarum tanto studio incénsus fuit, ut, quemcúmque vidéret áliqua virtútis laude excellentem, illum imitari studéret.\par
\switchcolumn
\EN
\noindent Anthony was an Egyptian, the child of noble and Christian parents, whom he lost while yet very young. On one occasion he entered a Church, and heard these words of the Gospel, (Matth. xix. 21,) If thou wilt be perfect, go and sell that thou hast, and give to the poor. He took these words as if they were addressed to himself personally, for this was the obedience which he thought every word of the Lord Christ should meet with. He therefore sold his whole possessions, and gave the price to the poor. Being thus delivered from worldly entanglements, he set himself to lead on earth the life of an angel. Finding himself, as it were, about to enter the field of battle against Satan, he thought it wisest to add to the shield of faith, which he already possessed, all the rest of the armour of God, wherefore he observed all those who were eminent for any grace, and strove to copy them.\par
\switchcolumn*

\LA
\begin{Resp}\Rsign Honéstum fecit illum Dóminus, et custodívit eum ab inimícis, et a seductóribus tutávit illum: \Star{} Et dedit illi claritátem ætérnam.\end{Resp}
\begin{Resp}\Vsign Justum dedúxit Dóminus per vias rectas, et osténdit illi regnum Dei.\end{Resp}
\begin{Resp}\Rsign Et dedit illi claritátem ætérnam.\end{Resp}
\switchcolumn
\EN
\begin{Resp}\Rsign The Lord made him honourable, and defended him from his enemies, and kept him safe from those that lay in wait for him; \Star{} And gave him perpetual glory.\end{Resp}
\begin{Resp}\Vsign He went down with him into the pit, and left him not in bonds.\end{Resp}
\begin{Resp}\Rsign And gave him perpetual glory.\end{Resp}
\switchcolumn*

\LA
\LessonHead{Lectio V}
\switchcolumn
\EN
\LessonHead{Lesson V}
\switchcolumn*

\LA
\noindent Nihil ígitur eo continentius, nihil vigilantius erat. Patiéntia, mansuetúdine, misericórdia, humilitate, labóre, ac studio divinárum Scripturárum superábat omnes. Ab hæreticórum et schismaticórum hóminum, maxime Arianórum, congressu et colloquio sic abhorrebat, ut ne prope quidem ad eos accedéndum díceret. Humi jacebat, cum eum necessarius somnus occupasset. Jejunium autem adeo coluit, ut salem tantummodo ad panem adhibéret, sitim aqua exstingueret; neque se ante solis occásum cibo aut potu recreabat, sæpe étiam biduum cibo abstinebat; sæpíssime in oratióne pernoctábat. Cum talis tantusque Dei miles evasísset Antonius, sanctíssimum juvenem hostis humani generis variis tentatiónibus aggréditur, quas ille jejunio et oratióne vincebat. Nec vero frequens de sátana triumphus securum reddebat Antónium, qui diaboli innumerábiles artes nocéndi nóverat.\par
\switchcolumn
\EN
\noindent He was excelled by none in watchfulness and self-restraint. He surpassed all in long-suffering, meekness, tenderness, lowliness, perseverance, and continual study of the Holy Scriptures. He had such a loathing of the company and conversation of heretics and schismatics, especially Arians, that he used to say that a faithful Christian ought as far as possible never to come near any such. He took the sleep which was needful for the body lying on the ground. Such was his devotion to fasting, that he took nothing with his bread but salt, and drank only water; he never ate or drank before sunset; he often abstained from food altogether for two days at a time; and very often passed whole nights in prayer. Being so valiant a soldier of God, Anthony was attacked by the devil with diverse temptations, but he overcame them all by prayer and fasting. Nevertheless, these frequent triumphs over Satan did not lull Anthony into security, for he was well aware of the numberless arts of assault possessed by the evil one.\par
\switchcolumn*

\LA
\begin{Resp}\Rsign Amávit eum Dóminus, et ornávit eum: stolam glóriæ índuit eum, \Star{} Et ad portas paradísi coronávit eum.\end{Resp}
\begin{Resp}\Vsign Índuit eum Dóminus lorícam fídei, et ornávit eum.\end{Resp}
\begin{Resp}\Rsign Et ad portas paradísi coronávit eum.\end{Resp}
\switchcolumn
\EN
\begin{Resp}\Rsign The Lord loved him and beautified him; He clothed him with a robe of glory, \Star{} And crowned him at the gates of Paradise.\end{Resp}
\begin{Resp}\Vsign The Lord hath put on him the breast-plate of faith, and hath adorned him.\end{Resp}
\begin{Resp}\Rsign And crowned him at the gates of Paradise.\end{Resp}
\switchcolumn*

\LA
\LessonHead{Lectio VI}
\switchcolumn
\EN
\LessonHead{Lesson VI}
\switchcolumn*

\LA
\noindent Itaque cóntulit se in vastíssimam Ægypti solitúdinem: ubi quotídie ad christianam perfectiónem proficiens, dæmones (quorum tanto erant acriores ímpetus, quanto Antonius ad resisténdum fortior evadebat) ita contempsit, ut illis exprobraret imbecillitátem; ac sæpe discipulos suos éxcitans ad pugnándum contra diabolum, docensque quibus armis vincerétur: Mihi credite, dicebat, fratres, pertimescit sátanas piórum vigílias, oratiónes, jejunia, voluntáriam paupertátem, misericórdiam et humilitátem, maxime vero ardentem amórem in Christum Dóminum, cujus único sanctíssimæ crucis signo debilitátus áufugit. Sic autem dæmónibus erat formidolósus, ut multi per Ægyptum ab illis agitati, invocáto nómine Antonii liberaréntur: tántaque erat ejus fama sanctitátis, ut per litteras se ejus oratiónibus Constantinus Magnus et fílii commendarent. Qui aliquándo quintum et centésimum annum agens, cum innumerábiles sui institúti imitatores haberet, convocátis mónachis et ad perféctam christianæ vitæ regulam instructis, sanctitáte et miraculis clarus migrávit in cælum, décimo sexto Kalendas Februarii.\par
\switchcolumn
\EN
\noindent Then he betook himself into the vast deserts of Africa that lie near Egypt. Day by day he advanced on the path to perfection. Day by day the attacks of the fiends became more violent, but day by day his' strength grew greater to strive against them. At length he came to mock at the powerlessness of the devils, against whom he stirred up his disciples to fight, teaching them with what arms to combat. Believe me, my brethren, he used to say, Satan is afraid of good men's watchings, and prayers, and fasts, and voluntary poverty, and mercifulness, and lowliness, but above all, of their warm love for Christ our Lord, the mere sign of Whose most holy Cross is enough to undo him and put him to flight. He became such an object of dread to the devils, that many persons throughout Egypt who were tormented by them, were delivered by calling on his name moreover the fame of his holiness was so spread abroad, that Constantine the Great and his sons wrote to him to commend themselves to his prayers. In the hundred and fifth year of his age, and the fulness of his reputation for piety and miracles, having roused up great numbers to follow his example, he gathered his monks around him, and when he had exhorted them to strive after Christian perfection, he passed to heaven on the 17th day of January, (in the year of our Lord 356.)\par
\switchcolumn*

\LA
\begin{Resp}\Rsign Iste homo perfécit ómnia quæ locútus est ei Deus, et dixit ad eum: Ingrédere in réquiem meam: \Star{} Quia te vidi justum coram me ex ómnibus géntibus.\end{Resp}
\begin{Resp}\Vsign Iste est, qui contémpsit vitam mundi, et pervénit ad cæléstia regna.\end{Resp}
\begin{Resp}\Rsign Quia te vidi justum coram me ex ómnibus géntibus.\end{Resp}
\begin{Resp}\Vsign Glória Patri, et Fílio, \Star{} et Spirítui Sancto.\end{Resp}
\begin{Resp}\Rsign Quia te vidi justum coram me ex ómnibus géntibus.\end{Resp}
\switchcolumn
\EN
\begin{Resp}\Rsign This is he which did according unto all that God commanded him; and God said unto him: Enter thou into My rest, \Star{} For thee have I seen righteous before Me among all people.\end{Resp}
\begin{Resp}\Vsign This is he which loved not his life in this world, and is come unto an everlasting kingdom.\end{Resp}
\begin{Resp}\Rsign For thee have I seen righteous before Me among all people.\end{Resp}
\begin{Resp}\Vsign Glory be to the Father, and to the Son, \Star{} and to the Holy Ghost.\end{Resp}
\begin{Resp}\Rsign For thee have I seen righteous before Me among all people.\end{Resp}
\switchcolumn*

\end{paracol}

\par\needspace{10\baselineskip}\addvspace{1.2em}{\centering\small\textsc{Die 18 Januarii} · \textsc{18 January}\par}
\Office{Cathedræ S. Petri Romæ}{Chair of St. Peter at Rome}{Duplex majus}
\addcontentsline{toc}{subsection}{Die 18 Januarii · Cathedræ S. Petri Romæ\enspace\textit{Chair of St. Peter at Rome}}
\begin{paracol}{2}
\LA
\RefLine{Lectio IX: de S. Priscæ Virginis et Martyris (01-18r).}
\switchcolumn
\EN
\RefLine{Lesson IX: of St. Prisca, Virgin and Martyr (01-18r).}
\switchcolumn*

\LA
\LessonHead{Lectio I}
\switchcolumn
\EN
\LessonHead{Lesson I}
\switchcolumn*

\LA
\LessonTitle{Incipit Epístola prima beáti Petri Apóstoli}
\switchcolumn
\EN
\LessonTitle{Lesson from the first letter of St. Peter the Apostle}
\switchcolumn*

\LA
\Cite{1 Pet 1:1-5}
\switchcolumn
\EN
\Cite{1 Pet 1:1-5}
\switchcolumn*

\LA
\noindent Petrus Apóstolus Jesu Christi, eléctis ádvenis dispersiónis Ponti, Galátiæ, Cappadóciæ, Asiæ, et Bithýniæ, secúndum præsciéntiam Dei Patris, in sanctificatiónem Spíritus, in obediéntiam, et aspersiónem sánguinis Jesu Christi: Grátia vobis, et pax multiplicétur. Benedíctus Deus et Pater Dómini nostri Jesu Christi, qui secúndum misericórdiam suam magnam regenerávit nos in spem vivam, per resurrectiónem Jesu Christi ex mórtuis, in hereditátem incorruptíbilem, et incontaminátam, et immarcescíbilem, conservátam in cælis in vobis, qui in virtúte Dei custodímini per fidem in salútem, parátam revelári in témpore novíssimo.\par
\switchcolumn
\EN
\noindent Peter, an apostle of Jesus Christ, to the strangers dispersed through Pontus, Galatia, Cappadocia, Asia, and Bithynia, elect, According to the foreknowledge of God the Father, unto the sanctification of the Spirit, unto obedience and sprinkling of the blood of Jesus Christ: Grace unto you and peace be multiplied. Blessed be the God and Father of our Lord Jesus Christ, who according to his great mercy hath regenerated us unto a lively hope, by the resurrection of Jesus Christ from the dead, Unto an inheritance incorruptible, and undefiled, and that can not fade, reserved in heaven for you, Who, by the power of God, are kept by faith unto salvation, ready to be revealed in the last time.\par
\switchcolumn*

\LA
\begin{Resp}\Rsign Simon Petre, ántequam de navi vocárem te, novi te et super plebem meam príncipem te constítui: \Star{} Et claves regni cælórum trádidi tibi.\end{Resp}
\begin{Resp}\Vsign Quodcúmque ligáveris super terram erit ligátum et in cælis; et quodcúmque sólveris super terram, erit solútum et in cælis.\end{Resp}
\begin{Resp}\Rsign Et claves regni cælórum trádidi tibi.\end{Resp}
\switchcolumn
\EN
\begin{Resp}\Rsign Simon Peter, before I called thee out of the ship, I knew thee, and appointed thee for a ruler over My people. \Star{} And I have given unto thee the keys of the kingdom of heaven.\end{Resp}
\begin{Resp}\Vsign Whatsoever, thou shalt bind on earth, shall be bound in heaven; and whatsoever thou shalt loose on earth, shall be loosed in heaven.\end{Resp}
\begin{Resp}\Rsign And I have given unto thee the keys of the kingdom of heaven.\end{Resp}
\switchcolumn*

\LA
\LessonHead{Lectio II}
\switchcolumn
\EN
\LessonHead{Lesson II}
\switchcolumn*

\LA
\Cite{1 Pet 1:6-9}
\switchcolumn
\EN
\Cite{1 Pet 1:6-9}
\switchcolumn*

\LA
\noindent In quo exsultábitis, módicum nunc si opórtet contristári in váriis tentatiónibus: ut probátio vestræ fídei multo pretiósior auro (quod per ignem probátur) inveniátur in laudem, et glóriam, et honórem, in revelatióne Jesu Christi: quem cum non vidéritis, dilígitis: in quem nunc quoque non vidéntes créditis: credéntes autem exsultábitis lætítia inenarrábili, et glorificáta: reportántes finem fídei vestræ, salútem animárum.\par
\switchcolumn
\EN
\noindent Wherein you shall greatly rejoice, if now you must be for a little time made sorrowful in diverse temptations: That the trial of your faith much more precious than gold which is tried by the fire may be found unto praise and glory and honour at the appearing of Jesus Christ: Whom having not seen, you love: in whom also now, though you see him not, you believe: and believing shall rejoice with joy unspeakable and glorified; Receiving the end of your faith, even the salvation of your souls.\par
\switchcolumn*

\LA
\begin{Resp}\Rsign Si díligis me, Simon Petre, pasce oves meas. Dómine, tu nosti quia amo te, \Star{} Et ánimam meam pono pro te.\end{Resp}
\begin{Resp}\Vsign Si oportúerit me mori tecum, non te negábo.\end{Resp}
\begin{Resp}\Rsign Et ánimam meam pono pro te.\end{Resp}
\switchcolumn
\EN
\begin{Resp}\Rsign Simon Peter, if thou lovest Me, feed My sheep. Lord, Thou knowest that I love thee; \Star{} I will lay down my life for thy sake.\end{Resp}
\begin{Resp}\Vsign If I should die with thee, I will not deny thee.\end{Resp}
\begin{Resp}\Rsign I will lay down my life for thy sake.\end{Resp}
\switchcolumn*

\LA
\LessonHead{Lectio III}
\switchcolumn
\EN
\LessonHead{Lesson III}
\switchcolumn*

\LA
\Cite{1 Pet 1:10-12}
\switchcolumn
\EN
\Cite{1 Pet 1:10-12}
\switchcolumn*

\LA
\noindent De qua salúte exquisiérunt atque scrutáti sunt prophétæ, qui de futúra in vobis grátia prophetavérunt; scrutántes in quod vel quale tempus significáret in eis Spíritus Christi: prænúntians eas quæ in Christo sunt passiónes et posterióres glórias: quibus revelátum est, quia non sibimetípsis, vobis autem ministrábant ea, quæ nunc nuntiáta sunt vobis per eos, qui evangelizavérunt vobis, Spíritu Sancto misso de cælo, in quem desíderant Angeli prospícere.\par
\switchcolumn
\EN
\noindent Of which salvation the prophets have inquired and diligently searched, who prophesied of the grace to come in you. Searching what or what manner of time the Spirit of Christ in them did signify: when it foretold those sufferings that are in Christ, and the glories that should follow: To whom it was revealed, that not to themselves, but to you they ministered those things which are now declared to you by them that have preached the gospel to you, the Holy Ghost being sent down from heaven, on whom the angels desire to look.\par
\switchcolumn*

\LA
\begin{Resp}\Rsign Tu es Petrus, et super hanc petram ædificábo Ecclésiam meam, et portæ ínferi non prævalébunt advérsus eam: \Star{} Et tibi dabo claves regni cælórum.\end{Resp}
\begin{Resp}\Vsign Quodcúmque ligáveris super terram, erit ligátum et in cælis; et quodcúmque sólveris super terram, erit solútum et in cælis.\end{Resp}
\begin{Resp}\Rsign Et tibi dabo claves regni cælórum.\end{Resp}
\begin{Resp}\Vsign Glória Patri, et Fílio, \Star{} et Spirítui Sancto.\end{Resp}
\begin{Resp}\Rsign Et tibi dabo claves regni cælórum.\end{Resp}
\switchcolumn
\EN
\begin{Resp}\Rsign Thou art Peter, and upon this rock I will build My Church, and the gates of hell shall not prevail against it. \Star{} And I will give unto thee the keys of the kingdom of heaven.\end{Resp}
\begin{Resp}\Vsign Whatsoever thou shalt bind on earth, shall be bound in heaven; and whatsoever thou shalt loose on earth, shall be loosed in heaven.\end{Resp}
\begin{Resp}\Rsign And I will give unto thee the keys of the kingdom of heaven.\end{Resp}
\begin{Resp}\Vsign Glory be to the Father, and to the Son, \Star{} and to the Holy Ghost.\end{Resp}
\begin{Resp}\Rsign And I will give unto thee the keys of the kingdom of heaven.\end{Resp}
\switchcolumn*

\LA
\LessonHead{Lectio IV}
\switchcolumn
\EN
\LessonHead{Lesson IV}
\switchcolumn*

\LA
\LessonTitle{Sermo sancti Leónis Papæ}
\switchcolumn
\EN
\LessonTitle{From the Sermons of Pope St. Leo (the Great.)}
\switchcolumn*

\LA
\Source{Sermo 1 de Ss. Apost. Petro et Paulo, ante medium.}
\switchcolumn
\EN
\Source{First Sermon on the Holy Apostles Peter and Paul, before the middle.}
\switchcolumn*

\LA
\noindent Cum duódecim Apóstoli, accépta per Spíritum Sanctum ómnium locutióne linguárum, imbuéndum Evangélio mundum, distribútis sibi terrárum pártibus, suscepíssent, beatíssimus Petrus, princeps apostólici órdinis, ad arcem Románi destinátur impérii, ut lux veritátis, quæ in ómnium géntium revelabátur salútem, efficácius se ab ipso cápite per totum mundi corpus effúnderet. Cujus autem natiónis hómines in hac tunc urbe non essent? aut quæ usquam gentes ignorárent quod Roma didicísset?\par
\switchcolumn
\EN
\noindent When the twelve holy Apostles had received from the Holy Ghost the power to speak all languages, they divided the whole world into districts, which they severally allotted to themselves as fields for their Gospel labours. Then was Peter, the Prince of the Apostles, sent to the capital city of the Roman Empire, that he might cause the light to shine thence throughout the whole body of the civilized nations. At that time what nation was there that had no representative in Rome? When Rome had learnt, what people that did not learn too?\par
\switchcolumn*

\LA
\begin{Resp}\Rsign Tu es pastor óvium, princeps Apostolórum: tibi trádidit Deus ómnia regna mundi: \Star{} Et ídeo tibi tráditæ sunt claves regni cælórum.\end{Resp}
\begin{Resp}\Vsign Quodcúmque ligáveris super terram, erit ligátum et in cælis; et quodcúmque sólveris super terram, erit solútum et in cælis.\end{Resp}
\begin{Resp}\Rsign Et ídeo tibi tráditæ sunt claves regni cælórum.\end{Resp}
\switchcolumn
\EN
\begin{Resp}\Rsign Thou art the Shepherd of the sheep, and the Prince of the Apostles, and unto thee hath God given all the kingdoms of the world. \Star{} Therefore unto thee hath He given the keys of the kingdom of heaven.\end{Resp}
\begin{Resp}\Vsign Whatsoever thou shalt bind on earth shall be bound in heaven; and whatsoever thou shalt loose on earth shall be loosed in heaven.\end{Resp}
\begin{Resp}\Rsign Therefore unto thee hath He given the keys of the kingdom of heaven.\end{Resp}
\switchcolumn*

\LA
\LessonHead{Lectio V}
\switchcolumn
\EN
\LessonHead{Lesson V}
\switchcolumn*

\LA
\noindent Hic conculcándæ philosophíæ opiniónes, hic dissolvéndæ erant terrénæ sapiéntiæ vanitátes, hic confutándi dǽmonum cultus, hic ómnium sacrilegiórum impíetas destruénda, ubi diligentíssima superstitióne habebátur colléctum quidquid usquam fúerat vanis erróribus institútum. Ad hanc ergo urbem tu, beatíssime Petre Apóstole, veníre non métuis, et, consórte glóriæ tuæ Paulo Apóstolo, aliárum adhuc ecclesiárum ordinatiónibus occupáto, silvam istam freméntium bestiárum, et turbullentíssimæ profunditátis océanum, constántior quam cum supra mare graderéris, ingréderis.\par
\switchcolumn
\EN
\noindent In Rome were the dreams of an unbelieving philosophy to be destroyed, in Rome were the empty utterances of earthly wisdom to be confuted, in Rome was idolatry to be overcome, in Rome profanity to be put down, even in Rome, where the activity of superstition had gathered together from the whole earth every error which it could find. O most blessed Apostle Peter! this was the city to which thou didst not shrink to come. The Apostle Paul, thy comrade in glory, was yet occupied in founding the Churches, and thou didst enter alone into that forest of wild beasts roaring furiously; thou didst commit thyself to that stormy ocean, more boldly than when thou walkest upon the waters to come to Jesus.\par
\switchcolumn*

\LA
\begin{Resp}\Rsign Ego pro te rogávi, Petre, ut non defíciat fides tua: \Star{} Et tu aliquándo convérsus confírma fratres tuos.\end{Resp}
\begin{Resp}\Vsign Caro et sanguis non revelávit tibi, sed Pater meus, qui est in cælis.\end{Resp}
\begin{Resp}\Rsign Et tu aliquándo convérsus confírma fratres tuos.\end{Resp}
\switchcolumn
\EN
\begin{Resp}\Rsign Peter, I have prayed for thee, that thy faith fail not \Star{} And when thou art converted, strengthen thy brethren.\end{Resp}
\begin{Resp}\Vsign Flesh and blood hath not revealed it unto thee, but My Father Which is in heaven.\end{Resp}
\begin{Resp}\Rsign And when thou art converted, strengthen thy brethren.\end{Resp}
\switchcolumn*

\LA
\LessonHead{Lectio VI}
\switchcolumn
\EN
\LessonHead{Lesson VI}
\switchcolumn*

\LA
\noindent Jam pópulos, qui ex circumcisióne credíderant, erudíeras: jam Antiochénam ecclésiam, ubi primum christiáni nóminis dignitas est orta, fundáveras: jam Pontum, Galátiam, Cappadóciam, Asiam atque Bithýniam légibus evangélicæ prædicatiónis impléveras; nec aut dúbius de provéctu óperis, aut de spátio tuæ ignárus ætátis, trophǽum crucis Christi Románis árcibus inferébas, quo te divínis præordinatiónibus anteíbant, et honor potestátis, et glória passiónis.\par
\switchcolumn
\EN
\noindent Thou hadst already taught them of the circumcision who were converted; thou hadst founded the Church of Antioch, the first that bore the noble name of Christian; thou hadst published the law of the Gospel throughout Pontus, Galatia, Cappadocia, Asia, and Bithynia; and thou didst not fear for the hardness of thy work, nor turn back because of thine old age, but didst boldly set up the trophy of the cross of Christ upon those Roman walls, where the Providence of God had appointed the throne of thine honour, and the glorious -scene of thy passion.\par
\switchcolumn*

\LA
\begin{Resp}\Rsign Petre, amas me? Tu scis, Dómine, quia amo te. \Star{} Pasce oves meas.\end{Resp}
\begin{Resp}\Vsign Simon Joánnis, díligis me plus his? Tu scis, Dómine, quia amo te.\end{Resp}
\begin{Resp}\Rsign Pasce oves meas.\end{Resp}
\begin{Resp}\Vsign Glória Patri, et Fílio, \Star{} et Spirítui Sancto.\end{Resp}
\begin{Resp}\Rsign Pasce oves meas.\end{Resp}
\switchcolumn
\EN
\begin{Resp}\Rsign Peter, lovest thou Me? Lord, Thou knowest that I love thee. \Star{} Feed My sheep.\end{Resp}
\begin{Resp}\Vsign Simon, son of Jonas, lovest thou Me more than these? Lord, Thou knowest that I love thee.\end{Resp}
\begin{Resp}\Rsign Feed My sheep.\end{Resp}
\begin{Resp}\Vsign Glory be to the Father, and to the Son, \Star{} and to the Holy Ghost.\end{Resp}
\begin{Resp}\Rsign Feed My sheep.\end{Resp}
\switchcolumn*

\LA
\LessonHead{Lectio VII}
\switchcolumn
\EN
\LessonHead{Lesson VII}
\switchcolumn*

\LA
\LessonTitle{Léctio sancti Evangélii secúndum Matthǽum}
\switchcolumn
\EN
\LessonTitle{From the Holy Gospel according to Matthew}
\switchcolumn*

\LA
\Cite{Matt 16:13-19}
\switchcolumn
\EN
\Cite{Matt 16:13-19}
\switchcolumn*

\LA
\noindent In illo témpore: Venit Jesus in partes Cæsaréæ Philíppi, et interrogábat discipulos suos, dicens: Quem dicunt hómines esse Fílium hóminis? Et réliqua.\par
\switchcolumn
\EN
\noindent At that time, Jesus came into the quarters of Cesarea Philippi: and he asked his disciples, saying: Whom do men say that the Son of man is? And so on.\par
\switchcolumn*

\LA
\LessonTitle{Homilía sancti Hilarii Epíscopi}
\switchcolumn
\EN
\LessonTitle{Homily by St. Hilary, Bishop (of Poitiers.)}
\switchcolumn*

\LA
\Source{Comment. in Matth. can. 16, post initium}
\switchcolumn
\EN
\Source{Commentary on Matthew xvi.}
\switchcolumn*

\LA
\noindent Dóminus a discípulis requírit, quem se hómines esse dícerent; ut adjécit, hóminis fílium. Hæc enim confessiónis tenénda rátio est, ut sicut Dei Fílium, ita et fílium hóminis meminérimus: quia álterum sine áltero nihil spei tríbuit ad salútem. Editis ítaque, quæ diversæ de eo erant, hóminum opiniónibus, quid de se ipse séntiant quærit. Petrus respóndit: Tu es Christus Fílius Dei vivi. Sed Petrus conditiónes propositiónis expénderat. Dóminus enim díxerat: Quem me hómines esse dicunt, fílium hóminis? Et certe fílium hóminis contemplátio córporis præferébat. Sed addéndo, Quem me esse dicunt, significávit, præter id quod in se videbátur, esse áliud sentiéndum: erat enim hóminis fílius. Quod ígitur de se opinándi judícium desiderábat? Non illud arbitrámur, quod de se ipse conféssus est: sed occúltum erat de quo quærebátur, in quod se credéntium fides debébat exténdere.\par
\switchcolumn
\EN
\noindent The Lord asketh His disciples who men say that He is, and He addeth, He, the Son of Man. Let us ever remember to hold fast this truth of our profession, namely, that the Son of God is the Son of Man also. Were He one and not the other, then were He no Saviour for us. The Lord then, having heard the various opinions of men, asketh, But Who say ye that I am? And Simon Peter answered and said Thou art the Christ, the Son of the living God. Peter had weighed the questions. The Lord had asked, Who do men say that I, the Son of Man, am? That He was Son of Man was sufficiently evident to all who looked upon His Body. But when He spake of His whole Self, and asked, Who do ye say that I am? He showed that the mind had something to grasp beyond That Which was seen, for Son of Man He was manifestly. What judgment did He wish them to give? I think it was not that which He had owned concerning Himself. That something more, which He wished them to own, was a hidden thing, whereunto the faith of them that believed in Him was to reach.\par
\switchcolumn*

\LA
\begin{Resp}\Rsign Quem dicunt hómines esse Filium hóminis? dixit Jesus discípulis suis. Respóndens Petrus dixit: Tu es Christus Fílius Dei vivi. \Star{} Et ego dico tibi, quia tu es Petrus, et super hanc petram ædificábo Ecclésiam meam.\end{Resp}
\begin{Resp}\Vsign Beátus es, Simon Bar-Jona, quia caro et sanguis non revelávit tibi, sed Pater meus qui est in cælis.\end{Resp}
\begin{Resp}\Rsign Et ego dico tibi, quia tu es Petrus, et super hanc petram ædificábo Ecclésiam meam.\end{Resp}
\switchcolumn
\EN
\begin{Resp}\Rsign Jesus asked His disciples, saying: Who do men say that I, the Son of Man, am? Peter answered, and said: Thou art the Christ, the Son of the living God. \Star{} And I say unto thee, that thou art Peter, and upon this rock I will build My Church.\end{Resp}
\begin{Resp}\Vsign Blessed art thou, Simon Bar-Jona, for flesh and blood hath not revealed it unto thee, but My Father Which is in heaven.\end{Resp}
\begin{Resp}\Rsign And I say unto thee, that thou art Peter, and upon this rock I will build My Church.\end{Resp}
\switchcolumn*

\LA
\LessonHead{Lectio VIII}
\switchcolumn
\EN
\LessonHead{Lesson VIII}
\switchcolumn*

\LA
\noindent Et dignum plane conféssio Petri præmium consecúta est, quia Dei Fílium in hómine vidísset. Beátus hic est, qui ultra humánum óculos intendísse et vidísse laudátus est: non id quod ex carne et sánguine erat cóntuens, sed Dei Fílium cæléstis Patris revelatióne conspíciens; dignúsque judicátus, qui quod in Christo Dei esset, primus agnósceret. O in nuncupatióne novi nóminis felix Ecclésiæ fundaméntum, dígnaque ædificatióne illíus petra, quæ inférnas leges, et tártari portas, et ómnia mortis claustra dissólveret! O beátus cæli jánitor, cujus arbítrio claves ætérni áditus tradúntur, cujus terréstre judícium præjudicáta auctóritas sit in cælo!\par
\switchcolumn
\EN
\noindent Peters confession was followed by a proper reward for having seen the Son of God in the Son of Man. Blessed is this holy Apostle, in whose praise it is said that he saw with more than human eyes That Which was unseen, who gazed upon Flesh and Blood, and by the secret revelation of the Heavenly Father recognised the Eternal Son of God; who was the first thought worthy to acknowledge the Divinity of Christ. \{Here, if necessary, the Lesson is divided.) God bless thee, O Peter, thou who by uttering for the first time the title of Divine honour, didst lay the goodly foundation of the Church! God bless thee, thou worthy rock whereon she is built, for ever triumphant over the infernal powers, the gates of hell, and the bands of death! God bless thee, happy door-keeper of heaven, to whose keeping are given the keys of the everlasting mansions, whose sentences on earth are already confirmed in heaven so that whatsoever thou shalt bind on earth shall be bound in heaven, and whatsoever thou shalt loose on earth shall be loosed in heaven.\par
\switchcolumn*

\LA
\begin{Resp}\Rsign Elégit te Dóminus sacerdótem sibi, ad sacrificándum ei \Star{} Hóstiam laudis.\end{Resp}
\begin{Resp}\Vsign Immola Deo sacrifícium laudis, et redde Altíssimo vota tua.\end{Resp}
\begin{Resp}\Rsign Hóstiam laudis.\end{Resp}
\begin{Resp}\Vsign Glória Patri, et Fílio, \Star{} et Spirítui Sancto.\end{Resp}
\begin{Resp}\Rsign Hóstiam laudis.\end{Resp}
\switchcolumn
\EN
\begin{Resp}\Rsign The Lord hath chosen thee for a priest unto Himself, to offer up unto Him \Star{} The sacrifice of praise.\end{Resp}
\begin{Resp}\Vsign Offer unto God thanksgiving, and pay thy vows unto the Most High.\end{Resp}
\begin{Resp}\Rsign The sacrifice of praise.\end{Resp}
\begin{Resp}\Vsign Glory be to the Father, and to the Son, \Star{} and to the Holy Ghost.\end{Resp}
\begin{Resp}\Rsign The sacrifice of praise.\end{Resp}
\switchcolumn*

\end{paracol}

\par\needspace{10\baselineskip}\addvspace{1.2em}{\centering\small\textsc{Die 18 Januarii} · \textsc{18 January}\par}
\Office{S. Priscæ Virginis et Martyris}{St. Prisca, Virgin and Martyr}{Simplex}
\addcontentsline{toc}{subsection}{Die 18 Januarii · S. Priscæ Virginis et Martyris\enspace\textit{St. Prisca, Virgin and Martyr}}
\begin{paracol}{2}
\LA
\LessonHead{Lectio IX (pro commemoratione)}
\switchcolumn
\EN
\LessonHead{Lesson IX (when commemorated)}
\switchcolumn*

\LA
\LessonTitle{Commemoratio: S. Priscæ Virginis et Martyris}
\switchcolumn
\EN
\LessonTitle{Commemoration of St. Prisca, Virgin and Martyr.}
\switchcolumn*

\LA
\noindent Prisca, nobilis virgo Romana, trédecim annos nata, Cláudio imperatóre, christianæ fidei accusata, ejusdem jussu ducta ad Apóllinis templum, ut idolis immolaret, cum rem detestarétur, cólaphis cæsa, in carcerem traditur: atque inde emissa, cum in fidei constántia perseveraret, affecta verbéribus, ferventíque ádipe delibuta, rursus in carcerem includitur. Post triduum in amphitheatrum producta, leoni objícitur; qui suæ feritátis oblítus, humíliter se ad ejus pedes abjecit. Quæ póstea in ergastulo triduum inedia afflicta, in equuleo suspénditur, et ungulis férreis excarnificáta in rogum injícitur, unde étiam mirabíliter evasit incólumis. Dénique extra Urbem cápite abscisso, virginitátis palmam martyrii coróna cumulávit. Cujus corpus via Ostiénsi, décimo ab Urbe milliario, a Christiánis décimo quinto Kalendas Februarii sepelítur.\par
\switchcolumn
\EN
\noindent Prisca was a noble Roman maiden, who at thirteen years of age was accused of Christianity before the Emperor Claudius. By his command she was taken to the temple of Apollo to sacrifice there, and when she refused, was buffeted and sent to prison. She was taken out from thence again, but as she still held steadfastly to the faith, they flogged her, poured boiling tallow upon her, and sent her back a second time. She was at last thrown to a lion in the amphitheatre, but it quietly lay down at her feet. She was starved for three days in a slaves' prison house, and then tortured upon the rack. Pieces of flesh were next torn from her body with iron hooks, and she was thrown on a burning pile. She marvellously still remained alive, and was accordingly beheaded outside the city. Thus she added the crown of martyrdom to the palm of virginity. The Christians buried her body at the tenth milestone on the road from Rome to Ostia on the eighteenth of January.\par
\switchcolumn*

\LA
\TeDeum{Te Deum}
\switchcolumn
\EN
\TeDeum{Te Deum}
\switchcolumn*

\end{paracol}

\par\needspace{10\baselineskip}\addvspace{1.2em}{\centering\small\textsc{Die 19 Januarii} · \textsc{19 January}\par}
\Office{Ss. Marii, Marthæ, Audifacis, et Abachum Martyrum}{Ss. Marius, Martha, Audifax, and Abachum, Martyrs}{Simplex}
\addcontentsline{toc}{subsection}{Die 19 Januarii · Ss. Marii, Marthæ, Audifacis, et Abachum Martyrum\enspace\textit{Ss. Marius, Martha, Audifax, and Abachum, Martyrs}}
\begin{paracol}{2}
\LA
\RefLine{Lectiones I–II: de Scriptura occurrente.}
\switchcolumn
\EN
\RefLine{Lessons I–II: from the Scripture of the day.}
\switchcolumn*

\LA
\LessonHead{Lectio III (pro commemoratione)}
\switchcolumn
\EN
\LessonHead{Lesson III (when commemorated)}
\switchcolumn*

\LA
\noindent Marius Persa, nobili loco natus, cum Martha cónjuge pari nobilitate, et duobus fíliis Audíface et Abachum, Romam venit Cláudio imperatóre, ut Mártyrum sepúlcra venerarétur. Ibi Christiános in víncula conjectos fovébant, et ópera ac facultátibus suis sustentábant, et Sanctórum córpora sepeliebant. Quam ob rem comprehénsi omnes, cum nec impiórum minis nec terróre commoveréntur, ut diis sacrificarent; primum fustibus debilitati, deínde funibus attracti, tum admótis candéntibus, láminis combusti, et úngulis férreis excarnificáti sunt. Postremo præcisis mánibus, et ad collum alligatis, ducti per mediam urbem, via Cornelia ad tertium décimum ab Urbe milliárium, in eum locum, qui Nymphe dicebátur, necántur: ac primum Martha, quæ virum ac fílios ad supplícia pro Jesu Christi fide constanter perferenda, veheménter fuerat cohortáta; mox ceteris in eádem arenaria cervices abscindúntur, eorúmque córpora conjiciúntur in ignem. Quæ semiusta, Felícitas matrona Romana nobilis, colligenda et in suo prædio sepelienda curávit.\par
\switchcolumn
\EN
\noindent Maris was a Persian of high rank, who came to Rome in the reign of the Emperor Claudius, with his wife Martha, who was equally noble, and their two sons Audifax and Abachum, to pray at the graves of the Martyrs. Here they comforted the Christians who were in prison, and whom they relieved by their ministrations and alms, and buried the bodies of the Saints. For these acts they were all arrested, but no threats or terrors could move them to sacrifice to idols. They were accordingly mangled with clubs, and drawn with ropes, after which they were burnt by applying plates of red-hot metal to their bodies, and their flesh partly torn off with metal hooks. Lastly their hands were all cut off, and they were fastened together by the neck, in which state they were driven through the city to the thirteenth mile-stone on the Cornelian Way, a place now called Santa Ninfa, where they were to die. Martha addressed a moving exhortation to her husband and sons to hold out bravely to the last, for the love of Jesus Christ; and was then herself drowned. The other three martyrs were next beheaded in the same sand-pit. Their bodies were thrown into a fire. The lady Felicity of Rome collected the half-burnt remains, and caused them to be buried at her own farm.\par
\switchcolumn*

\LA
\TeDeum{Te Deum}
\switchcolumn
\EN
\TeDeum{Te Deum}
\switchcolumn*

\end{paracol}

\par\needspace{10\baselineskip}\addvspace{1.2em}{\centering\small\textsc{Die 20 Januarii} · \textsc{20 January}\par}
\Office{Ss. Fabiani et Sebastiani Martyrum}{Ss. Fabian and Sebastian, Martyrs}{Duplex}
\addcontentsline{toc}{subsection}{Die 20 Januarii · Ss. Fabiani et Sebastiani Martyrum\enspace\textit{Ss. Fabian and Sebastian, Martyrs}}
\begin{paracol}{2}
\LA
\RefLine{Lectiones I–III: de Scriptura occurrente.}
\switchcolumn
\EN
\RefLine{Lessons I–III: from the Scripture of the day.}
\switchcolumn*

\LA
\LessonHead{Lectio IV}
\switchcolumn
\EN
\LessonHead{Lesson IV}
\switchcolumn*

\LA
\noindent Fabiánus Románus, a Maximíno usque ad Décium regens Ecclésiam, septem diáconis regiónes divísit, qui páuperum curam habérent. Tótidem subdiáconos creávit, qui res gestas Mártyrum a septem notáriis scriptas collígerent. Idem státuit, ut quotánnis féria quinta in Cœna Dómini, vétere combústo, chrisma renovarétur. Dénique décimo tértio Kaléndas Februárii in persecutióne Décii martýrio coronátus, in cœmetério Callísti via Appia sepelítur, cum sedísset annos quíndecim, dies quátuor. Hic fecit ordinatiónes quinque mense Decémbri, quibus creávit presbýteros vigínti duos, diáconos septem, epíscopos per divérsa loca úndecim.\par
\switchcolumn
\EN
\noindent Fabian was a Roman, and sat as Pope from (the year of our Lord 236, in) the reign of the Emperor Maximin till (250, in) that of Decius. He appointed a Deacon to each of the seven districts of Rome to look after the poor. He likewise appointed the same number of Subdeacons to collect the acts of the Martyrs from the records kept by the seven district notaries. It was by him that it was ordained that every Maundy Thursday the old Chrism should be burnt and new consecrated. He was crowned with martyrdom upon the 20th of January, in the persecution of Decius, and buried in the cemetery of St. Kallistus on the Appian Way, having sat in the throne of Peter fifteen years and four days. He held five Advent ordinations, in which he ordained twenty - two Priests, seven Deacons, and eleven Bishops for diverse Sees.\par
\switchcolumn*

\LA
\begin{Resp}\Rsign Sancti tui, Dómine, mirábile consecúti sunt iter, serviéntes præcéptis tuis, ut inveniréntur illǽsi in aquis válidis: \Star{} Terra appáruit árida: et in Mari Rubro via sine impediménto.\end{Resp}
\begin{Resp}\Vsign Quóniam percússit petram, et fluxérunt aquæ, et torréntes inundavérunt.\end{Resp}
\begin{Resp}\Rsign Terra appáruit árida: et in Mari Rubro via sine impediménto.\end{Resp}
\switchcolumn
\EN
\begin{Resp}\Rsign thy Saints, O Lord, have passed a wonderful way, serving thy commandments, that they might be found without hurt in the midst of the mighty waters. \Star{} Dry land appeared, and, out of the Red Sea, a way without impediment.\end{Resp}
\begin{Resp}\Vsign He smote the rock, and the waters gushed out, and the streams overflowed.\end{Resp}
\begin{Resp}\Rsign Dry land appeared, and, out of the Red Sea, a way without impediment.\end{Resp}
\switchcolumn*

\LA
\LessonHead{Lectio V}
\switchcolumn
\EN
\LessonHead{Lesson V}
\switchcolumn*

\LA
\noindent Sebastiánus ex patre Narbonénsi, matre Mediolanénsi natus, ob géneris nobilitátem et virtútem Diocletiáno carus fuit. Dux primæ cohórtis, Christiános, quorum fidem clam colébat, ópera et facultátibus adjuvábat; et qui ex eis tormentórum vim reformidáre videbántur, cohortatióne sic confirmábat, ut pro Jesu Christo multi se ultro tortóribus offérrent. In illis fuére Marcus et Marcelliánus fratres, qui Romæ in custódia erant apud Nicóstratum: cujus uxor Zoë vocem, quam amíserat, Sebastiáni oratióne recuperávit. Quibus Diocletiáno delátis, Sebastiánum accérsit, et veheméntius objurgátum ómnibus artifíciis a Christi fide conátur avértere. Sed cum nihil nec pollicéndo, nec terréndo profíceret, ad palum alligátum sagíttis confígi jubet.\par
\switchcolumn
\EN
\noindent The father of Sebastian was of Narbonne, and his mother a Milanese. He was a great favourite of the Emperor Diocletian, both on account of his noble birth and his personal bravery, and was by him appointed captain of the first company of the Praetorian Guards. He was in secret a Christian, and often supported the others both by good offices and money. When some showed signs of yielding under persecution, he so successfully exhorted them, that, for Jesus Christ's sake, many offered themselves to the tormentors. Among these were the brothers Mark and Marcellian who were imprisoned at Rome in the house of Nicostratus. The wife of Nicostratus himself, named Zoe, had lost her voice, but it was restored to her at the prayer of Sebastian. These facts becoming known to Diocletian, he sent for Sebastian, and after violently rebuking him, used every means to turn him from his faith in Christ. But as neither promises nor threats availed, he ordered him to be tied to a post and shot to death with arrows.\par
\switchcolumn*

\LA
\begin{Resp}\Rsign Vérbera carníficum non timuérunt Sancti Dei, moriéntes pro Christi nómine: \Star{} Ut herédes fíerent in domo Dómini.\end{Resp}
\begin{Resp}\Vsign Tradidérunt córpora sua propter Deum ad supplícia.\end{Resp}
\begin{Resp}\Rsign Ut herédes fíerent in domo Dómini.\end{Resp}
\switchcolumn
\EN
\begin{Resp}\Rsign The Saints of God shrank not from the stripes of the executioners, but died for Christ's Name's sake \Star{} That they might be made joint-heirs in the house of the Lord.\end{Resp}
\begin{Resp}\Vsign They gave their bodies for God's sake to death.\end{Resp}
\begin{Resp}\Rsign That they might be made joint-heirs in the house of the Lord.\end{Resp}
\switchcolumn*

\LA
\LessonHead{Lectio VI}
\switchcolumn
\EN
\LessonHead{Lesson VI}
\switchcolumn*

\LA
\noindent Quem ómnium opinióne mórtuum noctu sancta múlier Iréne sepeliéndi grátia jussit auférri; sed vivum repértum domi suæ curávit. Itaque paulo post confirmáta valetúdine, Diocletiáno óbviam factus, ejus impietátem libérius accusávit. Cujus aspéctu cum ille primum obstupuísset, quod mórtuum créderet; rei novitáte et acri Sebastiáni reprehensióne excandéscens, eum támdiu virgis cædi imperávit, donec ánimam Deo rédderet. Ejus corpus in cloácam dejéctum Lucína, a Sebastiáno in somnis admónita ubi esset et quo loco humári vellet, ad Catacúmbas sepelívit, ubi sancti Sebastiáni nómine célebris ecclésia est ædificáta.\par
\switchcolumn
\EN
\noindent Sebastian was treated accordingly, and left for dead, but in the night the holy widow Irene sent for the body in order to bury it, and then found that he was still alive, and nursed him in her own house. As soon as his health was restored, he went out to meet Diocletian, and boldly rebuked him for his wickedness. The Emperor was first thunderstruck at the sight of a man whom he believed to have been some time dead, but afterwards, frenzied with rage at the reproaches of Sebastian, ordered him to be beaten to death with rods, under which torment the martyr yielded his blessed soul to God, (upon the 20th day of January, in the year of our Lord 288.) His body was thrown into a sewer, but he appeared in sleep to Lucina, and made known to her where it was, and where he would have it buried. She accordingly found it and laid it in those Catacombs, over which a famous Church hath since been built, called St. Sebastian's-without-the-Walls.\par
\switchcolumn*

\LA
\begin{Resp}\Rsign Tamquam aurum in fornáce probávit eléctos Dóminus, et quasi holocáusti hóstiam accépit illos: et in témpore erit respéctus illórum: \Star{} Quóniam donum et pax est eléctis Dei.\end{Resp}
\begin{Resp}\Vsign Qui confídunt in illum, intélligent veritátem, et fidéles in dilectióne acquiéscent illi.\end{Resp}
\begin{Resp}\Rsign Quóniam donum et pax est eléctis Dei.\end{Resp}
\begin{Resp}\Vsign Glória Patri, et Fílio, \Star{} et Spirítui Sancto.\end{Resp}
\begin{Resp}\Rsign Quóniam donum et pax est eléctis Dei.\end{Resp}
\switchcolumn
\EN
\begin{Resp}\Rsign As gold in the furnace hath the Lord tried His chosen ones, and received them as a burnt-offering, and yet a while, and they shall be regarded; \Star{} For the grace of God, and His peace, are with His chosen.\end{Resp}
\begin{Resp}\Vsign They that put their trust in Him shall understand the truth and such as be faithful in love shall abide with Him.\end{Resp}
\begin{Resp}\Rsign For the grace of God, and His peace, are with His chosen.\end{Resp}
\begin{Resp}\Vsign Glory be to the Father, and to the Son, \Star{} and to the Holy Ghost.\end{Resp}
\begin{Resp}\Rsign For the grace of God, and His peace, are with His chosen.\end{Resp}
\switchcolumn*

\LA
\LessonHead{Lectio VII}
\switchcolumn
\EN
\LessonHead{Lesson VII}
\switchcolumn*

\LA
\LessonTitle{Léctio sancti Evangélii secúndum Lucam}
\switchcolumn
\EN
\LessonTitle{From the Holy Gospel according to Luke}
\switchcolumn*

\LA
\Cite{Luc 6:17-23}
\switchcolumn
\EN
\Cite{Luke 6:17-23}
\switchcolumn*

\LA
\noindent In illo témpore: Descéndens Jesus de monte, stetit in loco campéstri, et turba discipulórum ejus, et multitúdo copiósa plebis ab omni Judǽa, et Jerúsalem, et marítima, et Tyri, et Sidónis. Et réliqua.\par
\switchcolumn
\EN
\noindent At that time, Jesus came down from the mountain, and stood in the plain, and the company of His disciples, and a great multitude of people out of all Judea, and Jerusalem, and from the sea coast of Tyre and Sidon. And so on.\par
\switchcolumn*

\LA
\LessonTitle{Homilía sancti Ambrósii Epíscopi}
\switchcolumn
\EN
\LessonTitle{Homily by St Ambrose, Bishop (of Milan.)}
\switchcolumn*

\LA
\Source{Lib. 5. in Lucam. Cap. 6., post init.}
\switchcolumn
\EN
\Source{Bk. v. on Luke vi.}
\switchcolumn*

\LA
\noindent Advérte ómnia diligénter, quómodo et cum Apóstolis ascéndat et descéndat ad turbas. Quómodo enim turba nisi in húmili Christum vidéret? Non séquitur ad excélsa, non ascéndit ad sublímia. Dénique ubi descéndit, invénit infírmos: in excélsis enim infírmi esse non possunt. Hinc étiam Matthǽus docet in inferióribus débiles esse sanátos. Prius enim unusquísque sanándus est, ut paulátim virtútibus procedéntibus ascéndere possit ad montem; et ídeo quemque in inferióribus sanat, hoc est, a libídine révocat, injúriam cæcitátis avértit. Ad vúlnera nostra descéndit: ut usu quodam et cópia suæ natúræ, compartícipes nos fáciat esse regni cæléstis.\par
\switchcolumn
\EN
\noindent Mark well how Jesus goeth upward with His disciples, and downward to the multitude. How should the multitude behold Christ, save in a lower place? Such go not up to the things which are above; such attain not to the things which are high. And when Jesus cometh down, He findeth such as are diseased for such like go not up to the heights. Hence also Matthew saith that there were there all sick people, (iv. 23.) Of these every man had need of healing, that, when he had received strength, by and by, he might go up into the mountain. And therefore, being Himself come down, He healeth them in the plain, that is to say, He calleth them away from their lust, and freeth them of their blindness. He cometh down to our wounds, to the end that by a certain use of His nature, and by the abundance thereof, He might make us joint-heirs of the kingdom of heaven.\par
\switchcolumn*

\LA
\begin{Resp}\Rsign Propter testaméntum Dómini, et leges patérnas, Sancti Dei perstitérunt in amóre fraternitátis: \Star{} Quia unus fuit semper spíritus in eis, et una fides.\end{Resp}
\begin{Resp}\Vsign Ecce quam bonum et quam jucúndum habitáre fratres in unum.\end{Resp}
\begin{Resp}\Rsign Quia unus fuit semper spíritus in eis, et una fides.\end{Resp}
\switchcolumn
\EN
\begin{Resp}\Rsign Because of the covenant of the Lord, and the laws of their fathers, the Saints of God abode in brotherly love, \Star{} For one spirit and one faith was ever in them.\end{Resp}
\begin{Resp}\Vsign Behold how good and how pleasant it is for brethren to dwell together in unity.\end{Resp}
\begin{Resp}\Rsign For one spirit and one faith was ever in them.\end{Resp}
\switchcolumn*

\LA
\LessonHead{Lectio VIII}
\switchcolumn
\EN
\LessonHead{Lesson VIII}
\switchcolumn*

\LA
\noindent Beáti páuperes, quia vestrum est regnum Dei. Quátuor tantum beatitúdines sanctus Lucas Domínicas posuit, octo vero sanctus Matthǽus: sed in illis octo istæ quátuor sunt, et in quátuor istis illæ octo. Hic enim quátuor velut virtútes ampléxus est cardináles: ille in illis octo mýsticum númerum reserávit. Pro octáva enim multi inscribúntur Psalmi: et mandátum áccipis octo illis partem dare fortásse benedictiónibus. Sicut enim spei nostræ octáva perféctio est, ita octáva summa virtútum est.\par
\switchcolumn
\EN
\noindent Blessed be ye poor, for your's is the kingdom of God. Saint Luke giveth us but four of the Lord's Beatitudes, and Saint Matthew eight but in those eight are contained these four, and in these four those eight. For in these four are embraced the cardinal virtues and in those eight they are set forth in a number full of mystery. It is written at the head of more than one of the Psalms that they are for the octave, and thou hast received the commandment Give a portion to seven, and also to eight to seven or eight what? Perchance degrees of blessedness. For as this eighth Beatitude doth name the most glorious realization of our hope the kingdom of Heaven so doth it also name the most royal exertion of our strength blessed are they which are persecuted.\par
\switchcolumn*

\LA
\begin{Resp}\Rsign Sancti mei, qui in carne pósiti, certámen habuístis: \Star{} Mercédem labóris ego reddam vobis.\end{Resp}
\begin{Resp}\Vsign Veníte benedícti Patris mei, percípite regnum.\end{Resp}
\begin{Resp}\Rsign Mercédem labóris ego reddam vobis.\end{Resp}
\begin{Resp}\Vsign Glória Patri, et Fílio, \Star{} et Spirítui Sancto.\end{Resp}
\begin{Resp}\Rsign Mercédem labóris ego reddam vobis.\end{Resp}
\switchcolumn
\EN
\begin{Resp}\Rsign O ye My Saints, who, being in the flesh, didst have striving \Star{} I will render unto you a reward of your labours.\end{Resp}
\begin{Resp}\Vsign Come, ye blessed of My Father, inherit the kingdom\end{Resp}
\begin{Resp}\Rsign I will render unto you a reward of your labours.\end{Resp}
\begin{Resp}\Vsign Glory be to the Father, and to the Son, \Star{} and to the Holy Ghost.\end{Resp}
\begin{Resp}\Rsign I will render unto you a reward of your labours.\end{Resp}
\switchcolumn*

\LA
\LessonHead{Lectio IX}
\switchcolumn
\EN
\LessonHead{Lesson IX}
\switchcolumn*

\LA
\noindent Sed prius quæ sunt amplióra videámus. Beáti, inquit, páuperes, quóniam vestrum est regnum Dei. Primam benedictiónem hanc utérque Evangelísta pósuit. Órdine enim prima est, et parens quædam generatióque virtútum: quia qui contempsérit sæculária, ipse merébitur sempitérna: nec potest quisquam méritum regni cæléstis adipísci, qui mundi cupiditáte pressus, emergéndi non habet facultátem.\par
\switchcolumn
\EN
\noindent But let us first consider the fuller of the forms of these Beatitudes. Blessed be ye poor, for your's is the kingdom of God. Both of the Evangelists give to this Beatitude the first place. Yea, surely, for poorness, at least in spirit, is the first in order, the mother, and procreatrix of virtues; since he that setteth no store by temporal things, winneth toward eternal things; neither is any man able to gain the kingdom of heaven, on whom the love of this present world doth so press, that he cannot rid himself thereof.\par
\switchcolumn*

\LA
\TeDeum{Te Deum}
\switchcolumn
\EN
\TeDeum{Te Deum}
\switchcolumn*

\end{paracol}

\par\needspace{10\baselineskip}\addvspace{1.2em}{\centering\small\textsc{Die 21 Januarii} · \textsc{21 January}\par}
\Office{S. Agnetis Virginis et Martyris}{S. Agnes, Virgin and Martyr}{Duplex}
\addcontentsline{toc}{subsection}{Die 21 Januarii · S. Agnetis Virginis et Martyris\enspace\textit{S. Agnes, Virgin and Martyr}}
\begin{paracol}{2}
\LA
\RefLine{Lectiones I–III: ex Commúni Unius Virginis et Martyris (C6), in 2 loco.}
\switchcolumn
\EN
\RefLine{Lessons I–III: from the Common of a Virgin Martyr (C6), set 2.}
\switchcolumn*

\LA
\RefLine{Lectiones VII–IX: ex Commúni Unius Virginis et Martyris (C6).}
\switchcolumn
\EN
\RefLine{Lessons VII–IX: from the Common of a Virgin Martyr (C6).}
\switchcolumn*

\LA
\LessonHead{Lectio IV}
\switchcolumn
\EN
\LessonHead{Lesson IV}
\switchcolumn*

\LA
\LessonTitle{Ex libro S. Ambrósii Epíscopi de Virgínibus}
\switchcolumn
\EN
\LessonTitle{From the Book of St. Ambrose, Bishop (of Milan,) on Virgins}
\switchcolumn*

\LA
\Source{Liber 1, post initium}
\switchcolumn
\EN
\Source{Bk. i, Chap. 2}
\switchcolumn*

\LA
\noindent Hódie natális est Vírginis, integritátem sequámur. Natális est Mártyris, hóstias immolémus. Natális est sanctæ Agnétis, miréntur viri, non despérent párvuli, stúpeant nuptæ, imiténtur innúptæ. Sed quid dignum ea loqui póssumus, cujus ne nomen quidem vácuum laudis est? Devótio supra ætátem, virtus supra natúram: ut mihi videátur non hóminis habuísse nomen, sed oráculum Mártyris, quod indicávit, quid esset futúra. Nomen Vírginis títulus est pudóris. Appellábo Mártyrem: prædicávi satis. Prolíxa laudátio est, quæ non quǽritur, sed tenétur. Nemo est laudabílior, quam qui ab ómnibus laudári potest. Quot hómines, tot præcónes, qui Mártyrem prædícant, dum loquúntur.\par
\switchcolumn
\EN
\noindent This is a virgin's birthday; let us then follow the example of her chastity. It is a Martyr's birthday; let us then offer sacrifices. It is the birthday of the holy Agnes; let men then be filled with wonder, little ones with hope, married women with awe, and the unmarried with emulation. But how shall I set forth the glory of her whose very name is an utterance of praise? It seemeth to me that this being, holy beyond her years, and strong beyond human nature, received the name of Agnes, not as an earthly designation, but as a revelation from God of what she was to be. For this name Agnes is from the Greek, and being interpreted, signifieth Pure. So that this saintly maiden is known by the very title of Chastity and when I have added thereto the word Martyr, I have said enough. She needeth not the praise which we could utter, but do not. None is more praiseworthy than she for whose praise all mouths are fitted. As many as name her, so many praise her, by the noble title of martyr.\par
\switchcolumn*

\LA
\begin{Resp}\Rsign Induit me Dóminus vestiménto salútis, et induménto lætítiæ circúmdedit me: \Star{} Et tamquam sponsam decorávit me coróna.\end{Resp}
\begin{Resp}\Vsign Tradidit áuribus meis inæstimábiles margarítas, circúmdedit me vernántibus atque coruscántibus gemmis.\end{Resp}
\begin{Resp}\Rsign Et tamquam sponsam decorávit me coróna.\end{Resp}
\switchcolumn
\EN
\begin{Resp}\Rsign The Lord hath clothed me with the garments of salvation, and covered me with the robe of righteousness: \Star{} And hath set a crown upon my head, as the crown of his bride.\end{Resp}
\begin{Resp}\Vsign He hath put pearls beyond price in mine ears, and hath crowned me with bright blossoms of the eternal springtime.\end{Resp}
\begin{Resp}\Rsign And hath set a crown upon my head, as the crown of his bride.\end{Resp}
\switchcolumn*

\LA
\LessonHead{Lectio V}
\switchcolumn
\EN
\LessonHead{Lesson V}
\switchcolumn*

\LA
\noindent Hæc trédecim annórum martýrium fecísse tráditur. Quo detestabílior crudélitas quæ nec minúsculæ pepércit ætáti; immo magna vis fídei, quæ étiam ab illa testimónium invénit ætáte. Fuítne in illo corpúsculo vúlneri locus? Et quæ non hábuit quo ferrum recíperet, hábuit quo ferrum vínceret. Hæc inter cruéntas carníficum impávida manus, hæc stridéntium grávibus immóbilis tráctibus catenárum, nunc furéntis mucróni mílitis totum offérre corpus, mori adhuc néscia, sed paráta, vel si ad aras invíta raperétur, téndere Christo inter ignes manus, atque in ipsis sacrílegis focis trophǽum Dómini signáre victóris: nunc ferrátis colla manúsque ambas insérere néxibus. Sed nullus tam tenúia membra póterat nexus inclúdere. Novum martýrii genus! Nondum idónea pœnæ, et jam matúra victóriæ; certáre diffícilis, fácilis coronári: magistérium virtútis implévit, quæ præjudícium vehébat ætátis.\par
\switchcolumn
\EN
\noindent We learn by tradition that this holy martyr testified in the thirteenth year of her age. We will pass by the foul cruelty which did not spare her tender years, to contemplate the great power of her faith, whereby she overcame the weakness of childhood, and witnessed a good confession. Her little body was hardly big enough to give play to the instruments of their cruelty, but if they could scarce sheathe their swords in her slight frame, they found in her that which laughed the power of the sword to scorn. She had no fear when she found herself grasped by the bloody hands of the executioners. She was unmoved when they dragged her with clanging chains. Hardly entered on life, she stood fully prepared to die. She quailed not when the weapons of the angry soldiery were pointed at her breast. If they forced her against her will to approach the altars of devils, she could stretch forth her hands to Christ amid the very flames which consumed the idolatrous offerings, and mark on the heathen shrine the victorious Cross of the Lord. She was ready to submit her neck and hands to the iron shackles, but they were too big to clasp her slender limbs. Behold a strange martyr! She is not grown of stature to fight the battle, but she is ripe for the triumph; too weak to run in the race, and yet clearly entitled to the prize; unable from her age to be aught but a learner, she is found a teacher.\par
\switchcolumn*

\LA
\begin{Resp}\Rsign Mel et lac ex ejus ore suscépi, \Star{} Et sanguis ejus ornávit genas meas.\end{Resp}
\begin{Resp}\Vsign Osténdit mihi thesáuros incomparábiles, quos mihi se datúrum repromísit.\end{Resp}
\begin{Resp}\Rsign Et sanguis ejus ornávit genas meas.\end{Resp}
\switchcolumn
\EN
\begin{Resp}\Rsign Honey and milk have I received from his mouth, \Star{} And his Blood is red upon my cheeks.\end{Resp}
\begin{Resp}\Vsign He hath shewn me countless treasures, and hath sworn to endow me therewith.\end{Resp}
\begin{Resp}\Rsign And his Blood is red upon my cheeks.\end{Resp}
\switchcolumn*

\LA
\LessonHead{Lectio VI}
\switchcolumn
\EN
\LessonHead{Lesson VI}
\switchcolumn*

\LA
\noindent Non sic ad thálamum nupta properáret, ut ad supplícii locum, læta succéssu, gradu festína Virgo procéssit. Flere omnes, ipsa sine fletu. Mirári pleríque, quod tam fácile vitæ suæ pródiga, quam nondum háuserat, jam quasi perfúncta donáret. Stupére univérsi, quod jam divinitátis testis exsísteret, quæ adhuc árbitra sui per ætátem esse non posset. Quanto terróre egit cárnifex ut timerétur, quantis blandítiis ut suadéret, quantórum vota, ut sibi ad núptias proveníret! At illa: Et hæc Sponsi injúria est, inquit, exspectáre placitúram. Qui me sibi prior elégit, accípiet: quid, percússor, moráris? Péreat corpus, quod amári potest óculis, quibus nolo. Stetit, orávit, cervícem infléxit. Cérneres trepidáre carníficem, quasi ipse addíctus fuísset, trémere percussóris déxteram, pallére ora aliéno timéntis perículo, cum puélla non timéret suo. Habétis ígitur in una hóstia duplex martýrium, pudóris et religiónis. Et virgo permánsit, et martýrium obtínuit.\par
\switchcolumn
\EN
\noindent She went to the place of execution a virgin, with more willing and joyful footsteps than she would have gone with to the nuptial chamber as a bride. The spectators were all in tears, and she alone did not weep. They beheld her with wonder, laying down that life of which she had hardly begun to taste the sweets, as freely as though she had drained it to the dregs and was weary of its burden. All men were amazed when they saw her whose years had not made her her own mistress, arise as a witness for the Deity. Consider how many threats her murderer used to excite her fears, how many arguments to shake her resolution, how many promises to bribe her to accept his offers of marriage. But she answered him It is an insult to Him Whom I have wedded to expect me to comply. He That first chose me, His will I be. Headsman, why waitest thou? Perish the body which draweth the admiration of eyes from which I would turn away. She stood, prayed, and then bent her neck for the stroke. Now mightest thou have seen the murderer trembling as though he himself were the criminal, the executioner's hand shake, and the faces of them that stood by turn white at the sight of her position, and all the while herself remain without fear. This one victim brought God a double offering, that of her purity, and that of her faith. She preserved virginity and achieved martyrdom.\par
\switchcolumn*

\LA
\begin{Resp}\Rsign Jam corpus ejus córpori meo sociátum est, et sanguis ejus ornávit genas meas: \Star{} Cujus mater virgo est, cujus pater féminam nescit.\end{Resp}
\begin{Resp}\Vsign Ipsi sum desponsáta, cui Angeli serviunt, cujus pulchritúdinem sol et luna mirántur.\end{Resp}
\begin{Resp}\Rsign Cujus mater virgo est, cujus pater féminam nescit.\end{Resp}
\begin{Resp}\Vsign Glória Patri, et Fílio, \Star{} et Spirítui Sancto.\end{Resp}
\begin{Resp}\Rsign Cujus mater virgo est, cujus pater féminam nescit.\end{Resp}
\switchcolumn
\EN
\begin{Resp}\Rsign His holy Body is already united to me, and his Blood is red upon my cheeks: \Star{} He it is whose Mother is a virgin, and whose Father hath not known a woman.\end{Resp}
\begin{Resp}\Vsign I am wedded to the Lord of Angels, at whose beauty the sun and the moon do wonder.\end{Resp}
\begin{Resp}\Rsign He it is whose Mother is a virgin, and whose Father hath not known a woman.\end{Resp}
\begin{Resp}\Vsign Glory be to the Father, and to the Son, \Star{} and to the Holy Ghost.\end{Resp}
\begin{Resp}\Rsign He it is whose Mother is a virgin, and whose Father hath not known a woman.\end{Resp}
\switchcolumn*

\end{paracol}

\par\needspace{10\baselineskip}\addvspace{1.2em}{\centering\small\textsc{Die 22 Januarii} · \textsc{22 January}\par}
\Office{Ss. Vincentii et Anastasii Martyrum}{Ss. Vincent and Anastasius, Martyrs}{Semiduplex}
\addcontentsline{toc}{subsection}{Die 22 Januarii · Ss. Vincentii et Anastasii Martyrum\enspace\textit{Ss. Vincent and Anastasius, Martyrs}}
\begin{paracol}{2}
\LA
\RefLine{Lectiones I–III: de Scriptura occurrente.}
\switchcolumn
\EN
\RefLine{Lessons I–III: from the Scripture of the day.}
\switchcolumn*

\LA
\RefLine{Lectiones VII–IX: ex Commúni Plurimorum Martyrum Pontificum (C3).}
\switchcolumn
\EN
\RefLine{Lessons VII–IX: from the Commune Plurimorum Martyrum Pontificum (C3).}
\switchcolumn*

\LA
\LessonHead{Lectio IV}
\switchcolumn
\EN
\LessonHead{Lesson IV}
\switchcolumn*

\LA
\noindent Vincentius, Oscæ in Hispánia citeriore natus, a prima ætate stúdiis déditus, sacras litteras a Valerio Cæsaraugustano episcopo didicit; cujus étiam partes suscépit prædicándi Evangélium, quod epíscopus, propter linguæ impediméntum, prædicatiónis officio fungi non poterat. Ea re ad Dacianum, provinciæ a Diocletiáno et Maximiáno præpositum, delata, Vincentius Cæsaraugustæ comprehénditur, et vinctus ad Dacianum Valéntiam adducitur. Ubi verbéribus et equuleo tortus, multis præséntibus, cum nulla aut tormentórum vi, aut acerbitate vel lenitate verbórum a proposito deterreri posset; in cratícula impositus, prunis ardéntibus suppositis, ac férreis unguibus excarnificatus, candentibúsque láminis exustus, íterum ducitur in carcerem stratum testáceis fragmentis, ut ejus nudum corpus, somno oppressum, a subjectis étiam testárum aculeis torquerétur.\par
\switchcolumn
\EN
\noindent Vincent was born at Huesca in Aragon in Spain. He was early turned to study, and learned sacred letters from Valerius, Bishop of Saragossa. He was accustomed to deliver discourses for this Prelate, who, owing to an impediment in his speech, was not able to preach himself. This coming to the ears of Dacian, Prefect of the province under Diocletian and Maximian, he caused Vincent to be arrested at Saragossa, and brought before him at Valencia in bonds. The saint was scourged, and afterwards tormented on the rack, in presence of numerous spectators, but neither torture, threats, nor fair words could bend his resolution. He was then laid on a grating over hot coals, his flesh mangled with iron hooks, and white-hot plates of metal applied to the wounds. The still breathing remains were taken back to a prison, and laid on broken potsherds, that the agony of his naked body might prevent his sleeping from exhaustion.\par
\switchcolumn*

\LA
\begin{Resp}\Rsign Sancti tui, Dómine, mirábile consecúti sunt iter, serviéntes præcéptis tuis, ut inveniréntur illǽsi in aquis válidis: \Star{} Terra appáruit árida: et in Mari Rubro via sine impediménto.\end{Resp}
\begin{Resp}\Vsign Quóniam percússit petram, et fluxérunt aquæ, et torréntes inundavérunt.\end{Resp}
\begin{Resp}\Rsign Terra appáruit árida: et in Mari Rubro via sine impediménto.\end{Resp}
\switchcolumn
\EN
\begin{Resp}\Rsign thy Saints, O Lord, have passed a wonderful way, serving thy commandments, that they might be found without hurt in the midst of the mighty waters. \Star{} Dry land appeared, and, out of the Red Sea, a way without impediment.\end{Resp}
\begin{Resp}\Vsign He smote the rock, and the waters gushed out, and the streams overflowed.\end{Resp}
\begin{Resp}\Rsign Dry land appeared, and, out of the Red Sea, a way without impediment.\end{Resp}
\switchcolumn*

\LA
\LessonHead{Lectio V}
\switchcolumn
\EN
\LessonHead{Lesson V}
\switchcolumn*

\LA
\noindent Verum, illo in tenebricosa incluso custódia, claríssimus splendor obortus totum carcerem illustrávit; quæ lux cum summa admiratióne omnes, qui aderant, affecísset, res a custode carceris ad Dacianum defertur. Qui eductum in molli cúlcitra collocat; et quem cruciátibus in suam senténtiam trahere non poterat, deliciis perducere conátur. Sed invictus Vincentii animus, Jesu Christi fide speque munítus, vicit ómnia; et ignis, ferri, tortórum immanitate superata, victor ad cælestem martyrii corónam advolávit, undecimo Kalendas Februarii. Cujus corpus, cum projectum esset inhumátum, corvus et a volucribus et a lupo únguibus, rostro, alis mirabíliter defendit. Qua re cógnita, Dacianus illud in altum mare demergi jubet: sed inde étiam divinitus ejectum ad littus, Christiáni sepeliunt.\par
\switchcolumn
\EN
\noindent As he lay in his dark cell, a glorious light suddenly filled the prison, to the astonishment of all who saw it. The gaoler informed Dacian, who caused the martyr to be brought out and cared for in a soft bed, hoping that though he had failed to move him by cruelty, he might seduce him by pretended kindness. But the indomitable soul of Vincent, armed with faith and hope in Christ Jesus, remained unconquered even to the end, and triumphing over the fire, the steel, and the cruelty of the tormentors, passed away to receive the victorious crown of martyrdom in heaven, on the 22nd day of January, (in the year of our Lord 304.) His body was thrown out unburied. A raven perched upon it and kept off with his beak, claws, and wings both the other birds and a wolf, which came to prey on it. Dacian then had it thrown into the sea, but by the will of God it was washed up again, and the Christians took and buried it.\par
\switchcolumn*

\LA
\begin{Resp}\Rsign Vérbera carníficum non timuérunt Sancti Dei, moriéntes pro Christi nómine: \Star{} Ut herédes fíerent in domo Dómini.\end{Resp}
\begin{Resp}\Vsign Tradidérunt córpora sua propter Deum ad supplícia.\end{Resp}
\begin{Resp}\Rsign Ut herédes fíerent in domo Dómini.\end{Resp}
\switchcolumn
\EN
\begin{Resp}\Rsign The Saints of God shrank not from the stripes of the executioners, but died for Christ's Name's sake \Star{} That they might be made joint-heirs in the house of the Lord.\end{Resp}
\begin{Resp}\Vsign They gave their bodies for God's sake to death.\end{Resp}
\begin{Resp}\Rsign That they might be made joint-heirs in the house of the Lord.\end{Resp}
\switchcolumn*

\LA
\LessonHead{Lectio VI}
\switchcolumn
\EN
\LessonHead{Lesson VI}
\switchcolumn*

\LA
\noindent Anastasius Persa, monachus, Heraclío imperatóre, cum sanctam Jerosolymórum terram visitasset, ad Cæsaréam Palæstinæ pro Christi religióne víncula et verbera constanter perpéssus est. Mox a Persis ob eamdem causam variis cruciátibus afféctus, a rege Chósroa una cum septuagínta aliis Christiánis secúri percútitur. Cujus relíquiæ primum Jerosolymam ad monastérium, in quo monasticam vitam proféssus erat, deínde Romam delátæ, collocátæ sunt in monasterio ad Aquas Sálvias.\par
\switchcolumn
\EN
\noindent Anastasius was a Persian monk who made a pilgrimage to the Holy Places at Jerusalem in the reign of the Emperor Heraclius, during which journey he endured bonds and stripes on account of his confession of Christ at Banias, then called Cassarea, in Palestine. Soon after his return, he was arrested by the Persians for the same cause, and, after enduring diverse torments, he and seventy other Christians were beheaded by order of King Chosroes. He testified upon the 22nd day of January in the year of our Lord 628. His reliques were first carried to Jerusalem, to the monastery in which he had made his monastic profession, and afterwards to Rome, where they were laid in the monastery of Saints Vincent and Anastasius.\par
\switchcolumn*

\LA
\begin{Resp}\Rsign Tamquam aurum in fornáce probávit eléctos Dóminus, et quasi holocáusti hóstiam accépit illos: et in témpore erit respéctus illórum: \Star{} Quóniam donum et pax est eléctis Dei.\end{Resp}
\begin{Resp}\Vsign Qui confídunt in illum, intélligent veritátem, et fidéles in dilectióne acquiéscent illi.\end{Resp}
\begin{Resp}\Rsign Quóniam donum et pax est eléctis Dei.\end{Resp}
\begin{Resp}\Vsign Glória Patri, et Fílio, \Star{} et Spirítui Sancto.\end{Resp}
\begin{Resp}\Rsign Quóniam donum et pax est eléctis Dei.\end{Resp}
\switchcolumn
\EN
\begin{Resp}\Rsign As gold in the furnace hath the Lord tried His chosen ones, and received them as a burnt-offering, and yet a while, and they shall be regarded; \Star{} For the grace of God, and His peace, are with His chosen.\end{Resp}
\begin{Resp}\Vsign They that put their trust in Him shall understand the truth and such as be faithful in love shall abide with Him.\end{Resp}
\begin{Resp}\Rsign For the grace of God, and His peace, are with His chosen.\end{Resp}
\begin{Resp}\Vsign Glory be to the Father, and to the Son, \Star{} and to the Holy Ghost.\end{Resp}
\begin{Resp}\Rsign For the grace of God, and His peace, are with His chosen.\end{Resp}
\switchcolumn*

\end{paracol}

\par\needspace{10\baselineskip}\addvspace{1.2em}{\centering\small\textsc{Die 23 Januarii} · \textsc{23 January}\par}
\Office{S. Raymundi de Peñafort Confessoris}{S. Raymond of Penafort, Confessor}{Semiduplex}
\addcontentsline{toc}{subsection}{Die 23 Januarii · S. Raymundi de Peñafort Confessoris\enspace\textit{S. Raymond of Penafort, Confessor}}
\begin{paracol}{2}
\LA
\RefLine{Lectiones I–III: de Scriptura occurrente.}
\switchcolumn
\EN
\RefLine{Lessons I–III: from the Scripture of the day.}
\switchcolumn*

\LA
\RefLine{Lectiones VII–VIII: ex Commúni Confessoris non pontificis (C5).}
\switchcolumn
\EN
\RefLine{Lessons VII–VIII: from the Common of a Confessor not a Bishop (C5).}
\switchcolumn*

\LA
\RefLine{Lectio IX: de S. Emerentianæ Virginis et Martyris (01-23o).}
\switchcolumn
\EN
\RefLine{Lesson IX: of St. Emerentiana, Virgin and Martyr (01-23o).}
\switchcolumn*

\LA
\LessonHead{Lectio IV}
\switchcolumn
\EN
\LessonHead{Lesson IV}
\switchcolumn*

\LA
\noindent Beátus Raymundus Barcinonénsis, ex nobili família de Péñafort, christianæ religiónis rudimentis imbutus, adhuc párvulus exímia animi et corporis índole magnum aliquid portendere visus est. Nam adoléscens humaniores litteras in pátria proféssus, Bonóniam se cóntulit, ubi pietátis officiis ac pontificio civilíque juri sedulo incumbens, et doctoris laurea insignítus, ibidem sacros cánones magna cum hóminum admiratióne est interpretatus. Ejus virtútum fama percrebrescente, Berengarius Barcinonénsis epíscopus, cum Roma suam ad ecclésiam rediret, eum conveniéndi causa Bonóniam iter instituit, et tandem summis precibus, ut secum in pátriam reverterétur, obtinuit. Mox ejusdem ecclésiæ canonicatu et præpositura ornatus, universo clero et pópulo, integritate, modestia, doctrina et morum suavitáte præfulsit, ac Deíparæ Vírginis, quam singulari pietátis afféctu venerabátur, honórem et cultum semper pro viribus auxit.\par
\switchcolumn
\EN
\noindent The blessed Raymond was born at Barcelona, (in the year of our Lord 1175,) and was of the noble family of the De Penafuerte. He was early instructed in the Christian religion, and even while he was still a little child, he showed such excellence of mind and body, as filled his friends with strong hopes of his future greatness. As a young man he taught letters in his native place. He afterwards went to Bologna, where he applied himself to works of godliness, and to the study of the Ecclesiastical and Civil Law. He took the degree of Doctor, and lectured with great applause upon the Canon Law. He attained so much celebrity that Berengar, Bishop of Barcelona, on his way from Rome to his own See, turned aside to visit the Saint at Bologna, and at length persuaded him after many entreaties to return with him to Spain. He was appointed to a Canonry and the Archdeaconry in the Church of Barcelona, in which offices he set both clergy and people a brilliant example of uprightness, modesty, learning, and meekness, and more especially strove, as far as in him lay, to increase the honour and reverence paid to the Virgin Mother of God, whom he venerated with an affection singularly devoted.\par
\switchcolumn*

\LA
\begin{Resp}\Rsign Honéstum fecit illum Dóminus, et custodívit eum ab inimícis, et a seductóribus tutávit illum: \Star{} Et dedit illi claritátem ætérnam.\end{Resp}
\begin{Resp}\Vsign Justum dedúxit Dóminus per vias rectas, et osténdit illi regnum Dei.\end{Resp}
\begin{Resp}\Rsign Et dedit illi claritátem ætérnam.\end{Resp}
\switchcolumn
\EN
\begin{Resp}\Rsign The Lord made him honourable, and defended him from his enemies, and kept him safe from those that lay in wait for him; \Star{} And gave him perpetual glory.\end{Resp}
\begin{Resp}\Vsign He went down with him into the pit, and left him not in bonds.\end{Resp}
\begin{Resp}\Rsign And gave him perpetual glory.\end{Resp}
\switchcolumn*

\LA
\LessonHead{Lectio V}
\switchcolumn
\EN
\LessonHead{Lesson V}
\switchcolumn*

\LA
\noindent Annum circiter quintum supra quadragesimum agens, in ordine fratrum Prædicatórum solemni emissa professióne, ut novus miles, in omni virtútum genere, sed præcipue in caritate erga egénos, et maxime captivos ab infidelibus detentos, se exercuit. Unde cum ejus hortatu sanctus Petrus Nolascus (cujus ipse confessiónes audiebat) suas opes piíssimo huic operi conferret, tum eidem, tum beato Raymundo, et Jacobo primo Aragoniæ regi apparens beatíssima Virgo, gratíssimum sibi et unigenito Fílio suo fore dixit, si in suum honórem instituerétur ordo religiosórum, quibus captivos ex infidelium tyrannide liberándi cura incumberet. Quare collátis inter se consíliis, ordinem beátæ Maríæ de mercede redemptiónis captivórum fundavérunt: cui beátus Raymundus certas vivéndi leges præscripsit ad ejusdem ordinis vocatiónem accommodatíssimas, quarum approbatiónem aliquot post annos a Gregório nono impetrávit, et dictum sanctum Petrum primum generalem ordinis magistrum suis ipse mánibus habitu eodem indutum creávit.\par
\switchcolumn
\EN
\noindent When he was about forty-five years of age he solemnly professed in the Order of Friars Preachers, and strove, as a new recruit, to perfect himself in all the duties of his calling, particularly in charity to the poor, and above all to those unhappy Christians who were slaves to the unbelievers. He was the Confessor of St. Peter Nolasco and of James I., King of Aragon, and by his advice St Peter Nolasco gave up his whole worldly possessions to ransom as many as possible of the wretched captives. At this moment the Most Blessed Virgin appeared simultaneously to St. Raymond, St. Peter Nolasco, and King James, and revealed to them the pleasure of her Only-begotten Son and herself, that they should establish in her honour an order of Religious persons whose work should be the redemption of Christian slaves from bondage among unbelievers. The three took counsel together, and then founded the Order of the Blessed Mary of Ransom, for the Redemption of Captives. The blessed Raymond himself composed a most appropriate code of rules for the new institution, for which he after some years obtained the express sanction of Pope Gregory IX., and himself (on the ioth day of August, in the year 1223,) with his own hands clothed St. Peter Nolasco in the habit, and constituted him the first Master General of the Order.\par
\switchcolumn*

\LA
\begin{Resp}\Rsign Amávit eum Dóminus, et ornávit eum: stolam glóriæ índuit eum, \Star{} Et ad portas paradísi coronávit eum.\end{Resp}
\begin{Resp}\Vsign Índuit eum Dóminus lorícam fídei, et ornávit eum.\end{Resp}
\begin{Resp}\Rsign Et ad portas paradísi coronávit eum.\end{Resp}
\switchcolumn
\EN
\begin{Resp}\Rsign The Lord loved him and beautified him; He clothed him with a robe of glory, \Star{} And crowned him at the gates of Paradise.\end{Resp}
\begin{Resp}\Vsign The Lord hath put on him the breast-plate of faith, and hath adorned him.\end{Resp}
\begin{Resp}\Rsign And crowned him at the gates of Paradise.\end{Resp}
\switchcolumn*

\LA
\LessonHead{Lectio VI}
\switchcolumn
\EN
\LessonHead{Lesson VI}
\switchcolumn*

\LA
\noindent Ab eodem Gregório Romam accersítus, et capellani ac pœnitentiarii et confessarii sui munere decoratus, ejusdem jussu Romanórum Pontificum decreta, in diversis conciliis et epistolis sparsa, in unum Decretalium volúmen redégit. Archiepiscopátum Tarraconensem ab ipso Pontifice sibi oblátum constantíssime recusávit, et totius ordinis Prædicatórum, generale magistérium, quod per biennium sanctíssime administraverat, sponte dimísit. Jacobo Aragoniæ regi sacræ Inquisitiónis officii suis in regnis instituéndi auctor fuit. Multa patrávit miracula, inter quæ illud claríssimum, quod ex insula Baleari Majori Barcinónem reversurus, strato super aquas pallio centum sexagínta milliaria sex horis confécerit et suum cœnobium januis clausis fúerit ingréssus. Tandem prope centenarius, virtútibus et meritis cumulátus, obdormívit in Dómino, anno salútis millesimo ducentésimo septuagesimo quinto. Quem Clemens octavus in Sanctórum númerum rétulit.\par
\switchcolumn
\EN
\noindent He was summoned to Rome by Gregory IX., (in the year 1230,) and appointed by him his Chaplain, Penitentiary, and Confessor, and by his orders collected into one volume of the Decretals the ordinances of the Roman Pontiffs, which up to that time were only to be found scattered among the records of diverse Councils and Churches. He firmly refused the Archbishopric of Tarascon, which was offered him by the Pope himself, and, having been chosen Master General of the whole order of Friars Preachers, he discharged the duties of that office in holiness for two years, and then resigned it. It was by his advice that James, King of Aragon, established the Office of the Holy Inquisition in his dominions. He was distinguished by many miracles, of which the chief which is narrated of him is that on one occasion being in the island of Majorca and wishing to go to Barcelona, he spread his cloak upon the sea, and passed over the waters on it, accomplishing the whole distance of sixty leagues in six hours, and finally entering his convent through the closed doors. He attained the age of nearly an hundred years, and fell asleep in the Lord (on the 6th day of January,) in the year of salvation 1275. His name was enrolled by Clement VIII. among those of the Saints.\par
\switchcolumn*

\LA
\begin{Resp}\Rsign Iste homo perfécit ómnia quæ locútus est ei Deus, et dixit ad eum: Ingrédere in réquiem meam: \Star{} Quia te vidi justum coram me ex ómnibus géntibus.\end{Resp}
\begin{Resp}\Vsign Iste est, qui contémpsit vitam mundi, et pervénit ad cæléstia regna.\end{Resp}
\begin{Resp}\Rsign Quia te vidi justum coram me ex ómnibus géntibus.\end{Resp}
\begin{Resp}\Vsign Glória Patri, et Fílio, \Star{} et Spirítui Sancto.\end{Resp}
\begin{Resp}\Rsign Quia te vidi justum coram me ex ómnibus géntibus.\end{Resp}
\switchcolumn
\EN
\begin{Resp}\Rsign This is he which did according unto all that God commanded him; and God said unto him: Enter thou into My rest, \Star{} For thee have I seen righteous before Me among all people.\end{Resp}
\begin{Resp}\Vsign This is he which loved not his life in this world, and is come unto an everlasting kingdom.\end{Resp}
\begin{Resp}\Rsign For thee have I seen righteous before Me among all people.\end{Resp}
\begin{Resp}\Vsign Glory be to the Father, and to the Son, \Star{} and to the Holy Ghost.\end{Resp}
\begin{Resp}\Rsign For thee have I seen righteous before Me among all people.\end{Resp}
\switchcolumn*

\end{paracol}

\par\needspace{10\baselineskip}\addvspace{1.2em}{\centering\small\textsc{Die 23 Januarii} · \textsc{23 January}\par}
\Office{S. Emerentianæ Virginis et Martyris}{St. Emerentiana, Virgin and Martyr}{Simplex}
\addcontentsline{toc}{subsection}{Die 23 Januarii · S. Emerentianæ Virginis et Martyris\enspace\textit{St. Emerentiana, Virgin and Martyr}}
\begin{paracol}{2}
\LA
\LessonHead{Lectio IX (pro commemoratione)}
\switchcolumn
\EN
\LessonHead{Lesson IX (when commemorated)}
\switchcolumn*

\LA
\LessonTitle{Commemoratio: S. Emerentianæ Virginis et Martyris}
\switchcolumn
\EN
\LessonTitle{Commemoration of: St. Emerentiana, Virgin and Martyr}
\switchcolumn*

\LA
\noindent Emerentiána virgo Romana, collactanea beátæ Agnetis, adhuc catechumena, fide et caritate flagrans, furéntes in Christiános idolórum cultores cum vehementius accusaret, a concitáta multitúdine lapidibus óbruta est. Quæ in cruciátibus orans ad sepúlcrum sanctæ Agnetis, proprio sánguine, quem pro Christo constanter effudit, baptizata, ánimam Deo réddidit.\par
\switchcolumn
\EN
\noindent Emerentiana was a Roman maiden, and the foster-sister of the blessed Agnes. While she was still a Catechumen she was inspired by her faith and love to rebuke the fury of the idol-worshippers against the Christians, whereupon a mob assembled, and stoned her so severely that she was only able to drag herself to the grave of holy Agnes, where, while she prayed, she gave up her soul to God, being baptized, not in water, but in her own blood, so freely shed for Christ.\par
\switchcolumn*

\LA
\TeDeum{Te Deum}
\switchcolumn
\EN
\TeDeum{Te Deum}
\switchcolumn*

\end{paracol}

\par\needspace{10\baselineskip}\addvspace{1.2em}{\centering\small\textsc{Die 24 Januarii} · \textsc{24 January}\par}
\Office{S. Timothei Episcopi et Martyris}{St. Timothy, Bishop and Martyr}{Duplex}
\addcontentsline{toc}{subsection}{Die 24 Januarii · S. Timothei Episcopi et Martyris\enspace\textit{St. Timothy, Bishop and Martyr}}
\begin{paracol}{2}
\LA
\RefLine{Lectiones I–III: de Scriptura occurrente.}
\switchcolumn
\EN
\RefLine{Lessons I–III: from the Scripture of the day.}
\switchcolumn*

\LA
\RefLine{Lectiones VII–IX: ex Commúni Unius Martyris Pontificis (C2).}
\switchcolumn
\EN
\RefLine{Lessons VII–IX: from the Common of a Single Martyr Bishop (C2).}
\switchcolumn*

\LA
\LessonHead{Lectio IV}
\switchcolumn
\EN
\LessonHead{Lesson IV}
\switchcolumn*

\LA
\noindent Timótheus, Lystris in Lycaónia natus ex patre Gentíli et matre Judǽa, christiánam colébat religiónem, cum in ea loca venit Paulus Apóstolus. Qui fama commótus, quæ de Timóthei sanctitáte percrebúerat, ipsum adhíbuit sócium suæ peregrinatiónis: sed propter Judǽos qui se ad Christum convérterant, sciéntes Timóthei patrem esse Gentílem, eum circumcídit. Cum ígitur ambo Ephesum veníssent, ibi ordinátus est epíscopus ab Apóstolo, ut eam ecclésiam gubernáret.\par
\switchcolumn
\EN
\noindent Timothy was a native of Lystra in Lycaonia, born of a Gentile father and a Jewish mother. He embraced the Christian religion when the Apostle Paul came into those parts. The holy Apostle was so struck with the fame of Timothy's sanctity, that he chose him to be the companion of his journeys, and caused him to be circumcised, in order to remove a stumbling-block from the way of those Jews who felt drawn to Christianity. When they came together to Ephesus, the Apostle consecrated him Bishop of that Church.\par
\switchcolumn*

\LA
\begin{Resp}\Rsign Honéstum fecit illum Dóminus, et custodívit eum ab inimícis, et a seductóribus tutávit illum: \Star{} Et dedit illi claritátem ætérnam.\end{Resp}
\begin{Resp}\Vsign Descendítque cum illo in fóveam, et in vínculis non derelíquit eum.\end{Resp}
\begin{Resp}\Rsign Et dedit illi claritátem ætérnam.\end{Resp}
\switchcolumn
\EN
\begin{Resp}\Rsign The Lord made him honourable, and defended him from his enemies, and kept him safe from those that lay in wait for him, \Star{} And gave him perpetual glory.\end{Resp}
\begin{Resp}\Vsign He went down with him into the pit, and left him not in bonds.\end{Resp}
\begin{Resp}\Rsign And gave him perpetual glory.\end{Resp}
\switchcolumn*

\LA
\LessonHead{Lectio V}
\switchcolumn
\EN
\LessonHead{Lesson V}
\switchcolumn*

\LA
\noindent Ad quem Apóstolus duas epístolas scripsit, álteram Laodicéa, álteram Roma. Quibus in pastorális offícii cura confirmátus, cum sacrifícium, quod uni Deo debétur fíeri dǽmonum simulácris, ferre non posset; pópulum Ephesínum, Diánæ in ejus celebritáte immolántem, ab illa impietáte removére conátus, lapídibus óbrutus est; ac pene mórtuus, a Christiánis eréptus et in montem óppido vicínum delátus, nono Kaléndas Februárii obdormívit in Dómino.\par
\switchcolumn
\EN
\noindent Two of the Apostle Paul's Epistles are addressed to this Saint, of which one was written from Laodicea, and the other from Rome. These sacred writings so stirred him up to the zealous discharge of his duties as a spiritual shepherd, that he strove to prevent the people of Ephesus from sacrificing to Artemis on her feastday, knowing that sacrifice is due to God alone. The heathens thereupon stoned him till he was well-nigh dead, and although he was rescued by the Christians, and carried to a mountain near the city, he then fell asleep in the Lord, on the 24th day of January.\par
\switchcolumn*

\LA
\begin{Resp}\Rsign Desidérium ánimæ ejus tribuísti ei Dómine, \Star{} Et voluntáte labiórum ejus non fraudásti eum.\end{Resp}
\begin{Resp}\Vsign Quóniam prævenísti eum in benedictiónibus dulcédinis: posuísti in cápite ejus corónam de lápide pretióso.\end{Resp}
\begin{Resp}\Rsign Et voluntáte labiórum ejus non fraudásti eum.\end{Resp}
\switchcolumn
\EN
\begin{Resp}\Rsign Lord, Thou hast given him his heart's desire. \Star{} And hast not withholden the request of his lips.\end{Resp}
\begin{Resp}\Vsign For Thou hast prevented him with the blessings of sweetness: Thou hast set a crown of precious stones upon his head.\end{Resp}
\begin{Resp}\Rsign And hast not withholden the request of his lips.\end{Resp}
\switchcolumn*

\LA
\LessonHead{Lectio VI}
\switchcolumn
\EN
\LessonHead{Lesson VI}
\switchcolumn*

\LA
\LessonTitle{Sermo sancti Augustíni Epíscopi}
\switchcolumn
\EN
\LessonTitle{From the Sermons of St Augustine, Bishop (of Hippo.)}
\switchcolumn*

\LA
\Source{Sermo 44 de Sanctis}
\switchcolumn
\EN
\Source{44 on the Saints.}
\switchcolumn*

\LA
\noindent Triumphális beáti Mártyris Timóthei dies hódie nobis anniversária celebritáte recúrrit: cujus glorificatióni sicut congáudet Ecclésia, sic ejus propónit sequénda vestígia. Si enim compátimur, et conglorificábimur. In cujus glorióso agóne duo nobis præcípue consideránda sunt: induráta vidélicet tortóris sævítia, et Mártyris invícta patiéntia. Sævítia tortóris, ut eam detestémur: patiéntia Mártyris, ut eam imitémur. Audi Psalmístam advérsus malítiam increpántem: Noli æmulári in malignántibus quóniam tamquam fœnum velóciter aréscent. Quod autem advérsus malignántes patiéntia exhibénda sit, audi Apóstolum suadéntem: Patiéntia vobis necessária est, ut reportétis promissiónes.\par
\switchcolumn
\EN
\noindent The illustrious day whereon the blessed Martyr Timothy conquered, doth this day come round to us again and as the Church doth rejoice with him in his glory, so doth she set before us his footsteps to be followed. For if we suffer, we shall also reign with him. In his glorious battle we have two things chiefly to consider the hardened cruelty of the tormentor, and the unconquered patience of the Martyr the cruelty of the tormentor, that we may abhor it; the patience of the Martyr, that we may imitate it. Hear what the Psalmist saith, complaining against sin: “Fret not thyself because of the evil-doers, for they shall soon dry up like the grass.” (xxxvi. 1.) But touching the patience which is to be shown against the evil-doers, hear the word wherewith the Apostle moveth us Ye have need of patience, that ye may receive the promise. (Heb. x. 36.)\par
\switchcolumn*

\LA
\begin{Resp}\Rsign Stola jucunditátis índuit eum Dóminus: \Star{} Et corónam pulchritúdinis pósuit super caput ejus.\end{Resp}
\begin{Resp}\Vsign Cibávit illum Dóminus pane vitæ et intelléctus: et aqua sapiéntiæ salutáris potávit illum.\end{Resp}
\begin{Resp}\Rsign Et corónam pulchritúdinis pósuit super caput ejus.\end{Resp}
\begin{Resp}\Vsign Glória Patri, et Fílio, \Star{} et Spirítui Sancto.\end{Resp}
\begin{Resp}\Rsign Et corónam pulchritúdinis pósuit super caput ejus.\end{Resp}
\switchcolumn
\EN
\begin{Resp}\Rsign The Lord hath put on him a robe of honour, \Star{} And put about his head a crown of joy.\end{Resp}
\begin{Resp}\Vsign With the bread of life and understanding hath the Lord fed him, and given him the water of health and wisdom to drink.\end{Resp}
\begin{Resp}\Rsign And put about his head a crown of joy.\end{Resp}
\begin{Resp}\Vsign Glory be to the Father, and to the Son, \Star{} and to the Holy Ghost.\end{Resp}
\begin{Resp}\Rsign And put about his head a crown of joy.\end{Resp}
\switchcolumn*

\end{paracol}

\par\needspace{10\baselineskip}\addvspace{1.2em}{\centering\small\textsc{Die 25 Januarii} · \textsc{25 January}\par}
\Office{In Conversione S. Pauli Apostoli}{Conversion of St. Paul the Apostle}{Duplex majus}
\addcontentsline{toc}{subsection}{Die 25 Januarii · In Conversione S. Pauli Apostoli\enspace\textit{Conversion of St. Paul the Apostle}}
\begin{paracol}{2}
\LA
\LessonHead{Lectio I}
\switchcolumn
\EN
\LessonHead{Lesson I}
\switchcolumn*

\LA
\LessonTitle{De Actibus Apostolórum}
\switchcolumn
\EN
\LessonTitle{Lesson from the Acts of the Apostles}
\switchcolumn*

\LA
\Cite{Act 9:1-5}
\switchcolumn
\EN
\Cite{Acts 9:1-5}
\switchcolumn*

\LA
\noindent Saulus autem adhuc spirans minárum et cædis in discípulos Dómini, accéssit ad príncipem sacerdótum, et pétiit ab eo epístolas in Damáscum ad synagógas: ut si quos invenísset hujus viæ viros ac mulíeres, vinctos perdúceret in Jerúsalem. Et cum iter fáceret, cóntigit ut appropinquáret Damásco: et súbito circumfúlsit eum lux de cælo. Et cadens in terram audívit vocem dicéntem sibi: Saule, Saule, quid me perséqueris? Qui dixit: Quis es, dómine? Et ille: Ego sum Jesus, quem tu perséqueris: durum est tibi contra stímulum calcitráre.\par
\switchcolumn
\EN
\noindent And Saul, as yet breathing out threatenings and slaughter against the disciples of the Lord, went to the high priest, And asked of him letters to Damascus, to the synagogues: that if he found any men and women of this way, he might bring them bound to Jerusalem. And as he went on his journey, it came to pass that he drew nigh to Damascus; and suddenly a light from heaven shined round about him. And falling on the ground, he heard a voice saying to him: Saul, Saul, why persecutest thou me? Who said: Who art thou, Lord? And he: I am Jesus whom thou persecutest. It is hard for thee to kick against the goad.\par
\switchcolumn*

\LA
\begin{Resp}\Rsign Qui operátus est Petro in apostolátum, operátus est et mihi inter gentes: \Star{} Et cognovérunt grátiam Dei, quæ data est mihi.\end{Resp}
\begin{Resp}\Vsign Grátia Dei in me vácua non fuit, sed grátia ejus semper in me manet.\end{Resp}
\begin{Resp}\Rsign Et cognovérunt grátiam Dei, quæ data est mihi.\end{Resp}
\switchcolumn
\EN
\begin{Resp}\Rsign He That wrought effectually in Peter to the Apostleship, the Same was mighty in me toward the Gentiles, \Star{} And they perceived the grace that was given unto me of the Lord Christ.\end{Resp}
\begin{Resp}\Vsign The grace of God which was bestowed upon me was not in vain, but His grace abideth ever in me.\end{Resp}
\begin{Resp}\Rsign And they perceived the grace that was given unto me of the Lord Christ.\end{Resp}
\switchcolumn*

\LA
\LessonHead{Lectio II}
\switchcolumn
\EN
\LessonHead{Lesson II}
\switchcolumn*

\LA
\Cite{Act 9:6-9}
\switchcolumn
\EN
\Cite{Acts 9:6-9}
\switchcolumn*

\LA
\noindent Et tremens ac stupens dixit: Dómine, quid me vis fácere? Et Dóminus ad eum: Surge, et ingrédere civitátem, et ibi dicétur tibi quid te opórteat fácere. Viri autem illi qui comitabántur cum eo, stabant stupefácti, audiéntes quidem vocem, néminem autem vidéntes. Surréxit autem Saulus de terra, apertísque óculis nihil vidébat. Ad manus autem illum trahéntes, introduxérunt Damáscum. Et erat ibi tribus diébus non videns, et non manducávit, neque bibit.\par
\switchcolumn
\EN
\noindent And he trembling and astonished, said: Lord, what wilt thou have me to do? And the Lord said to him: Arise, and go into the city, and there it shall be told thee what thou must do. Now the men who went in company with him, stood amazed, hearing indeed a voice, but seeing no man. And Saul arose from the ground; and when his eyes were opened, he saw nothing. But they leading him by the hands, brought him to Damascus. And he was there three days, without sight, and he did neither eat nor drink.\par
\switchcolumn*

\LA
\begin{Resp}\Rsign Bonum certámen certávi, cursum consummávi, fidem servávi: \Star{} Ideóque repósita est mihi coróna justítiæ.\end{Resp}
\begin{Resp}\Vsign Scio cui crédidi, et certus sum quia potens est depósitum meum serváre in illum diem.\end{Resp}
\begin{Resp}\Rsign Ideóque repósita est mihi coróna justítiæ.\end{Resp}
\switchcolumn
\EN
\begin{Resp}\Rsign I have fought a good fight, I have finished my course, I have kept the faith; \Star{} Henceforth there is laid up for me a crown of righteousness.\end{Resp}
\begin{Resp}\Vsign I know Whom I have believed, and am persuaded that He is able to keep that which I have committed unto Him against that day.\end{Resp}
\begin{Resp}\Rsign Henceforth there is laid up for me a crown of righteousness.\end{Resp}
\switchcolumn*

\LA
\LessonHead{Lectio III}
\switchcolumn
\EN
\LessonHead{Lesson III}
\switchcolumn*

\LA
\Cite{Act 9:10-16}
\switchcolumn
\EN
\Cite{Acts 9:10-16}
\switchcolumn*

\LA
\noindent Erat autem quidam discípulus Damásci, nómine Ananías: et dixit ad illum in visu Dóminus: Ananía. At ille ait: Ecce ego, Dómine. Et Dóminus ad eum: Surge, et vade in vicum qui vocátur Rectus: et quære in domo Judæ Saulum nómine Tarsénsem: ecce enim orat. (Et vidit virum Ananíam nómine, introeúntem, et imponéntem sibi manus ut visum recípiat.) Respóndit autem Ananías: Dómine, audívi a multis de viro hoc, quanta mala fécerit sanctis tuis in Jerúsalem: et hic habet potestátem a princípibus sacerdótum alligándi omnes qui ínvocant nomen tuum. Dixit autem ad eum Dóminus: Vade, quóniam vas electiónis est mihi iste, ut portet nomen meum coram géntibus, et régibus, et fíliis Israël. Ego enim osténdam illi quanta opórteat eum pro nómine meo pati.\par
\switchcolumn
\EN
\noindent Now there was a certain disciple at Damascus, named Ananias. And the Lord said to him in a vision: Ananias. And he said: Behold I am here, Lord. And the Lord said to him: Arise, and go into the street that is called Straight, and seek in the house of Judas, one named Saul of Tarsus. For behold he prayeth. And he saw a man named Ananias coming in, and putting his hands upon him, that he might receive his sight. But Ananias answered: Lord, I have heard by many of this man, how much evil he hath done to thy saints in Jerusalem. And here he hath authority from the chief priests to bind all that invoke thy name. And the Lord said to him: Go thy way; for this man is to me a vessel of election, to carry my name before the Gentiles, and kings, and the children of Israel. For I will shew him how great things he must suffer for my name's sake.\par
\switchcolumn*

\LA
\begin{Resp}\Rsign Repósita est mihi coróna justítiæ, \Star{} Quam reddet mihi Dóminus in illum diem justus judex.\end{Resp}
\begin{Resp}\Vsign Scio cui crédidi, et certus sum quia potens est depósitum meum serváre in illum diem.\end{Resp}
\begin{Resp}\Rsign Quam reddet mihi Dóminus in illum diem justus judex.\end{Resp}
\begin{Resp}\Vsign Glória Patri, et Fílio, \Star{} et Spirítui Sancto.\end{Resp}
\begin{Resp}\Rsign Quam reddet mihi Dóminus in illum diem justus judex.\end{Resp}
\switchcolumn
\EN
\begin{Resp}\Rsign There is laid up for me a crown of righteousness: \Star{} Which the Lord, the righteous Judge, shall give me at that day.\end{Resp}
\begin{Resp}\Vsign I know Whom I have believed, and am persuaded that He is able to keep that which I have committed unto Him against that day.\end{Resp}
\begin{Resp}\Rsign Which the Lord, the righteous Judge, shall give me at that day.\end{Resp}
\begin{Resp}\Vsign Glory be to the Father, and to the Son, \Star{} and to the Holy Ghost.\end{Resp}
\begin{Resp}\Rsign Which the Lord, the righteous Judge, shall give me at that day.\end{Resp}
\switchcolumn*

\LA
\LessonHead{Lectio IV}
\switchcolumn
\EN
\LessonHead{Lesson IV}
\switchcolumn*

\LA
\LessonTitle{Sermo sancti Augustíni Epíscopi}
\switchcolumn
\EN
\LessonTitle{From the Sermons of St. Augustine, Bishop (of Hippo.)}
\switchcolumn*

\LA
\Source{Sermo 14 de Sanctis}
\switchcolumn
\EN
\Source{14th on the Saints}
\switchcolumn*

\LA
\noindent Hódie de Actibus Apostolórum léctio hæc pronuntiáta est, ubi Paulus Apóstolus, ex persecutóre Christianórum, annuntiátor factus est Christi. Prostrávit enim Christus persecutórem, ut fáceret Ecclésiæ doctórem: percútiens eum, et sanans; occídens, et vivíficans: occísus agnus a lupis, et fáciens agnos de lupis. Ita enim in præclára prophetía cum Jacob Patriárcha benedíceret fíliis suis (præséntes tangens, futúra prospíciens) prædíctum erat quod in Paulo cóntigit. Erat autem Paulus, sicut ipse testátur, de tribu Bénjamin. Cum autem Jacob, benedícens fílios suos, venísset ad benedicéndum Bénjamin, ait de illo: Bénjamin lupus rapax.\par
\switchcolumn
\EN
\noindent We have this day heard read out of the Acts of the Apostles how that the Apostle Paul, from being a persecutor of the Christians, was changed into a preacher of Christ. Christ laid low the persecutor, that He might raise him up a teacher of His Church. He smote and healed him, slew him and made him alive again. For the Lord Christ is that Lamb That was Himself slain by the wolves, and That now turneth the wolves into lambs. Now was fulfilled in Paul that which was clearly spoken in prophecy by the Patriarch Jacob, when he blessed his children, laying hands indeed on them which then were, but looking forward to the things which were yet for to come. Paul beareth witness of himself that he was of the tribe of Benjamin, (Rom. xi. 1,) and when Jacob blessed his sons, and came to bless Benjamin, he said: Benjamin shall ravin as a wolf. (Gen. xlix. 27.)\par
\switchcolumn*

\LA
\begin{Resp}\Rsign Tu es vas electiónis, sancte Paule Apóstole, prædicator veritátis in univérso mundo, \Star{} Per quem omnes gentes cognovérunt grátiam Dei.\end{Resp}
\begin{Resp}\Vsign Intercéde pro nobis ad Deum, qui te elégit.\end{Resp}
\begin{Resp}\Rsign Per quem omnes gentes cognovérunt grátiam Dei.\end{Resp}
\switchcolumn
\EN
\begin{Resp}\Rsign O Holy Apostle Paul, thou art a chosen vessel unto God, to preach the Gospel throughout the whole world \Star{} Through whom all nations have known the grace of God.\end{Resp}
\begin{Resp}\Vsign Pray for us to God Who hath chosen thee.\end{Resp}
\begin{Resp}\Rsign Through whom all nations have known the grace of God.\end{Resp}
\switchcolumn*

\LA
\LessonHead{Lectio V}
\switchcolumn
\EN
\LessonHead{Lesson V}
\switchcolumn*

\LA
\noindent Quid ergo? lupus rapax semper? Absit; sed qui mane rapit prædam, ad vésperam dívidit escas. Hoc in Apóstolo Paulo implétum est, quia et de illo dictum erat. Jam, si placet, audiámus illum mane rapiéntem, ad vésperam escas dividéntem. Mane et véspere pósita sunt pro eo, ac si dicerétur, prius et póstea. Sic ergo accipiámus: Prius rápiet, póstea dívidet escas. Atténdite raptórem: Saulus, inquit, accéptis epístolis a princípibus sacerdótum, ibat, ut ubicúmque inveníret Christiános, ad sacerdótes attráheret et addúceret, útique puniéndos.\par
\switchcolumn
\EN
\noindent What then? Is Benjamin a wolf that shall ravin for ever? God forbid. In the morning he shall devour the prey, and at night he shall divide the spoil. This is exactly what was fulfilled in the Apostle Paul. If it please you, we will now consider how in the morning he devoured the prey, and at night divided the spoil. Here morning and evening are put for the beginning and the end. So we may read, In the beginning he shall devour the prey, and at the end he shall divide the spoil. First, then, in the beginning, he devoured the prey. So it is written that he received letters from the chief priests and went forth, that wheresoever he should find any Christians, he might bring them bound unto the priests, that they might be punished.\par
\switchcolumn*

\LA
\begin{Resp}\Rsign Grátia Dei sum id quod sum: \Star{} Et grátia ejus in me vacua non fuit, sed semper in me manet.\end{Resp}
\begin{Resp}\Vsign Qui operátus est Petro in apostolátum, operátus est et mihi inter gentes.\end{Resp}
\begin{Resp}\Rsign Et grátia ejus in me vacua non fuit, sed semper in me manet.\end{Resp}
\switchcolumn
\EN
\begin{Resp}\Rsign By the grace of God I am what I am. \Star{} And His grace which was bestowed upon me was not in vain, but abideth ever in me.\end{Resp}
\begin{Resp}\Vsign He that wrought effectually in Peter to the Apostleship, the Same was mighty in me toward the Gentiles.\end{Resp}
\begin{Resp}\Rsign And His grace which was bestowed upon me was not in vain, but abideth ever in me.\end{Resp}
\switchcolumn*

\LA
\LessonHead{Lectio VI}
\switchcolumn
\EN
\LessonHead{Lesson VI}
\switchcolumn*

\LA
\noindent Ibat spirans et anhélans cædes: hoc est, mane rápiens. Nam et quando lapidátus est Stéphanus primus Martyr pro nómine Christi, evidéntius áderat et Saulus; et sic áderat lapidántibus, ut non ei suffíceret, si tantum suis mánibus lapidáret. Ut enim esset in ómnium lapidántium mánibus, ipse ómnium vestiménta servábat; magis sǽviens omnes adjuvándo, quam suis mánibus lapidándo. Audívimus, mane rápiet: videámus ad vésperam quáliter dívidat escas. Voce Christi prostrátus de cælo, et accípiens désuper interdíctum jam sæviéndi, cécidit in fáciem suam, prius prosternéndus, póstea erigéndus; prius percutiéndus, póstea sanándus.\par
\switchcolumn
\EN
\noindent He went breathing out threatenings and slaughter, yea, truly, devouring the prey. When also they stoned Stephen, the first Martyr that laid down his life for Christ's name's sake, Saul was consenting unto his death, and, as though it contented him not to stone him, he kept the clothes of all them that did it, urging them on more than if he had joined them. So in the morning he devoured the prey. How in the evening did he divide the spoil? Struck down by the voice of Christ from heaven, ravining no more, he falleth upon his face, cast down to be raised up, smitten to be healed.\par
\switchcolumn*

\LA
\begin{Resp}\Rsign Saulus, qui et Paulus, magnus prædicátor, \Star{} A Deo confortátus convalescébat, et confundébat Judæos.\end{Resp}
\begin{Resp}\Vsign Osténdens quia hic est Christus, Fílius Dei.\end{Resp}
\begin{Resp}\Rsign A Deo confortátus convalescébat, et confundébat Judæos.\end{Resp}
\begin{Resp}\Vsign Glória Patri, et Fílio, \Star{} et Spirítui Sancto.\end{Resp}
\begin{Resp}\Rsign A Deo confortátus convalescébat, et confundébat Judæos.\end{Resp}
\switchcolumn
\EN
\begin{Resp}\Rsign Saul, who also is called Paul, was made a great preacher \Star{} And being of God increased the more in strength, he confounded the Jews.\end{Resp}
\begin{Resp}\Vsign Proving that This is very Christ, the Son of God.\end{Resp}
\begin{Resp}\Rsign And being of God increased the more in strength, he confounded the Jews.\end{Resp}
\begin{Resp}\Vsign Glory be to the Father, and to the Son, \Star{} and to the Holy Ghost.\end{Resp}
\begin{Resp}\Rsign And being of God increased the more in strength, he confounded the Jews.\end{Resp}
\switchcolumn*

\LA
\LessonHead{Lectio VII}
\switchcolumn
\EN
\LessonHead{Lesson VII}
\switchcolumn*

\LA
\LessonTitle{Léctio sancti Evangélii secúndum Matthǽum}
\switchcolumn
\EN
\LessonTitle{From the Holy Gospel according to Matthew}
\switchcolumn*

\LA
\Cite{Matt 19:27-29}
\switchcolumn
\EN
\Cite{Matt 19:27-29}
\switchcolumn*

\LA
\noindent In illo témpore: Dixit Petrus ad Jesum: Ecce nos relíquimus ómnia, et secúti sumus te: quid ergo erit nobis? Et réliqua.\par
\switchcolumn
\EN
\noindent In that time, Peter said to Jesus said to him: Behold we have left all things, and have followed thee: what therefore shall we have? And so on.\par
\switchcolumn*

\LA
\LessonTitle{Homilía sancti Bedæ Venerábilis Presbýteri}
\switchcolumn
\EN
\LessonTitle{Sermon of the Venerable Bede, Priest (at Jarrow)}
\switchcolumn*

\LA
\Source{Homilia in Natali S. Benedicti Ep.}
\switchcolumn
\EN
\Source{Homily for St. Benedict's Birth-day}
\switchcolumn*

\LA
\noindent Perféctus ille est, qui ábiens vendit ómnia quæ habet, et dat paupéribus, ac véniens séquitur Christum; habébit enim thesáurum non deficiéntem in cælis. Unde bene, interrogánte Petro, dixit tálibus Jesus: Amen dico vobis, quod vos qui secúti estis me, in regeneratióne, cum séderit Fílius hóminis in sede majestátis suæ, sedébitis et vos super sedes duódecim, judicántes duódecim tribus Israël. In hac quippe vita pro ejus nómine laborántes, in ália prǽmium speráre dócuit, id est, in regeneratióne; cum vidélicet in vitam immortálem fuérimus resurgéndo regeneráti, qui in vitam cadúcam mortáliter erámus géniti.\par
\switchcolumn
\EN
\noindent If thou wilt be perfect, saith Christ, go and sell that thou hast, and give to the poor, and come and follow Me and thou shalt have treasure in heaven. (Matth. xix. 21.) Yea, treasure that passeth not away! Unto such saith Jesus, at the questioning of Peter Amen I say unto you, that ye which have followed Me, in the regeneration, when the Son of Man shall sit in the throne of His glory, ye also shall sit upon twelve thrones, judging the twelve tribes of Israel. He taught them, which work for His Name's sake in this life, to look for their reward in another life that is, in the regeneration. In the regeneration! when we who have been born dying creatures into a dying life, shall in the resurrection be born again into an undying life.\par
\switchcolumn*

\LA
\begin{Resp}\Rsign Sancte Paule Apóstole, prædicátor veritátis et doctor géntium, \Star{} Intercéde pro nobis ad Deum, qui te elégit, ut digni efficiámur grátia Dei.\end{Resp}
\begin{Resp}\Vsign Tu es vas electiónis, sancte Paule Apostole, prædicátor veritátis.\end{Resp}
\begin{Resp}\Rsign Intercéde pro nobis ad Deum, qui te elégit, ut digni efficiámur grátia Dei.\end{Resp}
\switchcolumn
\EN
\begin{Resp}\Rsign O Holy Apostle Paul, Preacher of the truth, and teacher of the Gentiles, \Star{} Pray for us to God, Who hath chosen thee, that we may be made worthy of the grace of God.\end{Resp}
\begin{Resp}\Vsign O Holy Apostle Paul, thou art a chosen vessel unto God, and a Preacher of the truth.\end{Resp}
\begin{Resp}\Rsign Pray for us to God Who hath chosen thee, that we may be made worthy of the grace of God.\end{Resp}
\switchcolumn*

\LA
\LessonHead{Lectio VIII}
\switchcolumn
\EN
\LessonHead{Lesson VIII}
\switchcolumn*

\LA
\noindent Et justa prorsus retribútio, ut, qui hic pro Christo humánæ glóriam celsitúdinis neglexérunt, illic a Christo júdices glorificáti singuláriter cum eo assídeant qui a sequéndis ejus vestígiis nulla ratióne póterant avélli. Nemo autem putet duódecim tantum Apóstolos, quia pro Juda prævaricánte Matthías eléctus est, tunc esse judicatúros; sicut nec duódecim solæ sunt tribus Israël judicándæ: alióquin tribus Levi, quæ tertiadécima est, injudicáta recédet.\par
\switchcolumn
\EN
\noindent And soothly, it is a just retribution, that they, who, while they were yet here, have for Christ's sake set no store by being great among men, should there by Christ be singularly glorified to be the assessors of His judgment-seat, even they whom nothing here could turn aside from being the followers of His footsteps. Nevertheless, let there be no man that believeth that the twelve Apostles only, among whom Matthias holdeth that place from which Judas by transgression fell, (Acts i. 25,) that they only shall judge, even as the twelve tribes of Israel shall not alone be judged for then were the tribe of Levi, which is the thirteenth, unjudged.\par
\switchcolumn*

\LA
\begin{Resp}\Rsign Damasci, præpósitus gentis Aretæ regis vóluit me comprehéndere: \Star{} Et a frátribus per murum demíssus sum in sporta, * Et sic evasi manus ejus in nómine Dómini.\end{Resp}
\begin{Resp}\Vsign Deus et Pater Dómini nostri Jesu Christi scit quia non mentior.\end{Resp}
\begin{Resp}\Rsign Et a frátribus per murum demissus sum in sporta.\end{Resp}
\begin{Resp}\Vsign Glória Patri, et Fílio, \Star{} et Spirítui Sancto.\end{Resp}
\begin{Resp}\Rsign Et sic evasi manus ejus in nómine Dómini.\end{Resp}
\switchcolumn
\EN
\begin{Resp}\Rsign In Damascus the governor under Aretas the king was desirous to apprehend me \Star{} And by the brethren in a basket was I let down by the wall, and so escaped I his hands, in the name of the Lord.\end{Resp}
\begin{Resp}\Vsign The God and Father of our Lord Jesus Christ knoweth that I lie not.\end{Resp}
\begin{Resp}\Rsign And by the brethren in a basket was I let down by the wall.\end{Resp}
\begin{Resp}\Vsign Glory be to the Father, and to the Son, \Star{} and to the Holy Ghost.\end{Resp}
\begin{Resp}\Rsign And by the brethren in a basket was I let down by the wall.\end{Resp}
\switchcolumn*

\LA
\LessonHead{Lectio IX}
\switchcolumn
\EN
\LessonHead{Lesson IX}
\switchcolumn*

\LA
\noindent Et Paulus, qui tertiusdécimus est Apóstolus, judicándi sorte privábitur? cum ipse dicat: Nescítis quóniam ángelos judicábimus, quanto magis sæculária? Sciéndum namque est, omnes, qui ad exémplum Apostolórum, sua reliquérunt ómnia et secúti sunt Christum, júdices cum eo ventúros, sicut étiam omne mortálium genus esse judicándum. Quia enim duodenário sæpe número solet in Scriptúris univérsitas designári, per duódecim sedes Apostolórum ómnium numerósitas judicántium, et, per duódecim tribus Israël, univérsitas eórum qui judicándi sunt, osténditur.\par
\switchcolumn
\EN
\noindent Moreover, then, were Paul, who is the thirteenth Apostle, deprived of all part in the judgment; whereas he saith of himself: Know ye not that we shall judge angels? How much more things that pertain to this life? But it behoveth us to know that every one who hath forsaken all and followed Christ, as did the Apostles, shall also come with Him to judgment, even as every man shall stand at His judgment seat. And the Scriptures use often to signify all by this number twelve; by the twelve thrones of the Apostles are signified the thrones of all them that shall judge; and by the twelve tribes of Israel, the whole number of them that shall be judged.\par
\switchcolumn*

\LA
\TeDeum{Te Deum}
\switchcolumn
\EN
\TeDeum{Te Deum}
\switchcolumn*

\end{paracol}

\par\needspace{10\baselineskip}\addvspace{1.2em}{\centering\small\textsc{Die 26 Januarii} · \textsc{26 January}\par}
\Office{S. Polycarpi Episcopi et Martyris}{St. Polycarp, Bishop and Martyr}{Duplex}
\addcontentsline{toc}{subsection}{Die 26 Januarii · S. Polycarpi Episcopi et Martyris\enspace\textit{St. Polycarp, Bishop and Martyr}}
\begin{paracol}{2}
\LA
\RefLine{Lectiones I–III: de Scriptura occurrente.}
\switchcolumn
\EN
\RefLine{Lessons I–III: from the Scripture of the day.}
\switchcolumn*

\LA
\LessonHead{Lectio IV}
\switchcolumn
\EN
\LessonHead{Lesson IV}
\switchcolumn*

\LA
\LessonTitle{Ex libro sancti Hierónymi Presbýteri de Scriptóribus ecclesiásticis}
\switchcolumn
\EN
\LessonTitle{From the Book on Ecclesiastical Writers, composed by St. Jerome, Priest (at Bethlehem.)}
\switchcolumn*

\LA
\Source{Cap. 17}
\switchcolumn
\EN
\Source{Chapter 17}
\switchcolumn*

\LA
\noindent Polycárpus, Joánnis Apóstoli discípulus et ab eo Smyrnæ epíscopus ordinátus, totíus Asiæ princeps fuit, quippe qui nonnúllos Apostolórum et eórum qui víderant Dóminum, magístros habúerit et víderit. Hic propter quasdam super die Paschæ quæstiónes, sub imperatóre Antoníno Pio, Ecclésiam in Urbe regénte Anicéto, Romam venit; ubi plúrimos credéntium, Marciónis et Valentíni persuasióne decéptos, redúxit ad fidem. Cumque ei fortúito óbviam fuísset Márcion et díceret: Cognóscis nos? respóndit: Cognósco primogénitum diáboli. Póstea vero, regnánte Marco Antoníno et Lúcio Aurélio Cómmodo, quarta post Nerónem persecutióne, Smyrnæ, sedénte procónsule et univérso pópulo in amphitheátro advérsus eum personánte, igni tráditus est. Scripsit ad Philippénses valde útilem epístolam, quæ usque hódie in Asiæ convéntu légitur.\par
\switchcolumn
\EN
\noindent Polycarp was a disciple of the Apostle John, and was consecrated by him Bishop of Smyrna. He was reckoned the chief of all the Christians of Asia, because he had been taught by several of the Apostles, and other persons who had seen the Lord. During the reign of the Emperor Antoninus Pius, and while Anicetus presided over the Church of Rome, Polycarp came thither to discuss some questions regarding the time for observing Easter. He found some heretics at Rome, who had been led astray by the doctrine of Marcion and Valentine, and brought back many of them to the faith. One day Marcion met him by accident, and said to him: Do you recognise me? whereto he replied: I recognise the devil's eldest son. Some time after, in the reign of Mark Antonine and Lucius Aurelius Commodus, during the fourth persecution since Nero, when the Pro-consul was ruling in Smyrna, the whole population being assembled in the theatre, clamoured against Polycarp, and to please them he was burnt. He wrote an extremely useful Epistle to the Philippians, which is publicly read in the Churches of Asia even to this day.\par
\switchcolumn*

\LA
\begin{Resp}\Rsign Honéstum fecit illum Dóminus, et custodívit eum ab inimícis, et a seductóribus tutávit illum: \Star{} Et dedit illi claritátem ætérnam.\end{Resp}
\begin{Resp}\Vsign Descendítque cum illo in fóveam, et in vínculis non derelíquit eum.\end{Resp}
\begin{Resp}\Rsign Et dedit illi claritátem ætérnam.\end{Resp}
\switchcolumn
\EN
\begin{Resp}\Rsign The Lord made him honourable, and defended him from his enemies, and kept him safe from those that lay in wait for him, \Star{} And gave him perpetual glory.\end{Resp}
\begin{Resp}\Vsign He went down with him into the pit, and left him not in bonds.\end{Resp}
\begin{Resp}\Rsign And gave him perpetual glory.\end{Resp}
\switchcolumn*

\LA
\LessonHead{Lectio V}
\switchcolumn
\EN
\LessonHead{Lesson V}
\switchcolumn*

\LA
\LessonTitle{De Expositione sancti Ambrósii Epíscopi in Psalmum centesimum decimum octavum.}
\switchcolumn
\EN
\LessonTitle{From the Exposition of the hundred-and-eighteenth Psalm by St. Ambrose, Bishop of Milan.}
\switchcolumn*

\LA
\Source{Lectio v. Serm. 21.}
\switchcolumn
\EN
\Source{21st Sermon.}
\switchcolumn*

\LA
\noindent Principes persecuti sunt me gratis: et a verbis tuis trepidavit cor meum. Bene hoc Martyr dicit, quod injuste persecutionum tormenta sustineat: qui nihil rapuerit, nullum violentus oppresserit, nullius sanguinem fuderit, nullius torum putaverit esse violandum: qui nihil legibus debeat, et graviora latronum sustinere cogatur supplicia: qui loquatur juste, et non audiatur: qui loquatur plena salutis, et impugnetur: ut possit dicere, Cum loquebar illis, impugnabant me gratis. Gratis igitur persecutionem patitur, qui impugnatur sine crimine impugnatur ut noxius, cum sit in tali confessione laudabilis: impugnatur quasi veneficus, qui in nomine Domini gloriatur, cum pietas virtutum omnium fundamentum sit.\par
\switchcolumn
\EN
\noindent Princes have persecuted me I without a cause; but my heart standeth in awe of thy word. These are rightly the words of a martyr, who beareth unjustly the torments of the persecutors, who hath robbed no man, who hath violently oppressed no man, who hath shed the blood of no man, who hath imagined to defile the bed of no man, who is debtor to the laws in nothing, and who is punished more grievously than if he were a robber: who speaketh righteousness, and there is none that will hear: who speaketh salvation, and all men fight against him: who is able to say: “When I spoke unto them, they fought against me without a cause.” (Ps. cxix. 7.) They fight against him without a cause, who can lay no sin to his charge; they fight against him as an evildoer, who is by their own acknowledgment righteous: they fight against him as a warlock, who glorieth in the name of the Lord, and who doeth all things well because he doeth all things for God's sake.\par
\switchcolumn*

\LA
\begin{Resp}\Rsign Desidérium ánimæ ejus tribuísti ei Dómine, \Star{} Et voluntáte labiórum ejus non fraudásti eum.\end{Resp}
\begin{Resp}\Vsign Quóniam prævenísti eum in benedictiónibus dulcédinis: posuísti in cápite ejus corónam de lápide pretióso.\end{Resp}
\begin{Resp}\Rsign Et voluntáte labiórum ejus non fraudásti eum.\end{Resp}
\switchcolumn
\EN
\begin{Resp}\Rsign Lord, Thou hast given him his heart's desire. \Star{} And hast not withholden the request of his lips.\end{Resp}
\begin{Resp}\Vsign For Thou hast prevented him with the blessings of sweetness: Thou hast set a crown of precious stones upon his head.\end{Resp}
\begin{Resp}\Rsign And hast not withholden the request of his lips.\end{Resp}
\switchcolumn*

\LA
\LessonHead{Lectio VI}
\switchcolumn
\EN
\LessonHead{Lesson VI}
\switchcolumn*

\LA
\noindent Vere frustra impugnatur, qui apud impios et infidos impietatis arcessitur, cum fidei sit magister. Verum qui gratis impugnatur, fortis debet esse et constans, quomodo ergo subtexuit: Et a verbis tuis trepidavit cor meum? Trepidare infirmitatis est, timoris atque formidinis. Sed est étiam infirmitas ad salutem, est étiam timor sanctorum. Timete Dominum omnes sancti ejus: et, Beatus vir, qui timet Dominum. Qua ratione beatus? Quia in mandatis ejus cupit nimis.\par
\switchcolumn
\EN
\noindent They fight against him in vain who is accused of ungodliness among the ungodly and the unfaithful, because he teacheth Faith. Verily, him that is fought against without a cause it behoveth to be strong and patient. Wherefore then saith he: My heart standeth in awe of thy word? Awe is the mark of the weak, the timid, and the fearful But there is also a weakness unto salvation, there is a fear which is an holy fear. O fear the Lord, all ye His Saints. (Ps. xxxiii. 10.) And again: Blessed is the man that feareth the Lord. (Ps. cxi. i.) And wherefore is he blessed? because he delighteth greatly in His commandments.\par
\switchcolumn*

\LA
\begin{Resp}\Rsign Stola jucunditátis índuit eum Dóminus: \Star{} Et corónam pulchritúdinis pósuit super caput ejus.\end{Resp}
\begin{Resp}\Vsign Cibávit illum Dóminus pane vitæ et intelléctus: et aqua sapiéntiæ salutáris potávit illum.\end{Resp}
\begin{Resp}\Rsign Et corónam pulchritúdinis pósuit super caput ejus.\end{Resp}
\begin{Resp}\Vsign Glória Patri, et Fílio, \Star{} et Spirítui Sancto.\end{Resp}
\begin{Resp}\Rsign Et corónam pulchritúdinis pósuit super caput ejus.\end{Resp}
\switchcolumn
\EN
\begin{Resp}\Rsign The Lord hath put on him a robe of honour, \Star{} And put about his head a crown of joy.\end{Resp}
\begin{Resp}\Vsign With the bread of life and understanding hath the Lord fed him, and given him the water of health and wisdom to drink.\end{Resp}
\begin{Resp}\Rsign And put about his head a crown of joy.\end{Resp}
\begin{Resp}\Vsign Glory be to the Father, and to the Son, \Star{} and to the Holy Ghost.\end{Resp}
\begin{Resp}\Rsign And put about his head a crown of joy.\end{Resp}
\switchcolumn*

\LA
\LessonHead{Lectio VII}
\switchcolumn
\EN
\LessonHead{Lesson VII}
\switchcolumn*

\LA
\LessonTitle{Léctio sancti Evangélii secúndum Matthǽum}
\switchcolumn
\EN
\LessonTitle{The Lesson is taken from the Holy Gospel according to Matthew}
\switchcolumn*

\LA
\Cite{Matt 10:26-32}
\switchcolumn
\EN
\Cite{Matt 10:26-32}
\switchcolumn*

\LA
\noindent In illo témpore: Dixit Jesus discípulis suis: Nihil est opértum, quod non revelábitur; et occúltum, quod non sciétur. Et réliqua.\par
\switchcolumn
\EN
\noindent At that time, Jesus said unto His disciples: There is nothing covered, that shall not be revealed, and hid, that shall not be known. And so on.\par
\switchcolumn*

\LA
\LessonTitle{Homilía sancti Hilárii Epíscopi}
\switchcolumn
\EN
\LessonTitle{Homily by St Hilary, Bishop (of Poitiers.)}
\switchcolumn*

\LA
\Source{Comment. in Matthaeum, can. 10 post medium}
\switchcolumn
\EN
\Source{Comm. on Matt Chap. 10}
\switchcolumn*

\LA
\noindent Dóminus diem judícii osténdit, quæ abstrúsam voluntátis nostræ consciéntiam prodet: et ea quæ nunc occúlta existimántur, luce cognitiónis públicæ déteget. Ígitur non minas, non consília, non potestátes insectántium monet esse metuéndas: quia dies judícii nulla hæc fuísse, atque inánia revelábit. Et quod dico vobis in ténebris, dícite in lúmine: et quod in aure audítis, prædicáte super tecta. Non légimus Dóminum sólitum fuísse nóctibus sermocinári, et doctrínam in ténebris tradidísse: sed quia omnis sermo ejus carnálibus ténebræ sunt, et verbum ejus infidélibus nox est.\par
\switchcolumn
\EN
\noindent The Lord pointeth to the day of judgment, that day wherein the hidden counsels of the hearts shall be made manifest, and those things which are dark now shall be the subject of all men's knowledge. Therefore He warneth us not to fear threats, nor persuasions, nor the power of such as fight against us since in the day of judgment it will be manifest that all these things are null and void. “And what I tell you in darkness, that speak ye in light; and what ye hear in the ear, that preach ye upon the house-tops.” We read not that the Lord's use was to speak by night, or to tell His doctrine in darkness, but that to the carnal all His words were darkness, and to the unbelieving all His discourse night.\par
\switchcolumn*

\LA
\begin{Resp}\Rsign Coróna áurea super caput ejus, \Star{} Expréssa signo sanctitátis, glória honóris, et opus fortitúdinis.\end{Resp}
\begin{Resp}\Vsign Quóniam prævenísti eum in benedictiónibus dulcédinis, posuísti in cápite ejus corónam de lápide pretióso.\end{Resp}
\begin{Resp}\Rsign Expréssa signo sanctitátis, glória honóris, et opus fortitúdinis.\end{Resp}
\switchcolumn
\EN
\begin{Resp}\Rsign A crown of gold upon his head, \Star{} Wherein is engraved Holiness, an ornament of honour, a costly work.\end{Resp}
\begin{Resp}\Vsign For Thou hast prevented him with the blessings of sweetness, Thou hast set a crown of precious stones upon his head.\end{Resp}
\begin{Resp}\Rsign Wherein is engraved Holiness, an ornament of honour, a costly work.\end{Resp}
\switchcolumn*

\LA
\LessonHead{Lectio VIII}
\switchcolumn
\EN
\LessonHead{Lesson VIII}
\switchcolumn*

\LA
\noindent Ítaque id quod a se dictum est, cum libertáte fídei et confessiónis vult esse loquéndum. Idcírco quæ in ténebris dicta sunt prædicári jussit in lúmine: ut quæ secréto áurium commíssa sunt, super tecta, id est, excélso loquéntium præcónio audiántur. Constánter enim Dei ingerénda cognítio est, et profúndum doctrínæ evangélicæ secrétum in lúmine prædicatiónis apostólicæ revelándum: non timéntes eos, quibus cum sit licéntia in córpora, tamen in ánimam jus nullum est: sed timéntes pótius Deum, cui perdéndæ in gehénna et ánimæ et córporis sit potéstas.\par
\switchcolumn
\EN
\noindent Therefore willeth He that that which He hath spoken, should be freely proclaimed in faith and in confession. Therefore commandeth He that that which He hath told in darkness shall be spoken in light, and that that which He hath made to be heard in the ear should be preached upon the house-tops, that is, with loud and high words. For it behoveth us ever to make God known, and to speak in the light of Apostolic preaching the dark things of the Gospel message, having no fear of them which have power over bodies, but none over our souls, but rather fearing God, Which is able to destroy both body and soul in hell.\par
\switchcolumn*

\LA
\begin{Resp}\Rsign Hic est vere Martyr, qui pro Christi nómine sánguinem suum fudit: \Star{} Qui minas júdicum non tímuit, nec terrénæ dignitátis glóriam quæsívit, sed ad cæléstia regna pervénit.\end{Resp}
\begin{Resp}\Vsign Justum dedúxit Dóminus per vias rectas, et osténdit illi regnum Dei.\end{Resp}
\begin{Resp}\Rsign Qui minas júdicum non tímuit, nec terrénæ dignitátis glóriam quæsívit, sed ad cæléstia regna pervénit.\end{Resp}
\begin{Resp}\Vsign Glória Patri, et Fílio, \Star{} et Spirítui Sancto.\end{Resp}
\begin{Resp}\Rsign Qui minas júdicum non tímuit, nec terrénæ dignitátis glóriam quæsívit, sed ad cæléstia regna pervénit.\end{Resp}
\switchcolumn
\EN
\begin{Resp}\Rsign This is a Martyr indeed, who shed his blood for Christ's Name's sake \Star{} Who feared not for the threats of judges, nor sought to be great with the glory of this world, but pressed on unto the kingdom of heaven.\end{Resp}
\begin{Resp}\Vsign The Lord guided the righteous in right paths, and showed him the kingdom of God.\end{Resp}
\begin{Resp}\Rsign Who feared not for the threats of judges, nor sought to be great with the glory of this world, but pressed on unto the kingdom of heaven.\end{Resp}
\begin{Resp}\Vsign Glory be to the Father, and to the Son, \Star{} and to the Holy Ghost.\end{Resp}
\begin{Resp}\Rsign Who feared not for the threats of judges, nor sought to be great with the glory of this world, but pressed on unto the kingdom of heaven.\end{Resp}
\switchcolumn*

\LA
\LessonHead{Lectio IX}
\switchcolumn
\EN
\LessonHead{Lesson IX}
\switchcolumn*

\LA
\noindent Nolíte timére eos qui occídunt corpus. Nullus ígitur córporum nostrórum casus est pertimescéndus, neque ullus interiméndæ carnis admitténdus est dolor: quando pro natúræ suæ atque oríginis conditióne resolúta, in substántiam spirituális ánimæ refundátur. Et quia doctrínis tálibus confirmátos opórtet líberam confiténdi Dei habére constántiam, étiam conditiónem, qua tenerémur, adjécit: negatúrum se eum Patri in cælis, qui se homínibus in terra negásset: eum porro qui conféssus coram homínibus se fuísset, a se in cælis confiténdum: qualésque nos nóminis sui testes homínibus fuissémus, tali nos apud Deum Patrem testimónio ejus usúros.\par
\switchcolumn
\EN
\noindent Fear not them which kill the body. Therefore we need fear nothing which may chance to our bodies, nor sorrow because of the destruction of the flesh, when, according to the laws of our nature and that from whence we are taken, we are unclothed upon, and become a pure spirit. And, since it behoveth us who are rooted in such a doctrine, freely and constantly to confess God, even were it only because of the alternative whereby we are bound, He saith further: “Whosoever shall confess Me before men, him will I confess also before My Father, Which is in heaven. But whosoever shall deny Me before men, him will I also deny before My Father, Which is in heaven.” Such witnesses as He hath seen us to have been here to His name before men, such a Witness shall we find Him to be hereafter to our names before His Father Which is in heaven.\par
\switchcolumn*

\LA
\TeDeum{Te Deum}
\switchcolumn
\EN
\TeDeum{Te Deum}
\switchcolumn*

\end{paracol}

\par\needspace{10\baselineskip}\addvspace{1.2em}{\centering\small\textsc{Die 27 Januarii} · \textsc{27 January}\par}
\Office{S. Joannis Chrysostomi Episcopi Confessoris et Ecclesiæ Doctoris}{St. John Chrysostom, Bishop, Confessor, and Doctor of the Church}{Duplex}
\addcontentsline{toc}{subsection}{Die 27 Januarii · S. Joannis Chrysostomi Episcopi Confessoris et Ecclesiæ Doctoris\enspace\textit{St. John Chrysostom, Bishop, Confessor, and Doctor of the Church}}
\begin{paracol}{2}
\LA
\RefLine{Lectiones I–III: de Scriptura occurrente.}
\switchcolumn
\EN
\RefLine{Lessons I–III: from the Scripture of the day.}
\switchcolumn*

\LA
\LessonHead{Lectio IV}
\switchcolumn
\EN
\LessonHead{Lesson IV}
\switchcolumn*

\LA
\noindent Joánnes Antiochénus, propter áureum eloquéntiæ flumen cognoménto Chrysóstomus, a forensibus et sæcularibus stúdiis ad divinas litteras summa cum ingenii et industriæ laude se cóntulit. Itaque sacris initiátus ac presbyter Antiochenæ ecclésiæ factus, mortuo Nectario, Arcadii imperatóris opera, invitus Constantinopolitanæ ecclésiæ præfícitur. Quo suscepto pastorali munere, depravatos mores, et nobiliórum hóminum vivéndi licéntiam vehementius objurgare cœpit: qua ex libertáte magnam multórum subiit invidiam. Apud Eudoxiam etiam, quod eam propter Callítropæ víduæ pecúniam, et alterius víduæ agrum reprehendísset, gráviter offendit.\par
\switchcolumn
\EN
\noindent John of Antioch, who, on account of the golden stream of his eloquence, is called by the Greeks Chrysostomos, or, the golden-mouthed, was a lawyer and man of the world of much eminence, before he turned his great intellect and wonderful industry to the study of things sacred. He took orders, and was ordained a priest of the Church of Antioch, (in the year of our Lord 386,) and after the death of Nectarius, was forced by the Emperor Arcadius to accept, though sorely against his own will, the Archbishopric of Constantinople. Having received the burden of a shepherd's office, (upon the 26th day of February,) in the year 398, he set himself zealously to do his duty, struggling against the degradation of public morality and the loose lives of the nobility, and thereby drew upon himself the ill-will of many enemies, especially the Empress Eudoxia, whom he had rebuked on account of the money of the widow Callitropa, and the land of another widow.\par
\switchcolumn*

\LA
\begin{Resp}\Rsign Invéni David servum meum, óleo sancto meo unxi eum: \Star{} Manus enim mea auxiliábitur ei.\end{Resp}
\begin{Resp}\Vsign Nihil profíciet inimícus in eo, et fílius iniquitátis non nocébit ei.\end{Resp}
\begin{Resp}\Rsign Manus enim mea auxiliábitur ei.\end{Resp}
\switchcolumn
\EN
\begin{Resp}\Rsign I have found David My servant, with My holy oil have I anointed him \Star{} For My hand shall help him.\end{Resp}
\begin{Resp}\Vsign The enemy shall prevail nothing against him, nor the son of wickedness afflict him.\end{Resp}
\begin{Resp}\Rsign For My hand shall help him.\end{Resp}
\switchcolumn*

\LA
\LessonHead{Lectio V}
\switchcolumn
\EN
\LessonHead{Lesson V}
\switchcolumn*

\LA
\noindent Quare aliquot episcopórum acto Chalcédone convéntu, quo ipse vocátus ire nóluit, quod nec legítimum concílium nec publicum esse diceret, niténte in primis ipsa contra Chrysóstomum Eudoxia, ejícitur in exsílium: sed paulo post propter ejus desidérium seditióne pópuli facta, admirábili civitátis plausu ab exsílio revocátur. Verum cum pérditos mores increpare non desísteret, et ad argenteam Eudoxiæ státuam in foro sanctæ Sophíæ ludos fíeri prohibéret; conspiratióne inimicórum episcopórum íterum exsulare cógitur, viduis et egéntibus ómnibus communis parentis ejectiónem lugéntibus. In exsílio Chrysóstomus incredibile est et quanta mala perpéssus sit, et quam multos ad Jesu Christi fidem convérterit.\par
\switchcolumn
\EN
\noindent Come Bishops being assembled in a Council at Chalcedon, which Council the Saint held to be neither lawful, nor public, although he was commanded to go there, he refused. Whereupon Eudoxia, striving earnestly against him, caused him to be sent into exile. Soon after, however, the people of the city rose, and demanded his recall, and he was then brought back again amid great public rejoicings. Nevertheless he ceased not to war against vice, and absolutely forbade the celebration of public games round the silver statue of Eudoxia in the square outside the Church of the Eternal Wisdom. Upon this, a party of Bishops, who were enemies to him, banded together, and obtained that he should be banished again, which was done accordingly, (on the 20th day of June, 404,) amid the lamentations of widows and the poor, who felt as if they were being deprived of a common father. During this exile, it almost passeth belief how much Chrysostom suffered, and how many souls he turned to the faith which is in Christ Jesus.\par
\switchcolumn*

\LA
\begin{Resp}\Rsign Pósui adjutórium super poténtem, et exaltávi eléctum de plebe mea: \Star{} Manus enim mea auxiliábitur ei.\end{Resp}
\begin{Resp}\Vsign Invéni David servum meum, óleo sancto meo unxi eum.\end{Resp}
\begin{Resp}\Rsign Manus enim mea auxiliábitur ei.\end{Resp}
\switchcolumn
\EN
\begin{Resp}\Rsign I have laid help upon one that is mighty, and have exalted one chosen out of My people \Star{} For My hand shall help him.\end{Resp}
\begin{Resp}\Vsign I have found David My servant, with My holy oil have I anointed him.\end{Resp}
\begin{Resp}\Rsign For My hand shall help him.\end{Resp}
\switchcolumn*

\LA
\LessonHead{Lectio VI}
\switchcolumn
\EN
\LessonHead{Lesson VI}
\switchcolumn*

\LA
\noindent Verum dum, concílio Romæ hábito, decreto Innocentii primi Pontificis restitúitur, a milítibus, qui eum custodiébant, miris in itinere malis et calamitátibus affícitur. Cumque per Armeniam ducerétur, sanctus Basiliscus Martyr, in cujus templo antea oraverat, noctu sic eum affátus est: Joánnes frater, crástinus dies nos loco conjunget. Quare postridie sumpto Eucharistiæ sacraménto seque crucis signo muniens, ánimam Deo réddidit, décimo octavo Kalendas Octobris. Quo mortuo, horribilis grando Constantinópoli cécidit, et quatríduo Augusta cessit e vita. Ejus corpus insígni pompa et hóminum multitúdine celebrátum, Theodosius Arcadii fílius Constantinópolim portándum, et honorifice sepeliéndum curávit sexto Kalendas Februarii: cujus étiam relíquias veneratus, paréntum suórum véniam pétiit: quod deínde Romam translátum in basilica Vaticana cónditum est. Multitúdinem, pietátem, ac splendórem conciónum ceterorúmque ejus scriptórum, interpretándi étiam ratiónem, et inhærentem senténtiæ sacrórum librórum explanatiónem omnes admirántur; dignumque existimant, cui Paulus Apóstolus, quem ille mirifice coluit, scribénti et prædicanti multa dictasse videátur. Hunc vero præclaríssimum univérsæ Ecclésiæ Doctórem Pius décimus Pontifex maximus cælestem oratórum sacrórum patronum declarávit atque constítuit.\par
\switchcolumn
\EN
\noindent At this time a Council was assembled at Rome, wherein Chrysostom's restoration to his See was decreed by Pope Innocent I., but meanwhile, he was suffering great hardships and cruelties on his journey at the hands of the soldiers who had him in charge. As he passed through Armenia he prayed in the Church of the holy martyr Basiliscus, and the same night that blessed conqueror appeared to him in a vision and said Brother John, to-morrow thou shalt be with me. On the next day, therefore, he received the Sacrament of the Eucharist, and, arming himself with the sign of the cross, resigned his soul to God, it being the 14th of September, (in the year of salvation, 407.) As soon as he was dead a furious hailstorm took place at Constantinople, and after four days the Empress died. The Emperor Theodosius, the son of Arcadius, brought the body of John Chrysostom to Constantinople with great state, and numerously attended, and on the 27th of January, (438,) laid it with magnificent honours in the grave, beside which he prayed for the forgiveness of his own father and mother. The holy body was afterwards taken to Rome, and is now buried in the Vatican Basilica. The number, devoutness, and brilliance of St. John Chrysostom's sermons and other writings, his acuteness in exposition, and the close aptness of his explanations of Holy Scripture, have been and are the object of universal wonder and admiration, and often seem not unworthy to have been dictated to him by the Apostle Paul, for whom he entertained a wonderful devotion.\par
\switchcolumn*

\LA
\begin{Resp}\Rsign Iste est, qui ante Deum magnas virtútes operátus est, et omnis terra doctrína ejus repléta est: \Star{} Ipse intercédat pro peccátis ómnium populórum.\end{Resp}
\begin{Resp}\Vsign Iste est, qui contémpsit vitam mundi, et pervénit ad cæléstia regna.\end{Resp}
\begin{Resp}\Rsign Ipse intercédat pro peccátis ómnium populórum.\end{Resp}
\begin{Resp}\Vsign Glória Patri, et Fílio, \Star{} et Spirítui Sancto.\end{Resp}
\begin{Resp}\Rsign Ipse intercédat pro peccátis ómnium populórum.\end{Resp}
\switchcolumn
\EN
\begin{Resp}\Rsign This is he which wrought great wonders before God, and the whole earth is full of his teaching \Star{} May he pray for all people, that their sins may be forgiven unto them\end{Resp}
\begin{Resp}\Vsign This is he which loved not his life in this world, and hath attained unto the kingdom of heaven.\end{Resp}
\begin{Resp}\Rsign May he pray for all people, that their sins may be forgiven unto them\end{Resp}
\begin{Resp}\Vsign Glory be to the Father, and to the Son, \Star{} and to the Holy Ghost.\end{Resp}
\begin{Resp}\Rsign May he pray for all people, that their sins may be forgiven unto them\end{Resp}
\switchcolumn*

\LA
\LessonHead{Lectio VII}
\switchcolumn
\EN
\LessonHead{Lesson VII}
\switchcolumn*

\LA
\LessonTitle{Léctio sancti Evangélii secúndum Matthǽum}
\switchcolumn
\EN
\LessonTitle{From the Holy Gospel according to Matthew}
\switchcolumn*

\LA
\Cite{Matt 5:13-19}
\switchcolumn
\EN
\Cite{Matt 5:13-19}
\switchcolumn*

\LA
\noindent In illo témpore: Dixit Jesus discípulis suis: Vos estis sal terræ. Quod si sal evanúerit, in quo saliétur? Et réliqua.\par
\switchcolumn
\EN
\noindent At that time: Jesus said unto His disciples: Ye are the salt of the earth. But if the salt have lost his savour, wherewith shall it be salted? And so on.\par
\switchcolumn*

\LA
\LessonTitle{Homilía sancti Joánnis Chrysóstomi}
\switchcolumn
\EN
\LessonTitle{Homily by St John Chrysostom, Patriarch (of Constantinople.)}
\switchcolumn*

\LA
\Source{Homilia 15 in Matthæum, sub medium}
\switchcolumn
\EN
\Source{15/A on Matth.}
\switchcolumn*

\LA
\noindent Atténdite quid díxerit, Vos estis sal terræ: per quod osténdit quam necessário ista præcípiat. Non enim de vestra, inquit, tantúmmodo vita, sed de univérso orbe vobis ratio reddénda est. Non ad duas quippe urbes, aut decem, aut vigínti, neque ad unam gentem vos mitto, sicut mittébam prophétas: sed ad omnem terram prorsus ac mare, totúmque mundum, et hunc váriis crimínibus oppréssum.\par
\switchcolumn
\EN
\noindent Consider how that the Lord saith: “Ye are the salt of the earth,” by the which figure He showeth what a necessary of life is the Gospel. By this figure, He hath us to know that they unto whom He spake have an account to render, not of their own life only, but for the whole world. Not unto two cities, saith the Lord, nor unto ten, nor unto twenty, nor unto one people, as I sent the Prophets, send I you. But I send you unto every land and sea, even unto the whole world, lying groaning, as it is, under the burden of diverse sins.\par
\switchcolumn*

\LA
\begin{Resp}\Rsign Amávit eum Dóminus, et ornávit eum: stolam glóriæ índuit eum, \Star{} Et ad portas paradísi coronávit eum.\end{Resp}
\begin{Resp}\Vsign Induit eum Dóminus lorícam fídei, et ornávit eum.\end{Resp}
\begin{Resp}\Rsign Et ad portas paradísi coronávit eum.\end{Resp}
\switchcolumn
\EN
\begin{Resp}\Rsign The Lord loved him and beautified him He clothed him with a robe of glory \Star{} And crowned him at the gates of Paradise.\end{Resp}
\begin{Resp}\Vsign The Lord hath put on him the breast-plate of faith, and hath adorned him.\end{Resp}
\begin{Resp}\Rsign And crowned him at the gates of Paradise.\end{Resp}
\switchcolumn*

\LA
\LessonHead{Lectio VIII}
\switchcolumn
\EN
\LessonHead{Lesson VIII}
\switchcolumn*

\LA
\noindent Dicéndo enim, Vos estis sal terræ, osténdit univérsam hóminum infatuátam esse natúram, et peccatórum vi corrúptam: et idcírco illas ab eis virtútes requírit, quæ máxime ad multórum salútem procurándam necessáriæ sunt atque útiles. Nam qui mansuétus est ac modéstus, et miséricors et justus, non intra se tantúmmodo hæc recte facta conclúdit, verum in aliórum quoque utilitátem præcláros hos fáciet efflúere fontes. Igitur qui corde mundo est atque pacíficus, et persecutiónem pro veritáte pátitur, nihilóminus in commúne cómmodum vitam instítuit.\par
\switchcolumn
\EN
\noindent These words, “Ye are the salt of the earth,” show unto us the whole nature of man as savourless and stinking with the strong corruption of sin. And therefore demandeth He of His Apostles such qualities as are most needful and useful to the furthering the salvation of many. He that is gentle and lowly, tender and just, shutteth not up all these good things in his own heart, but openeth these bright fountains that they may gush forth for the use of his neighbour. He whose heart is pure, and who seeketh peace, suffering persecution for the truth's sake, doth still lead a life for the good of the commonwealth.\par
\switchcolumn*

\LA
\begin{Resp}\Rsign In médio Ecclésiæ apéruit os ejus, \Star{} Et implévit eum Dóminus spíritu sapiéntiæ et intelléctus.\end{Resp}
\begin{Resp}\Vsign Jucunditátem et exsultatiónem thesaurizávit super eum.\end{Resp}
\begin{Resp}\Rsign Et implévit eum Dóminus spíritu sapiéntiæ et intelléctus.\end{Resp}
\begin{Resp}\Vsign Glória Patri, et Fílio, \Star{} et Spirítui Sancto.\end{Resp}
\begin{Resp}\Rsign Et implévit eum Dóminus spíritu sapiéntiæ et intelléctus.\end{Resp}
\switchcolumn
\EN
\begin{Resp}\Rsign In the midst of the congregation did the Lord open his mouth. \Star{} And filled him with the spirit of wisdom and understanding.\end{Resp}
\begin{Resp}\Vsign He made him rich with joy and gladness.\end{Resp}
\begin{Resp}\Rsign And filled him with the spirit of wisdom and understanding.\end{Resp}
\begin{Resp}\Vsign Glory be to the Father, and to the Son, \Star{} and to the Holy Ghost.\end{Resp}
\begin{Resp}\Rsign And filled him with the spirit of wisdom and understanding.\end{Resp}
\switchcolumn*

\LA
\LessonHead{Lectio IX}
\switchcolumn
\EN
\LessonHead{Lesson IX}
\switchcolumn*

\LA
\noindent Ne ígitur putétis, inquit, ad lévia vos ducéndos esse certámina, neque exiguárum rerum vobis ineúndam esse ratiónem. Vos estis sal terræ. Quid ígitur? Ipsine putrefácta medicáti sunt? Nequáquam: neque enim fíeri potest, ut ea quæ jam corrúpta sunt, salis perfricatióne reparéntur. Non ergo hoc fecérunt, sed ante renováta, sibíque trádita, atque ab illa jam putrédine liberáta, aspergébant sale, et in ea novitáte conservábant, quam a Dómino suscéperant. Liberáre quippe a putrédine peccatórum, Christi virtútis est: ut autem ad illa íterum non revertántur, Apostolórum curæ est ac labóris.\par
\switchcolumn
\EN
\noindent Think not, saith the Lord, that the struggle is easy whereunto ye shall be led, neither shall your reckoning be of light matters. Ye are the salt of the earth. Have ye then salted that which is corrupted? Nay, for it is impossible that that which is once corrupted can be made sound again by the rubbing it with salt. This it is not asked of them to do. But their work is to sprinkle with salt, and to keep fresh thereafter, such things as the Lord hath given over into their charge, and which He Himself hath made new, and freed from all taint, before giving them. To make sound after the corruption of sin, is the work of Christ's power alone; to preserve from falling away again, is the duty and the toil commanded to the Apostles.\par
\switchcolumn*

\LA
\TeDeum{Te Deum}
\switchcolumn
\EN
\TeDeum{Te Deum}
\switchcolumn*

\end{paracol}

\par\needspace{10\baselineskip}\addvspace{1.2em}{\centering\small\textsc{Die 28 Januarii} · \textsc{28 January}\par}
\Office{S. Petri Nolasci Confessoris}{St. Peter Nolasco, Confessor}{Duplex}
\addcontentsline{toc}{subsection}{Die 28 Januarii · S. Petri Nolasci Confessoris\enspace\textit{St. Peter Nolasco, Confessor}}
\begin{paracol}{2}
\LA
\RefLine{Lectiones I–III: de Scriptura occurrente.}
\switchcolumn
\EN
\RefLine{Lessons I–III: from the Scripture of the day.}
\switchcolumn*

\LA
\RefLine{Lectiones VII–VIII: ex Commúni Confessoris non pontificis (C5).}
\switchcolumn
\EN
\RefLine{Lessons VII–VIII: from the Common of a Confessor not a Bishop (C5).}
\switchcolumn*

\LA
\RefLine{Lectio IX: de S. Agnetis Virginis et Martyris secundo (01-28t).}
\switchcolumn
\EN
\RefLine{Lesson IX: of S. Agnes Virgin and Martyr, second (01-28t).}
\switchcolumn*

\LA
\LessonHead{Lectio IV}
\switchcolumn
\EN
\LessonHead{Lesson IV}
\switchcolumn*

\LA
\noindent Petrus Nolascus, Recáudi prope Carcasonam in Gállia nobili genere natus, singulari erga próximum caritate excelluit. Cujus virtútis præsagium fuit, quod, cum adhuc in cunábulis vagiret infans, exámen apum ad eum convolávit et favum mellis in ejus déxtera construxit. Adoléscens paréntibus orbatus, Albigensium hæresim, quæ tunc in Gállia grassabátur, éxsecrans, divéndito patrimonio, in Hispániam secessit; et apud beátam Vírginem montis Serráti votum, quo pridem se obstrinxerat, exsolvit. Tum Barcinonam pergens, cum Christi fidelibus ab hostium servitute rediméndis omnem pecúniam consumpsísset, seípsum pro iis liberandis venum ire, aut in illórum víncula súffici cúpere dictitabat.\par
\switchcolumn
\EN
\noindent Peter Nolasco was born of noble parents at Recaudun near Carcassonne in France \textasciitilde{}(about the year 1189,) and is chiefly distinguished for his great love toward his neighbour. It was considered a foreshadowing of this virtue, that when he was a little child in his cradle, a swarm of bees settled on his right hand, and began to make an honey-comb there. He lost his parents while still young, and in consequence of his horror of the Albigensian heresy, with which France was then plagued, he sold his property there and emigrated to Spain. Here he first discharged a vow which he had made at the sanctuary of the Blessed Virgin of Monserrat, and afterwards went to Barcelona. Here he was so affected by the miserable state of the Christians who were in slavery to the Moors, that he expended his whole fortune in ransoming as many of them as possible, and used to say that he wished he could be sold himself to ransom more, or could himself change places with them.\par
\switchcolumn*

\LA
\begin{Resp}\Rsign Honéstum fecit illum Dóminus, et custodívit eum ab inimícis, et a seductóribus tutávit illum: \Star{} Et dedit illi claritátem ætérnam.\end{Resp}
\begin{Resp}\Vsign Justum dedúxit Dóminus per vias rectas, et osténdit illi regnum Dei.\end{Resp}
\begin{Resp}\Rsign Et dedit illi claritátem ætérnam.\end{Resp}
\switchcolumn
\EN
\begin{Resp}\Rsign The Lord made him honourable, and defended him from his enemies, and kept him safe from those that lay in wait for him; \Star{} And gave him perpetual glory.\end{Resp}
\begin{Resp}\Vsign He went down with him into the pit, and left him not in bonds.\end{Resp}
\begin{Resp}\Rsign And gave him perpetual glory.\end{Resp}
\switchcolumn*

\LA
\LessonHead{Lectio V}
\switchcolumn
\EN
\LessonHead{Lesson V}
\switchcolumn*

\LA
\noindent Quam gratum Deo fúerit hoc sancti viri desidérium, súbsequens declarávit eventus; nam noctu oranti, et de Christianórum in captivitáte degéntium subsidio multa animo volvénti beáta Virgo apparens, Fílio suo sibíque acceptíssimum fore suggessit, si ad sui honórem religiosórum ordo instituerétur, quibus præcipue esset cura captivos ab infidelium tyrannide liberare. Huic cælésti mónito íllico obtemperans, una cum sancto Raymundo de Péñafort, et Jacobo primo, rege Aragoniæ, de eádem re a Dei Genitrice ipsa nocte præmónitis, religiónem beátæ Maríæ de Mercede redemptiónis captivórum instituit; sodálibus suis quarto voto obstrictis, manéndi in pignus sub paganórum potestate, si pro Christianórum liberatióne opus fúerit.\par
\switchcolumn
\EN
\noindent It came to pass that God showed how agreeable to Him was the charitable zeal of Peter. One night when he was praying, and his mind was much exercised on the means of succouring the enslaved Christians, the Blessed Virgin appeared to him in a vision, and gave him to understand that it would be most pleasing to her Son and herself, if he would found in her honour an order of religious men, whose chief duty it should be to effect the redemption of Christian bondsmen out of the hand of the unbelievers. In conformity to this revelation, which had likewise on the same night been made to St. Raymond de Penafort and King James I. of Aragon, he founded the Religious Order of the Blessed Mary of Ransom, for the redemption of captives. The members of this order add a fourth vow to the three essential ones of Poverty, Chastity, and Obedience, namely, that they will be ready if need be to remain as hostages in 'the hand of the unbelievers for the liberation of others.\par
\switchcolumn*

\LA
\begin{Resp}\Rsign Amávit eum Dóminus, et ornávit eum: stolam glóriæ índuit eum, \Star{} Et ad portas paradísi coronávit eum.\end{Resp}
\begin{Resp}\Vsign Índuit eum Dóminus lorícam fídei, et ornávit eum.\end{Resp}
\begin{Resp}\Rsign Et ad portas paradísi coronávit eum.\end{Resp}
\switchcolumn
\EN
\begin{Resp}\Rsign The Lord loved him and beautified him; He clothed him with a robe of glory, \Star{} And crowned him at the gates of Paradise.\end{Resp}
\begin{Resp}\Vsign The Lord hath put on him the breast-plate of faith, and hath adorned him.\end{Resp}
\begin{Resp}\Rsign And crowned him at the gates of Paradise.\end{Resp}
\switchcolumn*

\LA
\LessonHead{Lectio VI}
\switchcolumn
\EN
\LessonHead{Lesson VI}
\switchcolumn*

\LA
\noindent Edito virginitátis voto, illibatam perpetuo castitátem servávit. Patiéntia, humilitate, abstinéntia, ceterisque virtútibus mirabíliter enituit. Prophetíæ dono illustris, futura prædixit; inter quæ maxime celebrátur, quod Jacobus rex Valéntiam a Mauris occupatam expugnáverit, accepta prius ab eo obtinendæ victóriæ securitate. Angeli Custodis ac Deíparæ Vírginis frequénti apparitióne recreabátur. Senio tandem confectus, de imminénti morte certior factus, in morbum íncidit, sanctisque refectus sacramentis, fratres suos ad caritátem erga captivos cohortatus, et Psalmum, Confitébor tibi, Dómine, in toto corde meo, devotíssime récitans, ad illa verba, Redemptiónem misit Dóminus pópulo suo, spíritum Deo réddidit media nocte Vigíliæ Nativitátis Dómini, anno millesimo ducentésimo quinquagesimo sexto. Ejus autem festum Alexander septimus ad universam exténdit Ecclésiam.\par
\switchcolumn
\EN
\noindent After he took the vow of virginity he remained with his purity quite unsullied all his life, and was at the same time a bright pattern of long-suffering, lowliness, temperance, and other virtues. God was pleased to adorn him with the gift of Prophecy, whereby he foretold things to come. Among others, he prophesied to King James that he would take the city of Valencia from the Moors, which he afterwards did. He was refreshed by frequent apparitions of his Guardian Angel and of the Virgin Mother of God. He had lived to a great age, when being quite worn out, and falling into a grievous sickness, he perceived that his end was at hand. He then received the holy Sacraments, and, gathering his brethren around him, exhorted them for the last time to show pity to slaves. After this he began to repeat with great emotion the ex. Psalm, I will praise thee, O Lord, with my whole heart, and when he had uttered the words He sent redemption unto His people, he resigned his soul to God. This happened at midnight between the 23rd and 24th of December, 1256. Alexander VII. ordered that his feast should be celebrated on the 31st of January.\par
\switchcolumn*

\LA
\begin{Resp}\Rsign Iste homo perfécit ómnia quæ locútus est ei Deus, et dixit ad eum: Ingrédere in réquiem meam: \Star{} Quia te vidi justum coram me ex ómnibus géntibus.\end{Resp}
\begin{Resp}\Vsign Iste est, qui contémpsit vitam mundi, et pervénit ad cæléstia regna.\end{Resp}
\begin{Resp}\Rsign Quia te vidi justum coram me ex ómnibus géntibus.\end{Resp}
\begin{Resp}\Vsign Glória Patri, et Fílio, \Star{} et Spirítui Sancto.\end{Resp}
\begin{Resp}\Rsign Quia te vidi justum coram me ex ómnibus géntibus.\end{Resp}
\switchcolumn
\EN
\begin{Resp}\Rsign This is he which did according unto all that God commanded him; and God said unto him: Enter thou into My rest, \Star{} For thee have I seen righteous before Me among all people.\end{Resp}
\begin{Resp}\Vsign This is he which loved not his life in this world, and is come unto an everlasting kingdom.\end{Resp}
\begin{Resp}\Rsign For thee have I seen righteous before Me among all people.\end{Resp}
\begin{Resp}\Vsign Glory be to the Father, and to the Son, \Star{} and to the Holy Ghost.\end{Resp}
\begin{Resp}\Rsign For thee have I seen righteous before Me among all people.\end{Resp}
\switchcolumn*

\end{paracol}

\par\needspace{10\baselineskip}\addvspace{1.2em}{\centering\small\textsc{Die 28 Januarii} · \textsc{28 January}\par}
\Office{S. Agnetis Virginis et Martyris secundo}{S. Agnes Virgin and Martyr, second}{Simplex}
\addcontentsline{toc}{subsection}{Die 28 Januarii · S. Agnetis Virginis et Martyris secundo\enspace\textit{S. Agnes Virgin and Martyr, second}}
\begin{paracol}{2}
\LA
\LessonHead{Lectio IX (pro commemoratione)}
\switchcolumn
\EN
\LessonHead{Lesson IX (when commemorated)}
\switchcolumn*

\LA
\LessonTitle{Commemoratio: S. Agnetis Virginis et Martyris secundo}
\switchcolumn
\EN
\LessonTitle{Commemoration of: S. Agnes Virgin and Martyr, second}
\switchcolumn*

\LA
\noindent Beáta Agnes, paréntibus ad ejus sepúlcrum assidue vigilántibus, quadam nocte choro vírginum comitáta cum apparuísset, sic ad eos locuta dícitur: Ne me, paréntes, mortuam lugeátis; nam una cum his virgínibus vivo apud illum in cælis, quem in terris tota mente diléxi. Cum aliquot post annos Constantia, Constantini imperatóris fília, insanábili úlceri medicinam quærens, nondum christiana ad idem sepúlcrum accessísset; obdórmiscens sibi audíre visa est vocem hanc Agnetis: Constanter age, Constántia: crede in Jesum Christum Fílium Dei, qui te sanam fáciet. Quæ sanata, paulo post una cum multis ex família imperatóris baptismum suscípiens, ibi ecclésiam nómine beátæ Agnétis ædificávit.\par
\switchcolumn
\EN
\noindent The night when the parents of the blessed Agnes were watching at her grave, she appeared to them in company with a band of virgins, and said to them Father and Mother, weep not for me as though I were dead; for now these virgins and I live together in Him Whose love was my whole life upon earth. Some years afterwards, Constance, the daughter of the Emperor Constantine, being sick of an incurable ulcer, betook herself to the said grave, although she was not yet a Christian, and as she lay by it and slept, she seemed to hear the voice of Agnes, saying to her Constance, be of good courage believe in Jesus Christ the Son of God, and He will make thee whole. The Princess, being healed, was baptized, along with many others of the Emperor's family and household, and afterwards built over the grave of the blessed Agnes a Church named in her honour.\par
\switchcolumn*

\LA
\TeDeum{Te Deum}
\switchcolumn
\EN
\TeDeum{Te Deum}
\switchcolumn*

\end{paracol}

\par\needspace{10\baselineskip}\addvspace{1.2em}{\centering\small\textsc{Die 29 Januarii} · \textsc{29 January}\par}
\Office{S. Francisci Salesii Episcopi Confessoris et Ecclesiæ Doctoris}{St. Francis de Sales, Bishop, Confessor, and Doctor of the Church}{Duplex}
\addcontentsline{toc}{subsection}{Die 29 Januarii · S. Francisci Salesii Episcopi Confessoris et Ecclesiæ Doctoris\enspace\textit{St. Francis de Sales, Bishop, Confessor, and Doctor of the Church}}
\begin{paracol}{2}
\LA
\RefLine{Lectiones I–III: de Scriptura occurrente.}
\switchcolumn
\EN
\RefLine{Lessons I–III: from the Scripture of the day.}
\switchcolumn*

\LA
\RefLine{Lectiones VII–IX: ex Commúni Doctoris Pontificis (C4a).}
\switchcolumn
\EN
\RefLine{Lessons VII–IX: from the Common of a Doctor Bishop (C4a).}
\switchcolumn*

\LA
\LessonHead{Lectio IV}
\switchcolumn
\EN
\LessonHead{Lesson IV}
\switchcolumn*

\LA
\noindent Franciscus in oppido Salesio, unde famíliæ cognomen, piis et nobílibus paréntibus natus, a téneris annis futuræ sanctitátis indícia præbuit, morum innocéntia et gravitáte. Adoléscens liberálibus disciplinis eruditus, mox philosophíæ ac theologíæ Parisiis operam dedit; et, ne quid sibi deésset ad animi culturam, juris utriusque lauream summa cum laude Patavii obtinuit. In sacra æde Lauretana perpetuæ virginitátis votum, quo pridem Parisiis se obstrinxerat, innovávit; a cujus virtútis proposito nullis umquam dæmonum fraudibus, nullis sensuum illécebris pótuit dimoveri.\par
\switchcolumn
\EN
\noindent Francis was born of godly and noble parents, in the town of Sales, from which his family take their name of de Sales, (upon the 21st day of August, in the year of our Lord 1567.) In his childish years his staid and godly demeanour gave promise of his future sanctity. He received a liberal education as he grew up, and afterwards studied Philosophy and Theology at Paris. In order to the complete furnishing of his mind, he took the degree of Doctor of Laws, both Civil and Ecclesiastical, at Padua, with much distinction. He had already bound himself with a vow of perpetual virginity at Paris, and he renewed the same in the Holy House of Loreto. From this path of virtue, neither the temptations of the devil nor the allurements of the world ever induced him\par
to\par
\switchcolumn*

\LA
\begin{Resp}\Rsign Invéni David servum meum, óleo sancto meo unxi eum: \Star{} Manus enim mea auxiliábitur ei.\end{Resp}
\begin{Resp}\Vsign Nihil profíciet inimícus in eo, et fílius iniquitátis non nocébit ei.\end{Resp}
\begin{Resp}\Rsign Manus enim mea auxiliábitur ei.\end{Resp}
\switchcolumn
\EN
\begin{Resp}\Rsign I have found David My servant, with My holy oil have I anointed him \Star{} For My hand shall help him.\end{Resp}
\begin{Resp}\Vsign The enemy shall prevail nothing against him, nor the son of wickedness afflict him.\end{Resp}
\begin{Resp}\Rsign For My hand shall help him.\end{Resp}
\switchcolumn*

\LA
\LessonHead{Lectio V}
\switchcolumn
\EN
\LessonHead{Lesson V}
\switchcolumn*

\LA
\noindent Recusáta in Sabáudiæ senatu amplíssima dignitate, clericali milítiæ nomen dedit. Tum sacerdotio initiátus et Genevénsis ecclésiæ præposituram adeptus, ejus muneris partes adeo perfecte explevit, ut eum Granerius epíscopus vindicándis ab hæresi Calviniana Chaballicensibus aliisque Genevæ finitimis pópulis, divini verbi præcónem destinarit. Quam expeditiónem álacri animo suscípiens, asperrima quæque perpéssus est, sæpe ab hæréticis conquisítus ad necem, variisque calumniis et insídiis vexatus. Sed inter tot discrímina et agones, insuperábilis ejus constántia semper enituit; Deique ope protectus, septuagínta duo míllia hæreticórum ad catholicam fidem reduxísse dícitur, inter quos multi nobilitate et doctrina insignes numerántur.\par
\switchcolumn
\EN
\noindent He refused to be made Counsellor of the Parliament of Chambery, for which his family had obtained for him patents from the Duke of Savoy, and determined to become a clergyman. He was appointed to the Provostship of the Church of Geneva, and, being shortly afterwards ordained Priest, discharged so admirably the duties of his position, that he was sent by Granier, his Bishop, to preach the word of God in Chablais, and other places in the outskirts of the diocese, where the inhabitants had embraced the heresy of Calvin. He joyfully undertook this mission, in which he suffered much, being often hunted by the Protestants to murder him, and assailed by many calumnies and plots. Amid all these dangers and struggles his constancy remained invincible, and under the blessing and care of God he is said to have recalled seventy-two thousand of these heretics to the Faith of Christ's Universal Church, among whom were many distinguished by rank and learning.\par
\switchcolumn*

\LA
\begin{Resp}\Rsign Pósui adjutórium super poténtem, et exaltávi eléctum de plebe mea: \Star{} Manus enim mea auxiliábitur ei.\end{Resp}
\begin{Resp}\Vsign Invéni David servum meum, óleo sancto meo unxi eum.\end{Resp}
\begin{Resp}\Rsign Manus enim mea auxiliábitur ei.\end{Resp}
\switchcolumn
\EN
\begin{Resp}\Rsign I have laid help upon one that is mighty, and have exalted one chosen out of My people \Star{} For My hand shall help him.\end{Resp}
\begin{Resp}\Vsign I have found David My servant, with My holy oil have I anointed him.\end{Resp}
\begin{Resp}\Rsign For My hand shall help him.\end{Resp}
\switchcolumn*

\LA
\LessonHead{Lectio VI}
\switchcolumn
\EN
\LessonHead{Lesson VI}
\switchcolumn*

\LA
\noindent Mortuo Granerio, qui eum sibi coadjutórem decerni curáverat, epíscopus consecratus, sanctitátis suæ radios circumquaque diffudit, zelo ecclesiásticæ disciplinæ, pacis studio, misericórdia in páuperes, omnique virtúte conspicuus. Ad divini cultus augméntum novum ordinem sanctimoniálium instituit, a Visitatióne beátæ Vírginis nuncupátum, sub regula sancti Augustíni; cui addidit constitutiónes sapiéntia, discretióne et suavitáte mirábiles. Suis étiam scriptis cælésti doctrina refertis Ecclésiam illustrávit, quibus iter ad christianam perfectiónem tutum et planum demonstrat. Annum denique agens quinquagesimum quintum, dum e Gállia Annesium regréditur, post Sacrum in die sancti Joánnis Evangelistæ Lugduni celebrátum, gravi morbo correptus, sequénti die migrávit in cælum, anno Dómini millesimo sexcentésimo vigesimo secundo. Ejus corpus Annesium delátum, in ecclésia moniálium dicti ordinis honorifice cónditum fuit, cœpitque statis miraculis claréscere. Quibus rite probátis, ab Alexandro septimo Pontifice máximo in Sanctórum númerum relátus est, assignáta ejus festivitáti die vigesima nona Januarii; et a summo Pontifice Pio nono, ex sacrórum Rituum Congregatiónis consulto, universalis Ecclésiæ Doctor fuit declaratus.\par
\switchcolumn
\EN
\noindent After the death of Bishop Granier, who had procured his appointment as Coadjutor, he was consecrated Bishop, (upon the 3rd day of December, 1602.) In that office he was truly a burning and a shining light, showing all around a bright example of godliness, zeal for the discipline of the Church, ardent love of peace, tenderness to the poor, and, indeed, of all graces. For the greater ornament of God's worship he established a new Order of Nuns, which is named from the Visitation of the Blessed Virgin. These nuns follow the Rule of St. Augustine, but Francis added thereto several additional constitutions distinguished by wisdom, prudence, and tenderness. He enlightened the Church by writings full of heavenly teaching, and pointing out a safe and simple road to Christian perfection. In the 55th year of his age, while on his way from France to Annecy, after saying mass at Lyons on the Feast of St. John the Evangelist, he was seized with fatal illness, and on the next day passed from earth to heaven, in the year of our Lord 1622. His body was carried to Annecy and honourably buried in the Church of the nuns of the Visitation, where it soon began to be distinguished for miracles. The truth of these having been proved, the Supreme Pontiff, Alexander VI L, enrolled his name among those of the Saints, and appointed for his Feast-day the 29th of January. And the Supreme Pontiff, Pius IX., on the advice of the Congregation of Sacred Rites, declared him a Doctor of the Universal Church.\par
\switchcolumn*

\LA
\begin{Resp}\Rsign Iste est, qui ante Deum magnas virtútes operátus est, et omnis terra doctrína ejus repléta est: \Star{} Ipse intercédat pro peccátis ómnium populórum.\end{Resp}
\begin{Resp}\Vsign Iste est, qui contémpsit vitam mundi, et pervénit ad cæléstia regna.\end{Resp}
\begin{Resp}\Rsign Ipse intercédat pro peccátis ómnium populórum.\end{Resp}
\begin{Resp}\Vsign Glória Patri, et Fílio, \Star{} et Spirítui Sancto.\end{Resp}
\begin{Resp}\Rsign Ipse intercédat pro peccátis ómnium populórum.\end{Resp}
\switchcolumn
\EN
\begin{Resp}\Rsign This is he which wrought great wonders before God, and the whole earth is full of his teaching \Star{} May he pray for all people, that their sins may be forgiven unto them\end{Resp}
\begin{Resp}\Vsign This is he which loved not his life in this world, and hath attained unto the kingdom of heaven.\end{Resp}
\begin{Resp}\Rsign May he pray for all people, that their sins may be forgiven unto them\end{Resp}
\begin{Resp}\Vsign Glory be to the Father, and to the Son, \Star{} and to the Holy Ghost.\end{Resp}
\begin{Resp}\Rsign May he pray for all people, that their sins may be forgiven unto them\end{Resp}
\switchcolumn*

\end{paracol}

\par\needspace{10\baselineskip}\addvspace{1.2em}{\centering\small\textsc{Die 30 Januarii} · \textsc{30 January}\par}
\Office{S. Martinæ Virginis et Martyris}{St. Martina, Virgin and Martyr}{Semiduplex}
\addcontentsline{toc}{subsection}{Die 30 Januarii · S. Martinæ Virginis et Martyris\enspace\textit{St. Martina, Virgin and Martyr}}
\begin{paracol}{2}
\LA
\RefLine{Lectiones I–III: de Scriptura occurrente.}
\switchcolumn
\EN
\RefLine{Lessons I–III: from the Scripture of the day.}
\switchcolumn*

\LA
\RefLine{Lectiones VII–IX: ex Commúni Unius Virginis et Martyris (C6).}
\switchcolumn
\EN
\RefLine{Lessons VII–IX: from the Common of a Virgin Martyr (C6).}
\switchcolumn*

\LA
\LessonHead{Lectio IV}
\switchcolumn
\EN
\LessonHead{Lesson IV}
\switchcolumn*

\LA
\noindent Martina virgo Romana, patre consulari, illustri genere nata, teneris, adhuc annis orbáta paréntibus, christianæ pietátis ardore succensa, divítias, quibus affluebat, mira in páuperes liberalitate distríbuit. Sub Alexandro príncipe, cum deos inánes cólere juberétur, immane fácinus summa libertáte detestátur. Quapropter íterum atque íterum affecta verbéribus, uncis, úngulis férreis, testárum fragmentis lacerata, acutíssimis gládiis membratim concisa, ádipe fervénti peruncta, demum in amphitheatro damnátur ad bestias; a quibus illǽsa divinitus evadens, in ardentem rogum injecta, incólumis pari beneficio servátur.\par
\switchcolumn
\EN
\noindent Martina was a maiden of a most illustrious Roman family, daughter of a Consul. She lost her parents while still very young, and, being inflamed with Christian zeal, she distributed her wealth, whereof she had abundance, with great profusion among the poor. Under the Emperor Alexander, she was commanded to sacrifice to the imaginary gods, and refused with much boldness to commit this great wickedness. Upon this she was again and again scourged, and mangled with iron prongs and hooks, and pieces of broken pottery. Her limbs were cut off piece by piece with sharp swords, and boiling tallow poured upon the living trunk. Lastly she was sent to be eaten publicly by the wild beasts in the amphitheatre, but by the will of God they would not touch her, and she was then thrown upon a burning pile, but still remained alive.\par
\switchcolumn*

\LA
\begin{Resp}\Rsign Propter veritátem, et mansuetúdinem, et justítiam: \Star{} Et dedúcet te mirabíliter déxtera tua.\end{Resp}
\begin{Resp}\Vsign Spécie tua et pulchritúdine tua inténde, próspere procéde, et regna.\end{Resp}
\begin{Resp}\Rsign Et dedúcet te mirabíliter déxtera tua.\end{Resp}
\switchcolumn
\EN
\begin{Resp}\Rsign Because of truth, and meekness, and righteousness; \Star{} And thy right hand shall lead thee wonderfully.\end{Resp}
\begin{Resp}\Vsign In thy comeliness, and thy beauty, go forward, fare prosperously, and reign.\end{Resp}
\begin{Resp}\Rsign And thy right hand shall lead thee wonderfully.\end{Resp}
\switchcolumn*

\LA
\LessonHead{Lectio V}
\switchcolumn
\EN
\LessonHead{Lesson V}
\switchcolumn*

\LA
\noindent Ex ejus tortóribus nonnulli, miraculi novitate correpti, Dei aspirante grátia, Christi fidem ampléxi, post cruciátus gloriósam martyrii palmam cápitis abscissióne promeruére. Ad ejusdem preces nunc terræmótibus exórtis, nunc ignibus e cælo tonante delapsis, deórum templa prostráta sunt et simulacra consumpta. Interdum ex vulnéribus lac cum sánguine erupit, splendorque nitidíssimus ac suavíssimus odor e corpore emanávit: intérdum sublimis regia in sede divinis laudibus una cum cælítibus interesse visa est.\par
\switchcolumn
\EN
\noindent Come of her tormentors were so moved by the spectacle, that they repented, and, by the grace of God confessing the faith of Christ, through which she remained constant, were themselves tortured and beheaded. Martina herself lay praying, with a brightness on her face, while a matter like milk oozed from her body along with the blood, emitting a soft, sweet smell. She was as it were unconscious of an earthquake and most violent thunderstorm which arose and was raging, and while the lightning struck temples, and melted statues, she seemed in spirit rather to be seated above on a queenly throne, praising God in heaven among the Blessed.\par
\switchcolumn*

\LA
\begin{Resp}\Rsign Dilexísti justítiam, et odísti iniquitátem: \Star{} Proptérea unxit te Deus, Deus tuus, óleo lætítiæ.\end{Resp}
\begin{Resp}\Vsign Propter veritátem, et mansuetúdinem, et justítiam.\end{Resp}
\begin{Resp}\Rsign Proptérea unxit te Deus, Deus tuus, óleo lætítiæ.\end{Resp}
\switchcolumn
\EN
\begin{Resp}\Rsign Thou hast loved righteousness, and hated iniquity; \Star{} Therefore God, thy God, hath anointed thee with the oil of gladness.\end{Resp}
\begin{Resp}\Vsign Because of truth, and meekness, and righteousness.\end{Resp}
\begin{Resp}\Rsign Therefore God, thy God, hath anointed thee with the oil of gladness.\end{Resp}
\switchcolumn*

\LA
\LessonHead{Lectio VI}
\switchcolumn
\EN
\LessonHead{Lesson VI}
\switchcolumn*

\LA
\noindent Hisce prodigiis ejusque in primis constantia, acriter permotus judex, caput Virgini amputari præcepit: qua perempta, auditáque de cælo voce, qua ad Súperos evocabátur, Urbs tota contrémuit, ac multi idolórum cultores ad Christi fidem convérsi sunt. Sacrum Martinæ corpus, sedénte sancto Urbano primo, martyrio afféctum, Urbano octavo Pontifice máximo, in pervetusta ejusdem ecclésia, ad Mamertinum carcerem, in Capitolini clivi radicibus, cum sanctórum Mártyrum Concordii, Epiphanii et Sociórum corpóribus repertum, eodem loco in meliórem formam redacto, atque decentius ornato, magno pópuli concursu, totius Urbis lætítia, solemni ritu ac pompa repósitum est.\par
\switchcolumn
\EN
\noindent The judge being infuriated at what had taken place, and chiefly at her unbending firmness, ordered the head of the martyr to be cut off. At the moment this was done, a peal which shook the city was heard, like a voice calling her home, and so great was the consternation, that it was made the means of conversion to many idolaters. The holy body of Martina wherein she had suffered in the Pontificate of Urban I. was discovered in the time of Urban VIII., in the very old Church called after her, situated at the foot of the Capitoline Hill, near the Mamertine Prison, along with the bodies of the holy martyrs Concordius, Epiphanius, and others. The Church was then altered and restored and handsomely decorated, and then the body was replaced in it, amid public rejoicings, with a solemn ceremony and procession.\par
\switchcolumn*

\LA
\begin{Resp}\Rsign Afferéntur Regi vírgines post eam, próximæ ejus; \Star{} Afferéntur tibi in lætítia et exsultatióne.\end{Resp}
\begin{Resp}\Vsign Spécie tua et pulchritúdine tua; inténde, próspere procéde, et regna.\end{Resp}
\begin{Resp}\Rsign Afferéntur tibi in lætítia et exsultatióne.\end{Resp}
\begin{Resp}\Vsign Glória Patri, et Fílio, \Star{} et Spirítui Sancto.\end{Resp}
\begin{Resp}\Rsign Afferéntur tibi in lætítia et exsultatióne.\end{Resp}
\switchcolumn
\EN
\begin{Resp}\Rsign After her shall virgins be brought unto the King, her fellows \Star{} They shall be brought unto thee with gladness and rejoicing.\end{Resp}
\begin{Resp}\Vsign In thy comeliness and thy beauty, go forward, fare prosperously, and reign.\end{Resp}
\begin{Resp}\Rsign They shall be brought unto thee with gladness and rejoicing.\end{Resp}
\begin{Resp}\Vsign Glory be to the Father, and to the Son, \Star{} and to the Holy Ghost.\end{Resp}
\begin{Resp}\Rsign They shall be brought unto thee with gladness and rejoicing.\end{Resp}
\switchcolumn*

\end{paracol}

\par\needspace{10\baselineskip}\addvspace{1.2em}{\centering\small\textsc{Die 31 Januarii} · \textsc{31 January}\par}
\Office{S. Joannis Bosco Confessoris}{St. John Bosco, Confessor}{Duplex}
\addcontentsline{toc}{subsection}{Die 31 Januarii · S. Joannis Bosco Confessoris\enspace\textit{St. John Bosco, Confessor}}
\begin{paracol}{2}
\LA
\RefLine{Lectiones I–III: de Scriptura occurrente.}
\switchcolumn
\EN
\RefLine{Lessons I–III: from the Scripture of the day.}
\switchcolumn*

\LA
\LessonHead{Lectio IV}
\switchcolumn
\EN
\LessonHead{Lesson IV}
\switchcolumn*

\LA
\noindent Joánnes Bosco, humili loco natus apud Castum Novum in Astensibus atque sanctíssime, patre amisso, disciplína materna institútus, mira de se vel a prima ætate portendit. Ingenio enim mitis atque ad pietátem pronus, singulari se gerebat auctoritate inter æquales, quorum lites dirímere, faciles rixas compónere, turpia verba jocosque lascivos compéscere mature cœpit. Tum verbis eos jucundis ad se advocare, ludis preces insérere, quæ sacra elóquia in templo audivísset mirábili sermónis copia ac dulcédine referre, ad Pœniténtiæ et Eucharistiæ sacraménta rite suscipienda puerulos quam primum quamque creberrime inducere sategit. Oris quoque decor verecundus, morum suávitas, atque innocentíssimæ vitæ candor ómnium animos ad eum pertrahebant. Licet vero, familiáris rei angustia pressus labórum ærumnarúmque plenam adolescéntiam égerit, in id tamen unum hílaris ac Deo fidens contendit, ut sacerdotio augerétur.\par
\switchcolumn
\EN
\noindent John Bosco, born in the poor town of Castelnuovo d'Asti, and having lost his father at the age of two, was raised by his mother in a most saintly manner, and from his earliest years gave evidence of an extraordinary future. Docile and pious, he had a remarkable influence over those of his own age, whose fights he soon began to settle, and whose indecent words and improper jokes he stopped. Then he busied himself with drawing them to him by good stories, by including prayers in their games, by repeating in an attractive way the complete sermons he had heard in church, and with persuading them to receive the sacraments of Penance and of the Holy Eucharist without delay and frequently. His unassuming manner, his affability, and his innocence drew everyone to him. Although pressed with difficulties at home, and forced to work hard in his youth, he ardently desired with trust in God to become a priest.\par
\switchcolumn*

\LA
\begin{Resp}\Rsign Honéstum fecit illum Dóminus, et custodívit eum ab inimícis, et a seductóribus tutávit illum: \Star{} Et dedit illi claritátem ætérnam.\end{Resp}
\begin{Resp}\Vsign Justum dedúxit Dóminus per vias rectas, et osténdit illi regnum Dei.\end{Resp}
\begin{Resp}\Rsign Et dedit illi claritátem ætérnam.\end{Resp}
\switchcolumn
\EN
\begin{Resp}\Rsign The Lord made him honourable, and defended him from his enemies, and kept him safe from those that lay in wait for him; \Star{} And gave him perpetual glory.\end{Resp}
\begin{Resp}\Vsign He went down with him into the pit, and left him not in bonds.\end{Resp}
\begin{Resp}\Rsign And gave him perpetual glory.\end{Resp}
\switchcolumn*

\LA
\LessonHead{Lectio V}
\switchcolumn
\EN
\LessonHead{Lesson V}
\switchcolumn*

\LA
\noindent Voti tandem compos effectus, Cheriensem civitátem primum, ac dein Augustam Taurinórum pétiit, quo acrius, beato Josepho Cafasso magistro, et in sciéntia sanctórum profíceret et ad sacram morum doctrinam addiscéndam ánimum adjungeret. Ibi autem, cum voluntátis inclinatióne tum superno instinctu incitatus, suum in adolescéntulos ánimum convértit, ut prima iis christianæ sapiéntiæ traderet rudimenta. Quorum cum fieret in dies major numerus, sedem ad eos coadunándos stábilem ac firmam, haud sine cælésti afflatu, ásperis et diuturnis difficultátibus superatis, in illa urbis parte collocávit, quæ vulgo Valdocco appellátur, in eamque rem totus incúbuit. Paulo vero post, Vírgine Deípara auxiliante, quæ ei puérulo per visum in somnis futura innúerat, Joánnes Salesianórum Societátem institúere decrevit, cujus esset præsertim juveníles ánimas Christo lucrifácere; item novam familiam suscépit constituéndam sacrárum vírginum, quæ, ab Auxiliatrice Dei Matre nuncupátæ, adolescéntulas dirígerent in vias Dómini; quibus demum pium Cooperatórum cœtum adjecit ad Salesianórum ópera studio ac pietáte fovenda. Itaque brevi factum est, ut permagnam et christianæ et civíli societáti utilitátem afferret.\par
\switchcolumn
\EN
\noindent His wish was fulfilled, and he went first to Chieri, and then to Turin, where under the direction of Blessed Joseph Cafasso, he made rapid progress in the science of the Saints and in the learning of moral theology. There moved by divine grace and personal liking he began to take an interest in the youths, whom he taught the rudiments of the Christian religion. Their number increased day by day, and notwithstanding great and persistent difficulties, under divine inspiration he made a foundation for them in that section of the city called the Valdocco, on which he began to spend all his energy. Shortly after, with the help of the Blessed Virgin, who in a vision to him when a boy had revealed his future, John founded the Society of the Salesians, whose principal purpose was to be the saving of youthful souls for Christ. In like manner he founded a new family of nuns, who were called the daughters of St. Mary Auxiliatrix, and who would do for poor girls what the Salesians were doing for boys. To these he finally attached the Third Order of Salesian Cooperators, who by their piety and zeal were to assist in the educational work of the Salesians. And so in a short time he made great contributions both to the Church and to the State.\par
\switchcolumn*

\LA
\begin{Resp}\Rsign Amávit eum Dóminus, et ornávit eum: stolam glóriæ índuit eum, \Star{} Et ad portas paradísi coronávit eum.\end{Resp}
\begin{Resp}\Vsign Índuit eum Dóminus lorícam fídei, et ornávit eum.\end{Resp}
\begin{Resp}\Rsign Et ad portas paradísi coronávit eum.\end{Resp}
\switchcolumn
\EN
\begin{Resp}\Rsign The Lord loved him and beautified him; He clothed him with a robe of glory, \Star{} And crowned him at the gates of Paradise.\end{Resp}
\begin{Resp}\Vsign The Lord hath put on him the breast-plate of faith, and hath adorned him.\end{Resp}
\begin{Resp}\Rsign And crowned him at the gates of Paradise.\end{Resp}
\switchcolumn*

\LA
\LessonHead{Lectio VI}
\switchcolumn
\EN
\LessonHead{Lesson VI}
\switchcolumn*

\LA
\noindent Animárum enim studio flagrans, nulli pepércit labori nullique impensæ, ut festórum diérum asceteria pro adolescentulis, pupillórum hospitia, pusiónum operariórum scholas, ædes púeris alendis, instituéndis, templa Deo longe lateque per orbem excitaret. Simul Christi fidem in Subalpinis verbo et exemplo fovére, per totam Italiam optimos libros conficiendo, edendo, divulgando tutari, Evangélii præcones ad gentes infidéles sæpius mitténdo propagare non desiit. Simplex ac rectus homo Dei, ad omne opus bonum instructus, omnígenis virtútibus flóruit, quas incensíssimæ caritátis ardor alebat. Mente in Deum constanter erecta ac supernis charismátibus cumulátus, nullis sanctíssimus vir, nec minis terreri, nec labóribus fatigari, nec curis ópprimi, neque rebus adversis perturbari videbátur. Tria autem pietátis offícia suis maxime commendávit: ut quam sæpíssime ad sacram exhomologésim sacramque synáxim accederent, ut Mariam Auxiliatricem peramanter cólerent, ut Pontifici máximo ceu fílii addictíssimi obsequeréntur. Nec silentio prætereundum est eum, in difficilimis rerum adjunctis, præsto non semel Romano Pontifici adfuísse, ut mala ex légibus contra Ecclésiam eo témpore latis deriváta temperaret. Vitæ cursum, tot tantisque opéribus ac labóribus refertum, confecit pridie Kalendas Februarias anno salútis millesimo octingentésimo octogesimo octavo, ætátis septuagesimo tertio. Quem multis clarum miraculis Pius undecimus, Pontifex maximus, anno millesimo nongentésimo vigesimo nono Beatórum, quinquennio post, die solemni Paschæ, décimo nono exeunte sæculo a peracta humani generis Redemptióne, géntibus ex orbe universo in Urbem confluéntibus, Sanctórum ordínibus adserebat.\par
\switchcolumn
\EN
\noindent Filled with zeal for souls, he spared himself no labour and no expense to build recreational centres for the young, orphanages, schools for working children, schools and homes for the training of the young, and churches far and wide throughout the world. At the same time he did not stop spreading the Faith throughout the Subalpine country by word and by example, and throughout the whole of Italy, by writing and editing good books and by distributing the same, and in the foreign missions to which he sent numerous preachers. A simple and upright man, bent on every good work, he shone with all manner of virtue, which was fostered by his intense and ardent charity. With his mind always on God, and showered with heavenly gifts, this holy man of God was not disturbed by threats, nor tired by work, nor overwhelmed by care, nor upset by adversity. He recommended three works of piety to his followers: to receive as frequently as possible the sacraments of Penance and of Holy Eucharist, to cultivate a devotion to St. Mary Auxiliatrix, and to be the most loyal children of the Sovereign Pontiff. It should also be mentioned that John Bosco in very difficult circumstances went to the Pope more than once to console him in the evils coming from laws at that time passed against the Church. With a life of such accomplishments he died on the 31st day of January, 1888. Illustrious for his many miracles, the Supreme Pontiff, Pius XI, beatified him in 1929. Five years later, in the nineteenth centenary of the anniversary of our redemption, he was canonized among a vast gathering come to the Eternal City from every part of the world.\par
\switchcolumn*

\LA
\begin{Resp}\Rsign Iste homo perfécit ómnia quæ locútus est ei Deus, et dixit ad eum: Ingrédere in réquiem meam: \Star{} Quia te vidi justum coram me ex ómnibus géntibus.\end{Resp}
\begin{Resp}\Vsign Iste est, qui contémpsit vitam mundi, et pervénit ad cæléstia regna.\end{Resp}
\begin{Resp}\Rsign Quia te vidi justum coram me ex ómnibus géntibus.\end{Resp}
\begin{Resp}\Vsign Glória Patri, et Fílio, \Star{} et Spirítui Sancto.\end{Resp}
\begin{Resp}\Rsign Quia te vidi justum coram me ex ómnibus géntibus.\end{Resp}
\switchcolumn
\EN
\begin{Resp}\Rsign This is he which did according unto all that God commanded him; and God said unto him: Enter thou into My rest, \Star{} For thee have I seen righteous before Me among all people.\end{Resp}
\begin{Resp}\Vsign This is he which loved not his life in this world, and is come unto an everlasting kingdom.\end{Resp}
\begin{Resp}\Rsign For thee have I seen righteous before Me among all people.\end{Resp}
\begin{Resp}\Vsign Glory be to the Father, and to the Son, \Star{} and to the Holy Ghost.\end{Resp}
\begin{Resp}\Rsign For thee have I seen righteous before Me among all people.\end{Resp}
\switchcolumn*

\LA
\LessonHead{Lectio VII}
\switchcolumn
\EN
\LessonHead{Lesson VII}
\switchcolumn*

\LA
\LessonTitle{Léctio sancti Evangélii secúndum Matthǽum}
\switchcolumn
\EN
\LessonTitle{From the Holy Gospel according to Matthew}
\switchcolumn*

\LA
\Cite{Matt 18:1-5}
\switchcolumn
\EN
\Cite{Matt 18:1-5}
\switchcolumn*

\LA
\noindent In illo témpore: Accessérunt discípuli ad Jesum, dicéntes: Quis, putas, major est in regno cælórum? Et réliqua.\par
\switchcolumn
\EN
\noindent At that time: Came the disciples unto Jesus, saying: Who is the greatest in the kingdom of heaven? And so on, and that which followeth.\par
\switchcolumn*

\LA
\LessonTitle{Homilía sancti Joánnis Chrysostomi}
\switchcolumn
\EN
\LessonTitle{Homily of St. John Chrysostom}
\switchcolumn*

\LA
\Source{Homilia 60 in Matth. cap. 18}
\switchcolumn
\EN
\Source{In cap. 18 Matt., Hom. 60}
\switchcolumn*

\LA
\noindent Vide quot modis nos inducat Dóminus ad curam vel minórum fratrum? Ne ítaque dixeris: Ærarius est ille, aut calceórum sutor, agricola, insípiens, ut ideo despícias illum. Ne enim in illud mali íncidas, perpende quot modis te inducat, ut modeste agas et eórum curam geras. Puerum in médio státuit et ait: Efficiámini ut párvuli: et: Quicúmque suscéperit párvulum talem, me suscipit: et: Qui scandalizaverit, extrema patiétur. Si ergo Deus ita gaudet de párvulo qui repertus est, cur tu despícias eos quos Deus tam sollicite curat, cum oporteret ánimam ipsam tradere pro uno ex párvulis istis? Tanta quippe est Deo cura de anima, ut ne Fílio quidem suo pepércerit. Quare, óbsecro, primo dilúculo cum a domo exiérimus, hunc unum scopum habeámus et hanc præcipue sollicitúdinem, ut periclitántem eripiámus. Non loquor hic de sensibili periculo; hoc enim ne periculum quidem est; sed de periculo ánimæ, quod homínibus parat diabolus.\par
\switchcolumn
\EN
\noindent Do ye see in how many ways the Lord leadeth us on to the care of our lesser brethren? Therefore do not say: He is only a taxpayer, or a shoemaker, or a farmer, or that he is foolish, that ye may in that way look down on him. So that ye will not fall into such evil, reflect on how many ways he leadeth you to act humbly and take care of them. He placed a child in their midst and said: Become like little children; and: Whosoever receiveth one such little child for my sake, receiveth me; and: Who causeth one to sin, will suffer the severest penalties. If therefore God so rejoiceth over a little one that hath been found, why do ye despise those for whom God is so solicitous, when ye should trade your own life for one of those little ones? So great indeed is God's care for a soul that he did not spare his own Son. Wherefore, I entreat you, the first thing in the morning when we go out of the house, let us have in mind this purpose and this concern, to save someone in danger. I do not speak here of a visible danger, for this is not danger at all, but of danger to the soul, which the devil prepareth for men.\par
\switchcolumn*

\LA
\begin{Resp}\Rsign Iste est qui ante Deum magnas virtútes operátus est, et de omni corde suo laudávit Dóminum: \Star{} Ipse intercédat pro peccátis ómnium populórum.\end{Resp}
\begin{Resp}\Vsign Ecce homo sine queréla, verus Dei cultor, ábstinens se ab omni ópere malo, et pérmanens in innocéntia sua.\end{Resp}
\begin{Resp}\Rsign Ipse intercédat pro peccátis ómnium populórum.\end{Resp}
\switchcolumn
\EN
\begin{Resp}\Rsign This is he which wrought great wonders before God, and praised the Lord with all his heart. \Star{} May he pray for all people, that their sins may be forgiven unto them.\end{Resp}
\begin{Resp}\Vsign Behold a man without blame, a worshipper of God in truth, keeping himself clean from every evil work, and abiding still in his innocency.\end{Resp}
\begin{Resp}\Rsign May he pray for all people, that their sins may be forgiven unto them.\end{Resp}
\switchcolumn*

\LA
\LessonHead{Lectio VIII}
\switchcolumn
\EN
\LessonHead{Lesson VIII}
\switchcolumn*

\LA
\noindent Improbus, inquis, difficile tolerátur. Atque ideo debes illi amóre jungi, ut eum a vítio removeas, ut convértas et ad virtútem redúcas. At non obtemperat, inquis, neque consílium admittit. Unde hoc nosti? An exhortátus es et emendare studuísti? Hortátus sæpe sum, inquies. Quóties? Sæpius: semel et íterum. Idne sæpius vocas? Etiamsi per totam vitam id fecisses, nec deficere, nec desperáre oportebat. Non vides quómodo nos Deus semper hortátur per Prophétas, per Apóstolos, per Evangelistas? Quid ígitur? Num recte operamur? Num in ómnibus obtemperamus? Minime. Num ideo finem fecit admonendi?\par
\switchcolumn
\EN
\noindent Ye say it is hard to tolerate the bad. Ye should be joined to him in love, to lead him away from vice, to convert him and lead him back to virtue. But he doth not follow, say ye, nor taketh he advice. How know ye this? Have ye exhorted him and have ye tried to correct him? I have often exhorted, ye will say. How often? Once or twice. Do ye call that very often? Even if ye were to do it for your whole life, ye should neither stop nor despair. Do ye not see how God always encourageth us through the Prophets, through the Apostles, and through the Evangelists? What followeth? Do we act right? Do we obey in everything? Not at all. Should there be an end, then, to admonishing?\par
\switchcolumn*

\LA
\begin{Resp}\Rsign Sint lumbi vestri præcíncti, et lucérnæ ardéntes in mánibus vestris: \Star{} Et vos símiles homínibus exspectántibus dóminum suum, quando revertátur a núptiis.\end{Resp}
\begin{Resp}\Vsign Vigiláte ergo, quia nescítis qua hora Dóminus vester ventúrus sit.\end{Resp}
\begin{Resp}\Rsign Et vos símiles homínibus exspectántibus dóminum suum, quando revertátur a núptiis.\end{Resp}
\begin{Resp}\Vsign Glória Patri, et Fílio, \Star{} et Spirítui Sancto.\end{Resp}
\begin{Resp}\Rsign Et vos símiles homínibus exspectántibus dóminum suum, quando revertátur a núptiis.\end{Resp}
\switchcolumn
\EN
\begin{Resp}\Rsign Let your loins be girded about, and your lights burning; \Star{} And ye yourselves like unto men that wait for their lord, when he will return from the wedding.\end{Resp}
\begin{Resp}\Vsign Watch therefore, for ye know not what hour your Lord doth come.\end{Resp}
\begin{Resp}\Rsign And ye yourselves like unto men that wait for their lord, when he will return from the wedding.\end{Resp}
\begin{Resp}\Vsign Glory be to the Father, and to the Son, \Star{} and to the Holy Ghost.\end{Resp}
\begin{Resp}\Rsign And ye yourselves like unto men that wait for their lord, when he will return from the wedding.\end{Resp}
\switchcolumn*

\LA
\LessonHead{Lectio IX}
\switchcolumn
\EN
\LessonHead{Lesson IX}
\switchcolumn*

\LA
\noindent Nihil quippe tam pretiósum est quam ánima: quid enim prodest homini, si mundum univérsum lucrétur, ánimæ vero suæ detriméntum patiátur? Verum ómnia pervértit et dejécit pecuniárum amor, Deique timórem decussit, sicut tyrannus arcem sic ánimas occupans. Idcirco et filiórum et nostram neglígimus salútem. Hinc magna insipiéntia; hinc líberi servis vilióres fiunt. Ecquid de servis loquor? Mulum si quis habeat, multum curat ut agasónem illi optimum provideat, non ímprobum, non furacem, non temuléntum, non artis suæ imperítum: si autem fílio pædagogum dare opus sit, casu et sine delectu obvium quemque excipimus; etsi hac arte nulla sit major. Quid par illi arti, quæ dirigendæ juvenis menti et índoli incumbit? Qui tali instructus est facultate, plus diligéntiæ exhibeat oportet, quam quivis pictor aut statuarius.\par
\switchcolumn
\EN
\noindent There is indeed nothing as precious as a soul. For what doth it profit a man if he gain the whole world, and suffer the loss of his soul? But love of money hath destroyed and cast down everything, it hath thrust aside the fear of God, taking possession of the soul as a tyrant occupieth a fortress. And so we neglect our own salvation and that of our children. Great is that folly, and our children are worse than servants. Why do I speak of servants? If ye have a mule, ye take care to give it the best groom, one who is not worthless, nor a thief, nor a drunk, and one is not inexperienced in his work. If, however, it is necessary to have a preceptor for your son, ye take anyone ye may meet by chance and ye give no thought to selection, although there is no profession greater than this one. What is equal to that profession which is concerned with directing the soul and forming the mind and character of the young? He who hath such a task should shew more diligence than any painter or sculptor.\par
\switchcolumn*

\LA
\TeDeum{Te Deum}
\switchcolumn
\EN
\TeDeum{Te Deum}
\switchcolumn*

\end{paracol}

\SeasonTitle{Mensis Februarius}{February}

\addcontentsline{toc}{section}{Mensis Februarius\enspace\textit{February}}

\par\needspace{10\baselineskip}\addvspace{1.2em}{\centering\small\textsc{Die 1 Februarii} · \textsc{1 February}\par}
\Office{S. Ignatii Episcopi et Martyris}{St. Ignatius, Bishop and Martyr}{Duplex}
\addcontentsline{toc}{subsection}{Die 1 Februarii · S. Ignatii Episcopi et Martyris\enspace\textit{St. Ignatius, Bishop and Martyr}}
\begin{paracol}{2}
\LA
\RefLine{Lectiones I–III: de Scriptura occurrente.}
\switchcolumn
\EN
\RefLine{Lessons I–III: from the Scripture of the day.}
\switchcolumn*

\LA
\LessonHead{Lectio IV}
\switchcolumn
\EN
\LessonHead{Lesson IV}
\switchcolumn*

\LA
\LessonTitle{Ex libro sancti Hierónymi Presbýteri de Scriptóribus ecclesiásticis}
\switchcolumn
\EN

\switchcolumn*

\LA
\Source{Cap. 16}
\switchcolumn
\EN
\Source{From the book of saint Jerome Presbyter on the Ecclesiastical writers.}
\switchcolumn*

\LA
\noindent Ignátius, Antiochénæ ecclésiæ tértius post Petrum Apóstolum epíscopus, commovénte persecutiónem Trajáno, damnátus ad béstias, Romam vinctus míttitur. Cumque návigans Smyrnam venísset, ubi Polycárpus, audítor Joánnis, epíscopus erat, scripsit unam epístolam ad Ephésios, álteram ad Magnesiános, tértiam ad Trallénses, quartam ad Romános. Et inde egrédiens scripsit ad Philadélphios, et ad Smyrnǽos, et própriam ad Polycárpum, comméndans illi Antiochénsem ecclésiam; in qua et de Evangélio, quod nuper a me translátum est, super persóna Christi ponit testimónium.\par
\switchcolumn
\EN
\noindent Ignatius was the third Bishop of Antioch after the Apostle Peter. When Trajan stirred up his persecution, he was condemned to be devoured by wild beasts, and sent to Rome in chains. When on his journey thither he arrived at Smyrna, where Polycarp, the disciple of John, was Bishop, he wrote an Epistle to the Ephesians, another to the Magnesians, a third to the Trallians, and a fourth to the Romans and after leaving Smyrna, he addressed a further Epistle to the Philadelphians, and another to the Smyrnians, along with a private Epistle to Polycarp, to whose care he commended the Church of Antioch. In this last he quoteth a passage regarding the Person of Christ from the Gospel, which I have recently translated.\par
\switchcolumn*

\LA
\begin{Resp}\Rsign Honéstum fecit illum Dóminus, et custodívit eum ab inimícis, et a seductóribus tutávit illum: \Star{} Et dedit illi claritátem ætérnam.\end{Resp}
\begin{Resp}\Vsign Descendítque cum illo in fóveam, et in vínculis non derelíquit eum.\end{Resp}
\begin{Resp}\Rsign Et dedit illi claritátem ætérnam.\end{Resp}
\switchcolumn
\EN
\begin{Resp}\Rsign The Lord made him honourable, and defended him from his enemies, and kept him safe from those that lay in wait for him, \Star{} And gave him perpetual glory.\end{Resp}
\begin{Resp}\Vsign He went down with him into the pit, and left him not in bonds.\end{Resp}
\begin{Resp}\Rsign And gave him perpetual glory.\end{Resp}
\switchcolumn*

\LA
\LessonHead{Lectio V}
\switchcolumn
\EN
\LessonHead{Lesson V}
\switchcolumn*

\LA
\noindent Dignum autem vidétur, quia tanti viri fécimus mentiónem, et de epístola ejus, quam ad Romános scribit, pauca pónere: De Sýria usque ad Romam pugno ad béstias in mari et in terra, nocte diéque ligátus cum decem leopárdis, hoc est milítibus, qui me custódiunt; quibus et cum beneféceris, pejóres fiunt. Iníquitas autem eórum mea doctrína est; sed non idcírco justificátus sum. Utinam fruar béstiis, quæ mihi sunt præparátæ; quas et oro mihi velóces esse ad intéritum et ad supplícia, et állici ad comedéndum me, ne, sicut et aliórum Mártyrum, non áudiant corpus attíngere. Quod si veníre nolúerint, ego vim fáciam, ego me urgébo, ut dévorer. Ignóscite mihi, filíoli; quid mihi prosit, ego scio.\par
\switchcolumn
\EN
\noindent It is fitting that, as we have made mention of a man of so much importance, we should also note briefly the Epistle which he addressed to the Romans. I am on my way, saith he, from Syria to Rome, and am already fighting with beasts on sea and on land all the way. I may say I am chained day and night to ten leopards, for indeed the soldiers, who have charge of me, are no better. The more courteous I am to them, the worse they use me. But still their wickedness is good schooling for me, though I know that my mere sufferings cannot in themselves gain me justification. I earnestly wish for the beasts which are to devour me; at any rate, I pray they may put me out of pain quickly, and fly on me willingly, that I be not like some other Martyrs, whose bodies the animals have refused to touch. If I find that they will not come on, I will run at them as quick as I can, to make them devour me. Let me be, my little children I know what is good for me.\par
\switchcolumn*

\LA
\begin{Resp}\Rsign Desidérium ánimæ ejus tribuísti ei Dómine, \Star{} Et voluntáte labiórum ejus non fraudásti eum.\end{Resp}
\begin{Resp}\Vsign Quóniam prævenísti eum in benedictiónibus dulcédinis: posuísti in cápite ejus corónam de lápide pretióso.\end{Resp}
\begin{Resp}\Rsign Et voluntáte labiórum ejus non fraudásti eum.\end{Resp}
\switchcolumn
\EN
\begin{Resp}\Rsign Lord, Thou hast given him his heart's desire. \Star{} And hast not withholden the request of his lips.\end{Resp}
\begin{Resp}\Vsign For Thou hast prevented him with the blessings of sweetness: Thou hast set a crown of precious stones upon his head.\end{Resp}
\begin{Resp}\Rsign And hast not withholden the request of his lips.\end{Resp}
\switchcolumn*

\LA
\LessonHead{Lectio VI}
\switchcolumn
\EN
\LessonHead{Lesson VI}
\switchcolumn*

\LA
\noindent Nunc incípio Christi esse discípulus, nihil de his quæ vidéntur, desíderans, ut Jesum Christum invéniam. Ignis, crux, béstiæ, confráctio ossium, membrórum divísio, et totíus córporis contrítio, et tota torménta diáboli in me véniant; tantum ut Christo fruar. Cumque jam damnátus esset ad béstias, et ardóre patiéndi rugiéntes audíret leónes, ait: Fruméntum Christi sum; déntibus bestiárum molar, ut panis mundus invéniar. Passus est anno undécimo Trajáni. Relíquiæ córporis ejus Antiochíæ jacent extra portam Daphníticam, in cœmetério.\par
\switchcolumn
\EN
\noindent I feel now that I am beginning to be Christ's disciple; I desire none of those things which are seen, if so be I may find Christ Jesus. I care not that there come upon me fire, or cross, or wild beasts, or breaking of my bones, or sundering of my members, or destruction of my whole body, yea, or all the torments of the devil, if only so be I may win Christ. When he was brought condemned to the theatre, and heard the roaring of the beasts which were to devour him, he felt so strong an eagerness to suffer, that he cried out I am Christ's wheat, and so let the beasts' teeth be my mill, that I may be ground, and be found to make good bread. He suffered in the eleventh year of Trajan. What was left of his body lieth at Antioch, in the graveyard outside the gate which leadeth toward Daphne.\par
\switchcolumn*

\LA
\begin{Resp}\Rsign Stola jucunditátis índuit eum Dóminus: \Star{} Et corónam pulchritúdinis pósuit super caput ejus.\end{Resp}
\begin{Resp}\Vsign Cibávit illum Dóminus pane vitæ et intelléctus: et aqua sapiéntiæ salutáris potávit illum.\end{Resp}
\begin{Resp}\Rsign Et corónam pulchritúdinis pósuit super caput ejus.\end{Resp}
\begin{Resp}\Vsign Glória Patri, et Fílio, \Star{} et Spirítui Sancto.\end{Resp}
\begin{Resp}\Rsign Et corónam pulchritúdinis pósuit super caput ejus.\end{Resp}
\switchcolumn
\EN
\begin{Resp}\Rsign The Lord hath put on him a robe of honour, \Star{} And put about his head a crown of joy.\end{Resp}
\begin{Resp}\Vsign With the bread of life and understanding hath the Lord fed him, and given him the water of health and wisdom to drink.\end{Resp}
\begin{Resp}\Rsign And put about his head a crown of joy.\end{Resp}
\begin{Resp}\Vsign Glory be to the Father, and to the Son, \Star{} and to the Holy Ghost.\end{Resp}
\begin{Resp}\Rsign And put about his head a crown of joy.\end{Resp}
\switchcolumn*

\LA
\LessonHead{Lectio VII}
\switchcolumn
\EN
\LessonHead{Lesson VII}
\switchcolumn*

\LA
\LessonTitle{Léctio sancti Evangélii secúndum Joánnem}
\switchcolumn
\EN
\LessonTitle{From the Holy Gospel according to John}
\switchcolumn*

\LA
\Cite{Joannes 12:24-26}
\switchcolumn
\EN
\Cite{John 12:24-26}
\switchcolumn*

\LA
\noindent In illo témpore: Dixit Jesus discípulis suis: Amen, amen dico vobis, nisi granum fruménti cadens in terram, mortuum fúerit, ipsum solum manet. Et réliqua.\par
\switchcolumn
\EN
\noindent In that time Jesus said to his disciples: Amen, amen I say to you, unless the grain of wheat falling into the ground die, Itself remaineth alone. And so on.\par
\switchcolumn*

\LA
\LessonTitle{Homilía sancti Augustíni Epíscopi}
\switchcolumn
\EN
\LessonTitle{Homily by St. Augustine, Bishop (of Hippo.)}
\switchcolumn*

\LA
\Source{Tract. 51 in Joánn., sub med.}
\switchcolumn
\EN
\Source{Tract 51 on John.}
\switchcolumn*

\LA
\noindent Ipse Dóminus Jesus erat granum mortificándum et multiplicándum; mortificándum infidelitate Judæórum, multiplicándum fide populórum. Jam vero exhortans ad passiónis suæ sectanda vestígia, Qui amat, inquit, ánimam suam, perdet eam. Quod duobus modis intelligi potest. Qui amat, perdet; id est, si amas, perdes. Si cupis vitam tenére in Christo, noli mortem timére pro Christo. Item alio modo: Qui amat ánimam suam, perdet eam; noli amare, ne perdas; noli amare in hac vita, ne perdas in æterna vita.\par
\switchcolumn
\EN
\noindent The Lord Jesus was Himself a corn of wheat that was to die and bring forth much fruit; to die by the unbelief of the Jews, and to bring forth much fruit in the faith of the Gentiles. He, exhorting men to follow His steps, saith He that loveth his life shall lose it. Now, these words may be understood in two ways. First: he that loveth his life shall lose it, that is, If thou love life, thou wilt lose it; if thou wilt live for ever in Christ, refuse not to die for Christ. Or secondly: he that loveth his life shall lose it; love not then that which thou shalt lose; love not this present life, so that thou be thereby in jeopardy of losing life eternal.\par
\switchcolumn*

\LA
\begin{Resp}\Rsign Coróna áurea super caput ejus, \Star{} Expréssa signo sanctitátis, glória honóris, et opus fortitúdinis.\end{Resp}
\begin{Resp}\Vsign Quóniam prævenísti eum in benedictiónibus dulcédinis, posuísti in cápite ejus corónam de lápide pretióso.\end{Resp}
\begin{Resp}\Rsign Expréssa signo sanctitátis, glória honóris, et opus fortitúdinis.\end{Resp}
\switchcolumn
\EN
\begin{Resp}\Rsign A crown of gold upon his head, \Star{} Wherein is engraved Holiness, an ornament of honour, a costly work.\end{Resp}
\begin{Resp}\Vsign For Thou hast prevented him with the blessings of sweetness, Thou hast set a crown of precious stones upon his head.\end{Resp}
\begin{Resp}\Rsign Wherein is engraved Holiness, an ornament of honour, a costly work.\end{Resp}
\switchcolumn*

\LA
\LessonHead{Lectio VIII}
\switchcolumn
\EN
\LessonHead{Lesson VIII}
\switchcolumn*

\LA
\noindent Hoc autem, quod posterius dixi, magis habére vidétur evangélicus sensus; séquitur enim: Et qui odit ánimam suam in hoc mundo, in vitam ætérnam custódit eam. Ergo, quod supra dictum est, Qui amat, subintellígitur in hoc mundo, ipse útique perdet; qui autem odit, útique in hoc mundo, in vitam ætérnam ipse custódit eam. Magna et mira senténtia, quemádmodum sit hóminis in ánimam suam amor ut pereat, ódium ne pereat. Si male amáveris, tunc odísti; si bene óderis, tunc amásti. Felices, qui odérunt custodiéndo, ne perdant amándo.\par
\switchcolumn
\EN
\noindent What this second interpretation is the meaning of the Gospel, appeareth most probably from the words which follow And he that hateth his life in this world, shall keep it unto life eternal. From which we may suppose the sense of the first words to be He that loveth his life in this world shall lose it unto life eternal. This is a great and marvellous saying, showing how a man may so love life as to lose life, and so hate life as to keep life. If thou love it too well, then dost thou hate it if thou hate it with an holy hatred, then dost thou love it. Blessed are they that, lest they should so love it as to lose it, so hate it as to keep it.\par
\switchcolumn*

\LA
\begin{Resp}\Rsign Hic est vere Martyr, qui pro Christi nómine sánguinem suum fudit: \Star{} Qui minas júdicum non tímuit, nec terrénæ dignitátis glóriam quæsívit, sed ad cæléstia regna pervénit.\end{Resp}
\begin{Resp}\Vsign Justum dedúxit Dóminus per vias rectas, et osténdit illi regnum Dei.\end{Resp}
\begin{Resp}\Rsign Qui minas júdicum non tímuit, nec terrénæ dignitátis glóriam quæsívit, sed ad cæléstia regna pervénit.\end{Resp}
\begin{Resp}\Vsign Glória Patri, et Fílio, \Star{} et Spirítui Sancto.\end{Resp}
\begin{Resp}\Rsign Qui minas júdicum non tímuit, nec terrénæ dignitátis glóriam quæsívit, sed ad cæléstia regna pervénit.\end{Resp}
\switchcolumn
\EN
\begin{Resp}\Rsign This is a Martyr indeed, who shed his blood for Christ's Name's sake \Star{} Who feared not for the threats of judges, nor sought to be great with the glory of this world, but pressed on unto the kingdom of heaven.\end{Resp}
\begin{Resp}\Vsign The Lord guided the righteous in right paths, and showed him the kingdom of God.\end{Resp}
\begin{Resp}\Rsign Who feared not for the threats of judges, nor sought to be great with the glory of this world, but pressed on unto the kingdom of heaven.\end{Resp}
\begin{Resp}\Vsign Glory be to the Father, and to the Son, \Star{} and to the Holy Ghost.\end{Resp}
\begin{Resp}\Rsign Who feared not for the threats of judges, nor sought to be great with the glory of this world, but pressed on unto the kingdom of heaven.\end{Resp}
\switchcolumn*

\LA
\LessonHead{Lectio IX}
\switchcolumn
\EN
\LessonHead{Lesson IX}
\switchcolumn*

\LA
\noindent Sed vide ne tibi subrépat ut teípsum velis interímere, sic intelligendo, quod debes odisse in hoc mundo ánimam tuam. Hinc enim quidam malígni atque perversi hómines, et in seipsis crudeliores et sceleratiores homicídæ flammis se donant, aquis se præfócant, præcipítio se collídunt et pereunt. Hoc Christus non docuit; immo étiam diabolo præcipitium suggerénti respóndit: Redi retro, sátana: scriptum est, Non tentábis Dóminum, Deum tuum. Petro autem dixit, signíficans qua morte clarificatúrus erat Deum: cum esses júnior, cingébas te, et ambulabas quo volébas; cum autem senúeris, alter te cinget et feret quo tu non vis. Ubi satis expressit, non a seipso, sed ab alio debére occídi, qui vestígia sequitur Christi.\par
\switchcolumn
\EN
\noindent Beware lest thou take these words He that hateth his life in this world shall keep it unto life eternal as some do, for an approval of suicide. Some evil and perverse men, bloody and guilty murderers of themselves, do indeed throw themselves into the fire, drown themselves in water, and cast themselves down precipices, and so perish. This is not the teaching of Christ, Who, when the devil would have Him cast Himself down from an high place, answered Get thee behind Me, Satan. It is written, Thou shalt not tempt the Lord thy God. (Matth. iv. 5-7.) Who also said to Peter, signifying by what death he should glorify God When thou wast young thougirdedst thyself and walkedst whither thou wouldest; but when thou shalt be old, another shall gird thee, and carry thee whither thou wouldest not. (John xxi. 18.) From which it is evident that he that would follow Christ's footsteps, must be slain, not by himself, but by another.\par
\switchcolumn*

\LA
\TeDeum{Te Deum}
\switchcolumn
\EN
\TeDeum{Te Deum}
\switchcolumn*

\end{paracol}

\par\needspace{10\baselineskip}\addvspace{1.2em}{\centering\small\textsc{Die 2 Februarii} · \textsc{2 February}\par}
\Office{In Purificatione Beatæ Mariæ Virginis}{In Purificatione Beatæ Mariæ Virginis}{Duplex II classis}
\addcontentsline{toc}{subsection}{Die 2 Februarii · In Purificatione Beatæ Mariæ Virginis\enspace\textit{In Purificatione Beatæ Mariæ Virginis}}
\begin{paracol}{2}
\LA
\LessonHead{Lectio I}
\switchcolumn
\EN
\LessonHead{Lesson I}
\switchcolumn*

\LA
\LessonTitle{De libro Exodi}
\switchcolumn
\EN
\LessonTitle{From the book of Exodus}
\switchcolumn*

\LA
\Cite{Exod 13:1-3; 13:11-13}
\switchcolumn
\EN
\Cite{Exod 13:1-3; 13:11-13}
\switchcolumn*

\LA
\noindent Locútus est Dóminus ad Móysen, dicens: Sanctífica mihi omne primogénitum quod áperit vulvam in fíliis Israël, tam de homínibus quam de juméntis; mea sunt enim ómnia. Et ait Móyses ad pópulum: Cum introdúxerit te Dóminus in terram Chananǽi, sicut jurávit tibi et pátribus tuis, et déderit tibi eam; separábis omne quod áperit vulvam Dómino, et quod primitívum est in pecóribus tuis; quidquid habúeris masculíni sexus, consecrábis Dómino. Primogénitum ásini mutábis ove; quod, si non redémeris, interfícies. Omne autem primogénitum hóminis de fíliis tuis, prétio rédimes.\par
\switchcolumn
\EN
\noindent And the Lord spoke to Moses, saying: Sanctify unto me every firstborn that openeth the womb among the children of Israel, as well of men as of beasts: for they are all mine. And Moses said to the people: When the Lord shall have brought thee into the land of the Chanaanite, as he swore to thee and thy fathers, and shall give it thee: Thou shalt set apart all that openeth the womb for the Lord, and all that is first brought forth of thy cattle: whatsoever thou shalt have of the male sex, thou shalt consecrate to the Lord. The firstborn of an ass thou shalt change for a sheep: and if thou do not redeem it, thou shalt kill it. And every firstborn of men thou shalt redeem with a price.\par
\switchcolumn*

\LA
\begin{Resp}\Rsign Adórna thálamum tuum, Sion, et súscipe Regem Christum: \Star{} Quem virgo concépit, virgo péperit, virgo post partum, quem génuit, adorávit.\end{Resp}
\begin{Resp}\Vsign Accípiens Símeon púerum in mánibus, grátias agens benedíxit Dóminum.\end{Resp}
\begin{Resp}\Rsign Quem virgo concépit, virgo péperit, virgo post partum, quem génuit, adorávit.\end{Resp}
\switchcolumn
\EN
\begin{Resp}\Rsign Make ready thy chamber, O Zion, to receive Christ thy King \Star{} Even that Child Whom a Virgin bore, and remained a Virgin as before the fruit of her womb and the God of her soul.\end{Resp}
\begin{Resp}\Vsign Simeon took the Child up in his arms, and blessed God.\end{Resp}
\begin{Resp}\Rsign Even that Child Whom a Virgin bore, and remained a Virgin as before the fruit of her womb and the God of her soul.\end{Resp}
\switchcolumn*

\LA
\LessonHead{Lectio II}
\switchcolumn
\EN
\LessonHead{Lesson II}
\switchcolumn*

\LA
\LessonTitle{De libro Levítici}
\switchcolumn
\EN
\LessonTitle{From the Book of Leviticus}
\switchcolumn*

\LA
\Cite{Lev 12:1-5}
\switchcolumn
\EN
\Cite{Lev 12:1-5}
\switchcolumn*

\LA
\noindent Locútus est Dóminus ad Móysen, dicens: Lóquere fíliis Israël, et dices ad eos: Múlier, si suscépto sémine pepérerit másculum, immúnda erit septem diébus juxta dies separatiónis ménstruæ. Et die octávo circumcidétur infántulus: ipsa vero trigínta tribus diébus manébit in sánguine purificatiónis suæ. Omne sanctum non tanget, nec ingrediétur in Sanctuárium donec impleántur dies purificatiónis suæ. Sin autem féminam pepérerit, immúnda erit duábus hebdomádibus juxta ritum fluxus ménstrui, et sexagínta sex diébus manébit in sánguine purificatiónis suæ.\par
\switchcolumn
\EN
\noindent And the Lord spoke to Moses, saying: Speak to the children of Israel, and thou shalt say to them: If a woman having received seed shall bear a man child, she shall be unclean seven days, according to the days of the separation of her flowers. And on the eighth day the infant shall be circumcised: But she shall remain three and thirty days in the blood of her purification. She shall touch no holy thing, neither shall she enter into the sanctuary, until the days of her purification be fulfilled. But if she shall bear a maid child, she shall be unclean two weeks, according to the custom of her monthly courses, and she shall remain in the blood of her purification sixty-six days.\par
\switchcolumn*

\LA
\begin{Resp}\Rsign Postquam impléti sunt dies purgatiónis Maríæ secúndum legem Móysi, tulérunt Jesum in Jerúsalem, ut sísterent eum Dómino, \Star{} Sicut scriptum est in lege Dómini: Quia omne masculínum adapériens vulvam, sanctum Dómino vocábitur.\end{Resp}
\begin{Resp}\Vsign Obtulérunt pro eo Dómino par túrturum aut duos pullos columbárum.\end{Resp}
\begin{Resp}\Rsign Sicut scriptum est in lege Dómini: Quia omne masculínum adapériens vulvam, sanctum Dómino vocábitur.\end{Resp}
\switchcolumn
\EN
\begin{Resp}\Rsign When the days of Mary's purification according to the law of Moses were accomplished, they brought Jesus to Jerusalem to present Him to the Lord; \Star{} As it is written in the law of the Lord: Every male that openeth the womb shall be called holy unto the Lord.\end{Resp}
\begin{Resp}\Vsign They offered for Him unto the Lord a pair of turtle-doves or two young pigeons.\end{Resp}
\begin{Resp}\Rsign As it is written in the law of the Lord: Every male that openeth the womb shall be called holy unto the Lord.\end{Resp}
\switchcolumn*

\LA
\LessonHead{Lectio III}
\switchcolumn
\EN
\LessonHead{Lesson III}
\switchcolumn*

\LA
\Cite{Lev 12:6-8}
\switchcolumn
\EN
\Cite{Lev 12:6-8}
\switchcolumn*

\LA
\noindent Cumque expléti fúerint dies purificatiónis suæ, pro fílio, sive pro fília, déferet agnum annículum in holocáustum, et pullum colúmbæ sive túrturem pro peccáto, ad óstium tabernáculi testimónii, et tradet sacerdóti, qui ófferet illa coram Dómino; et orábit pro ea, et sic mundábitur a proflúvio sánguinis sui. Ista est lex pariéntis másculum aut féminam. Quod si non invénerit manus ejus, nec potúerit offérre agnum, sumet duos túrtures vel duos pullos columbárum, unum in holocáustum, et álterum pro peccáto; orabítque pro ea sacérdos, et sic mundábitur.\par
\switchcolumn
\EN
\noindent And when the days of her purification are expired, for a son, or for a daughter, she shall bring to the door of the tabernacle of the testimony, a lamb of a year old for a holocaust, and a young pigeon or a turtle for sin, and shall deliver them to the priest: Who shall offer them before the Lord, and shall pray for her, and so she shall be cleansed from the issue of her blood. This is the law for her that beareth a man child or a maid child. And if her hand find not sufficiency, and she is not able to offer a lamb, she shall take two turtles, or two young pigeons, one for a holocaust, and another for sin: and the priest shall pray for her, and so she shall be cleansed.\par
\switchcolumn*

\LA
\begin{Resp}\Rsign Obtulérunt pro eo Dómino par túrturum aut duos pullos columbárum. \Star{} Sicut scriptum est in lege Dómini.\end{Resp}
\begin{Resp}\Vsign Postquam autem impléti sunt dies purgatiónis Maríæ secúndum legem Móysi, tulérunt illum in Jerúsalem, ut sísterent eum Dómino.\end{Resp}
\begin{Resp}\Rsign Sicut scriptum est in lege Dómini.\end{Resp}
\begin{Resp}\Vsign Glória Patri, et Fílio, \Star{} et Spirítui Sancto.\end{Resp}
\begin{Resp}\Rsign Sicut scriptum est in lege Dómini.\end{Resp}
\switchcolumn
\EN
\begin{Resp}\Rsign They offered for Him unto the Lord a pair of turtle-doves, or two young pigeons \Star{} As it is written in the law of the Lord.\end{Resp}
\begin{Resp}\Vsign And when the days of Mary's purification according to the law of Moses were accomplished, they brought Him to Jerusalem, to present Him to the Lord.\end{Resp}
\begin{Resp}\Rsign As it is written in the law of the Lord.\end{Resp}
\begin{Resp}\Vsign Glory be to the Father, and to the Son, \Star{} and to the Holy Ghost.\end{Resp}
\begin{Resp}\Rsign As it is written in the law of the Lord\end{Resp}
\switchcolumn*

\LA
\LessonHead{Lectio IV}
\switchcolumn
\EN
\LessonHead{Lesson IV}
\switchcolumn*

\LA
\LessonTitle{Sermo sancti Augustíni Epíscopi}
\switchcolumn
\EN
\LessonTitle{From the Sermons of St. Augustine, Bishop (of Hippo.)}
\switchcolumn*

\LA
\Source{Sermo 13 de Tempore, post initium}
\switchcolumn
\EN
\Source{13th on the Season}
\switchcolumn*

\LA
\noindent Sic olim prædíctum est: Mater Sion dicit: Homo et homo factus est in ea; et ipse fundávit eam Altíssimus. O omnipoténtia nascéntis! o magnificéntia de cælo ad terram descendéntis! Adhuc in útero portabátur, et ex útero matris a Joánne Baptísta salutabátur. In templo præsentabátur, et a Simeóne sene famóso, annóso, probáto, coronáto agnoscebátur. Tunc cognóvit, tunc adorávit, tunc dixit: Nunc, Dómine, dimíttis servum tuum in pace: quia vidérunt óculi mei salutáre tuum.\par
\switchcolumn
\EN
\noindent In old time it was written: And of Zion shall it not be said This and that man was born in her, and the Highest Himself shall establish her? Blessed be the omnipotence of Him That was born! Blessed the glory of Him That came from heaven to earth! While yet He was borne in His Mother's womb, He was saluted by John the Baptist; He was presented in the temple, and recognised by that famous, ancient, and glorious worthy, the old man Simeon. As soon as he knew Him he worshipped Him, and said Lord, now lettest Thou thy servant depart in peace for mine eyes have seen thy Salvation.\par
\switchcolumn*

\LA
\begin{Resp}\Rsign Símeon justus et timorátus exspectábat redemptiónem Israël, \Star{} Et Spíritus Sanctus erat in eo.\end{Resp}
\begin{Resp}\Vsign Respónsum accépit Símeon a Spíritu Sancto, non visúrum se mortem, nisi vidéret Christum Dómini.\end{Resp}
\begin{Resp}\Rsign Et Spíritus Sanctus erat in eo.\end{Resp}
\switchcolumn
\EN
\begin{Resp}\Rsign Simeon was just and devout, waiting for the redemption of Israel \Star{} And the Holy Ghost was upon him.\end{Resp}
\begin{Resp}\Vsign It was revealed unto him by the Holy Ghost, that he should not see death, before he had seen the Lord's Christ.\end{Resp}
\begin{Resp}\Rsign And the Holy Ghost was upon him.\end{Resp}
\switchcolumn*

\LA
\LessonHead{Lectio V}
\switchcolumn
\EN
\LessonHead{Lesson V}
\switchcolumn*

\LA
\noindent Differebátur exíre de sǽculo, ut vidéret natum, per quem cónditum est sǽculum. Agnóvit Infántem senex, factus est in Púero puer. Innovátus in ætáte, qui plenus erat pietáte. Símeon senex ferébat Christum infántem, Christus regébat Simeónis senectútem. Dictum ei fúerat a Dómino, quod non gustáret mortem, nisi vidéret Christum Dómini natum. Natus est Christus, et implétum est desidérium senis in mundi ipsíus senectúte. Ipse ad senem hóminem venit, qui mundum inveterátum invénit.\par
\switchcolumn
\EN
\noindent He lingered in the world to see the birth of Him Who made the world. The old man knew the Child, and in that Child became a child himself, for in the love wherewith he regarded the Father of all, he felt his own years to be but as of yesterday. The old man Simeon bore the new-born Christ, and all the while, Christ was the old man's Lord. It had been told him by the Lord that he should not taste of death before he had seen the birth of the Lord's Christ. Now Christ is born, and all the old man's wishes on earth are fulfilled. He That came to a decrepit world came to an old man.\par
\switchcolumn*

\LA
\begin{Resp}\Rsign Respónsum accépit Símeon a Spíritu Sancto, non visúrum se mortem, nisi vidéret Christum Dómini: \Star{} Et benedíxit Deum, et dixit: Nunc dimíttis servum tuum in pace, quia vidérunt óculi mei salutáre tuum, Dómine.\end{Resp}
\begin{Resp}\Vsign Cum indúcerent púerum Jesum paréntes ejus, ut fácerent secúndum consuetúdinem legis pro eo, ipse accépit eum in ulnas suas.\end{Resp}
\begin{Resp}\Rsign Et benedíxit Deum, et dixit: Nunc dimíttis servum tuum in pace, quia vidérunt óculi mei salutáre tuum, Dómine.\end{Resp}
\switchcolumn
\EN
\begin{Resp}\Rsign It was revealed unto Simeon by the Holy Ghost that he should not see death before he had seen the Lord's Christ; \Star{} And he blessed God, and said: Lord, now lettest Thou thy servant depart in peace, for mine eyes have seen thy salvation.\end{Resp}
\begin{Resp}\Vsign When His parents brought in the Child Jesus, to do for Him after the custom of the law, then took he Him up in his arms.\end{Resp}
\begin{Resp}\Rsign And blessed God, and said Lord, now lettest Thou thy servant depart in peace, for mine eyes have seen thy salvation.\end{Resp}
\switchcolumn*

\LA
\LessonHead{Lectio VI}
\switchcolumn
\EN
\LessonHead{Lesson VI}
\switchcolumn*

\LA
\noindent In isto quidem sǽculo diu esse nolébat, et Christum in hoc sǽculo vidére cupiébat, cantans cum prophéta et dicens: Osténde nobis, Dómine, misericórdiam tuam, et salutáre tuum da nobis. Dénique, ut novéritis ita esse istíus lætítiam, conclúsit dicens: Nunc dimíttis servum tuum in pace: quia vidérunt óculi mei salutáre tuum. Prophétæ cecinérunt Conditórem cæli et terræ in terra cum homínibus futúrum; Angelus nuntiávit Creatórem carnis et spíritus in carne ventúrum; salutávit Joánnes ex útero in útero Salvatórem; Símeon senex Deum agnóvit infántem.\par
\switchcolumn
\EN
\noindent He wished not to remain long in the world, but he longed to see Christ in the world, singing with the Prophet, and saying: Show us thy mercy, O Lord, and grant us thy salvation. (Ps. lxxxiv. 8.) And now at last, that ye may know that the cause of his joy was that this prayer was granted, he saith: Now lettest Thou thy servant depart in peace, for mine eyes have seen thy salvation. The Prophets have sung that the Maker of heaven and earth would converse on earth with men, an angel hath declared that the Creator of flesh and spirit would come in the flesh, the unborn John, yet in the womb, hath saluted the unborn Saviour yet in the womb. The old man Simeon hath seen God, a little Child.\par
\switchcolumn*

\LA
\begin{Resp}\Rsign Cum indúcerent púerum Jesum paréntes ejus in templum, ut fácerent secúndum consuetúdinem legis pro eo, accépit eum Símeon in ulnas suas, et benedíxit Deum, dicens: \Star{} Nunc dimíttis, Dómine, servum tuum in pace.\end{Resp}
\begin{Resp}\Vsign Suscípiens Símeon Púerum in mánibus, exclamávit, dicens.\end{Resp}
\begin{Resp}\Rsign Nunc dimíttis, Dómine, servum tuum in pace.\end{Resp}
\begin{Resp}\Vsign Glória Patri, et Fílio, \Star{} et Spirítui Sancto.\end{Resp}
\begin{Resp}\Rsign Nunc dimíttis, Dómine, servum tuum in pace.\end{Resp}
\switchcolumn
\EN
\begin{Resp}\Rsign When His parents brought the Child Jesus into the temple, to do for Him after the custom of the law, Simeon took Him up in his arms, and blessed God, and said: \Star{} Lord, now lettest Thou thy servant depart in peace.\end{Resp}
\begin{Resp}\Vsign Simeon took up the Child in his arms, and cried out, and said:\end{Resp}
\begin{Resp}\Rsign Lord, now lettest Thou thy servant depart in peace.\end{Resp}
\begin{Resp}\Vsign Glory be to the Father, and to the Son, \Star{} and to the Holy Ghost.\end{Resp}
\begin{Resp}\Rsign Lord, now lettest Thou thy servant depart in peace.\end{Resp}
\switchcolumn*

\LA
\LessonHead{Lectio VII}
\switchcolumn
\EN
\LessonHead{Lesson VII}
\switchcolumn*

\LA
\LessonTitle{Léctio sancti Evangélii secúndum Lucam}
\switchcolumn
\EN
\LessonTitle{From the Holy Gospel according to Luke}
\switchcolumn*

\LA
\Cite{Luc 2:22-32}
\switchcolumn
\EN
\Cite{Luke 2:22-32}
\switchcolumn*

\LA
\noindent In illo témpore: Postquam impléti sunt dies purgatiónis Maríæ secúndum legem Móysi, tulérunt Jesum in Jerúsalem, ut sísterent eum Dómino, sicut scriptum est in lege Dómini. Et réliqua.\par
\switchcolumn
\EN
\noindent In that time, after the days of her purification according to the law of Moses were accomplished, they carried him to Jerusalem to present him to the Lord, as it is written in the law of the Lord. And so on.\par
\switchcolumn*

\LA
\LessonTitle{Homilía sancti Ambrósii Epíscopi}
\switchcolumn
\EN
\LessonTitle{Homily on this passage by St. Ambrose, Bishop (of Milan.)}
\switchcolumn*

\LA
\Source{Liber 2 Comment. in Lucæ cap. 2, post initium}
\switchcolumn
\EN
\Source{Bk. 2, Comm. on Luke II}
\switchcolumn*

\LA
\noindent Et ecce homo erat in Jerúsalem, cui nomen Símeon, et homo iste justus et timorátus, exspéctans consolatiónem Israël. Non solum ab Angelis, et prophétis, et a pastóribus, sed étiam a senióribus et justis generátio Dómini accépit testimónium. Omnis ætas, et utérque sexus, eventorúmque mirácula fidem ástruunt. Virgo génerat, stérilis parit, mutus lóquitur, Elísabeth prophétat, Magus adórat, útero clausus exsúltat, vídua confitétur, justus exspéctat.\par
\switchcolumn
\EN
\noindent And, behold, there was a man in Jerusalem, whose name was Simeon, and the same man was just and devout, waiting for the consolation of Israel. The birth of the Lord is attested not only by Angels and Prophets, and shepherds, but also by elders and just men. Every age, and both sexes, as well as the miracles of the events themselves, are here to strengthen our faith. A virgin conceiveth, a barren woman beareth, a dumb man speaketh, Elizabeth prophesieth, the wise man worshippeth, the unborn child leapeth, the widow praiseth, and the just man waiteth.\par
\switchcolumn*

\LA
\begin{Resp}\Rsign Suscípiens Jesum in ulnas suas Símeon, exclamávit, et dixit: \Star{} Tu es vere lumen ad illuminatiónem géntium, et glóriam plebis tuæ Israël.\end{Resp}
\begin{Resp}\Vsign Cum indúcerent púerum Jesum paréntes ejus, et ipse accépit eum in ulnas suas, et benedíxit Deum, et dixit.\end{Resp}
\begin{Resp}\Rsign Tu es vere lumen ad illuminatiónem géntium, et glóriam plebis tuæ Israël.\end{Resp}
\switchcolumn
\EN
\begin{Resp}\Rsign Simeon took Jesus up in his arms, and cried out, and said: \Star{} Verily Thou art a light to lighten the Gentiles, and the Glory of thy people Israel.\end{Resp}
\begin{Resp}\Vsign When His parents brought in the Child Jesus, then took he Him up in his arms, and blessed God, and said:\end{Resp}
\begin{Resp}\Rsign Verily Thou art a light to lighten the Gentiles, and the Glory of thy people Israel.\end{Resp}
\switchcolumn*

\LA
\LessonHead{Lectio VIII}
\switchcolumn
\EN
\LessonHead{Lesson VIII}
\switchcolumn*

\LA
\noindent Et bene justus, qui non suam, sed pópuli grátiam requirébat, cúpiens ipse corpóreæ vínculis fragilitátis exsólvi, sed exspéctans vidére promíssum: sciébat enim quia beáti óculi qui eum vidérent. Et ipse accépit eum in ulnas suas, et benedíxit Deum, et dixit: Nunc dimítte servum tuum, Dómine, secúndum verbum tuum in pace. Vide justum, velut corpóreæ cárcere molis inclúsum, velle dissólvi, ut incípiat esse cum Christo. Dissólvi enim, et cum Christo esse, multo mélius est.\par
\switchcolumn
\EN
\noindent Well is he called just, who looked not for favour for himself, but for consolation for his people. He desired to be set free from the bondage of this frail body, but he waited to see the Promised One for he knew that blessed are the eyes that see Him. Then took he Him up in his arms, and blessed God, and said: Lord, now lettest Thou thy servant depart in peace, according to thy word. Behold a just man, confined in the weary prison of the body, desiring to be dissolved and to begin to be with Christ. For to be dissolved and to be with Christ is much better. (Phil. i. 23.)\par
\switchcolumn*

\LA
\begin{Resp}\Rsign Senex Púerum portábat, Puer autem senem regébat: \Star{} Quem virgo concépit, virgo péperit, virgo post partum, quem génuit, adorávit.\end{Resp}
\begin{Resp}\Vsign Accípiens Símeon Púerum in mánibus, grátias agens benedíxit Dóminum.\end{Resp}
\begin{Resp}\Rsign Quem virgo concépit, virgo péperit, virgo post partum, quem génuit, adorávit.\end{Resp}
\begin{Resp}\Vsign Glória Patri, et Fílio, \Star{} et Spirítui Sancto.\end{Resp}
\begin{Resp}\Rsign Quem virgo concépit, virgo péperit, virgo post partum, quem génuit, adorávit.\end{Resp}
\switchcolumn
\EN
\begin{Resp}\Rsign The old man bore the Child, but the Child was the old man's King; \Star{} Even that Child Whom a virgin bore, and remained a virgin as before; and when that virgin had brought Him into the world, she fell down and worshipped Him.\end{Resp}
\begin{Resp}\Vsign Simeon took the Child up in his arms, and gave thanks, and blessed the Lord.\end{Resp}
\begin{Resp}\Rsign Even that Child Whom a virgin bore, and remained a virgin as before; and when that virgin had brought Him into the world, she fell down and worshipped Him.\end{Resp}
\begin{Resp}\Vsign Glory be to the Father, and to the Son, \Star{} and to the Holy Ghost.\end{Resp}
\begin{Resp}\Rsign Even that Child Whom a virgin bore, and remained a virgin as before; and when that virgin had brought Him into the world, she fell down and worshipped Him.\end{Resp}
\switchcolumn*

\LA
\LessonHead{Lectio IX}
\switchcolumn
\EN
\LessonHead{Lesson IX}
\switchcolumn*

\LA
\noindent Sed qui vult dimítti, véniat in templum, véniat in Jerúsalem, exspéctet Christum Dómini, accípiat in mánibus Verbum Dei, complectátur opéribus velut quibúsdam suæ fídei brácchiis: tunc dimittétur, ut non vídeat mortem, qui víderit Vitam. Vides úberem in omnes grátiam, Dómini generatióne diffúsam, et prophetíam incrédulis negátam esse, non justis. Ecce et Símeon prophétat, in ruínam et resurrectiónem plurimórum venísse Dóminum Jesum Christum, ut justórum iniquorúmque mérita discérnat; et pro nostrórum qualitáte factórum, judex verus et justus aut supplícia decérnat aut prǽmia.\par
\switchcolumn
\EN
\noindent Whosoever will be dissolved and be with Christ, let him come into the Temple, let him come to Jerusalem, let him wait for the Lord's Christ, let him take hold on the Word of God, let him embrace it with good works, as it were with arms of faith and then let him depart in peace, for he shall not see death, who hath seen life. Behold how the Lord's Birth doth overflow with abounding grace for all, and prophecy is not denied to the just, but to the unbelieving. Behold, Simeon prophesieth that the Lord Jesus Christ is come for the fall and rising again of many; yea, He shall separate the just from the unjust by their deserts, and according as our work shall be, so shall the true and righteous Judge command us to be punished or rewarded.\par
\switchcolumn*

\LA
\TeDeum{Te Deum}
\switchcolumn
\EN
\TeDeum{Te Deum}
\switchcolumn*

\end{paracol}

\par\needspace{10\baselineskip}\addvspace{1.2em}{\centering\small\textsc{Die 3 Februarii} · \textsc{3 February}\par}
\Office{S. Blasii Episcopi et Martyris}{St. Blase, Bishop and Martyr}{Simplex}
\addcontentsline{toc}{subsection}{Die 3 Februarii · S. Blasii Episcopi et Martyris\enspace\textit{St. Blase, Bishop and Martyr}}
\begin{paracol}{2}
\LA
\RefLine{Lectiones I–II: de Scriptura occurrente.}
\switchcolumn
\EN
\RefLine{Lessons I–II: from the Scripture of the day.}
\switchcolumn*

\LA
\LessonHead{Lectio III (pro commemoratione)}
\switchcolumn
\EN
\LessonHead{Lesson III (when commemorated)}
\switchcolumn*

\LA
\noindent Blásius, Sebáste in Arménia cum virtútum laude floréret, ejúsdem civitátis epíscopus elígitur. Qui, quo témpore Diocletiánus insatiábilem crudelitátem in Christiános exercébat, se in spelúncam ábdidit montis Argǽi; ubi támdiu látuit, dum ab Agricolái prǽsidis milítibus venántibus deprehénsus et ad prǽsidem ductus, ejus jussu conjéctus est in víncula. Quo in loco multos ægrótos sanávit, qui ad Blásium, ejus fama sanctitátis addúcti, deferebántur. In illis puer fuit, qui, desperáta a médicis salúte, transvérsa spina fáucibus inhærénte, ánimam agébat. Prodúctus autem ad prǽsidem Blásius semel et íterum, cum nec blandítiis nec minis addúci posset, ut diis sacrificáret, primum virgis cæsus, deínde in equúleo férreis pectínibus dilaniátus est; postrémo, dempto cápite, illústre fídei testimónium Christo Dómino dedit, tértio Nonas Februárii.\par
\switchcolumn
\EN
\noindent This Blaise was chosen Bishop of the city of Sebaste in Armenia, in which place he enjoyed a great reputation for virtue. When Diocletian began to make the Christians the objects of his insatiable cruelty, the Saint hid himself in a cave on Mount Argaeus, where he lay till he was found by some of the soldiers of Agricolaus the governor, who were out hunting. He was brought before the governor, who commanded him to be thrown into irons. While he was in prison, Blaise healed many of the sick, who were brought to him on account of his reputation of saintliness, and among others a boy who had been despaired of by the physicians, and who was at the point of death, from a fish bone which had become fixed in his throat. Blaise appeared twice before the governor, but neither cajolements nor threats could induce him to sacrifice to the gods. He was first beaten with rods, and afterwards put on the rack, where his flesh was mangled with iron combs. At last his head was cut off, whereby he finished a noble testimony to the faith which is in Christ our Lord. He bore witness on the 3rd day of February, (in the year of salvation 316).\par
\switchcolumn*

\LA
\TeDeum{Te Deum}
\switchcolumn
\EN
\TeDeum{Te Deum}
\switchcolumn*

\end{paracol}

\par\needspace{10\baselineskip}\addvspace{1.2em}{\centering\small\textsc{Die 4 Februarii} · \textsc{4 February}\par}
\Office{S. Andreæ Corsini Episcopi et Confessoris}{St. Andrew Corsini, Bishop and Confessor}{Duplex}
\addcontentsline{toc}{subsection}{Die 4 Februarii · S. Andreæ Corsini Episcopi et Confessoris\enspace\textit{St. Andrew Corsini, Bishop and Confessor}}
\begin{paracol}{2}
\LA
\RefLine{Lectiones I–III: de Scriptura occurrente.}
\switchcolumn
\EN
\RefLine{Lessons I–III: from the Scripture of the day.}
\switchcolumn*

\LA
\RefLine{Lectiones VII–IX: ex Commúni unius Confessoris Pontificis (C4).}
\switchcolumn
\EN
\RefLine{Lessons VII–IX: from the Common of a Confessor Bishop (C4).}
\switchcolumn*

\LA
\LessonHead{Lectio IV}
\switchcolumn
\EN
\LessonHead{Lesson IV}
\switchcolumn*

\LA
\noindent Andream, Floréntiæ ex nobili Corsinórum família natum, paréntes precibus a Deo impetrarunt et beátæ Virgini spopondérunt. Qualis autem futurus esset, divino præsagio, antequam nascerétur, ostensum est: nam mater gravida sibi visa est per quietem lupum edidisse, qui ad Carmelitárum ædem pergens, in ipso templi vestibulo statim in agnum convérsus est. Adoléscens pie et ingenue educatus, cum sensim ad vítia declináret, sæpe a matre increpátus fuit. Ubi autem cognóvit se paréntum voto Deíparæ Virgini dicátum fuísse, Dei amóre succénsus, dique visu matris admónitus, Carmelitárum institutum amplexus est; in quo variis tentatiónibus a dæmone vexatus, numquam tamen pótuit a religiónis proposito dimoveri. Mox Lutétiam missus, emenso studiórum curriculo, et laurea donatus, in pátriam revocátur, suique ordinis regímini in Etruria præfícitur.\par
\switchcolumn
\EN
\noindent This Andrew was born at Florence, of the noble family of Corsini, (upon the 30th day of November, in the year 1302.) His birth was a special answer to prayer, and his parents vowed him to the Blessed Virgin. God fore-showed even before his birth what he was to be. While his mother was great with child she dreamt that she brought forth a wolf, which ran to the Carmelite Church and was changed into a lamb as soon as it reached the porch. The lad was brought up in godliness and learning becoming his rank, but turned to bad courses; wherefore his mother often rebuked him. Nevertheless, when he knew how his parents had vowed him to the Maiden Mother of God, the love of God touched his heart, and the vision of his mother moving him, he betook himself to the Institute of the Carmelites. In that place the devil exercised him with many and diverse temptations, but could not break him off from his determination to profess as a friar. He was soon after sent to Paris, where he finished his studies at the University, and took his degree; after which he returned to his own country, and was set over the houses of his order in Tuscany.\par
\switchcolumn*

\LA
\begin{Resp}\Rsign Invéni David servum meum, óleo sancto meo unxi eum: \Star{} Manus enim mea auxiliábitur ei.\end{Resp}
\begin{Resp}\Vsign Nihil profíciet inimícus in eo, et fílius iniquitátis non nocébit ei.\end{Resp}
\begin{Resp}\Rsign Manus enim mea auxiliábitur ei.\end{Resp}
\switchcolumn
\EN
\begin{Resp}\Rsign I have found David My servant, with My holy oil have I anointed him \Star{} For My hand shall help him.\end{Resp}
\begin{Resp}\Vsign The enemy shall prevail nothing against him, nor the son of wickedness afflict him.\end{Resp}
\begin{Resp}\Rsign For My hand shall help him.\end{Resp}
\switchcolumn*

\LA
\LessonHead{Lectio V}
\switchcolumn
\EN
\LessonHead{Lesson V}
\switchcolumn*

\LA
\noindent Intérea Fesulana ecclésia, suo viduáta pastore, eum sibi episcopum elégit: quo munere se indignum æstimans, diu látuit ignotus, donec púeri voce mirabíliter loquentis próditus et extra urbem invéntus, ne divinæ contradíceret voluntati, episcopátum suscépit. Ea dignitate auctus, humilitati, quam semper colúerat, impensius incúbuit; et pastorali solicitudini misericórdiam in páuperes, liberalitátem, oratiónis assiduitátem, vigílias, aliasque virtútes adjunxit, et spíritu étiam prophetico clarus fuit; adeo ut ejus sanctitas ab ómnibus celebrarétur.\par
\switchcolumn
\EN
\noindent The Bishop of Fiesole being dead, the Church in that place chose Andrew Corsini for his successor. He held himself altogether unworthy of that office, and for a long time lay hidden and unknown, till he was betrayed by the voice of a child marvellously speaking, and found outside the city. Then, lest he should seem to resist the Will of God, he took the Bishopric, (in the year 1360.) Being dignified with this office, he set himself to a more perfect exercise of the. virtue of lowliness, whereof he was already a diligent practiser. He was eminent in watchfulness over the flock committed to his charge, joining thereto great tenderness and liberality towards the poor. He continued instant in prayer and watching. Thus was he so adorned with these and many other virtues, and even with the gift of prophecy, that the fame of his holy life was in the mouths of all men.\par
\switchcolumn*

\LA
\begin{Resp}\Rsign Pósui adjutórium super poténtem, et exaltávi eléctum de plebe mea: \Star{} Manus enim mea auxiliábitur ei.\end{Resp}
\begin{Resp}\Vsign Invéni David servum meum, óleo sancto meo unxi eum.\end{Resp}
\begin{Resp}\Rsign Manus enim mea auxiliábitur ei.\end{Resp}
\switchcolumn
\EN
\begin{Resp}\Rsign I have laid help upon one that is mighty, and have exalted one chosen out of My people \Star{} For My hand shall help him.\end{Resp}
\begin{Resp}\Vsign I have found David My servant, with My holy oil have I anointed him.\end{Resp}
\begin{Resp}\Rsign For My hand shall help him.\end{Resp}
\switchcolumn*

\LA
\LessonHead{Lectio VI}
\switchcolumn
\EN
\LessonHead{Lesson VI}
\switchcolumn*

\LA
\noindent His permotus, Urbanus quintus, ad sedandas Bononiæ turbas, Andream legátum misit. Quo in munere multa perpéssus, civium odia, quæ ad interneciónem exárserant, summa prudéntia restinxit; tum restituta tranquillitate, ad propria reversus est. Nec multo post assiduis labóribus et voluntária carnis maceratióne confectus, óbitus die a beáta Vírgine sibi prædicto, ad cæléstia regna migrávit, anno Dómini millesimo tercentésimo septuagesimo tertio, ætátis suæ septuagesimo primo. Quem Urbanus octavus, multis magnisque miraculis clarum, Sanctórum número adscripsit. Ejus corpus Floréntiæ in ecclésia sui ordinis quiescit, et maxima civium veneratióne cólitur, quibus non semel in præsénti discrimine præsidio fuit.\par
\switchcolumn
\EN
\noindent Urban V., moved by the fame of his godly conversation, sent him as his Legate to quiet disturbances at Bologna. He endured much in the discharge of this duty, calming with great wisdom the angry passions of the citizens, who had broken out into civil war, and when peace was restored, he returned home. Shortly after, he received from the Blessed Virgin a warning of his approaching death, and being worn out with his unceasing toil, and the rigour of his voluntary mortifications, he passed to the kingdom of heaven, (upon the 6th day of January,) in the year of our Lord 1373, and the 71st of his own age. His name became illustrious for many and great miracles, and Urban VIII. enrolled him in the number of the Saints. His body resteth at Florence in the Church of his Order, and is looked on with great reverence by the citizens, to whom, even in these days, he hath more than once shown himself a protector.\par
\switchcolumn*

\LA
\begin{Resp}\Rsign Iste est, qui ante Deum magnas virtútes operátus est, et omnis terra doctrína ejus repléta est: \Star{} Ipse intercédat pro peccátis ómnium populórum.\end{Resp}
\begin{Resp}\Vsign Iste est, qui contémpsit vitam mundi, et pervénit ad cæléstia regna.\end{Resp}
\begin{Resp}\Rsign Ipse intercédat pro peccátis ómnium populórum.\end{Resp}
\begin{Resp}\Vsign Glória Patri, et Fílio, \Star{} et Spirítui Sancto.\end{Resp}
\begin{Resp}\Rsign Ipse intercédat pro peccátis ómnium populórum.\end{Resp}
\switchcolumn
\EN
\begin{Resp}\Rsign This is he which wrought great wonders before God, and the whole earth is full of his teaching \Star{} May he pray for all people, that their sins may be forgiven unto them\end{Resp}
\begin{Resp}\Vsign This is he which loved not his life in this world, and hath attained unto the kingdom of heaven.\end{Resp}
\begin{Resp}\Rsign May he pray for all people, that their sins may be forgiven unto them\end{Resp}
\begin{Resp}\Vsign Glory be to the Father, and to the Son, \Star{} and to the Holy Ghost.\end{Resp}
\begin{Resp}\Rsign May he pray for all people, that their sins may be forgiven unto them\end{Resp}
\switchcolumn*

\end{paracol}

\par\needspace{10\baselineskip}\addvspace{1.2em}{\centering\small\textsc{Die 5 Februarii} · \textsc{5 February}\par}
\Office{S. Agathæ Virginis et Martyris}{St. Agatha, Virgin and Martyr}{Duplex}
\addcontentsline{toc}{subsection}{Die 5 Februarii · S. Agathæ Virginis et Martyris\enspace\textit{St. Agatha, Virgin and Martyr}}
\begin{paracol}{2}
\LA
\RefLine{Lectiones I–III: ex Commúni Unius Virginis et Martyris (C6), in 2 loco.}
\switchcolumn
\EN
\RefLine{Lessons I–III: from the Common of a Virgin Martyr (C6), set 2.}
\switchcolumn*

\LA
\RefLine{Lectiones VII–IX: ex Commúni Unius Virginis et Martyris (C6), in 3 loco.}
\switchcolumn
\EN
\RefLine{Lessons VII–IX: from the Common of a Virgin Martyr (C6), set 3.}
\switchcolumn*

\LA
\LessonHead{Lectio IV}
\switchcolumn
\EN
\LessonHead{Lesson IV}
\switchcolumn*

\LA
\noindent Agatha virgo, in Sicília nobílibus paréntibus nata, quam Panormitáni et Catanénses civem suam esse dicunt, in persecutióne Decii imperatóris Cátanæ gloriósi martyrii corónam consecúta est. Nam cum pari pulchritúdinis et castitátis laude commendarétur, Quintiánus Sicíliæ prætor ejus amóre captus est. Sed cum, tentáta modis ómnibus ejus pudicítia, Agatham in suam senténtiam perdúcere non posset, christiánæ superstitiónis nómine comprehénsam, Aphrodísiæ cuidam mulíeri depravándam tradit. Quæ Aphrodísiæ consuetúdine cum de constántia coléndæ christiánæ fídei et servándæ virginitátis removéri non posset, núntiat illa Quintiáno se in Agatha óperam pérdere. Quare ille ad se Vírginem addúci jubet; et, Nonne, inquit, te pudet nóbili génere natam, húmilem et servílem Christianórum vitam ágere? Cui Agatha: Multo præstántior est christiána humílitas et sérvitus, regum ópibus ac supérbia.\par
\switchcolumn
\EN
\noindent The Maiden Agatha was a Sicilian of noble birth. The citizens of Palermo and Catania dispute as to which city had the honour of being her birthplace. It was at Catania that, during the persecution under the Emperor Decius, she won the crown of a glorious martyrdom. She was equally celebrated for her beauty and her chastity, and Quintianus, Praetor of Sicily, conceived a passion for her. He tried every sort of device to overcome her modesty, and when he found it impossible to make her consent to his wishes, he caused her to be arrested on a charge of Christian superstition, and handed over to a woman named Aphrodisia to be corrupted. The company, however, of this woman had no effect in shaking her constancy in the Christian worship, nor her settled determination to preserve her purity. Aphrodisia therefore reported to Quintianus that she was only throwing away her pains on Agatha. He ordered her to be brought before him. Thou, said he, art the daughter of a noble family dost thou feel no shame in living the degraded and slavish life of a Christian? Agatha answered him, The lowliness and bondage of a Christian are far nobler than the estate and pride of a king.\par
\switchcolumn*

\LA
\begin{Resp}\Rsign Ego autem adjúta a Dómino, perseverábo in confessióne ejus, qui me salvam fecit, \Star{} Et consolátus est me.\end{Resp}
\begin{Resp}\Vsign Grátias tibi ago, Dómine Jesu Christe, qui misísti ad me Apóstolum tuum curáre vúlnera mea.\end{Resp}
\begin{Resp}\Rsign Et consolátus est me.\end{Resp}
\switchcolumn
\EN
\begin{Resp}\Rsign But by the Lord's help I will continue to acknowledge Him Who hath saved me \Star{} And strengthened me.\end{Resp}
\begin{Resp}\Vsign I thank thee, O my Lord Jesus Christ, because Thou hast sent thine Apostle unto me to heal my wounds.\end{Resp}
\begin{Resp}\Rsign And strengthened me.\end{Resp}
\switchcolumn*

\LA
\LessonHead{Lectio V}
\switchcolumn
\EN
\LessonHead{Lesson V}
\switchcolumn*

\LA
\noindent Quam ob rem irátus prætor hanc ei optiónem dat, velítne pótius venerári deos, an vim tormentórum subíre. At illa constans in fide, primum cólaphis cæsa míttitur in cárcerem; unde postrídie edúcta, cum in senténtia permanéret, admótis candéntibus láminis in equúleo torquétur. Tum ei mamílla abscínditur; quo in vúlnere Quintiánum appéllans Virgo, Crudélis, inquit, tyránne, non te pudet amputáre in fémina, quod ipse in matre suxísti? Mox conjécta in víncula, sequénti nocte a sene quodam, qui se Christi Apóstolum esse dicébat, sanáta est. Rursum evocáta a prætóre, et in Christi confessióne persevérans, in acutis téstulis et candéntibus carbónibus ei subjéctis volutátur.\par
\switchcolumn
\EN
\noindent Then the Praetor, being incensed against her, gave her the alternative of either sacrificing to the gods, or being submitted to the torture; and as she remained firm in the faith, she was buffeted and sent back to prison. The next day she was brought forth, and, because her resolution was still unshaken, she was stretched on the rack and tortured with pieces of white-hot metal. Then her breasts were cut off. When Agatha received this injury she cried out to Quintianus, Cruel tyrant, art thou not ashamed to do this to me, having thyself sucked at a mother's breast? She was remanded again to prison and put in irons. That night an old man, who called himself an Apostle of Christ, came to her, and healed her wounds. The following day she was brought for the last time before the Praetor. Her constancy was unmoved, and she was rolled on sharp potsherds and live embers.\par
\switchcolumn*

\LA
\begin{Resp}\Rsign Ipse me curávit, qui per Apóstolum Petrum in custódia me confortávit, pro eo quod jussa sum suspéndi in equúleo. \Star{} Propter fidem castitátis ádjuva me, Dómine, Deus meus, in tortúra mamillárum meárum.\end{Resp}
\begin{Resp}\Vsign Ipse me dignátus est ab omni plaga curáre, et mamíllam meam meo péctori restitúere.\end{Resp}
\begin{Resp}\Rsign Propter fidem castitátis ádjuva me, Dómine, Deus meus, in tortúra mamillárum mearum.\end{Resp}
\switchcolumn
\EN
\begin{Resp}\Rsign He hath cured me Who hath sent His Apostle Peter to the prison to strengthen me, after I was ordered to be put on the rack. \Star{} Because of my chastity, help me, O Lord my God, for they are torturing my breasts.\end{Resp}
\begin{Resp}\Vsign He hath vouchsafed to heal me of all my wounds, and to put new paps on my breasts.\end{Resp}
\begin{Resp}\Rsign Because of my chastity, help me, O Lord my God, for they are torturing my breasts.\end{Resp}
\switchcolumn*

\LA
\LessonHead{Lectio VI}
\switchcolumn
\EN
\LessonHead{Lesson VI}
\switchcolumn*

\LA
\noindent Quo témpore, ingénti terræmótu urbs tota contrémuit, ac duo paríetes, corruéntes, Silvínum et Falcónium íntimos prætóris familiáres oppressérunt. Quare veheménter commóta civitáte, véritus pópuli tumúltum Quintiánus, Agatham semimórtuam clam redúci ímperat in cárcerem. Quæ sic Deum precáta: Dómine, qui me custodísti ab infántia, qui abstulísti a me amórem sæculi, qui me carníficum torméntis superiórem præstitísti, áccipe ánimam meam. Ea in oratióne migrávit in cælum, Nonis Februárii: cujus corpus a Christiánis sepelítur.\par
\switchcolumn
\EN
\noindent At that time the whole city was shaken with a great earthquake, and two of the Praetor's dearest friends, Silvinus and Falconius, were killed by falling walls. The townspeople were in an uproar, and Quintianus, in fear of a riot, ordered Agatha, who was half dead, to be carried back to prison quietly. Then she made the following prayer O Lord, Who hast been my Keeper from my childhood, Who hast taken from me all love for this present world, Who hast strengthened me so that I am more than conqueror over the cruelty of the executioners, receive my spirit, and with these words she passed to heaven. She finished her testimony on the 5th day of February, (in the year of our Lord 251.) Her body was buried by the Christians.\par
\switchcolumn*

\LA
\begin{Resp}\Rsign Vidísti, Dómine, et spectásti agónem meum, quómodo pugnávi in stadio; sed, quia nólui obedíre mandátis príncipum, \Star{} Jussa sum in mamílla torquéri.\end{Resp}
\begin{Resp}\Vsign Propter veritátem, et mansuetúdinem, et justítiam.\end{Resp}
\begin{Resp}\Rsign Jussa sum in mamílla torquéri.\end{Resp}
\begin{Resp}\Vsign Glória Patri, et Fílio, \Star{} et Spirítui Sancto.\end{Resp}
\begin{Resp}\Rsign Jussa sum in mamílla torquéri.\end{Resp}
\switchcolumn
\EN
\begin{Resp}\Rsign O Lord, Thou hast seen and known how I have fought, and how I have run in the race; but, because I would not obey the magistrates \Star{} They ordered me to be tortured in the breasts.\end{Resp}
\begin{Resp}\Vsign Because of truth, and meekness, and righteousness,\end{Resp}
\begin{Resp}\Rsign They ordered me to be tortured in the breasts.\end{Resp}
\begin{Resp}\Vsign Glory be to the Father, and to the Son, \Star{} and to the Holy Ghost.\end{Resp}
\begin{Resp}\Rsign They ordered me to be tortured in the breasts.\end{Resp}
\switchcolumn*

\end{paracol}

\par\needspace{10\baselineskip}\addvspace{1.2em}{\centering\small\textsc{Die 6 Februarii} · \textsc{6 February}\par}
\Office{S. Titi Episcopi et Confessoris}{St. Titus, Bishop and Confessor}{Duplex}
\addcontentsline{toc}{subsection}{Die 6 Februarii · S. Titi Episcopi et Confessoris\enspace\textit{St. Titus, Bishop and Confessor}}
\begin{paracol}{2}
\LA
\RefLine{Lectiones I–III: de Scriptura occurrente.}
\switchcolumn
\EN
\RefLine{Lessons I–III: from the Scripture of the day.}
\switchcolumn*

\LA
\RefLine{Lectio IX: de S. Dorotheæ Virginis et Martyris (02-06o).}
\switchcolumn
\EN
\RefLine{Lesson IX: of St. Dorothy, Virgin and Martyr (02-06o).}
\switchcolumn*

\LA
\LessonHead{Lectio IV}
\switchcolumn
\EN
\LessonHead{Lesson IV}
\switchcolumn*

\LA
\noindent Titus Cretensium episcopum, vix Pauli Apóstoli verbo christianæ fidei sacramentis mysteriísque excultum, ea sanctitátis luce Ecclésiæ tunc vagiénti effulsisse compertum est, ut inter ejusdem Doctoris Géntium discipulos merúerit cooptari. Ascitus in partem óneris prædicatiónis, ádeo evangelizándi ardore, et fidelitate Paulo éxstitit carus, ut ipse, cum venísset Tróadem propter Evangélium Christi, testátus sit non habuísse requiem spiritui suo, eo quod Titum fratrem suum ibi non invenerit. Et paulo post Macedóniam petens, rursus suam in eum caritátem ita éxprimit: Sed qui consolátur húmiles, consolátus est nos Deus in adventu Titi.\par
\switchcolumn
\EN
\noindent The Apostle Paul had no sooner by his preaching gathered the Cretans to the Sacraments and doctrines of the Christian Faith, than he made Titus their Bishop. It is agreed that the holiness of Titus shone with so bright a light before the infant Church, that he earned fellowship with the disciples of the Teacher of the Gentiles. Being taken to share in the work of preaching, he so endeared himself to Paul by his faithfulness and zeal in declaring the Gospel, that the Apostle saith When I came to Troas to preach Christ's Gospel, and a door was opened unto me of the Lord, I had no rest in my spirit, because I found not Titus my brother but taking my leave of them, I went from thence into Macedonia. (2 Cor. ii. 12, 13.) And again he saith: When we were come into Macedonia, our flesh had no rest, but we were troubled on every side; without were fightings, within were fears. Nevertheless, God, that comforteth those that are cast down, comforted us by the coming of Titus. (vii. 5, 6.)\par
\switchcolumn*

\LA
\begin{Resp}\Rsign Invéni David servum meum, óleo sancto meo unxi eum: \Star{} Manus enim mea auxiliábitur ei.\end{Resp}
\begin{Resp}\Vsign Nihil profíciet inimícus in eo, et fílius iniquitátis non nocébit ei.\end{Resp}
\begin{Resp}\Rsign Manus enim mea auxiliábitur ei.\end{Resp}
\switchcolumn
\EN
\begin{Resp}\Rsign I have found David My servant, with My holy oil have I anointed him \Star{} For My hand shall help him.\end{Resp}
\begin{Resp}\Vsign The enemy shall prevail nothing against him, nor the son of wickedness afflict him.\end{Resp}
\begin{Resp}\Rsign For My hand shall help him.\end{Resp}
\switchcolumn*

\LA
\LessonHead{Lectio V}
\switchcolumn
\EN
\LessonHead{Lesson V}
\switchcolumn*

\LA
\noindent Quam ob rem Corinthum ab Apóstolo missus, ea sapiéntia et lenitate legatiónis hujus munere functus est, quæ præsertim de fidelium pietáte eleemosynas colligendas ad sublevandam Ecclésiæ Hebræórum inópiam spectabat, ut Corinthios non solum in Christi fide continúerit, sed étiam desidérium, fletum, æmulatiónem inter eos pro Paulo, qui illos primus instituit, excitaverit. Ad effundéndum interim inter gentes linguis locisque distinctas divini verbi semen, plúribus terra marique itinéribus reléctis, magnaque animi firmitáte pro crucis trophæo curis laboribúsque exantlátis, una cum duce Paulo Cretæ insulam áppulit. Cum porro huic ecclésiæ epíscopus ab ipso Apóstolo deléctus esset, dubitándum non est quin in eo munere ita versátus sit, ut, juxta ipsíus Pauli præcéptoris mónita, seípsum præbúerit exemplum bonórum operum in doctrina, in integritate, in gravitáte.\par
\switchcolumn
\EN
\noindent It was this affection of Paul toward Titus, which had induced him to send him to Corinth upon a Mission which mainly concerned the collection of alms from the charity of the faithful for the relief of the poor Hebrew Saints at Jerusalem. This mission Titus discharged with such wisdom and gentleness, that he not only strengthened the Corinthians in the faith, but also stirred up in them an earnest desire, a mourning, a fervent mind toward Paul, their earliest teacher. \textasciitilde{}(vii. 7.) Many were the other journeys by land and sea which Titus undertook in order to sow the seed of God's word among men of diverse nations, tongues, and countries. Filled with bold loyalty to the banner of the Cross, he went with Paul to the island of Crete. Of the Church of Crete the Apostle himself made him the first Bishop; and we may not doubt that, as such, he was what his Teacher bade him be, in all things showing himself a pattern of good works, in doctrine, in uncorruptness, in gravity. (Tit 2:7)\par
\switchcolumn*

\LA
\begin{Resp}\Rsign Pósui adjutórium super poténtem, et exaltávi eléctum de plebe mea: \Star{} Manus enim mea auxiliábitur ei.\end{Resp}
\begin{Resp}\Vsign Invéni David servum meum, óleo sancto meo unxi eum.\end{Resp}
\begin{Resp}\Rsign Manus enim mea auxiliábitur ei.\end{Resp}
\switchcolumn
\EN
\begin{Resp}\Rsign I have laid help upon one that is mighty, and have exalted one chosen out of My people \Star{} For My hand shall help him.\end{Resp}
\begin{Resp}\Vsign I have found David My servant, with My holy oil have I anointed him.\end{Resp}
\begin{Resp}\Rsign For My hand shall help him.\end{Resp}
\switchcolumn*

\LA
\LessonHead{Lectio VI}
\switchcolumn
\EN
\LessonHead{Lesson VI}
\switchcolumn*

\LA
\noindent Itaque tamquam lucérna inter eos, qui in idololatríæ et mendaciórum ténebris véluti in umbra mortis sedébant, religiónis jubar diffudit. Traditur eum inter Dálmatas, ut crucis vexíllum explicaret, strenue consudasse. Tandem meritórum et diérum plenus, quarto supra nonagesimum anno, pridie Nonas Januarii, pretiósa justórum morte obdormívit in Dómino, et sepúltus est in ecclésia ubi ab Apóstolo minister fuerat constitútus. Hujus nomen, a sancto Joánne Chrysostomo et a sancto Hieronymo præcipue commendátum, Martyrologio Romano eadem die inscriptum legitur: ejus autem festum summus Pontifex Pius nonus ab univérsa Ecclésia celebrari præcepit.\par
\switchcolumn
\EN
\noindent Like a candle, he gave forth the light of faith in the midst of men sitting in the darkness of idolatry and falsehood, as in the shadow of death. He is said to have sweated mightily to unfurl the banner of the Cross among the Dalmatians. He was full of days and good works, when, upon a 4th of January, in the 94th year of his age, he died one of those deaths which are precious in the sight of the Lord. He was buried in the Church of which the Apostle had made him the minister. His praises have been mostly written by St. John Chrysostom and St Jerome. The 4th of January is the day upon which his name is read in the Roman Martyrology, but Pope Pius IX. assigned for his Festival, to be kept with an Office and Mass by the clergy secular and regular throughout the Catholic world, the first free day afterwards.\par
\switchcolumn*

\LA
\begin{Resp}\Rsign Iste est, qui ante Deum magnas virtútes operátus est, et omnis terra doctrína ejus repléta est: \Star{} Ipse intercédat pro peccátis ómnium populórum.\end{Resp}
\begin{Resp}\Vsign Iste est, qui contémpsit vitam mundi, et pervénit ad cæléstia regna.\end{Resp}
\begin{Resp}\Rsign Ipse intercédat pro peccátis ómnium populórum.\end{Resp}
\begin{Resp}\Vsign Glória Patri, et Fílio, \Star{} et Spirítui Sancto.\end{Resp}
\begin{Resp}\Rsign Ipse intercédat pro peccátis ómnium populórum.\end{Resp}
\switchcolumn
\EN
\begin{Resp}\Rsign This is he which wrought great wonders before God, and the whole earth is full of his teaching \Star{} May he pray for all people, that their sins may be forgiven unto them\end{Resp}
\begin{Resp}\Vsign This is he which loved not his life in this world, and hath attained unto the kingdom of heaven.\end{Resp}
\begin{Resp}\Rsign May he pray for all people, that their sins may be forgiven unto them\end{Resp}
\begin{Resp}\Vsign Glory be to the Father, and to the Son, \Star{} and to the Holy Ghost.\end{Resp}
\begin{Resp}\Rsign May he pray for all people, that their sins may be forgiven unto them\end{Resp}
\switchcolumn*

\LA
\LessonHead{Lectio VII}
\switchcolumn
\EN
\LessonHead{Lesson VII}
\switchcolumn*

\LA
\LessonTitle{Léctio sancti Evangélii secúndum Lucam}
\switchcolumn
\EN
\LessonTitle{From the Holy Gospel according to Luke}
\switchcolumn*

\LA
\Cite{Luc 10:1-9}
\switchcolumn
\EN
\Cite{Luke 10:1-9}
\switchcolumn*

\LA
\noindent In illo témpore: Designávit Dóminus et álios septuagínta duos: et misit illos binos ante fáciem suam, in omnem civitátem et locum, quo erat ipse ventúrus. Et réliqua.\par
\switchcolumn
\EN
\noindent At that time, the Lord appointed other seventy-two also, and sent them two and two before His face into every city and place, whither He Himself would come. And so on.\par
\switchcolumn*

\LA
\LessonTitle{Homilía sancti Gregórii Papæ}
\switchcolumn
\EN
\LessonTitle{Homily by Pope St. Gregory (the Great.)}
\switchcolumn*

\LA
\Source{Homilia 17 in Evangelia}
\switchcolumn
\EN
\Source{17th Homily on the Gospels}
\switchcolumn*

\LA
\noindent Dóminus et Salvátor noster, fratres caríssimi, aliquándo nos sermónibus, aliquándo vero opéribus ádmonet. Ipsa étenim facta ejus præcépta sunt: quia dum áliquid tácitus facit, quid ágere debeámus innotéscit. Ecce enim binos in prædicatiónem discípulos mittit: quia duo sunt præcépta caritátis, Dei vidélicet amor, et próximi: et minus quam inter duos cáritas habéri non potest. Nemo enim próprie ad semetípsum habére caritátem dícitur: sed diléctio in álterum tendit, ut cáritas esse possit.\par
\switchcolumn
\EN
\noindent Dearly beloved brethren, our Lord and Saviour doth sometimes admonish us by words, and sometimes by works. Yea, His very works do themselves teach us for that which He doth silently His example still moveth us to copy. Behold how He sendeth forth His disciples to preach by two and two since there are two commandments to love, that is, a commandment to love God, and a commandment to love our neighbour and where there are not two, the one, being alone, hath not whereon to do the Lord's commandment. And no man can properly be said to love himself: for love tendeth outward toward our neighbour, if it be the love whereto the Gospel doth oblige us.\par
\switchcolumn*

\LA
\begin{Resp}\Rsign Amávit eum Dóminus, et ornávit eum: stolam glóriæ índuit eum, \Star{} Et ad portas paradísi coronávit eum.\end{Resp}
\begin{Resp}\Vsign Induit eum Dóminus lorícam fídei, et ornávit eum.\end{Resp}
\begin{Resp}\Rsign Et ad portas paradísi coronávit eum.\end{Resp}
\switchcolumn
\EN
\begin{Resp}\Rsign The Lord loved him and beautified him He clothed him with a robe of glory \Star{} And crowned him at the gates of Paradise.\end{Resp}
\begin{Resp}\Vsign The Lord hath put on him the breast-plate of faith, and hath adorned him.\end{Resp}
\begin{Resp}\Rsign And crowned him at the gates of Paradise.\end{Resp}
\switchcolumn*

\LA
\LessonHead{Lectio VIII}
\switchcolumn
\EN
\LessonHead{Lesson VIII}
\switchcolumn*

\LA
\noindent Ecce enim binos ad prædicándum discípulos Dóminus mittit: quátenus hoc nobis tácitus ínnuat, quia qui caritátem erga álterum non habet, prædicatiónis offícium suscípere nullátenus debet. Bene autem dícitur, quia misit eos ante fáciem suam in omnem civitátem et locum, quo erat ipse ventúrus. Prædicatóres enim suos Dóminus séquitur: quia prædicátio prǽvenit, et tunc ad mentis nostræ habitáculum Dóminus venit, quando verba exhortatiónis præcúrrunt: atque per hoc véritas in mente suscípitur.\par
\switchcolumn
\EN
\noindent Behold, the Lord sendeth forth His disciples to preach by two and two and thus doing, He doth silently teach us that whosoever loveth not his neighbour, such a one it behoveth not to take upon him the office of a preacher. Well also is it said that He sent them before His face into every city and place whither He Himself would come. The Lord followeth His preachers first cometh preaching, and then the Lord Himself cometh to the house of our mind, whither the word of exhortation hath come before and so cometh the truth into our mind.\par
\switchcolumn*

\LA
\begin{Resp}\Rsign Sint lumbi vestri præcíncti, et lucérnæ ardéntes in mánibus vestris: \Star{} Et vos símiles homínibus exspectántibus dóminum suum, quando revertátur a núptiis.\end{Resp}
\begin{Resp}\Vsign Vigiláte ergo, quia nescítis qua hora Dóminus vester ventúrus sit.\end{Resp}
\begin{Resp}\Rsign Et vos símiles homínibus exspectántibus dóminum suum, quando revertátur a núptiis.\end{Resp}
\begin{Resp}\Vsign Glória Patri, et Fílio, \Star{} et Spirítui Sancto.\end{Resp}
\begin{Resp}\Rsign Et vos símiles homínibus exspectántibus dóminum suum, quando revertátur a núptiis.\end{Resp}
\switchcolumn
\EN
\begin{Resp}\Rsign Let your loins be girded about, and your lights burning; \Star{} And ye yourselves like unto men that wait for their lord, when he will return from the wedding.\end{Resp}
\begin{Resp}\Vsign Watch therefore, for ye know not what hour your Lord doth come.\end{Resp}
\begin{Resp}\Rsign And ye yourselves like unto men that wait for their lord, when he will return from the wedding.\end{Resp}
\begin{Resp}\Vsign Glory be to the Father, and to the Son, \Star{} and to the Holy Ghost.\end{Resp}
\begin{Resp}\Rsign And ye yourselves like unto men that wait for their lord, when he will return from the wedding.\end{Resp}
\switchcolumn*

\end{paracol}

\par\needspace{10\baselineskip}\addvspace{1.2em}{\centering\small\textsc{Die 6 Februarii} · \textsc{6 February}\par}
\Office{S. Dorotheæ Virginis et Martyris}{St. Dorothy, Virgin and Martyr}{Simplex}
\addcontentsline{toc}{subsection}{Die 6 Februarii · S. Dorotheæ Virginis et Martyris\enspace\textit{St. Dorothy, Virgin and Martyr}}
\begin{paracol}{2}
\LA
\LessonHead{Lectio IX (pro commemoratione)}
\switchcolumn
\EN
\LessonHead{Lesson IX (when commemorated)}
\switchcolumn*

\LA
\LessonTitle{Commemoratio: S. Dorotheæ Virginis et Martyris}
\switchcolumn
\EN
\LessonTitle{Commemoration of: St. Dorothy, Virgin and Martyr}
\switchcolumn*

\LA
\noindent Doróthea virgo, ex Cæsaréa Cappadociæ, propter Christi confessiónem ab Apricio præside comprehensa, Chrystæ et Callistæ soróribus, quæ a fide defécerant, tradita est, ut eam a propósito removérent. Sed contra factum est; nam eas Doróthea ad cultum christianæ religiónis redúxit, propter quam étiam martyrium suscepérunt. Quare Virgo equuleo diu torta et palmis cæsa, ad extremum cápitis damnata, duplicatam virginitátis et martyrii palmam accepit.\par
\switchcolumn
\EN
\noindent The maiden Dorothy, of Caesarea in Cappadocia, was betrayed to Apricius the President by her two sisters, Chrysta and Callista, who had denied the faith, in the hope that he would induce her to do likewise. She was arrested, but it came not to pass as they hoped. On the contrary, she brought them back to the Christian worship, and they received martyrdom. She was long tormented upon the rack, and scourged with palm branches, and in the end was beheaded, receiving the double palm of virginity and martyrdom.\par
\switchcolumn*

\LA
\TeDeum{Te Deum}
\switchcolumn
\EN
\TeDeum{Te Deum}
\switchcolumn*

\end{paracol}

\par\needspace{10\baselineskip}\addvspace{1.2em}{\centering\small\textsc{Die 7 Februarii} · \textsc{7 February}\par}
\Office{S. Romualdi Abbatis}{St. Romuald, Abbot}{Duplex}
\addcontentsline{toc}{subsection}{Die 7 Februarii · S. Romualdi Abbatis\enspace\textit{St. Romuald, Abbot}}
\begin{paracol}{2}
\LA
\RefLine{Lectiones I–III: de Scriptura occurrente.}
\switchcolumn
\EN
\RefLine{Lessons I–III: from the Scripture of the day.}
\switchcolumn*

\LA
\LessonHead{Lectio IV}
\switchcolumn
\EN
\LessonHead{Lesson IV}
\switchcolumn*

\LA
\noindent Romuáldus, Ravénnæ, Sérgio patre, nóbili génere natus, adoléscens in propínquum monastérium Classénse pœniténtiæ causa secéssit: ubi religiósi hóminis sermóne ad pietátis stúdium veheméntius incénsus, viso étiam semel et íterum per noctem in ecclésia beáto Apollinári, quod Dei servus illi futúrum promíserat, mónachus effícitur. Mox ad Marínum, vitæ sanctitáte ac severióre disciplína in fínibus Venetórum eo témpore célebrem, se cóntulit, ut ad arctam et sublímem perfectiónis viam eo magístro ac duce uterétur.\par
\switchcolumn
\EN
\noindent The holy Abbot Romuald was the son of one Sergius, of a noble family of Ravenna. While he was still very young, he went to a neighbouring monastery at Classis to do penance. While he was there he heard a discourse by a monk, which stirred him up strongly to aim at godliness of living; and he had afterwards in the Church by night two visions in which the blessed servant of God Apollinaris foretold to him that he should become a monk himself. He accordingly did so; and soon afterwards betook himself to one Marinus, whose holy life and strict discipline were then much noised about in all the coasts of the Venetians, that he might by his teaching and guidance attain towards the hard and lofty point of perfection.\par
\switchcolumn*

\LA
\begin{Resp}\Rsign Honéstum fecit illum Dóminus, et custodívit eum ab inimícis, et a seductóribus tutávit illum: \Star{} Et dedit illi claritátem ætérnam.\end{Resp}
\begin{Resp}\Vsign Justum dedúxit Dóminus per vias rectas, et osténdit illi regnum Dei.\end{Resp}
\begin{Resp}\Rsign Et dedit illi claritátem ætérnam.\end{Resp}
\switchcolumn
\EN
\begin{Resp}\Rsign The Lord made him honourable, and defended him from his enemies, and kept him safe from those that lay in wait for him; \Star{} And gave him perpetual glory.\end{Resp}
\begin{Resp}\Vsign He went down with him into the pit, and left him not in bonds.\end{Resp}
\begin{Resp}\Rsign And gave him perpetual glory.\end{Resp}
\switchcolumn*

\LA
\LessonHead{Lectio V}
\switchcolumn
\EN
\LessonHead{Lesson V}
\switchcolumn*

\LA
\noindent Multis sátanæ insídiis et hóminum invídia oppugnátus, tanto humílior se assídue jejúniis et oratiónibus exercébat, et rerum cæléstium meditatióne, vim lacrimárum profúndens, fruebátur: vultu tamen ádeo læto semper erat, ut intuéntes exhilaráret. Magno apud príncipes et reges in honóre fuit; multíque ejus consílio, mundi illécebris abjéctis, solitúdinem petiérunt. Martýrii quoque cupiditáte flagrávit, cujus causa dum in Pannóniam proficíscitur, morbo, quo afflictabátur cum progrederétur, levabátur cum recéderet, revérti cógitur.\par
\switchcolumn
\EN
\noindent The more he was assailed by the wiles of Satan and the unkindness of men, the more did he exercise himself in lowliness, with continual fasting and prayer, and rejoice in thinking of heavenly things, with abundance of tears. And all the while he bore so bright a face as gladdened all who looked on him. He was held in great honour by princes and kings, and his counsel moved many to leave the blandishments of the world and withdraw to the desert. He had such a burning desire to obtain the crown of martyrdom that he set out for Pannonia on purpose to seek it, but, falling into sickness whenever he went forward though growing strong again whenever he drew back, he behoved to return home.\par
\switchcolumn*

\LA
\begin{Resp}\Rsign Amávit eum Dóminus, et ornávit eum: stolam glóriæ índuit eum, \Star{} Et ad portas paradísi coronávit eum.\end{Resp}
\begin{Resp}\Vsign Índuit eum Dóminus lorícam fídei, et ornávit eum.\end{Resp}
\begin{Resp}\Rsign Et ad portas paradísi coronávit eum.\end{Resp}
\switchcolumn
\EN
\begin{Resp}\Rsign The Lord loved him and beautified him; He clothed him with a robe of glory, \Star{} And crowned him at the gates of Paradise.\end{Resp}
\begin{Resp}\Vsign The Lord hath put on him the breast-plate of faith, and hath adorned him.\end{Resp}
\begin{Resp}\Rsign And crowned him at the gates of Paradise.\end{Resp}
\switchcolumn*

\LA
\LessonHead{Lectio VI}
\switchcolumn
\EN
\LessonHead{Lesson VI}
\switchcolumn*

\LA
\noindent In vita et post mortem miráculis clarus, spíritu étiam prophetíæ non cáruit. Scalam a terra cælum pertingéntem, in similitúdinem Jacob Patriárchæ, per quam hómines in veste cándida ascendébant et descendébant, per visum conspéxit; eóque Camaldulénses mónachos, quorum institúti auctor fuit, designári mirabíliter agnóvit. Dénique cum annos centum et vigínti ágeret, et centum ipsos in summa vitæ asperitáte Deo servísset, ad eum migrávit, anno salútis millésimo vigésimo séptimo. Ejus corpus quinquénnio postquam sepúltum fúerat, íntegrum repértum, Fabriáni in ecclésia sui órdinis honorífice cónditum est.\par
\switchcolumn
\EN
\noindent God worked miracles by him both during his life and after his death, and likewise gave him the gift of prophecy. Like the Patriarch Jacob, he saw a ladder reaching from earth to heaven, and men in white garments ascending and descending upon it, in whom he marvellously knew were represented the monks of the Camaldolese Institute, of which he was the founder. At the age of 120 years, of which he had spent 100 in serving God in great hardness, he passed into His Presence, in the year of Salvation 1027. Five years after his death his body was found incorrupt, and laid in a magnificent grave in the Church of his order at Fabriano.\par
\switchcolumn*

\LA
\begin{Resp}\Rsign Iste homo perfécit ómnia quæ locútus est ei Deus, et dixit ad eum: Ingrédere in réquiem meam: \Star{} Quia te vidi justum coram me ex ómnibus géntibus.\end{Resp}
\begin{Resp}\Vsign Iste est, qui contémpsit vitam mundi, et pervénit ad cæléstia regna.\end{Resp}
\begin{Resp}\Rsign Quia te vidi justum coram me ex ómnibus géntibus.\end{Resp}
\begin{Resp}\Vsign Glória Patri, et Fílio, \Star{} et Spirítui Sancto.\end{Resp}
\begin{Resp}\Rsign Quia te vidi justum coram me ex ómnibus géntibus.\end{Resp}
\switchcolumn
\EN
\begin{Resp}\Rsign This is he which did according unto all that God commanded him; and God said unto him: Enter thou into My rest, \Star{} For thee have I seen righteous before Me among all people.\end{Resp}
\begin{Resp}\Vsign This is he which loved not his life in this world, and is come unto an everlasting kingdom.\end{Resp}
\begin{Resp}\Rsign For thee have I seen righteous before Me among all people.\end{Resp}
\begin{Resp}\Vsign Glory be to the Father, and to the Son, \Star{} and to the Holy Ghost.\end{Resp}
\begin{Resp}\Rsign For thee have I seen righteous before Me among all people.\end{Resp}
\switchcolumn*

\LA
\LessonHead{Lectio VII}
\switchcolumn
\EN
\LessonHead{Lesson VII}
\switchcolumn*

\LA
\LessonTitle{Léctio sancti Evangélii secúndum Matthǽum}
\switchcolumn
\EN
\LessonTitle{From the Holy Gospel according to Matthew}
\switchcolumn*

\LA
\Cite{Matt 19:27-29}
\switchcolumn
\EN
\Cite{Matt 19:27-29}
\switchcolumn*

\LA
\noindent In illo témpore: Dixit Petrus ad Jesum: Ecce nos relíquimus ómnia, et secúti sumus te: quid ergo erit nobis? Et réliqua.\par
\switchcolumn
\EN
\noindent In that time, Peter said to Jesus said to him: Behold we have left all things, and have followed thee: what therefore shall we have? And so on.\par
\switchcolumn*

\LA
\LessonTitle{Homilía sancti Bedæ Venerábilis Presbýteri}
\switchcolumn
\EN
\LessonTitle{Sermon of the Venerable Bede, Priest (at Jarrow)}
\switchcolumn*

\LA
\Source{Homilia in Natali S. Benedicti Ep.}
\switchcolumn
\EN
\Source{Homily for St. Benedict's Birth-day}
\switchcolumn*

\LA
\noindent Perféctus ille est, qui ábiens vendit ómnia quæ habet, et dat paupéribus, ac véniens séquitur Christum; habébit enim thesáurum non deficiéntem in cælis. Unde bene, interrogánte Petro, dixit tálibus Jesus: Amen dico vobis, quod vos qui secúti estis me, in regeneratióne, cum séderit Fílius hóminis in sede majestátis suæ, sedébitis et vos super sedes duódecim, judicántes duódecim tribus Israël. In hac quippe vita pro ejus nómine laborántes, in ália prǽmium speráre dócuit, id est, in regeneratióne; cum vidélicet in vitam immortálem fuérimus resurgéndo regeneráti, qui in vitam cadúcam mortáliter erámus géniti.\par
\switchcolumn
\EN
\noindent If thou wilt be perfect, saith Christ, go and sell that thou hast, and give to the poor, and come and follow Me and thou shalt have treasure in heaven. (Matth. xix. 21.) Yea, treasure that passeth not away! Unto such saith Jesus, at the questioning of Peter Amen I say unto you, that ye which have followed Me, in the regeneration, when the Son of Man shall sit in the throne of His glory, ye also shall sit upon twelve thrones, judging the twelve tribes of Israel. He taught them, which work for His Name's sake in this life, to look for their reward in another life that is, in the regeneration. In the regeneration! when we who have been born dying creatures into a dying life, shall in the resurrection be born again into an undying life.\par
\switchcolumn*

\LA
\begin{Resp}\Rsign Iste est qui ante Deum magnas virtútes operátus est, et de omni corde suo laudávit Dóminum: \Star{} Ipse intercédat pro peccátis ómnium populórum.\end{Resp}
\begin{Resp}\Vsign Ecce homo sine queréla, verus Dei cultor, ábstinens se ab omni ópere malo, et pérmanens in innocéntia sua.\end{Resp}
\begin{Resp}\Rsign Ipse intercédat pro peccátis ómnium populórum.\end{Resp}
\switchcolumn
\EN
\begin{Resp}\Rsign This is he which wrought great wonders before God, and praised the Lord with all his heart. \Star{} May he pray for all people, that their sins may be forgiven unto them.\end{Resp}
\begin{Resp}\Vsign Behold a man without blame, a worshipper of God in truth, keeping himself clean from every evil work, and abiding still in his innocency.\end{Resp}
\begin{Resp}\Rsign May he pray for all people, that their sins may be forgiven unto them.\end{Resp}
\switchcolumn*

\LA
\LessonHead{Lectio VIII}
\switchcolumn
\EN
\LessonHead{Lesson VIII}
\switchcolumn*

\LA
\noindent Et justa prorsus retribútio, ut, qui hic pro Christo humánæ glóriam celsitúdinis neglexérunt, illic a Christo júdices glorificáti singuláriter cum eo assídeant qui a sequéndis ejus vestígiis nulla ratióne póterant avélli. Nemo autem putet duódecim tantum Apóstolos, quia pro Juda prævaricánte Matthías eléctus est, tunc esse judicatúros; sicut nec duódecim solæ sunt tribus Israël judicándæ: alióquin tribus Levi, quæ tertiadécima est, injudicáta recédet.\par
\switchcolumn
\EN
\noindent And soothly, it is a just retribution, that they, who, while they were yet here, have for Christ's sake set no store by being great among men, should there by Christ be singularly glorified to be the assessors of His judgment-seat, even they whom nothing here could turn aside from being the followers of His footsteps. Nevertheless, let there be no man that believeth that the twelve Apostles only, among whom Matthias holdeth that place from which Judas by transgression fell, (Acts i. 25,) that they only shall judge, even as the twelve tribes of Israel shall not alone be judged for then were the tribe of Levi, which is the thirteenth, unjudged.\par
\switchcolumn*

\LA
\begin{Resp}\Rsign Sint lumbi vestri præcíncti, et lucérnæ ardéntes in mánibus vestris: \Star{} Et vos símiles homínibus exspectántibus dóminum suum, quando revertátur a núptiis.\end{Resp}
\begin{Resp}\Vsign Vigiláte ergo, quia nescítis qua hora Dóminus vester ventúrus sit.\end{Resp}
\begin{Resp}\Rsign Et vos símiles homínibus exspectántibus dóminum suum, quando revertátur a núptiis.\end{Resp}
\begin{Resp}\Vsign Glória Patri, et Fílio, \Star{} et Spirítui Sancto.\end{Resp}
\begin{Resp}\Rsign Et vos símiles homínibus exspectántibus dóminum suum, quando revertátur a núptiis.\end{Resp}
\switchcolumn
\EN
\begin{Resp}\Rsign Let your loins be girded about, and your lights burning; \Star{} And ye yourselves like unto men that wait for their lord, when he will return from the wedding.\end{Resp}
\begin{Resp}\Vsign Watch therefore, for ye know not what hour your Lord doth come.\end{Resp}
\begin{Resp}\Rsign And ye yourselves like unto men that wait for their lord, when he will return from the wedding.\end{Resp}
\begin{Resp}\Vsign Glory be to the Father, and to the Son, \Star{} and to the Holy Ghost.\end{Resp}
\begin{Resp}\Rsign And ye yourselves like unto men that wait for their lord, when he will return from the wedding.\end{Resp}
\switchcolumn*

\LA
\LessonHead{Lectio IX}
\switchcolumn
\EN
\LessonHead{Lesson IX}
\switchcolumn*

\LA
\noindent Et Paulus, qui tertiusdécimus est Apóstolus, judicándi sorte privábitur? cum ipse dicat: Nescítis quóniam ángelos judicábimus, quanto magis sæculária? Sciéndum namque est, omnes, qui ad exémplum Apostolórum, sua reliquérunt ómnia et secúti sunt Christum, júdices cum eo ventúros, sicut étiam omne mortálium genus esse judicándum. Quia enim duodenário sæpe número solet in Scriptúris univérsitas designári, per duódecim sedes Apostolórum ómnium numerósitas judicántium, et, per duódecim tribus Israël, univérsitas eórum qui judicándi sunt, osténditur.\par
\switchcolumn
\EN
\noindent Moreover, then, were Paul, who is the thirteenth Apostle, deprived of all part in the judgment; whereas he saith of himself: Know ye not that we shall judge angels? How much more things that pertain to this life? But it behoveth us to know that every one who hath forsaken all and followed Christ, as did the Apostles, shall also come with Him to judgment, even as every man shall stand at His judgment seat. And the Scriptures use often to signify all by this number twelve; by the twelve thrones of the Apostles are signified the thrones of all them that shall judge; and by the twelve tribes of Israel, the whole number of them that shall be judged.\par
\switchcolumn*

\LA
\TeDeum{Te Deum}
\switchcolumn
\EN
\TeDeum{Te Deum}
\switchcolumn*

\end{paracol}

\par\needspace{10\baselineskip}\addvspace{1.2em}{\centering\small\textsc{Die 8 Februarii} · \textsc{8 February}\par}
\Office{S. Joannis de Matha Confessoris}{St. John of Matha, Confessor}{Duplex}
\addcontentsline{toc}{subsection}{Die 8 Februarii · S. Joannis de Matha Confessoris\enspace\textit{St. John of Matha, Confessor}}
\begin{paracol}{2}
\LA
\RefLine{Lectiones I–III: de Scriptura occurrente.}
\switchcolumn
\EN
\RefLine{Lessons I–III: from the Scripture of the day.}
\switchcolumn*

\LA
\LessonHead{Lectio IV}
\switchcolumn
\EN
\LessonHead{Lesson IV}
\switchcolumn*

\LA
\noindent Joánnes de Matha, ordinis sanctíssimæ Trinitátis redemptiónis captivórum institutor, Falcone in Provincia natus est paréntibus pietáte et nobilitate conspicuis. Studiórum causa Aquas Sextias, mox Parisios profectus, confectóque theologíæ curriculo, magisterii lauream adeptus, doctrinæ et virtútum splendore enituit. Quibus motus Parisiénsis antístes, ad sacrum presbyterátus ordinem præ humilitate reluctántem promovit, eo consílio, ut in ea civitáte cómmorans, sapiéntia et móribus studiosæ juventúti prælucéret. Cum autem in sacello ejusdem epíscopi, ipso cum aliis astante, primum Deo Sacrum offerret, cælésti favore meruit recreari. Nam Angelus cándida et fulgénti veste indútus, cui in péctore crux rúbei et cærúlei coloris assuta erat, brácchiis cancellátis et super duos captivos ad látera positos, Christianum unum, álterum Maurum, exténsis, appáruit. Qua visióne in éxtasim raptus, intelléxit prótinus vir Dei, se ad redimendos ab infidelibus captivos destinari.\par
\switchcolumn
\EN
\noindent John de la Mata, the founder of the Order of the Most Holy Trinity for the Ransom of Prisoners, was born at Faucon, in Provence, (upon Midsummer's Day, in the year 1169,) and was the child of parents equally distinguished for their rank and their godly life. He went for his education first to Aix and then to Paris. At the University of Paris, where he went through the course of Divinity and took the degree of Doctor, he became eminent for learning and virtue. For this reason the Bishop of Paris ordained him Priest, an honour from which his lowliness caused him to shrink, in the hope that he should induce him to remain at Paris, and be a bright example of wisdom and manners to the students who resorted thither. He offered up the Holy Sacrifice to God for the first time in the private Chapel of the Bishop, and in the presence of that Prelate and diverse other persons. In the midst of the ceremony, a vision from God appeared to John. There appeared to him an angel, clad in raiment white and glistering; having sewn on his breast a cross of red and blue. His arms were crossed before him, and his hands were upon the heads of two slaves, one a Christian and the other a Moor. And immediately the man of God was in the spirit, and knew that he was called to the work of ransoming bondsmen from the power of the unbelievers.\par
\switchcolumn*

\LA
\begin{Resp}\Rsign Honéstum fecit illum Dóminus, et custodívit eum ab inimícis, et a seductóribus tutávit illum: \Star{} Et dedit illi claritátem ætérnam.\end{Resp}
\begin{Resp}\Vsign Justum dedúxit Dóminus per vias rectas, et osténdit illi regnum Dei.\end{Resp}
\begin{Resp}\Rsign Et dedit illi claritátem ætérnam.\end{Resp}
\switchcolumn
\EN
\begin{Resp}\Rsign The Lord made him honourable, and defended him from his enemies, and kept him safe from those that lay in wait for him; \Star{} And gave him perpetual glory.\end{Resp}
\begin{Resp}\Vsign He went down with him into the pit, and left him not in bonds.\end{Resp}
\begin{Resp}\Rsign And gave him perpetual glory.\end{Resp}
\switchcolumn*

\LA
\LessonHead{Lectio V}
\switchcolumn
\EN
\LessonHead{Lesson V}
\switchcolumn*

\LA
\noindent Quo vero matúrius in re tanti moménti procéderet, in solitúdinem secessit, ibique divino nutu factum est, ut Felicem Valesium in ipsa erémo jam multis annis degentem repérerit: cum quo ínita societáte, se per triennium in oratióne, et contemplatióne, omniúmque virtútum studio exercuit. Cóntigit autem, ut, dum secum de rebus divinis prope fontem colloqueréntur, cervus ad eos accesserit, crucem inter córnua gerens rúbei et cærúlei coloris. Cumque Felix ob rei novitátem mirarétur, narrávit ei Joánnes visiónem in prima Missa hábitam; et exínde ferventius oratióni incumbéntes, ter in somnis admoniti, Romam proficísci decrevérunt, ut a summo Pontifice novi ordinis pro rediméndis captivis institutiónem impetrarent. Electus fuerat eo témpore Innocentius tertius; qui, illis benígne acceptis, dum secum de re proposita deliberat, in festo sanctæ Agnetis secundo, Laterani intra Missárum solemnia ad sacræ Hostiæ elevatiónem, Angelus ei cándida veste, cruce bicolori, spécie redimentis captivos appáruit. Quo viso, Pontifex institutum approbávit, et novum ordinem sanctíssimæ Trinitátis redemptiónis captivórum vocari jussit, ejusque professóribus albas vestes cum cruce rúbei et cærúlei coloris præbuit.\par
\switchcolumn
\EN
\noindent That he might set himself with due forethought to the carrying out of his work, he withdrew into a certain desert, and there, by the will of God, he found Felix de Valois, who had already spent many years in that place. With him he joined company, and they passed three years together in continual prayer, meditation, and all spiritual exercises. It came to pass, one day, when they were sitting on the bank of a spring, that there came to them a stag having between his horns a cross of red and blue. Felix cried out in wonder at that sight, and John then told him of the vision that had appeared to him when he was saying his first Mass. Thenceforth they gave themselves with redoubled fervour to prayer, and, being three times warned in sleep, they determined to go to Rome, and pray the Pope to institute an Order for the ransom of prisoners. They arrived at the time of the election of Innocent III., who received them courteously, and entertained in his mind their petition. While he was in consideration, he went to the Lateran Cathedral, on the second Feast of St. Agnes, and there, while Mass was being solemnly sung, at the moment of the elevation of the Sacred Host, there appeared to him an angel, clad in raiment white and glistering, having sewn on his breast a cross of red and blue, and making as though he would free prisoners. Thereupon the Pope founded the Order, commanding that it should be called the Order of the Most Holy Trinity for the Ransom of prisoners, and that they who professed in it should be clad in white raiment, having sewn on their breasts a cross of red and blue.\par
\switchcolumn*

\LA
\begin{Resp}\Rsign Amávit eum Dóminus, et ornávit eum: stolam glóriæ índuit eum, \Star{} Et ad portas paradísi coronávit eum.\end{Resp}
\begin{Resp}\Vsign Índuit eum Dóminus lorícam fídei, et ornávit eum.\end{Resp}
\begin{Resp}\Rsign Et ad portas paradísi coronávit eum.\end{Resp}
\switchcolumn
\EN
\begin{Resp}\Rsign The Lord loved him and beautified him; He clothed him with a robe of glory, \Star{} And crowned him at the gates of Paradise.\end{Resp}
\begin{Resp}\Vsign The Lord hath put on him the breast-plate of faith, and hath adorned him.\end{Resp}
\begin{Resp}\Rsign And crowned him at the gates of Paradise.\end{Resp}
\switchcolumn*

\LA
\LessonHead{Lectio VI}
\switchcolumn
\EN
\LessonHead{Lesson VI}
\switchcolumn*

\LA
\noindent Sic stabilito ordine, sancti fundatores in Gálliam rediérunt, primoque cœnobio Cervi Frígidi in diœcesi Meldénsi constructo, ad ejus regimen Felix remánsit; et Joánnes Romam cum aliquot sociis reversus est, ubi Innocentius domum, ecclésiam et hospitale sancti Thomæ de Formis in monte Cælio eis donávit cum multis redítibus et possessiónibus. Datis quoque litteris ad Miramolínum regem Marochii, opus redemptiónis felici auspicio inchoátum fuit. Tum ad Hispánias, sub jugo Saracenórum magna ex parte oppressas, Joánnes profectus est, regumque, príncipum atque aliórum fidelium animos ad captivórum et páuperum commiseratiónem commovit. Monasteria ædificávit, hospitalia eréxit, magnoque lucro animárum plures captivos redémit. Romam tandem reversus, sanctisque opéribus incumbens, assiduis labóribus attritus et morbo confectus, ardentíssimo Dei et próximi amóre exæstuans, ad extremum devenit. Quare frátribus convocatis, eisque ad opus redemptiónis cælitus præmonstrátum efficaciter cohortátis, obdormívit in Dómino sextodecimo Kalendas Januarii, anno salútis millesimo ducentésimo décimo tertio; ejusque corpus in ipsa ecclésia sancti Thomæ de Formis condigno honóre tumulátum fuit.\par
\switchcolumn
\EN
\noindent The Order being thus established, the holy Founders returned into France, and built their first Convent at Cerfroid, in the diocese of Meaux. Felix remained in charge of this house, and John went back to Rome with several companions. To them Innocent gave the house, Church, and hospital of St. Thomas de Formis on the Coelian Mount, with great endowments and property. Moreover he gave them a letter of introduction to Miramolin, King of Morocco, and they began with bright hopes the work of ransoming prisoners. John next betook himself to Spain, great part of which was then in the hands of the Saracens, and stirred up the hearts of the kings, princes, and all the faithful to have pity on slaves and the poor. He built Convents, founded Hospitals, and ransomed many bondsmen, to the great gain of souls. At last he returned to Rome, still busied in good works, but worn out by unceasing toil, and weakened by sickness. As he drew near the end of his earthly pilgrimage, his burning love for God and for his neighbour suffered no diminution. He called together his brethren, and earnestly exhorted them to go on with that work of ransom which had been pointed out to them from heaven, and then fell asleep in the Lord, on the 21st day of December, 1213. His body was buried with due honour in the Church of St. Thomas de Formis.\par
\switchcolumn*

\LA
\begin{Resp}\Rsign Iste homo perfécit ómnia quæ locútus est ei Deus, et dixit ad eum: Ingrédere in réquiem meam: \Star{} Quia te vidi justum coram me ex ómnibus géntibus.\end{Resp}
\begin{Resp}\Vsign Iste est, qui contémpsit vitam mundi, et pervénit ad cæléstia regna.\end{Resp}
\begin{Resp}\Rsign Quia te vidi justum coram me ex ómnibus géntibus.\end{Resp}
\begin{Resp}\Vsign Glória Patri, et Fílio, \Star{} et Spirítui Sancto.\end{Resp}
\begin{Resp}\Rsign Quia te vidi justum coram me ex ómnibus géntibus.\end{Resp}
\switchcolumn
\EN
\begin{Resp}\Rsign This is he which did according unto all that God commanded him; and God said unto him: Enter thou into My rest, \Star{} For thee have I seen righteous before Me among all people.\end{Resp}
\begin{Resp}\Vsign This is he which loved not his life in this world, and is come unto an everlasting kingdom.\end{Resp}
\begin{Resp}\Rsign For thee have I seen righteous before Me among all people.\end{Resp}
\begin{Resp}\Vsign Glory be to the Father, and to the Son, \Star{} and to the Holy Ghost.\end{Resp}
\begin{Resp}\Rsign For thee have I seen righteous before Me among all people.\end{Resp}
\switchcolumn*

\LA
\LessonHead{Lectio VII}
\switchcolumn
\EN
\LessonHead{Lesson VII}
\switchcolumn*

\LA
\LessonTitle{Léctio sancti Evangélii secúndum Lucam}
\switchcolumn
\EN
\LessonTitle{From the Holy Gospel according to Luke}
\switchcolumn*

\LA
\Cite{Luc 12:35-40}
\switchcolumn
\EN
\Cite{Luke 12:35-40}
\switchcolumn*

\LA
\noindent In illo témpore: Dixit Jesus discípulis suis: Sint lumbi vestri præcincti, et lucernæ ardentes in mánibus vestris. Et réliqua.\par
\switchcolumn
\EN
\noindent At that time, Jesus said unto His disciples: Let your loins be girded about, and your lights burning. And so on.\par
\switchcolumn*

\LA
\LessonTitle{Homilía sancti Gregórii Papæ}
\switchcolumn
\EN
\LessonTitle{Homily by Pope St. Gregory (the Great.)}
\switchcolumn*

\LA
\Source{Homilia 13 in Evang.}
\switchcolumn
\EN
\Source{13th on the Gospels.}
\switchcolumn*

\LA
\noindent Sancti Evangélii, fratres caríssimi, apérta vobis est léctio recitata. Sed ne alíquibus ipsa ejus planíties alta fortásse videátur, eam sub brevitáte transcúrrimus, quátenus ejus exposítio ita nesciéntibus fiat cógnita, ut tamen sciéntibus non sit onerósa. Dóminus dicit: Sint lumbi vestri præcíncti. Lumbos enim præcíngimus, cum carnis luxúriam per continéntiam coarctámus. Sed quia minus est, mala non ágere, nisi étiam quisque stúdeat, et bonis opéribus insudáre, prótinus ádditur: Et lucérnæ ardéntes in mánibus vestris. Lucérnas quippe ardéntes in mánibus tenémus, cum per bona ópera próximis nostris lucis exémpla monstrámus. De quibus profécto opéribus Dóminus dicit: Lúceat lux vestra coram homínibus, ut vídeant ópera vestra bona, et gloríficent Patrem vestrum qui in cælis est.\par
\switchcolumn
\EN
\noindent Dearly beloved brethren, the words of the Holy Gospel, which have just been read, lie open before you, and, lest their very plainness should make them seem to some to be hard, we will go through them with such shortness as that neither may they which understand not remain unenlightened, nor they which understand be wearied. The Lord saith Let your loins be girded about. Now, we gird our loins about, when by continency we master the lustful inclination of the flesh. But, forasmuch as it sufficeth not for a man to abstain from evil deeds, if he strive not to join thereto the earnest doing of good works, it is immediately added And your lights burning. Our lights burn when, by good works, we give bright example to our neighbour; concerning which works the Lord saith Let your light so shine before men, that they may see your good works, and glorify your Father Which is in heaven. (Matth. v. 16.)\par
\switchcolumn*

\LA
\begin{Resp}\Rsign Iste est qui ante Deum magnas virtútes operátus est, et de omni corde suo laudávit Dóminum: \Star{} Ipse intercédat pro peccátis ómnium populórum.\end{Resp}
\begin{Resp}\Vsign Ecce homo sine queréla, verus Dei cultor, ábstinens se ab omni ópere malo, et pérmanens in innocéntia sua.\end{Resp}
\begin{Resp}\Rsign Ipse intercédat pro peccátis ómnium populórum.\end{Resp}
\switchcolumn
\EN
\begin{Resp}\Rsign This is he which wrought great wonders before God, and praised the Lord with all his heart. \Star{} May he pray for all people, that their sins may be forgiven unto them.\end{Resp}
\begin{Resp}\Vsign Behold a man without blame, a worshipper of God in truth, keeping himself clean from every evil work, and abiding still in his innocency.\end{Resp}
\begin{Resp}\Rsign May he pray for all people, that their sins may be forgiven unto them.\end{Resp}
\switchcolumn*

\LA
\LessonHead{Lectio VIII}
\switchcolumn
\EN
\LessonHead{Lesson VIII}
\switchcolumn*

\LA
\noindent Duo autem sunt, quæ jubéntur, et lumbos restríngere, et lucérnas tenére: ut et mundítia sit castitátis in córpore, et lumen veritátis in operatióne. Redemptóri étenim nostro unum sine áltero placére nequáquam potest: si aut is qui bona agit, adhuc luxúriæ inquinaménta non déserit: aut is qui castitáte præéminet, necdum se per bona ópera exércet. Nec cástitas ergo magna est sine bono ópere, nec opus bonum est áliquod sine castitáte. Sed et si utrúmque ágitur, restat, ut quisquis ille est, spe ad supérnam pátriam tendat, et nequáquam se a vítiis pro mundi hujus honestáte contíneat.\par
\switchcolumn
\EN
\noindent Here, then, are two commandments, to gird our loins about, and to keep our lights burning the cleanness of purity in our body, and the light of the truth in our works. Whoso hath the one and not the other, pleaseth not thereby our Redeemer; that is, he pleaseth Him not which doth good works, but bridleth not himself from the pollutions of lust, neither he which is eminent in chastity, but exerciseth not himself in good works. Neither is chastity a great thing without good works, nor good works anything without chastity. And if any man do both, it remaineth that he must look by hope toward our Fatherland above, and not have for his reason wherethrough he turneth himself away from vice, the love of honour in this present world.\par
\switchcolumn*

\LA
\begin{Resp}\Rsign Sint lumbi vestri præcíncti, et lucérnæ ardéntes in mánibus vestris: \Star{} Et vos símiles homínibus exspectántibus dóminum suum, quando revertátur a núptiis.\end{Resp}
\begin{Resp}\Vsign Vigiláte ergo, quia nescítis qua hora Dóminus vester ventúrus sit.\end{Resp}
\begin{Resp}\Rsign Et vos símiles homínibus exspectántibus dóminum suum, quando revertátur a núptiis.\end{Resp}
\begin{Resp}\Vsign Glória Patri, et Fílio, \Star{} et Spirítui Sancto.\end{Resp}
\begin{Resp}\Rsign Et vos símiles homínibus exspectántibus dóminum suum, quando revertátur a núptiis.\end{Resp}
\switchcolumn
\EN
\begin{Resp}\Rsign Let your loins be girded about, and your lights burning; \Star{} And ye yourselves like unto men that wait for their lord, when he will return from the wedding.\end{Resp}
\begin{Resp}\Vsign Watch therefore, for ye know not what hour your Lord doth come.\end{Resp}
\begin{Resp}\Rsign And ye yourselves like unto men that wait for their lord, when he will return from the wedding.\end{Resp}
\begin{Resp}\Vsign Glory be to the Father, and to the Son, \Star{} and to the Holy Ghost.\end{Resp}
\begin{Resp}\Rsign And ye yourselves like unto men that wait for their lord, when he will return from the wedding.\end{Resp}
\switchcolumn*

\LA
\LessonHead{Lectio IX}
\switchcolumn
\EN
\LessonHead{Lesson IX}
\switchcolumn*

\LA
\noindent Et vos símiles homínibus exspectántibus dóminum suum, quando revertátur a núptiis: ut cum vénerit et pulsáverit, conféstim apériant ei. Venit quippe Dóminus, cum ad judícium próperat: pulsat vero, cum jam per ægritúdinis moléstias esse mortem vicínam desígnat. Cui conféstim aperímus, si hunc cum amóre suscípimus. Aperíre enim júdici pulsánti non vult, qui exíre de córpore trépidat: et vidére eum, quem contempsísse se méminit, júdicem formídat. Qui autem de sua spe et operatióne secúrus est, pulsanti conféstim áperit, quia lætus júdicem sústinet; et, cum tempus propínquæ mortis advénerit, de glória retributiónis hilaréscit.\par
\switchcolumn
\EN
\noindent And ye yourselves like unto men that wait for their lord, when he will return from the wedding that, when he cometh and knocketh, they may open unto him immediately. The Lord cometh at the hour of judgment He knocketh when, by the pains of sickness, He biddeth us know that death is nigh. To Him open we immediately, if we receive Him in love. Whoso feareth to leave this body, will not open to the Judge when He knocketh, for he dreadeth to see that Judge, Whom he knoweth that he hath despised. But whosoever knoweth that his hope and works are built upon a good foundation, when he heareth the Judge knock, openeth to Him immediately, for to such an one that coming is blessed, yea, when the hour of death is at hand, such an one haileth with gladness a glorious reward.\par
\switchcolumn*

\LA
\TeDeum{Te Deum}
\switchcolumn
\EN
\TeDeum{Te Deum}
\switchcolumn*

\end{paracol}

\par\needspace{10\baselineskip}\addvspace{1.2em}{\centering\small\textsc{Die 9 Februarii} · \textsc{9 February}\par}
\Office{S. Cyrilli Episc. Alexandrini Confessoris et Ecclesiæ Doctoris}{St. Cyril of Alexandria, Bishop, Confessor, and Doctor of the Church}{Duplex}
\addcontentsline{toc}{subsection}{Die 9 Februarii · S. Cyrilli Episc. Alexandrini Confessoris et Ecclesiæ Doctoris\enspace\textit{St. Cyril of Alexandria, Bishop, Confessor, and Doctor of the Church}}
\begin{paracol}{2}
\LA
\RefLine{Lectiones I–III: de Scriptura occurrente.}
\switchcolumn
\EN
\RefLine{Lessons I–III: from the Scripture of the day.}
\switchcolumn*

\LA
\RefLine{Lectiones VII–VIII: ex Commúni Doctoris Pontificis (C4a).}
\switchcolumn
\EN
\RefLine{Lessons VII–VIII: from the Common of a Doctor Bishop (C4a).}
\switchcolumn*

\LA
\RefLine{Lectio IX: de S. Apolloniæ Virginis et Martyris (02-09o).}
\switchcolumn
\EN
\RefLine{Lesson IX: of S. Apollonia, Virgin and Martyr (02-09o).}
\switchcolumn*

\LA
\LessonHead{Lectio IV}
\switchcolumn
\EN
\LessonHead{Lesson IV}
\switchcolumn*

\LA
\noindent Cyríllus Alexandrínus, cujus præcónia non uníus tantum vel altérius sunt comprobáta testimónio, sed étiam œcumenicórum conciliórum Ephesíni et Chalcedonénsis actis celebráta, claris ortus paréntibus, ac Theóphili epíscopi Alexandríni nepos, adhuc adoléscens præcelléntis ingénii clara specímina dedit. Lítteris ac sciéntiis egrégie imbútus, ad Joánnem epíscopum, Jerosolymitánum se cóntulit, ut in Christiána fide perficerétur. Alexandríam deínde cum rediísset, Theóphilo vita functo, ad illíus sedem evéctus est: quo in múnere ita óptimi pastóris formam ab Apóstolo definítam constánter præ se tulit, ut sanctíssimi præsulis glóriam mérito sit adéptus.\par
\switchcolumn
\EN
\noindent The praises of Cyril of Alexandria have been celebrated not only by one writer or another, but have even been registered in the acts of the Ecumenical Councils of Ephesus and Chalcedon. He was born of distinguished parents, and was the nephew of Theophilus, Pope of Alexandria. While he was still young he displayed marks of his excellent understanding. After giving a deep study to letters and science he betook himself to John, Bishop of Jerusalem, to be perfected in the Christian faith. After his return to Alexandria, and the death of Theophilus, he was raised to that see. In this office he kept ever before his eyes the type of the Shepherd of souls as it had been laid down by the Apostle; and by ever adhering thereto deservedly earned the glory of an holy Bishop.\par
\switchcolumn*

\LA
\begin{Resp}\Rsign Invéni David servum meum, óleo sancto meo unxi eum: \Star{} Manus enim mea auxiliábitur ei.\end{Resp}
\begin{Resp}\Vsign Nihil profíciet inimícus in eo, et fílius iniquitátis non nocébit ei.\end{Resp}
\begin{Resp}\Rsign Manus enim mea auxiliábitur ei.\end{Resp}
\switchcolumn
\EN
\begin{Resp}\Rsign I have found David My servant, with My holy oil have I anointed him \Star{} For My hand shall help him.\end{Resp}
\begin{Resp}\Vsign The enemy shall prevail nothing against him, nor the son of wickedness afflict him.\end{Resp}
\begin{Resp}\Rsign For My hand shall help him.\end{Resp}
\switchcolumn*

\LA
\LessonHead{Lectio V}
\switchcolumn
\EN
\LessonHead{Lesson V}
\switchcolumn*

\LA
\noindent Salútis animárum zelo incénsus curas omnes inténdit, ut sibi commíssum gregem in fídei et morum integritáte serváret, atque a venenátis infidélium et hæreticórum páscuis defénderet. Hinc tum Nováti ásseclas e civitáte expélli, tum Judæos, qui furóre acti in cædem Christianórum conspiráverant, juxta leges puníri satégit. Singuláre vero Cyrílli pro cathólicæ fídei incolumitáte enítuit stúdium contra Nestórium Constantinopolitánum epíscopum, asseréntem Jesum Christum ex María Vírgine hóminem tantum et non Deum natum, eíque divinitátem pro méritis esse collátam; cujus emendatiónem cum frustra tentásset, eum sancto Cælestíno Pontífici máximo denuntiávit.\par
\switchcolumn
\EN
\noindent Zeal for the salvation of souls was kindled in him, and he undertook all care to keep in the faith and in soundness of life the flock unto him committed, and to preserve them from the poisonous pastures of infidelity and heresy; hence, in accordance with the laws, he caused the followers of Novatus to be expelled from the city, and those Jews to be punished who had been induced by rage to plan a massacre of the Christians. His eminent care for the preservation of the Catholic faith pure and undenled shone forth especially in his controversy against Nestorius, Patriarch of Constantinople, who asserted that Jesus Christ had been born of the Virgin Mary as man only and not as God, and that the Godhead had been bestowed upon Him because of His merits. Cyril first attempted to convert Nestorius, but when he found this hopeless he denounced him to the Supreme Pontiff the holy Celestine.\par
\switchcolumn*

\LA
\begin{Resp}\Rsign Pósui adjutórium super poténtem, et exaltávi eléctum de plebe mea: \Star{} Manus enim mea auxiliábitur ei.\end{Resp}
\begin{Resp}\Vsign Invéni David servum meum, óleo sancto meo unxi eum.\end{Resp}
\begin{Resp}\Rsign Manus enim mea auxiliábitur ei.\end{Resp}
\switchcolumn
\EN
\begin{Resp}\Rsign I have laid help upon one that is mighty, and have exalted one chosen out of My people \Star{} For My hand shall help him.\end{Resp}
\begin{Resp}\Vsign I have found David My servant, with My holy oil have I anointed him.\end{Resp}
\begin{Resp}\Rsign For My hand shall help him.\end{Resp}
\switchcolumn*

\LA
\LessonHead{Lectio VI}
\switchcolumn
\EN
\LessonHead{Lesson VI}
\switchcolumn*

\LA
\noindent Cælestíni delegáta auctoritáte, concílio Ephesíno præfuit, in quo hæresis Nestoriána pénitus proscrípta est, damnátus Nestórius et a sua Sede dejéctus, ac dogma cathólicum de una in Christo eáque divína persóna, et divína gloriósæ Vírginis Maríæ maternitáte assértum; plaudénte pópulo univérso, qui incredíbili gáudio géstiens, collucéntibus fácibus domum dedúxit epíscopus. Sed hac de causa Cyríllus calúmniis, injúriis et persecutiónibus, plúrimis a Nestório ejúsque fautóribus impetítus fuit; quas ipse patientíssime tulit, ita ut, de sola fide sollícitus, quidquid advérsus eum effutiébant ac moliebántur hærétici, pro níhilo habéret. Tandem pro Ecclésia Dei máximis perfúnctus labóribus, plurimísque scriptis éditis tum ad éthnicos et hæréticos confutándos, tum ad sacras Scriptúras et cathólica explanánda dógmata, sancto fine quiévit anno quadringentésimo quadragésimo quarto, episcopátus trigésimo secúndo. Leo décimus tértius Póntifex máximus Offícium et Missam præclaríssimi hujus fídei cathólicæ propugnatóris et Orientális ecclésiæ lúminis, ad Ecclésiam univérsam exténdit.\par
\switchcolumn
\EN
\noindent As delegate of Pope Celestine, Cyril presided at the Council of Ephesus where the Nestorian heresy was condemned; Nestorius deprived of his see; and the Catholic doctrine as to the unity of Person in Christ and the divine Motherhood of the glorious Virgin Mary was laid down amid the rejoicings of all the people, who escorted the bishops to their lodgings with a torch-light procession. For this reason Nestorius and his followers made Cyril the object of slanders, insults, and persecutions which he bore with profound patience, having all his care for the purity of the faith, and taking no heed to what the heretics might say or try against him. At length he died a holy death, in the year of salvation 444 and of his own papacy the 32nd. After vast work for the Church of God, and leaving behind him diverse writings directed either against heathens and heretics or to the exposition of the holy Scriptures and of Catholic doctrine, the Supreme Pontiff Leo XIII. extended to the Universal Church the Office and Mass of this most eminent champion of the Catholic faith, and light of the Eastern Church.\par
\switchcolumn*

\LA
\begin{Resp}\Rsign Iste est, qui ante Deum magnas virtútes operátus est, et omnis terra doctrína ejus repléta est: \Star{} Ipse intercédat pro peccátis ómnium populórum.\end{Resp}
\begin{Resp}\Vsign Iste est, qui contémpsit vitam mundi, et pervénit ad cæléstia regna.\end{Resp}
\begin{Resp}\Rsign Ipse intercédat pro peccátis ómnium populórum.\end{Resp}
\begin{Resp}\Vsign Glória Patri, et Fílio, \Star{} et Spirítui Sancto.\end{Resp}
\begin{Resp}\Rsign Ipse intercédat pro peccátis ómnium populórum.\end{Resp}
\switchcolumn
\EN
\begin{Resp}\Rsign This is he which wrought great wonders before God, and the whole earth is full of his teaching \Star{} May he pray for all people, that their sins may be forgiven unto them\end{Resp}
\begin{Resp}\Vsign This is he which loved not his life in this world, and hath attained unto the kingdom of heaven.\end{Resp}
\begin{Resp}\Rsign May he pray for all people, that their sins may be forgiven unto them\end{Resp}
\begin{Resp}\Vsign Glory be to the Father, and to the Son, \Star{} and to the Holy Ghost.\end{Resp}
\begin{Resp}\Rsign May he pray for all people, that their sins may be forgiven unto them\end{Resp}
\switchcolumn*

\end{paracol}

\par\needspace{10\baselineskip}\addvspace{1.2em}{\centering\small\textsc{Die 9 Februarii} · \textsc{9 February}\par}
\Office{S. Apolloniæ Virginis et Martyris}{S. Apollonia, Virgin and Martyr}{Simplex}
\addcontentsline{toc}{subsection}{Die 9 Februarii · S. Apolloniæ Virginis et Martyris\enspace\textit{S. Apollonia, Virgin and Martyr}}
\begin{paracol}{2}
\LA
\LessonHead{Lectio IX (pro commemoratione)}
\switchcolumn
\EN
\LessonHead{Lesson IX (when commemorated)}
\switchcolumn*

\LA
\LessonTitle{Commemoratio: S. Apolloniæ Virginis et Martyris}
\switchcolumn
\EN
\LessonTitle{Commemoration of: S. Apollonia, Virgin and Martyr}
\switchcolumn*

\LA
\noindent Apollónia virgo Alexandrína, sub Décio imperatóre, cum ingravescénte jam ætáte ad idóla sisterétur, ut eis veneratiónem adhibéret; illis contémptis, Jesum Christum verum Deum coléndum esse prædicábat. Quam ob rem omnes ei contúsi sunt et evúlsi dentes; ac, nisi Christum detestáta deos cóleret, accénso rogo combustúros vivam mináti sunt ímpii carnífices. Quibus illa, se quamvis mortem pro Jesu Christi fide subitúram, respóndit. Itaque comprehénsa ut comburerétur, cum paulísper, quasi delíberans quid agéndum esset, stetísset, ex illórum mánibus elápsa, álacris in ignem sibi parátum, majóri Spíritus Sancti flamma intus accénsa, se injécit. Unde brevi, consúmpto córpore, puríssimus spíritus in cælum ad sempitérnam martyrii corónam evolávit.\par
\switchcolumn
\EN
\noindent Apollonia was an aged virgin of Alexandria, who, (in the year of salvation 249,) in the reign of the Emperor Decius, was brought before the idols to worship them, but refused, declaring that Christ Jesus is True God, and that to Him worship is due. The cruel executioners beat and pulled out all her teeth, and threatened to burn her alive if she would not deny Christ. To whom she answered, that for Christ Jesus' sake she was ready to die. Being taken to the place of execution she stood for a few moments as if in doubt, and then, the fire of the Holy Ghost burning up in her heart, she broke from those that held her, and leapt of her own accord into the flames. Her body was quickly consumed, and her soul departed pure to obtain the eternal crown of martyrdom.\par
\switchcolumn*

\LA
\TeDeum{Te Deum}
\switchcolumn
\EN
\TeDeum{Te Deum}
\switchcolumn*

\end{paracol}

\par\needspace{10\baselineskip}\addvspace{1.2em}{\centering\small\textsc{Die 10 Februarii} · \textsc{10 February}\par}
\Office{S. Scholasticæ Virginis}{St. Scholastica, Virgin}{Duplex}
\addcontentsline{toc}{subsection}{Die 10 Februarii · S. Scholasticæ Virginis\enspace\textit{St. Scholastica, Virgin}}
\begin{paracol}{2}
\LA
\RefLine{Lectiones I–III: de Scriptura occurrente.}
\switchcolumn
\EN
\RefLine{Lessons I–III: from the Scripture of the day.}
\switchcolumn*

\LA
\RefLine{Lectiones VII–IX: ex Commúni Virginum tantum (C6a).}
\switchcolumn
\EN
\RefLine{Lessons VII–IX: from the Common of a Virgin not a Martyr (C6a).}
\switchcolumn*

\LA
\LessonHead{Lectio IV}
\switchcolumn
\EN
\LessonHead{Lesson IV}
\switchcolumn*

\LA
\LessonTitle{Ex libro Dialogórum sancti Gregórii Papæ}
\switchcolumn
\EN
\LessonTitle{From the Second Book of the Dialogues of Pope St. Gregory (the Great.)}
\switchcolumn*

\LA
\Source{Lib. 2, Cap. 33.}
\switchcolumn
\EN
\Source{Ch. 33}
\switchcolumn*

\LA
\noindent Scholástica, venerábilis patris Benedícti soror, omnipoténti Dómino ab ipso infántiæ témpore dedicáta, ad eum semel per annum veníre consuéverat; ad quam vir Dei non longe extra jánuam in possessióne monastérii descendébat. Quadam vero die venit ex more, atque ad eam, cum discípulis venerábilis ejus descéndit frater; qui totum diem in Dei láudibus sacrísque collóquiis ducéntes, incumbéntibus jam noctis ténebris, simul accepérunt cibum. Cumque adhuc ad mensam sedérent, et inter sacra collóquia tárdior se hora protráheret, éadem sanctimoniális fémina soror ejus eum rogávit, dicens: Quæso te, ut ista nocte me non déseras, ut usque mane de cæléstis vitæ gáudiis loquámur. Cui ille respóndit: Quid est quod lóqueris, soror? manére extra cellam nullátenus possum. Tanta vero erat cæli serénitas, ut nulla in áëre nubes apparéret. Sanctimoniális autem fémina, cum verba fratris negántis pósuit, et caput in mánibus omnipoténtem Dóminum rogatúra declinávit. Cumque leváret de mensa caput, tanta coruscatiónis et tonítrui virtus, tántaque inundátio plúviæ erúpit, ut neque venerábilis Benedíctus, neque fratres, qui cum eo áderant, extra loci limen, quo conséderant, pedem movére potúerint.\par
\switchcolumn
\EN
\noindent The worshipful Scholastica, the sister of our Father Benedict, was hallowed unto the Lord Almighty from a child. Her custom was to come to see her brother once every year. And when she came, the man of God went down unto her, not far from the gate, but, as it were, within the borders of his monastery. And there was a day when she came, as her custom was, and her worshipful brother went down to her, and his disciples with him. Then they passed the whole day together, praising God, and speaking one to the other of spiritual things. And when the night came, they brake bread together. And while they were yet at table, and conversed together on spiritual things, the hour was late. Then the holy woman his sister besought him, saying Leave me not, I pray thee, this night, but let us speak even until morning of the gladness of the eternal life. He answered her: What is it that thou sayest, my sister? I can by no means remain out of my cell. Now the firmament was so clear that there were no clouds in the sky. Then the holy nun, when she had heard the words of her brother, that he would not abide with her, clasped her hands on the table, and laid her face on her hands, and besought the Lord Almighty. And it came to pass that when she lifted up her head from the table, there were great thunderings and lightnings, and a flood of rain, insomuch that neither the worshipful Benedict nor the brethren that were with him could move as much as a foot over the threshold of the place where they sat.\par
\switchcolumn*

\LA
\begin{Resp}\Rsign Propter veritátem, et mansuetúdinem, et justítiam: \Star{} Et dedúcet te mirabíliter déxtera tua.\end{Resp}
\begin{Resp}\Vsign Spécie tua et pulchritúdine tua inténde, próspere procéde, et regna.\end{Resp}
\begin{Resp}\Rsign Et dedúcet te mirabíliter déxtera tua.\end{Resp}
\switchcolumn
\EN
\begin{Resp}\Rsign Because of truth, and meekness, and righteousness; \Star{} And thy right hand shall lead thee wonderfully.\end{Resp}
\begin{Resp}\Vsign In thy comeliness, and thy beauty, go forward, fare prosperously, and reign.\end{Resp}
\begin{Resp}\Rsign And thy right hand shall lead thee wonderfully.\end{Resp}
\switchcolumn*

\LA
\LessonHead{Lectio V}
\switchcolumn
\EN
\LessonHead{Lesson V}
\switchcolumn*

\LA
\noindent Sanctimoniális quippe fémina caput in mánibus declínans, lacrimárum flúvium in mensam fúderat, per quas serenitátem áëris ad plúviam traxit. Nec paulo tárdius post oratiónem inundátio illa secúta est; sed tanta fuit conveniéntia oratiónis et inundatiónis, ut de mensa caput jam cum tonítruo leváret; quátenus unum idémque esset moméntum, et leváre caput, et plúviam depónere. Tunc vir Dei inter corúscos et tonítruos atque ingéntis plúviæ inundatiónem, videns se ad monastérium non posse remeáre, cœpit cónqueri contristátus, dicens: Parcat tibi omnípotens Deus, soror: quid est quod fecísti? Cui illa respóndit: Ecce rogávi te et audíre me noluísti; rogávi Deum meum, et audívit me: modo ergo, si potes, egrédere et, me dimíssa, ad monastérium recéde. Ipse autem exíre extra tectum non valens, qui remanére sponte nóluit in loco, mansit invítus. Sicque factum est, ut totam noctem pervígilem dúcerent, atque per sacra spiritális vitæ collóquia, sese vicária relatióne satiárent.\par
\switchcolumn
\EN
\noindent Now when the holy woman laid her head in her hands upon the table, she wept bitterly, and as she wept, the clearness of the sky was turned to a tempest. As she prayed, immediately the flood followed. And the time was so, that she lifted up her head when it thundered, and when she had lifted up her head, the rain came. When the man of God saw that he could not return to his monastery, because of the lightnings, and thunderings, and the great rain, he was sorrowful and grieved, saying Almighty God forgive thee, my sister; what is this that thou hast done? She answered him Behold, I besought thee, and thou wouldest not hear; I besought my God, and He hath heard me; if, therefore, thou wilt, go forth, leave me alone, and go thy way to thy monastery. But he could not, and so he tarried in the same place, not willingly, but of necessity. And so it came to pass that they slept not all that night, but fed one another with discourse on spiritual things.\par
\switchcolumn*

\LA
\begin{Resp}\Rsign Dilexísti justítiam, et odísti iniquitátem: \Star{} Proptérea unxit te Deus, Deus tuus, óleo lætítiæ.\end{Resp}
\begin{Resp}\Vsign Propter veritátem, et mansuetúdinem, et justítiam.\end{Resp}
\begin{Resp}\Rsign Proptérea unxit te Deus, Deus tuus, óleo lætítiæ.\end{Resp}
\switchcolumn
\EN
\begin{Resp}\Rsign Thou hast loved righteousness, and hated iniquity; \Star{} Therefore God, thy God, hath anointed thee with the oil of gladness.\end{Resp}
\begin{Resp}\Vsign Because of truth, and meekness, and righteousness.\end{Resp}
\begin{Resp}\Rsign Therefore God, thy God, hath anointed thee with the oil of gladness.\end{Resp}
\switchcolumn*

\LA
\LessonHead{Lectio VI}
\switchcolumn
\EN
\LessonHead{Lesson VI}
\switchcolumn*

\LA
\noindent Cumque die áltero éadem venerábilis fémina ad cellam própriam recessísset, vir Dei ad monastérium rédiit. Cum ecce post tríduum, in cella consístens, elevátis in áëra óculis vidit ejúsdem soróris suæ ánimam de córpore egréssam in colúmbæ spécie cæli secréta penetráre. Qui, tantæ ejus glóriæ congáudens, omnipoténti Deo in hymnis et láudibus grátias réddidit, ejúsque óbitum frátribus denuntiávit. Quos étiam prótinus misit, ut ejus corpus ad monastérium deférrent, atque in sepúlcro, quod sibi ipsi paráverat, pónerent. Quo facto cóntigit, ut, quorum mens una semper in Deo fúerat, eórum quoque córpora nec sepultúra separáret.\par
\switchcolumn
\EN
\noindent And when the morning was come, the worshipful woman arose, and went unto her own cell, and the man of God went back to his monastery. And, behold, after three days he was sitting in his cell, and he lifted up his eyes to heaven, and saw the soul of his sister, delivered from the body, fly to heaven in a bodily shape like a dove. Wherefore he rejoiced because of the glory that was revealed in her, and gave thanks to Almighty God in hymns and praises, and made known to the brethren that she was dead. He commanded them also to go and take up her body, and bring it to his monastery, and lay it in the grave which he had made ready for himself. Whereby it came to pass that they twain who had ever been of one mind in the Lord, even in death were not divided.\par
\switchcolumn*

\LA
\begin{Resp}\Rsign Afferéntur Regi vírgines post eam, próximæ ejus; \Star{} Afferéntur tibi in lætítia et exsultatióne.\end{Resp}
\begin{Resp}\Vsign Spécie tua et pulchritúdine tua; inténde, próspere procéde, et regna.\end{Resp}
\begin{Resp}\Rsign Afferéntur tibi in lætítia et exsultatióne.\end{Resp}
\begin{Resp}\Vsign Glória Patri, et Fílio, \Star{} et Spirítui Sancto.\end{Resp}
\begin{Resp}\Rsign Afferéntur tibi in lætítia et exsultatióne.\end{Resp}
\switchcolumn
\EN
\begin{Resp}\Rsign After her shall virgins be brought unto the King, her fellows \Star{} They shall be brought unto thee with gladness and rejoicing.\end{Resp}
\begin{Resp}\Vsign In thy comeliness and thy beauty, go forward, fare prosperously, and reign.\end{Resp}
\begin{Resp}\Rsign They shall be brought unto thee with gladness and rejoicing.\end{Resp}
\begin{Resp}\Vsign Glory be to the Father, and to the Son, \Star{} and to the Holy Ghost.\end{Resp}
\begin{Resp}\Rsign They shall be brought unto thee with gladness and rejoicing.\end{Resp}
\switchcolumn*

\end{paracol}

\par\needspace{10\baselineskip}\addvspace{1.2em}{\centering\small\textsc{Die 11 Februarii} · \textsc{11 February}\par}
\Office{In Apparitione Beatæ Mariæ Virginis Immaculatæ}{In Apparitione Beatæ Mariæ Virginis Immaculatæ}{Duplex majus}
\addcontentsline{toc}{subsection}{Die 11 Februarii · In Apparitione Beatæ Mariæ Virginis Immaculatæ\enspace\textit{In Apparitione Beatæ Mariæ Virginis Immaculatæ}}
\begin{paracol}{2}
\LA
\LessonHead{Lectio I}
\switchcolumn
\EN
\LessonHead{Lesson I}
\switchcolumn*

\LA
\LessonTitle{De Parábolis Salomónis}
\switchcolumn
\EN
\LessonTitle{Lesson from the book of Proverbs}
\switchcolumn*

\LA
\Cite{Prov 8:12-17}
\switchcolumn
\EN
\Cite{Prov 8:12-17}
\switchcolumn*

\LA
\noindent Ego sapiéntia hábito in consílio et erudítis intérsum cogitatiónibus. Timor Dómini odit malum: arrogántiam, et supérbiam, et viam pravam, et os bilíngue detéstor. Meum est consílium et ǽquitas, mea est prudéntia, mea est fortitúdo. Per me reges regnant, et legum conditóres justa decérnunt; per me príncipes ímperant, et poténtes decérnunt justítiam. Ego diligéntes me díligo; et qui mane vígilant ad me, invénient me.\par
\switchcolumn
\EN
\noindent I wisdom dwell in counsel, and am present in learned thoughts. The fear of the Lord hateth evil: I hate arrogance, and pride, and every wicked way, and a mouth with a double tongue. Counsel and equity is mine, prudence is mine, strength is mine. By me kings reign, and lawgivers decree just things, By me princes rule, and the mighty decree justice. I love them that love me: and they that in the morning early watch for me, shall find me.\par
\switchcolumn*

\LA
\begin{Resp}\Rsign Sapiéntia quæ attíngit a fine usque ad finem fórtiter, et dispónit ómnia suáviter, ædificávit sibi domum: \Star{} Ecce tabernáculum Dei cum homínibus.\end{Resp}
\begin{Resp}\Vsign Vidi sanctam civitátem Jerúsalem novam, parátam sicut sponsam ornátam viro suo.\end{Resp}
\begin{Resp}\Rsign Ecce tabernáculum Dei cum homínibus.\end{Resp}
\switchcolumn
\EN
\begin{Resp}\Rsign Wisdom which reacheth from end to end mightily, and ordereth all things sweetly, hath built herself a house: \Star{} Behold the tabernacle of God with men.\end{Resp}
\begin{Resp}\Vsign I saw the holy city, the new Jerusalem, prepared as a bride adorned for her husband.\end{Resp}
\begin{Resp}\Rsign Behold the tabernacle of God with men.\end{Resp}
\switchcolumn*

\LA
\LessonHead{Lectio II}
\switchcolumn
\EN
\LessonHead{Lesson II}
\switchcolumn*

\LA
\Cite{Prov 8:18-25}
\switchcolumn
\EN
\Cite{Prov 8:18-25}
\switchcolumn*

\LA
\noindent Mecum sunt divítiæ et glória, opes supérbæ et justítia. Mélior est enim fructus meus auro et lápide pretióso, et genímina mea argénto elécto. In viis justítiæ ámbulo, in médio semitárum judícii, ut ditem diligéntes me et thesáuros eórum répleam. Dóminus possidébit me in inítio viárum suárum, ántequam quidquam fáceret a princípio. Ab ætérno ordináta sum et ex antíquis, ántequam terra fíeret. Nondum erant abýssi, et ego jam concépta eram; necdum fontes aquárum erúperant, necdum montes gravi mole constíterant; ante colles ego parturiébar.\par
\switchcolumn
\EN
\noindent With me are riches and glory, glorious riches and justice. For my fruit is better than gold and the precious stone, and my blossoms than choice silver. I walk in the way of justice, in the midst of the paths of judgment, That I may enrich them that love me, and may fill their treasures. The Lord possessed me in the beginning of his ways, before he made any thing from the beginning. I was set up from eternity, and of old before the earth was made. The depths were not as yet, and I was already conceived. neither had the fountains of waters as yet sprung out: The mountains with their huge bulk had not as yet been established: before the hills I was brought forth:\par
\switchcolumn*

\LA
\begin{Resp}\Rsign Quasi arcus refúlgens inter nébulas, et quasi flos rosárum in diébus vernis, et quasi lília in tránsitu aquæ, \Star{} Sic fulget Virgo immaculáta.\end{Resp}
\begin{Resp}\Vsign Arcum meum ponam in núbibus, et erit signum fœ́deris mei vobíscum.\end{Resp}
\begin{Resp}\Rsign Sic fulget Virgo immaculáta.\end{Resp}
\switchcolumn
\EN
\begin{Resp}\Rsign As the rainbow giving light in the clouds, and as the flower of roses in the days of the spring, and as the lilies on the brink of the water, \Star{} Thus shineth the Virgin immaculate.\end{Resp}
\begin{Resp}\Vsign I will set my bow in the clouds, and it shall be a sign of my covenant with you.\end{Resp}
\begin{Resp}\Rsign Thus shineth the Virgin immaculate.\end{Resp}
\switchcolumn*

\LA
\LessonHead{Lectio III}
\switchcolumn
\EN
\LessonHead{Lesson III}
\switchcolumn*

\LA
\Cite{Prov 8:34-36; 9:1-5}
\switchcolumn
\EN
\Cite{Prov 8:34-36; 9:1-5}
\switchcolumn*

\LA
\noindent Beátus homo qui audit me, et qui vígilat ad fores meas cotídie, et obsérvat ad postes óstii mei. Qui me invénerit, invéniet vitam, et háuriet salútem a Dómino; qui autem in me peccáverit, lædet ánimam suam. Omnes, qui me odérunt, díligunt mortem. Sapiéntia ædificávit sibi domum, excídit colúmnas septem. Immolávit víctimas suas, míscuit vinum et propósuit mensam suam. Misit ancíllas suas ut vocárent ad arcem et ad mœ́nia civitátis: Si quis est párvulus, véniat ad me. Et insipiéntibus locúta est: Veníte, comédite panem meum, et bíbite vinum quod míscui vobis.\par
\switchcolumn
\EN
\noindent Blessed is the man that heareth me, and that watcheth daily at my gates, and waiteth at the posts of my doors. He that shall find me, shall find life, and shall have salvation from the Lord: But he that shall sin against me, shall hurt his own soul. All that hate me love death. Wisdom hath built herself a house, she hath hewn her out seven pillars. She hath slain her victims, mingled her wine, and set forth her table. She hath sent her maids to invite to the tower, and to the walls of the city: Whosoever is a little one, let him come to me. And to the unwise she said: Come, eat my bread, and drink the wine which I have mingled for you.\par
\switchcolumn*

\LA
\begin{Resp}\Rsign Surge, amíca mea, speciósa mea, et veni, colúmba mea: \Star{} Osténde mihi fáciem tuam, sonet vox tua in áuribus meis.\end{Resp}
\begin{Resp}\Vsign Vox túrturis audíta est in terra nostra.\end{Resp}
\begin{Resp}\Rsign Osténde mihi fáciem tuam, sonet vox tua in áuribus meis.\end{Resp}
\begin{Resp}\Vsign Glória Patri, et Fílio, \Star{} et Spirítui Sancto.\end{Resp}
\begin{Resp}\Rsign Osténde mihi fáciem tuam, sonet vox tua in áuribus meis.\end{Resp}
\switchcolumn
\EN
\begin{Resp}\Rsign Arise, my love, my beautiful one, and come, my dove: \Star{} Shew me thy face, let thy voice sound in my ears.\end{Resp}
\begin{Resp}\Vsign The voice of the turtle is heard in our land.\end{Resp}
\begin{Resp}\Rsign Shew me thy face, let thy voice sound in my ears\end{Resp}
\begin{Resp}\Vsign Glory be to the Father, and to the Son, \Star{} and to the Holy Ghost.\end{Resp}
\begin{Resp}\Rsign Shew me thy face, let thy voice sound in my ears\end{Resp}
\switchcolumn*

\LA
\LessonHead{Lectio IV}
\switchcolumn
\EN
\LessonHead{Lesson IV}
\switchcolumn*

\LA
\noindent Anno quarto a dogmática definitióne de immaculáto beátæ Vírginis Concéptu, ad Gavi flúminis oram prope óppidum Lourdes diœcésis Tarbiénsis in Gállia, ipsa Virgo in rupis sinu super specum Massabiélle puéllæ cuidam, vernácula lingua Bernadétte nuncupátæ, paupérrimæ quidem sed ingénuæ ac piæ, plúries se conspiciéndam óbtulit. Immaculáta Virgo juveníli ac benígno videbátur aspéctu, nívea veste niveóque pállio contécta, ac zona cærúlea succíncta; nudos pedes áurea rosa ornábat. Primo apparitiónis die, qui fuit undécimus Februárii anno millésimo octingentésimo quinquagésimo octávo, puéllam signum crucis rite piéque faciéndum edócuit, atque ad sacri rosárii recitatiónem, exémplo suo, corónam, quæ prius ex brácchio demíssa pendébat, manu advólvens, excitávit: quod in céteris étiam apparitiónibus prǽstitit. Altero autem apparitiónis die, puélla in simplicitáte cordis sui, diabólicam fraudem timens, lustrálem aquam in Vírginem effúdit; sed beáta Virgo, léniter arrídens, benigniórem illi vultum osténdit. Cum vero tértio apparuísset, puéllam ad specum per quíndecim dies invitávit. Exínde eam sǽpius est allocúta, ac pro peccatóribus oráre, terram deosculári, pœnitentiámque ágere est hortáta; deínde imperávit, ut sacerdótibus edíceret, ædificándum ibi esse sacéllum, solemnísque supplicatiónibus more illo accedéndum. Mandávit ínsuper ut e fonte, qui sub aréna adhuc latébat sed mox erat eruptúrus, aquam bíberet, eáque se abstérgeret. Dénique die festo Annuntiatiónis, percontánti eníxe puéllæ illíus nomen, cujus aspéctu tóties dignáta fúerat, Virgo, admótis péctori mánibus elatísque in cælum óculis, respóndit: Immaculáta Concéptio ego sum.\par
\switchcolumn
\EN
\noindent In the fourth year from the dogmatic definition of the Immaculate Conception of the Blessed Virgin, at the bank of the river Gave near the town of Lourdes of the diocese of Tarbes in France, the Virgin herself in a bend of the rock above the grotto of Massabiele, often shewed herself to a certain girl, called in the vernacular tongue Bernadette, indeed most poor but noble and pious, to be seen. The Immaculate Virgin appeared with a young and kind appearance, clothed with a white garment and white veil, and girt with a blue girdle; she adorned her bare feet with a golden rose. On the first day of the apparition, which was the eleventh of February in the one thousand eight hundred fifty-eighth year, she taught the girl the sign of the cross to be duly and piously made, and she incited her, by her example, to the recitation of the sacred rosary, turning over with her hand the chaplet, which before was hanging down from her arm: which she supplied also in the other apparitions. And on the second day of the apparition, the girl in the simplicity of her heart, fearing diabolic fraud, flung holy water on the Virgin; but the blessed Virgin, smiling gently, shewed her face more kindly to her. And when she appeared a third time, she invited the girl to the grotto for fifteen days. Thence she often addressed her; then commanded, that it might be declared to the priests, a chapel to be built there, and to be approached there for supplications in a manner of solemnity. Moreover she commanded that from the spring, which thus far was hidden under the sand but now was about to erupt, she might drink the water, and by it cleanse herself. Finally on the feast day of the Annunciation, to the girl earnestly inquiring her name, the Virgin, for whose appearance she was so often deemed worthy, the Virgin, her hands having been moved to her breast and her eyes raised to heaven, answered: I am the Immaculate Conception.\par
\switchcolumn*

\LA
\begin{Resp}\Rsign Quæ est ista, quæ progréditur quasi auróra consúrgens, \Star{} Pulchra ut luna, elécta ut sol?\end{Resp}
\begin{Resp}\Vsign Ipsa est colúmba mea, perfécta mea, immaculáta mea.\end{Resp}
\begin{Resp}\Rsign Pulchra ut luna, elécta ut sol.\end{Resp}
\switchcolumn
\EN
\begin{Resp}\Rsign Who is she, that cometh forth as the morning rising, \Star{} Fair as the moon, bright as the sun?\end{Resp}
\begin{Resp}\Vsign She is my dove, my perfect one, my undefiled.\end{Resp}
\begin{Resp}\Rsign Fair as the moon, bright as the sun?\end{Resp}
\switchcolumn*

\LA
\LessonHead{Lectio V}
\switchcolumn
\EN
\LessonHead{Lesson V}
\switchcolumn*

\LA
\noindent Percrebrescénte fama beneficiórum, quæ in sacro specu recepísse fidéles dicebántur, augebátur in dies hóminum concúrsus, quos loci relígio ad specum advocábat. Itaque prodigiórum fama puellǽque candóre mortus Tarbiénsis epíscopus, quarto ab enarrátis anno, post jurídiam factórum inquisitiónem, supernaturáles esse apparitiónis notas sua senténtia probávit, cultúmque Vírginis immaculátæ in eódem specu permísit. Mox ædificátum sacéllum: ex illa die pene innúmeræ fidélium turbæ, voti ac supplicatiónis causa, ex Gállia, Bélgio, Itália, Hispánia ceterísque Európæ provínciis necnon ex longínquis Américæ regiónibus quovis anno illuc advéniunt, noménque Immaculátæ de Lourdes ubíque terrárum inclaréscit. Fontis aqua, in cunctas orbis partes deláta, ægris sanitátem restítuit. Orbis vero cathólicus tantórum memor benefactórum, ædes sacras mirábili ópere ibi exstrúxit. Vexílla innúmera, acceptórum beneficiórum véluti monuménta, illuc a civitátibus ac géntibus missa, ædem Vírginis miro ornátu décorant. In hac sua véluti sede Immaculáta Virgo júgiter cólitur; intérdiu quidem précibus, religióso cantu solemnibúsque áliis cæremóniis; noctu vero sacris illis supplicatiónibus, quibus infinítæ propémodum peregrinántium turbæ céreis facibúsque accénsis procédunt et laudes beátæ Vírginis cóncinunt.\par
\switchcolumn
\EN
\noindent The fame of favours becoming widespread, which the faithful were said to have received in the sacred grotto, in time the assembly of men increased, whom the sanctity of the place called to the grotto. Therefore the bishop of Tarbes moved by the fame of the wonders and the purity of the girl, after a juridical inquiry of the happenings, approved by his decree the signs of the apparition to be supernatural, and permitted the worship of the immaculate Virgin in the same grotto. Soon a chapel is built: from that day the almost innumerable crowds of the faithful, for the purpose of prayer and supplication, from France, Belgium, Italy, Spain, and the other provinces of Europe and also from the remote regions of America arrived there in every year, and the name of the Immaculate of Lourdes becometh famous of countries everywhere. The water of the spring, having been carried into all parts of the world, restored health to the sick. And the Catholic world mindful of such great favours, built holy shrines there with wonderful work. Innumerable banners, as if monuments of received benefits, sent there by cities and nations, decorate the shrine of the Virgin with wonderful adornment. In this as if her own seat the immaculate Virgin is continually honoured: indeed by day with prayers, religious chant and other solemn ceremonies; and by night with those solemn supplications, by which the almost infinite crowds of pilgrims proceed with candles and lit torches and sing the praises of the blessed Virgin.\par
\switchcolumn*

\LA
\begin{Resp}\Rsign Erit in novíssimis diébus mons præparátus Vírgini Maríæ in vértice móntium, et elevábitur super cælos, et ibunt pópuli multi et dicent: \Star{} Veníte et ascendámus ad montem.\end{Resp}
\begin{Resp}\Vsign Sicut lætántium ómnium habitátio est in te.\end{Resp}
\begin{Resp}\Rsign Veníte et ascendámus ad montem.\end{Resp}
\switchcolumn
\EN
\begin{Resp}\Rsign In the last days the mountain a house for the Virgin Mary shall be prepared on the top of mountains, and it shall be exalted above the heavens, and many people shall go and say: \Star{} Come and let us go up to the mountain.\end{Resp}
\begin{Resp}\Vsign The dwelling in thee is as it were of all rejoicing.\end{Resp}
\begin{Resp}\Rsign Come and let us go up to the mountain.\end{Resp}
\switchcolumn*

\LA
\LessonHead{Lectio VI}
\switchcolumn
\EN
\LessonHead{Lesson VI}
\switchcolumn*

\LA
\noindent Peregrinatiónes hujúsmodi fidem, frigescénte sǽculo, excitásse, ánimum ad christiánam legem profiténdam addidísse, cultúmque Vírginis immaculátæ mirum in modum auxísse, ómnibus compértum est. In qua mirábili fídei professióne christiánus pópulus sacerdótes véluti duces habet, qui illuc suas plebes addúcunt. Ipse étiam Sacrórum antístites sanctum locum frequénter ádeunt, peregrinatiónibus præsunt, solemnioribúsque festis intérsunt. Nec ádeo rarum est ipsos Románæ Ecclésiæ purpurátos patres húmili peregrinórum more accedéntes conspícere. Ipsi quoque Románi Pontífices, pro sua erga Immaculátam de Lourdes pietáte, sacram ædem donis nobilíssimis cumulárunt. Pius nonus, sacris indulgéntiis, archiconfraternitátis privilégio ac minóris Basílicæ título ipsam insignívit; ac Deíparæ imáginem íbidem cultam, solémni ritu per legátum suum apostólicum in Gállia, diadémate distínctam vóluit. Leo vero décimus tértius innúmera étiam cóntulit benefícia, indulgéntias ad modum jubilǽi vigésimo quinto Apparitiónis anno verténte concéssit, peregrinatiónes sua auctoritáte verbóque provéxit, ac solémnem Ecclésiæ sub título Rosárii dedicatiónem suo nómine péragi curávit. Quorum beneficiórum amplitúdinem cumulávit, cum, plúrium episcopórum rogátu, solémne festum sub título Apparitiónis beátæ Maríæ Vírginis immaculátæ, próprio Offício et própria Missa celebrándum benígne concéssit. Tandem Pius décimus Póntifex máximus pro sua erga Deíparam pietáte, ac plurimórum votis ánnuens Sacrórum antístitum, idem festum ad Ecclésiam univérsam exténdit.\par
\switchcolumn
\EN
\noindent It is known by all for certain the pilgrimages of this sort to have stirred up faith, in a world becoming cold, the soul to have increased to profess the Christian law, and to have spread the worship of the immaculate Virgin in a wonderful manner. In which wonderful profession of faith the Christian people have the priests as leaders, who lead their people thence. The holy Bishops themselves also frequently visit the holy place, lead pilgrimages, and take part in more solemn feasts. Nor indeed is it rare to observe these empurpled fathers of the Roman Church approaching humble in the manner of pilgrims. Even these Roman Pontiffs, for their piety towards the Immaculate of Lourdes, enhanced the sacred shrine with most noble gifts. Pius the ninth, with sacred indulgences, distinguished it with the privilege of an archconfraternity and title of a minor Basilica; and willed the image honoured there of the Godbearer, with solemn rite through his apostolic legate in France, to be adorned with a crown. And Leo the thirteenth also conferred innumerable benefits, granted indulgences in the manner of a jubilee in the twenty-fifth year of the Apparition occurring, promoted pilgrimages by his authority and word, and arranged to complete by his name the solemn dedication of the Church under the title of the Rosary. He accumulated the extent of which benefits, when, with the asking of many bishops, he kindly granted a solemn feast to be celebrated under the title of the Apparition of the blessed Virgin Mary immaculate with a proper Office and proper Mass. Finally supreme Pontiff Pius the tenth for his piety towards the Godbearer, and agreeing to the wishes of many holy Bishops, extended the same feast to the universal Church.\par
\switchcolumn*

\LA
\begin{Resp}\Rsign Prævenísti eam, Dómine, in benedictiónibus dulcédinis, posuísti in cápite ejus \Star{} Corónam de lápide pretióso.\end{Resp}
\begin{Resp}\Vsign Magna est glória ejus in salutári tuo, glóriam et magnum decórem impónes super eam.\end{Resp}
\begin{Resp}\Rsign Corónam de lápide pretióso.\end{Resp}
\begin{Resp}\Vsign Glória Patri, et Fílio, \Star{} et Spirítui Sancto.\end{Resp}
\begin{Resp}\Rsign Corónam de lápide pretióso.\end{Resp}
\switchcolumn
\EN
\begin{Resp}\Rsign Thou hast prevented her, O Lord, with blessings of sweetness, thou hast set on her head \Star{} A crown of precious stones.\end{Resp}
\begin{Resp}\Vsign Her glory is great in thy salvation, glory and great beauty shalt thou lay upon her.\end{Resp}
\begin{Resp}\Rsign A crown of precious stones.\end{Resp}
\begin{Resp}\Vsign Glory be to the Father, and to the Son, \Star{} and to the Holy Ghost.\end{Resp}
\begin{Resp}\Rsign A crown of precious stones.\end{Resp}
\switchcolumn*

\LA
\LessonHead{Lectio VII}
\switchcolumn
\EN
\LessonHead{Lesson VII}
\switchcolumn*

\LA
\LessonTitle{Léctio sancti Evangélii secúndum Lucam}
\switchcolumn
\EN
\LessonTitle{From the Holy Gospel according to Luke}
\switchcolumn*

\LA
\Cite{Luc 1:26-28}
\switchcolumn
\EN
\Cite{Luke 1:26-28}
\switchcolumn*

\LA
\noindent In illo témpore: Missus est Angelus Gábriel a Deo in civitátem Galilææ, cui nomen Názareth, ad Vírginem desponsátam viro, cui nomen erat Joseph, de domo David, et nomen Vírginis María. Et réliqua.\par
\switchcolumn
\EN
\noindent In that time the angel Gabriel was sent from God into a city of Galilee, called Nazareth, To a virgin espoused to a man whose name was Joseph, of the house of David; and the virgin's name was Mary. And so on.\par
\switchcolumn*

\LA
\LessonTitle{Homilía sancti Bernárdi Abbátis}
\switchcolumn
\EN
\LessonTitle{Homily by St. Bernard, Abbot (of Clairvaux.)}
\switchcolumn*

\LA
\Source{Ex Homilía 2 super Missus}
\switchcolumn
\EN
\Source{2nd on this text.}
\switchcolumn*

\LA
\noindent Lætáre, pater Adam, sed magis tu, o Heva mater, exsúlta, qui sicut ómnium paréntes, ita ómnium fuístis peremptóres; et quod infelícius est, prius peremptóres quam paréntes. Ambo, inquam, consolámini super fília, et tali fília, sed illa ámplius de qua malum ortum est prius, cujus oppróbrium in omnes pertransívit mulíeres. Instat namque tempus, quo jam tollátur oppróbrium, nec hábeat vir quid causétur advérsus féminam: qui útique, dum se imprudénter excusáre conarétur, crudéliter illam accusáre non cunctátus est, dicens: Múlier quam dedísti mihi, dedit mihi de ligno, et comédi. Proptérea curre Heva ad Maríam; curre mater ad fíliam; fília pro matre respóndeat: ipsa matris oppróbrium áuferat; ipsa patri pro matre satisfáciat: quia ecce si vir cécidit per féminam, jam non erígitur nisi per féminam.\par
\switchcolumn
\EN
\noindent Rejoice, father Adam, and yet more thou mother Eve, ye that are the source of all, and the ruin of all, and the unhappy cause of their ruin before ye gave them birth. Be comforted both in your daughter, and such a daughter; but chiefly thou, O woman, of whom the first evil came, and who hast cast thy slur upon all women. The time is come for the slur to be taken away, and for the man to have nothing to say against the woman. At the first, when he unwisely began to make excuse, he scrupled not to throw the blame upon her, saying, The woman whom Thou gavest to be with me, she gave me of the tree, and I did eat. Wherefore, O Eve, betake thyself to Mary Mother, betake thyself to thy daughter let the daughter answer for the mother let her take away her mother's reproach; let her make up to her father for her mother's fault for if man be fallen by means of woman, it is by means of woman that he is raised up again.\par
\switchcolumn*

\LA
\begin{Resp}\Rsign Tu ergo ínvoca Dóminum, lóquere Regi pro nobis, \Star{} Et líbera nos de morte.\end{Resp}
\begin{Resp}\Vsign Omnes sitiéntes, veníte ad aquas, et hauriétis salútem a Dómino.\end{Resp}
\begin{Resp}\Rsign Et líbera nos de morte.\end{Resp}
\switchcolumn
\EN
\begin{Resp}\Rsign Therefore do thou call upon the Lord, speak to the King for us, \Star{} And deliver us from death.\end{Resp}
\begin{Resp}\Vsign All you that thirst, come to the waters, and draw salvation from the Lord.\end{Resp}
\begin{Resp}\Rsign And deliver us from death.\end{Resp}
\switchcolumn*

\LA
\LessonHead{Lectio VIII}
\switchcolumn
\EN
\LessonHead{Lesson VIII}
\switchcolumn*

\LA
\noindent Quid dicébas, o Adam? Múlier quam dedísti mihi, dedit mihi de ligno, et comédi. Verba malítiæ sunt hæc quibus magis áugeas quam déleas culpam. Verúmtamen Sapiéntia vicit malítiam, cum occasiónem véniæ, quam a te Deus interrogándo elícere tentávit, sed non pótuit, in thesáuro indeficiéntis suæ pietátis invénit. Rédditur nempe fémina pro fémina, prudens pro fátua, húmilis pro supérba, quæ pro ligno mortis gustum tibi pórrigat vitæ, et pro venenóso cibo illo amaritúdinis, dulcédinem páriat fructus ætérni. Muta ergo iníquæ excusatiónis verbum in vocem gratiárum actiónis, et dic: Dómine, múlier, quam dedísti mihi, dedit mihi de ligno vitæ, et comédi; et dulce factum est super mel ori meo, quia in ipso vivificásti me. Ecce enim ad hoc missus est Angelus ad Vírginem. O admirándam et omni honóre digníssimam Vírginem; o féminam singuláriter venerándam, super omnes féminas admirábilem, paréntum reparatrícem, posterórum vivificatrícem.\par
\switchcolumn
\EN
\noindent What didst thou say, O Adam? The woman whom Thou gavest to be with me, she gave me of the tree, and I did eat. These are wrathful words, by the which thou dost rather magnify than diminish thine offence. Nevertheless, Wisdom hath defeated thy malice. God asked thee that He might find in thee an occasion of pardon, but, in that He found it not, He hath sought and found it in the Treasure of His Own mercy. One woman answereth for another; the wise for the foolish; the lowly for the proud; for her that gave thee of the tree of death, another that giveth thee to taste of the tree of life; for her that brought thee the bitter food of sin, another that giveth thee of the sweet fruits of righteousness. Wherefore accuse the woman no more, but speak in thanksgiving, and say, Lord, the woman whom Thou hast given me, she hath given me of the tree of life, and I have eaten; and it is in my mouth sweeter than honey, for thereby hast Thou quickened me. (Ps. cxviii. 103, 93.) Behold, it was for this that the angel Gabriel was sent to the Virgin, to the most worshipful of women, a woman more wonderful than all women, the restorer of them that went before, and the quickener of them that come after her.\par
\switchcolumn*

\LA
\begin{Resp}\Rsign Plantávit Dóminus Deus paradísum voluptátis, produxítque lignum vitæ in médio ejus: \Star{} Et flúvius egrediebátur de loco voluptátis.\end{Resp}
\begin{Resp}\Vsign Emissiónes tuæ paradisus, Virgo María.\end{Resp}
\begin{Resp}\Rsign Et flúvius egrediebátur de loco voluptátis.\end{Resp}
\begin{Resp}\Vsign Glória Patri, et Fílio, \Star{} et Spirítui Sancto.\end{Resp}
\begin{Resp}\Rsign Et flúvius egrediebátur de loco voluptátis.\end{Resp}
\switchcolumn
\EN
\begin{Resp}\Rsign The Lord God planted a paradise of pleasure, and brought forth the tree of life in its midst: \Star{} And a river went out of the place of pleasure.\end{Resp}
\begin{Resp}\Vsign Thy plants are a paradise, O Virgin Mary.\end{Resp}
\begin{Resp}\Rsign And a river went out of the place of pleasure.\end{Resp}
\begin{Resp}\Vsign Glory be to the Father, and to the Son, \Star{} and to the Holy Ghost.\end{Resp}
\begin{Resp}\Rsign And a river went out of the place of pleasure.\end{Resp}
\switchcolumn*

\LA
\LessonHead{Lectio IX}
\switchcolumn
\EN
\LessonHead{Lesson IX}
\switchcolumn*

\LA
\noindent Quam tibi áliam prædixísse Deus vidétur, quando ad serpéntem ait: Inimicítias ponam inter te et mulíerem? Et si adhuc dúbitas, quod de María díxerit, audi quod séquitur: Ipsa cónteret caput tuum. Cui hæc serváta victória est, nisi Maríæ? Ipsa procul dúbio caput contrívit venenátum, quæ omnimódam malígni suggestiónem tam de carnis illécebra, quam de mentis supérbia dedúxit ad níhilum. Quam vero áliam Sálomon requirébat, cum dicébat: Mulíerem fortem quis invéniet? Nóverat quippe vir sápiens hujus sexus infirmitátem, frágile corpus, lúbricam mentem. Quia tamen et Deum légerat promisísse, et ita vidébat congrúere, ut qui vícerat per féminam, vincerétur per ipsam, veheménter admírans ajébat: Mulíerem fortem quis invéniet? Quod est dícere: si ita de manu féminæ pendet et nostra ómnium salus, et innocéntiæ restitútio, et de hoste victória: fortis omníno necésse est ut provideátur, quæ ad tantum opus possit esse idónea.\par
\switchcolumn
\EN
\noindent Has it not of this thy daughter, O Adam, that God spake when He said unto the serpent, I will put enmity between thee and the woman And if thou wilt still doubt that He speaketh of Mary, hear what followeth She shall bruise thy head. Who won this conquest but Mary? She brought to nought the whole wiles of Satan, whether for the pollution of her body or the injury of her soul. Was it not of her that Solomon spake, where he saith, Who shall find a virtuous woman? (Prov. xxxi. 10.) The wise man knew the weaknesses of women, how frail they are in body, and how changeable in mind. But he had read that God had promised that the enemy, who had prevailed by means of a woman, was by a woman to be overthrown, and he believed. But he wondered greatly, and said, Who shall find a virtuous woman? that is to say If our salvation, and the bringing back of that which is lost, and the final triumph over the enemy, is in the hand of a woman, it must needs be that a virtuous woman be found, meet to work in that matter\par
\switchcolumn*

\LA
\TeDeum{Te Deum}
\switchcolumn
\EN
\TeDeum{Te Deum}
\switchcolumn*

\end{paracol}

\par\needspace{10\baselineskip}\addvspace{1.2em}{\centering\small\textsc{Die 12 Februarii} · \textsc{12 February}\par}
\Office{Ss. Septem Fundatorum Ordinis Servorum B. M. V.}{Seven Holy Founders of the Order of the Servants of the Blessed Virgin Mary}{Duplex}
\addcontentsline{toc}{subsection}{Die 12 Februarii · Ss. Septem Fundatorum Ordinis Servorum B. M. V.\enspace\textit{Seven Holy Founders of the Order of the Servants of the Blessed Virgin Mary}}
\begin{paracol}{2}
\LA
\RefLine{Lectiones I–III: de Scriptura occurrente.}
\switchcolumn
\EN
\RefLine{Lessons I–III: from the Scripture of the day.}
\switchcolumn*

\LA
\LessonHead{Lectio IV}
\switchcolumn
\EN
\LessonHead{Lesson IV}
\switchcolumn*

\LA
\noindent Sǽculo tértio décimo cum Frideríci secúndi diro schísmate cruentísque factiónibus cultióres Itáliæ pópuli scinderéntur, próvidens Dei misericórdia, præter álios sanctitáte illústres, septem e Florentína nobilitáte viros suscitávit, qui in caritáte conjúncti, præclárum fratérnæ dilectiónis præbérent exémplum. Hi, nimírum Bonfílius Monáldius, Bonajúncta Manéttus, Manéttus Antellénsis, Amidéus de Amidéis, Ugúccio Ugucciónum, Sostenéus de Sostenéis et Aléxius Falconérius, cum, anno trigésimo tértio ejus sǽculi, die sacra Vírgini cælo recéptæ, in quodam piórum hóminum convéntu, Laudántium nuncupáto, fervéntius orárent; ab eádem Deípara síngulis apparénte sunt admóniti, ut sánctius perfectiúsque vitæ genus amplecteréntur. Re ítaque prius cum Florentíno prǽsule colláta, hi septem viri, géneris nobilitáte divitiísque posthábitis, sub vilíssimis detritísque véstibus cilício indúti, octáva die Septémbris in rurálem quamdam ædículam secessére, ut ea die primórdia vitæ sanctióris auspicaréntur, qua ipsa Dei Génetrix mortálibus orta sanctíssimam vitam incéperat.\par
\switchcolumn
\EN
\noindent In the thirteenth century, when the more cultured parts of Italy were rent by the dread dissension of the Emperor Frederick the Second and by bloody civil wars, the mercy of God set forth diverse men eminent for holiness, and among others raised up seven nobles of Florence, who were bound one to another in charity and gave an illustrious example of brotherly love. Their names were Bonfiglio Monaldi, Bonajuncta Manetti, Manetto Antalli, Amadeo de' Amidei, Uguccio de' Uguccioni, Sosteneo de' Sostenei, and Alexis de' Falconieri. Upon the holiday of the Assumption of the Virgin into heaven in the year\par
1233 They were praying in the oratory of a guild called the Guild of Praise, when the same Mother of God appeared to each one of them, and bade them embrace a life of greater holiness and perfection. These seven men discussed the matter with the Bishop of Florence, and then, considering neither the nobility of their birth nor their wealth, and clad in haircloth under vile and worn - out garments, withdrew into a little house in the country upon the th day of September, that they might begin their holier life upon the same day whereon the Mother of God herself had by her birth begun her life of holiness upon earth.\par
\switchcolumn*

\LA
\begin{Resp}\Rsign Honéstum fecit illum Dóminus, et custodívit eum ab inimícis, et a seductóribus tutávit illum: \Star{} Et dedit illi claritátem ætérnam.\end{Resp}
\begin{Resp}\Vsign Justum dedúxit Dóminus per vias rectas, et osténdit illi regnum Dei.\end{Resp}
\begin{Resp}\Rsign Et dedit illi claritátem ætérnam.\end{Resp}
\switchcolumn
\EN
\begin{Resp}\Rsign The Lord made him honourable, and defended him from his enemies, and kept him safe from those that lay in wait for him; \Star{} And gave him perpetual glory.\end{Resp}
\begin{Resp}\Vsign He went down with him into the pit, and left him not in bonds.\end{Resp}
\begin{Resp}\Rsign And gave him perpetual glory.\end{Resp}
\switchcolumn*

\LA
\LessonHead{Lectio V}
\switchcolumn
\EN
\LessonHead{Lesson V}
\switchcolumn*

\LA
\noindent Hoc vitæ institútum quam sibi foret accéptum Deus miráculo osténdit. Nam cum paulo deínceps hi septem viri per Florentínam urbem ostiátim eleemósynam emendicárent, áccidit, ut repénte infántium voce, quos inter fuit sanctus Philíppus Benítius quintum ætátis mensem vix ingréssus, beátæ Maríæ Servi acclamaréntur: quo deínde nómine semper appelláti sunt. Quare, vitándi pópuli occúrsus ac solitúdinis amóre ducti, in Senárii montis recéssu omnes convenére, ibíque cæléste quoddam vitæ genus aggréssi sunt. Victitábant enim in spelúncis, sola aqua herbísque conténti; vigíliis aliísque asperitátibus corpus atterébant, Christi passiónem ac mæstíssimæ ejúsdem Genitrícis dolóres assídue meditántes. Quod cum olim sacra Parascéves die impénsius exsequeréntur, ipsa beáta Virgo illis iteráto appárens, lúgubrem vestem, quam indúerent, osténdit; sibíque acceptíssimum fore significávit, ut novum in Ecclésia regulárem órdinem excitárent, qui jugem recóleret ac promovéret memóriam dolórum, quos ipsa pértulit sub cruce Dómini. Hæc sanctus Petrus, ínclytus órdinis Prædicatórum Martyr, ex familiári cum sanctis illis viris consuetúdine ac peculiári étiam Deíparæ visióne cum didicísset; iis auctor fuit, ut órdinem regulárem sub appellatióne Servórum beátæ Vírginis institúerent: qui póstea ab Innocéntio quarto Pontífice máximo approbátus fuit.\par
\switchcolumn
\EN
\noindent She showed by a miracle how acceptable in His sight should be their manner of life, for a short while after, when these seven men were begging alms from door to door through the city of Florence, it came to pass that some children, among whom was holy Philip Benizi, who had then scarcely entered the fifth month of his age, called them blessed Mary's servants, by the which name they were called ever after. To avoid meeting people, and in the desire to be alone, they all withdrew together to the solitude of Monte Senario, and there began a kind of heavenly life. They lived in caves and upon herbs and water only, while they wore out their bodies with watching and other hardships, while they contemplated unweariedly the sufferings of Christ and the woes of His most sorrowful Mother. One Good Friday, when their thoughts were fixed thereon more than ever, the Blessed Virgin appeared to them twice, and showed them her garments of mourning as those wherein they should clothe themselves. She bade them know that she would take it right well that they should raise up in the Church a new order to recall the memory of the sorrows which she bore beneath the Cross of the Lord. Holy Peter, the illustrious martyr of the Order of Friars Preachers, learnt this not only from his familiar converse with these holy men, but also from a special vision of the Mother of God, and it was on his incitement that they founded the regular Order called that of the Servites, or servants of the Blessed Virgin, the which Order was afterward approved by the Supreme Pontiff Innocent IV.\par
\switchcolumn*

\LA
\begin{Resp}\Rsign Amávit eum Dóminus, et ornávit eum: stolam glóriæ índuit eum, \Star{} Et ad portas paradísi coronávit eum.\end{Resp}
\begin{Resp}\Vsign Índuit eum Dóminus lorícam fídei, et ornávit eum.\end{Resp}
\begin{Resp}\Rsign Et ad portas paradísi coronávit eum.\end{Resp}
\switchcolumn
\EN
\begin{Resp}\Rsign The Lord loved him and beautified him; He clothed him with a robe of glory, \Star{} And crowned him at the gates of Paradise.\end{Resp}
\begin{Resp}\Vsign The Lord hath put on him the breast-plate of faith, and hath adorned him.\end{Resp}
\begin{Resp}\Rsign And crowned him at the gates of Paradise.\end{Resp}
\switchcolumn*

\LA
\LessonHead{Lectio VI}
\switchcolumn
\EN
\LessonHead{Lesson VI}
\switchcolumn*

\LA
\noindent Porro sancti illi viri, cum plures sibi sócios adjunxíssent, Itáliæ civitátes atque óppida, præsértim Etrúriæ excúrrere cœpérunt, prædicántes ubíque Christum crucifíxum, civíles discórdias compescéntes et innúmeros fere dévios ad virtútis sémitam revocántes. Neque Itáliam modo, sed et Gálliam, Germániam ac Polóniam suis evangélicis labóribus excoluérunt. Dénique cum bonum Christi odórem longe latéque diffudíssent, portentórum quoque glória illústres, migrárunt ad Dóminum. Sed quos unus veræ fraternitátis ac religiónis amor in vita sociáverat, unum páriter demórtuos contéxit sepúlcrum, únaque pópuli venerátio prosecúta est. Quaprópter Clemens undécimus et Benedíctus décimus tértius Pontífices máximi delátum iísdem a plúribus sǽculis indivíduum cultum confírmarunt; ac Leo décimus tértius, approbátis ántea miráculis, post indúltam veneratiónem ad collectívam eorúmdem invocatiónem a Deo patrátis, eósdem anno quinquagésimo sacerdótii sui Sanctórum honóribus cumulávit, eorúmque memóriam Offício et Missa in univérsa Ecclésia quotánnis recoléndam instítuit.\par
\switchcolumn
\EN
\noindent These holy men, when they had gathered to themselves some companions, began to go through the cities and towns of Italy, and especially of Tuscany, everywhere preaching Christ crucified, stilling contests among the citizens, and calling back almost countless backsliders into the path of grace. Neither did they make Italy only the field of their Gospel labours, but also France, Germany, and Poland. They passed away to be ever with the Lord when they had spread far and wide a sweet savour of Christ, and were famous also for the glory of signs and wonders. As one love of brotherhood and of the monastic life had joined them together upon earth, so one grave held their dead bodies, and one honour was paid them by the people. For this reason the Supreme Pontiffs Clement XI. and Benedict XIII. confirmed the honour whicfl had for centuries been paid to them individually, and Leo XIII., after proof of their miracles which had been wrought by God on the common invocation of these saints, after their veneration had been sanctioned in the jubilee year of his priesthood, decreed to them the honours paid to Saints, and ordered that their memory should every year be kept throughout the universal Church with an office and Mass.\par
\switchcolumn*

\LA
\begin{Resp}\Rsign Iste homo perfécit ómnia quæ locútus est ei Deus, et dixit ad eum: Ingrédere in réquiem meam: \Star{} Quia te vidi justum coram me ex ómnibus géntibus.\end{Resp}
\begin{Resp}\Vsign Iste est, qui contémpsit vitam mundi, et pervénit ad cæléstia regna.\end{Resp}
\begin{Resp}\Rsign Quia te vidi justum coram me ex ómnibus géntibus.\end{Resp}
\begin{Resp}\Vsign Glória Patri, et Fílio, \Star{} et Spirítui Sancto.\end{Resp}
\begin{Resp}\Rsign Quia te vidi justum coram me ex ómnibus géntibus.\end{Resp}
\switchcolumn
\EN
\begin{Resp}\Rsign This is he which did according unto all that God commanded him; and God said unto him: Enter thou into My rest, \Star{} For thee have I seen righteous before Me among all people.\end{Resp}
\begin{Resp}\Vsign This is he which loved not his life in this world, and is come unto an everlasting kingdom.\end{Resp}
\begin{Resp}\Rsign For thee have I seen righteous before Me among all people.\end{Resp}
\begin{Resp}\Vsign Glory be to the Father, and to the Son, \Star{} and to the Holy Ghost.\end{Resp}
\begin{Resp}\Rsign For thee have I seen righteous before Me among all people.\end{Resp}
\switchcolumn*

\LA
\LessonHead{Lectio VII}
\switchcolumn
\EN
\LessonHead{Lesson VII}
\switchcolumn*

\LA
\LessonTitle{Léctio sancti Evangélii secúndum Matthǽum.}
\switchcolumn
\EN
\LessonTitle{From the Holy Gospel according to Matthew}
\switchcolumn*

\LA
\Cite{Matt 19:27-29}
\switchcolumn
\EN
\Cite{Matt 19:27-29}
\switchcolumn*

\LA
\noindent In illo témpore: Dixit Petrus ad Jesum: Ecce nos relíquimus ómnia, et secúti sumus te: quid ergo erit nobis? Et réliqua.\par
\switchcolumn
\EN
\noindent At that time, Peter said unto Jesus: Behold, we have forsaken all, and followed thee what shall we have therefore? And so on.\par
\switchcolumn*

\LA
\LessonTitle{Homilía sancti Hierónymi Presbýteri.}
\switchcolumn
\EN
\LessonTitle{Homily by St. Jerome, Priest (at Bethlehem.)}
\switchcolumn*

\LA
\Source{Lib. 3. in Matth. cap. 19.}
\switchcolumn
\EN
\Source{Bk. iii. on Matth xix.}
\switchcolumn*

\LA
\noindent Grandis fidúcia! Petrus piscátor erat, dives non fúerat, cibos manu et arte quærébat: et tamen lóquitur confidénter: Relíquimus ómnia. Et quia non súfficit tantum relínquere, jungit quod perféctum est: Et secúti sumus te. Fécimus quod jussísti: quid ígitur nobis dabis prǽmii? Jesus autem dixit illis: Amen dico vobis, quod vos qui secúti estis me, in regeneratióne, cum séderit Fílius hóminis in sede majestátis suæ, sedébitis et vos super sedes duódecim, judicántes duódecim tribus Ísraël. Non dixit: Qui reliquístis ómnia; hoc enim et Crates fecit philósophus, et multi álii divítias contempsérunt: sed, Qui secúti estis me: quod próprie Apostolórum est, atque credéntium.\par
\switchcolumn
\EN
\noindent Peter was a fisherman, he was not rich, he earned his bread by his hand and skill, and nevertheless he is thus bold, and saith confidently: We have forsaken all. And because it sufficeth not to forsake only, he addeth that which to do is to be perfect: and followed thee. We have done that which Thou hast commanded us, what reward therefore wilt Thou give us? And Jesus said unto them Amen I say unto you, that ye which have followed Me, in the regeneration, when the Son of Man shall sit in the throne of His glory, ye also shall sit upon twelve thrones, judging the twelve tribes of Israel. He said not, Ye which have forsaken all, for this did even Crates the philosopher, and they which have set nothing by riches are many, but, Ye which have followed Me. This did the Apostles, and this do believers do.\par
\switchcolumn*

\LA
\begin{Resp}\Rsign Iste est qui ante Deum magnas virtútes operátus est, et de omni corde suo laudávit Dóminum: \Star{} Ipse intercédat pro peccátis ómnium populórum.\end{Resp}
\begin{Resp}\Vsign Ecce homo sine queréla, verus Dei cultor, ábstinens se ab omni ópere malo, et pérmanens in innocéntia sua.\end{Resp}
\begin{Resp}\Rsign Ipse intercédat pro peccátis ómnium populórum.\end{Resp}
\switchcolumn
\EN
\begin{Resp}\Rsign This is he which wrought great wonders before God, and praised the Lord with all his heart. \Star{} May he pray for all people, that their sins may be forgiven unto them.\end{Resp}
\begin{Resp}\Vsign Behold a man without blame, a worshipper of God in truth, keeping himself clean from every evil work, and abiding still in his innocency.\end{Resp}
\begin{Resp}\Rsign May he pray for all people, that their sins may be forgiven unto them.\end{Resp}
\switchcolumn*

\LA
\LessonHead{Lectio VIII}
\switchcolumn
\EN
\LessonHead{Lesson VIII}
\switchcolumn*

\LA
\noindent In regeneratióne, cum séderit Fílius hóminis in sede majestátis suæ (quando et mórtui de corruptióne resúrgent incorrúpti) sedébitis et vos in sóliis judicántium, condemnántes duódecim tribus Ísraël: quia vobis credéntibus, illi crédere noluérunt. Et omnis qui relíquerit domum, vel fratres, aut soróres, aut patrem, aut matrem, aut uxórem, aut fílios, aut agros propter nomen meum, céntuplum accípiet, et vitam ætérnam possidébit. Locus iste cum illa senténtia cóngruit, in qua Salvátor lóquitur: Non veni pacem míttere, sed gládium. Veni enim separáre hóminem a patre suo, et matrem a fília, et nurum a socru: et inimíci hóminis doméstici ejus. Qui ergo propter fidem Christi, et prædicatiónem Evangélii, omnes afféctus contémpserint, atque divítias et sǽculi voluptátes: isti céntuplum recípient, et vitam ætérnam possidébunt.\par
\switchcolumn
\EN
\noindent In the regeneration, when the Son of Man shall sit in the throne of His glory, and when the dead shall rise again from corruption incorruptible, (i Cor. xv. 53,) ye also shall sit upon twelve thrones of judgment, condemning the twelve tribes of Israel, because, when ye believed in Me, they would not. (John iii. 18.) And every one that hath forsaken houses, or brethren, or sisters, or father, or mother, or wife, or children, or lands, for My Name's sake, shall receive an hundredfold, and shall inherit everlasting life. This place agreeth well with that other where the Saviour saith I came not to send peace, but a sword. For I am come to set a man at variance against his father, and the daughter against her mother, and the daughter-in-law against her mother-in-law; and a man's foes shall be they of his own household. (Matth. x. 34.) Every one, therefore, that hath set no store by affection, and riches, and the pleasures of the world, for Christ's faith's sake, and the preaching of the Gospel, shall receive an hundred-fold, and shall inherit everlasting life.\par
\switchcolumn*

\LA
\begin{Resp}\Rsign Sint lumbi vestri præcíncti, et lucérnæ ardéntes in mánibus vestris: \Star{} Et vos símiles homínibus exspectántibus dóminum suum, quando revertátur a núptiis.\end{Resp}
\begin{Resp}\Vsign Vigiláte ergo, quia nescítis qua hora Dóminus vester ventúrus sit.\end{Resp}
\begin{Resp}\Rsign Et vos símiles homínibus exspectántibus dóminum suum, quando revertátur a núptiis.\end{Resp}
\begin{Resp}\Vsign Glória Patri, et Fílio, \Star{} et Spirítui Sancto.\end{Resp}
\begin{Resp}\Rsign Et vos símiles homínibus exspectántibus dóminum suum, quando revertátur a núptiis.\end{Resp}
\switchcolumn
\EN
\begin{Resp}\Rsign Let your loins be girded about, and your lights burning; \Star{} And ye yourselves like unto men that wait for their lord, when he will return from the wedding.\end{Resp}
\begin{Resp}\Vsign Watch therefore, for ye know not what hour your Lord doth come.\end{Resp}
\begin{Resp}\Rsign And ye yourselves like unto men that wait for their lord, when he will return from the wedding.\end{Resp}
\begin{Resp}\Vsign Glory be to the Father, and to the Son, \Star{} and to the Holy Ghost.\end{Resp}
\begin{Resp}\Rsign And ye yourselves like unto men that wait for their lord, when he will return from the wedding.\end{Resp}
\switchcolumn*

\LA
\LessonHead{Lectio IX}
\switchcolumn
\EN
\LessonHead{Lesson IX}
\switchcolumn*

\LA
\noindent Ex occasióne hujus senténtiæ quidam introdúcunt mille annos post resurrectiónem, dicéntes, tunc nobis céntuplum ómnium rerum quas dimísimus, et vitam ætérnam esse reddéndam: non intelligéntes, quod, si in céteris digna sit repromíssio, in uxóribus appáreat turpitúdo: ut qui unam pro Dómino dimíserit centum recípiat in futúro. Sensus ergo iste est: Qui carnália pro Salvatóre dimíserit, spirituália recípiet: quæ comparatióne et mérito sui ita erunt, quasi si parvo número centenárius númerus comparétur.\par
\switchcolumn
\EN
\noindent By reason of these words, an hundredfold, some will have that there shall be a thousand years after the resurrection, wherein they that have forsaken all things shall receive an hundredfold of those things which they have forsaken, and shall inherit everlasting life. Such men consider not that though in other things this were worthy, as touching wives it is unseemly for it becometh us not to think that he that hath forsaken one wife in this world, shall receive an hundred wives in that which is to come. But the meaning is this, that every one that for the Saviour's sake hath forsaken earthly things, shall receive spiritual things which things, being rightly weighed against earthly things, are as though an hundredfold were weighed against one.\par
\switchcolumn*

\LA
\TeDeum{Te Deum}
\switchcolumn
\EN
\TeDeum{Te Deum}
\switchcolumn*

\end{paracol}

\par\needspace{10\baselineskip}\addvspace{1.2em}{\centering\small\textsc{Die 14 Februarii} · \textsc{14 February}\par}
\Office{S. Valentini Presbyteri et Martyris}{St. Valentine, Priest and Martyr}{Simplex}
\addcontentsline{toc}{subsection}{Die 14 Februarii · S. Valentini Presbyteri et Martyris\enspace\textit{St. Valentine, Priest and Martyr}}
\begin{paracol}{2}
\LA
\RefLine{Lectiones I–II: de Scriptura occurrente.}
\switchcolumn
\EN
\RefLine{Lessons I–II: from the Scripture of the day.}
\switchcolumn*

\LA
\LessonHead{Lectio III (pro commemoratione)}
\switchcolumn
\EN
\LessonHead{Lesson III (when commemorated)}
\switchcolumn*

\LA
\LessonTitle{Sermo sancti Augustíni Epíscopi}
\switchcolumn
\EN
\LessonTitle{From the Sermons of St Augustine, Bishop (of Hippo.)}
\switchcolumn*

\LA
\Source{Sermo 44 de Sanctis}
\switchcolumn
\EN
\Source{44 on the Saints.}
\switchcolumn*

\LA
\noindent Triumphális beáti Mártyris Valentini dies hódie nobis anniversária celebritáte recúrrit: cujus glorificatióni sicut congáudet Ecclésia, sic ejus propónit sequénda vestígia. Si enim compátimur, et conglorificábimur. In cujus glorióso agóne duo nobis præcípue consideránda sunt: induráta vidélicet tortóris sævítia, et Mártyris invícta patiéntia. Sævítia tortóris, ut eam detestémur: patiéntia Mártyris, ut eam imitémur. Audi Psalmístam advérsus malítiam increpántem: Noli æmulári in malignántibus quóniam tamquam fœnum velóciter aréscent. Quod autem advérsus malignántes patiéntia exhibénda sit, audi Apóstolum suadéntem: Patiéntia vobis necessária est, ut reportétis promissiónes.\par
\switchcolumn
\EN
\noindent The illustrious day whereon the blessed Martyr Valentine conquered, doth this day come round to us again and as the Church doth rejoice with him in his glory, so doth she set before us his footsteps to be followed. For if we suffer, we shall also reign with him. In his glorious battle we have two things chiefly to consider the hardened cruelty of the tormentor, and the unconquered patience of the Martyr the cruelty of the tormentor, that we may abhor it; the patience of the Martyr, that we may imitate it. Hear what the Psalmist saith, complaining against sin: “Fret not thyself because of the evil-doers, for they shall soon dry up like the grass.” (xxxvi. 1.) But touching the patience which is to be shown against the evil-doers, hear the word wherewith the Apostle moveth us Ye have need of patience, that ye may receive the promise. (Heb. x. 36.)\par
\switchcolumn*

\LA
\TeDeum{Te Deum}
\switchcolumn
\EN
\TeDeum{Te Deum}
\switchcolumn*

\end{paracol}

\par\needspace{10\baselineskip}\addvspace{1.2em}{\centering\small\textsc{Die 15 Februarii} · \textsc{15 February}\par}
\Office{SS. Faustini et Jovitæ}{Sts. Faustinus and Jovita, Martyrs}{Simplex}
\addcontentsline{toc}{subsection}{Die 15 Februarii · SS. Faustini et Jovitæ\enspace\textit{Sts. Faustinus and Jovita, Martyrs}}
\begin{paracol}{2}
\LA
\RefLine{Lectiones I–II: de Scriptura occurrente.}
\switchcolumn
\EN
\RefLine{Lessons I–II: from the Scripture of the day.}
\switchcolumn*

\LA
\LessonHead{Lectio III (pro commemoratione)}
\switchcolumn
\EN
\LessonHead{Lesson III (when commemorated)}
\switchcolumn*

\LA
\noindent Faustinus et Jovita fratres, nobiles Brixiani, in multis Italiæ urbibus, quo vincti sæviente Trajani persecutione ducebantur, acerbissima supplicia perpessi, fortes in christianæ fidei confessione perstiterunt. Nam Brixiæ diu vinculis constricti, feris étiam objecti in ignemque conjecti, et a bestiis et a flamma íntegri et incolumes servati sunt. Inde vero eisdem catenis colligati Mediolanum venerunt; ubi eorum fides tentata exquisitissimis tormentis, tamquam igne aurum, in cruciatibus magis enituit. Postea Romam missi, ab Evaristo Pontifice confirmati, ibi quoque crudelissime torquentur. Denique perducti Neápolim, in ea étiam urbe varie cruciati, vinctis manibus pedibusque in mare demerguntur; unde per Angelos mirabiliter erepti sunt. Quare multos et constantia in tormentis et miraculorum virtute ad Christi fidem converterunt. Postrémo reducti Brixiam initio suscepti ab Hadriano imperii, securi percussi, illustrem martyrii coronam acceperunt.\par
\switchcolumn
\EN
\noindent Faustinus and Jovita were brothers, born of a noble family at Brescia. While Trajan's persecution was raging, they were taken about in chains from one city of Italy to another, and exhibited in torture in each. This cruelty utterly failed to silence their confession of Christ, Whom they preached by their sufferings in every place where they were shown. They were afterwards kept for a long time at Brescia, where they were exhibited with wild beasts, and tormented with fire. Being both still alive, they were brought to Milan, without their chains having ever been taken off. At Milan they were tortured again with every invention of cruelty that could be devised. Nevertheless the great power of their faith made them more than conquerors, shining even as gold tried in the furnace. From Milan they were brought to Rome, where they were confirmed by Pope Evaristus, and where they were put to the torture again with extreme barbarity. They were afterwards shown in public at Naples, where the tormentors displayed their skill in diverse ways upon them. Here they were thrown chained into the sea, but the angels delivered them. Their stations of suffering, by their God-like patience, and the wonderful Power displayed in them, had now turned many souls to Jesus. In the end they were carried back to Brescia, and, when Hadrian took the empire, they were put to death by the axe at that place. The crown of martyrdom which they won is glorious.\par
\switchcolumn*

\LA
\TeDeum{Te Deum}
\switchcolumn
\EN
\TeDeum{Te Deum}
\switchcolumn*

\end{paracol}

\par\needspace{10\baselineskip}\addvspace{1.2em}{\centering\small\textsc{Die 18 Februarii} · \textsc{18 February}\par}
\Office{S. Simeonis Episcopi et Martyris}{St. Simeon, Bishop and Martyr}{Simplex}
\addcontentsline{toc}{subsection}{Die 18 Februarii · S. Simeonis Episcopi et Martyris\enspace\textit{St. Simeon, Bishop and Martyr}}
\begin{paracol}{2}
\LA
\RefLine{Lectiones I–II: de Scriptura occurrente.}
\switchcolumn
\EN
\RefLine{Lessons I–II: from the Scripture of the day.}
\switchcolumn*

\LA
\LessonHead{Lectio III (pro commemoratione)}
\switchcolumn
\EN
\LessonHead{Lesson III (when commemorated)}
\switchcolumn*

\LA
\noindent Simeon, fílius Cléophæ, post Jacobum próximus Jerosolymis ordinátus epíscopus, Trajano imperatóre apud Atticum consularem est accusatus, quod christianus esset et Christi propinquus. Comprehendebántur enim omnes eo témpore, quicúmque ex genere David orti essent. Quare multis cruciátus tormentis eodem passiónis genere, quod Salvátor noster subíerat, affícitur: mirántibus ómnibus, quod homo ætate confectus (erat enim centum et viginti annórum) acerbíssimos crucis dolóres fortiter constanterque paterétur.\par
\switchcolumn
\EN
\noindent Simeon, the son of Cleophas, (Matth. xiii. 55,) was chosen the second Bishop of Jerusalem, (in the year 62,) being the first after James. Under the Emperor Trajan he was accused before the Pro-Consul Atticus, as being both a Christian and a relation of Christ, this being the time when all were arrested that were of the lineage of David. He underwent with great suffering the same things that were inflicted on our Saviour, and all men marvelled to see with how great boldness and firmness he endured the grievous torment of the cross, at his great age, for he was an hundred and twenty years old. (a.d. 107 or 116.)\par
\switchcolumn*

\LA
\TeDeum{Te Deum}
\switchcolumn
\EN
\TeDeum{Te Deum}
\switchcolumn*

\end{paracol}

\par\needspace{10\baselineskip}\addvspace{1.2em}{\centering\small\textsc{Die 22 Februarii} · \textsc{22 February}\par}
\Office{In Cathedra S. Petri Apostoli Antiochiæ}{Chair of St. Peter at Antioch}{Duplex majus}
\addcontentsline{toc}{subsection}{Die 22 Februarii · In Cathedra S. Petri Apostoli Antiochiæ\enspace\textit{Chair of St. Peter at Antioch}}
\begin{paracol}{2}
\LA
\LessonHead{Lectio I}
\switchcolumn
\EN
\LessonHead{Lesson I}
\switchcolumn*

\LA
\LessonTitle{Incipit Epístola prima beáti Petri Apóstoli}
\switchcolumn
\EN
\LessonTitle{Lesson from the first letter of St. Peter the Apostle}
\switchcolumn*

\LA
\Cite{1 Pet 1:1-5}
\switchcolumn
\EN
\Cite{1 Pet 1:1-5}
\switchcolumn*

\LA
\noindent Petrus Apóstolus Jesu Christi, eléctis ádvenis dispersiónis Ponti, Galátiæ, Cappadóciæ, Asiæ, et Bithýniæ, secúndum præsciéntiam Dei Patris, in sanctificatiónem Spíritus, in obediéntiam, et aspersiónem sánguinis Jesu Christi: Grátia vobis, et pax multiplicétur. Benedíctus Deus et Pater Dómini nostri Jesu Christi, qui secúndum misericórdiam suam magnam regenerávit nos in spem vivam, per resurrectiónem Jesu Christi ex mórtuis, in hereditátem incorruptíbilem, et incontaminátam, et immarcescíbilem, conservátam in cælis in vobis, qui in virtúte Dei custodímini per fidem in salútem, parátam revelári in témpore novíssimo.\par
\switchcolumn
\EN
\noindent Peter, an apostle of Jesus Christ, to the strangers dispersed through Pontus, Galatia, Cappadocia, Asia, and Bithynia, elect, According to the foreknowledge of God the Father, unto the sanctification of the Spirit, unto obedience and sprinkling of the blood of Jesus Christ: Grace unto you and peace be multiplied. Blessed be the God and Father of our Lord Jesus Christ, who according to his great mercy hath regenerated us unto a lively hope, by the resurrection of Jesus Christ from the dead, Unto an inheritance incorruptible, and undefiled, and that can not fade, reserved in heaven for you, Who, by the power of God, are kept by faith unto salvation, ready to be revealed in the last time.\par
\switchcolumn*

\LA
\begin{Resp}\Rsign Simon Petre, ántequam de navi vocárem te, novi te et super plebem meam príncipem te constítui: \Star{} Et claves regni cælórum trádidi tibi.\end{Resp}
\begin{Resp}\Vsign Quodcúmque ligáveris super terram erit ligátum et in cælis; et quodcúmque sólveris super terram, erit solútum et in cælis.\end{Resp}
\begin{Resp}\Rsign Et claves regni cælórum trádidi tibi.\end{Resp}
\switchcolumn
\EN
\begin{Resp}\Rsign Simon Peter, before I called thee out of the ship, I knew thee, and appointed thee for a ruler over My people. \Star{} And I have given unto thee the keys of the kingdom of heaven.\end{Resp}
\begin{Resp}\Vsign Whatsoever, thou shalt bind on earth, shall be bound in heaven; and whatsoever thou shalt loose on earth, shall be loosed in heaven.\end{Resp}
\begin{Resp}\Rsign And I have given unto thee the keys of the kingdom of heaven.\end{Resp}
\switchcolumn*

\LA
\LessonHead{Lectio II}
\switchcolumn
\EN
\LessonHead{Lesson II}
\switchcolumn*

\LA
\Cite{1 Pet 1:6-9}
\switchcolumn
\EN
\Cite{1 Pet 1:6-9}
\switchcolumn*

\LA
\noindent In quo exsultábitis, módicum nunc si opórtet contristári in váriis tentatiónibus: ut probátio vestræ fídei multo pretiósior auro (quod per ignem probátur) inveniátur in laudem, et glóriam, et honórem, in revelatióne Jesu Christi: quem cum non vidéritis, dilígitis: in quem nunc quoque non vidéntes créditis: credéntes autem exsultábitis lætítia inenarrábili, et glorificáta: reportántes finem fídei vestræ, salútem animárum.\par
\switchcolumn
\EN
\noindent Wherein you shall greatly rejoice, if now you must be for a little time made sorrowful in diverse temptations: That the trial of your faith much more precious than gold which is tried by the fire may be found unto praise and glory and honour at the appearing of Jesus Christ: Whom having not seen, you love: in whom also now, though you see him not, you believe: and believing shall rejoice with joy unspeakable and glorified; Receiving the end of your faith, even the salvation of your souls.\par
\switchcolumn*

\LA
\begin{Resp}\Rsign Si díligis me, Simon Petre, pasce oves meas. Dómine, tu nosti quia amo te, \Star{} Et ánimam meam pono pro te.\end{Resp}
\begin{Resp}\Vsign Si oportúerit me mori tecum, non te negábo.\end{Resp}
\begin{Resp}\Rsign Et ánimam meam pono pro te.\end{Resp}
\switchcolumn
\EN
\begin{Resp}\Rsign Simon Peter, if thou lovest Me, feed My sheep. Lord, Thou knowest that I love thee; \Star{} I will lay down my life for thy sake.\end{Resp}
\begin{Resp}\Vsign If I should die with thee, I will not deny thee.\end{Resp}
\begin{Resp}\Rsign I will lay down my life for thy sake.\end{Resp}
\switchcolumn*

\LA
\LessonHead{Lectio III}
\switchcolumn
\EN
\LessonHead{Lesson III}
\switchcolumn*

\LA
\Cite{1 Pet 1:10-12}
\switchcolumn
\EN
\Cite{1 Pet 1:10-12}
\switchcolumn*

\LA
\noindent De qua salúte exquisiérunt atque scrutáti sunt prophétæ, qui de futúra in vobis grátia prophetavérunt; scrutántes in quod vel quale tempus significáret in eis Spíritus Christi: prænúntians eas quæ in Christo sunt passiónes et posterióres glórias: quibus revelátum est, quia non sibimetípsis, vobis autem ministrábant ea, quæ nunc nuntiáta sunt vobis per eos, qui evangelizavérunt vobis, Spíritu Sancto misso de cælo, in quem desíderant Angeli prospícere.\par
\switchcolumn
\EN
\noindent Of which salvation the prophets have inquired and diligently searched, who prophesied of the grace to come in you. Searching what or what manner of time the Spirit of Christ in them did signify: when it foretold those sufferings that are in Christ, and the glories that should follow: To whom it was revealed, that not to themselves, but to you they ministered those things which are now declared to you by them that have preached the gospel to you, the Holy Ghost being sent down from heaven, on whom the angels desire to look.\par
\switchcolumn*

\LA
\begin{Resp}\Rsign Tu es Petrus, et super hanc petram ædificábo Ecclésiam meam, et portæ ínferi non prævalébunt advérsus eam: \Star{} Et tibi dabo claves regni cælórum.\end{Resp}
\begin{Resp}\Vsign Quodcúmque ligáveris super terram, erit ligátum et in cælis; et quodcúmque sólveris super terram, erit solútum et in cælis.\end{Resp}
\begin{Resp}\Rsign Et tibi dabo claves regni cælórum.\end{Resp}
\begin{Resp}\Vsign Glória Patri, et Fílio, \Star{} et Spirítui Sancto.\end{Resp}
\begin{Resp}\Rsign Et tibi dabo claves regni cælórum.\end{Resp}
\switchcolumn
\EN
\begin{Resp}\Rsign Thou art Peter, and upon this rock I will build My Church, and the gates of hell shall not prevail against it. \Star{} And I will give unto thee the keys of the kingdom of heaven.\end{Resp}
\begin{Resp}\Vsign Whatsoever thou shalt bind on earth, shall be bound in heaven; and whatsoever thou shalt loose on earth, shall be loosed in heaven.\end{Resp}
\begin{Resp}\Rsign And I will give unto thee the keys of the kingdom of heaven.\end{Resp}
\begin{Resp}\Vsign Glory be to the Father, and to the Son, \Star{} and to the Holy Ghost.\end{Resp}
\begin{Resp}\Rsign And I will give unto thee the keys of the kingdom of heaven.\end{Resp}
\switchcolumn*

\LA
\LessonHead{Lectio IV}
\switchcolumn
\EN
\LessonHead{Lesson IV}
\switchcolumn*

\LA
\LessonTitle{Sermo sancti Augustíni Epíscopi}
\switchcolumn
\EN
\LessonTitle{From the Sermons of St. Augustine, Bishop (of Hippo.)}
\switchcolumn*

\LA
\Source{Sermo 15 de Sanctis}
\switchcolumn
\EN
\Source{16th on the Saints}
\switchcolumn*

\LA
\noindent Institútio solemnitátis hodiérnæ a senióribus nostris Cáthedræ nomen accépit, ídeo quod primus Apostolórum Petrus hódie episcopátus cáthedram suscepísse referátur. Recte ergo ecclésiæ natálem sedis illíus colunt, quam Apóstolus pro ecclesiárum salúte suscépit, dicénte Dómino: Tu es Petrus, et super hanc petram ædificábo Ecclésiam meam.\par
\switchcolumn
\EN
\noindent The solemn Feast of today received from our forefathers the name of that of St Peter's Chair at Antioch, because there is a tradition that it was on this day that Peter, first of the Apostles, was enthroned in a Bishop's Chair. Rightly, therefore, do the Churches observe the first day of that Chair, the right to which the Apostle received for the salvation of the Churches from the Lord of the Churches Himself, with the words: Thou art Peter, and upon this rock I will build My Church.\par
\switchcolumn*

\LA
\begin{Resp}\Rsign Tu es pastor óvium, princeps Apostolórum: tibi trádidit Deus ómnia regna mundi: \Star{} Et ídeo tibi tráditæ sunt claves regni cælórum.\end{Resp}
\begin{Resp}\Vsign Quodcúmque ligáveris super terram, erit ligátum et in cælis; et quodcúmque sólveris super terram, erit solútum et in cælis.\end{Resp}
\begin{Resp}\Rsign Et ídeo tibi tráditæ sunt claves regni cælórum.\end{Resp}
\switchcolumn
\EN
\begin{Resp}\Rsign Thou art the Shepherd of the sheep, and the Prince of the Apostles, and unto thee hath God given all the kingdoms of the world. \Star{} Therefore unto thee hath He given the keys of the kingdom of heaven.\end{Resp}
\begin{Resp}\Vsign Whatsoever thou shalt bind on earth shall be bound in heaven; and whatsoever thou shalt loose on earth shall be loosed in heaven.\end{Resp}
\begin{Resp}\Rsign Therefore unto thee hath He given the keys of the kingdom of heaven.\end{Resp}
\switchcolumn*

\LA
\LessonHead{Lectio V}
\switchcolumn
\EN
\LessonHead{Lesson V}
\switchcolumn*

\LA
\noindent Petrum ítaque fundaméntum Ecclésiæ Dóminus nominávit: et ídeo digne fundaméntum hoc Ecclésia colit, supra quod ecclesiástici ædifícii altitúdo consúrgit. Unde conveniénter psalmus, qui lectus est, dicit: Exáltent eum in ecclésia plebis, et in cáthedra seniórum laudent eum. Benedíctus Deus, qui beátum Petrum Apóstolum in Ecclésia exaltári præcépit: quia dignum est ut fundaméntum hoc in Ecclésia honorétur, per quod ad cælum conscénditur.\par
\switchcolumn
\EN
\noindent It was the Lord Himself Who called Peter the foundation of the Church, and therefore it is right that the Church should reverence this foundation whereon her mighty structure riseth. Justly is it written in the Psalm which we have just heard: Let them exalt him in the congregation of the people, and praise him in the assembly of the elders. Blessed be God, Who hath commanded that the Blessed Apostle Peter should be exalted in the congregation! Worthy to be honoured by the Church is that foundation from which her goodly towers rise, pointing to heaven!\par
\switchcolumn*

\LA
\begin{Resp}\Rsign Ego pro te rogávi, Petre, ut non defíciat fides tua: \Star{} Et tu aliquándo convérsus confírma fratres tuos.\end{Resp}
\begin{Resp}\Vsign Caro et sanguis non revelávit tibi, sed Pater meus, qui est in cælis.\end{Resp}
\begin{Resp}\Rsign Et tu aliquándo convérsus confírma fratres tuos.\end{Resp}
\switchcolumn
\EN
\begin{Resp}\Rsign Peter, I have prayed for thee, that thy faith fail not \Star{} And when thou art converted, strengthen thy brethren.\end{Resp}
\begin{Resp}\Vsign Flesh and blood hath not revealed it unto thee, but My Father Which is in heaven.\end{Resp}
\begin{Resp}\Rsign And when thou art converted, strengthen thy brethren.\end{Resp}
\switchcolumn*

\LA
\LessonHead{Lectio VI}
\switchcolumn
\EN
\LessonHead{Lesson VI}
\switchcolumn*

\LA
\noindent Quod natális ergo Cáthedræ hódie cólitur, sacerdotále honorátur offícium. Sibi hoc Ecclésiæ ínvicem præstant, quia tanto necésse plus habet Ecclésia dignitatis, quanto sacerdotále offícium plus honóris. Cum solemnitátem hanc Ecclésiis mérito religiósa observátio introdúxerit, miror, cur apud quosdam infidéles hódie tam perniciósus error incréverit, ut super túmulos defunctórum cibos et vina conférant, quasi egréssæ de corpóribus ánimæ carnáles cibos requírant.\par
\switchcolumn
\EN
\noindent In the honour which is this day paid to the inauguration of the first Bishop's throne, an honour is paid to the office of all Bishops. The Churches testify one to another, that, the greater the Church's dignity, the greater the reverence due to her priests. While I confess how rightly godly custom hath exalted this Feast in the estimation of all the Churches, the more do I wonder at the growth of that unhealthy error which at this day causeth some unbelievers to lay food and wine upon the graves of the dead, as if souls once rid of the body had any longer any need of bodily refreshment.\par
\switchcolumn*

\LA
\begin{Resp}\Rsign Petre, amas me? Tu scis, Dómine, quia amo te. \Star{} Pasce oves meas.\end{Resp}
\begin{Resp}\Vsign Simon Joánnis, díligis me plus his? Tu scis, Dómine, quia amo te.\end{Resp}
\begin{Resp}\Rsign Pasce oves meas.\end{Resp}
\begin{Resp}\Vsign Glória Patri, et Fílio, \Star{} et Spirítui Sancto.\end{Resp}
\begin{Resp}\Rsign Pasce oves meas.\end{Resp}
\switchcolumn
\EN
\begin{Resp}\Rsign Peter, lovest thou Me? Lord, Thou knowest that I love thee. \Star{} Feed My sheep.\end{Resp}
\begin{Resp}\Vsign Simon, son of Jonas, lovest thou Me more than these? Lord, Thou knowest that I love thee.\end{Resp}
\begin{Resp}\Rsign Feed My sheep.\end{Resp}
\begin{Resp}\Vsign Glory be to the Father, and to the Son, \Star{} and to the Holy Ghost.\end{Resp}
\begin{Resp}\Rsign Feed My sheep.\end{Resp}
\switchcolumn*

\LA
\LessonHead{Lectio VII}
\switchcolumn
\EN
\LessonHead{Lesson VII}
\switchcolumn*

\LA
\LessonTitle{Léctio sancti Evangélii secúndum Matthǽum}
\switchcolumn
\EN
\LessonTitle{From the Holy Gospel according to Matthew}
\switchcolumn*

\LA
\Cite{Matt 16:13-19}
\switchcolumn
\EN
\Cite{Matt 16:13-19}
\switchcolumn*

\LA
\noindent In illo témpore: Venit Jesus in partes Cæsaréæ Philíppi, et interrogábat discípulos suos, dicens: Quem dicunt hómines esse Fílium hóminis? Et réliqua.\par
\switchcolumn
\EN
\noindent At that time, Jesus came into the coasts of Caesarea Philippi, and He asked His disciples, saying: Who do men say that I, the Son of Man, am? And so on.\par
\switchcolumn*

\LA
\LessonTitle{Homilía sancti Leónis Papæ}
\switchcolumn
\EN
\LessonTitle{Homily by Pope St. Leo (the Great.)}
\switchcolumn*

\LA
\Source{Sermo 3 in annivers. assumpt. suæ, post initium}
\switchcolumn
\EN
\Source{3rd on the Anniversary of his own election}
\switchcolumn*

\LA
\noindent Apóstolos Dóminus, quid de se hómines opinéntur, intérrogat: et tamdiu sermo respondéntium commúnis est, quámdiu humánæ intellegéntiæ ambigúitas explicátur. At ubi quid hábeat sensus discipulórum exígitur, primus est in Dómini confessióne, qui primus est in apostólica dignitáte. Qui cum dixísset: Tu es Christus Fílius Dei vivi: respóndit ei Jesus: Beátus es Simon Bar-Jona, quia caro et sanguis non revelávit tibi, sed Pater meus, qui in cælis est. Id est, Ideo beátus es, quia Pater meus te dócuit: nec terréna opínio te feféllit, sed inspirátio cæléstis instrúxit: et non caro et sanguis, sed ille me tibi, cujus sum unigénitus Fílius, indicávit.\par
\switchcolumn
\EN
\noindent The Lord asked His disciples Who men said that He was, and their answers were human as long as they were the answers of human reason, unilluminated by Divine light. At last, when the glimmerings of earthly conjecture were spoken, he whose Apostleship is the first in dignity, was the first to confess his Lord. And Simon Peter answered and said: Thou art the Christ, the Son of the living God. And Jesus answered and said unto him: Blessed art thou, Simon Bar-Jona, for flesh and blood hath not revealed it unto thee, but My Father Which is in heaven. That is to say, For this cause art thou blessed, because My Father Himself hath taught thee; the opinions of men have not beguiled thee, the voices of angels have not taught thee, not flesh and blood, but He, Whose Only begotten Son I am, hath revealed Me unto thee.\par
\switchcolumn*

\LA
\begin{Resp}\Rsign Quem dicunt hómines esse Filium hóminis? dixit Jesus discípulis suis. Respóndens Petrus dixit: Tu es Christus Fílius Dei vivi. \Star{} Et ego dico tibi, quia tu es Petrus, et super hanc petram ædificábo Ecclésiam meam.\end{Resp}
\begin{Resp}\Vsign Beátus es, Simon Bar-Jona, quia caro et sanguis non revelávit tibi, sed Pater meus qui est in cælis.\end{Resp}
\begin{Resp}\Rsign Et ego dico tibi, quia tu es Petrus, et super hanc petram ædificábo Ecclésiam meam.\end{Resp}
\switchcolumn
\EN
\begin{Resp}\Rsign Jesus asked His disciples, saying: Who do men say that I, the Son of Man, am? Peter answered, and said: Thou art the Christ, the Son of the living God. \Star{} And I say unto thee, that thou art Peter, and upon this rock I will build My Church.\end{Resp}
\begin{Resp}\Vsign Blessed art thou, Simon Bar-Jona, for flesh and blood hath not revealed it unto thee, but My Father Which is in heaven.\end{Resp}
\begin{Resp}\Rsign And I say unto thee, that thou art Peter, and upon this rock I will build My Church.\end{Resp}
\switchcolumn*

\LA
\LessonHead{Lectio VIII}
\switchcolumn
\EN
\LessonHead{Lesson VIII}
\switchcolumn*

\LA
\noindent Et ego, inquit, dico tibi: hoc est, sicut Pater meus tibi manifestávit divinitátem meam, ita et ego tibi notam fácio excelléntiam tuam, quia tu es Petrus: id est, cum ego sim inviolábilis petra, ego lapis anguláris, qui fácio útraque unum, ego fundaméntum præter quod nemo potest áliud pónere: tamen tu quoque petra es, quia mea virtúte solidáris, ut quæ mihi potestáte sunt própria, sint tibi mecum participatióne commúnia. Et super hanc petram ædificábo Ecclésiam meam, et portæ ínferi non prævalébunt advérsus eam. Super hanc, inquit, fortitúdinem ætérnum éxstruam templum: et Ecclésiæ meæ cælo inserénda sublímitas, in hujus fídei firmitáte consúrget.\par
\switchcolumn
\EN
\noindent Thus saith the Lord unto Simon Peter: And I say also unto thee, That thou art Peter. That is to say, Even as My Father hath revealed unto thee concerning Me that I am God, even so now will I also reveal unto thee that thou art Peter; I am the sure Rock of defence, the Corner Stone, Who make both one, (Eph. ii. 20, 15,) I am the Foundation, beside Which other can no man lay, (1 Cor. iii. 11,) and thou also art a rock, in My Strength made hard, and those things whereof I by right am Lord, into thy hand do I give them, that thou mayst bear rule over them, for Me, and with Me. And upon this rock I will build My Church, and the gates of hell shall not prevail against it. Upon this strength of thine, whereof I am the Strength, I will build My eternal temple, and upon the truth of thy confession of Me I will make to rise My glorious Church whose spires shall pierce to heaven.\par
\switchcolumn*

\LA
\begin{Resp}\Rsign Elégit te Dóminus sacerdótem sibi, ad sacrificándum ei \Star{} Hóstiam laudis.\end{Resp}
\begin{Resp}\Vsign Immola Deo sacrifícium laudis, et redde Altíssimo vota tua.\end{Resp}
\begin{Resp}\Rsign Hóstiam laudis.\end{Resp}
\begin{Resp}\Vsign Glória Patri, et Fílio, \Star{} et Spirítui Sancto.\end{Resp}
\begin{Resp}\Rsign Hóstiam laudis.\end{Resp}
\switchcolumn
\EN
\begin{Resp}\Rsign The Lord hath chosen thee for a priest unto Himself, to offer up unto Him \Star{} The sacrifice of praise.\end{Resp}
\begin{Resp}\Vsign Offer unto God thanksgiving, and pay thy vows unto the Most High.\end{Resp}
\begin{Resp}\Rsign The sacrifice of praise.\end{Resp}
\begin{Resp}\Vsign Glory be to the Father, and to the Son, \Star{} and to the Holy Ghost.\end{Resp}
\begin{Resp}\Rsign The sacrifice of praise.\end{Resp}
\switchcolumn*

\LA
\LessonHead{Lectio IX}
\switchcolumn
\EN
\LessonHead{Lesson IX}
\switchcolumn*

\LA
\noindent Hanc confessiónem portæ ínferi non tenébunt, mortis víncula non ligábunt: vox enim ista vox vitæ est. Et sicut confessóres suos in cæléstia próvehit, ita negatóres ad inférna demérgit. Propter quod dicit beatíssimo Petro: Tibi dabo claves regni cælórum: Et quæcúmque ligáveris super terram, erunt ligáta et in cælis: et quæcúmque sólveris super terram, erunt solúta et in cælis. Transívit quidem étiam in álios Apóstolos vis potestátis istíus, et ad omnes Ecclésiæ príncipes decréti hujus constitútio commeávit: sed non frustra uni commendátur, quod ómnibus intimátur. Petro enim ídeo hoc singuláriter créditur, quia cunctis Ecclésiæ rectóribus Petri forma præpónitur. Manet ergo Petri privilégium, ubicúmque ex ipsíus fertur æquitáte judícium. Nec nímia est vel sevéritas vel remíssio, ubi nihil erit ligátum, nihil solútum, nisi quod beátus Petrus aut sólverit, aut ligáverit.\par
\switchcolumn
\EN
\noindent Against this confession the gates of hell shall never prevail, neither shall the bonds of death take hold upon it. Thus saith He That is faithful and true. And as this confession hath power to lift up to heaven them that make it, so is it able to thrust down to hell them that gainsay it. Wherefore it is said unto the most blessed Peter: And I will give unto thee the keys of the kingdom of heaven and whatsoever thou shalt bind on earth, shall be bound in heaven; and whatsoever thou shalt loose on earth, shall be loosed in heaven. This power passed indeed to the other Apostles also; this the Lord's will had effect in them; but it is not in vain that it is written that that was given to one which passed from him to all. To Peter alone were the keys given, and Peter is set as the pattern for all them that bear rule in the Church to follow. There remaineth therefore the right of Peter, wheresoever his judgment decreeth justice. Neither is there anything too hard, or too lax, where there is nothing bound and nothing loosed, save when Peter bindeth or looseth.\par
\switchcolumn*

\LA
\TeDeum{Te Deum}
\switchcolumn
\EN
\TeDeum{Te Deum}
\switchcolumn*

\end{paracol}

\par\needspace{10\baselineskip}\addvspace{1.2em}{\centering\small\textsc{Die 23 Februarii} · \textsc{23 February}\par}
\Office{In Vigilia S. Matthiæ Ap.}{Vigil of St. Matthias, Apostle}{Vigilia}
\addcontentsline{toc}{subsection}{Die 23 Februarii · In Vigilia S. Matthiæ Ap.\enspace\textit{Vigil of St. Matthias, Apostle}}
\begin{paracol}{2}
\LA
\RefLine{Lectiones I–III: de Scriptura occurrente.}
\switchcolumn
\EN
\RefLine{Lessons I–III: from the Scripture of the day.}
\switchcolumn*

\end{paracol}

\par\needspace{10\baselineskip}\addvspace{1.2em}{\centering\small\textsc{Die 23 Februarii} · \textsc{23 February}\par}
\Office{S. Petri Damiani Episcopi Confessoris et Ecclesiæ Doctoris}{St. Peter Damian, Confessor and Doctor of the Church}{Duplex}
\addcontentsline{toc}{subsection}{Die 23 Februarii · S. Petri Damiani Episcopi Confessoris et Ecclesiæ Doctoris\enspace\textit{St. Peter Damian, Confessor and Doctor of the Church}}
\begin{paracol}{2}
\LA
\RefLine{Lectiones I–III: ex Commúni Doctoris Pontificis (C4a).}
\switchcolumn
\EN
\RefLine{Lessons I–III: from the Common of a Doctor Bishop (C4a).}
\switchcolumn*

\LA
\RefLine{Lectiones VII–VIII: ex Commúni Doctoris Pontificis (C4a).}
\switchcolumn
\EN
\RefLine{Lessons VII–VIII: from the Common of a Doctor Bishop (C4a).}
\switchcolumn*

\LA
\RefLine{Lectio IX: de Scriptura occurrente.}
\switchcolumn
\EN
\RefLine{Lesson IX: from the Scripture of the day.}
\switchcolumn*

\LA
\LessonHead{Lectio IV}
\switchcolumn
\EN
\LessonHead{Lesson IV}
\switchcolumn*

\LA
\noindent Petrus, Ravennæ honestis paréntibus natus, adhuc lactens a matre, numerosæ prolis pertæsa, abjícitur; sed domesticæ mulíeris ópera semivivus exceptus ac recreatus, genitrici ad humanitátis sensum revocátæ redditur. Utroque orbátus parente, tamquam vile mancipium sub aspera fratris tutela duram servitútem exercuit. Religiónis in Deum ac pietátis erga patrem, egregium tunc specimen dedit; invéntum síquidem forte nummum, non propriæ inediæ sublevandæ, sed sacerdoti, qui divinum sacrifícium ad illíus expiatiónem offerret, erogávit. A Damiáno fratre, a quo, uti fertur, cognoméntum accepit, benígne recéptus, ejus cura litteris erudítur in quibus brevi tantum profecit, ut magistris admiratióni esset. Cum autem liberálibus sciéntiis floréret et nómine, eas cum laude docuit. Interim ut corpus ratióni súbderet, sub mollibus vestibus cilícium adhibuit; jejuniis, vigiliis et oratiónibus solerter insistens. Calénte juvénta dum carnis stimulis acriter urgerétur, insultántium libidinum faces rigéntibus fluvii mersus aquis noctu exstinguebat; tum venerabília quæque loca obire, totumque Psaltérium recitare consueverat. Ope assidua páuperes levabat, quibus frequenter pastis convivio, propriis ipse mánibus ministrabat.\par
\switchcolumn
\EN
\noindent The holy Doctor Peter Damian was born of respectable parents at Ravenna, (about the year of our Lord 988.) While he was still a suckling, his mother, overcome with the care of many children, cast him out to perish, but one of the women servants saved him when he was nigh to death, and fed him until natural affection appeared again in his mother, to whom she then gave him back. After the death of both his parents he lived with a brother who treated him like the lowest slave, and in whose house he underwent a hard bondage. Even while he was in this condition he gave a wonderful proof of his faith toward God, and his dutiful love toward his father. It chanced that one day he found a considerable sum of money, but instead of using it to relieve his own poverty, he gave it all to a priest to offer God's sacrifice for the forgiveness of his father's sins. He had happily another brother called Damian, the same from whom he seemeth afterwards to have taken his surname. By him he was affectionately adopted, and put in the way of being educated. He made such progress in learning as astonished his teachers, and when he had won an eminent name in letters, he began to teach on his own accord with general applause. Meanwhile, lest his body should get the better of his mind, he constantly wore a hair-shirt under his softer clothes, and exercised himself in fasting, watching, and prayer. In the spring-time of his age he was grievously tormented by the stings of the flesh; and sometimes, when the rebellions of lust seemed about to get the mastery over him at night, he threw himself into a freezing stream to check them. After this he would go about visiting consecrated places, and repeat the whole book of Psalms. He was most careful in relieving the poor, on whom he would wait with his own hands.\par
\switchcolumn*

\LA
\begin{Resp}\Rsign Invéni David servum meum, óleo sancto meo unxi eum: \Star{} Manus enim mea auxiliábitur ei.\end{Resp}
\begin{Resp}\Vsign Nihil profíciet inimícus in eo, et fílius iniquitátis non nocébit ei.\end{Resp}
\begin{Resp}\Rsign Manus enim mea auxiliábitur ei.\end{Resp}
\switchcolumn
\EN
\begin{Resp}\Rsign I have found David My servant, with My holy oil have I anointed him \Star{} For My hand shall help him.\end{Resp}
\begin{Resp}\Vsign The enemy shall prevail nothing against him, nor the son of wickedness afflict him.\end{Resp}
\begin{Resp}\Rsign For My hand shall help him.\end{Resp}
\switchcolumn*

\LA
\LessonHead{Lectio V}
\switchcolumn
\EN
\LessonHead{Lesson V}
\switchcolumn*

\LA
\noindent Perficiendæ magis vitæ causa, in Avellanénsi Eugubinæ diœcesis cœnobio, ordini monachórum sanctæ Crucis Fontis Avellanæ a beato Ludulpho, sancti Romualdi discipulo, fundato, nomen dedit. Non ita multo post in monastérium Pomposianum, mox in cœnobium sancti Vincentii Petræ Pertusæ ab abbate suo missus, utrumque ascetérium verbo sacro, præcláris institutiónibus et móribus excoluit. Ad suos revocatus, post præsidis obitum, Avellanitárum famíliæ præfícitur, quam, novis variis in locis exstructis domiciliis et sanctíssimis institútis ita auxit, ut alter ejus ordinis parens ac præcipuum ornaméntum jure sit habitus. Salutarem Petri sollicitúdinem alia quoque diversi institúti cœnobia, canonicórum convéntus, et pópuli sunt experti. Urbináti diœcesi non uno nómine profuit: Theuzóni episcopo in causa gravíssima assedit, ipsumque in recte administrando episcopatu consílio et ópera juvit. Divinórum contemplatióne, corporis maceratiónibus, ceterisque spectátæ sanctimoniæ exemplis excelluit. His motus Stephanus nonus Pontifex maximus eum, licet invitum et reluctántem, sanctæ Romanæ Ecclésiæ cardinalem creávit et Ostiensem episcopum. Quas Petrus dignitates splendidíssimis virtútibus et consentáneis episcopali ministerio opéribus gessit.\par
\switchcolumn
\EN
\noindent Desiring to attain to perfection of life he betook himself to the convent of Font-Avellano, in the diocese of Gubbio, in Umbria, a house founded by the blessed Ludolph, the disciple of St. Romuald, for the monks of the Holy Cross. He dwelt there not long before he was sent by his Abbat, first to the Abbey of Pomposia, and, secondly, to that of St. Vincent at Pietra Pertusa, both which brotherhoods he greatly profited by his godly exhortations, discreet rules, and grave manners. After his return home, and the death of his Superior, he was chosen to rule the brethren of Avellano. Here he founded diverse new hermitages, and made the community so to flourish under his saintly direction, that he is esteemed the second Father and chief ornament of that Order. This healthful care of Peter was made a blessing to convents of other Rules than his own, to houses of Canons, and to the people. He was many ways profitable to the diocese of Urbino. He sat with Theuzo the Bishop of that See to judge of a most weighty matter, and led him by his counsel and assistance rightly to administer his Bishopric. He was foremost in contemplation of the things of God, in severity toward his own body, and in other things whereby to set a bright example of godliness. In consideration of these things the Supreme Pontiff Stephen IX., (in the year 1057,) created him, in spite of his own unwillingness and objections, a Cardinal of the Holy Roman Church, and appointed him Bishop of Ostia. This dignity Peter bore with the highest reputation for piety, and adorned with works meet for a Bishop.\par
\switchcolumn*

\LA
\begin{Resp}\Rsign Pósui adjutórium super poténtem, et exaltávi eléctum de plebe mea: \Star{} Manus enim mea auxiliábitur ei.\end{Resp}
\begin{Resp}\Vsign Invéni David servum meum, óleo sancto meo unxi eum.\end{Resp}
\begin{Resp}\Rsign Manus enim mea auxiliábitur ei.\end{Resp}
\switchcolumn
\EN
\begin{Resp}\Rsign I have laid help upon one that is mighty, and have exalted one chosen out of My people \Star{} For My hand shall help him.\end{Resp}
\begin{Resp}\Vsign I have found David My servant, with My holy oil have I anointed him.\end{Resp}
\begin{Resp}\Rsign For My hand shall help him.\end{Resp}
\switchcolumn*

\LA
\LessonHead{Lectio VI}
\switchcolumn
\EN
\LessonHead{Lesson VI}
\switchcolumn*

\LA
\noindent Dificillimo témpore Romanæ Ecclésiæ, summisque Pontificibus doctrina, legatiónibus, aliisque susceptis labóribus mirifice adfuit. Adversus Nicolaitárum et simoníacam hæreses ad mortem usque strenue decertávit. Hujusmodi depulsis malis, Mediolanensem Ecclésiæ Romanæ conciliávit. Benedicto et Cadalóo falsis pontificibus fortiter restitit: Henricum quartum Germaniæ regem ab iníquo uxoris divortio deterruit: Ravennates ad débita Romano Pontifici obsequia revocatos sacris restituit: canonicos Veliternos ad sanctioris vitæ leges compósuit. In provincia præsertim Urbinate vix ulla fuit episcopalis ecclésia, de qua Petrus non sit bene meritus: Eugubínam, quam aliquándo creditam hábuit, multis levávit incommodis; alias álibi, quando oportuit, perinde curávit, ac suæ essent tutelæ commissæ. Cardinalatu et episcopali dignitate depositis, nihil de prístina juvándi próximos sedulitate remisit. Jejunium sextæ fériæ in honórem sanctæ Crucis Jesu Christi, horarias beátæ Dei Genitrícis preces, ejusque die Sabbato cultum propagávit. Inferendæ quoque sibi verberatiónis morem ad patratórum scelerum expiatiónem provéxit. Demum sanctitáte, doctrina, miraculis et præcláre actis illustris, dum e Ravennate legatióne rediret, Favéntiæ octavo Kalendas Martii migrávit ad Christum. Ejus corpus ibidem apud Cistercienses, multis miraculis clarum, frequénti populórum veneratióne cólitur. Ipsum Faventini, non semel in præsénti discrimine propitium apud Deum delegérunt. Leo vero duodecimus Pontifex maximus Offícium Missamque in ejus honórem tamquam Confessóris Pontificis, quæ aliquibus in diœcesibus atque in ordine Camalulensium jam celebrabántur, ex sacrórum Rituum Congregatiónis consulto, áddita Doctoris qualitate, ad universam exténdit Ecclésiam.\par
\switchcolumn
\EN
\noindent At the most anxious times he greatly sustained the Church of Rome and the Supreme Pontiffs by his teaching, by missions which he discharged, and by diverse other labours which he undertook on their behalf. He strove manfully even unto death against the heresies of the Nicola'itans and the Simoniacs, by putting down which evils he reconciled the Church of Milan to that of Rome. He was one of the stoutest opponents of the false Popes Benedict and Cadalous. He deterred Henry IV., King of Germany, from his wicked scheme for putting away his wife. He recalled the people of Ravenna to their bounden duty to the Bishop of Rome, and restored them to the communion of the Church. He reformed the Canons of Velletri, and brought them to lead more godly lives. There were hardly any Cathedral Churches, especially in the province of Urbino, of which he did not deserve well. In Gubbio, of which he had at one time the management, he abolished many things unseemly. He brought about improvements in many and diverse places, as if each were his special charge. (In 1062) he gave up his dignities of Cardinal and Bishop, but he allowed his love toward his neighbours to know no diminution. He was particularly zealous in spreading abroad four devout practices 1st, To fast every Friday in honour of the Holy Cross of Jesus Christ; 2nd, To recite the Hours of the Blessed Mother of God, called also her Little Office; 3rd, To sanctify Saturday in her honour; and 4th, and especially, to scourge oneself in punishment for sin committed. At length he departed to be with Christ, at Faenza, on his way back from his mission to Ravenna, on the 22nd of February, (in the year 1072,) at the height of his reputation for holiness, learning, miracles, and good works. His body is buried in the house of the Cistercians at Faenza, where the people resort often to his grave with great reverence. The citizens of Faenza, to whom he hath been found good at need even to this day, have chosen him for their Patron in the presence of God. The supreme Pontiff Leo XII., finding that an Office and Mass in memory of him, as a Confessor and Bishop, was in use in some dioceses, and in the Camaldolese Order, by advice of the Sacred Congregation of Rites, added the title of Doctor, and extended the use of the said Office and Mass to the whole Church.\par
\switchcolumn*

\LA
\begin{Resp}\Rsign Iste est, qui ante Deum magnas virtútes operátus est, et omnis terra doctrína ejus repléta est: \Star{} Ipse intercédat pro peccátis ómnium populórum.\end{Resp}
\begin{Resp}\Vsign Iste est, qui contémpsit vitam mundi, et pervénit ad cæléstia regna.\end{Resp}
\begin{Resp}\Rsign Ipse intercédat pro peccátis ómnium populórum.\end{Resp}
\begin{Resp}\Vsign Glória Patri, et Fílio, \Star{} et Spirítui Sancto.\end{Resp}
\begin{Resp}\Rsign Ipse intercédat pro peccátis ómnium populórum.\end{Resp}
\switchcolumn
\EN
\begin{Resp}\Rsign This is he which wrought great wonders before God, and the whole earth is full of his teaching \Star{} May he pray for all people, that their sins may be forgiven unto them\end{Resp}
\begin{Resp}\Vsign This is he which loved not his life in this world, and hath attained unto the kingdom of heaven.\end{Resp}
\begin{Resp}\Rsign May he pray for all people, that their sins may be forgiven unto them\end{Resp}
\begin{Resp}\Vsign Glory be to the Father, and to the Son, \Star{} and to the Holy Ghost.\end{Resp}
\begin{Resp}\Rsign May he pray for all people, that their sins may be forgiven unto them\end{Resp}
\switchcolumn*

\end{paracol}

\par\needspace{10\baselineskip}\addvspace{1.2em}{\centering\small\textsc{Die 24 Februarii} · \textsc{24 February}\par}
\Office{S. Matthiæ Apostoli}{St. Matthias the Apostle}{Duplex II classis}
\addcontentsline{toc}{subsection}{Die 24 Februarii · S. Matthiæ Apostoli\enspace\textit{St. Matthias the Apostle}}
\begin{paracol}{2}
\LA
\RefLine{Lectio IX: de Scriptura occurrente.}
\switchcolumn
\EN
\RefLine{Lesson IX: from the Scripture of the day.}
\switchcolumn*

\LA
\LessonHead{Lectio I}
\switchcolumn
\EN
\LessonHead{Lesson I}
\switchcolumn*

\LA
\LessonTitle{De Actibus Apostolórum}
\switchcolumn
\EN
\LessonTitle{From the Acts of the Apostles}
\switchcolumn*

\LA
\Cite{Acts 1:15-18}
\switchcolumn
\EN
\Cite{Acts 1:15-18}
\switchcolumn*

\LA
\noindent In diébus illis, exsúrgens Petrus in médio fratrum, dixit (erat autem turba hóminum simul, fere centum vigínti): Viri fratres, opórtet impléri Scriptúram quam prædíxit Spíritus Sanctus per os David de Juda, qui fuit dux eórum qui comprehendérunt Jesum: qui connumerátus erat in nobis, et sortítus est sortem ministérii hujus. Et hic quidem possédit agrum de mercéde iniquitátis, et suspénsus crépuit médius: et diffúsa sunt ómnia víscera ejus.\par
\switchcolumn
\EN
\noindent In those days Peter rising up in the midst of the brethren, said: now the number of persons together was about an hundred and twenty: Men, brethren, the scripture must needs be fulfilled, which the Holy Ghost spoke before by the mouth of David concerning Judas, who was the leader of them that apprehended Jesus: Who was numbered with us, and had obtained part of this ministry. And he indeed hath possessed a field of the reward of iniquity, and being hanged, burst asunder in the midst: and all his bowels gushed out.\par
\switchcolumn*

\LA
\begin{Resp}\Rsign Ecce ego mitto vos sicut oves in médio lupórum, dicit Dóminus: \Star{} Estóte ergo prudéntes sicut serpéntes, et símplices sicut colúmbæ.\end{Resp}
\begin{Resp}\Vsign Dum lucem habétis, crédite in lucem, ut fílii lucis sitis.\end{Resp}
\begin{Resp}\Rsign Estóte ergo prudéntes sicut serpéntes, et símplices sicut colúmbæ.\end{Resp}
\switchcolumn
\EN
\begin{Resp}\Rsign Behold, I send you as sheep in the midst of wolves, said the Lord: \Star{} Be ye, therefore, wise as serpents and simple as doves.\end{Resp}
\begin{Resp}\Vsign Whilst you have the light, believe in the light, that you may be the children of light.\end{Resp}
\begin{Resp}\Rsign Be ye therefore wise as serpents and simple as doves.\end{Resp}
\switchcolumn*

\LA
\LessonHead{Lectio II}
\switchcolumn
\EN
\LessonHead{Lesson II}
\switchcolumn*

\LA
\Cite{Acts 1:19-22}
\switchcolumn
\EN
\Cite{Acts 1:19-22}
\switchcolumn*

\LA
\noindent Et notum factum est ómnibus habitántibus Jerúsalem, ita ut appellarétur ager ille, lingua eórum, Hacéldama, hoc est, ager sánguinis. Scriptum est enim in libro Psalmórum: Fiat commorátio eórum desérta, et non sit qui inhábitet in ea: et episcopátum ejus accípiat alter. Opórtet ergo ex his viris qui nobíscum sunt congregáti in omni témpore quo intrávit et exívit inter nos Dóminus Jesus, incípiens a baptísmate Joánnis usque in diem qua assúmptus est a nobis, testem resurrectiónis ejus nobíscum fíeri unum ex istis.\par
\switchcolumn
\EN
\noindent And it became known to all the inhabitants of Jerusalem: so that the same field was called in their tongue, Haceldama, that is to say, The field of blood. For it is written in the book of Psalms: Let their habitation become desolate, and let there be none to dwell therein. And his bishopric let another take. Wherefore of these men who have companied with us all the time that the Lord Jesus came in and went out among us, Beginning from the baptism of John, until the day wherein he was taken up from us, one of these must be made a witness with us of his resurrection.\par
\switchcolumn*

\LA
\begin{Resp}\Rsign Tóllite jugum meum super vos, dicit Dóminus, et díscite a me, quia mitis sum et húmilis corde: \Star{} Jugum enim meum suáve est, et onus meum leve.\end{Resp}
\begin{Resp}\Vsign Et inveniétis réquiem animábus vestris.\end{Resp}
\begin{Resp}\Rsign Jugum enim meum suáve est, et onus meum leve.\end{Resp}
\switchcolumn
\EN
\begin{Resp}\Rsign Take up my yoke upon you, said the Lord, and learn of me, because I am meek, and humble of heart. \Star{} For my yoke is sweet and my burden light.\end{Resp}
\begin{Resp}\Vsign And you shall find rest to your souls.\end{Resp}
\begin{Resp}\Rsign For my yoke is sweet and my burden light.\end{Resp}
\switchcolumn*

\LA
\LessonHead{Lectio III}
\switchcolumn
\EN
\LessonHead{Lesson III}
\switchcolumn*

\LA
\Cite{Acts 1:23-26}
\switchcolumn
\EN
\Cite{Acts 1:23-26}
\switchcolumn*

\LA
\noindent Et statuérunt duos, Joseph, qui vocabátur Bársabas, qui cognominátus est Justus, et Matthíam. Et orántes dixérunt: Tu Dómine, qui corda nosti ómnium, osténde quem elégeris ex his duóbus unum, accípere locum ministérii hujus et apostolátus, de quo prævaricátus est Judas ut abíret in locum suum. Et dedérunt sortes eis, et cécidit sors super Matthíam: et annumerátus est cum úndecim Apóstolis.\par
\switchcolumn
\EN
\noindent And they appointed two, Joseph, called Barsabas, who was surnamed Justus, and Matthias. And praying, they said: Thou, Lord, who knowest the hearts of all men, shew whether of these two thou hast chosen, To take the place of this ministry and apostleship, from which Judas hath by transgression fallen, that he might go to his own place. And they gave them lots, and the lot fell upon Matthias, and he was numbered with the eleven apostles.\par
\switchcolumn*

\LA
\begin{Resp}\Rsign Dum stetéritis ante reges et prǽsides, nolíte cogitáre quómodo aut quid loquámini: \Star{} Dábitur enim vobis in illa hora, quid loquámini.\end{Resp}
\begin{Resp}\Vsign Non enim vos estis qui loquímini: sed Spíritus Patris vestri, qui lóquitur in vobis.\end{Resp}
\begin{Resp}\Rsign Dábitur enim vobis in illa hora, quid loquámini.\end{Resp}
\begin{Resp}\Vsign Glória Patri, et Fílio, \Star{} et Spirítui Sancto.\end{Resp}
\begin{Resp}\Rsign Dábitur enim vobis in illa hora, quid loquámini.\end{Resp}
\switchcolumn
\EN
\begin{Resp}\Rsign But when they shall deliver you up to the judges, take no thought how or what to speak: \Star{} for it shall be given you in that hour what to speak:\end{Resp}
\begin{Resp}\Vsign For it is not you that speak, but the spirit of your Father that speaketh in you.\end{Resp}
\begin{Resp}\Rsign For it shall be given you in that hour what to speak.\end{Resp}
\begin{Resp}\Vsign Glory be to the Father, and to the Son, \Star{} and to the Holy Ghost.\end{Resp}
\begin{Resp}\Rsign For it shall be given you in that hour what to speak.\end{Resp}
\switchcolumn*

\LA
\LessonHead{Lectio IV}
\switchcolumn
\EN
\LessonHead{Lesson IV}
\switchcolumn*

\LA
\LessonTitle{De Expositióne sancti Augustíni Epíscopi super psalmum octogésimum sextum}
\switchcolumn
\EN
\LessonTitle{From the Exposition of the Eighty-sixth Psalm by St. Augustine, Bishop (of Hippo.)}
\switchcolumn*

\LA
\Source{Ante medium}
\switchcolumn
\EN
\Source{Ante medium}
\switchcolumn*

\LA
\noindent Fundaménta ejus in móntibus sanctis: díligit Dóminus portas Sion. Quare sunt fundaménta Apóstoli et Prophétæ? Quia eórum auctóritas portat infirmitátem nostram. Quare sunt portæ? Quia per ipsos intrámus ad regnum Dei. Prǽdicant enim nobis: et, cum per ipsos intrámus, per Christum intrámus; ipse est enim jánua. Et cum dicúntur duódecim portæ Jerúsalem, et una porta Christus et duódecim portæ Christus, quia in duódecim portis Christus; et ídeo duodenárius númerus Apostolórum. Sacraméntum magnum hujus duodenárii significátio est númeri. Sedébitis, inquit, super duódecim sedes, judicántes duódecim tribus Israël.\par
\switchcolumn
\EN
\noindent Her foundation is in the holy mountains; the Lord loveth the gates of Zion. Wherefore hath the city twelve foundations, and in them the names of the Prophets and of the Apostles of the Lamb? Because their authority is the foundation whereon our weakness resteth. Wherefore are they the gates? Because through them we enter in unto the kingdom of God, since they have preached the same unto us, and when we enter in through their preaching, we enter in by Christ, Who is Himself The Door. (John x. 7.) And, whereas it is written that the city hath twelve gates, and, again, that Christ is the one Door, Christ is all the twelve, for He is in all the twelve and therefore were twelve Apostles chosen. There lieth a great mystery in the signification of this number twelve: Ye shall sit, said the Lord, upon twelve thrones, judging the twelve tribes of Israel.\par
\switchcolumn*

\LA
\begin{Resp}\Rsign Vidi conjúnctos viros, habéntes spléndidas vestes, et Ángelus Dómini locútus est ad me, dicens: \Star{} Isti sunt viri sancti facti amíci Dei.\end{Resp}
\begin{Resp}\Vsign Vidi Ángelum Dei fortem, volántem per médium cælum, voce magna clamántem et dicéntem.\end{Resp}
\begin{Resp}\Rsign Isti sunt viri sancti facti amíci Dei.\end{Resp}
\switchcolumn
\EN
\begin{Resp}\Rsign I saw men standing together, clad in shining raiment, and the Angel of the Lord spoke unto me, saying: \Star{} These men are holy, for they are the friends of God.\end{Resp}
\begin{Resp}\Vsign I saw a strong Angel of God fly into the midst of heaven, saying with a loud voice:\end{Resp}
\begin{Resp}\Rsign These men are holy, for they are the friends of God.\end{Resp}
\switchcolumn*

\LA
\LessonHead{Lectio V}
\switchcolumn
\EN
\LessonHead{Lesson V}
\switchcolumn*

\LA
\noindent Si duódecim sellæ ibi sunt, non est ubi sédeat tértius décimus Paulus Apóstolus, et non erit quómodo júdicet; et ipse se judicatúrum dixit, non hómines tantum, sed et ángelos. Quos ángelos, nisi apóstatas ángelos? Nescítis, inquit, quia ángelos judicábimus? Respondéret ergo turba: Quid te jactas judicatúrum? Ubi sedébis? Duódecim sedes dixit Dóminus duódecim Apóstolis, unus cécidit Judas, in locum ipsíus sanctus Matthías ordinátus est; implétus est duodenárius númerus sédium. Primo locum invéni, ubi sédeas; et sic te mináre judicatúrum. Duódecim ergo sedes quid sibi velint, videámus. Sacraméntum est cujúsdam universitátis; quia per totum orbem terrárum futúra erat Ecclésia, unde vocátur hoc ædifícium ad Christi compágem.\par
\switchcolumn
\EN
\noindent If then there be set there twelve thrones of judgment, (Ps. cxxi. 5,) Paul, in that he is the thirteenth Apostle, hath not where to sit, nor wherein to judge. Nevertheless, he hath said of himself that he will judge not men only, but angels. Know ye not, saith he, that we shall judge angels? (i Cor. vi. 3,) that is, the fallen angels. Then might they have answered him: Wherefore boastest thou thyself to be a judge? For where is thy seat? The Lord hath said that for the twelve Apostles there shall be twelve thrones; one of the twelve, even Judas, is indeed fallen, but holy Matthias is chosen into his place; for the twelve thrones there are still twelve to sit thereon first find whereon thou shalt sit, and afterward give thyself out for a judge. Let us see, then, what is the meaning of these twelve thrones. By them is signified in a mystery the whole world, since the Church shall be through all the earth, whence this building is called to be built up together in Christ.\par
\switchcolumn*

\LA
\begin{Resp}\Rsign Beáti estis, cum maledíxerint vobis hómines, et persecúti vos fúerint, et díxerint omne malum advérsum vos, mentiéntes, propter me: \Star{} Gaudéte et exsultáte, quóniam merces vestra copiósa est in cælis.\end{Resp}
\begin{Resp}\Vsign Cum vos óderint hómines, et cum separáverint vos, et exprobráverint, et ejécerint nomen vestrum tamquam malum propter Fílium hóminis.\end{Resp}
\begin{Resp}\Rsign Gaudéte et exsultáte, quóniam merces vestra copiósa est in cælis.\end{Resp}
\switchcolumn
\EN
\begin{Resp}\Rsign Blessed are ye when men shall revile you, and persecute you, and shall say all manner of evil against you falsely, for My sake; \Star{} Rejoice, and be exceeding glad, for great is your reward in heaven.\end{Resp}
\begin{Resp}\Vsign When men shall hate you, and when they shall separate you from their company, and shall reproach you, and cast out your name as evil, for the Son of Man's sake.\end{Resp}
\begin{Resp}\Rsign Rejoice, and be exceeding glad, for great is your reward in heaven.\end{Resp}
\switchcolumn*

\LA
\LessonHead{Lectio VI}
\switchcolumn
\EN
\LessonHead{Lesson VI}
\switchcolumn*

\LA
\noindent Et ídeo, quia úndique venítur ad judicándum, duódecim sedes sunt; sicut, quia úndique intrátur in illam civitátem, duódecim portæ sunt. Non solum ergo illi duódecim et Apóstolus Paulus, sed quotquot judicatúri sunt, propter significatiónem universitátis, ad sedes duódecim pértinent; quemádmodum quotquot intrábunt, ad duódecim portas pértinent. Partes enim mundi quátuor sunt, Oriens, Occidens, Aquilo et Merídies. Istæ quátuor partes assídue nominántur in Scriptúris. Ab istis quátuor ventis, sicut dicit Dóminus in Evangélio, a quátuor ventis se collectúrum eléctos suos; ab ómnibus ergo istis quátuor ventis vocátur Ecclésia. Quómodo vocátur? Undique in Trinitáte vocátur. Non vocátur nisi per baptísmum in nómine Patris, et Fílii, et Spíritus Sancti. Quátuor ergo ter ducta duódecim inveniúntur.\par
\switchcolumn
\EN
\noindent Therefore is it said that there shall be twelve thrones, because from all quarters shall there come men to be judged; even as it is said that the city hath twelve gates, because from all quarters shall the nations of them which are saved, enter into it. So, not the twelve only, and the Apostle Paul, but all, as many as shall judge, have part in these twelve thrones, this signifying, that they shall judge all men; even as all that enter into the city, have part in her twelve gates. For there are four quarters of the world, the East, and the West, and the North, and the South of which four quarters is mention often made in the Scriptures. From the four winds shall the elect be gathered together, as saith the Lord in the Gospel: And He shall send His Angels with a great sound of a trumpet; and they shall gather together His elect from the four winds, from one end of heaven to the other. (Matth. xxiv. 31.) From the four winds, therefore, is the Church called together; and how are they called? Everywhere are they called in the Trinity; for they are called no otherwise than by baptizing them in the Name of the Father, and of the Son, and of the Holy Ghost. (Matth. xxvii. 19.) Now four being multiplied by three is twelve.\par
\switchcolumn*

\LA
\begin{Resp}\Rsign Isti sunt triumphatóres et amíci Dei, qui contemnéntes jussa príncipum, meruérunt prǽmia ætérna: \Star{} Modo coronántur, et accípiunt palmam.\end{Resp}
\begin{Resp}\Vsign Isti sunt qui venérunt ex magna tribulatióne, et lavérunt stolas suas in sánguine Agni.\end{Resp}
\begin{Resp}\Rsign Modo coronántur, et accípiunt palmam.\end{Resp}
\begin{Resp}\Vsign Glória Patri, et Fílio, \Star{} et Spirítui Sancto.\end{Resp}
\begin{Resp}\Rsign Modo coronántur, et accípiunt palmam.\end{Resp}
\switchcolumn
\EN
\begin{Resp}\Rsign These are they which have conquered, and are become the friends of God, who recked not of the commandments of princes, and earned the everlasting reward. \Star{} And now have they crowns on their heads, and palms in their hands.\end{Resp}
\begin{Resp}\Vsign These are they which came out of great tribulation, and have washed their robes in the blood of the Lamb.\end{Resp}
\begin{Resp}\Rsign And now have they crowns on their heads, and palms in their hands.\end{Resp}
\begin{Resp}\Vsign Glory be to the Father, and to the Son, \Star{} and to the Holy Ghost.\end{Resp}
\begin{Resp}\Rsign And now have they crowns on their heads, and palms in their hands.\end{Resp}
\switchcolumn*

\LA
\LessonHead{Lectio VII}
\switchcolumn
\EN
\LessonHead{Lesson VII}
\switchcolumn*

\LA
\LessonTitle{Léctio sancti Evangélii secúndum Matthǽum}
\switchcolumn
\EN
\LessonTitle{From the holy Gospel according to Matthew}
\switchcolumn*

\LA
\Cite{Matt 11:25-30}
\switchcolumn
\EN
\Cite{Matt 11:25-30}
\switchcolumn*

\LA
\noindent In illo témpore: Respóndens Jesus dixit: Confíteor tibi, Pater, Dómine cæli et terræ, quia abscondísti hæc a sapiéntibus et prudéntibus, et revelásti ea párvulis. Et réliqua.\par
\switchcolumn
\EN
\noindent At that time, Jesus answered and said: I confess to thee, O Father, Lord of heaven and earth, because thou hast hid these things from the wise and prudent, and hast revealed them to the little ones. And so on.\par
\switchcolumn*

\LA
\LessonTitle{Homilía sancti Augustíni Epíscopi}
\switchcolumn
\EN
\LessonTitle{Homily by St. Augustine, Bishop (of Hippo.)}
\switchcolumn*

\LA
\Source{Sermo 10 de verbis Domini}
\switchcolumn
\EN
\Source{10th Sermon on the Words of the Lord}
\switchcolumn*

\LA
\noindent Veníte ad me, omnes qui laborátis. Quare enim omnes laborámus, nisi quia sumus hómines mortáles, frágiles, infírmi, lútea vasa portántes, quæ fáciunt ínvicem angústias? Sed, si angustiántur vasa carnis, dilaténtur spátia caritátis. Quid ergo dicit, Veníte ad me, omnes qui laborátis, nisi ut non laborétis? Dénique promíssio ejus in promptu est; quóniam laborántes vocávit, quærent forte, qua mercéde vocáti sunt. Et ego vos, inquit, refíciam. Tóllite jugum meum super vos, et díscite a me, non mundum fabricáre, non cuncta visibília et invisibília creáre, non in ipso mundo mirabília fácere et mórtuos suscitáre; sed, Quóniam mitis sum et húmilis corde.\par
\switchcolumn
\EN
\noindent Come unto Me, all ye that labour! And wherefore labour we all, but because we are frail, sickly, dying creatures, burdened with earthen vessels which distress us? But if these fleshly vessels be distressful, let the open expanse of love be free and wide. Come unto Me, all ye that labour! and why? That we may labour no more. His promise is an instant promise, for He calleth such as are labouring. Perchance they will ask Him what shall be their reward? And I, saith He, will give you rest. Take My yoke upon you, and learn of me not how to make the world, not how to create all things visible and invisible, not to work wonders in the earth, nor to raise the dead but for I am meek and lowly in heart.\par
\switchcolumn*

\LA
\begin{Resp}\Rsign Isti sunt qui vivéntes in carne, plantavérunt Ecclésiam sánguine suo: \Star{} Cálicem Dómini bibérunt, et amíci Dei facti sunt.\end{Resp}
\begin{Resp}\Vsign In omnem terram exívit sonus eórum, et in fines orbis terræ verba eórum.\end{Resp}
\begin{Resp}\Rsign Cálicem Dómini bibérunt, et amíci Dei facti sunt.\end{Resp}
\switchcolumn
\EN
\begin{Resp}\Rsign These are they who while yet they lived in the flesh, planted the Church in their own blood; \Star{} They drank of the Lord's cup, and became the friends of God.\end{Resp}
\begin{Resp}\Vsign Their sound is gone out through all the earth, and their words to the ends of the world.\end{Resp}
\begin{Resp}\Rsign They drank of the Lord's cup, and became the friends of God.\end{Resp}
\switchcolumn*

\LA
\LessonHead{Lectio VIII}
\switchcolumn
\EN
\LessonHead{Lesson VIII}
\switchcolumn*

\LA
\noindent Magnus esse vis? a mínimo íncipe. Cógitas magnam fábricam constrúere celsitúdinis? de fundaménto prius cógita humilitátis. Et quantam quisque vult et dispónit superimpónere molem ædifícii, quanto erit majus ædifícium, tanto áltius fodit fundaméntum. Et fábrica quidem cum constrúitur, in supérna consúrgit; qui autem fodit fundaméntum, ad ima deprímitur. Ergo et fábrica ante celsitúdinem humiliátur, et fastígium post humiliatiónem erígitur.\par
\switchcolumn
\EN
\noindent Wilt thou be great? Begin by being little. Dost thou think to raise up a lofty building? Then lay the foundations thereof in lowliness. The greater soever, and the more massy, be that which any man thinketh to build, so much the deeper doth he dig his foundation. And when the house is built, it towereth heavenward; but he which layeth the foundation goeth down into the earth. The building, therefore, is low before it is high, and, after it is low, it riseth high to the roof.\par
\switchcolumn*

\LA
\begin{Resp}\Rsign Isti sunt viri sancti, quos elégit Dóminus in caritáte non ficta, et dedit illis glóriam sempitérnam: \Star{} Quorum doctrína fulget Ecclésia, ut sole luna.\end{Resp}
\begin{Resp}\Vsign Sancti per fidem vicérunt regna: operáti sunt justítiam.\end{Resp}
\begin{Resp}\Rsign Quorum doctrína fulget Ecclésia, ut sole luna.\end{Resp}
\begin{Resp}\Vsign Glória Patri, et Fílio, \Star{} et Spirítui Sancto.\end{Resp}
\begin{Resp}\Rsign Quorum doctrína fulget Ecclésia, ut sole luna.\end{Resp}
\switchcolumn
\EN
\begin{Resp}\Rsign These men are saints, whom the Lord hath chosen in love unfeigned, and hath given them glory everlasting. These are they \Star{} By the light of whose teaching the Church is glorified, even as the moon is glorified by the light of the sun.\end{Resp}
\begin{Resp}\Vsign The saints through faith subdued kingdoms, wrought righteousness.\end{Resp}
\begin{Resp}\Rsign By the light of whose teaching the Church is glorified, even as the moon is glorified by the light of the sun.\end{Resp}
\begin{Resp}\Vsign Glory be to the Father, and to the Son, \Star{} and to the Holy Ghost.\end{Resp}
\begin{Resp}\Rsign By the light of whose teaching the Church is glorified, even as the moon is glorified by the light of the sun.\end{Resp}
\switchcolumn*

\LA
\LessonHead{Lectio IX}
\switchcolumn
\EN
\LessonHead{Lesson IX}
\switchcolumn*

\LA
\noindent Quod est fastígium construéndæ fábricæ, quam molímur? quo perventúrum est cacúmen ædifícii? Cito dico, usque ad conspéctum Dei. Vidétis quam excélsum est, quanta res est conspícere Deum. Qui desíderat, et quod dico et quod audit intéllegit. Promíttitur nobis conspéctus Dei, veri Dei, summi Dei. Hoc enim bonum est, vidéntem vidére. Nam, qui colunt falsos deos, fácile illos vident; sed eos vident, qui óculos habent et non vident. Nobis autem promíttitur vísio Dei vivéntis et vidéntis.\par
\switchcolumn
\EN
\noindent What is the roof of the house on which we labour? Whither do its spires rise? I answer you at once; to the presence of God. You see how high it is, yea, what it is to see God. He that will, understandeth what I say, and he heareth. What is promised you is to see God, God, the True, God, the Supreme. Blessed is he who seeth Him by Whom he is seen. Such as worship false gods see them easily, but they see them who have eyes and see not. But unto us it is promised that we shall see that God Who liveth and seeth. (Gen. xvi. 14.)\par
\switchcolumn*

\LA
\TeDeum{Te Deum}
\switchcolumn
\EN
\TeDeum{Te Deum}
\switchcolumn*

\end{paracol}

\par\needspace{10\baselineskip}\addvspace{1.2em}{\centering\small\textsc{Die 27 Februarii} · \textsc{27 February}\par}
\Office{S. Gabrielis a Virgine Perdolente Confessoris}{St. Gabriel of the Sorrowful Virgin, Confessor}{Duplex}
\addcontentsline{toc}{subsection}{Die 27 Februarii · S. Gabrielis a Virgine Perdolente Confessoris\enspace\textit{St. Gabriel of the Sorrowful Virgin, Confessor}}
\begin{paracol}{2}
\LA
\RefLine{Lectiones I–III: ex Commúni Confessoris non pontificis (C5), in 2 loco.}
\switchcolumn
\EN
\RefLine{Lessons I–III: from the Common of a Confessor not a Bishop (C5), set 2.}
\switchcolumn*

\LA
\RefLine{Lectio IX: de Scriptura occurrente.}
\switchcolumn
\EN
\RefLine{Lesson IX: from the Scripture of the day.}
\switchcolumn*

\LA
\LessonHead{Lectio IV}
\switchcolumn
\EN
\LessonHead{Lesson IV}
\switchcolumn*

\LA
\noindent Gabriel, Assisii in Umbria, honesto genere natus, et Franciscus ob seráphici civis memóriam vocatus, egregiam animi índolem a púero osténdit. Adoléscens, cum Spoleti litteris operam daret, inani sæculi spécie et pompa aliquantulum állici visus est. Sed miserentis Dei munere, qui eum ad perfectiónem christianæ vitæ jamdúdum invitabat, cum in morbum incidísset, sæculi vanitátem fastidire cœpit, atque immortalia dumtáxat bona appétere. Quo autem citius Deo vocanti obtemperaret, factum est, ut insignem illam beatíssimæ Vírginis Icónem, solemni pompa extra Spoletinæ ecclésiæ septa delatam intúitus, divini amoris flammam conciperet, simulque Institutum Clericórum a Passióne Jesu amplecti statúeret. Itaque non exíguas difficultates eluctatus, in recessu Morrovallénsi, lúgubrem vestem lætus induit, et Gabriel a Vírgine perdolénte maluit appellari; ad ejusdem gaudiórum et dolórum memóriam perpetuo recoléndam.\par
\switchcolumn
\EN
\noindent Gabriel, born at Assisi in Umbria of a respectable family, and called Francis in memory of his seraphic fellow-townsman, showed from boyhood an excellent disposition of soul. As a youth, when studying letters at Spoleto, he seemed for a time to be allured by the empty beauty and pomp of the world. But by the gift of the merciful God, who had already called him to the perfection of a Christian life when he had fallen sick, he began to tire of the vanity of the world, and to desire immortal treasures alone. But to quicken his obedience to the call of God, it happened that as he saw the celebrated Image of the Blessed Virgin being carried with solemn pomp outside the precincts of the church of Spoleto, he experienced the flame of divine love, and at the same time decided to enter the Institute of the Clerks of the Passion of Jesus. Therefore, after overcoming no slight difficulties, he joyfully donned the somber habit in the secluded place of Morrovalle, and chose to be called Gabriel of Our Lady of Sorrows, to recall forever the memory of her joys and griefs.\par
\switchcolumn*

\LA
\begin{Resp}\Rsign Honéstum fecit illum Dóminus, et custodívit eum ab inimícis, et a seductóribus tutávit illum: \Star{} Et dedit illi claritátem ætérnam.\end{Resp}
\begin{Resp}\Vsign Justum dedúxit Dóminus per vias rectas, et osténdit illi regnum Dei.\end{Resp}
\begin{Resp}\Rsign Et dedit illi claritátem ætérnam.\end{Resp}
\switchcolumn
\EN
\begin{Resp}\Rsign The Lord made him honourable, and defended him from his enemies, and kept him safe from those that lay in wait for him; \Star{} And gave him perpetual glory.\end{Resp}
\begin{Resp}\Vsign He went down with him into the pit, and left him not in bonds.\end{Resp}
\begin{Resp}\Rsign And gave him perpetual glory.\end{Resp}
\switchcolumn*

\LA
\LessonHead{Lectio V}
\switchcolumn
\EN
\LessonHead{Lesson V}
\switchcolumn*

\LA
\noindent In tirocinio, cum regulari observántia et ómnium exercitatióne virtútum cotidie magis eminéret, brevi eo pérvenit, ut absolutæ sanctimoniæ exemplar haberétur non modo a sodálibus, vel provectis, sed étiam ultra cœnobii septa, factus bonus odor Christi in omni loco. Dominicæ passiónis cultor assiduus, in ea meditanda dies noctesque insumebat. In divinam Eucharistiam, quæ ejusdem Passiónis memóriam prodit, incredibili quodam studio ferebátur; qua cum se refíceret, seraphico ardore flagrábat. Nihil autem insígnius quam ejus erga magnam Dei Parentem pietas fuit. Ipsam omni obsequii genere percólere consuevit; sed præsertim confectam afflictamque cruciátibus Jesu tam dolenter contemplari, ut vim lacrimárum profúnderet. Pérdolens Virgo quasi tota ei vitæ rátio fuit, adeptæque ab eo sanctitátis magistra; ita ut inter æquales una fúerit senténtia, ideo excitátum Dei famulum divinitus fuísse, ut cultus Maríæ perdolentis magnum exemplo ejus cáperet increméntum.\par
\switchcolumn
\EN
\noindent In the novitiate, day by day he became conspicuous for regular observance and for the exercise of all the virtues, and in a short time he came to be considered a pattern of perfect holiness, not only by his companions and his seniors, but also beyond the confines of the monastery; he became a sweet odour in Christ in every place. An assiduous devotee of the Lord's Passion, he spent days and nights meditating upon it. He was drawn by unbelievable zeal towards the Holy Eucharist, a memorial of that Passion; and when he nourished himself with it, he burned with seraphic ardour. There was nothing more noticeable than his filial piety towards the great Mother of God. He was accustomed to pay her honour for every type of devotion, but especially to contemplate her stricken and afflicted by the sufferings of Jesus, with such sorrow that he shed floods of tears. The sorrowful Virgin was, as it were, the whole reason of his being, and the teacher of the holiness that he had acquired. As a result all his associates shared the one opinion that this servant of God had been inspired from on high so that the cult of St. Mary of Sorrows through his example might receive a great increase.\par
\switchcolumn*

\LA
\begin{Resp}\Rsign Amávit eum Dóminus, et ornávit eum: stolam glóriæ índuit eum, \Star{} Et ad portas paradísi coronávit eum.\end{Resp}
\begin{Resp}\Vsign Índuit eum Dóminus lorícam fídei, et ornávit eum.\end{Resp}
\begin{Resp}\Rsign Et ad portas paradísi coronávit eum.\end{Resp}
\switchcolumn
\EN
\begin{Resp}\Rsign The Lord loved him and beautified him; He clothed him with a robe of glory, \Star{} And crowned him at the gates of Paradise.\end{Resp}
\begin{Resp}\Vsign The Lord hath put on him the breast-plate of faith, and hath adorned him.\end{Resp}
\begin{Resp}\Rsign And crowned him at the gates of Paradise.\end{Resp}
\switchcolumn*

\LA
\LessonHead{Lectio VI}
\switchcolumn
\EN
\LessonHead{Lesson VI}
\switchcolumn*

\LA
\noindent Inter ceteras virtútes christianam humilitátem et obœdiéntiam maxime diléxit: nam inter omnes se minimum exístimans, abjectíssima quæque ministeria domus cupide affectabat, et antistitum suórum non modo jussa, sed et optáta diligentíssime perficiebat. Idem, refrenátis sensibus et vitæ asperitáte usus, illibátum retinuit florem virginitátis ac plane mundo crucifixus unice Deo vixit, intima Dómini sui fruítus consuetúdine. Ita brevem vitæ cursum tot virtútibus nobilitátum conficiens, Insulæ in Aprutio, caritátis incendio verius quam vi morbi consumptus, divinæque Matris ope recreatus, placidíssimo éxitu ad Superos evolávit, anno millesimo octingentésimo sexagesimo secundo, ætátis suæ vigesimo quarto. Eum deinceps, a Deo miraculis illustrátum, Pius Papa décimus Cælitum beatórum número accensuit. Benedíctus vero décimus quintus, Pontifex maximus, anno millesimo nongentésimo vigesimo, post cónditum Institutum a Passióne ducentésimo, in solemnitate Ascensiónis Dómini, beato júveni Sanctórum honores decrevit; et Pius undecimus ejus Offícium et Missam ad universam Ecclésiam extendit.\par
\switchcolumn
\EN
\noindent Among other virtues, he especially loved Christian humility and obedience; for he considered himself the least of all. He therefore strove eagerly to do all the most menial work of the house, and he most diligently performed, not only the direct commands, but even the unexpressed wishes of his superiors. Curbing his senses, and accustoming himself to a life of austerity, he retained unfaded the flower of his virginity, and completely crucified to the world, he lived to God alone, enjoying an intimate familiarity with his Lord. And so, at Isola in the Abruzzi, filling the short span of his life with so many noble virtues, consumed by the fire of charity rather than by disease, and refreshed by the aid of the Mother of God, his soul flew to heaven in a most peaceful journey in the year 1862, at the age of twenty-four. Then, as he had been made illustrious by God through miracles, Pope Pius X added him to the number of the Blessed in heaven. Likewise, the Supreme Pontiff Benedict XV in 1920, two hundred years after the foundation of the Institute of the Passion, on the feast of the Ascension of the Lord, decreed the honours of the Saints to the blessed youth; and Pius XI extended his Office and Mass to the Universal Church.\par
\switchcolumn*

\LA
\begin{Resp}\Rsign Iste homo perfécit ómnia quæ locútus est ei Deus, et dixit ad eum: Ingrédere in réquiem meam: \Star{} Quia te vidi justum coram me ex ómnibus géntibus.\end{Resp}
\begin{Resp}\Vsign Iste est, qui contémpsit vitam mundi, et pervénit ad cæléstia regna.\end{Resp}
\begin{Resp}\Rsign Quia te vidi justum coram me ex ómnibus géntibus.\end{Resp}
\begin{Resp}\Vsign Glória Patri, et Fílio, \Star{} et Spirítui Sancto.\end{Resp}
\begin{Resp}\Rsign Quia te vidi justum coram me ex ómnibus géntibus.\end{Resp}
\switchcolumn
\EN
\begin{Resp}\Rsign This is he which did according unto all that God commanded him; and God said unto him: Enter thou into My rest, \Star{} For thee have I seen righteous before Me among all people.\end{Resp}
\begin{Resp}\Vsign This is he which loved not his life in this world, and is come unto an everlasting kingdom.\end{Resp}
\begin{Resp}\Rsign For thee have I seen righteous before Me among all people.\end{Resp}
\begin{Resp}\Vsign Glory be to the Father, and to the Son, \Star{} and to the Holy Ghost.\end{Resp}
\begin{Resp}\Rsign For thee have I seen righteous before Me among all people.\end{Resp}
\switchcolumn*

\LA
\LessonHead{Lectio VII}
\switchcolumn
\EN
\LessonHead{Lesson VII}
\switchcolumn*

\LA
\LessonTitle{Léctio sancti Evangélii secúndum Marcum}
\switchcolumn
\EN
\LessonTitle{From the Holy Gospel according to Mark}
\switchcolumn*

\LA
\Cite{Marc 10:13-21}
\switchcolumn
\EN
\Cite{Mark 10:13-21}
\switchcolumn*

\LA
\noindent In illo témpore: Offerébant Jesu párvulos, ut tángeret illos: discípuli autem comminabántur offeréntibus. Et réliqua.\par
\switchcolumn
\EN
\noindent At that time: They brought young children to Jesus, that he should touch them: and his disciples rebuked those that brought them. And so on, and that which followeth.\par
\switchcolumn*

\LA
\LessonTitle{Homilía sancti Bedæ Venerábilis Presbýteri}
\switchcolumn
\EN
\LessonTitle{A Homily by St. Venerable Bede the Priest}
\switchcolumn*

\LA
\Source{Commentárium in Marcum 10:13-21}
\switchcolumn
\EN
\Source{Commentarium in Marcum 10:13-21}
\switchcolumn*

\LA
\noindent Ait discípulis Jesus: Sínite párvulos veníre ad me, et ne prohibuéritis eos; talium enim est regnum Dei. Significanter dixit: Talium est; non: Istórum; ut osténderet, non ætátem regnare, sed mores; et his, qui similem haberent innocéntiam et simplicitátem, præmium repromitti: Apóstolo quoque in eandem senténtiam congruénte: Fratres, nolíte fíeri púeri sensibus; sed malítia párvuli estote, sensu autem ut perfecti sitis. Amen, dico vobis: quisquis non recéperit regnum Dei velut párvulus, non intrábit in illud. Sicut puer non persevérat in iracúndia, non læsus méminit, non videns pulchram mulierem delectátur, non aliud cogitat, aliud lóquitur; sic et vos, nisi talem habuéritis innocéntiam et animi puritátem, regnum cælórum non poteritis intrare. Aliter, regnum Dei, id est doctrinam Evangélii, sicut párvuli accípere jubemur; quia quómodo párvulus in discéndo non contradicit doctóribus neque ratiónes et verba componit advérsum eos resistens, sed fideliter suscipit quod docétur et cum metu obtemperat et quiescit; ita et nos, in obœdiéndo simpliciter et sine ulla rectractatióne verbis Dómini, fácere debemus. Et complexans eos, et impónens manus super illos, benedicebat eos. Complexus benedicit párvulos, ut húmiles spíritu sua benedictióne, grátia et dilectióne dignos esse significet.\par
\switchcolumn
\EN
\noindent Jesus saith unto his disciples: Suffer the little children to come unto me, and forbid them not: for of such is the kingdom of God. It is noteworthy that he saith: Of such; not: Of these. That is, he is concerned with conduct, not with age; for the reward is promised to all such as are like unto children, not in age, but in innocence and simplicity; as saith the Apostle: Be not children in understanding; howbeit, in malice be ye children, but in understanding be men. Verily I say unto you, Whosoever shall not receive the kingdom of God as a little child, he shall not enter therein. A child doth not for long remain angry or remember an injury; he taketh no lustful pleasure in looking at a beautiful woman; he doth not think one thing and say another. In like manner, ye also, unless ye have this innocence and purity of soul, cannot enter the kingdom of heaven. We are commanded to receive the kingdom of God (that is, the teaching of the Gospel) like as children, because they, when learning, do not contradict their teacher, nor adduce reasons and arguments against them, but trustingly accept what they are taught in respectful silence and obedience. And the Lord took them up in his arms, put his hands upon them, and blessed them. He blesseth the children by taking them up in his arms, to shew that such as are poor in spirit do merit his blessing, grace, and love.\par
\switchcolumn*

\LA
\begin{Resp}\Rsign Iste est qui ante Deum magnas virtútes operátus est, et de omni corde suo laudávit Dóminum: \Star{} Ipse intercédat pro peccátis ómnium populórum.\end{Resp}
\begin{Resp}\Vsign Ecce homo sine queréla, verus Dei cultor, ábstinens se ab omni ópere malo, et pérmanens in innocéntia sua.\end{Resp}
\begin{Resp}\Rsign Ipse intercédat pro peccátis ómnium populórum.\end{Resp}
\switchcolumn
\EN
\begin{Resp}\Rsign This is he which wrought great wonders before God, and praised the Lord with all his heart. \Star{} May he pray for all people, that their sins may be forgiven unto them.\end{Resp}
\begin{Resp}\Vsign Behold a man without blame, a worshipper of God in truth, keeping himself clean from every evil work, and abiding still in his innocency.\end{Resp}
\begin{Resp}\Rsign May he pray for all people, that their sins may be forgiven unto them.\end{Resp}
\switchcolumn*

\LA
\LessonHead{Lectio VIII}
\switchcolumn
\EN
\LessonHead{Lesson VIII}
\switchcolumn*

\LA
\noindent Et, cum egréssus esset in viam, procurrens quidam, genu flexo ante eum, rogábat eum: Magister bone, quid fáciam, ut vitam ætérnam percipiam? Audierat, credo, iste quæsítor vitæ ætérnæ a Dómino, tantum eos, qui parvulórum velint esse similes, dignos esse introítu regni cæléstis: atque ideo curam gerens tractátus certioris, poscit sibi non per parábolas, sed aperte, quibus óperum meritis vitam ætérnam cónsequi possit, exponi. Jesus autem dixit ei: Præcepta nosti. Hæc est puerílis innocéntiæ castitas, quæ nobis imitanda propónitur, si regnum Dei volumus intrare. At ille respondens, ait ille: Magister, hæc ómnia observávi a juventúte mea. Non est putandus homo iste vel voto tentántis (ut quidam putavére) Dóminum interrogasse, vel de sua esse vita mentitus, cum se legis mandáta custodisse dicebat; sed, simpliciter, ut víxerit esse conféssus. Quia si mendacii aut simulatiónis noxa reus tenerétur, nequáquam intuitus arcana cordis ejus, eum diligere dicerétur Jesus.\par
\switchcolumn
\EN
\noindent And as he was going forth on his journey, a certain man came running up, and fell on his knees before him, and asked him: Good Master, what shall I do that I may inherit eternal life? Doubtless this seeker after eternal life was one of those who had heard the Lord say that whoso shall not receive the kingdom of God as little child shall not enter therein; and therefore, taking the safer course, he asketh for himself not in parables, but simply and openly, for teaching concerning the meritorious works by which eternal life is to be gained. Whereupon Christ said: Thou knowest the commandments. This is the chastity of childlike innocence, which is proposed for our imitation, if we wish to enter the kingdom of God. But he answering, said unto him: Master, all these have I observed from my youth. This man should not be considered as having asked the Lord with a view to tempting him (as some have thought), or to have lied about his life when he said he had kept the commandments of the law; but that rather he was trying to explain how he had hitherto lived. Because, if he be considered guilty of lies and pretence, it would not have been said of Jesus, who looketh on the secrets of the heart, that he, beholding him, loved him.\par
\switchcolumn*

\LA
\begin{Resp}\Rsign Sint lumbi vestri præcíncti, et lucérnæ ardéntes in mánibus vestris: \Star{} Et vos símiles homínibus exspectántibus dóminum suum, quando revertátur a núptiis.\end{Resp}
\begin{Resp}\Vsign Vigiláte ergo, quia nescítis qua hora Dóminus vester ventúrus sit.\end{Resp}
\begin{Resp}\Rsign Et vos símiles homínibus exspectántibus dóminum suum, quando revertátur a núptiis.\end{Resp}
\begin{Resp}\Vsign Glória Patri, et Fílio, \Star{} et Spirítui Sancto.\end{Resp}
\begin{Resp}\Rsign Et vos símiles homínibus exspectántibus dóminum suum, quando revertátur a núptiis.\end{Resp}
\switchcolumn
\EN
\begin{Resp}\Rsign Let your loins be girded about, and your lights burning; \Star{} And ye yourselves like unto men that wait for their lord, when he will return from the wedding.\end{Resp}
\begin{Resp}\Vsign Watch therefore, for ye know not what hour your Lord doth come.\end{Resp}
\begin{Resp}\Rsign And ye yourselves like unto men that wait for their lord, when he will return from the wedding.\end{Resp}
\begin{Resp}\Vsign Glory be to the Father, and to the Son, \Star{} and to the Holy Ghost.\end{Resp}
\begin{Resp}\Rsign And ye yourselves like unto men that wait for their lord, when he will return from the wedding.\end{Resp}
\switchcolumn*

\LA
\LessonHead{Lectio IX}
\switchcolumn
\EN
\LessonHead{Lesson IX}
\switchcolumn*

\LA
\noindent Díligit enim Dóminus eos, qui mandáta legis, quamvis minora, custódiunt: sed nihilóminus, quod in lege minus fuerat, iis qui perfecti esse desiderant, osténdit, quia non venit solvere legem aut Prophétas, sed adimplere. Ad quam profecto adimpletiónem pertinet, quod hic consequenter adjungitur: Vade, quæcúmque habes vende et da paupéribus, et habebis thesáurum in cælo, et veni, séquere me. Quicúmque perfectus esse volúerit, debet vendere quæ habet; et non ex parte véndere, sicut Ananías et Saphira, sed totum véndere: et cum vendiderit, dare omne paupéribus, et sic sibi præparáre thesáurum in regno cælórum. Nec hoc ad perfectiónem sufficit, nisi, post contemptas divítias, Salvatórem sequátur; id est relictis malis, fáciat bona. Facílius enim sæculum contemnitur quam volúntas. Multi divítias relinquéntes, Dóminum non sequúntur. Sequitur autem Dóminum, qui imitator ejus est et per vestígia illíus graditur. Qui enim dicit se in Christo credere, debet, quómodo ille ambulávit, et ipse ambuláre.\par
\switchcolumn
\EN
\noindent For the love of the Lord goeth out to all those who not only keep the commandments of the law, but have regard even for things which are not of strict precept. Which latter things he who came not to destroy the Law and the Prophets, but to fulfill them, showeth forth unto such as would be perfect, saying: Sell whatsoever thou hast, and give to the poor, and thou shalt have treasure in heaven; and come, follow me. Whosoever would be perfect must sell all (not merely in part, as did Ananias and Sapphira), and thereafter give all to the poor, and thus prepare for himself a treasure in the kingdom of heaven. But contempt for riches is not enough to achieve perfection. He must also follow the Saviour, leaving evil and doing good. For it is easier to spurn the world than the will; and many spurn their riches, but fail thereafter to follow the Lord. Whosoever would follow him must be his imitator, and walk in his footsteps. For whoso saith he believeth in Christ ought also to walk as Christ walked.\par
\switchcolumn*

\LA
\TeDeum{Te Deum}
\switchcolumn
\EN
\TeDeum{Te Deum}
\switchcolumn*

\end{paracol}

\SeasonTitle{Mensis Martius}{March}

\addcontentsline{toc}{section}{Mensis Martius\enspace\textit{March}}

\par\needspace{10\baselineskip}\addvspace{1.2em}{\centering\small\textsc{Die 4 Martii} · \textsc{4 March}\par}
\Office{S. Casimiri Confessoris}{St. Casimir, Confessor}{Semiduplex}
\addcontentsline{toc}{subsection}{Die 4 Martii · S. Casimiri Confessoris\enspace\textit{St. Casimir, Confessor}}
\begin{paracol}{2}
\LA
\RefLine{Lectiones I–III: ex Commúni Confessoris non pontificis (C5), in 2 loco.}
\switchcolumn
\EN
\RefLine{Lessons I–III: from the Common of a Confessor not a Bishop (C5), set 2.}
\switchcolumn*

\LA
\RefLine{Lectiones VII–VIII: ex Commúni Confessoris non pontificis (C5).}
\switchcolumn
\EN
\RefLine{Lessons VII–VIII: from the Common of a Confessor not a Bishop (C5).}
\switchcolumn*

\LA
\RefLine{Lectio IX: de Scriptura occurrente.}
\switchcolumn
\EN
\RefLine{Lesson IX: from the Scripture of the day.}
\switchcolumn*

\LA
\LessonHead{Lectio IV}
\switchcolumn
\EN
\LessonHead{Lesson IV}
\switchcolumn*

\LA
\noindent Casimírus, patre Casimiro, matre Elisabetha Austriaca, Poloniæ regibus ortus, a pueritia sub optimis magistris, pietate et bonis artibus instructus, juveniles artus aspero domabat cilicio, et assiduis extenuabat jejuniis. Regii spreta lecti mollitie, dura cubabat humo, et clam intempesta nocte, præ foribus templorum pronus in terra divinam exorabat clementiam. In Christi contemplanda passione assiduus, Missarum solemniis adeo erecta in Deum mente solebat adesse, ut extra se rapi videretur.\par
\switchcolumn
\EN
\noindent This Casimir was the son of Casimir III, King of Poland, by Elizabeth of Austria, his wife, (and was born upon the 3rd day of October, in the year 1458.) From his childhood he was taught by the best masters, and was trained in all godliness and good learning. While he was still a boy he wore rough haircloth, and chastened himself with much fasting. He forsook the softness of his princely bed, and lay upon the hard ground, and on stormy nights he would go out secretly and prostrate himself before the doors of the churches, crying to God for mercy. He was unwearied in contemplating the Passion of Christ, and when he was present at Mass, so profound was his recollection, that he seemed to be altogether beside himself.\par
\switchcolumn*

\LA
\begin{Resp}\Rsign Honéstum fecit illum Dóminus, et custodívit eum ab inimícis, et a seductóribus tutávit illum: \Star{} Et dedit illi claritátem ætérnam.\end{Resp}
\begin{Resp}\Vsign Justum dedúxit Dóminus per vias rectas, et osténdit illi regnum Dei.\end{Resp}
\begin{Resp}\Rsign Et dedit illi claritátem ætérnam.\end{Resp}
\switchcolumn
\EN
\begin{Resp}\Rsign The Lord made him honourable, and defended him from his enemies, and kept him safe from those that lay in wait for him; \Star{} And gave him perpetual glory.\end{Resp}
\begin{Resp}\Vsign He went down with him into the pit, and left him not in bonds.\end{Resp}
\begin{Resp}\Rsign And gave him perpetual glory.\end{Resp}
\switchcolumn*

\LA
\LessonHead{Lectio V}
\switchcolumn
\EN
\LessonHead{Lesson V}
\switchcolumn*

\LA
\noindent Catholicam promovere fidem summopere studuit, et Ruthenorum schisma abolere: quapropter Casimírum patrem induxit, ut legem ferret, ne schismatici nova templa construerent, nec vetera collabentia restaurarent. Erga pauperes et calamitatibus oppressos beneficus et misericors, patris et defensoris egenorum nomen obtinuit. Virginitatem, quam ab incunabulis servavit illæsam, sub extremo vitæ termino fortiter asseruit, dum gravi pressus infirmitate, mori potius quam castitatis jacturam ex medicorum consilio subire, constanter decrevit.\par
\switchcolumn
\EN
\noindent He made the propagation of the Catholic faith one of the chief works of his life, and strove hard against the schism in Ruthenia. He persuaded his father to forbid by law that the schismatics should build any new churches, or repair the existing ones when they fell into decay. So great was his liberality and tenderness toward the needy and the afflicted, that he came to be called the father and guardian of the poor. From his infancy he never soiled his purity, and in his last illness, when his physicians advised him to seek for relief from his grievous sufferings by the sacrifice of his chastity, he cheerfully determined rather to die.\par
\switchcolumn*

\LA
\begin{Resp}\Rsign Amávit eum Dóminus, et ornávit eum: stolam glóriæ índuit eum, \Star{} Et ad portas paradísi coronávit eum.\end{Resp}
\begin{Resp}\Vsign Índuit eum Dóminus lorícam fídei, et ornávit eum.\end{Resp}
\begin{Resp}\Rsign Et ad portas paradísi coronávit eum.\end{Resp}
\switchcolumn
\EN
\begin{Resp}\Rsign The Lord loved him and beautified him; He clothed him with a robe of glory, \Star{} And crowned him at the gates of Paradise.\end{Resp}
\begin{Resp}\Vsign The Lord hath put on him the breast-plate of faith, and hath adorned him.\end{Resp}
\begin{Resp}\Rsign And crowned him at the gates of Paradise.\end{Resp}
\switchcolumn*

\LA
\LessonHead{Lectio VI}
\switchcolumn
\EN
\LessonHead{Lesson VI}
\switchcolumn*

\LA
\noindent Consummatus in brevi, virtutibus et meritis plenus, prænuntiato mortis die, inter sacerdotum et religiosorum choros' spiritum Deo reddidit, anno ætatis vigesimo quinto. Corpus Vilnam delatum multis claret miraculis. Etenim præterquam quod puella defuncta vitam, cæci visum, claudi gressum, et varii infirmi sanitatem ad ejus sepulchrum recuperarunt; Lithuanis exiguo numero ad potentissimi hostis insperatam irruptionem trepidantibus in aëre apparens, insignem tribuit victoriam. Quibus permotus Leo decimus eumdem Sanctorum catalogo adscripsit.\par
\switchcolumn
\EN
\noindent Being made perfect in a short space, and full of piety and good works, he foretold the day of his own death, and, gathering round him a choir of priests and monks, he rendered his soul into the hands of God Whom they were praising, \textasciitilde{}(upon the 4th day of March, in the year of our Lord 1482, and) the 25th of his own age. His body was carried to Wilna, where many miracles are reputed to have been wrought around it. At his grave a dead girl is said to have received her life again, blind men their sight, cripples the power of walking, and many sick folk health. Moreover, on an occasion when the Lithuanians in scanty numbers were exposed to the shock of a powerful enemy, they believed that he appeared in the air, and gave them the signal victory which they won. On the assurance of these things, Leo X. was moved to add his name to those of the Saints.\par
\switchcolumn*

\LA
\begin{Resp}\Rsign Iste homo perfécit ómnia quæ locútus est ei Deus, et dixit ad eum: Ingrédere in réquiem meam: \Star{} Quia te vidi justum coram me ex ómnibus géntibus.\end{Resp}
\begin{Resp}\Vsign Iste est, qui contémpsit vitam mundi, et pervénit ad cæléstia regna.\end{Resp}
\begin{Resp}\Rsign Quia te vidi justum coram me ex ómnibus géntibus.\end{Resp}
\begin{Resp}\Vsign Glória Patri, et Fílio, \Star{} et Spirítui Sancto.\end{Resp}
\begin{Resp}\Rsign Quia te vidi justum coram me ex ómnibus géntibus.\end{Resp}
\switchcolumn
\EN
\begin{Resp}\Rsign This is he which did according unto all that God commanded him; and God said unto him: Enter thou into My rest, \Star{} For thee have I seen righteous before Me among all people.\end{Resp}
\begin{Resp}\Vsign This is he which loved not his life in this world, and is come unto an everlasting kingdom.\end{Resp}
\begin{Resp}\Rsign For thee have I seen righteous before Me among all people.\end{Resp}
\begin{Resp}\Vsign Glory be to the Father, and to the Son, \Star{} and to the Holy Ghost.\end{Resp}
\begin{Resp}\Rsign For thee have I seen righteous before Me among all people.\end{Resp}
\switchcolumn*

\end{paracol}

\par\needspace{10\baselineskip}\addvspace{1.2em}{\centering\small\textsc{Die 6 Martii} · \textsc{6 March}\par}
\Office{Ss. Perpetuæ et Felicitatis Martyrum}{Sts. Perpetua and Felicity, Martyrs}{Duplex}
\addcontentsline{toc}{subsection}{Die 6 Martii · Ss. Perpetuæ et Felicitatis Martyrum\enspace\textit{Sts. Perpetua and Felicity, Martyrs}}
\begin{paracol}{2}
\LA
\RefLine{Lectiones I–III: ex Commúni plurimarum non Virginum Martyrum (C7b).}
\switchcolumn
\EN
\RefLine{Lessons I–III: from the Common of multiple widowed Martyrs (C7b).}
\switchcolumn*

\LA
\RefLine{Lectiones VII–VIII: ex Commúni plurimarum non Virginum Martyrum (C7b).}
\switchcolumn
\EN
\RefLine{Lessons VII–VIII: from the Common of multiple widowed Martyrs (C7b).}
\switchcolumn*

\LA
\RefLine{Lectio IX: de Scriptura occurrente.}
\switchcolumn
\EN
\RefLine{Lesson IX: from the Scripture of the day.}
\switchcolumn*

\LA
\LessonHead{Lectio IV}
\switchcolumn
\EN
\LessonHead{Lesson IV}
\switchcolumn*

\LA
\noindent Perpetua et Felicitas, in persecutione Severi imperatoris, in Africa, una cum Revocato, Saturnino et Secundulo comprehensæ sunt et in tenebricosum carcerem detrusæ; quibus ultra adjunctus est Satyrus. Erant adhuc catechumenæ, sed paulo post baptizatæ sunt. Paucis diebus interjectis, e carcere ad forum deductæ cum sociis, post gloriosam confessionem, ab Hilarione procuratore damnantur ad bestias. Inde hilares descendunt ad carcerem, ubi variis visionibus recreantur et ad martyrii palmam accenduntur. Perpetuam, nec patris senio confecti iteratæ preces et lacrimæ, nec erga filium infantem pendentem ad ubera maternus amor, nec supplicii atrocitas, a Christi fide dimovere umquam potuerunt\par
\switchcolumn
\EN
\noindent Perpetua and Felicity were arrested during the persecution of the emperor Severus in Africa, together with Revocatus, Saturninus and Secundulus, and were cast into a dark dungeon, where Satyrus was added to their company. They were as yet catechumens but shortly afterwards they were baptized. After a few days they and their companions were led forth from prison to the court, and, after a glorious profession of faith, were condemned by the procurator Hilarion to the beasts. Thereupon they went down to the prison rejoicing, and there they were refreshed with many visions, and fired with longing for the martyr's palm. Neither the repeated prayers and tears of Perpetua's father, a man of extreme old age, nor her motherly love for her baby son still at the breast, nor the atrocity of the torture, could shake her faith in Christ.\par
\switchcolumn*

\LA
\begin{Resp}\Rsign Propter veritátem, et mansuetúdinem, et justítiam: \Star{} Et dedúcet te mirabíliter déxtera tua.\end{Resp}
\begin{Resp}\Vsign Spécie tua et pulchritúdine tua inténde, próspere procéde, et regna.\end{Resp}
\begin{Resp}\Rsign Et dedúcet te mirabíliter déxtera tua.\end{Resp}
\switchcolumn
\EN
\begin{Resp}\Rsign Because of truth, and meekness, and righteousness; \Star{} And thy right hand shall lead thee wonderfully.\end{Resp}
\begin{Resp}\Vsign In thy comeliness, and thy beauty, go forward, fare prosperously, and reign.\end{Resp}
\begin{Resp}\Rsign And thy right hand shall lead thee wonderfully.\end{Resp}
\switchcolumn*

\LA
\LessonHead{Lectio V}
\switchcolumn
\EN
\LessonHead{Lesson V}
\switchcolumn*

\LA
\noindent Felicitas vero instante spectaculo die, cum octo jam menses prægnans esset, in magno erat luctu, ne differretur; leges quippe vetabant prægnantes supplicio affici. At precibus commartyrum accelerato partu, enixa est filiam. Cumque in partu laborans doleret, ait illi quidam de custodibus: Quæ sic modo doles, quid facies objecta bestiis? Cui illa: Modo ego patior; illic autem alius erit in me, qui patietur pro me, quia et ego pro illo passura sum.\par
\switchcolumn
\EN
\noindent As the day of the spectacle came close, Felicity was afflicted with great sorrow lest it should be put off, since she was eight months with child; and the law forbade expectant mothers to be put to the torture. But at the prayers of her fellow-martyrs her delivery was hastened, and she gave birth to a daughter. While she was groaning in the pains of childbirth, one of the jailers said to her: What wilt thou do when thrown to the beasts, if thou groanest thus now? She replied: Now it is I who suffer; but then Another will be within me, who will suffer on my behalf, seeing that it is for him that I am to suffer.\par
\switchcolumn*

\LA
\begin{Resp}\Rsign Dilexísti justítiam, et odísti iniquitátem: \Star{} Proptérea unxit te Deus, Deus tuus, óleo lætítiæ.\end{Resp}
\begin{Resp}\Vsign Propter veritátem, et mansuetúdinem, et justítiam.\end{Resp}
\begin{Resp}\Rsign Proptérea unxit te Deus, Deus tuus, óleo lætítiæ.\end{Resp}
\switchcolumn
\EN
\begin{Resp}\Rsign Thou hast loved righteousness, and hated iniquity; \Star{} Therefore God, thy God, hath anointed thee with the oil of gladness.\end{Resp}
\begin{Resp}\Vsign Because of truth, and meekness, and righteousness.\end{Resp}
\begin{Resp}\Rsign Therefore God, thy God, hath anointed thee with the oil of gladness.\end{Resp}
\switchcolumn*

\LA
\LessonHead{Lectio VI}
\switchcolumn
\EN
\LessonHead{Lesson VI}
\switchcolumn*

\LA
\noindent In amphitheatrum, toto inspectante populo, producuntur tandem generosæ mulieres, Nonis Martii, ac primum flagellis cæduntur. Tunc a ferocissima vacca aliquamdiu jactatæ, plagis concisæ et in terram elisæ sunt. Demum cum sociis, qui a variis bestiis vexati fuerant, gladiorum ictibus conficiuntur. Harum sanctarum Mártyrum festum Pius decimus Pontifex maximus ad ritum duplicem pro universa Ecclesia evexit ac diei sextæ Martii adsignari mandavit.\par
\switchcolumn
\EN
\noindent At length the noble-hearted women were brought into the amphitheatre, in the sight of all the people, on the 7th day of March, and were first beaten with scourges. Then they were tossed for some time by a ferocious cow, beaten with lashes, and dashed on the ground. Lastly, together with their companions, who had been attacked by various beasts, they were slain with blows of the sword. Pope Pius X raised the feast of these holy Martyrs to the rite of a double for the Universal Church, and ordered it to be kept on March 6th.\par
\switchcolumn*

\LA
\begin{Resp}\Rsign Fallax grátia, et vana est pulchritúdo: \Star{} Múlier timens Dóminum, ipsa laudábitur.\end{Resp}
\begin{Resp}\Vsign Date ei de fructu mánuum suárum, et laudent eam in portis ópera ejus.\end{Resp}
\begin{Resp}\Rsign Múlier timens Dóminum, ipsa laudábitur.\end{Resp}
\begin{Resp}\Vsign Glória Patri, et Fílio, \Star{} et Spirítui Sancto.\end{Resp}
\begin{Resp}\Rsign Múlier timens Dóminum, ipsa laudábitur.\end{Resp}
\switchcolumn
\EN
\begin{Resp}\Rsign Favour is deceitful, and beauty is vain \Star{} A woman that feareth God she shall be praised.\end{Resp}
\begin{Resp}\Vsign Give her of the fruit of her hands, and let her own works praise her in the gates.\end{Resp}
\begin{Resp}\Rsign A woman that feareth God she shall be praised.\end{Resp}
\begin{Resp}\Vsign Glory be to the Father, and to the Son, \Star{} and to the Holy Ghost.\end{Resp}
\begin{Resp}\Rsign A woman that feareth God she shall be praised.\end{Resp}
\switchcolumn*

\end{paracol}

\par\needspace{10\baselineskip}\addvspace{1.2em}{\centering\small\textsc{Die 7 Martii} · \textsc{7 March}\par}
\Office{S. Thomæ de Aquino Confessoris et Ecclesiæ Doctoris}{St. Thomas Aquinas, Confessor and Doctor of the Church}{Duplex}
\addcontentsline{toc}{subsection}{Die 7 Martii · S. Thomæ de Aquino Confessoris et Ecclesiæ Doctoris\enspace\textit{St. Thomas Aquinas, Confessor and Doctor of the Church}}
\begin{paracol}{2}
\LA
\RefLine{Lectiones I–III: ex Commúni Doctoris non Pontificis (C5a).}
\switchcolumn
\EN
\RefLine{Lessons I–III: from the Common of a Doctor not a Bishop (C5a).}
\switchcolumn*

\LA
\RefLine{Lectiones VII–VIII: ex Commúni Doctoris non Pontificis (C5a).}
\switchcolumn
\EN
\RefLine{Lessons VII–VIII: from the Common of a Doctor not a Bishop (C5a).}
\switchcolumn*

\LA
\RefLine{Lectio IX: de Scriptura occurrente.}
\switchcolumn
\EN
\RefLine{Lesson IX: from the Scripture of the day.}
\switchcolumn*

\LA
\LessonHead{Lectio IV}
\switchcolumn
\EN
\LessonHead{Lesson IV}
\switchcolumn*

\LA
\noindent Præclárum christiáni orbis decus et Ecclésiæ lumen, beatíssimus vir Thomas, Landulpho comite Aquinate et Theodora Neapolitana, nobílibus paréntibus natus, futuræ in Deiparam devotiónis afféctum adhuc infantulus osténdit. Nam chártulam ab eo invéntam, in qua salutátio Angelica scripta erat, frustra adniténte nutrice, compressa manu valide retinuit, et a matre per vim abreptam, ploratu et gestu repétiit, ac mox redditam deglutívit. Quintum annum agens, monachis sancti Benedicti Cassinátibus custodiendus traditur. Inde Neapolim studiórum causa missus, jam adoléscens Fratrum Prædicatórum Ordinem suscépit. Sed matre ac frátribus id indigne feréntibus, Lutétiam Parisiórum mittitur. Quem fratres in itinere per vim raptum, in arcem castri sancti Joánnis perducunt: ubi varie exagitatus, ut sanctum proponsitum mutaret, mulierem etiam, quæ ad labefactándam ejus constantiam introducta fuerat, titióne fugávit. Mox beátus juvenis, flexis genibus ante signum crucis orans, ibique somno correptus, per quietem sentire visus est sibi ab Angelis constringi lumbos; quo ex témpore omni póstea libidinis sensu cáruit. Soróribus, quæ ut eum a pio consílio removerent, in castrum venerant, persuasit, ut, contemptis curis sæcularibus, ad exercitatiónem cæléstis vitæ se conferrent.\par
\switchcolumn
\EN
\noindent That splendid adornment of the Christian world and light of the Church, blessed Thomas of Aquino, was the son of Landulph, Earl of Aquino, and Theodora of Naples, his wife, being nobly descended on both sides. (He was born in the year of salvation 1226,) and even as an infant gave token of the love which he afterwards bore to the Mother of God. He found a little bit of paper upon which was written the Angelic Salutation, and held it firm in his hand in spite of the efforts of his wet-nurse; his mother took it away by force, but he cried and stretched out for it, and when she gave it back to him, he swallowed it. When he was only four years old, he was given into the keeping of the Benedictine monks of Monte Cassino. He was thence sent to Naples to study, and there, while very young, entered the Order of Friars Preachers. This displeased his mother and brothers, and he left Naples for Paris. When he was on his journey his brothers met him, and carried him off by force to the castle of Monte San Giovanni, where they imprisoned him in the keep. Here they used every means to break him of his intention, and at last brought a woman into his room to try to overcome his purity. The lad drove her out with a fire-brand. When he was alone he knelt down before the figure of the Cross, and there he fell asleep. As he slept, it seemed to him that angels came and girded his loins and from this time he never felt the least sexual inclination. His sisters came to the castle to beseech him to give up his purpose of leaving the world, but he so worked on them by his godly exhortations, that both of them ever after set no value on earthly things, and busied themselves rather with heavenly.\par
\switchcolumn*

\LA
\begin{Resp}\Rsign Honéstum fecit illum Dóminus, et custodívit eum ab inimícis, et a seductóribus tutávit illum: \Star{} Et dedit illi claritátem ætérnam.\end{Resp}
\begin{Resp}\Vsign Justum dedúxit Dóminus per vias rectas, et osténdit illi regnum Dei.\end{Resp}
\begin{Resp}\Rsign Et dedit illi claritátem ætérnam.\end{Resp}
\switchcolumn
\EN
\begin{Resp}\Rsign The Lord made him honourable, and defended him from his enemies, and kept him safe from those that lay in wait for him; \Star{} And gave him perpetual glory.\end{Resp}
\begin{Resp}\Vsign He went down with him into the pit, and left him not in bonds.\end{Resp}
\begin{Resp}\Rsign And gave him perpetual glory.\end{Resp}
\switchcolumn*

\LA
\LessonHead{Lectio V}
\switchcolumn
\EN
\LessonHead{Lesson V}
\switchcolumn*

\LA
\noindent Emissus e castro per fenestram, Neapolim redúcitur; unde Romam, póstea Parisium a fratre Joánne Theutónico, ordinis Prædicatórum generali magistro, ductus, Alberto Magno doctore, philosophíæ ac theologíæ operam dedit. Viginti quinque annos natus, magister est appellatus, publiceque philósophos ac Theólogos summa cum laude est interpretatus. Numquam se lectióni aut scriptióni dedit, nisi post oratiónem. In difficultátibus locórum sacræ Scripturæ, ad oratiónem jejunium adhibebat. Quin étiam sodali suo fratri Regínaldo dicere solebat, quidquid sciret non tam studio aut labóre suo peperisse, quam divínitus tráditum accepisse. Neapoli cum ad imáginem Crucifixi vehementius oraret, hanc vocem audívit: Bene scripsísti de me, Thoma; quam ergo mercedem accípies? Cui ille: Non aliam, Dómine, nisi teípsum. Collatiónes Patrum assidue pervolutábat; et nullum fuit scriptórum genus, in quo non esset diligentíssime versatus. Scripta ejus et multitúdine, et varietáte, et facilitate explicándi res difficiles adeo excellunt, ut uberrima atque incorrupta illíus doctrina, cum revelátis veritátibus mire consentiens, aptíssima sit ad ómnium témporum erróres pervincendos.\par
\switchcolumn
\EN
\noindent Being let down from a window, Thomas escaped out of the castle of Monte San Giovanni, and returned to Naples. Thence he went first to Rome, and then to Paris, in company of Brother John the German, then Master-General of the Friars Preachers. At Paris he studied Philosophy and Theology under Albert the Great Doctor. At the age of twenty-five years he took the degree of Master, and gave public disquisitions on the Philosophers and Theologians with great distinction. He never set himself to read or write till he had first prayed, and when he was about to take in hand a hard passage of the Holy Scriptures, he fasted also. Hence he was wont to say to Brother Reginald his comrade, that whatever he knew, he had learnt, not so much from his own labour and study, as from the inspiration of God. At Naples he was once kneeling in very earnest prayer before an image of Christ Crucified, when he heard a voice which said Thomas, thou hast written well of Me what reward wilt thou that I give thee? He answered: Lord, thyself. He studied most carefully the works of the Fathers, and there was no kind of author in which he was not well read. His own writings are so wonderful, both because of their number, their variety, and the clearness of his explanations of hard things, that his rich and pure teaching, marvellously consonant with revealed truth, is an admirable antidote for the errors of all times.\par
\switchcolumn*

\LA
\begin{Resp}\Rsign Amávit eum Dóminus, et ornávit eum: stolam glóriæ índuit eum, \Star{} Et ad portas paradísi coronávit eum.\end{Resp}
\begin{Resp}\Vsign Índuit eum Dóminus lorícam fídei, et ornávit eum.\end{Resp}
\begin{Resp}\Rsign Et ad portas paradísi coronávit eum.\end{Resp}
\switchcolumn
\EN
\begin{Resp}\Rsign The Lord loved him and beautified him; He clothed him with a robe of glory, \Star{} And crowned him at the gates of Paradise.\end{Resp}
\begin{Resp}\Vsign The Lord hath put on him the breast-plate of faith, and hath adorned him.\end{Resp}
\begin{Resp}\Rsign And crowned him at the gates of Paradise.\end{Resp}
\switchcolumn*

\LA
\LessonHead{Lectio VI}
\switchcolumn
\EN
\LessonHead{Lesson VI}
\switchcolumn*

\LA
\noindent A summo Pontifice Urbano quarto Romam vocatus, ejus jussu ecclesiásticum lucubrávit Offícium in Corporis Christi solemnitate celebrándum: oblatos vero honores et Neapolitanum archiepiscopátum, étiam deferénte Cleménte quarto, recusávit. A prædicatióne divini verbi non desistebat; quod cum fáceret per octavam Paschæ in basilica sancti Petri, mulierem, quæ ejus fimbriam tetígerat, a fluxu sánguinis liberávit. Missus a beato Gregório décimo ad concílium Lugdunense, in monasterio Fossæ Novæ in morbum íncidit, ubi ægrotus Cantica canticórum explanávit. Ibidem obiit quinquagenarius, anno salútis millesimo ducentésimo septuagesimo quarto, Nonis Martii. Miraculis étiam mórtuus cláruit; quibus probátis, a Joánne vigesimo secundo in Sanctórum númerum relátus est, anno millesimo tercentésimo vigesimo tertio, translato póstea ejus corpore Tolosam, ex mandato beáti Urbani quinti. Cum sanctis angelicis spirítibus non minus innocéntia quam ingenio comparatus, Doctoris Angelici nomen jure est adeptus, eidem auctoritate sancti Pii quinti confirmátum. Leo autem décimus tertius, libentíssime excípiens postulatiónes et vota ómnium pene Sacrórum antistitum orbis catholici, ad tot præcipue philosophicórum systématum a veritáte aberrántium luem propulsandam, ad increménta scientiárum et communem humani generis utilitátem, eum, ex sacrórum Rituum Congregatiónis consulto, per apostolicas litteras cælestem patronum scholárum ómnium catholicárum declarávit et instituit.\par
\switchcolumn
\EN
\noindent The Supreme Pontiff Urban IV. sent for him to Rome, and at his command he composed the Church Office for the feast of Corpus Christi. The Pope could not persuade him to accept any dignity. Pope Clement IV. also offered him the Archbishopric of Naples, but he refused it. He did not neglect the preaching of the Word of God. Once while he was giving a course of sermons in the Basilica of St Peter, during the octave of Easter, a woman who had an issue of blood was healed by touching the hem of his garment. He was sent by blessed Gregory X. to the Council of Lyons, but fell sick on his way to the Abbey of Fossa Nuovo, and there during his illness he made an exposition of the Song of Songs. There he died on the 7th day of March, in the year of salvation 1274, aged fifty years. He was distinguished for miracles even after his death, and on proof of these Pope John XXII. added his name to those of the Saints in the year 1323. His body was afterwards carried to Toulouse by command of blessed Urban V. He has been compared to an angel, both on account of his innocency and of his intellectual power, and has hence been deservedly termed the Angelic Doctor. The use of which title as applied to him was approved by the authority of holy Pius V. Leo XIII. cheerfully agreeing to the prayers and wishes of nearly all the bishops of the Catholic world, and in conformity with a vote of the Congregation of Sacred Rites, by his Apostolic letters declared and recognized Thomas of Aquino as the patron in heaven of all Catholic schools, as an antidote to the plague of so many false systems, especially of philosophy, for the increase of scientific knowledge, and for the common good of all mankind.\par
\switchcolumn*

\LA
\begin{Resp}\Rsign Iste homo perfécit ómnia quæ locútus est ei Deus, et dixit ad eum: Ingrédere in réquiem meam: \Star{} Quia te vidi justum coram me ex ómnibus géntibus.\end{Resp}
\begin{Resp}\Vsign Iste est, qui contémpsit vitam mundi, et pervénit ad cæléstia regna.\end{Resp}
\begin{Resp}\Rsign Quia te vidi justum coram me ex ómnibus géntibus.\end{Resp}
\begin{Resp}\Vsign Glória Patri, et Fílio, \Star{} et Spirítui Sancto.\end{Resp}
\begin{Resp}\Rsign Quia te vidi justum coram me ex ómnibus géntibus.\end{Resp}
\switchcolumn
\EN
\begin{Resp}\Rsign This is he which did according unto all that God commanded him; and God said unto him: Enter thou into My rest, \Star{} For thee have I seen righteous before Me among all people.\end{Resp}
\begin{Resp}\Vsign This is he which loved not his life in this world, and is come unto an everlasting kingdom.\end{Resp}
\begin{Resp}\Rsign For thee have I seen righteous before Me among all people.\end{Resp}
\begin{Resp}\Vsign Glory be to the Father, and to the Son, \Star{} and to the Holy Ghost.\end{Resp}
\begin{Resp}\Rsign For thee have I seen righteous before Me among all people.\end{Resp}
\switchcolumn*

\end{paracol}

\par\needspace{10\baselineskip}\addvspace{1.2em}{\centering\small\textsc{Die 8 Martii} · \textsc{8 March}\par}
\Office{S. Joannis de Deo Confessoris}{St. John of God, Confessor}{Duplex}
\addcontentsline{toc}{subsection}{Die 8 Martii · S. Joannis de Deo Confessoris\enspace\textit{St. John of God, Confessor}}
\begin{paracol}{2}
\LA
\RefLine{Lectiones I–III: ex Commúni Confessoris non pontificis (C5).}
\switchcolumn
\EN
\RefLine{Lessons I–III: from the Common of a Confessor not a Bishop (C5).}
\switchcolumn*

\LA
\RefLine{Lectio IX: de Scriptura occurrente.}
\switchcolumn
\EN
\RefLine{Lesson IX: from the Scripture of the day.}
\switchcolumn*

\LA
\LessonHead{Lectio IV}
\switchcolumn
\EN
\LessonHead{Lesson IV}
\switchcolumn*

\LA
\noindent Joánnes de Deo, ex catholicis piisque paréntibus in oppido Montis Majoris Junióris, regni Lusitaniæ, natus, quam sublimiter in sortem Dómini fúerit eléctus, insuetus splendor super ejus domo refulgens, sonitúsque æris campani sua sponte emissus, ab ipso nativitátis témpore non obscure prænuntiárunt. A laxiori vivéndi ratióne, divina operante virtúte, revocatus, magnæ sanctitátis exhibere specimen cœpit, et ob audítam prædicatiónem verbi Dei sic ad melióra se excitátum sensit, ut jam ab ipso sanctioris vitæ rudiménto consummátum aliquid perfectumque visus sit attigisse. Bonis ómnibus in páuperes carcéribus inclusos erogatis, admirábilis pœniténtiæ suique ipsíus contemptus cuncto pópulo spectaculum factus, a plerisque, ceu demens, gráviter afflíctus, in carcerem améntibus destinátum conjícitur. At Joánnes cælésti caritate magis incénsus, gemino atque amplo valetudinario ex piórum eleemosynis in civitáte Granaténsi exstructo, jactoque novi ordinis fundaménto, Ecclésiam nova prole fœcundávit fratrum Hospitalitátis, infirmis præcláro animárum corporúmque profectu inserviéntium, et longe lateque per orbem diffusórum.\par
\switchcolumn
\EN
\noindent John of God was born of Catholic and godly parents in the town of Montemor in Portugal, (in the year 1495.) The lot to which God had elected him was foreshown at his birth by a light shining over the house, and by the ringing of a bell untouched by human hands. He fell at one time into a loose habit of life, but was recalled by the grace of God, and began to show tokens of true reformation. By hearing the Word of God, he so felt himself stirred up to strive after nobler things, that he considered not that to which he had already attained, and yearned to be perfect, as our Father in heaven is perfect. He gave away all his property to the poor and prisoners, and became a gazing-stock to all that knew him, by the strength of his repentance, and the depth of his selfcontempt. On this account he was commonly supposed to be mad, and was once shut up in a lunatic asylum. He was only the more filled with schemes of charity, and collected, by begging, funds sufficient to build a large double Hospital in the town ot Granada. Here he founded the new Order of Hospital Brethren with which he enriched the Church. These Brethren are now spread throughout all parts of the world, and engaged in ministering to the souls and bodies of the sick.\par
\switchcolumn*

\LA
\begin{Resp}\Rsign Honéstum fecit illum Dóminus, et custodívit eum ab inimícis, et a seductóribus tutávit illum: \Star{} Et dedit illi claritátem ætérnam.\end{Resp}
\begin{Resp}\Vsign Justum dedúxit Dóminus per vias rectas, et osténdit illi regnum Dei.\end{Resp}
\begin{Resp}\Rsign Et dedit illi claritátem ætérnam.\end{Resp}
\switchcolumn
\EN
\begin{Resp}\Rsign The Lord made him honourable, and defended him from his enemies, and kept him safe from those that lay in wait for him; \Star{} And gave him perpetual glory.\end{Resp}
\begin{Resp}\Vsign He went down with him into the pit, and left him not in bonds.\end{Resp}
\begin{Resp}\Rsign And gave him perpetual glory.\end{Resp}
\switchcolumn*

\LA
\LessonHead{Lectio V}
\switchcolumn
\EN
\LessonHead{Lesson V}
\switchcolumn*

\LA
\noindent Paupéribus ægrotis, quos propriis quandoque humeris domum deferebat, nulla re ad ánimæ corporisque salútem profícua déerit. Effusa quoque extra nosocomíum caritate indigéntibus muliéribus viduis, et præcipue virgínibus periclitántibus clam aliménta subministrabat, curamque indefessam adhibebat, ut carnis concupiscéntiam a próximis hujusmodi vítio inquinátis exterminaret. Cum autem maximum in regio Granaténsi valetudinario excitátum fuísset incendium, Joánnes impávidus prosiliit in ignem, huc illuc discurrens, quoúsque tum infirmos humeris exportatos, tum lectulos e fenestris projectos ab igne vindicávit, ac per dimidiam horam inter flammas, jam in immensum succrescéntes, versatus, exínde divínitus incólumis, univérsis civibus admirántibus exívit; in schola caritátis édocens, segniórem in eum fuísse ignem qui foris ússerat, quam qui intus accenderat.\par
\switchcolumn
\EN
\noindent He strove to get for the sick poor, whom he sometimes brought to the Hospital on his own shoulders, whatever was needful for their souls or bodies. His charity was extended to the poor outside of his institution, and he used to supply food privately to necessitous widows, and more so to young women whose virtue was tempted on account of their poverty. He was most careful in encouraging the virtue of purity in all whom he knew. On one occasion when there was a great fire in the hospital at Granada, John bravely entered the burning house, ran from one part of it to another, carried out the sick on his shoulders, and threw the beds out of the windows, and finally, after passing half-an-hour in the midst of the flames, which were now raging with great violence, by the mercy of God left the building uninjured, to the great wonder of all the citizens; thereby to teach all them that love God that the fire which burnt in his heart gave him strength to risk the fire which threatened him from without.\par
\switchcolumn*

\LA
\begin{Resp}\Rsign Amávit eum Dóminus, et ornávit eum: stolam glóriæ índuit eum, \Star{} Et ad portas paradísi coronávit eum.\end{Resp}
\begin{Resp}\Vsign Índuit eum Dóminus lorícam fídei, et ornávit eum.\end{Resp}
\begin{Resp}\Rsign Et ad portas paradísi coronávit eum.\end{Resp}
\switchcolumn
\EN
\begin{Resp}\Rsign The Lord loved him and beautified him; He clothed him with a robe of glory, \Star{} And crowned him at the gates of Paradise.\end{Resp}
\begin{Resp}\Vsign The Lord hath put on him the breast-plate of faith, and hath adorned him.\end{Resp}
\begin{Resp}\Rsign And crowned him at the gates of Paradise.\end{Resp}
\switchcolumn*

\LA
\LessonHead{Lectio VI}
\switchcolumn
\EN
\LessonHead{Lesson VI}
\switchcolumn*

\LA
\noindent Multiplici asperitátum genere, demissíssima obediéntia, extrema paupertáte, orándi studio, rerum divinárum contemplatióne, ac in beátam Vírginem pietáte mirifice excelluit, et lacrimárum dono enituit. Denique gravi morbo correptus, ómnibus Ecclésiæ sacramentis rite sancteque refectus, viribus licet destitútus, propriis indútus vestibus, e léctulo surgens ac provolutus in genua, manu et corde Christum Dóminum e cruce pendentem perstringens, octavo Idus Martii anno millesimo quingentésimo quinquagesimo obiit in osculo Dómini; quem étiam mórtuus ténuit, nec dimísit, et in eádem corporis constitutióne sex circiter horas, quoúsque inde dimotus fuísset, tota civitáte inspectante, mirabíliter permansit, odórem mire fragrántem diffúndens. Quem ante et post óbitum plurimis miraculis clarum, Alexander octavus Pontifex maximus in Sanctórum númerum rétulit; et Leo décimus tertius, ex Sacrórum catholici orbis antístitum voto ac Rituum Congregatiónis consulto, cælestem ómnium hospitalium et infirmórum ubíque degéntium patronum declarávit, ipsiúsque nomen in agonizántium litaníis invocari præcepit.\par
\switchcolumn
\EN
\noindent He was a marked example of every kind of austerity, of the most lowly obedience, of the deepest voluntary poverty, of the most constant prayer, of ghostly contemplation, and of love towards the Blessed Virgin. He was distinguished for the gift of tears. Being at last seized by deadly sickness, he duly received, with saintly affection, all the Sacraments of the Church. After all strength seemed to have left him, he got out of his bed, put on his own clothes, and knelt down before an image of the Lord Christ hanging on the Cross. Round it he threw his arms and pressed it against his heart, and in this position, as it were in the kiss of the Lord, he died, on the 8th day of March 1550. After his death his body did not leave its grip of the crucifix until it was forcibly taken away, six hours after. During these six hours all the inhabitants of the city came to see it, and noticed a savour of strange sweetness proceeding from it. His name was illustrious as a worker of miracles both before and after his death, and the Supreme Pontiff Alexander VIII added it to those of the Saints, and Leo XIII, at the desire of the Bishops of the Catholic world, and in accordance with a vote of the Congregation of Rites, declared him the patron in heaven of all the sick and those who nurse them, wheresoever dwelling, and ordered that his name should be called upon in the Litany for the dying.\par
\switchcolumn*

\LA
\begin{Resp}\Rsign Iste homo perfécit ómnia quæ locútus est ei Deus, et dixit ad eum: Ingrédere in réquiem meam: \Star{} Quia te vidi justum coram me ex ómnibus géntibus.\end{Resp}
\begin{Resp}\Vsign Iste est, qui contémpsit vitam mundi, et pervénit ad cæléstia regna.\end{Resp}
\begin{Resp}\Rsign Quia te vidi justum coram me ex ómnibus géntibus.\end{Resp}
\begin{Resp}\Vsign Glória Patri, et Fílio, \Star{} et Spirítui Sancto.\end{Resp}
\begin{Resp}\Rsign Quia te vidi justum coram me ex ómnibus géntibus.\end{Resp}
\switchcolumn
\EN
\begin{Resp}\Rsign This is he which did according unto all that God commanded him; and God said unto him: Enter thou into My rest, \Star{} For thee have I seen righteous before Me among all people.\end{Resp}
\begin{Resp}\Vsign This is he which loved not his life in this world, and is come unto an everlasting kingdom.\end{Resp}
\begin{Resp}\Rsign For thee have I seen righteous before Me among all people.\end{Resp}
\begin{Resp}\Vsign Glory be to the Father, and to the Son, \Star{} and to the Holy Ghost.\end{Resp}
\begin{Resp}\Rsign For thee have I seen righteous before Me among all people.\end{Resp}
\switchcolumn*

\LA
\LessonHead{Lectio VII}
\switchcolumn
\EN
\LessonHead{Lesson VII}
\switchcolumn*

\LA
\LessonTitle{Léctio sancti Evangélii secúndum Matthǽum}
\switchcolumn
\EN
\LessonTitle{From the Holy Gospel according to Matthew}
\switchcolumn*

\LA
\Cite{Matt 22:34-46}
\switchcolumn
\EN
\Cite{Matt 22:34-46}
\switchcolumn*

\LA
\noindent In illo témpore: Accessérunt ad Jesum pharisæi, et interrogávit eum unus ex eis legis doctor tentans eum: Magister, quod est mandátum magnum in lege? Et réliqua.\par
\switchcolumn
\EN
\noindent At that time The Pharisees came together: And one of them, a doctor of the law, asking him, tempting him: Master, which is the greatest commandment in the law? And so on.\par
\switchcolumn*

\LA
\LessonTitle{Homilía sancti Joánnis Chrysóstomi}
\switchcolumn
\EN
\LessonTitle{Homily by St. John Chrysostom, Patriarch (of Constantinople.)}
\switchcolumn*

\LA
\Source{Homilia 72 in Matthæum}
\switchcolumn
\EN
\Source{72nd on Matthew.}
\switchcolumn*

\LA
\noindent Sadducæis confúsis, pharisæi rursus aggrediúntur; cumque quiéscere oporteret, decertare voluérunt: et legis perítiam profitentem præmittunt, non discere, sed tentare cupiéntes; ac ita intérrogant: Quodnam primum mandátum in lege sit. Nam cum primum illud sit, Diliges Dóminum Deum tuum: putántes causas sibi allatúrum ad mandátum hoc corrigéndum, aliquid addendo, quóniam Deum se faciebat, hoc modo intérrogant. Quid ígitur Christus? Ut ostendat idcirco ad hæc eos devenisse, quia nulla in eis esset caritas, sed invidiæ livore tabéscerent: Diliges, inquit, Dóminum Deum tuum: hoc primum et magnum mandátum est. Secúndum autem simile huic: Diliges próximum tuum sicut teípsum.\par
\switchcolumn
\EN
\noindent When the Pharisees had heard that Christ had put the Sadducees to silence, they gathered themselves together for a fresh attack; just when it behoved them to be quiet, they willed to contend; and so they put forward one of themselves, who professed skill in the law, not wishing to learn, but to lay a snare. This person therefore proposed the question Which is the great commandment in the law? The first and great commandment is Thou shalt love the Lord thy God, but they expected that He would make some exception or addition to this in His Own case, since He made Himself God. (John x. 33.) With this expectation they asked Him the question: But what said Christ? To show that they had adopted this course, because they were loveless, and sick with envy, He answered Thou shalt love the Lord thy God with all thy heart, and with all thy soul, and with all thy mind. This is the first and great commandment. And the second is like unto it Thou shalt love thy neighbour as thyself.\par
\switchcolumn*

\LA
\begin{Resp}\Rsign Iste est qui ante Deum magnas virtútes operátus est, et de omni corde suo laudávit Dóminum: \Star{} Ipse intercédat pro peccátis ómnium populórum.\end{Resp}
\begin{Resp}\Vsign Ecce homo sine queréla, verus Dei cultor, ábstinens se ab omni ópere malo, et pérmanens in innocéntia sua.\end{Resp}
\begin{Resp}\Rsign Ipse intercédat pro peccátis ómnium populórum.\end{Resp}
\switchcolumn
\EN
\begin{Resp}\Rsign This is he which wrought great wonders before God, and praised the Lord with all his heart. \Star{} May he pray for all people, that their sins may be forgiven unto them.\end{Resp}
\begin{Resp}\Vsign Behold a man without blame, a worshipper of God in truth, keeping himself clean from every evil work, and abiding still in his innocency.\end{Resp}
\begin{Resp}\Rsign May he pray for all people, that their sins may be forgiven unto them.\end{Resp}
\switchcolumn*

\LA
\LessonHead{Lectio VIII}
\switchcolumn
\EN
\LessonHead{Lesson VIII}
\switchcolumn*

\LA
\noindent Quam ob rem simile est huic? Quóniam hoc illud inducit, et ab illo rursus munítur. Quicúmque enim male agit, ódio habet lucem, et non venit ad lucem. Et rursus: Dixit insípiens in corde suo, Non est Deus. Deínde sequitur: Corrúpti sunt, et abominábiles facti sunt in stúdiis suis. Et íterum: Radix ómnium malórum avarítia est; quam quidam appeténtes, erravérunt a fide. Et, Qui díligit me, mandáta mea servábit: quorum caput et radix est: Diliges Dóminum Deum tuum, et próximum tuum sicut teípsum.\par
\switchcolumn
\EN
\noindent Mighty is this second commandment like unto the first? Because the first is the second's source and sanction. For every one that doeth evil hateth the light, neither cometh to the light. (John. iii. 20.) And again The fool hath said in his heart There is no God and there folio weth They are corrupt, and become abominable in their works. (Ps. xiii. I.) And yet again: The love of money is the root of all evil; which while some coveted after, they have erred from the faith. (i Tim. vi. 10.) And yet once more: If ye love Me, keep My commandments. (John xiv. 15,) of which commandments the head and root is Thou shalt love the Lord thy God; and thy neighbour as thyself.\par
\switchcolumn*

\LA
\begin{Resp}\Rsign Sint lumbi vestri præcíncti, et lucérnæ ardéntes in mánibus vestris: \Star{} Et vos símiles homínibus exspectántibus dóminum suum, quando revertátur a núptiis.\end{Resp}
\begin{Resp}\Vsign Vigiláte ergo, quia nescítis qua hora Dóminus vester ventúrus sit.\end{Resp}
\begin{Resp}\Rsign Et vos símiles homínibus exspectántibus dóminum suum, quando revertátur a núptiis.\end{Resp}
\begin{Resp}\Vsign Glória Patri, et Fílio, \Star{} et Spirítui Sancto.\end{Resp}
\begin{Resp}\Rsign Et vos símiles homínibus exspectántibus dóminum suum, quando revertátur a núptiis.\end{Resp}
\switchcolumn
\EN
\begin{Resp}\Rsign Let your loins be girded about, and your lights burning; \Star{} And ye yourselves like unto men that wait for their lord, when he will return from the wedding.\end{Resp}
\begin{Resp}\Vsign Watch therefore, for ye know not what hour your Lord doth come.\end{Resp}
\begin{Resp}\Rsign And ye yourselves like unto men that wait for their lord, when he will return from the wedding.\end{Resp}
\begin{Resp}\Vsign Glory be to the Father, and to the Son, \Star{} and to the Holy Ghost.\end{Resp}
\begin{Resp}\Rsign And ye yourselves like unto men that wait for their lord, when he will return from the wedding.\end{Resp}
\switchcolumn*

\LA
\LessonHead{Lectio IX}
\switchcolumn
\EN
\LessonHead{Lesson IX}
\switchcolumn*

\LA
\noindent Si ergo diligere Deum, diligere proximum est (nam si diligis me, o Petre, inquit, pasce oves meas): si étiam dilectio proximi facit ut mandata custodias: merito ait in his totam legem et prophetas pendere. Et quemadmodum in superioribus, cum de resurrectione interrogaretur, plus docuit quam tentantes petebant: sic in hoc loco de primo interrogatus mandato, secundum étiam non valde quam primum inferius, sponte attulit: secundum enim est primo simile. Ita occulte insinuavit, odio illos ad quærendum incitari. Caritas enim, inquit, non gemulatur.\par
\switchcolumn
\EN
\noindent If therefore, to love God is to love our neighbour also, as it appeareth where it is written Simon, son of Jonas, lovest thou Me? And he said unto Him Lord, Thou knowest all things; Thou knowest that I love thee. Jesus saith unto him Feed My sheep, (John xxi. 17,) and if love is the fulfilling of the law, \textasciitilde{}(Rom. xiii. 10,) justly doth the Lord say that on these two commandments hang all the law and the Prophets. And even as when, before this, (23-32,) being interrogated about the Resurrection, He answered them more than they asked, so, now, being interrogated concerning the first and great commandment, He answereth them, of His own accord, touching that second one also, which is little lower than the first, for the second is like unto it. Herein He would have them understand that it was hatred stirred them up to question Him. For Charity, saith the Apostle, envieth not. (1 Cor. xiii. 4.)\par
\switchcolumn*

\LA
\TeDeum{Te Deum}
\switchcolumn
\EN
\TeDeum{Te Deum}
\switchcolumn*

\end{paracol}

\par\needspace{10\baselineskip}\addvspace{1.2em}{\centering\small\textsc{Die 9 Martii} · \textsc{9 March}\par}
\Office{S. Franciscæ Romanæ Viduæ}{St. Frances of Rome, Widow}{Duplex}
\addcontentsline{toc}{subsection}{Die 9 Martii · S. Franciscæ Romanæ Viduæ\enspace\textit{St. Frances of Rome, Widow}}
\begin{paracol}{2}
\LA
\RefLine{Lectiones I–III: ex Commúni non Virginum non Martyrum (C7a).}
\switchcolumn
\EN
\RefLine{Lessons I–III: from the Common of Widows (C7a).}
\switchcolumn*

\LA
\RefLine{Lectiones VII–VIII: ex Commúni non Virginum non Martyrum (C7a).}
\switchcolumn
\EN
\RefLine{Lessons VII–VIII: from the Common of Widows (C7a).}
\switchcolumn*

\LA
\RefLine{Lectio IX: de Scriptura occurrente.}
\switchcolumn
\EN
\RefLine{Lesson IX: from the Scripture of the day.}
\switchcolumn*

\LA
\LessonHead{Lectio IV}
\switchcolumn
\EN
\LessonHead{Lesson IV}
\switchcolumn*

\LA
\noindent Francisca nobilis matrona Romana, ab ineunte ætate illustria dedit virtutum exempla: etenim pueriles ludos et illecebras mundi respuens, solitudine et oratione magnopere delectabatur. Undecim annos nata, virginitatem suam Deo consecrare, et monasterium ingredi proposuit: parentum tamen voluntati humiliter obtemperans, Laurentio de Pontianis, juveni æque diviti ac nobili nupsit. In matrimonio arctioris vitæ propositum, quantum licuit, semper retinuit: a spectaculis, conviviis, aliisque hujusmodi oblectamentis abhorrens, lanea ac vulgari veste utens, et quidquid a domesticis curis supererat temporis, orationi aut proximorum utilitati tribuens: in id vero maxima sollicitudine incumbens, ut matronas Romanas a pompis sæculi et ornatus vanitate revocaret. Quapropter domum Oblatarum, sub regula sancti Benedicti, congregationis Montis Oliveti, adhuc viro alligata, in Urbe instituit. Viri exsilium, bonorum jacturam, ac universæ domus mœrorem non modo constantissime toleravit: sed gratias agens cum beato Job, illud frequenter usurpabat: Dominus dedit, Dominus abstulit: sit nomen Domini benedictum.\par
\switchcolumn
\EN
\noindent The noble Roman matron Frances A (was born in the year 1384, and) was a pattern of godliness from her earliest years. As a child she shrank from games, and set no store by the amusements of the world, but delighted to be continually alone and engaged in prayer. At the age of eleven years she desired to consecrate her virginity to God, and to enter a convent, but humbly yielded obedience to the wishes of her parents, and was married to Lawrence de' Pontiani, a young man whose rank was equal to his wealth. As a wife she persevered, as far as she lawfully could, in her determination to lead an austere life; she abstained as much as possible from going to shows, feasts, and such like amusements, dressed plainly in woollen stuffs, and spent in prayer or the service of her neighbour whatever time she did not occupy with her duties as mistress of her husband's house. She strove earnestly to wean the married women of Rome from the vanities of the world and the frivolities of dress. To this end she founded during her husband's lifetime the Sisterhood of the Oblates, under the rule of the Benedictine congregation called of the Mount of Olives. When it pleased God, (in the year 1413,) that her husband should be banished, all her goods taken away, and her home ruined, she meekly bowed down before His holy will, often repeating the words of the blessed Job The Lord gave, and the Lord hath taken away; blessed be the name of the Lord. (i. 21.)\par
\switchcolumn*

\LA
\begin{Resp}\Rsign Propter veritátem, et mansuetúdinem, et justítiam: \Star{} Et dedúcet te mirabíliter déxtera tua.\end{Resp}
\begin{Resp}\Vsign Spécie tua et pulchritúdine tua inténde, próspere procéde, et regna.\end{Resp}
\begin{Resp}\Rsign Et dedúcet te mirabíliter déxtera tua.\end{Resp}
\switchcolumn
\EN
\begin{Resp}\Rsign Because of truth, and meekness, and righteousness; \Star{} And thy right hand shall lead thee wonderfully.\end{Resp}
\begin{Resp}\Vsign In thy comeliness, and thy beauty, go forward, fare prosperously, and reign.\end{Resp}
\begin{Resp}\Rsign And thy right hand shall lead thee wonderfully.\end{Resp}
\switchcolumn*

\LA
\LessonHead{Lectio V}
\switchcolumn
\EN
\LessonHead{Lesson V}
\switchcolumn*

\LA
\noindent Viro defuncto, ad prædictam Oblatarum domum convolans, nudis pedibus, fune ad collum alligato, humi prostrata, multis cum lacrimis earum numero adscribi suppliciter postulavit. Voti compos facta, licet esset omnium mater, non alio tamen quam ancillæ, vilissimæque feminæ, et immunditiæ vasculi titulo gloriabatur. Quam vilem sui existimationem et verbo declaravit et exemplo: sæpe enim e suburbana vinea revertens, et lignorum fascem proprio capiti impositum deferens, vel eisdem onustum agens per Urbem asellum, pauperibus subveniebat, in quos étiam largas eleemosynas erogabat: ægrotantesque in xenodochiis visitans, non corporali tantum cibo, sed salutaribus monitis recreabat. Corpus suum vigiliis, jejuniis, cilicio, ferreo cingulo, crebrisque flagellis in servitutem redigere jugiter satagebat. Cibum illi semel in die herbæ et legumina, aqua potum præbuit. Hos tamen corporis cruciatus aliquando confessarii mandato, a cujus ore nutuque pendebat, modice temperavit.\par
\switchcolumn
\EN
\noindent On her husband's death she (in 1437) betook herself immediately to the house of the Oblates, and, with her feet bare and a rope round her neck, threw herself down on the threshold, entreating the sisters with tears to receive her into their number. When she obtained her wish, although she was the mother of them all, she would be among them only as one that served, glorying rather to be called the most degraded of women and a vessel of uncleanness. Her lowly esteem of herself was shown both by her word and example. She passed often through the city from a vineyard in the country carrying a bundle of sticks on her head, or driving an ass laden with faggots; she succoured the needy, for whom she collected large alms, and visited the sick in the hospitals, ministering to them both food for the body and exhortations healthful for their souls. She strove continually to bring her body into subjection by watchings, fastings, haircloth, the wearing of an iron girdle, and the often use of a scourge. She never ate but once a day, and then only vegetables, and she took no drink but water. These severities she however sometimes relaxed, in obedience to her confessor, on whose word and wishes she framed her customs.\par
\switchcolumn*

\LA
\begin{Resp}\Rsign Dilexísti justítiam, et odísti iniquitátem: \Star{} Proptérea unxit te Deus, Deus tuus, óleo lætítiæ.\end{Resp}
\begin{Resp}\Vsign Propter veritátem, et mansuetúdinem, et justítiam.\end{Resp}
\begin{Resp}\Rsign Proptérea unxit te Deus, Deus tuus, óleo lætítiæ.\end{Resp}
\switchcolumn
\EN
\begin{Resp}\Rsign Thou hast loved righteousness, and hated iniquity; \Star{} Therefore God, thy God, hath anointed thee with the oil of gladness.\end{Resp}
\begin{Resp}\Vsign Because of truth, and meekness, and righteousness.\end{Resp}
\begin{Resp}\Rsign Therefore God, thy God, hath anointed thee with the oil of gladness.\end{Resp}
\switchcolumn*

\LA
\LessonHead{Lectio VI}
\switchcolumn
\EN
\LessonHead{Lesson VI}
\switchcolumn*

\LA
\noindent Divina mysteria, præsertim vero Christi Domini passionem, tanto mentis ardore, tantaque lacrimarum vi contemplabatur, ut præ doloris magnitudine pene confici videretur. Sæpe étiam cum oraret, maxime sumpto sanctissimæ Eucharistiæ sacramento, spiritu in Deum elevata, ac cælestium contemplatione rapta, immobilis permanebat. Quapropter humani generis hostis variis eam contumeliis ac verberibus a proposito dimovere conabatur: quem tamen illa imperterrita semper elusit; Angeli præsertim præsidio, cujus familiari consuetudine gloriosum de eo triumphum reportavit. Gratia curationum et prophetiæ dono enituit, quo et futura prædixit, et cordium secreta penetravit. Non semel aquæ, vel per rivum decurrentes, vel e cælo labentes, intactam prorsus, dum Deo vacaret, reliquerunt. Modica panis fragmenta, quæ vix tribus sororibus reficiendis fuíssent satis, sic ejus precibus Dominus multiplicavit, ut quindecim inde exsaturatis, tantum superfuerit, ut canistrum impleverit: et aliquando earumdem sororum, extra Urbem mense Januario ligna parantium, sitim recentis uvæ racemis ex vitæ in arbore pendentibus mirabiliter obtentis, abunde expleverit. Denique meritis et miraculis clara, migravit ad Dominum, anno ætatis suæ quinquagesimo sexto. Quam Paulus quintus Pontifex Maximus in Sanctorum numerum retulit.\par
\switchcolumn
\EN
\noindent So great was her mental realisation of the things of God, and chiefly of the sufferings of the Lord Christ, and so abundant her tears in contemplating them, that she seemed sometimes about to sink under her grief. Often when she was engaged in prayer, and principally after she had received the Most Holy Sacrament of the Eucharist, her spirit became altogether lifted up to God, and she remained motionless, carried away by the thought of heavenly things. The enemy of man assailed her with diverse reproaches and buffetings to break her off her intent, but she feared him not, and with the help of an Angel whom God gave her to be her familiar friend, she won a noble victory over the tempter. God glorified her with the gifts of healing and of prophecy, whereby she foretold things to come, and saw the secrets of the hearts of men. More than once while her thoughts were busy in God she remained unwet by streams or rain. When there was left only bread enough for three sisters, the Lord at her prayers was pleased so to multiply it, that fifteen had enough, and the basket was filled again with the fragments. In the month of January also, when the sisters were gathering sticks in the country, and were thirsty, she satisfied them abundantly with bunches of fresh grapes from a tree. She departed to be with the Lord, famous for good works and miracles, in the fifty-sixth year of her age, (upon the 9th day of March, in the year of our Lord 1440.) The Supreme Pontiff Paul V. caused her to be numbered among the saints.\par
\switchcolumn*

\LA
\begin{Resp}\Rsign Fallax grátia, et vana est pulchritúdo: \Star{} Múlier timens Dóminum, ipsa laudábitur.\end{Resp}
\begin{Resp}\Vsign Date ei de fructu mánuum suárum, et laudent eam in portis ópera ejus.\end{Resp}
\begin{Resp}\Rsign Múlier timens Dóminum, ipsa laudábitur.\end{Resp}
\begin{Resp}\Vsign Glória Patri, et Fílio, \Star{} et Spirítui Sancto.\end{Resp}
\begin{Resp}\Rsign Múlier timens Dóminum, ipsa laudábitur.\end{Resp}
\switchcolumn
\EN
\begin{Resp}\Rsign Favour is deceitful, and beauty is vain \Star{} A woman that feareth God she shall be praised.\end{Resp}
\begin{Resp}\Vsign Give her of the fruit of her hands, and let her own works praise her in the gates.\end{Resp}
\begin{Resp}\Rsign A woman that feareth God she shall be praised.\end{Resp}
\begin{Resp}\Vsign Glory be to the Father, and to the Son, \Star{} and to the Holy Ghost.\end{Resp}
\begin{Resp}\Rsign A woman that feareth God she shall be praised.\end{Resp}
\switchcolumn*

\end{paracol}

\par\needspace{10\baselineskip}\addvspace{1.2em}{\centering\small\textsc{Die 10 Martii} · \textsc{10 March}\par}
\Office{Ss. Quadraginta Martyrum}{Forty Holy Martyrs}{Semiduplex}
\addcontentsline{toc}{subsection}{Die 10 Martii · Ss. Quadraginta Martyrum\enspace\textit{Forty Holy Martyrs}}
\begin{paracol}{2}
\LA
\RefLine{Lectiones I–III: ex Commúni Plurimorum Martyrum Pontificum (C3).}
\switchcolumn
\EN
\RefLine{Lessons I–III: from the Commune Plurimorum Martyrum Pontificum (C3).}
\switchcolumn*

\LA
\RefLine{Lectio IX: de Scriptura occurrente.}
\switchcolumn
\EN
\RefLine{Lesson IX: from the Scripture of the day.}
\switchcolumn*

\LA
\LessonHead{Lectio IV}
\switchcolumn
\EN
\LessonHead{Lesson IV}
\switchcolumn*

\LA
\noindent Licinio imperatore, et Agricolao præside, ad Sebasten Armeniæ urbem, quadraginta militum fides in Jesum Christum, et fortitudo in cruciatibus perferendis enituit. Qui sæpius in horribilem carcerem detrusi, vinculisque constricti, cum ora ipsorum lapidibus contusa fuíssent, hiemis tempore frigidissimo nudi sub aperto aëre supra stagnum rigens pernoctare jussi sunt, ut frigore congelati necarentur. Una autem erat omnium oratio: Quadraginta in stadium ingressi sumus, quadraginta item Domine corona donemur, ne una quidem huic numero desit. Est in honore hic numerus, quem tu quadraginta dierum jejunio decorasti, per quem divina lex ingressa est in orbem terrarum. Elias quadraginta dierum jejunio Deum quærens, ejus visionem consecutus est. Et hæc quidem illorum erat oratio.\par
\switchcolumn
\EN
\noindent While Licinius was Emperor and Agricolaus President, (in the year of our Lord 320,) forty soldiers at Sebaste, a city of Armenia, gave a singular instance of faith in Jesus Christ, and bravery under suffering. After being often remanded to an horrid prison-house, bound in fetters, and their mouths bruised with stones, they were ordered out in the depth of winter, stripped naked, and put upon a frozen pool, to die of cold during the night. The prayer of them all was the same O Lord, forty of us have begun to run in the race, grant that all forty may receive the crown, let not one be wanting at the last. Behold, is it not an honourable number in thy sight, Who didst bless the fast of forty days, and at the end thy Divine Law came forth to the earth? When also Elias sought thee, Thou, O God, didst reveal thyself unto him when he had fasted for forty days. Even so was their petition.\par
\switchcolumn*

\LA
\begin{Resp}\Rsign Sancti tui, Dómine, mirábile consecúti sunt iter, serviéntes præcéptis tuis, ut inveniréntur illǽsi in aquis válidis: \Star{} Terra appáruit árida: et in Mari Rubro via sine impediménto.\end{Resp}
\begin{Resp}\Vsign Quóniam percússit petram, et fluxérunt aquæ, et torréntes inundavérunt.\end{Resp}
\begin{Resp}\Rsign Terra appáruit árida: et in Mari Rubro via sine impediménto.\end{Resp}
\switchcolumn
\EN
\begin{Resp}\Rsign thy Saints, O Lord, have passed a wonderful way, serving thy commandments, that they might be found without hurt in the midst of the mighty waters. \Star{} Dry land appeared, and, out of the Red Sea, a way without impediment.\end{Resp}
\begin{Resp}\Vsign He smote the rock, and the waters gushed out, and the streams overflowed.\end{Resp}
\begin{Resp}\Rsign Dry land appeared, and, out of the Red Sea, a way without impediment.\end{Resp}
\switchcolumn*

\LA
\LessonHead{Lectio V}
\switchcolumn
\EN
\LessonHead{Lesson V}
\switchcolumn*

\LA
\noindent Ceteris autem custodibus somno deditis, solus vigilabat janitor, qui et illos orantes, et luce circumfusos, et quosdam e cælo descendentes Angelos tamquam a Rege missos, qui coronas triginta novem militibus distribuerent intuens ita secum loquebatur: Quadraginta hi sunt: quadragesimi corona ubi est? Quæ dum cogitaret, unus ex illo numero, cui animus ad frigus ferendum defecerat in proximum tepefactum balneum desiliens, sanctos illos summo dolore affecit. Verum Deus illorum preces irritas esse non est passus: nam rei eventum admiratus janitor, mox custodibus e somno excitatis, detractisque sibi vestibus, ac se Christianum esse clara voce professus, Martyribus se adjunxit. Cum vero præsidis satellites janitorem quoque Christianum esse cognovissent bacillis comminuta omnium eorum crura fregerunt.\par
\switchcolumn
\EN
\noindent When the keepers were all asleep and the watchman only was awake, he heard them praying and saw a light shining round about them, and Angels coming down from heaven, as the messengers of the King, bearing nine-and-thirty crowns, and distributing them to the soldiers. Then he said within himself: Are not forty here? Where is the crown of the fortieth? And as he looked he saw one of them whose courage could not bear the cold, come and leap into a warm bath that stood by; and the Saints were grievously afflicted. Nevertheless God suffered not that their prayer should return unto them void; for the watchman wondered, and called the keepers, and stripped himself of his clothes; and, when with a loud voice he had confessed himself a Christian, he joined the Martyrs. When the servants of the President knew that the watchman also was a Christian, they brake the legs of them all with staves.\par
\switchcolumn*

\LA
\begin{Resp}\Rsign Vérbera carníficum non timuérunt Sancti Dei, moriéntes pro Christi nómine: \Star{} Ut herédes fíerent in domo Dómini.\end{Resp}
\begin{Resp}\Vsign Tradidérunt córpora sua propter Deum ad supplícia.\end{Resp}
\begin{Resp}\Rsign Ut herédes fíerent in domo Dómini.\end{Resp}
\switchcolumn
\EN
\begin{Resp}\Rsign The Saints of God shrank not from the stripes of the executioners, but died for Christ's Name's sake \Star{} That they might be made joint-heirs in the house of the Lord.\end{Resp}
\begin{Resp}\Vsign They gave their bodies for God's sake to death.\end{Resp}
\begin{Resp}\Rsign That they might be made joint-heirs in the house of the Lord.\end{Resp}
\switchcolumn*

\LA
\LessonHead{Lectio VI}
\switchcolumn
\EN
\LessonHead{Lesson VI}
\switchcolumn*

\LA
\noindent In eo supplicio mortui sunt omnes, præter Melithonem natu minimum. Quem cum præsens mater ejus, fractis cruribus, adhuc viventem vidisset, sic cohortata est: Fili, paulisper sustine; ecce Christus ad januam stat adjuvans te. Cum vero reliquorum corpora plaustris imponi cerneret, ut in rogum inférrentur, ac filium suum relinqui, quod speraret impia turba puerum, si vixisset, ad idolorum cultum revocari posse; ipso in humeros sublato, sancta mater vehicula Mártyrum corporibus onusta strenue prosequebatur: in cujus amplexu Melithon spiritum Deo reddidit, ejusque corpus in eumdem illum ceterorum Mártyrum rogum pia mater injecit: ut qui fide et virtute conjunctissimi fuerant, funeris étiam societate copulati, una in cælum pervenirent. Combustis illis, eorum reliquiæ projectæ in profluentem, cum mirabiliter in unum confluxissent locum, salvæ et integræ repertæ, honorifico sepulchro conditæ sunt.\par
\switchcolumn
\EN
\noindent Under this torment died they all, saving Melithon, who was the youngest. Now, his mother stood by, and when she saw that his legs were broken, but that he was yet alive, she cried, and said My son, have patience but a little longer. Behold how Christ standeth at the door to help thee. When she saw the bodies of all the others put upon carts and taken away to be burned, and that her son was left behind, because the multitude wickedly hoped that being but a lad, if he lived, he might yet be drawn to commit idolatry, the holy mother took him on her own shoulders and bravely followed behind the carts laden with the bodies of the Martyrs. In her arms Melithon gave up his soul to God, and the mother who loved him so well laid his body with her own hands upon the pile, with those of the other Martyrs, that, as they had all been one in faith and strength, in death they might not be divided, and might enter heaven together. After the burning, what remained of them was thrown into a running stream, but the ashes were all washed together into one place, and being found and rescued, they were laid in an honourable sepulchre.\par
\switchcolumn*

\LA
\begin{Resp}\Rsign Tamquam aurum in fornáce probávit eléctos Dóminus, et quasi holocáusti hóstiam accépit illos: et in témpore erit respéctus illórum: \Star{} Quóniam donum et pax est eléctis Dei.\end{Resp}
\begin{Resp}\Vsign Qui confídunt in illum, intélligent veritátem, et fidéles in dilectióne acquiéscent illi.\end{Resp}
\begin{Resp}\Rsign Quóniam donum et pax est eléctis Dei.\end{Resp}
\begin{Resp}\Vsign Glória Patri, et Fílio, \Star{} et Spirítui Sancto.\end{Resp}
\begin{Resp}\Rsign Quóniam donum et pax est eléctis Dei.\end{Resp}
\switchcolumn
\EN
\begin{Resp}\Rsign As gold in the furnace hath the Lord tried His chosen ones, and received them as a burnt-offering, and yet a while, and they shall be regarded; \Star{} For the grace of God, and His peace, are with His chosen.\end{Resp}
\begin{Resp}\Vsign They that put their trust in Him shall understand the truth and such as be faithful in love shall abide with Him.\end{Resp}
\begin{Resp}\Rsign For the grace of God, and His peace, are with His chosen.\end{Resp}
\begin{Resp}\Vsign Glory be to the Father, and to the Son, \Star{} and to the Holy Ghost.\end{Resp}
\begin{Resp}\Rsign For the grace of God, and His peace, are with His chosen.\end{Resp}
\switchcolumn*

\LA
\LessonHead{Lectio VII}
\switchcolumn
\EN
\LessonHead{Lesson VII}
\switchcolumn*

\LA
\LessonTitle{Léctio sancti Evangélii secúndum Lucam}
\switchcolumn
\EN
\LessonTitle{From the Holy Gospel according to Luke}
\switchcolumn*

\LA
\Cite{Luc 6:17-23}
\switchcolumn
\EN
\Cite{Luke 6:17-23}
\switchcolumn*

\LA
\noindent In illo témpore: Descéndens Jesus de monte, stetit in loco campéstri, et turba discipulórum ejus, et multitúdo copiósa plebis ab omni Judǽa, et Jerúsalem, et marítima, et Tyri, et Sidónis. Et réliqua.\par
\switchcolumn
\EN
\noindent At that time, Jesus came down from the mountain, and stood in the plain, and the company of His disciples, and a great multitude of people out of all Judea, and Jerusalem, and from the sea coast of Tyre and Sidon. And so on.\par
\switchcolumn*

\LA
\LessonTitle{Homilía sancti Ambrósii Epíscopi}
\switchcolumn
\EN
\LessonTitle{Homily by St Ambrose, Bishop (of Milan.)}
\switchcolumn*

\LA
\Source{Lib. 5. in Lucam. Cap. 6., post init.}
\switchcolumn
\EN
\Source{Bk. v. on Luke vi.}
\switchcolumn*

\LA
\noindent Advérte ómnia diligénter, quómodo et cum Apóstolis ascéndat et descéndat ad turbas. Quómodo enim turba nisi in húmili Christum vidéret? Non séquitur ad excélsa, non ascéndit ad sublímia. Dénique ubi descéndit, invénit infírmos: in excélsis enim infírmi esse non possunt. Hinc étiam Matthǽus docet in inferióribus débiles esse sanátos. Prius enim unusquísque sanándus est, ut paulátim virtútibus procedéntibus ascéndere possit ad montem; et ídeo quemque in inferióribus sanat, hoc est, a libídine révocat, injúriam cæcitátis avértit. Ad vúlnera nostra descéndit: ut usu quodam et cópia suæ natúræ, compartícipes nos fáciat esse regni cæléstis.\par
\switchcolumn
\EN
\noindent Mark well how Jesus goeth upward with His disciples, and downward to the multitude. How should the multitude behold Christ, save in a lower place? Such go not up to the things which are above; such attain not to the things which are high. And when Jesus cometh down, He findeth such as are diseased for such like go not up to the heights. Hence also Matthew saith that there were there all sick people, (iv. 23.) Of these every man had need of healing, that, when he had received strength, by and by, he might go up into the mountain. And therefore, being Himself come down, He healeth them in the plain, that is to say, He calleth them away from their lust, and freeth them of their blindness. He cometh down to our wounds, to the end that by a certain use of His nature, and by the abundance thereof, He might make us joint-heirs of the kingdom of heaven.\par
\switchcolumn*

\LA
\begin{Resp}\Rsign Propter testaméntum Dómini, et leges patérnas, Sancti Dei perstitérunt in amóre fraternitátis: \Star{} Quia unus fuit semper spíritus in eis, et una fides.\end{Resp}
\begin{Resp}\Vsign Ecce quam bonum et quam jucúndum habitáre fratres in unum.\end{Resp}
\begin{Resp}\Rsign Quia unus fuit semper spíritus in eis, et una fides.\end{Resp}
\switchcolumn
\EN
\begin{Resp}\Rsign Because of the covenant of the Lord, and the laws of their fathers, the Saints of God abode in brotherly love, \Star{} For one spirit and one faith was ever in them.\end{Resp}
\begin{Resp}\Vsign Behold how good and how pleasant it is for brethren to dwell together in unity.\end{Resp}
\begin{Resp}\Rsign For one spirit and one faith was ever in them.\end{Resp}
\switchcolumn*

\LA
\LessonHead{Lectio VIII}
\switchcolumn
\EN
\LessonHead{Lesson VIII}
\switchcolumn*

\LA
\noindent Beáti páuperes, quia vestrum est regnum Dei. Quátuor tantum beatitúdines sanctus Lucas Domínicas posuit, octo vero sanctus Matthǽus: sed in illis octo istæ quátuor sunt, et in quátuor istis illæ octo. Hic enim quátuor velut virtútes ampléxus est cardináles: ille in illis octo mýsticum númerum reserávit. Pro octáva enim multi inscribúntur Psalmi: et mandátum áccipis octo illis partem dare fortásse benedictiónibus. Sicut enim spei nostræ octáva perféctio est, ita octáva summa virtútum est.\par
\switchcolumn
\EN
\noindent Blessed be ye poor, for your's is the kingdom of God. Saint Luke giveth us but four of the Lord's Beatitudes, and Saint Matthew eight but in those eight are contained these four, and in these four those eight. For in these four are embraced the cardinal virtues and in those eight they are set forth in a number full of mystery. It is written at the head of more than one of the Psalms that they are for the octave, and thou hast received the commandment Give a portion to seven, and also to eight to seven or eight what? Perchance degrees of blessedness. For as this eighth Beatitude doth name the most glorious realization of our hope the kingdom of Heaven so doth it also name the most royal exertion of our strength blessed are they which are persecuted.\par
\switchcolumn*

\LA
\begin{Resp}\Rsign Sancti mei, qui in carne pósiti, certámen habuístis: \Star{} Mercédem labóris ego reddam vobis.\end{Resp}
\begin{Resp}\Vsign Veníte benedícti Patris mei, percípite regnum.\end{Resp}
\begin{Resp}\Rsign Mercédem labóris ego reddam vobis.\end{Resp}
\begin{Resp}\Vsign Glória Patri, et Fílio, \Star{} et Spirítui Sancto.\end{Resp}
\begin{Resp}\Rsign Mercédem labóris ego reddam vobis.\end{Resp}
\switchcolumn
\EN
\begin{Resp}\Rsign O ye My Saints, who, being in the flesh, didst have striving \Star{} I will render unto you a reward of your labours.\end{Resp}
\begin{Resp}\Vsign Come, ye blessed of My Father, inherit the kingdom\end{Resp}
\begin{Resp}\Rsign I will render unto you a reward of your labours.\end{Resp}
\begin{Resp}\Vsign Glory be to the Father, and to the Son, \Star{} and to the Holy Ghost.\end{Resp}
\begin{Resp}\Rsign I will render unto you a reward of your labours.\end{Resp}
\switchcolumn*

\end{paracol}

\par\needspace{10\baselineskip}\addvspace{1.2em}{\centering\small\textsc{Die 12 Martii} · \textsc{12 March}\par}
\Office{S. Gregorii Papæ Confessoris et Ecclesiæ Doctoris}{S. Gregory the Great, Pope, Confessor and Doctor of the Church}{Duplex}
\addcontentsline{toc}{subsection}{Die 12 Martii · S. Gregorii Papæ Confessoris et Ecclesiæ Doctoris\enspace\textit{S. Gregory the Great, Pope, Confessor and Doctor of the Church}}
\begin{paracol}{2}
\LA
\RefLine{Lectiones I–III: ex Commúni Doctoris Pontificis (C4a).}
\switchcolumn
\EN
\RefLine{Lessons I–III: from the Common of a Doctor Bishop (C4a).}
\switchcolumn*

\LA
\RefLine{Lectio IX: de Scriptura occurrente.}
\switchcolumn
\EN
\RefLine{Lesson IX: from the Scripture of the day.}
\switchcolumn*

\LA
\LessonHead{Lectio IV}
\switchcolumn
\EN
\LessonHead{Lesson IV}
\switchcolumn*

\LA
\noindent Gregórius Magnus, Románus, Gordiáni senatóris fílius, adoléscens philosophíæ óperam dedit; et prætório offício functus, patre mórtuo, sex monastéria in Sicília ædificávit, Romæ séptimum sancti Andréæ nómine in suis ǽdibus, prope basílicam sanctórum Joánnis et Pauli ad Clivum Scauri, ubi Hilarióne ac Maximiáno magístris, monáchi vitam proféssus, póstea abbas fuit. Mox diáconus cardinális creátus, Constantinópolim a Pelágio Pontífice ad Tibérium Constantínum imperatórem legátus míttitur: apud quem memorábile étiam illud effécit, quod Eutýchium patriárcham, qui scrípserat contra veram ac tractábilem córporum resurrectiónem, ita convícit, ut ejus librum imperátor in ignem injíceret. Quare Eutýchius paulo post cum in morbum incidísset, instánte morte, pellem manus suæ tenébat, multis præséntibus, dicens: Confíteor quia omnes in hac carne resurgémus.\par
\switchcolumn
\EN
\noindent Gregory the Great was a Roman, the son of Gordian the Senator, (and was born about the year of our Lord 540.) As a young man he studied philosophy, and afterwards discharged the office of Praetor. After his father's death he built six monasteries in Sicily, and a seventh in honour of St. Andrew, in his own house at Rome, hard by the Church of Saints John and Paul at the ascent of the hill Scaurus. In this monastery of St. Andrew, he and his masters, Hilarion and Maximian, professed themselves monks, and Gregory was afterwards Abbot. Later on, he was created a Cardinal Deacon, and sent to Constantinople as legate from Pope Pelagius to the Emperor Tiberius Constantine. Before the Emperor he so successfully disputed against the Patriarch Eutychius, who had denied that our bodies shall verily and indeed rise again, that the Prince threw the book of the said Patriarch into the fire. Eutychius himself also, soon after fell sick, and when he felt death coming on him, he took hold of the skin of his own hand and said in the hearing of many that stood by: I acknowledge that we shall all rise again in this flesh.\par
\switchcolumn*

\LA
\begin{Resp}\Rsign Invéni David servum meum, óleo sancto meo unxi eum: \Star{} Manus enim mea auxiliábitur ei.\end{Resp}
\begin{Resp}\Vsign Nihil profíciet inimícus in eo, et fílius iniquitátis non nocébit ei.\end{Resp}
\begin{Resp}\Rsign Manus enim mea auxiliábitur ei.\end{Resp}
\switchcolumn
\EN
\begin{Resp}\Rsign I have found David My servant, with My holy oil have I anointed him \Star{} For My hand shall help him.\end{Resp}
\begin{Resp}\Vsign The enemy shall prevail nothing against him, nor the son of wickedness afflict him.\end{Resp}
\begin{Resp}\Rsign For My hand shall help him.\end{Resp}
\switchcolumn*

\LA
\LessonHead{Lectio V}
\switchcolumn
\EN
\LessonHead{Lesson V}
\switchcolumn*

\LA
\noindent Romam rédiens, Pelágio pestiléntia subláto, summo ómnium consénsu Póntifex elígitur: quem honórem ne accíperet, quámdiu pótuit, recusávit; nam aliéno vestítu in spelúnca delítuit: ubi deprehénsus indício ígnæa colúmnæ, ad sanctum Petrum consecrátur. In pontificátu multa successóribus doctrínæ ac sanctitátis exémpla relíquit. Peregrínos quotídie ad mensam adhibébat: in quibus et Ángelum, et Dóminum Angelórum peregríni fácie accépit. Páuperes et urbános et extérnos, quorum númerum descríptum habébat, benígne sustentábat. Cathólicam fidem multis locis labefactátam restítuit. Nam Donatístas in Africa, Ariános in Hispánia représsit: Agnoítas Alexandría ejécit. Pállium Syágrio Augustodunénsi epíscopo dare nóluit, nisi neóphytos hæréticos expélleret ex Gállia. Gothos hǽresim Ariánam relínquere coégit. Missis in Británniam doctis et sanctis viris Augustíno et áliis mónachis, ínsulam ad Jesu Christi fidem convértit, vere a Beda presbýtero Angliæ vocátus Apóstolus. Joánnis patriárchæ Constantinopolitáni audáciam fregit, qui sibi universális Ecclésiæ epíscopi nomen arrogábat. Maurítium imperatórem, eos qui mílites fuíssent, mónachos fíeri prohibéntem, a senténtia detérruit.\par
\switchcolumn
\EN
\noindent Gregory returned to Rome, and, Pelagius being dead of a plague, he was unanimously chosen Pope. This honour he refused as long as he could. He disguised himself and took refuge in a cave, but was betrayed by a fiery pillar. Being discovered and overruled, he was consecrated at the grave of St. Peter, \textasciitilde{}(upon the 3rd day of September, in the year 590.) He left behind him many examples of doctrine and holiness to them that have followed him in the Papacy. Every day he brought pilgrims to his table, and among them he entertained not an Angel only, but the very Lord of Angels in the guise of a pilgrim. He tenderly cared for the poor, of whom he kept a list, as well without as within the city. He restored the Catholic faith in many places where it had been overthrown. He fought successfully against the Donatists in Africa and the Arians in Spain. He cleansed Alexandria of the Agnoites. He refused to give the Pall to Syagrius, Bishop of Autun, unless he would expel the Neophyte heretics from Gaul. He caused the Goths to abandon the Arian heresy. He sent into Britain Augustine and diverse other learned and holy monks, who brought the inhabitants of that island to believe in Jesus Christ. Hence Gregory is justly called by Bede, the Priest of Jarrow, the Apostle of England. He rebuked the presumption of John, Patriarch of Constantinople, who had taken to himself the title of Bishop of the Universal Church, and he dissuaded the Emperor Maurice from forbidding soldiers to become monks.\par
\switchcolumn*

\LA
\begin{Resp}\Rsign Pósui adjutórium super poténtem, et exaltávi eléctum de plebe mea: \Star{} Manus enim mea auxiliábitur ei.\end{Resp}
\begin{Resp}\Vsign Invéni David servum meum, óleo sancto meo unxi eum.\end{Resp}
\begin{Resp}\Rsign Manus enim mea auxiliábitur ei.\end{Resp}
\switchcolumn
\EN
\begin{Resp}\Rsign I have laid help upon one that is mighty, and have exalted one chosen out of My people \Star{} For My hand shall help him.\end{Resp}
\begin{Resp}\Vsign I have found David My servant, with My holy oil have I anointed him.\end{Resp}
\begin{Resp}\Rsign For My hand shall help him.\end{Resp}
\switchcolumn*

\LA
\LessonHead{Lectio VI}
\switchcolumn
\EN
\LessonHead{Lesson VI}
\switchcolumn*

\LA
\noindent Ecclésiam ornávit sanctíssimis institútis et légibus. Apud sanctum Petrum coácta sýnodo, multa constítuit: in iis, Ut in Missa Kýrie eléison nóvies repeterétur: ut extra id tempus, quod continétur Septuagésima et Pascha, Alle­lúja dicerétur: ut adderétur in Cánone, Diésque nostros in tua pace dispónas. Litanías, Statiónes, et ecclesiásticum offícium auxit. Quátuor concíliis, Nicǽno, Constantinopolitáno, Ephesíno, Chalcedonénsi, tamquam quátuor Evangéliis honórem habéri vóluit. Epíscopis Sicíliæ, qui ex antíqua ecclesiárum consuetúdine Romam síngulis triénniis conveniébant, quinto quoque anno semel veníre indúlsit. Multos libros confécit: quos cum dictáret, testátus est Petrus diáconus, se Spíritum Sanctum colúmbæ spécie in ejus cápite sæpe vidísse. Admirabília sunt quæ dixit, fecit, scripsit, decrévit, præsértim infírma semper et ægra valetúdine. Qui dénique multis éditis miráculis, pontificátus anno décimo tértio, mense sexto, die décimo, quarto Idus Mártii, qui dies festus a Græcis étiam propter insígnem hujus Pontíficis sapiéntiam ac sanctitátem præcípuo honóre celebrátur, ad cæléstem beatitúdinem evocátus est. Cujus corpus sepúltum est in basílica sancti Petri, prope Secretárium.\par
\switchcolumn
\EN
\noindent Gregory adorned the Church with holy customs and laws. He called together a Synod in the Church of St. Peter, and therein ordained many things; among others, the ninefold repetition of the words Kyrie eleison in the Mass, the saying of the word Alle­luja in the Church service except between Septuagesima inclusive and Easter exclusive, and the addition to the Canon of the Mass of the words Do Thou order all our days in thy peace. He increased the Litanies, the number of the Churches where is held the observance called a Station; and the length of the Church Service. He would that the four Councils of Nice, Constantinople, Ephesus, and Chalcedon should be honoured like four Gospels. He released the Sicilian Bishops from visiting Rome every three years, willing them to come instead once every five years. He was the author of many books, and Peter the Deacon declareth that he often saw the Holy Ghost on his head in the form of a dove when he was dictating them. It is a marvel how much he spoke, did, wrote, and legislated, suffering all the while from a weak and sickly body. He worked many miracles. At last God called him away to be blessed for ever in heaven, in the thirteenth year, sixth month, and tenth day of his Pontificate, being the 12 th day of March, (in the year of salvation 604.) This day is observed by the Greeks, as well as by us, as a festival, on account of the eminent wisdom and holiness of this Pope. His body was buried in the Church of St. Peter, hard by the Private Chapel.\par
\switchcolumn*

\LA
\begin{Resp}\Rsign Iste est, qui ante Deum magnas virtútes operátus est, et omnis terra doctrína ejus repléta est: \Star{} Ipse intercédat pro peccátis ómnium populórum.\end{Resp}
\begin{Resp}\Vsign Iste est, qui contémpsit vitam mundi, et pervénit ad cæléstia regna.\end{Resp}
\begin{Resp}\Rsign Ipse intercédat pro peccátis ómnium populórum.\end{Resp}
\begin{Resp}\Vsign Glória Patri, et Fílio, \Star{} et Spirítui Sancto.\end{Resp}
\begin{Resp}\Rsign Ipse intercédat pro peccátis ómnium populórum.\end{Resp}
\switchcolumn
\EN
\begin{Resp}\Rsign This is he which wrought great wonders before God, and the whole earth is full of his teaching \Star{} May he pray for all people, that their sins may be forgiven unto them\end{Resp}
\begin{Resp}\Vsign This is he which loved not his life in this world, and hath attained unto the kingdom of heaven.\end{Resp}
\begin{Resp}\Rsign May he pray for all people, that their sins may be forgiven unto them\end{Resp}
\begin{Resp}\Vsign Glory be to the Father, and to the Son, \Star{} and to the Holy Ghost.\end{Resp}
\begin{Resp}\Rsign May he pray for all people, that their sins may be forgiven unto them\end{Resp}
\switchcolumn*

\LA
\LessonHead{Lectio VII}
\switchcolumn
\EN
\LessonHead{Lesson VII}
\switchcolumn*

\LA
\LessonTitle{Léctio sancti Evangélii secúndum Matthǽum}
\switchcolumn
\EN
\LessonTitle{From the Holy Gospel according to Matthew}
\switchcolumn*

\LA
\Cite{Matt 16:13-19}
\switchcolumn
\EN
\Cite{Matt 16:13-19}
\switchcolumn*

\LA
\noindent In illo témpore: Venit Jesus in partes Cæsaréæ Philíppi, et interrogábat discípulos suos, dicens: Quem dicunt hómines esse Fílium hóminis? Et réliqua.\par
\switchcolumn
\EN
\noindent At that time: When Jesus came into the coasts of Caesarea Philippi, he asked his disciples, saying, Whom do men say that I the Son of man am? And so on.\par
\switchcolumn*

\LA
\LessonTitle{Homilía sancti Leónis Papæ}
\switchcolumn
\EN
\LessonTitle{A Homily by St. Leo the Pope}
\switchcolumn*

\LA
\Source{Sermo 2 in anniversario assumpt. suæ, ante medium}
\switchcolumn
\EN
\Source{Sermo 2 in anniversario assumpt. suæ ante medium}
\switchcolumn*

\LA
\noindent Cum, sicut evangélica lectióne reserátum est, interrogásset Dóminus discípulos, quem ipsum (multis divérsa opinántibus) créderent; respondissétque beátus Petrus, dicens: Tu es Christus Fílius Dei vivi; Dóminus ait: Beátus es, Simon Bar-Jona, quia caro et sanguis non revelávit tibi, sed Pater meus, qui in cælis est: et ego dico tibi, quia tu es Petrus, et super hanc petram ædificábo Ecclésiam meam, et portæ ínferi non prævalébunt advérsus eam. Et tibi dabo claves regni cælórum: et quodcúmque ligáveris super terram, erit ligátum et in cælis: et quodcúmque sólveris super terram, erit solútum et in cælis. Manet ergo disposítio veritátis, et beátus Petrus, in accépta fortitúdine petræ persevérans, suscépta Ecclésiæ gubernácula non relíquit.\par
\switchcolumn
\EN
\noindent When the Lord, as we read in the Gospel, asked his disciples who did men, amid their diverse speculations, believe him the Son of Man to be, blessed Peter answered and said: Thou art the Christ, the Son of the living God. And the Lord answered and said unto him: Blessed art thou, Simon Bar-Jona: for flesh and blood hath not revealed it unto thee, but my Father, which is in heaven: and I say also unto thee: That thou art Peter, and upon this rock I will build my Church, and the gates of hell shall not prevail against it; and I will give unto thee the keys of the kingdom of heaven; and whatsoever thou shalt bind on earth shall be bound in heaven; and whatsoever thou shalt loose on earth shall be loosed in heaven. But the dispensation of truth perdures, and blessed Peter, persevering in the strength of the rock which he hath received, hath not relinquished the position he assumed at the helm of the Church.\par
\switchcolumn*

\LA
\begin{Resp}\Rsign Amávit eum Dóminus, et ornávit eum: stolam glóriæ índuit eum, \Star{} Et ad portas paradísi coronávit eum.\end{Resp}
\begin{Resp}\Vsign Induit eum Dóminus lorícam fídei, et ornávit eum.\end{Resp}
\begin{Resp}\Rsign Et ad portas paradísi coronávit eum.\end{Resp}
\switchcolumn
\EN
\begin{Resp}\Rsign The Lord loved him and beautified him He clothed him with a robe of glory \Star{} And crowned him at the gates of Paradise.\end{Resp}
\begin{Resp}\Vsign The Lord hath put on him the breast-plate of faith, and hath adorned him.\end{Resp}
\begin{Resp}\Rsign And crowned him at the gates of Paradise.\end{Resp}
\switchcolumn*

\LA
\LessonHead{Lectio VIII}
\switchcolumn
\EN
\LessonHead{Lesson VIII}
\switchcolumn*

\LA
\noindent In univérsa namque Ecclésia, Tu es Christus Fílius Dei vivi, quotídie Petrus dicit; et omnis lingua, quæ confitétur Dóminum, magistério hujus vocis imbúitur. Hæc fides diábolum vincit et captivórum ejus víncula dissólvit. Hæc érutos mundo, ínserit cælo, et portæ ínferi advérsus eam prævalére non possunt. Tanta enim divínitus soliditáte muníta est, ut eam neque hærética umquam corrúmpere právitas, nec pagána potúerit superáre perfídia. His ítaque modis, dilectíssimi, rationábile obséquio celebrétur hodiérna festívitas: ut in persóna humilitátis meæ ille intelligátur, ille honorétur, in quo et ómnium pastórum sollicitúdo, cum commendatárum sibi óvium custódia persevérat, et cujus étiam dígnitas in indígno heréde non déficit.\par
\switchcolumn
\EN
\noindent In the universal Church it is as if Peter were still saying every day: Thou art the Christ, the Son of the living God. For every tongue which confesseth the Lord is taught that confession by the teaching of Peter. This is the Faith that overcometh the devil and looseth the bonds of his prisoners. This is the Faith which maketh men free of the world and bringeth them to heaven, and the gates of hell are impotent to prevail against it. This is the rock which God hath fortified with such ramparts of salvation, that the contagion of heresy will never be able to infect it, nor idolatry and unbelief to overcome it. And therefore, dearly beloved, we celebrate today’s festival with reasonable obedience, that in my humble person he may be acknowledged and honoured who doth continue to care for all the shepherds as well as sheep entrusted unto him, and who doth lose none of his dignity even in an unworthy successor.\par
\switchcolumn*

\LA
\begin{Resp}\Rsign In médio Ecclésiæ apéruit os ejus, \Star{} Et implévit eum Dóminus spíritu sapiéntiæ et intelléctus.\end{Resp}
\begin{Resp}\Vsign Jucunditátem et exsultatiónem thesaurizávit super eum.\end{Resp}
\begin{Resp}\Rsign Et implévit eum Dóminus spíritu sapiéntiæ et intelléctus.\end{Resp}
\begin{Resp}\Vsign Glória Patri, et Fílio, \Star{} et Spirítui Sancto.\end{Resp}
\begin{Resp}\Rsign Et implévit eum Dóminus spíritu sapiéntiæ et intelléctus.\end{Resp}
\switchcolumn
\EN
\begin{Resp}\Rsign In the midst of the congregation did the Lord open his mouth. \Star{} And filled him with the spirit of wisdom and understanding.\end{Resp}
\begin{Resp}\Vsign He made him rich with joy and gladness.\end{Resp}
\begin{Resp}\Rsign And filled him with the spirit of wisdom and understanding.\end{Resp}
\begin{Resp}\Vsign Glory be to the Father, and to the Son, \Star{} and to the Holy Ghost.\end{Resp}
\begin{Resp}\Rsign And filled him with the spirit of wisdom and understanding.\end{Resp}
\switchcolumn*

\end{paracol}

\PartTitle{Commune Sanctorum}{Common of the Saints}

\addcontentsline{toc}{chapter}{Commune Sanctorum\enspace\textit{Common of the Saints}}

\Office{Commune Apostolorum}{Commune Apostolorum}{}
\addcontentsline{toc}{subsection}{Commune Apostolorum\enspace\textit{Commune Apostolorum}}
\begin{paracol}{2}
\LA
\LessonHead{Lectio I}
\switchcolumn
\EN
\LessonHead{Lesson I}
\switchcolumn*

\LA
\LessonTitle{De Epístola prima beáti Pauli Apóstoli ad Corínthios}
\switchcolumn
\EN
\LessonTitle{Lesson from the first letter of St. Paul the Apostle to the Corinthians}
\switchcolumn*

\LA
\Cite{1 Cor 4:1-5}
\switchcolumn
\EN
\Cite{1 Cor 4:1-5}
\switchcolumn*

\LA
\noindent Sic nos exístimet homo ut minístros Christi, et dispensatóres mysteriórum Dei. Hic jam quǽritur inter dispensatóres ut fidélis quis inveniátur. Mihi autem pro mínimo est ut a vobis júdicer, aut ab humáno die: sed neque meípsum júdico. Nihil enim mihi cónscius sum, sed non in hoc justificátus sum: qui autem júdicat me, Dóminus est. Ítaque nolíte ante tempus judicáre, quoadúsque véniat Dóminus: qui et illuminábit abscóndita tenebrárum, et manifestábit consília córdium: et tunc laus erit unicuíque a Deo.\par
\switchcolumn
\EN
\noindent Let a man so account of us as of the ministers of Christ, and the dispensers of the mysteries of God. Here now it is required among the dispensers, that a man be found faithful. But to me it is a very small thing to be judged by you, or by man's day; but neither do I judge my own self. For I am not conscious to myself of any thing, yet am I not hereby justified; but he that judgeth me, is the Lord. Therefore judge not before the time; until the Lord come, who both will bring to light the hidden things of darkness, and will make manifest the counsels of the hearts; and then shall every man have praise from God.\par
\switchcolumn*

\LA
\begin{Resp}\Rsign Ecce ego mitto vos sicut oves in médio lupórum, dicit Dóminus: \Star{} Estóte ergo prudéntes sicut serpéntes, et símplices sicut colúmbæ.\end{Resp}
\begin{Resp}\Vsign Dum lucem habétis, crédite in lucem, ut fílii lucis sitis.\end{Resp}
\begin{Resp}\Rsign Estóte ergo prudéntes sicut serpéntes, et símplices sicut colúmbæ.\end{Resp}
\switchcolumn
\EN
\begin{Resp}\Rsign Behold, I send you as sheep in the midst of wolves, said the Lord: \Star{} Be ye, therefore, wise as serpents and simple as doves.\end{Resp}
\begin{Resp}\Vsign Whilst you have the light, believe in the light, that you may be the children of light.\end{Resp}
\begin{Resp}\Rsign Be ye therefore wise as serpents and simple as doves.\end{Resp}
\switchcolumn*

\LA
\LessonHead{Lectio II}
\switchcolumn
\EN
\LessonHead{Lesson II}
\switchcolumn*

\LA
\Cite{1 Cor 4:6-9}
\switchcolumn
\EN
\Cite{1 Cor 4:6-9}
\switchcolumn*

\LA
\noindent Hæc autem, fratres, transfigurávi in me et Apóllo, propter vos: ut in nobis discátis, ne supra quam scriptum est, unus advérsus álterum inflétur pro álio. Quis enim te discérnit? Quid autem habes quod non accepísti? Si autem accepísti, quid gloriáris quasi non accéperis? Jam saturáti estis, jam dívites facti estis: sine nobis regnátis: et útinam regnétis, ut et nos vobíscum regnémus. Puto enim quod Deus nos Apóstolos novíssimos osténdit, tamquam morti destinátos: quia spectáculum facti sumus mundo, et Ángelis, et homínibus.\par
\switchcolumn
\EN
\noindent But these things, brethren, I have in a figure transferred to myself and to Apollo, for your sakes; that in us you may learn, that one be not puffed up against the other for another, above that which is written. For who distinguisheth thee? Or what hast thou that thou hast not received? And if thou hast received, why dost thou glory, as if thou hadst not received it? You are now full; you are now become rich; you reign without us; and I would to God you did reign, that we also might reign with you. For I think that God hath set forth us apostles, the last, as it were men appointed to death: we are made a spectacle to the world, and to angels, and to men.\par
\switchcolumn*

\LA
\begin{Resp}\Rsign Tóllite jugum meum super vos, dicit Dóminus, et díscite a me, quia mitis sum et húmilis corde: \Star{} Jugum enim meum suáve est, et onus meum leve.\end{Resp}
\begin{Resp}\Vsign Et inveniétis réquiem animábus vestris.\end{Resp}
\begin{Resp}\Rsign Jugum enim meum suáve est, et onus meum leve.\end{Resp}
\switchcolumn
\EN
\begin{Resp}\Rsign Take up my yoke upon you, said the Lord, and learn of me, because I am meek, and humble of heart. \Star{} For my yoke is sweet and my burden light.\end{Resp}
\begin{Resp}\Vsign And you shall find rest to your souls.\end{Resp}
\begin{Resp}\Rsign For my yoke is sweet and my burden light.\end{Resp}
\switchcolumn*

\LA
\LessonHead{Lectio III}
\switchcolumn
\EN
\LessonHead{Lesson III}
\switchcolumn*

\LA
\Cite{1 Cor 4:10-15}
\switchcolumn
\EN
\Cite{1 Cor 4:10-15}
\switchcolumn*

\LA
\noindent Nos stulti propter Christum, vos autem prudéntes in Christo: nos infírmi, vos autem fortes: vos nóbiles, nos autem ignóbiles. Usque in hanc horam et esurímus, et sitímus, et nudi sumus, et cólaphis cǽdimur, et instábiles sumus, et laborámus operántes mánibus nostris: maledícimur, et benedícimus: persecutiónem pátimur, et sustinémus: blasphemámur, et obsecrámus: tamquam purgaménta hujus mundi facti sumus, ómnium peripséma usque adhuc. Non ut confúndam vos, hæc scribo, sed ut fílios meos caríssimos móneo. Nam si decem míllia pædagogórum habeátis in Christo, sed non multos patres, nam in Christo Jesu per Evangélium ego vos génui.\par
\switchcolumn
\EN
\noindent We are fools for Christ's sake, but you are wise in Christ; we are weak, but you are strong; you are honourable, but we without honour. Even unto this hour we both hunger and thirst, and are naked, and are buffeted, and have no fixed abode; And we labour, working with our own hands: we are reviled, and we bless; we are persecuted, and we suffer it. We are blasphemed, and we entreat; we are made as the refuse of this world, the offscouring of all even until now. I write not these things to confound you; but I admonish you as my dearest children. For if you have ten thousand instructors in Christ, yet not many fathers. For in Christ Jesus, by the gospel, I have begotten you.\par
\switchcolumn*

\LA
\begin{Resp}\Rsign Dum stetéritis ante reges et prǽsides, nolíte cogitáre quómodo aut quid loquámini: \Star{} Dábitur enim vobis in illa hora, quid loquámini.\end{Resp}
\begin{Resp}\Vsign Non enim vos estis qui loquímini: sed Spíritus Patris vestri, qui lóquitur in vobis.\end{Resp}
\begin{Resp}\Rsign Dábitur enim vobis in illa hora, quid loquámini.\end{Resp}
\begin{Resp}\Vsign Glória Patri, et Fílio, \Star{} et Spirítui Sancto.\end{Resp}
\begin{Resp}\Rsign Dábitur enim vobis in illa hora, quid loquámini.\end{Resp}
\switchcolumn
\EN
\begin{Resp}\Rsign But when they shall deliver you up to the judges, take no thought how or what to speak: \Star{} for it shall be given you in that hour what to speak:\end{Resp}
\begin{Resp}\Vsign For it is not you that speak, but the spirit of your Father that speaketh in you.\end{Resp}
\begin{Resp}\Rsign For it shall be given you in that hour what to speak.\end{Resp}
\begin{Resp}\Vsign Glory be to the Father, and to the Son, \Star{} and to the Holy Ghost.\end{Resp}
\begin{Resp}\Rsign For it shall be given you in that hour what to speak.\end{Resp}
\switchcolumn*

\end{paracol}

\Office{Commune Evangelistarum}{Commune Evangelistarum}{}
\addcontentsline{toc}{subsection}{Commune Evangelistarum\enspace\textit{Commune Evangelistarum}}
\begin{paracol}{2}
\LA
\LessonHead{Lectio I}
\switchcolumn
\EN
\LessonHead{Lesson I}
\switchcolumn*

\LA
\LessonTitle{Incipit liber Ezechiélis Prophétæ}
\switchcolumn
\EN
\LessonTitle{Lesson from the book of Ezekiel}
\switchcolumn*

\LA
\Cite{Ezech 1:1-4}
\switchcolumn
\EN
\Cite{Ezek 1:1-4}
\switchcolumn*

\LA
\noindent Et factum est in trigésimo anno, in quarto, in quinta mensis, cum essem in médio captivórum juxta flúvium Chobar, apérti sunt cæli, et vidi visiónes Dei. In quinta mensis, ipse est annus quintus transmigratiónis regis Jóachim, factum est verbum Dómini ad Ezechiélem fílium Buzi sacerdótem, in terra Chaldæórum, secus flumen Chobar: et facta est super eum ibi manus Dómini. Et vidi, et ecce ventus túrbinis veniébat ab Aquilóne, et nubes magna, et ignis invólvens, et splendor in circúitu ejus: et de médio ejus, quasi spécies eléctri, id est, de médio ignis:\par
\switchcolumn
\EN
\noindent Now it came to pass in the thirtieth year, in the fourth month, on the fifth day of the month, when I was in the midst of the captives by the river Chobar, the heavens were opened, and I saw the visions of God. On the fifth day of the month, the same was the fifth year of the captivity of king Joachin, The word of the Lord came to Ezechiel the priest the son of Buzi in the land of the Chaldeans, by the river Chobar: and the hand of the Lord was there upon him. And I saw, and behold a whirlwind came out of the north: and a great cloud, and a fire infolding it, and brightness was about it: and out of the midst thereof, that is, out of the midst of the fire, as it were the resemblance of amber.\par
\switchcolumn*

\LA
\begin{Resp}\Rsign Ecce ego mitto vos sicut oves in médio lupórum, dicit Dóminus: \Star{} Estóte ergo prudéntes sicut serpéntes, et símplices sicut colúmbæ.\end{Resp}
\begin{Resp}\Vsign Dum lucem habétis, crédite in lucem, ut fílii lucis sitis.\end{Resp}
\begin{Resp}\Rsign Estóte ergo prudéntes sicut serpéntes, et símplices sicut colúmbæ.\end{Resp}
\switchcolumn
\EN
\begin{Resp}\Rsign Behold, I send you as sheep in the midst of wolves, said the Lord: \Star{} Be ye, therefore, wise as serpents and simple as doves.\end{Resp}
\begin{Resp}\Vsign Whilst you have the light, believe in the light, that you may be the children of light.\end{Resp}
\begin{Resp}\Rsign Be ye therefore wise as serpents and simple as doves.\end{Resp}
\switchcolumn*

\LA
\LessonHead{Lectio II}
\switchcolumn
\EN
\LessonHead{Lesson II}
\switchcolumn*

\LA
\Cite{Ezech 1:5-9}
\switchcolumn
\EN
\Cite{Ezek 1:5-9}
\switchcolumn*

\LA
\noindent Et in médio ejus similitúdo quátuor animálium, et hic aspéctus eórum, similitúdo hóminis in eis. Quátuor fácies uni, et quátuor pennæ uni. Pedes eórum, pedes recti, et planta pedis eórum quasi planta pedis vítuli: et scintíllæ quasi aspéctus æris candéntis. Et manus hóminis sub pennis eórum, in quátuor pártibus: et fácies et pennas per quátuor partes habébant, junctǽque erant pennæ eórum altérius ad álterum. Non revertebántur cum incéderent, sed unumquódque ante fáciem suam gradiebátur.\par
\switchcolumn
\EN
\noindent And in the midst thereof the likeness of four living creatures: and this was their appearance: there was the likeness of a man in them. Every one had four faces, and every one four wings. Their feet were straight feet, and the sole of their foot was like the sole of a calf's foot, and they sparkled like the appearance of glowing brass. And they had the hands of a man under their wings on their four sides: and they had faces, and wings on the four sides, And the wings of one were joined to the wings of another. They turned not when they went: but every one went straight forward.\par
\switchcolumn*

\LA
\begin{Resp}\Rsign Tóllite jugum meum super vos, dicit Dóminus, et díscite a me, quia mitis sum et húmilis corde: \Star{} Jugum enim meum suáve est, et onus meum leve.\end{Resp}
\begin{Resp}\Vsign Et inveniétis réquiem animábus vestris.\end{Resp}
\begin{Resp}\Rsign Jugum enim meum suáve est, et onus meum leve.\end{Resp}
\switchcolumn
\EN
\begin{Resp}\Rsign Take up my yoke upon you, said the Lord, and learn of me, because I am meek, and humble of heart. \Star{} For my yoke is sweet and my burden light.\end{Resp}
\begin{Resp}\Vsign And you shall find rest to your souls.\end{Resp}
\begin{Resp}\Rsign For my yoke is sweet and my burden light.\end{Resp}
\switchcolumn*

\LA
\LessonHead{Lectio III}
\switchcolumn
\EN
\LessonHead{Lesson III}
\switchcolumn*

\LA
\Cite{Ezech 1:10-12}
\switchcolumn
\EN
\Cite{Ezek 1:10-12}
\switchcolumn*

\LA
\noindent Similitúdo autem vultus eórum, fácies hóminis et fácies leónis a dextris ipsórum quátuor, fácies autem bovis a sinístris ipsórum quátuor, et fácies áquilæ désuper ipsórum quátuor. Fácies eórum et pennæ eórum exténtæ désuper: duæ pennæ singulórum jungebántur, et duæ tegébant córpora eórum: et unumquódque eórum coram fácie sua ambulábat: ubi erat ímpetus spíritus, illuc gradiebántur, nec revertebántur cum ambulárent.\par
\switchcolumn
\EN
\noindent And as for the likeness of their faces: there was the face of a man, and the face of a lion on the right side of all the four: and the face of an ox, on the left side of all the four: and the face of an eagle over all the four. And their faces, and their wings were stretched upward: two wings of every one were joined, and two covered their bodies: And every one of them went straight forward: whither the impulse of the spirit was to go, thither they went: and they turned not when they went.\par
\switchcolumn*

\LA
\begin{Resp}\Rsign Dum stetéritis ante reges et prǽsides, nolíte cogitáre quómodo aut quid loquámini: \Star{} Dábitur enim vobis in illa hora, quid loquámini.\end{Resp}
\begin{Resp}\Vsign Non enim vos estis qui loquímini: sed Spíritus Patris vestri, qui lóquitur in vobis.\end{Resp}
\begin{Resp}\Rsign Dábitur enim vobis in illa hora, quid loquámini.\end{Resp}
\begin{Resp}\Vsign Glória Patri, et Fílio, \Star{} et Spirítui Sancto.\end{Resp}
\begin{Resp}\Rsign Dábitur enim vobis in illa hora, quid loquámini.\end{Resp}
\switchcolumn
\EN
\begin{Resp}\Rsign But when they shall deliver you up to the judges, take no thought how or what to speak: \Star{} for it shall be given you in that hour what to speak:\end{Resp}
\begin{Resp}\Vsign For it is not you that speak, but the spirit of your Father that speaketh in you.\end{Resp}
\begin{Resp}\Rsign For it shall be given you in that hour what to speak.\end{Resp}
\begin{Resp}\Vsign Glory be to the Father, and to the Son, \Star{} and to the Holy Ghost.\end{Resp}
\begin{Resp}\Rsign For it shall be given you in that hour what to speak.\end{Resp}
\switchcolumn*

\end{paracol}

\Office{Commune Unius Martyris Pontificis}{Common of a Single Martyr Bishop}{}
\addcontentsline{toc}{subsection}{Commune Unius Martyris Pontificis\enspace\textit{Common of a Single Martyr Bishop}}
\begin{paracol}{2}
\LA
\LessonHead{Lectio VII}
\switchcolumn
\EN
\LessonHead{Lesson VII}
\switchcolumn*

\LA
\LessonTitle{Léctio sancti Evangélii secúndum Lucam}
\switchcolumn
\EN
\LessonTitle{From the Holy Gospel according to Luke}
\switchcolumn*

\LA
\Cite{Luc 14:26-33}
\switchcolumn
\EN
\Cite{Luke 14:26-33}
\switchcolumn*

\LA
\noindent In illo témpore: Dixit Jesus turbis: Si quis venit ad me, et non odit patrem suum, et matrem, et uxórem, et fílios, et fratres, et soróres, adhuc autem et ánimam suam, non potest meus esse discípulus. Et réliqua.\par
\switchcolumn
\EN
\noindent At that time, Jesus said unto the multitudes: If any man come to Me, and hate not his father, and mother, and wife, and children, and brethren, and sisters, yea, and his own life also, he cannot be My disciple. And so on.\par
\switchcolumn*

\LA
\LessonTitle{Homilía sancti Gregórii Papæ}
\switchcolumn
\EN
\LessonTitle{Homily by Pope St Gregory the Great}
\switchcolumn*

\LA
\Source{Homil. 37 in Evangelia}
\switchcolumn
\EN
\Source{Homilia 37th on the Gospels.}
\switchcolumn*

\LA
\noindent Si considerémus fratres caríssimi, quæ et quanta sunt, quæ nobis promittúntur in cælis, viléscunt ánimo ómnia quæ habéntur in terris. Terréna namque substántia supérnæ felicitáti comparáta, pondus est, non subsídium. Temporális vita ætérnæ vitæ comparáta, mors est pótius dicénda quam vita. Ipse enim quotidiánus deféctus corruptiónis quid est áliud, quam quædam prolíxitas mortis? Quæ autem lingua dícere, vel quis intelléctus cápere súfficit, illa supérnæ civitátis quanta sint gáudia; Angelórum choris interésse, cum beatíssimis spirítibus glóriæ Conditóris assístere, præséntem Dei vultum cérnere, incircumscríptum lumen vidére, nullo mortis metu áffici, incorruptiónis perpétuæ múnere lætári?\par
\switchcolumn
\EN
\noindent Dearly beloved brethren, if we consider what and how great things are promised unto us in heaven, all things which are upon earth grow poor to our mind. For when this world's goods are reckoned against the gladness above, they are found to be a clog rather than an help. This present life being compared to life eternal, ought rather to be called death than life. For what is the daily failing of our corruption but, as it were, a creeping death But what tongue is there that can tell, or what understanding that can comprehend how great is the rejoicing in the city above, where they have part with the choirs of Angels, where they stand with the most blessed spirits before the glory of the Creator, where they see the face of God present, where they behold the Incomprehensible Light, where they have no fear of death, and where they rejoice eternally incorruptible\par
\switchcolumn*

\LA
\begin{Resp}\Rsign Coróna áurea super caput ejus, \Star{} Expréssa signo sanctitátis, glória honóris, et opus fortitúdinis.\end{Resp}
\begin{Resp}\Vsign Quóniam prævenísti eum in benedictiónibus dulcédinis, posuísti in cápite ejus corónam de lápide pretióso.\end{Resp}
\begin{Resp}\Rsign Expréssa signo sanctitátis, glória honóris, et opus fortitúdinis.\end{Resp}
\switchcolumn
\EN
\begin{Resp}\Rsign A crown of gold upon his head, \Star{} Wherein is engraved Holiness, an ornament of honour, a costly work.\end{Resp}
\begin{Resp}\Vsign For Thou hast prevented him with the blessings of sweetness, Thou hast set a crown of precious stones upon his head.\end{Resp}
\begin{Resp}\Rsign Wherein is engraved Holiness, an ornament of honour, a costly work.\end{Resp}
\switchcolumn*

\LA
\LessonHead{Lectio VIII}
\switchcolumn
\EN
\LessonHead{Lesson VIII}
\switchcolumn*

\LA
\noindent Sed ad hæc audíta inardéscit ánimus, jamque illic cupit assístere, ubi se sperat sine fine gaudére. Sed ad magna prǽmia perveníri non potest, nisi per magnos labóres. Unde et Paulus egrégius prædicátor dicit: Non coronábitur, nisi qui legítime certáverit. Deléctet ergo mentem magnitúdo præmiórum, sed non detérreat certámen labórum. Unde ad se veniéntibus Véritas dicit: Si quis venit ad me, et non odit patrem suum, et matrem, et uxórem, et fílios, et fratres, et soróres, adhuc autem et ánimam suam, non potest meus esse discípulus.\par
\switchcolumn
\EN
\noindent Then we hear these things our hearts burn within us and we long to be already there, where we hope to rejoice for ever. But we cannot attain unto great rewards, save through great labour. Therefore saith the excellent preacher Paul: “He is not crowned, except he strive lawfully.” (2 Tim. ii. 5.) The greatness of the reward doth delight our mind let not the throes of the struggle dishearten us. Therefore the Truth saith unto every one that cometh unto Him “If any man come to Me, and hate not his father and mother, and wife, and children, and brethren, and sisters, yea, and his own life also, he cannot be My disciple.”\par
\switchcolumn*

\LA
\begin{Resp}\Rsign Hic est vere Martyr, qui pro Christi nómine sánguinem suum fudit: \Star{} Qui minas júdicum non tímuit, nec terrénæ dignitátis glóriam quæsívit, sed ad cæléstia regna pervénit.\end{Resp}
\begin{Resp}\Vsign Justum dedúxit Dóminus per vias rectas, et osténdit illi regnum Dei.\end{Resp}
\begin{Resp}\Rsign Qui minas júdicum non tímuit, nec terrénæ dignitátis glóriam quæsívit, sed ad cæléstia regna pervénit.\end{Resp}
\begin{Resp}\Vsign Glória Patri, et Fílio, \Star{} et Spirítui Sancto.\end{Resp}
\begin{Resp}\Rsign Qui minas júdicum non tímuit, nec terrénæ dignitátis glóriam quæsívit, sed ad cæléstia regna pervénit.\end{Resp}
\switchcolumn
\EN
\begin{Resp}\Rsign This is a Martyr indeed, who shed his blood for Christ's Name's sake \Star{} Who feared not for the threats of judges, nor sought to be great with the glory of this world, but pressed on unto the kingdom of heaven.\end{Resp}
\begin{Resp}\Vsign The Lord guided the righteous in right paths, and showed him the kingdom of God.\end{Resp}
\begin{Resp}\Rsign Who feared not for the threats of judges, nor sought to be great with the glory of this world, but pressed on unto the kingdom of heaven.\end{Resp}
\begin{Resp}\Vsign Glory be to the Father, and to the Son, \Star{} and to the Holy Ghost.\end{Resp}
\begin{Resp}\Rsign Who feared not for the threats of judges, nor sought to be great with the glory of this world, but pressed on unto the kingdom of heaven.\end{Resp}
\switchcolumn*

\LA
\LessonHead{Lectio IX}
\switchcolumn
\EN
\LessonHead{Lesson IX}
\switchcolumn*

\LA
\noindent Sed percontári libet, quómodo paréntes et carnáliter propínquos præcípimur odísse, qui jubémur et inimícos dilígere? Et certe Véritas de uxóre dicit: Quod Deus conjúnxit, homo non séparet. Et Paulus ait: Viri, dilígite uxóres vestras, sicut et Christus Ecclésiam. Ecce, discípulus uxórem diligéndam prǽdicat, cum magíster dicat: Qui uxórem non odit, non potest meus esse discípulus. Numquid áliud judex núntiat, áliud præco clamat? An simul et odísse póssumus, et dilígere? Sed si vim præcépti perpéndimus, utrúmque ágere per discretiónem valémus: ut uxórem, et eos, qui nobis carnis cognatióne conjúncti sunt, et quos próximos nóvimus, diligámus; et quos adversários in via Dei pátimur, odiéndo et fugiéndo nesciámus.\par
\switchcolumn
\EN
\noindent Or it may be asked how we are commanded in one place to hate our parents, and them that are near us in the flesh, and in another place to love even our enemies. And, verily, the Truth hath said, as touching a wife “What God hath joined together, let not man put asunder.” (Matth. xix. 6.) And Paul saith “ Husbands, love your wives, even as Christ also loved the Church.” (Eph. v. 25.) Behold, the disciple commandeth a man to love his wife, and the Master saith “If any man hate not his wife, he cannot be My disciple.” Doth the judge, then, order one proclamation, and the crier make another? or can the man both love and hate? If we consider well the force of the commandment, we shall be able in wisdom to do both. Let us love wife, and kindred, and neighbour, as touching their nearness in the flesh but as touching the way of God, if they withstand us therein, let us not know them, but hate them and flee from them.\par
\switchcolumn*

\LA
\TeDeum{Te Deum}
\switchcolumn
\EN
\TeDeum{Te Deum}
\switchcolumn*

\end{paracol}

\Office{Commune unius Summi Pontificis et Martyris}{Common of a Single Martyr Pope}{}
\addcontentsline{toc}{subsection}{Commune unius Summi Pontificis et Martyris\enspace\textit{Common of a Single Martyr Pope}}
\begin{paracol}{2}
\LA
\LessonHead{Lectio VII}
\switchcolumn
\EN
\LessonHead{Lesson VII}
\switchcolumn*

\LA
\LessonTitle{Léctio sancti Evangélii secúndum Matthǽum}
\switchcolumn
\EN
\LessonTitle{From the Holy Gospel according to Matthew}
\switchcolumn*

\LA
\Cite{Matt 16:13-19}
\switchcolumn
\EN
\Cite{Matt 16:13-19}
\switchcolumn*

\LA
\noindent In illo témpore: Venit Jesus in partes Cæsaréæ Philíppi, et interrogábat discípulos suos, dicens: Quem dicunt hómines esse Fílium hóminis? Et réliqua.\par
\switchcolumn
\EN
\noindent At that time: When Jesus came into the coasts of Caesarea Philippi, he asked his disciples, saying, Whom do men say that I the Son of man am? And so on.\par
\switchcolumn*

\LA
\LessonTitle{Homilía sancti Leónis Papæ}
\switchcolumn
\EN
\LessonTitle{A Homily by St. Leo the Pope}
\switchcolumn*

\LA
\Source{Sermo 2 in anniversario assumpt. suæ, ante medium}
\switchcolumn
\EN
\Source{Sermo 2 in anniversario assumpt. suæ ante medium}
\switchcolumn*

\LA
\noindent Cum, sicut evangélica lectióne reserátum est, interrogásset Dóminus discípulos, quem ipsum (multis divérsa opinántibus) créderent; respondissétque beátus Petrus, dicens: Tu es Christus Fílius Dei vivi; Dóminus ait: Beátus es, Simon Bar-Jona, quia caro et sanguis non revelávit tibi, sed Pater meus, qui in cælis est: et ego dico tibi, quia tu es Petrus, et super hanc petram ædificábo Ecclésiam meam, et portæ ínferi non prævalébunt advérsus eam. Et tibi dabo claves regni cælórum: et quodcúmque ligáveris super terram, erit ligátum et in cælis: et quodcúmque sólveris super terram, erit solútum et in cælis. Manet ergo disposítio veritátis, et beátus Petrus, in accépta fortitúdine petræ persevérans, suscépta Ecclésiæ gubernácula non relíquit.\par
\switchcolumn
\EN
\noindent When the Lord, as we read in the Gospel, asked his disciples who did men, amid their diverse speculations, believe him the Son of Man to be, blessed Peter answered and said: Thou art the Christ, the Son of the living God. And the Lord answered and said unto him: Blessed art thou, Simon Bar-Jona: for flesh and blood hath not revealed it unto thee, but my Father, which is in heaven: and I say also unto thee: That thou art Peter, and upon this rock I will build my Church, and the gates of hell shall not prevail against it; and I will give unto thee the keys of the kingdom of heaven; and whatsoever thou shalt bind on earth shall be bound in heaven; and whatsoever thou shalt loose on earth shall be loosed in heaven. But the dispensation of truth perdures, and blessed Peter, persevering in the strength of the rock which he hath received, hath not relinquished the position he assumed at the helm of the Church.\par
\switchcolumn*

\LA
\begin{Resp}\Rsign Coróna áurea super caput ejus, \Star{} Expréssa signo sanctitátis, glória honóris, et opus fortitúdinis.\end{Resp}
\begin{Resp}\Vsign Quóniam prævenísti eum in benedictiónibus dulcédinis, posuísti in cápite ejus corónam de lápide pretióso.\end{Resp}
\begin{Resp}\Rsign Expréssa signo sanctitátis, glória honóris, et opus fortitúdinis.\end{Resp}
\switchcolumn
\EN
\begin{Resp}\Rsign A crown of gold upon his head, \Star{} Wherein is engraved Holiness, an ornament of honour, a costly work.\end{Resp}
\begin{Resp}\Vsign For Thou hast prevented him with the blessings of sweetness, Thou hast set a crown of precious stones upon his head.\end{Resp}
\begin{Resp}\Rsign Wherein is engraved Holiness, an ornament of honour, a costly work.\end{Resp}
\switchcolumn*

\LA
\LessonHead{Lectio VIII}
\switchcolumn
\EN
\LessonHead{Lesson VIII}
\switchcolumn*

\LA
\noindent In univérsa namque Ecclésia, Tu es Christus Fílius Dei vivi, quotídie Petrus dicit; et omnis lingua, quæ confitétur Dóminum, magistério hujus vocis imbúitur. Hæc fides diábolum vincit et captivórum ejus víncula dissólvit. Hæc érutos mundo, ínserit cælo, et portæ ínferi advérsus eam prævalére non possunt. Tanta enim divínitus soliditáte muníta est, ut eam neque hærética umquam corrúmpere právitas, nec pagána potúerit superáre perfídia. His ítaque modis, dilectíssimi, rationábile obséquio celebrétur hodiérna festívitas: ut in persóna humilitátis meæ ille intelligátur, ille honorétur, in quo et ómnium pastórum sollicitúdo, cum commendatárum sibi óvium custódia persevérat, et cujus étiam dígnitas in indígno heréde non déficit.\par
\switchcolumn
\EN
\noindent In the universal Church it is as if Peter were still saying every day: Thou art the Christ, the Son of the living God. For every tongue which confesseth the Lord is taught that confession by the teaching of Peter. This is the Faith that overcometh the devil and looseth the bonds of his prisoners. This is the Faith which maketh men free of the world and bringeth them to heaven, and the gates of hell are impotent to prevail against it. This is the rock which God hath fortified with such ramparts of salvation, that the contagion of heresy will never be able to infect it, nor idolatry and unbelief to overcome it. And therefore, dearly beloved, we celebrate today’s festival with reasonable obedience, that in my humble person he may be acknowledged and honoured who doth continue to care for all the shepherds as well as sheep entrusted unto him, and who doth lose none of his dignity even in an unworthy successor.\par
\switchcolumn*

\LA
\begin{Resp}\Rsign Hic est vere Martyr, qui pro Christi nómine sánguinem suum fudit: \Star{} Qui minas júdicum non tímuit, nec terrénæ dignitátis glóriam quæsívit, sed ad cæléstia regna pervénit.\end{Resp}
\begin{Resp}\Vsign Justum dedúxit Dóminus per vias rectas, et osténdit illi regnum Dei.\end{Resp}
\begin{Resp}\Rsign Qui minas júdicum non tímuit, nec terrénæ dignitátis glóriam quæsívit, sed ad cæléstia regna pervénit.\end{Resp}
\begin{Resp}\Vsign Glória Patri, et Fílio, \Star{} et Spirítui Sancto.\end{Resp}
\begin{Resp}\Rsign Qui minas júdicum non tímuit, nec terrénæ dignitátis glóriam quæsívit, sed ad cæléstia regna pervénit.\end{Resp}
\switchcolumn
\EN
\begin{Resp}\Rsign This is a Martyr indeed, who shed his blood for Christ's Name's sake \Star{} Who feared not for the threats of judges, nor sought to be great with the glory of this world, but pressed on unto the kingdom of heaven.\end{Resp}
\begin{Resp}\Vsign The Lord guided the righteous in right paths, and showed him the kingdom of God.\end{Resp}
\begin{Resp}\Rsign Who feared not for the threats of judges, nor sought to be great with the glory of this world, but pressed on unto the kingdom of heaven.\end{Resp}
\begin{Resp}\Vsign Glory be to the Father, and to the Son, \Star{} and to the Holy Ghost.\end{Resp}
\begin{Resp}\Rsign Who feared not for the threats of judges, nor sought to be great with the glory of this world, but pressed on unto the kingdom of heaven.\end{Resp}
\switchcolumn*

\LA
\LessonHead{Lectio IX}
\switchcolumn
\EN
\LessonHead{Lesson IX}
\switchcolumn*

\LA
\noindent Cum ergo cohortatiónes nostras áuribus vestræ sanctitátis adhibémus, ipsum vobis, cujus vice fúngimur, loqui crédite: quia et illíus vos afféctu monémus, et non áliud vobis, quam quod dócuit, prædicámus; obsecrántes, ut succíncti lumbos mentis vestræ, castam et sóbriam vitam in Dei timóre ducátis. Coróna mea, sicut Apóstolus ait, et gáudium vos estis, si fides vestra, quæ ab inítio Evangélii in univérso mundo prædicáta est, in dilectióne et sanctitáte permánserit. Nam licet omnem Ecclésiam, quæ in toto est orbe terrárum, cunctis opórteat florére virtútibus; vos tamen præcípue inter céteros pópulos decet méritis pietátis excéllere, quos in ipsa apostólicæ petræ arce fundátos, et Dóminus noster Jesus Christus cum ómnibus redémit, et beátus Apóstolus Petrus præ ómnibus erudívit.\par
\switchcolumn
\EN
\noindent When, therefore, we address our exhortations to your godly ears, believe ye that ye are hearing him speak whose office we are discharging. Yea, it is with his love for you that we warn you. And we preach unto you no other thing than that which he taught, entreating you as did he: Gird up the loins of your mind; be sober; be ye holy in all manner of living; pass the time of your sojourning here in the fear of God. My disciples, dearly beloved, ye are to me as the disciples of the Apostle Paul were to him, namely: My crown and joy; if so be that your faith, abide, still in all lowliness and holiness, like unto the first times of the Gospel. For although the whole Church, which is in all the world, should indeed abound in all the virtues, it becometh especially you among all others to excel in acts of piety, founded as ye be on the very citadel of the Apostolic Rock ye who have not only been redeemed with the rest of men by our Lord Jesus Christ, but who have been instructed by the blessed Apostle Peter far beyond all others.\par
\switchcolumn*

\LA
\TeDeum{Te Deum}
\switchcolumn
\EN
\TeDeum{Te Deum}
\switchcolumn*

\end{paracol}

\Office{Commune Plurimorum Martyrum Pontificum}{Commune Plurimorum Martyrum Pontificum}{}
\addcontentsline{toc}{subsection}{Commune Plurimorum Martyrum Pontificum\enspace\textit{Commune Plurimorum Martyrum Pontificum}}
\begin{paracol}{2}
\LA
\LessonHead{Lectio I}
\switchcolumn
\EN
\LessonHead{Lesson I}
\switchcolumn*

\LA
\LessonTitle{De Epístola beáti Pauli Apóstoli ad Romános}
\switchcolumn
\EN
\LessonTitle{Lesson from the letter of St Paul the Apostle to the Romans}
\switchcolumn*

\LA
\Cite{Rom 8:12-19}
\switchcolumn
\EN
\Cite{Rom 8:12-19}
\switchcolumn*

\LA
\noindent Fratres: Debitóres sumus non carni, ut secúndum carnem vivámus. Si enim secúndum carnem vixéritis, moriémini: si autem spíritu facta carnis mortificavéritis, vivétis. Quicúmque enim spíritu Dei agúntur, ii sunt fílii Dei. Non enim accepístis spíritum servitútis íterum in timóre, sed accepístis spíritum adoptiónis filiórum, in quo clamámus: Abba (Pater). Ipse enim Spíritus testimónium reddit spirítui nostro quod sumus fílii Dei. Si autem fílii, et herédes: herédes, quidem Dei, coherédes autem Christi: si tamen compátimur ut et conglorificémur. Exístimo enim quod non sunt condígnæ passiónes hujus témporis ad futúram glóriam, quæ revelábitur in nobis. Nam exspectátio creatúræ revelatiónem filiórum Dei exspéctat.\par
\switchcolumn
\EN
\noindent Therefore, brethren, we are debtors, not to the flesh, to live according to the flesh. For if you live according to the flesh, you shall die: but if by the Spirit you mortify the deeds of the flesh, you shall live. For whosoever are led by the Spirit of God, they are the sons of God. For you have not received the spirit of bondage again in fear; but you have received the spirit of adoption of sons, whereby we cry: Abba, Father. For the Spirit himself giveth testimony to our spirit, that we are the sons of God. And if sons, heirs also; heirs indeed of God, and joint heirs with Christ: yet so, if we suffer with him, that we may be also glorified with him. For I reckon that the sufferings of this time are not worthy to be compared with the glory to come, that shall be revealed in us. For the expectation of the creature waiteth for the revelation of the sons of God\par
\switchcolumn*

\LA
\begin{Resp}\Rsign Abstérget Deus omnem lácrimam ab óculis Sanctórum: et jam non erit ámplius neque luctus, neque clamor, sed nec ullus dolor: \Star{} Quóniam prióra transiérunt.\end{Resp}
\begin{Resp}\Vsign Non esúrient, neque sítient ámplius, neque cadet super illos sol, neque ullus æstus.\end{Resp}
\begin{Resp}\Rsign Quóniam prióra transiérunt.\end{Resp}
\switchcolumn
\EN
\begin{Resp}\Rsign God shall wipe away all tears from the eyes of His Saints, and there shall be no more sorrow, there shall be no more death, neither sorrow, nor crying, neither shall there be any more pain; \Star{} For the former things are passed away.\end{Resp}
\begin{Resp}\Vsign They shall hunger no more, neither thirst any more, neither shall the sun light on them, nor any heat.\end{Resp}
\begin{Resp}\Rsign For the former things are passed away.\end{Resp}
\switchcolumn*

\LA
\LessonHead{Lectio II}
\switchcolumn
\EN
\LessonHead{Lesson II}
\switchcolumn*

\LA
\Cite{Rom 8:28-34}
\switchcolumn
\EN
\Cite{Rom 8:28-34}
\switchcolumn*

\LA
\noindent Scimus autem quóniam diligéntibus Deum ómnia cooperántur in bonum iis qui secúndum propósitum vocáti sunt sancti. Nam quos præscívit et prædestinávit confórmes fíeri imáginis Fílii sui ut sit ipse primogénitus in multis frátribus. Quos autem prædestinávit, hos et vocávit: et quos vocávit, hos et justificávit: quos autem justificávit, illos et glorificávit. Quid ergo dicémus ad hæc? Si Deus pro nobis, quis contra nos? Qui étiam próprio Fílio suo non pepércit, sed pro nobis ómnibus trádidit illum, quómodo non étiam cum illo ómnia nobis donávit? Quis accusábit advérsus eléctos Dei? Deus qui justíficat, quis est qui condémnet? Christus Jesus, qui mórtuus est, immo qui et resurréxit, qui est ad déxteram Dei, qui étiam interpéllat pro nobis.\par
\switchcolumn
\EN
\noindent And we know that to them that love God, all things work together unto good, to such as, according to his purpose, are called to be saints. For whom he foreknew, he also predestinated to be made conformable to the image of his Son; that he might be the firstborn amongst many brethren. And whom he predestinated, them he also called. And whom he called, them he also justified. And whom he justified, them he also glorified. What shall we then say to these things? If God be for us, who is against us? He that spared not even his own Son, but delivered him up for us all, how hath he not also, with him, given us all things? Who shall accuse against the elect of God? God that justifieth. Who is he that shall condemn? Christ Jesus that died, yea that is risen also again; who is at the right hand of God, who also maketh intercession for us.\par
\switchcolumn*

\LA
\begin{Resp}\Rsign Viri sancti gloriósum sánguinem fudérunt pro Dómino, amavérunt Christum in vita sua, imitáti sunt eum in morte sua: \Star{} Et ídeo corónas triumpháles meruérunt.\end{Resp}
\begin{Resp}\Vsign Unus spíritus, et una fides erat in eis.\end{Resp}
\begin{Resp}\Rsign Et ídeo corónas triumpháles meruérunt.\end{Resp}
\switchcolumn
\EN
\begin{Resp}\Rsign These men are holy, who have gloriously shed their blood for the Lord's sake, yea, who loved Christ in their lives, and were made like unto Him in their flesh, \Star{} And therefore they have earned crowns of victory.\end{Resp}
\begin{Resp}\Vsign One spirit, and one faith was in them.\end{Resp}
\begin{Resp}\Rsign And therefore they have earned crowns of victory.\end{Resp}
\switchcolumn*

\LA
\LessonHead{Lectio III}
\switchcolumn
\EN
\LessonHead{Lesson III}
\switchcolumn*

\LA
\Cite{Rom 8:35-39}
\switchcolumn
\EN
\Cite{Rom 8:35-39}
\switchcolumn*

\LA
\noindent Quis ergo nos separábit a caritáte Christi? tribulátio? an angústia? an fames? an núditas? an perículum? an persecútio? an gládius? (sicut scriptum est: Quia propter te mortificámur tota die: æstimáti sumus sicut oves occisiónis.) Sed in his ómnibus superámus propter eum qui diléxit nos. Certus sum enim quia neque mors, neque vita, neque Ángeli, neque Principátus, neque Virtútes, neque instántia, neque futúra, neque fortitúdo, neque altitúdo, neque profúndum, neque creatúra ália póterit nos separáre a caritáte Dei, quæ est in Christo Jesu Dómino nostro.\par
\switchcolumn
\EN
\noindent Who then shall separate us from the love of Christ? Shall tribulation? or distress? or famine? or nakedness? or danger? or persecution? or the sword? As it is written: For thy sake we are put to death all the day long. We are accounted as sheep for the slaughter. But in all these things we overcome, because of him that hath loved us. For I am sure that neither death, nor life, nor angels, nor principalities, nor powers, nor things present, nor things to come, nor might, Nor height, nor depth, nor any other creature, shall be able to separate us from the love of God, which is in Christ Jesus our Lord.\par
\switchcolumn*

\LA
\begin{Resp}\Rsign Tradidérunt córpora sua propter Deum ad supplícia: \Star{} Et meruérunt habére corónas perpétuas.\end{Resp}
\begin{Resp}\Vsign Isti sunt qui venérunt ex magna tribulatióne, et lavérunt stolas suas in sánguine Agni.\end{Resp}
\begin{Resp}\Rsign Et meruérunt habére corónas perpétuas.\end{Resp}
\begin{Resp}\Vsign Glória Patri, et Fílio, \Star{} et Spirítui Sancto.\end{Resp}
\begin{Resp}\Rsign Et meruérunt habére corónas perpétuas.\end{Resp}
\switchcolumn
\EN
\begin{Resp}\Rsign They gave their bodies for God’s sake to death \Star{} And gained the everlasting crown.\end{Resp}
\begin{Resp}\Vsign These are they which came out of great tribulation, and have washed their robes in the Blood of the Lamb.\end{Resp}
\begin{Resp}\Rsign And gained the everlasting crown.\end{Resp}
\begin{Resp}\Vsign Glory be to the Father, and to the Son, \Star{} and to the Holy Ghost.\end{Resp}
\begin{Resp}\Rsign And gained the everlasting crown.\end{Resp}
\switchcolumn*

\LA
\LessonHead{Lectio VII}
\switchcolumn
\EN
\LessonHead{Lesson VII}
\switchcolumn*

\LA
\LessonTitle{Léctio sancti Evangélii secúndum Lucam}
\switchcolumn
\EN
\LessonTitle{From the Holy Gospel according to Luke}
\switchcolumn*

\LA
\Cite{Luc 21:9-19}
\switchcolumn
\EN
\Cite{Luke 21:9-19}
\switchcolumn*

\LA
\noindent In illo témpore: Dixit Jesus discípulis suis: Cum audiéritis prǽlia, et seditiónes, nolíte terréri: opórtet primum hæc fíeri, sed nondum statim finis. Et réliqua.\par
\switchcolumn
\EN
\noindent At that time, Jesus said unto His disciples: When ye shall hear of wars and commotions, be not terrified for these things must first come to pass; but the end is not by and by. And so on.\par
\switchcolumn*

\LA
\LessonTitle{Homilía sancti Gregórii Papæ}
\switchcolumn
\EN
\LessonTitle{Homily by Pope St Gregory (the Great.)}
\switchcolumn*

\LA
\Source{Homilia 35 in Evangel.}
\switchcolumn
\EN
\Source{35th on the Gospels.}
\switchcolumn*

\LA
\noindent Dóminus ac Redémptor noster peritúri mundi præcurréntia mala denúntiat, ut eo minus pertúrbent veniéntia, quo fúerint præscíta. Minus enim jácula fériunt, quæ prævidéntur: et nos tolerabílius mundi mala suscípimus, si contra hæc per præsciéntiæ clípeum munímur. Ecce enim dicit: Cum audiéritis prǽlia et seditiónes, nolíte terréri: opórtet enim primum hæc fíeri, sed nondum statim finis. Pensánda sunt verba Redemptóris nostri, per quæ nos áliud intérius, áliud extérius passúros esse denúntiat; bella quippe ad hostes pértinent, seditiónes ad cives. Ut ergo nos índicet intérius exteriúsque turbári, áliud nos fatétur ab hóstibus, áliud a frátribus pérpeti.\par
\switchcolumn
\EN
\noindent Our Lord and Redeemer willeth us to know what shall be the signs that the end of the world is at hand, to the end that ye may be the less terrified, when that cometh whereof ye have already had warning. Darts strike less which are seen coming: and the plagues of the earth will be to us more bearable, if we are harnessed against them with the shield of foreknowledge. Behold, how He saith When ye shall hear of wars and commotions be not terrified for these things must first come to pass; but the end is not by and by. It behoveth us to ponder these words of our Redeemer, wherein He warneth us of suffering, from without, and from within. Wars are the work of a foreign enemy, commotions of the citizens. Therefore, that He may let us know that we shall be troubled from within and from without, He showeth that our wrestling shall be in part against strangers, and in part against our brethren.\par
\switchcolumn*

\LA
\begin{Resp}\Rsign Propter testaméntum Dómini, et leges patérnas, Sancti Dei perstitérunt in amóre fraternitátis: \Star{} Quia unus fuit semper spíritus in eis, et una fides.\end{Resp}
\begin{Resp}\Vsign Ecce quam bonum et quam jucúndum habitáre fratres in unum.\end{Resp}
\begin{Resp}\Rsign Quia unus fuit semper spíritus in eis, et una fides.\end{Resp}
\switchcolumn
\EN
\begin{Resp}\Rsign Because of the covenant of the Lord, and the laws of their fathers, the Saints of God abode in brotherly love, \Star{} For one spirit and one faith was ever in them.\end{Resp}
\begin{Resp}\Vsign Behold how good and how pleasant it is for brethren to dwell together in unity.\end{Resp}
\begin{Resp}\Rsign For one spirit and one faith was ever in them.\end{Resp}
\switchcolumn*

\LA
\LessonHead{Lectio VIII}
\switchcolumn
\EN
\LessonHead{Lesson VIII}
\switchcolumn*

\LA
\noindent Sed his malis præveniéntibus, quia non statim finis sequátur, adjúngit: Surget gens contra gentem, et regnum advérsus regnum: et terræmótus magni erunt per loca, et pestiléntiæ, et fames, terrorésque de cælo: et signa magna erunt. Última tribulátio multis tribulatiónibus prævenítur: et per crebra mala, quæ prævéniunt, indicántur mala perpétua, quæ subsequéntur. Et ídeo post bella et seditiónes non statim finis: quia multa debent mala præcúrrere, ut malum váleant sine fine nuntiáre.\par
\switchcolumn
\EN
\noindent When these woes come, the end is not by and by. And He saith further Nation shall rise against nation, and kingdom against kingdom; and great earthquakes shall be in diverse places, and pestilences, and famines, and fearful sights and great signs shall there be from heaven. Before the last tribulation cometh, shall come many other tribulations: and, by the many woes which shall come first, shall be foreshadowed the everlasting woe which shall come in the end. And therefore, after wars and commotions, the end is not yet by and by many woes must come first, to give warning of the woe that hath no end.\par
\switchcolumn*

\LA
\begin{Resp}\Rsign Sancti mei, qui in carne pósiti, certámen habuístis: \Star{} Mercédem labóris ego reddam vobis.\end{Resp}
\begin{Resp}\Vsign Veníte benedícti Patris mei, percípite regnum.\end{Resp}
\begin{Resp}\Rsign Mercédem labóris ego reddam vobis.\end{Resp}
\begin{Resp}\Vsign Glória Patri, et Fílio, \Star{} et Spirítui Sancto.\end{Resp}
\begin{Resp}\Rsign Mercédem labóris ego reddam vobis.\end{Resp}
\switchcolumn
\EN
\begin{Resp}\Rsign O ye My Saints, who, being in the flesh, didst have striving \Star{} I will render unto you a reward of your labours.\end{Resp}
\begin{Resp}\Vsign Come, ye blessed of My Father, inherit the kingdom\end{Resp}
\begin{Resp}\Rsign I will render unto you a reward of your labours.\end{Resp}
\begin{Resp}\Vsign Glory be to the Father, and to the Son, \Star{} and to the Holy Ghost.\end{Resp}
\begin{Resp}\Rsign I will render unto you a reward of your labours.\end{Resp}
\switchcolumn*

\LA
\LessonHead{Lectio IX}
\switchcolumn
\EN
\LessonHead{Lesson IX}
\switchcolumn*

\LA
\noindent Sed cum tot signa perturbatiónis dicta sint, opórtet, ut eórum consideratiónem bréviter per síngula perstringámus: quia necésse est, ut ália e cælo, ália e terra, ália ab eleméntis, ália ab homínibus patiámur. Ait enim: Surget gens contra gentem, ecce perturbátio hóminum: erunt terræmótus magni per loca, ecce respéctus iræ désuper: erunt pestiléntiæ ecce inæquálitas córporum: erit fames, ecce sterílitas terræ: terrorésque de cælo et tempestátes, ecce inæquálitas áëris. Quia ergo ómnia consummánda sunt, ante consummatiónem ómnia perturbántur: et qui in cunctis delíquimus, in cunctis ferímur: ut impleátur quod dícitur, Et pugnábit pro eo orbis terrárum contra insensátos.\par
\switchcolumn
\EN
\noindent But, forasmuch as the signs and troubles whereof the Lord speaketh are so manifold, we must needs shortly consider each for, of necessity, we must suffer some things from heaven, some from the earth, some from the powers of nature, and some from men. For where He saith Nation shall rise against nation He speaketh concerning the troubling of men where great earthquakes shall be in diverse places concerning wrath from above: where: and pestilences concerning the frailty of the body where and famines concerning the barrenness of the earth: where fearful signs from heaven, and tempests concerning commotions of the air. As, then, all things shall have an end, so, before the end, shall all things be troubled and we who have sinned and come short in all things, shall in all things be afflicted, that it may be fulfilled that is written: and the world shall fight with Him against the unwise. (Wisd v. 21.)\par
\switchcolumn*

\LA
\TeDeum{Te Deum}
\switchcolumn
\EN
\TeDeum{Te Deum}
\switchcolumn*

\end{paracol}

\Office{Commune unius Confessoris Pontificis}{Common of a Confessor Bishop}{}
\addcontentsline{toc}{subsection}{Commune unius Confessoris Pontificis\enspace\textit{Common of a Confessor Bishop}}
\begin{paracol}{2}
\LA
\LessonHead{Lectio VII}
\switchcolumn
\EN
\LessonHead{Lesson VII}
\switchcolumn*

\LA
\LessonTitle{Léctio sancti Evangélii secúndum Matthǽum}
\switchcolumn
\EN
\LessonTitle{From the Holy Gospel according to Matthew}
\switchcolumn*

\LA
\Cite{Matt 25:14-23}
\switchcolumn
\EN
\Cite{Matt 25:14-23}
\switchcolumn*

\LA
\noindent In illo témpore: Dixit Jesus discípulis suis parábolam hanc: Homo péregre proficíscens, vocávit servos suos, et trádidit illis bona sua. Et réliqua.\par
\switchcolumn
\EN
\noindent At that time Jesus spake unto His disciples this parable A man, travelling into a far country, called his own servants, and delivered unto them his goods. And so on.\par
\switchcolumn*

\LA
\LessonTitle{Homilía sancti Gregórii Papæ}
\switchcolumn
\EN
\LessonTitle{Homily by Pope St Gregory (the Great.)}
\switchcolumn*

\LA
\Source{Homilia 9 in Evangelia}
\switchcolumn
\EN
\Source{9th on the Gospels.}
\switchcolumn*

\LA
\noindent Léctio sancti Evangélii, fratres caríssimi, solícite consideráre nos ádmonet, ne nos, qui plus céteris in hoc mundo accepísse áliquid cérnimur, ab Auctóre mundi grávius inde judicémur. Cum enim augéntur dona, ratiónes étiam crescunt donórum. Tanto ergo esse humílior atque ad serviéndum Deo prómptior quisque debet ex múnere, quanto se obligatiórem esse cónspicit in reddénda ratióne. Ecce homo, qui péregre proficíscitur, servos suos vocat, eisque ad negótium talénta partítur. Post multum vero témporis positúrus ratiónem revértitur: bene operántes pro apportáto lucro remúnerat, servum vero a bono ópere torpéntem damnat.\par
\switchcolumn
\EN
\noindent Dearly beloved brethren, this Lesson from the Holy Gospel moveth us to take good heed lest we, who are seen in this world to have received more than others, should thereby bring ourselves into greater condemnation from the Maker of this world. To whom much is given, of the same is much required. Therefore, let him that receiveth much, strive to be all the more lowly, and all the more ready to do God service, for his very gifts’ sake, knowing that he will be obliged to give account thereof. Behold, a man, travelling into a far country, calleth his own servants, and delivereth unto them talents, to the end that they may trade therewith. After a long time, the lord of those servants cometh, and reckoneth with them, and to them that have done well He rendereth a reward of their labours, but that servant which was careless of his master’s work He condemneth.\par
\switchcolumn*

\LA
\begin{Resp}\Rsign Amávit eum Dóminus, et ornávit eum: stolam glóriæ índuit eum, \Star{} Et ad portas paradísi coronávit eum.\end{Resp}
\begin{Resp}\Vsign Induit eum Dóminus lorícam fídei, et ornávit eum.\end{Resp}
\begin{Resp}\Rsign Et ad portas paradísi coronávit eum.\end{Resp}
\switchcolumn
\EN
\begin{Resp}\Rsign The Lord loved him and beautified him He clothed him with a robe of glory \Star{} And crowned him at the gates of Paradise.\end{Resp}
\begin{Resp}\Vsign The Lord hath put on him the breast-plate of faith, and hath adorned him.\end{Resp}
\begin{Resp}\Rsign And crowned him at the gates of Paradise.\end{Resp}
\switchcolumn*

\LA
\LessonHead{Lectio VIII}
\switchcolumn
\EN
\LessonHead{Lesson VIII}
\switchcolumn*

\LA
\noindent Quis ítaque iste homo est qui péregre proficíscitur, nisi Redémptor noster, qui, in ea carne quam assúmpserat, ábiit in cælum? Carnis enim locus próprius terra est; quæ quasi ad peregrína dúcitur, dum per Redemptórem nostrum in cælo collocátur. Sed homo iste, péregre proficíscens, servis suis bona sua trádidit; quia fidélibus suis spirituália dona concéssit. Et uni quidem quinque talénta, álii duo, álii vero commísit unum. Quinque étenim sunt córporis sensus, vidélicet: visus, audítus, gustus, odorátus et tactus. Quinque ergo taléntis, donum quinque sénsuum, id est, exteriórum sciéntia, exprímitur. Duóbus vero, intelléctus et operátio designátur. Uníus autem talénti nómine, intelléctus tantúmmodo designátur.\par
\switchcolumn
\EN
\noindent What other, then, is that man travelling into a far country but our Redeemer, Who is gone up from us into heaven in that Flesh Which He had taken into Himself? For the earth is the home of the Flesh, Which travelleth into a far country when our Redeemer giveth It a place in heaven. But that man travelling into a far country delivered unto his servants his goods and so doth our Redeemer give spiritual gifts unto His faithful people. “And unto one he gave five talents, to another two, and to another one.” There are five bodily senses that is, sight, hearing, taste, smell, and touch. By the five talents therefore are signified the five senses, that is, outward knowledge. By the two, wit and work. And by the figure of the one talent, understanding, which is alone.\par
\switchcolumn*

\LA
\begin{Resp}\Rsign Sint lumbi vestri præcíncti, et lucérnæ ardéntes in mánibus vestris: \Star{} Et vos símiles homínibus exspectántibus dóminum suum, quando revertátur a núptiis.\end{Resp}
\begin{Resp}\Vsign Vigiláte ergo, quia nescítis qua hora Dóminus vester ventúrus sit.\end{Resp}
\begin{Resp}\Rsign Et vos símiles homínibus exspectántibus dóminum suum, quando revertátur a núptiis.\end{Resp}
\begin{Resp}\Vsign Glória Patri, et Fílio, \Star{} et Spirítui Sancto.\end{Resp}
\begin{Resp}\Rsign Et vos símiles homínibus exspectántibus dóminum suum, quando revertátur a núptiis.\end{Resp}
\switchcolumn
\EN
\begin{Resp}\Rsign Let your loins be girded about, and your lights burning; \Star{} And ye yourselves like unto men that wait for their lord, when he will return from the wedding.\end{Resp}
\begin{Resp}\Vsign Watch therefore, for ye know not what hour your Lord doth come.\end{Resp}
\begin{Resp}\Rsign And ye yourselves like unto men that wait for their lord, when he will return from the wedding.\end{Resp}
\begin{Resp}\Vsign Glory be to the Father, and to the Son, \Star{} and to the Holy Ghost.\end{Resp}
\begin{Resp}\Rsign And ye yourselves like unto men that wait for their lord, when he will return from the wedding.\end{Resp}
\switchcolumn*

\LA
\LessonHead{Lectio IX}
\switchcolumn
\EN
\LessonHead{Lesson IX}
\switchcolumn*

\LA
\noindent Sed is, qui quinque talénta accéperat, ália quinque lucrátus est: quia sunt nonnúlli, qui, etsi intérna ac mýstica penetráre nésciunt, pro intentióne tamen supérnæ pátriæ docent recta quos possunt; de ipsis exterióribus, quæ accepérunt, duplum taléntum portant; dumque se a carnis petulántia et a terrenárum rerum ámbitu atque a visibílium voluptáte custódiunt, ab his étiam álios admonéndo compéscunt. Et sunt nonnúlli, qui, quasi duóbus taléntis ditáti, intelléctum atque operatiónem percípiunt, subtília de intérnis intélligunt, mira in exterióribus operántur; cumque et intelligéndo et operándo áliis prǽdicant, quasi duplicátum de negótio lucrum repórtant.\par
\switchcolumn
\EN
\noindent And so he that had received five talents, gained other five talents for some there be who, while yet they are not able to go on unto things inward and mystic, do yet so desire our Fatherland which is above, that they teach well all whom they can, and of those very outward things which they have received make gain double. These are they which keep themselves clean from the unruly motions of the flesh, and from the lust of the world, and from the delight of things which are seen, and, by their preaching, keep other men also clean from all these things. And some there are who receive, as their two talents, the power to think and the power to work. These are they which inwardly understand dark things, and outwardly work wonders. And these, since they preach unto others, both through their understanding and their works, gain, as it were, double, for the talents which they have received.\par
\switchcolumn*

\LA
\TeDeum{Te Deum}
\switchcolumn
\EN
\TeDeum{Te Deum}
\switchcolumn*

\end{paracol}

\Office{Commune Doctoris Pontificis}{Common of a Doctor Bishop}{}
\addcontentsline{toc}{subsection}{Commune Doctoris Pontificis\enspace\textit{Common of a Doctor Bishop}}
\begin{paracol}{2}
\LA
\LessonHead{Lectio I}
\switchcolumn
\EN
\LessonHead{Lesson I}
\switchcolumn*

\LA
\LessonTitle{De libro Ecclesiástici}
\switchcolumn
\EN
\LessonTitle{Lesson from the book of Ecclesiasticus}
\switchcolumn*

\LA
\Cite{Sir 39:1-5}
\switchcolumn
\EN
\Cite{Sir 39:1-5}
\switchcolumn*

\LA
\noindent Sapiéntiam ómnium antiquórum exquíret sápiens, et in prophétis vacábit. Narratiónem virórum nominatórum conservábit, et in versútias parabolárum simul introíbit. Occúlta proverbiórum exquíret, et in abscónditis parabolárum conversábitur. In médio magnatórum ministrábit, et in conspéctu prǽsidis apparébit. In terram alienigenárum géntium pertránsiet; bona enim et mala in homínibus tentábit.\par
\switchcolumn
\EN
\noindent The wise men will seek out the wisdom of all the ancients, and will be occupied in the prophets. He will keep the sayings of renowned men, and will enter withal into the subtilties of parables. He will search out the hidden meanings of proverbs, and will be conversant in the secrets of parables. He shall serve among great men, and: appear before the governor. He shall pass into strange countries: for he shall try good and evil among men.\par
\switchcolumn*

\LA
\begin{Resp}\Rsign Euge, serve bone et fidélis, quia in pauca fuísti fidélis, supra multa te constítuam: \Star{} Intra in gáudium Dómini tui.\end{Resp}
\begin{Resp}\Vsign Dómine, quinque talénta tradidísti mihi, ecce ália quinque superlucrátus sum.\end{Resp}
\begin{Resp}\Rsign Intra in gáudium Dómini tui.\end{Resp}
\switchcolumn
\EN
\begin{Resp}\Rsign Well done, thou good and faithful servant, thou hast been faithful over a few things. I will make thee ruler over many things: \Star{} Enter thou into the joy of thy Lord.\end{Resp}
\begin{Resp}\Vsign Lord, thou deliveredst unto me five talents; behold, I have gained beside them five talents more.\end{Resp}
\begin{Resp}\Rsign Enter thou into the joy of thy Lord.\end{Resp}
\switchcolumn*

\LA
\LessonHead{Lectio II}
\switchcolumn
\EN
\LessonHead{Lesson II}
\switchcolumn*

\LA
\Cite{Sir 39:6-10}
\switchcolumn
\EN
\Cite{Sir 39:6-10}
\switchcolumn*

\LA
\noindent Cor suum tradet ad vigilándum dilúculo ad Dóminum, qui fecit illum, et in conspéctu Altíssimi deprecábitur. Apériet os suum in oratióne, et pro delíctis suis deprecábitur. Si enim Dóminus magnus volúerit, spíritu intellegéntiæ replébit illum: et ipse tamquam imbres mittet elóquia sapiéntiæ suæ, et in oratióne confitébitur Dómino: et ipse díriget consílium ejus, et disciplínam, et in abscónditis suis consiliábitur.\par
\switchcolumn
\EN
\noindent He will give his heart to resort early to the Lord that made him, and he will pray in the sight of the most High. He will open his mouth in prayer, and will make supplication for his sins. For if it shall please the great Lord, he will fill him with the spirit of understanding: And he will pour forth the words of his wisdom as showers, and in his prayer he will confess to the Lord. And he shall direct his counsel, and his knowledge, and in his secrets shall he meditate.\par
\switchcolumn*

\LA
\begin{Resp}\Rsign Ecce sacérdos magnus, qui in diébus suis plácuit Deo: \Star{} Ideo jurejurándo fecit illum Dóminus créscere in plebem suam.\end{Resp}
\begin{Resp}\Vsign Benedictiónem ómnium géntium dedit illi, et testaméntum suum confirmávit super caput ejus.\end{Resp}
\begin{Resp}\Rsign Ideo jurejurándo fecit illum Dóminus créscere in plebem suam.\end{Resp}
\switchcolumn
\EN
\begin{Resp}\Rsign Behold an high priest, who in his days pleased God \Star{} Therefore the Lord assured him by an oath that He would multiply his seed among His people.\end{Resp}
\begin{Resp}\Vsign He hath made him a blessing unto all nations, and hath established His covenant upon his head.\end{Resp}
\begin{Resp}\Rsign Therefore the Lord assured him by an oath that He would multiply his seed among His people\end{Resp}
\switchcolumn*

\LA
\LessonHead{Lectio III}
\switchcolumn
\EN
\LessonHead{Lesson III}
\switchcolumn*

\LA
\Cite{Sir 39:11-14}
\switchcolumn
\EN
\Cite{Sir 39:11-14}
\switchcolumn*

\LA
\noindent Ipse palam fáciet disciplínam doctrínæ suæ, et in lege testaménti Dómini gloriábitur. Collaudábunt multi sapiéntiam ejus, et usque in sǽculum non delébitur. Non recédet memória ejus, et nomen ejus requirétur a generatióne in generatiónem. Sapiéntiam ejus enarrábunt gentes, et laudem ejus enuntiábit Ecclésia.\par
\switchcolumn
\EN
\noindent He shall show forth the discipline he hath learned, and shall glory in the law of the covenant of the Lord. Many shall praise his wisdom, and it shall never be forgotten. The memory of him shall not depart away, and his name shall be in request from generation to generation. Nations shall declare his wisdom, and the church shall show forth his praise.\par
\switchcolumn*

\LA
\begin{Resp}\Rsign Jurávit Dóminus, et non pœnitébit eum: \Star{} Tu es sacérdos in ætérnum secúndum órdinem Melchísedech.\end{Resp}
\begin{Resp}\Vsign Dixit Dóminus Dómino meo, Sede a dextris meis.\end{Resp}
\begin{Resp}\Rsign Tu es sacérdos in ætérnum secúndum órdinem Melchísedech.\end{Resp}
\begin{Resp}\Vsign Glória Patri, et Fílio, \Star{} et Spirítui Sancto.\end{Resp}
\begin{Resp}\Rsign Tu es sacérdos in ætérnum secúndum órdinem Melchísedech.\end{Resp}
\switchcolumn
\EN
\begin{Resp}\Rsign The Lord hath sworn and will not repent \Star{} Thou art a Priest for ever after the order of Melchisedek.\end{Resp}
\begin{Resp}\Vsign The Lord said unto my Lord Sit Thou at My right hand. R.* Thou art a Priest for ever after the order of Melchisedek.\end{Resp}
\begin{Resp}\Vsign Glory be to the Father, and to the Son, \Star{} and to the Holy Ghost.\end{Resp}
\begin{Resp}\Rsign Thou art a Priest for ever after the order of Melchisedek.\end{Resp}
\switchcolumn*

\LA
\LessonHead{Lectio VII}
\switchcolumn
\EN
\LessonHead{Lesson VII}
\switchcolumn*

\LA
\LessonTitle{Léctio sancti Evangélii secúndum Matthǽum}
\switchcolumn
\EN
\LessonTitle{From the Holy Gospel according to Matthew}
\switchcolumn*

\LA
\Cite{Matt 5:13-19}
\switchcolumn
\EN
\Cite{Matt 5:13-19}
\switchcolumn*

\LA
\noindent In illo témpore: Dixit Jesus discípulis suis: Vos estis sal terræ. Quod si sal evanúerit, in quo saliétur? Et réliqua.\par
\switchcolumn
\EN
\noindent At that time: Jesus said unto his disciples: Ye are the salt of the earth: But if the salt have lost his savour, wherewith shall it be salted? And so on, and that which followeth.\par
\switchcolumn*

\LA
\LessonTitle{Homilía sancti Augustíni Epíscopi}
\switchcolumn
\EN
\LessonTitle{A Homily by St. Augustine the Bishop}
\switchcolumn*

\LA
\Source{Lib. 1 de Sermone Domini in monte, cap. 6}
\switchcolumn
\EN
\Source{Lib. 1 de Sermóne Dómini in monte, cap. 6}
\switchcolumn*

\LA
\noindent Osténdit Dóminus fátuos esse judicándos, qui temporálium bonórum vel cópiam sectántes vel inópiam metuéntes, amíttunt ætérna, quæ nec dari possunt ab homínibus nec auférri. Itaque, si sal infatuátum fúerit, in quo saliétur? Id est, si vos, per quos condiéndi sunt quodámmodo pópuli, metu persecutiónum temporálium amiséritis regna cælórum; qui erunt hómines, per quos a vobis error auferátur, cum vos elégerit Deus, per quos errórem áuferat ceterórum?\par
\switchcolumn
\EN
\noindent The Lord would have us understand how that men do lose their power of savouring others with righteousness when they are willing to place their eternal welfare in jeopardy for the sake of any témporal advantage, like as attainment of ease or luxury, or escape from suffering or toil. For that which is eternal, unlike things of this world, can neither be bestowed by men, nor by them taken away. Hence, when he asketh: If the salt have lost his savour, wherewith shall it be salted? he would have us understand the question to be: If ye, by whom mankind is preserved from corruption, be willing to lose the kingdom of heaven so as to escape trials or persecutions in this world, who is there to preserve you from corruption, seeing ye are they that God hath chosen to preserve all others from corruption?\par
\switchcolumn*

\LA
\begin{Resp}\Rsign Amávit eum Dóminus, et ornávit eum: stolam glóriæ índuit eum, \Star{} Et ad portas paradísi coronávit eum.\end{Resp}
\begin{Resp}\Vsign Induit eum Dóminus lorícam fídei, et ornávit eum.\end{Resp}
\begin{Resp}\Rsign Et ad portas paradísi coronávit eum.\end{Resp}
\switchcolumn
\EN
\begin{Resp}\Rsign The Lord loved him and beautified him He clothed him with a robe of glory \Star{} And crowned him at the gates of Paradise.\end{Resp}
\begin{Resp}\Vsign The Lord hath put on him the breast-plate of faith, and hath adorned him.\end{Resp}
\begin{Resp}\Rsign And crowned him at the gates of Paradise.\end{Resp}
\switchcolumn*

\LA
\LessonHead{Lectio VIII}
\switchcolumn
\EN
\LessonHead{Lesson VIII}
\switchcolumn*

\LA
\noindent Ergo ad níhilum valet sal infatuátum, nisi ut mittátur foras et calcétur ab homínibus. Non ítaque calcátur ab homínibus qui pátitur persecutiónem; sed qui, persecutiónem timéndo, infatuátur. Calcári enim non potest nisi inférior; sed inférior non est, qui, quamvis córpore multa in terra sustíneat, corde tamen fixus in cælo est.\par
\switchcolumn
\EN
\noindent Those that should be the salt of the earth, but have lost their savour, are thenceforth good for nothing, but to be cast out, and to be trodden under foot of men. But no one that suffereth persecution is truly said to be trodden under foot of men. Rather, that one is truly trodden under foot of men who through fear of persecution hath lost the savour of righteousness. For no one can be trodden upon, unless he be beneath him which treadeth upon him. And certainly no one who hath his heart in heaven, no matter how grievously he doth suffer in his body on earth, is rightly said to be beneath anyone who misuseth him.\par
\switchcolumn*

\LA
\begin{Resp}\Rsign In médio Ecclésiæ apéruit os ejus, \Star{} Et implévit eum Dóminus spíritu sapiéntiæ et intelléctus.\end{Resp}
\begin{Resp}\Vsign Jucunditátem et exsultatiónem thesaurizávit super eum.\end{Resp}
\begin{Resp}\Rsign Et implévit eum Dóminus spíritu sapiéntiæ et intelléctus.\end{Resp}
\begin{Resp}\Vsign Glória Patri, et Fílio, \Star{} et Spirítui Sancto.\end{Resp}
\begin{Resp}\Rsign Et implévit eum Dóminus spíritu sapiéntiæ et intelléctus.\end{Resp}
\switchcolumn
\EN
\begin{Resp}\Rsign In the midst of the congregation did the Lord open his mouth. \Star{} And filled him with the spirit of wisdom and understanding.\end{Resp}
\begin{Resp}\Vsign He made him rich with joy and gladness.\end{Resp}
\begin{Resp}\Rsign And filled him with the spirit of wisdom and understanding.\end{Resp}
\begin{Resp}\Vsign Glory be to the Father, and to the Son, \Star{} and to the Holy Ghost.\end{Resp}
\begin{Resp}\Rsign And filled him with the spirit of wisdom and understanding.\end{Resp}
\switchcolumn*

\LA
\LessonHead{Lectio IX}
\switchcolumn
\EN
\LessonHead{Lesson IX}
\switchcolumn*

\LA
\noindent Vos estis lumen mundi. Quómodo dixit supérius sal terræ, sic nunc dicit lumen mundi. Nam, neque supérius ista terra accipiénda est, quam pédibus corpóreis calcámus; sed hómines, qui in terra hábitant, vel étiam peccatóres, quorum condiéndis et exstinguéndis putóribus apostólicum salem Dóminus misit. Et hic mundum non cælum et terram, sed hómines qui sunt in mundo vel díligunt mundum, opórtet intélligi; quibus illuminándis Apóstoli missi sunt. Non potest cívitas abscóndi super montem pósita; id est, fundáta super insígnem magnámque justítiam, quam signíficat étiam ipse mons, in quo dísputat Dóminus.\par
\switchcolumn
\EN
\noindent Ye are the light of the world. And we are to understand the word World in the same sense as the word Earth when he spoke above of the salt of the earth, that is, not that earth whereupon we walk with our bodily feet, but the men which dwell upon the earth; in other words, sinners, for the sweetening and correction of whose corruption, the Lord hath sent his Apostles, as it were, as so much salt. And so by the world we are to understand, not the heaven and the earth, but the men who are in the world and love the world, for the enlightening of whom the Apostles have been sent. A city that is set on a hill cannot be hid: that is, what is founded upon the heights of righteousness, whereof the mountain upon which the Lord gave this discourse was itself a figure, is magnificent in the eyes of all men.\par
\switchcolumn*

\LA
\TeDeum{Te Deum}
\switchcolumn
\EN
\TeDeum{Te Deum}
\switchcolumn*

\end{paracol}

\Office{Commune Summorum Pontificum Confessorum}{Common of Popes}{}
\addcontentsline{toc}{subsection}{Commune Summorum Pontificum Confessorum\enspace\textit{Common of Popes}}
\begin{paracol}{2}
\LA
\LessonHead{Lectio VII}
\switchcolumn
\EN
\LessonHead{Lesson VII}
\switchcolumn*

\LA
\LessonTitle{Léctio sancti Evangélii secúndum Matthǽum}
\switchcolumn
\EN
\LessonTitle{From the Holy Gospel according to Matthew}
\switchcolumn*

\LA
\Cite{Matt 16:13-19}
\switchcolumn
\EN
\Cite{Matt 16:13-19}
\switchcolumn*

\LA
\noindent In illo témpore: Venit Jesus in partes Cæsaréæ Philíppi, et interrogábat discípulos suos, dicens: Quem dicunt hómines esse Fílium hóminis? Et réliqua.\par
\switchcolumn
\EN
\noindent At that time: When Jesus came into the coasts of Caesarea Philippi, he asked his disciples, saying, Whom do men say that I the Son of man am? And so on.\par
\switchcolumn*

\LA
\LessonTitle{Homilía sancti Leónis Papæ}
\switchcolumn
\EN
\LessonTitle{A Homily by St. Leo the Pope}
\switchcolumn*

\LA
\Source{Sermo 2 in anniversario assumpt. suæ, ante medium}
\switchcolumn
\EN
\Source{Sermo 2 in anniversario assumpt. suæ ante medium}
\switchcolumn*

\LA
\noindent Cum, sicut evangélica lectióne reserátum est, interrogásset Dóminus discípulos, quem ipsum (multis divérsa opinántibus) créderent; respondissétque beátus Petrus, dicens: Tu es Christus Fílius Dei vivi; Dóminus ait: Beátus es, Simon Bar-Jona, quia caro et sanguis non revelávit tibi, sed Pater meus, qui in cælis est: et ego dico tibi, quia tu es Petrus, et super hanc petram ædificábo Ecclésiam meam, et portæ ínferi non prævalébunt advérsus eam. Et tibi dabo claves regni cælórum: et quodcúmque ligáveris super terram, erit ligátum et in cælis: et quodcúmque sólveris super terram, erit solútum et in cælis. Manet ergo disposítio veritátis, et beátus Petrus, in accépta fortitúdine petræ persevérans, suscépta Ecclésiæ gubernácula non relíquit.\par
\switchcolumn
\EN
\noindent When the Lord, as we read in the Gospel, asked his disciples who did men, amid their diverse speculations, believe him the Son of Man to be, blessed Peter answered and said: Thou art the Christ, the Son of the living God. And the Lord answered and said unto him: Blessed art thou, Simon Bar-Jona: for flesh and blood hath not revealed it unto thee, but my Father, which is in heaven: and I say also unto thee: That thou art Peter, and upon this rock I will build my Church, and the gates of hell shall not prevail against it; and I will give unto thee the keys of the kingdom of heaven; and whatsoever thou shalt bind on earth shall be bound in heaven; and whatsoever thou shalt loose on earth shall be loosed in heaven. But the dispensation of truth perdures, and blessed Peter, persevering in the strength of the rock which he hath received, hath not relinquished the position he assumed at the helm of the Church.\par
\switchcolumn*

\LA
\begin{Resp}\Rsign Amávit eum Dóminus, et ornávit eum: stolam glóriæ índuit eum, \Star{} Et ad portas paradísi coronávit eum.\end{Resp}
\begin{Resp}\Vsign Induit eum Dóminus lorícam fídei, et ornávit eum.\end{Resp}
\begin{Resp}\Rsign Et ad portas paradísi coronávit eum.\end{Resp}
\switchcolumn
\EN
\begin{Resp}\Rsign The Lord loved him and beautified him He clothed him with a robe of glory \Star{} And crowned him at the gates of Paradise.\end{Resp}
\begin{Resp}\Vsign The Lord hath put on him the breast-plate of faith, and hath adorned him.\end{Resp}
\begin{Resp}\Rsign And crowned him at the gates of Paradise.\end{Resp}
\switchcolumn*

\LA
\LessonHead{Lectio VIII}
\switchcolumn
\EN
\LessonHead{Lesson VIII}
\switchcolumn*

\LA
\noindent In univérsa namque Ecclésia, Tu es Christus Fílius Dei vivi, quotídie Petrus dicit; et omnis lingua, quæ confitétur Dóminum, magistério hujus vocis imbúitur. Hæc fides diábolum vincit et captivórum ejus víncula dissólvit. Hæc érutos mundo, ínserit cælo, et portæ ínferi advérsus eam prævalére non possunt. Tanta enim divínitus soliditáte muníta est, ut eam neque hærética umquam corrúmpere právitas, nec pagána potúerit superáre perfídia. His ítaque modis, dilectíssimi, rationábile obséquio celebrétur hodiérna festívitas: ut in persóna humilitátis meæ ille intelligátur, ille honorétur, in quo et ómnium pastórum sollicitúdo, cum commendatárum sibi óvium custódia persevérat, et cujus étiam dígnitas in indígno heréde non déficit.\par
\switchcolumn
\EN
\noindent In the universal Church it is as if Peter were still saying every day: Thou art the Christ, the Son of the living God. For every tongue which confesseth the Lord is taught that confession by the teaching of Peter. This is the Faith that overcometh the devil and looseth the bonds of his prisoners. This is the Faith which maketh men free of the world and bringeth them to heaven, and the gates of hell are impotent to prevail against it. This is the rock which God hath fortified with such ramparts of salvation, that the contagion of heresy will never be able to infect it, nor idolatry and unbelief to overcome it. And therefore, dearly beloved, we celebrate today’s festival with reasonable obedience, that in my humble person he may be acknowledged and honoured who doth continue to care for all the shepherds as well as sheep entrusted unto him, and who doth lose none of his dignity even in an unworthy successor.\par
\switchcolumn*

\LA
\begin{Resp}\Rsign Sint lumbi vestri præcíncti, et lucérnæ ardéntes in mánibus vestris: \Star{} Et vos símiles homínibus exspectántibus dóminum suum, quando revertátur a núptiis.\end{Resp}
\begin{Resp}\Vsign Vigiláte ergo, quia nescítis qua hora Dóminus vester ventúrus sit.\end{Resp}
\begin{Resp}\Rsign Et vos símiles homínibus exspectántibus dóminum suum, quando revertátur a núptiis.\end{Resp}
\begin{Resp}\Vsign Glória Patri, et Fílio, \Star{} et Spirítui Sancto.\end{Resp}
\begin{Resp}\Rsign Et vos símiles homínibus exspectántibus dóminum suum, quando revertátur a núptiis.\end{Resp}
\switchcolumn
\EN
\begin{Resp}\Rsign Let your loins be girded about, and your lights burning; \Star{} And ye yourselves like unto men that wait for their lord, when he will return from the wedding.\end{Resp}
\begin{Resp}\Vsign Watch therefore, for ye know not what hour your Lord doth come.\end{Resp}
\begin{Resp}\Rsign And ye yourselves like unto men that wait for their lord, when he will return from the wedding.\end{Resp}
\begin{Resp}\Vsign Glory be to the Father, and to the Son, \Star{} and to the Holy Ghost.\end{Resp}
\begin{Resp}\Rsign And ye yourselves like unto men that wait for their lord, when he will return from the wedding.\end{Resp}
\switchcolumn*

\LA
\LessonHead{Lectio IX}
\switchcolumn
\EN
\LessonHead{Lesson IX}
\switchcolumn*

\LA
\noindent Cum ergo cohortatiónes nostras áuribus vestræ sanctitátis adhibémus, ipsum vobis, cujus vice fúngimur, loqui crédite: quia et illíus vos afféctu monémus, et non áliud vobis, quam quod dócuit, prædicámus; obsecrántes, ut succíncti lumbos mentis vestræ, castam et sóbriam vitam in Dei timóre ducátis. Coróna mea, sicut Apóstolus ait, et gáudium vos estis, si fides vestra, quæ ab inítio Evangélii in univérso mundo prædicáta est, in dilectióne et sanctitáte permánserit. Nam licet omnem Ecclésiam, quæ in toto est orbe terrárum, cunctis opórteat florére virtútibus; vos tamen præcípue inter céteros pópulos decet méritis pietátis excéllere, quos in ipsa apostólicæ petræ arce fundátos, et Dóminus noster Jesus Christus cum ómnibus redémit, et beátus Apóstolus Petrus præ ómnibus erudívit.\par
\switchcolumn
\EN
\noindent When, therefore, we address our exhortations to your godly ears, believe ye that ye are hearing him speak whose office we are discharging. Yea, it is with his love for you that we warn you. And we preach unto you no other thing than that which he taught, entreating you as did he: Gird up the loins of your mind; be sober; be ye holy in all manner of living; pass the time of your sojourning here in the fear of God. My disciples, dearly beloved, ye are to me as the disciples of the Apostle Paul were to him, namely: My crown and joy; if so be that your faith, abide, still in all lowliness and holiness, like unto the first times of the Gospel. For although the whole Church, which is in all the world, should indeed abound in all the virtues, it becometh especially you among all others to excel in acts of piety, founded as ye be on the very citadel of the Apostolic Rock ye who have not only been redeemed with the rest of men by our Lord Jesus Christ, but who have been instructed by the blessed Apostle Peter far beyond all others.\par
\switchcolumn*

\LA
\TeDeum{Te Deum}
\switchcolumn
\EN
\TeDeum{Te Deum}
\switchcolumn*

\end{paracol}

\Office{Commune Confessoris non pontificis}{Common of a Confessor not a Bishop}{}
\addcontentsline{toc}{subsection}{Commune Confessoris non pontificis\enspace\textit{Common of a Confessor not a Bishop}}
\begin{paracol}{2}
\LA
\LessonHead{Lectio I}
\switchcolumn
\EN
\LessonHead{Lesson I}
\switchcolumn*

\LA
\LessonTitle{De libro Ecclesiástici}
\switchcolumn
\EN
\LessonTitle{Lesson from the book of Ecclesiasticus}
\switchcolumn*

\LA
\Cite{Sir 31:8-11}
\switchcolumn
\EN
\Cite{Sir 31:8-11}
\switchcolumn*

\LA
\noindent Beátus vir, qui invéntus est sine mácula, et qui post aurum non ábiit, nec sperávit in pecúnia et thesáuris. Quis est hic? et laudábimus eum: fecit enim mirabília in vita sua. Qui probátus est in illo, et perféctus est, erit illi glória ætérna: Qui pótuit tránsgredi, et non est transgréssus; fácere mala, et non fecit: ídeo stabilíta sunt bona illíus in Dómino, et eleemósynas illíus enarrábit omnis ecclésia sanctórum.\par
\switchcolumn
\EN
\noindent Blessed is the rich man that is found without blemish: and that hath not gone after gold, nor put his trust in money nor in treasures. Who is he, and we will praise him? for he hath done wonderful things in his life. Who hath been tried thereby, and made perfect, he shall have glory everlasting. He that could have transgressed, and hath not transgressed: and could do evil things, and hath not done them: Therefore are his goods established in the Lord, and all the church of the saints shall declare his alms\par
\switchcolumn*

\LA
\begin{Resp}\Rsign Euge, serve bone et fidélis, quia in pauca fuísti fidélis, supra multa te constítuam: \Star{} Intra in gáudium Dómini tui.\end{Resp}
\begin{Resp}\Vsign Dómine, quinque talénta tradidísti mihi, ecce ália quinque superlucrátus sum.\end{Resp}
\begin{Resp}\Rsign Intra in gáudium Dómini tui.\end{Resp}
\switchcolumn
\EN
\begin{Resp}\Rsign Well done, thou good and faithful servant, thou hast been faithful over a few things. I will make thee ruler over many things: \Star{} Enter thou into the joy of thy Lord.\end{Resp}
\begin{Resp}\Vsign Lord, thou deliveredst unto me five talents; behold, I have gained beside them five talents more.\end{Resp}
\begin{Resp}\Rsign Enter thou into the joy of thy Lord.\end{Resp}
\switchcolumn*

\LA
\LessonHead{Lectio II}
\switchcolumn
\EN
\LessonHead{Lesson II}
\switchcolumn*

\LA
\Cite{Sir 32:18-20; 32:28; 33:1-3}
\switchcolumn
\EN
\Cite{Sir 32:18-20; 32:28; 33:1-3}
\switchcolumn*

\LA
\noindent Qui timet Dóminum excípiet doctrínam ejus: et qui vigiláverint ad illum invénient benedictiónem. Qui quærit legem, replébitur ab ea: et qui insidióse agit, scandalizábitur in ea. Qui timent Dóminum invénient judícium justum, et justítias quasi lumen accéndent. Qui credit Deo atténdit mandátis: et qui confídit in illo non minorábitur. Timénti Dóminum non occúrrent mala: sed in tentatióne Deus illum conservábit, et liberábit a malis. Sápiens non odit mandáta et justítias, et non illidétur quasi in procélla navis. Homo sensátus credit legi Dei, et lex illi fidélis.\par
\switchcolumn
\EN
\noindent He that feareth the Lord, will receive his discipline: and they that will seek him early, shall find a blessing. He that seeketh the law, shall be filled with it: and he that dealeth deceitfully, shall meet with a stumblingblock therein. They that fear the Lord, shall find just judgment, and shall kindle justice as a light. He that believeth God, taketh heed to the commandments: and he that trusteth in him, shall fare never the worse. No evils shall happen to him that feareth the Lord, but in temptation God will keep him, and deliver him from evils. A wise man hateth not the commandments and justices, and he shall not be dashed in pieces as a ship in a storm. A man of understanding is faithful to the law of God, and the law is faithful to him.\par
\switchcolumn*

\LA
\begin{Resp}\Rsign Justus germinábit sicut lílium: \Star{} Et florébit in ætérnum ante Dóminum.\end{Resp}
\begin{Resp}\Vsign Plantátus in domo Dómini, in átriis domus Dei nostri.\end{Resp}
\begin{Resp}\Rsign Et florébit in ætérnum ante Dóminum.\end{Resp}
\switchcolumn
\EN
\begin{Resp}\Rsign The righteous shall grow as the lily; \Star{} Yea, he shall flourish in the presence of the Lord for ever.\end{Resp}
\begin{Resp}\Vsign Those that be planted in the house of the Lord, shall flourish in the courts of the house of our God.\end{Resp}
\begin{Resp}\Rsign Yea, he shall flourish in the presence of the Lord for ever.\end{Resp}
\switchcolumn*

\LA
\LessonHead{Lectio III}
\switchcolumn
\EN
\LessonHead{Lesson III}
\switchcolumn*

\LA
\Cite{Sir 34:14-20}
\switchcolumn
\EN
\Cite{Sir 34:14-20}
\switchcolumn*

\LA
\noindent Spíritus timéntium Deum quǽritur, et in respéctu illíus benedicétur. Spes enim illórum in salvántem illos, et óculi Dei in diligéntes se. Qui timet Dóminum nihil trepidábit: et non pavébit, quóniam ipse est spes ejus. Timéntis Dóminum, beáta est ánima ejus. Ad quem réspicit, et quis est fortitúdo ejus? Óculi Dómini super timéntes eum: protéctor poténtiæ, firmaméntum virtútis, tégimen ardóris, et umbráculum meridiáni: deprecátio offensiónis, et adjutórium casus: exáltans ánimam, et illúminans óculos, dans sanitátem, et vitam, et benedictiónem.\par
\switchcolumn
\EN
\noindent The spirit of those that fear God; is sought after, and by his regard shall be blessed. For their hope is on him that saveth them, and the eyes of God are upon them that love him. He that feareth the Lord shall tremble at nothing, and shall not be afraid for he is his hope. The soul of him that feareth the Lord is blessed. To whom doth he look, and who in his strength? The eyes of the Lord are upon them that fear him, he is their powerful protector, and strong stay, a defence from the heat, and a cover from the sun at noon, A preservation from stumbling, and a help from falling; he raiseth up the soul, and enlighteneth the eyes, and giveth health, and life, and blessing.\par
\switchcolumn*

\LA
\begin{Resp}\Rsign Iste cognóvit justítiam, et vidit mirabília magna, et exorávit Altíssimum: \Star{} Et invéntus est in número Sanctórum.\end{Resp}
\begin{Resp}\Vsign Iste est qui contémpsit vitam mundi, et pervénit ad cæléstia regna.\end{Resp}
\begin{Resp}\Rsign Et invéntus est in número Sanctórum.\end{Resp}
\begin{Resp}\Vsign Glória Patri, et Fílio, \Star{} et Spirítui Sancto.\end{Resp}
\begin{Resp}\Rsign Et invéntus est in número Sanctórum.\end{Resp}
\switchcolumn
\EN
\begin{Resp}\Rsign This is he which knew righteousness, and saw great wonders, and made his prayer unto the Most High; \Star{} And he is numbered among the Saints.\end{Resp}
\begin{Resp}\Vsign This is he which loved not his life in this world, and is come unto an everlasting kingdom.\end{Resp}
\begin{Resp}\Rsign And he is numbered among the Saints.\end{Resp}
\begin{Resp}\Vsign Glory be to the Father, and to the Son, \Star{} and to the Holy Ghost.\end{Resp}
\begin{Resp}\Rsign And he is numbered among the Saints.\end{Resp}
\switchcolumn*

\LA
\LessonHead{Lectio VII}
\switchcolumn
\EN
\LessonHead{Lesson VII}
\switchcolumn*

\LA
\LessonTitle{Léctio sancti Evangélii secúndum Lucam}
\switchcolumn
\EN
\LessonTitle{From the Holy Gospel according to Luke}
\switchcolumn*

\LA
\Cite{Luc 12:35-40}
\switchcolumn
\EN
\Cite{Luke 12:35-40}
\switchcolumn*

\LA
\noindent In illo témpore: Dixit Jesus discípulis suis: Sint lumbi vestri præcincti, et lucernæ ardentes in mánibus vestris. Et réliqua.\par
\switchcolumn
\EN
\noindent At that time, Jesus said unto His disciples: Let your loins be girded about, and your lights burning. And so on.\par
\switchcolumn*

\LA
\LessonTitle{Homilía sancti Gregórii Papæ}
\switchcolumn
\EN
\LessonTitle{Homily by Pope St. Gregory (the Great.)}
\switchcolumn*

\LA
\Source{Homilia 13 in Evang.}
\switchcolumn
\EN
\Source{13th on the Gospels.}
\switchcolumn*

\LA
\noindent Sancti Evangélii, fratres caríssimi, apérta vobis est léctio recitata. Sed ne alíquibus ipsa ejus planíties alta fortásse videátur, eam sub brevitáte transcúrrimus, quátenus ejus exposítio ita nesciéntibus fiat cógnita, ut tamen sciéntibus non sit onerósa. Dóminus dicit: Sint lumbi vestri præcíncti. Lumbos enim præcíngimus, cum carnis luxúriam per continéntiam coarctámus. Sed quia minus est, mala non ágere, nisi étiam quisque stúdeat, et bonis opéribus insudáre, prótinus ádditur: Et lucérnæ ardéntes in mánibus vestris. Lucérnas quippe ardéntes in mánibus tenémus, cum per bona ópera próximis nostris lucis exémpla monstrámus. De quibus profécto opéribus Dóminus dicit: Lúceat lux vestra coram homínibus, ut vídeant ópera vestra bona, et gloríficent Patrem vestrum qui in cælis est.\par
\switchcolumn
\EN
\noindent Dearly beloved brethren, the words of the Holy Gospel, which have just been read, lie open before you, and, lest their very plainness should make them seem to some to be hard, we will go through them with such shortness as that neither may they which understand not remain unenlightened, nor they which understand be wearied. The Lord saith Let your loins be girded about. Now, we gird our loins about, when by continency we master the lustful inclination of the flesh. But, forasmuch as it sufficeth not for a man to abstain from evil deeds, if he strive not to join thereto the earnest doing of good works, it is immediately added And your lights burning. Our lights burn when, by good works, we give bright example to our neighbour; concerning which works the Lord saith Let your light so shine before men, that they may see your good works, and glorify your Father Which is in heaven. (Matth. v. 16.)\par
\switchcolumn*

\LA
\begin{Resp}\Rsign Iste est qui ante Deum magnas virtútes operátus est, et de omni corde suo laudávit Dóminum: \Star{} Ipse intercédat pro peccátis ómnium populórum.\end{Resp}
\begin{Resp}\Vsign Ecce homo sine queréla, verus Dei cultor, ábstinens se ab omni ópere malo, et pérmanens in innocéntia sua.\end{Resp}
\begin{Resp}\Rsign Ipse intercédat pro peccátis ómnium populórum.\end{Resp}
\switchcolumn
\EN
\begin{Resp}\Rsign This is he which wrought great wonders before God, and praised the Lord with all his heart. \Star{} May he pray for all people, that their sins may be forgiven unto them.\end{Resp}
\begin{Resp}\Vsign Behold a man without blame, a worshipper of God in truth, keeping himself clean from every evil work, and abiding still in his innocency.\end{Resp}
\begin{Resp}\Rsign May he pray for all people, that their sins may be forgiven unto them.\end{Resp}
\switchcolumn*

\LA
\LessonHead{Lectio VIII}
\switchcolumn
\EN
\LessonHead{Lesson VIII}
\switchcolumn*

\LA
\noindent Duo autem sunt, quæ jubéntur, et lumbos restríngere, et lucérnas tenére: ut et mundítia sit castitátis in córpore, et lumen veritátis in operatióne. Redemptóri étenim nostro unum sine áltero placére nequáquam potest: si aut is qui bona agit, adhuc luxúriæ inquinaménta non déserit: aut is qui castitáte præéminet, necdum se per bona ópera exércet. Nec cástitas ergo magna est sine bono ópere, nec opus bonum est áliquod sine castitáte. Sed et si utrúmque ágitur, restat, ut quisquis ille est, spe ad supérnam pátriam tendat, et nequáquam se a vítiis pro mundi hujus honestáte contíneat.\par
\switchcolumn
\EN
\noindent Here, then, are two commandments, to gird our loins about, and to keep our lights burning the cleanness of purity in our body, and the light of the truth in our works. Whoso hath the one and not the other, pleaseth not thereby our Redeemer; that is, he pleaseth Him not which doth good works, but bridleth not himself from the pollutions of lust, neither he which is eminent in chastity, but exerciseth not himself in good works. Neither is chastity a great thing without good works, nor good works anything without chastity. And if any man do both, it remaineth that he must look by hope toward our Fatherland above, and not have for his reason wherethrough he turneth himself away from vice, the love of honour in this present world.\par
\switchcolumn*

\LA
\begin{Resp}\Rsign Sint lumbi vestri præcíncti, et lucérnæ ardéntes in mánibus vestris: \Star{} Et vos símiles homínibus exspectántibus dóminum suum, quando revertátur a núptiis.\end{Resp}
\begin{Resp}\Vsign Vigiláte ergo, quia nescítis qua hora Dóminus vester ventúrus sit.\end{Resp}
\begin{Resp}\Rsign Et vos símiles homínibus exspectántibus dóminum suum, quando revertátur a núptiis.\end{Resp}
\begin{Resp}\Vsign Glória Patri, et Fílio, \Star{} et Spirítui Sancto.\end{Resp}
\begin{Resp}\Rsign Et vos símiles homínibus exspectántibus dóminum suum, quando revertátur a núptiis.\end{Resp}
\switchcolumn
\EN
\begin{Resp}\Rsign Let your loins be girded about, and your lights burning; \Star{} And ye yourselves like unto men that wait for their lord, when he will return from the wedding.\end{Resp}
\begin{Resp}\Vsign Watch therefore, for ye know not what hour your Lord doth come.\end{Resp}
\begin{Resp}\Rsign And ye yourselves like unto men that wait for their lord, when he will return from the wedding.\end{Resp}
\begin{Resp}\Vsign Glory be to the Father, and to the Son, \Star{} and to the Holy Ghost.\end{Resp}
\begin{Resp}\Rsign And ye yourselves like unto men that wait for their lord, when he will return from the wedding.\end{Resp}
\switchcolumn*

\LA
\LessonHead{Lectio I (in 2 loco)}
\switchcolumn
\EN
\LessonHead{Lesson I (set 2)}
\switchcolumn*

\LA
\LessonTitle{De libro Sapiéntiæ}
\switchcolumn
\EN
\LessonTitle{Lesson from the book of Wisdom}
\switchcolumn*

\LA
\Cite{Sap 4:7-14}
\switchcolumn
\EN
\Cite{Wis 4:7-14}
\switchcolumn*

\LA
\noindent Justus si morte præoccupátus fúerit, in refrigério erit; Senéctus enim venerábilis est non diutúrna, neque annórum número computáta: cani autem sunt sensus hóminis, et ætas senectútis vita immaculáta. Placens Deo factus est diléctus, et vivens inter peccatóres translátus est. Raptus est, ne malítia mutáret intelléctum ejus, aut ne fíctio decíperet ánimam illíus. Fascinátio enim nugacitátis obscúrat bona, et inconstántia concupiscéntiæ transvértit sensum sine malítia. Consummátus in brevi, explévit témpora multa; plácita enim erat Deo ánima illíus: propter hoc properávit edúcere illum de médio iniquitátum.\par
\switchcolumn
\EN
\noindent But the just man, if he be prevented with death, shall be in rest. For venerable old age is not that of long time, nor counted by the number of years: but the understanding of a man is grey hairs. And a spotless life is old age. He pleased God and was beloved, and living among sinners he was translated. He was taken away lest wickedness should alter his understanding, or deceit beguile his soul. For the bewitching of vanity obscureth good things, and the wandering of concupiscence overturneth the innocent mind. Being made perfect in a short space, he fulfilled a long time: For his soul pleased God: therefore he hastened to bring him out of the midst of iniquities:\par
\switchcolumn*

\LA
\begin{Resp}\Rsign Euge, serve bone et fidélis, quia in pauca fuísti fidélis, supra multa te constítuam: \Star{} Intra in gáudium Dómini tui.\end{Resp}
\begin{Resp}\Vsign Dómine, quinque talénta tradidísti mihi, ecce ália quinque superlucrátus sum.\end{Resp}
\begin{Resp}\Rsign Intra in gáudium Dómini tui.\end{Resp}
\switchcolumn
\EN
\begin{Resp}\Rsign Well done, thou good and faithful servant, thou hast been faithful over a few things. I will make thee ruler over many things: \Star{} Enter thou into the joy of thy Lord.\end{Resp}
\begin{Resp}\Vsign Lord, thou deliveredst unto me five talents; behold, I have gained beside them five talents more.\end{Resp}
\begin{Resp}\Rsign Enter thou into the joy of thy Lord.\end{Resp}
\switchcolumn*

\LA
\LessonHead{Lectio II (in 2 loco)}
\switchcolumn
\EN
\LessonHead{Lesson II (set 2)}
\switchcolumn*

\LA
\Cite{Sap 4:14-19}
\switchcolumn
\EN
\Cite{Wis 4:14-19}
\switchcolumn*

\LA
\noindent Pópuli autem vidéntes, et non intelligéntes, nec ponéntes in præcórdiis tália, quóniam grátia Dei et misericórdia est in sanctos ejus, et respéctus in eléctos illíus. Condémnat autem justus mórtuus vivos ímpios, et juvéntus celérius consummáta longam vitam injústi. Vidébunt enim finem sapiéntis, et non intélligent quid cogitáverit de illo Deus, et quare muníerit illum Dóminus. Vidébunt, et contémnent eum; illos autem Dóminus irridébit. Et erunt post hæc decidéntes sine honóre, et in contumélia inter mórtuos in perpétuum: quóniam disrúmpet illos inflátos sine voce, et commovébit illos a fundaméntis, et usque ad suprémum desolabúntur.\par
\switchcolumn
\EN
\noindent but the people see this, and understand not, nor lay up such things in their hearts: That the grace of God, and his mercy is with his saints, and that he hath respect to his chosen. But the just that is dead, condemneth the wicked that are living, and youth soon ended, the long life of the unjust. For they shall see the end of the wise man, and shall not understand what God hath designed for him, and why the Lord hath set him in safety. They shall see him, and shall despise him: but the Lord shall laugh them to scorn. And they shall fall after this without honour, and be a reproach among the dead for ever: for he shall burst them puffed up and speechless, and shall shake them from the foundations, and they shall be utterly laid waste:\par
\switchcolumn*

\LA
\begin{Resp}\Rsign Justus germinábit sicut lílium: \Star{} Et florébit in ætérnum ante Dóminum.\end{Resp}
\begin{Resp}\Vsign Plantátus in domo Dómini, in átriis domus Dei nostri.\end{Resp}
\begin{Resp}\Rsign Et florébit in ætérnum ante Dóminum.\end{Resp}
\switchcolumn
\EN
\begin{Resp}\Rsign The righteous shall grow as the lily; \Star{} Yea, he shall flourish in the presence of the Lord for ever.\end{Resp}
\begin{Resp}\Vsign Those that be planted in the house of the Lord, shall flourish in the courts of the house of our God.\end{Resp}
\begin{Resp}\Rsign Yea, he shall flourish in the presence of the Lord for ever.\end{Resp}
\switchcolumn*

\LA
\LessonHead{Lectio III (in 2 loco)}
\switchcolumn
\EN
\LessonHead{Lesson III (set 2)}
\switchcolumn*

\LA
\Cite{Sap 4:19-20; 5:1-5}
\switchcolumn
\EN
\Cite{Wis 4:19-20; 5:1-5}
\switchcolumn*

\LA
\noindent Et erunt geméntes, et memória illórum períbit. Vénient in cogitatióne peccatórum suórum tímidi, et tradúcent illos ex advérso iniquitátes ipsórum. Tunc stabunt justi in magna constántia advérsus eos qui se angustiavérunt, et qui abstulérunt labóres eórum. Vidéntes turbabúntur timóre horríbili, et mirabúntur in subitatióne insperátæ salútis; dicéntes intra se, pœniténtiam agéntes, et præ angústia spíritus geméntes: Hi sunt quos habúimus aliquándo in derísum, et in similitúdinem impropérii. Nos insensáti, vitam illórum æstimabámus insániam, et finem illórum sine honóre; ecce quómodo computáti sunt inter fílios Dei, et inter sanctos sors illórum est.\par
\switchcolumn
\EN
\noindent they shall be in sorrow, and their memory shall perish. They shall come with fear at the thought of their sins, and their iniquities shall stand against them to convict them. Then shall the just stand with great constancy against those that have afflicted them, and taken away their labours. These seeing it, shall be troubled with terrible fear, and shall be amazed at the suddenness of their unexpected salvation. Saying within themselves, repenting, and groaning for anguish of spirit: These are they, whom we had some time in derision, and for a parable of reproach. We fools esteemed their life madness, and their end without honour. Behold how they are numbered among the children of God, and their lot is among the saints.\par
\switchcolumn*

\LA
\begin{Resp}\Rsign Iste cognóvit justítiam, et vidit mirabília magna, et exorávit Altíssimum: \Star{} Et invéntus est in número Sanctórum.\end{Resp}
\begin{Resp}\Vsign Iste est qui contémpsit vitam mundi, et pervénit ad cæléstia regna.\end{Resp}
\begin{Resp}\Rsign Et invéntus est in número Sanctórum.\end{Resp}
\begin{Resp}\Vsign Glória Patri, et Fílio, \Star{} et Spirítui Sancto.\end{Resp}
\begin{Resp}\Rsign Et invéntus est in número Sanctórum.\end{Resp}
\switchcolumn
\EN
\begin{Resp}\Rsign This is he which knew righteousness, and saw great wonders, and made his prayer unto the Most High; \Star{} And he is numbered among the Saints.\end{Resp}
\begin{Resp}\Vsign This is he which loved not his life in this world, and is come unto an everlasting kingdom.\end{Resp}
\begin{Resp}\Rsign And he is numbered among the Saints.\end{Resp}
\begin{Resp}\Vsign Glory be to the Father, and to the Son, \Star{} and to the Holy Ghost.\end{Resp}
\begin{Resp}\Rsign And he is numbered among the Saints.\end{Resp}
\switchcolumn*

\end{paracol}

\Office{Commune Doctoris non Pontificis}{Common of a Doctor not a Bishop}{}
\addcontentsline{toc}{subsection}{Commune Doctoris non Pontificis\enspace\textit{Common of a Doctor not a Bishop}}
\begin{paracol}{2}
\LA
\LessonHead{Lectio I}
\switchcolumn
\EN
\LessonHead{Lesson I}
\switchcolumn*

\LA
\LessonTitle{De libro Ecclesiástici}
\switchcolumn
\EN
\LessonTitle{Lesson from the book of Ecclesiasticus}
\switchcolumn*

\LA
\Cite{Sir 39:1-5}
\switchcolumn
\EN
\Cite{Sir 39:1-5}
\switchcolumn*

\LA
\noindent Sapiéntiam ómnium antiquórum exquíret sápiens, et in prophétis vacábit. Narratiónem virórum nominatórum conservábit, et in versútias parabolárum simul introíbit. Occúlta proverbiórum exquíret, et in abscónditis parabolárum conversábitur. In médio magnatórum ministrábit, et in conspéctu prǽsidis apparébit. In terram alienigenárum géntium pertránsiet; bona enim et mala in homínibus tentábit.\par
\switchcolumn
\EN
\noindent The wise men will seek out the wisdom of all the ancients, and will be occupied in the prophets. He will keep the sayings of renowned men, and will enter withal into the subtilties of parables. He will search out the hidden meanings of proverbs, and will be conversant in the secrets of parables. He shall serve among great men, and: appear before the governor. He shall pass into strange countries: for he shall try good and evil among men.\par
\switchcolumn*

\LA
\begin{Resp}\Rsign Euge, serve bone et fidélis, quia in pauca fuísti fidélis, supra multa te constítuam: \Star{} Intra in gáudium Dómini tui.\end{Resp}
\begin{Resp}\Vsign Dómine, quinque talénta tradidísti mihi, ecce ália quinque superlucrátus sum.\end{Resp}
\begin{Resp}\Rsign Intra in gáudium Dómini tui.\end{Resp}
\switchcolumn
\EN
\begin{Resp}\Rsign Well done, thou good and faithful servant, thou hast been faithful over a few things. I will make thee ruler over many things: \Star{} Enter thou into the joy of thy Lord.\end{Resp}
\begin{Resp}\Vsign Lord, thou deliveredst unto me five talents; behold, I have gained beside them five talents more.\end{Resp}
\begin{Resp}\Rsign Enter thou into the joy of thy Lord.\end{Resp}
\switchcolumn*

\LA
\LessonHead{Lectio II}
\switchcolumn
\EN
\LessonHead{Lesson II}
\switchcolumn*

\LA
\Cite{Sir 39:6-10}
\switchcolumn
\EN
\Cite{Sir 39:6-10}
\switchcolumn*

\LA
\noindent Cor suum tradet ad vigilándum dilúculo ad Dóminum, qui fecit illum, et in conspéctu Altíssimi deprecábitur. Apériet os suum in oratióne, et pro delíctis suis deprecábitur. Si enim Dóminus magnus volúerit, spíritu intellegéntiæ replébit illum: et ipse tamquam imbres mittet elóquia sapiéntiæ suæ, et in oratióne confitébitur Dómino: et ipse díriget consílium ejus, et disciplínam, et in abscónditis suis consiliábitur.\par
\switchcolumn
\EN
\noindent He will give his heart to resort early to the Lord that made him, and he will pray in the sight of the most High. He will open his mouth in prayer, and will make supplication for his sins. For if it shall please the great Lord, he will fill him with the spirit of understanding: And he will pour forth the words of his wisdom as showers, and in his prayer he will confess to the Lord. And he shall direct his counsel, and his knowledge, and in his secrets shall he meditate.\par
\switchcolumn*

\LA
\begin{Resp}\Rsign Justus germinábit sicut lílium: \Star{} Et florébit in ætérnum ante Dóminum.\end{Resp}
\begin{Resp}\Vsign Plantátus in domo Dómini, in átriis domus Dei nostri.\end{Resp}
\begin{Resp}\Rsign Et florébit in ætérnum ante Dóminum.\end{Resp}
\switchcolumn
\EN
\begin{Resp}\Rsign The righteous shall grow as the lily; \Star{} Yea, he shall flourish in the presence of the Lord for ever.\end{Resp}
\begin{Resp}\Vsign Those that be planted in the house of the Lord, shall flourish in the courts of the house of our God.\end{Resp}
\begin{Resp}\Rsign Yea, he shall flourish in the presence of the Lord for ever.\end{Resp}
\switchcolumn*

\LA
\LessonHead{Lectio III}
\switchcolumn
\EN
\LessonHead{Lesson III}
\switchcolumn*

\LA
\Cite{Sir 39:11-14}
\switchcolumn
\EN
\Cite{Sir 39:11-14}
\switchcolumn*

\LA
\noindent Ipse palam fáciet disciplínam doctrínæ suæ, et in lege testaménti Dómini gloriábitur. Collaudábunt multi sapiéntiam ejus, et usque in sǽculum non delébitur. Non recédet memória ejus, et nomen ejus requirétur a generatióne in generatiónem. Sapiéntiam ejus enarrábunt gentes, et laudem ejus enuntiábit Ecclésia.\par
\switchcolumn
\EN
\noindent He shall show forth the discipline he hath learned, and shall glory in the law of the covenant of the Lord. Many shall praise his wisdom, and it shall never be forgotten. The memory of him shall not depart away, and his name shall be in request from generation to generation. Nations shall declare his wisdom, and the church shall show forth his praise.\par
\switchcolumn*

\LA
\begin{Resp}\Rsign Iste cognóvit justítiam, et vidit mirabília magna, et exorávit Altíssimum: \Star{} Et invéntus est in número Sanctórum.\end{Resp}
\begin{Resp}\Vsign Iste est qui contémpsit vitam mundi, et pervénit ad cæléstia regna.\end{Resp}
\begin{Resp}\Rsign Et invéntus est in número Sanctórum.\end{Resp}
\begin{Resp}\Vsign Glória Patri, et Fílio, \Star{} et Spirítui Sancto.\end{Resp}
\begin{Resp}\Rsign Et invéntus est in número Sanctórum.\end{Resp}
\switchcolumn
\EN
\begin{Resp}\Rsign This is he which knew righteousness, and saw great wonders, and made his prayer unto the Most High; \Star{} And he is numbered among the Saints.\end{Resp}
\begin{Resp}\Vsign This is he which loved not his life in this world, and is come unto an everlasting kingdom.\end{Resp}
\begin{Resp}\Rsign And he is numbered among the Saints.\end{Resp}
\begin{Resp}\Vsign Glory be to the Father, and to the Son, \Star{} and to the Holy Ghost.\end{Resp}
\begin{Resp}\Rsign And he is numbered among the Saints.\end{Resp}
\switchcolumn*

\LA
\LessonHead{Lectio VII}
\switchcolumn
\EN
\LessonHead{Lesson VII}
\switchcolumn*

\LA
\LessonTitle{Léctio sancti Evangélii secúndum Matthǽum}
\switchcolumn
\EN
\LessonTitle{From the Holy Gospel according to Matthew}
\switchcolumn*

\LA
\Cite{Matt 5:13-19}
\switchcolumn
\EN
\Cite{Matt 5:13-19}
\switchcolumn*

\LA
\noindent In illo témpore: Dixit Jesus discípulis suis: Vos estis sal terræ. Quod si sal evanúerit, in quo saliétur? Et réliqua.\par
\switchcolumn
\EN
\noindent At that time: Jesus said unto his disciples: Ye are the salt of the earth: But if the salt have lost his savour, wherewith shall it be salted? And so on, and that which followeth.\par
\switchcolumn*

\LA
\LessonTitle{Homilía sancti Augustíni Epíscopi}
\switchcolumn
\EN
\LessonTitle{A Homily by St. Augustine the Bishop}
\switchcolumn*

\LA
\Source{Lib. 1 de Sermone Domini in monte, cap. 6}
\switchcolumn
\EN
\Source{Lib. 1 de Sermóne Dómini in monte, cap. 6}
\switchcolumn*

\LA
\noindent Osténdit Dóminus fátuos esse judicándos, qui temporálium bonórum vel cópiam sectántes vel inópiam metuéntes, amíttunt ætérna, quæ nec dari possunt ab homínibus nec auférri. Itaque, si sal infatuátum fúerit, in quo saliétur? Id est, si vos, per quos condiéndi sunt quodámmodo pópuli, metu persecutiónum temporálium amiséritis regna cælórum; qui erunt hómines, per quos a vobis error auferátur, cum vos elégerit Deus, per quos errórem áuferat ceterórum?\par
\switchcolumn
\EN
\noindent The Lord would have us understand how that men do lose their power of savouring others with righteousness when they are willing to place their eternal welfare in jeopardy for the sake of any témporal advantage, like as attainment of ease or luxury, or escape from suffering or toil. For that which is eternal, unlike things of this world, can neither be bestowed by men, nor by them taken away. Hence, when he asketh: If the salt have lost his savour, wherewith shall it be salted? he would have us understand the question to be: If ye, by whom mankind is preserved from corruption, be willing to lose the kingdom of heaven so as to escape trials or persecutions in this world, who is there to preserve you from corruption, seeing ye are they that God hath chosen to preserve all others from corruption?\par
\switchcolumn*

\LA
\begin{Resp}\Rsign Iste est qui ante Deum magnas virtútes operátus est, et de omni corde suo laudávit Dóminum: \Star{} Ipse intercédat pro peccátis ómnium populórum.\end{Resp}
\begin{Resp}\Vsign Ecce homo sine queréla, verus Dei cultor, ábstinens se ab omni ópere malo, et pérmanens in innocéntia sua.\end{Resp}
\begin{Resp}\Rsign Ipse intercédat pro peccátis ómnium populórum.\end{Resp}
\switchcolumn
\EN
\begin{Resp}\Rsign This is he which wrought great wonders before God, and praised the Lord with all his heart. \Star{} May he pray for all people, that their sins may be forgiven unto them.\end{Resp}
\begin{Resp}\Vsign Behold a man without blame, a worshipper of God in truth, keeping himself clean from every evil work, and abiding still in his innocency.\end{Resp}
\begin{Resp}\Rsign May he pray for all people, that their sins may be forgiven unto them.\end{Resp}
\switchcolumn*

\LA
\LessonHead{Lectio VIII}
\switchcolumn
\EN
\LessonHead{Lesson VIII}
\switchcolumn*

\LA
\noindent Ergo ad níhilum valet sal infatuátum, nisi ut mittátur foras et calcétur ab homínibus. Non ítaque calcátur ab homínibus qui pátitur persecutiónem; sed qui, persecutiónem timéndo, infatuátur. Calcári enim non potest nisi inférior; sed inférior non est, qui, quamvis córpore multa in terra sustíneat, corde tamen fixus in cælo est.\par
\switchcolumn
\EN
\noindent Those that should be the salt of the earth, but have lost their savour, are thenceforth good for nothing, but to be cast out, and to be trodden under foot of men. But no one that suffereth persecution is truly said to be trodden under foot of men. Rather, that one is truly trodden under foot of men who through fear of persecution hath lost the savour of righteousness. For no one can be trodden upon, unless he be beneath him which treadeth upon him. And certainly no one who hath his heart in heaven, no matter how grievously he doth suffer in his body on earth, is rightly said to be beneath anyone who misuseth him.\par
\switchcolumn*

\LA
\begin{Resp}\Rsign In médio Ecclésiæ apéruit os ejus, \Star{} Et implévit eum Dóminus spíritu sapiéntiæ et intelléctus.\end{Resp}
\begin{Resp}\Vsign Jucunditátem et exsultatiónem thesaurizávit super eum.\end{Resp}
\begin{Resp}\Rsign Et implévit eum Dóminus spíritu sapiéntiæ et intelléctus.\end{Resp}
\begin{Resp}\Vsign Glória Patri, et Fílio, \Star{} et Spirítui Sancto.\end{Resp}
\begin{Resp}\Rsign Et implévit eum Dóminus spíritu sapiéntiæ et intelléctus.\end{Resp}
\switchcolumn
\EN
\begin{Resp}\Rsign In the midst of the congregation did the Lord open his mouth. \Star{} And filled him with the spirit of wisdom and understanding.\end{Resp}
\begin{Resp}\Vsign He made him rich with joy and gladness.\end{Resp}
\begin{Resp}\Rsign And filled him with the spirit of wisdom and understanding.\end{Resp}
\begin{Resp}\Vsign Glory be to the Father, and to the Son, \Star{} and to the Holy Ghost.\end{Resp}
\begin{Resp}\Rsign And filled him with the spirit of wisdom and understanding.\end{Resp}
\switchcolumn*

\end{paracol}

\Office{Commune Abbatis}{Common of an Abbot}{}
\addcontentsline{toc}{subsection}{Commune Abbatis\enspace\textit{Common of an Abbot}}
\begin{paracol}{2}
\LA
\LessonHead{Lectio VII}
\switchcolumn
\EN
\LessonHead{Lesson VII}
\switchcolumn*

\LA
\LessonTitle{Léctio sancti Evangélii secúndum Matthǽum.}
\switchcolumn
\EN
\LessonTitle{From the Holy Gospel according to Matthew}
\switchcolumn*

\LA
\Cite{Matt 19:27-29}
\switchcolumn
\EN
\Cite{Matt 19:27-29}
\switchcolumn*

\LA
\noindent In illo témpore: Dixit Petrus ad Jesum: Ecce nos relíquimus ómnia, et secúti sumus te: quid ergo erit nobis? Et réliqua.\par
\switchcolumn
\EN
\noindent At that time, Peter said unto Jesus: Behold, we have forsaken all, and followed thee what shall we have therefore? And so on.\par
\switchcolumn*

\LA
\LessonTitle{Homilía sancti Hierónymi Presbýteri.}
\switchcolumn
\EN
\LessonTitle{Homily by St. Jerome, Priest (at Bethlehem.)}
\switchcolumn*

\LA
\Source{Lib. 3. in Matth. cap. 19.}
\switchcolumn
\EN
\Source{Bk. iii. on Matth xix.}
\switchcolumn*

\LA
\noindent Grandis fidúcia! Petrus piscátor erat, dives non fúerat, cibos manu et arte quærébat: et tamen lóquitur confidénter: Relíquimus ómnia. Et quia non súfficit tantum relínquere, jungit quod perféctum est: Et secúti sumus te. Fécimus quod jussísti: quid ígitur nobis dabis prǽmii? Jesus autem dixit illis: Amen dico vobis, quod vos qui secúti estis me, in regeneratióne, cum séderit Fílius hóminis in sede majestátis suæ, sedébitis et vos super sedes duódecim, judicántes duódecim tribus Ísraël. Non dixit: Qui reliquístis ómnia; hoc enim et Crates fecit philósophus, et multi álii divítias contempsérunt: sed, Qui secúti estis me: quod próprie Apostolórum est, atque credéntium.\par
\switchcolumn
\EN
\noindent Peter was a fisherman, he was not rich, he earned his bread by his hand and skill, and nevertheless he is thus bold, and saith confidently: We have forsaken all. And because it sufficeth not to forsake only, he addeth that which to do is to be perfect: and followed thee. We have done that which Thou hast commanded us, what reward therefore wilt Thou give us? And Jesus said unto them Amen I say unto you, that ye which have followed Me, in the regeneration, when the Son of Man shall sit in the throne of His glory, ye also shall sit upon twelve thrones, judging the twelve tribes of Israel. He said not, Ye which have forsaken all, for this did even Crates the philosopher, and they which have set nothing by riches are many, but, Ye which have followed Me. This did the Apostles, and this do believers do.\par
\switchcolumn*

\LA
\begin{Resp}\Rsign Iste est qui ante Deum magnas virtútes operátus est, et de omni corde suo laudávit Dóminum: \Star{} Ipse intercédat pro peccátis ómnium populórum.\end{Resp}
\begin{Resp}\Vsign Ecce homo sine queréla, verus Dei cultor, ábstinens se ab omni ópere malo, et pérmanens in innocéntia sua.\end{Resp}
\begin{Resp}\Rsign Ipse intercédat pro peccátis ómnium populórum.\end{Resp}
\switchcolumn
\EN
\begin{Resp}\Rsign This is he which wrought great wonders before God, and praised the Lord with all his heart. \Star{} May he pray for all people, that their sins may be forgiven unto them.\end{Resp}
\begin{Resp}\Vsign Behold a man without blame, a worshipper of God in truth, keeping himself clean from every evil work, and abiding still in his innocency.\end{Resp}
\begin{Resp}\Rsign May he pray for all people, that their sins may be forgiven unto them.\end{Resp}
\switchcolumn*

\LA
\LessonHead{Lectio VIII}
\switchcolumn
\EN
\LessonHead{Lesson VIII}
\switchcolumn*

\LA
\noindent In regeneratióne, cum séderit Fílius hóminis in sede majestátis suæ (quando et mórtui de corruptióne resúrgent incorrúpti) sedébitis et vos in sóliis judicántium, condemnántes duódecim tribus Ísraël: quia vobis credéntibus, illi crédere noluérunt. Et omnis qui relíquerit domum, vel fratres, aut soróres, aut patrem, aut matrem, aut uxórem, aut fílios, aut agros propter nomen meum, céntuplum accípiet, et vitam ætérnam possidébit. Locus iste cum illa senténtia cóngruit, in qua Salvátor lóquitur: Non veni pacem míttere, sed gládium. Veni enim separáre hóminem a patre suo, et matrem a fília, et nurum a socru: et inimíci hóminis doméstici ejus. Qui ergo propter fidem Christi, et prædicatiónem Evangélii, omnes afféctus contémpserint, atque divítias et sǽculi voluptátes: isti céntuplum recípient, et vitam ætérnam possidébunt.\par
\switchcolumn
\EN
\noindent In the regeneration, when the Son of Man shall sit in the throne of His glory, and when the dead shall rise again from corruption incorruptible, (i Cor. xv. 53,) ye also shall sit upon twelve thrones of judgment, condemning the twelve tribes of Israel, because, when ye believed in Me, they would not. (John iii. 18.) And every one that hath forsaken houses, or brethren, or sisters, or father, or mother, or wife, or children, or lands, for My Name's sake, shall receive an hundredfold, and shall inherit everlasting life. This place agreeth well with that other where the Saviour saith I came not to send peace, but a sword. For I am come to set a man at variance against his father, and the daughter against her mother, and the daughter-in-law against her mother-in-law; and a man's foes shall be they of his own household. (Matth. x. 34.) Every one, therefore, that hath set no store by affection, and riches, and the pleasures of the world, for Christ's faith's sake, and the preaching of the Gospel, shall receive an hundred-fold, and shall inherit everlasting life.\par
\switchcolumn*

\LA
\begin{Resp}\Rsign Sint lumbi vestri præcíncti, et lucérnæ ardéntes in mánibus vestris: \Star{} Et vos símiles homínibus exspectántibus dóminum suum, quando revertátur a núptiis.\end{Resp}
\begin{Resp}\Vsign Vigiláte ergo, quia nescítis qua hora Dóminus vester ventúrus sit.\end{Resp}
\begin{Resp}\Rsign Et vos símiles homínibus exspectántibus dóminum suum, quando revertátur a núptiis.\end{Resp}
\begin{Resp}\Vsign Glória Patri, et Fílio, \Star{} et Spirítui Sancto.\end{Resp}
\begin{Resp}\Rsign Et vos símiles homínibus exspectántibus dóminum suum, quando revertátur a núptiis.\end{Resp}
\switchcolumn
\EN
\begin{Resp}\Rsign Let your loins be girded about, and your lights burning; \Star{} And ye yourselves like unto men that wait for their lord, when he will return from the wedding.\end{Resp}
\begin{Resp}\Vsign Watch therefore, for ye know not what hour your Lord doth come.\end{Resp}
\begin{Resp}\Rsign And ye yourselves like unto men that wait for their lord, when he will return from the wedding.\end{Resp}
\begin{Resp}\Vsign Glory be to the Father, and to the Son, \Star{} and to the Holy Ghost.\end{Resp}
\begin{Resp}\Rsign And ye yourselves like unto men that wait for their lord, when he will return from the wedding.\end{Resp}
\switchcolumn*

\LA
\LessonHead{Lectio IX}
\switchcolumn
\EN
\LessonHead{Lesson IX}
\switchcolumn*

\LA
\noindent Ex occasióne hujus senténtiæ quidam introdúcunt mille annos post resurrectiónem, dicéntes, tunc nobis céntuplum ómnium rerum quas dimísimus, et vitam ætérnam esse reddéndam: non intelligéntes, quod, si in céteris digna sit repromíssio, in uxóribus appáreat turpitúdo: ut qui unam pro Dómino dimíserit centum recípiat in futúro. Sensus ergo iste est: Qui carnália pro Salvatóre dimíserit, spirituália recípiet: quæ comparatióne et mérito sui ita erunt, quasi si parvo número centenárius númerus comparétur.\par
\switchcolumn
\EN
\noindent By reason of these words, an hundredfold, some will have that there shall be a thousand years after the resurrection, wherein they that have forsaken all things shall receive an hundredfold of those things which they have forsaken, and shall inherit everlasting life. Such men consider not that though in other things this were worthy, as touching wives it is unseemly for it becometh us not to think that he that hath forsaken one wife in this world, shall receive an hundred wives in that which is to come. But the meaning is this, that every one that for the Saviour's sake hath forsaken earthly things, shall receive spiritual things which things, being rightly weighed against earthly things, are as though an hundredfold were weighed against one.\par
\switchcolumn*

\LA
\TeDeum{Te Deum}
\switchcolumn
\EN
\TeDeum{Te Deum}
\switchcolumn*

\end{paracol}

\Office{Commune Unius Virginis et Martyris}{Common of a Virgin Martyr}{}
\addcontentsline{toc}{subsection}{Commune Unius Virginis et Martyris\enspace\textit{Common of a Virgin Martyr}}
\begin{paracol}{2}
\LA
\LessonHead{Lectio VII}
\switchcolumn
\EN
\LessonHead{Lesson VII}
\switchcolumn*

\LA
\LessonTitle{Léctio sancti Evangélii secúndum Matthǽum}
\switchcolumn
\EN
\LessonTitle{From the Holy Gospel according to St. Matthew}
\switchcolumn*

\LA
\Cite{Matt 25:1-13}
\switchcolumn
\EN
\Cite{Matt 25:1}
\switchcolumn*

\LA
\noindent In illo témpore: Dixit Jesus discípulis suis parábolam hanc: Símile erit regnum cælórum decem virgínibus, quæ, accipiéntes lámpades suas, exiérunt óbviam sponso et sponsæ. Et réliqua.\par
\switchcolumn
\EN
\noindent At that time, Jesus said to His disciples: The Kingdom of heaven shall be likened unto ten virgins, which took their lamps, and went forth to meet the Bridegroom and the Bride. And so on.\par
\switchcolumn*

\LA
\LessonTitle{Homilía sancti Gregórii Papæ}
\switchcolumn
\EN
\LessonTitle{Homily by Pope St. Gregory the Great}
\switchcolumn*

\LA
\Source{Homilía 12 in Evangelia}
\switchcolumn
\EN
\Source{12th on the Gospels}
\switchcolumn*

\LA
\noindent Sæpe vos, fratres caríssimi, admóneo prava ópera fúgere, mundi hujus inquinaménta devitáre, sed hodiérna sancti Evangélii lectióne compéllor dícere, ut, et bona, quæ ágitis, cum magna cautéla teneátis; ne per hoc, quod a vobis rectum géritur, favor aut grátia humána requirátur; ne appetítus laudis subrépat, et, quod foris osténditur, intus a mercéde vacuétur. Ecce enim Redemptóris voce decem vírgines, et omnes dicúntur vírgines, et tamen intra beatitúdinis jánuam non omnes sunt recéptæ; quia eárum quædam, dum de virginitáte sua glóriam foris éxpetunt, in vasis suis óleum habére noluérunt.\par
\switchcolumn
\EN
\noindent Dearly beloved brethren oftentimes do I warn you to fly corrupt conversation, and to keep yourselves unspotted from the world. But the portion which is this day read from the Holy Gospel doth oblige me to say that even to these good things which ye do, ye must needs take all careful heed. Look ye well to it, that, when ye work righteousness, ye do it not as seeking the praise and admiration of men, for if the lust of praise do once creep in, that which seemeth so fair without, loseth its reward within. Behold how the Redeemer speaketh of these ten virgins. He calleth them all virgins, yet entered not all of them into the door of blessedness, for there were some of them who sought outwardly the honour of virginity, but would take no oil within their vessels with their lamps.\par
\switchcolumn*

\LA
\begin{Resp}\Rsign Hæc est Virgo sápiens, quam Dóminus vigilántem invénit, quæ accéptis lampádibus sumpsit secum óleum: \Star{} Et veniénte Dómino, introívit cum eo ad núptias.\end{Resp}
\begin{Resp}\Vsign Média nocte clamor factus est: Ecce sponsus venit, exíte óbviam ei.\end{Resp}
\begin{Resp}\Rsign Et veniénte Dómino, introívit cum eo ad núptias.\end{Resp}
\switchcolumn
\EN
\begin{Resp}\Rsign This is one of those wise virgins, whom the Lord found watching, for when she took her lamp, she took oil with her. \Star{} And when the Lord came, she went in with him to the marriage.\end{Resp}
\begin{Resp}\Vsign At midnight there was a cry made: Behold, the Bridegroom cometh, go ye out to meet him.\end{Resp}
\begin{Resp}\Rsign And when the Lord came, she went in with him to the marriage.\end{Resp}
\switchcolumn*

\LA
\LessonHead{Lectio VIII}
\switchcolumn
\EN
\LessonHead{Lesson VIII}
\switchcolumn*

\LA
\noindent Sed prius quæréndum nobis est quid sit regnum cælórum, aut cur decem virgínibus comparétur, quæ étiam vírgines prudéntes et fátuæ dicántur. Dum enim cælórum regnum constat quia reprobórum nullus ingréditur, étiam fátuis virgínibus cur símile esse perhibétur? Sed sciéndum nobis est quod sæpe in sacro elóquio regnum cælórum præséntis témporis Ecclésia dícitur. De quo álio in loco Dóminus dicit: Mittet Fílius hóminis Ángelos suos, et cólligent de regno ejus ómnia scándala. Neque enim in illo regno beatitúdinis, in quo pax summa est, inveníri scándala póterunt, quæ colligántur.\par
\switchcolumn
\EN
\noindent First of all, it is for us to ask What is the kingdom of Heaven? And wherefore shall the same be likened unto ten virgins, whereof, albeit five were wise, yet five were foolish For if the kingdom of heaven be such that there shall in no wise enter into it anything that defileth, neither whatsoever worketh abomination, or maketh a lie, (Apoc. xxi. 27,) how can it be like unto five virgins which were foolish? But we must know that, in the word of God, the kingdom of heaven doth oftentimes signify the Church as she now is, touching the which the Lord saith in another place “The Son of Man shall send forth His Angels, and they shall gather out of His kingdom all things that offend.” (Matth. xiii. 41.) In that kingdom of Blessedness, wherein peace shall have her perfect reign, there shall be nothing found that offendeth for the angels to gather out.\par
\switchcolumn*

\LA
\begin{Resp}\Rsign Média nocte clamor factus est: \Star{} Ecce sponsus venit, exíte óbviam ei.\end{Resp}
\begin{Resp}\Vsign Prudéntes vírgines, aptáte vestras lámpades.\end{Resp}
\begin{Resp}\Rsign Ecce sponsus venit, exíte óbviam ei.\end{Resp}
\begin{Resp}\Vsign Glória Patri, et Fílio, \Star{} et Spirítui Sancto.\end{Resp}
\begin{Resp}\Rsign Ecce sponsus venit, exíte óbviam ei.\end{Resp}
\switchcolumn
\EN
\begin{Resp}\Rsign At midnight there was a cry made: \Star{} Behold, the Bridegroom cometh, go ye out to meet him.\end{Resp}
\begin{Resp}\Vsign Trim your lamps, O ye wise virgins.\end{Resp}
\begin{Resp}\Rsign Behold, the Bridegroom cometh, go ye out to meet him.\end{Resp}
\begin{Resp}\Vsign Glory be to the Father, and to the Son, \Star{} and to the Holy Ghost.\end{Resp}
\begin{Resp}\Rsign Behold, the Bridegroom cometh, go ye out to meet him.\end{Resp}
\switchcolumn*

\LA
\LessonHead{Lectio IX}
\switchcolumn
\EN
\LessonHead{Lesson IX}
\switchcolumn*

\LA
\noindent In quinque autem córporis sénsibus unusquísque subsístit; geminátus autem quinárius denárium pérficit. Et, quia ex utróque sexu fidélium multitúdo collígitur, sancta Ecclésia decem virgínibus símilis esse denuntiátur. In qua quia mali cum bonis et réprobi cum eléctis admíxti sunt, recte símilis virgínibus prudéntibus et fátuis esse perhibétur. Sunt namque pleríque continéntes, qui ab appetítu se exterióri custódiunt et spe ad interióra rapiúntur, carnem mácerant, et toto desidério ad supérnam pátriam anhélant, ætérna prǽmia éxpetunt, pro labóribus suis recípere laudes humánas nolunt. Hi nimírum glóriam suam non in ore hóminum ponunt, sed intra consciéntiam cóntegunt. Et sunt pleríque, qui corpus per abstinéntiam afflígunt, sed de ipsa sua abstinéntia humános favóres éxpetunt.\par
\switchcolumn
\EN
\noindent The body of every man doth consist of five senses, and five being doubled, is ten. Forasmuch, therefore, as the whole body of the faithful doth consist of two sexes, the Holy Church is likened unto ten virgins. And forasmuch as in the Church the good are for the present mingled with the bad, and the reprobate with the elect, it is rightly said that, of the ten virgins, five are wise and five are foolish. There are many who have self-control, which do keep themselves from lusting after things outward, whose hope beareth them to things inward, who chastise the flesh, who long with intense home-sickness for their Fatherland which is in heaven, who seek an eternal reward, and who will not to receive for their labours the praise of men. These are they who reckon their glory, not in the mouths of men, but in the testimony of their own conscience. And many there be likewise who afflict the body by self-control, and yet who seek for their self-control applause from men.\par
\switchcolumn*

\LA
\TeDeum{Te Deum}
\switchcolumn
\EN
\TeDeum{Te Deum}
\switchcolumn*

\LA
\LessonHead{Lectio I (in 2 loco)}
\switchcolumn
\EN
\LessonHead{Lesson I (set 2)}
\switchcolumn*

\LA
\noindent Confitébor tibi, Dómine, Rex, et collaudábo te Deum Salvatórem meum. Confitébor nómini tuo: quóniam adjútor et protéctor factus es mihi, et liberásti corpus meum a perditióne, a láqueo linguæ iníquæ et a lábiis operántium mendácium, et in conspéctu astántium factus es mihi adjútor. Et liberásti me secúndum multitúdinem misericórdiæ nóminis tui a rugiéntibus, præparátis ad escam: de mánibus quæréntium ánimam meam, et de portis tribulatiónum, quæ circumdedérunt me; a pressúra flammæ, quæ circúmdedit me, et in médio ignis non sum æstuáta; de altitúdine ventris ínferi, et a lingua coinquináta, et a verbo mendácii, a rege iníquo, et a lingua injústa.\par
\switchcolumn
\EN
\noindent A prayer of Jesus the son of Sirach. I will give glory to thee, O Lord, O King, and I will praise thee, O God my Saviour. I will give glory to thy name: for thou hast been a helper and protector to me. And hast preserved my body from destruction, from the snare of an unjust tongue, and from the lips of them that forge lies, and in the sight of them that stood by, thou hast been my helper. And thou hast delivered me, according to the multitude of the mercy of thy name, from them that did roar, prepared to devour. Out of the hands of them that sought my life, and from the gates of afflictions, which compassed me about: From the oppression of the flame which surrounded me, and in the midst of the fire I was not burnt. From the depth of the belly of hell, and from an unclean tongue, and from lying words, from an unjust king, and from a slanderous tongue:\par
\switchcolumn*

\LA
\begin{Resp}\Rsign Veni Sponsa Christi, áccipe corónam, quam tibi Dóminus præparávit in ætérnum; pro cujus amóre sánguinem tuum fudísti, \Star{} Et cum Ángelis in paradísum introísti.\end{Resp}
\begin{Resp}\Vsign Veni, elécta mea, et ponam in te thronum meum; quia concupívit Rex spéciem tuam.\end{Resp}
\begin{Resp}\Rsign Et cum Ángelis in paradísum introísti.\end{Resp}
\switchcolumn
\EN

\switchcolumn*

\LA
\LessonHead{Lectio II (in 2 loco)}
\switchcolumn
\EN
\LessonHead{Lesson II (set 2)}
\switchcolumn*

\LA
\Cite{Sir 51:8-12}
\switchcolumn
\EN
\Cite{Sir 51:8-12}
\switchcolumn*

\LA
\noindent Laudábit usque ad mortem ánima mea Dóminum, et vita mea appropínquans erat in inférno deórsum. Circumdedérunt me úndique, et non erat qui adjuváret. Respíciens eram ad adjutórium hóminum, et non erat. Memoráta sum misericórdiæ tuæ, Dómine, et operatiónis tuæ, quæ a sǽculo sunt: quóniam éruis sustinéntes te, Dómine, et líberas eos de mánibus géntium.\par
\switchcolumn
\EN
\noindent My soul shall praise the Lord even to death. And my life was drawing near to hell beneath. They compassed me on every side, and there was no one that would help me. I looked for the succour of men, and there was none. I remembered thy mercy, O Lord, and thy works, which are from the beginning of the world. How thou deliverest them that wait for thee, O Lord, and savest them out of the hands of the nations.\par
\switchcolumn*

\LA
\begin{Resp}\Rsign Diffúsa est grátia in lábiis tuis, \Star{} Proptérea benedíxit te Deus in ætérnum.\end{Resp}
\begin{Resp}\Vsign Spécie tua et pulchritúdine tua inténde, próspere procéde, et regna.\end{Resp}
\begin{Resp}\Rsign Proptérea benedíxit te Deus in ætérnum.\end{Resp}
\switchcolumn
\EN
\begin{Resp}\Rsign Grace is poured into thy lips, therefore; \Star{} God hath blessed thee for ever.\end{Resp}
\begin{Resp}\Vsign In thy comeliness and thy beauty, go forward, fare prosperously, and reign.\end{Resp}
\begin{Resp}\Rsign God hath blessed thee for ever.\end{Resp}
\switchcolumn*

\LA
\LessonHead{Lectio III (in 2 loco)}
\switchcolumn
\EN
\LessonHead{Lesson III (set 2)}
\switchcolumn*

\LA
\Cite{Sir 51:13-17}
\switchcolumn
\EN
\Cite{Sir 51:13-17}
\switchcolumn*

\LA
\noindent Exaltásti super terram habitatiónem meam, et pro morte defluénte deprecáta sum. Invocávi Dóminum, Patrem Dómini mei, ut non derelínquat me in die tribulatiónis meæ, et in témpore superbórum, sine adjutório. Laudábo nomen tuum assídue, et collaudábo illud in confessióne, et exaudíta est orátio mea. Et liberásti me de perditióne, et eripuísti me de témpore iníquo. Proptérea confitébor, et laudem dicam tibi, et benedícam nómini Dómini.\par
\switchcolumn
\EN
\noindent Thou hast exalted my dwelling place upon the earth and I have prayed for death to pass away. I called upon the Lord, the father of my Lord, that he would not leave me in the day of my trouble, and in the time of the proud without help. I will praise thy name continually, and will praise it with thanksgiving, and my prayer was heard. And thou hast saved me from destruction, and hast delivered me from the evil time. Therefore I will give thanks, and praise thee, and bless the name of the Lord.\par
\switchcolumn*

\LA
\begin{Resp}\Rsign Spécie tua et pulchritúdine tua \Star{} Inténde, próspere procéde, et regna.\end{Resp}
\begin{Resp}\Vsign Diffúsa est grátia in lábiis tuis, proptérea benedíxit te Deus in ætérnum.\end{Resp}
\begin{Resp}\Rsign Inténde, próspere procéde, et regna.\end{Resp}
\begin{Resp}\Vsign Glória Patri, et Fílio, \Star{} et Spirítui Sancto.\end{Resp}
\begin{Resp}\Rsign Inténde, próspere procéde, et regna.\end{Resp}
\switchcolumn
\EN
\begin{Resp}\Rsign In thy comeliness and thy beauty; \Star{} Go forward, fare prosperously, and reign.\end{Resp}
\begin{Resp}\Vsign Grace is poured into thy lips, therefore God hath blessed thee for ever.\end{Resp}
\begin{Resp}\Rsign Go forward, fare prosperously, and reign.\end{Resp}
\begin{Resp}\Vsign Glory be to the Father, and to the Son, \Star{} and to the Holy Ghost.\end{Resp}
\begin{Resp}\Rsign Go forward, fare prosperously, and reign.\end{Resp}
\switchcolumn*

\LA
\LessonHead{Lectio VII (in 2 loco)}
\switchcolumn
\EN
\LessonHead{Lesson VII (set 2)}
\switchcolumn*

\LA
\LessonTitle{Léctio sancti Evangélii secúndum Matthǽum}
\switchcolumn
\EN
\LessonTitle{From the Holy Gospel according to Matthew}
\switchcolumn*

\LA
\Cite{Matt 13:44-52}
\switchcolumn
\EN
\Cite{Matt 13:44}
\switchcolumn*

\LA
\noindent In illo témpore: Dixit Jesus discípulis suis parábolam hanc: Simile est regnum cælórum thesáuro abscóndito in agro. Et réliqua.\par
\switchcolumn
\EN
\noindent At that time Jesus spake unto His disciples this parable The kingdom of heaven is like unto treasure hid in a field. And so on.\par
\switchcolumn*

\LA
\LessonTitle{Homilía sancti Gregórii Papæ}
\switchcolumn
\EN
\LessonTitle{Homily by Pope St Gregory the Great}
\switchcolumn*

\LA
\Source{Homilia 11 in Evangelia}
\switchcolumn
\EN
\Source{11th on the Gospels}
\switchcolumn*

\LA
\noindent Cælórum regnum, fratres caríssimi, idcírco terrénis rebus símile dícitur, ut, ex his quæ ánimus novit, surgat ad incógnita, quæ non novit: quátenus exémplo visibílium se ad invisibília rápiat, et, per ea quæ usu dídicit quasi confricátus incaléscat; ut per hoc, quod scit notum dilígere, discat et incógnita amáre. Ecce enim cælórum regnum thesáuro abscóndito in agro comparátur; quem, qui invénit homo, abscóndit, et præ gáudio illíus vadit, et vendit univérsa quæ habet, et emit agrum illum.\par
\switchcolumn
\EN
\noindent Dearly beloved brethren, the kingdom of heaven is likened unto the things of earth, to the end that by the mean of things which we know, our mind may rise to the contemplation of the things which we know not by the example of things which are seen, may fix her gaze on things which are not seen by the touch of things which she useth, may be warmed towards the things which she useth not; by things which she knoweth and loveth, to love also the things which she knoweth not. For, behold, “the kingdom of heaven is likened unto treasure hid in a field, the which when a man hath found, he hideth, and, for joy thereof, goeth and selleth all that he hath and buyeth that field.”\par
\switchcolumn*

\LA
\begin{Resp}\Rsign Hæc est Virgo sápiens, quam Dóminus vigilántem invénit, quæ accéptis lampádibus sumpsit secum óleum: \Star{} Et veniénte Dómino, introívit cum eo ad núptias.\end{Resp}
\begin{Resp}\Vsign Média nocte clamor factus est: Ecce sponsus venit, exíte óbviam ei.\end{Resp}
\begin{Resp}\Rsign Et veniénte Dómino, introívit cum eo ad núptias.\end{Resp}
\switchcolumn
\EN
\begin{Resp}\Rsign This is one of those wise virgins, whom the Lord found watching, for when she took her lamp, she took oil with her. \Star{} And when the Lord came, she went in with him to the marriage.\end{Resp}
\begin{Resp}\Vsign At midnight there was a cry made: Behold, the Bridegroom cometh, go ye out to meet him.\end{Resp}
\begin{Resp}\Rsign And when the Lord came, she went in with him to the marriage.\end{Resp}
\switchcolumn*

\LA
\LessonHead{Lectio VIII (in 2 loco)}
\switchcolumn
\EN
\LessonHead{Lesson VIII (set 2)}
\switchcolumn*

\LA
\noindent Qua in re hoc quoque notándum est, quod invéntus thesáurus abscónditur, ut servétur: quia stúdium cæléstis desidérii a malígnis spirítibus custodíre non súfficit, qui hoc ab humánis láudibus non abscóndit. In præsénti étenim vita, quasi in via sumus, qua ad pátriam pérgimus. Malígni autem spíritus iter nostrum quasi quidam latrúnculi óbsident. Deprædári ergo desíderat, qui thesáurum públice portat in via. Hoc autem dico, non ut próximi ópera nostra bona non vídeant, cum scriptum sit: Vídeant ópera vestra bona, et gloríficent Patrem vestrum qui in cælis est; sed, ut per hoc quod ágimus, laudes extérius non quærámus. Sic autem sit opus in público, quátenus inténtio máneat in occúlto; ut, et de bono ópere próximis præbeámus exémplum, et tamen per intentiónem, qua Deo soli placére quǽrimus, semper optémus secrétum.\par
\switchcolumn
\EN
\noindent And herein we must remark that the treasure, when once it hath been found, is hidden to keep it safe. He whose intimate yearnings after God are not hidden from the praise of men, is open thereby to the attacks of evil spirits. In this life we are, as it were, journeying homewards on a road beset by evil spirits who are like highwaymen. He therefore inviteth robbery who carrieth his treasure ostentatiously. Doubtless our neighbour should be able to see our good works, as it is written: “Let your light so shine before men that they may see your good works, and glorify your Father which is in heaven.” But this is not to be understood to mean that we are to seek the praise of men by what we do. Rather, let us in such wise work in the open that the inner intention of devotion is not advertised. So we shall give an example to our neighbour, and yet keep hidden, except from the sight of God, our purpose of pleasing him.\par
\switchcolumn*

\LA
\begin{Resp}\Rsign Média nocte clamor factus est: \Star{} Ecce sponsus venit, exíte óbviam ei.\end{Resp}
\begin{Resp}\Vsign Prudéntes vírgines, aptáte vestras lámpades.\end{Resp}
\begin{Resp}\Rsign Ecce sponsus venit, exíte óbviam ei.\end{Resp}
\begin{Resp}\Vsign Glória Patri, et Fílio, \Star{} et Spirítui Sancto.\end{Resp}
\begin{Resp}\Rsign Ecce sponsus venit, exíte óbviam ei.\end{Resp}
\switchcolumn
\EN
\begin{Resp}\Rsign At midnight there was a cry made: \Star{} Behold, the Bridegroom cometh, go ye out to meet him.\end{Resp}
\begin{Resp}\Vsign Trim your lamps, O ye wise virgins.\end{Resp}
\begin{Resp}\Rsign Behold, the Bridegroom cometh, go ye out to meet him.\end{Resp}
\begin{Resp}\Vsign Glory be to the Father, and to the Son, \Star{} and to the Holy Ghost.\end{Resp}
\begin{Resp}\Rsign Behold, the Bridegroom cometh, go ye out to meet him.\end{Resp}
\switchcolumn*

\LA
\LessonHead{Lectio IX (in 2 loco)}
\switchcolumn
\EN
\LessonHead{Lesson IX (set 2)}
\switchcolumn*

\LA
\noindent Thesáurus autem cæléste est desidérium; ager vero, in quo thesáurus abscónditur, disciplína stúdii cæléstis. Quem profécto agrum, vénditis ómnibus, cómparat, qui, voluptátibus carnis renúntians, cuncta sua terréna desidéria per disciplínæ cæléstis custódiam calcat: ut nihil jam quod caro blandítur, líbeat; nihil, quod carnálem vitam trucídat, spíritus perhorréscat.\par
\switchcolumn
\EN
\noindent The treasure is the desire for heaven, the field wherein it is hidden is the earnest observance wherewith this desire is surrounded. Whosoever turneth his back upon the enjoyments of the flesh, and by earnest striving heavenward, putteth all earthly lusts under the feet of discipline, so that he smileth back no more when the flesh smileth at him, and shuddereth no more at anything that can only kill the body whosoever doth thus, hath sold all that he had, and bought that field.\par
\switchcolumn*

\LA
\TeDeum{Te Deum}
\switchcolumn
\EN
\TeDeum{Te Deum}
\switchcolumn*

\LA
\LessonHead{Lectio VII (in 3 loco)}
\switchcolumn
\EN
\LessonHead{Lesson VII (set 3)}
\switchcolumn*

\LA
\LessonTitle{Léctio sancti Evangélii secúndum Matthǽum}
\switchcolumn
\EN
\LessonTitle{From the Holy Gospel according to Matthew}
\switchcolumn*

\LA
\Cite{Matt 19:3-12}
\switchcolumn
\EN
\Cite{Matt 19:3-12}
\switchcolumn*

\LA
\noindent In illo témpore: Accessérunt ad Jesum pharisæi, tentántes eum et dicéntes: Si licet hómini dimíttere uxórem suam quacúmque ex causa? Et réliqua.\par
\switchcolumn
\EN
\noindent In that time, there came to him the Pharisees tempting him, and saying: Is it lawful for a man to put away his wife for every cause? And so on.\par
\switchcolumn*

\LA
\LessonTitle{Homilía S. Joánnis Chrysostomi}
\switchcolumn
\EN
\LessonTitle{Homily by St. John Chrysostom, Patriarch of (Constantinople.)}
\switchcolumn*

\LA
\Source{Homilia 63 in Matth., sub med.}
\switchcolumn
\EN
\Source{63th on Matthew.}
\switchcolumn*

\LA
\noindent Dóminus noster, quia ad virginitátem hortári grave esse videbátur, a necessitáte legis de non solvéndo matrimónio, eos ad cupiditátem illíus tráhere studet. Deínde, ut eam esse possíbilem osténdat, sic inquit: Sunt eunúchi qui ex ventre matris ita nati sunt; et sunt qui ab homínibus eunúchi facti sunt, et sunt qui seípsos castravérunt propter regnum cælórum. Quibus verbis laténter eos ad eligéndam virginitátem indúcit, dum eam virtútem esse possíbilem ástruit.\par
\switchcolumn
\EN
\noindent Seeing that directly to exhort them unto virginity was well-nigh more than they could bear, our Lord seeketh to draw them to the desire thereof, taking occasion by the needfulness of a law against divorce. Then He showeth that virginity is possible, saying There are some eunuchs which were so born from their mother's womb; and there are some eunuchs which were made eunuchs of men; and there be eunuchs which have made themselves eunuchs for the kingdom of heaven's sake. In these words He persuadeth them indirectly to choose virginity, while He teacheth them that such a gift is not so good as to be impossible.\par
\switchcolumn*

\LA
\begin{Resp}\Rsign Hæc est Virgo sápiens, quam Dóminus vigilántem invénit, quæ accéptis lampádibus sumpsit secum óleum: \Star{} Et veniénte Dómino, introívit cum eo ad núptias.\end{Resp}
\begin{Resp}\Vsign Média nocte clamor factus est: Ecce sponsus venit, exíte óbviam ei.\end{Resp}
\begin{Resp}\Rsign Et veniénte Dómino, introívit cum eo ad núptias.\end{Resp}
\switchcolumn
\EN
\begin{Resp}\Rsign This is one of those wise virgins, whom the Lord found watching, for when she took her lamp, she took oil with her. \Star{} And when the Lord came, she went in with him to the marriage.\end{Resp}
\begin{Resp}\Vsign At midnight there was a cry made: Behold, the Bridegroom cometh, go ye out to meet him.\end{Resp}
\begin{Resp}\Rsign And when the Lord came, she went in with him to the marriage.\end{Resp}
\switchcolumn*

\LA
\LessonHead{Lectio VIII (in 3 loco)}
\switchcolumn
\EN
\LessonHead{Lesson VIII (set 3)}
\switchcolumn*

\LA
\noindent Id autem hoc fere modo confírmat: Cógita tecum, si aut a natúra talis esses, aut ab homínibus eam injúriam passus, quid fáceres, cum et ejúsmodi voluptáte caréres, et nullam, caréndo, mercédem consequeréris? Grátias ígitur nunc Deo agas, quod cum mercéde atque corónis idem sústines, quod illi sine corónis ac præmio tólerant. Immo vero non idem, sed multo lévius; tum quia spe erígeris et consciéntia recte facti, tum quia non ita ingéntibus concupiscéntiæ flúctibus jactáris.\par
\switchcolumn
\EN
\noindent This doctrine He establisheth somewhat thus. Supposing that thou hadst been born a eunuch by nature, or hadst been made a eunuch by the cruelty of men, so that thou hadst no sexual enjoyment, and hadst no credit for having none, what wouldest thou do? Give God thanks therefore, that thou dost, for a reward and a crown, what such others suffer with no reward and no crown yea, and that, rather a lighter burden than the same, and that not only because thou hast the joy of hope, and of knowing that thou doest well, but also because thou art not so battered by storms of desire as they are.\par
\switchcolumn*

\LA
\begin{Resp}\Rsign Média nocte clamor factus est: \Star{} Ecce sponsus venit, exíte óbviam ei.\end{Resp}
\begin{Resp}\Vsign Prudéntes vírgines, aptáte vestras lámpades.\end{Resp}
\begin{Resp}\Rsign Ecce sponsus venit, exíte óbviam ei.\end{Resp}
\begin{Resp}\Vsign Glória Patri, et Fílio, \Star{} et Spirítui Sancto.\end{Resp}
\begin{Resp}\Rsign Ecce sponsus venit, exíte óbviam ei.\end{Resp}
\switchcolumn
\EN
\begin{Resp}\Rsign At midnight there was a cry made: \Star{} Behold, the Bridegroom cometh, go ye out to meet him.\end{Resp}
\begin{Resp}\Vsign Trim your lamps, O ye wise virgins.\end{Resp}
\begin{Resp}\Rsign Behold, the Bridegroom cometh, go ye out to meet him.\end{Resp}
\begin{Resp}\Vsign Glory be to the Father, and to the Son, \Star{} and to the Holy Ghost.\end{Resp}
\begin{Resp}\Rsign Behold, the Bridegroom cometh, go ye out to meet him.\end{Resp}
\switchcolumn*

\LA
\LessonHead{Lectio IX (in 3 loco)}
\switchcolumn
\EN
\LessonHead{Lesson IX (set 3)}
\switchcolumn*

\LA
\noindent Cum ergo de illis eunúchis díxerit, qui, nisi étiam ipsi mente se contíneant, frustra et inutíliter tales sunt, ac de illis, qui, ut regnum cælórum consequántur, se cóntinent; subjúnxit rursus, dicens: Qui potest cápere, capiat; ut alacrióres effíciat eo ipso quod grave esse opus osténdit, et ineffábili benignitáte sua nóluit eam rem intra necessitátem legis conclúdere. Idque dicéndo, adhuc magis possíbilem esse demónstrat, ut voluntátis stúdium plus augeátur.\par
\switchcolumn
\EN
\noindent When, therefore, He had spoken of such as are eunuchs by nature, or by mutilation, and are eunuchs vainly and uselessly, unless they also bridle their thoughts, and of such as deny themselves for the kingdom of heaven's sake, He added He that is able to receive it, let him receive it, that He might make them the readier by showing the very sternness of the work, and, in His unspeakable goodness, He would not include any such precept within the requirements of the law, and, by saying this, showeth it to be the more possible, that He might increase the desire of freely choosing it.\par
\switchcolumn*

\LA
\TeDeum{Te Deum}
\switchcolumn
\EN
\TeDeum{Te Deum}
\switchcolumn*

\end{paracol}

\Office{Commune Virginum tantum}{Common of a Virgin not a Martyr}{}
\addcontentsline{toc}{subsection}{Commune Virginum tantum\enspace\textit{Common of a Virgin not a Martyr}}
\begin{paracol}{2}
\LA
\LessonHead{Lectio VII}
\switchcolumn
\EN
\LessonHead{Lesson VII}
\switchcolumn*

\LA
\LessonTitle{Léctio sancti Evangélii secúndum Matthǽum}
\switchcolumn
\EN
\LessonTitle{From the Holy Gospel according to St. Matthew}
\switchcolumn*

\LA
\Cite{Matt 25:1-13}
\switchcolumn
\EN
\Cite{Matt 25:1}
\switchcolumn*

\LA
\noindent In illo témpore: Dixit Jesus discípulis suis parábolam hanc: Símile erit regnum cælórum decem virgínibus, quæ, accipiéntes lámpades suas, exiérunt óbviam sponso et sponsæ. Et réliqua.\par
\switchcolumn
\EN
\noindent At that time, Jesus said to His disciples: The Kingdom of heaven shall be likened unto ten virgins, which took their lamps, and went forth to meet the Bridegroom and the Bride. And so on.\par
\switchcolumn*

\LA
\LessonTitle{Homilía sancti Gregórii Papæ}
\switchcolumn
\EN
\LessonTitle{Homily by Pope St. Gregory the Great}
\switchcolumn*

\LA
\Source{Homilía 12 in Evangelia}
\switchcolumn
\EN
\Source{12th on the Gospels}
\switchcolumn*

\LA
\noindent Sæpe vos, fratres caríssimi, admóneo prava ópera fúgere, mundi hujus inquinaménta devitáre, sed hodiérna sancti Evangélii lectióne compéllor dícere, ut, et bona, quæ ágitis, cum magna cautéla teneátis; ne per hoc, quod a vobis rectum géritur, favor aut grátia humána requirátur; ne appetítus laudis subrépat, et, quod foris osténditur, intus a mercéde vacuétur. Ecce enim Redemptóris voce decem vírgines, et omnes dicúntur vírgines, et tamen intra beatitúdinis jánuam non omnes sunt recéptæ; quia eárum quædam, dum de virginitáte sua glóriam foris éxpetunt, in vasis suis óleum habére noluérunt.\par
\switchcolumn
\EN
\noindent Dearly beloved brethren oftentimes do I warn you to fly corrupt conversation, and to keep yourselves unspotted from the world. But the portion which is this day read from the Holy Gospel doth oblige me to say that even to these good things which ye do, ye must needs take all careful heed. Look ye well to it, that, when ye work righteousness, ye do it not as seeking the praise and admiration of men, for if the lust of praise do once creep in, that which seemeth so fair without, loseth its reward within. Behold how the Redeemer speaketh of these ten virgins. He calleth them all virgins, yet entered not all of them into the door of blessedness, for there were some of them who sought outwardly the honour of virginity, but would take no oil within their vessels with their lamps.\par
\switchcolumn*

\LA
\begin{Resp}\Rsign Hæc est Virgo sápiens, quam Dóminus vigilántem invénit, quæ accéptis lampádibus sumpsit secum óleum: \Star{} Et veniénte Dómino, introívit cum eo ad núptias.\end{Resp}
\begin{Resp}\Vsign Média nocte clamor factus est: Ecce sponsus venit, exíte óbviam ei.\end{Resp}
\begin{Resp}\Rsign Et veniénte Dómino, introívit cum eo ad núptias.\end{Resp}
\switchcolumn
\EN
\begin{Resp}\Rsign This is one of those wise virgins, whom the Lord found watching, for when she took her lamp, she took oil with her. \Star{} And when the Lord came, she went in with him to the marriage.\end{Resp}
\begin{Resp}\Vsign At midnight there was a cry made: Behold, the Bridegroom cometh, go ye out to meet him.\end{Resp}
\begin{Resp}\Rsign And when the Lord came, she went in with him to the marriage.\end{Resp}
\switchcolumn*

\LA
\LessonHead{Lectio VIII}
\switchcolumn
\EN
\LessonHead{Lesson VIII}
\switchcolumn*

\LA
\noindent Sed prius quæréndum nobis est quid sit regnum cælórum, aut cur decem virgínibus comparétur, quæ étiam vírgines prudéntes et fátuæ dicántur. Dum enim cælórum regnum constat quia reprobórum nullus ingréditur, étiam fátuis virgínibus cur símile esse perhibétur? Sed sciéndum nobis est quod sæpe in sacro elóquio regnum cælórum præséntis témporis Ecclésia dícitur. De quo álio in loco Dóminus dicit: Mittet Fílius hóminis Ángelos suos, et cólligent de regno ejus ómnia scándala. Neque enim in illo regno beatitúdinis, in quo pax summa est, inveníri scándala póterunt, quæ colligántur.\par
\switchcolumn
\EN
\noindent First of all, it is for us to ask What is the kingdom of Heaven? And wherefore shall the same be likened unto ten virgins, whereof, albeit five were wise, yet five were foolish For if the kingdom of heaven be such that there shall in no wise enter into it anything that defileth, neither whatsoever worketh abomination, or maketh a lie, (Apoc. xxi. 27,) how can it be like unto five virgins which were foolish? But we must know that, in the word of God, the kingdom of heaven doth oftentimes signify the Church as she now is, touching the which the Lord saith in another place “The Son of Man shall send forth His Angels, and they shall gather out of His kingdom all things that offend.” (Matth. xiii. 41.) In that kingdom of Blessedness, wherein peace shall have her perfect reign, there shall be nothing found that offendeth for the angels to gather out.\par
\switchcolumn*

\LA
\begin{Resp}\Rsign Média nocte clamor factus est: \Star{} Ecce sponsus venit, exíte óbviam ei.\end{Resp}
\begin{Resp}\Vsign Prudéntes vírgines, aptáte vestras lámpades.\end{Resp}
\begin{Resp}\Rsign Ecce sponsus venit, exíte óbviam ei.\end{Resp}
\begin{Resp}\Vsign Glória Patri, et Fílio, \Star{} et Spirítui Sancto.\end{Resp}
\begin{Resp}\Rsign Ecce sponsus venit, exíte óbviam ei.\end{Resp}
\switchcolumn
\EN
\begin{Resp}\Rsign At midnight there was a cry made: \Star{} Behold, the Bridegroom cometh, go ye out to meet him.\end{Resp}
\begin{Resp}\Vsign Trim your lamps, O ye wise virgins.\end{Resp}
\begin{Resp}\Rsign Behold, the Bridegroom cometh, go ye out to meet him.\end{Resp}
\begin{Resp}\Vsign Glory be to the Father, and to the Son, \Star{} and to the Holy Ghost.\end{Resp}
\begin{Resp}\Rsign Behold, the Bridegroom cometh, go ye out to meet him.\end{Resp}
\switchcolumn*

\LA
\LessonHead{Lectio IX}
\switchcolumn
\EN
\LessonHead{Lesson IX}
\switchcolumn*

\LA
\noindent In quinque autem córporis sénsibus unusquísque subsístit; geminátus autem quinárius denárium pérficit. Et, quia ex utróque sexu fidélium multitúdo collígitur, sancta Ecclésia decem virgínibus símilis esse denuntiátur. In qua quia mali cum bonis et réprobi cum eléctis admíxti sunt, recte símilis virgínibus prudéntibus et fátuis esse perhibétur. Sunt namque pleríque continéntes, qui ab appetítu se exterióri custódiunt et spe ad interióra rapiúntur, carnem mácerant, et toto desidério ad supérnam pátriam anhélant, ætérna prǽmia éxpetunt, pro labóribus suis recípere laudes humánas nolunt. Hi nimírum glóriam suam non in ore hóminum ponunt, sed intra consciéntiam cóntegunt. Et sunt pleríque, qui corpus per abstinéntiam afflígunt, sed de ipsa sua abstinéntia humános favóres éxpetunt.\par
\switchcolumn
\EN
\noindent The body of every man doth consist of five senses, and five being doubled, is ten. Forasmuch, therefore, as the whole body of the faithful doth consist of two sexes, the Holy Church is likened unto ten virgins. And forasmuch as in the Church the good are for the present mingled with the bad, and the reprobate with the elect, it is rightly said that, of the ten virgins, five are wise and five are foolish. There are many who have self-control, which do keep themselves from lusting after things outward, whose hope beareth them to things inward, who chastise the flesh, who long with intense home-sickness for their Fatherland which is in heaven, who seek an eternal reward, and who will not to receive for their labours the praise of men. These are they who reckon their glory, not in the mouths of men, but in the testimony of their own conscience. And many there be likewise who afflict the body by self-control, and yet who seek for their self-control applause from men.\par
\switchcolumn*

\LA
\TeDeum{Te Deum}
\switchcolumn
\EN
\TeDeum{Te Deum}
\switchcolumn*

\end{paracol}

\Office{Commune non Virginum non Martyrum}{Common of Widows}{}
\addcontentsline{toc}{subsection}{Commune non Virginum non Martyrum\enspace\textit{Common of Widows}}
\begin{paracol}{2}
\LA
\LessonHead{Lectio I}
\switchcolumn
\EN
\LessonHead{Lesson I}
\switchcolumn*

\LA
\LessonTitle{De Parábolis Salomónis}
\switchcolumn
\EN
\LessonTitle{Lesson from the book of Proverbs}
\switchcolumn*

\LA
\Cite{Prov 31:10-17}
\switchcolumn
\EN
\Cite{Prov 31:10-17}
\switchcolumn*

\LA
\noindent Mulíerem fortem quis invéniet? procul et de últimis fínibus prétium ejus. Confídit in ea cor viri sui, et spóliis non indigébit. Reddet ei bonum, et non malum, ómnibus diébus vitæ suæ. Quæsívit lanam et linum, et operáta est consília mánuum suárum. Facta est quasi navis institóris, de longe portans panem suum. Et de nocte surréxit, dedítque prædam domésticis suis, et cibária ancíllis suis. Considerávit agrum, et emit eum; de fructu mánuum suárum plantávit víneam. Accínxit fortitúdine lumbos suos, et roborávit bráchium suum.\par
\switchcolumn
\EN
\noindent Who shall find a valiant woman? Far and from the uttermost coasts is the price of her. The heart of her husband trusteth in her, and he shall have no need of spoils. She will render him good, and not evil, all the days of her life. She hath sought wool and flax, and hath wrought by the counsel of her hands. She is like the merchant's ship, she bringeth her bread from afar. And she hath risen in the night, and given a prey to her household, and victuals to her maidens. She hath considered a field, and bought it: with the fruit of her hands she hath planted a vineyard. She hath girded her loins with strength, and hath strengthened her arm\par
\switchcolumn*

\LA
\begin{Resp}\Rsign Veni, elécta mea, et ponam in te thronum meum \Star{} Quia concupívit Rex spéciem tuam.\end{Resp}
\begin{Resp}\Vsign Spécie tua et pulchritúdine tua inténde, próspere procéde, et regna.\end{Resp}
\begin{Resp}\Rsign Quia concupívit Rex spéciem tuam.\end{Resp}
\switchcolumn
\EN
\begin{Resp}\Rsign Come, O My chosen one, and I will establish My throne in thee; \Star{} For the King hath greatly desired thy beauty.\end{Resp}
\begin{Resp}\Vsign In thy comeliness and thy beauty, go forward, fare prosperously, and reign.\end{Resp}
\begin{Resp}\Rsign For the King hath greatly desired thy beauty.\end{Resp}
\switchcolumn*

\LA
\LessonHead{Lectio II}
\switchcolumn
\EN
\LessonHead{Lesson II}
\switchcolumn*

\LA
\Cite{Prov 31:18-24}
\switchcolumn
\EN
\Cite{Prov 31:18-24}
\switchcolumn*

\LA
\noindent Gustávit, et vidit quia bona est negotiátio ejus; non extinguétur in nocte lucérna ejus. Manum suam misit ad fórtia, et dígiti ejus apprehendérunt fusum. Manum suam apéruit ínopi, et palmas suas exténdit ad páuperem. Non timébit dómui suæ a frigóribus nivis; omnes enim doméstici ejus vestíti sunt duplícibus. Stragulátam vestem fecit sibi; byssus et púrpura induméntum ejus. Nóbilis in portis vir ejus, quando séderit cum senatóribus terræ. Síndonem fecit, et véndidit, et cíngulum trádidit Chananǽo.\par
\switchcolumn
\EN
\noindent She hath tasted and seen that her traffic is good: her lamp shall not be put out in the night. She hath put out her hand to strong things, and her fingers have taken hold of the spindle. She hath opened her hand to the needy, and stretched out her hands to the poor. She shall not fear for her house in the cold of snow: for all her domestics are clothed with double garments. She hath made for herself clothing of tapestry: fine linen, and purple is her covering. Her husband is honourable in the gates, when he sitteth among the senators of the land. She made fine linen, and sold it, and delivered a girdle to the Chanaanite.\par
\switchcolumn*

\LA
\begin{Resp}\Rsign Diffúsa est grátia in lábiis tuis, \Star{} Proptérea benedíxit te Deus in ætérnum.\end{Resp}
\begin{Resp}\Vsign Spécie tua et pulchritúdine tua inténde, próspere procéde, et regna.\end{Resp}
\begin{Resp}\Rsign Proptérea benedíxit te Deus in ætérnum.\end{Resp}
\switchcolumn
\EN
\begin{Resp}\Rsign Grace is poured into thy lips, therefore; \Star{} God hath blessed thee for ever.\end{Resp}
\begin{Resp}\Vsign In thy comeliness and thy beauty, go forward, fare prosperously, and reign.\end{Resp}
\begin{Resp}\Rsign God hath blessed thee for ever.\end{Resp}
\switchcolumn*

\LA
\LessonHead{Lectio III}
\switchcolumn
\EN
\LessonHead{Lesson III}
\switchcolumn*

\LA
\Cite{Prov 31:25-31}
\switchcolumn
\EN
\Cite{Prov 31:25-31}
\switchcolumn*

\LA
\noindent Fortitúdo et decor induméntum ejus, et ridébit in die novíssimo. Os suum apéruit sapiéntiæ, et lex cleméntiæ in lingua ejus. Considerávit sémitas domus suæ, et panem otiósa non comédit. Surrexérunt fílii ejus, et beatíssimam prædicavérunt; vir ejus, et laudávit eam. Multæ fíliæ congregavérunt divítias; tu supergréssa es univérsas. Fallax grátia, et vana est pulchritúdo: múlier timens Dóminum, ipsa laudábitur. Date ei de fructu mánuum suárum, et laudent eam in portis ópera ejus.\par
\switchcolumn
\EN
\noindent Strength and beauty are her clothing, and she shall laugh in the latter day. She hath opened her mouth to wisdom, and the law of clemency is on her tongue. She hath looked well to the paths of her house, and hath not eaten her bread idle. Her children rose up, and called her blessed: her husband, and he praised her. Many daughters have gathered together riches: thou hast surpassed them all. Favour is deceitful, and beauty is vain: the woman that feareth the Lord, she shall be praised. Give her of the fruit of her hands: and let her works praise her in the gates.\par
\switchcolumn*

\LA
\begin{Resp}\Rsign Spécie tua et pulchritúdine tua \Star{} Inténde, próspere procéde, et regna.\end{Resp}
\begin{Resp}\Vsign Diffúsa est grátia in lábiis tuis, proptérea benedíxit te Deus in ætérnum.\end{Resp}
\begin{Resp}\Rsign Inténde, próspere procéde, et regna.\end{Resp}
\begin{Resp}\Vsign Glória Patri, et Fílio, \Star{} et Spirítui Sancto.\end{Resp}
\begin{Resp}\Rsign Inténde, próspere procéde, et regna.\end{Resp}
\switchcolumn
\EN
\begin{Resp}\Rsign In thy comeliness and thy beauty; \Star{} Go forward, fare prosperously, and reign.\end{Resp}
\begin{Resp}\Vsign Grace is poured into thy lips, therefore God hath blessed thee for ever.\end{Resp}
\begin{Resp}\Rsign Go forward, fare prosperously, and reign.\end{Resp}
\begin{Resp}\Vsign Glory be to the Father, and to the Son, \Star{} and to the Holy Ghost.\end{Resp}
\begin{Resp}\Rsign Go forward, fare prosperously, and reign.\end{Resp}
\switchcolumn*

\LA
\LessonHead{Lectio VII}
\switchcolumn
\EN
\LessonHead{Lesson VII}
\switchcolumn*

\LA
\LessonTitle{Léctio sancti Evangélii secúndum Matthǽum}
\switchcolumn
\EN
\LessonTitle{From the Holy Gospel according to Matthew}
\switchcolumn*

\LA
\Cite{Matt 13:44-52}
\switchcolumn
\EN
\Cite{Matt 13:44}
\switchcolumn*

\LA
\noindent In illo témpore: Dixit Jesus discípulis suis parábolam hanc: Simile est regnum cælórum thesáuro abscóndito in agro. Et réliqua.\par
\switchcolumn
\EN
\noindent At that time Jesus spake unto His disciples this parable The kingdom of heaven is like unto treasure hid in a field. And so on.\par
\switchcolumn*

\LA
\LessonTitle{Homilía sancti Gregórii Papæ}
\switchcolumn
\EN
\LessonTitle{Homily by Pope St Gregory the Great}
\switchcolumn*

\LA
\Source{Homilia 11 in Evangelia}
\switchcolumn
\EN
\Source{11th on the Gospels}
\switchcolumn*

\LA
\noindent Cælórum regnum, fratres caríssimi, idcírco terrénis rebus símile dícitur, ut, ex his quæ ánimus novit, surgat ad incógnita, quæ non novit: quátenus exémplo visibílium se ad invisibília rápiat, et, per ea quæ usu dídicit quasi confricátus incaléscat; ut per hoc, quod scit notum dilígere, discat et incógnita amáre. Ecce enim cælórum regnum thesáuro abscóndito in agro comparátur; quem, qui invénit homo, abscóndit, et præ gáudio illíus vadit, et vendit univérsa quæ habet, et emit agrum illum.\par
\switchcolumn
\EN
\noindent Dearly beloved brethren, the kingdom of heaven is likened unto the things of earth, to the end that by the mean of things which we know, our mind may rise to the contemplation of the things which we know not by the example of things which are seen, may fix her gaze on things which are not seen by the touch of things which she useth, may be warmed towards the things which she useth not; by things which she knoweth and loveth, to love also the things which she knoweth not. For, behold, “the kingdom of heaven is likened unto treasure hid in a field, the which when a man hath found, he hideth, and, for joy thereof, goeth and selleth all that he hath and buyeth that field.”\par
\switchcolumn*

\LA
\begin{Resp}\Rsign Os suum apéruit sapiéntiæ, et lex cleméntiæ in lingua ejus: considerávit sémitas domus suæ, \Star{} Et panem otiósa non comédit.\end{Resp}
\begin{Resp}\Vsign Gustávit et vidit quia bona est negotiátio ejus: non exstinguétur in nocte lucérna ejus.\end{Resp}
\begin{Resp}\Rsign Et panem otiósa non comédit.\end{Resp}
\switchcolumn
\EN
\begin{Resp}\Rsign She openeth her mouth with wisdom, and in her tongue is the law of kindness. She looketh well to the ways of her household \Star{} And eateth not the bread of idleness.\end{Resp}
\begin{Resp}\Vsign She tasteth and perceiveth that her merchandise is good. Her candle goeth not out by night.\end{Resp}
\begin{Resp}\Rsign And she eateth not the bread of idleness.\end{Resp}
\switchcolumn*

\LA
\LessonHead{Lectio VIII}
\switchcolumn
\EN
\LessonHead{Lesson VIII}
\switchcolumn*

\LA
\noindent Qua in re hoc quoque notándum est, quod invéntus thesáurus abscónditur, ut servétur: quia stúdium cæléstis desidérii a malígnis spirítibus custodíre non súfficit, qui hoc ab humánis láudibus non abscóndit. In præsénti étenim vita, quasi in via sumus, qua ad pátriam pérgimus. Malígni autem spíritus iter nostrum quasi quidam latrúnculi óbsident. Deprædári ergo desíderat, qui thesáurum públice portat in via. Hoc autem dico, non ut próximi ópera nostra bona non vídeant, cum scriptum sit: Vídeant ópera vestra bona, et gloríficent Patrem vestrum qui in cælis est; sed, ut per hoc quod ágimus, laudes extérius non quærámus. Sic autem sit opus in público, quátenus inténtio máneat in occúlto; ut, et de bono ópere próximis præbeámus exémplum, et tamen per intentiónem, qua Deo soli placére quǽrimus, semper optémus secrétum.\par
\switchcolumn
\EN
\noindent And herein we must remark that the treasure, when once it hath been found, is hidden to keep it safe. He whose intimate yearnings after God are not hidden from the praise of men, is open thereby to the attacks of evil spirits. In this life we are, as it were, journeying homewards on a road beset by evil spirits who are like highwaymen. He therefore inviteth robbery who carrieth his treasure ostentatiously. Doubtless our neighbour should be able to see our good works, as it is written: “Let your light so shine before men that they may see your good works, and glorify your Father which is in heaven.” But this is not to be understood to mean that we are to seek the praise of men by what we do. Rather, let us in such wise work in the open that the inner intention of devotion is not advertised. So we shall give an example to our neighbour, and yet keep hidden, except from the sight of God, our purpose of pleasing him.\par
\switchcolumn*

\LA
\begin{Resp}\Rsign Regnum mundi et omnem ornátum sǽculi contémpsi, propter amórem Dómini mei Jesu Christi: \Star{} Quem vidi, quem amávi, in quem crédidi, quem diléxi.\end{Resp}
\begin{Resp}\Vsign Eructávit cor meum verbum bonum: dico ego ópera mea Regi.\end{Resp}
\begin{Resp}\Rsign Quem vidi, quem amávi, in quem crédidi, quem diléxi.\end{Resp}
\begin{Resp}\Vsign Glória Patri, et Fílio, \Star{} et Spirítui Sancto.\end{Resp}
\begin{Resp}\Rsign Quem vidi, quem amávi, in quem crédidi, quem diléxi.\end{Resp}
\switchcolumn
\EN
\begin{Resp}\Rsign The kingdom of this world and all the beauty of life I have esteemed as nothing, for the excellency of the love of Jesus Christ my Lord \Star{} Whom, having seen, I loved Whom, having believed, I longed after.\end{Resp}
\begin{Resp}\Vsign My heart is overflowing with a good matter I speak of my works unto the King.\end{Resp}
\begin{Resp}\Rsign Whom, having seen, I loved Whom, having believed, I longed after.\end{Resp}
\begin{Resp}\Vsign Glory be to the Father, and to the Son, \Star{} and to the Holy Ghost.\end{Resp}
\begin{Resp}\Rsign Whom, having seen, I loved Whom, having believed, I longed after.\end{Resp}
\switchcolumn*

\end{paracol}

\Office{Commune plurimarum non Virginum Martyrum}{Common of multiple widowed Martyrs}{}
\addcontentsline{toc}{subsection}{Commune plurimarum non Virginum Martyrum\enspace\textit{Common of multiple widowed Martyrs}}
\begin{paracol}{2}
\LA
\LessonHead{Lectio I}
\switchcolumn
\EN
\LessonHead{Lesson I}
\switchcolumn*

\LA
\noindent Confitébor tibi, Dómine, Rex, et collaudábo te Deum Salvatórem meum. Confitébor nómini tuo: quóniam adjútor et protéctor factus es mihi, et liberásti corpus meum a perditióne, a láqueo linguæ iníquæ et a lábiis operántium mendácium, et in conspéctu astántium factus es mihi adjútor. Et liberásti me secúndum multitúdinem misericórdiæ nóminis tui a rugiéntibus, præparátis ad escam: de mánibus quæréntium ánimam meam, et de portis tribulatiónum, quæ circumdedérunt me; a pressúra flammæ, quæ circúmdedit me, et in médio ignis non sum æstuáta; de altitúdine ventris ínferi, et a lingua coinquináta, et a verbo mendácii, a rege iníquo, et a lingua injústa.\par
\switchcolumn
\EN
\noindent A prayer of Jesus the son of Sirach. I will give glory to thee, O Lord, O King, and I will praise thee, O God my Saviour. I will give glory to thy name: for thou hast been a helper and protector to me. And hast preserved my body from destruction, from the snare of an unjust tongue, and from the lips of them that forge lies, and in the sight of them that stood by, thou hast been my helper. And thou hast delivered me, according to the multitude of the mercy of thy name, from them that did roar, prepared to devour. Out of the hands of them that sought my life, and from the gates of afflictions, which compassed me about: From the oppression of the flame which surrounded me, and in the midst of the fire I was not burnt. From the depth of the belly of hell, and from an unclean tongue, and from lying words, from an unjust king, and from a slanderous tongue:\par
\switchcolumn*

\LA
\begin{Resp}\Rsign Veni Sponsa Christi, áccipe corónam, quam tibi Dóminus præparávit in ætérnum; pro cujus amóre sánguinem tuum fudísti, \Star{} Et cum Ángelis in paradísum introísti.\end{Resp}
\begin{Resp}\Vsign Veni, elécta mea, et ponam in te thronum meum; quia concupívit Rex spéciem tuam.\end{Resp}
\begin{Resp}\Rsign Et cum Ángelis in paradísum introísti.\end{Resp}
\switchcolumn
\EN

\switchcolumn*

\LA
\LessonHead{Lectio II}
\switchcolumn
\EN
\LessonHead{Lesson II}
\switchcolumn*

\LA
\Cite{Sir 51:8-12}
\switchcolumn
\EN
\Cite{Sir 51:8-12}
\switchcolumn*

\LA
\noindent Laudábit usque ad mortem ánima mea Dóminum, et vita mea appropínquans erat in inférno deórsum. Circumdedérunt me úndique, et non erat qui adjuváret. Respíciens eram ad adjutórium hóminum, et non erat. Memoráta sum misericórdiæ tuæ, Dómine, et operatiónis tuæ, quæ a sǽculo sunt: quóniam éruis sustinéntes te, Dómine, et líberas eos de mánibus géntium.\par
\switchcolumn
\EN
\noindent My soul shall praise the Lord even to death. And my life was drawing near to hell beneath. They compassed me on every side, and there was no one that would help me. I looked for the succour of men, and there was none. I remembered thy mercy, O Lord, and thy works, which are from the beginning of the world. How thou deliverest them that wait for thee, O Lord, and savest them out of the hands of the nations.\par
\switchcolumn*

\LA
\begin{Resp}\Rsign Diffúsa est grátia in lábiis tuis, \Star{} Proptérea benedíxit te Deus in ætérnum.\end{Resp}
\begin{Resp}\Vsign Spécie tua et pulchritúdine tua inténde, próspere procéde, et regna.\end{Resp}
\begin{Resp}\Rsign Proptérea benedíxit te Deus in ætérnum.\end{Resp}
\switchcolumn
\EN
\begin{Resp}\Rsign Grace is poured into thy lips, therefore; \Star{} God hath blessed thee for ever.\end{Resp}
\begin{Resp}\Vsign In thy comeliness and thy beauty, go forward, fare prosperously, and reign.\end{Resp}
\begin{Resp}\Rsign God hath blessed thee for ever.\end{Resp}
\switchcolumn*

\LA
\LessonHead{Lectio III}
\switchcolumn
\EN
\LessonHead{Lesson III}
\switchcolumn*

\LA
\Cite{Sir 51:13-17}
\switchcolumn
\EN
\Cite{Sir 51:13-17}
\switchcolumn*

\LA
\noindent Exaltásti super terram habitatiónem meam, et pro morte defluénte deprecáta sum. Invocávi Dóminum, Patrem Dómini mei, ut non derelínquat me in die tribulatiónis meæ, et in témpore superbórum, sine adjutório. Laudábo nomen tuum assídue, et collaudábo illud in confessióne, et exaudíta est orátio mea. Et liberásti me de perditióne, et eripuísti me de témpore iníquo. Proptérea confitébor, et laudem dicam tibi, et benedícam nómini Dómini.\par
\switchcolumn
\EN
\noindent Thou hast exalted my dwelling place upon the earth and I have prayed for death to pass away. I called upon the Lord, the father of my Lord, that he would not leave me in the day of my trouble, and in the time of the proud without help. I will praise thy name continually, and will praise it with thanksgiving, and my prayer was heard. And thou hast saved me from destruction, and hast delivered me from the evil time. Therefore I will give thanks, and praise thee, and bless the name of the Lord.\par
\switchcolumn*

\LA
\begin{Resp}\Rsign Spécie tua et pulchritúdine tua \Star{} Inténde, próspere procéde, et regna.\end{Resp}
\begin{Resp}\Vsign Diffúsa est grátia in lábiis tuis, proptérea benedíxit te Deus in ætérnum.\end{Resp}
\begin{Resp}\Rsign Inténde, próspere procéde, et regna.\end{Resp}
\begin{Resp}\Vsign Glória Patri, et Fílio, \Star{} et Spirítui Sancto.\end{Resp}
\begin{Resp}\Rsign Inténde, próspere procéde, et regna.\end{Resp}
\switchcolumn
\EN
\begin{Resp}\Rsign In thy comeliness and thy beauty; \Star{} Go forward, fare prosperously, and reign.\end{Resp}
\begin{Resp}\Vsign Grace is poured into thy lips, therefore God hath blessed thee for ever.\end{Resp}
\begin{Resp}\Rsign Go forward, fare prosperously, and reign.\end{Resp}
\begin{Resp}\Vsign Glory be to the Father, and to the Son, \Star{} and to the Holy Ghost.\end{Resp}
\begin{Resp}\Rsign Go forward, fare prosperously, and reign.\end{Resp}
\switchcolumn*

\LA
\LessonHead{Lectio VII}
\switchcolumn
\EN
\LessonHead{Lesson VII}
\switchcolumn*

\LA
\LessonTitle{Léctio sancti Evangélii secúndum Matthǽum}
\switchcolumn
\EN
\LessonTitle{From the Holy Gospel according to Matthew}
\switchcolumn*

\LA
\Cite{Matt 13:44-52}
\switchcolumn
\EN
\Cite{Matt 13:44}
\switchcolumn*

\LA
\noindent In illo témpore: Dixit Jesus discípulis suis parábolam hanc: Simile est regnum cælórum thesáuro abscóndito in agro. Et réliqua.\par
\switchcolumn
\EN
\noindent At that time Jesus spake unto His disciples this parable The kingdom of heaven is like unto treasure hid in a field. And so on.\par
\switchcolumn*

\LA
\LessonTitle{Homilía sancti Gregórii Papæ}
\switchcolumn
\EN
\LessonTitle{Homily by Pope St Gregory the Great}
\switchcolumn*

\LA
\Source{Homilia 11 in Evangelia}
\switchcolumn
\EN
\Source{11th on the Gospels}
\switchcolumn*

\LA
\noindent Cælórum regnum, fratres caríssimi, idcírco terrénis rebus símile dícitur, ut, ex his quæ ánimus novit, surgat ad incógnita, quæ non novit: quátenus exémplo visibílium se ad invisibília rápiat, et, per ea quæ usu dídicit quasi confricátus incaléscat; ut per hoc, quod scit notum dilígere, discat et incógnita amáre. Ecce enim cælórum regnum thesáuro abscóndito in agro comparátur; quem, qui invénit homo, abscóndit, et præ gáudio illíus vadit, et vendit univérsa quæ habet, et emit agrum illum.\par
\switchcolumn
\EN
\noindent Dearly beloved brethren, the kingdom of heaven is likened unto the things of earth, to the end that by the mean of things which we know, our mind may rise to the contemplation of the things which we know not by the example of things which are seen, may fix her gaze on things which are not seen by the touch of things which she useth, may be warmed towards the things which she useth not; by things which she knoweth and loveth, to love also the things which she knoweth not. For, behold, “the kingdom of heaven is likened unto treasure hid in a field, the which when a man hath found, he hideth, and, for joy thereof, goeth and selleth all that he hath and buyeth that field.”\par
\switchcolumn*

\LA
\begin{Resp}\Rsign Os suum apéruit sapiéntiæ, et lex cleméntiæ in lingua ejus: considerávit sémitas domus suæ, \Star{} Et panem otiósa non comédit.\end{Resp}
\begin{Resp}\Vsign Gustávit et vidit quia bona est negotiátio ejus: non exstinguétur in nocte lucérna ejus.\end{Resp}
\begin{Resp}\Rsign Et panem otiósa non comédit.\end{Resp}
\switchcolumn
\EN
\begin{Resp}\Rsign She openeth her mouth with wisdom, and in her tongue is the law of kindness. She looketh well to the ways of her household \Star{} And eateth not the bread of idleness.\end{Resp}
\begin{Resp}\Vsign She tasteth and perceiveth that her merchandise is good. Her candle goeth not out by night.\end{Resp}
\begin{Resp}\Rsign And she eateth not the bread of idleness.\end{Resp}
\switchcolumn*

\LA
\LessonHead{Lectio VIII}
\switchcolumn
\EN
\LessonHead{Lesson VIII}
\switchcolumn*

\LA
\noindent Qua in re hoc quoque notándum est, quod invéntus thesáurus abscónditur, ut servétur: quia stúdium cæléstis desidérii a malígnis spirítibus custodíre non súfficit, qui hoc ab humánis láudibus non abscóndit. In præsénti étenim vita, quasi in via sumus, qua ad pátriam pérgimus. Malígni autem spíritus iter nostrum quasi quidam latrúnculi óbsident. Deprædári ergo desíderat, qui thesáurum públice portat in via. Hoc autem dico, non ut próximi ópera nostra bona non vídeant, cum scriptum sit: Vídeant ópera vestra bona, et gloríficent Patrem vestrum qui in cælis est; sed, ut per hoc quod ágimus, laudes extérius non quærámus. Sic autem sit opus in público, quátenus inténtio máneat in occúlto; ut, et de bono ópere próximis præbeámus exémplum, et tamen per intentiónem, qua Deo soli placére quǽrimus, semper optémus secrétum.\par
\switchcolumn
\EN
\noindent And herein we must remark that the treasure, when once it hath been found, is hidden to keep it safe. He whose intimate yearnings after God are not hidden from the praise of men, is open thereby to the attacks of evil spirits. In this life we are, as it were, journeying homewards on a road beset by evil spirits who are like highwaymen. He therefore inviteth robbery who carrieth his treasure ostentatiously. Doubtless our neighbour should be able to see our good works, as it is written: “Let your light so shine before men that they may see your good works, and glorify your Father which is in heaven.” But this is not to be understood to mean that we are to seek the praise of men by what we do. Rather, let us in such wise work in the open that the inner intention of devotion is not advertised. So we shall give an example to our neighbour, and yet keep hidden, except from the sight of God, our purpose of pleasing him.\par
\switchcolumn*

\LA
\begin{Resp}\Rsign Regnum mundi et omnem ornátum sǽculi contémpsi, propter amórem Dómini mei Jesu Christi: \Star{} Quem vidi, quem amávi, in quem crédidi, quem diléxi.\end{Resp}
\begin{Resp}\Vsign Eructávit cor meum verbum bonum: dico ego ópera mea Regi.\end{Resp}
\begin{Resp}\Rsign Quem vidi, quem amávi, in quem crédidi, quem diléxi.\end{Resp}
\begin{Resp}\Vsign Glória Patri, et Fílio, \Star{} et Spirítui Sancto.\end{Resp}
\begin{Resp}\Rsign Quem vidi, quem amávi, in quem crédidi, quem diléxi.\end{Resp}
\switchcolumn
\EN
\begin{Resp}\Rsign The kingdom of this world and all the beauty of life I have esteemed as nothing, for the excellency of the love of Jesus Christ my Lord \Star{} Whom, having seen, I loved Whom, having believed, I longed after.\end{Resp}
\begin{Resp}\Vsign My heart is overflowing with a good matter I speak of my works unto the King.\end{Resp}
\begin{Resp}\Rsign Whom, having seen, I loved Whom, having believed, I longed after.\end{Resp}
\begin{Resp}\Vsign Glory be to the Father, and to the Son, \Star{} and to the Holy Ghost.\end{Resp}
\begin{Resp}\Rsign Whom, having seen, I loved Whom, having believed, I longed after.\end{Resp}
\switchcolumn*

\end{paracol}

\Office{Sanctæ Mariæ Sabbato (a Nativitate usque ad Purificationem)}{Our Lady's Saturday (from Christmas to the Purification)}{Simplex}
\addcontentsline{toc}{subsection}{Sanctæ Mariæ Sabbato (a Nativitate usque ad Purificationem)\enspace\textit{Our Lady's Saturday (from Christmas to the Purification)}}
\begin{paracol}{2}
\LA
\RefLine{Lectiones I–II: de Scriptura occurrente.}
\switchcolumn
\EN
\RefLine{Lessons I–II: from the Scripture of the day.}
\switchcolumn*

\LA
\LessonHead{Lectio III (mense Januario)}
\switchcolumn
\EN
\LessonHead{Lesson III (in January)}
\switchcolumn*

\LA
\LessonTitle{Ex Epístola sancti Ambrósii Epíscopi ad Sirícium Papam}
\switchcolumn
\EN
\LessonTitle{From the letter of Blessed Ambrose, bishop of Milan, to Pope Siricius}
\switchcolumn*

\LA
\Source{Epist. 81, alias 7, post initium}
\switchcolumn
\EN
\Source{Epist. 81, alias 7}
\switchcolumn*

\LA
\noindent De via perversitátis prodúntur dícere: Virgo concépit, sed non virgo generávit. Pótuit ergo virgo concípere, non pótuit virgo generáre, cum semper concéptus præcédat, partus sequátur? Sed, si doctrínis non créditur sacerdótum, credátur oráculis Christi; credátur mónitis Angelórum dicéntium: Quia non est impossíbile Deo omne verbum; credátur Sýmbolo Apostolórum, quod Ecclésia Romána intemerátum semper custódit et servat. Audívit María vocem Angeli, et, quæ ante díxerat: Quómodo fiet istud; non de fide generatiónis intérrogans, respóndit póstea: Ecce ancílla Dómini, contíngat mihi secúndum verbum tuum.\par
\switchcolumn
\EN
\noindent How great is the madness of their dismal barkings, that the same persons should say that Christ could not be born of a virgin, and yet assert that women, after having given birth to human pledges, remain virgins? But if they will not believe the doctrines of the Clergy, let them believe the oracles of Christ, let them believe the admonitions of Angels who say, For with God nothing shall be impossible. Let them give credit to the Creed of the Apostles, which the Roman Church has always kept and preserved undefiled. Mary heard the voice of the Angel, and she who before had said “How shall this be?”, not asking from want of faith in the mode of generation, afterwards replied: “Behold the handmaid of the Lord, be it unto me according to thy word.”\par
\switchcolumn*

\LA
\TeDeum{Te Deum}
\switchcolumn
\EN
\TeDeum{Te Deum}
\switchcolumn*

\end{paracol}

\Office{Sanctæ Mariæ Sabbato (post Purificationem usque ad Dominicam Palmarum)}{Our Lady's Saturday (from the Purification to Palm Sunday)}{Simplex}
\addcontentsline{toc}{subsection}{Sanctæ Mariæ Sabbato (post Purificationem usque ad Dominicam Palmarum)\enspace\textit{Our Lady's Saturday (from the Purification to Palm Sunday)}}
\begin{paracol}{2}
\LA
\RefLine{Lectiones I–II: de Scriptura occurrente.}
\switchcolumn
\EN
\RefLine{Lessons I–II: from the Scripture of the day.}
\switchcolumn*

\LA
\LessonHead{Lectio III (mense Februario)}
\switchcolumn
\EN
\LessonHead{Lesson III (in February)}
\switchcolumn*

\LA
\LessonTitle{Ex libro sancti Hierónymi Presbýteri advérsus Joviniánum}
\switchcolumn
\EN
\LessonTitle{From the book of St. Jerome, Priest, against Jovian}
\switchcolumn*

\LA
\Source{Apologia ad Pammachium pro lib. adversus Jovinianum, in fine}
\switchcolumn
\EN
\Source{Apologia to Panmachius to the Book against Jovian, at end}
\switchcolumn*

\LA
\noindent Christus virgo, Mater vírginis nostri virgo perpétua, mater et virgo. Jesus enim clausis ingréssus est óstiis; et in sepúlcro ejus, quod novum et in petra duríssima fúerat excísum, nec ántea quis nec póstea pósitus est. Hortus conclúsus, fons signátus. De quo fonte ille flúvius manat, juxta Joël, qui írrigat torréntem vel fúnium vel spinárum: fúnium, peccatórum, quibus ante alligabámur; spinárum, quæ suffócant seméntem patrisfamílias. Hæc est porta orientális, ut ait Ezéchiel, semper clausa et lúcida, opériens in se vel ex se próferens Sancta sanctórum; per quam Sol justítiæ et Póntifex noster secúndum órdinem Melchísedech ingréditur et egréditur.\par
\switchcolumn
\EN
\noindent Christ is virgin, and the mother of our virgin Christ remains forever virgin, the Virgin Mother. For Jesus entered the world, as he entered the upper room, “the doors being closed”, and in his sepulcher, which had been hewn out of the hard rock, “no one had yet been laid”, nor was anyone else laid in it afterward. “An enclosed garden a fountain sealed” was Mary. And from that fountain, according to Joel, flows the stream that waters the torrent either of “ropes” or of “thorns”. The ropes would be the sins by which we had been bound down; the thorns would be those which choke the householder's seed. Mary is the eastern gate, of which Ezechiel speaks, always closed and luminous, whether concealing in itself or bringing forth from itself the Holy of holies. She is the gate through which the Sun of justice and our High Priest according to the order of Melchisedech goes in and out.\par
\switchcolumn*

\LA
\TeDeum{Te Deum}
\switchcolumn
\EN
\TeDeum{Te Deum}
\switchcolumn*

\LA
\LessonHead{Lectio III (mense Martio)}
\switchcolumn
\EN
\LessonHead{Lesson III (in March)}
\switchcolumn*

\LA
\LessonTitle{Ex libro sancti Irenǽi Epíscopi et Mártyris advérsus hǽreses}
\switchcolumn
\EN
\LessonTitle{From the book of St. Irenaeus, Bishop and Martyr, Against Heretics}
\switchcolumn*

\LA
\Source{Lib. 5, c. 19}
\switchcolumn
\EN
\Source{Book 5, ch. 19}
\switchcolumn*

\LA
\noindent In sua própria veniénte Dómino et sua própria eum bajulánte conditióne, quæ bajulátur ab ipso, et recapitulatiónem ejus quæ in ligno fuit inobediéntiæ, per eam quæ in ligno est obediéntiam, faciénte, et seductióne illa solúta, qua sedúcta est male illa, quæ jam viro destináta erat, virgo Heva; per veritátem evangelizáta est bene ab Angelo jam sub viro Virgo María. Quemádmodum enim illa per Angélicum sermónem sedúcta est, ut effúgeret Deum, prævaricáta verbum ejus: ita et hæc per Angélicum sermónem evangelizáta est, ut portáret Deum, obédiens ejus verbo. Et sicut illa sedúcta est, ut effúgeret Deum; sic hæc suása est obedíre Deo, ut vírginis Hevæ Virgo María fíeret advocáta. Et quemádmodum adstríctum est morte genus humánum per vírginem, solvátur per Vírginem: æqua lance dispósita virginális inobediéntia per Virginitátem obediéntiam.\par
\switchcolumn
\EN
\noindent That the Lord should come “unto His own” and that His own creation should bear Him by whom itself was born; that He should make amends through His obedience on the tree of the Cross for the disobedience committed at the tree in Paradise; that He should annul the effects of the seduction by which the virgin Eve, already destined for a husband, was led into evil - all this good news, this true revelation was fittingly brought by an Angel to the Virgin Mary, already espoused to a husband. For, just as Eve was seduced by the words of a fallen angel into turning away from God by disobeying His command, so Mary was evangelized by the words of an Angel, who persuaded her to give birth to God in obedience to His word. One has lured into fleeing from God, the other was persuaded to obey God, that the Virgin Mary might became an advocate for the virgin Eve. Through a virgin mankind was tied to death; so also through a virgin the bonds loosed; the Virgin's obedience balanced the virgin's disobedience.\par
\switchcolumn*

\LA
\TeDeum{Te Deum}
\switchcolumn
\EN
\TeDeum{Te Deum}
\switchcolumn*

\end{paracol}
\backmatter
\IndexOfOffices
\end{document}
